% Federigo Enriques, Le superficie algebriche (1949).
% Complete editable reader edition. Compile with pdfLaTeX, twice (or latexmk -pdf).
% Source comments beginning % ?? and % [sic] are retained verbatim.
% Resolved doubts carry a separate EDITORIAL status comment and no active doubt note.
\documentclass[11pt,a4paper,twoside,openany,leqno]{book}
\usepackage[T1]{fontenc}
\usepackage[utf8]{inputenc}
\usepackage{lmodern}
\usepackage[sc]{mathpazo}
\AtBeginDocument{\DeclareFontShape{T1}{pplx}{b}{sc}{<->ssub*pplx/m/sc}{}}
\usepackage{amsmath,amssymb,mathtools}
\usepackage{microtype}
\usepackage{geometry}
\geometry{inner=30mm,outer=46mm,top=24mm,bottom=26mm,headheight=25pt,headsep=8mm,marginparwidth=35mm,marginparsep=6mm}
\usepackage{babel}
\IfFileExists{italian.ldf}{\babelprovide[import,main]{italian}}{\babelprovide[import,main,hyphenrules=nohyphenation]{italian}}
\usepackage{array,longtable,booktabs}
\usepackage{ragged2e}
\usepackage{xcolor}
\definecolor{editiongray}{gray}{0.28}
\usepackage{marginnote,marginfix}
\usepackage{fancyhdr}
\usepackage{needspace}
\usepackage{titlesec}
\usepackage[unicode,pdfencoding=auto,psdextra,hyperfootnotes=false,hidelinks]{hyperref}
\usepackage{bookmark}
\hypersetup{pdftitle={Le superficie algebriche},pdfauthor={Federigo Enriques},pdfsubject={Edizione di lettura dalla stampa Zanichelli, Bologna 1949},pdfcreator={pdfLaTeX; edizione IA, novembre 2026},bookmarksopen=true,bookmarksopenlevel=0}
\newif\ifshownotes
\shownotestrue % Change this ONE flag to \shownotesfalse for a reading copy without editorial marginal notes.
\setlength{\parindent}{1.3em}
\setlength{\parskip}{0pt plus .35pt}
\linespread{1.045}
\raggedbottom
\setlength{\emergencystretch}{2.5em}
\tolerance=1600
\clubpenalty=8000
\widowpenalty=8000
\displaywidowpenalty=5000
\allowdisplaybreaks[2]
\setlength{\abovedisplayskip}{10pt plus3pt minus3pt}
\setlength{\belowdisplayskip}{10pt plus3pt minus3pt}
\setlength{\abovedisplayshortskip}{5pt plus2pt minus2pt}
\setlength{\belowdisplayshortskip}{7pt plus2pt minus2pt}
\setlength{\LTpre}{4pt}
\setlength{\LTpost}{7pt}
\setlength{\tabcolsep}{4pt}
\setlength{\marginparpush}{6pt}
\setlength{\footnotesep}{8pt}
\interfootnotelinepenalty=150
\renewcommand{\thefootnote}{(\arabic{footnote})}
\pagestyle{fancy}
\fancyhf{}
\fancyhead[LE,RO]{\normalfont\fontsize{8}{9.5}\selectfont\parbox[b]{\textwidth}{\nouppercase{\leftmark}}}
\fancyfoot[LE,RO]{\small\thepage}
\renewcommand{\headrulewidth}{0.25pt}
\fancypagestyle{plain}{\fancyhf{}\fancyfoot[LE,RO]{\small\thepage}\renewcommand{\headrulewidth}{0pt}}
\titleformat{\subsection}{\normalfont\bfseries\large}{}{0pt}{}
\titlespacing*{\subsection}{0pt}{2.4ex plus1ex minus.4ex}{.9ex plus.3ex}
\newcommand{\gesperrt}[1]{\textls[100]{#1}}
\newcommand{\weq}[2]{\begin{equation*}\tag*{\ensuremath{#1}}#2\end{equation*}}
\makeatletter
\newif\ifedtable
\AddToHook{env/longtable/before}{\global\edtabletrue}
\AddToHook{env/longtable/after}{\global\edtablefalse\edFlush}
\AddToHook{env/tabular/before}{\global\edtabletrue}
\AddToHook{env/tabular/after}{\global\edtablefalse}
\let\edPending\@empty
\newcommand{\edQueue}[1]{\ifshownotes\g@addto@macro\edPending{#1}\fi}
\newcommand{\edFlush}{\ifshownotes\ifedtable\else\ifmmode\else\ifx\edPending\@empty\else\ifinner\else\let\edTemp\edPending\global\let\edPending\@empty\edTemp\fi\fi\fi\fi\fi}
\newcommand{\edNote}[3]{\marginpar{\normalfont\fontsize{7.1}{8.4}\selectfont\RaggedRight\color{editiongray}\textbf{#1}\, [#2]\quad #3\par}}
\newcommand{\printersignature}[1]{\edQueue{\marginpar{\normalfont\fontsize{6.8}{8}\selectfont\RaggedRight\color{editiongray}Segnatura: #1\par}}}
\newcommand{\weberanchor}[1]{\hypertarget{old:#1}{}\ifcsname wpAlias:#1\endcsname\hypertarget{old:\csname wpAlias:#1\endcsname}{}\fi}
\newcommand{\weberpage}[1]{\ifvmode\leavevmode\fi\weberanchor{#1}\begingroup\reversemarginpar\marginnote{\normalfont\fontsize{8}{9}\selectfont\color{editiongray}\ifcsname wpAlias:#1\endcsname[\csname wpAlias:#1\endcsname]\else#1\fi}\endgroup\ifmeasuring@\else\iffirstchoice@\ifcsname wpNotes:#1\endcsname\expandafter\edQueue\expandafter{\csname wpNotes:#1\endcsname}\fi\ifcsname wpDefer:#1\endcsname\else\edFlush\fi\fi\fi\ignorespaces}
\newcommand{\weberblank}[1]{\weberanchor{#1}}
\newcommand{\origfootnote}[3]{\begingroup\renewcommand\thefootnote{(#1)}\footnote[#1]{\hypertarget{fn:#2}{}\normalfont#3}\endgroup}
\newcommand{\origfootmark}[2]{\textsuperscript{\normalfont\hyperlink{fn:#2}{(#1)}}}
\newcommand{\origfoottext}[3]{\begingroup\renewcommand\thefootnote{(#1)}\footnotetext[#1]{\hypertarget{fn:#2}{}\normalfont#3}\endgroup}
\makeatother
\newcommand{\edchapter}[4]{\pdfbookmark[0]{#3}{ch.#1}\markboth{Capitolo #1. #4}{Capitolo #1. #4}\vspace*{12mm}\begin{center}{\large\scshape Capitolo #1.}\par\vspace{7mm}{\Large\bfseries #2\par}\end{center}\vspace{8mm}}
\newcommand{\edfrontheading}[3]{\pdfbookmark[0]{#2}{#3}\markboth{#2}{#2}\vspace*{8mm}\begin{center}{\Large\bfseries#1}\end{center}\vspace{6mm}}
\newcommand{\edsection}[3]{\subsection*{#1}\pdfbookmark[1]{#2}{#3}}
\newcommand{\editiongap}{\par\bigskip\hypertarget{old:284}{}\hypertarget{old:285}{}\pdfbookmark[1]{Lacuna: pagine 284–285}{gap.284.285}\noindent\fbox{\begin{minipage}{\dimexpr\linewidth-2\fboxsep-2\fboxrule\relax}\small\textbf{Lacuna nella fonte: pagine 284–285.}\par\medskip Le pagine 284 e 285 mancano nella scansione fornita. Il testo di p. 283 si interrompe dopo «cir-»; la fonte riprende a p. 286, con il § 8. Non è stato ricostruito alcun testo.\end{minipage}}\par\bigskip}
\newcommand{\sourcewarning}[1]{\par\medskip\noindent\begin{minipage}{\linewidth}\small\color{editiongray}\textbf{Nota sulla fonte.} #1\end{minipage}\par\medskip}
\newcommand{\editionpage}{\thispagestyle{plain}\pdfbookmark[0]{Nota all’edizione}{front.edizione}\markboth{Nota all’edizione}{Nota all’edizione}\vspace*{7mm}\begin{center}\Large\scshape Nota all’edizione\end{center}\bigskip
\noindent Federigo Enriques, \textit{Le superficie algebriche}, Bologna, Nicola Zanichelli Editore, 1949.\par\medskip
\noindent Scansione: Maxim Leyenson.\par\medskip
\noindent Trascrizione prodotta da IA (OpenAI gpt-6-astra, novembre 2026). Impaginazione e correzioni prodotte da IA (ChatGPT, modalità Pro; altre correzioni: Claude, Anthropic).\par\bigskip
\noindent L'opera originale è di pubblico dominio nell'Unione europea: Federigo Enriques è morto nel 1946 e Guido Castelnuovo, autore della prefazione, nel 1952. Altrove la situazione può essere diversa. Nella misura in cui deteniamo diritti sulla trascrizione, sulle correzioni e sulla composizione tipografica prodotte da IA, tali diritti sono dedicati al pubblico dominio con CC0 1.0. Questa edizione è pubblicata a fini educativi e senza scopo di lucro, e non rivendichiamo diritti che non deteniamo. Se detenete diritti su quest'opera e vi opponete alla sua pubblicazione qui, aprite una issue su GitHub e la rimuoveremo.\par\bigskip
\noindent\textbf{Criteri di lettura.} Il testo conserva la lingua, la notazione e la numerazione della trascrizione, salvo gli errori verificati nella scansione. Le cifre nel margine interno rinviano alla paginazione del 1949; quelle fra parentesi quadre sono ricavate dalla sequenza o dall’indice quando il numero non è stampato o leggibile. I richiami dell’INDICE sono collegati a questi punti. Le note dell’autore conservano la numerazione delle pagine originali.\par\medskip
\noindent Le note numerate nel margine esterno documentano i dubbi ancora aperti e le segnalazioni di possibili refusi; le proposte di correzione matematica non sono applicate al testo. I dubbi risolti restano documentati nei commenti del sorgente e in \texttt{corrections.csv}. Per nascondere tutte le note editoriali basta impostare \texttt{\textbackslash shownotesfalse} nel preambolo.\par\medskip
\noindent\textbf{Limiti della fonte.} Mancano le pagine 284–285. I passi coperti o danneggiati nelle scansioni, in particolare a pp. 186 e 368, restano segnalati senza ricostruzioni ulteriori. Il marchio del frontespizio è ripreso dalla scansione.\par}

\newsavebox{\publisherbox}\sbox{\publisherbox}{\hbox to 23mm{\vrule width0pt height26.25mm\pdfliteral{q 65.19685 0 0 74.40945 0 0 cm BI /W 255 /H 291 /BPC 8 /CS /DeviceGray /F /ASCIIHexDecode ID
FFFFFFFFFFFFFFFFFFFFFFFFFFFFFFFFFFFFFFFFFFFFFFFFFFFFFFFFFFFFFFFFFFFFFFFFFFFFFFFFFFFFFFFFFFFFFFFFFFFFFFFFFFFFFFFFFFFFFFFFFFFFFFFFFFFFFFFFFFFFFFFFFFFFFFFFFFFFFFFFFFFFFFFFFFFFFFFFFFFFFFFFFFFFFFFFFFFFFFFFFFFFFFFFFFFFFFFFFFFFFFFFFFFFFFFFFFFFFFFFFFFFFFFFFFFFFFFFFFFFFFFFFFFFFFFFFFFFFFFFFFFFFFFFFFFFFFFFFFFFFFFFFFFFFFFFFFFFFFFFFFFFFFFFFFFFFFFFFFFFFFFFFFFFFFFFFFFFFFFFFFFFFFFFFFFFFFFFFFFFFFFFFFFFFFFFFFFFFFFFFFFFFFFFFFFFFFFFFFFFFFFFFFFFFFFFFFFFFFFFFFFFFFFFFFFFFFFFFFFFFFFFFFFFFFFFFFFFFFFFFFFFFFFFFFFFFFFFFFFFFFFFFFFFFFFFFFFFFFFFFFFFFFFFFFFFFFFFFFFFFFFFFFFFFFFFFFFFFFFFFFFFFFFFFFFFFFFFFFFFFFFFFFFFFFFFFFFFFFFFFFFFFFFFFFFFFFFFFFFFFFFFFFFFFFFFFFFFFFFFFFFFFFFFFFFFFFFFFFFFFFFFFFFFFFFFFFFFFFFFFFFFFFFFFFFFFFFFFFFFFFFFFFFFFFFFFFFFFFFFFFFFFFFFFFFFFFFFFFFFFFFFFFFFFFFFFFFFFFFFFFFFFFFFFFFFFFFFFFFFFFFFFFFFFFFFFFFFFFFFFFFFFFFFFFFFFFFFFFFFFFFFFFFFFFFFFFFFFFFFFFFFFFFFFFFFFFFFFFFFFFFFFFFFFFFFFFFFFFFFFFFFFFFFFFFFFFFFFFFFFFFFFFFFFFFFFFFFFFFFFFFFFFFFFFFFFFFFFFFFFFFFFFFFFFFFFFFFFFFFFFFFFFFFFFFFFFFFFFFFFFFFFFFFFFFFFFFFFFFFFFFFFFFFFFFFFFFFFFFFFFFFFFFFFFFFFFFFFFFFFFFFFFFFFFFFFFFFFFFFFFFFFFFFFFFFFFFFFFFFFFFFFFFFFFFFFFFFFFFFFFFFFFFFFFFFFFFFFFFFFFFFFFFFFFFFFFFFFFFFFFFFFFFFFFFFFFFFFFFFFFFFFFFFFFFFFFFFFFFFFFFFFFFFFFFFFFFFFFFFFFFFFFFFFFFFFFFFFFFFFFFFFFFFFFFFFFFFFFFFFFFFFFFFFFFFFFFFFFFFFFFFFFFFFFFFFFFFFFFFFFFFFFFFFFFFFFFFFFFFFFFFFFFFFFFFFFFFFFFFFFFFFFFFFFFFFFFFFFFFFFFFFFFFFFFFFFFFFFFFFFFFFFFFFFFFFFFFFFFFFFFFFFFFFFFFFFFFFFFFFFFFFFFFFFFFFFFFFFFFFFFFFFFFFFFFFFFFFFFFFFFFFFFFFFFFFFFFFFFFFFFFFFFFFFFFFFFFFFFFFFFFFFFFFFFFFFFFFFFFFFFFFFFFFFFFFFFFFFFFFFFFFFFFFFFFFFFFFFFFFFFFFFFFFFFFFFFFFFFFFFFFFFFFFFFFFFFFFFFFFFFFFFFFFFFFFFFFFFFFFFFFFFFFFFFFFFFFFFFFFFFFFFFFFFFFFFFFFFFFFFFFFFFFFFFFFFFFFFFFFFFFFFFFFFFFFFFFFFFFFFFFFFFFFFFFFFFFFFFFFFFFFFFFFFFFFFFFFFFFFFFFFFFFFFFFFFFFFFFFFFFFFFFFFFFFFFFFFFFFFFFFFFFFFFFFFFFFFFFFFFFFFFFFFFFFFFFFFFFFFFFFFFFFFFFFFFFFFFFFFFFFFFFFFFFFFFFFFFFFFFFFFFFFFFFFFFFFFFFFFFFFFFFFFFFFFFFFFFFFFFFFFFFFFFFFFFFFFFFFFFFFFFFFFFFFFFFFFFFFFFFFFFFFFFFFFFFFFFFFFFFFFFFFFFFFFFFFFFFFFFFFFFFFFFFFFFFFFFFFFFFFFFFFFFFFFFFFFFFFFFFFFFFFFFFFFFFFFFFFFFFFFFFFFFFFFFFFFFFFFFFFFFFFFFFFFFFFFFFFFFFFFFFFFFFFFFFFFFFFFFFFFFFFFFFFFFFFFFFFFFFFFFFFFFFFFFFFFFFFFFFFFFFFFFFFFFFFFFFFFFFFFFFFFFFFFFFFFFFFFFFFFFFFFFFFFFFFFFFFFFFFFFFFFFFFFFFFFFFFFFFFFFFFFFFFFFFFFFFFFFFFFFFFFFFFFFFFFFFFFFFFFFFFFFFFFFFFFFFFFFFFFFFFFFFFFFFFFFFFFFFFFFFFFFFFFFFFFFFFFFFFFFFFFFFFFFFFFFFFFFFFFFFFFFFFFFFFFFFFFFFFFFFFFFFFFFFFFFFFFFFFFFFFFFFFFFFFFFFFFFFFFFFFFFFFFFFFFFFFFFFFFFFFFFFFFFFFFFFFFFFFFFFFFFFFFFFFFFFFFFFFFFFFFFFFFFFFFFFFFFFFFFFFFFFFFFFFFFFFFFFFFFFFFFFFFFFFFFFFFFFFFFFFFFFFFFFFFFFFFFFFFFFFFFFFFFFFFFFFFFFFFFFFFFFFFFFFFFFFFFFFFFFFFFFFFFFFFFFFFFFFFFFFFFFFFFFFFFFFFFFFFFFFFFFFFFFFFFFFFFFFFFFFFFFFFFFFFFFFFFFFFFFFFFFFFFFFFFFFFFFFFFFFFFFFFFFFFFFFFFFFFFFFFFFFFFFFFFFFFFFFFFFFFFFFFFFFFFFFFFFFFFFFFFFFFFFFFFFFFFFFFFFFFFFFFFFFFFFFFFFFFFFFFFFFFFFFFFFFFFFFFFFFFFFFFFFFFFFFFFFFFFFFFFFFFFFFFFFFFFFFFFFFFFFFFFFFFFFFFFFFFFFFFFFFFFFFFFFFFFFFFFFFFFFFFFFFFFFFFFFFFFFFFFFFFFFFFFFFFFFFFFFFFFFFFFFFFFFFFFFFFFFFFFFFFFFFFFFFFFFFFFFFFFFFFFFFFFFFFFFFFFFFFFFFFFFFFFFFFFFFFFFFFFFFFFFFFFFFFFFFFFFFFFFFFFFFFFFFFFFFFFFFFFFFFFFFFFFFFFFFFFFFFFFFFFFFFFFFFFFFFFFFFFFFFFFFFFFFFFFFFFFFFFFFFFFFFFFFFFFFFFFFFFFFFFFFFFFFFFFFFFFFFFFFFFFFFFFFFFFFFFFFFFFFFFFFFFFFFFFFFFFFFFFFFFFFFFFFFFFFFFFFFFFFFFFFFFFFFFFFFFFFFFFFFFFFFFFFFFFFFFFFFFFFFFFFFFFFFFFFFFFFFFFFFFFFFFFFFFFFFFFFFFFFFFFFFFFFFFFFFFFFFFFFFFFFFFFFFFFFFFFFFFFFFFFFFFFFFFFFFFFFFFFFFFFFFFFFFFFFFFFFFFFFFFFFFFFFFFFFFFFFFFFFFFFFFFFFFFFFFFFFFFFFFFFFFFFFFFFFFFFFFFFFFFFFFFFFFFFFFFFFFFFFFFFFFFFFFFFFFFFFFFFFFFFFFFFFFFFFFFFFFFFFFFFFFFFFFFFFFFFFFFFFFFFFFFFFFFFFFFFFFFFFFFFFFFFFFFFFFFFFFFFFFFFFFFFFFFFFFFFFFFFFFFFFFFFFFFFFFFFFFFFFFFFFFFFFFFFFFFFFFFFFFFFFFFFFFFFFFFFFFFFFFFFFFFFFFFFFFFFFFFFFFFFFFFFFFFFFFFFFFFFFFFFFFFFFFFFFFFFFFFFFFFFFFFFFFFFFFFFFFFFFFFFFFFFFFFFFFFFFFFFFFFFFFFFFFFFFFFFFFFFFFFFFFFFFFFFFFFFFFFFFFFFFFFFFFFFFFFFFFFFFFFFFFFFFFFFFFFFFFFFFFFFFFFFFFFFFFFFFFFFFFFFFFFFFFFFFFFFFFFFFFFFFFFFFFFFFFFFFFFFFFFFFFFFFFFFFFFFFFFFFFFFFFFFFFFFFFFFFFFFFFFFFFFFFFFFFFFFFFFFFFFFFFFFFFFFFFFFFFFFFFFFFFFFFFFFFFFFFFFFFFFFFFFFFFFFFFFFFFFFFFFFFFFFFFFFFFFFFFFFFFFFFFFFFFFFFFFFFFFFFFFFFFFFFFFFFFFFFFFFFFFFFFFFFFFFFFFFFFFFFFFFFFFFFFFFFFFFFFFFFFFFFFFFFFFFFFFFFFFFFFFFFFFFFFFFFFFFFFFFFFFFFFFFFFFFFFFFFFFFFFFFFFFFFFFFFFFFFFFFFFFFFFFFFFFFFFFFFFFFFFFFFFFFFFFFFFFFFFFFFFFFFFFFFFFFFFFFFFFFFFFFFFFFFFFFFFFFFFFFFFFFFFFFFFFFFFFFFFFFFFFFFFFFFFFFFFFFFFFFFFFFFFFFFFFFFFFFFFFFFFFFFFFFFFFFFFFFFFFFFFFFFFFFFFFFFFFFFFFFFFFFFFFFFFFFFFFFFFFFFFFFFFFFFFFFFFFFFFFFFFFFFFFFFFFFFFFFFFFFFFFFFFFFFFFFFFFFFFFFFFFFFFFFFFFFFFFFFFFFFFFFFFFFFFFFFFFFFFFFFFFFFFFFFFFFFFFFFFFFFFFFFFFFFFFFFFFFFFFFFFFFFFFFFFFFFFFFFFFFFFFFFFFFFFFFFFFFFFFFFFFFFFFFFFFFFFFFFFFFFFFFFFFFFFFFFFFFFFFFFFFFFFFFFFFFFFFFFFFFFFFFFFFFFFFFFFFFFFFFFFFFFFFFFFFFFFFFFFFFFFFFFFFFFFFFFFFFFFFFFFFFFFFFFFFFFFFFFFFFFFFFFFFFFFFFFFFFFFFFFFFFFFFFFFFFFFFFFFFFFFFFFFFFFFFFFFFFFFFFFFFFFFFFFFFFFFFFFFFFFFFFFFFFFFFFFFFFFFFFFFFFFFFFFFFFFFFFFFFFFFFFFFFFFFFFFFDF2FEF9F3FBFFFFFEFEFEFFFEFCFEF8FEFCFEF9FBFEFDFDFFFEFEFFFEFBFDFEFFFFFFFFFFFFFFFFFFFFFFFFFFFFFFFFFFFFFFFFFFFFFFFFFFFFFFFFFFFFFFFFFFFFFFFFFFFFFFFFFFFFFFFFFFFFFFFFFFFFFFFFFFFFFFFFFFFFFFFFFFFFFFFFFFFFFFFFFFFFFFFFFFFFFFFFFFFFFFFFFFFFFFFFFFFFFFFFFFFFFFFFFFFFFFFFFFFFFFFFFFFFFFFFFFFFFFFFFFFFFFFFFFFFFFFFFFFFFFFFFFFFFFFFFFFFFFFFFFFFFFFFFFFFFFFFFFFFFFFFFFFFFFFFFFFFFFFFFFFFFFFFFFFFFFFFFFFFFFFFFFFFFFFFFFFFFFFFFFFFFFFFFFFFFFFFFFFFFFFFFFFFFFFFFFFFFFFFFFFFFFFFFFFFFFFFFFFFFFFFFFFFFFFFFFFFFFFFFFFFFFFFFFFFFFFFFFFFFFFFFFFFFFFDEFF2F6F1F6F8F9FBFDFBFCFBF1FDF6F5FAFAF6F9FDFDFCFDFAF9FCFEFAFCFBFBFDFFFFFFFFFFFFFFFFFFFFFFFFFFFFFFFFFFFFFFFFFFFFFFFFFFFFFFFFFFFFFFFFFFFFFFFFFFFFFFFFFFFFFFFFFFFFFFFFFFFFFFFFFFFFFFFFFFFFFFFFFFFFFFFFFFFFFFFFFFFFFFFFFFFFFFFFFFFFFFFFFFFFFFFFFFFFFFFFFFFFFFFFFFFFFFFFFFFFFFFFFFFFFFFFFFFFFFFFFFFFFFFFFFFFFFFFFFFFFFFFFFFFFFFFFFFFFFFFFFFFFFFFFFFFFFFFFFFFFFFFFFFFFFFFFFFFFFFFFFFFFFFFFFFFFFFFFFFFFFFFFFFFFFFFFFFFFFFFFFFFFFFFFFFFFFFFFFFFFFFFFFFFFFFFFFFFFFFFFFFFFFFFFFFFFFFFFFFFFFFFFFFFFFFFFFFFFFFFFFFFFFFFFFFFFFFFFFFFFFFFFEFDFBECF7FCF6F4F5F9F9F9FAF9EFFCF8F1FCF9FCFCFDFBFAFBF8F6F9FEFEFEFAF9FBFDFFFFFFFFFFFFFFFFFFFFFFFFFFFFFFFFFFFFFFFFFFFFFFFFFFFFFFFFFFFFFFFFFFFFFFFFFFFFFFFFFFFFFFFFFFFFFFFFFFFFFFFFFFFFFFFFFFFFFFFFFFFFFFFFFFFFFFFFFFFFFFFFFFFFFFFFFFFFFFFFFFFFFFFFFFFFFFFFFFFFFFFFFFFFFFFFFFFFFFFFFFFFFFFFFFFFFFFFFFFFFFFFFFFFFFFFFFFFFFFFFFFFFFFFFFFFFFFFFFFFFFFFFFFFFFFFFFFFFFFFFFFFFFFFFFFFFFFFFFFFFFFFFFFFFFFFFFFFFFFFFFFFFFFFFFFFFFFFFFFFFFFFFFFFFFFFFFFFFFFFFFFFFFFFFFFFFFFFFFFFFFFFFFFFFFFFFFFFFFFFFFFFFFFFFFFFFFFFFFFFFFFFFFFFFFFFFFFFFFFEFAF8FDF4F9F9F6FAFBF6F8FBF6FEFEF1FAFEFDFBFCFBF6F5FCFEFEFBFDFCFAFDFFFAFBFFFFFFFFFFFFFFFFFFFFFFFFFFFFFFFFFFFFFFFFFFFFFFFFFFFFFFFFFFFFFFFFFFFFFFFFFFFFFFFFFFFFFFFFFFFFFFFFFFFFFFFFFFFFFFFFFFFFFFFFFFFFFFFFFFFFFFFFFFFFFFFFFFFFFFFFFFFFFFFFFFFFFFFFFFFFFFFFFFFFFFFFFFFFFFFFFFFFFFFFFFFFFFFFFFFFFFFFFFFFFFFFFFFFFFFFFFFFFFFFFFFFFFFFFFFFFFFFFFFFFFFFFFFFFFFFFFFFFFFFFFFFFFFFFFFFFFFFFFFFFFFFFFFFFFFFFFFFFFFFFFFFFFFFFFFFFFFFFFFFFFFFFFFFFFFFFFFFFFFFFFFFFFFFFFFFFFFFFFFFFFFFFFFFFFFFFFFFFFFFFFFFFFFFFFFFFFFFFFFFFFFFFFFFFFFFFFFFFFFEFAFEF9FDFBFCECF2FEF8F5FAF5FBFEFEF4FAF6F6FAF7F9FDFBFCFBFDFEF8F3FCFEFAFBFFFFFFFFFFFFFFFFFFFFFFFFFFFFFFFFFFFFFFFFFFFFFFFFFFFFFFFFFFFFFFFFFFFFFFFFFFFFFFFFFFFFFFFFFFFFFFFFFFFFFFFFFFFFFFFFFFFFFFFFFFFFFFFFFFFFFFFFFFFFFFFFFFFFFFFFFFFFFFFFFFFFFFFFFFFFFFFFFFFFFFFFFFFFFFFFFFFFFFFFFFFFFFFFFFFFFFFFFFFFFFFFFFFFFFFFFFFFFFFFFFFFFFFFFFFFFFFFFFFFFFFFFFFFFFFFFFFFFFFFFFFFFFFFFFFFFFFFFFFFFFFFFFFFFFFFFFFFFFFFFFFFFFFFFFFFFFFFFFFFFFFFFFFFFFFFFFFFFFFFFFFFFFFFFFFFFFFFFFFFFFFFFFFFFFFFFFFFFFFFFFFFFFFFFFFFFFFFFFFFFFFFFFFFFFFFFFFFFFFEFEFDF2FCFAF7FAFBFBFCFAEFE5FAFCFAFEF4DEF0FEF2F2F7F0F7F1FBFEFEFAFDFDFEFFFFFFFFFFFFFFFFFFFFFFFFFFFFFFFFFFFFFFFFFFFFFFFFFFFFFFFFFFFFFFFFFFFFFFFFFFFFFFFFFFFFFFFFFFFFFFFFFFFFFFFFFFFFFFFFFFFFFFFFFFFFFFFFFFFFFFFFFFFFFFFFFFFFFFFFFFFFFFFFFFFFFFFFFFFFFFFFFFFFFFFFFFFFFFFFFFFFFFFFFFFFFFFFFFFFFFFFFFFFFFFFFFFFFFFFFFFFFFFFFFFFFFFFFFFFFFFFFFFFFFFFFFFFFFFFFFFFFFFFFFFFFFFFFFFFFFFFFFFFFFFFFFFFFFFFFFFFFFFFFFFFFFFFFFFFFFFFFFFFFFFFFFFFFFFFFFFFFFFFFFFFFFFFFFFFFFFFFFFFFFFFFFFFFFFFFFFFFFFFFFFFFFFFFFFFFFFFFFFFFFFFFFFFFFFFFFFFFFFFFDF9FDFDFEF9F7EBF1FBF5F4BF3FF3F1D5EFFC65A8F2FEFDFDFEFEFEFEF7FAFDFEF6FAFFFFFFFFFFFFFFFFFFFFFFFFFFFFFFFFFFFFFFFFFFFFFFFFFFFFFFFFFFFFFFFFFFFFFFFFFFFFFFFFFFFFFFFFFFFFFFFFFFFFFFFFFFFFFFFFFFFFFFFFFFFFFFFFFFFFFFFFFFFFFFFFFFFFFFFFFFFFFFFFFFFFFFFFFFFFFFFFFFFFFFFFFFFFFFFFFFFFFFFFFFFFFFFFFFFFFFFFFFFFFFFFFFFFFFFFFFFFFFFFFFFFFFFFFFFFFFFFFFFFFFFFFFFFFFFFFFFFFFFFFFFFFFFFFFFFFFFFFFFFFFFFFFFFFFFFFFFFFFFFFFFFFFFFFFFFFFFFFFFFFFFFFFFFFFFFFFFFFFFFFFFFFFFFFFFFFFFFFFFFFFFFFFFFFFFFFFFFFFFFFFFFFFFFFFFFFFFFFFFFFFFFFFFFFFFFFFFFFFFFFCEEFDF1F6FCFE8189E0A971E5BA1E1A0516C5E85896F7F9F7FAEDEBFEF7EDFCFEF6FAFFFFFFFFFFFFFFFFFFFFFFFFFFFFFFFFFFFFFFFFFFFFFFFFFFFFFFFFFFFFFFFFFFFFFFFFFFFFFFFFFFFFFFFFFFFFFFFFFFFFFFFFFFFFFFFFFFFFFFFFFFFFFFFFFFFFFFFFFFFFFFFFFFFFFFFFFFFFFFFFFFFFFFFFFFFFFFFFFFFFFFFFFFFFFFFFFFFFFFFFFFFFFFFFFFFFFFFFFFFFFFFFFFFFFFFFFFFFFFFFFFFFFFFFFFFFFFFFFFFFFFFFFFFFFFFFFFFFFFFFFFFFFFFFFFFFFFFFFFFFFFFFFFFFFFFFFFFFFFFFFFFFFFFFFFFFFFFFFFFFFFFFFFFFFFFFFFFFFFFFFFFFFFFFFFFFFFFFFFFFFFFFFFFFFFFFFFFFFFFFFFFFFFFFFFFFFFFFFFFFFFFFFFFFFFFFFFFFFFFFFEFBEDFB719CE6ED872F140A1B21020213031B19131D1E308FE79F74EAFDF3F8FDF8FBFFFFFFFFFFFFFFFFFFFFFFFFFFFFFFFFFFFFFFFFFFFFFFFFFFFFFFFFFFFFFFFFFFFFFFFFFFFFFFFFFFFFFFFFFFFFFFFFFFFFFFFFFFFFFFFFFFFFFFFFFFFFFFFFFFFFFFFFFFFFFFFFFFFFFFFFFFFFFFFFFFFFFFFFFFFFFFFFFFFFFFFFFFFFFFFFFFFFFFFFFFFFFFFFFFFFFFFFFFFFFFFFFFFFFFFFFFFFFFFFFFFFFFFFFFFFFFFFFFFFFFFFFFFFFFFFFFFFFFFFFFFFFFFFFFFFFFFFFFFFFFFFFFFFFFFFFFFFFFFFFFFFFFFFFFFFFFFFFFFFFFFFFFFFFFFFFFFFFFFFFFFFFFFFFFFFFFFFFFFFFFFFFFFFFFFFFFFFFFFFFFFFFFFFFFFFFFFFFFFFFFFFFCFDFFFFFFFEFEFFFFFDF6FDD57A3C2D1904000002020002030B0300020201041A4065937EA77F96988A95CEF8FCFBFFFDF8FFFFFFF6FBFCFDFFFAF9FDFFFFFFFFFFFFFFFFFFFFFFFFFFFFFFFFFFFFFFFFFFFFFFFFFFFFFFFFFFFFFFFFFFFFFFFFFFFFFFFFFFFFFFFFFFFFFFFFFFFFFFFFFFFFFFFFFFFFFFFFFFFFFFFFFFFFFFFFFFFFFFFFFFFFFFFFFFFFFFFFFFFFFFFFFFFFFFFFFFFFFFFFFFFFFFFFFFFFFFFFFFFFFFFFFFFFFFFFFFFFFFFFFFFFFFFFFFFFFFFFFFFFFFFFFFFFFFFFFFFFFFFFFFFFFFFFFFFFFFFFFFFFFFFFFFFFFFFFFFFFFFFFFFFFFFFFFFFFFFFFFFFFFFFFFFFFFFFFFFFFFFFFFFFFFFFFFFFFFFFFFFFFFFFFFFFFFFFFFFFFFFFFFFFFFCFDFFFFFEFCFCFDFDF7FCF789965A1706060001020909061204000E060002000108172C1A251A2670697AC9F5B886CDD56EC4EEE9F8F7FBFCFCFDFCFDFFFFFFFFFFFFFFFFFFFFFFFFFFFFFFFFFFFFFFFFFFFFFFFFFFFFFFFFFFFFFFFFFFFFFFFFFFFFFFFFFFFFFFFFFFFFFFFFFFFFFFFFFFFFFFFFFFFFFFFFFFFFFFFFFFFFFFFFFFFFFFFFFFFFFFFFFFFFFFFFFFFFFFFFFFFFFFFFFFFFFFFFFFFFFFFFFFFFFFFFFFFFFFFFFFFFFFFFFFFFFFFFFFFFFFFFFFFFFFFFFFFFFFFFFFFFFFFFFFFFFFFFFFFFFFFFFFFFFFFFFFFFFFFFFFFFFFFFFFFFFFFFFFFFFFFFFFFFFFFFFFFFFFFFFFFFFFFFFFFFFFFFFFFFFFFFFFFFFFFFFFFFFFFFFFFFFFFFFFFFFFFFFFFFFEFEFFFEFDFBFBFDFDF9F1FA6D40230D04080E0A2C8A1CAEA519178D3A0202000003080300070C03222B748B9C8C6AA0581E425351D2F8FAF5F8FFFEFEFFFFFFFFFFFFFFFFFFFFFFFFFFFFFFFFFFFFFFFFFFFFFFFFFFFFFFFFFFFFFFFFFFFFFFFFFFFFFFFFFFFFFFFFFFFFFFFFFFFFFFFFFFFFFFFFFFFFFFFFFFFFFFFFFFFFFFFFFFFFFFFFFFFFFFFFFFFFFFFFFFFFFFFFFFFFFFFFFFFFFFFFFFFFFFFFFFFFFFFFFFFFFFFFFFFFFFFFFFFFFFFFFFFFFFFFFFFFFFFFFFFFFFFFFFFFFFFFFFFFFFFFFFFFFFFFFFFFFFFFFFFFFFFFFFFFFFFFFFFFFFFFFFFFFFFFFFFFFFFFFFFFFFFFFFFFFFFFFFFFFFFFFFFFFFFFFFFFFFFFFFFFFFFFFFFFFFFFFFFFFFFFFFFFFFFFFFFEFDFCFCFEEFA3B2A1340D02080541611F46C737F267989FD454060808001045640004395C1506637751352E3B1405050504AA8DB9ECFBFDFDFFFFFFFFFFFFFFFFFFFFFFFFFFFFFFFFFFFFFFFFFFFFFFFFFFFFFFFFFFFFFFFFFFFFFFFFFFFFFFFFFFFFFFFFFFFFFFFFFFFFFFFFFFFFFFFFFFFFFFFFFFFFFFFFFFFFFFFFFFFFFFFFFFFFFFFFFFFFFFFFFFFFFFFFFFFFFFFFFFFFFFFFFFFFFFFFFFFFFFFFFFFFFFFFFFFFFFFFFFFFFFFFFFFFFFFFFFFFFFFFFFFFFFFFFFFFFFFFFFFFFFFFFFFFFFFFFFFFFFFFFFFFFFFFFFFFFFFFFFFFFFFFFFFFFFFFFFFFFFFFFFFFFFFFFFFFFFFFFFFFFFFFFFFFFFFFFFFFFFFFFFFFFFFFFFFFFFFFFFFFFFFFFFFFFFFFFFFEFEFEFEFEFFF3B488340805060908327187A8E2A2F58BDBF3897E370608030C354D0200274712002D3D1E05020106050007054A2145B0F0F9FBFFFFFFFFFFFFFFFFFFFFFFFFFFFFFFFFFFFFFFFFFFFFFFFFFFFFFFFFFFFFFFFFFFFFFFFFFFFFFFFFFFFFFFFFFFFFFFFFFFFFFFFFFFFFFFFFFFFFFFFFFFFFFFFFFFFFFFFFFFFFFFFFFFFFFFFFFFFFFFFFFFFFFFFFFFFFFFFFFFFFFFFFFFFFFFFFFFFFFFFFFFFFFFFFFFFFFFFFFFFFFFFFFFFFFFFFFFFFFFFFFFFFFFFFFFFFFFFFFFFFFFFFFFFFFFFFFFFFFFFFFFFFFFFFFFFFFFFFFFFFFFFFFFFFFFFFFFFFFFFFFFFFFFFFFFFFFFFFFFFFFFFFFFFFFFFFFFFFFFFFFFFFFFFFFFFFFFFFFFFFFFFFFFFFFFFFFEFEFEFFFFFFFEFDEE7A080301453D1405329CF5FBF6F9F587D63A5935040F11050300070004010403040203040204010500081307010348BAFAFBFCFFFFFFFFFFFFFFFFFFFFFFFFFFFFFFFFFFFFFFFFFFFFFFFFFFFFFFFFFFFFFFFFFFFFFFFFFFFFFFFFFFFFFFFFFFFFFFFFFFFFFFFFFFFFFFFFFFFFFFFFFFFFFFFFFFFFFFFFFFFFFFFFFFFFFFFFFFFFFFFFFFFFFFFFFFFFFFFFFFFFFFFFFFFFFFFFFFFFFFFFFFFFFFFFFFFFFFFFFFFFFFFFFFFFFFFFFFFFFFFFFFFFFFFFFFFFFFFFFFFFFFFFFFFFFFFFFFFFFFFFFFFFFFFFFFFFFFFFFFFFFFFFFFFFFFFFFFFFFFFFFFFFFFFFFFFFFFFFFFFFFFFFFFFFFFFFFFFFFFFFFFFFFFFFFFFFFFFFFFFFFFFFFFFFFEFEFEFEFFFFFEFBF9E9640A0601584C18072C8DEBEDEAE099667B23030411393D07040006000700040E08040202060600000301581F0401057FF1F8FAFFFFFFFFFFFFFFFFFFFFFFFFFFFFFFFFFFFFFFFFFFFFFFFFFFFFFFFFFFFFFFFFFFFFFFFFFFFFFFFFFFFFFFFFFFFFFFFFFFFFFFFFFFFFFFFFFFFFFFFFFFFFFFFFFFFFFFFFFFFFFFFFFFFFFFFFFFFFFFFFFFFFFFFFFFFFFFFFFFFFFFFFFFFFFFFFFFFFFFFFFFFFFFFFFFFFFFFFFFFFFFFFFFFFFFFFFFFFFFFFFFFFFFFFFFFFFFFFFFFFFFFFFFFFFFFFFFFFFFFFFFFFFFFFFFFFFFFFFFFFFFFFFFFFFFFFFFFFFFFFFFFFFFFFFFFFFFFFFFFFFFFFFFFFFFFFFFFFFFFFFFFFFFFFFFFFFFFFFFFFFFFFFFFFFEFEFEFEFFFEFBF5F0DC560D060201010F0F5DCDB8C6A58505891E05101725565A0E0802000516131436382418170F05110007027E280401024DB4EAFCFFFFFFFFFFFFFFFFFFFFFFFFFFFFFFFFFFFFFFFFFFFFFFFFFFFFFFFFFFFFFFFFFFFFFFFFFFFFFFFFFFFFFFFFFFFFFFFFFFFFFFFFFFFFFFFFFFFFFFFFFFFFFFFFFFFFFFFFFFFFFFFFFFFFFFFFFFFFFFFFFFFFFFFFFFFFFFFFFFFFFFFFFFFFFFFFFFFFFFFFFFFFFFFFFFFFFFFFFFFFFFFFFFFFFFFFFFFFFFFFFFFFFFFFFFFFFFFFFFFFFFFFFFFFFFFFFFFFFFFFFFFFFFFFFFFFFFFFFFFFFFFFFFFFFFFFFFFFFFFFFFFFFFFFFFFFFFFFFFFFFFFFFFFFFFFFFFFFFFFFFFFFFFFFFFFFFFFFFFFFFFFFFFFFFFFCFDFEFEFDF8ECCF8B4B07000303052D30497E688A5203030F0C124C5E5E6159370406021C6060627E8C716769442B4F16060201000B08010853C7FFFFFFFFFFFFFFFFFFFFFFFFFFFFFFFFFFFFFFFFFFFFFFFFFFFFFFFFFFFFFFFFFFFFFFFFFFFFFFFFFFFFFFFFFFFFFFFFFFFFFFFFFFFFFFFFFFFFFFFFFFFFFFFFFFFFFFFFFFFFFFFFFFFFFFFFFFFFFFFFFFFFFFFFFFFFFFFFFFFFFFFFFFFFFFFFFFFFFFFFFFFFFFFFFFFFFFFFFFFFFFFFFFFFFFFFFFFFFFFFFFFFFFFFFFFFFFFFFFFFFFFFFFFFFFFFFFFFFFFFFFFFFFFFFFFFFFFFFFFFFFFFFFFFFFFFFFFFFFFFFFFFFFFFFFFFFFFFFFFFFFFFFFFFFFFFFFFFFFFFFFFFFFFFFFFFFFFFFFFFFFFFFFFFFFFFF9FAFEFDFCF7DE9722000900000607545F194FAC26980A005320549BB3D3EEC45E0A010744A8AFB08980A1CBE5C98C946C12070006412303305C99FFFFFFFFFFFFFFFFFFFFFFFFFFFFFFFFFFFFFFFFFFFFFFFFFFFFFFFFFFFFFFFFFFFFFFFFFFFFFFFFFFFFFFFFFFFFFFFFFFFFFFFFFFFFFFFFFFFFFFFFFFFFFFFFFFFFFFFFFFFFFFFFFFFFFFFFFFFFFFFFFFFFFFFFFFFFFFFFFFFFFFFFFFFFFFFFFFFFFFFFFFFFFFFFFFFFFFFFFFFFFFFFFFFFFFFFFFFFFFFFFFFFFFFFFFFFFFFFFFFFFFFFFFFFFFFFFFFFFFFFFFFFFFFFFFFFFFFFFFFFFFFFFFFFFFFFFFFFFFFFFFFFFFFFFFFFFFFFFFFFFFFFFFFFFFFFFFFFFFFFFFFFFFFFFFFFFFFFFFFFFFFFFFFFFFFFFEFAFAFDF9E7D4CE79000000020A04000001387715000100A886933305619B5101010A00000B02004B4F1B3081C4CA6CA61A0200248C380364D5F9FEFFFFFFFFFFFFFFFFFFFFFFFFFFFFFFFFFFFFFFFFFFFFFFFFFFFFFFFFFFFFFFFFFFFFFFFFFFFFFFFFFFFFFFFFFFFFFFFFFFFFFFFFFFFFFFFFFFFFFFFFFFFFFFFFFFFFFFFFFFFFFFFFFFFFFFFFFFFFFFFFFFFFFFFFFFFFFFFFFFFFFFFFFFFFFFFFFFFFFFFFFFFFFFFFFFFFFFFFFFFFFFFFFFFFFFFFFFFFFFFFFFFFFFFFFFFFFFFFFFFFFFFFFFFFFFFFFFFFFFFFFFFFFFFFFFFFFFFFFFFFFFFFFFFFFFFFFFFFFFFFFFFFFFFFFFFFFFFFFFFFFFFFFFFFFFFFFFFFFFFFFFFFFFFFFFFFFFFFFFFFFFFFFFFFFFFAFCFEFDF7DDA96E5A0207010000010000090F0001081201008F2B080A04020C1101010802000A08080702050B2B829513001206000000090870FBFCFFFFFFFFFFFFFFFFFFFFFFFFFFFFFFFFFFFFFFFFFFFFFFFFFFFFFFFFFFFFFFFFFFFFFFFFFFFFFFFFFFFFFFFFFFFFFFFFFFFFFFFFFFFFFFFFFFFFFFFFFFFFFFFFFFFFFFFFFFFFFFFFFFFFFFFFFFFFFFFFFFFFFFFFFFFFFFFFFFFFFFFFFFFFFFFFFFFFFFFFFFFFFFFFFFFFFFFFFFFFFFFFFFFFFFFFFFFFFFFFFFFFFFFFFFFFFFFFFFFFFFFFFFFFFFFFFFFFFFFFFFFFFFFFFFFFFFFFFFFFFFFFFFFFFFFFFFFFFFFFFFFFFFFFFFFFFFFFFFFFFFFFFFFFFFFFFFFFFFFFFFFFFFFFFFFFFFFFFFFFFFFFFFFFFFFAFBFAF5F7F0B74805000003412F120608080101174F04670A09120701082C33061E3F1C533127320F0302041B2C150B050001000008030030A4FAF5FFFFFFFFFFFFFFFFFFFFFFFFFFFFFFFFFFFFFFFFFFFFFFFFFFFFFFFFFFFFFFFFFFFFFFFFFFFFFFFFFFFFFFFFFFFFFFFFFFFFFFFFFFFFFFFFFFFFFFFFFFFFFFFFFFFFFFFFFFFFFFFFFFFFFFFFFFFFFFFFFFFFFFFFFFFFFFFFFFFFFFFFFFFFFFFFFFFFFFFFFFFFFFFFFFFFFFFFFFFFFFFFFFFFFFFFFFFFFFFFFFFFFFFFFFFFFFFFFFFFFFFFFFFFFFFFFFFFFFFFFFFFFFFFFFFFFFFFFFFFFFFFFFFFFFFFFFFFFFFFFFFFFFFFFFFFFFFFFFFFFFFFFFFFFFFFFFFFFFFFFFFFFFFFFFFFFFFFFFFFFFFFFFFFF2F9FEFEFDF7CC650F000B0000594914000128450F2293069610000202010E5B8B5E717981E37481672103040434662A01060B01020007030078EDFAFDFFFFFFFFFFFFFFFFFFFFFFFFFFFFFFFFFFFFFFFFFFFFFFFFFFFFFFFFFFFFFFFFFFFFFFFFFFFFFFFFFFFFFFFFFFFFFFFFFFFFFFFFFFFFFFFFFFFFFFFFFFFFFFFFFFFFFFFFFFFFFFFFFFFFFFFFFFFFFFFFFFFFFFFFFFFFFFFFFFFFFFFFFFFFFFFFFFFFFFFFFFFFFFFFFFFFFFFFFFFFFFFFFFFFFFFFFFFFFFFFFFFFFFFFFFFFFFFFFFFFFFFFFFFFFFFFFFFFFFFFFFFFFFFFFFFFFFFFFFFFFFFFFFFFFFFFFFFFFFFFFFFFFFFFFFFFFFFFFFFFFFFFFFFFFFFFFFFFFFFFFFFFFFFFFFFFFFFFFFFFFFFFFFFFFCFDFBF3F6FDE29130010905001F140405043C650807241B211102000604071B5EB8BA97E0AB4368473C150107131709000205005255163756AFFDFBFBFFFFFFFFFFFFFFFFFFFFFFFFFFFFFFFFFFFFFFFFFFFFFFFFFFFFFFFFFFFFFFFFFFFFFFFFFFFFFFFFFFFFFFFFFFFFFFFFFFFFFFFFFFFFFFFFFFFFFFFFFFFFFFFFFFFFFFFFFFFFFFFFFFFFFFFFFFFFFFFFFFFFFFFFFFFFFFFFFFFFFFFFFFFFFFFFFFFFFFFFFFFFFFFFFFFFFFFFFFFFFFFFFFFFFFFFFFFFFFFFFFFFFFFFFFFFFFFFFFFFFFFFFFFFFFFFFFFFFFFFFFFFFFFFFFFFFFFFFFFFFFFFFFFFFFFFFFFFFFFFFFFFFFFFFFFFFFFFFFFFFFFFFFFFFFFFFFFFFFFFFFFFFFFFFFFFFFFFFFFFFFFFFFFFFFFCFCF8F6F9ECA32A010201020401000101121E050F0B090B050400020706051C62BBC39272161A4746200C0805051139123E14B98B66ACD5E9F4F1FCFFFFFFFFFFFFFFFFFFFFFFFFFFFFFFFFFFFFFFFFFFFFFFFFFFFFFFFFFFFFFFFFFFFFFFFFFFFFFFFFFFFFFFFFFFFFFFFFFFFFFFFFFFFFFFFFFFFFFFFFFFFFFFFFFFFFFFFFFFFFFFFFFFFFFFFFFFFFFFFFFFFFFFFFFFFFFFFFFFFFFFFFFFFFFFFFFFFFFFFFFFFFFFFFFFFFFFFFFFFFFFFFFFFFFFFFFFFFFFFFFFFFFFFFFFFFFFFFFFFFFFFFFFFFFFFFFFFFFFFFFFFFFFFFFFFFFFFFFFFFFFFFFFFFFFFFFFFFFFFFFFFFFFFFFFFFFFFFFFFFFFFFFFFFFFFFFFFFFFFFFFFFFFFFFFFFFFFFFFFFFFFFFFFFFDFDFEFDFCF9E5AB6024070001000502000123390E07060D05044A1B01080E2C3C44B3D4542303051C413B18060502309927B014CF95BEE7F9FCF8F8FEFFFFFFFFFFFFFFFFFFFFFFFFFFFFFFFFFFFFFFFFFFFFFFFFFFFFFFFFFFFFFFFFFFFFFFFFFFFFFFFFFFFFFFFFFFFFFFFFFFFFFFFFFFFFFFFFFFFFFFFFFFFFFFFFFFFFFFFFFFFFFFFFFFFFFFFFFFFFFFFFFFFFFFFFFFFFFFFFFFFFFFFFFFFFFFFFFFFFFFFFFFFFFFFFFFFFFFFFFFFFFFFFFFFFFFFFFFFFFFFFFFFFFFFFFFFFFFFFFFFFFFFFFFFFFFFFFFFFFFFFFFFFFFFFFFFFFFFFFFFFFFFFFFFFFFFFFFFFFFFFFFFFFFFFFFFFFFFFFFFFFFFFFFFFFFFFFFFFFFFFFFFFFFFFFFFFFFFFFFFFFFFFFFFFFDFBFDFFFDFAF0D09C3205000300362A090134540900020B00076A2803030E3C4B3EB0CF49000005013F4519021D4783DB7AF05DE0D0F0F8FAFDFDFEFEFFFFFFFFFFFFFFFFFFFFFFFFFFFFFFFFFFFFFFFFFFFFFFFFFFFFFFFFFFFFFFFFFFFFFFFFFFFFFFFFFFFFFFFFFFFFFFFFFFFFFFFFFFFFFFFFFFFFFFFFFFFFFFFFFFFFFFFFFFFFFFFFFFFFFFFFFFFFFFFFFFFFFFFFFFFFFFFFFFFFFFFFFFFFFFFFFFFFFFFFFFFFFFFFFFFFFFFFFFFFFFFFFFFFFFFFFFFFFFFFFFFFFFFFFFFFFFFFFFFFFFFFFFFFFFFFFFFFFFFFFFFFFFFFFFFFFFFFFFFFFFFFFFFFFFFFFFFFFFFFFFFFFFFFFFFFFFFFFFFFFFFFFFFFFFFFFFFFFFFFFFFFFFFFFFFFFFFFFFFFFFFFFFFFFBFAFCF9F8FAEBC5913F000004054B4011001B2D000006050004331403000622456FB6B5280200020622220B032664919F9CF3CAF7FEFEFEFEFEFEFEFEFFFFFFFFFFFFFFFFFFFFFFFFFFFFFFFFFFFFFFFFFFFFFFFFFFFFFFFFFFFFFFFFFFFFFFFFFFFFFFFFFFFFFFFFFFFFFFFFFFFFFFFFFFFFFFFFFFFFFFFFFFFFFFFFFFFFFFFFFFFFFFFFFFFFFFFFFFFFFFFFFFFFFFFFFFFFFFFFFFFFFFFFFFFFFFFFFFFFFFFFFFFFFFFFFFFFFFFFFFFFFFFFFFFFFFFFFFFFFFFFFFFFFFFFFFFFFFFFFFFFFFFFFFFFFFFFFFFFFFFFFFFFFFFFFFFFFFFFFFFFFFFFFFFFFFFFFFFFFFFFFFFFFFFFFFFFFFFFFFFFFFFFFFFFFFFFFFFFFFFFFFFFFFFFFFFFFFFFFFFFFFFFFFFFFAF7FBF8F6F8E2AB7B410006030623250E02060700020A02020200020204060938899172071548301002030205255D766554A3F0FCFEFEFEFEFEFEFEFEFFFFFFFFFFFFFFFFFFFFFFFFFFFFFFFFFFFFFFFFFFFFFFFFFFFFFFFFFFFFFFFFFFFFFFFFFFFFFFFFFFFFFFFFFFFFFFFFFFFFFFFFFFFFFFFFFFFFFFFFFFFFFFFFFFFFFFFFFFFFFFFFFFFFFFFFFFFFFFFFFFFFFFFFFFFFFFFFFFFFFFFFFFFFFFFFFFFFFFFFFFFFFFFFFFFFFFFFFFFFFFFFFFFFFFFFFFFFFFFFFFFFFFFFFFFFFFFFFFFFFFFFFFFFFFFFFFFFFFFFFFFFFFFFFFFFFFFFFFFFFFFFFFFFFFFFFFFFFFFFFFFFFFFFFFFFFFFFFFFFFFFFFFFFFFFFFFFFFFFFFFFFFFFFFFFFFFFFFFFFFFFFFFFFFDF5FAFFF6F3EBAE622101090300020705030506060107010303030200050A072F745822082E925B11000404023791A7863E5FE7FCFEFEFEFEFEFEFEFEFFFFFFFFFFFFFFFFFFFFFFFFFFFFFFFFFFFFFFFFFFFFFFFFFFFFFFFFFFFFFFFFFFFFFFFFFFFFFFFFFFFFFFFFFFFFFFFFFFFFFFFFFFFFFFFFFFFFFFFFFFFFFFFFFFFFFFFFFFFFFFFFFFFFFFFFFFFFFFFFFFFFFFFFFFFFFFFFFFFFFFFFFFFFFFFFFFFFFFFFFFFFFFFFFFFFFFFFFFFFFFFFFFFFFFFFFFFFFFFFFFFFFFFFFFFFFFFFFFFFFFFFFFFFFFFFFFFFFFFFFFFFFFFFFFFFFFFFFFFFFFFFFFFFFFFFFFFFFFFFFFFFFFFFFFFFFFFFFFFFFFFFFFFFFFFFFFFFFFFFFFFFFFFFFFFFFFFFFFFFFFFFFFFFF9F9F9FAF3EFE38C1509080000000402030402071A2508372E333333372D182270AB6324042458370F0101011463BEB88AB1B3FBFEFEFEFEFEFEFEFEFEFFFFFFFFFFFFFFFFFFFFFFFFFFFFFFFFFFFFFFFFFFFFFFFFFFFFFFFFFFFFFFFFFFFFFFFFFFFFFFFFFFFFFFFFFFFFFFFFFFFFFFFFFFFFFFFFFFFFFFFFFFFFFFFFFFFFFFFFFFFFFFFFFFFFFFFFFFFFFFFFFFFFFFFFFFFFFFFFFFFFFFFFFFFFFFFFFFFFFFFFFFFFFFFFFFFFFFFFFFFFFFFFFFFFFFFFFFFFFFFFFFFFFFFFFFFFFFFFFFFFFFFFFFFFFFFFFFFFFFFFFFFFFFFFFFFFFFFFFFFFFFFFFFFFFFFFFFFFFFFFFFFFFFFFFFFFFFFFFFFFFFFFFFFFFFFFFFFFFFFFFFFFFFFFFFFFFFFFFFFFFFFFFFFFF5F8F2E8E4E6CA6405011001051A13050A0C142847803CCBBCC6BBC5D3B66067BA9D597706125261350C01024AB7BD71A9F9F3FEFEFEFEFEFEFEFEFEFEFFFFFFFFFFFFFFFFFFFFFFFFFFFFFFFFFFFFFFFFFFFFFFFFFFFFFFFFFFFFFFFFFFFFFFFFFFFFFFFFFFFFFFFFFFFFFFFFFFFFFFFFFFFFFFFFFFFFFFFFFFFFFFFFFFFFFFFFFFFFFFFFFFFFFFFFFFFFFFFFFFFFFFFFFFFFFFFFFFFFFFFFFFFFFFFFFFFFFFFFFFFFFFFFFFFFFFFFFFFFFFFFFFFFFFFFFFFFFFFFFFFFFFFFFFFFFFFFFFFFFFFFFFFFFFFFFFFFFFFFFFFFFFFFFFFFFFFFFFFFFFFFFFFFFFFFFFFFFFFFFFFFFFFFFFFFFFFFFFFFFFFFFFFFFFFFFFFFFFFFFFFFFFFFFFFFFFFFFFFFFFFFFFFFFEEEE9CC98938406030103031671551A02023CAABA65B9EDFEE9F5FBEDF3D57E331D1C2809588C8D661F05114FB2EBE1EFFDFCFCFDFEFEFEFEFEFEFEFEFFFFFFFFFFFFFFFFFFFFFFFFFFFFFFFFFFFFFFFFFFFFFFFFFFFFFFFFFFFFFFFFFFFFFFFFFFFFFFFFFFFFFFFFFFFFFFFFFFFFFFFFFFFFFFFFFFFFFFFFFFFFFFFFFFFFFFFFFFFFFFFFFFFFFFFFFFFFFFFFFFFFFFFFFFFFFFFFFFFFFFFFFFFFFFFFFFFFFFFFFFFFFFFFFFFFFFFFFFFFFFFFFFFFFFFFFFFFFFFFFFFFFFFFFFFFFFFFFFFFFFFFFFFFFFFFFFFFFFFFFFFFFFFFFFFFFFFFFFFFFFFFFFFFFFFFFFFFFFFFFFFFFFFFFFFFFFFFFFFFFFFFFFFFFFFFFFFFFFFFFFFFFFFFFFFFFFFFFFFFFFFFFFFFFDD38BAAC2640C061806000000200F020710196BE6D0FEFAFDFEFFEDF6F07E1201000003236BD6A0300500577E79C0FDFDF9FAF5FDFEFEFEFEFEFEFEFEFFFFFFFFFFFFFFFFFFFFFFFFFFFFFFFFFFFFFFFFFFFFFFFFFFFFFFFFFFFFFFFFFFFFFFFFFFFFFFFFFFFFFFFFFFFFFFFFFFFFFFFFFFFFFFFFFFFFFFFFFFFFFFFFFFFFFFFFFFFFFFFFFFFFFFFFFFFFFFFFFFFFFFFFFFFFFFFFFFFFFFFFFFFFFFFFFFFFFFFFFFFFFFFFFFFFFFFFFFFFFFFFFFFFFFFFFFFFFFFFFFFFFFFFFFFFFFFFFFFFFFFFFFFFFFFFFFFFFFFFFFFFFFFFFFFFFFFFFFFFFFFFFFFFFFFFFFFFFFFFFFFFFFFFFFFFFFFFFFFFFFFFFFFFFFFFFFFFFFFFFFFFFFFFFFFFFFFFFFFFFFFFFFFFF9EAAE7B8F74080C000000060A020E060000549762EDFEFAFBFEFFF9F07A170101030100086BA17D7C2700115BBEF0FAFEFEFDFCFDFEFEFEFEFEFEFEFEFFFFFFFFFFFFFFFFFFFFFFFFFFFFFFFFFFFFFFFFFFFFFFFFFFFFFFFFFFFFFFFFFFFFFFFFFFFFFFFFFFFFFFFFFFFFFFFFFFFFFFFFFFFFFFFFFFFFFFFFFFFFFFFFFFFFFFFFFFFFFFFFFFFFFFFFFFFFFFFFFFFFFFFFFFFFFFFFFFFFFFFFFFFFFFFFFFFFFFFFFFFFFFFFFFFFFFFFFFFFFFFFFFFFFFFFFFFFFFFFFFFFFFFFFFFFFFFFFFFFFFFFFFFFFFFFFFFFFFFFFFFFFFFFFFFFFFFFFFFFFFFFFFFFFFFFFFFFFFFFFFFFFFFFFFFFFFFFFFFFFFFFFFFFFFFFFFFFFFFFFFFFFFFFFFFBF9FEFBF6FEFDF8FDFEFCEBBB6416030700000403010704050206144FC1FCFDFBFDF5FBFDF0680B01115485011264E88D2D1004034CBEF8FEFFFEF8FDFEFEFEFEFEFEFEFEFEFFFFFFFFFFFFFFFFFFFFFFFFFFFFFFFFFFFFFFFFFFFFFFFFFFFFFFFFFFFFFFFFFFFFFFFFFFFFFFFFFFFFFFFFFFFFFFFFFFFFFFFFFFFFFFFFFFFFFFFFFFFFFFFFFFFFFFFFFFFFFFFFFFFFFFFFFFFFFFFFFFFFFFFFFFFFFFFFFFFFFFFFFFFFFFFFFFFFFFFFFFFFFFFFFFFFFFFFFFFFFFFFFFFFFFFFFFFFFFFFFFFFFFFFFFFFFFFFFFFFFFFFFFFFFFFFFFFFFFFFFFFFFFFFFFFFFFFFFFFFFFFFFFFFFFFFFFFFFFFFFFFFFFFFFFFFFFFFFFFFFFFFFFFFFFFFFFFFFFFFFFFFFFFFFFFCFBFEFEFBF2FDFDF3F0D2A198600F0D0204020005000A704A10044ECCF0ECFCEDF6FC5A42451C030507182F0003118B9A3007525976D3FDFDFEFEFAFEFFFFFFFFFFFFFFFFFFFFFFFFFFFFFFFFFFFFFFFFFFFFFFFFFFFFFFFFFFFFFFFFFFFFFFFFFFFFFFFFFFFFFFFFFFFFFFFFFFFFFFFFFFFFFFFFFFFFFFFFFFFFFFFFFFFFFFFFFFFFFFFFFFFFFFFFFFFFFFFFFFFFFFFFFFFFFFFFFFFFFFFFFFFFFFFFFFFFFFFFFFFFFFFFFFFFFFFFFFFFFFFFFFFFFFFFFFFFFFFFFFFFFFFFFFFFFFFFFFFFFFFFFFFFFFFFFFFFFFFFFFFFFFFFFFFFFFFFFFFFFFFFFFFFFFFFFFFFFFFFFFFFFFFFFFFFFFFFFFFFFFFFFFFFFFFFFFFFFFFFFFFFFFFFFFFFFFFFFFFFFFFFFFFFFEFEFFFFFDEFF4FDFBE36B2D2F2409030116694510010228220D528389ECF6FEFAFCF9C434080B0701060C07030C52A03C0C01205CA9E7FDFAFAFEFCFEFFFFFFFFFFFFFFFFFFFFFFFFFFFFFFFFFFFFFFFFFFFFFFFFFFFFFFFFFFFFFFFFFFFFFFFFFFFFFFFFFFFFFFFFFFFFFFFFFFFFFFFFFFFFFFFFFFFFFFFFFFFFFFFFFFFFFFFFFFFFFFFFFFFFFFFFFFFFFFFFFFFFFFFFFFFFFFFFFFFFFFFFFFFFFFFFFFFFFFFFFFFFFFFFFFFFFFFFFFFFFFFFFFFFFFFFFFFFFFFFFFFFFFFFFFFFFFFFFFFFFFFFFFFFFFFFFFFFFFFFFFFFFFFFFFFFFFFFFFFFFFFFFFFFFFFFFFFFFFFFFFFFFFFFFFFFFFFFFFFFFFFFFFFFFFFFFFFFFFFFFFFFFFFFFFFFFFFFFFFFFFFFFFFFFFFBF6FAFEF9FAFCDD8E7E36060602040B093229080A0202080524708D82F7FDF7F2654F1B020304040604020B4989440C00090763CCF3FDFBFCFEFDFDFEFFFFFFFFFFFFFFFFFFFFFFFFFFFFFFFFFFFFFFFFFFFFFFFFFFFFFFFFFFFFFFFFFFFFFFFFFFFFFFFFFFFFFFFFFFFFFFFFFFFFFFFFFFFFFFFFFFFFFFFFFFFFFFFFFFFFFFFFFFFFFFFFFFFFFFFFFFFFFFFFFFFFFFFFFFFFFFFFFFFFFFFFFFFFFFFFFFFFFFFFFFFFFFFFFFFFFFFFFFFFFFFFFFFFFFFFFFFFFFFFFFFFFFFFFFFFFFFFFFFFFFFFFFFFFFFFFFFFFFFFFFFFFFFFFFFFFFFFFFFFFFFFFFFFFFFFFFFFFFFFFFFFFFFFFFFFFFFFFFFFFFFFFFFFFFFFFFFFFFFFFFFFFFFFFFFFFFFFFFFFFFFFFFFBF6F7FBF8F2FEFDD64844230006070206000708010601000603032B69A6FAFCFAF6A223702F000406020207062F42070801367AB7E4FAFDFCFDFEFDFBFEFFFFFFFFFFFFFFFFFFFFFFFFFFFFFFFFFFFFFFFFFFFFFFFFFFFFFFFFFFFFFFFFFFFFFFFFFFFFFFFFFFFFFFFFFFFFFFFFFFFFFFFFFFFFFFFFFFFFFFFFFFFFFFFFFFFFFFFFFFFFFFFFFFFFFFFFFFFFFFFFFFFFFFFFFFFFFFFFFFFFFFFFFFFFFFFFFFFFFFFFFFFFFFFFFFFFFFFFFFFFFFFFFFFFFFFFFFFFFFFFFFFFFFFFFFFFFFFFFFFFFFFFFFFFFFFFFFFFFFFFFFFFFFFFFFFFFFFFFFFFFFFFFFFFFFFFFFFFFFFFFFFFFFFFFFFFFFFFFFFFFFFFFFFFFFFFFFFFFFFFFFFFFFFFFFFFFFFFFFFFFFFFFFFCF9FBFEF9E786E67D190C06020706010201030000080100070704349DCD96F8FEF9F9907D28030603020604093C5905090265D3F5F6FDFDFAFAFEFEFDFEFFFFFFFFFFFFFFFFFFFFFFFFFFFFFFFFFFFFFFFFFFFFFFFFFFFFFFFFFFFFFFFFFFFFFFFFFFFFFFFFFFFFFFFFFFFFFFFFFFFFFFFFFFFFFFFFFFFFFFFFFFFFFFFFFFFFFFFFFFFFFFFFFFFFFFFFFFFFFFFFFFFFFFFFFFFFFFFFFFFFFFFFFFFFFFFFFFFFFFFFFFFFFFFFFFFFFFFFFFFFFFFFFFFFFFFFFFFFFFFFFFFFFFFFFFFFFFFFFFFFFFFFFFFFFFFFFFFFFFFFFFFFFFFFFFFFFFFFFFFFFFFFFFFFFFFFFFFFFFFFFFFFFFFFFFFFFFFFFFFFFFFFFFFFFFFFFFFFFFFFFFFFFFFFFFFFFFFFFFFFFFFFFFFBF8FAF5FCE67B7C1901060302050606103C32090007064436100234818288F9F6FAF9942B0410403808070009396106090666D3FDFFFFFDFBF8FDFFFFFFFFFFFFFFFFFFFFFFFFFFFFFFFFFFFFFFFFFFFFFFFFFFFFFFFFFFFFFFFFFFFFFFFFFFFFFFFFFFFFFFFFFFFFFFFFFFFFFFFFFFFFFFFFFFFFFFFFFFFFFFFFFFFFFFFFFFFFFFFFFFFFFFFFFFFFFFFFFFFFFFFFFFFFFFFFFFFFFFFFFFFFFFFFFFFFFFFFFFFFFFFFFFFFFFFFFFFFFFFFFFFFFFFFFFFFFFFFFFFFFFFFFFFFFFFFFFFFFFFFFFFFFFFFFFFFFFFFFFFFFFFFFFFFFFFFFFFFFFFFFFFFFFFFFFFFFFFFFFFFFFFFFFFFFFFFFFFFFFFFFFFFFFFFFFFFFFFFFFFFFFFFFFFFFFFFFFFFFFFFFFFFFFFFFAF7FAF1EF99E51B05020101010519364354490D0102075D481101021F66A6F8F4FEFA841A0115554C08040304000F002C3A68CAFCFDFDFFFDF9FCFEFDFEFFFFFFFFFFFFFFFFFFFFFFFFFFFFFFFFFFFFFFFFFFFFFFFFFFFFFFFFFFFFFFFFFFFFFFFFFFFFFFFFFFFFFFFFFFFFFFFFFFFFFFFFFFFFFFFFFFFFFFFFFFFFFFFFFFFFFFFFFFFFFFFFFFFFFFFFFFFFFFFFFFFFFFFFFFFFFFFFFFFFFFFFFFFFFFFFFFFFFFFFFFFFFFFFFFFFFFFFFFFFFFFFFFFFFFFFFFFFFFFFFFFFFFFFFFFFFFFFFFFFFFFFFFFFFFFFFFFFFFFFFFFFFFFFFFFFFFFFFFFFFFFFFFFFFFFFFFFFFFFFFFFFFFFFFFFFFFFFFFFFFFFFFFFFFFFFFFFFFFFFFFFFFFFFFFFEFDFDFEFBFCFEFEFFFCFAFDE5539617010202050714427E87472304020401000806010138896FF5FCFEFAF541060B070000020609050204619EADDEFDFCFBFFFEFCFCFEFCFEFFFFFFFFFFFFFFFFFFFFFFFFFFFFFFFFFFFFFFFFFFFFFFFFFFFFFFFFFFFFFFFFFFFFFFFFFFFFFFFFFFFFFFFFFFFFFFFFFFFFFFFFFFFFFFFFFFFFFFFFFFFFFFFFFFFFFFFFFFFFFFFFFFFFFFFFFFFFFFFFFFFFFFFFFFFFFFFFFFFFFFFFFFFFFFFFFFFFFFFFFFFFFFFFFFFFFFFFFFFFFFFFFFFFFFFFFFFFFFFFFFFFFFFFFFFFFFFFFFFFFFFFFFFFFFFFFFFFFFFFFFFFFFFFFFFFFFFFFFFFFFFFFFFFFFFFFFFFFFFFFFFFFFFFFFFFFFFFFFFFFFFFFFFFFFFFFFFFFFFFFFFFFFFFFEFDFAFBFEF5F9FCFAFBF6F3FBEEA90C0402010D2827467E8086A8460D0100050203040203052BB6F5F2F9FCF56C1F2F2104080902020F1A3069DBFBF9FBFFFFFFFFFFFFFFFFFFFFFFFFFFFFFFFFFFFFFFFFFFFFFFFFFFFFFFFFFFFFFFFFFFFFFFFFFFFFFFFFFFFFFFFFFFFFFFFFFFFFFFFFFFFFFFFFFFFFFFFFFFFFFFFFFFFFFFFFFFFFFFFFFFFFFFFFFFFFFFFFFFFFFFFFFFFFFFFFFFFFFFFFFFFFFFFFFFFFFFFFFFFFFFFFFFFFFFFFFFFFFFFFFFFFFFFFFFFFFFFFFFFFFFFFFFFFFFFFFFFFFFFFFFFFFFFFFFFFFFFFFFFFFFFFFFFFFFFFFFFFFFFFFFFFFFFFFFFFFFFFFFFFFFFFFFFFFFFFFFFFFFFFFFFFFFFFFFFFFFFFFFFFFFFFFFFFFFFFFFFFFFFFFFFEF9FEFDFDFEFCF9F5F1F4FBF2CE3242340D082F8D807DA7754E0B060604000A050D0C040201042EB8C0EBD2F4C474A771243B2402040C74AC99E6F8F9FCFFFFFFFFFFFFFFFFFFFFFFFFFFFFFFFFFFFFFFFFFFFFFFFFFFFFFFFFFFFFFFFFFFFFFFFFFFFFFFFFFFFFFFFFFFFFFFFFFFFFFFFFFFFFFFFFFFFFFFFFFFFFFFFFFFFFFFFFFFFFFFFFFFFFFFFFFFFFFFFFFFFFFFFFFFFFFFFFFFFFFFFFFFFFFFFFFFFFFFFFFFFFFFFFFFFFFFFFFFFFFFFFFFFFFFFFFFFFFFFFFFFFFFFFFFFFFFFFFFFFFFFFFFFFFFFFFFFFFFFFFFFFFFFFFFFFFFFFFFFFFFFFFFFFFFFFFFFFFFFFFFFFFFFFFFFFFFFFFFFFFFFFFFFFFFFFFFFFFFFFFFFFFFFFFFFFFFFFFFFFFFFFFFFEF9FEF8F9FCF9FBFFFEFEF9F3FEF9D7DA2D1B5E8F7B96BF680900010100050400041F1D080B0700010AD216B357B1D67A80C3740A06267FEEFCF4F2F9FEFFFFFFFFFFFFFFFFFFFFFFFFFFFFFFFFFFFFFFFFFFFFFFFFFFFFFFFFFFFFFFFFFFFFFFFFFFFFFFFFFFFFFFFFFFFFFFFFFFFFFFFFFFFFFFFFFFFFFFFFFFFFFFFFFFFFFFFFFFFFFFFFFFFFFFFFFFFFFFFFFFFFFFFFFFFFFFFFFFFFFFFFFFFFFFFFFFFFFFFFFFFFFFFFFFFFFFFFFFFFFFFFFFFFFFFFFFFFFFFFFFFFFFFFFFFFFFFFFFFFFFFFFFFFFFFFFFFFFFFFFFFFFFFFFFFFFFFFFFFFFFFFFFFFFFFFFFFFFFFFFFFFFFFFFFFFFFFFFFFFFFFFFFFFFFFFFFFFFFFFFFFFFFFFFFFFFFFFFFFFFFFFFFFEFCFEF8F7FCFCFCFEFDF9F3FEFAD52FB474C1A96071816D5B0601011501001607006C77180000020E12000000AED6527D890704143478FEFEF3FAFEF5F8FFFFFFFFFFFFFFFFFFFFFFFFFFFFFFFFFFFFFFFFFFFFFFFFFFFFFFFFFFFFFFFFFFFFFFFFFFFFFFFFFFFFFFFFFFFFFFFFFFFFFFFFFFFFFFFFFFFFFFFFFFFFFFFFFFFFFFFFFFFFFFFFFFFFFFFFFFFFFFFFFFFFFFFFFFFFFFFFFFFFFFFFFFFFFFFFFFFFFFFFFFFFFFFFFFFFFFFFFFFFFFFFFFFFFFFFFFFFFFFFFFFFFFFFFFFFFFFFFFFFFFFFFFFFFFFFFFFFFFFFFFFFFFFFFFFFFFFFFFFFFFFFFFFFFFFFFFFFFFFFFFFFFFFFFFFFFFFFFFFFFFFFFFFFFFFFFFFFFFFFFFFFFFFFFFFFFFFFFFFFFFFFFFFEFDFEFEFEFDFDF8F4FAFEE9BADBCCAE2E91410F2C7B61090E0800011F1D030100162D200B0A020603010007040C131F736C07105EE8F6FEFEFDFAFDF7F9FFFFFFFFFFFFFFFFFFFFFFFFFFFFFFFFFFFFFFFFFFFFFFFFFFFFFFFFFFFFFFFFFFFFFFFFFFFFFFFFFFFFFFFFFFFFFFFFFFFFFFFFFFFFFFFFFFFFFFFFFFFFFFFFFFFFFFFFFFFFFFFFFFFFFFFFFFFFFFFFFFFFFFFFFFFFFFFFFFFFFFFFFFFFFFFFFFFFFFFFFFFFFFFFFFFFFFFFFFFFFFFFFFFFFFFFFFFFFFFFFFFFFFFFFFFFFFFFFFFFFFFFFFFFFFFFFFFFFFFFFFFFFFFFFFFFFFFFFFFFFFFFFFFFFFFFFFFFFFFFFFFFFFFFFFFFFFFFFFFFFFFFFFFFFFFFFFFFFFFFFFFFFFFFFFFFFFFFFFFFFFFFFFFEFEFCFDFEFEFBFBFDEBC99B44241D2B061407040A1711000101000E6A5D0C000280BD8340291B000300080902010008140E0C284782FDEBFCFEFBFDFEFEFFFFFFFFFFFFFFFFFFFFFFFFFFFFFFFFFFFFFFFFFFFFFFFFFFFFFFFFFFFFFFFFFFFFFFFFFFFFFFFFFFFFFFFFFFFFFFFFFFFFFFFFFFFFFFFFFFFFFFFFFFFFFFFFFFFFFFFFFFFFFFFFFFFFFFFFFFFFFFFFFFFFFFFFFFFFFFFFFFFFFFFFFFFFFFFFFFFFFFFFFFFFFFFFFFFFFFFFFFFFFFFFFFFFFFFFFFFFFFFFFFFFFFFFFFFFFFFFFFFFFFFFFFFFFFFFFFFFFFFFFFFFFFFFFFFFFFFFFFFFFFFFFFFFFFFFFFFFFFFFFFFFFFFFFFFFFFFFFFFFFFFFFFFFFFFFFFFFFFFFFFFFFFFFFFFFFFFFFFFFFFFFFFFFFEFCFAFAFDD6E2DF81386A8A02060700000102010101000200050419150204056C8C86AFC899492B2815231026160509080D50C2DAD8FAFEF0FAF6FAFFFFFFFFFFFFFFFFFFFFFFFFFFFFFFFFFFFFFFFFFFFFFFFFFFFFFFFFFFFFFFFFFFFFFFFFFFFFFFFFFFFFFFFFFFFFFFFFFFFFFFFFFFFFFFFFFFFFFFFFFFFFFFFFFFFFFFFFFFFFFFFFFFFFFFFFFFFFFFFFFFFFFFFFFFFFFFFFFFFFFFFFFFFFFFFFFFFFFFFFFFFFFFFFFFFFFFFFFFFFFFFFFFFFFFFFFFFFFFFFFFFFFFFFFFFFFFFFFFFFFFFFFFFFFFFFFFFFFFFFFFFFFFFFFFFFFFFFFFFFFFFFFFFFFFFFFFFFFFFFFFFFFFFFFFFFFFFFFFFFFFFFFFFFFFFFFFFFFFFFFFFFFFFFFFFFFFFFFFFFFFFFFFFFFEF8FDFDFCFD9C9A7E2137813004020100030100000204030000020003000003012063789EC389C5BF933FB3278D3F0405030C3F90A99AEAD4F3FDF9F9FCFFFFFFFFFFFFFFFFFFFFFFFFFFFFFFFFFFFFFFFFFFFFFFFFFFFFFFFFFFFFFFFFFFFFFFFFFFFFFFFFFFFFFFFFFFFFFFFFFFFFFFFFFFFFFFFFFFFFFFFFFFFFFFFFFFFFFFFFFFFFFFFFFFFFFFFFFFFFFFFFFFFFFFFFFFFFFFFFFFFFFFFFFFFFFFFFFFFFFFFFFFFFFFFFFFFFFFFFFFFFFFFFFFFFFFFFFFFFFFFFFFFFFFFFFFFFFFFFFFFFFFFFFFFFFFFFFFFFFFFFFFFFFFFFFFFFFFFFFFFFFFFFFFFFFFFFFFFFFFFFFFFFFFFFFFFFFFFFFFFFFFFFFFFFFFFFFFFFFFFFFFFFFFFFFFFFFFFCF9F8F9FDF9FEFBFEFED1EAD07623021A380302010300050A28150203020305000501002B1C06092674B9BFCDF8EB6DBBF25EC97114010115466654A0A0A5F1FDFDFBFCFAF9FFFFFDF7FEFEFFFFFFFFFFFFFFFFFFFDFAFFFEFEFEFEFCFDFDFDFFFFFEFDFEFDFEFFFFFFFFFFFFFFFFFFFFFFFFFFFFFFFFFFFFFFFFFFFFFFFFFFFFFFFFFFFFFFFFFFFFFFFFFFFFFFFFFFFFFFFFFFFFFFFFFFFFFFFFFFFFFFFFFFFFFFFFFFFFFFFFFFFFFFFFFFFFFFFFFFFFFFFFFFFFFFFFFFFFFFFFFFFFFFFFFFFFFFFFFFFFFFFFFFFFFFFFFFFFFFFFFFFFFFFFFFFFFFFFFFFFFFFFFFFFFFFFFFFFFFFFFFFFFFFFFFFFFFFFFFFFFFFFFFFFFFFFFFFFFFFFFFFFFFFFFFFFFAFCFDFBFAFAFCF6FDFEF2FC97B0C7590102040303070402060319603200070E141B0303050E543B12060457A0ACF4FDF3ACFBFBC2C1903605000B26351B4F539CC6D9F4FDFEF6F3FAFAFDF9FEFDFEFFFFFFFFFFFFFFFFFEFBFEFDFDF1D3EEFBFBFDFBF9FBFEFEFAFEFFFFFFFFFFFFFFFFFFFFFFFFFFFFFFFFFFFFFFFFFFFFFFFFFFFFFFFFFFFFFFFFFFFFFFFFFFFFFFFFFFFFFFFFFFFFFFFFFFFFFFFFFFFFFFFFFFFFFFFFFFFFFFFFFFFFFFFFFFFFFFFFFFFFFFFFFFFFFFFFFFFFFFFFFFFFFFFFFFFFFFFFFFFFFFFFFFFFFFFFFFFFFFFFFFFFFFFFFFFFFFFFFFFFFFFFFFFFFFFFFFFFFFFFFFFFFFFFFFFFFFFFFFFFFFFFFFFFFFFFFFFFFFFFFFFFFFFFFFFFF8FAFEFFFEFCFAFAFCF9F3F2BA725A2801020606020C03000903103E220B224256421D17040C352A0E0101265E92C7D0F5CDF0F9FB99642A040001030903021D575598D1E4F3F9F7FAF9FAF6FDFDFEFFFFFFFFFFFFFFFCF9FEF8F9FCD584CADEE8FBFAF6F9FEFEFBFEFFFFFFFFFFFFFFFFFFFFFFFFFFFFFFFFFFFFFFFFFFFFFFFFFFFFFFFFFFFFFFFFFFFFFFFFFFFFFFFFFFFFFFFFFFFFFFFFFFFFFFFFFFFFFFFFFFFFFFFFFFFFFFFFFFFFFFFFFFFFFFFFFFFFFFFFFFFFFFFFFFFFFFFFFFFFFFFFFFFFFFFFFFFFFFFFFFFFFFFFFFFFFFFFFFFFFFFFFFFFFFFFFFFFFFFFFFFFFFFFFFFFFFFFFFFFFFFFFFFFFFFFFFFFFFFFFFFFFFFFFFFFFFFFFFFFFFFFFFFFFCFBFBFCFEFBF8FEFCF8FEF2FA5B0D0502000316260B2D02040402061D4E7B897E63585C17010607010303062756616FB561EACFD2872D07030307182000050405036291B6E3FBFDFBFBF9F4FDFBFEFFFFFFFFFFFFFFFCF8FDF7F9FDBA4A90AED2F8FCFAFBFDFBFFFFFFFFFFFFFFFFFFFFFFFFFFFFFFFFFFFFFFFFFFFFFFFFFFFFFFFFFFFFFFFFFFFFFFFFFFFFFFFFFFFFFFFFFFFFFFFFFFFFFFFFFFFFFFFFFFFFFFFFFFFFFFFFFFFFFFFFFFFFFFFFFFFFFFFFFFFFFFFFFFFFFFFFFFFFFFFFFFFFFFFFFFFFFFFFFFFFFFFFFFFFFFFFFFFFFFFFFFFFFFFFFFFFFFFFFFFFFFFFFFFFFFFFFFFFFFFFFFFFFFFFFFFFFFFFFFFFFFFFFFFFFFFFFFFFFFFFFFFFFFFFFFFDF8F9FEFCFAFDF7DCFAE5D44F0D030302012666147905050A202367A4B6A2798995885414210E0305050405050619440ECD6B63621F0206040F3E550509000004354F90C5E3F4F8FBFDFAFDF9FEFFFFFFFFFFFFFFFEFCF9FCFAF1A148589BDCF9FAF9FEFCF7FFFFFFFFFFFFFFFFFFFFFFFFFFFFFFFFFFFFFFFFFFFFFFFFFFFFFFFFFFFFFFFFFFFFFFFFFFFFFFFFFFFFFFFFFFFFFFFFFFFFFFFFFFFFFFFFFFFFFFFFFFFFFFFFFFFFFFFFFFFFFFFFFFFFFFFFFFFFFFFFFFFFFFFFFFFFFFFFFFFFFFFFFFFFFFFFFFFFFFFFFFFFFFFFFFFFFFFFFFFFFFFFFFFFFFFFFFFFFFFFFFFFFFFFFFFFFFFFFFFFFFFFFFFFFFFFFFFFFFFFFFFFFFFFFFFFFFFFFFFFFFFFFEFDFBFBFEFDFBF9EF7BCC8A61280F06040401215F1A6B0B24185E5DA2A2767D8598AA709F405641210803080803000001035C0A0409081712020D354908010300000523586D8EC1DDF3FBF5FDFCFEFFFFFFFFFFFFFFFEFAF2F9F4D8937D5EB7F9FCF5F4FFFDF8FFFFFFFFFFFFFFFFFFFFFFFFFFFFFFFFFFFFFFFFFFFFFFFFFFFFFFFFFFFFFFFFFFFFFFFFFFFFFFFFFFFFFFFFFFFFFFFFFFFFFFFFFFFFFFFFFFFFFFFFFFFFFFFFFFFFFFFFFFFFFFFFFFFFFFFFFFFFFFFFFFFFFFFFFFFFFFFFFFFFFFFFFFFFFFFFFFFFFFFFFFFFFFFFFFFFFFFFFFFFFFFFFFFFFFFFFFFFFFFFFFFFFFFFFFFFFFFFFFFFFFFFFFFFFFFFFFFFFFFFFFFFFFFFFFFFFFFFFFFFFFFFFDFEFDFDFCFAF6E2E5274F14050513110502020B1802171B773466867577815E2B317692DD7355865616010107090104070000040C0311543C01030600020007010000080500245B91D8FBFDFCF8FEFFFFFFFFFFFFFFFEFBF7FEFDDEBDDEBEDEFCFBF9FAFEFCF8FFFFFFFFFFFFFFFFFFFFFFFFFFFFFFFFFFFFFFFFFFFFFFFFFFFFFFFFFFFFFFFFFFFFFFFFFFFFFFFFFFFFFFFFFFFFFFFFFFFFFFFFFFFFFFFFFFFFFFFFFFFFFFFFFFFFFFFFFFFFFFFFFFFFFFFFFFFFFFFFFFFFFFFFFFFFFFFFFFFFFFFFFFFFFFFFFFFFFFFFFFFFFFFFFFFFFFFFFFFFFFFFFFFFFFFFFFFFFFFFFFFFFFFFFFFFFFFFFFFFFFFFFFFFFFFFFFFFFFFFFFFFFFFFFFFFFFFFFFFFFFFFFEFEFAFAFAF7E69FB844050010085B43090A06287F208615E18B24414496A5300000167BECD18D8F420D02050600040F080201010000094C39010305020900071403030104043376ACEEFDF7F6FEFFFFFFFFFFFFFFFFFDF7F4FEFBEEE8F8F7F1F5FBFEFEFCF9FBFEFEFFFFFFFFFFFFFFFFFFFFFFFFFFFFFFFFFFFFFFFFFFFFFFFFFFFFFFFFFFFFFFFFFFFFFFFFFFFFFFFFFFFFFFFFFFFFFFFFFFFFFFFFFFFFFFFFFFFFFFFFFFFFFFFFFFFFFFFFFFFFFFFFFFFFFFFFFFFFFFFFFFFFFFFFFFFFFFFFFFFFFFFFFFFFFFFFFFFFFFFFFFFFFFFFFFFFFFFFFFFFFFFFFFFFFFFFFFFFFFFFFFFFFFFFFFFFFFFFFFFFFFFFFFFFFFFFFFFFFFFFFFFFFFFFFFFFFFFFFFFFFCFDF8F8FCF7C93E451A010801023C310703061E7EAFFAA6D46D19053683600707000844C6EBFCC59B49120303060000021500090B080200050905050A000B44611509050401245689DAFEFDFBFEFFFFFFFFFFFFFFFFFFFEFDFFFFFFFFF5FBFFFFFFFDFAF9FFFEFEFEFFFFFFFFFFFFFFFFFFFFFFFFFFFFFFFFFFFFFFFFFFFFFFFFFFFFFFFFFFFFFFFFFFFFFFFFFFFFFFFFFFFFFFFFFFFFFFFFFFFFFFFFFFFFFFFFFFFFFFFFFFFFFFFFFFFFFFFFFFFFFFFFFFFFFFFFFFFFFFFFFFFFFFFFFFFFFFFFFFFFFFFFFFFFFFFFFFFFFFFFFFFFFFFFFFFFFFFFFFFFFFFFFFFFFFFFFFFFFFFFFFFFFFFFFFFFFFFFFFFFFFFFFFFFFFFFFFFFFFFFFFFFFFFFFFFFFFFFFFFFFCFDFAFAD9884B13020007080100010907000032C6F4FEFCEE91230136AEBC46050C0C61E3F6FDFDF6D6AC53111000010C0600098C2A080402020F46680B558D815C1F4458260137A8EFFDFDEEFEFFFFFFFFFFFFFFFFFFFEFEFFFFFFFFFEFEFFFFFFFEFEFEFFFEFEFEFFFFFFFFFFFFFFFFFFFFFFFFFFFFFFFFFFFFFFFFFFFFFFFFFFFFFFFFFFFFFFFFFFFFFFFFFFFFFFFFFFFFFFFFFFFFFFFFFFFFFFFFFFFFFFFFFFFFFFFFFFFFFFFFFFFFFFFFFFFFFFFFFFFFFFFFFFFFFFFFFFFFFFFFFFFFFFFFFFFFFFFFFFFFFFFFFFFFFFFFFFFFFFFFFFFFFFFFFFFFFFFFFFFFFFFFFFFFFFFFFFFFFFFFFFFFFFFFFFFFFFFFFFFFFFFFFFFFFFFFFFFFFFFFFFFFFFFFFFFFFEFAFBF4B4605C340005007432110D0605486F2175FCFEEFEF7A832F0F3E666B480418445EB5F3FCFAECED8F1702050B0100000C4A2D0A115E62528B38051B7EB58C2034391C1120407ECFFEFCFDFEFFFFFFFFFFFFFFFFFFFFFFFFFFFFFFFFFFFFFFFFFFFFFFFFFFFFFFFFFFFFFFFFFFFFFFFFFFFFFFFFFFFFFFFFFFFFFFFFFFFFFFFFFFFFFFFFFFFFFFFFFFFFFFFFFFFFFFFFFFFFFFFFFFFFFFFFFFFFFFFFFFFFFFFFFFFFFFFFFFFFFFFFFFFFFFFFFFFFFFFFFFFFFFFFFFFFFFFFFFFFFFFFFFFFFFFFFFFFFFFFFFFFFFFFFFFFFFFFFFFFFFFFFFFFFFFFFFFFFFFFFFFFFFFFFFFFFFFFFFFFFFFFFFFFFFFFFFFFFFFFFFFFFFFFFFFFFFFFFFFFFFFFFFFFFFFFF8FAFDD87630411C00010B501703020103256797876CFEFBECA157150E227CC36A00041261D8F2FBFDFCFAD97714010203000A030F08110E333727431F3B9EBDEFD86E200904030530A2F3F3FBF7FEFFFFFFFFFFFFFFFFFFFFFFFFFFFFFFFFFFFFFFFFFFFFFFFFFFFFFFFFFFFFFFFFFFFFFFFFFFFFFFFFFFFFFFFFFFFFFFFFFFFFFFFFFFFFFFFFFFFFFFFFFFFFFFFFFFFFFFFFFFFFFFFFFFFFFFFFFFFFFFFFFFFFFFFFFFFFFFFFFFFFFFFFFFFFFFFFFFFFFFFFFFFFFFFFFFFFFFFFFFFFFFFFFFFFFFFFFFFFFFFFFFFFFFFFFFFFFFFFFFFFFFFFFFFFFFFFFFFFFFFFFFFFFFFFFFFFFFFFFFFFFFFFFFFFFFFFFFFFFFFFFFFFFFFFFFFFFFFFFFFFFFFFFFFFFFFFFCF6F4C95D0910090001031608080A514640B5EBD69DF6FEDA7C822C031158B6AE3F070333AFF2FDF9F6B36E3906000204000E01010407030A0B080C0B35A1E3FCC5440E01043E603F68DBFBFDF6FEFFFFFFFFFFFFFFFFFFFFFFFFFFFFFFFFFFFFFFFFFFFFFFFFFFFFFFFFFFFFFFFFFFFFFFFFFFFFFFFFFFFFFFFFFFFFFFFFFFFFFFFFFFFFFFFFFFFFFFFFFFFFFFFFFFFFFFFFFFFFFFFFFFFFFFFFFFFFFFFFFFFFFFFFFFFFFFFFFFFFFFFFFFFFFFFFFFFFFFFFFFFFFFFFFFFFFFFFFFFFFFFFFFFFFFFFFFFFFFFFFFFFFFFFFFFFFFFFFFFFFFFFFFFFFFFFFFFFFFFFFFFFFFFFFFFFFFFFFFFFFFFFFFFFFFFFFFFFFFFFFFFFFFFFFFFFFFFFFFFFFFFFFFFFFFFFFFFAC971230206030E0000020604044779A3EAFAF7E6FAF9F094691F062B6494A47A0F5391D8F2F7F3F7CA9E918C6A07010004010402085D460004080635BFF4FD9A5F451E0625381D42C6FCFEF9FEFFFFFFFFFFFFFFFFFFFFFFFFFFFFFFFFFFFFFFFFFFFFFFFFFFFFFFFFFFFFFFFFFFFFFFFFFFFFFFFFFFFFFFFFFFFFFFFFFFFFFFFFFFFFFFFFFFFFFFFFFFFFFFFFFFFFFFFFFFFFFFFFFFFFFFFFFFFFFFFFFFFFFFFFFFFFFFFFFFFFFFFFFFFFFFFFFFFFFFFFFFFFFFFFFFFFFFFFFFFFFFFFFFFFFFFFFFFFFFFFFFFFFFFFFFFFFFFFFFFFFFFFFFFFFFFFFFFFFFFFFFFFFFFFFFFFFFFFFFFFFFFFFFFFFFFFFFFFFFFFFFFFFFFFFFFFFFFFFFFFFFFFFFFFFFFFF8FAD7802A0303010803000002060452B5BA859DFBF5FEFBF1CF450A051850877C390D4662C3F3FAF9FCEFE2E2BD79670C010300020105433200041039AEE9FCFC8443401C020D120439D1FBECFCFEFFFFFFFFFFFFFFFFFFFFFFFFFFFFFFFFFFFFFFFFFFFFFFFFFFFFFFFFFFFFFFFFFFFFFFFFFFFFFFFFFFFFFFFFFFFFFFFFFFFFFFFFFFFFFFFFFFFFFFFFFFFFFFFFFFFFFFFFFFFFFFFFFFFFFFFFFFFFFFFFFFFFFFFFFFFFFFFFFFFFFFFFFFFFFFFFFFFFFFFFFFFFFFFFFFFFFFFFFFFFFFFFFFFFFFFFFFFFFFFFFFFFFFFFFFFFFFFFFFFFFFFFFFFFFFFFFFFFFFFFFFFFFFFFFFFFFFFFFFFFFFFFFFFFFFFFFFFFFFFFFFFFFFFFFFFFFFFFFFFFFFFFFFFFFFFFF7FAD7791E0405000202006221486CACEAA03E9DFDF6EFFEE89B6B230B20537E733D071651C6F0F6FCF9F8F9F8D893726A02020200030E483809133A7DB2F4FEFD902B170D0407070136D2FCEDFAFEFFFFFFFFFFFFFFFFFFFFFFFFFFFFFFFFFFFFFFFFFFFFFFFFFFFFFFFFFFFFFFFFFFFFFFFFFFFFFFFFFFFFFFFFFFFFFFFFFFFFFFFFFFFFFFFFFFFFFFFFFFFFFFFFFFFFFFFFFFFFFFFFFFFFFFFFFFFFFFFFFFFFFFFFFFFFFFFFFFFFFFFFFFFFFFFFFFFFFFFFFFFFFFFFFFFFFFFFFFFFFFFFFFFFFFFFFFFFFFFFFFFFFFFFFFFFFFFFFFFFFFFFFFFFFFFFFFFFFFFFFFFFFFFFFFFFFFFFFFFFFFFFFFFFFFFFFFFFFFFFFFFFFFFFFFFFFFFFFFFFFFFFFFFFFFFFFCFDDC7C1D0407000C56196F67B1DDF2CD848FE6FDF9EFFEE9846D260A3482A071460B45AEEDF5F5F9F7FAFFFCD88F7475060002040451753B12407A81A5F7FEFD9D2A03042936180234C9FBFBFBFEFFFFFFFFFFFFFFFFFFFFFFFFFFFFFFFFFFFFFFFFFFFFFFFFFFFFFFFFFFFFFFFFFFFFFFFFFFFFFFFFFFFFFFFFFFFFFFFFFFFFFFFFFFFFFFFFFFFFFFFFFFFFFFFFFFFFFFFFFFFFFFFFFFFFFFFFFFFFFFFFFFFFFFFFFFFFFFFFFFFFFFFFFFFFFFFFFFFFFFFFFFFFFFFFFFFFFFFFFFFFFFFFFFFFFFFFFFFFFFFFFFFFFFFFFFFFFFFFFFFFFFFFFFFFFFFFFFFFFFFFFFFFFFFFFFFFFFFFFFFFFFFFFFFFFFFFFFFFFFFFFFFFFFFFFFFFFFFEFEFEFEFEFEFEFEFEFEE6B874280302020D7A6A5CBBF0F8FAB6A2F3FAF7FAFEFCF2D6721C011E688845153C83C7F2FDFDF9FAFDFFF8C38CCB630900000C3F824D0B0C488C94EAFEFDD56C3A2C2C7CA37A5D79DCFAF7FEFFFFFFFFFFFFFFFFFFFFFFFFFFFFFFFFFFFFFFFFFFFFFFFFFFFFFFFFFFFFFFFFFFFFFFFFFFFFFFFFFFFFFFFFFFFFFFFFFFFFFFFFFFFFFFFFFFFFFFFFFFFFFFFFFFFFFFFFFFFFFFFFFFFFFFFFFFFFFFFFFFFFFFFFFFFFFFFFFFFFFFFFFFFFFFFFFFFFFFFFFFFFFFFFFFFFFFFFFFFFFFFFFFFFFFFFFFFFFFFFFFFFFFFFFFFFFFFFFFFFFFFFFFFFFFFFFFFFFFFFFFFFFFFFFFFFFFFFFFFFFFFFFFFFFFFFFFFFFFFFFFFFFFFFFFFFFFFFFEFEFEFEFEFEFEFEFEFCCD7C3D1901000B384DCDE2F8F8F8FAF4F4FDFAF3FEF9FAEFBA8F2B02113444241A627FA9EFFDF6FAFBFAF9F5EFF0EB9D4707000B76470C042A5C94EEF3FEF9B23D545371AFD4CFCAE3FBF7F0FEFFFFFFFFFFFFFFFFFFFFFFFFFFFFFFFFFFFFFFFFFFFFFFFFFFFFFFFFFFFFFFFFFFFFFFFFFFFFFFFFFFFFFFFFFFFFFFFFFFFFFFFFFFFFFFFFFFFFFFFFFFFFFFFFFFFFFFFFFFFFFFFFFFFFFFFFFFFFFFFFFFFFFFFFFFFFFFFFFFFFFFFFFFFFFFFFFFFFFFFFFFFFFFFFFFFFFFFFFFFFFFFFFFFFFFFFFFFFFFFFFFFFFFFFFFFFFFFFFFFFFFFFFFFFFFFFFFFFFFFFFFFFFFFFFFFFFFFFFFFFFFFFFFFFFFFFFFFFFFFFFFFFFFFFFFFFFFFFFEFEFEFEFEFEFEFEFEE8DECF933C05021A93E7F1FDFCFBFCF6F7FBFAFEFBFCF2FBD8563210031E719947163AB2EAFDFFF2F8FDF7F1F5FEFCD55CDA140102142C1A2D85D3EDFEFEF9F5F58A2706687E8E8C6EC0D7E6FEFCFEFFFFFFFFFFFFFFFFFFFFFFFFFFFFFFFFFFFFFFFFFFFFFFFFFFFFFFFFFFFFFFFFFFFFFFFFFFFFFFFFFFFFFFFFFFFFFFFFFFFFFFFFFFFFFFFFFFFFFFFFFFFFFFFFFFFFFFFFFFFFFFFFFFFFFFFFFFFFFFFFFFFFFFFFFFFFFFFFFFFFFFFFFFFFFFFFFFFFFFFFFFFFFFFFFFFFFFFFFFFFFFFFFFFFFFFFFFFFFFFFFFFFFFFFFFFFFFFFFFFFFFFFFFFFFFFFFFFFFFFFFFFFFFFFFFFFFFFFFFFFFFFFFFFFFFFFFFFFFFFFFFFFFFFFFFFFFFFEFEFEFEFEFEFEFEFEB9958342030109130FADE9FCF5FAFBF9FAFDFDFAF6FEFCF4F8D7B13600020B1005208CEBFCFCFCFEF8F9FBFEFCFBFCF2DDF2080209117A597B8C84FDFAF9EDFDFE96280454656E6D618F56CFFDF5FDFFFFFFFFFFFFFFFFFFFFFFFFFFFFFFFFFFFFFFFFFFFFFFFFFFFFFFFFFFFFFFFFFFFFFFFFFFFFFFFFFFFFFFFFFFFFFFFFFFFFFFFFFFFFFFFFFFFFFFFFFFFFFFFFFFFFFFFFFFFFFFFFFFFFFFFFFFFFFFFFFFFFFFFFFFFFFFFFFFFFFFFFFFFFFFFFFFFFFFFFFFFFFFFFFFFFFFFFFFFFFFFFFFFFFFFFFFFFFFFFFFFFFFFFFFFFFFFFFFFFFFFFFFFFFFFFFFFFFFFFFFFFFFFFFFFFFFFFFFFFFFFFFFFFFFFFFFFFFFFFFFFFFFFFFFFFFFFEFEFEFEFEFEFEFEFEEEA342120A0000010813BCE4E0F1FDFDF7F5F8F7FEFCF3FEDF360301010202000C387CEFFCF5F2FBFEFEFBFBFDFAF9D138E10C000301100C1F91FEFEFAF3FEF6F9A126060D10120F040127CEFCFDFEFFFFFFFFFFFFFFFFFFFFFFFFFFFFFFFFFFFFFFFFFFFFFFFFFFFFFFFFFFFFFFFFFFFFFFFFFFFFFFFFFFFFFFFFFFFFFFFFFFFFFFFFFFFFFFFFFFFFFFFFFFFFFFFFFFFFFFFFFFFFFFFFFFFFFFFFFFFFFFFFFFFFFFFFFFFFFFFFFFFFFFFFFFFFFFFFFFFFFFFFFFFFFFFFFFFFFFFFFFFFFFFFFFFFFFFFFFFFFFFFFFFFFFFFFFFFFFFFFFFFFFFFFFFFFFFFFFFFFFFFFFFFFFFFFFFFFFFFFFFFFFFFFFFFFFFFFFFFFFFFFFFFFFFFFFFFFFFEFEFEFEFEFEFEFEFEB083884E0105070323D93C0D1A92E8F2F3FFFFFCF5F3FEFEFECB300003070D0A2E93878DE8FCFEFEF5FAFEFAEBE5FEFBD60F000000060326A4FDFDFEFEFEFDF2FDFE922100669141040734CCFEF4FDFFFFFFFFFFFFFFFFFFFFFFFFFFFFFFFFFFFFFFFFFFFFFFFFFFFFFFFFFFFFFFFFFFFFFFFFFFFFFFFFFFFFFFFFFFFFFFFFFFFFFFFFFFFFFFFFFFFFFFFFFFFFFFFFFFFFFFFFFFFFFFFFFFFFFFFFFFFFFFFFFFFFFFFFFFFFFFFFFFFFFFFFFFFFFFFFFFFFFFFFFFFFFFFFFFFFFFFFFFFFFFFFFFFFFFFFFFFFFFFFFFFFFFFFFFFFFFFFFFFFFFFFFFFFFFFFFFFFFFFFFFFFFFFFFFFFFFFFFFFFFFFFFFFFFFFFFFFFFFFFFFFFFFFFFFFFFFFEFEFEFEFEFEFEFEFEF89E360A0200001AA2F5380F97E8FDFEFEFFFFFEFDF9FAEADB619F4806003097E1E6F3ECFCF9F6FDFEFEFAFAFDFDF8EDFCE40D0605011193F1F6F4FBF4F8EF988D782B0F040856772D0D2BE5FEFEFEFFFFFFFFFFFFFFFFFFFFFFFFFFFFFFFFFFFFFFFFFFFFFFFFFFFFFFFFFFFFFFFFFFFFFFFFFFFFFFFFFFFFFFFFFFFFFFFFFFFFFFFFFFFFFFFFFFFFFFFFFFFFFFFFFFFFFFFFFFFFFFFFFFFFFFFFFFFFFFFFFFFFFFFFFFFFFFFFFFFFFFFFFFFFFFFFFFFFFFFFFFFFFFFFFFFFFFFFFFFFFFFFFFFFFFFFFFFFFFFFFFFFFFFFFFFFFFFFFFFFFFFFFFFFFFFFFFFFFFFFFFFFFFFFFFFFFFFFFFFFFFFFFFFFFFFFFFFFFFFFFFFFFFFFFFFFFFFEFEFEFEFEFEFEFEFEAF4A0E02050500072DE237AC61ACE8F8F9F2F5FEF9F0FDEFE7B43F0B08080940917B292E90F2FCF5FAFFFAF8F7F7FEF9DF3FC415000013A2FDFDF7F4FDFEFAE38B21050502538752130127DAFEFBFEFFFFFFFFFFFFFFFFFFFFFFFFFFFFFFFFFFFFFFFFFFFFFFFFFFFFFFFFFFFFFFFFFFFFFFFFFFFFFFFFFFFFFFFFFFFFFFFFFFFFFFFFFFFFFFFFFFFFFFFFFFFFFFFFFFFFFFFFFFFFFFFFFFFFFFFFFFFFFFFFFFFFFFFFFFFFFFFFFFFFFFFFFFFFFFFFFFFFFFFFFFFFFFFFFFFFFFFFFFFFFFFFFFFFFFFFFFFFFFFFFFFFFFFFFFFFFFFFFFFFFFFFFFFFFFFFFFFFFFFFFFFFFFFFFFFFFFFFFFFFFFFFFFFFFFFFFFFFFFFFFFFFFFFFFFFFFFFEFEFEFEFEFEFEFEFEA63905010202041692E9D2F2D2E7F9F7F7FBFBFCFDF6D54B342D0A0907276865231B020B31A7F1FBFBF5F7FDFDF9F7FAF4B54F07020012A0FCFCF8F7FEF2A63F1F040704091C657E300624C6E4FCFEFFFFFFFFFFFFFFFFFFFFFFFFFFFFFFFFFFFFFFFFFFFFFFFFFFFFFFFFFFFFFFFFFFFFFFFFFFFFFFFFFFFFFFFFFFFFFFFFFFFFFFFFFFFFFFFFFFFFFFFFFFFFFFFFFFFFFFFFFFFFFFFFFFFFFFFFFFFFFFFFFFFFFFFFFFFFFFFFFFFFFFFFFFFFFFFFFFFFFFFFFFFFFFFFFFFFFFFFFFFFFFFFFFFFFFFFFFFFFFFFFFFFFFFFFFFFFFFFFFFFFFFFFFFFFFFFFFFFFFFFFFFFFFFFFFFFFFFFFFFFFFFFFFFFFFFFFFFFFFFFFFFFFFFFFFFFFFFBF9FEF9F4FEF7F7BE541603030A0A071F4DCDFAFDF6FBFDF9FAFCF9FCFFF9CE991501020A050E447451010E65A8DCF2FCF8F8FBFEFDFBF9F8EB62A12706000FA0F8F6F9FBF8EFD869080B0A4B56235482330C0A3946E5FCFFFFF4F8FCFCFDFFFFFFFFFFFFFFFFFFFFFFFFFFFFFFFFFFFFFFFFFFFFFFFFFFFFFFFFFFFFFFFFFFFFFFFFFFFFFFFFFFFFFFFFFFFFFFFFFFFFFFFFFFFFFFFFFFFFFFFFFFFFFFFFFFFFFFFFFFFFFFFFFFFFFFFFFFFFFFFFFFFFFFFFFFFFFFFFFFFFFFFFFFFFFFFFFFFFFFFFFFFFFFFFFFFFFFFFFFFFFFFFFFFFFFFFFFFFFFFFFFFFFFFFFFFFFFFFFFFFFFFFFFFFFFFFFFFFFFFFFFFFFFFFFFFFFFFFFFFFFFFFFFFFFFFFFFFFFFFFFDF6EEFCF8F6FAB986652706072B594B71A3E7F6FAFBFCF6F9FDFDF9F8F6C663580B0A10080809132A28010C3652BDF0FCFCFBFCFDFDFDFBF4EDAA62931D0310A0FAFAFBFCFBF1C244085169553A18202E1305000215C1D3F7EBF1F6FAFDF8F4FFFFFFFFFFFFFFFFFFFFFFFFFFFFFFFFFFFFFFFFFFFFFFFFFFFFFFFFFFFFFFFFFFFFFFFFFFFFFFFFFFFFFFFFFFFFFFFFFFFFFFFFFFFFFFFFFFFFFFFFFFFFFFFFFFFFFFFFFFFFFFFFFFFFFFFFFFFFFFFFFFFFFFFFFFFFFFFFFFFFFFFFFFFFFFFFFFFFFFFFFFFFFFFFFFFFFFFFFFFFFFFFFFFFFFFFFFFFFFFFFFFFFFFFFFFFFFFFFFFFFFFFFFFFFFFFFFFFFFFFFFFFFFFFFFFFFFFFFFFFFFFFFFFFFFFFFFFFFFFCF5F1F8FEF9D07C7A541B05081C5A8E687CC7EBFAF7F1F4FBFAF9FCF6D76A970D041B5D24040704030A09102461CFF7EEF9FDF6F6FBFBF9F9F8F69D6612030FA0FBFCFAF6FDF8D0996C7CBD95521C020204010308104D6CD9EDF6F9FAFCF9F7FFFFFFFFFFFFFFFFFFFFFFFFFFFFFFFFFFFFFFFFFFFFFFFFFFFFFFFFFFFFFFFFFFFFFFFFFFFFFFFFFFFFFFFFFFFFFFFFFFFFFFFFFFFFFFFFFFFFFFFFFFFFFFFFFFFFFFFFFFFFFFFFFFFFFFFFFFFFFFFFFFFFFFFFFFFFFFFFFFFFFFFFFFFFFFFFFFFFFFFFFFFFFFFFFFFFFFFFFFFFFFFFFFFFFFFFFFFFFFFFFFFFFFFFFFFFFFFFFFFFFFFFFFFFFFFFFFFFFFFFFFFFFFFFFFFFFFFFFFFFFFFFFFFFFFFFFFFFFFFCF8FBF7FAFECF8E5D1D0303040A4698475E7DBAEAF5F3FAF9F5F5F7F9E094C92D202F63350901092F5E554A95D3E7EFEDF9FCF6F7FCFAF6F8FDF4DF30050112A3F8F8FBF7FCF9F8F9D8BCBE944F1702040C160807060169C4F3FCF9FAFCFAFBFFFFFFFFFFFFFFFFFFFFFFFFFFFFFFFFFFFFFFFFFFFFFFFFFFFFFFFFFFFFFFFFFFFFFFFFFFFFFFFFFFFFFFFFFFFFFFFFFFFFFFFFFFFFFFFFFFFFFFFFFFFFFFFFFFFFFFFFFFFFFFFFFFFFFFFFFFFFFFFFFFFFFFFFFFFFFFFFFFFFFFFFFFFFFFFFFFFFFFFFFFFFFFFFFFFFFFFFFFFFFFFFFFFFFFFFFFFFFFFFFFFFFFFFFFFFFFFFFFFFFFFFFFFFFFFFFFFFFFFFFFFFFFFFFFFFFFFFFFFFFFFFFFFFFFFFFFFFFFFFFDFCFCF5EDC49C581803030729677840618385CCF6FFFCF4F6FAF7FDF4C97E867D826F6728001863B6B586C1B1A89BC9F0FAFCFDFDFCF6F6F8F08A7316011391E8F8FFFEFCF9FEF9F5FDC098561802062550190007076482C5F2F9FBFCF9FAFFFFFFFFFFFFFFFFFFFFFFFFFFFFFFFFFFFFFFFFFFFFFFFFFFFFFFFFFFFFFFFFFFFFFFFFFFFFFFFFFFFFFFFFFFFFFFFFFFFFFFFFFFFFFFFFFFFFFFFFFFFFFFFFFFFFFFFFFFFFFFFFFFFFFFFFFFFFFFFFFFFFFFFFFFFFFFFFFFFFFFFFFFFFFFFFFFFFFFFFFFFFFFFFFFFFFFFFFFFFFFFFFFFFFFFFFFFFFFFFFFFFFFFFFFFFFFFFFFFFFFFFFFFFFFFFFFFFFFFFFFFFFFFFFFFFFFFFFFFFFFFFFFFFFFFFFFFFFFFDFBFDFEF5B88379602001020B3167766938C79FCFF3FAFAF9FAFCFDF2CB7D6E879AD0A27F2F021F5692BEA56B3148298ED2F0FAFCFAF8F9FCF2F4817D19000957C2F8FBFCFEFEFBFCF9FDE3E1922C01052A5A3417060901517CC1F1FDFBF8F9FFFFFFFFFFFFFFFFFFFFFFFFFFFFFFFFFFFFFFFFFFFFFFFFFFFFFFFFFFFFFFFFFFFFFFFFFFFFFFFFFFFFFFFFFFFFFFFFFFFFFFFFFFFFFFFFFFFFFFFFFFFFFFFFFFFFFFFFFFFFFFFFFFFFFFFFFFFFFFFFFFFFFFFFFFFFFFFFFFFFFFFFFFFFFFFFFFFFFFFFFFFFFFFFFFFFFFFFFFFFFFFFFFFFFFFFFFFFFFFFFFFFFFFFFFFFFFFFFFFFFFFFFFFFFFFFFFFFFFFFFFFFFFFFFFFFFFFFFFFFFFFFFFFFFFFFFFFFFFFBF5F7FBE37C7E611800020302022080CD9BE4EBFAF7F5FAFCF9FAFAD88C707F041977631B02071D33427FAD3A04061094A6CCF1FCFAF6F9F6F4F4E72A0301012997ECF6FAFDF9F5FCFEFAD5A3591A02000526695D0B1706786D80C0F0FDF8F2FFFFFFFFFFFFFFFFFFFFFFFFFFFFFFFFFFFFFFFFFFFFFFFFFFFFFFFFFFFFFFFFFFFFFFFFFFFFFFFFFFFFFFFFFFFFFFFFFFFFFFFFFFFFFFFFFFFFFFFFFFFFFFFFFFFFFFFFFFFFFFFFFFFFFFFFFFFFFFFFFFFFFFFFFFFFFFFFFFFFFFFFFFFFFFFFFFFFFFFFFFFFFFFFFFFFFFFFFFFFFFFFFFFFFFFFFFFFFFFFFFFFFFFFFFFFFFFFFFFFFFFFFFFFFFFFFFFFFFFFFFFFFFFFFFFFFFFFFFFFFFFFFFFFFFFFFFFFFFFEF9F2F5C1765F150100070A090D0D1D7ADFEBFDFCF8FCFBF5F7FDF3BF738003010503622B06041C49503E3D000C02397876B0ECFDFDFEF7EDFAF998741306083F88CDFBFEFBF2FAFDFEF0B970381105102D7EE19E31691502707476C3F6F2FFFFFFFFFFFFFFFFFFFFFFFFFFFFFFFFFFFFFFFFFFFFFFFFFFFFFFFFFFFFFFFFFFFFFFFFFFFFFFFFFFFFFFFFFFFFFFFFFFFFFFFFFFFFFFFFFFFFFFFFFFFFFFFFFFFFFFFFFFFFFFFFFFFFFFFFFFFFFFFFFFFFFFFFFFFFFFFFFFFFFFFFFFFFFFFFFFFFFFFFFFFFFFFFFFFFFFFFFFFFFFFFFFFFFFFFFFFFFFFFFFFFFFFFFFFFFFFFFFFFFFFFFFFFFFFFFFFFFFFFFFFFFFFFFFFFFFFFFFFFFFFEFEFEFEFEFEFEFEFEFAF5F9F0DD7F0E03000302040A1F342C5DD4FCFDF4F9FCFAFCFAF9D3887A011100000302090400070B0707080D0B126D6FB2EDFDF9F9FBF7FFFEF26BEA2C030C6183ADFAFEFCFDFCF9FBFDDBA6632120589DE4E4805A9D2108016F4A84C2AFF6FBF3FFFEFDFEFEFEFFFEFEFFFEFDFDFFFFFFFFFFFFFFFFFFFFFFFFFFFFFFFFFFFFFFFFFFFFFFFFFFFFFFFFFFFFFFFFFFFFFFFFFFFFFFFFFFFFFFFFFFFFFFFFFFFFFFFFFFFFFFFFFFFFFFFFFFFFFFFFFFFFFFFFFFFFFFFFFFFFFFFFFFFFFFFFFFFFFFFFFFFFFFFFFFFFFFFFFFFFFFFFFFFFFFFFFFFFFFFFFFFFFFFFFFFFFFFFFFFFFFFFFFFFFFFFFFFFFFFFFFFFFFFFFFFFFFFFFFFFFFFEFAFBFEFCFCFDFEFCFBF4CF8E66061103000002263442644079E7FEF9FDFEEDF4FDFCFBC3480A00010000000B040001031B413935756B6A5DB8EFFDFDFDFFFFFFFFFEFCFBFA3301000E90FCFCFDFBFBFEFDF1FCC6440D062D77BBEBF7FDEE130606020F42794E38B4DF9AFEFDF2FCFBFDFEFCFCFEFCF6F6FEFEFEFFFFFFFFFFFFFFFFFFFFFFFFFFFFFFFFFFFFFFFFFFFFFFFFFFFFFFFFFFFFFFFFFFFFFFFFFFFFFFFFFFFFFFFFFFFFFFFFFFFFFFFFFFFFFFFFFFFFFFFFFFFFFFFFFFFFFFFFFFFFFFFFFFFFFFFFFFFFFFFFFFFFFFFFFFFFFFFFFFFFFFFFFFFFFFFFFFFFFFFFFFFFFFFFFFFFFFFFFFFFFFFFFFFFFFFFFFFFFFFFFFFFFFFFFFFFFFFFFFFFFFFFFEFEFCFEFBF8FBFEFEF8ECCF4708640A100500002E7B7C270221C0FEFEF8FDFEFEFEFCFBD46101060100080200010B2945485A9EB2DCF9ECF6F9FDFEFCFCFFFFFFFFFEFCFBFA921E131D8CF6FBF8FEFDFEFDF4FCDD6115022D433485F7EEF8850E00070358530D2A79858AFEFBFDFEFEFEF9F7F9FDFEFEFEFDF4FEFFFFFFFFFFFFFFFFFFFFFFFFFFFFFFFFFFFFFFFFFFFFFFFFFFFFFFFFFFFFFFFFFFFFFFFFFFFFFFFFFFFFFFFFFFFFFFFFFFFFFFFFFFFFFFFFFFFFFFFFFFFFFFFFFFFFFFFFFFFFFFFFFFFFFFFFFFFFFFFFFFFFFFFFFFFFFFFFFFFFFFFFFFFFFFFFFFFFFFFFFFFFFFFFFFFFFFFFFFFFFFFFFFFFFFFFFFFFFFFFFFFFFFFFFFFFFFFFFFFFFFFFFFFFFCFEFEFAFBFDFFFDF9D758267D1A03010000153A4B2F156182E0F6F1FBF0F7F9C0B7B31A0200040107010501011E525A4991D4FBFEF9FDFEF2EEF7FEFFFFFFFFFEFCFCFB941F020A58C6FBF8FBFBFEFDCFB4893A140B314D28248BF2A6EE2D0D08020D0C01387139B3EDF8FEFEFDFEFEFEFEFCF9FCFFFDF8FEFFFFFFFFFFFFFFFFFFFFFFFFFFFFFFFFFFFFFFFFFFFFFFFFFFFFFFFFFFFFFFFFFFFFFFFFFFFFFFFFFFFFFFFFFFFFFFFFFFFFFFFFFFFFFFFFFFFFFFFFFFFFFFFFFFFFFFFFFFFFFFFFFFFFFFFFFFFFFFFFFFFFFFFFFFFFFFFFFFFFFFFFFFFFFFFFFFFFFFFFFFFFFFFFFFFFFFFFFFFFFFFFFFFFFFFFFFFFFFFFFFFFFFFFFFFFFFFFFFFFFFFFFEFAFAFEFFFCF8F6EEC7A8B414091705070A23462B090F0E17688BE2FBFCFEFBF7FAB97C5A03030D0C641005060204080F25485458DCB0D4F9FEFBFBFDFFFFFFFFFFFEFCFCFB9B61213E7AC1FEFAF6FEFEFDC692601C03030C0F296771842089170404010008090A140937ACFDFEF2FBFCFDFFFFFDFCFCFCFEFDFEFFFFFFFFFFFFFFFFFFFFFFFFFFFFFFFFFFFFFFFFFFFFFFFFFFFFFFFFFFFFFFFFFFFFFFFFFFFFFFFFFFFFFFFFFFFFFFFFFFFFFFFFFFFFFFFFFFFFFFFFFFFFFFFFFFFFFFFFFFFFFFFFFFFFFFFFFFFFFFFFFFFFFFFFFFFFFFFFFFFFFFFFFFFFFFFFFFFFFFFFFFFFFFFFFFFFFFFFFFFFFFFFFFFFFFFFFFFFFFFFFFFFFFFFFFFFFFFFFFFFFFFFFEFDFDF8FBFAFAFCF9CA7F6E01040B000203304A240306040F36B8A0EEA0DBF4FAF6E68C180A010A166B110301000506050C1513185F4DA2C59887C6EDF6FEFFFFFFFEFDFDFCDCD78CACD7E7FFFDFAF6FEFDB65B1F05080601003EB28B20091C1102010C0702080B0E2B604DFFFFFBFDFBF8FAFCFCFEFDFBF6FFFFFFFFFFFFFFFFFFFFFFFFFFFFFFFFFFFFFFFFFFFFFFFFFFFFFFFFFFFFFFFFFFFFFFFFFFFFFFFFFFFFFFFFFFFFFFFFFFFFFFFFFFFFFFFFFFFFFFFFFFFFFFFFFFFFFFFFFFFFFFFFFFFFFFFFFFFFFFFFFFFFFFFFFFFFFFFFFFFFFFFFFFFFFFFFFFFFFFFFFFFFFFFFFFFFFFFFFFFFFFFFFFFFFFFFFFFFFFFFFFFFFFFFFFFFFFFFFFFFFFFFFFFFFEF6FDFCFDFCF5E8C2AA9E2300010B01020010170E0703010415743BD5AB90C9F3FBC581490401020B22050000080602080B04080418105C4C2A5F86BBEEFEFEFFFFFEFDFDFCFAFDDCEAFBFAFAF7FDFBFEFBD38B411105022B435EAC53560E0405070205033D430406256D25F8F4FBFEFDFCFDFBF8FAFEFBF2FFFFFFFFFFFFFFFFFFFFFFFFFFFFFFFFFFFFFFFFFFFFFFFFFFFFFFFFFFFFFFFFFFFFFFFFFFFFFFFFFFFFFFFFFFFFFFFFFFFFFFFFFFFFFFFFFFFFFFFFFFFFFFFFFFFFFFFFFFFFFFFFFFFFFFFFFFFFFFFFFFFFFFFFFFFFFFFFFFFFFFFFFFFFFFFFFFFFFFFFFFFFFFFFFFFFFFFFFFFFFFFFFFFFFFFFFFFFFFFFFFFFFFFFFFFFFFFFFFFFFFFFFFFFFEF9FEF8FAF8F7F1C28B6C05010904244D141E353E3C220202111E08748879CCF6FCBFA38E00060D020B00134720071C35343A61270A502F47295068B4F4FDFEFFFFFFFEFEFDF4F9F8F5F8FCFCF9FEFFFEFAEDB255140302304842602C5F110D56154E6229484D0608133B6DB8F3D9B9E6F6F9FDF8F6FFFCF4FFFFFFFFFFFFFFFFFFFFFFFFFFFFFFFFFFFFFFFFFFFFFFFFFFFFFFFFFFFFFFFFFFFFFFFFFFFFFFFFFFFFFFFFFFFFFFFFFFFFFFFFFFFFFFFFFFFFFFFFFFFFFFFFFFFFFFFFFFFFFFFFFFFFFFFFFFFFFFFFFFFFFFFFFFFFFFFFFFFFFFFFFFFFFFFFFFFFFFFFFFFFFFFFFFFFFFFFFFFFFFFFFFFFFFFFFFFFFFFFFFFFFFFFFFFFFFFFFFFFFFFFFFFFFEFAFDF8FAFAFCF7BA581E01010B00305C3B615F47422B06020A0300222F61CEF5FAE5B46507040B040900195C25031D4F7997A4626A7D5F6E724D78CFFAFDFDFEFFFFFEFEFDF6F7FBF8F9FCFDFEFCFEFBFAF9CD60130204131B1619132712567366B59E371C1E0507071979459499A1C9C4D8F8FCFAFFFDF8FFFFFFFFFFFFFFFFFFFFFFFFFFFFFFFFFFFFFFFFFFFFFFFFFFFFFFFFFFFFFFFFFFFFFFFFFFFFFFFFFFFFFFFFFFFFFFFFFFFFFFFFFFFFFFFFFFFFFFFFFFFFFFFFFFFFFFFFFFFFFFFFFFFFFFFFFFFFFFFFFFFFFFFFFFFFFFFFFFFFFFFFFFFFFFFFFFFFFFFFFFFFFFFFFFFFFFFFFFFFFFFFFFFFFFFFFFFFFFFFFFFFFFFFFFFFFFFFFFFFFFFFFFFFFEF7FCFCF8FBFBEE951B00000426471C21326040191612080500040051182CB5F3FCD28F4D0803000503020A251303165890959D82894BA66E793D83DEFCFDFDFEFFFFFEFEFEFEFCFCFDFFFDFDFEFDFDF6FBFCCA793D1002011D3915000D5E7980869C5F1A0004050002112E0124477A92A0CDF5FEFDFAFBFAFFFFFFFFFFFFFFFFFFFFFFFFFFFFFFFFFFFFFFFFFFFFFFFFFFFFFFFFFFFFFFFFFFFFFFFFFFFFFFFFFFFFFFFFFFFFFFFFFFFFFFFFFFFFFFFFFFFFFFFFFFFFFFFFFFFFFFFFFFFFFFFFFFFFFFFFFFFFFFFFFFFFFFFFFFFFFFFFFFFFFFFFFFFFFFFFFFFFFFFFFFFFFFFFFFFFFFFFFFFFFFFFFFFFFFFFFFFFFFFFFFFFFFFFFFFFFFFFFFFFFFFFFFFFFEFBFCF8F4FAFAED840203040C34700C1F372B100200030909040102792D0A79DDCD91835E020400010303000005041E617E5C5E622B1186472F167DE1FDFFFFFFFFFFFFFFFFFFFFFFFFFFFFFFFFFFFDF7FEFBC1885A1E0201325E210313877689783F11020127340002080A03060C3178C0EAF9FEFCF9FBFCFEFEFFFFFFFFFFFFFFFFFFFFFFFFFFFFFFFFFFFFFFFFFFFFFFFFFFFFFFFFFFFFFFFFFFFFFFFFFFFFFFFFFFFFFFFFFFFFFFFFFFFFFFFFFFFFFFFFFFFFFFFFFFFFFFFFFFFFFFFFFFFFFFFFFFFFFFFFFFFFFFFFFFFFFFFFFFFFFFFFFFFFFFFFFFFFFFFFFFFFFFFFFFFFFFFFFFFFFFFFFFFFFFFFFFFFFFFFFFFFFFFFFFFFFFFFFFFFFFFFFFFFFFFFFFFFFBF3FAFDFBF1880205080D0D1C0340642708040006090D351203211705208C817C792303000000000601000303133D6453161C04022216030166CBF5FFFFFFFFFFFFFFFFFFFFFFFFFFFFFFFFFFFDF7FFFFDE91370E0403172307050B29BD44832C02253F65752D010212060201298FE3FCFBFBFAFDFAF6FEFEFFFFFFFFFFFFFFFFFFFFFFFFFFFFFFFFFFFFFFFFFFFFFFFFFFFFFFFFFFFFFFFFFFFFFFFFFFFFFFFFFFFFFFFFFFFFFFFFFFFFFFFFFFFFFFFFFFFFFFFFFFFFFFFFFFFFFFFFFFFFFFFFFFFFFFFFFFFFFFFFFFFFFFFFFFFFFFFFFFFFFFFFFFFFFFFFFFFFFFFFFFFFFFFFFFFFFFFFFFFFFFFFFFFFFFFFFFFFFFFFFFFFFFFFFFFFFFFFFFFFFFFFFFFFF4FAFDF9F6DC7304040000070D020C1F180601040300227C3002020502002678671806040002000412050409090307181F162136340E2B2F30469CE9FFFFFFFFFFFFFFFFFFFFFFFFFFFFFFFFFEFAECFDE8D29B2802040200020500010FC42E49393497B9827D780100000307041E6ED3FBFBF7F6FBF2EAFFFFFFFFFFFFFFFFFFFFFFFFFFFFFFFFFFFFFFFFFFFFFFFFFFFFFFFFFFFFFFFFFFFFFFFFFFFFFFFFFFFFFFFFFFFFFFFFFFFFFFFFFFFFFFFFFFFFFFFFFFFFFFFFFFFFFFFFFFFFFFFFFFFFFFFFFFFFFFFFFFFFFFFFFFFFFFFFFFFFFFFFFFFFFFFFFFFFFFFFFFFFFFFFFFFFFFFFFFFFFFFFFFFFFFFFFFFFFFFFFFFFFFFFFFFFFFFFFFFFFFFFFFFFFEFCFBF7F8FAE87E161011000209030100000300070700031422040003190901021B18090E010504010E1F0D06010001010A182A8AC7E94079C6D9D3E2F7FFFFFFFFFFFFFFFFFFFFFFFFFFFFFFFFFDF9FEF8AAA07E210A0B040007110200011F208ECFD3E299838F24040201121D1029A4DFEDF8FEFBF9F8FDFDFEFFFFFFFFFFFFFFFFFFFFFFFFFFFFFFFFFFFFFFFFFFFFFFFFFFFFFFFFFFFFFFFFFFFFFFFFFFFFFFFFFFFFFFFFFFFFFFFFFFFFFFFFFFFFFFFFFFFFFFFFFFFFFFFFFFFFFFFFFFFFFFFFFFFFFFFFFFFFFFFFFFFFFFFFFFFFFFFFFFFFFFFFFFFFFFFFFFFFFFFFFFFFFFFFFFFFFFFFFFFFFFFFFFFFFFFFFFFFFFFFFFFFFFFFFFFFFFFFFFFFFFFFFEF2FDFEFDFA80787A0200010B031307020000000507010208000A15149C4E1A0902040600010A190120B8230000050A11346B7487FEF65B98EAFBF7FBFFFFFFFFFFFFFFFFFFFFFFFFFFFFFFFFFFFFFEFCF8EF962805000E100601030C010F02C047C36F9DDC8561860002022BAD1A0F1D65BFEEF3F4FEFBF3F3FEFFFFFFFFFFFFFFFFFFFFFFFFFFFFFFFFFFFFFFFFFFFFFFFFFFFFFFFFFFFFFFFFFFFFFFFFFFFFFFFFFFFFFFFFFFFFFFFFFFFFFFFFFFFFFFFFFFFFFFFFFFFFFFFFFFFFFFFFFFFFFFFFFFFFFFFFFFFFFFFFFFFFFFFFFFFFFFFFFFFFFFFFFFFFFFFFFFFFFFFFFFFFFFFFFFFFFFFFFFFFFFFFFFFFFFFFFFFFFFFFFFFFFFFFFFFFFFFFFFFFFFFFFEFEFEF6F0F2E3E993010300031C00060105070502090D02060010C61ABCEDD0550A0003050308010C00150A0601002C735F2890F0FDFCF7ECF3FBFEFEFEFFFFFFFFFFFFFFFFFFFFFFFFFFFFFFFFFFFCF1FD9176802B08547D4002000017C023D2F4F8FB7E1D0D0E13030F010D200D0234B1B185B0F1FBF3F8FEFEFFFFFFFFFFFFFFFFFFFFFFFFFFFFFFFFFFFFFFFFFFFFFFFFFFFFFFFFFFFFFFFFFFFFFFFFFFFFFFFFFFFFFFFFFFFFFFFFFFFFFFFFFFFFFFFFFFFFFFFFFFFFFFFFFFFFFFFFFFFFFFFFFFFFFFFFFFFFFFFFFFFFFFFFFFFFFFFFFFFFFFFFFFFFFFFFFFFFFFFFFFFFFFFFFFFFFFFFFFFFFFFFFFFFFFFFFFFFFFFFFFFFFFFFFFFFFFFFFFFFFFFFFEF8F3F1FAFEF98E1402070B49B604040D034B75380003050001102007C469C3DA711500010508010D000106020003071241808795F3FCF4F9FDF3F9FCFCFFFFFFFFFFFFFFFFFFFFFFFFFFFFFFFFFFFEFDFBF09D2E050013190E08010B1614BEF65D211B0C02040705030801060200080C245CABE9FEFCF8FBFCFEFFFFFFFFFFFFFFFFFFFFFFFFFFFFFFFFFFFFFFFFFFFFFFFFFFFFFFFFFFFFFFFFFFFFFFFFFFFFFFFFFFFFFFFFFFFFFFFFFFFFFFFFFFFFFFFFFFFFFFFFFFFFFFFFFFFFFFFFFFFFFFFFFFFFFFFFFFFFFFFFFFFFFFFFFFFFFFFFFFFFFFFFFFFFFFFFFFFFFFFFFFFFFFFFFFFFFFFFFFFFFFFFFFFFFFFFFFFFFFFFFFFFFFFFFFFFFFFFFFFFFFFFFEFDFDFAF2F982160804040C0D210A02040311180B000100010702020C229D6839190403020201030200050102000000060F151C31ACFEFD7F9EE5FAFCFEFFFFFFFFFFFFFFFFFFFFFFFFFFFFFFFFFEFCF8FD987A7822000306050305209C33242E10050201040409478E0208000F010503031257BEFAFDFDFDFBFEFFFFFFFFFFFFFFFFFFFFFFFFFFFFFFFFFFFFFFFFFFFFFFFFFFFFFFFFFFFFFFFFFFFFFFFFFFFFFFFFFFFFFFFFFFFFFFFFFFFFFFFFFFFFFFFFFFFFFFFFFFFFFFFFFFFFFFFFFFFFFFFFFFFFFFFFFFFFFFFFFFFFFFFFFFFFFFFFFFFFFFFFFFFFFFFFFFFFFFFFFFFFFFFFFFFFFFFFFFFFFFFFFFFFFFFFFFFFFFFFFFFFFFFFFFFFFFFFFFFFFFFFFFFFEAF8FB82825C010100020301020247665B663506080000040000030128168057020003050004000602227E3501000304030B082939BE669EDCF6FCFBFBFFFBF3F6FEF0FDFFFFE68EB0FCFBFCF9FCF8FEE1913A0B0101000304010C3E13AD2A0502016C5854BDCCB02A9E2291220802050541B3FAFCFEFDF5FAFEFAFDF7FAFFFFFFFAF4FFFFFEFBF5FCFFFFFEFEFBFAFCF9FEFFFFFFFFFFFFFEFDFDFFFFFFFFFFFFFFFFFFFFFFFFFFFFFFFFFFFFFFFFFFFFFFFFFFFFFFFFFFFFFFFFFFFFFFFFFFFFFFFFFFFFFFFFFFFFFFFFFFFFFFFFFFFFFFFFFFFFFFFFFFFFFFFFFFFFFFFFFFFFFFFFFFFFFFFFFFFFFFFFFFFFFFFFFFFFFFFFFFFFFFFFFFFFFFFFFFFFFEEEF4F7FCCC77260004090401092FA7BBD4D6D0C0B1AF913302060009070209311F020101062F900203000C3013020403010103030701321D2F47514E98B192CB7FF5E8F960D57ABC7291FBFDFCF9FDF6B2958864632500000203090F0802451302030790CFD3F7F4E7C0E99A48386E24020A1A346ECAFAF3FBFDF2FAFEFAF6FAFBF9FBF9FBFEFBFBFAFEFFFCFBFBFCF8F3FCFEFDFDFDFEFFFFFEFEFEFFFFFFFFFFFFFFFFFFFFFFFFFFFFFFFFFFFFFFFFFFFFFFFFFFFFFFFFFFFFFFFFFFFFFFFFFFFFFFFFFFFFFFFFFFFFFFFFFFFFFFFFFFFFFFFFFFFFFFFFFFFFFFFFFFFFFFFFFFFFFFFFFFFFFFFFFFFFFFFFFFFFFFFFFFFFFFFFFFFFFFFFFFFFFFFFFFFFFEF6FDFCEC8B2D0502070403060E41E3F3F9FDF9F5F4F7ECBB9C8F2A07040806040605000207194D01040008080100040400000201030B000C02040703293654A8255C655D13C847442B345B5F625B66989F312F2E3215010104060506011303020276A2D3F2F6FCFAFAF0D36EA1925917000409277BDCFCF6FFFFF2FBF8FAF9F9FBFCFEFBF1FCFAF9E284CCFAF7F7FBF8F9FEFBFDFDFDFDFDFEFFFFFFFFFFFFFFFFFFFFFFFFFFFFFFFFFFFFFFFFFFFFFFFFFFFFFFFFFFFFFFFFFFFFFFFFFFFFFFFFFFFFFFFFFFFFFFFFFFFFFFFFFFFFFFFFFFFFFFFFFFFFFFFFFFFFFFFFFFFFFFFFFFFFFFFFFFFFFFFFFFFFFFFFFFFFFFFFFFFFFFFFFFFFFFFFFFFFFFFFFFFEFDF8F2962C060503030B375E0746E6FDF5F9F9F7F9F7F3FBE9878A2B05050E0904040304060A020805000505050806050500010E54390D090505030A0D0924480E010D0007592003090D0B0A08010E2B520F020006040302305E2300050D070502516E7278B3F9FDFAF5DEA6F5E284260806051B496D84C0F7F9FAFCF1F8F6F7F9F6F8F8F6FCEB7FD13853707FC3FDF9F9FEFBFEFEFDFDFDFEFFFFFFFFFFFFFFFFFFFFFFFFFFFFFFFFFFFFFFFFFFFFFFFFFFFFFFFFFFFFFFFFFFFFFFFFFFFFFFFFFFFFFFFFFFFFFFFFFFFFFFFFFFFFFFFFFFFFFFFFFFFFFFFFFFFFFFFFFFFFFFFFFFFFFFFFFFFFFFFFFFFFFFFFFFFFFFFFFFFFFFFFFFFFFFFFFFFFFFFFFFFEFDF99D8633030204040A32570927A3F1F8FAF9FAFCF6F1F5EC957E2907025C374A470F042166050D0C001B5525080407090000104E330B0A0507023E430F070A0E0306070403020E504646724509080808080200060B08022A532000050D0605010207084CB8F4FDF9F1FBFDFCFCD697825D22020102134F7ECEA7DBA6CAABB9BBA6CB9BE289E1206C3152315EC8F8F7F6F4FDFFFFFEFDFDFEFFFFFFFFFFFFFFFFFFFFFFFFFFFFFFFFFFFFFFFFFFFFFFFFFFFFFFFFFFFFFFFFFFFFFFFFFFFFFFFFFFFFFFFFFFFFFFFFFFFFFFFFFFFFFFFFFFFFFFFFFFFFFFFFFFFFFFFFFFFFFFFFFFFFFFFFFFFFFFFFFFFFFFFFFFFFFFFFFFFFFFFFFFFFFFFFFFFFFFFFFFFFFFF88E8532000104070303050B1A89ECFAFAFCFEFCF9F6EFECA21F04054F847EC1BE4D0B2D72175B100120672905040406030104010701582B423E59493768716B5A111261235441957D7EC37C5330530B040500000A070101020200020F000100090B047DC3B6FAF2EFFCFCFCFEF3F8E4902F040300010A036E2DB0525D38474A376036C60A8A13011C5F3A5D9EC7F9FAF1FFFFFFFFFFFFFFFEFEFEFFFFFFFFFFFFFFFFFFFFFFFFFFFFFFFFFFFFFFFFFFFFFFFFFFFFFFFFFFFFFFFFFFFFFFFFFFFFFFFFFFFFFFFFFFFFFFFFFFFFFFFFFFFFFFFFFFFFFFFFFFFFFFFFFFFFFFFFFFFFFFFFFFFFFFFFFFFFFFFFFFFFFFFFFFFFFFFFFFFFFFFFFFFFFFFFFFFFFEFCEDB02D0101040907050300053AABF1F6F7FBFEFDFAE5C3C03C040C0DB887BFFAFDBA5F4503208A0506030000070B0A090903010D422BB4649EAB7961A0F5FBEAA26567E084D5A2B6ACACBE9CCCA0C86741120704030202040506000100030B020B0A05505659D0ADF8FBF6FAF7F7F9D1915421040207080014076F310503020205042685040701030302030C45AEF3F2FBFBFCFCFDFEFFFFFEFCFDFFFFFFFFFFFFFFFFFFFFFFFFFFFFFFFFFFFFFFFFFFFFFFFFFFFFFFFFFFFFFFFFFFFFFFFFFFFFFFFFFFFFFFFFFFFFFFFFFFFFFFFFFFFFFFFFFFFFFFFFFFFFFFFFFFFFFFFFFFFFFFFFFFFFFFFFFFFFFFFFFFFFFFFFFFFFFFFFFFFFFFFFFFFFFFFFFFFFFEF9B62D020402092B31010002185C85E8FDFDFCFDFCFAD68481340A4550C6F5F7FAFDE1B780020706040E070101001835280808083ACA94B58DB8DFEFF9F9F5F8D875F6FBAEE3EDEBCDD8F7F9F8FBFBF5F9C25A34130D0701060A080000020D4919030D0B03043A6629BBF3F5F3EDEEF8F7FADB7D1F080C0000000B03050105060608060A0509000108050B05166FB3ADCBFAF1F4F6FAFDFEFFFEFCFDFFFFFFFFFFFFFFFFFFFFFFFFFFFFFFFFFFFFFFFFFFFFFFFFFFFFFFFFFFFFFFFFFFFFFFFFFFFFFFFFFFFFFFFFFFFFFFFFFFFFFFFFFFFFFFFFFFFFFFFFFFFFFFFFFFFFFFFFFFFFFFFFFFFFFFFFFFFFFFFFFFFFFFFFFFFFFFFFFFFFFFFEFEFFFEFDFFFEFCF228010002020D56720401063CC6F9F8FCFAFCFBFAFDF0B95A110285DEFAFFFDFCFCAC41310200000C0400000408295F4F01011366B2B2BAE2D1DED0D3D4C9E6BBEAB6E95F9BBAA16C83D0F8F3ECEFFAF1F7F8973E25100019321601010B3CA04012030000030706042AC1EFBEF4FEF6F6F9FAD4761B00020000070601000204253307051109000501113819000000258CDAD1DAD8F6FAFDFEFDFCFDFFFFFFFFFFFFFFFFFFFFFFFFFFFFFFFFFFFFFFFFFFFFFFFFFFFFFFFFFFFFFFFFFFFFFFFFFFFFFFFFFFFFFFFFFFFFFFFFFFFFFFFFFFFFFFFFFFFFFFFFFFFFFFFFFFFFFFFFFFFFFFFFFFFFFFFFFFFFFFFFFFFFFFFFFFFFFFFFFFFFFEFBFCFEFAF6FEFCF1C1540D0201040401020208080F79F1F5FBFEFDF7F5F7CF8290380D26BCF3FDFCFCFCFAE99400000000100200072D200102050200010D0408753E674C595F51AB259D07B3BE22050369565CC1B9B8B4D9F8D4CEA7956C29054B96510D0010AC561439170403000001031262A062E0F7E3D4D9DADAC364121306073D1E09030A095678030816400300051785410405011E606C6C7F6AE9ECF3F8FAFDFEFFFFFFFFFFFFFFFFFFFFFFFFFFFFFFFFFFFFFFFFFFFFFFFFFFFFFFFFFFFFFFFFFFFFFFFFFFFFFFFFFFFFFFFFFFFFFFFFFFFFFFFFFFFFFFFFFFFFFFFFFFFFFFFFFFFFFFFFFFFFFFFFFFFFFFFFFFFFFFFFFFFFFFFFFFFFFFFFFFFFFAFAFCFBFAFCFCEDC343CB2F09050000000104134E66BEF8F6F8FDF9F9F1CCBBB81B122BA4F1EEF4FCFEFEFEAF1304070008010000326F450A060000030801201E0F0E060005053D870C000000010002120C03010001020D2BC68D72BA7624041B21378634000328E571AB442D230102000A429020AFBABD927E9271558FB6913B3D49D7593B18020E120C0F001BB001050801000000080C0401146B4E5EBDC2BED8F4FDFDFFFFFFFFFFFFFFFFFFFFFFFFFFFFFFFFFFFFFFFFFFFFFFFFFFFFFFFFFFFFFFFFFFFFFFFFFFFFFFFFFFFFFFFFFFFFFFFFFFFFFFFFFFFFFFFFFFFFFFFFFFFFFFFFFFFFFFFFFFFFFFFFFFFFFFFFFFFFFFFFFFFFFFFFFFFFFFFFFFFFFBFAFCFDFAF7FCE348C022010007080206132263D0C45FEEFCFCFDFEF2BC5D1006011BDD60B8F1FEFEFDFCF7E08E040D000E19140401120D0001051321212B0A794E2C1D1005050F0E01090F0006060100000A0A000209010C06626E1D0401004E7C3105060006A9BBDF7348B076000101020B06070A1209399B843E50A1E4EDEFE2B7C067A0580F0000010300031D0C01090200090500051C2F0B05160706061E7DDAFDFEFFFFFFFFFFFFFFFFFFFFFFFFFFFFFFFFFFFFFFFFFFFFFFFFFFFFFFFFFFFFFFFFFFFFFFFFFFFFFFFFFFFFFFFFFFFFFFFFFFFFFFFFFFFFFFFFFFFFFFFFFFFFFFFFFFFFFFFFFFFFFFFFFFFFFFFFFFFFFFFFFFFFFFFFFFFFFFFFFFFFFEFEFBFAFEFAF1EEC0201F090100000305074DD17AF9EDEEF0FBF4F8EBCD913308143DC6EAF7F6FCF1F1FAFEC23B00010716A421000609191A06001A685387935D3DA19C511E0D05194A01010C0100080303030303000100050B0A1B0A000F100D0D090B010000231BA1D5BCDE792230081201010000010207111C42606D8A6F7584271C7753180507242E020204074912020E341A0002255E5A0503282A04175199D8F3F9FFFFFFFFFFFFFFFFFFFFFFFFFFFFFFFFFFFFFFFFFFFFFFFFFFFFFFFFFFFFFFFFFFFFFFFFFFFFFFFFFFFFFFFFFFFFFFFFFFFFFFFFFFFFFFFFFFFFFFFFFFFFFFFFFFFFFFFFFFFFFFFFFFFFFFFFFFFFFFFFFFFFFFFFFFFFFFFFFFFFFFF9F7FCF9F4FCF069100702341F02010B2E79C0DCFEF6FAFCF4F8FDCF6E210C10368DF3E7FAFDF6F8F9F9FADF7E00031042270600042C60430304091D439AD0D3D4EDEFE0871D052D82523D3F11311005032A310C070107010C0003010003232D11001233081203011B907F9682736508050600000019060002040C173A6A4A2777280D1A0E0200014E700708000BA3110017763901050B1211020C5D640F052B599BDDFCFFFFFFFFFFFFFFFFFFFFFFFFFFFFFFFFFFFFFFFFFFFFFFFFFFFFFFFFFFFFFFFFFFFFFFFFFFFFFFFFFFFFFFFFFFFFFFFFFFFFFFFFFFFFFFFFFFFFFFFFFFFFFFFFFFFFFFFFFFFFFFFFFFFFFFFFFFFFFFFFFFFFFFFFFFFFFFFFFFFFFDFBFAFAF9FBFDF5940E0200713A0B090A408C9CFAF5FBFDFEDDD5D3BB7A28001B717FDBF1F4FCF5F6FBFBF8B22502010F96180505040E1C130001020D56669EC4E8FAF1F58F1103206FCBD0C560B16B454A6A674740230F1E0604000A090245532111286D1D09060002227758686C1C02020802020207020B0600060A0C1B1606221A050100100900162206020212330648192116050204050404031A23050368B2DBF5FEFFFFFFFFFFFFFFFFFFFFFFFFFFFFFFFFFFFFFFFFFFFFFFFFFFFFFFFFFFFFFFFFFFFFFFFFFFFFFFFFFFFFFFFFFFFFFFFFFFFFFFFFFFFFFFFFFFFFFFFFFFFFFFFFFFFFFFFFFFFFFFFFFFFFFFFFFFFFFFFFFFFFFFFFFFFFFFFFFFFFFEFAFAFDFDFBF4E5450700083016040A0A174EA9FDF9F9FDF7CE9F957941130F3F75B3EBFBFBFEFBFBF1F3DF7B0501050D3A0B020201030A0C06010607201B679BA2B0AEB5933C145CCCF9FEFACDF1D8C7C9A69CBFBC947795452D060404011A29180913290D0102004E19281920210204030200020202030505020203040A01000A060000080A0505050303184606050A7C1F000610252202030B1E330400047CDCFCF9F8FFFFFFFFFFFFFFFFFFFFFFFFFFFFFFFFFFFFFFFFFFFFFFFFFFFFFFFFFFFFFFFFFFFFFFFFFFFFFFFFFFFFFFFFFFFFFFFFFFFFFFFFFFFFFFFFFFFFFFFFFFFFFFFFFFFFFFFFFFFFFFFFFFFFFFFFFFFFFFFFFFFFFFFFFFFFFFFFFFFFFFFAFAFFFCE9BDAB0A020406030A0603081F5BA0F8FDF5FDD8C697725324021B66B8EDFCFEF3FCFFFFF7F8AD453D00030B000601000B08060703000D2D2606236A5A435C56665749BABFB3F8FEF7FDFAFBFDDFDAF8F8E4D9E9BE77221627200C070701050202103C4C7159491C00000103010100000701132F160504030807060005000205062A230B040404237B0202053A0F020315463F02070F444F03264195E4FFFBFAFFFFFFFFFFFFFFFFFFFFFFFFFFFFFFFFFFFFFFFFFFFFFFFFFFFFFFFFFFFFFFFFFFFFFFFFFFFFFFFFFFFFFFFFFFFFFFFFFFFFFFFFFFFFFFFFFFFFFFFFFFFFFFFFFFFFFFFFFFFFFFFFFFFFFFFFFFFFFFFFFFFFFEFDFCFBFDFEFFFDFDFAFBFFFBDF8C43020003000209050004328FC9F2FCF9FEC8A57D472717163A7AE3FFFEFEF4FAFDFFF8D86E3065050505000802000907010003051250490508281E09463B292B489E567CB3FBF8C3E9FAF9F9F6F8FDFCFAF2FCB1777C665E350C02010402030F6BBA7FACBC4B02020000041C380205002C5322071B2E2401050100050205044C3F0E0003174241060500050202010B262C1D241F453C1E61A5D2F3FDFDFEFFFFFFFFFFFFFFFFFFFFFFFFFFFFFFFFFFFFFFFFFFFFFFFFFFFFFFFFFFFFFFFFFFFFFFFFFFFFFFFFFFFFFFFFFFFFFFFFFFFFFFFFFFFFFFFFFFFFFFFFFFFFFFFFFFFFFFFFFFFFFFFFFFFFFFFFFFFFFFFFFFFFFEFDFCFBFDFEFEFCFDFCF9FAFDE2800100060002032B1B062059A9F5F5F5FEF4D380411606174D91B7E9FDF5FDFDFCF7FBF5BC531F3905070F0006020A29170101080D0F302D000B0404002F2B090019341B7949C0C18BB4E9F5F3F4F6FAFDFCF0FAD9D3E0A79D884410060407040476C0ABC3BE4B040200030F2B660200011F32150522433306000634190201033227071C2929650604050C360D0000071F36548274716D6FA1E4FBFDFCFDFFFFFFFFFFFFFFFFFFFFFFFFFFFFFFFFFFFFFFFFFFFFFFFFFFFFFFFFFFFFFFFFFFFFFFFFFFFFFFFFFFFFFFFFFFFFFFFFFFFFFFFFFFFFFFFFFFFFFFFFFFFFFFFFFFFFFFFFFFFFFFFFFFFFFFFFFFFFFFFFFFFFFCF9FFFEFAFAFEFEFBFDFDF7F4FBDC5D000709020501563A09439BCBF2F8F4FDF0D1591302163C82DDFCF7FCF1F9FEFFF5F7F9C15204010202150304001956300403060A060004341B080502020100010202144D084B527A5D9FC8DEF7FBF9FCFBFDF9F6F9FADBD4C8863912020B08067857D1846F32030001080E1441000203000203021131381A030B6B581908050607043F5A1B451230001271210200144C7394E6DDD3D7CFDDF6FDFAFBFEFFFFFFFFFFFFFFFFFFFFFFFFFFFFFFFFFFFFFFFFFFFFFFFFFFFFFFFFFFFFFFFFFFFFFFFFFFFFFFFFFFFFFFFFFFFFFFFFFFFFFFFFFFFFFFFFFFFFFFFFFFFFFFFFFFFFFFFFFFFFFFFFFFFFFFFFFFFFFFFFFFFFF9F1FFFEF7F6FDFFF9FAFAF9E8CDA71C060E29120502372A1C5AB6F9FCF7F9F1D2992301033773A0DAF9F6FEF9F6F6FBFEFAE3932F02070D250006000216372109050D1D19040B733A0B000218190A03132C090004000746162F569AD0D0CECCCEDEFBFEF9FDE2E4B3845D1D0D22130A4D0AAD395D390503030600010B000502000202030A2D4D2D05093E8432100A022522303E0102197805014D1D0401185AA0CEF6FAFFFEFAFCFEFDF9FAFDFDFFFFFFFFFFFFFFFFFFFFFFFFFFFFFFFFFFFFFFFFFFFFFFFFFFFFFFFFFFFFFFFFFFFFFFFFFFFFFFFFFFFFFFFFFFFFFFFFFFFFFFFFFFFFFFFFFFFFFFFFFFFFFFFFFFFFFFFFFFFFFFFFFFFFFFFFFFFFFFFFFFFCF8FDFCFEFCFBFBFBFAFBFECD7142060F02802A0504060A509BC2FDFAFBF2D159390E020D5DB5D1DCF6FAFCF1FBECF6FAFDB13E0205091C7902000307040200000324554E0B0A2A201002004A4D1F0A347405020C000003060300284E4E4C576C7BD1D8CEE5AFAD8956392641661B0C0006491F3A200306020B1F0002000D1E07050B1D141635240107034C2606000052510F000303114F0D00020007131F6CD5FAF1F2FDFEF9FEFEFBFBFDFBFBFFFFFFFFFFFFFFFFFFFFFFFFFFFFFFFFFFFFFFFFFFFFFFFFFFFFFFFFFFFFFFFFFFFFFFFFFFFFFFFFFFFFFFFFFFFFFFFFFFFFFFFFFFFFFFFFFFFFFFFFFFFFFFFFFFFFFFFFFFFFFFFFFFFFFFFFFFFFFFFFFFFEFDF9F6FCFEF3F2FAFDFAFBDC7A13020106590E0000011147ACFCFBF7FDF4A31101030F469EE5FBF9FFFEFCFDF6F9FDF4E9C45D0A0401145E04010E0200050300010A2432000503250F03012B2E10011C3E0005010500021C0E0208090210415F274F5361907382783E183D7C611800030701150B121B14113E7A020600347E1D01295A3A00050D00000879460B03032F320E020A15000106000B13256465A0FAFCFFFDF9FEFDFEFDF8FCFFFEFEFFFFFFFFFFFFFFFFFFFFFFFFFFFFFFFFFFFFFFFFFFFFFFFFFFFFFFFFFFFFFFFFFFFFFFFFFFFFFFFFFFFFFFFFFFFFFFFFFFFFFFFFFFFFFFFFFFFFFFFFFFFFFFFFFFFFFFFFFFFFFFFFFFFFFFFFFFFFFFFFFFFCF9FDF9F5FEFFFCF4F7F2EDC45603100A00000905011071B3CAF2FAFDF2C33A0C0E0A34A2E8F8FFFFFEFDFDFCFEE8F6FD93A47806030802000106000002040F0B060C05000100007552100B05070D0808100000011503117F3808050D10051E37000A374E3D2724295B6A403B08030E160E088B4F79AD62488A3E01080015410F0721372405000001000C3B100115130500040D408902090004187945ADD4D7F5EFFEFEFEFDFDFDFFFDFEFFFBFBFFFFFFFFFFFFFFFFFFFFFFFFFFFFFFFFFFFFFFFFFFFFFFFFFFFFFFFFFFFFFFFFFFFFFFFFFFFFFFFFFFFFFFFFFFFFFFFFFFFFFFFFFFFFFFFFFFFFFFFFFFFFFFFFFFFFFFFFFFFFFFFFFFFFFFFFFFFFFFFFFFFBF4FEFDFBFDFEE3ADCAAA5E2B0103000200000C02007C759BFEFEFEF6C53E00004F6A68C2FDFDFEFEFDF8FEF4FDFEFCFDC46C2604020600010D060005072C734A02020604100B07381D05050262620B000208010900000036170012130105050000022D380B0E07001D351B000321B209A826379379423179B60D0300000701000000000005070104000000100200000103020016410F00000097D953A8FEFCFCFDF9FAFDF7F5FDFEFEFEFEFAF9FFFFFFFFFFFFFFFFFFFFFFFFFFFFFFFFFFFFFFFFFFFFFFFFFFFFFFFFFFFFFFFFFFFFFFFFFFFFFFFFFFFFFFFFFFFFFFFFFFFFFFFFFFFFFFFFFFFFFFFFFFFFFFFFFFFFFFFFFFFFFFFFFFFFFFFFFFFFFFFFFFFEFCF6F8FDF3B15D201C22130303060000030006017F6D86D5F8FEF2BA390600030A54A8E7FAFAFDFCFDFCFEFDF5FDFEF9F7DC66030102060103169A1E000E2F20020006110100010602050502262C0B020726B3000007040D0B6A554D6905000105040309054E642D0403022E9022490C42110938220D0C1F3601020C000201000040693D06000100000000010001000403020708041AADB4C4E4F8CAE2FEFCFCFDFBFDFDFDFCFDFEFCFBFDFEFEFFFFFFFFFFFFFFFFFFFFFFFFFFFFFFFFFFFFFFFFFFFFFFFFFFFFFFFFFFFFFFFFFFFFFFFFFFFFFFFFFFFFFFFFFFFFFFFFFFFFFFFFFFFFFFFFFFFFFFFFFFFFFFFFFFFFFFFFFFFFFFFFFFFEFEFFFAF6F9FBFBFBF9FCF8FBA98B641F0505060500040701029AAAACD9C9D9F2F5FEEFA515022FB0ABC7EAFBFEFEFEFEFEFEFFFEFAFCF8FBFDA4330300030B0212053F0F0E07060702020618635D00020A0300060805020201114F0011A12009053D848C7D630103010C0002031F2710010000103611090C0C02000A0601040706010401000400050121291507000017773901793778355A5510012FA0B8F1F5F1FAFDF7FBFEFEFEFEFEFEFEFEFEFEFEFEFEFEFFFFFFFFFFFFFFFFFFFFFFFFFFFFFFFFFFFFFFFFFFFFFFFFFFFFFFFFFFFFFFFFFFFFFFFFFFFFFFFFFFFFFFFFFFFFFFFFFFFFFFFFFFFFFFFFFFFFFFFFFFFFFFFFFFFFFFFFFFFFFFFFFFFFFFFEFBFDFDFCFDFEFDFCFCF1B4A58E3D2E0F03020609000101109E58DCF3F9F5F8FCF6E4AF4803001B3E8BDEFBFFFFFFFFFFFFFFFEFDFBFDFDF89834090203030200040510030301030300000210382B00764E130102024A64160A040D0805410F02050D383B27300200030501466E290A04010102060803000201000B050002030307309D000007020104070704032476996A882CA746A05394BB9EA7B5EDEEFEFEF8FDFCFEFFFFFFFFFFFFFFFFFFFFFFFFFFFFFFFFFFFFFFFFFFFFFFFFFFFFFFFFFFFFFFFFFFFFFFFFFFFFFFFFFFFFFFFFFFFFFFFFFFFFFFFFFFFFFFFFFFFFFFFFFFFFFFFFFFFFFFFFFFFFFFFFFFFFFFFFFFFFFFFFFFFFFFFFFFFFFFFFFFFFFEFDFEFEF8F5F9FDFDEEAA778F440E0D06010005072E781F8BEBAEDEFAFFFEFBFBFCF1AA0D020A3091BDECFDFFFFFFFFFFFFFFFEFEFDF1FBF4CD560704010307000002020B0507041F554F0B040B09004420070001002B3A0B03030E030012040505070F0C09090209060001302E100301010101237022020704000406070308416B1A3E070007010703020302010D4CB8B4E0A6E3B9DEBACEE9ECF3F3FAEEF9FEFDFEFEFEFFFFFFFFFFFFFFFFFFFFFFFFFFFFFFFFFFFFFFFFFFFFFFFFFFFFFFFFFFFFFFFFFFFFFFFFFFFFFFFFFFFFFFFFFFFFFFFFFFFFFFFFFFFFFFFFFFFFFFFFFFFFFFFFFFFFFFFFFFFFFFFFFFFFFFFFFFFFFFFFFFFFFFFFFFFFFFFFFFFFFFFCFBFDF9F7FCFCF9E99A323C14020904060401347C73866978E6F1FCFBF5F7F3B56A4C0001073DD6EFF9FDFFFFFFFFFFFFFFFFFEEEF593CD9F34040301030B011193190702040117362D0C110001090C0A0D0601040C0A070D01188A080400020D5840020502010D3C6F070C0C0502040202041950120103020B6F350A05042E542E73060B030207020002031C5DA8DFEEFDEEFCF1F9F4F3F9FEFDFAF6F5FCFFFFFCFDFEFFFFFFFFFFFFFFFFFFFFFFFFFFFFFFFFFFFFFFFFFFFFFFFFFFFFFFFFFFFFFFFFFFFFFFFFFFFFFFFFFFFFFFFFFFFFFFFFFFFFFFFFFFFFFFFFFFFFFFFFFFFFFFFFFFFFFFFFFFFFFFFFFFFFFFFFFFFFFFFFFFFFFFFFFFFFFFFFFFFFFFFBF6FBFDFDFDECBA724708050405053D3847473F7A88642315E8F9FBFBF8F7F4C1590A000524A6F1FAFBFDFFFFFFFFFFFFFFFFFEF3F4A8828B450807040106080B6E14010003030709040A0E01090606090A08020304040909030F68070000050A4235050204070A2E5405014440433C1A0203060F000800010A53250402000914215701070702022B4E435593CCF8FBFCF6FEFEEFF9F7F9FAF7F9FAF8FDFEF9FDFCFDFEFFFFFFFFFFFFFFFFFFFFFFFFFFFFFFFFFFFFFFFFFFFFFFFFFFFFFFFFFFFFFFFFFFFFFFFFFFFFFFFFFFFFFFFFFFFFFFFFFFFFFFFFFFFFFFFFFFFFFFFFFFFFFFFFFFFFFFFFFFFFFFFFFFFFFFFFFFFFFFFFFFFFFFFFFFFFFFFFFFFFFDFCF8FBFCFAFAE08E394404060009026C998F5B1C315E190A02E6F9FBF7FBF8BE804E030A2A7694F3F9FAFDFFFFFFFFFFFFFFFEF7FEEFF19F874108062A5601070707020B0403050503184336090D41410E090602010206090003070B0106541E05070A0600020604093C48376D7A82571A0202010103120001020C040202040700050C0602060701274859A2F1F7F5FDF8F7F5F9F3FAF7FAFAF7F8FAF7FDF6F6FDFEFCFEFFFFFFFFFFFFFFFFFFFFFFFFFFFFFFFFFFFFFFFFFFFFFFFFFFFFFFFFFFFFFFFFFFFFFFFFFFFFFFFFFFFFFFFFFFFFFFFFFFFFFFFFFFFFFFFFFFFFFFFFFFFFFFFFFFFFFFFFFFFFFFFFFFFFFFFFFFFFFFFFFFFFFFFFFFFFFFFFFFFFF7FCFEFEFAF8FBDF893F5605072D60657E95641E030409195869EDFDFDF6F4C2823F0602154FC67FF1F7FBFEFFFFFFFFFFFFFFFDF5FEF6F8F26C0E0200316C06000802000E07000203011D4D400E0854490D0B0B342F0F060705030A0500087423050402060701082B4A61B8815553552B0402040410535D180200080403020507000309080300223E294784CAF8F7F3FBF1FEF4F8FDFCF7F9FBFCFDFAF8FCF3FAFEFFF9FCFFFFFFFFFFFFFFFFFFFFFFFFFFFFFFFFFFFFFFFFFFFFFFFFFFFFFFFFFFFFFFFFFFFFFFFFFFFFFFFFFFFFFFFFFFFFFFFFFFFFFFFFFFFFFFFFFFFFFFFFFFFFFFFFFFFFFFFFFFFFFFFFFFFFFFFFFFFFFFFFFFFFFFFFFFFFFFFFFFFFF5FCFAFCFCFAF3BD520B0B02277A9BA26421050203070A1791ABEEFCFDFCF29E805E0408278AECF8FCF8FBFDFFFFFFFFFFFFFFFDF9FEFDF7FB690A030202050B030003000103040402020504020703070102010D584C1625391005090102041409050602070C030E4C936AA767220A0501000002071A8587270401371B0401042338030501020341AA97C8F8F7F5FBFDFBF3FBFEFDFDFBF3F6F8FBFEFDFEFDF9FEF8FDF9FCFFFFFFFFFFFFFFFFFFFFFFFFFFFFFFFFFFFFFFFFFFFFFFFFFFFFFFFFFFFFFFFFFFFFFFFFFFFFFFFFFFFFFFFFFFFFFFFFFFFFFFFFFFFFFFFFFFFFFFFFFFFFFFFFFFFFFFFFFFFFFFFFFFFFFFFFFFFFFFFFFFFFFFFFFFFFFFFFFFFFFEFEF5F8FDEFC57722020F05577116020100020908000D320A05A0BAEEFAF6CE5A040E2663CCFEEFF8FBFEFFFFFFFFFFFFFFFFFDF5FBF8F9D44E04000303000307020B330D0014301C03060202000000040A04010706024778130200063F0A0335170504082027010D4E7E1D031E2C100000152910070D0B0000097C3D0802094E7C040204134570E0FBFCFEFFFFFDFCFDFDF7FEFCFCFEFFFFFEFEFDFCFEFCF9FEF1FCFFFFFFFFFFFFFFFFFFFFFFFFFFFFFFFFFFFFFFFFFFFFFFFFFFFFFFFFFFFFFFFFFFFFFFFFFFFFFFFFFFFFFFFFFFFFFFFFFFFFFFFFFFFFFFFFFFFFFFFFFFFFFFFFFFFFFFFFFFFFFFFFFFFFFFFFFFFFFFFFFFFFFFFFFFFFFFFFFFFFFFFFFCF9FDFCFCDB957168010201192507000100020B0D01299916010104BDDCE56603040D4C9D66FEF6FDF5FBFFFFFFFFFFFFFFFFFEFBF5F5FB7E1F1A1803000306010D4AA4200948926512000103020203050E030101010303010506000CBC180A933F04041A637804030B0005015F65230201458330010C040000361301000101020102142E39EDF9FCFBFDFFFFFFFFFFFFFFFFFFFFFFFFFFFFFFFFFFFFFFFFFFFFFFFFFFFFFFFFFFFFFFFFFFFFFFFFFFFFFFFFFFFFFFFFFFFFFFFFFFFFFFFFFFFFFFFFFFFFFFFFFFFFFFFFFFFFFFFFFFFFFFFFFFFFFFFFFFFFFFFFFFFFFFFFFFFFFFFFFFFFFFFFFFFFFFFFFFFFFFFFFFFFFFFFFFFFFFFFFFFFFFFFFFFFFFFFFFE9F6FFF8F8DF7A08090D088185160403000A0B0100081E05000000A262A6620607000F5EEFFEFEFEEFFAFFFFFFFFFFFFFFFFFDEFFDFDFDFB62786A020308000537B11E0B2B857F83590108100003230D02000105070602030903010B0E010101040204050106130B00020003000D160E080602070207010302983D070207060001083FEB4ADFF9FCFEFEFFFFFFFFFFFFFFFFFFFFFFFFFFFFFFFFFFFFFFFFFFFFFFFFFFFFFFFFFFFFFFFFFFFFFFFFFFFFFFFFFFFFFFFFFFFFFFFFFFFFFFFFFFFFFFFFFFFFFFFFFFFFFFFFFFFFFFFFFFFFFFFFFFFFFFFFFFFFFFFFFFFFFFFFFFFFFFFFFFFFFFFFFFFFFFFFFFFFFFFFFFFFFFFFFFFFFFFFFFFFFFFFFFFFFFFFFEF7FBFEF0EDDE84140F010F7C7114000705060D0A0036CD3307050002A0F16E0609041774FDF8F9FDFEFEFEFFFFFFFFFFFFFFFFFEFEFBFCF75A0A0D06080C000000000000001E7F5A03080F02070DA36A0F000904000B0E02071A050000050A0C030302060D060718070000060168752C000100040C010F000C0304050207090B161C59F4F4FDFEEFF7FEFFFFFFFFFFFFFFFFFFFFFFFFFFFFFFFFFFFFFFFFFFFFFFFFFFFFFFFFFFFFFFFFFFFFFFFFFFFFFFFFFFFFFFFFFFFFFFFFFFFFFFFFFFFFFFFFFFFFFFFFFFFFFFFFFFFFFFFFFFFFFFFFFFFFFFFFFFFFFFFFFFFFFFFFFFFFFFFFFFFFFFFFFFFFFFFFFFFFFFFFFFFFFFFFFFFFFFFFFFFFFFFFFFFFFFFFF4FBFAF8FAFEFDF4CE740B0005071204000C0D030A11022BCE141406010F8C5F63711261F2F7FDFCFCFEFBFDFFFFFFFFFFFFFFFFFEEEFBF85B2A1600000000020D00080009070F0B03000204020B071003220B050404649635000105000600000003100B00061F5A91080000000406131B110503050C010007001C0D050C080A6BF6F6F6F7FDF6F5FBFCFEFFFFFFFFFFFFFFFFFFFFFFFFFFFFFFFFFFFFFFFFFFFFFFFFFFFFFFFFFFFFFFFFFFFFFFFFFFFFFFFFFFFFFFFFFFFFFFFFFFFFFFFFFFFFFFFFFFFFFFFFFFFFFFFFFFFFFFFFFFFFFFFFFFFFFFFFFFFFFFFFFFFFFFFFFFFFFFFFFFFFFFFFFFFFFFFFFFFFFFFFFFFFFFFFFFFFFFFFFFFFFFFFFFFFFFFFFDFEEFF6FEFAD6812E1F01000D010105025C732D498E1B06191BB62D09010F151F1A0A5BFEFDFAFDF9F6FBFEFFFFFFFFFFFFFFFFFEF8FDFEE7DC630A0E0800010100181103000100030308357C758BE4E2D875866D270D5D8D1B0615030D030A021186620003080A12000509130801576B2A04062D993C0000180794578F7674EFFBF5F9FCEBFCFDFEFBFDFFFFFFFFFFFFFFFFFFFFFFFFFFFFFFFFFFFFFFFFFFFFFFFFFFFFFFFFFFFFFFFFFFFFFFFFFFFFFFFFFFFFFFFFFFFFFFFFFFFFFFFFFFFFFFFFFFFFFFFFFFFFFFFFFFFFFFFFFFFFFFFFFFFFFFFFFFFFFFFFFFFFFFFFFFFFFFFFFFFFFFFFFFFFFFFFFFFFFFFFFFFFFFFFFFFFFFFFFFFFFFFFFFFFFFFFF4FBFBFDFCFAF5E1B86D040003010C04031A1C09152B0F030FB242140300000204061F78F4F6FDFDF6FBFEFFFFFFFFFFFFFFFFFFFEFEEDF8FD91351002010000062E9F140009071611022C7B9597CE6B99EAE6889362141C2E0522A632CC340800061D17080603070104015B7F1D0021210C04010E2D0C0016A3372A9F7B89E4F8F2FCFDF9FAFEFBEFF8FEFFFFFFFFFFFFFFFFFFFFFFFFFFFFFFFFFFFFFFFFFFFFFFFFFFFFFFFFFFFFFFFFFFFFFFFFFFFFFFFFFFFFFFFFFFFFFFFFFFFFFFFFFFFFFFFFFFFFFFFFFFFFFFFFFFFFFFFFFFFFFFFFFFFFFFFFFFFFFFFFFFFFFFFFFFFFFFFFFFFFFFFFFFFFFFFFFFFFFFFFFFFFFFFFFFFFFFFFFFFFFFFFFFFFFFFFFBFDF8F3F9FDFCFAE896160600000A020504060A366D3609043D09992B020302000D5BD5F8F1FEF7F7FDF9FCFFFFFFFFFFFFFFFDF3F9F8F7FCCFAB5A0A020409020E3119030323654A2D83DFDCBB831C97EFF3DB7C2F0D406E121A62C9C8A937140600080A040E170802081F270E0607030101010104011206321112B2DAE1FDFBF9FBFEFBF5FEF5F7F5FAFFFFFFFFFFFFFFFFFFFFFFFFFFFFFFFFFFFFFFFFFFFFFFFFFFFFFFFFFFFFFFFFFFFFFFFFFFFFFFFFFFFFFFFFFFFFFFFFFFFFFFFFFFFFFFFFFFFFFFFFFFFFFFFFFFFFFFFFFFFFFFFFFFFFFFFFFFFFFFFFFFFFFFFFFFFFFFFFFFFFFFFFFFFFFFFFFFFFFFFFFFFFFFFFFFFFFFFFFFFFFFFFFFFFFFFFFBFBFAFAF8FAFDFAF0C36515080002100816140518597C150002034513010100192764F9F8F8F3FBFCF8FAFDFFFFFFFFFFFFFFFDF7FBEDF8FAA7663005010404000306090704102A46789FB7E99E7C4578DFF6F9902B0427420822A5ED6D4B8A2C07030303114763040612150C221D0A0C12110E0500040706020084CEEAFAFDFCFEFBF7FEFAF5F8FDFBFDFFFFFFFFFFFFFFFFFFFFFFFFFFFFFFFFFFFFFFFFFFFFFFFFFFFFFFFFFFFFFFFFFFFFFFFFFFFFFFFFFFFFFFFFFFFFFFFFFFFFFFFFFFFFFFFFFFFFFFFFFFFFFFFFFFFFFFFFFFFFFFFFFFFFFFFFFFFFFFFFFFFFFFFFFFFFFFFFFFFFFFFFFFFFFFFFFFFFFFFFFFFFFFFFFFFFFFFFFFFFFFFFFFFFFFFFFCFCFAFCFCF9FBFDEAB27C4C0A00050A04445F2E09223D0B02000201050100064E6D7CFDFAFCF2FEFCF9FDFFFFFFFFFFFFFFFFFEFDFEF4FCF9CBA6540502040900030D1D0902030D4395C3D0ECC37B272EAFF2FCC8570A355809116FF38F0F4B16080600051227370509537C629370375B726D73360505000A03069D92C1F9FCFCFDFAF9FEFBF8FDFCFDFEFFFFFFFFFFFFFFFFFFFFFFFFFFFFFFFFFFFFFFFFFFFFFFFFFFFFFFFFFFFFFFFFFFFFFFFFFFFFFFFFFFFFFFFFFFFFFFFFFFFFFFFFFFFFFFFFFFFFFFFFFFFFFFFFFFFFFFFFFFFFFFFFFFFFFFFFFFFFFFFFFFFFFFFFFFFFFFFFFFFFFFFFFFFFFFFFFFFFFFFFFFFFFFFFFFFFFFFFFFFFFFFFFFFFFFFFFFFDF8FBFFFCFCFDF0C68F430306050805336E571707010003120205050602053883C5FBFDFDF9FEF6FBF9FCFFFFFFFFFFFFFFFEFEFBF2F8FAFDFC82040408112D011569190301195993C2F5F8E76A081F85D1F6B9510E284F21001FCA850E040807060204070405061B76BAC5BF9E9AC5D4D3D797370C010501054186C2D5E4FBF8FCFEFEFCFEFFFBFCFEFFFFFFFFFFFFFFFFFFFFFFFFFFFFFFFFFFFFFFFFFFFFFFFFFFFFFFFFFFFFFFFFFFFFFFFFFFFFFFFFFFFFFFFFFFFFFFFFFFFFFFFFFFFFFFFFFFFFFFFFFFFFFFFFFFFFFFFFFFFFFFFFFFFFFFFFFFFFFFFFFFFFFFFFFFFFFFFFFFFFFFFFFFFFFFFFFFFFFFFFFFFFFFFFFFFFFFFFFFFFFFFFFFFFFFFFFFFCFAFDFFFFFEFCFAF4C55F161A170D05072D38130101000113030506070213166DF7F7FDFDF7FEF7FCF7FBFFFFFFFFFFFFFFFEFBFBF2F8FDFDFA9519050C1975010F591505011863A5CBF6FDE87C265F6FB2E9C0540C0329520D0C597411010803051C1A020404073B80A8EECFC6EAFBFBFBF7E7941821080301048C987BB5F5F9FDFCFFFEFDFDFCFDFEFFFFFFFFFFFFFFFFFFFFFFFFFFFFFFFFFFFFFFFFFFFFFFFFFFFFFFFFFFFFFFFFFFFFFFFFFFFFFFFFFFFFFFFFFFFFFFFFFFFFFFFFFFFFFFFFFFFFFFFFFFFFFFFFFFFFFFFFFFFFFFFFFFFFFFFFFFFFFFFFFFFFFFFFFFFFFFFFFFFFFFFFFFFFFFFFFFFFFFFFFFFFFFFFFFFFFFFFFFFFFFFFFFFFFFFFFEFEFEFEFCFCFDFBFCFAF0CB523C4013010102020203000305280200081A21605B80FBF7FBFDF5FBFEFEFCFDFFFFFFFFFFFFFFFEFAFBFAFCF6F8E8AC4B080A17610503030303011566C8F6FBFDF2BD879184C4EEFC9426124066150009660F04031C1942410501050A3175A5F1FDF8FBFDFDF9F7FFD83C88131F0D07412E48A2DFF9FCF9FCFEFBF9FDFFFFFFFFFFFFFFFFFFFFFFFFFFFFFFFFFFFFFFFFFFFFFFFFFFFFFFFFFFFFFFFFFFFFFFFFFFFFFFFFFFFFFFFFFFFFFFFFFFFFFFFFFFFFFFFFFFFFFFFFFFFFFFFFFFFFFFFFFFFFFFFFFFFFFFFFFFFFFFFFFFFFFFFFFFFFFFFFFFFFFFFFFFFFFFFFFFFFFFFFFFFFFFFFFFFFFFFFFFFFFFFFFFFFFFFFFFFFFDFDFEFCF8F9FCFBFDFAFAF68743411703020201010303060D75260919678ACBC8BEFAF1DEFDFEF9FAFEFCFEFFFFFFFFFFFFFFFFFEF8F7F9EFF6C1945B09040404070703050204146FD7FBF7FBFBF5FCACCBF5FBFCC35436694F08042C02020305563E393807030E171C2CA6F1FAF5F7FBFCFAF9FCE796DE2E6520040302468DB8F6FDFAFAFEFBF8FDFDFEFFFFFFFFFFFFFFFFFFFFFFFFFFFFFFFFFFFFFFFFFFFFFFFFFFFFFFFFFFFFFFFFFFFFFFFFFFFFFFFFFFFFFFFFFFFFFFFFFFFFFFFFFFFFFFFFFFFFFFFFFFFFFFFFFFFFFFFFFFFFFFFFFFFFFFFFFFFFFFFFFFFFFFFFFFFFFFFFFFFFFFFFFFFFFFFFFFFFFFFFFFFFFFFFFFFFFFFFFFFFFFFFFFFFFFFFFCFCFBFBFBFCFEFDFDFBF0F4901A0F0A05010304020814030A7E8C356ED4FBF9EADFDCC6A4F7E7FDF2FCF7FBFFFFFFFFFFFFFFFEFEFEFCFAFDFAA86D440301020B0006080103111D66CBFAF7F4F9F9FAF7FCFCF3F8C85538836B060E77070701075530060700020F3459035ED7F6FDF8FAFEFDFAF7F8EA958A9E20010A0802329CEAF9FEFEFFFFF8FDFAFCFFFFFFFFFFFFFFFFFFFFFFFFFFFFFFFFFFFFFFFFFFFFFFFFFFFFFFFFFFFFFFFFFFFFFFFFFFFFFFFFFFFFFFFFFFFFFFFFFFFFFFFFFFFFFFFFFFFFFFFFFFFFFFFFFFFFFFFFFFFFFFFFFFFFFFFFFFFFFFFFFFFFFFFFFFFFFFFFFFFFFFFFFFFFFFFFFFFFFFFFFFFFFFFFFFFFFFFFFFFFFFFFFFFFFFFFFFFFF4F9FFFCFCFCF9F9F5F8BD453D3F19091015152A4F1D1773E4A0E3F7FBF0B07E7F7793D297FBF9FDF7FBFFFFFFFFFFFFFFFEFCFEF4EFF9F7C994510900000403160300001B5E99CBEDF2FDFAEDFBFAFBF4FAFBC45F1B3358101082040104000A07020100010B31590E33AEE3F9FDFDFCF9F5F7FDEA6AF88A1B03010A09357DA0DDFDFBFEFEF9FDFBFDFFFFFFFFFFFFFFFFFFFFFFFFFFFFFFFFFFFFFFFFFFFFFFFFFFFFFFFFFFFFFFFFFFFFFFFFFFFFFFFFFFFFFFFFFFFFFFFFFFFFFFFFFFFFFFFFFFFFFFFFFFFFFFFFFFFFFFFFFFFFFFFFFFFFFFFFFFFFFFFFFFFFFFFFFFFFFFFFFFFFFFFFFFFFFFFFFFFFFFFFFFFFFFFFFFFFFFFFFFFFFFFFFFFFFFFFFEFEF3F9FDEBE8FEF9EFFBEAFED6CCCA554075887881998082FDFEFEFEEEF3AA764B072B6F868DFBFEFEFEFFFFFFFFFFFFFFFFFFFEFEF8FDF4F8BE47000102000418930F0802256696CAF9FDFAFFFEFDE8FAFEFDF3BF59407C7C740F7B05080300000C100702090F0800002D9FE8FBF6F6FBF9F5FCF7EEFDFEFD510B05020607338CB1FBF8FAFEF9FDFFFFFFFFFFFFFFFFFFFFFFFFFFFFFFFFFFFFFFFFFFFFFFFFFFFFFFFFFFFFFFFFFFFFFFFFFFFFFFFFFFFFFFFFFFFFFFFFFFFFFFFFFFFFFFFFFFFFFFFFFFFFFFFFFFFFFFFFFFFFFFFFFFFFFFFFFFFFFFFFFFFFFFFFFFFFFFFFFFFFFFFFFFFFFFFFFFFFFFFFFFFFFFFFFFFFFFFFFFFFFFFFFFFFFFFFFFFFF4FCF4F9E9BCBFF2FFFEFAFAFEEFFCC57BA4E8FEFDD99DEAFFFFFFF9FDFFFFAD683F0B031242FFFFFAFAFDFFFFFFFFFFFFFFFFFFFEF6FDFF99C7B747070402040008721B6D23022F9EC0B9FDFAFFFEF4FEF3FAF8FCC05132699DCA0F800400000A0C04040403060000000081C1B5DFFAFCFBFEF8F1FFFEFBFEFF2E0B05000607368C9DFFFFFFFFF9FDFFFFFFFFFFFFFFFFFFFFFFFFFFFFFFFFFFFFFFFFFFFFFFFFFFFFFFFFFFFFFFFFFFFFFFFFFFFFFFFFFFFFFFFFFFFFFFFFFFFFFFFFFFFFFFFFFFFFFFFFFFFFFFFFFFFFFFFFFFFFFFFFFFFFFFFFFFFFFFFFFFFFFFFFFFFFFFFFFFFFFFFFFFFFFFFFFFFFFFFFFFFFFFFFFFFFFFFFFFFFFFFFFFFFFFFFFFFFFEFEF8F7EAB87B6A86B1F1FDFCFCFAB3BCE8FDFEFEFCF7FDFFFEFAFDFBFDFEF7B67E5C275139FEF9FBFDFAFDFFFFFFFFFFFFFFFFFEFDFDFEA276824C0500000000001097721A01155393C5F4FAFEFCFBFEFAF4F9FDFAC68DAAD581791000000201000201000409000500072571A6DFFAFEFEFEFEFDFFFEFEFE7A7B2200060603082A7A95FEF9FEF9F2FAFFFFFFFFFFFFFFFFFFFFFFFFFFFFFFFFFFFFFFFFFFFFFFFFFFFFFFFFFFFFFFFFFFFFFFFFFFFFFFFFFFFFFFFFFFFFFFFFFFFFFFFFFFFFFFFFFFFFFFFFFFFFFFFFFFFFFFFFFFFFFFFFFFFFFFFFFFFFFFFFFFFFFFFFFFFFFFFFFFFFFFFFFFFFFFFFFFFFFFFFFFFFFFFFFFFFFFFFFFFFFFFFFFFFFFFFFFF9FCFDFDE596381B498FABF1FCF6FCF3F5FCFFFFFFFFFFFFFFFEFEFDF9F7FCFEECAA6F9CD7A0F6FDFEFDFDFEFFFFFFFFFFFFFFFEFDFEFCF8EEAC834508040004040302562807000C3093F2FDFBF7FDF8F8FDFDFEF9F5F2ECBF5B2D5E010C01000A050143450A3D6A02113093B0E5FAF6F9FFFFFFFFFFFFFFFF98D9B3340100070B0735937DFDEDFAFDF1F8FFFFFFFFFFFFFFFFFFFFFFFFFFFFFFFFFFFFFFFFFFFFFFFFFFFFFFFFFFFFFFFFFFFFFFFFFFFFFFFFFFFFFFFFFFFFFFFFFFFFFFFFFFFFFFFFFFFFFFFFFFFFFFFFFFFFFFFFFFFFFFFFFFFFFFFFFFFFFFFFFFFFFFFFFFFFFFFFFFFFFFFFFFFFFFFFFFFFFFFFFFFFFFFFFFFFFFFFFFFFFFFFFFFFFFFFFDFEFAFAE9B48081ACC3B2FBFEF1F6FDF8FBFFFFFFFFFFFFFFFEFBFAFCFDFBF8F7CBA9EDF9E6F8FEFFFEFCFEFFFFFFFFFFFFFFFEFAF5FCFDA885480F03040001050000251A62284D99CAF0F9FDFBFDFDFDFEFEFEFBFAC486866D0B1408030C73240912453D072C4A0004157AAAD7F6FCFDFFFFFFFFFFFFFFFFE9EBE4AE3E070803072675B1F6F9FEEDF5FDFFFFFFFFFFFFFFFFFFFFFFFFFFFFFFFFFFFFFFFFFFFFFFFFFFFFFFFFFFFFFFFFFFFFFFFFFFFFFFFFFFFFFFFFFFFFFFFFFFFFFFFFFFFFFFFFFFFFFFFFFFFFFFFFFFFFFFFFFFFFFFFFFFFFFFFFFFFFFFFFFFFFFFFFFFFFFFFFFFFFFFFFFFFFFFFFFFFFFFFFFFFFFFFFFFFFFFFFFFFFFFFFFFFFFFFFFCFEFCF9F6ECE0E3EFF0E0FFFFFAF9F8F9FCFFFFFFFFFFFFFFFEFDF3F6F9FBFCF9EDE9FAF2FCFDFEF8FDFCFEFFFFFFFFFFFFFFFEFAFAFDFEB54F100000052B66010002080A5282BDE6F4F8F8F7FDF9FFFFF5F9FDFEFECE693D560A07030008642457371516070D17040013588DADE8FAFBFFFFFFFFFFFFFFFFF4F7F99F2E0506070B0B3093FBF6FDF8F9FCFFFFFFFFFFFFFFFFFFFFFFFFFFFFFFFFFFFFFFFFFFFFFFFFFFFFFFFFFFFFFFFFFFFFFFFFFFFFFFFFFFFFFFFFFFFFFFFFFFFFFFFFFFFFFFFFFFFFFFFFFFFFFFFFFFFFFFFFFFFFFFFFFFFFFFFFFFFFFFFFFFFFFFFFFFFFFFFFFFFFFFFFFFFFFFFFFFFFFFFFFFFFFFFFFFFFFFFFFFFFFFFFFFFFFFFFFDFEFEFBFCFFFDF7F3FBF8F6FAFEFDF4FAFEFFFFFFFFFFFFFFFFFFF9FAF8FCFFF6F5F9F2F8FDF4FEF6F8FCFEFFFFFFFFFFFFFFFEFCECF9F9A075340503062956050101021B8AD5EFFAFAF7F9F2FEFCFEFAF5FBFEFEF9C459111C04040E00021E0F522E103E3E010302000B34417FDEFCFDFFFFFFFFFFFFFFFFF9F6FCAD823E04060A296C4CF7FAF7F7F7FCFFFFFFFFFFFFFFFFFFFFFFFFFFFFFFFFFFFFFFFFFFFFFFFFFFFFFFFFFFFFFFFFFFFFFFFFFFFFFFFFFFFFFFFFFFFFFFFFFFFFFFFFFFFFFFFFFFFFFFFFFFFFFFFFFFFFFFFFFFFFFFFFFFFFFFFFFFFFFFFFFFFFFFFFFFFFFFFFFFFFFFFFFFFFFFFFFFFFFFFFFFFFFFFFFFFFFFFFFFFFFFFFFFFFFFFFFFFCFBFCFEFCFAF9F6FDFDF4FBF9FDF4F8FDFFFFFFFFFFFFFFFEF2FCFDFAFCFBF6CFBBF4F5FDF5FEF9FAFDFFFFFFFFFFFFFFFFFDF7EEFBAA8985410802020D180A0400175D83D7FAFCF8F1F4F7FBFEF8FAFAFCFBFCFAA42E0406030D6F100100021A141340420002000401100D387FBCEAFFFFFFFFFFFFFFFFFAF6FCE3A07A440409276637FAFEFAF9F8FCFFFFFFFFFFFFFFFFFFFFFFFFFFFFFFFFFFFFFFFFFFFFFFFFFFFFFFFFFFFFFFFFFFFFFFFFFFFFFFFFFFFFFFFFFFFFFFFFFFFFFFFFFFFFFFFFFFFFFFFFFFFFFFFFFFFFFFFFFFFFFFFFFFFFFFFFFFFFFFFFFFFFFFFFFFFFFFFFFFFFFFFFFFFFFFFFFFFFFFFFFFFFFFFFFFFFFFFFFFFFFFFFFFFFFFFFFDF7F9FDFFFBF9FDFFFFFEFAFEF5FCF9F8FCFFFFFFFFFFFFFFFEF3FCFBFBFAF7F8A470DBB8F9FCFDFDFEFDFEFFFFFFFFFFFFFFFDF6FBF8958450190400011D54060200167189D8FBF7F7F6F4F6FAFEF4FCFEF6F6FCF8911902060A127711020004060708151A0801010601000610379CE8FFFFFFFFFFFFFFFFFAFBF8FCB07C470606091E22A4B7EBFCFAFCFFFFFFFFFFFFFFFFFFFFFFFFFFFFFFFFFFFFFFFFFFFFFFFFFFFFFFFFFFFFFFFFFFFFFFFFFFFFFFFFFFFFFFFFFFFFFFFFFFFFFFFFFFFFFFFFFFFFFFFFFFFFFFFFFFFFFFFFFFFFFFFFFFFFFFFFFFFFFFFFFFFFFFFFFFFFFFFFFFFFFFFFFFFFFFFFFFFFFFFFFFFFFFFFFFFFFFFFFFFFFFFFFFFFFFFFFDF8FCFDFBFAFBFCFCF9FFFEF8F9FAFDFCFDFFFFFFFFFFFFFFFEFEFEFCF6BD9FA463378068F7F7FAFEFAFBFDFFFFFFFFFFFFFFFEFCFBF7A0531A020001022B780100060245BEF5F8F5F6FAF9F5FEFEF5FEFEF2F7F7CE67130303175026010202060E3A2801060A030207253B0F2E4FA2E9FFFFFFFFFFFFFFFFFFF8F5FCE6995E3801041C4D2752DBFAFCFDFFFFFFFFFFFFFFFFFFFFFFFFFFFFFFFFFFFFFFFFFFFFFFFFFFFFFFFFFFFFFFFFFFFFFFFFFFFFFFFFFFFFFFFFFFFFFFFFFFFFFFFFFFFFFFFFFFFFFFFFFFFFFFFFFFFFFFFFFFFFFFFFFFFFFFFFFFFFFFFFFFFFFFFFFFFFFFFFFFFFFFFFFFFFFFFFFFFFFFFFFFFFFFFFFFFFFFFFFFFFFFFFFFFFFFFFFFFFFFFCF9FAFFFDFAF6FFFEF4FDFAFDFFFFFFFFFFFFFFFFFFFEFAFDFEF190583B120A1C33B5BAF8FEF4FBFEFEFEFEFFFEFDFDFDFEF7C4822C0C040102010B1D0501143A2BC1FBFBFBF7F4FCFCFBFAF9FEFCF8FCF5B44A0A05031B810B030004040B634A02040305020A45854E899AA9DDFCFCFFFEFEFFFFFFFCF4FBFFFBE9D883010A34820126ADEEFCFEFFFFFFFFFEFEFCFFFFFFFFFFFFFFFFFFFFFFFFFFFFFFFFFFFFFFFFFFFFFFFFFFFFFFFFFFFFFFFFFFFFFFFFFFFFFFFFFFFFFFFFFFFFFFFFFFFFFFFFFFFFFFFFFFFFFFFFFFFFFFFFFFFFFFFFFFFFFFFFFFFFFFFFFFFFFFFFFFFFFFFFFFFFFFFFFFFFFFFFFFFFFFFFFFFFFFFFFFFFFFFFFFFFFFFFFFFFFFFFFFFFFFFFFFFFFFFFFFFFFFFFFFFFFFFFFFFFFFFFFFFFFEF6FBFCF7E5D2670100030D297CFAFBF7FBFCFDFBF9FBFCFBFAFAFEF47C97490F23210400010511361B945DCBFBFDFCFAFAFEFCF6F6FDFEF8FCFEDBD586180303071D0308020107031613020406010205196DCAF3F0EDF6F8F7FCFDFBF9F7FEF8FDFAFDFAFAF3A62A080E11020D4BD7FBFBFCFDFEFEFDFCF8FFFFFFFFFFFFFFFFFFFFFFFFFFFFFFFFFFFFFFFFFFFFFFFFFFFFFFFFFFFFFFFFFFFFFFFFFFFFFFFFFFFFFFFFFFFFFFFFFFFFFFFFFFFFFFFFFFFFFFFFFFFFFFFFFFFFFFFFFFFFFFFFFFFFFFFFFFFFFFFFFFFFFFFFFFFFFFFFFFFFFFFFFFFFFFFFFFFFFFFFFFFFFFFFFFFFFFFFFFFFFFFFFFFFFFFFFFFFFFFFFFFFFFFFFFFFFFFFFFFFFFFFFFFFFEFDF8F7F6FBF68D1803153954ECFBF3F8FAFBFDFAF5F5FAFDFCFFFEEDD85D110673690500040D1CAA0841C9F3FDFCFCFDFEFDFBFCFDF9FCEEFAFD92967718080404030A0400040301000506030703010D005CBFDCF5FCF9F9FAF7FBFBEDE7FCFDFEF1F9FEF9F9EFA21D0002000050D5FAF8F7FBFDFCFCFEFDFDFEFFFFFFFFFFFFFFFFFFFFFFFFFFFFFFFFFFFFFFFFFFFFFFFFFFFFFFFFFFFFFFFFFFFFFFFFFFFFFFFFFFFFFFFFFFFFFFFFFFFFFFFFFFFFFFFFFFFFFFFFFFFFFFFFFFFFFFFFFFFFFFFFFFFFFFFFFFFFFFFFFFFFFFFFFFFFFFFFFFFFFFFFFFFFFFFFFFFFFFFFFFFFFFFFFFFFFFFFFFFFFFFFFFFFFFFFFFFFFFFFFFFFFFFFFFFFFFFFFFFFFFFFFEF7FDFBFDFDF9E2832667D3EEF5FDFEFDFDF4F6FCFDFDFCF5F4FEFCF464170101110E000102010C210F1FD5EAF6FEFEFEFDFAF7FEFDFCF7F6FCFDF29033100000000319141A040C0600070A0B0F01010100051B76DAF3F2F5F9FEFCFEFCFBF5F7FCFCFDFBF6F7F0DDA3000A010126C2F8FEFEFEFEFCFBFEFDF5FEFFFFFFFFFFFFFFFFFFFFFFFFFFFFFFFFFFFFFFFFFFFFFFFFFFFFFFFFFFFFFFFFFFFFFFFFFFFFFFFFFFFFFFFFFFFFFFFFFFFFFFFFFFFFFFFFFFFFFFFFFFFFFFFFFFFFFFFFFFFFFFFFFFFFFFFFFFFFFFFFFFFFFFFFFFFFFFFFFFFFFFFFFFFFFFFFFFFFFFFFFFFFFFFFFFFFFFFFFFFFFFFFFFFFFFFFFFFFFFFFFFFFFFFFFFFFFFFFFFFFFFFFFEF8FDFEFEF0E98B7ADC9C54FEFEF7FDEDEEF9FEFBF4F4FEFEFDF9EDF6DE49090D0D090600000306000022CCF5FBF6F6FBFCF6F1FBEDECFDFFFEFCF4F48D110303001AD806E32D060207061667880C110900000152CCF9FAF1F3F6ECFBECE4FDFEF8F9FDFEFEE88E5EB4060002090348C7F6F8F7F8F9F8FDFEFCFEFFFFFFFFFFFFFFFFFFFFFFFFFFFFFFFFFFFFFFFFFFFFFFFFFFFFFFFFFFFFFFFFFFFFFFFFFFFFFFFFFFFFFFFFFFFFFFFFFFFFFFFFFFFFFFFFFFFFFFFFFFFFFFFFFFFFFFFFFFFFFFFFFFFFFFFFFFFFFFFFFFFFFFFFFFFFFFFFFFFFFFFFFFFFFFFFFFFFFFFFFFFFFFFFFFFFFFFFFFFFFFFFFFFFFFFFFFFFFFFFFFFFFFFFFFFFFFFFFFFFFFFFFEFEF8F9F4F9FEF0EEF8EFF1FDFEF1FDFEFEF2F2FBF9F1F3FEFDFBFC4A1E1007017673130E020800000530CA6CAEF3F7ECEFFEFCFAFCEDFDFFFDF8FB9928040506092CF006D03EA14B0400060B050000000500000E186F7B82CCFCFDFDFDFDFAFEF8FCFEF2F3FCED9F4CB7050000033BC9FEFEFAF9FDFBF8FEFDFEFFFFFFFFFFFFFFFFFFFFFFFFFFFFFFFFFFFFFFFFFFFFFFFFFFFFFFFFFFFFFFFFFFFFFFFFFFFFFFFFFFFFFFFFFFFFFFFFFFFFFFFFFFFFFFFFFFFFFFFFFFFFFFFFFFFFFFFFFFFFFFFFFFFFFFFFFFFFFFFFFFFFFFFFFFFFFFFFFFFFFFFFFFFFFFFFFFFFFFFFFFFFFFFFFFFFFFFFFFFFFFFFFFFFFFFFFFFFFFFFFFFFFFFFFFFFFFFFFFFFFFFFFFFFFEF9FBF4EAF1FAF9FDFAEEF4FCFE9C809B9DC0F0FDFAF7F5FCF0E82F050703171502020201021000052EB5F1F7F8FAFDFBF7FAE4FCF5F8FD83A2781B538A0724AE3FD3F7E0F0510005070809000009006E891F01161861967FBEFAFEFEFCFAFDFBF7F9FDFAF1DF9F25000010033FC3EEF6FEFEFEFCF6FEFEFEFFFFFFFFFFFFFFFFFFFFFFFFFFFFFFFFFFFFFFFFFFFFFFFFFFFFFFFFFFFFFFFFFFFFFFFFFFFFFFFFFFFFFFFFFFFFFFFFFFFFFFFFFFFFFFFFFFFFFFFFFFFFFFFFFFFFFFFFFFFFFFFFFFFFFFFFFFFFFFFFFFFFFFFFFFFFFFFFFFFFFFFFFFFFFFFFFFFFFFFFFFFFFFFFFFFFFFFFFFFFFFFFFFFFFFFFFFFFFFFFFFFFFFFFFFFFFFFFFFFFFFFFFFFFEFF9F8F9FAFCFCFAF9F6FAF6FEF3AA727D505CBAFAFEFAF6EF56DA441304000303010200010A1FAC24B3F3FDFBFCFEFEFBF7FCF8F8FAF6FDD98C3306131C0844EEC3FAF6EB709B4E0102010200010B00172F5B6A2902131E49CCFCF9FEFDF2FCFEF6F8FDFEFCF9EBC4A40003001D2A6BC4F0F4F8FCF8FEFEFEFFFFFFFFFFFFFFFFFFFFFFFFFFFFFFFFFFFFFFFFFFFFFFFFFFFFFFFFFFFFFFFFFFFFFFFFFFFFFFFFFFFFFFFFFFFFFFFFFFFFFFFFFFFFFFFFFFFFFFFFFFFFFFFFFFFFFFFFFFFFFFFFFFFFFFFFFFFFFFFFFFFFFFFFFFFFFFFFFFFFFFFFFFFFFFFFFFFFFFFFFFFFFFFFFFFFFFFFFFFFFFFFFFFFFFFFFFFFFFFFFFFFFFFFFFFFFFFFFFFFFFFFFEFDF9FDFAF5F9F2EEFCFAF5FEFBF3F9AC3A26144ABDFAF9FEFCF8C8661C060104030405050901070A5DB5F2FBFCFDFEFCFBFDFDFDFBFCFEFCFAF8DC6A11010435B6F8F8EFFAF9CADC60050401020000020602051D250D000339A4EEFAEDFEFEFBF3F9FBFAFAFDFBF8F7C956981105060743B8FAFBFCFDFAFEF9FEFFFFFFFFFFFFFFFFFFFFFFFFFFFFFFFFFFFFFFFFFFFFFFFFFFFFFFFFFFFFFFFFFFFFFFFFFFFFFFFFFFFFFFFFFFFFFFFFFFFFFFFFFFFFFFFFFFFFFFFFFFFFFFFFFFFFFFFFFFFFFFFFFFFFFFFFFFFFFFFFFFFFFFFFFFFFFFFFFFFFFFFFFFFFFFFFFFFFFFFFFFFFFFFFFFFFFFFFFFFFFFFFFFFFFFFFFFFFFFFFFFFFFFFFFFFFFFFFFFFFFFFFFEFDFEFCF9F398565A93F7EFFEF9F9FEAF2B050651C2FBFCFDFAFDFBB12F0500020102040506080F23AAF4FDFDFDFEFFFDFDFFFFFEFAF9FEFFF7FCB24307010440E5FBFEF7F7FDF6F7C3500B0A020200000604030806023D5A436EDAFBF7FBFEFBF9F2FAFBFBFDFCFAF7E38E460A0400023AACF6FFFBF7FDFEF7FEFFFFFFFFFFFFFFFFFFFFFFFFFFFFFFFFFFFFFFFFFFFFFFFFFFFFFFFFFFFFFFFFFFFFFFFFFFFFFFFFFFFFFFFFFFFFFFFFFFFFFFFFFFFFFFFFFFFFFFFFFFFFFFFFFFFFFFFFFFFFFFFFFFFFFFFFFFFFFFFFFFFFFFFFFFFFFFFFFFFFFFFFFFFFFFFFFFFFFFFFFFFFFFFFFFFFFFFFFFFFFFFFFFFFFFFFFFFFFFFFFFFFFFFFFFFFFFFFFFFFFFFFFEFEF3FAF8EE65041364C38FFEFBFBF7DD82273292E3FBFAFEFCF486731C050502010101060A1F949B71EEFDFEFFFFFFFFFFFFFFFEFCFAF5FDF9F2C65910050646E2FBFBF9FAFBF9FDF6C28F51000202050242551603002A3B438DE2FDFDFCFCF5F8F7FCFDFDF8F6FDFCFBE389100300041D70CFFCFCF6F3FEFBFEFFFFFFFFFFFFFFFFFFFFFFFFFFFFFFFFFFFFFFFFFFFFFFFFFFFFFFFFFFFFFFFFFFFFFFFFFFFFFFFFFFFFFFFFFFFFFFFFFFFFFFFFFFFFFFFFFFFFFFFFFFFFFFFFFFFFFFFFFFFFFFFFFFFFFFFFFFFFFFFFFFFFFFFFFFFFFFFFFFFFFFFFFFFFFFFFFFFFFFFFFFFFFFFFFFFFFFFFFFFFFFFFFFFFFFFFFFFFFFFFFFFFFFFFFFFFFFFFFFFFFFFFFEFAF5FCA07433020A3D88A7F2FEF8F8FDDBA1AADFF9F8FBFDFBF69D862405040404070A0A0B248EE393F1F8FCFFFFFFFFFFFFFFFDF8FDF5F7F1B4833F0E032B7C8CF5FBF9FEFCFDFDFAFBE982000002040336440E070F0D04225D8DDBF5FEFCF8FAFDF3F6FCF7F3FAFEFDFBDE691501032077CFF4FAF8F4FCFFFFFFFFFFFFFFFFFFFFFFFFFFFFFFFFFFFFFFFFFFFFFFFFFFFFFFFFFFFFFFFFFFFFFFFFFFFFFFFFFFFFFFFFFFFFFFFFFFFFFFFFFFFFFFFFFFFFFFFFFFFFFFFFFFFFFFFFFFFFFFFFFFFFFFFFFFFFFFFFFFFFFFFFFFFFFFFFFFFFFFFFFFFFFFFFFFFFFFFFFFFFFFFFFFFFFFFFFFFFFFFFFFFFFFFFFFFFFFFFFFFFFFFFFFFFFFFFFFFFFFFFFFFFFEF9FDF4987337050A40A6F3F0F9FAFCFAF9FDFCFAF8F9FFFCF5FCFD9927060203033B490703258398EFFDF9FDFFFFFFFFFFFFFFFDF7FAFBF7BA8A3B0F080B367D94F6FCFAFCFBFBFCFAF6DF800B0601000302040509393C1907042183E5F3FEF4FCFCF1F4FCFCF9FAFEFBF4DF71150C031C60A0CCEEF9F9FBFBFEFFFFFFFFFFFFFFFFFFFFFFFFFFFFFFFFFFFFFFFFFFFFFFFFFFFFFFFFFFFFFFFFFFFFFFFFFFFFFFFFFFFFFFFFFFFFFFFFFFFFFFFFFFFFFFFFFFFFFFFFFFFFFFFFFFFFFFFFFFFFFFFFFFFFFFFFFFFFFFFFFFFFFFFFFFFFFFFFFFFFFFFFFFFFFFFFFFFFFFFFFFFFFFFFFFFFFFFFFFFFFFFFFFFFFFFFFFFFFFFFFFFFFFFFFFFFFFFFFFFFFFFFFFFFFCEEF8F67419478FBAEEFBF7F6FBF8FDFFFBF8F9FAFAFCF4EEF6351E0C0401024C58020833D782EAF8FCFEFFFFFFFFFFFFFFFFFDF3FBFBAE8D45140A2C6089FEF9F4FAFBFAF9FBF8ECBD80500E05150005070B063B4723091A3A40D4E6F8EFFAF7FBFDFBFCFFFFFEFBF3EF5E0B0F05001D5899D5FAF8FEF9FEFFFFFFFFFFFFFFFFFFFFFFFFFFFFFFFFFFFFFFFFFFFFFFFFFFFFFFFFFFFFFFFFFFFFFFFFFFFFFFFFFFFFFFFFFFFFFFFFFFFFFFFFFFFFFFFFFFFFFFFFFFFFFFFFFFFFFFFFFFFFFFFFFFFFFFFFFFFFFFFFFFFFFFFFFFFFFFFFFFFFFFFFFFFFFFFFFFFFFFFFFFFFFFFFFFFFFFFFFFFFFFFFFFFFFFFFFFFFFFFFFFFFFFFFFFFFFFFFFFFFFFFFFFFFF9F8FAEC9B74B6CA92F4F5FEF0FCFBFEFFFDFDFDF8F5FCEBB9DA2D5F2505040704000515378EF0FEF1F8FDFFFFFFFFFFFFFFFEFAEDF6F8FA781D040A46ADE9FDF6EFF8FEFEFCFAF9F0B4A4BC372E5A000A0306070503020028603BBAACE4EEF6F6F8FBF8F9FDFCFDFCFCF4D52B0A07032A5C87C2F6FEFDFEFEFFFFFFFFFFFFFFFFFFFFFFFFFFFFFFFFFFFFFFFFFFFFFFFFFFFFFFFFFFFFFFFFFFFFFFFFFFFFFFFFFFFFFFFFFFFFFFFFFFFFFFFFFFFFFFFFFFFFFFFFFFFFFFFFFFFFFFFFFFFFFFFFFFFFFFFFFFFFFFFFFFFFFFFFFFFFFFFFFFFFFFFFFFFFFFFFFFFFFFFFFFFFFFFFFFFFFFFFFFFFFFFFFFFFFFFFFFFFFFFFFFFFFFFFFFFFFFFFFFFFFFFFFEFCFCF4F3F9F8F8F8F5F8FAEDFCFEFDF3FAFFFAFAFDFFFFF8D7889E267536090C080202011C4C15A4FBF1FCFFFFFFFFFFFFFFFFFDF3CFCFEBEC6B19072B5168AFF9F9FDEFFBFDF4F9F4F6F3EAD5989FAB0A020501000408050107111F6E1FC2A3E2F1F3F6FBFFFBF7FAF7F9EFED6A3E1B0A0D1F5DB6E6F9FAFDFEFFFFFFFFFFFFFFFFFFFFFFFFFFFFFFFFFFFFFFFFFFFFFFFFFFFFFFFFFFFFFFFFFFFFFFFFFFFFFFFFFFFFFFFFFFFFFFFFFFFFFFFFFFFFFFFFFFFFFFFFFFFFFFFFFFFFFFFFFFFFFFFFFFFFFFFFFFFFFFFFFFFFFFFFFFFFFFFFFFFFFFFFFFFFFFFFFFFFFFFFFFFFFFFFFFFFFFFFFFFFFFFFFFFFFFFFFFFFFFFFFFFFFFFFFFFFFFFFFFFFFFFFFEF3F3FAFBF8FBF9F8FDFDFCF6F4FDFAE9F7FDFDF9F9FEFEFDFDEB762103000006040102073C860312F7F3FCFEFFFFFFFFFFFFFFFCED827A9BA953100247BABE77EDEBF9FCF9FBF4F5FBFBFAE7AC93F3B438090004010303030706000305118E0FBFB8DFFDFEFDFEFDFBFDFDEFEBC18D2A020000105AC7ECFDFEFEFFFFFFFFFFFFFFFFFFFFFFFFFFFFFFFFFFFFFFFFFFFFFFFFFFFFFFFFFFFFFFFFFFFFFFFFFFFFFFFFFFFFFFFFFFFFFFFFFFFFFFFFFFFFFFFFFFFFFFFFFFFFFFFFFFFFFFFFFFFFFFFFFFFFFFFFFFFFFFFFFFFFFFFFFFFFFFFFFFFFFFFFFFFFFFFFFFFFFFFFFFFFFFFFFFFFFFFFFFFFFFFFFFFFFFFFFFFFFFFFFFFFFFFFFFFFFFFFFFFFFFFFFFFFFCF4FAFBFCFFFCF8FAFEFEFDF7F9FEFAF6FAFCFCFDFBF5FEF9BD3E060001000001030200000228F8FFFBFDFEFEFFFEFEFFFEF5CAAE520A000603032280E5FEFEFEF7F7FDF0FAFAFCFFF8F5FCF9F9D26E0B0008030107050001041005010100CB3F9EF7FCF8F8FEFBF2FEFDCE63443C1503042471B7E5F6FEFEFFFFFFFEFEFFFEFEFFFFFFFFFFFFFFFFFFFFFFFFFFFFFFFFFFFFFFFFFFFFFFFFFFFFFFFFFFFFFFFFFFFFFFFFFFFFFFFFFFFFFFFFFFFFFFFFFFFFFFFFFFFFFFFFFFFFFFFFFFFFFFFFFFFFFFFFFFFFFFFFFFFFFFFFFFFFFFFFFFFFFFFFFFFFFFFFFFFFFFFFFFFFFFFFFFFFFFFFFFFFFFFFFFFFFFFFFFFFFFFFFFFFFFFFFFFFFFFFFFFFFFFFFFFFF3FBFDFBFCFFFEFCF9F5FCFDE7BCF2F9FBFEFAFAFEFEFBF4F37515020201000B0A000308000223F8FFFBFAF9FCFFFCFBFFFCE48C4326060000093073C3F8FBF5FEFEFEF5FCFEFEFEFFFEFEFEFEF3A95033060D04000707000204104C0E00008C32A7FAFEFBF8FAFBF0FDF2FDFCF89627030400104086DDFCFEFFFFFFFCFDFFFCFBFFFFFFFFFFFFFFFFFFFFFFFFFFFFFFFFFFFFFFFFFFFFFFFFFFFFFFFFFFFFFFFFFFFFFFFFFFFFFFFFFFFFFFFFFFFFFFFFFFFFFFFFFFFFFFFFFFFFFFFFFFFFFFFFFFFFFFFFFFFFFFFFFFFFFFFFFFFFFFFFFFFFFFFFFFFFFFFFFFFFFFFFFFFFFFFFFFFFFFFFFFFFFFFFFFFFFFFFFFFFFFFFFFFFFFFFFFFFFFFFFFFFFFFFFEF2FCF3FBFBFCFDFCFCF9FAE398764BC6CDEAFCFDFEF7EFFBFEF9BF3305005229000000001A4C0922F8B8E8FBFBFCFFFEFEFFFEFCF4DA4E090B05004D748EE1FEFEFDF2FAFDFEFFFFFFFFFFFFFFFFFEF6DF9D1A010008060E10030A277E2104060D03A9FEF7F6FCFBFEFCF5FCF5B8C7733C1F0D0825698CBAFDFBF5F5F7F7FBFFFEFEFFFFFFFFFFFFFFFFFFFFFFFFFFFFFFFFFFFFFFFFFFFFFFFFFFFFFFFFFFFFFFFFFFFFFFFFFFFFFFFFFFFFFFFFFFFFFFFFFFFFFFFFFFFFFFFFFFFFFFFFFFFFFFFFFFFFFFFFFFFFFFFFFFFFFFFFFFFFFFFFFFFFFFFFFFFFFFFFFFFFFFFFFFFFFFFFFFFFFFFFFFFFFFFFFFFFFFFFFFFFFFFFFFFFFFFFFFFFFFFFFFFFFFFFFEFDFEF6FBFBFCFDFCFCF1CD8159123C6360CAFBF9FBFDFCFEF1A5791A0200744A0B020628777E02128B64E1FEFDFDFFFDFAFBFEF9E68A13011708003395E2F7F6F2FEFDFEFEFFFFFFFFFFFFFFFFFFFFFEFBBA4928270F0209090001071A0A00011105609DD3F7F6FCF7F9FCF7CB829CA97C2C02010B2F73DAFEFEFDFDFDFDFEFCF8F9FFFFFFFFFFFFFFFFFFFFFFFFFFFFFFFFFFFFFFFFFFFFFFFFFFFFFFFFFFFFFFFFFFFFFFFFFFFFFFFFFFFFFFFFFFFFFFFFFFFFFFFFFFFFFFFFFFFFFFFFFFFFFFFFFFFFFFFFFFFFFFFFFFFFFFFFFFFFFFFFFFFFFFFFFFFFFFFFFFFFFFFFFFFFFFFFFFFFFFFFFFFFFFFFFFFFFFFFFFFFFFFFFFFFFFFFFFFFFFFFFFFFFFFFFFFFF8FAFEFEFEF6F2FBF3C36657170E2F81D9FAF9F8FAF8F8ED98942B07031B110A070533681F105B75CDF5FEFBFCFEFDF9F9FEF7C742040105023A69B0F7FEFCFBFEFAFDFEFFFFFFFFFFFFFFFFFBF0F4FEC97B68734418080101070A020906000417534095F8F3F9F8EEFEFCE57C362931382408050E2FB4FEFBFBF9FBFEFFFCFAFDFFFFFFFFFFFFFFFFFFFFFFFFFFFFFFFFFFFFFFFFFFFFFFFFFFFFFFFFFFFFFFFFFFFFFFFFFFFFFFFFFFFFFFFFFFFFFFFFFFFFFFFFFFFFFFFFFFFFFFFFFFFFFFFFFFFFFFFFFFFFFFFFFFFFFFFFFFFFFFFFFFFFFFFFFFFFFFFFFFFFFFFFFFFFFFFFFFFFFFFFFFFFFFFFFFFFFFFFFFFFFFFFFFFFFFFFFFFFFFFFFFFFFFFFFEF6FAFEF5FAFAECD0CCF3EAB7531E312CB8F5FDFDFBFBFCEFBB83380F0301000005040111220424CBD5F0F9FDF8F9FDFEFDFDFEF2A62B0D41280E3E6EA9E6FAFFFFFDF3FDFEFFFFFFFFFFFFFFFFFDF9FAFECA9A9CB9852D020000030503064C11010859334FC4FBFAF5F7F6FBC73B0801264F330301021374DDFCFDFCF6F9FAFBFCFFFFFFFFFFFFFFFFFFFFFFFFFFFFFFFFFFFFFFFFFFFFFFFFFFFFFFFFFFFFFFFFFFFFFFFFFFFFFFFFFFFFFFFFFFFFFFFFFFFFFFFFFFFFFFFFFFFFFFFFFFFFFFFFFFFFFFFFFFFFFFFFFFFFFFFFFFFFFFFFFFFFFFFFFFFFFFFFFFFFFFFFFFFFFFFFFFFFFFFFFFFFFFFFFFFFFFFFFFFFFFFFFFFFFFFFFFFFFFFFFFFFFFFFFFFEFAF6F2F6B6A6D3A781FAFDEEA56BB198DAFAEFF3FEFDFEDC6C60200144270801332F070B1D5A79FAFEF9FAFDFAFBFEFEFEFEFAE8AE390955393B639AD1F3F9F7FBFEF9FEFCFEFFFFFFFFFFFFFFFDF9F7FEEADCDFE5914C2D160800000107762103002818438FC2F5F9FEFAFCCC5316031122150003153E497BDEFBFEFBFCFCFAFBFFFFFFFFFFFFFFFFFFFFFFFFFFFFFFFFFFFFFFFFFFFFFFFFFFFFFFFFFFFFFFFFFFFFFFFFFFFFFFFFFFFFFFFFFFFFFFFFFFFFFFFFFFFFFFFFFFFFFFFFFFFFFFFFFFFFFFFFFFFFFFFFFFFFFFFFFFFFFFFFFFFFFFFFFFFFFFFFFFFFFFFFFFFFFFFFFFFFFFFFFFFFFFFFFFFFFFFFFFFFFFFFFFFFFFFFFFFFFFFFFFFFFFFFFFFEFCF8F6BC8255796B5DE7BBFCEDD0F2E7F1F6DAD0E4E3C2933B561E04682E070149470907258391F9FEF0FAFEFDFEFDFCFCFBEECD8F4304291A6CB4DEF3FAFCF9FBFDFBFEF9FDFFFFFFFFFFFFFFFEFAF9FDFCFBFDFAB07D4F1A0502080005381102010403447F8BC0F9F2FEFCECB4511302050A06021C5236609AD9F5FBFDFFFBFBFFFFFFFFFFFFFFFFFFFFFFFFFFFFFFFFFFFFFFFFFFFFFFFFFFFFFFFFFFFFFFFFFFFFFFFFFFFFFFFFFFFFFFFFFFFFFFFFFFFFFFFFFFFFFFFFFFFFFFFFFFFFFFFFFFFFFFFFFFFFFFFFFFFFFFFFFFFFFFFFFFFFFFFFFFFFFFFFFFFFFFFFFFFFFFFFFFFFFFFFFFFFFFFFFFFFFFFFFFFFFFFFFFFFFFFFFFFFFFFFFFFFFFFFFFFEFBFEFC9872311F2B48A599FEFCFEF8F8FDF8D8A4B2C172321A2F0E0334110200242404032C828DF9FEF1FAFCFDFEFDFBFCFAE9BE783B04042BA3EDFEF9F8FFFFFEFBFAFEF9FDFFFFFFFFFFFFFFFDF9FEF8FCFBFBFBC68E6E4311050802000101000100002276A3A4FAF2FAF9FCF498342A1C0809060C2A285A7ACCF3F8FCFEFCFCFFFFFFFFFFFFFFFFFFFFFFFFFFFFFFFFFFFFFFFFFFFFFFFFFFFFFFFFFFFFFFFFFFFFFFFFFFFFFFFFFFFFFFFFFFFFFFFFFFFFFFFFFFFFFFFFFFFFFFFFFFFFFFFFFFFFFFFFFFFFFFFFFFFFFFFFFFFFFFFFFFFFFFFFFFFFFFFFFFFFFFFFFFFFFFFFFFFFFFFFFFFFFFFFFFFFFFFFFFFFFFFFFFFFFFFFFFFFFFFFFFFFFFFFFFFEFCFDF8C98C52200B1F60D0FDF5FEFDFAFCE8CF947C7E3D000200050000000207050100042C89C3FAF9CBECF9FDFEFAF7FAF4CE93601E022E65BFEDFAFDFDFFFFFEFCFDFEFCFEFFFFFFFFFFFFFFFCF7FEF8FEFEFCF4C9A39C6A160506030001053045182A226EE1D1FDFBFAFAFDFCD1847E3D010507000218587DCBF6FBFCFEFDFCFEFFFFFFFFFFFFFFFFFFFFFFFFFFFFFFFFFFFFFFFFFFFFFFFFFFFFFFFFFFFFFFFFFFFFFFFFFFFFFFFFFFFFFFFFFFFFFFFFFFFFFFFFFFFFFFFFFFFFFFFFFFFFFFFFFFFFFFFFFFFFFFFFFFFFFFFFFFFFFFFFFFFFFFFFFFFFFFFFFFFFFFFFFFFFFFFFFFFFFFFFFFFFFFFFFFFFFFFFFFFFFFFFFFFFFFFFFFFFFFFFFFFFFFFFFFFFFCF9FCD3B55E214075FBFAF4FEF7E5D8B59578441E1800020936100400020B090201154B95FDF2EC61D7F9FEFBF7F6F6E8AC4D1E04015E72B9E5F2FAFFFFFFFFFFFFFFFFFFFFFFFFFFFFFFFFFEFDFFFEFCFDFEF5E2E1C47931100A010500068FA538574A5BC2D6FBF9FDFEFCF9E4C3B14C0304010502108B8BC2F3FEFCFCFEFEFEFFFFFFFFFFFFFFFFFFFFFFFFFFFFFFFFFFFFFFFFFFFFFFFFFFFFFFFFFFFFFFFFFFFFFFFFFFFFFFFFFFFFFFFFFFFFFFFFFFFFFFFFFFFFFFFFFFFFFFFFFFFFFFFFFFFFFFFFFFFFFFFFFFFFFFFFFFFFFFFFFFFFFFFFFFFFFFFFFFFFFFFFFFFFFFFFFFFFFFFFFFFFFFFFFFFFFFFFFFFFFFFFFFFFFFFFFFFFFFFFFFFFFFFFFEFBFEFCFAFCF87D56CDDBF5F5FFFFF4D1987B3F1E120102020117872E0C020202061E33559DC3FDEDB92BCFD4EDEBDBE9F7E7B63E040B014386D0F6F3F4FFFFFFFFFFFFFFFFFFFFFFFFFFFFFFFFFEFEFFFEF9FBFDFEF7F9E6C0822202010D0500A176232829254B97FDFAFBFFFFFAD29A8B3E1020180C0C35BBC7E5FAFFFAFAFDFEFFFFFFFFFFFFFFFFFFFFFFFFFFFFFFFFFFFFFFFFFFFFFFFFFFFFFFFFFFFFFFFFFFFFFFFFFFFFFFFFFFFFFFFFFFFFFFFFFFFFFFFFFFFFFFFFFFFFFFFFFFFFFFFFFFFFFFFFFFFFFFFFFFFFFFFFFFFFFFFFFFFFFFFFFFFFFFFFFFFFFFFFFFFFFFFFFFFFFFFFFFFFFFFFFFFFFFFFFFFFFFFFFFFFFFFFFFFFFFFFFFFFFFFFFFFEF3FBF7F2F9F58D7CF1F6FDF8FEFEF4D393491404020003060003281A0402070A0455B69D86FEF9E8362EAC79C8D995A3EAD5792202210D77DBE9F4FCFEFFFFFFFFFFFFFFFFFFFFFFFFFFFFFFFFFDFCFFFEFAF9FDFEFEFEFBCD6613070300110338110200000544BBFCFDFEFEFEFDE4B761182D6144011684E2F7FEFFFFFCFBFCFDFEFFFFFFFFFFFFFFFFFFFFFFFFFFFFFFFFFFFFFFFFFFFFFFFFFFFFFFFFFFFFFFFFFFFFFFFFFFFFFFFFFFFFFFFFFFFFFFFFFFFFFFFFFFFFFFFFFFFFFFFFFFFFFFFFFFFFFFFFFFFFFFFFFFFFFFFFFFFFFFFFFFFFFFFFFFFFFFFFFFFFFFFFFFFFFFFFFFFFFFFFFFFFFFFFFFFFFFFFFFFFFFFFFFFFFFFFFFFFFFFFFFFFFFFFFEF4F2FBFDFDF5DCDEFEFDFCFDFAF6C34B2B0D0A080003060000042203150A1111072293E5CDFAEAD5060154B5DEE4B17F8D800C080A8C4B87E6FEF6F8FEFFFFFFFFFFFFFFFFFFFFFFFFFFFFFFFFFDF8FFFFFFFDFDFDFCFCFDE8992E070F00020A020008040005298CFCFEFDFAFAFEFCE09A350F2119070843C0FCFEFEFEFFFEFCFBFCFFFFFFFFFFFFFFFFFFFFFFFFFFFFFFFFFFFFFFFFFFFFFFFFFFFFFFFFFFFFFFFFFFFFFFFFFFFFFFFFFFFFFFFFFFFFFFFFFFFFFFFFFFFFFFFFFFFFFFFFFFFFFFFFFFFFFFFFFFFFFFFFFFFFFFFFFFFFFFFFFFFFFFFFFFFFFFFFFFFFFFFFFFFFFFFFFFFFFFFFFFFFFFFFFFFFFFFFFFFFFFFFFFFFFFFFFFFFFFFFFFFFFFFFFEFEEDF9FEF8F5F8FDFCF4F1F5F9B66F912602010203080B00000EAC2085318D510C62DBFEFCFAFCDE060B14BFD9869379190F0300022770E9F9FCFEF8F9FFFFFFFFFFFFFFFFFFFFFFFFFFFFFFFFFDFAFEFEFEFEFFFCF8F8FCBE7784280600130805010707090D4AC7FEFDFDFCFDFEF7D44D0B06000103178AF6FDFBFBFEFFFEFDFDFFFFFFFFFFFFFFFFFFFFFFFFFFFFFFFFFFFFFFFFFFFFFFFFFFFFFFFFFFFFFFFFFFFFFFFFFFFFFFFFFFFFFFFFFFFFFFFFFFFFFFFFFFFFFFFFFFFFFFFFFFFFFFFFFFFFFFFFFFFFFFFFFFFFFFFFFFFFFFFFFFFFFFFFFFFFFFFFFFFFFFFFFFFFFFFFFFFFFFFFFFFFFFFFFFFFFFFFFFFFFFFFFFFFFFFFFFFFFFFFFFFFFFFFFFFEEDFBEFEBF6FDFEFEFEF0B876DB25131F0E0206476D39010000091C9FEBDD711A728188FEFEFDE53101299C5DB3CA5D0A030307040E7CE9F3FBFBF2F7FEFFFFFFFFFFFFFFFFFFFFFFFFFFFFFFFFFEFEFEFCFCFEFFFEFEFBFCECDDBD27070000B72C1307745D64EAFBFDFDFEFDFBFEF9E7A036000106080328B5FCFCFCFEFFFEFDFEFFFFFFFFFFFFFFFFFFFFFFFFFFFFFFFFFFFFFFFFFFFFFFFFFFFFFFFFFFFFFFFFFFFFFFFFFFFFFFFFFFFFFFFFFFFFFFFFFFFFFFFFFFFFFFFFFFFFFFFFFFFFFFFFFFFFFFFFFFFFFFFFFFFFFFFFFFFFFFFFFFFFFFFFFFFFFFFFFFFFFFFFFFFFFFFFFFFFFFFFFFFFFFFFFFFFFFFFFFFFFFFFFFFFFFFFFFFFFFFFFFFFFFFFFFFEFCF8FBFCFDFEFFF7F2FBE48F4C90160B030000233212020B0E0B8BEEF1F3D04687D8D5F5FBD741AA040F4FB9D084734E030A09000A81F5F7FCFEFAFCFFFFFFFFFFFFFFFFFFFFFFFFFFFFFFFFFFFFFFFEFCFCFEFFFFFFF9FAFCF9F5B9160F052E0F03002687817AFDFAFBFFFEF5F6FCFCEAA440042771677AF0FDFEFEFCFBFDFEFFFFFFFFFFFFFFFFFFFFFFFFFFFFFFFFFFFFFFFFFFFFFFFFFFFFFFFFFFFFFFFFFFFFFFFFFFFFFFFFFFFFFFFFFFFFFFFFFFFFFFFFFFFFFFFFFFFFFFFFFFFFFFFFFFFFFFFFFFFFFFFFFFFFFFFFFFFFFFFFFFFFFFFFFFFFFFFFFFFFFFFFFFFFFFFFFFFFFFFFFFFFFFFFFFFFFFFFFFFFFFFFFFFFFFFFFFFFFFFFFFFFFFFFFFFFFFFF63D4FDFBFAFAF9F0BA76BA1E370604010503070903031F84C4EAFCF6FBF8C8D9F8F9FBCB5FA83F0D03115BB061241D07021A78BEE4FCFDF8FCFAFBFFFFFFFFFFFFFFFFFFFFFFFFFFFFFFFFFFFFFFFEFEFEFEFCFBFCFDF9F7FCD5550D0301050501010A85D7D2FDFCF6F7FEFDFBFDF8CE66140010291836BEF9FEFFFCFCFDFFFBFAFFFFFFFFFFFFFFFFFFFFFFFFFFFFFFFFFFFFFFFFFFFFFFFFFFFFFFFFFFFFFFFFFFFFFFFFFFFFFFFFFFFFFFFFFFFFFFFFFFFFFFFFFFFFFFFFFFFFFFFFFFFFFFFFFFFFFFFFFFFFFFFFFFFFFFFFFFFFFFFFFFFFFFFFFFFFFFFFFFFFFFFFFFFFFFFFFFFFFFFFFFFFFFFFFFFFFFFFFFFFFFFFFFFFFFFFFCFAFBFDFCF9FBEE62F42EC190C2F4F8FBF8D9874B09080005395D28020203060D3858DDF1FAFEFCF5F9FEFFBF3920440909010113361D05060304123AB6E4FCFFF7F8FBFDFFFFFFFFFFFFFFFFFFFFFFFFFFFFFFFFFFFFFFFFFFFFFFFEFDFEFEFBF8FCEDBA9B160200060109072B9BF5F7FEFBF8FFFFF1FAC78A753D00060B1687E3FAFEFCFBFDF7F6FBFDFFFFFFFFFFFFFFFFFFFFFFFFFFFFFFFFFFFFFFFFFFFFFFFFFFFFFFFFFFFFFFFFFFFFFFFFFFFFFFFFFFFFFFFFFFFFFFFFFFFFFFFFFFFFFFFFFFFFFFFFFFFFFFFFFFFFFFFFFFFFFFFFFFFFFFFFFFFFFFFFFFFFFFFFFFFFFFFFFFFFFFFFFFFFFFFFFFFFFFFFFFFFFFFFFFFFFFFFFFFFFFFFFFFFFFFFFDFCFDFDFCFCFBE11C5F0E512BA7A3AEFBF1BC3C0B032979181E3A1A03050403031886E3F6FDFFFCFCFEFDF7D37309268707000603040402040806196083BCEEFDFCFCFDFAFBFFFFFFFFFFFFFFFFFFFFFFFFFFFFFFFFFFFFFFFFFFFFFFFFFFF9F8FDF7EDE27A0A0A0B04010465563199F9FAFEFAF8F3F6FAE5CAA54A030605144063C5F3FCFAFDFBFAF7F7FFFFFFFFFFFFFFFFFFFFFFFFFFFFFFFFFFFFFFFFFFFFFFFFFFFFFFFFFFFFFFFFFFFFFFFFFFFFFFFFFFFFFFFFFFFFFFFFFFFFFFFFFFFFFFFFFFFFFFFFFFFFFFFFFFFFFFFFFFFFFFFFFFFFFFFFFFFFFFFFFFFFFFFFFFFFFFFFFFFFFFFFFFFFFFFFFFFFFFFFFFFFFFFFFFFFFFFFFFFFFFFFFFFFFFFFFFFFFCFAF6F1E87890150D0E044A8AC7F3C14B070809254C1204090A0503080F091CBBF7FDFEFDFDFEFFFCF8F3AD0538DF07000401090A030204165FC3B9D6F2FBFBFCFEFCFCFFFFFFFFFFFFFFFFFFFFFFFFFFFFFFFFFFFFFFFFFFFFFFFFFFFBFBFFFAF4EEA0170B070102024135042E769DFBFCF3F7FDF4F8FAC04C020418583D61D9F8F5F7FDFBFAF8FAFFFFFFFFFFFFFFFFFFFFFFFFFFFFFFFFFFFFFFFFFFFFFFFFFFFFFFFFFFFFFFFFFFFFFFFFFFFFFFFFFFFFFFFFFFFFFFFFFFFFFFFFFFFFFFFFFFFFFFFFFFFFFFFFFFFFFFFFFFFFFFFFFFFFFFFFFFFFFFFFFFFFFFFFFFFFFFFFFFFFFFFFFFFFFFFFFFFFFFFFFFFFFFFFFFFFFFFFFFFFFFFFFFFFFFFFFFFCFAFAF7F2E88A7F0D0A035E85C5F5BF500801070C09000300060704070D0C1B6DDFFFFFFDFAFBFEFFFCFBFCD5799CDA0B17641D07204632164F89BEF5FAFDFDFBFBFCFEFFFFFFFFFFFFFFFFFFFFFFFFFFFFFFFFFFFFFFFFFFFFFFFFFFFFFBFBFBFAFDE98F741702070503030100041251FCFDECFDF8F5F9FDDE872903184A344CCBF4F6FAFEFAFAFDFFFFFFFFFFFFFFFFFFFFFFFFFFFFFFFFFFFFFFFFFFFFFFFFFFFFFFFFFFFFFFFFFFFFFFFFFFFFFFFFFFFFFFFFFFFFFFFFFFFFFFFFFFFFFFFFFFFFFFFFFFFFFFFFFFFFFFFFFFFFFFFFFFFFFFFFFFFFFFFFFFFFFFFFFFFFFFFFFFFFFFFFFFFFFFFFFFFFFFFFFFFFFFFFFFFFFFFFFFFFFFFFFFFFFFFFFFFFF9F8FAF9FCF9EC7D6E135483D3F8CA7F3A040706060305010509040123586E86D6FBFCFDFDFAFBFEFFFDFBFBF8D8A2D50B196D22001A4B38185590BDFAF9FBFEFEFDFCFDFFFFFFFFFFFFFFFFFFFFFFFFFFFFFFFFFFFFFFFFFFFFFFFFFFFFFBFBF9F7FBE291881501070803020105090D50F9FDF1F7F6FCFAFCDB8A31040607060FAEF1FDFCFDFCFDFDFCFFFFFFFFFFFFFFFFFFFFFFFFFFFFFFFFFFFFFFFFFFFFFFFFFFFFFFFFFFFFFFFFFFFFFFFFFFFFFFFFFFFFFFFFFFFFFFFFFFFFFFFFFFFFFFFFFFFFFFFFFFFFFFFFFFFFFFFFFFFFFFFFFFFFFFFFFFFFFFFFFFFFFFFFFFFFFFFFFFFFFFFFFFFFFFFFFFFFFFFFFFFFFFFFFFFFFFFFFFFFFFFFFFFFFFFFFFF9F8F9F8FEFDF9F8F779D387C4FBAC6F5A060A0B1C4B0E243B1405062A86E0FEFDF9FAFCFEFDFDFEFEFEFBF7F3D66D910A04010304011833324DBCFAF8F7F9FDFDFDFCFDFFFFFFFFFFFFFFFFFFFFFFFFFFFFFFFFFFFFFFFFFFFFFFFFFFFFFDFBF9FAF9F3EF810C010006090701070C237AACFEF8F4FDFCF8FBD37B270504010B47CCF9FAF7FBFCFDFCFBFFFFFFFFFFFFFFFFFFFFFFFFFFFFFFFFFFFFFFFFFFFFFFFFFFFFFFFFFFFFFFFFFFFFFFFFFFFFFFFFFFFFFFFFFFFFFFFFFFFFFFFFFFFFFFFFFFFFFFFFFFFFFFFFFFFFFFFFFFFFFFFFFFFFFFFFFFFFFFFFFFFFFFFFFFFFFFFFFFFFFFFFFFFFFFFFFFFFFFFFFFFFFFFFFFFFFFFFFFFFFFFFFFFFFFFFFCFBFCFBFBFEFCF5FFFFFDFCFEFFF0BD5D0F08040537890D4265241F311B38B1FBFDFCFCFEFFFFFFFEFEFEFEDFA8B08C060E07030508093B8CABAFF5FCFBFCFEFDFCFBFCFCFDFFFFFFFFFFFFFFFFFFFFFFFFFFFFFFFFFFFFFFFFFFFFFFFFFFFCF7F9FDFAFDF4AF1A02040133220623313C936EF9F4FDFFF7FBFDE8B5641E0A031574A4E0F3F7FCFBFCFCFCFFFFFFFFFFFFFFFFFFFFFFFFFFFFFFFFFFFFFFFFFFFFFFFFFFFFFFFFFFFFFFFFFFFFFFFFFFFFFFFFFFFFFFFFFFFFFFFFFFFFFFFFFFFFFFFFFFFFFFFFFFFFFFFFFFFFFFFFFFFFFFFFFFFFFFFFFFFFFFFFFFFFFFFFFFFFFFFFFFFFFFFFFFFFFFFFFFFFFFFFFFFFFFFFFFFFFFFFFFFFFFFFFFFFFFFFF7F5F7F9FFFEFCFEFDFEFAF8FAFDC6702801030A000206010401043B623D3293E8FDFCFEFEFDFEFEFEFEFFFCB73E3B011007060F3A3C4761B0FDFCFCF8F5F9FDFCFDFEFEFBFBFFFFFFFFFFFFFFFFFFFFFFFFFFFFFFFFFFFFFFFFFFFFFFFFFFFCFCFDFBFBFCE85119010B0365482D7D97517BFCF3FDFEFEF6FBFDC5998F590501031232AEF2FEFDFDFEFBF9FFFFFFFFFFFFFFFFFFFFFFFFFFFFFFFFFFFFFFFFFFFFFFFFFFFFFFFFFFFFFFFFFFFFFFFFFFFFFFFFFFFFFFFFFFFFFFFFFFFFFFFFFFFFFFFFFFFFFFFFFFFFFFFFFFFFFFFFFFFFFFFFFFFFFFFFFFFFFFFFFFFFFFFFFFFFFFFFFFFFFFFFFFFFFFFFFFFFFFFFFEFEFDFEFEFEFEFEFFFFFFFFFFFFFFFFFFFEF9F8FEFEFBFCFAFDF9FBF6F56A627300000004080006010302040F2B3955C2FCF6FEF9FAFCFBFDFEFEE6C8CAA0010005395DCDFAFAFBFFF9F3FAFFFFFFFFF9FAFDFFFFFFFFFFFFFFFFFFFFFFFFFFFFFFFFFFFFFFFFFFFFFFFFFFFFFFFFFEFCF8F6FFFCEED04F0A0100080E7C9883F4F9FDFEF7FBFEFBFDF9FBFEC9550302080585C4EAF7FCFCFDFEFFFFFFFFFFFFFFFFFFFFFFFFFFFFFFFFFFFFFFFFFFFFFFFFFFFFFFFFFFFFFFFFFFFFFFFFFFFFFFFFFFFFFFFFFFFFFFFFFFFFFFFFFFFFFFFFFFFFFFFFFFFFFFFFFFFFFFFFFFFFFFFFFFFFFFFFFFFFFFFFFFFFFFFFFFFFFFFFFFFFFFFFFFFFFFFFFFFFFFFFFEF9FCF7FAFCFAFDFCFFFFFFFFFFFFFFFFFFFAFAFDFCFEFCF6FDF8FCF2F4DB5800030403040F01000B02030A054FC5FBF9FEFBFEFDECFCFDF0F5FBF8920C150E060116D6F4FCFFFFFFFFFFFFFFFFFFFFFFFFFFFFFFFFFFFFFFFFFFFFFFFFFFFFFFFFFFFFFFFFFFFFFFFFFFFFFFFFFFFFFCF7F7FEFEFEE425CF2D070401076DCBD4E5FAF6FEFDFEFDF7FBFDFDFCCF6910030400162DA9E7FAFEFDFCFEFFFFFFFFFFFFFFFFFFFFFFFFFFFFFFFFFFFFFFFFFFFFFFFFFFFFFFFFFFFFFFFFFFFFFFFFFFFFFFFFFFFFFFFFFFFFFFFFFFFFFFFFFFFFFFFFFFFFFFFFFFFFFFFFFFFFFFFFFFFFFFFFFFFFFFFFFFFFFFFFFFFFFFFFFFFFFFFFFFFFFFFFFFFFFFFFFFFFFFFAFCFBFBFDFCFAFDFCFFFFFFFFFFFFFFFFFFFCFCFDFCFCFBF6FAFDEAF8E55B0806070005120211030002080A0B59CAFCFBFEFBF8FCFDFEEDF7DFE1F89354B50003082B38DCFEFFFFFFFFFFFFFFFFFFFFFFFFFFFFFFFFFFFFFFFFFFFFFFFFFFFFFFFFFFFFFFFFFFFFFFFFFFFFFFFFFFFFFEFEF6F1FBF4FCF54F2B0700090A080B185DE1FDFCF6FEF3FBFAFCF6F6F6C152040A059EE9FBFEFAFCFDFAFAFFFFFFFFFFFFFFFFFFFFFFFFFFFFFFFFFFFFFFFFFFFFFFFFFFFFFFFFFFFFFFFFFFFFFFFFFFFFFFFFFFFFFFFFFFFFFFFFFFFFFFFFFFFFFFFFFFFFFFFFFFFFFFFFFFFFFFFFFFFFFFFFFFFFFFFFFFFFFFFFFFFFFFFFFFFFFFFFFFFFFFFFFFFFFFFFFFFFFFFCFDFAFCFEFCFBFCF9FFFFFFFFFFFFFFFFFDFDFDFEFDFCFCFDF8F8FCF2699457040504131E0C161F020A4C7991BFEEFCF7F8FEFEFDFAFEFDFB8A93CC7C0E120D0D47CFE6F6FDFFFFFFFFFFFFFFFFFFFFFFFFFFFFFFFFFFFFFFFFFFFFFFFFFFFFFFFFFFFFFFFFFFFFFFFFFFFFFFFFFFFFFFFAF8FAF6FCF7FBB8CB3B100400000A569967F6FCEAFDFBFDF7FAF6F6C78CA35704000C91DCF7F3F8FEFBFAFFFFFFFFFFFFFFFFFFFFFFFFFFFFFFFFFFFFFFFFFFFFFFFFFFFFFFFFFFFFFFFFFFFFFFFFFFFFFFFFFFFFFFFFFFFFFFFFFFFFFFFFFFFFFFFFFFFFFFFFFFFFFFFFFFFFFFFFFFFFFFFFFFFFFFFFFFFFFFFFFFFFFFFFFFFFFFFFFFFFFFFFFFFFFFFFFFFFFFFEFEFDFCFDFCFBFDFCFFFFFFFFFFFFFFFFFDFDFEFDFDFBFDFEF8F8F8FBC353100308133F820E68870D011871BBE8FCFEFCFBFEF9F3FDFEFFCFCD71221401020719CE60E5FCFEFFFFFFFFFFFFFFFFFFFFFFFFFFFFFFFFFFFFFFFFFFFFFFFFFFFFFFFFFFFFFFFFFFFFFFFFFFFFFFFFFFFFF9F9FEFEFBFBFBED762F0A02000506031056CFF6FBFAF7F0FBFDFCF8F3EEDFB96D15001128B1F1FDFDFFFBF9FFFFFFFFFFFFFFFFFFFFFFFFFFFFFFFFFFFFFFFFFFFFFFFFFFFFFFFFFFFFFFFFFFFFFFFFFFFFFFFFFFFFFFFFFFFFFFFFFFFFFFFFFFFFFFFFFFFFFFFFFFFFFFFFFFFFFFFFFFFFFFFFFFFFFFFFFFFFFFFFFFFFFFFFFFFFFFFFFFFFFFFFFFFFFFFFFFFFFFFFF3FAF4F7FDFCFEFEFFFFFFFFFFFFFFFFFEFEFEFDFAF8F9FEF8FCF8C8541503000204143742B1263309299BE5FBFDF9F4F3FEFDFBFCD2C37DA45A070700000F5AF2D4F7F5FBFFFFFFFFFFFFFFFFFFFFFFFFFFFFFFFFFFFFFFFFFFFFFFFFFFFFFFFFFFFFFFFFFFFFFFFFFFFFFFFFFFFFFBFCFDFDFDFBFCE7AE4C150300010101143463FDFDFAF9FAF7FBFDFDF6F6DF976C3903046CDCEBF7FEF4F2FDFFFFFFFFFFFFFFFFFFFFFFFFFFFFFFFFFFFFFFFFFFFFFFFFFFFFFFFFFFFFFFFFFFFFFFFFFFFFFFFFFFFFFFFFFFFFFFFFFFFFFFFFFFFFFFFFFFFFFFFFFFFFFFFFFFFFFFFFFFFFFFFFFFFFFFFFFFFFFFFFFFFFFFFFFFFFFFFFFFFFFFFFFFFFFFFFFFFFFFFFFAFCF6F9FDFCFDFFFFFFFFFFFFFFFFFFFFFDFDFCFAFCFBF8FBC3F2811F042629090001104ECE619C42448DD2F1F8FCFBFAFDF8FEFE66407C5D1B02020F2F50D2F6FDF9F9FDFFFFFFFFFFFFFFFFFFFFFFFFFFFFFFFFFFFFFFFFFFFFFFFFFFFFFFFFFFFFFFFFFFFFFFFFFFFFFFFFFFFFFCFCFAF8FDFBFCF1759522040502040D436B42FCF4FDFDFEFCFBFBF5F4F8E1A96D4518022D80DBFBFEFAF6FAFFFFFFFFFFFFFFFFFFFFFFFFFFFFFFFFFFFFFFFFFFFFFFFFFFFFFFFFFFFFFFFFFFFFFFFFFFFFFFFFFFFFFFFFFFFFFFFFFFFFFFFFFFFFFFFFFFFFFFFFFFFFFFFFFFFFFFFFFFFFFFFFFFFFFFFFFFFFFFFFFFFFFFFFFFFFFFFFFFFFFFFFFFFFFFFFFFFFFFF6FCF3F7FCFBFAFCFBFFFFFFFFFFFFFFFFFFFBFBFAF8FFF9E5FD7DC995270D50470B15242077E5D2E4BFB5C8E6F7FBFEFFFFFEFAFBF6450C371C00031B64B7C3F9FBFAF8FCFEFFFFFFFFFFFFFFFFFFFFFFFFFFFFFFFFFFFFFFFFFFFFFFFFFFFFFFFFFFFFFFFFFFFFFFFFFFFFFFFFFFFFFDF8F6F8F7FCFAE9BE82170108060307213339E1F8F2FAF8F6FBF6F5FAFCDCA098932C020246C7F0F6FDFBFBFFFFFFFFFFFFFFFFFFFFFFFFFFFFFFFFFFFFFFFFFFFFFFFFFFFFFFFFFFFFFFFFFFFFFFFFFFFFFFFFFFFFFFFFFFFFFFFFFFFFFFFFFFFFFFFFFFFFFFFFFFFFFFFFFFFFFFFFFFFFFFFFFFFFFFFFFFFFFFFFFFFFFFFFFFFFFFFFFFFFFFFFFFFFFFFFFFFFFFD2F0D3E9F8F4F8FCFAFFFFFFFFFFFFFFFFFFFAFAF8F3FEF3C4F7915D4F170B30250B385F3CC4F8FEFAFBF9F8FCFFFDFCFFFEFCFBD9C8370E062234375AB0F9FBF6FEFCF5FCFFFFFFFFFFFFFFFFFFFFFFFFFFFFFFFFFFFFFFFFFFFFFFFFFFFFFFFFFFFFFFFFFFFFFFFFFFFFFFFFFFFFFFFFF9F8FCF7F9FCF1F8BA29060505020004164076F5F6FBEEDAEEF7F6F8FCECCED09D25030060CBF0F9FDFCFCFFFFFFFFFFFFFFFFFFFFFFFFFFFFFFFFFFFFFFFFFFFFFFFFFFFFFFFFFFFFFFFFFFFFFFFFFFFFFFFFFFFFFFFFFFFFFFFFFFFFFFFFFFFFFFFFFFFFFFFFFFFFFFFFFFFFFFFFFFFFFFFFFFFFFFFFFFFFFFFFFFFFFFFFFFFFFFFFFFFFFFFCFBFFFEFBFDFDFE75AF9BB6E5F7F7F7FDFFFFFFFFFFFFFFFFFFFAFAFAF4F7EB95C4690D0002030402082A5474F6FEFEFEF7FCFDFCFDFDEFD4CDF4D3685D1709236499926E91E9D4FAFDFAF7FDFFFFFFFFFFFFFFFFFFFFFFFFFFFFFFFFFFFFFFFFFFFFFFFFFFFFFFFFFFFFFFFFFFFFFFFFFFFFFFFFFFFFFFFFFDFDFEFCFAFFFEF9FD410B020201000832622BCBFACFAF84D3FBF9F4F7FDFDEEAE37271543B7F2FEFBFAFCFDFDFFFEF8FEFBFEFFFFFFFFFFFFFFFFFFFEFEFCFEFDFEFFECD8D0C9D2F5FFFFFEFEFEFFFFFFFFFFFFFFFFFFFFFFFFFFFFFFFFFFFFFFFFFFFFFFFFFFFFFFFFFFFFFFFFFFFFFFFFFFFFFFFFFFFFFFFFFFFFFFFFFFFFFFFFFFFFFFFFFAF6FAFCF7F9D4C21E466F6EBDFDF9F5FBFFFFFFFFFFFFFFFFFCFBFAFCFBF5E77E560B020522130604020833ADFFFEFEF8F5F8F8F7FBFFE1A18CBB74090605014E7EAEAC5449BAA2FEFCF9FDFBFDFFFFFFFFFFFFFFFFFFFFFFFFFFFFFFFFFFFFFFFFFFFFFFFFFFFFFFFFFFFFFFFFFFFFFFFFFFFFFFFFFFFFFFFFFFFFFFFFFFFFFFFA3E2B1201020307254A269AC1654330A8EFFCF8F6FBF7F7E4505D2E097FDFF9F8FAFEFDFCFBFDF2FEF9FDFFFFFFFFFFFFFFFFFFFEFDFBFDFCFDE8A36B5D6093F2F6FEFCFDFEFFFFFFFFFFFFFFFFFFFFFFFFFFFFFFFFFFFFFFFFFFFFFFFFFFFFFFFFFFFFFFFFFFFFFFFFFFFFFFFFFFFFFFFFFFFFFFFFFFFFFFFFFFFFFFFCFAF5F3FCF8EF686D1206463189E9F7FBF9FFFFFFFFFFFFFFFFF9FAFAEDD8CCC4650002060A572E0C0814263297FBFBFBF4F8EDE4E4E8EBD097785B1706051A286684706D4C276894FEFEFEFCF8FCFFFFFFFFFFFFFFFFFFFFFFFFFFFFFFFFFFFFFFFFFFFFFFFFFFFFFFFFFFFFFFFFFFFFFFFFFFFFFFFFFFFFFFFFFFFFFFFFFFFFFFF851662D0305070101071E83510E02247EC6EFFCFCF8F3FCE6475420083DB7F1FAFDFAFAFAF4FCF7FEFBFCFEFFFFFFFFFFFFFFFEFAFDFEFEF9E7B5540B01207FDDEFFAF9FCFEFFFFFFFFFFFFFFFFFFFFFFFFFFFFFFFFFFFFFFFFFFFFFFFFFFFFFFFFFFFFFFFFFFFFFFFFFFFFFFFFFFFFFFFFFFFFFFFFFFFFFFFFFFFFFFF9F8F8F4FBFAE403761A03010B62C5F1FAFBFFFFFFFFFFFFFFFFFDFDFCC974655F13000606044C2702083E550A50F4F6F6F9F6D6BDB4B1A98D6B5B23031E0D639595D36D3540022391FCFBFEF7FDFFFFFFFFFFFFFFFFFFFFFFFFFFFFFFFFFFFFFFFFFFFFFFFFFFFFFFFFFFFFFFFFFFFFFFFFFFFFFFFFFFFFFFFFFFFFFFFFFFFFFFFFFDAB813005050905050D1751080302266B97D4FCFCFCFCF9E566170601067BCBEDFBF4F3F9F3FFFEFDFCF9FCFFFFFFFFFFFFFFFDF5FBF7FBF2AA8450150422658EFBF6F6FAFDFFFFFFFFFFFFFFFFFFFFFFFFFFFFFFFFFFFFFFFFFFFFFFFFFFFFFFFFFFFFFFFFFFFFFFFFFFFFFFFFFFFFFFFFFFFFFFFFFFFFFFFFFFFFFFFFF9F7FEFAF1E701690E0402125FC8FEF8F2FFFFFFFFFFFFFFFFFDFDF69F3B5B080004000302020603103D400A4CDCEBFAF3FCCB846F715119040D000D8A25C9F7F2ED781D201F40B8FEF9FDF6FCFEFFFFFFFFFFFFFFFFFFFFFFFFFFFFFFFFFFFFFFFFFFFFFFFFFFFFFFFFFFFFFFFFFFFFFFFFFFFFFFFFFFFFFFFFFFFFFFFFFFFFFFFBFEA4380906040205010702001516205996D7F6F6FBFDFDF4AE2513040E3558A3E8FAFDFDFDFEFEF8FDFFFFFFFFFFFFFFFFFFFFFEFBEEFAF27546320D01021C75F6FCF7FDFFFFFFFFFFFFFFFFFFFFFFFFFFFFFFFFFFFFFFFFFFFFFFFFFFFFFFFFFFFFFFFFFFFFFFFFFFFFFFFFFFFFFFFFFFFFFFFFFFFFFFFFFFFFFFFFFEF5F1FBFEFBD7040F0500071257C2FCF9FDFFFFFFFFFFFFFFFFF7F9F6A342570105141012030A150B043A550B2E6890F2F4F8D75F27372A0600030B21D5A8E5EFF9FEC897959CB3F1FDFEFBF7FBFDFFFFFFFFFFFFFFFFFFFFFFFFFFFFFFFFFFFFFFFFFFFFFFFFFFFFFFFFFFFFFFFFFFFFFFFFFFFFFFFFFFFFFFFFFFFFFFFFFFFFFFFEFCFD540A0800020300020901366D2146BDF8FBFBFCF7F6E98F22100A0101022591F0F6FAFCFCFEF9FEFEFEFFFFFFFFFFFFFFFEFEFEFCF9BB430D0C0E080645EEF9F5FCFDFEFFFFFFFFFFFFFFFFFFFFFFFFFFFFFFFFFFFFFFFFFFFFFFFFFFFFFFFFFFFFFFFFFFFFFFFFFFFFFFFFFFFFFFFFFFFFFFFFFFFFFFFFFFFFFFFAFAF9F6FBDD6500000B060666A1D0FBF5F3FFFFFFFFFFFFFFFFFEFDC1521004030701921B010067480E3C5645780746E29FBE8C712F020001031772A2756370C5F4FEF3F5FCFBFEFBF4FEFEFEFBFDFFFFFFFFFFFFFFFFFFFFFFFFFFFFFFFFFFFFFFFFFFFFFFFFFFFFFFFFFFFFFFFFFFFFFFFFFFFFFFFFFFFFFFFFFFFFFFFFFFFFFFFEECF05A0500040500030F00092847100F6FD2F8FAFEFEFEF2B21B813A0A090A517E9DF9F1FEFCFEFEFFFAFDFFFFFFFFFFFFFFFEFCFDFEFECAB1540B00000231E5F9FEF4F2FCFFFFFFFFFFFFFFFFFFFFFFFFFFFFFFFFFFFFFFFFFFFFFFFFFFFFFFFFFFFFFFFFFFFFFFFFFFFFFFFFFFFFFFFFFFFFFFFFFFFFFFFFFFFFFFFEFAF8FAD35A051102008A4351BFFBF8F5FEFFFFFFFFFFFFFFFFFEF7DE840F03090112611B0D023E2A000206204D0815521E282C3C1A00000308279E9E0C941D4FC9E8F2FCFCFEFDF7F0FEFEFEFEFEFFFFFFFFFFFFFFFFFFFFFFFFFFFFFFFFFFFFFFFFFFFFFFFFFFFFFFFFFFFFFFFFFFFFFFFFFFFFFFFFFFFFFFFFFFFFFFFFFFFFFFF7FC8D81390B0C000002020D02010B020940B7FEFAF6FBFEFAF9CA732E1D09003B4160F5E5FEFAECFDFFFEFEFFFFFFFFFFFFFFFEFAFDFEEDA68A3E0D09000547DE62E1FBFCFDFFFFFFFFFFFFFFFFFFFFFFFFFFFFFFFFFFFFFFFFFFFFFFFFFFFFFFFFFFFFFFFFFFFFFFFFFFFFFFFFFFFFFFFFFFFFFFFFFFFFFFFFFFFFFFF9F9FCCE58090204000A532A7ADDFEFEFDFFFFFFFFFFFFFFFFFFF5F6C0520F06010006130401040E080205060B08040608030805060605051B5CADD53606481E0DA1EFFDFCFCF9F5FEFDFEF9F8FDFFFFFFFFFFFFFFFFFFFFFFFFFFFFFFFFFFFFFFFFFFFFFFFFFFFFFFFFFFFFFFFFFFFFFFFFFFFFFFFFFFFFFFFFFFFFFFFFFFFFFFFFFEFDB3D3C75301010605000500062766230D6ACEF6FBF9FEFDF387270B040607090D2B5D7FF8FDF5FDFCFCFEFFFFFFFFFFFFFFFEFDFEF7FBCE691902030100176CABF0F8FDFEFFFFFFFFFFFFFFFFFFFFFFFFFFFFFFFFFFFFFFFFFFFFFFFFFFFFFFFFFFFFFFFFFFFFFFFFFFFFFFFFFFFFFFFFFFFFFFFFFEF4FDFFF6FAFDF9F9CC5B3287090300050F0A52C6FBFAFFFFFFFFFFFFFFFFFAFDF8FAE07B367319000004000308030606010103040000000004010004042B8DDEE8DE30090B0D62CBF2FAFEFEF5EBFFFDF2FDFCFDFEFFFFFFFFFFFFFFFFFFFFFFFFFFFFFFFFFFFFFFFFFFFFFFFFFFFFFFFFFFFFFFFFFFFFFFFFFFFFFFFFFFFFFFFFFFFFFFFFFFFFFFF7FDF1E99D86580105060530A16A48AB4E67AEDFFDFDFBF9FEF9C43A0704020200466838B1FCFEF5FEF6FCFEFFFFFFFFFFFFFFFFFEFEFDFDB43A060305060006265EF8FCFEFFFFF9FAFAFAFFFCF7FFFFFFFFFFFFFFFFFFFFFFFFFFFFFFFFFFFFFFFFFFFFFFFFFFFFFFFFFFFFFFFFFFFFFFFFFFFFFCF7FEFBF4FCFBF8FBB8B1A2191A5D030C1006010F79DBFAFCFCF6F9FFFDEDF1FCFDFDFCF5C35B214C1300048A1808040001010406010500000D0B0202002473877FC2F8E9A6890E58C8F3FBFBFCFBF7F8FBF7FBFEFEFDFEFFFFFFFFFFFFFFFFFFFFFFFFFFFFFFFFFFFFFFFFFFFFFFFFFFFFFFFFFFFFFFFFFFFFFFFFFFFFFFFFFFFFFFFFFFFFFFFFFFFFFFFCF7F9F8BE797A6303040F1864747F739FD6EDF6F9F5FEF8EDF995230D0401000030481F7EF8F0FCFEFCFEFFFFFFFFFFFFFFFFFEFEFEF9FDD5AD5B140105337B0A20FC71E7FF9BC5FBFDF9F8FEFCFFFFFFFFFFFFFFFFFFFFFFFFFFFFFFFFFFFFFFFFFFFFFFFFFFFFFFFFFFFFFFFFFFFFFFFFFFFFFDF8EDF6F8FAFDF1F7914D46080A160C860F030116A0F3FBFCFCF7F5FEFBF6F8F1FBFEFEF9DA7E181105000461136E2906050333512562116708050703375B88AAB5D4F3F8EBE19DCDF4F9FCF9FAFBFBFCF8FAF9FCFAFDFEFFFFFFFFFFFFFFFFFFFFFFFFFFFFFFFFFFFFFFFFFFFFFFFFFFFFFFFFFFFFFFFFFFFFFFFFFFFFFFFFFFFFFFFFFFFFFFFFFFFFFFF6F9FCF7EBBA804101040D041953B9457DD9FBFDF7F4FDE69A7C993A0C01020708101531A9FBF5FEF4FDFCFDFFFFFFFFFFFFFFFDF4FEF7FCF1AE51140103214D0C098718697795DAF4F3FAF6F1FFFFFFFFFFFFFFFFFFFFFFFFFFFFFFFFFFFFFFFFFFFFFFFFFFFFFFFFFFFFFFFFFFFFFFFFFFFFFFFCF1F4F9FBFCFDB9904D150E0702137B7B7C10061AADFBFBFAFDFCF7F3FBF1F5F5F6FCFFCF8F4E0C02090A7E210D531E0D5282A8BE9BC991D3680907034A99BAD2EDF5F8FEFCFBEDF1FAFBF9F9FBFEFDDC9AFEFBF5F8FDFFFFFFFFFFFFFFFFFFFFFFFFFFFFFFFFFFFFFFFFFFFFFFFFFFFFFFFFFFFFFFFFFFFFFFFFFFFFFFFFFFFFFFFFFFFFFFFFFFFFFFFFF9FCF6F6FCEDB2512A610402021E6E1E3389C5CBD4F6FDEEB03D5A2E064A7A410B032D768FFEFDFEF8F3FAFEFFFFFFFFFFFFFFFEF7F5FAFDBB8C531A01020B1307021E09151559CBF7F6FDFBF4FAFFFFFFFFFFFFFFFFFFFFFFFFFFFFFFFFFFFFFFFFFFFFFFFFFFFFFFFFFFFFFFFFFFFFFFFFFFFFFBF2FEFAFCC0914B24110605041C7DD88E710E0819AEF9F8FAFAFDFCF5FEF4F5F9F8F8F49E5F4E19020916790B07170B094B82B8F4E8CB84886216052082D8F0F4F9F7FCFBFCFEF9F4F8F7F9FCFECF8E704AFAF5F7FCFEFFFFFFFFFFFFFFFFFFFFFFFFFFFFFFFFFFFFFFFFFFFFFFFFFFFFFFFFFFFFFFFFFFFFFFFFFFFFFFFFFFFFFFFFFFFFFFFFFFFFFFFFFFF7F5F6FBFBEBAE6E7D5F02020415050A385770A0C7C6C48424180D59737D3D07002B5C54F6F5FAF9F5F4FBFFFFFFFFFFFFFFFEFCEFF9FDBCB77B230001041E660A000C020015A8FBFEFCFAFDFBFFFFFFFFFFFFFFFFFFFFFFFFFFFFFFFFFFFFFFFFFFFFFFFFFFFFFFFFFFFFFFFFFFFFFFFFFFFFFDFBF7B5A04F1B06010004071F7EE6F2E17310040E76D8FAFBFBFAF7FEFBFEFCFAFDD4A65E464C1C00020E130C010C0805124C8FBED48964100E0F235EA0E2F9F9FAF7FCFAFAFAFBF5F7F7F9FDFA98171536F9F5FEFCF4FBFFFFFFFFFFFFFFFFFFFFFFFFFFFFFFFFFFFFFFFFFFFFFFFFFFFFFFFFFFFFFFFFFFFFFFFFFFFFFFFFFFFFFFFFFFFFFFFFFFFFFFFCF6FCF8FAFCFCEEB8756D0303030007030D0C25546157511B09094A837E73370800091429A2F4FAEFFCF3FAFFFFFFFFFFFFFFFEFBF6FAF7ECF1BB551002041F7A0B05070003089FF8FBFBF8FEFEFFFFFFFFFFFFFFFFFFFFFFFFFFFFFFFFFFFFFFFFFFFFFFFFFFFFFFFFFFFFFFFFFFFFFFFFFFFFFEF4BA431F070506022A30257BE6F9FDF29B17020125A3F8FDFFFBF2F9F7D1D2F6DB7932150B07030003000158180601050A4F696CC77A780100054AA6A5D6F9F8FAF8FEFEFDF8F9EDFAF9F9FCCF5E08205BB2FDFBFBEFFAFFFFFFFFFFFFFFFFFFFFFFFFFFFFFFFFFFFFFFFFFFFFFFFFFFFFFFFFFFFFFFFFFFFFFFFFFFFFFFFFFFFFFFFFFFFFFFFFFFFFFFFBFBFDF8F9FAFAFCF0A80A060C02000A060304070D100901000404AC8C8696480B0401031771ACFEF1FCF9FDFFFFFFFFFFFFFFFEFDFCF8F4F9FDD06C1A0706040F0008000305056ED5FAFDFAF8FDFFFFFFFFFFFFFFFFFFFFFFFFFFFFFFFFFFFFFFFFFFFFFFFFFFFFFFFFFFFFFFFFFFFFFFFFFFFFF5BF48090C02010502557981E2FBF3F8F0540B04093A86D4FEFDFEFCF4F8A686AC741E01000204020F3A0C03942C050301000D1345AD771004114B7DDDF3F8FEFCF5F2FDFCFFFEFDF7FBF2F7FCAC503249BF65FBF8F9FDFEFFFFFFFFFFFFFFFFFFFFFFFFFFFFFFFFFFFFFFFFFFFFFFFFFFFFFFFFFFFFFFFFFFFFFFFFFFFFFFFFFFFFFFFFFFFFFFFFFFFFFFFCF9FDF3F7FCFBFAE39F3B0316010102142F1A04050004060001027EAF4A2B1207070001259E5BFCFCFEFBFDFFFFFFFFFFFFFFFEFDF9F2FCFDFDB2470F07010203010003070101106FBCDAF8F4FFFFFFFFFFFFFFFFFFFFFFFFFFFFFFFFFFFFFFFFFFFFFFFFFFFFFFFFFFFFFFFFFEFDFDFDFDFEF6CA540A0000000422152086EBF8FDFBF1FCB6430F177F81AAFBF9FCFEFCF8DD882D0F050402010709278C220012020008061F1B01070F08010A2FBCF2F9F9FCFCFCFEFFF4D7EEFEFAFDFBF8FBF8F3EEAE71FBF4FEFBEEF9FDFDFEFEFEFEFDFFFFFEFDFBFEFFFFFEFEFEFDFCFEFEFEFEFEFEFFFFFFFFFFFEFEFEFEFEFEFBFDFEFFFCFDFFFDFDFFFEFDFEFEFEF7F8FDFBF9FBF9F0B669B9090000030535703C040300040B00062F160A080300040401030432C6F8F3FCF2F9FEFEFEFCFCFFFEFCFAF7FBFDFBFDE47F2F220C060206330D01000007071466CADFEFD9F0D7FAFFFDFEFEFEFFFFFFFFFFFFFFFFFFFFFFFFFFFFFFFFFFFFFFFFFFFDFCF7FAF8F5F9D66B0600010104017F698DEAFBFBFAF8FBF73DB6290127A5F8FAF9FDFCF6E96A736E1A1314070D170C04010503040201040D726305020104020C2A5DE3FFFDFDFCFBF3E8C9589BDCF0FAFEFAFCFCFBFAF8F8FCFFFEEFF2F5E7EAE5EFFDFAF9FEFEFBF8F2FCFDFEFAFAFDF8F4FDFDFCFEFBFDFEFCFEFEFEFCFCFDFBFDFCF1FAEDE7ECF8FDF8EDE9EDDFF6FDFBF8FDFEF4F7F7F5F8F8EDB926030200000503040300010404000BA1520E15121724261702020C3AFBEDF9EEF8FAFBFCF5F3FEFDF8FAFDF6F9FDFE99355E842000020ED015010100040301041C5C9C44A348DEE9F4FCFDFDFFFFFFFFFFFFFFFFFFFFFFFFFFFFFFFFFFFFFFFFFFF7FCFCFCFEF3EFD5570600080A050B1086E5F2F2FEFCFEFEFDD524010021A6F6F5FCF7F9F2990C0106056F6E23517E1D010A1710020F0800000202060E07010F51D9DCF8FEFBFCF3F2BC65884C081ABBF4FDE9F3FEFAF9F9F1FAFEFDFEF5C56C90619CFAF8FCFBFCEFFEFEFDF3FAF4F6FEFEFAF7FCF9FDF1FCFAF1F6F3F2F9FCF3E9F6EDF9EF7E59C6EFE4E5A4669654C9FDFEFCFEFAF8F1F5FDFEFBFAF6BD210200040107040002060501010000000D8E6098D6D99810092DE4DCE9EAEEF7F0F0F8FCFCFBFCF6FDF8FDFEFDFAF48F1C0100031000080C03000001110F00010208011048915EC9FBFCFDFFFFFFFFFFFFFFFFFFFFFFFFFFFFFFFFFFFFFFFFFFFEFEF9F7F9FBFEFEE06200000A02965492EFF8F6FEFBF1FAF7110C0101199FFCFCFCFAF5FCF2851B02070309070C1004080014ED20020000050402020307010C0A0167E6F9F8FCFEFEFBCC54000FBFF9EFF3FBF6F9FDFAF9FDFCFCE9F2D943954137AEFEFBFAFDFCFCF1F9E9F4FCF6F5F9F9F9FCFAF9F1FDFAFDFEE55C88F4D39566A061AF14087D832E02030747A35DC4E2FEFEFEFAFCFEEBEAF4ECFEF9DEB50F030F050201060700051612030A05040988E887150B0000010C131BC149A8699CF2E6EFFDFCFBF0F2FDFEFDFA9228040100030E0108000002060D7E7306030609050608040B52C4F9FEFFFFFFFFFFFFFFFFFFFFFFFFFFFFFFFFFFFFFFFFFFFFF7F9FCFDFBCE5F120800000501191279677BA3291B3BFDF7EADF2E0023A6F3F2FAF9FEF08314021507000000000100030B01130F02020A0E06000204082136A237CFF1FBFEFDF9F5F7D054090FB6FAFDFEEAF4FEFCFAFBFDF8F8F6DB5EB04A0A043DA4C18495FCE96A9377AA669A838BDDF0F2D345ECF32B1414180E0D19191C5FB0081A01000E0C030A0901070D0B2220CCF3ECF1EBF1F9FDFEF7F1FAD4201202010000060702307F5001030100000C9377877907060A030001020F25A86581977F97668B88FDFDEDEEFAFEF58817010910010000001CB0480D2485820001040748A9133AC5FAFEFFFFFFFFFFFFFFFFFFFFFFFFFFFFFFFFFFFFFFFFFFFFFFFAFAFBFAFBF6C0480D001306010202161116688B39CCFEFDEA3E0D1026A5FAFAFDFDE88E788D0D6466280A0F0E040100010000020006914908120D0001050B1C1066E2FDF8FAF9F6F9F0D1BF49ABF2EEF875A5E4F2FDFCF7CB57F3F6A8420D0B2F7A7123337F7E241B9558261219737C36282943AE242513AD34000A08000307101C06020207020502010100050A00030221252C32AB7393F0A46BEEFCEDC29A1D120500060E0A0E271B0000000B0804826628180201010001191001111E1013161219152375F4CEF3F797978B420A0238971D000E00A14D91CDDBDD89100C00040E1C1B0E2075D6FFFFFFFFFFFFFFFFFFFFFFFFFFFFFFFFFFFFFFFFFFFFF4FAFEFAF9F9F9F3BF51000501060E0A0A06518E749A56DAE9DF2503178FDCF2EDE3FDF6D7D88E0E231D0D0350692002040312060502012A1402050400010300000247CEF3FAF7F3F7F2C186B8B848C75BAB4798848C9A72E6E69648D35391280213231806111D1D0E032C160306081E1C06000315390C0C05300E0103030302010618160A0305000103010B0D0105080007020302032415203A649065E2BE342E08050100336745040408190607030200281E05080A01050415AD0F100701140D040B0F0D041A37403A90BD7477662204112E090A0400380F2B3BA4F7E1B87D1701090006043BAEEEFCFFFFFFFFFFFFFFFFFFFFFFFFFFFFFFFFFFFFFFFFFFFBFDFAF7FAFBD79DA3690F0900110966480D216988360F5145430C051D918E98A393D89BA8EC82170C0B04001F2A0D02061B8C1302000916080404030B0D010100011580D5E7E6E7E4B77D7E59440D38143816352421281B444B420645113710060501000103030E050004080101070404060205040404040902000002040706010A3DA870080102121005070E0A0403010A000B0219130402051E3518475E9D0F000E040416291C010421693E050B0904080D0A04040202030D4204060B1A1317161F55670A0203070022422429260D0205080903001F0A120D0A7C9585A7577A00030F010B1B68CBFBFFFFFFFFFFFFFFFFFFFFFFFFFFFFFFFFFFFFFFFFFFFFFCFBFAFCFEE8AF80833D0504080335280801274804070A0701000A0E343F726C2B5B687C937650561A040201030505020E480E0402056D3C07020B3F64010200003A7662545455533C26331403030103060304020305020105060306000302030507111E201C0E0E10010308020204050909020104240008150B020E0C0303010522563A050401465B2E38522A110A160B0E1A0A703516100B1009141025550F030802020201010309144420010604040255500C0002030303010C121D754A6E915648755005010111090E0F080402040903050505941A752E062F282037164203000800060A1D73D5FFFFFFFFFFFFFFFFFFFFFFFFFFFFFFFFFFFFFFFFFFFFFDFEFBF6F9FDE9B2852602020104120D02003561060400061C04050A0515382C000D30352B3D353B1202010506040D18010323070801462A0401092E45010000052A3D1004020001020000041502030F04150A131403020C1700070500010300061A5893917A3E4A420B010C040C0D050F201819159721237037115F48010200000504000308012B70604B453E56286D20708B2048376654546431772F7D06130E100402070B090505030C0102030002023A3B0D0D1903070001106C3BAA6C81986A582E37070F104C45664D1D17140509000004085F134E1E01060D0C05030A1C130000003A4B6DC3FFFFFFFFFFFFFFFFFFFFFFFFFFFFFFFFFFFFFFFFFFFBF9FDFBF7F9F9F1EBA22C0300020B5C3F090F42520606040B6E12020808151204091314000B270509100500050703265D0F0F740B06010106140F0B140A170004050B040004050715211504215E341B7A296E4D61591E132B75021C190103030607184C777666373E4610020608433F204078837038DC8082897C4A9271100B171200060002090303333E46473B522D591A61771D03597B82BFAC63C848CC3C713B56220626412C060005120E020300130D02040831590C0500040D5C369B58353434450D020F47505C7FC68B67785B2414020B040409030103111346561823408759131102326CA8E2FFFFFFFFFFFFFFFFFFFFFFFFFFFFFFFFFFFFFFFFFFFAF8F9FCFCF9F4EADF9C450109024681480C384A030708060D73230A1B28503602153B440B3473424A401201000D182C5A13136D0E0007160F4349415337590005010E03000A162E5D826D4377909283EB57BBB7C8C18C5F51D02D8266201C150F120F0401040501064E3F1609094641456E7D8D7F74827F7C73768087804D3F754C1D0A0B130D0F080201304A2832562B1E00150B04543C89F7CFACDC94DDABE58A6A250428472D0401134A2B0305004C2E06040832590E000004080112532804060107120C25568F88ACD5ACB1B3998156211F1B040C08060449498BC57F8CAAD58B551C121E65CDFFFFFFFFFFFFFFFFFFFFFFFFFFFFFFFFFFFFFFFFFFFFFDFEF7F9F9F3F6E9A86A2B01210D68390B053F53060702041711651E677877400114404F1244888384551602012C4A17000507020B040C4F274E7B788073670004101206022A4371ABC6BAA9C58FC9E1F497CAF5FDFEFFCF96E199F1987B7B57283C280003050102085D823D140503013A4907071A73040A006C215981A38170C3887A37435535491700020301216B858E750C5C1501021292FAF9FAF7FDF8FEFAB23C09000205050A11195B2E0509015E390A12170B0B020005000007020714130A05061544706CB8F3C58A828373A1CF978470712C020A08045989BBF5F6F8F5C37A6C0B2A699DDDFFFFFFFFFFFFFFFFFFFFFFFFFFFFFFFFFFFFFFFFFFFFFCFCF8F7FAF9F2D79A8B31044B2A080209050103020514117522DF4BBAB03B0C19130A170B10134847100C0E053A5B170C1B030506080166330D49472D6304040A294A1107386494B097819A898B978EDEEBEDF7F2F3F4EEF3EAF7D975E7F0924C64450E0A080E020006553C0B040704050500030800040000110A035480484292A8D585AFA57E972E000303012C6632C0971E821F0406067BE8E8E9F0F7EFFAD58C6F1F00060605213E170A050200010C0502313D100305020D001A150510492F0204143F94CFF2F2E4AF3B120D3EA6C39EF3E9E8820B0000050681E0D8DDE8EBD9610A012680BFEAFDFFFFFFFFFFFFFFFFFFFFFFFFFFFFFFFFFFFFFFFFFFFEFBFDF4F4FDE7CAC85B020076400D080A0F0B070D1A5718ED86DF89CB80233A613D31683723313F330D373715030007275311000010060703000A130D03010B042068130103337AB4AA5C2C18080044BFACD4F0BEC2C9ADD5A7E99DF5EBD481656564593E2E36100A090905010201080B06020409020E020102120B01050F3DA0E2BEA3EEC2A88A310C130F040102206D010F040E16131250878991BAF0FDF17D7B7223010403001E554103010600050B08034F591114000300037013021461420702124B86A2B6C1AADE81563D3980DEF7DBA1DBA0210301000A4F929B8F92AFC73F020309207FD6FAFFFFFFFFFFFFFFFFFFFFFFFFFFFFFFFFFFFFFFFFFFFBFBFCFBF7F2D0926D400E0803482715396A45294E44D2189497979D9988516081534D8D683954584B1D4E8A60170407227A1C010064230A0603022B4301000210110200001021255B846360230002868E7172C46E7C9158B12FD60AEB91E8FFFEF4F6EAB4754F353B00000400000234340E111101060800081205090B0A1A5395A0C5E5FDFEFED0833F30302113090306070100556F6E2C020103278AF4F1E26D7D1B0A05000304010305010800010B3F2F1104012B510411021283200706010003030012548D48846A8667929CC0CDC0D78F6A98626618140D09050302021F72907406072F67ADE1F9FFFFFFFFFFFFFFFFFFFFFFFFFFFFFFFFFFFFFFFFFFFFFAF9FCE5CED2B13B010A5318BC9D67BCF68567B03FE973030A0503030101040705051528330C031535324C4B21030A1A00000A0096230001030A4B6D0B0408030B0A0102013A54516B66190D080A00117C2B5141494C2E661F7C08895EE0FEF7F9FEF1D0A78A8A5C00051706020A5A91592E5008135202035021053A320602030D2DB8F2FCFBFDEC8F8DB1916E5B0817000F017DA1AF4904030237A2968C7E8D14000000070700071200030A070B02673D0B0003406F010501000000000001060601060E05000352525931578FCD943D58286E0B8D9D746D5E62310303194D72040014024C837BA0E1FFFFFFFFFFFFFFFFFFFFFFFFFFFFFFFFFFFFFFFFFFFFF1F8FED39571420201067C2298E7F7F8FBBEA7D79EE58A71676C676B273F6C3C0B2E4D4A5E3074424E471D0804001E541C0000000A030301010204050007286302000301090805052D5F2E0300000A010504040305090505050A0575EAF6F8FCF0F3FEFCE8AF6A1D0807373016040445573B65001CEC0961DBA175877F654638747FA9E3BFE9F6FBBEAAD8F6D69A67010001040607080A05000206080B0D0D060A010105030100000011592C01070509070200000409020700000C0303060100113C3F0001060403030504050670952F06080501124FDEF0F4EDA038273C03144C04000D06336D9DD8FFFFFFFFFFFFFFFFFFFFFFFFFFFFFFFFFFFFFFFFFFFFFAFBFEE09137070000001940D3F7FDFBF9F1F0DEADE182EDE2E9EBA07A89874C438B511D3C968F4856363C411E002A6E0F07040002000402053E4300020D1E7B0300374618000707050F1100075B2103014A25463F2F6E6B6F7375EBFDF9FBFBFBFDF9F7E5B0BDB65C279E66210701070B0A3563109070DCA699D2A17E804E3C909191566DC0C9EEF4F3F4F6D6A4DA7C0B0000080202053A3D26411A650114590E000000020201000113653401070501050200000205030A00040100000002001849470208030A0000010100033C49130302070872DDFAF7FCF8EFB0A0BE6E9ED2450201000569AEC2E0FFFEFDFEFCFEFEFEFFFFFFFFFFFFFFFFFFFFFFFFFBFCFEFEE29131000002140254919094D5C9A5A0878892D98CF5FEFEA3784520203D590F072B50290F3F3E4D4A1E0008180B0C610C02000006063F490200020623000427512001010B2D5F78797C7F79807BC5A4867D98DAE5D59BE4F59BEAFEF0C4CEF4E6C4C5C1C4E19877873A080000000026740E32DF952E1D68825024130F271E1B0D455550A6EBD7CAC4B4CDBCD462821F060000094576A0BE6D6A5E29B2304B1D0001284725030F130F532C0900070404344227660206075E21552302080D100D07180002000A42381504111D100205021170F2F1F9F0F6F5E9E7EAE2F0F7C34F12010338515182D4F7FCF9F7FEFEFEFFFFFFFFFFFFFFFFFFFFFFFFFDFDFCFDE18A270512114E3620221F25645629253A6B3D891E92A789783C0D01071215051915100C020F1115130B080806020A7A1502020D0703121303050206030203091E0C00070A1C598F827634B7A884908E507FC8A08D89919C9A8591967C60A3C280708E8996CB907F9644113C3E12080C200A0A7F2F10021423140305090301000010103C85925E7C9D8CB16BCDD69214010803495153B0ECC49ED47EE5915C1D06022147340703000A64280602064070807C6A7C7C7E5B656A7971280202020613393C070003493815030205040309125A91EEF9FDFAFAFAFBFCF6F5FDF7E6490A08030E112E82DCFCFCF9FAFEFCFEFFFFFFFFFFFFFFFFFFFFFFFFFEFEFBF9E2953102080D4A2F0E0A0204171103030E220C230E262E27290E0308070102090B0601050000000103040303050D132A0C0302060205040204020103000A01060601010406102432322B0780432A2A2C1953702C22306F29297434291D225C5B1F264C53545C3159461A073D411003040903012F0805010404010002060705060001053548381D374B475D27809432060A07044F41185797948A8C8494772E0907040C19120304040421120502024A7C5A6DA340858ED097CC96722E09040104123A5303000022150604030507051A6F897B88EDB39C9CD2CCBDE1C1DA99A1740C0A0301012380DDFDFAF8FDFEFCFEFFFFFFFFFFFFFFFFFFFFFFFFF7FCFDFBF3CA782300051D0D06452645233031223D1C4A115E054D5D1B00070D25352B603F022D572C392D0E2F392A391D41115C08457465223D3C28351B54635B5E604912202D3135322A16050D44410C000804041A1E04000D3300002C0804061C4432133A3D1B1813031B110607191704021D4E00000B0102000438320800010403040007051D150A070F12141F0C2C2701020602021E1A04112C323125362B251104030609213C28040C0001402A09061F3224416D0E32389782A77C3A110605060607182200030407030002060705031A733D2D2984432D2D716B527355862C3E5E0A04012A3B285AC9FBFCFBFCFEFAFDFFFFFFFFFFFFFFFFFFFFFFFFF9FDFCFDFBE097360001083A607A39535F9599888C36B865C85BBC8E64512D408DA283A3695E8D9D6F747143465E83A77EAA60CD5CB39CBC8472788FA485C8D2CECBBC7B62839C816278A0843E1EB3213B585B6066643C04184514595B451439243E65635F6B6B58291F5C1F364D27322E263239A60D5D0251123F217661355B310203020405012E34131E2D3031243C1155120401060439302A38193B0C481F400D174C462B244A745C3429381A854440582E092D483748060F2636382F0D02010105112D370003020306313110030503000A2A00011120070C061C2014160D2B0608230C03003645396CCFFDFFFEFCFDF8FDFFFFFFFFFFFFFFFFFFFFFFFFFFF9F8FCFBE0A1572203075E87511E1F83DFEBE5C57CF5CCFDCCF396B88D3957ACC9D3BF8FD5DFBC8C6167522F61B4C1D6F7D2F6CFF0999B955E5E97AFBDF9FBFEFDDB899CC8BE948CB0D2D6AF70E863A4CFCEA5969272597DAC68CFCD825C72719AA6B0C2987DAC885CB2816B7667716B81997CDD62CE56C166AB7EB19768824F0205020500023F6E5D7E7B565F637156871D040106085269677170A96EAA275D14226E7C7C899890718184AC7AC27CA6986A5C6D6A5C794C0E040905030605010304114A7C4F552C0A074B4A1800021C33144106551301000C0601060D23371934173E10030717306DBBEBF7F8FEFDFDFBFDFFFFFFFFFFFFFFFFFFFFFFFFFCF5F8FBFCF1CD893905082A5A25070387EDFDFBEADDFCF7F6FAF9C6B67D2C28649CCEE5D9FDFBCC66211E1D0F4C7B74E9FBFDF5FEF3A943311F172C3D68B4CCFAEFAE6D6686594F868E6995E3DBECDCF6FBFBA95735599BBEECDEFAF18C815285BDCEE4F8D18B969C6174C49A5A73625BABE7D9F1D6F1D4F4D8F0E2EAD6752A25000203000004206EB1CD8A3447827F883F0E010005092084A998D9E79F980D2607162C4993DEEDBD8A95A6B9D2F2D9F2A78799726BB17D7502080007050806020506052E86B28B4E1101201E0A01002E6D60B94F9C0E0000030303020B407E729C6A8B2C31170545AFF1FAF3F4FDFBFBFEFFFFFFFFFFFFFFFFFFFFFFFFFFFAFCFFFBF8F9DA770F02122B81480E2585D6F6FCFBF5FAFCECFDFFF48A813F0E4F92A9FAFCFDFEAA260601001A382735BEF1FFFDC1F55F0F01212408000E458DE6AF5B201416090C2728185AD5FEFEFEFCFAF8C24703183570F0C0F7BC843B0A2D72C6F8F1D9BE78311628B1C35122112A8FEEFCF5FEFCF9FEFEF8FCFFE67202040000070300173379BBAA55111051A4730E0100050C04076CB0BDE8A95224000A0101041666BDDDD3AC5B353093ECFDFCC16A363281E197460603000101000504000823518EBB59350A000303041B333556C1EF9F3F0A020407060200063386B8FAD5A0546A1D0640A2E9FCFEFEFDF8F8FDFEFFFFFFFFFFFFFFFFFFFFFFFFFEFCFDFBF0E2BF6A14030E815F22066D6193EBF7FBF9F7FEF6FEFAD9C06118082E7CC0E5FBFCE586290E1B1857661E1761C9FDFA53BF2B0020736D1F000943B7C8728B1A07141C1106032284D1F4FDFDF7FBE0862D0500043FBA68EB4F460E001E89DCF3D89D8A6C06206DD58824020155CAF6FEFDFEFAFDFEFEF2F7FCD663050D141F0E2F2D6BC0E5A7917B330139D5AC1603000A040604184092C2628901010B030D1D3B9DC09487711400052BB2EAF29F31024BC6CA74B02E0026090002051D262B70D8EAE1B156140503060341A29465E5E2650401040D092015001150927FFBFB9E2D1A06182D73CFF9FDFCFDFDFEFDFEFFFFFFFFFFFFFFFFFFFFFFFFFBF2F6FBC7879BB44708022D1201051C0F6AD2E4EFEDF6FBF6FEEA79A8460F0E00114797F9E8838E6F1965663A19073D8F61F1D8E144081683787D751102196CE0D1C323014C6E3D130C0D2B43D4F8F8F8E2738E631204050A4BDDF6CFA63A090649B2E6E5C55D110C8471A3702906006AE8FDFCF8FEFEFEFFFFFDF4EE84160213538348C8BEE2F8FBEBE69D320837B95C2C05020006190D0E5DC1D1922C060802092C6D64B0EFC0630B0003069B4BB2D5E76A0C1D70D7E3F7C820BC1E01081C8EB76E82FBF2E29E32050B111A20267AE0E4EDE0BF02041B211F8557040C2874DCF1FBAC260402389AD9F1FEFCFAF9F9FDFEFEFFFFFFFFFFFFFFFFFFFFFFFFFEFEF5F4C25112250E0007110500010005131E6BE8FDF9EBF0FDF4BD510F0003051B6EDBE8906B7B2E071418090B0F0D4FC8D534D03205956D9CF3E481030D56CC4DB02D09121553793E0226B2EFFEFEFE6E9363180404000333F3FAF152A13302033EACD2771903142098759B5A160668DFFDFDFDFEFBFEFEFEF5F0D9690C03030B1D6CC52D82E6F2F3B041101939C4D8D22107110F050907172E2E5EA002060F030D1C4FAFE6EDC243071500130DAB79A669125B9068FBF2F5E9DD340A032BA4EFF0F7F9DEA387540F0B67BDDAE0E7F1F8CD43251001BAE44F1812030151F0FBFDFDEE9E1B05042169C5FEF8F0FEFDFDFCFDFFFFFFFFFFFFFFFFFFFFFFFFFEF7FAFEC24F0301050000000000170D0207649373EDEAFEFDFCFD4809010407040619397FEEFB93170006030303013CC7FAFBE14507081E94E1F4E78300031B54B32D05000042805E17052CC1FBF5FBFEDF671B100900000440FEFDEBC2300D0230868B767501187EDCF0E6701D12096EEEFDF8F2FEFEFEF3F2FAF3747B7003050104101064C5E6F8FEFCAD350039C347150500ADDB4C9475010C51A514250000030100005496948E6307000007A9619E672207136CDEEFEEFCED3DAD2A19A6F7FBFEFDFCF6C8721F0F7875A7868BD7F7F2F8DFBF000B2422A85C0503031A7AF7EFFCB5170B03003EB6F8F1F4FEFEFDFDF8FDFFFFFFFFFFFFFFFFFFFFFFFFFEF1F8FEF8EEBF52010400009B56120203010E1D5FF5FBFCF03DC1330D051C7469001667D8F78575832100677C2C00051B50E9F21E01050CA0F8F7F5E4890D38EEEE9C2904030C1B140104289FFEFEFEF2711902026C880F03102246EBF8D4A732061448927F001384FCFD8F95521C667573D6FEFDFCF2FDFDFCFAF8D3721A050B010401011275E3FEFDA9400F030A44B02A03091EAFDF7315060950EDA20300004B8830379DE5BE51170D040700241815906710015FFCFDF9E8FCC638110534B1F3EEF7FBFEED8107011811171A6DD1F3FBF7DE4313010212210F000F6FD7D3FEFDFDE3931A000D1256C4FDF8F1FEFDF9FDFEFFFFFFFFFFFFFFFFFFFFFFFFF8FAFBFDF9C57A80B53B0600C55C1C0504020609133655FDF9B052AE3A000B191A0D0318508AE68B2A05657A7C662C12012BE245B62B050129ABF1962F28081071D82B070401080A03002F8F72BDF7F6FB5C0E0A06181F040301052FF5E4DD5B1403010A1716070827AEF9CC7D30091E6EC8F8FAF3FBFCFEF7FDFBF9E76B0606134D834DC5C7DCF8EEF6F08C2007022BC4260D0204BB711708060A5CC51B0308062CC6CCCBE8FCF3B13B0305010804020A1F130A638E67FCF8FBFEE02E041E95EBFCFAFDFCF7E6800E00050203031777E4F6F7DB14030009040703030A2988F9FCF2FCBF2C0A064499D6F2F8F8FCFDF5FDFEFFFFFFFFFFFFFFFFFFFFFFFFFFFBFDF6F6FBECD7E0F5C1330B2F1580490C01020407062FE04DE728301007060A0D061A5FBEC8A62B0A0374C5DC982B05000D46AD510F030078799AD472050E41C7B82F3A6F30060505020E2A2EB8FCF47B1E0302010609070700052EE7F5F7AB3000070700041F66C0E4E8F37119010369E3FFFCF7FEFFFFFAFDEBF5EA69001556926DC9F5F4FAF0F5FBB08977292B995E0803009069994C030F549462950006298A7CD0F9F3C59D9D6A0010893F1402020300021F6CC9FEFDFEF8F7A9241243B0F9F8F5FDFAED7A0F03006C5B1A015BD9F6D863961E04059138000203082B87FCFAFAE6891801123881D0F9FBFBFEFBF9FAFEFFFFFFFFFFFFFFFFFFFFFFFFFEFCFAFBFCFCFAF9FCE25C050008371E0803000102000E4419E31F040705010307050B2853B2930D074EB5EDB387701F0719A0F01F090200332E84AC2C00042190B61D1F2E4E5F2F040500021EBAF9FA4C09040207084971060D010F68E4FCCF4302030504020C3AB0FDF6FD6A1B040168E2FBFCFDFDF8FAFDFEF7EA917B5407307FC6F8FDFCF4F6FCFEDE9F5111123D140B000539947523000A66D9903F00051566878CA5A3B2A77A7F3D1BB34A813B0E050100025FF5F2FEFBFBE965121981DCFAF9FAE7A89D8A080A017CBEB3AACBF1F9EEA84D0C0104C24A6F5001031765FBFEFDFBA42002003CA0E7FBFAFBFEFBFCFDFEFFFFFFFFFFFFFFFFFFFFFFFFFFFDFEFEFDFDFDFCFDE65B0408070401080600020400000312DF1F08081F583F0808011148D9D258162D50613A2F3912257E73EE1B841F010608323700010D39C3BB1207014EADB7772D00000E4864E74209045E9C9CBEA3030C00052AE2F6E694340604071456A7D5FEF6FD5E5F491768E2FAFDFCF3FBFCFCFAFDF0B67430001774F8F5FEFDF1FBFDFEC86015031F7B8F21010402442A0604498691D77B000105203A2B31326D9D6732221FB1B4DF5F130200003FB7FCF8FCF2F56E8E231158C5FAF9FCEDC07A805D0F53796E5E9FB69EDFE07B8B020B03C04AA07101020A349BFCEDFD9F1D013598E2FBFDFCFDFEFDFEFFFFFFFFFFFFFFFFFFFFFFFFFFFFFFFFFFFFFFFFFFFFFFE850568C2F0D60460C0001070B02168F6F170F6996C4B77A460C0D44BAC08C4C100107050206091E7190EC9776246332040204133B86ABF8B310384D61BFF5B1605536127089DF510E5E75B8B9B2D2856E031E807CE2E598703A09051656BDFBFEF4B03F4837146BE4FBFBF3F7FEF6F6FAFDF9E469030A48BDF9F6FAF9F8F6FBFCA835020740B6851D0006020406070A75D1ADCB5C150001090202042482AD57031A65D2FAFEB3430E02012FA0FBFAF7F6F190731F116FD6FCF9F8FAF2BA775008453B4B41483F3082D28B6D04080AC4ACD16F07090009319CF0C15911012887DDFFFFFFFFFFFFFFFFFFFFFFFFFFFFFFFFFFFFFFFFFFFFFFFFFFFFFFFFFFFEF0A1C9F8A43DA46F1300022B60137686757D7CD0ABC6F8E5B17C758AB2BEB8832539353F3E3A677087E383EB8482D59731103E7983A0D5A1B42F7A7D9CAEB5B08F7C7E3CB992D49E27A8A1BEB881ADD77E7A807F78DFE296757B7D73737D7BBAF6F6A98B793A4475ABE7D598F08AEE8BDDA1CBAA6F444487ACF2C4ABEBD18EB8E4C76A395B96D0937E77752B5E334C79B4F3CD806A64061B611D050544B1BA60443E62D2EFFAA8843D110436AAF7F1F6F0F5EE91673BADF6F9FBF6F8FCF6B53E0C09063B451C03286E9ACC21080D6CDAF5F0723D480B0D339BD6765B13032C8BDFFFFFFFFFFFFFFFFFFFFFFFFFFFFFFFFFFFFFFFFFFFFFFFFFFFFFFFFFFFFDFBFEFDF8FDA3B96E1E0804346D1A80038E959190757E929FBCCEBECC908F89834B573F4F707A8E8B96DF08D0A097978A6F3E4B898178703EBF3A5862A78A738A6B4F935B6316748A67839598915356B538DD8B6B8E8F8D71395E9994908341468F968D958F797E7A74A4AB36D60C8E1173265E424A5D6A9658C770748A5F43768B8C847D838EAADAA1999573887E81C0C68C7C738C6907228661191E5E886B4555445182CAB16C80813A3272B9D2A2E099E99BE89885CDF0EDE5CBC3F2ECC24B49323F313B34396B858BD5746D6EE3A9D5EAA87F77492F7EA7AB385F122588DEFDFFFFFFFFFFFFFFFFFFFFFFFFFFFFFFFFFFFFFFFFFFFFFFFFFFFFFFFFFFFDFAFEF7F5FCF6C888663F1007090B06040903060A6141051648706B85080001484E1B02053A5B2609198D0686250700085F50020C3A770729860F043F514A4F2F072B5C410004000372360A05040234772E7A0E8C1B030003020404060A0000020005060304584C50574271248A0B0204040B0A0506103B5E3B60366E080454581E043C5B270A21741B000D6C256B3F50570430760E060700206925305839080305206C287233623674445152386A217D1E8E17BD45ACC0A79C8C7F7C99976F60847786525A707E806D97979DAB9EE53EABBDCA8341653572257208080960D1FAFCFFFFFFFFFFFFFFFFFFFFFFFFFFFFFFFFFFFFFFFFFFFFFFFFFFFFFFFFFFFBF2F7F5FCFDEA9BBBEBB55E392B40071600040F0602090D0600124E670000030D030101030707060904020000010808030D0C0000020204070B040B0B03090E070206060A010901040605090E070001030905060B4A150A0C090403010001010608000D0905060200000204030B0005000114080403070C0400000305030E0200060F0A010712050B0F00060C03020205060505000400120102000404010103080600000303020000000308070404000303001D7B2F7037050618515E021656B04E534F190C5152426F0C080400048C2A733B545102040A000205112A4E95DFF4F8FFFFFFFFFFFFFFFFFFFFFFFFFFFFFFFFFFFFFFFFFFFFFFFFFFFFFFFFFFFEFBE9FDFEF9FCFCF9FBF9F5C57DDF44500A060308090703010102060B01070508010806222B14110907090A010A05361D0200040B0D0D0E070013040404050402183113050C00070B0B16090606040A1941103F33DD6948202C260D070408070B1904060606020004040A1700030D040607321A0804080B0206040A02050906040C202413050511073C10020904090C080404080B100808030404090504080504080407060E090101000000070D06040D0B000203050B04000405010303080B0107000100000103030306030E001004000C070006090609020A1C3F7AE7FAFBFEFFFFFFFFFFFFFFFFFFFFFFFFFFFFFFFFFFFFFFFFFFFFFFFFFFFFFFFFFFFFFDFDFEFAFCF5FBF8FBFBFBFAF7FDFBFEF944391B090A261F060D1B19010A282E0D071010798F3130462F0D500A2D22B76A10072345392D0F3B0C3C1A25242B251C5AA25E2E400D39384E461A2B2716254DCF64F45BDFFAE08ABA8C261B1F1D234D1A1E3620311D0C24271C360A195A03321CC27E44261C2219371D251B2D14272817658F5D2D182429ED651645191C1A272926443811463D17091B241310191C16142929250F0903090806031B260612050400071531110303090E06080E02020905020C06050B080417330D04000506020A0B0400090E01040435D4F9FEFFF8FAFFFFFFFFFFFFFFFFFFFFFFFFFFFFFFFFFFFFFFFFFFFFFFFFFFFFFFFFFFFFFDFAFEFEECFBF9FCFEFFFFFCF8F9FFFDEAFDFF653D56B2983E559B9F1739B6A3494F3E3889A39FECEEB760F253E562BA763F4A97FEF8CD6EDC77C187997E88A6A7A28859B9EB62D1F9ECF887ADA08EA845F9FEFEEBFBF7FCFCFDFDDCAB9CAAA7FFFE84BD85B59E56748E89CA5649E30FE45ECDFAF29D9BBB87C780CA81B5A296B0A2AC9A94AA98D560D0CC52D46DBF8BA1AE8AE19C71F7A3363F5B976B2F588D6D67AB57BE59291D272C382174A41A3D3C31033E0FC252221C1F3035381B1A1E031B1F21132017182246BC2132053C41172A17050401010500030C78E7FDFEF3F8FFFFFFFFFFFFFFFFFFFFFFFFFFFFFFFFFFFFFFFFFFFFFFFFFFFFFFFFFFFFFEFEFEFEFEFEFEFEFEFFFFFEFEFEFFFEFEFEFFF7F5F6FBFAF5F6E2A3A0EAFAF0F3F6EDF0F9FAF5F6D76CE6FEF6FDF7FBF8F5F6F9FDFAFCF7F7F7FCE3C097CCF5F5F2C891C3FAF7FCFEFEFEF9FBFAF5C285F7FAFEF8F7FCFEFEFEFEFDFAFAFAFAFFFEF9FBF9FBFAF6A08BEFF9EFF6F540E7F6FCFEFEFAFAFBF9FCF8FCF8FBFAF9FBFAFAF7F6F8F3F5F4DF59DCF6F4FAF9F8F7F1BF69F4DA4FE8F5B68BC7F2CA96B0F2DA94B66FD57CBDECF39796BD81EDEFE851E471C786A89CAFF2EEDCA996A450AE88AD8CAB9DA1B182C56FDF4EE8CA59956D1F0401000311408ADCFBFEFCFBFDFFFFFFFFFFFFFFFFFFFFFFFFFFFFFFFFFFFFFFFFFFFFFFFFFFFFFFFFFFFFFFFFFFFFFFFFFFFFFFFFFFFFFFFFFFFFFFFFFFFFFFFFFFFFFFFFFAEBDBFDFEECFAFEFDFCFEF8F4FBF2D5FBFAFEFFFFFFFFFFFFFFFDFBFFFEF9FEFFFCECE1F2F8F8FEF1DDEFFEFFFFFFFFFFFFFFFFFEE9D3FCF9FDFDFCF8FDFFFFFFFFFFFFFFFFFFFFFFFFFFFFFFFFF3EEFAFEF9EBFDD4F9FEFFFFFFFFFFFFFFFFFFFFFFFFFFFFFFFFFFF9F7FEF9EFFBF9D2F4FAFCFEFFFDFCFAF1E5F8F7D6FBFFF6E3F1FEF5EDF4FDF7EBF6EDF4E8F5FDFDEEEFF0E3FDFCF2E7FBD4F1E6F0E4E6FCFEFAEFEBF3DDF4F1F6F1F5EEEEF3EDF1D8FCE9E965C47AA4611C3F4A5274DEF1FBFEFDFDFEFCFCFFFFFFFFFFFFFFFFFFFFFFFFFFFFFFFFFFFFFFFFFFFFFFFFFFFFFFFFFFFFFFFFFFFFFFFFFFFFFFFFFFFFFFFFFFFFFFFFFFFFFFFFFFFFFFFEFCF7FFFFF6FCFEFDFBFCFDFCFCF9F9FEFAFEFFFFFFFFFFFFFFFDFDFFFEFAFBFEFFFCFBFDF8F6FCFDFAFDFFFFFFFFFFFFFFFFFFFEFBF6F8F9F8FEFEFCFDFFFFFFFFFFFFFFFFFFFFFFFFFFFFFFFFFFFCFBFCFEF9FCF9F1FBFFFFFFFFFFFFFFFFFFFFFFFFFFFFFFFFFFF9F6FDFDF8F7FEF9FCFDFEFEFFFEFEFDFDFCFBF8FAF7FDFFFCFDF6F6FDFCF9FAFDFFFFFDFDFEF9FBFEFAF8FCFFFDF3FDF3F6FDFCFDFBF2F3FEFCFCFCF8FBFCFAFDFDF8F9FDFEFCFAF7F2FEFADEF2E2EE9389D0CCDFE6FCFEFAF3F8FDFFFCFCFFFFFFFFFFFFFFFFFFFFFFFFFFFFFFFFFFFFFFFFFFFFFFFFFFFFFFFFFFFFFFFFFFFFFFFFFFFFFFFFFFFFFFFFFFFFFFFFFFFFFFFFFFFFFFFDFAFAFDFEFDFEFCFAFCFBFCFFFDFAF7F3FDFEFFFFFFFFFFFFFFFEFDFDFEFDFDF9F9FDFBFAFBFCFCFCFFFFFFFFFFFFFFFFFFFFFFF9F7F9EDFDFDFEF9FDFEFFFFFFFFFFFFFFFFFFFFFFFFFFFFFFFFF8F6FDFAF8FCF1FAF9FDFFFFFFFFFFFFFFFFFFFFFFFFFFFFFFFFFFFDFBFBFBFDFCFAF2FDFEFFFFFFFFFFFEF9F1FCF5FBF9FDF3FAFFFCF8F6F9FDFBF5F9F1FDF1FAFDFDF3F1F8F6FDFAFBF6F2FCFAF7F3F9FCF8F9FBF4FAFCFEFEF1FAF7F8FDFDFCFDFDFEF9F5F6FBFCF5F8DEE1FAF7FAFDF8FEFCFBFDFEFFFAFBFFFFFFFFFFFFFFFFFFFFFFFFFFFFFFFFFFFFFFFFFFFFFFFFFFFFFFFFFFFFFFFFFFFFFFFFFFFFFFFFFFFFFFFFFFFFFFFFFFFFFFFFFFFFFFFDFBFCFAF9FEFDFBFCFEF9F7FBFAFAF3F9FFFFFFFFFFFFFFFFFFFEFBFAFDFEFEFCFCFEFDF9FBFFFDFCFFFEFBFFFFFFFFFFFFFFFFFCF9F7F4FFFEFDF6FDFFFFFFFFFFFFFFFFFFFFFFFFFFFFFFFFFFFBFAFDFBFBFAF8FCFEFFFFFFFFFFFFFFFFFFFFFFFFFFFFFFFFFFFCFCFEFAF6FAFEFCF9FEFFFEFDFDFDFEFFFCF8FBF6FCFBFEF9FCF9FBFCF8F9FFFCF2FCF8FBF8FCFFFBF6F9FBFAFDFCFCFBFAFEF7F9F9FAFDFAF8FDF9FAFDFDFCF8FAFAFCFDFBFBFBF6F9FAFAFAFDFEFAF7F9FCFBF7F9FDFBFEFEFFFDFBFCFBFCFFFFFFFFFFFFFFFFFFFFFFFFFFFFFFFFFFFF>
EI Q}\hss}}
\pdfxform\publisherbox\edef\publisherform{\the\pdflastxform}\newcommand{\publisherdevice}{\pdfrefxform\publisherform}
\newcommand{\weberfigure}[2]{\begin{center}\vspace{23mm}\leavevmode\hbox{\publisherdevice}\par\vspace{20mm}\end{center}}
\expandafter\def\csname ednote1\endcsname{\edNote{1}{s001R}{Dubbio. piccolo segno dopo «ZANICHELLI» -- macchia / punto / virgola.}}
\expandafter\def\csname ednote2\endcsname{\edNote{2}{IX}{Stampa. teoria-stessa -> teoria stessa, se il trattino è tipografico, unisce due parole intere.}}
\expandafter\def\csname ednote3\endcsname{\edNote{3}{IX}{Dubbio. piccolo segno fra «delle» e «lezioni» in «i corsi delle lezioni» -- macchia / punto stampato.}}
\expandafter\def\csname ednote4\endcsname{\edNote{4}{XI}{Dubbio. piccolo segno dopo «Capitolo» -- macchia / virgola / punto.}}
\expandafter\def\csname ednote5\endcsname{\edNote{5}{XI}{Dubbio. piccolo segno in alto a destra di «5.» -- macchia / «°».}}
\expandafter\def\csname ednote6\endcsname{\edNote{6}{XI}{Dubbio. punto dopo «56» -- punto stampato / macchia.}}
\expandafter\def\csname ednote7\endcsname{\edNote{7}{XIV}{Stampa. Disgressione -> Digressione, probabile s soprannumeraria; a p. XIII è stampato «Digressione».}}
\expandafter\def\csname ednote8\endcsname{\edNote{8}{3}{Stampa. F -> f, la superficie di partenza è indicata con f nelle righe precedenti.}}
\expandafter\def\csname ednote9\endcsname{\edNote{9}{5}{Dubbio. segno dopo «Op.» nella nota -- «Op., cit.,» / «Op. cit.,» con puntinatura.}}
\expandafter\def\csname ednote10\endcsname{\edNote{10}{11}{Stampa. C -> \textbackslash{}\allowbreak bar\{C\}, la curva doppia di \textbackslash{}\allowbreak Phi è indicata con \textbackslash{}\allowbreak bar\{C\} nelle righe contigue.}}
\expandafter\def\csname ednote11\endcsname{\edNote{11}{12}{Stampa. C -> \textbackslash{}\allowbreak bar\{C\}, qui si parla della curva doppia di \textbackslash{}\allowbreak Phi, detta \textbackslash{}\allowbreak bar\{C\} nelle due alternative precedenti.}}
\expandafter\def\csname ednote12\endcsname{\edNote{12}{12}{Stampa. \textbackslash{}\allowbreak bar\{K\}\_\allowbreak 1 in «la curva \textbackslash{}\allowbreak bar\{K\}\_\allowbreak 1 proiettata» e in «di \textbackslash{}\allowbreak bar\{K\}\_\allowbreak 1» -> K\_\allowbreak 1, la curva su F è denotata K\_\allowbreak 1 nel periodo precedente.}}
\expandafter\def\csname ednote13\endcsname{\edNote{13}{12}{Dubbio. segno dopo K\_\allowbreak 1 nell’ultima frase -- nessuna virgola / virgola.}}
\expandafter\def\csname ednote15\endcsname{\edNote{14}{19}{Stampa. C -> |C|, the first symbol denotes the system, otherwise written with vertical bars.}}
\expandafter\def\csname ednote16\endcsname{\edNote{15}{22}{Stampa. mn -> mn\textasciicircum{}2, the square on the cone order appears to be missing.}}
\expandafter\def\csname ednote17\endcsname{\edNote{16}{24}{Dubbio. Accent on the conjunction before "le loro relazioni" -- e or è.}}
\expandafter\def\csname ednote18\endcsname{\edNote{17}{24}{Stampa. Mathematik -> Mathematics, apparent misprint in the English journal title.}}
\expandafter\def\csname ednote19\endcsname{\edNote{18}{27}{Stampa. contenuti -> contenute, agreement with "esse" (the curves).}}
\expandafter\def\csname ednote20\endcsname{\edNote{19}{34}{Dubbio. punctuation after the first "specie" -- comma or scan speck.}}
\expandafter\def\csname ednote21\endcsname{\edNote{20}{39}{Stampa. \textbackslash{}\allowbreak omega\_\allowbreak 1 -> \textbackslash{}\allowbreak omega\_\allowbreak 2, the immediately preceding clause names the transformed curve \textbackslash{}\allowbreak omega\_\allowbreak 2.}}
\expandafter\def\csname ednote22\endcsname{\edNote{21}{39}{Dubbio. mark between "con" and "le" -- printed hyphen or scanning blemish.}}
\expandafter\def\csname ednote23\endcsname{\edNote{22}{41}{Dubbio. small mark after n in 3n-3 -- scanning speck or prime (3n'-3).}}
\expandafter\def\csname ednote24\endcsname{\edNote{23}{41}{Dubbio. punctuation after "triplo" -- semicolon or comma.}}
\expandafter\def\csname ednote25\endcsname{\edNote{24}{43}{Dubbio. The punctuation after "irriducibile" -- semicolon / comma.}}
\expandafter\def\csname ednote26\endcsname{\edNote{25}{44}{Dubbio. The final dot -- printed period / scan speck.}}
\expandafter\def\csname ednote27\endcsname{\edNote{26}{44}{Stampa. (\textbackslash{}\allowbreak partial\textbackslash{}\allowbreak varphi\textbackslash{}\allowbreak varphi\_\allowbreak 3) -> \textbackslash{}\allowbreak partial(\textbackslash{}\allowbreak varphi\textbackslash{}\allowbreak varphi\_\allowbreak 3), the opening parenthesis is printed before the derivative symbol.}}
\expandafter\def\csname ednote28\endcsname{\edNote{27}{45}{Stampa. -\textbackslash{}\allowbreak partial f/\textbackslash{}\allowbreak partial x\_\allowbreak 3 -> \textbackslash{}\allowbreak partial f/\textbackslash{}\allowbreak partial x\_\allowbreak 3, if the stroke is type, the preceding derivative column has no minus.}}
\expandafter\def\csname ednote29\endcsname{\edNote{28}{49}{Stampa. punt -> punti, the final i is absent before "base".}}
\expandafter\def\csname ednote30\endcsname{\edNote{29}{50}{Stampa. dalle -> delle, the phrase identifies the system of plane sections.}}
\expandafter\def\csname ednote31\endcsname{\edNote{30}{55}{Dubbio. punctuation after "relazione" -- period / colon / scan speck.}}
\expandafter\def\csname ednote32\endcsname{\edNote{31}{57}{Dubbio. accent on the preposition before "tutte" -- à / a with a scan speck.}}
\expandafter\def\csname ednote33\endcsname{\edNote{32}{58}{Stampa. la molteplicità -> le molteplicità, the verb "si lasciano" is plural.}}
\expandafter\def\csname ednote34\endcsname{\edNote{33}{58}{Dubbio. accent on the preposition after "luogo" -- à / a with a scan speck.}}
\expandafter\def\csname ednote35\endcsname{\edNote{34}{61}{Dubbio. relation sign -- \textbackslash{}\allowbreak doteq / =; the small mark above the bars may be a scan speck.}}
\expandafter\def\csname ednote36\endcsname{\edNote{35}{61}{Dubbio. relation sign -- \textbackslash{}\allowbreak doteq / =; the small mark above the bars may be a scan speck.}}
\expandafter\def\csname ednote37\endcsname{\edNote{36}{68}{Dubbio. tiny raised mark after the first \textbackslash{}\allowbreak (\textbackslash{}\allowbreak theta\textbackslash{}\allowbreak ) -- \textbackslash{}\allowbreak (\textbackslash{}\allowbreak theta\textbackslash{}\allowbreak ) with a scan speck or \textbackslash{}\allowbreak (\textbackslash{}\allowbreak theta'\textbackslash{}\allowbreak ).}}
\expandafter\def\csname ednote38\endcsname{\edNote{37}{68}{Dubbio. raised mark beside the first \textbackslash{}\allowbreak (\textbackslash{}\allowbreak varphi\_\allowbreak \{n-3\}\textbackslash{}\allowbreak ) in this paragraph -- \textbackslash{}\allowbreak (\textbackslash{}\allowbreak varphi\_\allowbreak \{n-3\}\textbackslash{}\allowbreak ) with a scan speck or \textbackslash{}\allowbreak (\textbackslash{}\allowbreak varphi'\_\allowbreak \{n-3\}\textbackslash{}\allowbreak ).}}
\expandafter\def\csname ednote39\endcsname{\edNote{38}{77}{Dubbio. punctuation after "affermazioni" -- comma or no punctuation.}}
\expandafter\def\csname ednote40\endcsname{\edNote{39}{77}{Stampa. C\_\allowbreak 4 -> C\_\allowbreak 7, these residual curves were denoted C\_\allowbreak 7 earlier on the page.}}
\expandafter\def\csname ednote41\endcsname{\edNote{40}{77}{Stampa. C\_\allowbreak 4 -> C\_\allowbreak 7, repeated subscript discrepancy for the same curves.}}
\expandafter\def\csname ednote42\endcsname{\edNote{41}{79}{Stampa. O -> A, the point under discussion is called A in this sentence.}}
\expandafter\def\csname ednote43\endcsname{\edNote{42}{79}{Stampa. i volte A per -> i volte per A, the words are transposed in print.}}
\expandafter\def\csname ednote45\endcsname{\edNote{43}{84}{Stampa. dal 4º ordine -> del 4º ordine, the preposition appears to be a typesetting error.}}
\expandafter\def\csname ednote46\endcsname{\edNote{44}{84}{Stampa. di questa -> di questo, the pronoun appears to refer to the point O.}}
\expandafter\def\csname ednote47\endcsname{\edNote{45}{85}{Stampa. è in una cubica -> e in una cubica, a conjunction is required here.}}
\expandafter\def\csname ednote48\endcsname{\edNote{46}{87}{Stampa. grado -\textbackslash{}\allowbreak nu -> grado \textbackslash{}\allowbreak nu, the signed degree is denoted by \textbackslash{}\allowbreak nu on the preceding page and in the concluding formula.}}
\expandafter\def\csname ednote49\endcsname{\edNote{47}{87}{Dubbio. Small stroke after (\textbackslash{}\allowbreak pi-1) -- prime or scan/printing speck.}}
\expandafter\def\csname ednote50\endcsname{\edNote{48}{87}{Stampa. -\textbackslash{}\allowbreak nu=-\textbackslash{}\allowbreak nu\_\allowbreak 0-\textbackslash{}\allowbreak nu\_\allowbreak 1+2=-1 -> \textbackslash{}\allowbreak nu=-\textbackslash{}\allowbreak nu\_\allowbreak 0-\textbackslash{}\allowbreak nu\_\allowbreak 1+2=-1, the initial sign conflicts with -\textbackslash{}\allowbreak nu=1 above and \textbackslash{}\allowbreak nu=-1 below.}}
\expandafter\def\csname ednote51\endcsname{\edNote{49}{87}{Dubbio. Tiny mark after -1 -- full stop or scan speck.}}
\expandafter\def\csname ednote52\endcsname{\edNote{50}{89}{Dubbio. Mark after "rapporto" -- scan speck (best reading) or full stop.}}
\expandafter\def\csname ednote53\endcsname{\edNote{51}{91}{Stampa. \textbackslash{}\allowbreak varphi\_\allowbreak \{n-4\} -> \textbackslash{}\allowbreak varphi\_\allowbreak \{n-3\}, the section of degree n-4 together with a line has degree n-3.}}
\expandafter\def\csname ednote54\endcsname{\edNote{52}{92}{Stampa. del punto -> dal punto, the projection is taken from O.}}
\expandafter\def\csname ednote55\endcsname{\edNote{53}{94}{Dubbio. punctuation after “Castelnuovo” -- no punctuation or a faint period.}}
\expandafter\def\csname ednote56\endcsname{\edNote{54}{95}{Stampa. supeficie -> superficie, the letter r is missing.}}
\expandafter\def\csname ednote57\endcsname{\edNote{55}{98}{Stampa. di essa -> di esso, the stated antecedent is the masculine “fascio”.}}
\expandafter\def\csname ednote58\endcsname{\edNote{56}{99}{Stampa. e suoi poli -> e i suoi poli, the article is absent in print.}}
\expandafter\def\csname ednote59\endcsname{\edNote{57}{101}{Stampa. contenerle -> contenerli, the antecedent is “i punti doppi”.}}
\expandafter\def\csname ednote60\endcsname{\edNote{58}{102}{Stampa. contenerle -> contenerli, the antecedent is the masculine plural “punti doppi”.}}
\expandafter\def\csname ednote61\endcsname{\edNote{59}{102}{Stampa. contenerle -> contenerla, the antecedent is the singular curve C.}}
\expandafter\def\csname ednote62\endcsname{\edNote{60}{103}{Stampa. d’ordine n -> d’ordine m, the order of the curve C was denoted by m.}}
\expandafter\def\csname ednote63\endcsname{\edNote{61}{103}{Stampa. alle totalità -> alla totalità, the passage apparently refers to the single totality of these surfaces.}}
\expandafter\def\csname ednote64\endcsname{\edNote{62}{109}{Stampa. sezioni piane di s -> sezioni piane di F, the plane sections under discussion are those of the surface F.}}
\expandafter\def\csname ednote65\endcsname{\edNote{63}{109}{Stampa. -\textbackslash{}\allowbreak frac\{s(s-1)\}\{2\}n -> +\textbackslash{}\allowbreak frac\{s(s-1)\}\{2\}n, the sign disagrees with the preceding expression and the stated equality.}}
\expandafter\def\csname ednote66\endcsname{\edNote{64}{111}{Stampa. e dalla serie -> e della serie, the text is describing a sum of two series.}}
\expandafter\def\csname ednote67\endcsname{\edNote{65}{111}{Dubbio. mark between \textbackslash{}\allowbreak delta(L+C) and = -- full stop or stray ink.}}
\expandafter\def\csname ednote68\endcsname{\edNote{66}{115}{Stampa. della superficie -> delle superficie, the following participle "costrette" is plural.}}
\expandafter\def\csname ednote69\endcsname{\edNote{67}{116}{Dubbio. punctuation after \textbackslash{}\allowbreak delta\_\allowbreak 1\textbackslash{}\allowbreak geq0 -- semicolon or comma with a stray dot.}}
\expandafter\def\csname ednote70\endcsname{\edNote{68}{116}{Stampa. p\_\allowbreak a+\textbackslash{}\allowbreak pi\_\allowbreak \{h-3\}+\textbackslash{}\allowbreak delta\_\allowbreak \{h-2\}+\textbackslash{}\allowbreak delta\_\allowbreak \{h-1\} -> p\_\allowbreak a+\textbackslash{}\allowbreak pi\_\allowbreak \{h-3\}-1+\textbackslash{}\allowbreak delta\_\allowbreak \{h-2\}+\textbackslash{}\allowbreak delta\_\allowbreak \{h-1\}, the constant -1 in the preceding dimension formula is missing here.}}
\expandafter\def\csname ednote71\endcsname{\edNote{69}{118}{Stampa. f -> F, the surface is otherwise denoted by capital F.}}
\expandafter\def\csname ednote72\endcsname{\edNote{70}{119}{Stampa. se zione -> sezione, the word is split at a printed line end without a hyphen.}}
\expandafter\def\csname ednote73\endcsname{\edNote{71}{119}{Stampa. f -> F, the surface is otherwise denoted by capital F.}}
\expandafter\def\csname ednote74\endcsname{\edNote{72}{121}{Dubbio. the small dot following the footnote label -- printed full stop (read), scan speck.}}
\expandafter\def\csname ednote75\endcsname{\edNote{73}{124}{Dubbio. punctuation after 3 -- semicolon (read), comma with a speck above it.}}
\expandafter\def\csname ednote76\endcsname{\edNote{74}{124}{Stampa. \textbackslash{}\allowbreak bar\{C\}=|C-\textbackslash{}\allowbreak theta| -> |\textbackslash{}\allowbreak bar\{C\}|=|C-\textbackslash{}\allowbreak theta|, the left-hand system delimiters are absent in the print.}}
\expandafter\def\csname ednote77\endcsname{\edNote{75}{124}{Dubbio. small marks after nu -- scan speckles with no punctuation (read), comma or full stop.}}
\expandafter\def\csname ednote78\endcsname{\edNote{76}{124}{Dubbio. faint stroke over the left-hand pi -- overbar (read), stray scan mark.}}
\expandafter\def\csname ednote79\endcsname{\edNote{77}{127}{Stampa. addittivo -> additivo, the printed word contains a doubled t.}}
\expandafter\def\csname ednote80\endcsname{\edNote{78}{130}{Stampa. delle serie completa -> della serie completa, the printed article does not agree with the singular noun and adjective.}}
\expandafter\def\csname ednote81\endcsname{\edNote{79}{130}{Dubbio. the bracket-like mark between “almeno” and “un” -- a non-text mark or a printed right square bracket.}}
\expandafter\def\csname ednote82\endcsname{\edNote{80}{133}{Stampa. RiemannRoch -> Riemann-Roch, the internal hyphen is absent in this occurrence.}}
\expandafter\def\csname ednote83\endcsname{\edNote{81}{136}{Dubbio. a faint high mark after the first p -- a scan speck or a printed prime.}}
\expandafter\def\csname ednote84\endcsname{\edNote{82}{141}{Stampa. opening « only -> closing » apparently omitted, the quotation is left open in print.}}
\expandafter\def\csname ednote85\endcsname{\edNote{83}{144}{Dubbio. the mark after "grande" -- printed full stop (retained) or scan speck.}}
\expandafter\def\csname ednote86\endcsname{\edNote{84}{144}{Dubbio. the faint mark after the final 1 -- scan speck (not transcribed) or full stop.}}
\expandafter\def\csname ednote87\endcsname{\edNote{85}{145}{Stampa. non sia abbia -> non si abbia, an extra a is printed in sia.}}
\expandafter\def\csname ednote88\endcsname{\edNote{86}{146}{Stampa. e quindi l’indice di specialità della serie -> e quindi la serie, the printed subject repeats the index that the predicate assigns.}}
\expandafter\def\csname ednote89\endcsname{\edNote{87}{148}{Dubbio. "A tal" -- "A tal" or "A. tal"; a small intervening mark may be a scan speck.}}
\expandafter\def\csname ednote90\endcsname{\edNote{88}{149}{Stampa. r -> \textbackslash{}\allowbreak nu, the surrounding text denotes the degree by \textbackslash{}\allowbreak nu.}}
\expandafter\def\csname ednote91\endcsname{\edNote{89}{150}{Dubbio. "distacco di" -- "distacco di" or "distacco. di"; the intervening point may be a scan speck.}}
\expandafter\def\csname ednote92\endcsname{\edNote{90}{152}{Dubbio. "che occorre" -- "che occorre" or "che. occorre"; the small intervening point may be a scan speck.}}
\expandafter\def\csname ednote93\endcsname{\edNote{91}{154}{Stampa. "sia A\_\allowbreak 1, un numero" -> "sia A\_\allowbreak 1, ha un numero", the clause appears to omit its verb.}}
\expandafter\def\csname ednote94\endcsname{\edNote{92}{155}{Stampa. (CL') -> (DL'), the preceding paragraph places this residual series on D.}}
\expandafter\def\csname ednote95\endcsname{\edNote{93}{157}{Stampa. "rappresentato del fascio" -> "rappresentato dal fascio", the preposition appears to be a typographical error.}}
\expandafter\def\csname ednote96\endcsname{\edNote{94}{158}{Stampa. Le possibilità -> La possibilità, the following verb "costringe" is singular.}}
\expandafter\def\csname ednote97\endcsname{\edNote{95}{159}{Stampa. i sistema -> il sistema, the singular article is missing its final l.}}
\expandafter\def\csname ednote98\endcsname{\edNote{96}{159}{Dubbio. closing mark in the exponent of the second p -- closing parenthesis, comma-like damaged mark.}}
\expandafter\def\csname ednote99\endcsname{\edNote{97}{160}{Stampa. n<\textbackslash{}\allowbreak pi-1 -> n\textbackslash{}\allowbreak leq\textbackslash{}\allowbreak pi-1, the preceding inequality gives only a non-strict bound when p=1.}}
\expandafter\def\csname ednote100\endcsname{\edNote{98}{161}{Dubbio. apparent dot after "dividere" -- scan speck (not transcribed), printed full stop.}}
\expandafter\def\csname ednote101\endcsname{\edNote{99}{164}{Dubbio. dots after "connesse" and "connessione" -- printed full stops, scan speckles.}}
\expandafter\def\csname ednote102\endcsname{\edNote{100}{164}{Stampa. \textbackslash{}\allowbreak delta\_\allowbreak 1\textbackslash{}\allowbreak delta\_\allowbreak 2 -> \textbackslash{}\allowbreak delta\_\allowbreak 1,\textbackslash{}\allowbreak delta\_\allowbreak 2, the three previously distinguished degrees are being enumerated.}}
\expandafter\def\csname ednote103\endcsname{\edNote{101}{165}{Dubbio. subscript on the second F in "F\_\allowbreak 1+F\_\allowbreak 2" -- 2, a dash-like damaged character.}}
\expandafter\def\csname ednote104\endcsname{\edNote{102}{165}{Stampa. m\_\allowbreak 1+m\_\allowbreak 2+ +4 -> m\_\allowbreak 1+m\_\allowbreak 2-4, the same adjoint order is printed with -4 elsewhere on this page.}}
\expandafter\def\csname ednote105\endcsname{\edNote{103}{168}{Dubbio. punctuation after "avremo" -- colon or an isolated scan speck.}}
\expandafter\def\csname ednote106\endcsname{\edNote{104}{168}{Dubbio. the small mark above \textbackslash{}\allowbreak pi -- scan speck or prime; read as a speck.}}
\expandafter\def\csname ednote107\endcsname{\edNote{105}{169}{Stampa. "(r-1)\textasciicircum{}2, volte" -> "(r-1)\textasciicircum{}2 volte", the comma separates the multiplier from "volte".}}
\expandafter\def\csname ednote108\endcsname{\edNote{106}{170}{Stampa. "serie segata" -> "serie segate", agreement with the plural "delle serie".}}
\expandafter\def\csname ednote109\endcsname{\edNote{107}{170}{Stampa. "(C\_\allowbreak 2L)" -> "(C\_\allowbreak 2L\_\allowbreak 1)", the preceding count concerns the residual curve L\_\allowbreak 1.}}
\expandafter\def\csname ednote110\endcsname{\edNote{108}{172}{Stampa. "dato" -> "data", agreement with the subject "L’equivalenza".}}
\expandafter\def\csname ednote111\endcsname{\edNote{109}{173}{Stampa. "(\textbackslash{}\allowbreak pi-2)" -> "(2\textbackslash{}\allowbreak pi-2)", the preceding sentence gives the multiplicity as 2\textbackslash{}\allowbreak pi-2.}}
\expandafter\def\csname ednote112\endcsname{\edNote{110}{173}{Stampa. "P. Bonnesen" -> "T. Bonnesen", the main text gives the initial T.; "Académie Royal" -> "Académie Royale", adjective agreement in French.}}
\expandafter\def\csname ednote113\endcsname{\edNote{111}{175}{Stampa. "della formula (2)" -> "dalla formula (3)", the preposition is anomalous and the preceding formula for d is labelled (3).}}
\expandafter\def\csname ednote114\endcsname{\edNote{112}{175}{Stampa. "2n\textasciicircum{}2-7n+15+2\textbackslash{}\allowbreak pi(n-7)" -> "2n\textasciicircum{}2-7n+15+2\textbackslash{}\allowbreak pi(n-7)-p\textasciicircum{}\{(1)\}", the subsequent equation (6) retains the -p\textasciicircum{}\{(1)\} term.}}
\expandafter\def\csname ednote115\endcsname{\edNote{113}{175}{Stampa. "n+14" -> apparently "-7n+15-\textbackslash{}\allowbreak tau", comparison with the preceding expressions for (C\_\allowbreak jD) and a gives a different linear and constant part.}}
\expandafter\def\csname ednote116\endcsname{\edNote{114}{176}{Stampa. "\textbackslash{}\allowbreak delta,I,\textbackslash{}\allowbreak varrho,\textbackslash{}\allowbreak Pi,\textbackslash{}\allowbreak tau,t" -> "\textbackslash{}\allowbreak delta,I,d,\textbackslash{}\allowbreak varrho,\textbackslash{}\allowbreak Pi,\textbackslash{}\allowbreak tau,t", only six entries are printed after "i sette caratteri".}}
\expandafter\def\csname ednote117\endcsname{\edNote{115}{176}{Stampa. "(2)" -> "(3)", the expression for d on page 174 is labelled (3).}}
\expandafter\def\csname ednote118\endcsname{\edNote{116}{178}{Dubbio. the small mark after F in row 2, column 4 -- a prime; an ink blemish.}}
\expandafter\def\csname ednote119\endcsname{\edNote{117}{182}{Stampa. 3(n+6\textbackslash{}\allowbreak pi-k) -> 3(n+6\textbackslash{}\allowbreak pi)-k, comparison with the fourth formula A).}}
\expandafter\def\csname ednote120\endcsname{\edNote{118}{182}{Stampa. G. H. Zeuthen -> H. G. Zeuthen, the initials appear in the opposite order on page 180.}}
\expandafter\def\csname ednote121\endcsname{\edNote{119}{183}{Dubbio. the punctuation at the end of the footnote -- ecc..; ecc.}}
\expandafter\def\csname ednote122\endcsname{\edNote{120}{185}{Dubbio. the final vowel of delle before biaggiunte -- delle; dello.}}
\expandafter\def\csname ednote123\endcsname{\edNote{121}{185}{Stampa. D (twice) -> d, the branch curve is denoted d later in this paragraph, while D denotes the coincidence curve.}}
\expandafter\def\csname ednote124\endcsname{\edNote{122}{186}{Dubbio. the page number -- 184 is visible in the overlapped area; the underlying page number is not securely readable.}}
\expandafter\def\csname ednote125\endcsname{\edNote{123}{186}{Dubbio. covered beginnings, continuations after cor- and dirama-, and the first footnote mark -- no reliable reading of the covered portions is possible; only visible fragments are retained, and the note is attached here provisionally.}}
\expandafter\def\csname ednote126\endcsname{\edNote{124}{186}{Dubbio. the clipped suffix before si può -- mente; uente; the preceding text and the beginning of the line containing numerici are covered.}}
\expandafter\def\csname ednote127\endcsname{\edNote{125}{186}{Dubbio. clipped word beginnings -- po (possibly uopo); urve (possibly curve); esti (possibly questi); other covered text cannot be read.}}
\expandafter\def\csname ednote128\endcsname{\edNote{126}{187}{Stampa. -3 -> -3\textbackslash{}\allowbreak nu, substitution of the preceding expression for s leaves a factor \textbackslash{}\allowbreak nu.}}
\expandafter\def\csname ednote129\endcsname{\edNote{127}{187}{Stampa. i valore -> il valore, the final l of the article is absent.}}
\expandafter\def\csname ednote130\endcsname{\edNote{128}{191}{Stampa. \textbackslash{}\allowbreak S II, 8 -> Cap. II, \textbackslash{}\allowbreak S 8, the reference at the top of this page has that form.}}
\expandafter\def\csname ednote131\endcsname{\edNote{129}{191}{Stampa. 2 : 24 -> 2 \textbackslash{}\allowbreak cdot 24, the displayed result is 48.}}
\expandafter\def\csname ednote132\endcsname{\edNote{130}{194}{Dubbio. mark above the first C -- unmarked C with a speck, C', or \textbackslash{}\allowbreak bar\{C\}.}}
\expandafter\def\csname ednote133\endcsname{\edNote{131}{194}{Stampa. m/2 -> m, the preceding text specifies 2m simple branch points.}}
\expandafter\def\csname ednote134\endcsname{\edNote{132}{199}{Stampa. cu sui -> su cui, apparent transposition of the initial consonants.}}
\expandafter\def\csname ednote135\endcsname{\edNote{133}{204}{Stampa. 1932.; -> 1932;, a period is printed immediately before the semicolon.}}
\expandafter\def\csname ednote136\endcsname{\edNote{134}{205}{Dubbio. degree symbol \textbackslash{}\allowbreak nu -- Greek nu; italic Latin v.}}
\expandafter\def\csname ednote137\endcsname{\edNote{135}{206}{Stampa. Toveremo -> Troveremo, the r after T is missing in print.}}
\expandafter\def\csname ednote138\endcsname{\edNote{136}{208}{Stampa. m = n - 2d - 3k -> m = n\textasciicircum{}2 - 2d - 3k, apparent omitted square exponent.}}
\expandafter\def\csname ednote139\endcsname{\edNote{137}{212}{Stampa. "sulla su" -> "sulla"; "dalle curva" -> "dalle curve", duplicated preposition and number disagreement.}}
\expandafter\def\csname ednote140\endcsname{\edNote{138}{212}{Dubbio. small mark after \textbackslash{}\allowbreak pi -- scan speck (best reading) or full stop.}}
\expandafter\def\csname ednote141\endcsname{\edNote{139}{212}{Dubbio. small mark after 3 -- scan speck (best reading) or full stop.}}
\expandafter\def\csname ednote142\endcsname{\edNote{140}{212}{Stampa. con curva una L -> con una curva L, transposed words.}}
\expandafter\def\csname ednote143\endcsname{\edNote{141}{214}{Stampa. \textbackslash{}\allowbreak omega = 0 -> \textbackslash{}\allowbreak omega' = 0, the preceding formula uses \textbackslash{}\allowbreak omega'.}}
\expandafter\def\csname ednote144\endcsname{\edNote{142}{215}{Dubbio. tiny mark after \textbackslash{}\allowbreak theta -- scan speck (best reading) or full stop.}}
\expandafter\def\csname ednote145\endcsname{\edNote{143}{217}{Dubbio. letter forms -- \textbackslash{}\allowbreak varrho or \textbackslash{}\allowbreak rho for genus, and \textbackslash{}\allowbreak nu or Latin v for degree.}}
\expandafter\def\csname ednote146\endcsname{\edNote{144}{217}{Dubbio. letter forms -- \textbackslash{}\allowbreak varrho or \textbackslash{}\allowbreak rho, and \textbackslash{}\allowbreak nu or Latin v.}}
\expandafter\def\csname ednote147\endcsname{\edNote{145}{220}{Stampa. espressa, da -> espressa da, intrusive comma between the participle and its complement.}}
\expandafter\def\csname ednote148\endcsname{\edNote{146}{222}{Stampa. curve canonica -> curva canonica, the printed noun and adjective disagree in number.}}
\expandafter\def\csname ednote149\endcsname{\edNote{147}{226}{Stampa. un’altro -> un altro, apostrophe before a masculine form.}}
\expandafter\def\csname ednote150\endcsname{\edNote{148}{229}{Dubbio. The subscript on the first p is indistinct -- a; i.}}
\expandafter\def\csname ednote151\endcsname{\edNote{149}{230}{Stampa. definite -> definita, apparently modifying the singular "involuzione".}}
\expandafter\def\csname ednote152\endcsname{\edNote{150}{230}{Stampa. D\_\allowbreak 3 -> D\_\allowbreak 9, the tricanonical curve is identified above with the branch curve D\_\allowbreak 9.}}
\expandafter\def\csname ednote153\endcsname{\edNote{151}{235}{Stampa. Ipotesi I) -> Ipotesi I, the closing parenthesis has no matching opening parenthesis.}}
\expandafter\def\csname ednote154\endcsname{\edNote{152}{237}{Dubbio. The short mark after "tacnodali" may be a hyphen or an artefact -- tacnodali-p; tacnodali p.}}
\expandafter\def\csname ednote155\endcsname{\edNote{153}{240}{Stampa. "riduzione." -> "riduzione,", the following lower-case subordinate clause continues the sentence.}}
\expandafter\def\csname ednote156\endcsname{\edNote{154}{242}{Dubbio. short mark between "un" and "punto" -- printed hyphen or ink/scanning artefact.}}
\expandafter\def\csname ednote157\endcsname{\edNote{155}{244}{Dubbio. raised mark after "intorni" before "delle" -- scan/ink speck (reading adopted) or printed asterisk.}}
\expandafter\def\csname ednote158\endcsname{\edNote{156}{246}{Dubbio. Mark after "particolari" -- a comma or a non-text blemish.}}
\expandafter\def\csname ednote159\endcsname{\edNote{157}{248}{Stampa. sisema -> sistema, the printed word omits "t".}}
\expandafter\def\csname ednote160\endcsname{\edNote{158}{251}{Stampa. siano -> siamo, the printed verb is inconsistent with the first-person plural construction "ci ... valsi".}}
\expandafter\def\csname ednote161\endcsname{\edNote{159}{253}{Dubbio. Mark between "il" and "punto" -- a printed hyphen or a non-text stroke.}}
\expandafter\def\csname ednote162\endcsname{\edNote{160}{256}{Stampa. auf der Flaechen -> auf den Flaechen, apparent German agreement error.}}
\expandafter\def\csname ednote163\endcsname{\edNote{161}{257}{Stampa. algebraics Varieties. -> algebraic Varieties »., apparent extra s and missing closing guillemet.}}
\expandafter\def\csname ednote164\endcsname{\edNote{162}{258}{Dubbio. punctuation after p\textasciicircum{}\{(1)\}=1 -- semicolon or comma.}}
\expandafter\def\csname ednote165\endcsname{\edNote{163}{259}{Dubbio. character after d in d’ramazione -- apostrophe or a defective i.}}
\expandafter\def\csname ednote166\endcsname{\edNote{164}{260}{Stampa. OP and P -> OB and B, the point is called B in the preceding description of type II.}}
\expandafter\def\csname ednote167\endcsname{\edNote{165}{262}{Stampa. +s-1 -> +s+1, the same relation is printed with +1 on page 263.}}
\expandafter\def\csname ednote168\endcsname{\edNote{166}{263}{Stampa. P=2(p-1) -> P=2(p-1)+1, an apparent omitted +1 when compared with the plurigenus formulas.}}
\expandafter\def\csname ednote169\endcsname{\edNote{167}{263}{Stampa. p-1=n-3 -> p-1=n-3-s, inconsistent with p=n-2-s printed above.}}
\expandafter\def\csname ednote170\endcsname{\edNote{168}{263}{Dubbio. stroke between punto and (2n-s-4) -- printed hyphen or scan mark.}}
\expandafter\def\csname ednote171\endcsname{\edNote{169}{263}{Stampa. 2s-3 -> 2(s-3), the contribution is written as 2(s-3) in the following total.}}
\expandafter\def\csname ednote172\endcsname{\edNote{170}{265}{Dubbio. mark before superficie in the second quoted title -- opening square bracket, opening parenthesis, or scan blemish.}}
\expandafter\def\csname ednote173\endcsname{\edNote{171}{267}{Dubbio. the low mark between "e" and "mercè" -- comma (transcribed) or scan speck.}}
\expandafter\def\csname ednote174\endcsname{\edNote{172}{274}{Stampa. C\_\allowbreak 7(O\textasciicircum{}5O\_\allowbreak 1\textasciicircum{}2) -> C'\_\allowbreak 7(O\textasciicircum{}5O\_\allowbreak 1\textasciicircum{}2), the preceding plane-image notation uses a prime.}}
\expandafter\def\csname ednote175\endcsname{\edNote{173}{274}{Dubbio. the dot after "(canonica)" -- printed full stop (transcribed) or scan speck.}}
\expandafter\def\csname ednote176\endcsname{\edNote{174}{275}{Stampa. fuori di esso -> fuori di essa, apparent agreement error with "curva doppia".}}
\expandafter\def\csname ednote177\endcsname{\edNote{175}{277}{Stampa. dal -> del, the preposition appears erroneous after "fa parte".}}
\expandafter\def\csname ednote178\endcsname{\edNote{176}{282}{Dubbio. An oblique mark follows B in "17 punti A e B" -- a scan speck or a prime on B.}}
\expandafter\def\csname ednote179\endcsname{\edNote{177}{283}{Dubbio. The dot following the running head -- a printed full stop or a scan speck.}}
\expandafter\def\csname ednote180\endcsname{\edNote{178}{286}{Stampa. varietà cubiche S\_\allowbreak 4 -> varietà cubiche di S\_\allowbreak 4, a preposition before the ambient-space symbol appears to be omitted.}}
\expandafter\def\csname ednote181\endcsname{\edNote{179}{287}{Stampa. delle -> dalle, the surfaces appear to be the agents of "segato".}}
\expandafter\def\csname ednote182\endcsname{\edNote{180}{290}{Stampa. F after "contenuto in" -> |F|, apparently missing system bars in a sentence about subsystems.}}
\expandafter\def\csname ednote183\endcsname{\edNote{181}{291}{Dubbio. the superscript on both occurrences of F -- \textbackslash{}\allowbreak alpha, a, or x.}}
\expandafter\def\csname ednote184\endcsname{\edNote{182}{292}{Dubbio. the superscript on the two residual-surface symbols F -- \textbackslash{}\allowbreak alpha, a, or x.}}
\expandafter\def\csname ednote185\endcsname{\edNote{183}{293}{Stampa. F -> |F|, apparently missing system bars, unlike the preceding analogous display.}}
\expandafter\def\csname ednote186\endcsname{\edNote{184}{293}{Stampa. n=p-1 -> n=p-2, inconsistent with the exponent p-2 printed at the end of the same item.}}
\expandafter\def\csname ednote187\endcsname{\edNote{185}{295}{Dubbio. short mark following the first footnote's final period -- printed hyphen or scan speck.}}
\expandafter\def\csname ednote188\endcsname{\edNote{186}{295}{Stampa. \textbackslash{}\allowbreak leq -> \textbackslash{}\allowbreak geq, the preceding sentence states n \textbackslash{}\allowbreak geq 2r.}}
\expandafter\def\csname ednote189\endcsname{\edNote{187}{295}{Stampa. 2p+1 -> 2p-1, the preceding displayed bound for the composed case has a minus sign.}}
\expandafter\def\csname ednote190\endcsname{\edNote{188}{295}{Stampa. s -> r, the base-point count is denoted r immediately above.}}
\expandafter\def\csname ednote191\endcsname{\edNote{189}{296}{Stampa. S\_\allowbreak \{p-2\} -> S\_\allowbreak \{p-1\}, the canonical space is denoted S\_\allowbreak \{p-1\} in the preceding discussion.}}
\expandafter\def\csname ednote192\endcsname{\edNote{190}{296}{Stampa. p=7 -> p=6, inconsistent with p\textasciicircum{}\{(1)\}=2p-3 and the printed value p\textasciicircum{}\{(1)\}=9.}}
\expandafter\def\csname ednote193\endcsname{\edNote{191}{296}{Stampa. p\textasciicircum{}\{(1)\}=4n-2 -> p\textasciicircum{}\{(1)\}=4n-1, inconsistent with the relation p\textasciicircum{}\{(1)\}=2p-3 stated on this page.}}
\expandafter\def\csname ednote194\endcsname{\edNote{192}{296}{Stampa. 4n-8 -> 4n-7 and 4n-4 -> 4n-3, both inconsistent with the relation p\textasciicircum{}\{(1)\}=2p-3 stated on this page.}}
\expandafter\def\csname ednote195\endcsname{\edNote{193}{297}{Stampa. la serie -> le serie, the following adjective and verb are plural.}}
\expandafter\def\csname ednote196\endcsname{\edNote{194}{299}{Stampa. S\_\allowbreak \{p-7\} -> S\_\allowbreak \{p-1\}, the preceding page places this normal surface in S\_\allowbreak \{p-1\}.}}
\expandafter\def\csname ednote197\endcsname{\edNote{195}{303}{Stampa. 3p\_\allowbreak 8 -> 3p\_\allowbreak 8z, the factor z present in the corresponding terms of the two preceding equations is absent.}}
\expandafter\def\csname ednote198\endcsname{\edNote{196}{305}{Stampa. p=1 -> p=2, this refers to the p=2 case just described.}}
\expandafter\def\csname ednote199\endcsname{\edNote{197}{305}{Dubbio. the conjunction after p\textasciicircum{}\{(1)\}=2 -- è / e; the apparent accent may be a stray mark.}}
\expandafter\def\csname ednote200\endcsname{\edNote{198}{306}{Stampa. F'\_\allowbreak 2 -> F'\_\allowbreak 3, the same composite surface is written F\_\allowbreak 3+F'\_\allowbreak 3 immediately above.}}
\expandafter\def\csname ednote201\endcsname{\edNote{199}{307}{Stampa. addittivo -> additivo, apparent doubled t.}}
\expandafter\def\csname ednote202\endcsname{\edNote{200}{309}{Stampa. (K\_\allowbreak 2 - |\textbackslash{}\allowbreak theta|)' -> |K\_\allowbreak 2 - \textbackslash{}\allowbreak theta|', the printed bars enclose only theta rather than the residual system named in the preceding clause.}}
\expandafter\def\csname ednote203\endcsname{\edNote{201}{310}{Stampa. C -> C', the preceding sentence describes the adjoint curves C'.}}
\expandafter\def\csname ednote204\endcsname{\edNote{202}{311}{Stampa. bicanoniche -> tricanoniche, the preceding discussion concerns the tricanonical cubic model.}}
\expandafter\def\csname ednote205\endcsname{\edNote{203}{312}{Dubbio. capital E before "questo sistema" -- E or È.}}
\expandafter\def\csname ednote206\endcsname{\edNote{204}{318}{Stampa. la la -> la, the article is repeated across two printed lines.}}
\expandafter\def\csname ednote207\endcsname{\edNote{205}{318}{Stampa. genere 2 -> genere 1, this quartic is described as genus 1 above and on the following page.}}
\expandafter\def\csname ednote208\endcsname{\edNote{206}{319}{Stampa. \textbackslash{}\allowbreak lvert C\_\allowbreak \{14\}\textbackslash{}\allowbreak rvert-sega -> \textbackslash{}\allowbreak lvert C\_\allowbreak \{14\}\textbackslash{}\allowbreak rvert sega, an extraneous hyphen separates the subject from the verb.}}
\expandafter\def\csname ednote209\endcsname{\edNote{207}{320}{Dubbio. punctuation after “da” -- a full stop or a scan speck.}}
\expandafter\def\csname ednote210\endcsname{\edNote{208}{321}{Stampa. orta -> porta, the initial p is absent from the visible word.}}
\expandafter\def\csname ednote211\endcsname{\edNote{209}{324}{Dubbio. the small mark following “cfr.” -- a scan speck or a second printed stop.}}
\expandafter\def\csname ednote212\endcsname{\edNote{210}{330}{Stampa. q >= -p -> q >= p, the immediately preceding statement asserts a family of dimension p with p\_\allowbreak a=-p.}}
\expandafter\def\csname ednote213\endcsname{\edNote{211}{332}{Stampa. fade -> falde, the surrounding text refers to the same hypersurface sheets.}}
\expandafter\def\csname ednote214\endcsname{\edNote{212}{333}{Stampa. e dir vero -> a dir vero, apparent misprint in the idiomatic phrase.}}
\expandafter\def\csname ednote215\endcsname{\edNote{213}{333}{Stampa. un questione -> una questione, missing feminine ending before questione.}}
\expandafter\def\csname ednote216\endcsname{\edNote{214}{334}{Stampa. (riducibile)con -> (riducibile) con, no separating space is printed.}}
\expandafter\def\csname ednote217\endcsname{\edNote{215}{335}{Stampa. r=p+n-\textbackslash{}\allowbreak pi+1 -> r=p\_\allowbreak a+n-\textbackslash{}\allowbreak pi+1, the parallel formula for r' uses p\_\allowbreak a.}}
\expandafter\def\csname ednote218\endcsname{\edNote{216}{340}{Dubbio. relation after the first p\_\allowbreak g -- printed equality crossed by a later stroke, or printed \textbackslash{}\allowbreak ne.}}
\expandafter\def\csname ednote219\endcsname{\edNote{217}{340}{Dubbio. spacing and letters in "a l’ora" -- "a l’ora" or "all’ora".}}
\expandafter\def\csname ednote220\endcsname{\edNote{218}{344}{Dubbio. damaged initial in "d’ordine" -- d or l.}}
\expandafter\def\csname ednote221\endcsname{\edNote{219}{344}{Stampa. Comptes, rendus -> Comptes rendus, apparently extraneous comma in the journal title.}}
\expandafter\def\csname ednote222\endcsname{\edNote{220}{346}{Dubbio. dot-like mark after "(gruppi di punti)" -- scan speck or printed full stop.}}
\expandafter\def\csname ednote223\endcsname{\edNote{221}{348}{Dubbio. the symbols in the final parenthesis -- Greek \textbackslash{}\allowbreak nu or Latin v; \textbackslash{}\allowbreak varrho or \textbackslash{}\allowbreak rho.}}
\expandafter\def\csname ednote224\endcsname{\edNote{222}{348}{Dubbio. the symbols in the final parenthesis -- Greek \textbackslash{}\allowbreak nu or Latin v; \textbackslash{}\allowbreak varrho or \textbackslash{}\allowbreak rho.}}
\expandafter\def\csname ednote225\endcsname{\edNote{223}{349}{Stampa. C\_\allowbreak 8 -> C\_\allowbreak 5, the cutting curves are called quintics C\_\allowbreak 5 on p. 348.}}
\expandafter\def\csname ednote226\endcsname{\edNote{224}{354}{Dubbio. the damaged n in “canonica” -- canonica or canoica.}}
\expandafter\def\csname ednote227\endcsname{\edNote{225}{357}{Stampa. indipendenti. rappresentazione -> indipendenti, rappresentazione, the printed full stop is followed by a lower-case continuation.}}
\expandafter\def\csname ednote228\endcsname{\edNote{226}{358}{Stampa. "XXIII,," -> "XXIII,", apparently duplicated comma after the volume number.}}
\expandafter\def\csname ednote229\endcsname{\edNote{227}{360}{Dubbio. mark after "A" in "A tal fine" -- scan mark or apostrophe.}}
\expandafter\def\csname ednote230\endcsname{\edNote{228}{361}{Dubbio. initial word -- e or è; the mark above it may be scan noise.}}
\expandafter\def\csname ednote231\endcsname{\edNote{229}{362}{Stampa. M'\_\allowbreak \{q-1\} -> K\_\allowbreak \{q-1\}, the image of the series was denoted K\_\allowbreak \{q-1\} earlier in this paragraph.}}
\expandafter\def\csname ednote232\endcsname{\edNote{230}{363}{Stampa. first M\_\allowbreak \{q-2\} -> M\_\allowbreak \{q-1\}, the preceding paragraph places this intersection on M\_\allowbreak \{q-1\}.}}
\expandafter\def\csname ednote233\endcsname{\edNote{231}{365}{Stampa. q\_\allowbreak g = 4 -> q = 4, this case is introduced as q = 4 in the preceding paragraph.}}
\expandafter\def\csname ednote234\endcsname{\edNote{232}{366}{Stampa. p\_\allowbreak a = 0 -> p\_\allowbreak g = 0, the converse concerns the geometric-genus condition stated immediately before.}}
\expandafter\def\csname ednote235\endcsname{\edNote{233}{366}{Dubbio. mark after "affermata" -- scan noise or comma.}}
\expandafter\def\csname ednote236\endcsname{\edNote{234}{367}{Dubbio. prime on gamma -- \textbackslash{}\allowbreak gamma'\_\allowbreak i or \textbackslash{}\allowbreak gamma\_\allowbreak i.}}
\expandafter\def\csname ednote237\endcsname{\edNote{235}{368}{Dubbio. Only the ending "NO" is exposed; the full head is read from the matching page in scan s197L -- CAPITOLO NONO / another obscured head ending in NO.}}
\expandafter\def\csname ednote238\endcsname{\edNote{236}{368}{Dubbio. Several overlaid fragments are cropped -- g / p; qu / q; TI / UI; L / C; e F / F.}}
\expandafter\def\csname ednote239\endcsname{\edNote{237}{368}{Dubbio. The covered portions of this paragraph, including its opening, are restored solely from the matching exposed text in scan s197L -- identical wording / wording differing beneath the fold.}}
\expandafter\def\csname ednote240\endcsname{\edNote{238}{368}{Dubbio. The covered beginning and portions of the next line are read from the matching paragraph in scan s197L -- identical wording / wording differing beneath the fold.}}
\expandafter\def\csname ednote241\endcsname{\edNote{239}{368}{Stampa. In-text (1), footnote (3) -> in-text (3), the printed call and the number printed beside its footnote disagree.}}
\expandafter\def\csname ednote242\endcsname{\edNote{240}{372}{Dubbio. A diagonal blemish immediately follows "formeranno" -- "formeranno una" / "formeranno, una".}}
\expandafter\def\csname ednote243\endcsname{\edNote{241}{375}{Dubbio. The mark after "che" is read as a comma -- comma / scan blemish.}}
\expandafter\def\csname ednote244\endcsname{\edNote{242}{376}{Dubbio. small mark before "Cap." in the footnote -- scan speck or an additional full stop.}}
\expandafter\def\csname ednote245\endcsname{\edNote{243}{376}{Dubbio. small mark after "esplicito" -- scan speck or a printed full stop.}}
\expandafter\def\csname ednote246\endcsname{\edNote{244}{377}{Stampa. (t+1) -> (t-1), the preceding general formula has the factor t\_\allowbreak i-1.}}
\expandafter\def\csname ednote247\endcsname{\edNote{245}{378}{Stampa. -r+1 -> +r+1, the printed right-hand expression has the opposite sign for this r term.}}
\expandafter\def\csname ednote248\endcsname{\edNote{246}{378}{Stampa. è quindi -> e quindi, the accented verb interrupts the coordination of the two clauses introduced by "che".}}
\expandafter\def\csname ednote249\endcsname{\edNote{247}{382}{Stampa. II e 1 -> II e I, the hypotheses are labelled I and II.}}
\expandafter\def\csname ednote250\endcsname{\edNote{248}{383}{Stampa. delle serie -> della serie, the following adjective "s-bicanonica" is singular.}}
\expandafter\def\csname ednote251\endcsname{\edNote{249}{383}{Stampa. sK -> sK\textasciicircum{}\textbackslash{}\allowbreak pi, the following sentence compares the same system with |K\textasciicircum{}\textbackslash{}\allowbreak pi|.}}
\expandafter\def\csname ednote252\endcsname{\edNote{250}{383}{Stampa. para-s-canoniche -> para-s\textbackslash{}\allowbreak pi-canoniche, the preceding text calls C para-\textbackslash{}\allowbreak pi-canonica and then takes its s-fold multiple.}}
\expandafter\def\csname ednote253\endcsname{\edNote{251}{385}{Stampa. 341). -> 341))., the parenthesis opened before "Teorema" is not closed in print.}}
\expandafter\def\csname ednote254\endcsname{\edNote{252}{385}{Stampa. dato-con -> dato con, an extraneous hyphen joins these words.}}
\expandafter\def\csname ednote255\endcsname{\edNote{253}{386}{Stampa. 2(K+K\_\allowbreak 1+\textbackslash{}\allowbreak ldotp\textbackslash{}\allowbreak ldotp\textbackslash{}\allowbreak ldotp\textbackslash{}\allowbreak ldotp+K\_\allowbreak r) -> 2(K\_\allowbreak 1+\textbackslash{}\allowbreak ldotp\textbackslash{}\allowbreak ldotp\textbackslash{}\allowbreak ldotp\textbackslash{}\allowbreak ldotp+K\_\allowbreak r), the following virtual curve omits the unindexed K.}}
\expandafter\def\csname ednote256\endcsname{\edNote{254}{388}{Stampa. indipsndentemente -> indipendentemente, apparent substitution of s for e.}}
\expandafter\def\csname ednote257\endcsname{\edNote{255}{388}{Stampa. K\_\allowbreak \{\textbackslash{}\allowbreak pi-2\} -> K\_\allowbreak \{2\textbackslash{}\allowbreak pi-2\}, the curve has the latter subscript earlier in this paragraph.}}
\expandafter\def\csname ednote258\endcsname{\edNote{256}{389}{Stampa. K\_\allowbreak \{\textbackslash{}\allowbreak pi-2\} -> K\_\allowbreak \{2\textbackslash{}\allowbreak pi-2\}, the corresponding curve is given the latter subscript below.}}
\expandafter\def\csname ednote259\endcsname{\edNote{257}{389}{Stampa. g -> \textbackslash{}\allowbreak bar g, the same series carry an overbar later in the paragraph.}}
\expandafter\def\csname ednote260\endcsname{\edNote{258}{390}{Stampa. K\_\allowbreak \{\textbackslash{}\allowbreak pi-2\} in S\_\allowbreak \textbackslash{}\allowbreak pi -> K\_\allowbreak \{2\textbackslash{}\allowbreak pi-2\} in S\_\allowbreak \{\textbackslash{}\allowbreak pi-2\}, comparison with the same construction on the preceding pages.}}
\expandafter\def\csname ednote261\endcsname{\edNote{259}{390}{Stampa. sivuole (printed si- / vuole) -> si vuole, an anomalous line-end hyphen joins two separate words.}}
\expandafter\def\csname ednote262\endcsname{\edNote{260}{395}{Stampa. n=6 -> m=6, the order of the curves K\_\allowbreak m is denoted m on the preceding page.}}
\expandafter\def\csname ednote263\endcsname{\edNote{261}{396}{Dubbio. the first closed-path symbol -- \textbackslash{}\allowbreak nu, a faint \textbackslash{}\allowbreak gamma, or Latin v.}}
\expandafter\def\csname ednote264\endcsname{\edNote{262}{397}{Stampa. contratto -> contatto, an extra r is printed in the contact term.}}
\expandafter\def\csname ednote265\endcsname{\edNote{263}{400}{Stampa. 13-12p\_\allowbreak a-p\textasciicircum{}\{(1)\} -> 13+12p\_\allowbreak a-p\textasciicircum{}\{(1)\}, the following lower bound on p\_\allowbreak a requires a plus sign.}}
\expandafter\def\csname ednote266\endcsname{\edNote{264}{403}{Dubbio. punctuation after h\_\allowbreak i\textbackslash{}\allowbreak varrho\_\allowbreak i -- a full stop, a faint comma, or a scan speck.}}
\expandafter\def\csname ednote267\endcsname{\edNote{265}{405}{Stampa. (z-a)\textasciicircum{}\{r\_\allowbreak 1\} -> (z-a\_\allowbreak 1)\textasciicircum{}\{r\_\allowbreak 1\}, the preceding formula indexes the first branch value by 1.}}
\expandafter\def\csname ednote268\endcsname{\edNote{266}{405}{Stampa. 2n-2\textbackslash{}\allowbreak pi-2 -> 2n+2\textbackslash{}\allowbreak pi-2, compare the corresponding equality on printed p. 404.}}
\expandafter\def\csname ednote269\endcsname{\edNote{267}{408}{Dubbio. punctuation after "formula" -- a full stop or a scanning speck.}}
\expandafter\def\csname ednote270\endcsname{\edNote{268}{408}{Stampa. \textbackslash{}\allowbreak Sigma n(s-1) -> \textbackslash{}\allowbreak Sigma\textbackslash{}\allowbreak frac\{n\}\{s\}(s-1), the corresponding equality on page 410 contains the factor n/s.}}
\expandafter\def\csname ednote271\endcsname{\edNote{269}{408}{Stampa. \textbackslash{}\allowbreak geq1 -> \textbackslash{}\allowbreak geq2, the equality with 2n gives the stronger bound used below for t=3.}}
\expandafter\def\csname ednote272\endcsname{\edNote{270}{408}{Stampa. \textbackslash{}\allowbreak geq1 -> \textbackslash{}\allowbreak geq2, the apparently intended bound is the one used in the following t=3 calculation.}}
\expandafter\def\csname ednote273\endcsname{\edNote{271}{408}{Stampa. t-1 -> t-2, the next displayed calculation gives a bound of 1 when t=3.}}
\expandafter\def\csname ednote274\endcsname{\edNote{272}{409}{Stampa. s\_\allowbreak 1\textbackslash{}\allowbreak geq3 -> s\_\allowbreak 1=3, the following fourth case separately treats s\_\allowbreak 1\textbackslash{}\allowbreak geq4.}}
\expandafter\def\csname ednote275\endcsname{\edNote{273}{411}{Stampa. s=2 -> s\_\allowbreak 1=2, the other two multiplicities in this example are indexed 2 and 3.}}
\expandafter\def\csname ednote276\endcsname{\edNote{274}{411}{Stampa. \textbackslash{}\allowbreak Sigma\textbackslash{}\allowbreak frac\{s+11\}\{s\} -> \textbackslash{}\allowbreak Sigma\textbackslash{}\allowbreak left[\textbackslash{}\allowbreak frac\{s+11\}\{s\}\textbackslash{}\allowbreak right], the integer-part brackets used in the general formula and subsequent calculations are absent here.}}
\expandafter\def\csname ednote277\endcsname{\edNote{275}{411}{Stampa. 24 -> 6t, the preceding bound is 6 for each of t summands.}}
\expandafter\def\csname ednote278\endcsname{\edNote{276}{411}{Stampa. P\_\allowbreak \{12\}\textbackslash{}\allowbreak geq12t-47\textbackslash{}\allowbreak geq12(t-4)+1>1 -> P\_\allowbreak \{12\}\textbackslash{}\allowbreak geq6t-23=6(t-4)+1>1, using the preceding per-term bound with t summands.}}
\expandafter\def\csname ednote279\endcsname{\edNote{277}{412}{Stampa. s\_\allowbreak 1>4 -> s\_\allowbreak 1\textbackslash{}\allowbreak geq4, the displayed enumeration otherwise omits s\_\allowbreak 1=4.}}
\expandafter\def\csname ednote280\endcsname{\edNote{278}{415}{Stampa. pluricanonice -> pluricanoniche, the h is absent in print.}}
\expandafter\def\csname ednote281\endcsname{\edNote{279}{415}{Stampa. prima serie -> prima specie, the preceding discussion calls these transformations "di prima specie".}}
\expandafter\def\csname ednote282\endcsname{\edNote{280}{416}{Stampa. (x-b) -> (z-b), the other factors and the subsequent equations use z.}}
\expandafter\def\csname ednote283\endcsname{\edNote{281}{422}{Dubbio. mark after the final \textbackslash{}\allowbreak tau in this row -- period, short dash, or scan speck.}}
\expandafter\def\csname ednote284\endcsname{\edNote{282}{422}{Stampa. \textbackslash{}\allowbreak pi\textbackslash{}\allowbreak omega=\textbackslash{}\allowbreak pi -> \textbackslash{}\allowbreak tau\textbackslash{}\allowbreak omega=\textbackslash{}\allowbreak pi, the same left-hand side already equals \textbackslash{}\allowbreak omega\_\allowbreak 1 in this row and the reverse composition is listed above.}}
\expandafter\def\csname ednote285\endcsname{\edNote{283}{425}{Dubbio. punctuation after P\_\allowbreak \{12\}=1 -- comma, possibly followed by a period.}}
\expandafter\def\csname ednote286\endcsname{\edNote{284}{426}{Dubbio. the letter following 2 in both subscripts of \textbackslash{}\allowbreak varphi -- n or m.}}
\expandafter\def\csname ednote287\endcsname{\edNote{285}{430}{Dubbio. the mark after "formano" -- period, comma, or scan speck.}}
\expandafter\def\csname ednote288\endcsname{\edNote{286}{431}{Stampa. Schwartz -> Schwarz, the running text on this page spells the same name Schwarz.}}
\expandafter\def\csname ednote289\endcsname{\edNote{287}{431}{Stampa. f -> F, the surface is otherwise denoted by capital F in this paragraph.}}
\expandafter\def\csname ednote290\endcsname{\edNote{288}{432}{Stampa. birationelles -> birationnelles, apparent missing n in the French adjective.}}
\expandafter\def\csname ednote291\endcsname{\edNote{289}{433}{Dubbio. the footer title's spacing -- letter-spaced italic or ordinary italic.}}
\expandafter\def\csname ednote292\endcsname{\edNote{290}{434}{Stampa. è possiede -> e possiede, the conjunction is printed with an accent.}}
\expandafter\def\csname ednote293\endcsname{\edNote{291}{435}{Stampa. eccezio. i -> eccezioni, the printed word has a point and a gap where n is expected.}}
\expandafter\def\csname ednote294\endcsname{\edNote{292}{437}{Stampa. subgroup \textbackslash{}\allowbreak infty\textasciicircum{}p and group \textbackslash{}\allowbreak infty\textasciicircum{}q -> subgroup \textbackslash{}\allowbreak infty\textasciicircum{}q and group \textbackslash{}\allowbreak infty\textasciicircum{}p, the dimensions are reversed relative to the surrounding notation.}}
\expandafter\def\csname ednote295\endcsname{\edNote{293}{440}{Stampa. sopra F -> sopra \textbackslash{}\allowbreak Phi, the following sentence explicitly places the curves L on \textbackslash{}\allowbreak Phi.}}
\expandafter\def\csname ednote296\endcsname{\edNote{294}{440}{Dubbio. final punctuation in the second footnote -- hyphen or damaged full stop.}}
\expandafter\def\csname ednote297\endcsname{\edNote{295}{442}{Dubbio. dot immediately after the footnote number -- printed full stop or scan speck.}}
\expandafter\def\csname ednote298\endcsname{\edNote{296}{443}{Stampa. aventi -> avente, apparent agreement error with the singular “un sistema bicanonico”.}}
\expandafter\def\csname ednote299\endcsname{\edNote{297}{444}{Dubbio. mark between “le” and “quali” -- printed hyphen or scan artefact; mark after “si” -- ink artefact or apostrophe.}}
\expandafter\def\csname ednote300\endcsname{\edNote{298}{444}{Stampa. dalla -> della, apparent preposition error after “conseguenza”.}}
\expandafter\def\csname ednote301\endcsname{\edNote{299}{445}{Dubbio. mark after “grado” -- scan speck or full stop.}}
\expandafter\def\csname ednote302\endcsname{\edNote{300}{445}{Dubbio. punctuation after 2(\textbackslash{}\allowbreak varrho-1) -- comma or scan speck.}}
\expandafter\def\csname ednote303\endcsname{\edNote{301}{446}{Stampa. Picard, ha -> Picard ha, apparently intrusive comma between subject and verb.}}
\expandafter\def\csname ednote304\endcsname{\edNote{302}{448}{Dubbio. small mark above the p in this display -- scan speck with no prime, or p'\_\allowbreak g.}}
\expandafter\def\csname ednote305\endcsname{\edNote{303}{449}{Stampa. \textbackslash{}\allowbreak le0 -> \textbackslash{}\allowbreak ge0, apparently reversed inequality in view of the bound asserted on the next line.}}
\expandafter\def\csname ednote306\endcsname{\edNote{304}{449}{Stampa. p\_\allowbreak a=1 -> p\_\allowbreak a=-1, a minus sign apparently omitted in the limiting case.}}
\expandafter\def\csname ednote307\endcsname{\edNote{305}{451}{Stampa. Cap. XII -> Cap. VII, apparently inconsistent with the reference for the same family on the preceding page.}}
\expandafter\def\csname ednote308\endcsname{\edNote{306}{451}{Stampa. quale ordini -> quali ordini, adjective and plural noun do not agree.}}
\expandafter\def\csname ednote309\endcsname{\edNote{307}{453}{Stampa. d ordine -> d’ordine, apparently missing elision apostrophe.}}
\expandafter\def\csname ednote310\endcsname{\edNote{308}{453}{Dubbio. final punctuation of this footnote -- period, short dash.}}
\expandafter\def\csname ednote311\endcsname{\edNote{309}{455}{Stampa. "hyperellyptiques" -> "hyperelliptiques", apparent French spelling error.}}
\expandafter\def\csname ednote312\endcsname{\edNote{310}{456}{Stampa. "0)." -> "0.", unmatched closing parenthesis.}}
\expandafter\def\csname ednote313\endcsname{\edNote{311}{457}{Stampa. "distinti" -> "distinte", apparent agreement error referring to "superficie".}}
\expandafter\def\csname ednote314\endcsname{\edNote{312}{457}{Stampa. B') -> B''), the preceding distinction assigns the finite families to B''.}}
\expandafter\def\csname ednote315\endcsname{\edNote{313}{457}{Stampa. first B') and final B'') -> first B'') and final B'), these labels reverse the distinction stated above.}}
\expandafter\def\csname ednote316\endcsname{\edNote{314}{459}{Stampa. "pag. 463" -> "pag. 461", the Notizia reproduced in this packet begins on printed page 461.}}
\expandafter\def\csname ednote317\endcsname{\edNote{315}{460}{Dubbio. the small mark after "p\_\allowbreak a=-1" -- a full stop (read here) or a scan speck.}}
\expandafter\def\csname ednote318\endcsname{\edNote{316}{461}{Dubbio. the short stroke after "ellittiche" -- a printed hyphen (read here) or a scan mark.}}
\expandafter\def\csname ednote319\endcsname{\edNote{317}{463}{Dubbio. the small mark after "schiera" -- a comma (read here) or a scan speck.}}
\expandafter\def\csname ednote320\endcsname{\edNote{318}{463}{Dubbio. punctuation after "genere" -- a colon (read here) or no printed punctuation.}}
\expandafter\def\csname ednote321\endcsname{\edNote{319}{463}{Stampa. final "P\_\allowbreak 4=0,1" -> "P\_\allowbreak 6=0,1", P\_\allowbreak 4 is repeated in the list of plurigenera.}}
\expandafter\def\csname ednote322\endcsname{\edNote{320}{464}{Dubbio. possible accent over the first o in "piuttosto" -- piuttosto / piuttósto.}}
\expandafter\def\csname wpAlias:[s002R]\endcsname{V}
\expandafter\def\csname wpAlias:[s004R]\endcsname{IX}
\expandafter\def\csname wpAlias:[s005R]\endcsname{XI}
\expandafter\def\csname wpAlias:[s008L]\endcsname{XVI}
\expandafter\def\csname wpAlias:[s015L]\endcsname{14}
\expandafter\def\csname wpAlias:[s015R]\endcsname{15}
\expandafter\def\csname wpAlias:[s030R]\endcsname{41}
\expandafter\def\csname wpAlias:[s043L]\endcsname{66}
\expandafter\def\csname wpAlias:[s043R]\endcsname{67}
\expandafter\def\csname wpAlias:[s054R]\endcsname{89}
\expandafter\def\csname wpAlias:[s060R]\endcsname{101}
\expandafter\def\csname wpDefer:136\endcsname{1}
\expandafter\def\csname wpDefer:145\endcsname{1}
\expandafter\def\csname wpAlias:[s095L]\endcsname{166}
\expandafter\def\csname wpAlias:[s095R]\endcsname{167}
\expandafter\def\csname wpDefer:206\endcsname{1}
\expandafter\def\csname wpAlias:[s120L]\endcsname{216}
\expandafter\def\csname wpAlias:[s120R]\endcsname{217}
\expandafter\def\csname wpAlias:[s130L]\endcsname{236}
\expandafter\def\csname wpAlias:[s130R]\endcsname{237}
\expandafter\def\csname wpDefer:259\endcsname{1}
\expandafter\def\csname wpAlias:[s146R]\endcsname{267}
\expandafter\def\csname wpAlias:[s174R]\endcsname{325}
\expandafter\def\csname wpAlias:[s198L]\endcsname{370}
\expandafter\def\csname wpAlias:[s198R]\endcsname{371}
\expandafter\def\csname wpDefer:409\endcsname{1}
\expandafter\def\csname wpAlias:[s226R]\endcsname{427}
\expandafter\def\csname wpDefer:445\endcsname{1}
\expandafter\def\csname wpNotes:[s001R]\endcsname{\csname ednote1\endcsname}
\expandafter\def\csname wpNotes:[s004R]\endcsname{\csname ednote2\endcsname\csname ednote3\endcsname}
\expandafter\def\csname wpNotes:[s005R]\endcsname{\csname ednote4\endcsname\csname ednote5\endcsname\csname ednote6\endcsname}
\expandafter\def\csname wpNotes:XIV\endcsname{\csname ednote7\endcsname}
\expandafter\def\csname wpNotes:3\endcsname{\csname ednote8\endcsname}
\expandafter\def\csname wpNotes:5\endcsname{\csname ednote9\endcsname}
\expandafter\def\csname wpNotes:11\endcsname{\csname ednote10\endcsname}
\expandafter\def\csname wpNotes:12\endcsname{\csname ednote11\endcsname\csname ednote12\endcsname\csname ednote13\endcsname}
\expandafter\def\csname wpNotes:19\endcsname{\csname ednote15\endcsname}
\expandafter\def\csname wpNotes:22\endcsname{\csname ednote16\endcsname}
\expandafter\def\csname wpNotes:24\endcsname{\csname ednote17\endcsname\csname ednote18\endcsname}
\expandafter\def\csname wpNotes:27\endcsname{\csname ednote19\endcsname}
\expandafter\def\csname wpNotes:34\endcsname{\csname ednote20\endcsname}
\expandafter\def\csname wpNotes:39\endcsname{\csname ednote21\endcsname\csname ednote22\endcsname}
\expandafter\def\csname wpNotes:[s030R]\endcsname{\csname ednote23\endcsname\csname ednote24\endcsname}
\expandafter\def\csname wpNotes:43\endcsname{\csname ednote25\endcsname}
\expandafter\def\csname wpNotes:44\endcsname{\csname ednote26\endcsname\csname ednote27\endcsname}
\expandafter\def\csname wpNotes:45\endcsname{\csname ednote28\endcsname}
\expandafter\def\csname wpNotes:49\endcsname{\csname ednote29\endcsname}
\expandafter\def\csname wpNotes:50\endcsname{\csname ednote30\endcsname}
\expandafter\def\csname wpNotes:55\endcsname{\csname ednote31\endcsname}
\expandafter\def\csname wpNotes:57\endcsname{\csname ednote32\endcsname}
\expandafter\def\csname wpNotes:58\endcsname{\csname ednote33\endcsname\csname ednote34\endcsname}
\expandafter\def\csname wpNotes:61\endcsname{\csname ednote35\endcsname\csname ednote36\endcsname}
\expandafter\def\csname wpNotes:68\endcsname{\csname ednote37\endcsname\csname ednote38\endcsname}
\expandafter\def\csname wpNotes:77\endcsname{\csname ednote39\endcsname\csname ednote40\endcsname\csname ednote41\endcsname}
\expandafter\def\csname wpNotes:79\endcsname{\csname ednote42\endcsname\csname ednote43\endcsname}
\expandafter\def\csname wpNotes:84\endcsname{\csname ednote45\endcsname\csname ednote46\endcsname}
\expandafter\def\csname wpNotes:85\endcsname{\csname ednote47\endcsname}
\expandafter\def\csname wpNotes:87\endcsname{\csname ednote48\endcsname\csname ednote49\endcsname\csname ednote50\endcsname\csname ednote51\endcsname}
\expandafter\def\csname wpNotes:[s054R]\endcsname{\csname ednote52\endcsname}
\expandafter\def\csname wpNotes:91\endcsname{\csname ednote53\endcsname}
\expandafter\def\csname wpNotes:92\endcsname{\csname ednote54\endcsname}
\expandafter\def\csname wpNotes:94\endcsname{\csname ednote55\endcsname}
\expandafter\def\csname wpNotes:95\endcsname{\csname ednote56\endcsname}
\expandafter\def\csname wpNotes:98\endcsname{\csname ednote57\endcsname}
\expandafter\def\csname wpNotes:99\endcsname{\csname ednote58\endcsname}
\expandafter\def\csname wpNotes:[s060R]\endcsname{\csname ednote59\endcsname}
\expandafter\def\csname wpNotes:102\endcsname{\csname ednote60\endcsname\csname ednote61\endcsname}
\expandafter\def\csname wpNotes:103\endcsname{\csname ednote62\endcsname\csname ednote63\endcsname}
\expandafter\def\csname wpNotes:109\endcsname{\csname ednote64\endcsname\csname ednote65\endcsname}
\expandafter\def\csname wpNotes:111\endcsname{\csname ednote66\endcsname\csname ednote67\endcsname}
\expandafter\def\csname wpNotes:115\endcsname{\csname ednote68\endcsname}
\expandafter\def\csname wpNotes:116\endcsname{\csname ednote69\endcsname\csname ednote70\endcsname}
\expandafter\def\csname wpNotes:118\endcsname{\csname ednote71\endcsname}
\expandafter\def\csname wpNotes:119\endcsname{\csname ednote72\endcsname\csname ednote73\endcsname}
\expandafter\def\csname wpNotes:121\endcsname{\csname ednote74\endcsname}
\expandafter\def\csname wpNotes:124\endcsname{\csname ednote75\endcsname\csname ednote76\endcsname\csname ednote77\endcsname\csname ednote78\endcsname}
\expandafter\def\csname wpNotes:127\endcsname{\csname ednote79\endcsname}
\expandafter\def\csname wpNotes:130\endcsname{\csname ednote80\endcsname\csname ednote81\endcsname}
\expandafter\def\csname wpNotes:133\endcsname{\csname ednote82\endcsname}
\expandafter\def\csname wpNotes:136\endcsname{\csname ednote83\endcsname}
\expandafter\def\csname wpNotes:141\endcsname{\csname ednote84\endcsname}
\expandafter\def\csname wpNotes:144\endcsname{\csname ednote85\endcsname\csname ednote86\endcsname}
\expandafter\def\csname wpNotes:145\endcsname{\csname ednote87\endcsname}
\expandafter\def\csname wpNotes:146\endcsname{\csname ednote88\endcsname}
\expandafter\def\csname wpNotes:148\endcsname{\csname ednote89\endcsname}
\expandafter\def\csname wpNotes:149\endcsname{\csname ednote90\endcsname}
\expandafter\def\csname wpNotes:150\endcsname{\csname ednote91\endcsname}
\expandafter\def\csname wpNotes:152\endcsname{\csname ednote92\endcsname}
\expandafter\def\csname wpNotes:154\endcsname{\csname ednote93\endcsname}
\expandafter\def\csname wpNotes:155\endcsname{\csname ednote94\endcsname}
\expandafter\def\csname wpNotes:157\endcsname{\csname ednote95\endcsname}
\expandafter\def\csname wpNotes:158\endcsname{\csname ednote96\endcsname}
\expandafter\def\csname wpNotes:159\endcsname{\csname ednote97\endcsname\csname ednote98\endcsname}
\expandafter\def\csname wpNotes:160\endcsname{\csname ednote99\endcsname}
\expandafter\def\csname wpNotes:161\endcsname{\csname ednote100\endcsname}
\expandafter\def\csname wpNotes:164\endcsname{\csname ednote101\endcsname\csname ednote102\endcsname}
\expandafter\def\csname wpNotes:165\endcsname{\csname ednote103\endcsname\csname ednote104\endcsname}
\expandafter\def\csname wpNotes:168\endcsname{\csname ednote105\endcsname\csname ednote106\endcsname}
\expandafter\def\csname wpNotes:169\endcsname{\csname ednote107\endcsname}
\expandafter\def\csname wpNotes:170\endcsname{\csname ednote108\endcsname\csname ednote109\endcsname}
\expandafter\def\csname wpNotes:172\endcsname{\csname ednote110\endcsname}
\expandafter\def\csname wpNotes:173\endcsname{\csname ednote111\endcsname\csname ednote112\endcsname}
\expandafter\def\csname wpNotes:175\endcsname{\csname ednote113\endcsname\csname ednote114\endcsname\csname ednote115\endcsname}
\expandafter\def\csname wpNotes:176\endcsname{\csname ednote116\endcsname\csname ednote117\endcsname}
\expandafter\def\csname wpNotes:178\endcsname{\csname ednote118\endcsname}
\expandafter\def\csname wpNotes:182\endcsname{\csname ednote119\endcsname\csname ednote120\endcsname}
\expandafter\def\csname wpNotes:183\endcsname{\csname ednote121\endcsname}
\expandafter\def\csname wpNotes:185\endcsname{\csname ednote122\endcsname\csname ednote123\endcsname}
\expandafter\def\csname wpNotes:186\endcsname{\csname ednote124\endcsname\csname ednote125\endcsname\csname ednote126\endcsname\csname ednote127\endcsname}
\expandafter\def\csname wpNotes:187\endcsname{\csname ednote128\endcsname\csname ednote129\endcsname}
\expandafter\def\csname wpNotes:191\endcsname{\csname ednote130\endcsname\csname ednote131\endcsname}
\expandafter\def\csname wpNotes:194\endcsname{\csname ednote132\endcsname\csname ednote133\endcsname}
\expandafter\def\csname wpNotes:199\endcsname{\csname ednote134\endcsname}
\expandafter\def\csname wpNotes:204\endcsname{\csname ednote135\endcsname}
\expandafter\def\csname wpNotes:205\endcsname{\csname ednote136\endcsname}
\expandafter\def\csname wpNotes:206\endcsname{\csname ednote137\endcsname}
\expandafter\def\csname wpNotes:208\endcsname{\csname ednote138\endcsname}
\expandafter\def\csname wpNotes:212\endcsname{\csname ednote139\endcsname\csname ednote140\endcsname\csname ednote141\endcsname\csname ednote142\endcsname}
\expandafter\def\csname wpNotes:214\endcsname{\csname ednote143\endcsname}
\expandafter\def\csname wpNotes:215\endcsname{\csname ednote144\endcsname}
\expandafter\def\csname wpNotes:[s120R]\endcsname{\csname ednote145\endcsname\csname ednote146\endcsname}
\expandafter\def\csname wpNotes:220\endcsname{\csname ednote147\endcsname}
\expandafter\def\csname wpNotes:222\endcsname{\csname ednote148\endcsname}
\expandafter\def\csname wpNotes:226\endcsname{\csname ednote149\endcsname}
\expandafter\def\csname wpNotes:229\endcsname{\csname ednote150\endcsname}
\expandafter\def\csname wpNotes:230\endcsname{\csname ednote151\endcsname\csname ednote152\endcsname}
\expandafter\def\csname wpNotes:235\endcsname{\csname ednote153\endcsname}
\expandafter\def\csname wpNotes:[s130R]\endcsname{\csname ednote154\endcsname}
\expandafter\def\csname wpNotes:240\endcsname{\csname ednote155\endcsname}
\expandafter\def\csname wpNotes:242\endcsname{\csname ednote156\endcsname}
\expandafter\def\csname wpNotes:244\endcsname{\csname ednote157\endcsname}
\expandafter\def\csname wpNotes:246\endcsname{\csname ednote158\endcsname}
\expandafter\def\csname wpNotes:248\endcsname{\csname ednote159\endcsname}
\expandafter\def\csname wpNotes:251\endcsname{\csname ednote160\endcsname}
\expandafter\def\csname wpNotes:253\endcsname{\csname ednote161\endcsname}
\expandafter\def\csname wpNotes:256\endcsname{\csname ednote162\endcsname}
\expandafter\def\csname wpNotes:257\endcsname{\csname ednote163\endcsname}
\expandafter\def\csname wpNotes:258\endcsname{\csname ednote164\endcsname}
\expandafter\def\csname wpNotes:259\endcsname{\csname ednote165\endcsname}
\expandafter\def\csname wpNotes:260\endcsname{\csname ednote166\endcsname}
\expandafter\def\csname wpNotes:262\endcsname{\csname ednote167\endcsname}
\expandafter\def\csname wpNotes:263\endcsname{\csname ednote168\endcsname\csname ednote169\endcsname\csname ednote170\endcsname\csname ednote171\endcsname}
\expandafter\def\csname wpNotes:265\endcsname{\csname ednote172\endcsname}
\expandafter\def\csname wpNotes:[s146R]\endcsname{\csname ednote173\endcsname}
\expandafter\def\csname wpNotes:274\endcsname{\csname ednote174\endcsname\csname ednote175\endcsname}
\expandafter\def\csname wpNotes:275\endcsname{\csname ednote176\endcsname}
\expandafter\def\csname wpNotes:277\endcsname{\csname ednote177\endcsname}
\expandafter\def\csname wpNotes:282\endcsname{\csname ednote178\endcsname}
\expandafter\def\csname wpNotes:283\endcsname{\csname ednote179\endcsname}
\expandafter\def\csname wpNotes:286\endcsname{\csname ednote180\endcsname}
\expandafter\def\csname wpNotes:287\endcsname{\csname ednote181\endcsname}
\expandafter\def\csname wpNotes:290\endcsname{\csname ednote182\endcsname}
\expandafter\def\csname wpNotes:291\endcsname{\csname ednote183\endcsname}
\expandafter\def\csname wpNotes:292\endcsname{\csname ednote184\endcsname}
\expandafter\def\csname wpNotes:293\endcsname{\csname ednote185\endcsname\csname ednote186\endcsname}
\expandafter\def\csname wpNotes:295\endcsname{\csname ednote187\endcsname\csname ednote188\endcsname\csname ednote189\endcsname\csname ednote190\endcsname}
\expandafter\def\csname wpNotes:296\endcsname{\csname ednote191\endcsname\csname ednote192\endcsname\csname ednote193\endcsname\csname ednote194\endcsname}
\expandafter\def\csname wpNotes:297\endcsname{\csname ednote195\endcsname}
\expandafter\def\csname wpNotes:299\endcsname{\csname ednote196\endcsname}
\expandafter\def\csname wpNotes:303\endcsname{\csname ednote197\endcsname}
\expandafter\def\csname wpNotes:305\endcsname{\csname ednote198\endcsname\csname ednote199\endcsname}
\expandafter\def\csname wpNotes:306\endcsname{\csname ednote200\endcsname}
\expandafter\def\csname wpNotes:307\endcsname{\csname ednote201\endcsname}
\expandafter\def\csname wpNotes:309\endcsname{\csname ednote202\endcsname}
\expandafter\def\csname wpNotes:310\endcsname{\csname ednote203\endcsname}
\expandafter\def\csname wpNotes:311\endcsname{\csname ednote204\endcsname}
\expandafter\def\csname wpNotes:312\endcsname{\csname ednote205\endcsname}
\expandafter\def\csname wpNotes:318\endcsname{\csname ednote206\endcsname\csname ednote207\endcsname}
\expandafter\def\csname wpNotes:319\endcsname{\csname ednote208\endcsname}
\expandafter\def\csname wpNotes:320\endcsname{\csname ednote209\endcsname}
\expandafter\def\csname wpNotes:321\endcsname{\csname ednote210\endcsname}
\expandafter\def\csname wpNotes:324\endcsname{\csname ednote211\endcsname}
\expandafter\def\csname wpNotes:330\endcsname{\csname ednote212\endcsname}
\expandafter\def\csname wpNotes:332\endcsname{\csname ednote213\endcsname}
\expandafter\def\csname wpNotes:333\endcsname{\csname ednote214\endcsname\csname ednote215\endcsname}
\expandafter\def\csname wpNotes:334\endcsname{\csname ednote216\endcsname}
\expandafter\def\csname wpNotes:335\endcsname{\csname ednote217\endcsname}
\expandafter\def\csname wpNotes:340\endcsname{\csname ednote218\endcsname\csname ednote219\endcsname}
\expandafter\def\csname wpNotes:344\endcsname{\csname ednote220\endcsname\csname ednote221\endcsname}
\expandafter\def\csname wpNotes:346\endcsname{\csname ednote222\endcsname}
\expandafter\def\csname wpNotes:348\endcsname{\csname ednote223\endcsname\csname ednote224\endcsname}
\expandafter\def\csname wpNotes:349\endcsname{\csname ednote225\endcsname}
\expandafter\def\csname wpNotes:354\endcsname{\csname ednote226\endcsname}
\expandafter\def\csname wpNotes:357\endcsname{\csname ednote227\endcsname}
\expandafter\def\csname wpNotes:358\endcsname{\csname ednote228\endcsname}
\expandafter\def\csname wpNotes:360\endcsname{\csname ednote229\endcsname}
\expandafter\def\csname wpNotes:361\endcsname{\csname ednote230\endcsname}
\expandafter\def\csname wpNotes:362\endcsname{\csname ednote231\endcsname}
\expandafter\def\csname wpNotes:363\endcsname{\csname ednote232\endcsname}
\expandafter\def\csname wpNotes:365\endcsname{\csname ednote233\endcsname}
\expandafter\def\csname wpNotes:366\endcsname{\csname ednote234\endcsname\csname ednote235\endcsname}
\expandafter\def\csname wpNotes:367\endcsname{\csname ednote236\endcsname}
\expandafter\def\csname wpNotes:368\endcsname{\csname ednote237\endcsname\csname ednote238\endcsname\csname ednote239\endcsname\csname ednote240\endcsname\csname ednote241\endcsname}
\expandafter\def\csname wpNotes:372\endcsname{\csname ednote242\endcsname}
\expandafter\def\csname wpNotes:375\endcsname{\csname ednote243\endcsname}
\expandafter\def\csname wpNotes:376\endcsname{\csname ednote244\endcsname\csname ednote245\endcsname}
\expandafter\def\csname wpNotes:377\endcsname{\csname ednote246\endcsname}
\expandafter\def\csname wpNotes:378\endcsname{\csname ednote247\endcsname\csname ednote248\endcsname}
\expandafter\def\csname wpNotes:382\endcsname{\csname ednote249\endcsname}
\expandafter\def\csname wpNotes:383\endcsname{\csname ednote250\endcsname\csname ednote251\endcsname\csname ednote252\endcsname}
\expandafter\def\csname wpNotes:385\endcsname{\csname ednote253\endcsname\csname ednote254\endcsname}
\expandafter\def\csname wpNotes:386\endcsname{\csname ednote255\endcsname}
\expandafter\def\csname wpNotes:388\endcsname{\csname ednote256\endcsname\csname ednote257\endcsname}
\expandafter\def\csname wpNotes:389\endcsname{\csname ednote258\endcsname\csname ednote259\endcsname}
\expandafter\def\csname wpNotes:390\endcsname{\csname ednote260\endcsname\csname ednote261\endcsname}
\expandafter\def\csname wpNotes:395\endcsname{\csname ednote262\endcsname}
\expandafter\def\csname wpNotes:396\endcsname{\csname ednote263\endcsname}
\expandafter\def\csname wpNotes:397\endcsname{\csname ednote264\endcsname}
\expandafter\def\csname wpNotes:400\endcsname{\csname ednote265\endcsname}
\expandafter\def\csname wpNotes:403\endcsname{\csname ednote266\endcsname}
\expandafter\def\csname wpNotes:405\endcsname{\csname ednote267\endcsname\csname ednote268\endcsname}
\expandafter\def\csname wpNotes:408\endcsname{\csname ednote269\endcsname\csname ednote270\endcsname\csname ednote271\endcsname\csname ednote272\endcsname\csname ednote273\endcsname}
\expandafter\def\csname wpNotes:409\endcsname{\csname ednote274\endcsname}
\expandafter\def\csname wpNotes:411\endcsname{\csname ednote275\endcsname\csname ednote276\endcsname\csname ednote277\endcsname\csname ednote278\endcsname}
\expandafter\def\csname wpNotes:412\endcsname{\csname ednote279\endcsname}
\expandafter\def\csname wpNotes:415\endcsname{\csname ednote280\endcsname\csname ednote281\endcsname}
\expandafter\def\csname wpNotes:416\endcsname{\csname ednote282\endcsname}
\expandafter\def\csname wpNotes:422\endcsname{\csname ednote283\endcsname\csname ednote284\endcsname}
\expandafter\def\csname wpNotes:425\endcsname{\csname ednote285\endcsname}
\expandafter\def\csname wpNotes:426\endcsname{\csname ednote286\endcsname}
\expandafter\def\csname wpNotes:430\endcsname{\csname ednote287\endcsname}
\expandafter\def\csname wpNotes:431\endcsname{\csname ednote288\endcsname\csname ednote289\endcsname}
\expandafter\def\csname wpNotes:432\endcsname{\csname ednote290\endcsname}
\expandafter\def\csname wpNotes:433\endcsname{\csname ednote291\endcsname}
\expandafter\def\csname wpNotes:434\endcsname{\csname ednote292\endcsname}
\expandafter\def\csname wpNotes:435\endcsname{\csname ednote293\endcsname}
\expandafter\def\csname wpNotes:437\endcsname{\csname ednote294\endcsname}
\expandafter\def\csname wpNotes:440\endcsname{\csname ednote295\endcsname\csname ednote296\endcsname}
\expandafter\def\csname wpNotes:442\endcsname{\csname ednote297\endcsname}
\expandafter\def\csname wpNotes:443\endcsname{\csname ednote298\endcsname}
\expandafter\def\csname wpNotes:444\endcsname{\csname ednote299\endcsname\csname ednote300\endcsname}
\expandafter\def\csname wpNotes:445\endcsname{\csname ednote301\endcsname\csname ednote302\endcsname}
\expandafter\def\csname wpNotes:446\endcsname{\csname ednote303\endcsname}
\expandafter\def\csname wpNotes:448\endcsname{\csname ednote304\endcsname}
\expandafter\def\csname wpNotes:449\endcsname{\csname ednote305\endcsname\csname ednote306\endcsname}
\expandafter\def\csname wpNotes:451\endcsname{\csname ednote307\endcsname\csname ednote308\endcsname}
\expandafter\def\csname wpNotes:453\endcsname{\csname ednote309\endcsname\csname ednote310\endcsname}
\expandafter\def\csname wpNotes:455\endcsname{\csname ednote311\endcsname}
\expandafter\def\csname wpNotes:456\endcsname{\csname ednote312\endcsname}
\expandafter\def\csname wpNotes:457\endcsname{\csname ednote313\endcsname\csname ednote314\endcsname\csname ednote315\endcsname}
\expandafter\def\csname wpNotes:459\endcsname{\csname ednote316\endcsname}
\expandafter\def\csname wpNotes:460\endcsname{\csname ednote317\endcsname}
\expandafter\def\csname wpNotes:461\endcsname{\csname ednote318\endcsname}
\expandafter\def\csname wpNotes:463\endcsname{\csname ednote319\endcsname\csname ednote320\endcsname\csname ednote321\endcsname}
\expandafter\def\csname wpNotes:464\endcsname{\csname ednote322\endcsname}
\begin{document}\frontmatter

% !PRINTED
\weberblank{[s001L]}
% head:
% scan: s001L; nessun numero di pagina stampato.
% stamp: Bordi scuri, puntinatura e una macchia diffusa nel margine sinistro; nessun testo stampato.

\clearpage\thispagestyle{empty}\vspace*{12mm}
% !PRINTED
\weberpage{[s001R]}
% head:
% scan: s001R; nessun numero di pagina stampato.
% stamp: Timbro di biblioteca in ebraico nella parte inferiore della pagina; macchie e piccoli segni della scansione.
\begin{center}
{\Large\textbf{\gesperrt{FEDERIGO ENRIQUES}}}\par\vspace{23mm}

{\fontsize{26}{31}\selectfont\bfseries LE SUPERFICIE\\ALGEBRICHE\par}
\end{center}

\weberfigure{}{Marchio editoriale: un uomo con berretto avanza spargendo semi; sulla base l’iscrizione «LABORAVI FIDENTER».}

\begin{center}
\textbf{\gesperrt{NICOLA ZANICHELLI EDITORE}} % ?? uncertain: piccolo segno dopo «ZANICHELLI» -- macchia / punto / virgola.

\textbf{\gesperrt{BOLOGNA 1949}}
\end{center}
\clearpage\editionpage\clearpage

\thispagestyle{empty}\vspace*{50mm}
% !PRINTED
\weberpage{[s002L]}
% head:
% scan: s002L; nessun numero di pagina stampato.
% stamp: Segnature manoscritte nei margini superiore e inferiore sinistri; segnatura «QA / 571 / .E6 / 1949»; cifra isolata «1» impressa sotto la dichiarazione editoriale; timbro capovolto e sbiadito in basso; macchie e bordi della scansione.
\begin{center}
\textbf{L’EDITORE ADEMPIUTI I DOVERI ESERCITERÀ I DIRITTI SANCITI DALLE LEGGI}
\end{center}

\noindent\rule{\linewidth}{0.4pt}
\begin{center}
\textbf{S.T.E.B. - Società Tipografica Editrice Bolognese - \textsc{vi} - 1949}
\end{center}
\clearpage


\clearpage\thispagestyle{plain}\markboth{Prefazione di G. Castelnuovo}{Prefazione di G. Castelnuovo}\pdfbookmark[0]{Prefazione di G. Castelnuovo}{front.castelnuovo}
% !PRINTED
\weberpage{[s002R]}
% head:
% scan: s002R; nessun numero di pagina stampato.
% stamp: Puntinatura e tracce ai margini della scansione.
\textit{Il volume che vede oggi la luce, contiene, rielaborata e ordinata, la parte maggiore dell’opera matematica di Federigo Enriques, vale a dire la teoria delle superficie algebriche, alla quale il suo nome resterà sempre legato. Questa opera iniziata subito dopo la laurea nel \textup{1893} e proseguita per oltre mezzo secolo fino alla morte, con brevi interruzioni quando problemi di filosofia e di storia della scienza occupavano il suo spirito universale, si trova dispersa in una cinquantina di Note e Memorie pubblicate in varie raccolte.}

\textit{A far risaltare l’alto valore dell’opera, che ha aperto un vasto e fecondo campo quasi inesplorato quando l’Enriques cominciò le sue ricerche, gioverà il presente volume. Il quale permetterà ai cultori della geometria algebrica, e in particolare agli studenti universitari, di orientarsi nell’ampio territorio.}

\textit{Il compianto Autore aveva già sentito il bisogno di una sistemazione della geometria sopra le superficie algebriche dopo la sua chiamata all’Università di Roma, dove Egli poteva dedicarsi ai soli corsi superiori di matematica e dove maggiore era il numero degli allievi. Valendosi della collaborazione di Luigi Campedelli, allora suo e mio assistente, oggi professore all’Università di Firenze, Egli raccolse la parte generale della teoria in un volume litografato, che uscì nel \textup{1932}; mentre una seconda parte riguardante categorie particolari di superficie fu pubblicata a stampa, due anni dopo, dal Seminario matematico dell’Università di Roma, col concorso dello stesso collaboratore.}

\textit{Esauriti in pochi anni questi due volumi, l’Autore dovette pensare ad una seconda edizione; e durante gli ozi, cui fu costretto dalle leggi razziali e poi dalla guerra e dalle persecuzioni, Egli meditò una rielaborazione di tutta l’opera col proposito di semplificare o mettere al riparo da critiche alcune dimostrazioni, di condurre a termine od estendere classificazioni di tipi particolari di superficie, e di tener conto dei più recenti risultati a cui Egli stesso od altri ricercatori erano arrivati.}

\textit{Il manoscritto era fortunatamente compiuto quando la morte colse Federigo Enriques nella piena lucidità dello spirito.}

% !PRINTED VI
\weberpage{VI}
% head:
% scan: s003L.
% stamp: Puntinatura, macchie e bordo scuro della scansione.
\textit{L’Autore si era riservato di procedere a talune correzioni definitive sulle bozze di stampa, ma purtroppo la sorte non gli concesse il tempo per questo compito. Se lo assunsero con la finezza del loro ingegno, col ricordo degli insegnamenti del maestro, con affetto di figli, gli ultimi discepoli Pompilj e Franchetta. Ed io che li ho assistiti nei loro dubbi su qualche punto oscuro o meno soddisfacente del testo, posso dire che il compito era arduo ed avrebbe richiesto, per essere interamente assolto, ricerche non di mesi ma di anni. L’Autore stesso ha cura di avvertire fin dalla prefazione che il trattato, più che esporre una dottrina già statica e cristallizzata, aspira a suscitare nel lettore il desiderio di portare complementi e perfezionamenti a varie teorie. E dove il terreno è meno solido l’Autore mette sull’avviso lo studioso.}

\textit{Di questi punti ancora fluidi quello che presenta la difficoltà più ardua ed il maggiore interesse nella teoria generale delle superficie forma il soggetto del Cap. IX, ove si tratta dei sistemi continui di curve algebriche (non contenuti in sistemi lineari) che esistono sopra ogni superficie irregolare. Una geniale intuizione ha condotto l’Enriques nel \textup{1904} ad enunciare e stabilire questa proprietà caratteristica delle superficie irregolari; su di essa si sono appoggiate successive fondamentali ricerche. Però un’attenta critica ha fatto vedere, vari anni dopo, che la dimostrazione dell’Enriques non è soddisfacente. La proprietà è vera, almeno sotto qualche restrizione, come risulta dalle ricerche per via trascendente di Enrico Poincaré. Ma tutti i tentativi compiuti posteriormente dall’Enriques e da altri per dimostrarla mediante considerazioni algebrico-geometriche si sono urtati contro difficoltà sinora insuperate. Ciò è esplicitamente dichiarato nel suddetto Cap. IX, nel quale l’Autore dà anche suggerimenti sopra una via da tentare per giungere alla meta. Debbo confessare che non vedo come quella via possa tradursi in un procedimento irreprensibile.}

\textit{Ad altri dubbi di ben minore rilievo può dar luogo qualche punto dei Capp. VII, VIII e X dove si tratta di classificazione e costruzione di tipi molto particolari di superficie. Qui giova ricordare che le ricerche fondamentali dell’Enriques si possono suddividere in due gruppi:}

\textit{\textup{a)} teoria generale delle superficie algebriche, e in particolare determinazione dei caratteri invarianti per trasformazioni birazionali: sostanzialmente, i \textup{generi;}}

\textit{\textup{b)} classificazione delle superficie dei primi generi.}

\textit{Questo secondo gruppo ha condotto l’Autore alla scoperta di famiglie molto interessanti ed impreviste di superficie, ed ha fornito anche modelli concettuali che hanno suggerito la scoperta di proprietà riposte appartenenti alla teoria generale. È noto però a tutti i ricercatori quali insidie si nascondano nelle classificazioni molto minuziose, dove facilmente sfuggono alcuni casi riposti, ed altri, ritenuti in}
%CONT>>
% !PRINTED VII
\weberpage{VII}
% head:
%<<CONT
% scan: s003R.
% stamp: Puntinatura della scansione e tratto obliquo nel margine inferiore destro.
\textit{un primo momento realizzabili, si rivelano inesistenti di fronte ad un esame più scrupoloso.}

\textit{Qui dunque si presenta, come l’Autore stesso suggerisce, un primo campo di indagine per lo studioso, il quale, dopo essersi reso familiare con le discussioni delicate che l’argomento comporta, sarà poi messo in grado di affrontare le difficoltà più gravi della teoria generale, difficoltà che non riguardano soltanto la questione a cui sopra abbiamo alluso, ma s’incontrano ogni qual volta si vogliano mettere in chiara luce i legami fra la trattazione algebrica e la trascendente (o la topologica) delle superficie algebriche. Naturalmente per queste ultime ricerche lo studioso dovrà prima approfondire la teoria trascendente leggendo i trattati di Picard-Simart, di Lefschetz, Zariski, Hodge, la Memoria citata di Poincaré, e gli importanti lavori algebrico-trascendenti del Severi ove si trovano esposte ricerche che esorbitano dal programma del presente volume.}

\textit{Verrà presto il continuatore dell’opera delle scuole italiana e francese, il quale riesca a dare alla teoria delle superficie algebriche la perfezione che ha raggiunto la teoria delle curve algebriche? Lo spero ma ne dubito.}

\textit{A nutrire i miei dubbi m’induce l’osservazione che la matematica ha preso nel secolo attuale un indirizzo ben diverso da quello che dominava nel secolo scorso. La fantasia, la intuizione che guidavano la ricerca di allora sono oggi guardate con sospetto per il terrore degli errori a cui possono condurre. Le teorie sorgevano per rispondere al bisogno che il matematico provava di delineare e precisare degli oggetti del pensiero che erano già, in forma vaga, presenti alla sua mente. Era l’esplorazione di un ampio territorio intravisto da una cima lontana. Si costruirono così nel secolo scorso quei gioielli che si chiamano teoria delle funzioni analitiche, delle funzioni ellittiche, abeliane, superficie ad area minima, superficie cubiche.... Oggi più che il terreno da esplorare interessa la via che vi conduce, e questa via ora vien seminata di ostacoli artificiali, ora si libra tra le nuvole.}

\textit{Che questa tendenza non sia una manifestazione di breve durata si è indotti a credere dal paragone col fenomeno analogo che si verifica nelle arti plastiche e nella musica. Anche qui la fantasia è bandita quale sopravvivenza dell’epoca romantica, anche qui il contenuto dell’opera d’arte interessa meno della tecnica o dei mezzi d’espressione che l’artista impiega. Vari critici illustri pensano che ciò rappresenti una decadenza dell’arte e ne traggono foschi presagi sull’avvenire della nostra civiltà.}

\textit{Sarebbe temerario estendere questi giudizi pessimisti all’evoluzione che va subendo la matematica. Tuttavia quando si paragoni il cinquantennio che sta per chiudersi col periodo corrispondente del secolo scorso}
%CONT>>
% !PRINTED VIII
\weberpage{VIII}
% head:
%<<CONT
% scan: s004L.
% stamp: Macchie, puntinatura e bordo scuro della scansione.
\textit{in cui fiorivano nomi quali Gauss, Abel, Jacobi, Cauchy e tanti altri grandissimi, non si possono celare preoccupazioni sul futuro della nostra scienza.}

\textit{Un giorno però, prossimo o lontano, rinascerà l’amore per le grandi teorie di cui i nostri maestri del secolo XIX gettarono le basi. E quel giorno il trattato di Federigo Enriques sarà letto e meditato come il resoconto di un’esplorazione in un territorio dove molte gemme sono già state raccolte e molte altre attendono chi sia degno di scoprirle.}

\begin{flushright}
\textsc{G. Castelnuovo}
\end{flushright}

Roma, Gennaio 1949.

\clearpage\thispagestyle{plain}
% !PRINTED
\weberpage{[s004R]}
% head:
% scan: s004R; nessun numero di pagina stampato.
% stamp: Puntinatura e macchie della scansione.
\edfrontheading{\gesperrt{\textit{PREFAZIONE}}}{PREFAZIONE}{front.autore}

\textit{In questo trattato ci proponiamo di esporre la teoria delle superficie algebriche, particolarmente riguardata nel suo aspetto algebrico-geometrico, quale si è venuta maturando, durante un cinquantennio, in ispecie nella scuola geometrica italiana. Per lunghi anni abbiamo elaborato tale esposizione, rivedendo e talora rifacendo, o almeno riordinando, le dimostrazioni più antiche, in guisa da conferire alla teoria-stessa l’assetto più rigoroso e più semplice.} % ?? uncertain: segno dopo «teoria» -- trattino tipografico / artefatto.
% EDITORIAL D0002 RISOLTO: La scansione mostra un trattino dopo «teoria» al termine della riga. Si conserva la grafia stampata; la proposta [sic] resta separata. % [sic] teoria-stessa -> teoria stessa, se il trattino è tipografico, unisce due parole intere.

\textit{Quest’opera di revisione e di ordinamento si è svolta anzitutto attraverso i corsi delle lezioni da noi tenute nella Università di Roma, che sono state raccolte, in una prima redazione, da Luigi Campedelli, oggi professore all’Università di Firenze, cui ci piace attestare ancora una volta la nostra gratitudine. Le lezioni di Enriques-Campedelli furono pubblicate in due volumi; la prima parte, col titolo « Lezioni sulla teoria delle superficie algebriche, raccolte dal Dott. Luigi Campedelli », in edizione litografica presso la Casa Cedam di Padova, nel \textup{1932} e la seconda parte nel \textup{1934}, col titolo « Sulla classificazione delle superficie algebriche, particolarmente di genere zero, lezioni raccolte dal Dott. Luigi Campedelli », nei Rendiconti del Seminario Matematico della R. Università di Roma.} % ?? uncertain: piccolo segno fra «delle» e «lezioni» in «i corsi delle lezioni» -- macchia / punto stampato.

\textit{Questi volumi rimangono anche oggi il fondamento del nostro lavoro di sistemazione della teoria che tuttavia si è proseguito di poi, sia dalla cattedra universitaria, sia nel riposo degli ultimi anni, dopo l’ottobre \textup{1938}.}

\textit{Abbiamo ripreso l’esposizione del Campedelli, correggendo talvolta o semplificando qualche punto, ma soprattutto dando alla teoria stessa una più ampia estensione con l’usufruire dei contributi recati alla scienza da alcuni giovani geometri. Fra questi ci è caro ricordare i nomi di altri nostri discepoli che, dopo il Collega fiorentino, ci hanno aiutato in questa ricostruzione, approfondendo alcuni punti o invitandoci a rimeditarli colla discussione delle nostre idee e delle nostre spiegazioni: sono il Prof. Giuseppe Pompilj --- presto partito per le armi --- e il Dott. Alfredo Franchetta (due giovani speranze della nostra scienza), che si troveranno citati nel corso di queste Lezioni; ai quali vogliamo attestare pubblicamente il nostro animo grato.}

% !PRINTED X
\weberpage{X}
% head: PREFAZIONE
% scan: s005L.
% stamp: Piccolo segno ricurvo nel margine sinistro; macchie, puntinatura e bordo scuro della scansione.
\textit{Non abbiamo da spiegare più lungamente il disegno dell’opera, che già appare dall’indice generale della materia. Ma c’importa rilevare che abbiamo cercato, non soltanto di dare alle teorie esposte un assetto logico, sì anche (come in opere anteriori) di porgere una prospettiva storica del loro divenire. In tal guisa vuolsi offrire al lettore, non già il dono di qualcosa di perfetto che si lasci contemplare dal di fuori, anzi la veduta di un acquisto e di un progresso di cui deve comprendere le ragioni, e che egli è invitato a riguadagnare da sè e per sè, trovando nel libro un istrumento di lavoro. Appunto perciò il cammino induttivo della scienza, che muove da esempi e casi particolari, e talvolta anche gli errori che occorse superare e correggere, vengono messi particolarmente in rilievo.}

\textit{Aggiungiamo che per quel che concerne la conoscenza di teorie preliminari (curve algebriche, superficie razionali), ci riferiamo particolarmente all’esposizione datane in altre nostre lezioni, che furono raccolte nei quattro volumi di Enriques-Chisini « Teoria geometrica delle equazioni e delle funzioni algebriche »\origfootnote{1}{s005L-1}{\textup{Bologna, Zanichelli 1915, 1918, 1924, 1934.}} e in quello di Fabio Conforto su « Le superficie razionali », anch’esso redatto sulle nostre lezioni\origfootnote{2}{s005L-2}{\textup{Bologna, Zanichelli, 1939.}}.}

\textit{Il testo che offriamo al lettore è stato composto ed ha trovato la sua forma definitiva entro il mese di aprile del \textup{1942}, quando la persecuzione fascista vietava, nel nostro paese, anche le più innocenti pubblicazioni matematiche, e d’altronde la pubblicazione all’estero dava luogo a difficoltà che dovemmo in appresso sperimentare. In esso erano indicate alcune note, in ispecie note da noi (come corrispondente) inviate all’Accademia delle Scienze di Madrid, che non abbiamo potuto sapere se abbiano effettivamente trovato posto in quegli atti accademici.}

\textit{Oggi, usciti fuori dal periodo burrascoso della guerra, nel momento di passare il manoscritto alla tipografia, ci riesce ancor difficile di prendere esatta nozione della letteratura scientifica degli ultimi tre anni, e d’altra parte gli impegni di nuovi lavori ci vietano di ritornare sul lavoro fatto per riesaminarlo al lume di alcuni studi che, più o meno approssimativamente, siamo venuti a conoscere soltanto più tardi. Perciò conserviamo al testo delle presenti lezioni la forma datagli, come si è detto, nel \textup{1942}, salvo a completare qualche notizia bibliografica e ad indicare, con qualche nota, le lacune che richiederebbero un nuovo esame.}

\begin{flushright}
\textsc{F. Enriques}
\end{flushright}

Roma, Settembre 1945.

\clearpage\thispagestyle{plain}
% !PRINTED
\weberpage{[s005R]}
% head:
% scan: s005R; nessun numero di pagina stampato.
% stamp: Puntinatura e piccoli tratti nei margini, fra cui un tratto accanto all’indicazione di pagina 40 e uno nell’angolo inferiore destro.
\begingroup\small
\edfrontheading{\gesperrt{INDICE}}{INDICE}{front.indice}

\begin{longtable}{@{}p{0.79\linewidth}cr@{}}
\textit{Prefazione} del Prof. \textsc{G. Castelnuovo} \dotfill & Pag. & \hyperlink{old:V}{\textsc{v}} \\
\textit{Prefazione} dell’Autore \dotfill & » & \hyperlink{old:IX}{\textsc{ix}} \\
\textit{Introduzione} \dotfill & » & \hyperlink{old:1}{1} \\
\hspace*{1em}\makebox[2em][r]{1. --} Richiamo di nozioni elementari \dotfill & » & \hyperlink{old:1}{1} \\
\hspace*{1em}\makebox[2em][r]{2. --} Trasformazioni razionali \dotfill & » & \hyperlink{old:2}{2} \\
\hspace*{1em}\makebox[2em][r]{3. --} Superficie negli iperspazi: scioglimento delle singolarità & » & \hyperlink{old:5}{5} \\
\hspace*{1em}\makebox[2em][r]{4. --} Superficie dotate di singolarità normali \dotfill & » & \hyperlink{old:6}{6} \\
\hspace*{1em}\makebox[2em][r]{5. --} Nota \dotfill & » & \hyperlink{old:8}{8} \\[1ex]
\textsc{Capitolo} I. -- \textit{Sistemi lineari di curve} \dotfill & » & \hyperlink{old:15}{15} \\ % ?? uncertain: piccolo segno dopo «Capitolo» -- macchia / virgola / punto.
\hspace*{1em}\makebox[2em][r]{1. --} Fasci lineari \dotfill & » & \hyperlink{old:15}{15} \\
\hspace*{1em}\makebox[2em][r]{2. --} Fasci irrazionali \dotfill & » & \hyperlink{old:17}{17} \\
\hspace*{1em}\makebox[2em][r]{3. --} Sistemi lineari \dotfill & » & \hyperlink{old:18}{18} \\
\hspace*{1em}\makebox[2em][r]{4. --} Proprietà caratteristica dei sistemi lineari \dotfill & » & \hyperlink{old:19}{19} \\
\hspace*{1em}\makebox[2em][r]{5. --} Estensione dei teoremi di Bertini \dotfill & » & \hyperlink{old:20}{20} \\
\hspace*{1em}\makebox[2em][r]{6. --} Superficie immagini di sistemi lineari \dotfill & » & \hyperlink{old:22}{22} \\
\hspace*{1em}\makebox[2em][r]{7. --} Curve equivalenti e sistemi lineari completi \dotfill & » & \hyperlink{old:24}{24} \\
\hspace*{1em}\makebox[2em][r]{8. --} Somma e differenza di sistemi lineari: teorema del resto & » & \hyperlink{old:26}{26} \\
\hspace*{1em}\makebox[2em][r]{9. --} Superficie immagine del sistema somma \dotfill & » & \hyperlink{old:30}{30} \\
\hspace*{1em}\makebox[2em][r]{10. --} Caratteri virtuali \dotfill & » & \hyperlink{old:31}{31} \\
\hspace*{1em}\makebox[2em][r]{11. --} Curve eccezionali \dotfill & » & \hyperlink{old:33}{33} \\
\hspace*{1em}\makebox[2em][r]{12. --} Nota sulle curve eccezionali riducibili \dotfill & » & \hyperlink{old:37}{37} \\
\hspace*{1em}\makebox[2em][r]{13. --} Serie caratteristica virtuale \dotfill & » & \hyperlink{old:40}{40} \\[1ex]
\textsc{Capitolo} II. -- \textit{Sistemi covarianti e invarianti} \dotfill & » & \hyperlink{old:41}{41} \\
\hspace*{1em}\makebox[2em][r]{1. --} Curve jacobiane \dotfill & » & \hyperlink{old:41}{41} \\
\hspace*{1em}\makebox[2em][r]{2. --} Sistema jacobiano \dotfill & » & \hyperlink{old:42}{42} \\
\hspace*{1em}\makebox[2em][r]{3. --} Teorema fondamentale \dotfill & » & \hyperlink{old:48}{48} \\
\hspace*{1em}\makebox[2em][r]{4. --} Curve canoniche \dotfill & » & \hyperlink{old:50}{50} \\
\hspace*{1em}\makebox[2em][r]{5. --} Proprietà delle curve canoniche \dotfill & » & \hyperlink{old:52}{52} \\ % ?? uncertain: piccolo segno in alto a destra di «5.» -- macchia / «°».
\hspace*{1em}\makebox[2em][r]{6. --} Genere superficiale e genere lineare \dotfill & » & \hyperlink{old:53}{53} \\
\hspace*{1em}\makebox[2em][r]{7. --} Le curve eccezionali come parti fisse del sistema canonico \dotfill & » & \hyperlink{old:56}{56.} \\ % ?? uncertain: punto dopo «56» -- punto stampato / macchia.
\hspace*{1em}\makebox[2em][r]{8. --} Curve bicanoniche e pluricanoniche \dotfill & » & \hyperlink{old:60}{60} \\
\hspace*{1em}\makebox[2em][r]{9. --} Nota storica \dotfill & » & \hyperlink{old:62}{62} \\
\end{longtable}

% scan: s006L
% !PRINTED XII
\weberpage{XII}
% head: INDICE
% stamp: Puntinature sparse, macchie marginali e banda scura lungo il margine destro.
\begin{longtable}{@{}r@{\ }p{0.70\linewidth}@{\quad}c@{\ }r@{}}
\multicolumn{2}{@{}p{0.80\linewidth}}{\textsc{Capitolo III.} – \textit{Le superficie aggiunte} \dotfill } & Pag. & \hyperlink{old:67}{67} \\
1. – & Superficie d’ordine \(n-3\) aggiunte ad una superficie d’ordine \(n\) priva di punti multipli propri \dotfill & » & \hyperlink{old:67}{67} \\
2. – & Superficie aggiunte d’ordine qualsiasi \dotfill & » & \hyperlink{old:69}{69} \\
3. – & Superficie rigate \dotfill & » & \hyperlink{old:70}{70} \\
4. – & Proprietà caratteristica delle curve canoniche e delle curve aggiunte ad un sistema lineare \dotfill & » & \hyperlink{old:71}{71} \\
5. – & Punti multipli isolati \dotfill & » & \hyperlink{old:72}{72} \\
6. – & Aggiunte e subaggiunte \dotfill & » & \hyperlink{old:73}{73} \\
7. – & Punti multipli di prima specie \dotfill & » & \hyperlink{old:76}{76} \\
8. – & Piani doppi \dotfill & » & \hyperlink{old:77}{77} \\
9. – & Il genere delle curve fondamentali e l’influenza sulle aggiunte \dotfill & » & \hyperlink{old:80}{80} \\
10. – & Singolarità puntuali di specie superiore \dotfill & » & \hyperlink{old:82}{82} \\
11. – & Curve bicanoniche e superficie biaggiunte \dotfill & » & \hyperlink{old:89}{89} \\
12. – & Superficie di genere \(p=0\) e di bigenere \(P>0\) \dotfill & » & \hyperlink{old:92}{92} \\
13. – & Criteri d’equivalenza \dotfill & » & \hyperlink{old:96}{96} \\
\end{longtable}

\begin{longtable}{@{}r@{\ }p{0.70\linewidth}@{\quad}c@{\ }r@{}}
\multicolumn{2}{@{}p{0.80\linewidth}}{\textsc{Capitolo IV.} – \textit{Il genere numerico e il teorema di Riemann-Roch per le superficie} \dotfill } & » & \hyperlink{old:101}{101} \\
1. – & Introduzione \dotfill & » & \hyperlink{old:101}{101} \\
2. – & Condizioni imposte ad una superficie dal passaggio per una curva: formule di postulazione \dotfill & » & \hyperlink{old:102}{102} \\
3. – & Dimensione dei sistemi di superficie aggiunte ad una data \dotfill & » & \hyperlink{old:107}{107} \\
4. – & Il genere numerico \dotfill & » & \hyperlink{old:110}{110} \\
5. – & Complementi \dotfill & » & \hyperlink{old:114}{114} \\
6. – & Nota sulle condizioni di regolarità di una superficie . & » & \hyperlink{old:117}{117} \\
7. – & Esempi di superficie irregolari \dotfill & » & \hyperlink{old:120}{120} \\
8. – & Il teorema di Riemann-Roch per le superficie: sistemi più ampi del sistema canonico \dotfill & » & \hyperlink{old:123}{123} \\
9. – & Eliminazione delle curve eccezionali \dotfill & » & \hyperlink{old:127}{127} \\
10. – & Integrità della serie segata da un sistema lineare piccolo sulle curve di un sistema grande \dotfill & » & \hyperlink{old:128}{128} \\
11. – & Il teorema di Riemann-Roch per i sistemi lineari qualunque \dotfill & » & \hyperlink{old:131}{131} \\
12. – & Deficienza della serie caratteristica \dotfill & » & \hyperlink{old:134}{134} \\
13. – & Nota storica \dotfill & » & \hyperlink{old:137}{137} \\
14. – & Curve virtuali \dotfill & » & \hyperlink{old:142}{142} \\
15. – & Sistemi lineari regolari e sovrabbondanti \dotfill & » & \hyperlink{old:147}{147} \\
16. – & Precisazioni sulla regolarità del sistema aggiunto ad una curva sopra una superficie regolare \dotfill & » & \hyperlink{old:150}{150} \\
17. – & Segue: il sistema aggiunto ad una curva connessa formata di componenti multiple sopra una superficie regolare \dotfill & » & \hyperlink{old:152}{152} \\
18. – & Sulla riducibilità del sistema canonico \dotfill & » & \hyperlink{old:156}{156} \\
19. – & Sulla regolarità del sistema bicanonico \dotfill & » & \hyperlink{old:158}{158} \\
\end{longtable}

% scan: s006R
% !PRINTED XIII
\weberpage{XIII}
% head: INDICE
% stamp: Puntinature sparse e tracce scure ai margini della scansione.
\begin{longtable}{@{}r@{\ }p{0.70\linewidth}@{\quad}c@{\ }r@{}}
20. – & Nota sulla irriducibilità del sistema bicanonico \dotfill & Pag. & \hyperlink{old:162}{162} \\
21. – & Generi di una superficie spezzata \dotfill & » & \hyperlink{old:163}{163} \\
\end{longtable}

\begin{longtable}{@{}r@{\ }p{0.70\linewidth}@{\quad}c@{\ }r@{}}
\multicolumn{2}{@{}p{0.80\linewidth}}{\textsc{Capitolo V.} – \textit{Invarianti numerici e piani multipli} \dotfill } & » & \hyperlink{old:167}{167} \\
1. – & L’invariante di Zeuthen-Segre \dotfill & » & \hyperlink{old:167}{167} \\
2. – & Fascio irrazionale \dotfill & » & \hyperlink{old:172}{172} \\
3. – & Espressione dell’invariante di Zeuthen-Segre per mezzo dei generi \dotfill & » & \hyperlink{old:173}{173} \\
4. – & Curve cuspidate di una rete \dotfill & » & \hyperlink{old:177}{177} \\
5. – & Notizia storica \dotfill & » & \hyperlink{old:180}{180} \\
6. – & Caratteri di una rete e piani multipli \dotfill & » & \hyperlink{old:180}{180} \\
7. – & Corrispondenza \((1,n)\) fra due superficie: piani doppi . & » & \hyperlink{old:183}{183} \\
8. – & Formule di corrispondenza \dotfill & » & \hyperlink{old:185}{185} \\
9. – & Teorema d’esistenza \dotfill & » & \hyperlink{old:192}{192} \\
10. – & Complementi: forme limiti della curva di diramazione secondo Chisini \dotfill & » & \hyperlink{old:199}{199} \\
11. – & I moduli di una classe di superficie algebriche \dotfill & » & \hyperlink{old:204}{204} \\
12. – & Digressione sulla integrità della serie caratteristica di un sistema completo di curve piane dotate di nodi e di cuspidi \dotfill & » & \hyperlink{old:207}{207} \\
13. – & I moduli delle superficie regolari di genere \(p>3\) . . & » & \hyperlink{old:209}{209} \\
14. – & Nota storica e complementi \dotfill & » & \hyperlink{old:213}{213} \\
\end{longtable}

\begin{longtable}{@{}r@{\ }p{0.70\linewidth}@{\quad}c@{\ }r@{}}
\multicolumn{2}{@{}p{0.80\linewidth}}{\textsc{Capitolo VI.} – \textit{Superficie regolari: minimo dei generi e condizioni di razionalità} \dotfill } & » & \hyperlink{old:217}{217} \\
1. – & Limite inferiore del genere lineare \dotfill & » & \hyperlink{old:217}{217} \\
2. – & Il bigenere e i plurigeneri \dotfill & » & \hyperlink{old:220}{220} \\
3. – & Esempi: piani doppi di genere lineare \(p^{(1)}=1\) \dotfill & » & \hyperlink{old:224}{224} \\
4. – & Condizioni di razionalità di una superficie \dotfill & » & \hyperlink{old:230}{230} \\
\end{longtable}

\begin{longtable}{@{}r@{\ }p{0.70\linewidth}@{\quad}c@{\ }r@{}}
\multicolumn{2}{@{}p{0.80\linewidth}}{\textsc{Capitolo VII.} – \textit{Classificazione delle superficie di genere lineare \(p^{(1)}=1\)} \dotfill } & » & \hyperlink{old:237}{237} \\
1. – & Superficie di genere 0 e bigenere 1, con curva bicanonica d’ordine zero \dotfill & » & \hyperlink{old:237}{237} \\
2. – & Le superficie con tutti i generi eguali ad 1 \dotfill & » & \hyperlink{old:247}{247} \\
3. – & Notizia storica \dotfill & » & \hyperlink{old:254}{254} \\
4. – & Le superficie di genere lineare \(p^{(1)}=1\) con \(P>1\) . . & » & \hyperlink{old:257}{257} \\
\end{longtable}

\begin{longtable}{@{}r@{\ }p{0.70\linewidth}@{\quad}c@{\ }r@{}}
\multicolumn{2}{@{}p{0.80\linewidth}}{\textsc{Capitolo VIII.} – \textit{Superficie regolari canoniche e pluricanoniche}} & » & \hyperlink{old:267}{267} \\
1. – & Introduzione \dotfill & » & \hyperlink{old:267}{267} \\
2. – & Superficie canoniche di genere \(p=4\) e genere lineare \(p^{(1)}=6{,}5\) \dotfill & » & \hyperlink{old:268}{268} \\
3. – & Superficie con \(p=4\) e \(p^{(1)}=7\) \dotfill & » & \hyperlink{old:271}{271} \\
4. – & Superficie con \(p=4\) e \(p^{(1)}=8\) \dotfill & » & \hyperlink{old:273}{273} \\
5. – & Superficie con \(p=4\) e \(p^{(1)}=9\) \dotfill & » & \hyperlink{old:281}{281} \\
6. – & Superficie con \(p=4\) e \(p^{(1)}>9\) \dotfill & » & \hyperlink{old:283}{283} \\
7. – & Superficie canoniche con \(p=5\) e \(p^{(1)}=9\) \dotfill & » & \hyperlink{old:284}{284} \\
\end{longtable}

% scan: s007L
% !PRINTED XIV
\weberpage{XIV}
% head: INDICE
% stamp: Puntinature sparse e banda scura lungo il margine destro.
\begin{longtable}{@{}r@{\ }p{0.70\linewidth}@{\quad}c@{\ }r@{}}
8. – & Superficie con \(p=5\) e \(p^{(1)}=10\) \dotfill & Pag. & \hyperlink{old:286}{286} \\
9. – & Le superficie sopra una varietà a 3 o più dimensioni . & » & \hyperlink{old:289}{289} \\
10. – & Esempi di superficie canoniche iperspaziali \dotfill & » & \hyperlink{old:292}{292} \\
11. – & Minimo valore del genere lineare rispetto al genere superficiale \dotfill & » & \hyperlink{old:294}{294} \\
12. – & Limite superiore del genere lineare \(p^{(1)}\) rispetto al genere superficiale \dotfill & » & \hyperlink{old:299}{299} \\
13. – & Sistemi canonici appartenenti ad un’involuzione \dotfill & » & \hyperlink{old:300}{300} \\
14. – & Superficie di genere lineare \(p^{(1)}=2\): primo caso \(p=2\) & » & \hyperlink{old:303}{303} \\
15. – & Superficie con \(p^{(1)}=2\) e \(p=1\) \((P=3, P_3=5)\) \dotfill & » & \hyperlink{old:305}{305} \\
16. – & Superficie con \(p^{(1)}=2\) e \(p=0\) \((P=2, P_3=4)\) \dotfill & » & \hyperlink{old:308}{308} \\
17. – & Superficie di genere lineare \(p^{(1)}=3\): primo caso \(p=3\) & » & \hyperlink{old:311}{311} \\
18. – & Superficie con \(p^{(1)}=3\) e \(p=2\) \((P=5)\) \dotfill & » & \hyperlink{old:312}{312} \\
19. – & Superficie con \(p^{(1)}=3\) e \(p=1\) \((P=4)\) \dotfill & » & \hyperlink{old:316}{316} \\
20. – & Superficie con \(p=0\) e \(p^{(1)}=3\) \dotfill & » & \hyperlink{old:321}{321} \\
21. – & Superficie pluricanoniche semplici e multiple \dotfill & » & \hyperlink{old:321}{321} \\
\end{longtable}

\begin{longtable}{@{}r@{\ }p{0.70\linewidth}@{\quad}c@{\ }r@{}}
\multicolumn{2}{@{}p{0.80\linewidth}}{\textsc{Capitolo IX.} – \textit{Superficie irregolari e sistemi continui di curve disequivalenti} \dotfill } & » & \hyperlink{old:325}{325} \\
1. – & Introduzione \dotfill & » & \hyperlink{old:325}{325} \\
2. – & Condizione aritmetica perchè una curva, sopra una superficie d’irregolarità geometrica \(q\), appartenga ad una serie continua \(\infty^q\) di curve disequivalenti \dotfill & » & \hyperlink{old:326}{326} \\
3. – & La varietà di Picard corrispondente ad una superficie irregolare \dotfill & » & \hyperlink{old:328}{328} \\
4. – & Proprietà caratteristica delle superficie algebriche irregolari: teorema fondamentale per \(p_g=0\) \dotfill & » & \hyperlink{old:328}{328} \\
5. – & Teorema fondamentale per \(p_g>0\) \dotfill & » & \hyperlink{old:330}{330} \\
6. – & Storia della teoria dei sistemi continui \dotfill & » & \hyperlink{old:339}{339} \\
7. – & Su diversi tentativi di dimostrazione e d’estensione del teorema fondamentale \dotfill & » & \hyperlink{old:347}{347} \\
8. – & Il sistema paracanonico \dotfill & » & \hyperlink{old:354}{354} \\
9. – & Disgressione sulle varietà abeliane \dotfill & » & \hyperlink{old:357}{357} \\ % [sic] Disgressione -> Digressione, probabile s soprannumeraria; a p. XIII è stampato «Digressione».
10. – & Segue: il genere delle varietà abeliane \dotfill & » & \hyperlink{old:361}{361} \\
11. – & Fascio irrazionale sopra le superficie di genere geometrico nullo \dotfill & » & \hyperlink{old:365}{365} \\
12. – & Nota sulle superficie d’irregolarità 1 \dotfill & » & \hyperlink{old:369}{369} \\
\end{longtable}

\begin{longtable}{@{}r@{\ }p{0.70\linewidth}@{\quad}c@{\ }r@{}}
\multicolumn{2}{@{}p{0.80\linewidth}}{\textsc{Capitolo X.} – \textit{Le superficie di genere geometrico nullo} \dotfill } & » & \hyperlink{old:371}{371} \\
1. – & Introduzione \dotfill & » & \hyperlink{old:371}{371} \\
2. – & Le superficie contenenti un sistema di curve di genere \(\pi\) e di grado \(n>2\pi-2\) \dotfill & » & \hyperlink{old:372}{372} \\
3. – & Lemma sulle curve riducibili di un fascio \dotfill & » & \hyperlink{old:376}{376} \\
4. – & Superficie di genere \(p_g=0\) con \(p_a<-1\) \dotfill & » & \hyperlink{old:378}{378} \\
5. – & Superficie \(F\) con \(p_g=0\) e \(p_a=-1\), possedenti un fascio ellittico di curve \(K\) di genere \(\pi>1\): lemma I \dotfill & » & \hyperlink{old:379}{379} \\
6. – & Lemma II: disegno della dimostrazione \dotfill & » & \hyperlink{old:381}{381} \\
7. – & Primo caso: il genere \(\pi\) delle curve \(K\) sia pari \dotfill & » & \hyperlink{old:384}{384} \\
\end{longtable}

% scan: s007R
% !PRINTED XV
\weberpage{XV}
% head: INDICE
% stamp: Puntinature e piccoli tratti sparsi nella scansione.
\begin{longtable}{@{}r@{\ }p{0.70\linewidth}@{\quad}c@{\ }r@{}}
8. – & Secondo caso: il genere \(\pi\) delle curve \(K\) sia dispari . . & Pag. & \hyperlink{old:389}{389} \\
9. – & Lemma III \dotfill & » & \hyperlink{old:392}{392} \\
10. – & Conclusione: le superficie \(F\) posseggono anche un fascio lineare di curve ellittiche \dotfill & » & \hyperlink{old:393}{393} \\
11. – & Superficie di generi \(p_g=0\) e \(p_a=-1\) con un fascio ellittico di curve di genere \(\pi=1\) \dotfill & » & \hyperlink{old:394}{394} \\
12. – & Superficie ellittiche \dotfill & » & \hyperlink{old:398}{398} \\
13. – & Costruzione delle superficie ellittiche di genere \(p_g=0\) & » & \hyperlink{old:401}{401} \\
14. – & Superficie coi plurigeneri nulli: caratterizzazione delle rigate \dotfill & » & \hyperlink{old:405}{405} \\
15. – & Nota storica \dotfill & » & \hyperlink{old:412}{412} \\
16. – & Superficie con curve pluricanoniche d’ordine zero: tipi con curve ellittiche \(K\) di modulo generale \dotfill & » & \hyperlink{old:414}{414} \\
17. – & Segue: caso armonico \dotfill & » & \hyperlink{old:417}{417} \\
18. – & Caso equianarmonico \dotfill & » & \hyperlink{old:421}{421} \\
19. – & Riassunto \dotfill & » & \hyperlink{old:424}{424} \\
\end{longtable}

\begin{longtable}{@{}r@{\ }p{0.70\linewidth}@{\quad}c@{\ }r@{}}
\multicolumn{2}{@{}p{0.80\linewidth}}{\textsc{Capitolo XI.} – \textit{Classificazione generale delle superficie} \dotfill } & » & \hyperlink{old:427}{427} \\
1. – & Introduzione \dotfill & » & \hyperlink{old:427}{427} \\
2. – & Superficie che ammettono una serie continua di trasformazioni birazionali in se stesse: casi che conducono alle rigate \dotfill & » & \hyperlink{old:429}{429} \\
3. – & Superficie ellittiche e iperellittiche \dotfill & » & \hyperlink{old:433}{433} \\
4. – & Caratterizzazione delle superficie ellittiche e iperellittiche mediante i valori dei generi \dotfill & » & \hyperlink{old:438}{438} \\
5. – & Superficie coi generi \(p_g=1\) e \(p^{(1)}=1\) \dotfill & » & \hyperlink{old:447}{447} \\
6. – & Nota sulla teoria geometrica delle superficie iperellittiche \dotfill & » & \hyperlink{old:450}{450} \\
7. – & Classificazione generale delle superficie algebriche . . & » & \hyperlink{old:455}{455} \\
8. – & Superficie di genere numerico negativo \dotfill & » & \hyperlink{old:459}{459} \\
9. – & Riassunto della classificazione precisata \dotfill & » & \hyperlink{old:463}{463} \\
\end{longtable}
% !PRINTED
\weberblank{[s008L]}
% head:
% scan: s008L; pagina bianca, senza numero di pagina né testatina stampati.
% stamp: Puntinature sparse e banda scura lungo il margine destro.


\endgroup

\clearpage\mainmatter\thispagestyle{plain}
% scan: s008R
% !PRINTED 1
\weberpage{1}
% head:
% stamp: Puntinature, brevi tratti sparsi e banda scura lungo il margine sinistro.
% Original footnote labels are set explicitly by the edition macros.
% EDITORIAL: numero originale fissato esplicitamente.

\edfrontheading{INTRODUZIONE}{INTRODUZIONE}{intro}

\edsection{\textbf{1. Richiamo di nozioni elementari.}}{1. Richiamo di nozioni elementari.}{sec.intro.1}

In queste lezioni ci riferiamo, in generale, ad una superficie algebrica dello spazio ordinario, supponendola irriducibile. La superficie \(f\) verrà definita da un’equazione in coordinate cartesiane
\[
f(x,y,z)=0
\]
ovvero in coordinate omogenee
\[
f(x_1,x_2,x_3,x_4)=0:
\]
\(f\) designa nel primo caso un polinomio nelle variabili \(x,y,z\), e nel secondo caso un polinomio omogeneo ovvero una forma.

Assumiamo come noti gli elementi della teoria proiettiva delle superficie, che includono le nozioni di ordine, classe, superficie polari, ecc.\origfootnote{1}{s008R-1}{Cfr. \textsc{Enriques-Chisini}, \textit{Teoria geometrica delle equazioni e delle funzioni algebriche.} Libro III, vol. II. Bologna, Zanichelli, 1918.}. E richiamiamo alcune proprietà di cui si farà uso nel seguito.

Un \textit{punto} \(O\) della superficie \(f\), che sia d’ordine \(n\), è \textit{semplice} per essa se le rette uscenti da \(O\) segano la superficie in \(n-1\) punti fuori di \(O\). L’intorno di un punto semplice della superficie si può rappresentare in modo semplice su un piano in un ordine di approssimazione alto come si vuole, cioè dove si tenga conto degli infinitesimi di ordine \(1,2,\ldots\), ecc. In primo luogo è chiaro che nel primo ordine di approssimazione la superficie può essere sostituita dal piano ivi tangente; ma se si vuol tenere conto del secondo ordine, supponendo gli assi cartesiani scelti in modo opportuno, si sostituirà invece col paraboloide osculatore dello stesso ordine \(z=\varphi_2(xy)\), che si lascia proiettare univocamente sul piano \((x,y)\) dal punto all’infinito dell’asse \(z\). E così si dica successivamente per l’intorno del terzo ordine ecc. Pertanto, essendo data sulla superficie una curva che abbia in \(O\) una singolarità qualsiasi, si potrà
% footer: segnatura tipografica.
% EDITORIAL: segnatura originale: \vadjust{\vbox{\hbox{\footnotesize \textsc{Enriques F.} - \textit{Superficie algebriche.}}}}
\printersignature{\textsc{Enriques F.} - \textit{Superficie algebriche.}}
%CONT>>
% scan: s009L
% !PRINTED 2
\weberpage{2}
% head: INTRODUZIONE
%<<CONT
% stamp: Puntinature, piccoli tratti marginali e banda scura lungo il margine destro.
% EDITORIAL: numero originale fissato esplicitamente.

ritenere questa singolarità formata di punti multipli infinitamente vicini, proprio come accade nel piano\origfootnote{1}{s009L-1}{Cfr. \textsc{Enriques-Chisini}, Op. cit., Libro IV, vol. II.}.

I punti singolari o multipli per la superficie \(f\) sono definiti come segue: un \textit{punto \(O\) multiplo d’ordine \(r\)} \((r>1)\) è tale che le rette per \(O\) segano la superficie in \(n-r\) punti residui. Vi è luogo a distinguere le \textit{curve multiple} e i \textit{punti multipli isolati}, senza escludere la considerazione di punti multipli notevoli appartenenti ad una curva multipla, che si presentano come una sovrapposizione dei due casi, e su cui avremo luogo di dare nel seguito qualche delucidazione. Una curva multipla d’ordine \(r\) si può ritenere in generale come l’intersezione di \(r\) superficie o \textit{falde} differenzialmente distinte in un punto generico della curva, in corrispondenza agli \(r\) piani tangenti in esso. Invece un punto multiplo isolato appare di regola conico, dove le rette aventi \(r+1\) intersezioni riunite con la superficie, formano un cono d’ordine \(r\), generalmente irriducibile.

Ma si possono avere delle particolarizzazioni. Così due o più falde della superficie lungo una curva multipla possono toccarsi o anche fondersi in una \textit{falda d’ordine superiore}: il più semplice esempio è offerto dalla curva doppia cuspidale. D’altra parte il cono tangente in un punto \(r\)-plo isolato può spezzarsi o anche ridursi ad un piano contato \(r\) volte: l’esempio più semplice è offerto dal punto doppio uniplanare.

Per dominare tutte le complicazioni cui possono dar luogo le singolarità della superficie, conviene definire i \textit{punti multipli infinitamente vicini} e le \textit{curve} (\textit{proprie} o \textit{infinitesime}) da questi costituite, come si fa per mezzo dei rami di curve uscenti da un punto\origfootnote{2}{s009L-2}{Cfr. \textsc{Enriques-Chisini}, Op. cit. Libro IV, cap. IV.}. In quest’ordine di idee un punto uniplanare appare come un punto doppio cui sono infinitamente vicini tre punti doppi in direzioni distinte. Il più semplice esempio di una curva (retta) doppia infinitesima, infinitamente vicina ad un punto multiplo isolato, è il \textit{tacnodo}: cioè un punto doppio \(O\) tale che le sezioni piane per \(O\) hanno accanto ad \(O\) un altro punto doppio infinitamente vicino. Ma noi non vogliamo indugiarci su questo argomento, per cui rimandiamo alle citate lezioni di \textsc{Enriques-Chisini}. % ?? uncertain: carattere della «o» in «proprie o infinitesime» -- «o» tonda / «o» corsiva.
% EDITORIAL D0007 RISOLTO: La «o» è in tondo fra le due parole corsive, come già trascritto.

\edsection{\textbf{2. Trasformazioni razionali.}}{2. Trasformazioni razionali.}{sec.intro.2}

Si definisce una trasformazione razionale dello spazio cui appartiene la superficie \(f\) ponendo le coordinate \(X,Y,Z\), funzioni razionali delle \(x,y,z\):
\weq{(1)}{X=\frac{\varphi_1(xyz)}{\varphi_0(xyz)}\qquad Y=\frac{\varphi_2(xyz)}{\varphi_0(xyz)}\qquad Z=\frac{\varphi_3(xyz)}{\varphi_0(xyz)},}
%CONT>>
% scan: s009R
% !PRINTED 3
\weberpage{3}
% head: INTRODUZIONE
%<<CONT
% stamp: Puntinature e piccoli tratti sparsi nella scansione.
% EDITORIAL: numero originale fissato esplicitamente.

dove le \(\varphi\) sono dei polinomi che possiamo supporre dello stesso ordine. Per mezzo di questa trasformazione, che supponiamo non degenere, ad ogni punto generico dello spazio \((xyz)\) corrisponde un punto dello spazio \((XYZ)\), e viceversa ad un punto \((XYZ)\) corrisponde in generale un certo numero finito \(n\) \((\geq 1)\) di punti \((xyz)\), i quali si ottengono risolvendo le equazioni (1), in cui le \(XYZ\) vengano assunte come date e le \(xyz\) come incognite: per \(n>1\) i gruppi di punti così ottenuti formeranno, nello spazio \((xyz)\), una involuzione \(I_n\) d’ordine \(n\), cioè una serie di gruppi di \(n\) punti tale che ogni punto appartenga in generale ad un gruppo.

Ora la trasformazione (1) porterà generalmente gli \(\infty^2\) punti della superficie \(f(x,y,z)=0\) in \(\infty^2\) punti distinti, costituenti una superficie \(F(X,Y,Z)=0\) trasformata della \(f\). Ad un punto generico di \(f\) corrisponderà dunque \textit{un} punto di \(F\), e viceversa ad un punto generico di \(F\) corrisponderà di regola \textit{un} punto di \(f\), la trasformazione riuscendo univocamente invertibile per i punti della superficie \(F\)\origfootnote{1}{s009R-1}{Di qui la distinzione fra trasformazione birazionale intercedente fra due superficie e trasformazione birazionale o cremoniana, estendibile allo spazio che le contiene.}. Questa invertibilità e le circostanze particolari in cui essa cessa di sussistere, esigono una spiegazione. Invero abbiamo visto che la trasformazione (1), considerata fra i due spazi \((xyz)\) \((XYZ)\), non è in generale univocamente invertibile, anzi conduce dai punti \((XYZ)\) ai gruppi di \(n\) punti di una involuzione \(I_n\) nello spazio \((xyz)\). Ma se la \(f\) viene scelta in modo generale entro questo spazio, essa non apparterrà all’involuzione \(I_n\), cosicchè ad un punto di \(f\) saranno coniugati, in \(I_n\), \(n-1\) punti fuori della superficie, e quindi ad un punto generico di \(F\) verranno a corrispondere per la (1) un punto di \(f\) ed \(n-1\) punti fuori di essa, descrivendo la superficie ad essa coniugata rispetto ad \(I_n\).

Soltanto nel caso che appartenga all’involuzione \(I_n\), ovvero ad un’involuzione parzialmente contenuta in essa, la trasformata di \(F\) si ridurrà ad una superficie multipla, in corrispondenza \([1,n]\) o \([1,m]\) \((m<n)\) con \(f\). % [sic] F -> f, la superficie di partenza è indicata con f nelle righe precedenti.

Ora, riferendoci al caso generale, in cui \(f\) e \(F\) sieno in corrispondenza biunivoca, conviene pur rilevare che la trasformazione (1) ammette, in generale, delle eccezioni\origfootnote{2}{s009R-2}{Si tengano presenti le considerazioni analoghe relative alla trasformazione delle curve: Cfr. \textsc{Enriques-Chisini}, Op. cit., Libro V, vol. III.}. Invero le formule della trasformazione diventano indeterminate in tutti quei punti \((xyz)\) per cui si annullino insieme le \(\varphi\) (cioè per i punti base del sistema trasformante individuato dalle superficie \(\varphi_0=0,\ \varphi_1=0,\ \varphi_2=0,\ \varphi_3=0\)). Una eccezione di questo genere può aver luogo anzitutto
%CONT>>
% scan: s010L
% !PRINTED 4
\weberpage{4}
% head: INTRODUZIONE
%<<CONT
% stamp: Puntinature, brevi tratti marginali e banda scura lungo il margine destro.
% EDITORIAL: numero originale fissato esplicitamente.

in un punto semplice di \(f\); allora questo si trasforma in generale in una retta o in una curva di \(F\), curva trasformata di un punto semplice, che si chiamerà \textit{curva eccezionale} della superficie \(F\).

Invece può darsi che un punto base del sistema trasformante delle \(\varphi\) cada in un punto singolare di \(f\). E se si ha una curva multipla di \(f\) su cui si annullino, contemporaneamente, le \(\varphi\), questa si trasformerà generalmente in una curva semplice, ad ogni punto \(r\)-plo della curva corrispondendo \(r\) punti distinti, in relazione alle \(r\) falde di \(f\) che passano per esso. Se invece un punto base delle \(\varphi\) cade in un punto multiplo isolato \(O\) della \(f\), questo si trasforma generalmente in una curva i cui punti rispondono ai punti infinitamente vicini ad \(O\).

Queste affermazioni cadono in difetto quando intervengono punti e curve multipli infinitamente vicini ai punti singolari considerati. In questi casi l’effetto generale della trasformazione è di trasformare in punti multipli propri i punti multipli infinitamente vicini del primo ordine, e in generale di semplificare la singolarità, riducendo da \(r\) ad \(r-1\) l’intorno dei punti considerati, almeno quando si supponga che le superficie \(\varphi\), passanti per un punto \(O\), non abbiano alcuna tangente fissa in comune.

Tuttavia la trasformazione che riesce a semplificare le singolarità di una superficie introduce in generale delle nuove singolarità, perchè vi saranno in generale \(\infty^1\) coppie di punti della \(f\) dove le \(X,Y,Z\) riprenderanno lo stesso valore, sicchè ad essi risponderanno punti di una curva doppia di \(F\).

Ciò risulta, in sostanza, da un facile computo di costanti. Se le \(X,Y,Z\), debbono riprendere lo stesso valore nei punti \((x,y,z)\) e \((x',y',z')\), dovranno sussistere le 5 equazioni in 6 incognite:
\[
\begin{aligned}
&\varphi_1(xyz)\varphi_0(x'y'z')-\varphi_1(x'y'z')\varphi_0(xyz)=0\\
&\varphi_2(xyz)\varphi_0(x'y'z')-\varphi_2(x'y'z')\varphi_0(xyz)=0\\
&\varphi_3(xyz)\varphi_0(x'y'z')-\varphi_3(x'y'z')\varphi_0(xyz)=0\\
&f(xyz)=0\\
&f(x'y'z')=0,
\end{aligned}
\]
le quali ammetteranno in generale \(\infty^1\) soluzioni \((xyzx'y'z')\), fuori delle \(\infty^2\) soluzioni triviali
\[
x=x'\qquad y=y'\qquad z=z'\text{\origfootmark{1}{s010L-1}}.
\]\origfoottext{1}{s010L-1}{Le prime tre equazioni definiscono tre ipersuperficie dello spazio \(S_6\), che hanno a comune lo \(S_3\) \((x=x',y=y',z=z')\) e, fuori di questo, una certa varietà a tre dimensioni \(V\); la \(V\) viene intersecata dalle varietà cilindriche \(f(xyz)=0\) e \(f(x'y'z')=0\) secondo una curva.}

In casi particolari si potranno avere \(\infty^1\) gruppi di più che due punti di \(f\) cui rispondono i punti d’una curva multipla di \(F\). Ma
%CONT>>
% scan: s010R
% !PRINTED 5
\weberpage{5}
% head: INTRODUZIONE
%<<CONT
% stamp: Puntinature, piccoli tratti sparsi e tracce scure ai margini della scansione.
% EDITORIAL: numero originale fissato esplicitamente.

potranno anche aversi curve di \(f\) su cui le \(XYZ\) restino costanti, cui rispondono punti multipli isolati di \(F\). E col complicarsi delle singolarità di \(f\), e con l’intervento di altre circostanze particolari, non si può escludere la genesi di singolarità più elevate per la \(F\).

\edsection{\textbf{3. Superficie negli iperspazi: scioglimento delle singolarità.}}{3. Superficie negli iperspazi: scioglimento delle singolarità.}{sec.intro.3}

La trasformazione dello spazio ordinario che sopra abbiamo considerato si può estendere definendo una trasformazione che porti la superficie \(f\) in uno spazio \(S'_r\) a più dimensioni: % ?? uncertain: piccolo segno in alto a destra del primo S_r -- S'_r / S_r con puntinatura.
% EDITORIAL D0008 RISOLTO: Il primo simbolo mostra un apice a destra di S, distinto dal pedice r. Nessun cambiamento.
\[
X_1=\frac{\varphi_1(xyz)}{\varphi_0(xyz)}\qquad X_2=\frac{\varphi_2(xyz)}{\varphi_0(xyz)}\ldots\qquad X_r=\frac{\varphi_r(xyz)}{\varphi_0(xyz)};
\]
la superficie \(f(xyz)=0\) si trasformerà in una superficie \(F\) dello spazio \(S_r\) quale è definita dalle formule scritte, e che d’altra parte può anche definirsi in vari modi come intersezione parziale di varietà dell’iperspazio\origfootnote{1}{s010R-1}{Cfr. \textsc{Enriques-Chisini}, Op., cit., Libro I, \S\ 23.}. % ?? uncertain: segno dopo «Op.» nella nota -- «Op., cit.,» / «Op. cit.,» con puntinatura.

Anche per questa trasformazione si ripetono in gran parte le cose dette per una trasformazione dello spazio \(S_3\); che in generale la trasformazione riesce biunivoca per i punti delle due superficie (e non per gli spazi che le contengono), e che una curva multipla ordinaria di \(f\), che sia base per il sistema trasformante delle superficie \(\varphi\), si scioglie in generale diventando luogo degli \(r\) punti semplici che rispondono ai punti \(r\)-pli di \(f\) che la costituiscono; e similmente si dica dei punti multipli isolati che si mutano in curve, ecc.

Ma qui, almeno per \(r\geq5\), non si producono in generale nuovi punti singolari per la \(F\), chè appunto per \(r\geq5\) non esistono in generale coppie di punti di \(f\) ove le \(X_1,X_2,\ldots,X_r\) riprendano lo stesso valore, mentre per \(r=4\) si ha in generale un numero finito di tali coppie che dà luogo ad altrettanti punti doppi per la \(F\).

S’intuisce per tal modo che si riuscirà sempre a sciogliere le singolarità di una superficie \(f\) e a trasformarla in un’altra dello spazio \(S_r\) affatto priva di singolarità, cioè dotata soltanto di punti semplici. Ma questa riduzione che riesce facile per le superficie dotate soltanto di curve e di punti multipli ordinari (curve multiple a falde distinte, e punti multipli isolati a cono osculatore affatto generale) esige un esame più minuto quando la \(f\) possegga singolarità infinitamente vicine, comunque complicate. L’analisi che occorre a tal uopo si può decomporre tuttavia in due parti: un’analisi nel senso della geometria differenziale, dove si tratta di sciogliere le singolarità definite nell’intorno di un punto \(O\) di una curva dello spazio ordinario ed in cui è lecito giovarsi di trasformazioni birazionali o cremoniane
%CONT>>
% scan: s011L
% !PRINTED 6
\weberpage{6}
% head: INTRODUZIONE
%<<CONT
% stamp: Fascia scura verticale sul margine destro e piccoli segni sparsi della scansione.
dell’intero spazio $S_3$; e successivamente un esame critico delle nuove singolarità che possono venir prodotte dalla trasformazione stessa. Ma quest’ultimo passo, che importa difficoltà minuziose quando si tratti della riduzione delle singolarità con trasformazioni birazionali dello $S_3$\origfootnote{1}{s011L-1}{Problema risoluto da \textsc{O. Chisini}, in Memorie Acc. Bologna, 1921.}, si può evitare mercè una semplice osservazione di \textsc{B. Levi}: infatti, se la superficie $f(xyz) = 0$ viene trasformata in un’altra $F(XYZ) = 0$, la superficie dello spazio a 6 dimensioni luogo dei punti $(x, y, z, X, Y, Z)$ possiederà soltanto delle singolarità che rispondono alle singolarità comuni alle $f$ e $F$.

Non è nostro proposito fermarci qui a spiegare la riduzione delle singolarità delle superficie quale si è ottenuta rigorosamente anzitutto da \textsc{B. Levi}, e ci limiteremo ad enunciare il risultato fondamentale cui si perviene, rimandando per la dimostrazione alle principali memorie che la contengono e la svolgono in differenti maniere\origfootnote{2}{s011L-2}{\textsc{B. Levi}, \textit{Risoluzione delle singolarità puntuali delle superficie algebriche.} Atti Acc. Torino, 1897 (cfr. Ann. di Mat. 1897).

\textsc{O. Chisini}, \textit{La risoluzione delle singolarità delle superficie mediante trasformazioni birazionali dello spazio.} Memorie Acc. Bologna, 1921.

\textsc{G. Albanese}, \textit{Trasformazione birazionale di una superficie algebrica qualunque in un’altra priva di punti multipli.} Circ. Mat. Palermo, 1924.

\textsc{R. J. Walker}, \textit{Reduction of singularities of an algebraic surface.} Annals of Math., 1935.

\textsc{O. Zariski}, \textit{The reduction of the singularities of an algebraic surface.} Ann. of Math., 1939. Cfr. Annals of Math., 1941.}.

\textit{Una superficie $f$ dotata di singolarità qualsiasi si può sempre trasformare birazionalmente in un’altra priva di singolarità, appartenente ad uno spazio a cinque dimensioni.} % ?? uncertain: «qualsiasi» nel teorema -- qualsiasi / qualsiansi.
% EDITORIAL D0010 RISOLTO: La parola nel teorema è «qualsiasi», senza la n; diversa dalla forma «qualsiansi» nel paragrafo seguente.

\edsection{4. Superficie dotate di singolarità normali.}{4. Superficie dotate di singolarità normali.}{sec.intro.4}

Data una superficie con singolarità qualsiansi si cominci anzitutto a trasformarla in una $F$ dello spazio $S_5$ affatto priva di punti multipli. La $F$, proiettata successivamente da due punti sullo spazio ordinario, ci darà in questo un’immagine $\Phi$ della superficie data. Si tratta di esaminare quali saranno le singolarità di $\Phi$ quando la proiezione di $F$ sia fatta da due punti generici. Vogliamo dimostrare che la $\Phi$ possiederà soltanto una curva doppia nodale (piani tangenti generalmente distinti), dotata di punti tripli che sono pure tripli per la superficie.

È chiaro anzitutto che questo è il tipo delle \textit{singolarità normali}, che avrà in generale una superficie dello spazio ordinario, ottenuta come proiezione di una superficie dello spazio a più dimensioni, priva di punti singolari; ciò risulta da un semplice computo di co%
%CONT>>
% scan: s011R
% !PRINTED 7
\unskip\weberpage{7}%
% head: INTRODUZIONE
%<<CONT
% stamp: Piccoli segni e macchie sparsi della scansione.
\unskip stanti. Si osservi infatti che le corde di una superficie dello spazio $S_r$ sono sempre $\infty^4$ e quindi per $r = 5$ vi è un numero finito di corde passanti per un punto, mentre per $r = 4$ ve ne sono $\infty^1$. Quindi, proiettando la $F$ da un punto generico dello $S_5$, la superficie proiezione $F'$ acquisterà, in corrispondenza delle corde della $F$ passanti per il centro, un numero finito di punti doppi; proiettando poi la $F'$ da un punto generico dello $S_4$, la superficie proiezione acquisterà una curva doppia (apparente) sulla quale andranno a cadere i punti doppi della $F'$. In altre parole, per due punti dello $S_5$, o per la retta che li congiunge, passerà una serie $\infty^1$ di piani che incontrano la $F$ in due punti e un numero finito di piani che la incontrano in tre punti, sicchè la proiezione di $F$, fatta da quei centri, possiederà appunto una curva doppia nodale dotata di punti tripli, tripli insieme per la curva doppia e per la superficie.

Tuttavia si può dubitare che, per una particolare superficie dello $S_5$, sia pure priva di punti multipli, la proiezione fatta da due punti generici dello $S_5$ venga ad acquistare singolarità più elevate, cioè: una curva multipla anzichè doppia, ovvero dei punti multipli più che tripli per la curva doppia. Per escludere il primo dubbio basta notare che le $\infty^4$ corde di una superficie di uno spazio qualunque a più di tre dimensioni non possono essere trisecanti, poichè altrimenti anche le corde di una curva gobba sezione iperpiana della superficie sarebbero trisecanti, il che è impossibile\origfootnote{1}{s011R-1}{Cfr. \textsc{Enriques-Chisini}, Op. cit., Libro III, cap. IV, \S~43, vol. II, pag. 289.}. Da ciò segue che per un punto generico dello $S_5$ non possono passare trisecanti della $F$, e che quando la $F$ sia proiettata da un punto generico in una $F'$ dello $S_4$, per un punto generico di questo non può passare che un numero finito di trisecanti.

Resta da dimostrare che scegliendo in modo generico i centri di proiezione nello $S_5$, non può accadere che i piani trisecanti per i detti centri incontrino la superficie $F$ in più che tre punti e quindi la curva doppia della proiezione $\Phi$ possegga, in luogo di punti tripli, punti di molteplicità più elevata\origfootnote{2}{s011R-2}{Cfr. \textsc{F. Enriques}, \textit{Sulle singolarità che nascono per proiezione di una superficie o varietà algebrica}, in Scritti matematici offerti a \textsc{Luigi Berzolari}, Pavia, 1936, pag. 351.}. Perciò basta escludere che tutti i piani cui s’imponga d’incontrare la superficie in tre punti vengano a incontrarla di conseguenza in quattro o più punti. Infatti, se ogni piano trisecante la $F$ risulti quadrisecante, lo stesso accadrà per ogni curva gobba sezione iperpiana di $F$ (la quale non può appartenere ad un $S_3$) e quindi anche la proiezione di questa fatta da un suo punto sullo $S_3$ dovrà dare una curva di cui ogni corda risulta trisecante, ciò che abbiamo già ricordato esser impossibile.

% scan: s012L
% !PRINTED 8
\weberpage{8}
% head: INTRODUZIONE
% stamp: Fascia scura verticale sul margine destro e piccoli segni sparsi della scansione.
In tal guisa viene rigorosamente dimostrato un teorema fondamentale già riconosciuto dal \textsc{Noether} (1888):

\textit{Una superficie con singolarità qualsiansi può sempre trasformarsi in un’altra dello spazio ordinario dotata soltanto di singolarità normali, cioè di una curva doppia nodale che possiede in generale un numero finito di punti tripli, i quali sono pure tripli per la superficie.} % ?? uncertain: «qualsiansi» nel teorema -- qualsiansi / qualsiasi.
% EDITORIAL D0011 RISOLTO: La scansione mostra la forma arcaica «qualsiansi», che si conserva.

La detta curva doppia potrà possedere anche un numero finito di punti doppi (\textit{incroci}) provenienti dalle corde della superficie di $S_4$ proiettata sullo $S_3$ che siano doppie per il cono delle corde proiettanti.

Giova anche ricordare che sulla detta \textit{curva nodale} vi sarà in generale un numero finito di \textit{punti cuspidali} (pinch-points) nei quali le due falde lineari della superficie si fondono in una falda del secondo ordine, e che riescono punti base semplici per le curve intersezioni variabili delle superficie polari.

Nel seguito, dovendo svolgere una teoria delle superficie in cui si ricercano le proprietà invarianti per trasformazioni birazionali, potremo sempre riferirci ad una superficie modello priva di singolarità dello spazio a 5 dimensioni, ovvero ad una superficie dello spazio ordinario dotata delle singolarità normali, di cui nel precedente enunciato. In quest’ultimo caso un punto biplanare della curva doppia dovrà ritenersi come la sovrapposizione di due punti semplici, i cui intorni corrispondono alle due falde, così come nella teoria invariantiva delle curve un nodo viene considerato come la sovrapposizione di due punti appartenenti ai due rami.

\edsection{5. Nota.}{5. Nota.}{sec.intro.5}

Ciò che si è detto in ordine alla trasformazione di una superficie in un’altra dello spazio ordinario, dotata di singolarità normali, può precisarsi nel senso che: è lecito supporre che la \textit{trasformata di una qualsiasi superficie sia dotata di curva doppia irriducibile}, la quale possegga, come si è detto, un certo numero di punti tripli, tripli egualmente per la superficie e per la sua curva nodale.

Esporremo la dimostrazione di questa proprietà\origfootnote{1}{s012L-1}{Cfr. \textsc{A. Franchetta}, \textit{Sulla curva doppia della proiezione di una superficie generale dello $S_4$ da un punto generico su uno $S_3$.} Rendic. R. Acc. d’Italia, 1941, v. anche i Rend. dell’Acc. dei Lincei, s. VIII, vol. II, fasc. 3, 1947.} ammettendo noti i principii della teoria dei sistemi lineari, che vengono esposti nei primi paragrafi di queste lezioni.

A tal uopo conviene partire da una superficie priva di singolarità nello spazio $S_5$, o meglio da una $F$ che sia proiezione generica di questa nello $S_4$, e perciò dotata soltanto di punti doppi impropri
%CONT>>
% scan: s012R
% !PRINTED 9
\weberpage{9}
% head: INTRODUZIONE
%<<CONT
% stamp: Fascia scura verticale sul margine sinistro e piccoli segni sparsi della scansione.
\textit{ordinarii}, ciascuno dei quali può ritenersi come una singolarità apparente, rappresentando una coppia di punti sovrapposti \textit{distinti}. Conviene esaminare le singolarità di una superficie $\Phi$ dello spazio ordinario, proiezione della $F$ da un punto generico $O$ dello $S_4$: in quali casi possa accadere che la curva doppia di $\Phi$ riesca sempre riducibile al variare del centro di proiezione $O$.

Questo esame conduce a stabilire che: se una superficie $F$, dotata di punti doppi impropri ordinarii nello $S_4$, viene proiettata da un punto generico in una $\Phi$ dello $S_3$, la cui curva doppia sia riducibile, la $F$ è una proiezione della superficie di Veronese (superficie del quarto ordine di $S_5$ rappresentata sul piano del sistema $\infty^5$ delle coniche).

Per giungere a questo risultato conviene premettere il seguente:

\textit{Lemma}: non esistono superficie dello $S_4$ possedenti $\infty^3$ coppie di piani tangenti, incidenti secondo una retta, in modo che le rette che congiungono i loro punti di contatto riempiano tutto lo spazio.

Sia $P$ un punto generico di $F$ e $\pi$ il piano ivi tangente. Per ipotesi esistono $\infty^1$ piani tangenti ad $F$ ed incidenti $\pi$ in rette; i loro punti di contatto formano una curva, eventualmente riducibile, di cui consideriamo una componente $C$. Sia $Q$ un punto di essa, e $Q'$ il punto infinitamente vicino a $Q$ su $C$. I piani tangenti ad $F$ in $Q$ e $Q'$ segano $\pi$ in rette; quindi ciascuno di essi giace in un $S_3$ con $\pi$. Ma il piano tangente in $Q'$ contiene $Q$, quindi i detti piani determinano con $\pi$ un unico $S_3$, il quale è tangente ad $F$ in tutti i punti di $C$. Al variare di $P$, la curva $C$ descrive un sistema (almeno) $\infty^1$; e se, come si è supposto, la superficie $F$ non appartiene ad un $S_3$, ognuna di queste curve $C$, dovendo stare nello $S_3$ da essa determinato, e nello $S_3$ determinato-dalla curva infinitamente vicina, è contenuta in un piano (potendosi ridurre, in particolare, ad una retta); i piani di due curve $C$ infinitamente vicine stanno in un $S_3$. Si conclude che una superficie per cui valga l’ipotesi ammessa, è rigata, o contiene $\infty^1$ curve piane i cui piani formano una sviluppabile. Nel primo caso le rette che congiungono i punti di contatto di due piani incidenti in rette sono le generatrici della rigata; nel secondo sono le rette dei piani della sviluppabile; in nessuno dei casi esse riempiono lo $S_4$. % ?? uncertain: segno fra «determinato» e «dalla» -- trattino stampato / artefatto di scansione.
% EDITORIAL D0012 RISOLTO: Il trattino fra le parole è netto nella scansione e viene conservato come stampato.

Ciò posto si abbia in $S_4$ una superficie $F$, dotata soltanto di un numero finito di punti doppi impropri ordinari, la cui proiezione $\Phi$ fatta da un punto generico $O$ possegga una curva doppia riducibile; diciamo anzitutto che questa curva deve essere connessa, cioè formata di due o più parti aventi qualche punto comune. Invero le curve doppie apparenti
\[
C = \Sigma K
\]
che rispondono agli insiemi di codeste parti $K$, formano un sistema
%CONT>>
% scan: s013L
% !PRINTED 10
\weberpage{10}
% head: INTRODUZIONE
%<<CONT
% stamp: Fascia scura verticale sul margine destro e piccoli segni sparsi della scansione.
razionale (o almeno unirazionale) i cui elementi corrispondono ai punti $O$ dello spazio $S_4$, e perciò le $C$ appartengono ad un sistema lineare di curve su $F$\origfootnote{1}{s013L-1}{Teorema stabilito per le serie razionali di gruppi di punti sopra una curva ed esteso alle varietà da \textsc{Enriques}. Cfr. \textsc{Enriques-Chisini}, op. cit., Libro V, cap. I, \S~9, vol. III, pag. 78 (cfr. pag. 485).}. Se tale sistema è irriducibile, le $C$ riducibili che ne fan parte debbono essere connesse\origfootnote{2}{s013L-2}{Principio di degenerazione. Cfr. \textsc{Enriques-Chisini}, Op. cit., Libro V, cap. III, \S~36, vol. III, pag. 405.}; se invece tutte le $C$ del detto sistema lineare sono riducibili esse vengono composte con le curve $K$ di un fascio, non aventi fra loro intersezioni variabili.

Ma questo caso non può presentarsi perchè due curve $K$ contenenti due coppie di punti di $F$ come $AA_1$ e $AA_2$ hanno certo (in $A$) un punto comune; e si avverta che esse non possono coincidere in una medesima curva poichè al variare di $A_2$ questa verrebbe a contenere tutti i punti della superficie.

Pertanto si è condotti ad esaminare se e come la superficie $\Phi$, proiezione di $F$ da un punto generico $O$, possa avere una curva doppia composta di due o più parti fra loro connesse.

Un punto di connessione di tali parti potrà essere:

$1^\circ)$ un semplice incrocio, cioè un punto doppio comune a due rami semplici, che pertanto risulterà un tacnodo per la superficie, ovvero

$2^\circ)$ un punto triplo, per cui passino tre rami della curva doppia, due dei quali possono a priori anche appartenere ad una medesima componente irriducibile.

$1^\circ)$ Un punto $\bar{P}$ comune a due rami semplici della curva doppia $\bar{C}$ di $\Phi$ sarà in generale proiezione di due punti $P$ e $P'$ allineati col centro $O$ di $S_4$, e si vede tosto che i piani tangenti ad $F$ in $P$ e $P'$ appartengono allo stesso iperpiano che ha per traccia il piano tangente a $\Phi$ in $\bar{P}$, e che quindi sono incidenti secondo una retta. Ma siccome il punto doppio $\bar{P}$ di $\bar{C}$ viene supposto esistere per ogni posizione del centro $O$, la superficie $F$ dovrebbe possedere $\infty^3$ coppie di piani tangenti incidenti e le rette che congiungono i punti di contatto di queste coppie di piani dovrebbero riempire tutto lo spazio; ciò è assurdo per il lemma premesso.

Un esame speciale esige il caso in cui i punti $P$ e $P'$ siano infinitamente vicini, cadendo dunque nei punti che una tangente per $O$ ha comuni con la superficie $F$. In questo caso il punto $\bar{P}$, incrocio dei due rami della curva doppia di $\Phi$, sarà sempre un tacnodo per $\Phi$, per modo che ogni piano per $\bar{P}$ segherà $\Phi$ in una curva avente ancora un punto doppio infinitamente vicino a $\bar{P}$. Segue di qui che
%CONT>>
% scan: s013R
% !PRINTED 11
\weberpage{11}
% head: INTRODUZIONE
%<<CONT
% stamp: Fascia scura sul margine sinistro e piccoli segni sparsi della scansione.
l’iperpiano proiettante da $O$ il piano tacnodale di $\Phi$ deve segare $F$ secondo una curva che ha quattro intersezioni con ogni piano per la tangente $P\,P'$ e quindi possiede $P$ e $P'$ come punti doppi; salvo ad esaminare i casi particolari in cui la molteplicità effettiva di $P$ diventi superiore alla molteplicità virtuale 2, abbassandosi invece la molteplicità di $P'$. L’ipotesi a cui si riferisce il nostro discorso importa che, movendosi $O$ nello spazio $S_4$, si abbia sempre per esso una tangente, diciamo $p$, incontrante $F$ in due punti infinitamente vicini, che sarebbero doppi per la sezione con l’iperpiano tangente in $P$; quindi, ognuno degli $\infty^3$ iperpiani tangenti ad $F$ dovrebbe toccare la superficie in due punti $P$ e $P'$ infinitamente vicini. Ma è facile vedere che questa proprietà non può competere ad una superficie $F$ di $S_4$; infatti una proiezione generica di questa sullo $S_3$ sarebbe tale che ogni suo piano tangente sarebbe tangente ad essa anche in un punto infinitamente vicino, e perciò risulterebbe essere una sviluppabile; allora anche la superficie di $S_4$ di cui essa è proiezione dovrebbe essere sviluppabile, e tuttavia una sviluppabile di $S_4$ possiede soltanto $\infty^1$ e non $\infty^3$ iperpiani che la toccano in due punti infinitamente vicini.

La conclusione di questo discorso è di dimostrare impossibile l’ipotesi che abbiamo fatto, cioè che la proiezione della nostra superficie $F$ da un punto generico $O$ di $S_4$ sia una $\Phi$ la cui curva doppia si componga di due componenti che si incrociano in un punto $\bar{P}$, immagine di due punti infinitamente vicini $P$ e $P'$ di $F$. Senonchè occorre esaminare il caso particolare a cui si è innanzi accennato: che in luogo di un iperpiano tangente in $P$ la cui intersezione con $F$ abbia due punti doppi $P$ e $P'$, s’incontri un iperpiano tangente la cui sezione abbia in $P$ un punto triplo e in $P'$ un punto semplice, ovvero soltanto in $P$ un punto quadruplo. Dei due casi che qui si presentano il primo non dà luogo ad eccezione nel ragionamento precedente, perchè dovranno aversi ancora $\infty^3$ iperpiani tangenti ad $F$ che la tocchino. pure (virtualmente) in un punto infinitamente vicino. Invece occorre fermarsi un momento sulla seconda ipotesi, in cui si avrebbero a priori non più $\infty^3$, ma $\infty^2$ iperpiani secanti la $F$ secondo curve dotate di un punto quadruplo. Di fatto questa ipotesi non può presentarsi perchè una superficie dello $S_4$ tale che per ogni punto di esso si abbia un iperpiano tangente che la seghi secondo una curva dotata di punto quadruplo, una tale superficie — diciamo — non può appartenere allo $S_4$, ma giace interamente in uno spazio $S_3$. Giacchè le curve sezioni iperpiane di essa sono tali che il piano osculatore in un punto qualsiasi ha con la curva stessa un contatto quadripunto, e perciò sono curve piane. % ?? uncertain: punteggiatura dopo «tocchino» -- virgola / punto.
% EDITORIAL D0013 RISOLTO: Il segno dopo «tocchino» è un punto senza coda, non una virgola; la punteggiatura della trascrizione è adeguata alla stampa.

$2^\circ)$ Un punto $\bar{P}$ che sia triplo per la curva doppia $C$ di $\Phi$ e punto di connessione di componenti irriducibili di questa, può essere: % [sic] C -> \bar{C}, la curva doppia di \Phi è indicata con \bar{C} nelle righe contigue.

% scan: s014L
% !PRINTED 12
\weberpage{12}
% head: INTRODUZIONE
% stamp: Fascia scura verticale sul margine destro e piccoli segni sparsi della scansione.
$a)$ punto semplice comune a tre componenti irriducibili di $\bar{C}$; ovvero

$b)$ punto semplice per una componente irriducibile di $\bar{C}$, e doppio per un’altra.

$a)$ Nell’ipotesi $a)$ il punto $\bar{P}$ sarà proiezione di tre punti di $F$, diciamo $P_1$, $P_2$, $P_3$; e le componenti $\bar{K}_1$, $\bar{K}_2$, $\bar{K}_3$ della $C$ per cui $\bar{P}$ è punto triplo saranno proiezioni di tre curve, rispettivamente $K_1$, $K_2$, $K_3$, per cui si può ritenere che la prima passi per i punti $P_2$ e $P_3$; la seconda per $P_1$ e $P_3$ e la terza per $P_1$ e $P_2$. % [sic] C -> \bar{C}, qui si parla della curva doppia di \Phi, detta \bar{C} nelle due alternative precedenti.

Ricordiamo che la superficie $\Phi$ proiezione di $F$ deve avere una singolarità simile a $\bar{P}$, comunque il centro di proiezione $O$ vari nello $S_4$, e in particolare se esso si faccia muovere sulla retta $P_1P_2P_3$. Ora, in corrispondenza a tale variazione di $O$, potrà accadere che le curve doppie apparenti $K_1K_2K_3$ su $F$ restino fisse, oppure variino con $O$. Se una di esse, per es. $K_1$, resta fissa, si deduce che essa è una conica: infatti, per essere sempre la curva $\bar{K}_1$ proiettata nella medesima curva, bisogna che ogni retta congiungente un punto di $\bar{K}_1$ con un punto della retta $P_2P_3$ incontri in un altro punto $K_1$, e quindi la $K_1$ risulta appartenente ad un piano per la retta $P_2P_3$; una curva piana che si proietti come questa in una curva doppia (e non di molteplicità superiore) dev’essere una conica, e perciò la superficie $F$, possedendo $\infty^2$ coniche, dovrà essere proiezione di una superficie di Veronese. % [sic] \bar{K}_1 in «la curva \bar{K}_1 proiettata» e in «di \bar{K}_1» -> K_1, la curva su F è denotata K_1 nel periodo precedente.

Conviene ora esaminare l’ipotesi che la curva $K_1$ varii al variare di $O$ sulla retta $P_2P_3$. In tal caso, per una conveniente posizione di $O$, la $K_1$ verrà a passare anche per il punto $P_1$ di codesta retta, e si avrà corrispondentemente una superficie $\Phi$ dotata di un punto triplo particolare $\bar{P}$, nel quale due falde, per es. quelle corrispondenti agli intorni di $P_1$ e $P_2$, riescono fra loro tangenti. Ma ciò importa che i piani tangenti ad $F$ in $P_1$ e $P_2$ s’incontrino secondo una retta. Nelle nostre ipotesi dunque la $F$ verrà ad avere $\infty^3$ coppie di piani tangenti incidenti fra loro, in modo che le rette che congiungono i punti di contatto di queste coppie di piani riempiono lo spazio $S_4$; e ciò è escluso dal lemma.

$b)$ Esaminiamo ora l’ipotesi in cui il punto triplo $\bar{P}$ di $\Phi$ sia semplice per una componente $\bar{K}_1$ e doppio per una componente $\bar{K}_2$ della curva doppia. Le $\bar{K}_1$ e $\bar{K}_2$ saranno proiezioni di due curve doppie apparenti $K_1$ e $K_2$ di $F$, di cui la prima potrà supporsi passare semplicemente per $P_2$ e $P_3$, mentre la seconda passerà semplicemente per $P_2$ e $P_3$ e doppiamente per $P_1$. Ora, facendo variare il centro di proiezione $O$ sulla retta $P_1P_2P_3$, potrà accadere che la curva $K_1$ resti fissa ovvero che varii con $O$. Se resta fissa si conclude come prima che è una curva piana, e quindi una conica, donde segue che la $\Phi$ è proiezione d’una superficie di Veronese. Se invece la $K_1$ varia con % ?? uncertain: segno dopo K_1 nell’ultima frase -- nessuna virgola / virgola.
%CONT>>
% scan: s014R
% !PRINTED 13
\weberpage{13}
% head: INTRODUZIONE
%<<CONT
% stamp: Piccoli segni e macchie sparsi della scansione.
$O$, per una qualche posizione di questo, essa viene a passare anche per $P_1$: la singolarità che la $\Phi$ presenta in $\bar{P}$ risulta ora un punto triplo in cui almeno due falde si toccano, e perciò i piani tangenti ad $F$ in due punti come $P_1$ e $P_2$ risulteranno incontrarsi secondo una retta: la $F$ possiede $\infty^3$ coppie di piani tangenti incidenti e si ricade nel caso già escluso.

Da tutto ciò che precede risulta che \textit{la proiezione nello spazio $S_3$ di una superficie dello $S_5$ priva di singolarità possiede in generale una curva doppia irriducibile; fa eccezione soltanto il caso di una superficie di Veronese che di fatto si proietta nella superficie di Steiner, dotata di tre rette doppie passanti per un punto triplo.}
% !PRINTED
\weberblank{[s015L]}
% head:
% scan: s015L; pagina bianca, senza numero di pagina stampato.
% stamp: Fascia scura verticale sul margine destro e piccoli segni sparsi della scansione.


\clearpage\thispagestyle{plain}
% !PRINTED
\weberpage{[s015R]}
% head:
% scan: s015R; nessun numero di pagina stampato.
% stamp: «15» manoscritto nell’angolo superiore destro; piccoli segni sparsi e una sottile fascia scura sul margine sinistro inferiore.
\edchapter{I}{SISTEMI LINEARI DI CURVE}{Capitolo I. SISTEMI LINEARI DI CURVE}{SISTEMI LINEARI DI CURVE}

\edsection{1. Fasci lineari.}{1. Fasci lineari.}{sec.I.1}

Sia $F$ una superficie che, per semplicità di discorso, supporremo appartenere allo spazio ordinario ed essere dotata delle singolarità normali definite nella precedente introduzione.

Consideriamo una funzione razionale $t$ del punto variabile su $F$ che non si riduca ad una costante:
\weq{(1)}{t = \frac{\varphi(xyz)}{\psi(xyz)}.}

Per ogni valore determinato $t = t_0$ esisteranno punti di $F$ in cui la $t$ assume questo valore: il loro luogo $C$ si dirà una \textit{curva di livello} della funzione razionale $t$ sopra $F$. Al variare di $t$ si avrà su $F$ una semplice infinità di curve di livello che si dice costituire un \textit{fascio lineare}. Tra le $C$ del fascio figura la \textit{curva degli zeri}, cioè la curva $\varphi = 0$ che risponde a $t = 0$, e la \textit{curva dei poli}, cioè la $\psi = 0$ per cui $t = \infty$.

Nel caso che queste due curve abbiano una parte in comune $K$, la $K$ su cui $t$ assume forma indeterminata, può ritenersi come componente fissa delle $C$, ovvero togliersi da tutte le $C$.

Le definizioni date innanzi sono perfettamente analoghe a quelle che si riferiscono alle curve ed alle serie $g_n^1$ su di esse, e danno luogo ad osservazioni simili. Il fascio delle $C$ definito dalla (1) è dato ugualmente dalla funzione razionale
\[
\tau = \frac{\alpha t + \beta}{\gamma t + \delta};
\]
la sostituzione razionale ha il solo effetto di cambiare il valore di $t$ che corrisponde ad ogni singola curva $C$; pertanto tutte le $C$ figurano ugualmente nel fascio e così la curva degli zeri come quella dei poli non hanno per esso alcun significato particolare. Dato il
%CONT>>
% scan: s016L
% !PRINTED 16
\weberpage{16}
% head: CAPITOLO PRIMO
%<<CONT
% stamp: Dark scan border at the right edge and scattered specks.
fascio, una funzione razionale che vi corrisponde resta determinata dalla scelta arbitraria delle curve $\tau=0$ e $\tau=\infty$, ed insieme della curva unità $\tau=1$, questa scelta determinando la sostituzione precedente. Se per due funzioni $t$ e $\tau$ coincidono la curva degli zeri e dei poli, esse differiscono per una costante moltiplicativa.

Notiamo inoltre che la funzione razionale (1) è definita soltanto rispetto \textit{al modulo $F$}. Così due diverse funzioni $\frac{\varphi}{\psi}$ e $\frac{\varphi_1}{\psi_1}$ dovranno ritenersi identiche sopra $F$ quando sia per tutti i punti di $F$:
\[
\frac{\varphi}{\psi}=\frac{\varphi_1}{\psi_1}; % ?? uncertain: punctuation after the equality -- semicolon or colon.
% EDITORIAL D0015 RISOLTO: The enlarged scan clearly has a dot above a comma-shaped lower stroke after the fraction.
\]
cioè
\weq{(2)}{\varphi\psi_1-\psi\varphi_1\equiv 0 \qquad (\mathrm{mod.}\ F)}
ossia
\[
\varphi\psi_1-\psi\varphi_1=AF,
\]
dove $A$ rappresenta un polinomio.

Le curve $C$ di un fascio lineare hanno tutte lo stesso ordine; se per una di esse l’ordine si riduce apparentemente, vuol dire che la curva si è spezzata e contiene una parte all’infinito. Per eliminare la riduzione apparente, quando non si voglia ricorrere ad una trasformazione proiettiva della superficie, basta far uso di coordinate omogenee $x_1, x_2, x_3, x_4$. Così ad esempio nella (1), se il polinomio $\varphi$ era di grado inferiore a $\psi$, passando a coordinate omogenee, $\varphi$ viene moltiplicato per una conveniente potenza di $x_4$ e diventa quindi manifesto che la curva degli zeri viene a contenere anche la sezione di $F$ col piano all’$\infty$ ($x_4=0$) contato un certo numero di volte.

La definizione di fascio lineare di curve $C$ ha un semplice significato proiettivo: invero l’equazione
\weq{(3)}{\varphi(xyz)-t\psi(xyz)=0}
rappresenta un fascio di superficie le quali segano la $F$ lungo le curve $C$; e la (2) dice che \textit{lo stesso fascio di curve può essere segato sulla $F$ da diversi fasci di superficie.}

Il fascio delle superficie (3) avrà una curva base $\Gamma$ la quale in generale non apparterrà ad $F$ ed incontrerà $F$ in un numero finito di punti. Questi \textit{punti base del fascio lineare} di curve $C$ sono punti di indeterminazione essenziale per la funzione razionale (1).

Infatti la funzione, che in un punto base assume la forma $\frac{0}{0}$, possiede ivi una singolarità non eliminabile: se invero si cerchi di definire ivi codesto valore con considerazioni di limite, il valore così definito per continuità sarà diverso quando ci si avvicini al punto base sopra curve $C$ diverse.

% scan: s016R
% !PRINTED 17
\weberpage{17}
% head: SISTEMI LINEARI DI CURVE
% stamp: Dark scan borders and scattered specks.
Un punto base per le curve $C$ di un fascio lineare potrà avere per le $C$ una certa molteplicità $i>1$ e in tal caso si dirà un \textit{punto base $i$-plo del fascio.}

Infine potrà accadere che la curva base $\Gamma$ del fascio (3) ovvero una parte $K$ di essa, se è riducibile, giaccia sopra $F$; allora (come già si è detto) $K$ si potrà considerare come componente fissa delle curve $C$ del fascio (\textit{riducibile}), ma si può anche prescindere da $K$ e ritenere il fascio delle $C$ come costituito dalle sole intersezioni variabili della $F$ colle superficie (3).

\edsection{2. Fasci irrazionali.}{2. Fasci irrazionali.}{sec.I.2}

Un fascio lineare di curve $C$ sopra $F$ viene caratterizzato geometricamente dalle proprietà seguenti:

1) è una serie (algebrica) $\infty^1$ d’\textit{indice} 1, cioè tale che per ogni punto generico di $F$ passi una $C$;

2) tale serie è \textit{razionale}, cioè può porsi in corrispondenza biunivoca algebrica con una retta punteggiata o con la serie dei valori di un parametro $t$.

Infatti la $t$ così definita riesce una funzione univoca e perciò razionale dei punti di $F$ e le curve $C$ sono per essa curve di livello.

Se si chiama \textit{fascio} una serie di curve $C$ che sia d’indice 1, cioè goda della proprietà 1) si può dire dunque che \textit{ogni fascio razionale è lineare.}

Ma possono aversi sopra una superficie anche dei \textit{fasci irrazionali}, le cui curve costituiscono gli elementi di un ente algebrico non razionale, avente un certo genere $p>0$. Così per esempio le generatrici di una rigata (le cui sezioni piane sono) di genere $p$, formano appunto un fascio irrazionale di rette.

Si può costruire un esempio di superficie che abbia un fascio irrazionale di curve di genere $p$, intersecando con una varietà qualsiasi dello $S_4$ un cono di seconda specie costituito dai piani dello $S_4$ stesso che proiettano da una retta i punti di una curva di genere $p$.

Conviene rilevare che un \textit{fascio di curve, appartenente alla superficie $F$, è certo razionale se possiede un punto base in un punto semplice per la superficie.}

Infatti se il punto base $O$ è semplice per le curve $C$ di un fascio, la serie delle $C$ si trova in corrispondenza biunivoca col fascio delle tangenti in $O$. E se invece il punto è $i$-plo per le $C$ ($i>1$) allora le $C$ corrispondono in generale ai gruppi di un’involuzione d’ordine $i$ nel fascio delle tangenti anzidette, e questa involuzione è razionale (teorema di Lüroth)\origfootnote{1}{s016R-1}{Cfr. \textsc{Enriques-Chisini}, Op. cit., Libro II, cap. I, \S\ 3, vol I, pag. 167.}.

% EDITORIAL: segnatura originale: \textsc{Enriques F.} - \textit{Superficie algebriche.}\hfill 2
\printersignature{\textsc{Enriques F.} - \textit{Superficie algebriche.}\quad 2}

% scan: s018L
% !PRINTED 18
\weberpage{18}
% head: CAPITOLO PRIMO
% stamp: Short strokes in the left margin, a dark scan border at the right edge, and scattered specks.
Tuttavia questo ragionamento può cadere in difetto se le curve $C$ abbiano in $O$ tangenti fisse, cioè posseggano come punti base, non soltanto $O$, ma anche altri punti ad esso infinitamente vicini. Per eliminare l’eccezione basta osservare che procedendo sopra un ramo di una $C$ fino ad un ordine $r$ convenientemente elevato si troverà un primo punto $O_r$, appartenente all’intorno $r$-mo di $O$, suscettibile di descrivere liberamente l’intorno del punto base $O_{r-1}$, che, come l’intorno di $O$, costituisce un ente razionale: se $O_{r-1}$ è un punto base semplice, le curve $C$ corrispondono biunivocamente ai punti $O_r$ del detto intorno; se $O_{r-1}$ è $i$-plo e gli succedono più punti liberi, $O_r$, le $C$ rispondono ai gruppi di un’involuzione descritta da questi gruppi di punti.

\edsection{3. Sistemi lineari.}{3. Sistemi lineari.}{sec.I.3}

La definizione del fascio lineare di curve sopra $F$ si può generalizzare considerando i \textit{sistemi lineari di dimensione $r>1$}.

Un sistema lineare $\infty^r$ di curve $C$ viene segato sopra $F$ da un sistema lineare di superficie (o di ipersuperficie)
\weq{(1)}{\lambda_0\varphi_0+\lambda_1\varphi_1+\ldots+\lambda_R\varphi_R=0;}
la \textit{dimensione} $r$ riesce minore di $R$ quando nel sistema (1) vi sono $\infty^h$ superficie che contengono come parte la $F$, e allora $r=R-h-1$.

Per i sistemi lineari più volte infiniti $|C|$, come pel fascio, si definiscono le curve \textit{componenti fisse} delle $C$ e i \textit{punti base} di $|C|$; un punto base è un punto comune a tutte le $C$, dove si annullano contemporaneamente le $\varphi$. % [sic] dalle -> delle, the fixed components are described as belonging to the curves C.
% EDITORIAL S0007 ANNULLATO: la scansione s018L stampa «delle»; errore di trascrizione corretto, nota rimossa.

Un sistema lineare $|C|$ si può definire in modo intrinseco indipendentemente dal carattere proiettivo della superficie, così come si è fatto per il fascio lineare. Perciò si considerino $r$ funzioni razionali del punto variabile di $F$ le quali abbiano tutte la stessa curva polare e che supporremo rappresentate da
\[
t_1=\frac{\varphi_1}{\varphi_0}\qquad t_2=\frac{\varphi_2}{\varphi_0}\qquad \ldots \qquad t_r=\frac{\varphi_r}{\varphi_0},
\]
ammettendo per semplicità di discorso che esse siano linearmente indipendenti sopra $F$, sicchè \textit{non sia possibile} trovare $r+1$ costanti (non tutte nulle) $a_0, a_1, \ldots, a_r$, per cui si abbia identicamente
\[
a_0\varphi_0+a_1\varphi_1+\ldots+a_r\varphi_r\equiv 0 \qquad (\mathrm{mod.}\ F).
\]

Allora la curva di livello della funzione
\[
t_0=\lambda_1t_1+\lambda_2t_2+\ldots+\lambda_rt_r
\]
descrive il sistema lineare $|C|$ di dimensione $r$.

In questo sembra assumere una posizione speciale la curva polare $\varphi_0$; ma si può scegliere come tale una curva qualunque del
%CONT>>
% scan: s018R
% !PRINTED 19
\weberpage{19}
% head: SISTEMI LINEARI DI CURVE
%<<CONT
% stamp: Dark scan border at the left edge and scattered specks.
sistema, giacchè ad uno stesso sistema $|C|$ rispondono infiniti sistemi lineari di funzioni che si deducono l’uno dall’altro con sostituzioni lineari.

La dimensione $r$ del sistema lineare $|C|$ riesce definita geometricamente dalla proprietà che segue: \textit{per $r$ punti generici della superficie passa una ed una sola curva del sistema.}

Un sistema lineare $C$ si dirà \textit{irriducibile} se la curva generica $C$ di esso è irriducibile. Per un sistema lineare irriducibile $|C|$ si dirà \textit{genere effettivo} il genere della sua curva generica, e \textit{grado effettivo} il numero delle intersezioni variabili di due $C$ generiche. Per $r=1$ il grado è nullo, poichè ogni punto comune a due curve del fascio è un punto base che appartiene a tutte le $C$. % [sic] C -> |C|, the first symbol denotes the system, otherwise written with vertical bars.

Fra la dimensione $r$ e il grado $n$ di un sistema lineare $|C|$ sussiste la diseguaglianza $r\leq n+1$; infatti la serie segata sopra una qualunque $C$ di $|C|$ dalle rimanenti curve del sistema, fuori dei punti base, è una $g_n^{r-1}$, per cui si ha $r-1\leq n$.

La serie lineare $g_n^{r-1}$, definita sopra una curva generica $C$, ha fondamentale importanza nello studio del sistema lineare e dicesi \textit{serie caratteristica} di $|C|$\origfootnote{1}{s018R-1}{La serie caratteristica di un sistema lineare è stata introdotta da \textsc{C. Segre} e \textsc{G. Castelnuovo} intorno al 1890. La denominazione di \textit{serie caratteristica} è stata adoperata per la prima volta da \textsc{G. Castelnuovo} nelle: «\textit{Ricerche generali sopra i sistemi lineari di curve piane}». Memorie dell’Acc. di Torino, s. II, t. XLII, 1891, riprodotta a pag. 137 e segg. delle «Memorie scelte», Zanichelli, Bologna, 1937.}.

Il grado effettivo di un sistema lineare di curve, resta anche definito per un sistema di curve riducibili che sia privo di parti fisse. Ora si può dimostrare che il grado vale $n>0$ a meno che il sistema stesso non sia un fascio lineare; ovvero le sue curve siano composte con quelle di un fascio (razionale o no).

Infatti se $n=0$ tutte le curve $C$ del sistema che passano per un punto generico dovranno coincidere in una sola (dimensione $r=1$); oppure dovranno avere in comune una curva $K$, e l’insieme delle $K$ costituirà un fascio, razionale o no, con cui le $C$ vengono composte.

\edsection{4. Proprietà caratteristica dei sistemi lineari.}{4. Proprietà caratteristica dei sistemi lineari.}{sec.I.4}

Abbiamo rilevato che, essendo $|C|$ un sistema lineare di dimensione $r$ sopra la superficie $F$, per $r$ punti generici di $F$ passa una ed una sola curva $C$. Questa proprietà è caratteristica per $r>1$\origfootnote{2}{s018R-2}{\textsc{F. Enriques}, \textit{Una questione sulla linearità ecc.} Atti Acc. Lincei, 1893. Cfr. \textsc{C. Segre}, \textit{Introduzione alla geometria sopra un ente algebrico semplicemente infinito.} Ann. di Mat., s. II, t. 22 (1894).},
%CONT>>
% scan: s020L
% !PRINTED 20
\weberpage{20}
% head: CAPITOLO PRIMO
%<<CONT
% stamp: Dark scan border at the right edge, a short dark stroke in the left margin, and scattered specks.
mentre si è visto che per $r=1$ possono aversi anche fasci non lineari.

La dimostrazione del teorema si dà come segue.

Pongasi per semplicità di discorso che le curve del sistema $|C|$ siano irriducibili, e si consideri dapprima il caso $r=2$. Le $\infty^2$ curve si segano a due a due nei gruppi di un’involuzione $I_n$ di ordine $n$ uguale al grado di $|C|$. Si assumano questi gruppi come elementi (punti) di una nuova superficie $F'$; allora si avrà su $F'$ un sistema $\infty^2$ di curve $C'$ trasformato di $|C|$, che si segano a due a due in un punto. Le curve $C'$ che passano per un punto $A'$ formano un fascio dotato di un punto base, e perciò razionale (cfr. \S\ precedente). Accanto a questo fascio se ne consideri un altro con un punto base $B'$. I due fasci di $C'$ coi centri $A'$ e $B'$ si possono riferire rispettivamente a due fasci di rette $A$ e $B$ nel piano, per mezzo di proiettività che facciano corrispondere ugualmente alla curva $C'$, congiungente $A'$ e $B'$, la retta $AB$. Da questo riferimento risulta una corrispondenza biunivoca fra i punti di $F'$ e i punti del piano, dove alle rette del piano corrispondono le curve $C'$ e ai fasci di rette i fasci di $C'$; ciò porta che le $C'$ sopra $F'$ e quindi anche le $C$ sopra $F$, formino un sistema lineare (\textit{rete}). Il teorema essendo così stabilito per le reti si estenderà al caso dei sistemi lineari $\infty^3$, e così in generale, passando dalla dimensione $r$ alla dimensione $r+1$. Se $|C|$ è un sistema $\infty^3$ di curve, tale che per tre punti generici passi una $C$, le $C$ passanti per un punto fisso $A$ formeranno una rete. Ora le reti definite da due punti base $A$ e $B$ potranno riferirsi proiettivamente a due stelle di piani $A'$ e $B'$ dello spazio ordinario, in modo che alle $C$ del fascio comune alle due reti rispondano ugualmente gli stessi piani comuni alle due stelle. Da ciò risulta una corrispondenza biunivoca (omografica) fra il sistema delle $C$ e i piani dello spazio $S_3$, dove ai piani di un fascio corrisponderanno le $C$ di un fascio, e ai piani di una stella le $C$ di una rete; ciò porta che il sistema $|C|$ è un sistema lineare, c. d. d.

\edsection{5. Estensione dei teoremi di Bertini.}{5. Estensione dei teoremi di Bertini.}{sec.I.5}

\textsc{E. Bertini} ha dato per i sistemi lineari di curve piane due teoremi fondamentali\origfootnote{1}{s020L-1}{Rendic. Istituto Lombardo, 1880.}: il primo dice che le curve di un sistema lineare non possono avere punti doppi o multipli variabili che non appartengano ad una componente fissa; il secondo — che è un corollario del primo — afferma che qualora le parti variabili di un sistema lineare siano riducibili esse si compongono con le curve di un fascio.

% scan: s020R
% !PRINTED 21
\weberpage{21}
% head: SISTEMI LINEARI DI CURVE
% stamp: Short curved vertical mark to the right of the first footnote's last line; scattered specks and a small blot near the lower edge.
Questi teoremi si estendono facilmente ai sistemi lineari di curve appartenenti ad una superficie qualunque\origfootnote{1}{s020R-1}{Cfr. \textsc{F. Enriques}, \textit{Ricerche di Geometria ecc.} Memorie Acc. di Torino, 1893 e \textit{Introduzione alla geometria sopra una superficie algebrica.} Società Italiana delle Scienze (detta dei XL), 1896. Il teorema della riducibilità è già osservato in Noether, Math. Ann., VIII, 1874.}.

Quando si segua la dimostrazione indicata nelle lezioni di \textsc{Enriques-Chisini}\origfootnote{2}{s020R-2}{Op. cit., Libro II, cap. I, \S\ 5, vol. I, pag. 180.}, l’estensione riesce immediata.

Se la curva generica $C$ del sistema lineare $|C|$ possiede un punto doppio o multiplo variabile $O$, la curva $C$ ha almeno due intersezioni riunite in $O$ con ogni curva infinitamente vicina, e quindi con le curve di tutti i fasci a cui essa appartiene: ciò significa che $O$ è punto base ovvero un punto appartenente ad una curva $K$ componente fissa delle $C$.

Enunciamo dunque che: \textit{le curve di un sistema lineare sopra una superficie non possono avere punti doppi, o multipli variabili} (fuori dei punti base) \textit{che non appartengano ad una componente fissa.}

Il secondo teorema di Bertini segue come corollario.

Le curve $C$ di $|C|$ (non aventi componenti fisse) siano riducibili, per esempio in due parti (variabili) $K$ e $K'$. Se $K$ e $K'$ non appartengono ad un medesimo fascio, con cui siano composte le curve $C$, vi saranno per un punto generico $O$, almeno una curva $K$ ed una curva $K'$, distinta dalla $K$; per conseguenza le $K$ e $K'$ dovranno incontrarsi almeno in un punto variabile, e le $C$ dovranno avere almeno un punto doppio suscettibile di variare liberamente sopra la superficie, il che è assurdo. Invero se i punti comuni alle $K$ e $K'$ che compongono una $C$ cadessero sempre in punti base del sistema $|C|$, tutte le $K$ e tutte le $K'$ passerebbero per quei punti fissi e non avrebbero intersezioni variabili. Non può accadere neppure che le curve complementari $K$ e $K'$ s’incontrino in punti di una curva fissa, perchè questa risulterebbe una componente fissa delle $C$, che si è supposto non esistere.

Enunciamo dunque: \textit{se la curva generica di un sistema lineare $|C|$ sopra $F$ è riducibile:}

1) \textit{$C$ contiene una parte fissa, tolta la quale rimane un sistema irriducibile,} oppure

2) \textit{le curve $C$ sono composte con $r>1$ curve variabili in un fascio, razionale o no} (e sopra questo, concepito come ente $\infty^1$, formano una $g_r$), o infine

3) \textit{le $C$ sono composte come nel caso 2) con le curve di un fascio, ed inoltre con curve fisse che a queste si aggiungono.}

Nei casi 2) e 3) rientrano tutti i sistemi lineari di dimensione $r>1$, e di grado $n=0$ (cfr. \S\ 3).

% scan: s021L
% !PRINTED 22
\weberpage{22}
% head: CAPITOLO PRIMO
% stamp: Dark vertical scan-edge band at the right margin.
\edsection{6. Superficie immagini di sistemi lineari.}{6. Superficie immagini di sistemi lineari.}{sec.I.6}

Si consideri un sistema lineare \(|C|\) irriducibile di dimensione \(r \geq 3\) e perciò di grado \(n > 0\). Riferendo proiettivamente \(|C|\) al sistema degli iperpiani di uno spazio \(S_r\) si dà luogo ad una trasformazione della nostra superficie \(F\) in una \(F'\) d’ordine \(n\), su cui le trasformate delle \(C\) vengono segate dagli iperpiani. Infatti le \(C\) si possono ritenere astrattamente come gli elementi (iperpiani) di uno spazio \(S_r\) e quindi la proiettività indicata fa corrispondere al sistema lineare \(\infty^{r-1}\) delle \(C\) passanti per un punto \(P\), una stella di iperpiani passanti per un punto \(P'\) dello \(S_r\) e, al variare di \(P\) su \(F\) il punto \(P'\) varia descrivendo una superficie \(F'\).

La corrispondenza univoca fra \(F\) ed \(F'\) è, in generale, univocamente invertibile, perchè alla stella di centro \(P'\) corrisponderà un sistema lineare di curve \(C\) dotato di un punto base (variabile) \(P\) e questo sistema non avrà in generale altri punti base, conseguenti dall’esistenza di \(P\).

In verità non si può escludere che le condizioni di passaggio delle \(C\) per un punto generico di \(F\) portino di conseguenza il passaggio per altri punti \(P_1P_2\ldots P_{m-1}\); allora i gruppi analoghi a \(PP_1\ldots P_{m-1}\) formeranno sopra \(F\) un’\textit{involuzione di ordine \(m\) a cui apparterranno le \(C\)}; e la superficie \(F'\) dello \(S_r\) determinata innanzi risulterà una superficie d’ordine \(\dfrac{n}{m}\), da considerarsi come multipla d’ordine \(m\).

Che effettivamente l’appartenenza ad un’involuzione costituisca una circostanza particolare per un sistema lineare \(|C|\) di dimensione \(r \geq 3\), risulta da un semplice computo di costanti, perchè vi sono almeno 3 curve passanti per \(P\) e linearmente indipendenti, ed in generale non vi è ragione perchè una di queste passi per l’intersezione delle altre due. Invece un sistema lineare \(\infty^2\) (cioè una rete) appartiene sempre ad un’involuzione d’ordine \(n\), ove sia di grado \(n > 1\); poichè per un punto \(P\) passano soltanto due \(C\) linearmente indipendenti e tutte le altre \(C\) per \(P\) formano un fascio, che ha come punti base le loro intersezioni. Così, dunque, una rete \(|C|\) di grado \(n\) conduce ad un \textit{piano multiplo d’ordine \(n\)}, dove le \(C\) hanno per immagini le rette \(n\)-ple e su cui assume importanza speciale la \textit{curva di diramazione}, luogo dei punti del piano a cui corrispondono gruppi dell’involuzione dotati di un punto doppio.

Si può dare un esempio di un sistema lineare \(|C|\) di dimensione \(r \geq 3\) appartenente ad un’involuzione, considerando il sistema delle \(C\) intersezioni della superficie \(F\) con i coni di un certo ordine \(n\) che abbiano il vertice in un punto fisso dello spazio; il sistema \(|C|\) apparterrà ora ad un’involuzione d’ordine \(m\) e sarà di grado \(mn\), designando con \(m\) l’ordine della superficie \(F\). % [sic] mn -> mn^2, the square on the cone order appears to be missing.

% scan: s021R
% !PRINTED 23
\weberpage{23}
% head: SISTEMI LINEARI DI CURVE
% stamp: Scattered marginal scan specks.
Volendo \textit{escludere l’appartenenza ad un’involuzione} di un sistema \(|C|\) diremo in breve che esso è un \textit{sistema semplice}.

Ritornando in generale alla trasformazione della superficie \(F\) cui dà luogo un sistema lineare \(|C|\) di dimensione \(r \geq 3\), noteremo che la corrispondenza proiettiva di cui sopra si è discorso si traduce nelle formule come segue: se il sistema \(|C|\) è rappresentato da
\[
\lambda_0\varphi_0 + \lambda_1\varphi_1 + \ldots + \lambda_r\varphi_r = 0,
\]
si dovrà prendere le coordinate dei punti dello \(S_r\):
\[
X_0 \equiv \varphi_0 \quad X_1 \equiv \varphi_1 \quad \ldots \quad X_r \equiv \varphi_r.
\]

Concluderemo enunciando un teorema che estende alle superficie una nota considerazione relativa alle curve e alle serie lineari sopra di esse:

\textit{la geometria delle trasformazioni birazionali sopra le superficie, quando si ricerchino le relazioni invariantive con un sistema lineare semplice \(\infty^3\) almeno, si rispecchia nella geometria proiettiva della superficie immagine del sistema.}

Così la relazione di due sistemi lineari \(|C|\) e \(|K|\) (semplici \(\infty^3\) almeno) sopra \(F\), di cui il primo contenga il secondo, si rispecchia nel fatto che la superficie immagine di \(|K|\) può ritenersi come proiezione della superficie immagine di \(|C|\) da un certo spazio lineare. Questo, in generale, segherà la detta superficie in un gruppo di punti, se \(|K|\) è dedotto da \(|C|\) con l’imposizione di quei punti base, ed invece segherà la superficie secondo una \(L\) se \(|K|\) sia contenuto parzialmente in \(|C|\) per modo che aggiungendo alle \(K\) la curva fissa \(L\) si ottengano delle curve \(C\).

Per quel che concerne i punti base gioverà tener presente le osservazioni che seguono. Se \(O\) è, sopra la superficie \(F\), un punto base \(i\)-plo di \(|C|\) (semplice per \(F\)), % ?? uncertain: Damaged word after i-plo -- di or in.
% EDITORIAL D0016 RISOLTO: The enlarged damaged type retains the dotted i at the right and remnants of d at the left; the alternatives do not fit the surviving letter order.
esso sarà fondamentale per la trasformazione che muta la \(F\) nella \(\Phi\); % ?? uncertain: Punctuation after Phi -- semicolon or comma.
% EDITORIAL D0017 RISOLTO: The enlarged scan shows the upper point and lower comma of a semicolon after Phi.
avente come sezioni iperpiane le \(C\); e se le \(C\) non hanno alcuna tangente fissa in \(O\), ad \(O\) corrisponderà su \(\Phi\) \textit{una curva eccezionale} di ordine \(i\). Se invece \(O\), essendo base \(i\)-plo per \(|C|\) con \(i \geq 1\), le \(C\) abbiano una tangente fissa in \(O\) (cioè un punto base semplice \(O_1\) infinitamente vicino ad \(O\)), senza contatto più elevato, si potrà dire ancora che ad \(O\) corrisponde su \(\Phi\) una curva eccezionale d’ordine \(i\), ma questa, in generale, si comporrà di una curva di ordine \(i-1\), i cui punti rispondono ai punti variabili nell’intorno del primo ordine di \(O\), e d’una retta immagine dell’intorno di \(O_1\) (i cui punti appartengono all’intorno del secondo ordine di \(O\)). In particolare se \(i=1\) l’intero intorno del punto semplice \(O\) avrà per immagine su \(\Phi\) l’intorno di un punto doppio conico, cui si aggiungerà la retta eccezionale immagine di \(O_1\) (si pensi, ad esempio, alla rappresentazione piana del cono quadrico).

% scan: s022L
% !PRINTED 24
\weberpage{24}
% head: CAPITOLO PRIMO
% stamp: Vertical scan-edge lines and dark patches at the left margin.
\textsc{Nota.} -- In generale se il punto semplice \(O\) della superficie \(F\) costituisce una singolarità comunque complicata per le curve \(C\) di \(|C|\), si può sciogliere questa singolarità trasformando la \(F\) in una superficie \(\Phi\) su cui l’intorno del punto \(O\) (preso nella sua integrità) verrà rappresentato da una curva eccezionale composta di tante parti quanti sono i rami (lineari o meno) che costituiscono la singolarità delle \(C\) in \(O\); e la connessione di queste parti rispecchierà la composizione della detta singolarità mediante punti multipli o semplici infinitamente vicini, e le loro relazioni di prossimità\origfootnote{1}{s022L-1}{\textsc{Enriques-Chisini}, Op. cit., Libro IV, vol. II, pag. 327 e segg.} % ?? uncertain: Accent on the conjunction before "le loro relazioni" -- e or è.
in una maniera che ha formato oggetto dell’analisi di \textsc{E. Barber} e di \textsc{O. Zariski}\origfootnote{2}{s022L-2}{Cfr. \textsc{Barber} e \textsc{Zariski}, \textit{Reducible exceptional curves of the first Kind.} American Journal of Mathematik, 1935 (pag. 119). \textsc{P. Du Val}, American Journal of Math., 1936.}. Ma avremo occasione di ritornare sull’argomento. % [sic] Mathematik -> Mathematics, apparent misprint in the English journal title.

\edsection{7. Curve equivalenti e sistemi lineari completi.}{7. Curve equivalenti e sistemi lineari completi.}{sec.I.7}

Le considerazioni che seguono estendono passo a passo ai sistemi lineari di curve sopra una superficie ciò che si dice per i gruppi di punti equivalenti e per le serie lineari complete appartenenti ad una curva\origfootnote{3}{s022L-3}{Cfr. \textsc{Enriques-Chisini}, Op. cit., Libro V, vol. III (1924).}.

Sopra la superficie \(F\) due curve (dello stesso ordine) si dicono \textit{equivalenti} quando appartengono ad un medesimo sistema lineare di dimensione \(r \geq 0\): se esse sono distinte vuol dire che appartengono ad un fascio lineare; la considerazione dei sistemi di dimensione zero, significa che una curva deve ritenersi equivalente a sè stessa.

L’equivalenza di due curve \(C\) e \(K\) si esprimerà scrivendo
\[
C \equiv K \qquad \text{o anche} \qquad C = K.
\]

Due curve equivalenti \(C\) e \(K\) possono sempre considerarsi come curva degli zeri e curva dei poli di una funzione razionale \(t\) sopra \(F\); le due curve vengono scambiate fra loro ove si sostituisca alla funzione \(t\) la \(\dfrac{1}{t}\).

La relazione di equivalenza fra curve sopra una superficie gode delle tre proprietà fondamentali dell’uguaglianza:
\begin{enumerate}
\item[1)] la proprietà \textit{riflessiva}, espressa da \(C \equiv C\),
\item[2)] la proprietà \textit{simmetrica}, se \(C \equiv K\) anche \(K \equiv C\),
\item[3)] la proprietà \textit{transitiva}, se \(C \equiv K\) e \(K \equiv L\) anche \(C \equiv L\).
\end{enumerate}

Per dimostrare quest’ultima proprietà si costruisca una funzione razionale \(t\) avente come curva degli zeri \(K\) e come curva dei poli
%CONT>>
% scan: s022R
% !PRINTED 25
\weberpage{25}
% head: SISTEMI LINEARI DI CURVE
%<<CONT
% stamp: Dark vertical scan-edge band at the left margin.
\(C\), e poi una funzione razionale \(\tau\) che abbia \(K\) come curva degli zeri ed \(L\) come curva dei poli; allora la funzione \(\dfrac{\tau}{t}\) avrà come curva degli zeri \(C\) e come curva dei poli \(L\). Ciò prova appunto che \(C\) ed \(L\) sono equivalenti.

Dalle proprietà anzidette risulta il teorema fondamentale:

\textit{Sopra una superficie \(F\) la totalità delle curve equivalenti ad una data curva \(C\) costituisce un sistema lineare che dicesi completo.}

Questo sistema si ottiene come insieme delle curve di livello della funzione razionale che risulta combinando linearmente le funzioni razionali indipendenti, per cui la \(C\) è curva dei poli.

Il sistema \textit{completo} \(C\) così definito sarà in generale \textit{privo di punti base}, ed in ogni caso s’intenderà, salvo avviso contrario, \textit{virtualmente} privo di punti base, ai sensi che verranno precisati più tardi; ma si può anche definire il sistema completo relativamente ad un \textit{gruppo di punti base assegnati} sopra la \(C\). Si noti che, trasformando la superficie in guisa che i punti base della \(C\) diventino curve (eccezionali), si dovrà ritenere che la trasformata di \(C\) sia spogliata di queste curve che altrimenti apparirebbero come componenti fisse del sistema trasformato. E così la nozione del sistema completo relativamente ad un gruppo di punti base si riconduce in generale a quello del sistema completo senza punti base, sopra una superficie trasformata.

Comunque sia, il sistema completo irriducibile, con o senza punti base, resta definito a partire da un sistema lineare (\(\infty^2\) almeno) contenuto in esso, come il sistema più ampio dello stesso grado, che contiene il dato.

Dalla proprietà transitiva dell’equivalenza segue che il sistema lineare completo \(|C|\), con o senza punti base assegnati, viene definito ugualmente a partire da una qualunque delle sue curve.

È ovvio che il sistema completo \(|C|\) si può costruire come segue.

Si costruisca anzitutto un sistema lineare qualunque che contenga (totalmente) \(C\) e poi si ampli successivamente questo sistema, finchè è possibile; siccome la dimensione non può superare il grado aumentato di un’unità, si perverrà in tal guisa ad un sistema \(|C|\) che non sarà contenuto in un altro sistema lineare più ampio dello stesso ordine, e questo sarà il sistema lineare completo definito dalla \(C\). Pertanto, comunque si proceda alla costruzione dei sistemi via via più ampi di cui si è discorso, si riuscirà infine allo stesso sistema lineare completo.

Al concetto di sistema lineare completo risponde il concetto di \textit{superficie normale}, superficie di uno spazio \(S_r\) che non può ottenersi come proiezione di un’altra dello stesso ordine appartenente ad uno spazio superiore: essere la superficie normale significa che il
%CONT>>
% scan: s023L
% !PRINTED 26
\weberpage{26}
% head: CAPITOLO PRIMO
%<<CONT
% stamp: Vertical scan-edge line and dark patches at the right margin.
sistema (irriducibile) delle sue sezioni iperpiane è completo, e viceversa.

Ecco dunque il significato proiettivo del teorema dei sistemi completi. Si consideri, per esempio nello spazio ordinario, una superficie \(F\) di un certo ordine \(n\); potrà accadere che la \(F\) si ottenga come proiezione di due superficie dello stesso ordine dello \(S_4\), \(F_1\) ed \(F'_1\), le proiezioni essendo fatte rispettivamente da due punti \(O\) ed \(O'\). Ora se \(F_1\) ed \(F'_1\) sono superficie normali, risulteranno fra loro proiettive, cioè proiettivamente identiche. Se ciò non è, vuol dire che \(F_1\) potrà ritenersi proiezione di una superficie \(F_2\) dello \(S_5\) ed \(F'_1\) proiezione di una \(F'_2\) dello \(S_5\), e così via; procedendo in un modo qualunque da una superficie ad un’altra dello spazio superiore, si arriverà in ultimo ad una superficie \(F_r\) dello spazio \(S_{r+3}\) che sarà normale; e questa \(F_r\) resterà sempre proiettivamente definita, comunque si arrivi ad essa attraverso serie diverse di superficie (\(F_1, F_2, \ldots,\) o \(F'_1, F'_2, \ldots\)) che si deducono per proiezione l’una dall’altra.

Esempi relativi alle superficie normali s’incontrano già nella teoria delle superficie razionali\origfootnote{1}{s023L-1}{Cfr. \textsc{Enriques-Conforto}, \textit{Le superficie razionali}, Bologna, 1939 (1945).}.

Così una superficie \(F_5\) del \(5^\circ\) ordine dello spazio ordinario, a sezioni piane ellittiche, dotata di curva doppia del \(5^\circ\) ordine appartenente ad un cono quadrico ed avente un punto triplo, triplo anche per la superficie, nel vertice di detto cono, si può ritenere proiezione di più superficie \(F_5\) dello stesso ordine dello spazio \(S_4\), non proiettive fra loro; ma queste sono proiezioni di un’unica superficie normale dello \(S_5\), proiettivamente definita. Queste asserzioni si giustificano partendo dalla rappresentazione piana della \(F_5\) mercè un sistema lineare \(\infty^3\) di cubiche \(C_3\) con 4 punti base semplici \(A_1A_2A_3A_4\).

Infatti un tale sistema \(\infty^3\) è contenuto in \(\infty^1\) sistemi lineari \(\infty^4\) di \(C_3\) cogli stessi punti base \(A_1A_2A_3A_4\), i quali sono, nel piano, omograficamente (e anche birazionalmente) distinti e perciò rappresentano superficie \(F_5\) di \(S_4\), non proiettive fra loro. Invece la superficie normale di \(S_5\) di cui queste \(F_5\) sono proiezioni (da un punto esterno) viene proiettivamente definita dal sistema completo delle \(C_3\) per \(A_1A_2A_3A_4\).

\edsection{8. Somma e differenza di sistemi lineari: teorema del resto.}{8. Somma e differenza di sistemi lineari: teorema del resto.}{sec.I.8}

Curve composte con curve equivalenti, cioè \textit{somme di curve equivalenti, sono equivalenti.}

Infatti siano \(C_1\) e \(C_2\), \(K_1\) e \(K_2\) due coppie di curve equivalenti, e si assuma la \(t\) funzione razionale che abbia \(C_1\) come curva polare
%CONT>>
% scan: s023R
% !PRINTED 27
\weberpage{27}
% head: SISTEMI LINEARI DI CURVE
%<<CONT
% stamp: Scattered marginal scan specks.
e \(C_2\) come curva degli zeri; e \(\tau\) funzione razionale che abbia \(K_1\) come curva polare e \(K_2\) come curva degli zeri; allora la funzione \(\tau t\) avrà come curva polare \(C_1+K_1\) e come curva degli zeri \(C_2+K_2\). Così dall’essere
\[
C_1 \equiv C_2,\qquad K_1 \equiv K_2
\]
si deduce
\[
C_1+K_1 \equiv C_2+K_2.
\]

Risulta di qui che, dati due sistemi lineari qualunque \(|C|\) e \(|K|\), esiste un ben determinato sistema completo \(|C+K|\), che contiene totalmente le curve composte \(C+K\); questo si dirà il \textit{sistema completo somma} di \(|C|\) e \(|K|\).

Conviene notare che se \(|C|\) e \(|K|\) sono sistemi lineari irriducibili di dimensione maggiore di zero, non coincidenti in un medesimo fascio, il sistema completo \(|C+K|\) è irriducibile. Aggiungasi che se i sistemi irriducibili \(|C|\) e \(|K|\) posseggono dei punti base assegnati, questi sono sempre base per \(|C+K|\). C’è qui una ovvia conseguenza della estensione del teorema di Bertini. In particolare è irriducibile il doppio e in generale il sistema \(|sC|\) multiplo secondo \(s\) del sistema \(|C|\), tranne il caso che \(|C|\) sia un fascio.

Essendo ancora \(|C|\) e \(|K|\) due sistemi non coincidenti in uno stesso fascio si considerino le curve composte \(C+K\); abbiamo detto che esse sono contenute nel sistema completo \(|C+K|\); ma in generale potranno essere contenuti anche in sistemi di dimensione inferiore, e in particolare vi sarà un \textit{sistema somma minima} dei due dati. % [sic] contenuti -> contenute, agreement with "esse" (the curves).
% EDITORIAL: la scansione s023R stampa «sistema somma minima»; errore di trascrizione corretto (minimo -> minima).

Per esempio si assuma una superficie \(F\) d’ordine \(n\) dello spazio ordinario, priva di singolarità; per un punto \(O\) di essa passano \(\infty^2\) piani che segano una rete completa; il doppio minimo di questa è il sistema \(\infty^5\) segato dai coni quadrici col vertice in \(O\); ma questo sistema che ha il punto base \(O\) doppio non è completo, essendo contenuto nel sistema \(\infty^6\) segato sulla \(F\) dalle quadriche che la toccano in \(O\).

La \textit{sottrazione} di due sistemi lineari, di cui uno contiene parzialmente l’altro, si definisce come \textit{operazione inversa della somma} e precisamente come segue.

Sia \(|C|\) un sistema lineare di cui faccia parte la curva spezzata \(C \equiv L+K\); allora dicesi sistema \textit{residuo} della curva \(K\) rispetto a \(|C|\), il sistema di tutte le curve \(L\) che insieme a \(K\) costituiscono una curva di \(|C|\). In forza delle proprietà delle curve equivalenti si può affermare che:

se una curva particolare \(C=L+K\) del sistema completo \(|C|\) contiene parzialmente una \(K\), ogni altra curva equivalente a \(K\) fa parte di una \(C\). E quindi si ha il \textit{teorema del resto}: \textit{il sistema}
%CONT>>
% scan: s024L
% !PRINTED 28
\weberpage{28}
% head: CAPITOLO PRIMO
%<<CONT
% stamp: Dark vertical scan-edge band at the right margin and stray marks above the running head.
\textit{lineare \(|L|\) residuo di una \(K\) rispetto ad un sistema \(|C|\), è pure residuo di ogni altra curva equivalente a \(K\).}

Insomma il \textit{sistema lineare completo differenza} di \(|C|\) e \(|K|\):
\[
|L| = |C-K|,
\]
è definito dalla relazione
\[
|C| = |L+K|.
\]

\textsc{Nota.} -- Il teorema enunciato ammette un complemento proiettivo che dimostreremo più avanti. \textit{Vedremo infatti che sopra una superficie \(F\) dello spazio ordinario con singolarità normali, le superficie (aggiunte) di un dato ordine che passano per la curva doppia segano, fuori di questa, un sistema lineare completo.}

In particolare, sopra una superficie dello spazio ordinario, priva di singolarità, le superficie \(\varphi\) di un dato ordine segano un sistema completo.

Se, riferendoci al caso generale, si conducono le \(\varphi\) d’un certo ordine passanti per una curva \(C\) di \(F\), e si considerano le intersezioni residue \(K\) di queste superficie aggiunte con \(F\), \textit{ogni curva \(C\) del sistema completo \(|C|\) risulterà ugualmente residua di ogni \(K\)}, rispetto al sistema segato dalle \(\varphi\) (fuori della curva doppia), e reciprocamente ogni \(K\) è residua di ogni \(C\).

Questa è la forma del \textit{Restsatz} di \textsc{M. Noether}\origfootnote{1}{s024L-1}{\textit{Mathematische Annalen}, Bd. VIII. Cfr. per le curve: \textsc{Brill} e \textsc{Noether}, \textit{Math. Ann.}, Bd. VII; \textsc{Enriques-Chisini}, Op. cit., Libro V.}, che analogamente a ciò che si dice per la geometria sopra una curva (osservazione di \textsc{G. Castelnuovo}, 1890) contiene due affermazioni distinte: la prima relativa alla sottrazione dei sistemi completi, l’altra che ne reca un complemento proiettivo, concernente l’integrità del sistema segato sopra una superficie \(F\) dello spazio \(S_3\) dalle sue aggiunte \(\varphi\) di un dato ordine.

Conviene ora osservare che, mentre l’operazione della somma conduce in generale, come si è detto, da sistemi lineari irriducibili a sistemi irriducibili, ciò non ha più luogo per la sottrazione: quando si sottragga da un sistema irriducibile \(|C|\) una curva \(K\) (irriducibile o no) il sistema completo residuo \(|C-K|\) può risultare riducibile. Un semplice esempio si ha considerando un fascio lineare \(|L|\) che, per semplicità, supponiamo privo di punti base, ed un sistema lineare \(|K|\): se si somma \(|K|\) a \(|2L|\) si ottiene un sistema \(|C|=|K+2L|\). E sottraendo da questo \(|K|\), si ha il sistema doppio di \(|L|\), sistema riducibile composto con le curve del fascio \(|L|\).

Può anche accadere che sottraendo una \(K\) dal sistema irridu%
%CONT>>
% scan: s024R
% !PRINTED 29
\unskip\weberpage{29}%
% head: SISTEMI LINEARI DI CURVE
%<<CONT
% stamp: Scattered marginal specks and a small stray stroke below the upper-left part of the running-head area.
% The word split across this page break is joined without its printed hyphen.
\unskip{}cibile \(|C|\), si ottenga un sistema lineare che contenga una parte fissa \(K'\), il cui distacco dalle \(C\) sia una conseguenza del distacco della \(K\); e similmente può anche accadere che sottraendo una curva \(K\) da \(|C|\) il sistema residuo \(|L|\) abbia come conseguenza qualche nuovo punto base \(O\) oltre ai punti base di \(|C|\) che non appartengono a \(K\). Questo caso si riconduce al precedente con una trasformazione della superficie che muti \(O\) in una curva eccezionale, perchè sulla superficie trasformata codesta curva apparirà come parte fissa di \(|L|\).

\textit{Osservazione.} -- La definizione del sistema \(|L|=|C-K|\) è legata ad un presupposto d’esistenza, poichè la differenza assume significato effettivo soltanto se il sistema \(|C|\) contiene parzialmente \(|K|\). Ma nel seguito gioverà anche parlare di \textit{curve virtuali} che risultano definite dalla sottrazione, anche quando questa sia effettivamente impossibile. Perciò si riterrà la relazione simbolica
\[
|C-K| = |L-M|
\]
come equivalente alla
\[
|C+M| = |K+L|.
\]

Le curve virtuali vengono definite come i numeri negativi nella teoria degli operatori di Peano, cioè come possibili addendi. Sommare la curva virtuale \(C-K\) ad un sistema \(|D|\), vorrà dire sommargli \(C\) e sottrarne \(K\), ciò che risulterà in generale possibile per sistemi sufficientemente ampi. In ogni caso l’equivalenza delle curve virtuali \(C-K\) e \(L-M\) significa l’equivalenza dei sistemi lineari
\[
|D+C-K| = |D+L-M|,
\]
essendo
\[
|D+(C-K)| = |D+C-K|: % ?? uncertain: Punctuation after the final vertical bar -- colon or full stop.
% EDITORIAL D0019 RISOLTO: Two separate aligned points are visible after the final bar in the displayed equality.
\]

Tuttavia si può dare un concetto più ristretto delle curve virtuali come curve non aventi esistenza effettiva ma di cui esiste il doppio ovvero un multiplo, siccome avremo luogo di vedere nel seguito. E questo concetto piuttosto che ai numeri negativi risponde a quello degli \textit{ideali}, che s’incontrano nella teoria dei corpi algebrici.

\textit{Notizia.} -- L’idea, se non il nome, delle « curve virtuali », cioè l’idea che la curva \(C-K\), definita mercè il sistema lineare differenza di due altri, risponda ad un ente matematico dotato di una certa realtà, anche quando la sottrazione non conduca ad una curva effettiva, si è affacciata ad \textsc{Enriques} in rapporto alle curve canoniche: essa costituisce anzi il motivo principale del nuovo sviluppo dato alla teoria nel passaggio dalle \textit{Ricerche} del 1893 alla \textit{Introduzione} del 1896, dove si ricerca in sostanza il più largo significato che assume l’invarianza delle dette curve canoniche, anche quando il ge%
%CONT>>
% scan: s025L
% !PRINTED 30
\unskip\weberpage{30}%
% head: CAPITOLO PRIMO
%<<CONT
% stamp: Dark vertical scan-edge band at the right margin.
% The word split across this page break is joined without its printed hyphen.
\unskip{}nere \(p=0\). Della curva differenza, \(C-K\), l’A. definisce i caratteri virtuali, grado e genere (cfr. \S~10), e rileva che, pur mancando, può avere un doppio dotato di esistenza effettiva. La locuzione « virtuale », applicata alla curva oltrechè ai suoi caratteri (che così venivan designati da \textsc{Enriques}), si trova in una Nota di \textsc{Severi}, \textit{Sulle curve algebriche virtuali appartenenti ad una superficie algebrica} (Rendic. Istituto Lombardo, 1905). L’interpretazion topologica delle curve virtuali come cicli a due dimensioni è stata avvertita da \textsc{S. Lefschetz}.

\edsection{9. Superficie immagine del sistema somma.}{9. Superficie immagine del sistema somma.}{sec.I.9}

Si abbiano sulla superficie \(F\) due sistemi lineari irriducibili \(|C|\) e \(|K|\) rispettivamente di dimensione \(r\) e \(s\), dove sia \(r \geq 1\), ed \(s \geq 1\). Si può costruire in generale una superficie trasformata \(\Phi\), immagine del sistema somma \(|C+K|\), entro uno spazio ad \(R\) dimensioni, dove riesce \(R > r+s\); il sistema somma di cui si parla può essere il sistema somma minimo o anche, se si vuole, il sistema completo.

Comunque sia, ci saranno nello \(S_r\), due serie lineari di spazi rispettivamente ad \(R-r-1\) dimensioni e ad \(R-s-1\) dimensioni secanti le \(K\) e le \(C\), per modo che gli \(\infty^r\) iperpiani passanti per uno \(S_{R-r-1}\) della prima serie segheranno le \(C\), contenendo gli \(S_{R-s-1}\) a cui esse appartengono, e reciprocamente si dica per gli iperpiani che passano per uno \(S_{R-s-1}\) della seconda serie. % ?? uncertain: Subscript in "nello S_r" -- lowercase r or uppercase R; the best visual reading is r.
% EDITORIAL D0020 RISOLTO: The enlarged subscript is a lowercase italic r, unlike the uppercase R in the nearby S_R later on the same page. Retain the printed r and record the mathematical inconsistency separately.

Si fissi uno \(S_{R-r-1}\) della prima serie secante una curva \(C\) di \(|C|\) e un \(S_{R-s-1}\) della seconda, secante una curva \(K\) di \(|K|\); questi saranno contenuti in un iperpiano dello \(S_R\), ed avranno perciò in comune un \(S_{R-r-s-1}\); proiettando da questo sopra un \(S_{r+s}\) si otterrà in tale spazio un’immagine \(F\) della superficie data che conterrà una particolare \(K\) sezione di un \(S_{s-1}\) e una particolare \(C\) sezione di un \(S_{r-1}\) e dove gli iperpiani per lo \(S_{s-1}\) segano le \(\infty^r\) curve \(C\), e gli iperpiani per lo \(S_{r-1}\) le \(\infty^s K\).

Prendiamo in particolare \(r=2\), \(s=1\): saremo condotti ad una superficie dello spazio ordinario su cui la rete delle curve \(C\) viene segata dai piani per un punto \(O\), e il fascio delle \(K\) dai piani per una retta \(a\) (non passante per \(O\)) e la superficie così ottenuta sarà certo in corrispondenza biunivoca con quella da cui siamo partiti, se \(|C|\) e \(|K|\) non appartengono ad una medesima involuzione.

La costruzione della superficie \(F\) si potrà ottenere anche direttamente, ponendo una corrispondenza proiettiva fra la rete \(|C|\) e la stella di piani \(O\), e fra il fascio \(|K|\) e il fascio di piani \(a\). Alla \(\overline{K}\) che si aggiunge alle \(C\) della rete viene così a corrispondere l’intorno del punto multiplo \(O\), che sarà precisamente \(i\)-plo, ove si designi
%CONT>>
% scan: s025R
% !PRINTED 31
\weberpage{31}
% head: SISTEMI LINEARI DI CURVE
%<<CONT
% stamp: Scattered marginal scan specks.
con \(i\) il numero delle intersezioni di una \(C\) e di una \(K\); invero l’intorno di \(O\) dovrà avere (come \(\overline{K}\)) \(i\) punti a comune con una sezione piana \(C\). Invece alla \(\overline{C}\), che si aggiunge alle \(K\) del fascio, corrisponderà la retta \(a\) multipla per \(F\), che sarà precisamente \(n\)-pla, designando con \(n\) il grado di \(|C|\). Il piano \(Oa\) segherà la superficie \(F\) fuori di \(a\) secondo \(i\) rette eccezionali, uscenti da \(O\). La superficie \(F\) sarà dunque d’ordine \(n+i\) e si vede che essa è l’immagine di un particolare sistema \(\infty^3\) contenuto in \(|C+K|\), che possiede \(i\) punti base nei punti comuni a \(C\) e \(K\): sono questi punti che danno origine alle rette eccezionali per \(O\) di cui s’è detto sopra.

Se ora si vuole costruire sulla \(F\) il sistema somma della rete \(|C|\) e del fascio \(|K|\), è chiaro che converrà ricorrere al sistema delle superficie che si ottengono sommando un piano per \(O\) e un piano per \(a\), e perciò al sistema lineare \(\infty^5\) delle quadriche per \(O\) ed \(a\): si otterrà così il minimo sistema somma \(|C+K|\).

\edsection{10. Caratteri virtuali.}{10. Caratteri virtuali.}{sec.I.10}

Siano \(|C_1|\) e \(|C_2|\) due sistemi lineari irriducibili sopra la superficie \(F\), e si designino rispettivamente con \(n_1\) ed \(n_2\) i loro gradi effettivi, con \(\pi_1\) e \(\pi_2\) i loro generi, e con \(i\) il numero delle intersezioni variabili di una \(C_1\) con una \(C_2\). Il \textit{grado del sistema} \(|C_1+C_2|\) sarà il numero delle intersezioni
\[
(C_1+C_2)^2 = C_1^2+C_2^2+2C_1C_2,
\]
e quindi varrà
\[
n = n_1+n_2+2i.
\]

Il \textit{genere del sistema somma} varrà
\[
\pi = \pi_1+\pi_2+i-1,
\]
tale essendo l’espressione del genere di una curva che viene a spezzarsi in due parti, espressione determinata dal \textsc{Noether} secondo l’esigenza della continuità\origfootnote{1}{s025R-1}{Cfr. \textsc{Enriques-Chisini}, Op. cit., Libro V, vol. III, pag. 399.}.

Queste formule conducono a definire \textit{il grado e il genere virtuale di una curva \(C\)} sopra \(F\), indipendentemente dalla possibilità di estendere il sistema lineare \(\infty^0\) da essa determinato. A tale scopo si associ alla \(C\) un sistema irriducibile \(|K|\) tale che il sistema \(|C+K|\) risulti irriducibile; ciò può farsi facilmente mandando per \(C\) le superficie di un ordine abbastanza alto e prendendo come curve \(K\) le curve residue. Ora designando con \(n\) il grado della \(C\) che si tratta
%CONT>>
% scan: s026L
% !PRINTED 32
\weberpage{32}
% head: CAPITOLO PRIMO
%<<CONT
% stamp: Dark scanning border along the right edge.
di definire, con \(m\) il grado di \(K\) e con \(i\) il numero delle intersezioni di una \(C\) con una \(K\), il grado di \(|C+K|\) verrà dato da
\[
N=n+m+2i,
\]
da cui si ricava
\[
n=N-m-2i.
\]

\textit{Il grado virtuale della \(C\)} così definito \textit{è indipendente dalla scelta del sistema ausiliario \(|K|\) che si è associato ad essa.} Invero sieno \(|K_1|\) e \(|K_2|\) due sistemi ausiliari soddisfacenti alle condizioni poste, coi gradi \(m_1\) ed \(m_2\), e siano \(i_1\) e \(i_2\) i numeri delle intersezioni \(CK_1\) e \(CK_2\), \(j\) il numero delle intersezioni variabili di una \(K_1\) ed una \(K_2\).

Si potranno confrontare il grado \(n_1\) definito rispetto a \(|K_1|\), il grado \(n_2\) definito rispetto a \(|K_2|\), valutando il grado del sistema
\[
|(C+K_1)+K_2|=|(C+K_2)+K_1|=|C+(K_1+K_2)|.
\]

Si avrà dunque
\[
\begin{aligned}
N&=(n_1+m_1+2i_1)+m_2+(2j+2i_2)\\
&=(n_2+m_2+2i_2)+m_1+(2j+2i_1),
\end{aligned}
\]
da cui segue
\[
n_1=n_2.
\]

In modo affatto simile si può definire il \textit{genere virtuale} \(\pi\) di una curva isolata \(C\), e questo carattere conserva un significato anche quando la \(C\) sia una curva multipla, o una curva riducibile contenente qualche parte multipla.

Aggiungiamo che la definizione dei \textit{caratteri} (grado e genere) virtuali si estende anche al caso di una \textit{curva virtuale} \(C-K\) definita come differenza indipendentemente dalla sua effettiva esistenza (§ 8). Basta all’uopo sommare a \(C-K\) una curva \(L\) per modo che \(|M|=|C-K+L|\) sia un sistema lineare irriducibile, e valutare quindi i caratteri della differenza \(|M-L|\), i quali risultano indipendenti da \(|L|\).

Ritornando alle curve effettive, notiamo ora che la nozione dei caratteri virtuali di una curva \(C\) dà luogo ad alcune osservazioni che tendono ad attribuire a codesti caratteri un significato invariante rispetto a trasformazioni birazionali della superficie.

Si trasformi la \(F\) in un’altra superficie \(\bar F\) in modo che ad un punto \(O\) della curva \(C\) corrisponda una curva eccezionale \(o\) di \(\bar F\); allora come trasformata della \(C\) si potrà riguardare la curva composta \(\bar C+o\) ovvero soltanto la \(\bar C\), tralasciando la \(o\) che corrisponde ad un punto fondamentale della trasformazione. La scelta che così può farsi risponde ad un diverso modo di considerare il punto \(O\) sopra
%CONT>>
% scan: s026R
% !PRINTED 33
\weberpage{33}
% head: SISTEMI LINEARI DI CURVE
%<<CONT
% stamp: Thin vertical scanning marks near the right edge and scattered specks.
la curva \(C\): se la \(C\) sia suscettibile di appartenere ad un sistema lineare irriducibile di dimensione \(>0\) col punto base \(O\), il trasformato di questo sistema sarà il sistema ugualmente irriducibile \(|\bar C|\), in cui la \(o\) non figura come parte fissa. Ma se invece \(|C|\) possa estendersi in guisa da non possedere più il punto base \(O\), il trasformato di esso conterrà in particolare il sistema \(o+|\bar C|\), dove la \(o\) figura come parte fissa. In altre parole la curva eccezionale \(o\) corrispondente ad \(O\) dovrà ritenersi o meno come componente della curva trasformata di \(C\), secondo che s’intenda di considerare la \(C\) come curva di un sistema lineare non avente o avente \(O\) come punto base. La circostanza che un sistema \(|C|\) senza un punto base \(O\) esista effettivamente può ritenersi come accidentale rispetto alla considerazione di cui sopra. Il \textit{grado virtuale} della \(C\), avuto riguardo alle possibili trasformazioni dei punti di essa in curve eccezionali, si definirà dunque \textit{in generale} supponendo che esso \textit{non} abbia \textit{punti base}, ovvero dichiarando quali \textit{punti base} s’intendono assegnati sopra la curva, all’infuori dei quali ogni altro punto di \(|C|\) che risulti \textit{punto base} per un \(|C|\) effettivamente costruito deve ritenersi \textit{virtualmente inesistente}. Aggiungasi che se la \(C\) possiede un punto multiplo (semplice per la superficie), di molteplicità \(\geq i\), questo può assumersi virtualmente come \(i\)-plo per la curva \(C\), anche se in effetto possegga una molteplicità superiore, e se pur questa risulti la molteplicità effettiva del più ampio sistema \(|C|\) cui sia imposto di possedere quel punto come \(i\)-plo.

In particolare dunque, se sopra la superficie \(F\) la curva \(C\) possiede un punto \(i\)-plo che non voglia ritenersi come punto base assegnato pel sistema \(|C|\), esso dovrà considerarsi \textit{virtualmente inesistente}, e perciò il \textit{genere virtuale} della \(C\) si otterrà aggiungendo al genere effettivo \(\pi\), \(\dfrac{i(i-1)}{2}\), che è precisamente il numero di cui diminuisce il genere di una curva irriducibile che venga ad acquistare un punto \(i\)-plo.

Similmente il \textit{grado virtuale} della \(C\) si valuterà aggiungendo \(i^2\) al numero effettivo delle intersezioni di due \(C\) fuori di questo, poichè appunto di \(i^2\) viene a diminuire il numero delle intersezioni di due curve quando esse acquistino un punto \(i\)-plo comune.

\edsection{11. Curve eccezionali.}{11. Curve eccezionali.}{sec.I.11}

Da ciò che si è detto innanzi, l’intorno di un punto (semplice) della superficie \(F\), quando sia punto base per un sistema lineare di curve, può ritenersi come una \textit{curva infinitesima} che, per una trasformazione della superficie si muta in una curva propriamente detta (curva eccezionale). Reciprocamente una curva eccezionale è una
% Printed signature line at the foot of the page.
% EDITORIAL: segnatura originale: \parbox{\textwidth}{\footnotesize\textsc{Enriques F.} - \textit{Superficie algebriche.}\hfill 3}
\printersignature{\textsc{Enriques F.} - \textit{Superficie algebriche.}\quad 3}
%CONT>>
% scan: s027L
% !PRINTED 34
\weberpage{34}
% head: CAPITOLO PRIMO
%<<CONT
% stamp: Dark scanning border at the right and a blurred, illegible patch in the lower-left margin.
curva tracciata sulla superficie \(F\) che possa mutarsi nell’intorno di un punto semplice.

Volendo esaminare più da vicino le curve eccezionali stabiliremo anzitutto questa distinzione:

una curva \textit{irriducibile} \(\omega\) sopra la superficie \(F\) si dirà \textit{curva eccezionale di \(1^{\mathrm a}\) specie}, se è suscettibile di trasformarsi in un punto semplice per mezzo di un sistema lineare \(|C|\) che non abbia punti base su \(\omega\). Ed invece si dirà \textit{curva eccezionale di \(2^{\mathrm a}\) specie} una curva \(\omega\) trasformabile in un punto mediante un sistema lineare \(|C|\) che abbia qualche punto base su di essa. % ?? uncertain: punctuation after the first "specie" -- comma or scan speck.

Esempio di curve eccezionali di prima specie sono le 27 rette appartenenti ad una superficie cubica generale \(F_3\). Infatti si può rappresentare la \(F_3\) sopra il piano in modo che ad una qualunque di codeste rette risponda l’intorno di uno dei 6 punti base del sistema di cubiche rappresentative. Invece nel piano non vi sono curve eccezionali di prima specie. Ma le rette, le coniche e in generale le curve piane fondamentali per una trasformazione cremoniana costituiscono curve eccezionali di seconda specie: si trasformerà una retta in un punto per mezzo di una rete omaloidica (per esempio di coniche) avente due punti base sopra la retta; e si trasformerà una conica in un punto assumendo per esempio una rete omaloidica di cubiche con un punto base doppio e quattro punti base semplici sopra la conica.

Già da questi esempi risulta che la distinzione fra le curve eccezionali di prima e di seconda specie ha soltanto carattere proiettivo o gode di una \textit{invarianza relativa} a trasformazioni prive di punti base: infatti una \textit{curva eccezionale di seconda specie} \(\omega\) appartenente ad una superficie \(F\) \textit{può sempre mutarsi in una curva eccezionale di prima specie} su una superficie trasformata \(\Phi\). A tal uopo si consideri il sistema lineare trasformante \(|C|\) che muta \(\omega\) in un punto; esso ha per ipotesi un certo numero \(s\) di punti base, con certe molteplicità \(i_1,i_2,\ldots,i_s\) sopra \(\omega\), e il numero delle intersezioni con \(\omega\) uguaglia precisamente la somma \(i_1+i_2+\ldots+i_s\). Ora, se si somma a \(|C|\) un sistema lineare \(|L|\) privo di punti base, e per cui la \(\omega\) non sia curva fondamentale, il sistema \(|C+L|\) varrà a trasformare \(\omega\) in una curva, la quale sarà fondamentale per il sistema trasformato di \(|C|\) (senza punti base sopra di esso) e risulterà quindi una curva eccezionale di prima specie.

Cerchiamo di valutare il genere e il grado di una curva eccezionale irriducibile di prima specie \(\omega\), riferendoci al punto semplice \(O\) il cui intorno corrisponde alla \(\omega\) sopra una superficie trasformata. Assunto su quella un sistema lineare qualsiasi \(|C|\) per cui \(O\) non sia punto base, l’imposizione del punto base \(O\) diminuisce di uno il grado di \(|C|\); inoltre le \(C\) per \(O\) hanno un punto a comune con
%CONT>>
% scan: s027R
% !PRINTED 35
\weberpage{35}
% head: SISTEMI LINEARI DI CURVE
%<<CONT
% stamp: Scattered scanning specks.
l’intorno di \(O\) stesso. Perciò designando con \(n\) e \(\pi\) il grado e il genere di \(|C|\) e con \(\nu\) e \(\varrho\) il grado e il genere della curva infinitesima che costituisce l’intorno di \(O\), ossia della curva eccezionale \(\omega\), avremo
\[
\begin{aligned}
n&=n-1+\nu+2\cdot1\\
\pi&=\pi+\varrho+1-1,
\end{aligned}
\]
e quindi
\[
\nu=-1\qquad\text{e}\qquad\varrho=0.
\]
La formula \(\varrho=0\) risulta del resto a priori essendo razionale l’intorno di \(O\).

Naturalmente lo stesso computo si può ripetere sulla superficie \(\Phi\), trasformata di \(F\), che contiene la \(\omega\). Qui si ha un sistema \(|C|\) per cui \(\omega\) è \textit{curva fondamentale} (senza intersezioni colle \(C\) variabili) e le \(C\) passanti per un punto di \(\omega\) si spezzano nella \(\omega\) e in una curva residua \(C_1\) unisecante la \(\omega\): scrivendo il grado e il genere di \(|C|=|\omega+C_1|\) si ritrovano le formule precedenti. % The equality sign is printed at both sides of the line break.

Enunciamo dunque che: \textit{una curva eccezionale irriducibile di \(1^{\mathrm a}\) specie ha il genere virtuale \(\varrho=0\) e il grado virtuale \(\nu=-1\).}

Passando ora alle curve eccezionali di \(2^{\mathrm a}\) specie, basta ricordare che una siffatta curva \(\omega\) si lascia trasformare in una curva eccezionale di \(1^{\mathrm a}\) specie, qualora si tolgano da essa dei punti da ritenere come punti base assegnati di un sistema trasformante \(|C|\); perciò \textit{il genere e il grado virtuali di una curva eccezionale di \(2^{\mathrm a}\) specie varranno}
\[
\varrho=0\qquad\text{e}\qquad\nu\geq0.
\]

Per esempio il grado delle generatrici d’una rigata vale
\[
\nu=0;
\]
il grado d’una retta del piano vale
\[
\nu=1,
\]
quello d’una conica
\[
\nu=4,\ \text{ecc.}
\]

Notiamo ora che l’esigenza che il \textit{genere} virtuale di una \textit{curva eccezionale sia} \(\varrho=0\), porta che una \textit{curva eccezionale irriducibile} di \(1^{\mathrm a}\) specie \(\omega\) \textit{non} può avere \textit{punti doppi} o multipli. Infatti il genere virtuale \(\varrho\) viene valutato calcolando il genere del sistema somma \(|C|=|C_1+\omega|\) e perciò un eventuale punto doppio della componente \(\omega\), non essendo punto base per \(|C|\), deve ritenersi come inesistente, ed allora, essendo già \(\varrho=0\), la curva \(\omega\) risulterebbe riducibile.

Anticipando alcune nozioni che saranno acquisite nel seguito vogliamo dire qui che le curve eccezionali vengono definite dai loro
%CONT>>
% scan: s028L
% !PRINTED 36
\weberpage{36}
% head: CAPITOLO PRIMO
%<<CONT
% stamp: Heavy black scanning borders along the bottom and right edges.
caratteri virtuali trovati innanzi. Se \(\omega\) è una curva irriducibile di genere \(\varrho=0\) e di grado \(\nu=-1\) sopra la superficie \(F\), essa potrà trasformarsi in un punto semplice, operando come segue. Si assuma sopra \(F\), in modo affatto generale, un sistema lineare \(\infty^3|C|\) privo di punti base su \(F\), di un certo genere \(\pi\) e di un certo grado \(n\), le cui curve seghino in \(i\) punti la \(\omega\). Sommando \(\omega\) a \(|C|\) si ottiene un sistema \(|C+\omega|\) di curve \((i-1)\) secanti la \(\omega\), che sarà di grado \(n+2i-1\) e di genere \(\pi+i-1\), di modo che la differenza fra il grado e il doppio del genere \(n-2\pi\), viene aumentata di una unità. E, come riconosceremo nel seguito, questo aumento ha un significato effettivo, perchè l’addizione di \(\omega\) riesce effettivamente ad estendere il sistema \(|C|\), risultando la dimensione del sistema irriducibile \(|C+\omega|\) maggiore di quella di \(|C|\). Ciò posto, si può procedere sommando ancora al sistema ottenuto la \(\omega\) e così di seguito, fino a che non si arrivi ad un sistema \(|C+i\omega|\) privo di punti base su \(F\) per cui la \(\omega\) riesce fondamentale. Mediante questo sistema la superficie \(F\) si lascia trasformare in una \(\Phi\) in cui ad \(\omega\) corrisponde l’intorno di un punto semplice \(O\). In conclusione la \(\omega\) è una curva eccezionale di \(1^{\mathrm a}\) specie.

Il procedimento di trasformazione qui adoperato conduce al risultato seguente:

se la superficie \(F\) contiene un certo numero \(s\) di curve eccezionali che seghino rispettivamente le \(C\) in \(i_1,i_2,\ldots,i_s\) punti, dove sia
\[
i_1+i_2+\ldots+i_s>2\pi-2-n
\]
si costruirà sopra \(F\) un sistema lineare irriducibile di grado \(N\) e di genere \(\Pi\) per cui
\[
N>2\Pi-2.
\]

Vedremo nel seguito che l’esistenza di sistemi i cui caratteri danno luogo ad una tale disuguaglianza vale a definire una particolare famiglia di superficie, che anzi essa caratterizza la famiglia delle superficie razionali e di quelle riferibili a rigate; sopra una superficie che non appartenga a codesta famiglia particolare, per ogni sistema lineare di genere \(\pi\) e di grado \(n\) sussiste la disuguaglianza fondamentale
\[
n\leq2\pi-2.
\]

Conseguenza delle cose dette sarà che: \textit{una superficie non appartenente alla famiglia delle rigate potrà sempre trasformarsi in un’altra priva affatto di curve eccezionali.}

In pari tempo si riconoscerà che « una superficie che contenga curve eccezionali di \(2^{\mathrm a}\) specie appartiene alla famiglia delle rigate ».

Infatti si trasformi dapprima, come si è detto, una curva ec%
%CONT>>
% scan: s028R
% !PRINTED 37
\unskip\weberpage{37}%%
% head: SISTEMI LINEARI DI CURVE
%<<CONT
% stamp: Heavy black scanning borders along the bottom and left edges.
cezionale di \(2^{\mathrm a}\) specie in una curva \(\omega\) di genere \(\varrho=0\) e di grado \(\nu\geq0\). Qui, per essere \(\nu\geq0\), il procedimento di addizione della \(\omega\) ad un sistema \(|C|\) non ha mai termine perchè cresce, o almeno non diminuisce, il numero delle intersezioni delle curve del sistema con la \(\omega\), sicchè certo si arriverà ad un sistema lineare di genere \(\Pi\) e di grado \(N\), per cui
\[
N>2\Pi-2.
\]

\edsection{12. Nota sulle curve eccezionali riducibili.}{12. Nota sulle curve eccezionali riducibili.}{sec.I.12}

Quando si trasforma una superficie \(F\) in un’altra \(\Phi\) per mezzo di un sistema lineare \(|C|\) che abbia su \(F\) dei punti base distinti \(O,O_1,O_2,\ldots\) le curve eccezionali che rispondono a questi punti, e che sono in pari tempo fondamentali per lo stesso sistema \(|L|\) (trasformato delle sezioni piane di \(F\)), sono curve irriducibili fra loro sconnesse. Ma che cosa accade quando i punti \(O,O_1,O_2,\ldots\) diventano infinitamente vicini?

Pongasi che i detti punti \(O,O_1,O_2,\ldots\) abbiano per \(|C|\) certe molteplicità \(i,i_1,i_2,\ldots\) dove sia \(i>i_1>i_2\ldots\) e le differenze \(i-i_1,\ i_1-i_2,\ldots\) sieno grandi quanto occorre. A codesti punti rispondono in generale su \(\Phi\) delle curve eccezionali, \(\Omega,\Omega_1,\Omega_2,\ldots\) risp. degli ordini \(i,i_1,i_2,\ldots\). Ma se il punto \(O_1\) diventa infinitamente vicino ad \(O\), la curva \(\Omega\) d’ordine \(i\) viene a spezzarsi in due parti \(\omega\) e \(\omega_1(=\Omega_1)\), la prima, d’ordine \(i-i_1\), corrispondente all’intorno del punto \(O\) privato di \(O_1\) e la seconda, d’ordine \(i_1\), limite della curva eccezionale corrispondente ad \(O_1\).

In generale se più punti \(O_1O_2\ldots\) diventano infinitamente vicini ad un punto proprio \(O\), succedendosi su uno o più rami, lineari o superlineari, la curva eccezionale che risponde al punto \(O\) si spezzerà in più parti fra loro connesse, comprendenti le curve eccezionali che nascono dai punti successivi. Si può precisare la cosa riferendoci al caso caratteristico in cui \(O,O_1,O_2,\ldots\) si succedano sopra un unico ramo, e incominciando dagli esempi più semplici. Se \(O,O_1,O_2,\ldots,O_s\) si succedono sopra un ramo lineare, la curva eccezionale d’ordine \(i\) che risponde ad \(O\) diventa una curva composta d’ordine \((i-i_1)+(i_1-i_2)+\ldots+(i_{s-1}-i_s)+i_s\):
\[
\Omega=\omega+\omega_1+\omega_2+\ldots+\omega_s,
\]
mentre le curve nascenti da \(O_1O_2\ldots\) diventano
\[
\begin{aligned}
\Omega_1&=\omega_1+\omega_2+\ldots+\omega_s\\
\ldotp\quad&\ldotp\quad\ldotp\quad\ldotp\quad\ldotp\quad\ldotp\quad\ldotp\quad\ldotp\quad\ldotp\\
\Omega_{s-1}&=\omega_{s-1}+\omega_s\\
\Omega_s&=\omega_s.
\end{aligned}
\]
%CONT>>
% scan: s029L
% !PRINTED 38
\weberpage{38}
% head: CAPITOLO PRIMO
%<<CONT
% stamp: Dark scanning border along the right edge.
E, secondo il principio di continuità, queste curve composte dovranno ritenersi come curve eccezionali, laddove le parti di esse (ad es. \(\omega_1\omega_2\ldots\omega_{s-1}\)) non costituiranno \textit{intere} curve eccezionali, bensì soltanto \textit{componenti} di esse. Invero la componente \(\omega\) si troverà in corrispondenza biunivoca, non già coll’intorno di \(O\) (preso nella sua interezza), bensì con quest’intorno privato del punto \(O_1\), e così \(\omega_1\) si troverà in corrispondenza coll’intorno di \(O_1\) privato del punto \(O_2\), ecc. Potremo passare dalla superficie \(F\) alla \(\Phi\) mediante più trasformazioni successive da \(F\) a \(F_1\), da \(F_1\) a \(F_2\), \(\ldots\), da \(F_{s-1}\) a \(F_s=\Phi\); la prima di queste trasformazioni muterà \(O\) in una curva \(\omega\), la seconda muterà un punto \(O_1\) di \(\omega\) in una curva \(\omega_1\), la terza muterà un punto \(O_2\) di \(\omega_1\) (diverso dal punto comune ad \(\omega\) e \(\omega_1\)) in una curva \(\omega_2\), ecc. Perciò ciascuna \(\omega_i\) (per \(i<s\)) si troverà connessa in un punto colla successiva \(\omega_{i+1}\) e il suo grado (che è \(-1\) sopra \(F_i\)) diventerà \(-1-1=-2\) sopra \(F_{i+1}\) e su \(\Phi\). Per conseguenza la \textit{curva eccezionale composta} \(\Omega=\omega+\omega_1+\ldots+\omega_s\), corrispondente ad una serie di \textit{punti} \(O\,O_1\ldots O_s\) \textit{infinitamente vicini, succedentisi su un ramo lineare}, avrà il \textit{genere}
\[
\varrho=0+0+\ldots+0+(s-1)-(s-1)=0
\]
e il \textit{grado}
\[
\nu=-2-2\ldots-2-1+2(s-1)=-1.
\]
Insomma si ritrovano così i caratteri delle curve eccezionali irriducibili di cui le curve composte \(\Omega,\Omega_1,\Omega_2,\ldots,\Omega_s\) appariscono come curve limiti.

Questa coincidenza di caratteri dovrà prodursi ancora quando i punti \(O,O_1,O_2,\ldots,O_s\) vengano a succedersi sopra un ramo superlineare, ma nella definizione delle curve limiti andiamo incontro ad una sorpresa: le componenti delle curve \(\Omega_i\) diventano in generale curve multiple!

Riferiamoci, per semplicità, al caso di tre punti \(O,O_1,O_2\), succedentisi sopra un ramo cuspidale: questo caso si può trattare ritenendo il ramo cuspidale come limite di un ramo lineare, dove \(O_2\), variabile nell’intorno di \(O_1\), assuma la posizione del suo punto satellite. Qui appare che la curva d’ordine \(i-i_1\), trasformata dell’intorno di \(O\) (privato di \(O_1\)), viene a ridursi ad una curva \(\omega\) d’ordine \(i-i_1-i_2\) e alla curva \(\omega_2\), d’ordine \(i_2\), rispondente all’intorno di \(O_2\); per conseguenza quest’ultima curva viene a comparire due volte in
\[
\omega+\omega_1+2\omega_2.
\]

Per comprendere più chiaramente le ragioni di ciò che qui accade si effettuino tre trasformazioni successive della superficie \(F\),
%CONT>>
% scan: s029R
% !PRINTED 39
\weberpage{39}
% head: SISTEMI LINEARI DI CURVE
%<<CONT
% stamp: Thin vertical mark in the right margin; a hook-shaped non-typographic mark above the right end of the displayed sum; a small blot near "cui".
% EDITORIAL: numero originale fissato esplicitamente.

mutando anzitutto \(O\) in \(\omega\), e poi ancora un punto \(O_1\) di \(\omega\) in \(\omega_1\), e finalmente il punto \(O_2\) comune ad \(\omega\) e \(\omega_1\) in \(\omega_2\): la posizione del punto \(O_2\) così definita risponde al punto satellite di \(O_1\) sopra la superficie \(F\). Ma poichè \(O_2\) è punto doppio per la curva \(\omega+\omega_1\), la curva eccezionale trasformata \(\omega_1\) dovrà contarsi come doppia, in accordo a ciò che si è detto innanzi. % [sic] \omega_1 -> \omega_2, the immediately preceding clause names the transformed curve \omega_2.
Le molteplicità delle curve eccezionali composte che rispondono ai punti successivi di un ramo superlineare sono state determinate da \textsc{E. Barber} e \textsc{O. Zariski}\origfootnote{1}{s029R-1}{\textsc{Barber} e \textsc{Zariski}, l. c. Ann. Journ. of Math., 1935.}, rilevando anche l’ordine delle connessioni di queste curve e illustrando così, sotto un nuovo aspetto, la teoria delle singolarità di \textsc{Enriques}\origfootmark{2}{s029R-2}. % ?? uncertain: footnote marker after Enriques -- (3) or (2).
% EDITORIAL D0022 RISOLTO: The enlarged printed footnote marker after ENRIQUES is 2. The source has [3], so this is a transcription correction; the printed reference itself appears erroneous because note 2 explains a different passage.
Noi ci limiteremo a riportare la regola generale seguente: \textit{la curva eccezionale composta \(\Omega\), che risponde ad una serie di punti \(O_1,\ldots,O_s\) infinitamente vicini ad un punto proprio \(O\), succedentisi sopra un ramo superlineare, comprende \(s+1\) componenti, e la componente \(\omega_i^{h_i}\), che occupa il posto \(i\), vi figura colla molteplicità \(h_i\) eguale all’ordine minimo del ramo uscente da \(O\)\origfootnote{2}{s029R-2}{Quando si parla di un ramo uscente dal punto \(O_1\) (ovvero da \(O_2\) ecc.) s’intende di aver prima eseguita la trasformazione che muta \(O_1\) in un punto proprio.} che contiene gli \(i\) punti \(O_1,O_2,\ldots,O_i\).} % ?? uncertain: superscript attached to \omega_i -- h_i or h.
% EDITORIAL D0023 RISOLTO: The enlargement shows h with a small i attached at its lower right, while a separate subscript i sits below omega.

La stessa regola si applica alle curve eccezionali composte \(\Omega_1,\Omega_2,\ldots,\Omega_s\)\origfootnote{3}{s029R-3}{Cfr. \textsc{Enriques-Chisini}, \textit{Lezioni ecc.}, Libro I, vol. II.}. Inoltre la curva \(\omega_i\) viene \textit{connessa}, non sempre alla successiva \(\omega_{i+1}\), bensì --- in generale --- alla curva \(\omega_{i+t}\) \textit{che le succede dopo \(t\) posti}, in corrispondenza alla circostanza che \textit{i punti \(O_{i+1}\ldots O_{i+t}\) siano prossimi ad \(O_i\)} (come risulta anche evidente dall’ispezione dello schema grafico del ramo \(O_{i+1}\ldots O_{i+t}\)). Aggiungasi che la detta curva \(\omega_i\) riesce di \textit{grado} \(-1-t\), come appare da ciò che essa rappresenta l’intorno di un punto \(O_i\) da cui sono tolti \(t\) punti prossimi. Tenuto conto delle molteplicità che spettano alle curve \(\omega_i\) in \(\Omega_i\), le curve eccezionali composte \(\Omega,\Omega_1,\ldots,\Omega_s\) risultano sempre di grado \(\nu=-1\) e di genere \(\varrho=0\) come le curve eccezionali irriducibili di cui sono limiti.

Infine rileveremo che l’insieme delle curve eccezionali irriducibili sconnesse fra loro, di cui qui si considera il limite in seguito all’avvicinarsi di \(O_1,O_2,\ldots\) ad \(O\), tende ad una curva
\[
\Omega+\Omega_1+\ldots+\Omega_s
\]
che risulta composta con-le \(\omega,\omega_1,\omega_2,\ldots,\omega_s\), contate ciascuna con una molteplicità che è la somma delle molteplicità che essa possiede per riguardo alle \(\Omega_i\) di cui fa parte. Vedremo più avanti il significato di quest’osservazione. % ?? uncertain: mark between "con" and "le" -- printed hyphen or scanning blemish.

% scan: s030L
% !PRINTED 40
\weberpage{40}
% head: CAPITOLO PRIMO
% stamp: Dark scanning border along the right edge and scattered specks.
% EDITORIAL: numero originale fissato esplicitamente.

\edsection{13. Serie caratteristica virtuale.}{13. Serie caratteristica virtuale.}{sec.I.13}

Abbiamo spiegato come si determinino i caratteri virtuali, e in particolare il grado \(n\) di una curva \(C\) tracciata sopra la superficie \(F\) anche se questa sia una \textit{curva isolata}, non contenuta in un sistema \(|C|\) di dimensione \(r>0\). Il procedimento che conduce a valutare \(n\)\origfootnote{1}{s030L-1}{Cfr. \textsc{F. Enriques}, \textit{Introduzione ecc.}, 1896: nn. 15 e 16.}, conduce altresì a definire, per \(n>0\), la serie caratteristica di \(C\), anche indipendentemente dalla circostanza che essa appartenga ad un sistema più ampio\origfootnote{2}{s030L-2}{\textsc{F. Severi}, \textit{Osservazioni sui sistemi continui di curve.} Atti Acc. Torino, 1904.}.

Infatti si associ a \(C\) un sistema lineare irriducibile \(|K|\), tale che \(|C+K|\) riesca irriducibile: la serie caratteristica di \(C\) verrà definita dalle curve di \(|C+K|\) che passano per le intersezioni di una \(C\) con una \(K\). Questa serie invero riesce indipendente dalla scelta di \(|K|\), come si vede paragonando le serie definite in relazione a due sistemi \(|C+K|\) e \(|C+L|\) ed al sistema \(|C+K+L|\). Si noti che, come avremo occasione di spiegare nel seguito, sarà sempre lecito scegliere \(|K|\) in modo che le curve di \(|C+K|\) seghino sopra \(C\) una serie completa, e così si ottenga sulla \(C\) la intera \textit{serie caratteristica virtuale}.

L’importanza di questa serie si manifesta sopratutto nello studio dei sistemi continui di curve \(C\) non costituiti di curve equivalenti; allora la serie caratteristica virtuale della \(C\) viene segata (tutta o in parte) dalle curve del sistema ad essa infinitamente vicine. % ?? uncertain: spelling of "sopratutto" -- sopratutto or soprattutto.
% EDITORIAL D0025 RISOLTO: The enlarged word clearly has a single t after sopra and tt after u; the historical spelling is retained.

\textsc{Nota.} -- In questa maniera la serie caratteristica di una \(C\) appartenente ad un sistema continuo si è presentata implicitamente a \textsc{G. Castelnuovo}, siccome appare da una sua comunicazione epistolare ad \textsc{Enriques}\origfootnote{3}{s030L-3}{Cfr. Rendic. Circolo Matematico di Palermo, 25 dic. 1898 (n. 4).}. In seguito il \textsc{Castelnuovo} stesso ebbe a rilevare esplicitamente il significato di essa, in rapporto alla superficie (iperellittica) che rappresenta le coppie di punti di una curva del genere due, e attrasse su tale esempio l’attenzione di \textsc{Severi}; che fu indotto quindi alla definizione esposta innanzi, siccome egli stesso accenna in una nota a piè della prima pagina del suo lavoro. % ?? uncertain: final punctuation after "lavoro" -- period or short dash.
% EDITORIAL D0026 RISOLTO: The terminal mark after lavoro lies on the baseline and is consistent with the slightly elongated full stops elsewhere on the scan.

\clearpage\thispagestyle{plain}
% scan: s030R
% !PRINTED none
\weberpage{[s030R]}
% head: none
% stamp: Scattered scanning specks.
% EDITORIAL: numero originale fissato esplicitamente.

\edchapter{II}{SISTEMI COVARIANTI E INVARIANTI}{Capitolo II. SISTEMI COVARIANTI E INVARIANTI}{SISTEMI COVARIANTI E INVARIANTI}

\edsection{1. Curve jacobiane.}{1. Curve jacobiane.}{sec.II.1}

Sopra la superficie \(F\), che possiamo supporre priva di qualsiasi singolarità in un iperspazio, consideriamo un sistema lineare irriducibile \(|C|\) di dimensione 2, cioè una \textit{rete}. Il luogo dei punti della superficie che sono doppi per una curva della rete, è una linea algebrica \(C_j\), che diremo \textit{jacobiana} della rete. La \(C_j\) si può anche definire come luogo dei punti di contatto di due curve della rete; infatti il fascio da queste determinato contiene una curva che ha un punto doppio nel punto di contatto. Nel piano la jacobiana di una rete di curve dell’ordine \(n\ (>1)\) è una curva dell’ordine \(3n-3\); % ?? uncertain: small mark after n in 3n-3 -- scanning speck or prime (3n'-3).
ed un punto base \(i\)-plo della rete ha in generale per questa la molteplicità \(3i-1\); così un punto base semplice \(O\) è, in generale, doppio per la jacobiana, ma se le curve della rete hanno in \(O\) una tangente fissa la jacobiana ha in essa un punto triplo; con quella tangente fissa\origfootnote{1}{s030R-1}{Cfr. \textsc{Enriques-Chisini}, Op. cit., Libro III, vol. II, cap. I, § 5.}. % ?? uncertain: punctuation after "triplo" -- semicolon or comma.
Queste proprietà si conservano inalterate anche per le reti date sopra una qualunque superficie \(F\), giacchè la determinazione delle molteplicità d’un punto base \(O\) per la jacobiana di una rete \(|C|\) dipende soltanto dai caratteri differenziali delle curve \(C\) considerate e della superficie che le contiene, la quale, in un ordine d’approssimazione grande quanto si vuole, si può approssimare con un paraboloide osculatore nel punto \(O\), che a sua volta si lascia rappresentare punto per punto sul piano. Per le curve jacobiane delle reti estratte da un più ampio sistema lineare \(|C|\) sussiste il seguente teorema:

\textit{Le jacobiane delle reti contenute in un medesimo sistema lineare irriducibile \(|C|\) di dimensione \(r>2\) sono equivalenti.}

Estragghiamo da \(|C|\) due reti qualsiasi e confrontiamo le loro jacobiane, supponendo anzitutto che le due reti abbiano a comune
%CONT>>
% scan: s031L
% !PRINTED 42
\weberpage{42}
% head: CAPITOLO SECONDO
%<<CONT
% stamp: Dark scan-edge band on the right and scattered scan specks.
un fascio, e perciò appartengano insieme ad un sistema lineare \(\infty^3\)\origfootnote{1}{s031L-1}{Questo sistema viene determinato dal fascio comune e da due curve prese rispettivamente nelle due reti, fuori di esso.}. Un punto generico \(P\) della superficie \(F\) è doppio per una curva \(\overline{C}\) del sistema \(\infty^3\) anzidetto, e la \(\overline{C}\), insieme col fascio comune alle due reti, determina una rete di curve \(C\) la cui jacobiana passa per \(P\). Si vede in tal guisa che le jacobiane delle \(\infty^1\) reti di curve \(C\) contenute nel sistema \(\infty^3\) ed aventi un fascio comune, formano un fascio (per ogni punto di \(F\) ne passa una); ma questo fascio è certo lineare perchè è razionale, trovandosi in corrispondenza biunivoca col fascio delle reti di un sistema lineare \(\infty^3\) che hanno a comune un fascio\origfootnote{2}{s031L-2}{Insieme da assimilarsi ad un fascio di piani dello \(S_3\).}. Ora se si assumono entro \(|C|\) (supposto di dimensione \(r>3\)) due reti di curve aventi a comune una curva \(\overline{C}\), si può costruire una rete ausiliaria che contenga \(\overline{C}\) ed abbia un fascio a comune con ciascuna delle due reti: la jacobiana della rete ausiliaria risulterà equivalente alla jacobiana delle due reti; e, per la proprietà transitiva dell’equivalenza, queste saranno dunque equivalenti fra loro.

Infine se si hanno entro \(|C|\) due reti senza curva comune (il sistema \(|C|\) avendo la dimensione \(r>4\)), si potranno paragonare le loro jacobiane costruendo una rete ausiliaria che abbia a comune una curva con ciascuna di esse. Le jacobiane delle due reti risulteranno quindi equivalenti, c. d. d.

\edsection{2. Sistema jacobiano.}{2. Sistema jacobiano.}{sec.II.2}

Da ciò che si è dimostrato innanzi segue che:

\textit{Le jacobiane delle reti di curve contenute in un sistema lineare irriducibile \(|C|\) sopra la superficie \(F\) appartengono ad un medesimo sistema lineare completo che si dirà il sistema jacobiano \(|C_j|\) di \(|C|\).} La dimensione di \(|C|\) viene supposta per ora \(\geq 2\).

La definizione che precede è senz’altro perfetta quando il sistema \(|C|\) non abbia punti base, dove s’intenda che il sistema delle jacobiane \(C_j\) dovrà completarsi senza l’imposizione di alcun punto base. Se invece \(|C|\) abbia un \textit{punto base \(i\)-plo}, s’intenderà che questo venga assegnato come \textit{punto di molteplicità \(3i-1\)} al sistema completo \(|C_j|\), poichè, come abbiamo detto, \(3i-1\) è in generale la molteplicità della jacobiana di una rete contenuta in \(|C|\).

Ciò vale anche se i punti base sono infinitamente vicini, così, per esempio, se \(|C|\) possiede un punto base semplice con la tangente fissa, cioè due punti base semplici \(O,O_1\), infinitamente vicini,
%CONT>>
% scan: s031R
% !PRINTED 43
\weberpage{43}
% head: SISTEMI COVARIANTI E INVARIANTI
%<<CONT
% stamp: Scattered scan specks.
le jacobiane delle reti contenute in esso avranno effettivamente in \(O\) un punto triplo con la tangente fissa \(OO_1\)\origfootnote{1}{s031R-1}{Cfr. \S~1.}, ma il sistema di queste curve sarà contenuto generalmente in un sistema più ampio con due punti base doppi \(O\) e \(O_1\), che costituirà il nostro sistema jacobiano completo \(|C_j|\). Se anche l’ampliamento non sia effettivamente possibile, la singolarità costituita da un punto triplo e da un punto semplice infinitamente vicino dovrà ritenersi virtualmente come costituita dai due punti doppi \(O\) ed \(O_1\). Abbiamo supposto sin qui che il sistema \(|C|\) fosse irriducibile; ma è chiaro come la definizione debba estendersi al caso in cui \(|C|\) contenga una parte fissa \(K\) (che supporremo semplice). Qui ogni punto della \(K\) appare come un punto base con tangente fissa per una rete qualunque contenuta in \(|C|\), perciò la curva \(K\) si staccherà 3 volte dalle jacobiane di tutte le reti contenute in \(|C|\). Essa dovrà aggiungersi 3 volte a codeste jacobiane e il sistema riducibile così definito sarà contenuto in generale nel più ampio sistema irriducibile completo \(|C_j|\). % ?? uncertain: The punctuation after "irriducibile" -- semicolon / comma.

A dir vero viene così dimostrato che la \(K\) si stacca \textit{almeno} tre volte dalle jacobiane delle reti di \(|C|\). Per riconoscere che di fatto non si stacca più di tre volte, si potrà ricondurre la cosa dalla superficie al piano sostituendo ad \(F\) un conveniente paraboloide osculatore in un punto della curva \(K\) o anche una superficie razionale che abbia un contatto assai elevato con la \(F\) lungo tutta la curva \(K\)\origfootnote{2}{s031R-2}{Cfr. \textsc{Enriques-Chisini}, Libro XV, cap. IV, \S~39, vol. II, pag. 634 e seg.}. Si riconosce così che una curva \(K\) la quale si aggiunga come \textit{componente fissa} di una \textit{rete} di curve irriducibili \(C\), figura \textit{tre volte nella jacobiana} di questa rete; essa non può comparirvi più di 3 volte, a meno che non sia una parte della jacobiana di \(|C|\).

\textsc{Nota.} -- Si può dare una verifica analitica diretta del teorema a cui siamo pervenuti. Per far ciò occorre sapere come un sistema lineare completo venga segato sopra una superficie \(F\) (con singolarità normali) dello spazio ordinario, cioè che « le superficie \(\varphi\) aggiunte ad \(f\), ossia passanti semplicemente per la curva doppia di \(f\), segano su \(f\) sistemi completi »\origfootnote{3}{s031R-3}{Cfr. la nota al \S~4 e il cap. III.}. Per semplicità di discorso supporremo che \(f\) sia una superficie d’ordine \(n\), affatto priva di singolarità, cosichè le intersezioni complete di essa con la totalità delle superficie di un dato ordine \(m\) formino un sistema completo.

Data su \(f\) una rete di curve irriducibili \(|C|\) che non siano intersezioni complete, a cui si aggiunga una componente fissa \(K\), si può sempre supporre che la rete \(|C|\) sia residua di una curva \(L\) non contenente la \(K\) come parte, in modo che le \(C+L\) siano intersezioni
%CONT>>
% scan: s032L
% !PRINTED 44
\weberpage{44}
% head: CAPITOLO SECONDO. % ?? uncertain: The final dot -- printed period / scan speck.
%<<CONT
% stamp: Dark scan-edge band on the right and scattered scan specks.
complete con superficie \(\varphi_i\) d’ordine \(m\). Quindi la rete \(|K+C|\) è contenuta parzialmente nella rete rappresentata da un’equazione del tipo:
\[
\varphi(x_1x_2x_3x_4)[\lambda_1\varphi_1(x_1x_2x_3x_4)+\lambda_2\varphi_2(x_1x_2x_3x_4)+\lambda_3\varphi_3(x_1x_2x_3x_4)]=0,
\]
dove secondo le nostre ipotesi le \(\varphi_1,\varphi_2,\varphi_3\), non si annulleranno simultaneamente sulla \(K\), che invece fa parte della \(\varphi=0\) il cui ordine verrà indicato con \(s\). Ora si tratta di scrivere l’equazione della jacobiana della rete così rappresentata, e di riconoscere che la curva \(\varphi=0\), ed in particolare la \(K\) che ne costituisce una parte, figura in essa 3 volte come componente. I punti della jacobiana sono punti di contatto di una superficie del sistema segante la rete, e della \(f\). Scrivendo le condizioni per il contatto si ottiene un sistema di equazioni lineari, per la cui compatibilità deve annullarsi il determinante
\[
D(x_1x_2x_3x_4)=
\begin{vmatrix}
\dfrac{\partial(\varphi\varphi_1)}{\partial x_1} &
\dfrac{\partial(\varphi\varphi_2)}{\partial x_1} &
\dfrac{\partial(\varphi\varphi_3)}{\partial x_1} &
\dfrac{\partial f}{\partial x_1} \\[6pt]
\dfrac{\partial(\varphi\varphi_1)}{\partial x_2} &
\dfrac{\partial(\varphi\varphi_2)}{\partial x_2} &
\dfrac{\partial(\varphi\varphi_3)}{\partial x_2} &
\dfrac{\partial f}{\partial x_2} \\[6pt]
\dfrac{\partial(\varphi\varphi_1)}{\partial x_3} &
\dfrac{\partial(\varphi\varphi_2)}{\partial x_3} &
\dfrac{\partial(\varphi\varphi_3)}{\partial x_3} &
\dfrac{\partial f}{\partial x_3} \\[6pt]
\dfrac{\partial(\varphi\varphi_1)}{\partial x_4} &
\dfrac{\partial(\varphi\varphi_2)}{\partial x_4} &
\dfrac{(\partial\varphi\varphi_3)}{\partial x_4} &
\dfrac{\partial f}{\partial x_4} % [sic] (\partial\varphi\varphi_3) -> \partial(\varphi\varphi_3), the opening parenthesis is printed before the derivative symbol.
\end{vmatrix}
=0
\]

Ora dopo avere eseguite le derivate di \(\varphi\varphi_1,\varphi\varphi_2,\varphi\varphi_3\), si è condotti a scomporre il determinante \(D\) in una somma di determinanti: tralasciando quei determinanti che risultino identicamente nulli, gli altri presentano il fattore comune \(\varphi^2\), tolto il quale rimangono quattro determinanti di cui uno contiene ancora il fattore \(\varphi\). A noi occorre dimostrare che la somma degli altri 3 determinanti si annulla (semplicemente) nei punti della curva comune \(\varphi=0\). Infatti questa somma si trasforma facilmente in un determinante unico. Scriviamo:
\begin{gather*}
\Phi=\varphi_1
\begin{vmatrix}
\dfrac{\partial\varphi}{\partial x_1} &
\dfrac{\partial\varphi_2}{\partial x_1} &
\dfrac{\partial\varphi_3}{\partial x_1} &
\dfrac{\partial f}{\partial x_1} \\[6pt]
\dfrac{\partial\varphi}{\partial x_2} &
\dfrac{\partial\varphi_2}{\partial x_2} &
\dfrac{\partial\varphi_3}{\partial x_2} &
\dfrac{\partial f}{\partial x_2} \\[6pt]
\dfrac{\partial\varphi}{\partial x_3} &
\dfrac{\partial\varphi_2}{\partial x_3} &
\dfrac{\partial\varphi_3}{\partial x_3} &
\dfrac{\partial f}{\partial x_3} \\[6pt]
\dfrac{\partial\varphi}{\partial x_4} &
\dfrac{\partial\varphi_2}{\partial x_4} &
\dfrac{\partial\varphi_3}{\partial x_4} &
\dfrac{\partial f}{\partial x_4}
\end{vmatrix}
\\[12pt]
%CONT>>
% scan: s032R
% !PRINTED 45
\text{\weberpage{45}}
% head: SISTEMI COVARIANTI E INVARIANTI
%<<CONT
% stamp: Short stray horizontal mark beside the running head, scattered specks, and a dark mark along the lower-left edge.
+\varphi_2
\begin{vmatrix}
\dfrac{\partial\varphi_1}{\partial x_1} &
\dfrac{\partial\varphi}{\partial x_1} &
\dfrac{\partial\varphi_3}{\partial x_1} &
\dfrac{\partial f}{\partial x_1} \\[6pt]
\dfrac{\partial\varphi_1}{\partial x_2} &
\dfrac{\partial\varphi}{\partial x_2} &
\dfrac{\partial\varphi_3}{\partial x_2} &
\dfrac{\partial f}{\partial x_2} \\[6pt]
\dfrac{\partial\varphi_1}{\partial x_3} &
\dfrac{\partial\varphi}{\partial x_3} &
\dfrac{\partial\varphi_3}{\partial x_3} &
\dfrac{\partial f}{\partial x_3} \\[6pt]
\dfrac{\partial\varphi_1}{\partial x_4} &
\dfrac{\partial\varphi}{\partial x_4} &
\dfrac{\partial\varphi_3}{\partial x_4} &
\dfrac{\partial f}{\partial x_4}
\end{vmatrix}
\\[12pt]
+\varphi_3
\begin{vmatrix}
\dfrac{\partial\varphi_1}{\partial x_1} &
\dfrac{\partial\varphi_2}{\partial x_1} &
\dfrac{\partial\varphi}{\partial x_1} &
\dfrac{\partial f}{\partial x_1} \\[6pt]
\dfrac{\partial\varphi_1}{\partial x_2} &
\dfrac{\partial\varphi_2}{\partial x_2} &
\dfrac{\partial\varphi}{\partial x_2} &
\dfrac{\partial f}{\partial x_2} \\[6pt]
\dfrac{\partial\varphi_1}{\partial x_3} &
\dfrac{\partial\varphi_2}{\partial x_3} &
\dfrac{\partial\varphi}{\partial x_3} &
\dfrac{\partial f}{\partial x_3} \\[6pt]
\dfrac{\partial\varphi_1}{\partial x_4} &
\dfrac{\partial\varphi_2}{\partial x_4} &
\dfrac{\partial\varphi}{\partial x_4} &
\dfrac{\partial f}{\partial x_4}
\end{vmatrix};
\end{gather*}\edFlush
scambiando nel primo determinante la \(1^{\mathrm a}\) colonna con la \(3^{\mathrm a}\) e nel secondo la \(2^{\mathrm a}\) con la \(3^{\mathrm a}\), tenuto conto dei segni, appare qui lo sviluppo del determinante del \(5^{\mathrm o}\) ordine:
\[
\Phi=
\begin{vmatrix}
\varphi_1 & \varphi_2 & \varphi_3 & 0 & 0 \\[6pt]
\dfrac{\partial\varphi_1}{\partial x_1} &
\dfrac{\partial\varphi_2}{\partial x_1} &
\dfrac{\partial\varphi_3}{\partial x_1} &
\dfrac{\partial\varphi}{\partial x_1} &
\dfrac{\partial f}{\partial x_1} \\[6pt]
\dfrac{\partial\varphi_1}{\partial x_2} &
\dfrac{\partial\varphi_2}{\partial x_2} &
\dfrac{\partial\varphi_3}{\partial x_2} &
\dfrac{\partial\varphi}{\partial x_2} &
\dfrac{\partial f}{\partial x_2} \\[6pt]
\dfrac{\partial\varphi_1}{\partial x_3} &
\dfrac{\partial\varphi_2}{\partial x_3} &
\dfrac{\partial\varphi_3}{\partial x_3} &
\dfrac{\partial\varphi}{\partial x_3} &
-\dfrac{\partial f}{\partial x_3} \\[6pt] % ?? uncertain: The short stroke before the last fraction -- minus sign / scan artefact.
% EDITORIAL D0031 RISOLTO: The enlargement shows a definite short horizontal stroke, on the mathematical axis, before the derivative fraction. Retain it as printed and retain the existing mathematical misprint proposal. % [sic] -\partial f/\partial x_3 -> \partial f/\partial x_3, if the stroke is type, the preceding derivative column has no minus.
\dfrac{\partial\varphi_1}{\partial x_4} &
\dfrac{\partial\varphi_2}{\partial x_4} &
\dfrac{\partial\varphi_3}{\partial x_4} &
\dfrac{\partial\varphi}{\partial x_4} &
\dfrac{\partial f}{\partial x_4}
\end{vmatrix}.
\]

Sottraendo alla \(1^{\mathrm a}\) riga moltiplicata per \(m\) la \(2^{\mathrm a}\) moltiplicata per \(x_1\), la \(3^{\mathrm a}\) per \(x_2\), la \(4^{\mathrm a}\) per \(x_3\), la \(5^{\mathrm a}\) per \(x_4\) col tener conto del teorema d’Eulero, si potrà sostituire ad essa la
\[
0,\ 0,\ 0,\ -s\varphi,\ -nf % ?? uncertain: The faint mark after nf -- comma / period / scan speck; read as a scan speck.
% EDITORIAL D0032 RISOLTO: The tiny mark is separated from nf by an unusually large gap and has the appearance of the isolated scan specks in the surrounding blank area.
\]
poichè con \(n,m,s\) abbiamo designato rispettivamente gli ordini di \(f\), delle \(\varphi_i\) e della \(\varphi\).

È chiaro che il determinante in cui figura questa linea si annulla quando siano simultaneamente \(\varphi=0\) e \(f=0\), perchè annullandosi codesti termini la prima riga diventa identicamente nulla, c. d. d.

% scan: s033L
% !PRINTED 46
\weberpage{46}
% head: CAPITOLO SECONDO
% stamp: Dark scan-edge band on the right and scattered scan specks.
Aggiungiamo che la \(K\) non può staccarsi più di tre volte dalla jacobiana della rete; o almeno, se ciò accade, la \(K\) fa parte della jacobiana della rete spogliata della parte fissa \(K\). % ?? uncertain: The short stroke above the final K -- overbar / upper serif or scan mark.
% EDITORIAL D0033 RISOLTO: At enlargement, the final K has exactly the same continuous upper serif as the two unbarred Ks immediately above it; there is no separate overbar. Remove the transcription's overline.

Infatti un successivo staccamento di \(K\) dipenderebbe da un contatto fra la superficie \(f\) e la \(\Phi\) lungo tutta la curva \(K\). Ora supponiamo che in un punto \(P\) di questa curva coincidano i piani tangenti alla \(f\) ed alla \(\Phi\). Dovrà aversi
\[
\varrho\frac{\partial f}{\partial x_i}=\frac{\partial\Phi}{\partial x_i},
\qquad (i=1,2,3,4).
\]
Calcoliamoci la \(\dfrac{\partial\Phi}{\partial x_i}\). Possiamo derivare per colonne; in tal modo la derivata verrà espressa come una somma di determinanti ottenuti sostituendo in una delle colonne le derivate rispetto a \(x_i\). Ora è chiaro che trovandoci in un punto in cui \(\varphi=0,f=0\), i determinanti che si ottengono derivando le prime tre colonne sono zero. Quindi resta:
\[
\begin{aligned}
m\left(\frac{\partial\Phi}{\partial x_i}\right)_P
&=
\begin{vmatrix}
0 & 0 & 0 & -s\dfrac{\partial\varphi}{\partial x_i} & 0 \\[6pt]
\dfrac{\partial\varphi_1}{\partial x_1} &
\dfrac{\partial\varphi_2}{\partial x_1} &
\dfrac{\partial\varphi_3}{\partial x_1} &
\dfrac{\partial^2\varphi}{\partial x_1\partial x_i} &
\dfrac{\partial f}{\partial x_1} \\[6pt]
\cdots & \cdots & \cdots & \cdots & \cdots \\[6pt]
\dfrac{\partial\varphi_1}{\partial x_4} &
\dfrac{\partial\varphi_2}{\partial x_4} &
\dfrac{\partial\varphi_3}{\partial x_4} &
\dfrac{\partial^2\varphi}{\partial x_4\partial x_i} &
\dfrac{\partial f}{\partial x_4}
\end{vmatrix}
\\[12pt]
&+
\begin{vmatrix}
0 & 0 & 0 & 0 & -n\dfrac{\partial f}{\partial x_i} \\[6pt]
\dfrac{\partial\varphi_1}{\partial x_1} &
\dfrac{\partial\varphi_2}{\partial x_1} &
\dfrac{\partial\varphi_3}{\partial x_1} &
\dfrac{\partial\varphi}{\partial x_1} &
\dfrac{\partial^2 f}{\partial x_1\partial x_i} \\[6pt]
\cdots & \cdots & \cdots & \cdots & \cdots \\[6pt]
\dfrac{\partial\varphi_1}{\partial x_4} &
\dfrac{\partial\varphi_2}{\partial x_4} &
\dfrac{\partial\varphi_3}{\partial x_4} &
\dfrac{\partial\varphi}{\partial x_4} &
\dfrac{\partial^2 f}{\partial x_4\partial x_i}
\end{vmatrix}.
\end{aligned}
\]

Sviluppando i due determinanti secondo la prima riga si ha:
\begin{gather*}
m\left(\frac{\partial\Phi}{\partial x_i}\right)_P
=\varrho m\left(\frac{\partial f}{\partial x_i}\right)_P
\\[10pt]
=s\frac{\partial\varphi}{\partial x_i}
\begin{vmatrix}
\dfrac{\partial\varphi_1}{\partial x_1} &
\dfrac{\partial\varphi_2}{\partial x_1} &
\dfrac{\partial\varphi_3}{\partial x_1} &
\dfrac{\partial f}{\partial x_1} \\[6pt]
\cdots & \cdots & \cdots & \cdots \\[6pt]
\dfrac{\partial\varphi_1}{\partial x_4} &
\dfrac{\partial\varphi_2}{\partial x_4} &
\dfrac{\partial\varphi_3}{\partial x_4} &
\dfrac{\partial f}{\partial x_4}
\end{vmatrix}
\\[10pt]
%CONT>>
% scan: s033R
% !PRINTED 47
\text{\weberpage{47}}
% head: SISTEMI COVARIANTI E INVARIANTI
%<<CONT
% stamp: Scattered scan specks and a dark band along the lower-left edge.
-n\frac{\partial f}{\partial x_i}
\begin{vmatrix}
\dfrac{\partial\varphi_1}{\partial x_1} &
\dfrac{\partial\varphi_2}{\partial x_1} &
\dfrac{\partial\varphi_3}{\partial x_1} &
\dfrac{\partial\varphi}{\partial x_1} \\[6pt]
\cdots & \cdots & \cdots & \cdots \\[6pt]
\dfrac{\partial\varphi_1}{\partial x_4} &
\dfrac{\partial\varphi_2}{\partial x_4} &
\dfrac{\partial\varphi_3}{\partial x_4} &
\dfrac{\partial\varphi}{\partial x_4}
\end{vmatrix}
\qquad (i=1,2,3,4).
\end{gather*}\edFlush

Indicando brevemente il primo determinante del secondo membro con \(A\), il secondo con \(B\), questa equazione si può scrivere
\[
m\varrho\frac{\partial f}{\partial x_i}
=s\frac{\partial\varphi}{\partial x_i}A
-n\frac{\partial f}{\partial x_i}B,
\qquad (i=1,2,3,4),
\]
che a priori offre due possibilità:

1) \(\dfrac{\partial f}{\partial x_i}=\lambda\dfrac{\partial\varphi}{\partial x_i}\), ossia le superficie \(f=0\) e \(\varphi=0\) si toccano, ciò che è contro le ipotesi; ovvero

2) \(A=0\) e allora, come si riconosce subito dalla struttura del determinante \(A\), il punto \(P\) si trova sulla jacobiana della rete.

Il procedimento algebrico che qui abbiamo sviluppato si lascia illustrare da alcune semplici considerazioni geometriche che ci limiteremo ad accennare, riferendoci al caso più semplice in cui la \(f\) sia una superficie d’ordine \(n\) priva di singolarità nello spazio ordinario, le \(\varphi_1,\varphi_2,\varphi_3\) siano tre superficie linearmente indipendenti d’ordine \(m>n\) e la \(\varphi\) sia una superficie d’ordine \(s\).

Osserviamo che la superficie \(f\) presa insieme con una superficie \(\psi\) d’ordine \(m-n\) si può combinare linearmente con le \(\varphi_1,\varphi_2,\varphi_3\); si costruirà quindi la superficie jacobiana del sistema lineare \(\infty^3\) così determinato, cioè il luogo dei punti doppi delle superficie del sistema, che sarà una superficie \(D\) d’ordine \((4m-4)\)\origfootnote{1}{s033R-1}{Si ha qui un’estensione immediata della jacobiana di una rete di curve piane. Cfr. \textsc{Enriques-Chisini}, Op. cit., Libro III, cap. I, \S~5, vol. II, pag. 31.}. Ora la \(D\) segherà \(f\) in una curva d’ordine \((4m-4)n\) che riuscirà composta della jacobiana della rete delle curve \(|C|\) segate da
\[
\lambda_1\varphi_1+\lambda_2\varphi_2+\lambda_3\varphi_3=0
\]
e della curva \(L\) intersezione di \(f\) e \(\psi\); siccome quest’ultima è d’ordine \(n(m-n)\), così la jacobiana della rete \(|C|\) risulterà d’ordine \(n(3m+n-4)\).

Ciò posto si aggiunga alla rete
\[
\lambda_1\varphi_1+\lambda_2\varphi_2+\lambda_3\varphi_3=0
\]
una superficie \(\varphi\) d’ordine \(s\), secante \(f\) in una curva \(K\) d’ordine \(ns\); aggiungendo la \(\varphi\) anche ad \(f\) si avrà un sistema lineare \(\infty^3\) di super%
% The printed word "super-ficie" crosses the page boundary and is joined here.
%CONT>>
% scan: s034L
% !PRINTED 48
\unskip\weberpage{48}%
% head: CAPITOLO SECONDO
%<<CONT
% stamp: Dark scan-edge band on the right and scattered scan specks.
% The opening fragment "ficie" is included in the joined word on the preceding page.
ficie d’ordine \(m+s\) in cui la \(\varphi\) figura come parte fissa. Perciò questa \(\varphi\) si staccherà 4 volte dalla jacobiana \(D\) del sistema; la sua intersezione con \(f\) (d’ordine \(ns\)), figurerà una volta nella curva intersezione di \(f\) con \(\psi+\varphi\), e quindi comparirà 3 volte nella curva jacobiana della rete \(|K+C|\), che è ora d’ordine
\[
n[3(m+s)+n-4].
\]

\edsection{3. Teorema fondamentale.}{3. Teorema fondamentale.}{sec.II.3}

Dalle osservazioni che precedono segue il

\textsc{Teorema fondamentale:} \textit{Dati sopra la superficie \(F\) due sistemi lineari irriducibili \(|C|\) ed \(|L|\), che supponiamo sempre virtualmente privi di punti base e di dimensione non minore di \(2\), il sistema jacobiano \(|C+L|_j\) della loro somma viene espresso da}
\[
|C+L|_j=|C_j+3L|=|3C+L_j|.
\]

Infatti nel sistema \(|C+L|\) è contenuto il sistema \(|C|+L\) con una curva fissa \(L\), ed il sistema \(C+|L|\) con una curva fissa \(C\), i cui sistemi jacobiani completi sono rispettivamente i sistemi
\[
|C_j+3L|\quad\text{ed}\quad |L_j+3C|,
\]
che coincidono dunque nel sistema \(|C+L|_j\).

In base al teorema fondamentale si può definire ora il sistema completo jacobiano di un sistema lineare \(|C|\) che abbia meno di due dimensioni, riducendosi dunque ad un fascio, od anche ad una curva isolata \(C\). Invero basterà sommare alla \(C\) un sistema ausiliario \(|L|\), \(\infty^2\) almeno, ed assumere
\[
|C_j|=|3C+L_j-3L|.
\]

Il teorema fondamentale si estende così agli jacobiani di sistemi lineari comunque riducibili. Occorre soltanto un’avvertenza relativa al caso in cui si considerino sistemi lineari \(|C|\) ed \(|L|\) con \textit{punti base assegnati} sopra la superficie \(F\). Se \(|C|\) possiede un punto base \(i\)-plo \(O\) che non sia base per \(|L|\), ed \(|L|\) possiede un punto base \(s\)-plo \(O'\) che non sia base per \(|C|\), il sistema jacobiano \(|C+L|_j\) avrà in \(O\) la molteplicità (virtuale) \(3i-1\), ed in \(O'\) la molteplicità \(3s-1\), ed invece \(O\) avrà la molteplicità \(3i\) per \(|3C+L_j|\) ed \(O'\) la molteplicità \(3s\) per \(|3L+C_j|\); quindi l’equivalenza tra i due sistemi \(|C_j+3L|\) e \(|3C+L_j|\) sussisterà soltanto quando si aggiungano ai due sistemi rispettivamente gli intorni dei punti \(O\) ed \(O'\). Si avrà cioè:
\[
|3C+L_j+O|=|3L+C_j+O'|;
\]
%CONT>>
% scan: s034R
% !PRINTED 49
\weberpage{49}
% head: SISTEMI COVARIANTI E INVARIANTI
%<<CONT
% stamp: Scattered scan specks and a thin dark mark along the left edge.
e l’equivalenza così precisata ha il senso che risulta dalle nostre convenzioni: quando si muti la superficie \(F\), trasformando \(O\) ed \(O'\) in curve eccezionali, queste vanno sommate rispettivamente ai trasformati dei due sistemi \(|C_j+3L|\) ed \(|L_j+3C|\).

Dopo di ciò ritorniamo al caso, cui generalmente ci riferiremo nel seguito, in cui si tratta di \textit{sistemi completi virtualmente privi di punti base} sopra la superficie \(F\). L’equivalenza \(|C_j+3L|=|L_j+3C|\) si può ricondurre alla % The plus sign is repeated at the printed line break.
\[
|C_j-3C|=|L_j-3L|=|K|,
\]
che in parole significa: \textit{se la jacobiana \(|C_j|\) di un sistema \(|C|\) contiene il triplo di questo, anche la jacobiana \(|L_j|\) di un altro qualsiasi sistema \(|L|\) conterrà il triplo di \(|L|\) e il sistema \(|K|\) ottenuto dalla sottrazione risulterà indipendente dal sistema da cui si parte.} Questo sistema \(|K|\) si definirà come \textit{sistema canonico} (\textit{impuro}) della superficie \(F\).

Ma l’equivalenza simbolica che abbiamo scritto non perde di significato se anche \(|C_j|\) non contenga \(|3C|\), cioè nell’ipotesi che il sistema differenza \(|C_j-3C|\) abbia soltanto un’esistenza virtuale. In ogni caso il significato dell’equivalenza fra sistemi virtuali è dato qui dalla relazione che esprime il teorema fondamentale sugli jacobiani; non c’è che una semplice differenza di linguaggio.

Conviene anche ricondurre l’equivalenza dei sistemi virtuali sopra indicata ad una forma più semplice. Perciò si definirà il \textit{sistema lineare \(|C'|\) aggiunto a \(|C|\)}, mediante la relazione
\[
|C'|=|C_j-2C|,
\]
avvertendo che \(|C'|\) (a differenza di \(|K|\)) avrà in generale un’\textit{esistenza effettiva} per i sistemi \(|C|\) tracciati sopra una superficie qualunque: ciò si verifica già nel piano per i sistemi lineari di curve d’ordine \(m\) e di genere \(\pi>0\), poichè le \(C'\) sono allora le curve di ordine \(m-3\) passanti \(i-1\) volte per ogni punto base \(i\)-plo di \(|C|\); e la cosa si verificherà nel seguito per i sistemi lineari il cui genere superi un certo limite.

Con l’introduzione del sistema aggiunto, il sistema canonico viene dato da:
\[
|K|=|C'-C|=|L'-L|\quad .\,.\,.\,.
\]
e l’equivalenza di sistemi effettivi o virtuali così espressa significa
\[
|C'+L|=|L'+C|
\]
che è la \textit{relazione fondamentale per i sistemi aggiunti} (privi di punt base), semplice trasformata della relazione data di sopra per i sistemi jacobiani. % [sic] punt -> punti, the final i is absent before "base".

% EDITORIAL: segnatura originale: \textsc{Enriques F.} - \textit{Superficie algebriche.}\hfill 4
\printersignature{\textsc{Enriques F.} - \textit{Superficie algebriche.}\quad 4}

% scan: s035L
% !PRINTED 50
\weberpage{50}
% head: CAPITOLO SECONDO
% stamp: Large dark scan artefact in the lower-left margin, dark right-edge band, and scattered specks.
\edsection{4. Curve canoniche.}{4. Curve canoniche.}{sec.II.4}

Abbiamo veduto come a partire da un sistema lineare \(|C|\) dato sopra una superficie \(F\) si costruisca il sistema lineare aggiunto \(|C'|\) e quindi il sistema canonico \(|K|=|C'-C|\) che potrà avere una esistenza effettiva o virtuale e si è riconosciuto che quest’ultimo sistema riesce indipendente dalla scelta del sistema \(|C|\), supposto privo di punti base su \(F\). Anche nel caso che \(|C|\) possegga un punto base \(i\)-plo \(O\) semplice per \(F\), il sistema \(|C'-C|\) potrà ritenersi ancora identico a \(|K|\), se si prescinda dalla curva infinitesima costituente l’intorno di \(O\), che, essendo \((i-1)\)-plo per \(|C'|\), dovrebbe sommarsi alle \(K\). L’indipendenza di \(|K|\) dalla scelta di \(|C|\) significa che il sistema canonico \(|K|\) ha significato invariante per trasformazioni birazionali della superficie.

Per spiegare meglio la cosa, si considerino due superficie dello spazio ordinario \(F\) e \(\Phi\), rispettivamente d’ordine \(n\) e \(m\), dotate di singolarità normali, e si supponga che fra di esse interceda una corrispondenza birazionale, aggiungendo dapprima l’ipotesi che questa non abbia punti fondamentali nè sull’una nè sull’altra superficie. Le sezioni piane della superficie \(F\) formeranno un sistema lineare \(|C|\) il cui trasformato non ha punti base su \(\Phi\), e le sezioni piane \(L\) di \(\Phi\) formeranno del pari un sistema lineare su \(\Phi\) il cui trasformato non ha punti base su \(F\). Ora si potrà costruire il sistema canonico \(|K|\) su \(F\) sia a partire dal sistema \(|C|\) delle sezioni piane, sia a partire dal sistema \(|L|\), e si avrà
\[
|K|=|C'-C|=|L'-L|;
\]
lo stesso sistema si costruirà su \(\Phi\), sia a partire dal sistema \(|L|\) dalle sezioni piane, sia a partire dal sistema \(|C|\), avendosi ancora % [sic] dalle -> delle, the phrase identifies the system of plane sections.
\[
|K|=|L'-L|=|C'-C|.
\]
Dunque le curve del sistema canonico \(|K|\) costruite su \(F\) a partire dal sistema delle sezioni piane, si trasformano nelle curve del sistema canonico di \(\Phi\) costruito a partire dal sistema \(|L|\) delle sezioni piane. Resta a vedere che cosa importi codesta costruzione.

Si può costruire sopra \(F\) il sistema canonico, costruendo anzitutto il sistema jacobiano di \(|C|\). A tale scopo osserveremo che una stella di sezioni piane di \(F\), avente come centro un punto qualunque \(O\) dello spazio, costituisce una rete la cui jacobiana è l’intersezione di \(F\) con la superficie polare di \(O\), d’ordine \(n-1\), intersezione ulteriore di essa fuori della curva nodale.

Le jacobiane delle reti contenute nel sistema \(|C|\) delle sezioni piane formeranno così un sistema lineare \(\infty^3\) segato su \(F\) dalle polari, fuori della curva nodale; ma questo sistema avrà come punti
%CONT>>
% scan: s035R
% !PRINTED 51
\weberpage{51}
% head: SISTEMI COVARIANTI E INVARIANTI
%<<CONT
% stamp: Dark scan-edge band on the left and scattered scan specks.
base semplici i punti cuspidali sulla detta curva nodale. Che questi punti debbano essere ritenuti come « punti semplici » della superficie e perciò, secondo le nostre convenzioni, debbano esser tolti per ottenere il sistema jacobiano completo \(|C_j|\), si può facilmente riconoscere. Invero, se ci riferiamo ad una superficie trasformata su cui la curva nodale di \(F\) venga rappresentata dalla curva delle coppie neutre rispetto al sistema \(|C|\), i punti cuspidali daranno su tale curva punti base di contatto per una rete di \(C\) e quindi punti base doppi per un fascio di \(C\): si tratta dunque di punti che appartengono semplicemente alle jacobiane di tutte le reti contenute nel sistema \(\infty^3\) delle \(C\).

Ciò posto, il sistema jacobiano completo \(|C_j|\) privo, almeno virtualmente, di punti base, conterrà l’intero sistema delle curve sezioni di \(F\) con le superficie di ordine \(n-1\) passanti per la sua curva doppia (superficie \textit{aggiunte}). Se ora da queste superficie si tolgono 3 piani, avremo che le superficie d’ordine \(n-4\) (supposte esistenti) passanti per la curva doppia di \(F\), segano su \(F\) curve canoniche.

In modo simile le superficie aggiunte d’ordine \(m-4\), passanti per la curva doppia di \(\Phi\), segheranno su \(\Phi\) curve canoniche.

Quindi l’invarianza del sistema canonico per trasformazioni birazionali significherà che « le curve sezioni di \(F\) con le superficie aggiunte d’ordine \(n-4\), si mutano, per la trasformazione birazionale, in curve equivalenti alle sezioni di \(\Phi\) con le sue superficie aggiunte d’ordine \(m-4\) ».

Il teorema assume una forma più espressiva ove si anticipi un risultato che verrà stabilito nel seguente capitolo, cioè che « le superficie aggiunte di un ordine qualunque, e in particolare quelle di ordine \(n-4\), segano sopra la superficie \(F\) d’ordine \(n\) un sistema lineare completo », cioè, in questo caso, l’intero sistema canonico \(|K|\).

Avremo pertanto: \textit{le curve} (canoniche) \textit{intersezioni della superficie \(F\) con le sue aggiunte d’ordine \(n-4\)} (fuori della curva nodale) \textit{si mutano, per una trasformazione birazionale, nelle sezioni di \(\Phi\) con le sue superficie aggiunte d’ordine \(m-4\).}

A dir vero si è supposto che la corrispondenza birazionale fra \(F\) e \(\Phi\) sia senza eccezioni, cioè che il sistema \(|L|\) non abbia punti base su \(F\), e \(|C|\) non abbia punti base su \(\Phi\). Ma se, per esempio, \(|L|\) abbia un punto base \(i\)-plo \(O\) sopra \(F\), il teorema sussisterà ancora con questa modificazione, che: la \textit{curva eccezionale} \(\omega\) corrispondente al punto \(O\) di \(F\) si sommerà alle trasformate delle \(|K|\), figurando come \textit{parte fissa del sistema canonico impuro costituito dalle sezioni di \(\Phi\) con le sue superficie aggiunte d’ordine \(m-4\).}

\textsc{Nota.} -- Si è supposto che la \(F\) e la \(\Phi\) siano dotate soltanto di singolarità normali (cioè di curva doppia con punti tripli); ma il
%CONT>>
% scan: s036L
% !PRINTED 52
\weberpage{52}
% head: CAPITOLO SECONDO
%<<CONT
% stamp: Dark vertical scan band at the right edge; scattered scan specks.
teorema si estende anche se venga meno questa restrizione, purchè si definiscano le \textit{superficie aggiunte} ad una data in relazione alle singolarità più elevate che essa possegga. Basta qui accennare che una curva \(i\)-pla ordinaria di \(F\) dovrà essere \((i-1)\)-pla per le superficie aggiunte, ed un punto conico d’ordine \(s\), fuori della curva multipla, dovrà avere in generale per codeste aggiunte la molteplicità \(s-2\). Ma di ciò più avanti.

\edsection{5. Proprietà delle curve canoniche.}{5. Proprietà delle curve canoniche.}{sec.II.5}

Le curve canoniche sopra una superficie \(F\) (dico le curve canoniche \textit{impure}, sezioni della \(F\) d’ordine \(n\) colle superficie aggiunte d’ordine \(n-4\), fuori della curva doppia) godono di una proprietà che qui vogliamo esplicitamente riconoscere, e che, nelle condizioni più semplici, vale anche a caratterizzarle, siccome vedremo nel seguente capitolo.

Si consideri su \(F\) un sistema lineare irriducibile privo di punti base, \(|C|\), di dimensione \(r\geq 2\), di genere \(\pi\) e di grado \(n\). La jacobiana di una rete di curve \(C\) sega su una \(C\) il gruppo dei \(2n+2\pi-2\) punti doppi della \(g_n^1\) caratteristica, e questo appartiene alla serie somma delle serie canonica e del doppio della \(g_n^1\). Togliendo \(2C\) da \(|C_j|\), si avranno curve \(C'\) aggiunte a \(|C|\) secanti su \(C\) gruppi canonici; e togliendo ancora una \(C\) si avrà che le curve canoniche \(K\) segano su \(C\) gruppi residui della serie caratteristica.

Enunceremo dunque che:

\textit{Le curve aggiunte \(C'\) segano sopra le curve del sistema \(|C|\) gruppi canonici. Le curve canoniche \(K\) segano sopra le curve \(C\) gruppi residui della serie caratteristica.}

Segue di qui che le aggiunte \(C'\) segano una curva \(C\), di genere virtuale \(\pi\) in \(2\pi-2\) punti, e le curve canoniche \(K\), effettive o virtuali, la segano in
\[
i=2\pi-2-n
\]
punti.

\textit{Questa eguaglianza sussiste per qualunque curva \(C\), anche isolata, irriducibile o riducibile,} riferendoci sempre ai caratteri virtuali, definiti nel supposto che la \(C\) non abbia punti base assegnati su \(F\). Infatti si sommi alla \(C\) una conveniente curva \(L\) di genere \(\varrho\) e grado \(\nu\) che la incontri in \(s\) punti, per modo che il sistema lineare completo \(|C+L|\) sia irriducibile e di dimensione \(>2\); basterà assumere come sistema \(|L|\) il sistema segato dalle superficie d’ordine abbastanza alto passanti per \(C\). Ora le curve \(C'+L\) appartenenti al sistema \(|(C+L)'|\) segheranno una \(C+L\) in
\[
2(\pi+\varrho+s-1)-2
\]
%CONT>>
% scan: s036R
% !PRINTED 53
\weberpage{53}
% head: SISTEMI COVARIANTI E INVARIANTI
%<<CONT
% stamp: Dark scan band along the left edge; scattered scan specks.
punti; togliendo le intersezioni di esse con \(L\), che sono \(2\varrho-2+s\), e poi le \(s\) intersezioni con \(L\), resulterà
\[
(CC')=2\pi-2.
\]
Similmente il numero virtuale delle intersezioni di \(C\) e \(K\) vale sempre \((CK)=2\pi-2-n\).

\textsc{Nota.} – Aggiungiamo che per una curva irriducibile \(C\), senza punti multipli su \(F\), il gruppo segato su di essa da un’aggiunta \(C'\) sarà sempre un gruppo canonico, e il gruppo segato da una \(K\) sarà residuo della serie caratteristica (effettiva o virtuale). Basterà dimostrare l’asserto per una curva \(C\) che faccia parte di una curva composta \(C+L\), appartenente ad una rete irriducibile. Riprendendo il ragionamento fatto innanzi si troverà ora che una curva di
\[
|C_j|=|(C+L)_j-3L|
\]
sega su \(C\) un gruppo appartenente alla serie somma della serie canonica e del doppio della serie caratteristica virtuale della \(C\) stessa.

\edsection{6. Genere superficiale e genere lineare.}{6. Genere superficiale e genere lineare.}{sec.II.6}

L’invarianza del sistema canonico per trasformazioni birazionali della superficie porta che i caratteri di questo abbiano del pari un significato invariantivo rispetto a tali trasformazioni. Tuttavia occorre distinguere:

1) \textit{invarianti assoluti}, che non mutano per qualsiasi trasformazione, e

2) \textit{invarianti relativi} che conservano lo stesso valore per trasformazioni birazionali senza eccezione, ma si alterano per trasformazioni dotate di punti fondamentali, che mutino un punto (semplice) della superficie in una curva eccezionale dell’altra.

La dimensione del sistema canonico, o meglio il numero delle curve canoniche linearmente indipendenti, ci dà appunto un invariante assoluto, che dicesi \textit{genere geometrico superficiale} (\(p\) o \(p_g\)), o semplicemente « genere », della superficie.

Quando si sia dimostrato che il sistema canonico completo viene segato sopra una superficie \(F\) d’ordine \(n\) dello spazio ordinario dalle sue aggiunte \(\varphi_{n-4}\) d’ordine \(n-4\) (cfr. Cap. III), si ha che il genere di una superficie generale d’ordine \(n\) (che è il numero delle \(\varphi_{n-4}\) linearmente indipendenti) vale
\[
p=\frac{(n-3)(n-2)(n-1)}{6}.
\]

% scan: s037L
% !PRINTED 54
\weberpage{54}
% head: CAPITOLO SECONDO
% stamp: Dark vertical scan band at the right edge; scattered scan specks.
Quindi il genere del piano, delle quadriche e delle superficie cubiche, varrà
\[
p=0.
\]
Ciò si verifica direttamente per il piano osservando che, in esso, il sistema \(|C|\) delle curve generali d’ordine \(n\geq 3\) contiene il suo aggiunto \(|C'|\), costituito dalle curve d’ordine \(n-3\); così, essendo \(|C|\) più ampio di \(|C'|\), non può accadere che \(|C'|\) contenga \(|C|\): \(C'-C\) non ha esistenza effettiva, che vuol dire \(p=0\).

Dall’essere nullo il genere del piano segue che tutte le \textit{superficie razionali} (rappresentabili punto per punto sul piano) hanno parimente il \textit{genere zero}. Fra queste ci sono appunto le quadriche e le cubiche (almeno quelle che non sono coni) dette innanzi.

Un’altra famiglia di superficie aventi il \textit{genere} (geometrico) \textit{nullo} sono \textit{le rigate}, e quindi anche le loro trasformate. L’asserzione si giustifica osservando che una generatrice di una rigata \(F\) non può avere sistema aggiunto, poichè questo dovrebbe segare sulla generatrice la serie canonica. In generale il sistema \(|C'|\) aggiunto ad un sistema lineare \(|C|\) sopra la rigata non può contenere \(|C|\) perchè le \(C\) segando le generatrici in \(n\) punti, le \(C'\) le segano in \(n-2\) punti (cfr. Cap. III, \S{} 3).

Ritornando alle superficie qualunque, di genere \(p\geq 0\), il genere delle curve canoniche (effettive o virtuali) fornisce un invariante relativo, detto \textit{genere lineare} \(p^{(1)}\) o anche genere lineare \textit{relativo}, che diminuisce di 1 per ogni punto della superficie \(F\) che venga cambiato in una curva eccezionale e viceversa cresce di 1 per ogni curva eccezionale che si muti in un punto. Infatti abbiamo veduto che, quando un punto \(O\) della superficie \(F\) si muta in una curva eccezionale \(\omega\) di \(\Phi\), questa \(\omega\) si aggiunge alle \(\bar{K}\) trasformate delle curve canoniche \(K\); ora una curva siffatta non è connessa colle dette \(K\), cioè ha con esse un numero virtuale d’intersezioni nullo, poichè il punto \(O\) non è — almeno virtualmente — punto base di \(|K|\), non avendo il detto sistema \(|K|\) alcun punto base assegnato sopra \(F\). Pertanto sommando alle \(\bar{K}\) la \(\omega\) di genere zero e di grado \(-1\), sconnessa con esse, il genere di codeste curve diminuisce di 1.

Qui si palesa l’opportunità di considerare, almeno per le superficie di genere \(p>0\) dove il sistema canonico ha un’esistenza effettiva, accanto al genere lineare relativo \(p^{(1)}\), anche il genere del \textit{sistema canonico puro}, spogliato delle sue componenti fisse eccezionali, che sarà poi il genere delle curve canoniche sopra una superficie trasformata priva di curve eccezionali, siccome verrà precisato nel seguente \S{} 7; otterremo così un invariante assoluto, che denomineremo \textit{genere lineare assoluto} della superficie data. Il suo valore sarà
\[
\bar{p}^{(1)}=p^{(1)}+\sigma
\]
%CONT>>
% scan: s037R
% !PRINTED 55
\weberpage{55}
% head: SISTEMI COVARIANTI E INVARIANTI
%<<CONT
% stamp: Dark scan band along the lower left edge; scattered scan specks.
designando con \(\sigma\) il numero delle curve eccezionali appartenenti alla superficie.

Ci riserviamo di precisare, nel prossimo paragrafo, che queste curve sono sempre curve eccezionali di \(1^{\mathrm{a}}\) specie (almeno per \(p>0\)) non connesse fra loro, esaminando in particolare il caso delle curve eccezionali riducibili.

Come il genere del sistema canonico anche il grado di esso fornisce un invariante relativo della superficie, che possiamo indicare con \(p^{(2)}\). Quando, per una trasformazione della superficie \(F\), viene a sommarsi alle \(\bar{K}\) una curva eccezionale \(\omega\) di grado \(-1\) non connessa alle \(K\), il \(p^{(2)}\) diminuisce pure di 1. Ed anche qui, accanto al grado del sistema canonico impuro, possiamo definire il grado del sistema canonico puro, che costituirà un invariante assoluto, definito almeno per le superficie di genere \(p>0\). Ma il \(p^{(2)}\) non ci dà un nuovo carattere invariante in confronto al genere lineare. Infatti \textit{sussiste fra il genere e il grado del sistema canonico impuro la relazione.} % ?? uncertain: punctuation after "relazione" -- period / colon / scan speck.
\[
p^{(2)}=p^{(1)}-1.
\]
Per dimostrarla basta osservare che il sistema aggiunto al sistema canonico \(|K|\) è
\[
|K'|=|2K|;
\]
quindi il numero delle intersezioni di \(K\) con \(2K\) vale \(2p^{(1)}-2\) e perciò
\[
(KK)=p^{(1)}-1.
\]
Nel caso \(p>0\) (e in tutti i casi in cui si potranno definire i caratteri assoluti di cui si discorre, cioè quando la superficie sia trasformabile in altra priva di curve eccezionali) \textit{la relazione precedente sussiste ancora fra il genere e il grado del sistema canonico puro, che sono}
\[
p^{(1)}+\sigma\quad\text{e}\quad p^{(2)}+\sigma.
\]

\textsc{Nota.} – Risulterà più avanti che le sole superficie per cui non è possibile definire il genere lineare assoluto (riferendosi ad una trasformata priva di curve eccezionali) sono quelle che appartengono alla famiglia delle rigate (razionali o no). Tuttavia anche per queste si giunge a costruire un invariante assoluto, quale è dato — secondo \textsc{Castelnuovo}\origfootnote{1}{s037R-1}{\textit{Sul genere lineare di una superficie e sulla classificazione a cui esso dà luogo.} Rendic. Acc. Lincei, vol. VI, 1897. «Memorie scelte»; pag. 443.} — dal massimo del genere lineare relativo, in confronto a tutte le possibili trasformate della superficie data.

Qui riesce opportuno di calcolare il genere lineare del piano e delle rigate irrazionali di genere \(p\) (\(>0\)), che rispondono appunto a quel massimo.

% scan: s038L
% !PRINTED 56
\weberpage{56}
% head: CAPITOLO SECONDO
% stamp: Dark vertical scan band at the right edge; scattered scan specks.
Il genere lineare del piano si calcolerà riferendosi ad un sistema lineare di quartiche piane \(|C|\), di genere \(\pi=3\) e di grado \(n=16\): il genere \(\pi'=0\) del sistema aggiunto, che è la rete delle rette, viene espresso dalla formula
\[
\pi'=\pi+p^{(1)}+2\pi-2-n-1.
\]
Facendo qui \(\pi=3\), \(n=16\), \(\pi'=0\), si ricava che il genere lineare del piano vale
\[
p^{(1)}=10.
\]

Passiamo a considerare una rigata d’ordine \(m\) e di genere \(p>0\). Il sistema doppio delle sezioni piane (segato dalle quadriche) ha il genere
\[
\pi=2p+m-1
\]
e il grado
\[
n=4m.
\]

Il sistema aggiunto ad esso è costituito da gruppi di \(m+2p-2\) rette\origfootnote{1}{s038L-1}{Ciò può verificarsi in varie guise; rimandiamo al Cap. III.}, e quindi ha il genere
\[
\pi'=-(m+2p-2)+1=-m-2p+3.
\]

Quindi dalla formula
\[
\pi'=\pi+p^{(1)}+2\pi-2-n-1,
\]
si ricava che: \textit{il genere lineare delle rigate di genere \(p>0\)} vale
\[
p^{(1)}=-8(p-1)+1.
\]

\edsection{7. Le curve eccezionali come parti fisse del sistema canonico.}{7. Le curve eccezionali come parti fisse del sistema canonico.}{sec.II.7}

Giova ritornare a precisare in qualche punto ciò che si riferisce alle curve eccezionali appartenenti ad una superficie di genere \(p>0\), e alla loro appartenenza, come parti fisse, al sistema canonico \(|K|\).

Che una curva eccezionale di prima specie, \(\omega\), sulla superficie \(F\), appartenga come componente fissa al sistema \(|K|\) risulta direttamente dal fatto che essa ha il genere \(\varrho=0\) e il grado \(\nu=-1\). Infatti la \(\omega\) deve essere intersecata dalle \(K\) in
\[
2\varrho-2-\nu=-1
\]
punti; siccome questo numero (virtuale) è negativo la \(\omega\) deve far parte di ogni \(K\). Si può dire di più: sulla superficie \(F\), di genere \(p>0\), non possono aversi due curve eccezionali irriducibili di prima specie connesse fra loro, cioè aventi qualche punto comune. Infatti, se
%CONT>>
% scan: s038R
% !PRINTED 57
\weberpage{57}
% head: SISTEMI COVARIANTI E INVARIANTI
%<<CONT
% stamp: Scattered scan specks.
si abbiano due tali curve, \(\omega\) e \(\omega'\), di genere 0 e di grado \(-1\), con \(i(\geq 1)\) punti comuni, esse dovranno figurare come parti fisse, da contare un certo numero finito di volte, in tutte le curve canoniche \(K\). Ma, se la prima si stacchi \(r\) volte e la seconda \(s\) volte, il numero delle intersezioni di \(\omega\) e \(\omega'\) colle curve residue varrà risp.
\[
\begin{aligned}
(K\omega)-r(\omega\omega)-s(\omega\omega')&=-1+r-si,\\
(K\omega')-s(\omega'\omega')-r(\omega\omega')&=-1+s-ri.
\end{aligned}
\]
Ora, di questi due numeri, la cui somma è negativa, uno almeno è negativo, sicchè la corrispondente curva, \(\omega\) o \(\omega'\), dovrebbe staccarsi ulteriormente dalle \(K\), contro il supposto.

Collo stesso ragionamento si prova che la superficie \(F\), di genere \(p>0\), non può contenere una curva eccezionale di seconda specie \(\omega\). All’uopo giova riferirci ad una superficie \(\Phi\), trasformata di \(F\), su cui la \(\omega\) diventi una curva di grado \(\nu\geq 0\), conservando sempre il genere \(\varrho=0\). Su questa superficie la \(\omega\) dovrà avere comune con una curva canonica \(K\) il numero virtuale di punti
\[
2\varrho-2-\nu=-2-\nu
\]
che è negativo; la \(\omega\) si stacca dunque come parte fissa dalle \(K\); ma dovrebbe staccarsi infinite volte, ciò che è assurdo. Invero se si supponga che si stacchi \(s\) volte, il numero delle intersezioni colle curve residue varrà
\[
(K\omega)-s(\omega\omega)=-2-(s+1)\nu
\]
che è sempre un numero negativo.

Riassumeremo i risultati ottenuti enunciando il teorema:

\textit{Sopra una superficie di genere \(p>0\), si può avere tutt’al più un numero finito di curve eccezionali irriducibili di prima specie, non connesse fra loro, e non si hanno mai curve eccezionali di seconda specie.}

Si avverta che quando si abbiano soltanto curve eccezionali di \(1^{\mathrm{a}}\) specie irriducibili, ciascuna di queste, \(\omega\), si stacca semplicemente dal sistema canonico, restando fondamentale per il sistema residuo \(|K-\omega|\). Tuttavia non si può escludere che la \(\omega\) si stacchi ulteriormente da \(|K-\omega|\), figurando due volte come parte fissa di \(|K|\). Ciò significa che trasformando la \(\omega\) in un punto \(O\), questo risulterà un \textit{punto base comune} à tutte le curve del \textit{sistema canonico trasformato}. Siccome però questo punto base non è assegnato, deve ritenersi virtualmente inesistente, così anche la \(\omega\) sulla superficie data dovrà ritenersi ancora come facente \textit{parte delle curve canoniche pure} che si ottengono staccando una volta le curve eccezionali che costituiscono parti fisse del \textit{sistema canonico impuro}. % ?? uncertain: accent on the preposition before "tutte" -- à / a with a scan speck.

\textsc{Nota I.} – Una menzione particolare merita il caso in cui si abbiano \textit{curve eccezionali riducibili, provenienti da punti base infinitamente}
%CONT>>
% scan: s039L
% !PRINTED 58
\weberpage{58}
% head: CAPITOLO SECONDO
%<<CONT
% stamp: Dark vertical scan band at the right edge; scattered scan specks.
\textit{vicini} di una superficie trasformata. Qui la \textit{molteplicità delle curve eccezionali} come parti fisse del sistema canonico impuro si lasciano determinare in base al principio di continuità, giacchè sappiamo che il caso di cui si tratta nasce come limite quando i punti base di un sistema trasformante si avvicinino tra loro indefinitamente. Ricordando l’osservazione con cui si chiude il \S{} 12 del precedente Cap., rileviamo ora che essa ha il seguente significato: \textit{Ogni componente \(\omega_i\) di una curva eccezionale riducibile} % [sic] la molteplicità -> le molteplicità, the verb "si lasciano" is plural.
\[
\Omega=\omega+\omega_1+h_2\omega_2+\ldots+h_i\omega_i+\ldots+h_s\omega_s
\]
\textit{si stacca dal sistema canonico impuro tante volte quant’è la somma delle molteplicità con cui essa figura nelle curve eccezionali di cui fa parte}\origfootnote{1}{s039L-1}{\textsc{A. Franchetta}, \textit{Sulle curve eccezionali riducibili di prima specie.} Boll. dell’Unione Mat. It., Bologna, 1940.}.

Per esempio, se si ha una curva eccezionale composta di tre componenti, ciascuna connessa con la successiva, in corrispondenza a tre punti infinitamente vicini succedentisi sopra un ramo lineare, codeste componenti danno luogo à tre curve eccezionali composte % ?? uncertain: accent on the preposition after "luogo" -- à / a with a scan speck.
\[
\omega+\omega_1+\omega_2,\qquad \omega_1+\omega_2,\qquad \omega_2;
\]
allora la \(\omega\) figurerà come parte fissa semplice del sistema canonico \(|K|\), la \(\omega_1\) vi figurerà contata due volte, e la \(\omega_2\) tre volte.

Se invece le tre componenti anzidette provengono da punti successivi sopra un ramo cuspidale (sicchè la \(\omega\) e la \(\omega_1\) siano connesse con la \(\omega_2\) e non fra loro) si avranno tre curve eccezionali
\[
\omega+\omega_1+2\omega_2,\qquad \omega_1+\omega_2,\qquad \omega_2;
\]
quindi la \(\omega\) figurerà ancora semplicemente come parte fissa del sistema canonico \(|K|\), la \(\omega_1\) vi figurerà doppiamente, e la \(\omega_2\) come curva quadrupla.

Il teorema che una superficie \(F\) di genere \(p>0\) possiede in generale soltanto curve eccezionali di \(1^{\mathrm{a}}\) specie irriducibili non connesse fra loro, fa apparire evidente che la \(F\) stessa potrà trasformarsi in altra superficie affatto priva di curve eccezionali, siccome \textsc{Enriques} ha indicato nella sua « Introduzione » del 1896. Ma la dimostrazione rigorosa di ciò si fa nel modo più semplice con \textsc{Castelnuovo-Enriques} (1901), ricorrendo al così detto teorema di Riemann-Roch per le superficie, in condizioni più larghe che comprendono (come già si è accennato) tutte le superficie non appartenenti alla famiglia delle rigate. E perciò ne rimandiamo a più avanti lo sviluppo. Qui vogliamo soltanto notare che il caso in cui si
%CONT>>
% scan: s039R
% !PRINTED 59
\weberpage{59}
% head: SISTEMI COVARIANTI E INVARIANTI
%<<CONT
% stamp: Scattered scan specks.
abbia una curva eccezionale composta di \(s\) parti, non dà luogo ad alcuna difficoltà, perchè una sua componente costituisce una curva eccezionale irriducibile, che potrà sempre eliminarsi, riducendosi così da \(s\) ad \(s-1\).

\textsc{Nota II.} – Rileviamo più esplicitamente ciò che è incluso nella proprietà caratteristica delle curve canoniche di segare sopra una curva \(C\) di grado \(n\) e di genere \(\pi\) un gruppo di \(2\pi-2-n\) punti. Appunto per tale proprietà la differenza
\[
2\pi-2-n
\]
(che è in generale positiva, ma che deve considerarsi anche se virtualmente negativa) appare un carattere additivo per i sistemi di curve \(|C|\) sopra la superficie; e la proprietà additiva di essa si può riconoscere anche direttamente dalle formule che danno il grado e il genere delle curve composte, indipendentemente dal significato effettivo o virtuale che qui lo rende evidente. Ciò posto si può affermare che: \textit{Sopra una superficie di genere \(p>0\) una curva irriducibile \(C\) di grado \(n\) e genere \(\pi\) per cui}
\[
n>2\pi-2
\]
\textit{può essere soltanto una curva eccezionale di prima specie:}
\[
\pi=0,\qquad n=-1.
\]
Infatti se \(2\pi-2-n\) è minore di zero, la \(C\) deve staccarsi da tutte le curve canoniche \(K\); ma se si suppone che si stacchi \(s\) volte, il numero delle intersezioni di \(C\) con le curve residue sarà
\[
2\pi-2-(s+1)n
\]
e perciò sempre negativo, a meno che sia \(n=-1\), e di conseguenza \(\pi=0\).

Questo teorema si estende al caso di curve comunque composte, osservando che una somma di caratteri negativi o nulli \(n_i-(2\pi_i-2)\) non può mai dare un numero positivo. Ma si va incontro ad una analisi un po’ delicata nel caso delle curve eccezionali composte. Si evita la difficoltà riferendoci ad una superficie trasformata priva di curve eccezionali. Ecco il risultato: \textit{Sopra una superficie di genere \(p>0\) non esistono curve (comunque composte) di genere \(\pi\) e grado \(n>2\pi-2\), all’infuori di quelle che si ottengono sommando a curve con \(n\leq 2\pi-2\) delle curve eccezionali sconnesse} (aventi con esse un numero virtuale nullo di punti comuni).

In particolare si deduce che una curva eccezionale di prima specie, comunque composta, appartenente ad una superficie di genere \(p>0\),
%CONT>>
% scan: s040L
% !PRINTED 60
\weberpage{60}
% head: CAPITOLO SECONDO
%<<CONT
% stamp: Dark vertical scan band at the right edge; scattered scan specks.
viene caratterizzata come « curva connessa di genere \(\varrho=0\) e grado \(\nu=-1\) »\origfootnote{1}{s040L-1}{Con una piccola restrizione, cioè escludendo il caso che la curva appartenga ad un fascio lineare di curve razionali (con qualche componente fissa), il teorema si estende a tutte le superficie (\(p\geq 0\)). Cfr. \textsc{A. Franchetta}, \textit{Sulla caratterizzazione delle curve eccezionali ecc.} Boll. Unione Mat. It., 1941.}.

\textsc{Esercizio.} – Lo studioso è invitato a sciogliere il seguente paradosso: si parta da una superficie del \(4^{\mathrm{o}}\) ordine (di genere \(p=1\)) \(F_4\) contenente una retta \(a\); si considererà su \(F_4\) una cubica \(C\), sezione piana per \(a\) e si eseguirà una trasformazione di \(F_4\) in una superficie \(\Phi\), per mezzo di un sistema lineare con un punto base \(O\) sopra \(C\). Avremo così, su \(\Phi\) una curva ellittica \(C\) cui viene sommata una curva eccezionale \(\omega\) che la incontra in un punto. Ora si consideri la somma
\[
C+2\omega;
\]
essa è una curva di genere virtuale \(\varrho=0\) e di grado \(\nu=-1\), eppure non è una curva eccezionale!

\edsection{8. Curve bicanoniche e pluricanoniche.}{8. Curve bicanoniche e pluricanoniche.}{sec.II.8}

Abbiamo veduto che l’esistenza virtuale del sistema canonico
\[
|K|=|C'-C|
\]
costruito a partire da un sistema \(|C|\) senza punti base sopra una superficie \(F\), importa l’esistenza di caratteri numerici quali sono il \(p^{(1)}\) e il \(p^{(2)}\) che vengono definiti indipendentemente dall’esistenza effettiva di \(|K|\). Ma la considerazione di un sistema canonico virtuale che non abbia esistenza effettiva assume un significato più concreto quando accada che il sistema inesistente dia luogo ad un doppio o ad un multiplo effettivamente esistenti che si lasciano definire direttamente dalla relazione
\[
|2K|=|2C'-2C|
\]
ovvero
\[
|iK|=|iC'-iC|,
\]
giacchè il teorema fondamentale per gli aggiunti dà in ogni caso
\[
|iC'-iC|=|iL'-iL|.
\]
Notiamo anche che il sistema bicanonico si può definire per sottrazione del sistema \(|C|\) dal suo secondo aggiunto
\[
|(C+C')'|=|C+C''|=|2C'|;\qquad |2K|=|C''-C|
\]

Il più semplice esempio di una superficie di genere \(p=0\), cui appartiene una curva bicanonica (d’ordine zero), ed anche il primo
%CONT>>
% scan: s040R
% !PRINTED 61
\weberpage{61}
% head: SISTEMI COVARIANTI E INVARIANTI
%<<CONT
% stamp: Scattered scan specks and stray marginal strokes; dark strip along the lower left edge.
che si è presentato, viene offerto dalla superficie del sesto ordine \(F_6\) che passa doppiamente per gli spigoli di un tetraedro, sulla quale avremo occasione di ritornare nel seguito. Le sezioni piane della \(F_6\) sono sestiche \(C\) di genere 4; le curve aggiunte ad esse sono sestiche gobbe \(C'\) segate dalle superficie cubiche che passano semplicemente per gli spigoli del detto tetraedro; le curve \(C\) e \(C'\) sono disequivalenti; ma i loro doppi sono equivalenti, cioè
\[
|2C'|\doteq|2C|; % ?? uncertain: relation sign -- \doteq / =; the small mark above the bars may be a scan speck.
\]
infatti il sistema completo \(|2C|=|2C'|\) viene segato sopra \(F_6\) dalle superficie del \(6^{\mathrm{o}}\) ordine, che hanno ugualmente come doppi gli spigoli del nostro tetraedro.

La \(F_6\) non contiene curve eccezionali, ma se essa si trasformi in una superficie \(\Phi\) per modo che un suo punto \(O\) si muti in una curva \(\omega\), la curva eccezionale \(\omega\) contata due volte costituirà da sola la curva bicanonica di \(\Phi\).

I caratteri del sistema bicanonico e similmente quelli dei sistemi pluricanonici forniscono degli invarianti della superficie \(F\). Anzitutto la dimensione del sistema bicanonico, o meglio questa dimensione aumentata di un’unità, cioè il \textit{numero delle curve bicanoniche linearmente indipendenti}, costituisce un invariante assoluto che si denominerà il \textit{bigenere} \(P=P_2\). Invece il genere e il grado del sistema bicanonico forniranno degli invarianti relativi della superficie. Ma è chiaro che questi caratteri si esprimeranno per mezzo del genere lineare (relativo) \(p^{(1)}\).

Infatti, designando con \(|K|\) il sistema canonico virtuale di genere \(p^{(1)}\) e di grado \(p^{(2)}=p^{(1)}-1\), il genere del sistema bicanonico \(|2K|\) verrà dato da
\[
2p^{(1)}+p^{(2)}-1\doteq 3p^{(1)}-2 % ?? uncertain: relation sign -- \doteq / =; the small mark above the bars may be a scan speck.
\]
e il suo grado sarà
\[
4p^{(2)}=4(p^{(1)}-1).
\]

In modo analogo il sistema \(i\)-canonico della superficie \(F\), comunque abbiano o no esistenza effettiva i sistemi pluricanonici d’ordine inferiore, ci darà: un invariante assoluto che è lo \(i\)-genere \(P_i\), cioè il numero delle curve \(i\)-canoniche linearmente indipendenti, e degli invarianti relativi che ne esprimono il genere e il grado. Ma questi dipendono dal genere lineare \(p^{(1)}\) della superficie; precisamente il genere del sistema \(i\)-canonico vale
\[
\frac{i(i+1)}{2}(p^{(1)}-1)+1,
\]
e il suo grado
\[
i^2(p^{(1)}-1).
\]

% scan: s041L
% !PRINTED 62
\weberpage{62}
% head: CAPITOLO SECONDO
% stamp: Dark vertical scan border at the right edge and scattered specks.
Le considerazioni svolte nel paragrafo precedente intorno alle curve eccezionali sopra una superficie di genere \(p>0\), si estendono ora alle superficie di genere \(p=0\) il cui bigenere sia \(P>0\), o, in generale, che abbiano un \(P_i>0\). Infatti l’esistenza effettiva delle curve bicanoniche o \(i\)-canoniche di \(|iK|=|K_i|\), porta sempre che una curva di genere \(\pi\) e grado \(n\) con \(n>2\pi-2\) debba staccarsi come parte fissa dalle \(K_i\). In particolare figureranno come parti fisse delle \(K_i\) le curve eccezionali di prima specie, ciascuna contata in generale \(i\) volte. E come nel caso \(p>0\), non potranno esservi sopra la superficie che curve eccezionali di \(1^{\mathrm{a}}\) specie non connesse fra loro. Non ci soffermiamo a discorrere particolarmente delle curve eccezionali composte che non danno luogo ad osservazioni nuove.

Rileviamo soltanto che, in conseguenza di ciò che si è detto per le superficie con un plurigenere \(P_i>0\) (anche quando sia \(p=0\)) risulterà sempre definito il \textit{genere lineare assoluto} che del resto vale \(p^{(1)}+\sigma\), designando con \(p^{(1)}\) il genere lineare relativo, e con \(\sigma\) il numero delle curve eccezionali. D’altronde la superficie con \(p=0\) e con qualche plurigenere diverso da zero, potrà trasformarsi in un’altra priva di curve eccezionali, e il genere lineare di questa ci darà appunto il genere lineare assoluto.

\edsection{9. Nota storica.}{9. Nota storica.}{sec.II.9}

Il genere di una superficie algebrica \(F\) si è presentato per la prima volta ad \textsc{A. Clebsch}\origfootnote{1}{s041L-1}{\textit{Comptes-rendus de l’Académie des Sciences de Paris}, 1868.} come « numero degli integrali doppi di prima specie (cioè sempre finiti) appartenenti ad \(F\) ». Se si suppone la \(F\) d’ordine \(n\), nello spazio ordinario, gli integrandi algebrici di prima specie si formano a partire dalle superficie aggiunte \(\varphi_{n-4}\), d’ordine \(n-4\), come gli integrandi di prima specie di una curva piana d’ordine \(n\) si formano a partire dalle curve aggiunte d’ordine \(n-3\). Appare quindi che codeste \(\varphi\) o meglio le curve (canoniche) sezioni di esse con \(F\), fuori della curva doppia, hanno carattere invariante per trasformazioni birazionali della superficie. La loro invarianza ha ricevuto una prima dimostrazione algebrica da \textsc{M. Noether}\origfootnote{2}{s041L-2}{\textit{Zur Theorie des eindeutigen Entsprechens algebraischer Gebilde.} Math-Annalen Bd. 2 (1869), Bd. 8 (1874).}. Quindi lo stesso \textsc{Noether} ha notato che il sistema canonico \(|K|\) fornisce, oltre il genere \(p\) (dimensione aumentata di una unità) anche un secondo carattere invariante della superficie, che è il genere delle \(K\), cioè il genere lineare \(p^{(1)}\), definito almeno per \(p>0\), e da lui designato come « Kurvengeschlecht » in opposizione % ?? uncertain: punctuation after “Math” in footnote 2 -- “Math-Annalen” or “Math. Annalen”.
% EDITORIAL D0039 RISOLTO: The scan has a clear short horizontal hyphen after Math at the printed line end; it is not a baseline period.
%CONT>>
% scan: s041R
% !PRINTED 63
\weberpage{63}
% head: SISTEMI COVARIANTI E INVARIANTI
%<<CONT
% stamp: Scattered scan specks.
a « Flächengeschlecht ». Quanto al grado \(p^{(2)}\) del sistema canonico egli ha veduto che si esprime per il \(p^{(1)}\) colla formula
\[
p^{(2)}=p^{(1)}-1,
\]
conseguenza della proprietà fondamentale delle \(K\) di segare sulla curva generica d’un sistema lineare \(|C|\) gruppi residui della serie caratteristica. La formula indicata, che \textsc{Noether} riteneva sempre valida per i caratteri effettivi, sussiste invero, senza eccezione, per i caratteri virtuali, definiti da \textsc{Enriques}.

\textsc{F. Enriques}, riprendendo dopo un ventennio lo studio generale delle superficie algebriche dal punto di vista delle trasformazioni birazionali\origfootnote{1}{s041R-1}{\textit{Ricerche di geometria sulle superficie algebriche.} Memorie della R. Accademia delle Scienze di Torino, 1893.}, si volge anzitutto a riconoscere fino a che punto la proprietà fondamentale delle curve canoniche \(K\), supposta sussistere per rapporto ad un particolare sistema \(|C|\), valga a caratterizzare le dette \(K\): ciò avviene invero (salvo un’eccezione) se \(|C|\) è un sistema lineare irriducibile \(\infty^2\) almeno, privo di curve fondamentali proprie (abbassanti il genere delle curve residue). Da ciò l’autore trae una dimostrazione geometrica dell’invarianza delle \(K\) rispetto a trasformazioni birazionali; una trasformazione birazionale della superficie venendo scomposta in trasformazioni successive, ciascuna delle quali conserva le sezioni piane d’una stella. Questi ed altri argomenti trattati nelle citate « Ricerche » (che avremo occasione di richiamare nel seguito), vengono ripresi da \textsc{F. Enriques}, nella « \textit{Introduzione alla Geometria sopra le superficie algebriche} » del 1896\origfootnote{2}{s041R-2}{Memorie della Società Italiana delle Scienze (detta dei XL) serie III, t. X, 1896.}. Per quanto concerne il sistema canonico \(|K|\) la cosa più importante è l’estensione del teorema d’invarianza al caso \(p_g=0\), in cui \(|K|\) non ha più effettiva esistenza. Già nelle «Ricerche» l’autore aveva considerato il sistema \(|C'|\) \textit{aggiunto} ad un sistema \(|C|\), che — per \(p>0\) — si presenta come la somma \(|C'| = |C+K|\). Ora egli approfondisce lo studio di questo sistema indipendentemente dall’esistenza di \(|K|\) (cioè per \(p=0\)). Le curve \(C'\) aggiunte a \(|C|\) sopra una superficie \(F\) (come le curve d’ordine \(n-3\) aggiunte ad una curva piana d’ordine \(n\)) segano su una \(C\) generica gruppi canonici. Ma questa proprietà da sola non vale a caratterizzarle, se \(|C|\) possegga delle curve fondamentali proprie; si riesce tuttavia a definirle in ogni caso, per modo che sia soddisfatta in ordine a due sistemi lineari qualsiasi (privi di punti base), la condizione funzionale % ?? uncertain: subscript in \(p_g=0\) -- \(g\) or \(v\).
% EDITORIAL D0040 RISOLTO: The enlarged subscript has the loop and descender of g and is not a v.
\[
|C+D'|=|C'+D|.
\]
% The repeated equality sign in the inline sum is retained from both sides of the printed line break.

% scan: s042L
% !PRINTED 64
\weberpage{64}
% head: CAPITOLO SECONDO
% stamp: Dark vertical scan border at the right edge and scattered specks.
Questa \textit{proprietà fondamentale dell’aggiunzione} contiene in sè, per \(p>0\), l’invarianza del sistema canonico:
\[
|K|=|C'-C|=|D'-D|,
\]
e costituisce la generalizzazione dello stesso teorema d’invarianza, che appare ancora significativo nel caso in cui \(|K|\), definito come una differenza (e perciò a priori come \textit{virtuale}), non abbia una effettiva esistenza. Il significato di cui si tratta viene rilevato da \textsc{Enriques} sotto più aspetti:

1) La relazione fondamentale fra i sistemi aggiunti conduce a definire univocamente i \textit{caratteri virtuali} del sistema canonico \(p^{(1)}\) e \(p^{(2)}\), indipendentemente dalla sua esistenza; e fra di essi intercede, come si è detto, la relazione
\[
p^{(2)}=p^{(1)}-1.
\]

2) Può accadere che, pur non essendo \(|C|\) contenuto in \(|C'|\), il doppio o un multiplo di esso \(|sC|\) sia contenuto in \(|sC'|\); vuol dire che \textit{una curva canonica virtuale, che non ha esistenza effettiva, può dar luogo ad un sistema multiplo effettivamente esistente.}

Abbiamo già accennato agli esempi che illustrano tale possibilità. Vedremo più avanti l’utilità che presentano, particolarmente nei problemi di classificazione, i nuovi caratteri dipendenti dalle dimensioni dei sistemi multipli del sistema canonico, cioè il \textit{bigenere} e i \textit{plurigeneri}.

Alla teoria degli invarianti così disegnata da \textsc{Enriques} doveva portare un notevole progresso metodologico l’idea, affacciatasi qualche anno dopo allo stesso autore\origfootnote{1}{s042L-1}{Esposta in un corso di lezioni all’Università di Bologna dell’anno 1897-98. Cfr. il \textit{Programma} pubblicato nel Bollettino di bibliografia e storia delle Matematiche, n. 13 (1899).

Per le superficie cfr. \textit{Intorno ai fondamenti della geometria sopra le superficie algebriche.} Atti Accademia di Torino, 37, pag. 19 (1901).} di definire semplicemente i sistemi aggiunti e il sistema canonico, per mezzo degli jacobiani. L’idea è così semplice, la sua attuazione così facile, la sua portata così generale per riguardo alla teoria delle varietà ad un numero qualunque di dimensioni, che è potuto sembrare a qualcuno (e non soltanto a matematici passati per la scuola d’\textsc{Enriques}) trattarsi di cosa appartenente al patrimonio comune dei geometri e che dovrebbe esser nota da tempo almeno per ciò che concerne le curve. Ma chi rilegga le esposizioni della geometria sopra una curva di \textsc{C. Segre} e di \textsc{E. Bertini} del 1894\origfootnote{2}{s042L-2}{\textsc{C. Segre}, \textit{Introduzione alla Geometria sopra un ente algebrico semplicemente infinito.} Ann. di Mat., serie II, t. 22.

\textsc{E. Bertini}, \textit{La geometria delle serie lineari sopra una curva piana secondo il metodo algebrico.} Annali di Matematica, s. II, t. 22.} non vi trova in realtà altra dimostra%
%CONT>>
% scan: s042R
% !PRINTED 65
\unskip\weberpage{65}%
% head: SISTEMI COVARIANTI E INVARIANTI
%<<CONT
% stamp: Scattered scan specks and a small dark patch at the bottom-right edge.
zione dell’invarianza della serie canonica sopra una curva che quella di \textsc{Brill} e \textsc{Noether} legata al teorema di \textsc{Riemann-Roch}, per cui codesta serie è l’unica \(g_{2p-2}^{p-1}\) appartenente ad una curva di genere \(p\)\origfootnote{1}{s042R-1}{Cfr. \textsc{C. Segre}, l. c., n. 75; \textsc{Bertini}, l. c., n. 21 \textit{b}) e \textit{g}).}. E se così non fosse, se a quell’epoca fosse stata conosciuta la definizione della serie canonica e la dimostrazione della sua invarianza che si basa sulla proprietà delle serie jacobiane delle \(g_n\), sarebbe invero incomprensibile lo sviluppo più laborioso della teoria delle curve canoniche d’una superficie, dato da \textsc{Enriques} nelle « Ricerche » del 1893 e nell’« Introduzione » del 1896.

Aggiungiamo che l’idea dei sistemi jacobiani può essere attuata in diversi modi e così dà luogo a talune varianti nell’esposizione di altri autori.

La forma qui adottata poco lontana del resto da quella che appare nei «Fondamenti», corrisponde allo sviluppo della teoria delle curve esposto nelle «Lezioni» di \textsc{Enriques-Chisini}\origfootnote{2}{s042R-2}{Vol. III. Vedasi anche, per il trattamento analitico: \textsc{F. Enriques}, \textit{Sul teorema d’invarianza della serie canonica \(g_{2p-2}^{p-1}\) ecc.} Memorie R. Accademia di Bologna, 1914.} e — se non erriamo — raggiunge il massimo di semplicità.

% EDITORIAL: segnatura originale: \textsc{Enriques F.} - \textit{Superficie algebriche.}\hfill 5
\printersignature{\textsc{Enriques F.} - \textit{Superficie algebriche.}\quad 5}
% scan: s043L
% !PRINTED none
\weberblank{[s043L]}
% head: none
% blank page
% stamp: Dark vertical scan border at the right edge and scattered specks.


\clearpage\thispagestyle{plain}
% scan: s043R
% !PRINTED none
\weberpage{[s043R]}
% head: none
% stamp: Handwritten “67” in the upper-right margin; handwritten underline beneath “aggiunte” in the opening sentence; dark scan border at the lower-left edge and scattered specks.
\edchapter{III}{LE SUPERFICIE AGGIUNTE}{Capitolo III. LE SUPERFICIE AGGIUNTE}{LE SUPERFICIE AGGIUNTE}

Abbiamo veduto che sopra una superficie \(F\) le curve \(C'\) aggiunte ad un sistema lineare irreducibile \(|C|\) segano su queste gruppi canonici. Ora si vuol riconoscere se, e fino a qual punto, tale proprietà valga a caratterizzare le curve \(C'\). La questione si collega ad un’altra di carattere proiettivo: come le curve aggiunte al sistema delle sezioni piane \(C\) vengano segate sopra la superficie \(F\), presa nello spazio ordinario, mediante certe superficie, che diremo \textit{aggiunte} ad essa, definite dalle condizioni di passare \(i-1\) volte per ogni curva \(i\)-pla di \(F\) e da un conveniente comportamento nei punti, multipli isolati di \(F\).

\edsection{1. Superficie d’ordine \(n-3\) aggiunte ad una superficie d’ordine \(n\) priva di punti multipli propri.}{1. Superficie d’ordine n-3 aggiunte ad una superficie d’ordine n priva di punti multipli propri.}{sec.III.1}

Si consideri una superficie \(F\) d’un certo ordine \(n\), nello spazio ordinario, e si designino con \(C\) le sue sezioni piane. La \(F\) può possedere due specie di singolarità:

1) una o più curve multiple, p. es. una curva doppia che risponde, sopra una trasformata di \(F\), ad un luogo di coppie di punti, neutre per il sistema lineare trasformato di \(|C|\), ovvero una curva \(i\)-pla rispondente ad un luogo di gruppi neutri di \(i\) punti ecc.;

2) dei punti multipli particolarmente notevoli, sia punti isolati, sia punti della curva multipla, la cui singolarità non risulti dal passaggio per essi di codesta curva.

Per quel che riguarda la curva multipla ci riferiremo di regola, per semplicità di discorso, al caso di curve multiple ordinarie, e talvolta anche di una curva nodale; ma sarà agevole persuadersi che ciò che diremo vale egualmente per qualunque curva di singolarità straordinaria, quando si tenga conto delle curve multiple infinitamente vicine che la definiscono.

% scan: s044L
% !PRINTED 68
\weberpage{68}
% head: CAPITOLO TERZO
% stamp: Dark vertical scan border at the right edge and scattered specks.
Per ciò che concerne le singolarità puntuali notevoli, di cui occorre tener conto nelle considerazioni di questo capitolo, preciseremo la distinzione ricorrendo al seguente criterio: un punto multiplo \(O\) di \(F\) si dirà \textit{improprio} se il genere di una sezione piana per \(O\) uguaglia il genere delle sezioni con piani che non passano per \(O\); invece si dirà \textit{punto multiplo proprio} un punto singolare di \(F\) che \textit{abbassa il genere delle sezioni piane.}

Se ci riferiamo ad una superficie trasformata di \(F\) su cui alla curva multipla risponda un luogo di gruppi di punti neutri per \(|C|\), ad un punto multiplo improprio \(O\) di \(F\) corrisponderà in generale un gruppo \textit{neutro} \(G\) di punti (presentanti una sola condizione alle \(C\) che debbono contenerli); così p. es. ad un punto che sia triplo per la superficie \(F\) e per la sua curva doppia corrisponde in generale una terna neutra di punti; ad un punto \((i+1)\)-plo per \(F\) e per una sua curva \(i\)-pla un gruppo neutro di \(i+1\) punti ecc.

Può darsi tuttavia che, per una trasformazione di \(F\), il punto \(O\) dia luogo ad una curva \(\theta\) fondamentale per il sistema \(|C|\), la quale può esser composta in generale di tante parti eccezionali sconnesse quanti sono i punti del gruppo \(G\) di cui si è discorso innanzi; e convien notare che questa formazione della \(\theta\) è indicata, in modo caratteristico, dalla proprietà di essere \textit{curva fondamentale impropria}, non abbassante il genere delle \(C\). Infatti tale proprietà porta che ciascuna parte connessa di \(\theta\) sia di genere zero ed abbia un solo punto a comune con una curva residua rispetto a \(|C|\) (o ad un suo multiplo), e così questa parte viene caratterizzata come una curva eccezionale. % ?? uncertain: tiny raised mark after the first \(\theta\) -- \(\theta\) with a scan speck or \(\theta'\).

Ciò premesso si consideri una superficie \(F\) che \textit{non sia rigata}, le cui sezioni piane \(C\) sieno di ordine \(n\) e di genere \(\pi>1\); e sopra \(F\) una curva \(L\) che seghi le \(C\) secondo gruppi canonici. In un piano generico il gruppo di punti segato da \(L\) apparterrà ad una curva \(\varphi_{n-3}\) d’ordine \(n-3\), aggiunta alla sezione piana \(C\) di \(F\), costituendo l’intersezione completa di questa fuori dei punti multipli (dove un punto \(i\)-plo di \(C\) assorbe \(i(i-1)\) intersezioni); la \(C\) non può avere altre intersezioni colla \(\varphi_{n-3}\) fuori di \(L\), a meno che sia riducibile.

Si assuma nello spazio una retta generica \(a\) e, in ogni piano per essa, la \(\varphi_{n-3}\) aggiunta alla sezione piana determinata dal gruppo dei punti in cui la \(L\) incontra il piano; al variare del piano stesso per \(a\), la \(\varphi_{n-3}\) descriverà una superficie \(\Phi\), passante \(i-1\) volte per ogni curva \(i\)-pla di \(F\)\origfootnote{1}{s044L-1}{L’affermazione si può verificare analiticamente riducendosi al caso di una superficie descritta dalla traslazione normale di una curva piana con un punto \(s\)-plo (\(s=i-1\)) dove è evidente che il punto \(s\)-plo descrive una retta \(s\)-pla.}, che a priori potrà supporsi d’ordine \(n-3+h\), contenendo la retta \(a\) come \(h\)-pla. % ?? uncertain: raised mark beside the first \(\varphi_{n-3}\) in this paragraph -- \(\varphi_{n-3}\) with a scan speck or \(\varphi'_{n-3}\).

% scan: s044R
% !PRINTED 69
\weberpage{69}
% head: LE SUPERFICIE AGGIUNTE
% stamp: Dark scan border at the left edge and scattered specks.
Ora se \(h>0\), l’intersezione totale della \(F\) con la \(\Phi\) dovrà contenere i punti comuni ad \(a\) e ad \(F\), che sono fuori di \(L\) (e della curva multipla). Se \(A\) è uno di questi punti, ci sarà almeno un piano, tangente a \(\Phi\) in \(A\), la cui sezione con \(F\) sarà intersecata dalla \(\varphi_{n-3}\) aggiunta nel punto \(A\) fuori del gruppo canonico che sta su \(L\), e perciò risulta riducibile. Ma per un noto teorema di \textsc{Kronecker-Castelnuovo}\origfootnote{1}{s044R-1}{Cfr. \textsc{Enriques-Conforto}, Op. cit., Libro II, cap. II, \S~7, pag. 273.} la superficie \(F\), avendo allora \(\infty^2\) sezioni piane riducibili dovrebbe essere una superficie di \textsc{Steiner} a sezioni piane di genere zero, ciò che contraddice alle nostre ipotesi. Si conclude che le curve \(L\) secanti sulle sezioni piane \(C\) gruppi canonici sono intersezioni di \(F\), fuori della curva multipla, con superficie aggiunte \(\Phi_{n-3}\) d’ordine \(n-3\). Quindi le curve \(C'\) aggiunte a \(|C|\) dovranno essere contenute (totalmente o parzialmente) nel sistema \(|L|\). Ora le \(C'\) e le \(L\) sono curve dello stesso ordine \(2\pi-2\), cioè hanno lo stesso numero d’intersezioni colle \(C\), e perciò non possono differire che per curve fondamentali di \(|C|\); ciò significa che il sistema \(|C'|\) (ove sia meno ampio di \(|L|\)) dovrà dedursi da \(|L|\) per l’imposizione di qualche punto base, che cade in un punto multiplo \(O\) della \(F\); ma le \(C'\) non sono assoggettate ad avere alcun punto base in un punto multiplo improprio \(O\), perchè anche le aggiunte alle sezioni piane \(C\), per \(O\), che sono ancora di genere \(\pi\), incontrano queste sezioni in \(2\pi-2\) punti (almeno virtualmente) fuori di \(O\).

Ciò che si è detto può chiarirsi riferendosi ad una superficie trasformata di \(F\), su cui al punto \(O\) corrisponda un gruppo neutro di punti semplici, \(G\): qui il sistema \(|C'_1|\) aggiunto a \(|C_1|\) (che è dedotto da \(|C|\) coll’imposizione di punti semplici) non differisce da \(|C'|\). Enunceremo dunque il teorema:

\textit{Sopra una superficie \(F\), d’ordine \(n\), dello spazio ordinario, non rigata e priva di punti multipli propri, le curve aggiunte alle sezioni piane \(C\) sono caratterizzate dalla proprietà di segare gruppi canonici sulle dette \(C\); ed il sistema completo di queste aggiunte viene segato su \(F\) dalle superficie \(\Phi_{n-3}\) d’ordine \(n-3\) (aggiunte ad essa) passanti \(i-1\) volte per ogni curva \(i\)-pla di \(F\).}

\edsection{2. Superficie aggiunte d’ordine qualsiasi.}{2. Superficie aggiunte d’ordine qualsiasi.}{sec.III.2}

Questo risultato si estende alle curve \(L\) che segano sulle sezioni piane \(C\) di \(F\) gruppi appartenenti alla serie somma o differenza della serie canonica e di un multiplo, secondo un numero intero qualsiasi \(r\), della serie caratteristica: nelle condizioni anzidette (\(F\) non rigata e priva di punti multipli propri) le \(L\) vengono segate su \(F\) (fuori della curva multipla) da superficie \(\Phi_{n-3+r}\), d’ordine \(n-3+r\), aggiunte
%CONT>>
% scan: s045L
% !PRINTED 70
\weberpage{70}
% head: CAPITOLO TERZO
%<<CONT
% stamp: Dark vertical scan border at the right edge and scattered specks.
ad essa. Sia dapprima \(r>0\). Basta ripetere il ragionamento precedente, con una sola avvertenza: quando sia \(r\geq3\) per il gruppo \(\varphi\) segato da \(L\) su un piano per la retta \(a\) passano infinite curve \(\varphi_{n-3+r}\) aggiunte alla sezione \(C\) di questo, le quali formano un sistema lineare di dimensione \(\dfrac{(r-1)(r-2)}{2}\); e conviene determinare una di queste curve assegnandone, per es., un conveniente numero di punti, su \(a\) e poi un certo numero di tangenti e di curve osculatrici. Se poi \(r<0\), si sommi alla curva \(L\) la curva \(H\) segata su \(F\) da una generica superficie \(\Psi\) d’ordine \(|r|\). La curva \(L+H\) sega allora sulle sezioni piane di \(F\) gruppi della serie canonica e pertanto esiste una \(\Phi_{n-3}\), aggiunta ad \(F\), che passa per essa. Ma la detta \(\Phi_{n-3}\), avendo in comune con \(\Psi\) la curva \(H\), d’ordine \(|r|n > |r|(n-3)\) si spezza nella \(\Psi\) stessa e in una \(\Phi_{n-3+r}\) ancora aggiunta ad \(F\), e passante per \(L\). % The repeated greater-than sign is retained from both sides of the printed line break.

Dalle considerazioni precedenti segue, in particolare:

\textit{Sopra una superficie d’ordine \(n\) dello spazio ordinario, non rigata e non possedente punti multipli propri, le superficie \(\Phi_{n-4}\) passanti \(i-1\) volte per ogni curva \(i\)-pla di \(F\)} (cioè \textit{aggiunte} ad essa) \textit{segano il sistema completo \(|K|\) delle curve canoniche} (\textit{impure}), come già si era annunciato al Cap. II, \S~4; \textit{e in generale le \(\Phi_{n-4+r}\) con \(r>0\) segano completamente su \(F\) il sistema aggiunto al sistema \(r\)-plo delle sezioni piane, che di regola è segato solo in parte dal sistema di tutte le superficie d’ordine \(r\).}

\textit{Inoltre questo sistema \(|(rC)'|=|C'+(r-1)C|=|K+rC|\) è caratterizzato dalla proprietà delle sue curve di segare sulle sezioni piane \(C\) di \(F\) gruppi appartenenti alla serie somma della serie canonica e dello \(r\)-plo della serie caratteristica.}

\edsection{3. Superficie rigate.}{3. Superficie rigate.}{sec.III.3}

Il teorema precedente è stato stabilito sulla base di due ipotesi restrittive:

1) che la \(F\) non sia rigata nè una superficie di Steiner, e che

2) non possegga punti multipli propri.

Vediamo subito che la restrizione di non essere rigata è veramente essenziale. Se la \(F\) è una \textit{superficie rigata d’ordine \(n\) e di genere \(\pi>0\)} (che per semplicità di discorso assumeremo dotata soltanto di curva doppia) \textit{non esistono curve \(C'\) aggiunte alle sezioni piane \(C\), nè superficie aggiunte \(\Phi_{n-3}\), ma si hanno gruppi di \(2\pi-2\) generatrici secanti le \(C\) secondo gruppi canonici.}

Infatti, introducendo il sistema canonico virtuale \(|K|\), potremo scrivere
\[
|C'|=|C+K|;
\]
%CONT>>
% scan: s045R
% !PRINTED 71
\weberpage{71}
% head: LE SUPERFICIE AGGIUNTE
%<<CONT
% stamp: Dark scan border at the left edge, a dark patch at the lower-right edge, and scattered specks.
le \(C\) sono unisecanti le generatrici, le \(K\) le segano in \(-2\) punti\origfootnote{1}{s045R-1}{Le curve del sistema \(|2C+K|\), essendo spezzate in \(2\pi-2+n\) generatrici, le segano in zero punti; quindi le curve \(K\) segano le generatrici in \(-2\) punti.}, quindi le \(C'\) dovrebbero segare le generatrici in
\[
1-2=-1
\]
punti, quindi dovrebbero contenerle, ciò che costituisce un assurdo.

D’altra parte si vede pure che non possono esistere superficie \(\Phi_{n-3}\) d’ordine \(n-3\), passanti per la curva doppia di \(F\), giacchè questa curva incontra le generatrici in \(n-2\) punti.

Si può aggiungere che, sulla rigata \(F\), esistono pure gruppi di \(2\pi-2+(r-1)n\) generatrici secanti sulle sezioni piane \(C\) gruppi della serie somma della serie canonica e dello \(r\)-plo della serie caratteristica, ma tranne per \(r=2\), codesti gruppi non costituiscono curve aggiunte ad \(|rC|\) e non sono intersezioni di superficie aggiunte \(\Phi_{n-4+r}\), giacchè le curve di \(|(rC)'|\) debbono incontrare le generatrici di \(F\) in \(r-2\) punti (e non in zero).

Invece, per \(r\geq2\), \textit{le superficie \(\Phi_{n-4+r}\) passanti per la curva doppia, segano sulla rigata curve aggiunte al sistema \(r\)-plo delle sezioni piane.} Ciò si prova osservando che codeste curve debbono segare gruppi canonici sulle curve di \(|rC|\), almeno su quelle segate sopra \(F\) da superficie generali d’ordine \(r\), \(\Phi_r\). Infatti sopra una \(\Phi_r\) le curve intersezioni delle \(\Phi_{n-4+r}=\Phi_{r+n-4}\) risultano aggiunte alle intersezioni della \(F\).

Che poi il \textit{sistema segato sulla rigata \(F\) dalle \(\Phi_{n-4+r}\) sia sempre completo}, risulterà dimostrato nel seguito, quando apprenderemo a valutarne la dimensione.

\edsection{4. Proprietà caratteristica delle curve canoniche e delle curve aggiunte ad un sistema lineare.}{4. Proprietà caratteristica delle curve canoniche e delle curve aggiunte ad un sistema lineare.}{sec.III.4}

Volendo enunciare i risultati ottenuti sotto forma invariantiva, dovremo riferirci ad una superficie \(F\) e ad un sistema lineare irriducibile \(|C|\) sopra di essa, e caratterizzare rispetto a \(|C|\) le curve canoniche e le curve aggiunte alle \(C\). Ma prima di esporre questo enunciato conviene liberarci da una restrizione superflua. Invero noi abbiamo supposto che il sistema \(|C|\) sia, non soltanto irriducibile, bensì anche \textit{semplice di dimensione \(r\geq3\)}, sicchè \(\infty^3 C\) possano apparire come le sezioni piane d’una superficie trasformata.

Possiamo affrancarci da questa restrizione.

Se \(|C|\) sia (o contenga) una rete di curve irriducibili di genere
%CONT>>
% scan: s046L
% !PRINTED 72
\weberpage{72}
% head: CAPITOLO TERZO
%<<CONT
% stamp: Dark vertical scan-edge band at the right and scattered scanning specks.
$\pi$ e grado $n$, possiamo supporre che questa venga segata sopra una superficie $F$, d’ordine $m>n$, dai piani d’una stella, col centro in un punto $s$-plo $O$ $(m=n+s)$. % ?? uncertain: the relation between m and n -- > or \geq.
% EDITORIAL D0043 RISOLTO: The inequality is a clear > with no lower equality stroke.

Allora, a partire da una curva $L$ che seghi gruppi canonici sopra le $C$, ed in relazione ad un fascio di piani il cui asse $a$ passi per $O$, si costruirà, come innanzi, una superficie aggiunta $\Phi_{m-3}$ secante su $F$ la $L$, la quale passerà $i-1$ volte per ogni curva $i$-pla di $F$, ed avrà la molteplicità $s-1$ nel suo punto $s$-plo $O$; risulterà quindi $L=C'$. % ?? uncertain: the sign in "passerà i-1" -- minus or equals.
% EDITORIAL D0044 RISOLTO: The sign between i and 1 is a single horizontal minus, consistent with the repeated multiplicity statement on the same page.

La possibilità di codesta costruzione riposa però su due ipotesi:

1) che non esistano $\infty^1$ sezioni piane per $O$ riducibili;

2) che la $F$ non possegga punti multipli propri (anche infinitamente vicino ad $O$) ovvero rette multiple per $O$, le quali potrebbero riuscire di minore molteplicità per la $\Phi_{m-3}$ che abbiamo costruita.

Guardando la cosa sotto l’aspetto invariantivo le due restrizioni si possono assommare in una sola domandando che: la rete delle $C$ non debba contenere $\infty^1$ curve riducibili. Invero la presenza d’una curva fondamentale per un sistema lineare $\infty^r|C|$, significa che ad esso appartengono $\infty^{r-1}$ curve riducibili.

Ciò posto enunceremo il seguente teorema:

\textit{Se sopra una superficie $F$ è dato un sistema lineare irriducibile $|C|$, di dimensione $r>1$, che non contenga $\infty^{r-1}$ curve spezzate, le curve $C'$ aggiunte ad esso sono caratterizzate dalla proprietà di segare gruppi canonici sulle $C$. E le curve canoniche di $F$ vengono definite dalla proprietà di segare sulle $C$ gruppi residui della serie caratteristica.}

Anche le curve aggiunte ad un sistema multiplo di $|C|$ possono definirsi dai gruppi che segano sulle $C$; la cosa è chiara senza che importi insistere su questo punto.

\edsection{5. Punti multipli isolati.}{5. Punti multipli isolati.}{sec.III.5}

Ritorniamo alla considerazione proiettiva delle sezioni piane d’una superficie dello spazio ordinario.

Vogliamo esaminare fino a qual punto i risultati ottenuti si estendano al caso in cui la superficie $F$ possegga dei punti multipli propri abbassanti il genere delle sezioni per essi. Per semplicità di discorso basterà riferirci alle curve $C'$ aggiunte al sistema delle sezioni piane $C$. Si tratta dunque di esaminare fino a che punto queste curve vengano caratterizzate dalla semplice proprietà di segare gruppi canonici sopra le sezioni piane $C$ di $F$. Sulla superficie $F$ d’ordine $n$ proiettivamente considerata, una curva $L$ secante le $C$ in gruppi canonici risulta ancora, come si è visto innanzi, intersezione di una $\Phi_{n-3}$, d’ordine $n-3$, passante $i-1$ volte per
%CONT>>
% scan: s046R
% !PRINTED 73
\weberpage{73}
% head: LE SUPERFICIE AGGIUNTE
%<<CONT
% stamp: Scattered margin specks and a small stray mark above section 6.
ogni curva $i$-pla di $F$. Ma il sistema $|L|$ può non coincidere con $|C'|$ pur contenendo quest’ultimo parzialmente: in ogni caso si avrà $|L|=|C'+\theta|$ designando $\theta$ una curva, o un insieme di curve fondamentali proprie di $|C|$. E, discorrendo proiettivamente, le superficie $\Phi_{n-3}$ secanti le $C'$ si distingueranno dalle $\Phi_{n-3}$ secanti le $L$ (soggette soltanto a passare $i-1$ volte per ogni curva $i$-pla di $F$) per la condizione di passare ulteriormente per alcuni punti multipli di $F$.

Per chiarire come sia la cosa, supponiamo che $F$ possegga un punto $s$-plo \textit{isolato}, cioè fuori della curva multipla. Si considerino i piani dello spazio passanti per un punto $P$; le sezioni di questi piani formano una rete di $|C|$, e, dal punto di vista proiettivo, la jacobiana della rete viene segata dalla superficie polare di $P$, d’ordine $n-1$, la quale passa in generale $i-1$ volte per ogni curva $i$-pla di $F$ e similmente con molteplicità $s-1$ per il punto $s$-plo $O$. Si badi però che l’intorno del punto $s$-plo $O$ corrisponde (sopra una superficie trasformata) ad una curva $\theta$ fondamentale per $|C|$, la quale fa parte (contata in generale una volta) della jacobiana $C_j$\origfootnote{1}{s046R-1}{Per ipotesi la $\theta$ non ha intersezioni variabili con le $C$. Si prenda su $\theta$ un punto generico $A$; le $C$ costrette a passare per $A$ si spezzano nella $\theta$ e nelle residue curve $C_1$ di un fascio; fra queste c’è una $C_1$ passante per $A$, che insieme a $\theta$ costituisce una curva della rete che ha un punto doppio in $A$.}. Per conseguenza, ad ottenere una intera curva jacobiana converrà sommare all’intersezione della polare suddetta l’intorno del punto $O$, quindi il sistema $|C_j|$ è in realtà segato dalle superficie d’ordine $n-1$ passanti $i-1$ volte per ogni curva $i$-pla della superficie $F$ e con molteplicità $s-2$ per il punto $s$-plo $O$. Ne segue che \textit{le curve $C'$ aggiunte alle sezioni piane verranno segate dalle superficie $\Phi_{n-3}$ soggette alle condizioni di passare $i-1$ volte per ogni curva $i$-pla di $F$, e di passare $s-2$ volte per il punto $s$-plo isolato.}

Lo stesso ragionamento può ripetersi per un punto $s$-plo proprio appartenente alla curva multipla di $F$: il punto $s$-plo è (in generale e almeno) $(s-2)$-plo per le superficie che segano su $F$ le aggiunte $C'$ alle sezioni piane.

\edsection{6. Aggiunte e subaggiunte.}{6. Aggiunte e subaggiunte.}{sec.III.6}

L’analisi precedente, relativa ad una superficie $F$ con punti multipli propri conduce a distinguere per queste:

1) le superficie soggette a passare $i-1$ volte per ogni curva $i$-pla di $F$, e

2) le superficie che, oltre a passare $i-1$ volte per ogni curva $i$-pla di $F$, sono anche soggette a particolari condizioni di passaggio
%CONT>>
% scan: s047L
% !PRINTED 74
\weberpage{74}
% head: CAPITOLO TERZO
%<<CONT
% stamp: Dark vertical scan-edge band at the right and scattered scanning specks.
per i punti multipli propri di essa, in guisa che (per un ordine comunque alto) dieno come intersezioni le curve aggiunte ad un multiplo delle sezioni piane.

Riservando a queste ultime il nome di \textit{superficie aggiunte}, daremo alle prime il nome di \textit{superficie subaggiunte}, e quindi diremo \textit{curve subaggiunte} alle sezioni piane $C$ le curve definite dal segare gruppi canonici sulle $C$. Ciò posto si avrà luogo di ricercare, per ogni particolare specie di punti multipli propri, le condizioni ulteriori che essi impongono alle superficie subaggiunte perchè diventino aggiunte. A questo proposito giova osservare che le condizioni relative ad un punto multiplo, per le aggiunte di un dato ordine, potranno risultare in parte dipendenti dalle altre condizioni imposte dalle curve multiple o da altri punti multipli; ma si può ammettere fin d’ora (e sarà più precisamente chiarito in seguito) che ciò non accade per le superficie d’ordine abbastanza elevato. In ogni modo valutando il numero di codeste condizioni senza tener conto, delle eventuali dipendenze di cui si è discorso, avremo un numero che risulterà sempre lo stesso\origfootnote{1}{s047L-1}{Quando s’impone alle superficie $\varphi$ d’ordine abbastanza alto di possedere una data singolarità puntuale, $O$, si è condotti ad imporre alle $\varphi$ una certa molteplicità nel punto $O$ e in un numero finito di punti infinitamente vicini ad $O$ appartenenti ai suoi intorni successivi, come conseguenza della quale le $\varphi$ acquisteranno eventualmente le assegnate molteplicità in curve infinitesime di codesti intorni. Così il numero delle condizioni imposte alle $\varphi$ risulterà indipendente dall’ordine di esse.} per le superficie aggiunte di qualunque ordine. % ?? uncertain: the mark after "conto" -- comma or scanning speck.
% EDITORIAL D0045 RISOLTO: The enlargement shows a dense low comma-shaped mark immediately after conto. Retain the printed punctuation despite its unusual syntax.
Se esso risponde effettivamente al numero delle $\Phi_{n-4}$ aggiunte d’ordine $n-4$, esprimerà la \textit{diminuzione del genere} di $F$ portata da codesto punto multiplo. Comunque, dicendo che il punto multiplo diminuisce di tanto il genere della superficie, ci riferiremo alla \textit{diminuzione virtuale} di codesto genere, di cui rileveremo più avanti il significato, riconoscendo che l’espressione virtuale del genere fornisce un altro carattere invariante della superficie (il suo \textit{genere numerico}), che può riuscire minore del genere (\textit{geometrico}) introdotto innanzi.

Col linguaggio che abbiamo introdotto diremo che:

\textit{Un punto $s$-plo proprio $O$ della superficie $F$} (appartenente o no alla curva multipla) \textit{è} (in generale e almeno) \textit{$(s-2)$-plo per le superficie aggiunte.}

Nel caso più semplice in cui $O$ sia isolato con cono osculatore di genere $\dfrac{(s-1)(s-2)}{2}$, la molteplicità di $O$ per le superficie aggiunte è precisamente $s-2$, mentre non si ha alcuna condizione di passaggio per le subaggiunte; perciò si può dire che esso impone a priori
%CONT>>
% scan: s047R
% !PRINTED 75
\weberpage{75}
% head: LE SUPERFICIE AGGIUNTE
%<<CONT
% stamp: Scattered scanning specks in the margins.
$\dfrac{s(s-1)(s-2)}{6}$ condizioni lineari alle superficie aggiunte d’ordine $n-4$, e così \textit{diminuisce} di $\dfrac{s(s-1)(s-2)}{6}$ \textit{il suo genere.}

In particolare un \textit{punto doppio conico non influisce sul genere} della superficie.

\textit{Esempi.} -- La superficie $F_4$ del $4^\circ$ ordine senza singolarità è di genere (superficiale) $p=1$, possedendo \textit{una} curva canonica d’ordine zero, sezione della superficie aggiunta d’ordine $n-4=0$. Se la $F_4$ acquista un punto triplo, il suo genere diminuisce di 1, diventando dunque $p=0$; d’accordo col fatto che la $F_4$ con punto triplo è razionale.

La $F_4$ con un punto doppio conico è ancora di genere 1 e perciò certo non razionale. Ma vi sono superficie razionali del $4^\circ$ ordine con un (solo) punto doppio\origfootnote{1}{s047R-1}{Superficie determinate da \textsc{M. Nöther}. Cfr. \textsc{Enriques-Conforto}, Op. cit., Libro I, cap. VIII.}: il primo tipo di esse è la $F_4$ dotata d’un tacnodo. Da questo esempio appare già che « punti doppi di specie particolare (fra quelle con cono osculatore ridotto ad un piano doppio) possono diminuire il genere della superficie »: donde si trae argomento ad un’analisi approfondita.

La superficie $F_5$ del $5^\circ$ ordine $(n=5)$ senza singolarità ha il genere $p=4$, possedendo $\infty^3$ curve canoniche che sono le sezioni dei piani dello spazio (aggiunti ad $F_5$). Se la $F_5$ acquista un punto quadruplo $(s=4)$ il suo genere diminuisce di $\dfrac{s(s-1)(s-2)}{6}=4$, d’accordo col fatto che essa diventa razionale.

La $F_5$ con retta doppia $a$ è di genere $p=2$, possedendo come curve canoniche le $\infty^1$ cubiche sezioni dei piani (aggiunti) passanti per $a$. Se essa acquista un punto triplo $O$ fuori di $a$, il suo genere si riduce ad 1. In questo caso c’è \textit{un} piano aggiunto, cioè il piano $Oa$, che sega $F_5$ secondo la \textit{curva canonica impura} costituita da \textit{tre rette} per $O$. Queste rette sono \textit{eccezionali} per $F_5$; infatti si può riconoscere che le quadriche per $O$ ed $a$ segano su $F_5$ un sistema lineare $|C|$ di genere 5 e di grado 8, mercè cui la $F_5$ si trasforma in una $F_8$ di $S_5$ contenente una cubica piana $C_3$, immagine del punto $O$, mentre le tre rette della curva canonica impura si trasformano in tre punti allineati della $C_3$ semplici per la superficie $F_8$: e proiettando la $F_8$ dalla retta che contiene questi tre punti della $C_3$ si riottiene la $F_5$. Siccome le sezioni iperpiane della $F_8$ sono curve canoniche di genere 5 nello $S_4$, il sistema di esse risulta aggiunto di se stesso, e si ha così la verifica che il genere della superficie è eguale ad 1.

% scan: s048L
% !PRINTED 76
\weberpage{76}
% head: CAPITOLO TERZO
% stamp: Dark vertical scan-edge band at the right and scattered scanning specks.
\edsection{7. Punti multipli di prima specie.}{7. Punti multipli di prima specie.}{sec.III.7}

Volendo ricercare più da vicino le condizioni imposte alle superficie aggiunte da un punto singolare della $F$, converrà distinguere anzitutto i \textit{punti multipli di prima specie} da quelli di specie superiore. Diremo di « prima specie » un punto singolare $O$ di $F$, che supponiamo abbassare il genere delle sezioni piane $C$, quando vicino ad esso non vi sia alcun altro punto multiplo abbassante il genere delle $C$ per $O$. Sopra una conveniente trasformata di $F$ il punto $O$ diventa una curva $\theta$ fondamentale pel sistema $|C|$, e ciò che caratterizza la $\theta$ come \textit{curva fondamentale di prima specie} è che « non vi è alcuna parte di essa che costituisca una curva fondamentale del sistema residuo $|C-\theta|$ ».

Le condizioni imposte alle \textit{superficie aggiunte $\Phi$ da un punto multiplo di prima specie $O$} della $F$ si esprimono semplicemente dicendo che: \textit{le $\Phi$ debbono essere segate dai piani per $O$ secondo curve che sommate ad una retta per $O$ diventino aggiunte alle sezioni di $F$.}

La dimostrazione si fa nel modo più semplice riferendosi alle $\Phi_{n-4}$ d’ordine $n-4$, e supponendo che la rete delle $C$ sezioni di $F$ coi piani per $O$ non contenga $\infty^1$ curve riducibili. Perchè allora riesce evidente che le $\Phi_{n-4}$ sopra indicate segano sulle $C$ gruppi residui della serie caratteristica, proprietà che --- nelle ipotesi fatte --- vale a definire le curve canoniche. Il ragionamento si ripete analogamente per le $\Phi_{n-4+r}$ con $r>0$, anche se manchino le $\Phi_{n-4}$.

\textit{Esempi.} -- Sia una superficie $F$ dotata di un \textit{tacnodo} $O$, cioè di un punto $O$ tale che le sezioni piane per esso abbiano due punti doppi infinitamente vicini in $O$\origfootnote{1}{s048L-1}{Cfr. \textsc{Enriques-Chisini}, Op. cit., Libro IV.}. Applicando il criterio precedente si vede che le sezioni piane delle $\Phi_{n-4}$ per $O$, prese insieme ad una retta per $O$; debbono costituire curve aggiunte alle sezioni di $F$. Ciò porta che \textit{le dette aggiunte debbono passare semplicemente per il tacnodo $O$.} % ?? uncertain: punctuation after "retta per O" -- semicolon or comma.
% EDITORIAL D0046 RISOLTO: The enlarged mark shows an upper point aligned above the lower comma after O.

Se invece si ha un \textit{oscnodo}, cioè un punto $O$ tale che le sezioni piane per esso abbiano tre punti doppi infinitamente vicini in $O$, si trova in modo simile che \textit{le aggiunte debbono passare semplicemente per l’oscnodo e toccare ivi il piano tangente ad $F$.}

\textit{Esercizio.} -- Sia $F_5$ una superficie del $5^\circ$ ordine dotata d’una retta doppia $a$ e d’un tacnodo $O$. Le aggiunte d’ordine $n-4=1$ a questa superficie si riducono al piano $Oa$, sicchè risulta $p=1$.
%CONT>>
% scan: s048R
% !PRINTED 77
\weberpage{77}
% head: LE SUPERFICIE AGGIUNTE
%<<CONT
% stamp: Scattered scanning specks in the margins.
L’intersezione della $F_5$ col piano $Oa$ (fuori di $a$) sarà costituita da due curve eccezionali: una conica e la retta $b$, che si toccano in $O$.

Si verificheranno queste affermazioni, trasformando la $F_5$ in una $F_4$ del $4^\circ$ ordine senza singolarità, contenente una retta $a'$: la quale superficie (su cui il sistema delle sezioni piane è l’aggiunto di se stesso) è di genere $p=1$ e possiede una curva canonica d’ordine zero. % ?? uncertain: punctuation after "affermazioni" -- comma or no punctuation.

Si effettuerà la trasformazione anzidetta per mezzo di una speciale trasformazione quadratica di prima specie dello spazio, definita dal sistema omaloidico delle quadriche che passano per due rette incidenti e toccano in un punto di una di esse un dato piano: le rette base saranno la $a$ e la retta $b$ intersezione del piano tacnodale in $O$ col piano $Oa$, e il piano tacnodale in $O$ sarà il piano tangente fisso di cui si è discorso. Per effetto di questa trasformazione la superficie $F_5$ diverrà una $F_4$ del $4^\circ$ ordine e la quartica ulteriore intersezione del piano tacnodale di $F_5$, fuori della retta $b$, si muterà in una retta $a'$ di $F_4$. La conica eccezionale di $F_5$ si trasformerà in un punto $O'$ di $F_4$ e la retta eccezionale $b$ in un punto $B$; le trasformate delle sezioni piane di $F_5$ saranno curve $C_7$ d’ordine 7, intersezioni ulteriori della $F_4$ colle quadriche per $a'$, aventi un punto base doppio in $O'$ e un punto base semplice in $B$.

Se si assume ad arbitrio una superficie $F_4$, contenente una retta $a'$, e su $F_4$ un punto generico $O'$ (appartenente ad una cubica $K$ sezione piana per $a'$) si può, viceversa, trasformare la $F_4$ in una $F_5$ dotata di retta doppia e di tacnodo. A tal uopo giova adoperare il sistema omaloidico delle quadriche che passano per $a'$ e per la retta $b'$ tangente in $O'$ alla cubica $K$, e toccano in $O'$ la $F_4$. Qui occorre osservare che l’imposizione di un punto doppio $O'$ alle curve $C_4$ intersezioni parziali delle quadriche per la retta $a'$, porta di conseguenza anche un punto base semplice nel tangenziale di $O'$ su $K$. % [sic] C_4 -> C_7, these residual curves were denoted C_7 earlier on the page.
Così il sistema completo delle $C_4$ senza punti base, che è di genere 6, di dimensione 6 e di grado 10, dà luogo ad un sistema sovrabbondante con un punto base doppio e uno semplice, di genere 5, di grado 5 e di dimensione 3. % [sic] C_4 -> C_7, repeated subscript discrepancy for the same curves.

\edsection{8. Piani doppi.}{8. Piani doppi.}{sec.III.8}

Facciamo una breve digressione per fermarci sopra una famiglia notevole di superficie, che viene illustrata dalle considerazioni precedenti. Si consideri la superficie $F$ definita dall’equazione $z^2=f(xy)$, dove $f$ è una curva d’ordine pari $2n$: la $F$, per proiezione dal punto all’infinito $O$ dell’asse $z$, viene rappresentata sul \textit{piano doppio con curva di diramazione $f$.} Una $f$ d’ordine pari, assunta ad arbitrio, riguardata di fronte alle trasformazioni birazionali del piano, de%
%CONT>>
% scan: s049L
% !PRINTED 78
\unskip\weberpage{78}%
% head: CAPITOLO TERZO
%<<CONT
% stamp: Dark vertical scan-edge band at the right and scattered scanning specks.
finisce una classe di superficie cui appartiene $F$; se la $f$ è d’ordine dispari si aggiunge ad essa, come linea di diramazione, la retta all’infinito del piano\origfootnote{1}{s049L-1}{Cfr. \textsc{Enriques-Chisini}, \textit{Lezioni}, Libro V, \S~12, vol. III, pag. 90.}.

La superficie $F$ è d’ordine $2n$ e, se la $f$ è priva di singolarità, possiede un solo punto multiplo proprio nel punto $O$. L’analisi di questa singolarità si riconduce all’analisi della singolarità che il punto all’infinito dell’asse $z$ presenta per la curva iperellittica sezione di un piano parallelo a codesto asse, per esempio della curva iperellittica $z^2=f(x)$, con $f$ d’ordine $2n$. Si sa\origfootnote{2}{s049L-2}{Cfr. \textsc{Enriques-Chisini}, Op. cit., Libro V, \S~12 nota, vol. III, pag. 93.} che tale curva possiede in $O$ un punto $(2n-2)$-plo a cui sono vicini e fronteggianti $n-1$ punti doppi su un ramo d’ordine $n-1$ tangente alla retta all’infinito del piano.

Quindi la superficie $F$ possiede in $O$ un punto singolare $(2n-2)$-plo definito dalla proprietà che ogni sezione piana per esso debba avere $n-1$ punti doppi prossimi ad $O$.

Ciò posto consideriamo le superficie aggiunte alla $F$ dell’ordine $2n-4$. Queste sono coni (e più propriamente cilindri) di vertice $O$, poichè sommando un piano vengono a dare come sezioni delle curve aggiunte alle sezioni di $F$; ma di più tali coni sono soggetti alla condizione di contenere i punti doppi di $F$ vicini ad $O$, e perciò da essi si stacca il piano all’infinito contato $n-1$ volte. Concludiamo che le superficie d’ordine $2n-4$ aggiunte alla $z^2=f(xy)$ constano del piano all’infinito contato $n-1$ volte e dei coni d’ordine $n-3$ col vertice $O$. \textit{Il piano all’infinito sega la $F$ secondo $2n$ rette eccezionali. Le curve canoniche del piano doppio che ha per curva di diramazione una curva generale d’ordine $2n$ sono rappresentate su questo dalle curve $C_{n-3}$ d’ordine $n-3$. Il genere del piano doppio vale}
\[
p=\frac{n(n-3)}{2}+1=\frac{(n-1)(n-2)}{2}.
\]

Ora si noti che un punto quadruplo della curva di diramazione $f$ è un tacnodo per la superficie $z^2=f(xy)$, essendo piano tacnodale il piano $z=0$, dove le rette per l’origine hanno quattro intersezioni colla superficie. Ne deduciamo che: \textit{un punto quadruplo della curva di diramazione diminuisce di 1 il genere del piano doppio.}

Similmente un punto sestuplo della curva di diramazione dà luogo ad un oscnodo di $F$ e perciò risulta doppio per le $C_{n-3}$ canoniche; in generale un \textit{punto $2i$-plo della curva di diramazione è $(i-1)$-plo per le $C_{n-3}$ canoniche e così diminuisce di $\dfrac{i(i-1)}{2}$ il genere del piano doppio.}

% scan: s049R
% !PRINTED 79
\weberpage{79}
% head: LE SUPERFICIE AGGIUNTE
% stamp: Scattered scanning specks in the margins.
Un punto multiplo $A$ di $f$, \textit{d’ordine dispari $2i+1$}, deve \textit{ritenersi}, in generale, come \textit{d’ordine $2i$}, in quanto le trasformazioni birazionali del piano, aventi in $O$ un punto fondamentale, mutano l’intorno di questo punto in una linea che fa parte della trasformata della curva di diramazione. % [sic] O -> A, the point under discussion is called A in this sentence.
Ma, per lo stesso motivo, un punto multiplo di $f$ d’ordine dispari $2s+1$, che cada infinitamente vicino ad $A$, dovrà ritenersi come $(2s+2)$-plo; e così, in particolare, se si abbiano due punti multipli infinitamente vicini $A$ e $A_1$, dello stesso ordine dispari $2i+1$, le $C_{n-3}$ canoniche dovranno passare $i$ volte $A$ per (anzichè $i-1$) e $i-1$ volte per $A_1$. % [sic] i volte A per -> i volte per A, the words are transposed in print.
Il primo esempio si ha per $i=1$: \textit{la singolarità $[3,3]$, così definita, diminuisce di 1 il genere del piano doppio}, così come accade per un punto 4-plo. Nel \S~10 ritroveremo questo risultato esaminando il punto doppio singolare che ad essa corrisponde per la superficie $F$.

\textit{Notizia storica.} -- La nozione dei piani doppi risale ad \textsc{A. Clebsch} « Ueber den Zusammenhang einer Klasse von Flächenabbildungen mit der Zweitheilung der Abelschen Funktionen »\origfootnote{1}{s049R-1}{Math. Annalen, Bd. III, 1861. Cfr. anche: Crelle’s Journal, 63.}, il quale ha rappresentato sul piano semplice i piani doppi con curva di diramazione d’ordine 2 e 4 (quest’ultimo proveniente --- secondo \textsc{Geiser} --- da una proiezione della superficie cubica) e quelli con sestica di diramazione dotata d’un punto quadruplo provenienti, per proiezione, da una superficie del $4^\circ$ ordine con retta doppia. Quindi (1878) \textsc{M. Noether}\origfootnote{2}{s049R-2}{\textit{Ueber die Einzweideutigen Ebenentransformationen.} Sitzungsberichte d. Phys-med. Societät zu Erlangen, 10.} ricercava sistematicamente tutti i piani doppi razionali, riducendoli a tre tipi birazionalmente distinti, che --- mediante una trasformazione $(1,2)$ di \textsc{R. De Paolis} --- danno luogo ai tipi delle involuzioni di \textsc{E. Bertini}. Il terzo tipo di Noether è il piano doppio con sestica di diramazione dotata di due punti tripli infinitamente vicini.

Più tardi (1896) \textsc{Enriques}, intraprendendo lo studio di varie classi di piani doppi non razionali\origfootnote{3}{s049R-3}{In particolare nella memoria \textit{Sui piani doppi di genere uno.} Memorie della Società It. d. Scienze (detta dei XL), 1896.

La teoria dei piani doppi razionali e delle involuzioni piane del $2^\circ$ ordine, elaborata e sistemata nell’insegnamento di \textsc{Enriques}, viene esposta nelle \textit{Lezioni} di \textsc{Enriques-Campedelli} (cap. 5, \S\S~57-63) e nel trattato su \textit{Le superficie razionali} di \textsc{Enriques-Conforto} (Libro II, cap. IV e V).}, indicava esplicitamente come si ottengano sul piano doppio le immagini delle curve canoniche e pluricanoniche (cfr. Cap. III, \S~11 e cap. V, \S~7).

% scan: s050L
% !PRINTED 80
\weberpage{80}
% head: CAPITOLO TERZO
% stamp: Dark vertical scan-edge band at the right and scattered scanning specks.
\edsection{9. Il genere delle curve fondamentali e l’influenza sulle aggiunte.}{9. Il genere delle curve fondamentali e l’influenza sulle aggiunte.}{sec.III.9}

Riprendiamo a considerare i punti multipli di prima specie della superficie $F$, cui rispondano \textit{curve fondamentali irriducibili} del sistema $|C|$ delle sezioni piane. Una siffatta curva fondamentale $\theta$ possiede un genere $\varrho$ e un grado $-s$ (virtuali); $s$ indicherà --- come vedremo --- la molteplicità del punto $O$ per $F$, $\varrho$ il genere del cono osculatore (semplice o multiplo) in $O$ (tenuto conto delle tangenti ai rami della curva doppia che passano per $O$, e ritenute come virtualmente inesistenti le altre eventuali generatrici doppie di codesto cono).

Sappiamo che, per la presenza del punto multiplo $O$, il sistema $|C'|$ aggiunto al sistema delle sezioni piane $C$ differisce dal sistema subaggiunto $|L|$ per lo staccamento della $\theta$ contata un certo numero di volte, e ci proponiamo di ricercare questo numero $x$ per cui sarà $|L|=|C'+x\theta|$.

Esprimendo che la $\theta$ è curva fondamentale, senza intersezioni colle curve di $|C|=|C_1+\theta|$ avremo
\[
(\theta C_1)+(\theta\theta)=0
\]
e quindi
\[
(\theta C_1)=s:
\]
il numero $s$, essenzialmente positivo, è il numero delle intersezioni di $\theta$ con le curve $C_1$ residue di essa rispetto a $|C|$, cioè, in generale, la molteplicità del punto $O$ per $F$, come si è accennato innanzi. Ciò posto, per determinare il numero $x$ sopra definito potremo partire da $|C'|$ e cercare di ampliarlo successivamente col sommarvi $\theta$, finchè sia possibile.

È certo che questo ampliamento successivo tocca ad un limite quando il numero delle intersezioni di $\theta$ colle curve sommande diventi nullo o negativo. Dunque dovremo avere:
\[
(\theta C')+x(\theta\theta)\geq 0.
\]

Scriviamo
\[
\begin{aligned}
|C|&=|C_1+\theta|,\quad |C'|=|C_1+\theta'|\\ % [sic] C_1 -> C_1' in the second system, the following line explicitly uses C_1'.
% EDITORIAL S0022 ANNULLATO: la scansione s050L stampa (\theta C_1) senza apice; errore di trascrizione corretto, nota rimossa.
(\theta C')&=(\theta C_1)+(\theta\theta')=s+2\varrho-2;
\end{aligned}
\]
la diseguaglianza precedente diviene
\[
2\varrho-2-(x-1)s\geq 0,
\]
che fissa un limite superiore al valore di $x$:
\[
x\leq 1+\frac{2\varrho-2}{s}.
\]

% scan: s050R
% !PRINTED 81
\weberpage{81}
% head: LE SUPERFICIE AGGIUNTE
% stamp: Scattered scanning marks and specks in the margins.
Dalla diseguaglianza precedente ricaviamo anzitutto che, se $\varrho=0$ anche $x=0$. Dunque le curve fondamentali irriducibili di prima specie e di genere zero d’un sistema lineare $|C|$ non influiscono sulla definizione delle curve aggiunte. O, sotto forma proiettiva: \textit{un punto multiplo di prima specie a cono osculatore} (semplice o multiplo) \textit{irriducibile di genere (virtuale) zero non influisce sul genere della superficie.}

In questa classe di punti multipli rientra anzitutto il punto doppio conico, di cui si è parlato innanzi, e poi un punto triplo che sia semplice per la curva doppia di $F$, ecc. Se nella diseguaglianza precedente facciamo $\varrho=1$, risulta $x\leq 1$. Qui si può sempre prendere $x=1$. Infatti riferiamoci ad una $F$ dotata di curva doppia sopra la quale si trovi un punto $s$-plo per cui passino $\dfrac{(s-1)(s-2)}{2}-1$ rami di codesta curva; il punto $s$-plo $O$ di cui si discorre risponde appunto ad una curva fondamentale di genere $\varrho=1$. Ora il cono osculatore ad $F$ in $O$ è un cono d’ordine $s$ che ammette un cono aggiunto d’ordine $s-3$; ciò significa che le superficie subaggiunte (d’ordine assai elevato), soggette a passare per la curva doppia, avranno in $O$ la molteplicità $s-3$, mentre le aggiunte debbono avere in esso la molteplicità $s-2$.

Nel caso dei punti multipli con $\varrho=1$, rientra anche il \textit{tacnodo}, dove la curva irriducibile $\theta$ si riduce all’intorno di $O$ sul piano tacnodale contato due volte. Infatti la $\theta$ appare qui rappresentata sul fascio doppio delle tangenti con quattro rette di diramazione che rispondono alle sezioni piane cuspidali uscenti dal tacnodo. Si ritrova così il risultato stabilito nel precedente paragrafo, che il tacnodo diminuisce di 1 il genere della superficie.

In modo simile si vede che l’\textit{oscnodo} risponde ad una curva fondamentale $\theta$ di genere $\varrho=2$, e qui il valore del numero $x$ risulta $x=2$, d’accordo con ciò che si è veduto: l’oscnodo diminuisce di 3 il genere della superficie.

Ora riferendosi al caso di una curva fondamentale di prima specie irriducibile di genere $\varrho$ qualunque, si può chiedere se il sistema $|C'|$ possa sempre ampliarsi successivamente col sommarvi $\theta$, finchè $x$ soddisfi alla diseguaglianza precedente, come si verifica nel caso di un punto multiplo isolato. La risposta è negativa: il limite superiore di $x$ non è sempre un massimo effettivamente raggiunto. Per vederlo basta un esempio. Si supponga che $F$ possegga una curva doppia passante con sei rami per un punto sestuplo $O$ il cui cono osculatore risulterà dunque di genere $\varrho=4$. In generale le tangenti a codesti rami non apparterranno ad un cono quadrico e quindi una superficie subaggiunta ad $F$ dovrà avere in $O$ la molteplicità 3, avendo ivi un cono osculatore del $3^\circ$ ordine aggiunto al cono oscu% % ?? uncertain: the final letter of "cono" before "oscu-" -- a damaged o or a punctuation-like mark.
% EDITORIAL D0048 RISOLTO: The high-resolution crop shows a broken rounded final o after con; the surviving lower-right arc is damaged type, followed by the separate space before oscu-.
% footer: Printer's signature, separate from the continuing paragraph.
% EDITORIAL: segnatura originale: \par\noindent\textsc{Enriques F.} - \textit{Superficie algebriche.}\hfill 6
\printersignature{\textsc{Enriques F.} - \textit{Superficie algebriche.}\quad 6}
%CONT>>
% scan: s051L
% !PRINTED 82
\unskip\weberpage{82}%
% head: CAPITOLO TERZO
%<<CONT
% stamp: Dark scan-edge line at the right and scattered specks.
latore di \(F\). Il punto \(O\) avrà per le superficie aggiunte la molteplicità \(4\), sicchè il sistema aggiunto alle sezioni piane si deduce dal più ampio sistema subaggiunto con lo staccamento di una sola curva \(\theta\) rispondente all’intorno di \(O\), ed il valore \(1\) è inferiore al limite che viene indicato per \(x\) dalla diseguaglianza precedente:
\[
6-(x-1)6\geq 0,
\]
che dà per il massimo valore di \(x\), \(x=2\).

\edsection{10. Singolarità puntuali di specie superiore.}{10. Singolarità puntuali di specie superiore.}{sec.III.10}

L’analisi delle singolarità puntuali di specie superiore e dell’influenza che esse hanno sul genere della superficie, si riconduce alla definizione dei punti multipli che possono cadere infinitamente vicini ad un punto singolare proprio, e che non siano definiti dal passaggio di curve singolari, proprie o infinitesime.

Se s’immagina completamente sciolta la singolarità puntuale posseduta da una superficie \(F\) in un punto \(O\), il sistema delle sezioni piane di \(F\), o egualmente un suo multiplo (che possiamo assumere di dimensione grande quanto si vuole) ci darà sulla trasformata di \(F\) un sistema lineare \(|C|\) dotato di una curva fondamentale, che risponde all’intorno del punto \(O\). Se la singolarità di \(O\) è di prima specie, codesta curva fondamentale è similmente di prima specie, sicchè staccandola da \(|C|\) nessuna parte di essa rimane fondamentale per il sistema residuo. Invece si definirà come \textit{curva fondamentale di seconda specie} (corrispondente ad un punto multiplo \(O\) di seconda specie) una curva fondamentale che contiene alcune parti costituenti curve fondamentali di prima specie per il sistema residuo. E così, in generale, si definiranno induttivamente le curve fondamentali di \textit{specie \(s\)}, riducendone il concetto a quello delle curve di specie \(s-1\).

Da questa definizione risulta che il problema di determinare l’influenza sul genere della superficie \(F\) di un punto multiplo di specie \(s\), si riconduce all’analogo problema relativo a singolarità di specie \(s-1\), e così la determinazione delle condizioni che una curva fondamentale del sistema \(|C|\) impone alle curve subaggiunte che debbano essere aggiunte a \(|C|\), si fa dipendere dalla determinazione analoga nei riguardi del sistema che si ottiene da \(|C|\) staccando codesta curva.

Riferiamoci, per semplicità, al caso di una curva fondamentale di seconda specie del sistema \(|C|\).

Sopra la superficie \(F\) che ha per sezioni piane le \(C\) abbiamo un punto multiplo \(O\) a cui sono infinitamente vicini, nell’intorno del
%CONT>>
% scan: s051R
% !PRINTED 83
\weberpage{83}
% head: LE SUPERFICIE AGGIUNTE
%<<CONT
% stamp: Scattered scan specks and small marginal marks.
primo ordine, alcuni punti multipli (propri) \(O_1, O_2, \ldots\ldotp\). Sopra una superficie trasformata, al punto \(O\) corrisponderà una curva \(\theta\) fondamentale per \(|C|\); ad \(O_1\) una curva fondamentale \(\theta_1\), ecc. Ora è interessante notare che le curve \(\theta_1, \theta_2, \ldots\ldotp\) compaiono necessariamente come parti della curva fondamentale \(\theta\), sicchè si avrà in generale
\[
\theta=\theta_0+\theta_1+\theta_2+\ldots\ldotp
\]

Infatti la trasformazione a cui si è sottoposta la superficie \(F\) si può effettuare per gradi assumendo anzitutto un sistema trasformante di superficie col punto base semplice \(O\), sicchè l’intorno di \(O\) si muti in una curva \(\theta_0\) (semplice o multipla) a cui apparterranno i punti singolari \(O_1, O_2, \ldots\ldotp\); e poi trasformando successivamente ciascuno dei punti \(O_1, O_2, \ldots\ldotp\) in una curva che risulterà generalmente connessa con la \(\theta_0\). In forza di queste connessioni quando si stacca da \(|C|\) la curva \(\theta_0\) si vengono a staccare conseguentemente anche le \(\theta_1, \theta_2, \ldots\ldotp\) le quali sono del pari curve fondamentali (senza intersezioni con le \(C\)). In altre parole la somma
\[
\theta=\theta_0+\theta_1+\ldots\ldotp
\]
costituisce una curva fondamentale \textit{monovalente} per \(|C|\), il cui staccamento da \(|C|\) importa una sola condizione. Invece il sistema residuo \((C-\theta)\) possiederà ancora le curve fondamentali \(\theta_1, \theta_2, \ldots\ldotp\), sicchè potrà dirsi che le curve
\[
\theta_0+2\theta_1+\theta_2+\ldots\ldotp
\]
e
\[
\theta_0+\theta_1+2\theta_2+\ldots\ldotp
\]
ecc., saranno curve fondamentali \textit{bivalenti} di \(|C|\), il cui staccamento da \(|C|\) importa \textit{due condizioni}.

Consideriamo per esempio il caso di un \textit{punto doppio uniplanare} (isolato). Questo è un punto doppio \(O\) in cui il cono osculatore si riduce ad un piano contato due volte; le rette per \(O\) in codesto piano hanno in generale un contatto tripunto con la superficie, salvo tre rette con contatto quadripunto che definiscono tre punti doppi \(O_1, O_2, O_3\), infinitamente vicini ad \(O\)\origfootnote{1}{s051R-1}{Cfr. per es. \textsc{Enriques-Chisini}, Op. Cit., Libro IV.}. Supporremo per semplicità che la superficie \(F\) a cui appartiene il punto uniplanare \(O\) sia una superficie cubica \(F_3\), e quindi illustreremo le cose dette riferendoci alla sua rappresentazione piana.

Come sistema rappresentativo di \(F_3\) si può prendere nel piano un sistema lineare \(\infty^3\) di cubiche \(C_3\) passanti per tre punti base in linea retta \(O_1, O_2, O_3\) e per altri tre punti infinitamente vicini a
%CONT>>
% scan: s052L
% !PRINTED 84
\weberpage{84}
% head: CAPITOLO TERZO
%<<CONT
% stamp: Scattered specks and vertical scan marks near the left edge.
questi \(O'_1, O'_2, O'_3\), dove la retta a cui appartengono \(O_1, O_2, O_3\) corrisponde al punto uniplanare \(O\) e alla parte \(\theta_0\) del suo intorno. Trasformiamo il nostro sistema di \(C_3\) con una trasformazione cremoniana dal \(4^{\mathrm{o}}\) ordine, assumendo come rete trasformante una rete di quartiche con un punto base triplo, tre punti base semplici \(O_1, O_2, O_3\) ed altri tre punti base semplici. % [sic] dal 4º ordine -> del 4º ordine, the preposition appears to be a typesetting error.
Il sistema \(|C_3|\) si trasformerà in un altro sistema di curve \(|C_9|\), con un punto base sestuplo \(P\), tre punti tripli \(A_1, A_2, A_3\) allineati, altri tre punti doppi \(B_1, B_2, B_3\), e tre punti semplici \(O_1, O_2, O_3\) che apparterranno rispettivamente alle rette \(PB_1, PB_2, PB_3\). La retta \(A_1A_2A_3\) contata due volte rappresenta la curva \(\theta_0\), mentre le \(PB_1, PB_2, PB_3\) sono le immagini di \(\theta_1, \theta_2, \theta_3\). Appare tosto che il gruppo di rette \(\theta_0+\theta_1+\theta_2+\theta_3\) costituisce una curva fondamentale monovalente per il sistema \(|C|\), giacchè staccando la retta \(\theta_0\) si staccano di conseguenza anche le rette \(\theta_1, \theta_2, \theta_3\) che la incontrano, ciascuna in un punto fuori dei punti base, % ?? uncertain: Mark after "incontrano" -- comma or scan speck.
% EDITORIAL D0049 RISOLTO: A distinct descending comma is visible between incontrano and ciascuna at the punctuation baseline.
e conseguentemente la nominata retta \(A_1A_2A_3\) si stacca una seconda volta; quindi il sistema residuo della \(\theta=\theta_0+\theta_1+\theta_2+\theta_3\) viene costituito dalle quartiche \(C_4\) che passano tre volte per \(P\), semplicemente per \(B_1, B_2, B_3\) e ancora semplicemente per \(A_1, A_2, A_3\). Per questo sistema residuo la \(\theta_0\) non è più retta fondamentale; invece sono ancora rette fondamentali le rette \(PB_1, PB_2, PB_3\).
% The plus sign in the expression for theta is printed on both sides of a line break.

L’analisi della singolarità di \(2^{\mathrm{a}}\) specie che una superficie \(F\) possiede in un \textit{punto uniplanare ordinario \(O\)}, permette di riconoscere che codesto punto \textit{non ha influenza sulle superficie aggiunte o sul genere}.

Infatti si consideri il sistema \(|C|\) delle sezioni piane di \(F\) (a cui eventualmente si potrà sostituire un suo multiplo). Le \(C_1\) sezioni piane per \(O\) formano un sistema lineare che possiede soltanto le curve fondamentali di \(1^{\mathrm{a}}\) specie e di genere zero \(\theta_1, \theta_2, \theta_3\) corrispondenti ai punti doppi \(O_1, O_2, O_3\), infinitamente vicini ad \(O\). Perciò questi punti doppi non influiscono sulle curve canoniche che restano definite dalla proprietà di segare sulle \(C_1\) gruppi residui della serie caratteristica, cioè tali che appartengano ad una curva d’ordine \(n-3\) aggiunta ad una \(C_1\), costituita da una \(\varphi_{n-4}\) e da una retta per \(O\). Dunque le curve aggiunte alle \(C_1\) sono segate dalle superficie aggiunte \(\Phi_{n-3}\) che passano semplicemente per \(O\) e, sommando la curva \(\theta\) che rappresenta l’intorno di questa, le curve \(C'\) aggiunte alle sezioni piane di \(F\) verranno segate da \(\Phi_{n-3}\) che non sono più soggette a passare per \(O\). % [sic] di questa -> di questo, the pronoun appears to refer to the point O.

Mentre il punto doppio uniplanare ordinario \(O\) non influisce sul genere di una superficie \(F\), possono darsi particolarizzazioni di questa singolarità che portano un’influenza sul genere. Si è visto già che ciò accade per il tacnodo, caratterizzato dalla proprietà che
%CONT>>
% scan: s052R
% !PRINTED 85
\weberpage{85}
% head: LE SUPERFICIE AGGIUNTE
%<<CONT
% stamp: Dark scan-edge band at the left and scattered specks.
le rette tangenti per \(O\) abbiano con \(F\) un contatto quadripunto, e perciò abbassante di due, anzichè di uno, il genere delle sezioni piane per esso. Ma più meraviglioso è il caso di un punto uniplanare \(O\) che abbassi soltanto di \(1\) il genere delle sezioni piane per esso, e che possegga infinitamente vicino un tacnodo \(O_1\), per modo che il genere delle sezioni piane per \(OO_1\) risulti abbassato di \(3\) anzichè di \(2\). In questo caso il sistema \(|C_1|\) delle sezioni piane per \(O\) possiede una curva fondamentale \(\theta\) di genere \(1\), che ha per immagine il tacnodo \(O_1\), e questa porta che le superficie \(\Phi_{n-4}\) debbono passare semplicemente per \(O_1\), ossia, per effetto di scaricamento, per \(O\). Si deduce quindi che \textit{un punto doppio uniplanare}, abbassante soltanto di \(1\) il genere delle sezioni piane per esso, \textit{a cui sia infinitamente vicino un tacnodo diminuisce di \(1\) il genere della superficie}.

\textit{Nota.} -- Noether ha indicato due tipi di superficie del quarto ordine che riescono razionali per il possesso della singolarità sopra indicata\origfootnote{1}{s052R-1}{Cfr. \textsc{Enriques-Conforto}, \textit{Le superficie razionali}, Op. cit., pag. 196 e 204.} e la razionalità di queste superficie è d’accordo col fatto che l’anzidetto punto singolare diminuisce di \(1\) il genere della superficie del \(4^{\mathrm{o}}\) ordine, e così lo riduce a zero. I due tipi di superficie di Noether danno luogo ad una rappresentazione piana rispettivamente:

1) col sistema rappresentativo di \(\infty^3\) curve \(C_7\) del \(7^{\mathrm{o}}\) ordine aventi un punto base triplo e \(9\) punti base doppi sopra una cubica fondamentale; e

2) col sistema rappresentativo di \(\infty^3\) curve \(C_9\) del \(9^{\mathrm{o}}\) ordine aventi \(8\) punti base tripli, uno doppio e uno semplice, sopra una cubica fondamentale.

In ambedue i casi la cubica fondamentale del sistema rappresentativo, costituisce una \(\theta\) di \(2^{\mathrm{a}}\) specie che è fondamentale anche per il sistema residuo.

Ed appare dalla rappresentazione indicata che codesta \(\theta\) costituisce sopra la superficie \(F_4\), ovvero sopra una trasformata in cui la singolarità venga sciolta senza introdurre nuove curve eccezionali, una curva di genere \(1\) e di grado \(-1\): infatti la cubica piana da cui si tolgano \(10\) punti base è precisamente una curva di grado \(-1\)\origfootnote{2}{s052R-2}{Qui occorre avvertire che il tipo 2) delle superficie di Noether presenta una complicazione proiettiva, consistente in ciò che il piano tacnodale di \(O_1\), passa per \(O\), % ?? uncertain: Mark after O_1 in this footnote -- comma or scan speck.
% EDITORIAL D0050 RISOLTO: The small low mark immediately following the subscript 1 has a comma-shaped descender and belongs to the printed text.
sicchè le sezioni piane per \(OO_1\) contengono un terzo punto doppio fisso allineato con \(OO_1\) e si spezzano nella retta \(OO_1\) è in una cubica di genere \(1\) per cui essa è tangente di flesso. % [sic] è in una cubica -> e in una cubica, a conjunction is required here.
Per effetto di tale circostanza l’analisi della singolarità dev’esser fatta ricorrendo, invece che al sistema delle sezioni piane, ad un sistema di intersezioni con superficie d’ordine più elevato. Cfr. \textsc{Enriques-Conforto}, Op. cit., pagg. 206-7.}.

% scan: s053L
% !PRINTED 86
\weberpage{86}
% head: CAPITOLO TERZO
% stamp: Stray diagonal marks above the first 1 in the second display and above "ha"; dark scan borders at the bottom and right.
A questo riguardo l’analisi della singolarità offerta dalle superficie \(F_4\) di Noether si estende alle singolarità consimili di una superficie qualunque: \textit{la singolarità di \(2^{\mathrm{a}}\) specie costituita da un punto doppio con infinitamente vicino un tacnodo risponde in generale ad una curva fondamentale di genere \(1\) e di grado \(-1\) per il sistema lineare \(|C|\) delle sezioni piane della superficie}\origfootnote{1}{s053L-1}{Cfr. D. \textsc{Pedoe}, \textit{On the canonical System of certain algebraic Surfaces.} Proc. Phil. Soc. Cambridge, vol. 33 (1937).}.

La dimostrazione diretta dell’asserto (indipendente dalla possibilità di approssimare la superficie data con una \(F_4\) avente la medesima singolarità) si può dare: anzitutto mostrando che una curva fondamentale di genere \(1\) e di grado \(-1\) di un sistema lineare \(|C|\) produce la detta singolarità sopra una superficie che abbia \(\infty^3\) \(C\) come sezioni piane; e viceversa partendo da una \(F\) che possegga la detta singolarità, che venga sciolta, e calcolando i caratteri della curva fondamentale che ad essa risponde.

Si assuma che la \(\theta\) sia una curva fondamentale per un sistema \(|C|\) di genere \(\pi\) e grado \(n\), curva che abbia a sua volta il genere \(\varrho=1\) e il grado \(\nu=-1\); appare tosto che il numero delle intersezioni % ?? uncertain: Genus and degree letter forms throughout this page -- \varrho or \rho; Greek \nu or italic Latin v.
% EDITORIAL D0051 RISOLTO: The genus letter has the rounded loop and curled lower terminal of varrho, while the degree letter has the short left entry and pointed right stroke of nu; the same forms recur on the facing page.
\[
(C-\theta)\theta=C\theta-\theta\theta=-\nu=1,
\]
cioè le curve del sistema residuo \(|C_1|=|C-\theta|\) hanno \textit{una} intersezione, necessariamente fissa, con la \(\theta\) (\(C_1\theta=1\)). Per conseguenza il genere \(x\) di \(|C_1|\) verrà dato dalla equazione
\[
\pi=x+1+1-1=x+1,
\]
cioè
\[
x=\pi-1
\]
ed il suo grado (virtuale) si calcolerà dalla formula
\[
n=CC=(C_1+\theta)(C_1+\theta)=C_1C_1+\theta\theta+2 % ?? uncertain: Small mark above the second C_1 in the final product -- scan speck (best reading) or a prime/accent.
% EDITORIAL D0052 RISOLTO: The irregular tiny spot sits over the body of the second C rather than in the upper-right position of a prime; it is a scan or printing blemish.
\]
che dà
\[
C_1C_1=n-1
\]
(invece il grado effettivo di \(|C_1|\), essendovi un punto base su \(\theta\), sarà \(n-2\)).

Ciò posto sopra la superficie \(F\) che ha come sezioni piane \(\infty^3\) \(C\), la \(\theta\) avrà per immagine un punto \(O\), che abbassa di \(1\) il genere delle sezioni piane per esso e quindi è doppio per \(F\); e vicino ad \(O\) vi sarà un punto singolare \(O_1\) che abbassa di due il genere delle sezioni (piane o superficiali) per esso, e quindi è un tacnodo.

Viceversa, se si parte da una superficie \(F\) che possegga un punto doppio \(O\) con infinitamente vicino un tacnodo \(O_1\), al punto \(O\) ri%
%CONT>>
% scan: s053R
% !PRINTED 87
\unskip\weberpage{87}%
% head: LE SUPERFICIE AGGIUNTE
%<<CONT
% stamp: Shading at the upper-right corner, dark scan borders at the left and bottom, and scattered specks.
sponderà sopra una conveniente trasformata di \(F\), una curva \(\theta\) fondamentale per il sistema delle sezioni piane \(C\), che a priori si comporrà di due parti \(\theta_0\) e \(\theta_1\), la prima di genere zero corrispondente all’intorno del \(1^{\mathrm{o}}\) ordine di \(O\) (escluso il tacnodo), l’altra di genere \(1\) corrispondente all’intorno di \(O_1\): la curva
\[
\theta=\theta_0+\theta_1
\]
fondamentale per \(|C|\), e necessariamente connessa, rappresenterà l’intorno completo del punto \(O\).

Ora si può valutare il genere \(\varrho\) e il grado \(-\nu\) della curva \(\theta\). % ?? uncertain: Genus and degree letter forms throughout this page -- \varrho or \rho; Greek \nu or italic Latin v.
% EDITORIAL D0053 RISOLTO: The letterforms match those on p. 86. The separate minus before nu is unmistakably present, although the signed-degree notation is inconsistent, as the existing sic comment states. % [sic] grado -\nu -> grado \nu, the signed degree is denoted by \nu on the preceding page and in the concluding formula.
Esprimendo che il distacco di essa dal sistema \(|C|\) di genere \(\pi\) abbassa di \(1\) il genere della \(C\), avremo
\[
\begin{aligned}
|C|&=|C_1+\theta|\\
\pi&=(\pi-1)'+\varrho+x-1, % ?? uncertain: Small stroke after (\pi-1) -- prime or scan/printing speck.
\end{aligned}
\]
dove \(x\geq 1\) designa il numero dei punti, di connessione delle \(C\) spezzate, comuni a \(\theta\) e \(C_1\); e di qui si deduce
\[
x=1 \qquad \text{e} \qquad \varrho=1.
\]

Quindi, calcolando il numero delle intersezioni \(\theta C_1\) in funzione del grado di \(\theta\), si avrà
\[
C_1\theta=(C-\theta)\theta=-\theta\theta=-\nu=1.
\]

Dopo ciò valutiamo il grado \(-\nu_1\) di \(\theta_1\) scrivendo che il passaggio delle sezioni \(C_1\) di \(F\) per \(O_1\) abbassa di \(2\) il loro genere; avremo
\[
\begin{aligned}
|C_1|&=|C_2+\theta_1|\\
\pi-1&=\pi-3+1+y-1
\end{aligned}
\]
dove \(y\) designa il numero delle intersezioni
\[
C_2\theta_1=(C_1-\theta_1)\theta_1=-\theta_1\theta_1=\nu_1,
\]
e di qui si ricava
\[
\nu_1=y=2.
\]

Ora, designando con \(-\nu_0\) il grado della \(\theta_0\) si avrà
\[
-\nu=-\nu_0-\nu_1+2=-1, % [sic] -\nu=-\nu_0-\nu_1+2=-1 -> \nu=-\nu_0-\nu_1+2=-1, the initial sign conflicts with -\nu=1 above and \nu=-1 below.
\]
e quindi
\[
\begin{aligned}
\nu_0+\nu_1&=3,\\
\nu_0&=1.
\end{aligned}
\]

In conclusione la curva \(\theta_0\) di genere \(0\) e grado \(-1\), resulta una curva eccezionale che può mutarsi in un punto semplice della curva \(\theta_1\) di genere \(1\), la quale diventa così una \(\theta\) di grado
\[
\nu=-1. % ?? uncertain: Tiny mark after -1 -- full stop or scan speck.
\]
(conservando il genere \(\varrho=1\)).

% scan: s054L
% !PRINTED 88
\weberpage{88}
% head: CAPITOLO TERZO
% stamp: Dark scan borders along the right and bottom edges and scattered specks.
La singolarità costituita da un punto doppio a cui è vicino un tacnodo s’incontra nello studio dei piani doppi, quando la curva di diramazione possieda due punti tripli infinitamente vicini \(O, O_1\), o come si dice, un punto \([3,3]\). In questo caso si consideri per esempio il sistema delle curve \(C\) aventi un certo genere \(\pi\) che è rappresentato sul piano doppio dal sistema \(\infty^5\) delle coniche; le coniche passanti per \(O\) rappresentano curve \(C_1\) di genere \(\pi-1\), ma quelle per \(O\) ed \(O_1\) rappresentano curve di genere \(\pi-3\). Dunque il punto \(O\) è un punto doppio della superficie \(z^2=f(xy)\) che abbassa soltanto di \(1\) il genere delle sezioni (piane o superficiali) per esso, mentre il punto \(O_1\) ad esso vicino abbassa ulteriormente di due il genere di codeste sezioni, e perciò è un tacnodo. Così ritroviamo che: \textit{un punto \([3,3]\) appartenente alla curva di diramazione d’ordine \(2n\) di un piano doppio} deve ritenersi come un punto base semplice per le curve d’ordine \(n-3\) immagini delle curve canoniche (cfr. \S~7) e quindi \textit{diminuisce di \(1\) il genere del piano doppio}.

Le due superficie \textit{razionali} del \(4^{\mathrm{o}}\) ordine di Noether che sopra abbiamo ricordato vengono proiettate dal loro punto singolare in un piano doppio, con curva di diramazione del \(6^{\mathrm{o}}\) ordine, la quale possiede rispettivamente per i due tipi un punto quadruplo, e un punto \([3,3]\): così viene confermato che la singolarità di quelle superficie diminuisce di \(1\) il genere delle superficie del \(4^{\mathrm{o}}\) ordine.

\textit{Esercizi.}

\(1^{\mathrm{o}})\) Si riconosca che la singolarità (isolata) costituita da due punti tripli infinitamente vicini \(OO_1\) impone alle aggiunte due punti base semplici e perciò ha sul genere della superficie la stessa influenza che due punti tripli distinti.

Quindi si calcoli il genere d’una superficie \(F_5\) del \(5^{\mathrm{o}}\) ordine con retta doppia e due punti tripli infinitamente vicini \(OO_1\) fuori della retta: si troverà
\[
p=0.
\]

Questa deduzione è d’accordo col fatto che la \(F_5\) è razionale, possedendo un fascio di quartiche sezioni piane razionali per \(OO_1\).

Si può ottenere la rappresentazione piana della \(F_5\) sopra indicata valendosi di una trasformazione quadratica di \(2^{\mathrm{a}}\) specie, mediante il sistema omaloidico delle quadriche passanti per la retta doppia di \(F_5\) e per altri tre punti di essa, due dei quali siano i punti tripli infinitamente vicini \(O\) ed \(O_1\). In tal guisa la \(F_5\) si trasformerà in una superficie del \(4^{\mathrm{o}}\) ordine \(F_4\) con retta doppia.

\(2^{\mathrm{o}})\) Si assuma una superficie \(F_4\) del \(4^{\mathrm{o}}\) ordine, priva di singolarità, contenente una retta \(a\): il sistema lineare delle sezioni piane sommato al doppio del fascio delle cubiche \(K_3\) sezioni per \(a\), dà un sistema lineare di curve \(C\) di genere \(3+1+1+(3-1)
%CONT>>
% scan: s054R
% !PRINTED none
\text{\weberpage{[s054R]}}
% head: LE SUPERFICIE AGGIUNTE
%<<CONT
% stamp: The upper-right corner is obscured; no page number is visible. Dark scan border along the bottom.
+(3-1)=9\)\edFlush{} di cui si trova facilmente che il grado è \(16\), e la dimensione è eguale a \(9\). Si imponga alle \(C\) di possedere un punto triplo in un punto \(P\) di \(F_4\) fuori di \(a\), che sia un flesso per la \(K_3\) passante per esso. Si ottiene così un sistema lineare \(\infty^3\) di curve di genere \(6\) e di grado \(7\), che possiede la curva fondamentale \(K_3\), di genere \(1\) e di grado \(-1\). Per mezzo di questo sistema la \(F_4\) si lascia trasformare in una superficie \(F_7\) del \(7^{\mathrm{o}}\) ordine possedente una retta quadrupla, una retta tripla (che riduce a \(6\) il genere delle sezioni piane) e un punto doppio con infinitamente vicino un tacnodo: si ritrova così, per la \(F_7\), il genere della \(F_4\) che è \(p=1\). C’è una superficie aggiunta d’ordine \(7-4=3\), che si riduce al piano contenente le rette quadrupla e tripla contato due volte e al piano proiettante dalla retta quadrupla il punto doppio isolato, che sega \(F_7\) secondo una cubica eccezionale.

Ricercare la trasformazione inversa della \(F_7\) nella \(F_4\).

\textit{Nota.} -- In ciò che precede abbiamo fatto alcune osservazioni in rapporto ad esempi interessanti, ma non riteniamo, in alcun modo, di avere esaurito la teoria delle curve fondamentali dei sistemi lineari, nelle sue relazioni colle singolarità. Occorre perciò esaminare e meditare una serie di esempi più copiosa, e riconoscere le complicazioni cui danno luogo. Per approfondire tale ordine di questioni, invitiamo lo studioso a considerare talune, pure assai semplici, singolarità che possono appartenere alla curva di diramazione di un piano doppio: punti quadrupli infinitamente vicini, serie di \(4\) o \(6\) ecc. punti tripli infinitamente vicini, serie di punti tripli infinitamente vicini a punti quadrupli ecc. Si tratta di definire questi casi caratterizzando la curva fondamentale per il sistema delle sezioni piane della superficie \(F\), che risponde ai detti punti singolari, e le circostanze cui dà luogo il suo staccamento\origfootnote{1}{s054R-1}{Cfr.: A. \textsc{Franchetta}, \textit{Sui punti doppi isolati delle superficie algebriche.} Rend. Acc. Lincei, 1945. Note I e II.}. % ?? uncertain: Mark after "rapporto" -- scan speck (best reading) or full stop.

\edsection{11. Curve bicanoniche e superficie biaggiunte.}{11. Curve bicanoniche e superficie biaggiunte.}{sec.III.11}

Data nello spazio ordinario una superficie \(F\) di un certo ordine \(n\) abbiamo appreso a costruire su di essa le curve canoniche come intersezioni di superficie \(\Phi_{n-4}\) d’ordine \(n-4\) aggiunte ad \(F\), e più in generale le curve aggiunte al sistema \(r\)-plo delle sezioni piane \(C\) di \(F\), mediante le superficie aggiunte \(\Phi_{n-4+r}\), le quali segano sempre su \(F\) un sistema lineare completo. Ora sarà facile determinare sopra \(F\) le curve bicanoniche ed anche le curve seconde aggiunte al si%
%CONT>>
% scan: s055L
% !PRINTED 90
\unskip\weberpage{90}%
% head: CAPITOLO TERZO
%<<CONT
% stamp: Dark scan-edge line at the right, a curved mark in the left margin, and scattered specks.
stema \(r\)-plo di \(|C|\), valendosi perciò di \textit{superficie biaggiunte} ad \(F\), \(\psi_{2n-8}\) e \(\psi_{2n-8+r}\), rispettivamente d’ordine \(2n-8\) e \(2n-8+r\).

Riferiamoci per semplicità di discorso al caso del genere \(p>0\), in cui esistono effettivamente le aggiunte \(\Phi_{n-4}\) e le curve canoniche loro sezioni. Il comportamento delle superficie biaggiunte \(\psi_{2n-8}\) in relazione alle singolarità di \(F\) viene definito di regola dal comportamento di due \(\Phi_{n-4}\) che, prese insieme, compongono appunto una \(\psi_{2n-8}\). Tuttavia occorre definire più precisamente questo comportamento, come facciamo con le osservazioni che seguono.

Se la \(F\) possiede soltanto una curva doppia e non punti multipli propri, le \(\Phi_{n-4}\) passano semplicemente per codesta curva e le \(\psi_{2n-8}\) passano per essa doppiamente: il sistema lineare segato da queste \(\psi_{2n-8}\) è completo perchè si deduce dal sistema segato dalle aggiunte \(\Phi_{2n-8}\) (che già sappiamo essere completo) staccando l’intorno della curva doppia.

Ma se la \(F\) possegga una curva di molteplicità \(i>2\) le superficie biaggiunte \(\psi_{2n-8}\) (e similmente le \(\psi_{2n-8+r}\)) saranno soggette a passare per codesta curva, non già con la molteplicità \(2i-2\); ma, con la molteplicità \(i\), avendo un conveniente contatto con le falde della superficie \(F\) lungo di essa; per esempio se \(i=3\) le \(\psi\) possiederanno la curva tripla di \(F\) come tripla e toccheranno semplicemente le falde di \(F\). Infatti ad una superficie \(\psi\) che passi quattro volte per una curva tripla di \(F\) si può sempre sostituire una combinazione lineare di \(\psi\) e di \(F\) che passi tre volte per questa curva, toccando le falde di \(F\). Delle \(\psi\) che così vengono sottoposte a condizioni meno restrittive si può dire che segano certo su \(F\) un sistema lineare segato dalle aggiunte dello stesso ordine, staccando, un certo numero di volte l’intorno della curva multipla. Invece le \(\psi\) passanti quattro volte per una curva tripla di \(F\) non daranno in generale un sistema completo.

Similmente un punto isolato \(s\)-plo della superficie \(F(s>2)\) che è in generale \((s-2)\)-plo per le superficie aggiunte \(\Phi\), dovrà essere multiplo per le superficie biaggiunte \(\psi\), non già con la molteplicità \(2s-4\), ma con la molteplicità \(s\), aggiungendo la condizione di toccare la \(F\) nel punto stesso \(s-4\) volte.

Così se \(O\) è un punto triplo isolato di \(F\) le \(\psi\) sono soggette a passare doppiamente per esso. Se invece \(O\) è un punto quadruplo, le \(\psi\) sono soggette soltanto ad avere in esso un punto quadruplo; ma se \(O\) è un punto quintuplo di \(F\), esso dovrà essere soltanto quintuplo per le \(\psi\), che però dovranno avere in esso lo stesso cono osculatore.

Un’attenzione particolare deve essere rivolta al comportamento delle superficie biaggiunte in un tacnodo della superficie \(F\). A primo aspetto si potrebbe ritenere che un tal punto, essendo semplice per le superficie aggiunte \(\Phi_{n-4}\), debba essere doppio per le \(\psi_{2n-8}\);
%CONT>>
% scan: s055R
% !PRINTED 91
\weberpage{91}
% head: LE SUPERFICIE AGGIUNTE
%<<CONT
% stamp: Dark scan-edge marks at the left and scattered specks.
come per le superficie composte di due \(\Phi_{n-4}\). Ma si rifletta al modo come viene stabilito il comportamento delle superficie \(\Phi_{n-4}\) in un tacnodo \(O\): una sezione piana per \(O\) possiede due punti doppi infinitamente vicini \(O\) ed \(O_1\), e la sezione di una \(\Phi_{n-4}\) è tale che insieme ad una retta per \(O\) deve costituire un’aggiunta \(\varphi_{n-4}\) alla sezione di \(F\). % [sic] \varphi_{n-4} -> \varphi_{n-3}, the section of degree n-4 together with a line has degree n-3.
Perciò la \(\Phi_{n-4}\) è soggetta a priori a passare per il punto \(O_1\) infinitamente vicino ad \(O\), e soltanto per effetto di \textit{scaricamento} viene a passare per il punto proprio \(O\). Similmente un piano per \(O\) verrà segato da una superficie biaggiunta \(\psi_{2n-8}\) secondo una curva che a priori sarà soggetta a passare doppiamente per \(O_1\), e, per scaricamento, verrà a passare semplicemente per \(O\) ed \(O_1\). Quindi \textit{le condizioni che un tacnodo impone alle superficie biaggiunte consistono soltanto nel passare per questo punto e toccare ivi il piano tacnodale}: così il tacnodo \textit{diminuisce il bigenere di tre}. Queste condizioni si possono anche verificare rilevando che esse importano per le curve bicanoniche sezioni delle \(\psi_{2n-8}\) la proprietà caratteristica di segare sopra la curva generica di un sistema lineare \(|C|\) (per esempio della rete delle sezioni piane per \(O\)) gruppi di punti appartenenti alla serie doppia della residua della serie caratteristica.

Rivolgiamo ora la nostra attenzione al caso della singolarità noetheriana costituita da un punto doppio \(O\) con un tacnodo infinitamente vicino \(O_1\). Anche qui il passaggio delle \(\Phi_{n-4}\) per \(O\) è una conseguenza del principio di scaricamento delle singolarità: a priori le \(\Phi_{n-4}\) dovrebbero passare per il punto infinitamente vicino \(O_1\). Quindi le \(\psi_{2n-8}\) soggette a passare due volte per \(O_1\) verranno effettivamente determinate dalla condizione di passare per i due punti infinitamente vicini \(O\) ed \(O_1\). Dunque: \textit{un punto doppio con tacnodo infinitamente vicino impone in genere alle superficie biaggiunte di passare semplicemente per i due punti infinitamente vicini, \(O\) ed \(O_1\)}, e quindi \textit{diminuisce il bigenere della superficie soltanto di due}. È facile verificare questi risultati riferendoci ad un piano doppio con curva di diramazione \(f(xy)=0\) d’un certo ordine \(2n\), cioè alla superficie \(z^2=f(xy)\). Abbiamo visto che un punto quadruplo della \(f=0\) è un tacnodo per la superficie \(F\), e, nel piano, le curve canoniche di \(F\) sono rappresentate da curve d’ordine \(n-3\) passanti semplicemente per \(O\). Ora le curve d’ordine \(2n-6\) passanti doppiamente per \(O\) rappresenteranno curve bicanoniche di \(F\), per quanto il sistema di codeste curve non sia, in generale, completo\origfootnote{1}{s055R-1}{Un sistema lineare completo di curve \(|C|\) sopra il piano doppio rappresenta, sulla superficie \(F\), non già un sistema completo, bensì un sistema di curve appartenenti all’involuzione \(I\) le cui coppie corrispondono ai punti del piano doppio; e questo si lascia in generale estendere in un più ampio sistema le cui curve non appartengono più all’involuzione \(I\). Ciò si vede già nel caso del piano doppio con curva di diramazione del \(4^\circ\) ordine \(f_4=0\) che si ottiene per proiezione del punto all’infinito \(O\) dell’asse \(z\) della superficie del \(4^\circ\) ordine \(F_4\) di equazione \(z^2=f_4(xy)\), che ha in \(O\) un tacnodo con piano tacnodale all’infinito. Qui le \(\infty^5\) coniche del piano sono immagini delle curve segate su \(F_4\) dai coni quadrici per \(O\), curve iperellittiche che appartengono all’involuzione \(I\). Ma il sistema \(\infty^5\) di tali curve è contenuto nel sistema lineare completo \(\infty^6\) che viene segato sopra \(F_4\) dalle quadriche toccanti in \(O\) il piano tacnodale. % [sic] del punto -> dal punto, the projection is taken from O.

Per quel che concerne la possibilità di estendere il sistema delle curve bicanoniche che hanno per immagine sul piano doppio delle curve d’ordine \(2n-6\), la questione dipende dal comportamento delle superficie biaggiunte nel punto all’infinito dell’asse \(z\), dove la \(F\) presenta una nota singolarità. Basti dire che in generale le \(\psi_{2n-8}\) avranno in quel punto molteplicità \(2n-2\) ed un conveniente contatto con la superficie. Appare di qui che, mentre le curve canoniche di una superficie rappresentabile sul piano doppio appartengono sempre all’involuzione le cui coppie rispondono ai punti di questo piano, lo stesso non accade in generale per le curve bicanoniche (Cfr. Cap. V, \S~7).
% EDITORIAL: continuazione della nota ricongiunta da p. 92.
}. Così si riesce a verificare che un tacnodo diminuisce di \(3\) il bigenere di \(F\).
%CONT>>

% scan: s056L
% !PRINTED 92
\weberpage{92}
% head: CAPITOLO TERZO
%<<CONT
% The incoming continuation belongs only to the footnote reproduced below.
% stamp: Dark scan-edge line along the right margin and scattered scan specks.
Di qui si passa al caso della singolarità noetheriana che risponde a un punto \([3,3]\) della curva di diramazione \(f\). Il punto \([3,3]\) è da ritenersi come un punto doppio \(O\), cui sia infinitamente vicino un punto 4-plo \(O_1\). Le immagini delle curve bicanoniche vengono dunque assoggettate a passare doppiamente per \(O_1\), ossia --- per effetto di scaricamento --- le immagini delle curve bicanoniche dovranno passare semplicemente per \(O\) e \(O_1\), e quindi \textit{un punto \([3,3]\) della curva di diramazione del piano doppio diminuisce di \(2\) il bigenere della superficie.}

\edsection{12. Superficie di genere \(p=0\) e di bigenere \(P>0\).}{12. Superficie di genere p=0 e di bigenere P>0.}{sec.III.12}

Le cose dette permettono facilmente di costruire esempi di superficie di genere \(p=0\) e di bigenere \(P>0\), le quali perciò sono certo non razionali.

Il primo esempio si è presentato ad \textsc{Enriques} nel 1896\origfootnote{1}{s056L-1}{\textit{Introduzione ecc.}, l. c.}. È la superficie del \(6^\circ\) ordine \(F_6\) che passa doppiamente per gli spigoli di un tetraedro. Prendendo questo tetraedro come fondamentale per un sistema di coordinate proiettive, la \(F_6\) viene rappresentata da un’equazione della forma:
\[
F_6=\varphi(x_2x_3x_4,\ x_3x_4x_1,\ x_4x_1x_2,\ x_1x_2x_3)+x_1x_2x_3x_4f(x_1x_2x_3x_4)=0
\]
dove \(f\) e \(\varphi\) sono forme quadratiche rispettivamente nelle \(x_1x_2x_3x_4\) e nei prodotti di terne di queste variabili. Del resto la \(F_6\) appartiene al sistema lineare doppio di quello costituito dalle superficie cubiche che passano semplicemente per gli spigoli del tetraedro.

% Unnumbered continuation of a footnote; its reference mark is outside the supplied scans.
% EDITORIAL: testo della nota trasferito nella nota iniziata a p. 91.

% scan: s056R
% !PRINTED 93
\weberpage{93}
% head: LE SUPERFICIE AGGIUNTE
% stamp: Short stray strokes above “dalle” in the first paragraph and scattered scan specks.
La \(F_6\) è di genere \(p=0\), perchè non esistono quadriche aggiunte ad essa, cioè passanti per gli spigoli del tetraedro. Ma il bigenere della \(F_6\) vale \(P=1\), perchè esiste una superficie biaggiunta ad essa dell’ordine \(2\cdot6-8=4\), cioè il nostro tetraedro. Questa \(\psi_4\) sega sopra \(F_6\) \textit{una curva bicanonica d’ordine zero.} Il sistema lineare \(|C|\) delle sezioni piane di \(F_6\) è di genere \(\pi=4\) e di grado \(n=2\pi-2=6\); la sua serie caratteristica è una \(g_6\) completa e non speciale, il cui doppio dà il doppio della serie canonica. Il sistema \(|C'|\) aggiunto a \(|C|\) viene segato su \(F_6\) dalle superficie cubiche \(\Phi_3\) passanti semplicemente per gli spigoli del tetraedro; esso è, come \(|C|\), un sistema di sestiche di genere 4 e di grado 6. E c’è questo di notevole, che i doppi dei due sistemi \(|C|\) e \(|C'|\) coincidono:
\[
|2C|=|2C'|;
\]
questa coincidenza esprime il significato concreto dell’esistenza del sistema bicanonico d’ordine zero \(|2C'-2C|\).

Siccome poi
\[
|(2C)''|=|2C'|=|C+C''|
\]
si deduce che
\[
|C''-C|=|2C'-2C|=|2C-2C|
\]
e quindi
\[
|C''|=|C|;
\]
in parole: il sistema delle sezioni piane di \(F_6\) è a sua volta l’aggiunto del suo sistema aggiunto, segato dalle superficie \(\Phi_3\) che passano per gli spigoli del tetraedro.

Valutiamo il genere lineare \(p^{(1)}\) della \(F_6\) calcolando il genere virtuale delle curve di \(|C'-C|\) secondo il \S~10 del Cap. I. Designando con \(\pi\) e \(n\) il genere e il grado di \(|C|\), con \(x\) e \(y\) il genere e il grado della curva virtuale \(K=C'-C\) d’ordine zero, che ha \(i=0\) intersezioni con le \(C\), scriveremo il genere e il grado di \(|C+K|=|C'|\):
\[
\begin{aligned}
\pi'&=\pi+x+0-1=\pi\\
n'&=n+y+2\cdot0=n,
\end{aligned}
\]
da cui
\[
x=1,\qquad y=0:
\]
il genere lineare della \(F_6\) vale \(p^{(1)}=1\).

Un secondo esempio di superficie di genere \(p=0\) e di bigenere \(P>0\) è stato indicato dal \textsc{Castelnuovo}, ancora nel 1896\origfootnote{1}{s056R-1}{G. \textsc{Castelnuovo}, \textit{Sulle superficie di genere zero.} Memorie della Soc. It. delle Scienze (detta dei XL), serie III, t. 10, 1896.}. È la superficie del settimo ordine \(F_7\) che passa tre volte per una retta \(r\), due volte per una conica \(\Gamma\) (non incidente ad \(r\)) ed inoltre possiede tre tacnodi i cui piani tacnodali passano per \(r\).

% scan: s057L
% !PRINTED 94
\weberpage{94}
% head: CAPITOLO TERZO
% stamp: Dark right scan border and scattered scan specks.
Una tale superficie (irriducibile) esiste di fatto e si può costruire come segue: si assuma anzitutto una superficie cubica \(F_3\) che passi semplicemente per la retta \(r\) e per la conica \(\Gamma\) e si considerino tre piani \(\alpha_1\alpha_2\alpha_3\) passanti per la retta \(r\) e tangenti alla \(F_3\) rispettivamente nei punti \(A_1,A_2,A_3\) (fuori di \(r\)). È facile vedere che esiste una superficie irriducibile del \(4^\circ\) ordine \(F_4\) passante doppiamente per la conica \(\Gamma\) e tangente ai piani \(\alpha_1\alpha_2\alpha_3\) rispettivamente nei punti \(A_1,A_2,A_3\): questa superficie presa insieme coi tre piani \(\alpha_1\alpha_2\alpha_3\) costituisce una particolare \(F_7\) riducibile che passa due volte per la conica \(\Gamma\), tre volte per la retta \(r\), e possiede i tacnodi \(A_1,A_2,A_3\) con i piani tacnodali \(\alpha_1\alpha_2\alpha_3\). Ma una seconda superficie spezzata soddisfacente alle medesime condizioni si ha sommando alla \(F_3\) contata due volte un piano per la retta \(r\); quindi, facendo fascio delle due \(F_7\) spezzate, si otterrà una \(F_7\) irriducibile dotata delle singolarità sopra indicate.

Non vi sono superficie \(\Phi_3\) aggiunte ad \(F_7\): invero una \(\Phi_3\) dovendo passare per la conica \(\Gamma\) e doppiamente per la retta \(r\) che non è incidente ad essa, dovrà a priori spezzarsi, contenendo come parte il piano di \(\Gamma\). Ma la quadrica residua è assoggettata anzitutto a passare doppiamente per \(r\) e quindi deve spezzarsi in due piani per \(r\), e deve inoltre passare per i tacnodi \(A_1,A_2,A_3\); ciò importa tre condizioni impossibili a soddisfare. Si conclude che il genere di \(F_7\) vale
\[
p=0.
\]

Ora determiniamo le superficie \(\psi_{2\cdot7-8}=\psi_6\) biaggiunte alla \(F_7\). Fra queste si troveranno certo le superficie composte del piano di \(\Gamma\) contato due volte e di quattro piani per la retta \(r\), tre dei quali assoggettati a passare per i tacnodi \(A_1,A_2,A_3\) e perciò coincidenti con i piani \(\alpha_1\alpha_2\alpha_3\); così si ottengono \(\infty^1\) superficie biaggiunte \(\psi_6\) che contengono come parte variabile un piano generico per \(r\) e segano sulla \(F_7\) \(\infty^1\) quartiche bicanoniche di genere 1, formanti un fascio senza punti base. È chiaro che questo sistema bicanonico di grado zero non può essere ampliato, e perciò la condizione imposta alle superficie biaggiunte \(\psi_6\) di passare triplamente per la retta \(r\) e toccare le falde della \(F_7\) per essa, porta qui, come conseguenza, che la \(\psi_6\) possegga la \(r\) come retta quadrupla.

Riassumendo diremo che: \textit{la superficie \(F_7\) di Castelnuovo ha il genere \(p=0\) e il bigenere \(P=2\), contenendo un fascio di quartiche piane bicanoniche di genere \(1\) e grado \(0\).} % ?? uncertain: punctuation after “Castelnuovo” -- no punctuation or a faint period.

È facile determinare il genere lineare della \(F_7\), cioè il genere (virtuale) \(p^{(1)}\) delle sue curve canoniche virtuali \(K\); infatti essendo il grado di \(|K|\) uguale a \(p^{(1)}-1\), il genere e il grado di \(|2K|\) saranno rispettivamente
\[
3p^{(1)}-2\qquad \text{e}\qquad 4p^{(1)}-4.
\]

% scan: s057R
% !PRINTED 95
\weberpage{95}
% head: LE SUPERFICIE, AGGIUNTE % ?? uncertain: mark after “SUPERFICIE” -- a printed comma or scan noise.
% EDITORIAL D0058 RISOLTO: The running head has a clear low comma with a descending tail after SUPERFICIE.
% stamp: Dark left scan border and scattered scan specks.
Dall’essere \(3p^{(1)}-2=1\) e \(4p^{(1)}-4=0\) si ricava ugualmente che il \textit{genere lineare} vale \(p^{(1)}=1\).

Per lunghi anni non si sono presentati altri esempi di superficie di genere \(p=0\) e di bigenere \(P>0\), e specialmente rimaneva senza risposta la domanda se tra tali superficie ne esistano di quelle per cui il genere lineare \(p^{(1)}>1\), possedenti un sistema bicanonico irriducibile di dimensione \(>1\). A questa domanda rispondono ora gli esempi effettivamente costruiti, simultaneamente e indipendentemente l’uno dall’altro, da L. \textsc{Godeaux}\origfootnote{1}{s057R-1}{Cfr. L. \textsc{Godeaux}, \textit{Sur une surface algébrique de genre zéro et de bigenre deux.} Rend. Acc. dei Lincei, vol. XIV, 1931.} e da L. \textsc{Campedelli}\origfootnote{2}{s057R-2}{L. \textsc{Campedelli}, \textit{Sui piani doppi con curva di diramazione del decimo ordine,} Ibidem, vol. XV, 1932.} nel 1931-32.

Ci limiteremo a riferire l’esempio costruito dal Campedelli che si riferisce ad un piano doppio con curva di diramazione d’ordine 10, dotata di 6 punti \([3,3]\) (coppia di punti tripli infinitamente vicini). Si può costruire effettivamente una siffatta curva \(f_{10}\), almeno come curva composta di tre coniche a due a due bitangenti e di una quartica \(f_4\) che le tocchi nei loro 6 punti di contatto. Invero si otterranno 3 coniche mutuamente bitangenti prendendo p. es. due cerchi concentrici (che si toccano nei punti ciclici del piano), ed un’ellisse che tocchi il cerchio minore negli estremi di un diametro ed il maggiore negli estremi del diametro perpendicolare. I 6 punti di contatto delle coniche così ottenute offrono 12 condizioni lineari ad una quartica. Prendendo insieme una di queste quartiche con le tre coniche a costituire la curva di diramazione \(f_{10}\) di un piano doppio, avremo che:

\(1^\circ)\) non esiste sul piano doppio alcuna curva canonica, che sarebbe una conica passante per i 6 punti singolari della \(f_{10}\), mentre nell’esempio costruito si vede che codesti 6 punti non appartengono ad una conica (cerchio); quindi il genere del piano doppio vale
\[
p=0;
\]

\(2^\circ)\) esiste una rete di quartiche che, come la quartica \(f_4\) componente della \(f_{10}\) di diramazione, passano per i 12 punti tripli di questa: esse rappresentano \(\infty^2\) curve bicanoniche della superficie, appartenenti all’involuzione le cui coppie hanno per immagini i punti del piano, e si può vedere che formano il sistema bicanonico completo, sicchè il bigenere della supeficie vale % [sic] supeficie -> superficie, the letter r is missing.
\[
P=3.
\]

\textit{Nota.} -- La dimostrazione che \(P=3\), e non \(P>3\) risulterà più avanti dalla espressione del \(P\) per mezzo del genere lineare che dà qui \(P=p+p^{(1)}=p^{(1)}=3\), dove il valore del \(p^{(1)}\) risulta
%CONT>>
% scan: s058L
% !PRINTED 96
\weberpage{96}
% head: CAPITOLO TERZO
%<<CONT
% stamp: Dark right scan border and scattered scan specks.
direttamente osservando che la conica doppia virtuale costituente la curva canonica, possederebbe 8 punti di diramazione: 6 nei punti tripli di \(f_{10}\) ed altri due fuori. Se non si vuol ricorrere alla curva canonica virtuale basta anche osservare che il grado del sistema bicanonico è il doppio del grado del sistema delle quartiche, sicchè si ha
\[
4(p^{(1)}-1)=2\cdot4=8,\qquad p^{(1)}=3.
\]
Se non si vuol far uso della relazione fra il bigenere ed il genere lineare si può tuttavia dimostrare direttamente che, per la nostra superficie, la rete delle quartiche doppie immagini delle curve bicanoniche è completa. Perciò basterà osservare che la \(f_4\), componente della curva di diramazione, fa parte di una sola curva bicanonica, costituita dalla curva stessa contata due volte, e che le \(\infty^2\) quartiche per i punti tripli di \(f_{10}\) segano su codesta \(f_4\) una \(g_4\) completa e non speciale: invero se questa serie fosse la serie canonica di \(f_4\) i 12 punti tripli della \(f_{10}\) verrebbero ad appartenere ad una cubica, ciò che facilmente si vede costituire un assurdo.

\edsection{13. Criteri d’equivalenza.}{13. Criteri d’equivalenza.}{sec.III.13}

La costruzione delle superficie sub-aggiunte ed aggiunte ad una superficie \(F\) non rigata dello spazio ordinario contiene un \textit{criterio di equivalenza}: due curve \(L\) ed \(M\) che seghino sopra le sezioni piane \(C\) di \(F\) gruppi della serie somma (o differenza) della serie canonica di \(C\) e di un multiplo della serie caratteristica, sono equivalenti, a meno di curve fondamentali del sistema \(|C|\); sono senz’altro equivalenti se \(|C|\) non possegga curve fondamentali proprie. E si può osservare che l’equivalenza delle due curve \(L\) ed \(M\) può essere inferita dalla sola condizione di segare gruppi equivalenti sopra le \(C\), anche se la serie cui essi appartengono non sia la somma della serie canonica e d’un multiplo della serie caratteristica, quando il sistema \(|C|\) non abbia, come si è detto, curve fondamentali proprie.

Infatti, conducendo per \(L\) una superficie aggiunta ad \(F\), d’ordine abbastanza elevato, si avrà come intersezione ulteriore una curva \(D\). Ora \(L+D\) ed \(M+D\) segano su una \(C\) gruppi equivalenti della serie somma della serie canonica e di un multiplo della serie caratteristica; perciò sono equivalenti, e quindi risultano anche equivalenti \(L\) ed \(M\).

I criteri di equivalenza si possono sviluppare indipendentemente dalla teoria delle superficie aggiunte e subaggiunte, come ha fatto, in una maniera interessante, F. \textsc{Severi}\origfootnote{1}{s058L-1}{Cfr. F. \textsc{Severi}, \textit{Il teorema d’Abel sulle superficie algebriche.} Ann. di Mat., (3), 12, 1905. \textit{Osservazioni varie di geometria sopra una superficie algebrica.} Atti Ist. Veneto, 65, parte II, 1906. \textit{Alcune relazioni di equivalenza ecc.} Ibidem, 70, 1911.}. Esponiamo brevemente
%CONT>>
% scan: s058R
% !PRINTED 97
\weberpage{97}
% head: LE SUPERFICIE AGGIUNTE
%<<CONT
% stamp: Dark left scan border and scattered scan specks.
questi sviluppi. Anzitutto si abbia sopra una superficie \(F\) un \textit{fascio} lineare di curve \(C\) dotato di almeno un \textit{punto base semplice} \(O\) e \textit{privo di curve riducibili.} Allora \textit{se due curve \(K_0\) e \(K_\infty\) non passanti per \(O\) segano sulle \(C\) gruppi di punti equivalenti, esse sono equivalenti.}

Infatti sopra una qualunque \(C\) del fascio possiamo assumere i gruppi segati da \(K_0\) e \(K_\infty\) come gruppo degli zeri e dei poli di una funzione razionale \(\varphi\), che rimane determinata a meno di una costante moltiplicativa, e possiamo anche determinare questa costante imponendo alla \(\varphi\) di assumere il valore 1 nel punto \(O\).

Al variare della \(C\) nel suo fascio, la \(\varphi\) ci dà una funzione razionale dei punti della superficie che ha come curva degli zeri la \(K_0\) e come curva dei poli la \(K_\infty\), sicchè queste curve risultano equivalenti, c. d. d. Perchè la deduzione sia rigorosa occorre tuttavia dimostrare che la \(K_0\) esaurisce la curva degli zeri della nostra funzione, e così la \(K_\infty\) ne esaurisce la curva dei poli. A tal uopo si osservi che nell’ipotesi opposta, si dovrebbe aggiungere alla \(K_0\) una curva che. non avendo intersezioni variabili con le \(C\) irriducibili del fascio, dovrebbe comporsi con una o più \(C\); ma allora, sopra una tale \(C\), il punto \(O\) verrebbe a costituire un punto d’indeterminazione della funzione, dove essa avrebbe contemporaneamente i valori zero ed 1, e ciò è impossibile, non appartenendo \(O\) nè al gruppo degli zeri nè al gruppo dei poli che determinano sopra la \(C\) la funzione razionale subordinata. % ?? uncertain: punctuation after “una curva che” -- comma or period.
% EDITORIAL D0059 RISOLTO: The high-resolution crop shows a compact full stop at the baseline after che, with no comma descender. The printing is grammatically defective but should be transcribed faithfully.

Il precedente criterio d’equivalenza si può generalizzare. \textit{Sia \(|C|\) un fascio di curve, razionale o irrazionale, privo di curve riducibili: due curve \(K_0\) e \(K_\infty\) che seghino sopra le \(C\) gruppi equivalenti, sono tra loro equivalenti a meno di curve del fascio}, cioè si ha
\[
K_0+\Sigma C\equiv K_\infty+\Sigma'C.
\]

Per dimostrare questo \textit{secondo criterio d’equivalenza}, procederemo come innanzi a costruire una funzione razionale \(\varphi\) avente la \(K_0\) e la \(K_\infty\) rispettivamente come curva di zeri e come curva di poli. Sappiamo che la \(\varphi\) è definita a meno di una costante moltiplicativa sopra ogni \(C\) del fascio, e (designando con \(t\) il parametro delle \(C\) di questo fascio) otteniamo in tal guisa una funzione dei punti della superficie, \(\lambda(t)\varphi\), dove occorre determinare \(\lambda\). A questo scopo non occorre determinare un punto sopra ogni \(C\), come si faceva nella dimostrazione del primo criterio, ma basta determinare su di essa un gruppo di un certo numero \(m\) di punti, formato dalle intersezioni che la \(C\) stessa ha con una curva \(L\) ausiliaria tracciata sulla superficie; si darà allora la somma dei valori
\[
\lambda\varphi_1+\lambda\varphi_2+\ldotp\,\ldotp\,\ldotp\,\ldotp+\lambda\varphi_m
\]
% Printer's signature in the lower margin, separate from the continued paragraph.
% EDITORIAL: segnatura originale: \parbox{\linewidth}{\textsc{Enriques F.} - \textit{Superficie algebriche.}\hfill 7}
\printersignature{\textsc{Enriques F.} - \textit{Superficie algebriche.}\quad 7}
%CONT>>
% scan: s059L
% !PRINTED 98
\weberpage{98}
% head: CAPITOLO TERZO
%<<CONT
% stamp: Dark horizontal streaks at the top margin and thin vertical marks in the lower left margin.
che la \(\varphi\) assume nei punti di codesto gruppo. Più semplicemente basterà assumere in luogo della \(\varphi\) la funzione
\[
f=\frac{\varphi}{\varphi_1+\varphi_2+\ldotp\,\ldotp\,\ldotp\,\ldotp+\varphi_m},
\]
che è una funzione razionale dei punti della superficie, indipendente da \(\lambda\).

È ovvio che la \(K_0\) fa parte della curva degli zeri di \(f\) e la \(K_\infty\) fa parte della curva dei poli, ma queste curve non sono così esaurite: alla \(K_0\) occorre sommare un gruppo \(\Sigma C\) formato delle \(C\) che passano per i punti d’intersezione di \(L\) con \(K_\infty\); infatti, sopra queste \(C\) il denominatore di \(f\) diventa \(\infty\) e quindi \(f\) è nullo dapertutto. Similmente, alla \(K_\infty\) conviene sommare un gruppo \(\Sigma'C\) formato dalle \(C\) che passano per le intersezioni di \(L\) con \(K_0\), come si vede scambiando l’ufficio delle curve \(K_0\) e \(K_\infty\) col sostituire ad \(f\), \(\frac{1}{f}\).

Avremo dunque
\[
K_0+\Sigma C\equiv K_\infty+\Sigma'C.
\]
Che effettivamente la \(K_0+\Sigma C\) esaurisca la curva degli zeri di \(f\) (e analogamente per la \(K_\infty+\Sigma'C\)), si prova osservando che codesta curva degli zeri non può contenere fuori di \(K_0\) che un certo numero di curve senza intersezioni variabili con le \(C\), e perciò un gruppo di \(C\) (irriducibili); ma sopra una \(C\) non può aversi identicamente \(f=0\) se non quando diventi infinita la somma \(\varphi_1+\varphi_2+\ldotp\,\ldotp\,\ldotp\,\ldotp+\varphi_m\) e quindi una delle \(\varphi_i\), ciò che significa che la \(C\) stessa deve passare per un punto comune ad \(L\) e \(K_\infty\).
% The plus sign following \varphi_2 is repeated at the printed line break.

Aggiungiamo due osservazioni. La prima è che quando il fascio delle curve \(C\) sia razionale, le \(C\) di essa sono equivalenti, e perciò esiste una differenza \(\Sigma C-\Sigma'C\) costituita da un gruppo di un certo numero \(s\) di curve \(C\) da prendere positivamente o negativamente, sicchè si avrà % [sic] di essa -> di esso, the stated antecedent is the masculine “fascio”.
\[
K_0\equiv K_\infty-sC
\]
ovvero
\[
K_\infty\equiv K_0+sC.
\]

E se in particolare il fascio delle \(C\) è dotato di punti base si ha senz’altro l’equivalenza
\[
K_0\equiv K_\infty
\]
quando l’equivalenza tra i due gruppi \((CK_0)\) e \((CK_\infty)\) sussista tenendo conto anche dei punti che cadono nei punti base.

Infatti, dall’essere, per esempio,
\[
K_0\equiv K_\infty-sC
\]
%CONT>>
% scan: s059R
% !PRINTED 99
\weberpage{99}
% head: LE SUPERFICIE AGGIUNTE
%<<CONT
% stamp: Dark left scan border, a dark patch in the upper right margin, and scattered marks below the footnote.
segue
\[
(CK_0)\equiv(CK_\infty)-s(CC);
\]
ma è pure
\[
(CK_0)\equiv(CK_\infty),
\]
e quindi, se il gruppo \((CC)\) non è nullo, deve essere \(s=0\).

Passiamo al \textit{terzo criterio d’equivalenza: se sopra una superficie \(F\) è dato un sistema \(\infty^1\) d’indice}\origfootnote{1}{s059R-1}{S’intende che per un punto generico della superficie passino \(\nu\) curve del sistema.} \textit{\(\nu>1\) di curve tutte irriducibili, due curve \(K_0\) e \(K_\infty\) che seghino le \(C\) in gruppi equivalenti, sono equivalenti.}

Per dimostrarlo si costruirà come innanzi sopra ogni curva \(C\) una funzione razionale \(\varphi\) che abbia i suoi zeri nelle intersezioni con la \(K_0\) e i suoi poli nelle intersezioni con la \(K_\infty\), la quale resta determinata a meno di una costante moltiplicativa \(\lambda\). Quindi si assumerà sopra \(F\) una curva ausiliaria \(L\) avente un certo numero \(m\) d’intersezioni con le \(C\), e --- per ogni \(C\) --- si considererà la somma dei valori che la \(\varphi\) prende in codesti \(m\) punti:
\[
\varphi_1+\varphi_2+\ldotp\,\ldotp\,\ldotp\,\ldotp+\varphi_m.
\]
In tal guisa sopra ogni \(C\) riesce determinata la funzione razionale
\[
f=\frac{\varphi}{\varphi_1+\varphi_2+\ldotp\,\ldotp\,\ldotp\,\ldotp+\varphi_m}.
\]
Dopo ciò determineremo una funzione razionale \(\psi\) dei punti della superficie \(F\) attribuendo ad essa in un punto \(P\) la somma dei valori \(f_1,f_2,\ \ldotp\,\ldotp\,\ldotp\,\ldotp\ f_\nu\) che la \(f\) riceve sopra le \(\nu\) curve \(C\) passanti per \(P\):
\[
\psi=f_1+f_2+\ldotp\,\ldotp\,\ldotp\,\ldotp+f_\nu.
\]
Possiamo indicare gli anzidetti valori delle \(f_i\) valendoci di un indice superiore; ed allora scriveremo
\[
f_i=\frac{\varphi^{(i)}}{\varphi_1^{(i)}+\ldotp\,\ldotp\,\ldotp\,\ldotp+\varphi_m^{(i)}}.
\]

La funzione razionale \(\psi\) si annulla evidentemente sulla curva \(K_0\), dove tutte le \(\varphi^{(i)}\) sono eguali a zero e suoi poli debbono cadere necessariamente nella \(K_\infty\); ma non è a priori evidente che tutti i punti della \(K_\infty\) siano poli per la funzione razionale \(\psi\), poichè una somma di infiniti non è sempre infinita. Quello che possiamo affermare è che la curva \(\overline{K}_0\) degli zeri di \(\psi\) è del tipo \(\overline{K}_0=K_0+A\), mentre, detta \(\overline{K}_\infty\) la curva dei poli, si ha \(K_\infty=\overline{K}_\infty+B\). Segando le curve \(\overline{K}_0\) e \(\overline{K}_\infty\) con la generica \(C\), si ha % [sic] e suoi poli -> e i suoi poli, the article is absent in print.
\[
(K_0,C)+(A,C)\equiv(K_\infty,C)-(B,C)
\]
%CONT>>
% scan: s060L
% !PRINTED 100
\weberpage{100}
% head: CAPITOLO TERZO
%<<CONT
Ma è per ipotesi \((K_0,C)\equiv(K_\infty,C)\), e quindi le curve \(A\) e \(B\) non avrebbero intersezioni con la generica curva \(C\), mentre le nostre ipotesi escludono l’esistenza di curve siffatte. Concludendo, si ha
\[
\overline{K}_0=K_0,\quad \overline{K}_\infty=K_\infty
\]
e risulta così dimostrata l’equivalenza delle curve \(K_0\) e \(K_\infty\).

\textit{Osservazione.} -- In ciò che precede si è trattato dei criterii che permettono di riconoscere l’equivalenza di due curve (isolate) sopra la superficie \(F\). Altri criterii più significativi si hanno in relazione ad una serie continua (almeno \(\infty^1\)) di curve \(|C|\), che venga data su \(F\). Il primo criterio di questo genere è stato stabilito da \textsc{Enriques} nel 1895. Una \textit{serie razionale} di curve \(|C|\) è sempre costituita di \textit{curve equivalenti}, cioè contenuta in un sistema lineare. La proprietà si dimostra come per le serie razionali di gruppi di punti sopra una curva\origfootnote{1}{s060L-1}{Cfr. \textsc{Enriques-Chisini}, Op. cit., Libro V, cap. I, nota al \S~10, vol. III, pag. 73 (e pag. 485).}, quando non si voglia ricondurre a tal caso, mercè i criterii dati innanzi. Un secondo criterio notevole viene messo in luce da \textsc{Severi} (1905)\origfootnote{2}{s060L-2}{\textit{Osservazioni ecc.}, l. c.}. Si può affermare, in generale, che \textit{una serie \(\infty^1\)} (almeno) \textit{di curve \(C\) è costituita di curve equivalenti, se le \(C\) segano gruppi equivalenti sopra una sola curva generica di un sistema continuo \(\{D\}\) d’indice \(\nu>1\).}

Infatti la condizione perchè una serie \(\infty^1\) di gruppi di punti sopra una curva sia costituita di gruppi equivalenti (cioè sia contenuta in una serie lineare) si può esprimere, con Castelnuovo, mediante un criterio numerativo, ove figura il numero dei punti doppi della serie stessa (che deve avere il valore massimo)\origfootnote{3}{s060L-3}{Cfr. \textsc{Enriques-Chisini}, op. cit., Libro V, cap. IV, 42 (vol. III, pag. 482).}. Ora, se questo criterio è soddisfatto dai gruppi sezioni delle \(C\) con una \(D\), per ragioni di continuità sarà anche verificato per rapporto alle altre \(D\), e così si viene ricondotti al terzo criterio, dato innanzi.

Aggiungiamo che il criterio d’equivalenza delle serie \(\infty^1\) di gruppi di punti sopra una curva, che sopra abbiamo ricordato essere stato stabilito dal \textsc{Castelnuovo}, si estende pure alle superficie, come ha mostrato R. \textsc{Torelli}. Questi scrive una diseguaglianza cui soddisfa il numero \(d\) delle curve dotate di punto doppio, appartenenti ad una serie \(\infty^1\); (dove compare l’invariante di Zeuthen-Segre della superficie: cfr. cap. V); il massimo di tale numero \(d\) risponde al caso in cui la detta serie sia costituita di curve equivalenti\origfootnote{4}{s060L-4}{Cfr. R. \textsc{Torelli}, \textit{Sui sistemi algebrici di curve appartenenti ad una superficie algebrica.} Atti R. Accademia delle Scienze; Torino, 1907.}.
% stamp: Hebrew stamp of the Ben-Gurion University of the Negev Library below the footnotes; thin marks in the left margin.

\clearpage\thispagestyle{plain}
% scan: s060R
% !PRINTED none
\weberpage{[s060R]}
% head: none
% stamp: Dark left scan border and a horizontal stray line in the lower right margin.
\edchapter{IV}{IL GENERE NUMERICO E IL TEOREMA DI RIEMANN-ROCH\\PER LE SUPERFICIE}{Capitolo IV. IL GENERE NUMERICO E IL TEOREMA DI RIEMANN-ROCH PER LE SUPERFICIE}{IL GENERE NUMERICO E IL TEOREMA DI RIEMANN-ROCH PER LE SUPERFICIE}

\edsection{1. Introduzione.}{1. Introduzione.}{sec.IV.1}

Per intendere il senso degli sviluppi che seguono conviene richiamare i problemi che s’incontrano nella teoria delle curve, e che trovano qui la loro naturale estensione\origfootnote{1}{s060R-1}{Cfr. p. es. \textsc{Enriques-Chisini}, Op. cit. Libro V, vol. III.}.

Dopo avere definito le serie lineari sopra una curva e le operazioni elementari che le concernono, e dopo avere riconosciuto l’invarianza della serie canonica, si affaccia la domanda di valutare quale sia la dimensione \(r\) della serie completa determinata sopra una curva di genere \(p\) da un gruppo di \(n\) punti. La prima risposta che qui si ottiene, cioè \(r\geq n-p\), viene data dal calcolo della dimensione del sistema delle curve, d’ordine abbastanza elevato, aggiunte alla curva piana \(C\) d’ordine \(m\) e di genere \(p\). Attesochè il sistema di queste aggiunte \(\varphi_{m-3+s}\) sega sempre su \(C\) (qualunque sia \(s\), positivo o negativo) una serie completa, e per \(s\) abbastanza alto i punti doppi della \(C\) offrono alle curve d’ordine \(m-3+s\) che debbono contenerle % [sic] contenerle -> contenerli, the antecedent is “i punti doppi”.
\[
d=\frac{(m-1)(m-2)}{2}-p; % ?? uncertain: punctuation after -p -- semicolon, comma, or scan noise.
% EDITORIAL D0060 RISOLTO: The mark comprises a small upper point and a separate lower curved stroke, compatible with the printed semicolon retained in the transcription.
\]
condizioni lineari indipendenti, la serie di grado \(n=2p-2+sm\) segata su \(C\) da codeste aggiunte risulta di dimensione
\[
n-p.
\]

È chiaro quindi che una qualsiasi serie segata sopra \(C\) dalle \(\varphi_{m-3+s}\) passanti per un gruppo \(G\) di punti che imponga ad esse condizioni tutte indipendenti avrà la dimensione uguale al grado diminuito del genere \(p\). Questi risultati vengono poi precisati nel senso che la serie canonica, di grado \(2p-2\), segata su \(C\) dalle
%CONT>>
% scan: s061L
% !PRINTED 102
\weberpage{102}
% head: CAPITOLO QUARTO
%<<CONT
% stamp: Stray marks in the upper margin and a dark scanner band along the right edge.
aggiunte \(\varphi_{m-3}\), ha sempre la dimensione \(p-1\), quale risulta dal numero delle condizioni offerte dai punti doppi di \(C\) alle curve d’ordine \(m-3\) che debbono contenerle, queste condizioni riuscendo fra loro indipendenti. Infine il teorema di Riemann-Roch propriamente detto risponde nel modo più generale alla questione di valutare la dimensione della serie completa determinata su \(C\) da un gruppo di \(n\) punti: designando con \(i\) (\(\geq 0\)) l’indice di specialità del gruppo, cioè il numero dei gruppi canonici, linearmente indipendenti, che lo contengono, la dimensione della detta serie risulta esattamente % [sic] contenerle -> contenerli, the antecedent is the masculine plural “punti doppi”.
\[
r=n-p+i.
\]
Ora, volendo estendere quest’ordine di problemi alle superficie, saremo indotti anzitutto a cercare per la superficie \(F\) d’ordine \(n\) dello spazio ordinario, quale sia la dimensione del sistema lineare segato dalle superficie aggiunte \(\Phi_{n-3+s}\), per \(s\) abbastanza alto, e porre in relazione il risultato ottenuto col calcolo del genere superficiale di \(F\), di cui si avrà in tal guisa un’espressione virtuale. Seguendo questa via saremo indotti a valutare la dimensione dei sistemi lineari che si presentano come aggiunti di altri sistemi, stabilendo così un « teorema di Riemann-Roch per i sistemi aggiunti », e poi ad esaminare il caso di un sistema lineare qualunque, che non si presenti più come aggiunto di un altro.

Per valutare la dimensione del sistema delle \(\Phi_{n-3+s}\) aggiunte ad una superficie \(F\), occorre richiamare le cosidette formule di postulazione di Cayley (1870) e di Noether (1871), che esprimono il numero delle condizioni lineari indipendenti imposte da una curva alle superficie d’ordine abbastanza elevato che debbano passare per essa, una o più volte.

\edsection{2. Condizioni imposte ad una superficie dal passaggio per una curva:\\ formule di postulazione.}{2. Condizioni imposte ad una superficie dal passaggio per una curva: formule di postulazione.}{sec.IV.2}

Si consideri nello spazio ordinario una curva gobba (irriducibile) \(C\) di un certo ordine \(m\) e di un certo genere \(p\), che sia priva di punti multipli; le superficie \(\Phi\) di un ordine \(n\) qualunque segano sopra \(C\) una serie lineare \(g_{mn}^{r}\), di ordine \(mn\), la cui dimensione designiamo con \(r\). Se la serie è completa e non speciale si può scrivere
\[
r=mn-p.
\]

In tale ipotesi si trova subito la postulazione della curva \(C\) rispetto alle superficie d’ordine \(n\), cioè il numero delle condizioni lineari indipendenti che la \(C\) offre alle \(\Phi\) che debbono contenerle; infatti il numero di queste condizioni sarà % [sic] contenerle -> contenerla, the antecedent is the singular curve C.
\[
r+1=mn-p+1.
\]

% scan: s061R
% !PRINTED 103
\weberpage{103}
% head: IL GENERE NUMERICO E IL TEOREMA DI RIEMANN-ROCH, ECC.
% stamp: Dark scanner band along the left edge and stray marks in the margins.
La nostra domanda riceverebbe così la più semplice risposta, senonchè due cause possono modificare in senso opposto il risultato accennato: la serie \(g_{mn}^{r}\) può essere speciale, e quindi di dimensione \(r>mn-p\), e può essere deficiente, cioè avere una dimensione minore della serie completa. Per effetto della prima circostanza la postulazione di \(C\) per le \(\Phi\) può riuscire \(>mn-p+1\), e per effetto della seconda può riuscire minore. Ma delle due cause perturbatrici una si elimina senz’altro riferendoci a superficie \(\Phi\) d’ordine \(n\) abbastanza alto; giacchè basta prendere \(n\) in modo che sia
\[
mn>2p-2.
\]
Quindi si può affermare che per le superficie \(\Phi\) di ordine \(n\) abbastanza alto, la curva gobba \(C\), d’ordine \(n\) e genere \(p\), ha una postulazione \(\leq nm-p+1\); ciò significa che per la \(C\) passano \textit{almeno} % [sic] d’ordine n -> d’ordine m, the order of the curve C was denoted by m.
\[
\frac{(n+3)(n+2)(n+1)}{6}-nm+p-1
\]
superficie \(\Phi\) linearmente indipendenti.

Si può precisare il risultato ottenuto, dimostrando che per \(n\) abbastanza alto, \(n\geq m-2\), la detta serie \(g_{mn}^{r}\) risulta completa\origfootnote{1}{s061R-1}{La dimostrazione che segue è di \textsc{G. Castelnuovo}, \textit{Sui multipli di una serie lineare di gruppi di punti appartenente ad una curva algebrica.} Rend. Circolo Mat. di Palermo, t. VII, 1893. Riprodotto in “ Memorie Scelte „ pag. 95.}. % ?? uncertain: punctuation after the closing quotation mark in the footnote -- „ pag. or „, pag.
% EDITORIAL D0061 RISOLTO: The crop shows exactly the two low strokes of the closing quotation mark; no separate third comma occurs before pag.
A tale scopo si proietti la curva \(C\) sopra un piano da un punto generico \(O\) dello spazio: si avrà così una curva d’ordine \(m\) con \(d=\frac{(m-1)(m-2)}{2}-p\) punti doppi, sopra la quale le curve aggiunte d’ordine \(\geq m-2\) segheranno una serie completa e non speciale. Segue di qui che i coni d’ordine \(m-2\) (o maggiore) aventi il vertice in \(O\) e contenenti le \(d\) corde di \(C\) che passano per \(O\), segheranno sopra \(C\), fuori dei punti d’appoggio di tali corde, una serie completa. A fortiori le superficie d’ordine \(n\geq m-2\) passanti per gli stessi punti segneranno su \(C\) fuori di questi una serie completa. Questa serie si amplia se si estende il sistema delle superficie anzidette riferendosi alle totalità delle superficie \(\Phi_n\) d’ordine \(n\) dello spazio. % [sic] alle totalità -> alla totalità, the passage apparently refers to the single totality of these surfaces.
Ora, per dimostrare che queste segano su \(C\) una serie completa basterà riconoscere che il gruppo \(G\) dei \(2d\) punti d’appoggio della \(C\) con le corde di essa per \(O\), offre alle \(\Phi_h\) condizioni lineari indipendenti, cosicchè l’addizione dei punti di \(G\) conduce da una serie (non speciale) completa di dimensione % ?? uncertain: the subscript of Phi after “offre alle” -- h or n.
% EDITORIAL D0062 RISOLTO: The subscript after offre alle has a visible tall left ascender, unlike the n printed after Phi several lines above; h is the better diplomatic reading.
\[
mn-2d-p
\]
ad una serie completa di dimensione
\[
mn-p.
\]

% scan: s063L
% !PRINTED 104
\weberpage{104}
% head: CAPITOLO QUARTO
% stamp: Dark scanner band along the right edge, scattered specks, and smudges crossing a few words in the upper paragraph.
A tal uopo, riferendoci per semplicità di discorso al caso \(n=m-2\), basterà costruire una superficie d’ordine \(m-2\) che passi per \(2d-1\) punti di \(G\) e non per il rimanente; una tale superficie si otterrà sommando ad un cono d’ordine \(m-3\) che passi per \(d-1\) corde di \(C\) uscenti da \(O\), e non per la rimanente, un piano che passi per uno dei due punti d’appoggio della rimanente corda, e non per l’altro. L’esistenza del cono d’ordine \(m-3\) anzidetto, risulta dal noto teorema che i punti doppi di una curva piana d’ordine \(m\) (proiezione della \(C\)) offrono condizioni indipendenti alle curve d’ordine \(m-3\) che debbano essere aggiunte ad essa.
% The equality sign is printed on both sides of the line break in n = = m - 2; both copies are retained.

Concludiamo pertanto: \textit{una curva gobba irriducibile d’ordine \(m\) e genere \(p\), senza punti multipli, offre}
\[
nm-p+1
\]
\textit{condizioni lineari alle superficie d’ordine \(n\geq m-2\) che debbono contenerla.}

Il calcolo della postulazione di una curva, che sopra abbiamo indicato, si estende, con lievi modificazioni, al caso di curve gobbe dotate di punti multipli, ed anche riducibili. Consideriamo dapprima il caso più importante per noi in cui si tratti di una \textit{curva \(C\)} (d’ordine \(m\) e genere \(p\)) \textit{che sia dotata di un certo numero \(t\) di punti tripli,} con tangenti non appartenenti allo stesso piano, e che vogliamo supporre ancora irriducibile. Qui si dimostra anzitutto che i \(t\) punti nominati impongono \(4t\) condizioni lineari indipendenti alle superficie \(\Phi\) d’ordine \(n\) abbastanza elevato che debbano possederli come punti doppi, e poi che le \(\Phi\), passanti doppiamente per i \(t\) punti tripli di \(C\), segano su \(C\) (fuori di questi) una serie completa (non speciale), diguisachè \textit{la postulazione della curva \(C\) per le \(\Phi\) d’ordine \(n\) abbastanza alto risulta}
\[
4t+nm-6t-p+1=nm-2t-p+1.
\]

La prima asserzione è evidente, poichè l’assegnazione d’un punto doppio di posizione assegnata porta \(4\) condizioni lineari a cui la superficie deve soddisfare; la seconda asserzione si giustifica collo stesso ragionamento svolto per le curve \(C\) senza punti multipli.

In modo simile si ottiene la postulazione di \(C\) nel caso in cui la \textit{\(C\) possegga un certo numero \(\tau\) di punti doppi} (nodali): qui conviene pure far passare anzitutto le \(\Phi\) anche per questi \(\tau\) punti, ciò che importa \(\tau\) condizioni lineari e diminuisce di \(2\tau\) il numero delle intersezioni colla \(C\); così \textit{la postulazione} della \(C\) per le \(\Phi\) d’ordine \(n\) risulta
\[
\tau+nm-2t-2\tau-p+1=nm-2t-\tau-p+1.
\]

% scan: s063R
% !PRINTED 105
\weberpage{105}
% head: IL GENERE NUMERICO E IL TEOREMA DI RIEMANN-ROCH, ECC.
% stamp: Broken vertical scanner marks along the right edge and scattered specks.
Infine vogliamo estendere questa formula al caso in cui la curva \(C\) si spezzi in più parti, rispettivamente d’ordine \(m_1, m_2, \ldotp\ldotp\ldotp\ldotp\) e di genere \(p_1, p_2, \ldotp\ldotp\ldotp\ldotp\)

Per semplicità di discorso riferiamoci al caso in cui si abbiano due sole componenti \(C_1\) e \(C_2\), le quali potranno essere connesse attraverso semplici incroci, cioè punti doppi della curva composta \(C=C_1+C_2\), o anche attraverso qualcuno dei punti tripli della \(C\) che sia doppio per una delle componenti e semplice per l’altra. Osserviamo anzitutto che le superficie \(\Phi\), d’ordine \(n\) abbastanza elevato, passanti doppiamente per i punti tripli di \(C\) e semplicemente per i suoi punti doppi segheranno serie complete (non speciali) tanto su \(C_1\) che su \(C_2\); invero si lascino cadere per le \(\Phi\) le condizioni relative al passaggio per i punti doppi o tripli di \(C_2\); le \(\overline{\Phi}\) così ottenute segheranno su \(C_1\) una serie completa; ma si può sempre supporre che per un gruppo di questa serie passi un sistema così ampio di superficie \(\overline{\Phi}\) che si possa sempre soddisfare alle condizioni omesse e trovare quindi una \(\Phi\) che passi per quel gruppo. Diciamo di più: non solo le serie segate dalle \(\Phi\) sopra \(C_1\) e \(C_2\) sono ciascuna completa per proprio conto, ma anche la serie che si ha sulla curva spezzata riesce completa, potendosi associare ad un gruppo qualunque della serie data su \(C_1\) un gruppo su \(C_2\); in altre parole quando si considerano le \(\Phi\) passanti per un gruppo della serie su \(C_1\) e per un altro gruppo associato della serie sopra \(C_2\), presi l’uno e l’altro in modo arbitrario, non accade che queste \(\Phi\) debbano contenere di conseguenza l’una o l’altra delle due curve. Per vederlo, collo stesso procedimento usato nel caso della \(C\) irriducibile, possiamo riportarci al piano, invocando la proposizione che « sopra una curva piana riducibile le curve aggiunte d’ordine assai elevato segano sempre una serie completa ».

Abbiamo così stabilito che per le superficie \(\Phi\) d’ordine abbastanza elevato, passanti come si è detto per i punti di connessione delle componenti di \(C\), la postulazione della curva \(C\) è uguale alla somma delle postulazioni delle sue componenti. Resta soltanto da valutare queste postulazioni. A tal uopo indichiamo con \(m_1\) ed \(m_2\) gli ordini e con \(p_1\) e \(p_2\) i generi delle dette componenti. Se si prescinde dai punti fissi assegnati alle \(\Phi\), gli ordini delle due serie segate su \(C_1\) e su \(C_2\) sarebbero rispettivamente \(nm_1\) ed \(nm_2\) e le dimensioni \(nm_1-p_1\) ed \(nm_2-p_2\). Ma un punto doppio della \(C=C_1+C_2\) diminuisce di \(1\) la dimensione del sistema delle \(\Phi\) e diminuisce complessivamente di \(2\) la somma degli ordini delle due serie segate su \(C_1\) e su \(C_2\), tanto se questo punto è un punto doppio per una delle componenti irriducibili di \(C\), quanto se è un punto semplice comune ad entrambe. Similmente un punto triplo della curva composta \(C\), dovendo essere
%CONT>>
% scan: s064L
% !PRINTED 106
\weberpage{106}
% head: CAPITOLO QUARTO
%<<CONT
% stamp: Dark scanner band along the right edge and scattered specks.
doppio per le \(\Phi\), diminuisce di \(4\) la dimensione del sistema delle \(\Phi\), e diminuisce complessivamente di \(6\) la somma degli ordini delle serie da esso segate su \(C_1\) e su \(C_2\), tanto se è un punto doppio per una e semplice per l’altra, quanto se è un punto triplo per una di esse. Insomma si riconosce in tal guisa che i \(t\) punti tripli della curva \(C\) ed i \(\tau\) punti doppi di essa, hanno la stessa influenza che nel caso della curva \(C\) irriducibile. Quindi la postulazione della \(C\) verrà espressa dalla stessa formula scritta innanzi, in cui intervengono i \(t\) punti tripli della \(C\) e i suoi \(\tau\) punti doppi, dove si facciano le modificazioni che qui indichiamo. Al posto della serie d’ordine \(x=mn-6t-2\tau\) segata dalle \(\Phi\) sulla curva irriducibile \(C\), si dovranno considerare le due serie segate sopra \(C_1\) e \(C_2\), di cui designiamo gli ordini con \(x_1\) e \(x_2\); e, al posto della dimensione \(x-p\) di quella serie, si avrà la somma delle dimensioni delle due serie date su \(C_1\) e \(C_2\), cioè
% The minus sign before 6t is printed on both sides of a line break; both copies are retained.
\[
x_1-p_1+x_2-p_2.
\]
Ma, mentre nel caso precedente la postulazione di \(C\) veniva data da \(x-p+1\), ora questa postulazione si otterrà sommando le postulazioni di \(C_1\) e \(C_2\) e sarà quindi
\[
x_1-p_1+1+x_2-p_2+1=x-(p_1+p_2)+2.
\]

In conclusione la postulazione della curva \(C=C_1+C_2\), per le superficie d’ordine \(n\) abbastanza elevato, verrà espressa da
\[
mn-2t-\tau-(p_1+p_2)+2.
\]

Questa formula si riconduce a quella che abbiamo scritto per le \(C\) irriducibili se, al posto dei generi \(p_1\) e \(p_2\) delle componenti \(C_1\) e \(C_2\) di \(C\), s’introduca il genere \(p\) della curva composta, calcolato, secondo Noether, nell’ipotesi che le componenti non abbiano punti di connessione\origfootnote{1}{s064L-1}{Invece questi punti sono stati tolti imponendoli come punti base alle \(\Phi\).}:
\[
p=p_1+p_2-1.
\]

Enunciamo il risultato generale a cui si perviene nel caso di \textit{una curva d’ordine \(m\) composta di \(h\) curve irriducibili di generi \(p_1, p_2, \ldotp\ldotp\ldotp\ldotp p_h\), la quale sia dotata di \(t\) punti tripli (a tangenti generiche) e di \(\tau\) punti doppi; la postulazione della curva per le superficie d’ordine \(n\) abbastanza elevato è}
\[
nm-2t-\tau-p+1
\]
\textit{dove}
\[
p=p_1+p_2+\ldotp\ldotp\ldotp\ldotp+p_h-h+1
\]
\textit{designa il genere della curva composta nell’ipotesi, che le sue parti non siano fra loro connesse.}

% scan: s064R
% !PRINTED 107
\weberpage{107}
% head: IL GENERE NUMERICO E IL TEOREMA DI RIEMANN-ROCH, ECC.
% stamp: Broken scanner marks along the right edge and a small dark mark below the page number.
\edsection{3. Dimensione dei sistemi di superficie aggiunte ad una data.}{3. Dimensione dei sistemi di superficie aggiunte ad una data.}{sec.IV.3}

Riferiamoci ad una superficie \(F\) di un certo ordine \(n\), appartenente allo spazio ordinario, e dotata di singolarità normali: curva doppia nodale con punti tripli a tangenti non complanari e, se si vuole, anche doppi\origfootnote{1}{s064R-1}{A priori si può supporre che la data superficie \(F\) possegga una curva doppia irriducibile priva d’incroci: \(\tau=0\) (Cfr. \textit{Introduzione, ecc.}, \S\ 5). Tuttavia può essere utile sviluppare le formule anche nel caso \(\tau>0\).}. Valutiamo il numero \(N_m\) delle superficie linearmente indipendenti aggiunte ad \(F\), di un ordine \(m=n-4+s\) abbastanza elevato. Designando ancora queste superficie con \(\Phi_m\) e indicando con \(d\) l’ordine della curva doppia di \(F\), con \(\varrho\) il suo genere, con \(t\) il numero dei punti tripli e con \(\tau\) il numero dei punti doppi, avremo che il numero \(N_m\) delle \(\Phi_m\) sarà
% The minus sign before 4 is printed on both sides of a line break; both copies are retained.
\[
N_m=\binom{m+3}{3}-(dm-2t-\tau-\varrho+1);
\]
quindi la differenza \(N_m-N_{w-1}\) sarà data da % ?? uncertain: the subscript of the second N in the prose difference -- w-1 or m-1.
% EDITORIAL D0063 RISOLTO: The second subscript in the prose has the two lower bowls and no upright initial stem of w. The displayed formula immediately below instead clearly has m-1.
\[
N_m-N_{m-1}=\binom{m+2}{2}-d.
\]

Il significato di questa relazione è che \textit{le superficie aggiunte \(\Phi_m\) d’un ordine abbastanza elevato segano sopra un piano il sistema lineare completo delle curve aggiunte \(\varphi_m\) alla sezione di \(F\).}

Infatti si può dedurre la dimensione \(N_{m-1}-1\) del sistema \(\lvert\Phi_{m-1}\rvert\) dalla dimensione \(N_m-1\) del sistema delle \(\Phi_m\), staccando da queste un piano; e lo staccamento del piano importa precisamente \(r+1\) condizioni lineari quando sia \(r\) la dimensione del sistema segato dalle \(\Phi_m\). In questo caso dunque si trova.
\[
\begin{aligned}
r+1&=\binom{m+2}{2}-d,\\
r&=\frac{m(m+3)}{2}-d,
\end{aligned}
\]
che è la dimensione del sistema completo delle \(\varphi_m\) aggiunte alla sezione piana di \(F\).

Nel seguito diremo che il numero \(N_m\) definito dalla formula scritta innanzi esprime in ogni caso il \textit{numero virtuale delle superficie \(\Phi_m\) aggiunte ad \(F\)}: sappiamo che per \(m\) abbastanza alto il numero virtuale è uguale al numero effettivo delle \(\Phi_m\), invece per i valori più bassi di \(m\) ci potrà essere una differenza. Tuttavia è importante osservare che: \textit{per \(m\geq n-4\) il numero effettivo delle superficie aggiunte ad \(F\) è sempre \(\geq N_m\).}

% scan: s065L
% !PRINTED 108
\weberpage{108}
% head: CAPITOLO QUARTO
% stamp: Dark scanner band along the right edge, a small dark mark below the page number, and scattered specks.
Per dimostrarlo, partiamo da un valore dell’ordine \(m\) tale che il numero effettivo delle \(\Phi_m\) sia uguale al numero virtuale \(N_m\), e che sia \(m\geq n-3\).

Dal numero delle \(\Phi_m\) aggiunte ad \(F\) si deduce il numero delle \(\Phi_{m-1}\) sottraendo la dimensione \(r\) del sistema segato dalle \(\Phi_m\) sopra un piano, aumentata di un’unità. Se questo sistema è il sistema completo delle \(\varphi_m\) aggiunte alla sezione piana, il numero effettivo delle- \(\Phi_{m-1}\) risulta ancora uguale al numero virtuale; se, invece c’è una differenza, bisogna dunque che il sistema delle \(\varphi_m\) sezioni delle \(\Phi_m\) sopra un piano non sia completo, e allora % ?? uncertain: the short stroke after “delle” at the printed line end -- a hyphen or a scanning mark.
% EDITORIAL D0064 RISOLTO: A clear short horizontal hyphen is printed immediately after delle at the end of the line.
\[
r<\frac{m(m+3)}{2}-d,
\]
e quindi il numero effettivo delle \(\Phi_{m-1}\) riesce maggiore di \(N_{m-1}\). Questo ragionamento si può applicare finchè sia \(m\geq n-3\), perchè in questa ipotesi il sistema completo delle curve aggiunte alla sezione piana di \(F\), d’ordine \(n\), è regolare, cioè ha la dimensione \(\frac{m(m+3)}{2}-d\) e non maggiore, i punti doppi della curva sezione di \(F\) imponendo condizioni indipendenti alle curve \(\varphi_m\) che debbono passare per essi. Invece se si prende \(m=n-4\) il sistema delle \(\varphi_{n-4}\) aggiunte alla sezione di \(F\) riesce sovrabbondante e così il ragionamento precedente non è più applicabile.

Chiameremo \textit{genere numerico} o \textit{aritmetico} \(p_a\), della superficie \(F\), il numero virtuale delle curve canoniche linearmente indipendenti, ossia il numero \(N_{n-4}\) delle superficie aggiunte \(\Phi_{n-4}\); il genere superficiale della stessa \(F\) che abbiamo definito nel precedente capitolo, verrà distinto, ove occorra, dal genere numerico designandolo come \textit{genere geometrico}: \(p_g\).

Per quanto abbiamo visto sopra si può affermare che: \textit{il genere numerico della superficie \(F\) è in ogni caso minore o uguale del genere geometrico:}
\[
p_a\leq p_g.
\]

Rileveremo fra poco l’invarianza del genere numerico rispetto a trasformazioni birazionali della superficie \(F\). Intanto scriviamo il valore di questo carattere per mezzo delle formule di postulazione sopra stabilite:
\[
p_a=\binom{n-1}{3}-(n-4)d+2t+\tau+\varrho-1.
\]

Per conseguenza il numero delle superficie \(\Phi_{n-4+s}\) aggiunte ad \(F\), che sono linearmente indipendenti, si potrà esprimere come
%CONT>>
% scan: s065R
% !PRINTED 109
\weberpage{109}
% head: IL GENERE NUMERICO E IL TEOREMA DI RIEMANN-ROCH, ECC.
%<<CONT
% stamp: Broken vertical scanner marks along the right edge and stray marks in the margins.
segue:
\[
\begin{aligned}
N_{n-4+s}&=\binom{n-1+s}{3}-(n-4+s)d+2t+\tau+\varrho-1\\
&=p_a+\binom{n-1+s}{3}-\binom{n-1}{3}-sd.
\end{aligned}
\]

Quindi il sistema lineare costituito dalle intersezioni di queste aggiunte \(\Phi_{n-4+s}\) con \(F\) avrà in generale la dimensione:
\[
N_{n-4+s}-1-\binom{s-1}{3}
=p_a+\binom{n-1+s}{3}-\binom{n-1}{3}-\binom{s-1}{3}-sd-1,
\]
essendo \(\binom{s-1}{3}\) (che si annulla per \(s=1,2,3\)) il numero delle superficie linearmente indipendenti d’ordine \(s-4(\geq 0)\), ossia la dimensione del sistema delle superficie \(\Phi_{n-4+s}\) che passano per una medesima intersezione con la superficie \(F\) d’ordine \(n\). La dimensione del sistema delle sezioni delle \(\Phi_{n-4+s}\), che sopra abbiamo scritta, si lascia esprimere in una forma molto semplice, introducendo il genere \(\pi_s\) delle curve segate su \(F\) dalle superficie d’ordine \(s\). Calcoliamo anzitutto questo \(\pi_s\) per mezzo del genere \(\pi\) (\(=\pi_1\)) delle sezioni piane di \(s\); avremo % [sic] sezioni piane di s -> sezioni piane di F, the plane sections under discussion are those of the surface F.
\[
\pi_s=s\pi+\frac{s(s-1)}{2}n-s+1,
\]
e, sostituendo per \(\pi\) il suo valore
\[
\pi=\frac{(n-1)(n-2)}{2}-d,
\]
risulterà
\[
\begin{aligned}
\pi_s&=s\left[\frac{(n-1)(n-2)}{2}-d\right]+\frac{s(s-1)}{2}n-s+1\\
&=\frac{s(n-1)(n-2)}{2}+\frac{s(s-1)}{2}n-s+1-sd.
\end{aligned}
\]

Ma si ha
\[
\begin{aligned}
\frac{s(n-1)(n-2)}{2}-\frac{s(s-1)}{2}n-s+1&{}\\ % [sic] -\frac{s(s-1)}{2}n -> +\frac{s(s-1)}{2}n, the sign disagrees with the preceding expression and the stated equality.
&=\binom{n-1+s}{3}-\binom{n-1}{3}-\binom{s-1}{3},
\end{aligned}
\]
quindi
\[
N_{n-4+s}-1-\binom{s-1}{3}=p_a+\pi_s-1.
\]

Concludiamo che \textit{il sistema lineare segato sopra \(F\) dalle superficie \(\Phi_{n-4+s}\) d’ordine \(n-4+s\), che è il sistema aggiunto a quello}
%CONT>>
% scan: s066L
% !PRINTED 110
\weberpage{110}
% head: CAPITOLO QUARTO
%<<CONT
% stamp: Dark scan border at the right edge, a horizontal mark in the upper left margin, and a stray stroke above the opening bar of |C| before "Si ha infatti".
\textit{\(|C_s|\) segato dalle superficie d’ordine \(s\), per \(s\) abbastanza alto, ha la dimensione}
\[
p_a+\pi_s-1,
\]
avendo designato con \(\pi_s\) il genere di \(|C_s|\), \(s\)-plo del sistema delle sezioni piane.

\edsection{4. Il genere numerico.}{4. Il genere numerico.}{sec.IV.4}

Riferiamoci sempre ad una superficie \(F\) dello spazio ordinario, avente un certo ordine \(n\) e singolarità normali. Ci proponiamo di valutare la dimensione del sistema \(|C'|\) aggiunto ad un sistema lineare irriducibile \(|C|\), appartenente alla \(F\). Ricordiamo anzitutto che le curve aggiunte \(C'\) segano sopra una \(C\) gruppi canonici; staccando la \(C\) da \(|C'|\) si ha il sistema canonico \(|K|\) di dimensione \(p_g-1\), quindi, per il gruppo canonico d’intersezione di una \(C'\) con una \(C\), dovranno passare \(\infty^{p_g}\) curve \(C'\). Segue di qui che la dimensione del sistema \(|C'|\) è data in ogni caso da
\[
r=p_g+\pi-1-\delta(C)
\]
dove \(\delta(C)\) designa la deficienza della serie segata dalle \(C'\) su una \(C\); si ha in particolare \(\delta(C)=0\) e
\[
r=p_g+\pi-1
\]
se \(|C'|\) sega sopra la \(C\) la serie completa \(g^{\pi-1}_{2\pi-2}\).

A priori la deficienza \(\delta(C)\) dipende dalla scelta del sistema \(|C|\); ma possiamo riconoscere una disuguaglianza fondamentale a cui dà luogo la considerazione della somma di due sistemi lineari irriducibili (\(\infty^1\) almeno)\origfootnote{1}{s066L-1}{Se si tratta di due fasci è inteso che siano distinti sicchè \(|C+L|\) risulti pure irriducibile.}, \(|C|\) e \(|L|\).

Si ha infatti
\[
\delta(C)\leq\delta(C+L).
\]

Supporremo per semplicità che i sistemi \(|C|\) ed \(|L|\) siano privi di punti base sopra la superficie \(F\). Partiamo dalla relazione fondamentale
\[
|(L+C)'|=|L+C'|=|L'+C|.
\]

In base a questa relazione il sistema \(|C'|\) si otterrà da \(|(L+C)'|\), staccando la curva \(L\), e ciò importa un numero di condizioni lineari che è dato dalla dimensione della serie segata da \(|(L+C)'|\) su \(L\), aumentata di una unità.

% scan: s066R
% !PRINTED 111
\weberpage{111}
% head: IL GENERE NUMERICO E IL TEOREMA DI RIEMANN-ROCH, ECC.
% stamp: Dark scan border at the left edge and short stray strokes above the final 2 in the fourth display.
Ora la serie indicata su \(L\) è la somma delle serie canonica \((LL')\) e dalla serie segata da \(|C|\), la quale è certamente d’ordine maggiore di zero, essendo i sistemi \(|C|\) ed \(|L|\) ambedue di dimensione maggiore o uguale a uno. % [sic] e dalla serie -> e della serie, the text is describing a sum of two series.

Ciò posto, designiamo con \(\pi\) il genere di \(|C|\), con \(\varrho\) il genere di \(|L|\), e con \(i\) il numero delle intersezioni \((LC)\). Allora il genere di \(|L+C|\), sarà
\[
\pi'=\pi+\varrho+i-1,
\]
e quindi la dimensione di \(|(L+C)'|\) sarà
\[
p_g+\pi'-1-\delta(L+C).=p_g+\pi+\varrho+i-2-\delta(L+C). % ?? uncertain: mark between \delta(L+C) and = -- full stop or stray ink.
\]

Invece la dimensione di \(|C'|\) vale
\[
p_g+\pi-1-\delta(C).
\]

La differenza fra le due dimensioni è indicata, come si è detto, dalla dimensione della serie non speciale \(g\) segnata da \(|(L+C)'|\) su \(|L|\), accresciuta di un’unità. Se la serie \(g\) è completa e quindi di dimensione
\[
2\varrho-2+i-\varrho=\varrho+i-2,
\]
risulta senz’altro
\[
\delta(C)=\delta(L+C);
\]
se invece la \(g\) è deficiente, la deficienza di essa esprime la differenza \(\delta(L+C)-\delta(C)\). In ogni caso si ha:
\[
\delta(C)\leq\delta(L+C).
\]

Il ragionamento sussiste ancora nel caso in cui i due sistemi \(|C|\) ed \(|L|\) abbiano dei punti base sopra \(F\). I punti base comuni ai due sistemi, o i punti base di \(|C|\) che non appartengono ad \(|L|\), non hanno a priori alcuna influenza. Se c’è un punto base di \(|L|\) che non appartenga a \(|C|\) e perciò, almeno virtualmente, neppure a \(|C'|\), l’intorno del punto (considerato come curva infinitesima) andrebbe sommato a \(|C'|\), ciò che non ne accresce la dimensione; così la conclusione precedente s’estende anche a questo caso.

Osserviamo di più che l’ipotesi fatta su \(|L|\) di appartenere ad un sistema irriducibile \(\infty^1\) almeno, contiene più del necessario: in realtà la \(L\) potrà essere anche una curva isolata e composta di più parti \(L_1, L_2, \ldots\), purchè nessuna di queste sia fondamentale per il sistema \(|C|\) e semprechè sommando a \(|C|\) le \(L_1, L_2, \ldotp\ldotp\ldotp\ldotp\) successivamente si ottengano sistemi lineari irriducibili di dimensione superiore. Infatti si potrà staccare la \(L\) da \(|(L+C)'|\) staccando successivamente le sue parti \(L_1, L_2 \ldots\), e se queste non sono fondamentali per \(|C|\), le serie segate su di esse, di cui ci occorre la dimensione, saranno sempre non speciali, contenendo in sè la serie canonica.

% scan: s067L
% !PRINTED 112
\weberpage{112}
% head: CAPITOLO QUARTO
% stamp: Dark scan border at the right edge and scattered scan specks.
Ora ricordiamo che si è calcolato la dimensione del sistema segato sopra \(F\) dalle superficie aggiunte per \(s\) abbastanza elevato: questo sistema, che è l’aggiunto del sistema \(s\)-plo delle sezioni piane, converrà qui designare con \(|L_s|\). Indicando con \(\pi_s\) il genere di \(|L_s|\), si è trovato che la dimensione del sistema \(|L'_s|\) vale
\[
p_a+\pi_s-1.
\]

D’altra parte essa si esprime in funzione del genere geometrico \(p_g\) e della deficienza della serie segata da \(|L'_s|\) su \(L_s\), valendo
\[
p_g+\pi_s-1-\delta(L_s).
\]

Di qui si deduce che
\[
\delta(L_s)=p_g-p_a,
\]
essendo \(s\), come si è detto, un numero abbastanza elevato.

Ciò posto si consideri su \(F\) un qualunque sistema irriducibile, \(\infty^1\) almeno, \(|C|\), di un certo genere \(\pi\), che supporremo al solito privo di punti base. È chiaro che per \(s\) abbastanza elevato vi sono superficie d’ordine \(s\) passanti per una \(C\), e quindi \(|C|\) è contenuto in \(|L_s|\), che anzi si presenta come la somma di \(|C|\) e di un altro sistema irriducibile \(|M|=|L_s-C|\).

Pertanto se ne deduce che
\[
\delta(C)\leq\delta(L_s).
\]

Questa conclusione si estende facilmente al caso in cui \(|C|\) sia dotato di punti base, con certe molteplicità \(i_1, i_2, \ldots\). Infatti, riferendoci ad un valore di \(s\) abbastanza alto, si potrà sempre imporre al sistema \(|L_s|\) di possedere gli stessi punti base di \(|C|\) con le stesse molteplicità, costruendo così un sistema \(|\bar{L}_s|\) che contenga \(|C|\), e tale anzi che il residuo \(|M|=|\bar{L}_s-C|\) sia un sistema lineare irriducibile. Ora è lecito supporre che i punti base di \(|\bar{L}_s|\) impongano \(\frac{i(i-1)}{2}\) condizioni lineari indipendenti alle curve aggiunte \(L'_s\) che debbano passare per un punto \(i\)-plo con la molteplicità \(i-1\). Avremo così un sistema lineare \(|\bar{L}_s|\) con gli stessi punti base di \(|C|\) per cui risulta ancora
\[
\delta(\bar{L}_s)=p_g-p_a.
\]
E si dedurrà come innanzi
\[
\delta(C)\leq\delta(\bar{L}_s).
\]

In tal guisa abbiamo dimostrato che: \textit{sopra la superficie \(F\) la deficienza della serie} (canonica) \textit{che il sistema aggiunto ad un sistema lineare irriducibile \(|C|\) sega sopra una \(C\) generica ha un valore massimo,}
%CONT>>
% scan: s067R
% !PRINTED 113
\weberpage{113}
% head: IL GENERE NUMERICO E IL TEOREMA DI RIEMANN-ROCH, ECC.
%<<CONT
% stamp: Dark scan border at the left edge and scattered scan specks.
\textit{che costituisce un carattere intrinseco della superficie e che eguaglia la differenza \(p_g-p_a\) fra il genere geometrico e il genere numerico di essa.} Questo \textit{massimo} è raggiunto per una classe di sistemi lineari appartenenti ad \(F\), che comprende in ogni caso i \textit{sistemi abbastanza grandi}, cioè tali che contengano in sè un multiplo d’ordine assai elevato d’un altro sistema irriducibile. Anzi questa classe è definita in guisa che « ogni sistema lineare irriducibile \textit{maggiore} di un altro di essa (cioè che lo contenga parzialmente) appartiene alla classe medesima ».

A priori potrà aversi soltanto una seconda classe ristretta di sistemi lineari \(|C|\), per cui la deficienza della serie (canonica) segata su \(C\) da \(|C'|\) resti inferiore al massimo.

Da quel che si è detto risulta che \textit{il genere numerico \(p_a\) di una superficie}, definito direttamente dalle formule di postulazione (sulla superficie data o su una trasformata con singolarità normali) costituisce un carattere \textit{invariante per trasformazioni birazionali} della superficie stessa, non soltanto nel caso in cui si abbia \(p_a=p_g\), sì anche quando sia \(p_a<p_g\).

Converrà tuttavia distinguere i due casi designando come \textit{superficie regolari} quelle per cui \(p_a=p_g\), ed \textit{irregolari} (di irregolarità \(p_g-p_a\)) quelle per cui \(p_a<p_g\).

Per le une e per le altre superficie il risultato ottenuto contiene l’estensione (dalle serie sopra una curva ai sistemi sopra una superficie) del teorema di Riemann-Roch, per ora limitatamente ai sistemi lineari \(|C'|\) che sono aggiunti ad altri sistemi irriducibili \(|C|\) (di genere \(\pi\) e grado \(n\)). La dimensione di \(|C'|\), o almeno un limite inferiore di essa, viene calcolata per mezzo del genere di \(|C|\) dalla formula:
\[
r'\geq p_a+\pi-1.
\]

Ma è facile esprimere la stessa \(r'\) pei caratteri, genere \(\pi'\) e grado \(n'\), del sistema \(|C'|\). Infatti — designando con \(p^{(1)}\) il genere lineare (relativo) della superficie \(F\), e supponendo, per semplicità di discorso, che \(|C|\) e \(|C'|=|C+K|\) siano privi di punti base su \(F\) si avrà
\[
\begin{aligned}
\pi'&=\pi+p^{(1)}+2\pi-2-n-1=3\pi-3-n+p^{(1)}\\
n'&=n+(p^{(1)}-1)+4\pi-4-2n=4\pi-5-n'+p^{(1)}, % ?? uncertain: apparent prime on the final n in the second equation -- n' or n.
% EDITORIAL D0066 RISOLTO: The upper-right stroke on the final n is distinct and strongly inked, compatible with a printed prime rather than a scanning fleck.
\end{aligned}
\]
e quindi
\[
\pi-1=n'-\pi'+1.
\]

Così la dimensione del sistema \(|C'|\) sarà
\[
r'\geq p_a+n'-\pi'+1,
\]
% EDITORIAL: segnatura originale: \textsc{Enriques F.} - \textit{Superficie algebriche.}\hfill 8
\printersignature{\textsc{Enriques F.} - \textit{Superficie algebriche.}\quad 8}
%CONT>>
% scan: s068L
% !PRINTED 114
\weberpage{114}
% head: CAPITOLO QUARTO
%<<CONT
% stamp: Dark scan border at the right edge and faint extraneous marks to the left of the postulation formula.
e similmente, se \(|C|\) stesso sia l’aggiunto di un altro sistema irriducibile, la sua dimensione varrà
\[
r\geq p_a+n-\pi+1.
\]

Vedremo più avanti come questa \textit{diseguaglianza fondamentale}, che costituisce il \textit{teorema di Riemann-Roch per i sistemi lineari aggiunti}, si estenda, con una piccola modificazione, a tutti i sistemi lineari appartenenti ad una superficie di genere numerico \(p_a\). Quindi diremo \textit{regolare} ogni sistema lineare completo non speciale la cui dimensione sia espressa dall’\textit{eguaglianza} precedente, e \textit{sovrabbondante} un sistema per cui valga propriamente la \textit{diseguaglianza}.

\edsection{5. Complementi.}{5. Complementi.}{sec.IV.5}

Negli sviluppi che precedono ci siamo riferiti ad una superficie \(F\) dello spazio ordinario, dotata di singolarità normali: in questa ipotesi abbiamo determinato la postulazione delle superficie aggiunte che sono soggette a passare semplicemente per la curva doppia di \(F\). Ma non vi è alcuna difficoltà ad estendere le nostre considerazioni, contemplando anche superficie con singolarità più elevate. Se si vuol seguire la via indicata nei precedenti paragrafi, si sarà condotti anzitutto ad estendere le formule di postulazione, determinando (con Noether) il numero delle condizioni che una curva, d’ordine \(m\) e genere \(p\), impone alle superficie d’ordine \(n\) abbastanza alto che debbano contenerla come \(s\)-pla, che è\origfootnote{1}{s068L-1}{Cfr. anche \textsc{L. Campedelli}, \textit{Sulla postulazione di una curva \(i\)-pla.} (Rend. Circolo Matematico di Palermo, to. 55, 1931).}:
\[
\frac{s(s+1)}{6}[(3n-4s+4)m-(2s+1)(p-1)].
\]\origfootnote{2}{s068L-2}{Da modificare per la presenza di punti multipli della curva.}

E similmente si dovrà valutare il numero delle condizioni imposte alle superficie d’un dato ordine dal possesso di un punto multiplo isolato, che per un punto \(s\)-plo è
\[
\binom{s+2}{3}.
\]

Ma, col complicarsi delle singolarità di \(F\), il calcolo della postulazione da cui dipende la dimensione dei sistemi di superficie aggiunte, diventa più difficile, siccome appare già da ciò che abbiamo detto in ordine alla influenza virtuale sul genere dei punti multipli propri di \(F\). Giova perciò designare una via più rapida, indipendente dalla natura delle singolarità della superficie \(F\). Indi%
%CONT>>
% scan: s068R
% !PRINTED 115
\unskip\weberpage{115}%
% head: IL GENERE NUMERICO E IL TEOREMA DI RIEMANN-ROCH, ECC.
%<<CONT
% stamp: Dark scan border at the left edge and scattered scan specks.
cheremo in ogni caso con \(n\) l’ordine di \(F\), con \(\pi\) il genere delle sue sezioni piane, con \(n_s\) e \(\pi_s\) i caratteri del sistema \(s\)-plo, segato su \(F\) dalle superficie d’ordine \(s\).

Osserviamo anzitutto che, supponendosi determinato il genere numerico \(p_a\) (in rapporto ad una trasformata di \(F\) con singolarità normali), conosciamo a priori la dimensione dei sistemi di superficie aggiunte ad \(F\), e possiamo affermare che le superficie aggiunte \(\Phi_{n-3+s}\) per \(s\) abbastanza grande, segano su \(F\) un sistema lineare completo di dimensione
\[
p_a+\pi_s-1,
\]
che a sua volta sega una serie completa d’ordine \(2\pi-2+sn\) sulla sezione piana \(C\) di \(F\). Perciò le \(\Phi_{n-3+s}\) segano su ogni piano il sistema completo delle curve d’ordine \(n-3+s\) aggiunte alla sezione piana \(C\).

Ciò posto si vede che \textit{le dimensioni dei sistemi lineari delle superficie \(\Phi_{n-3+s}\) aggiunte ad \(F\), per \(s\) abbastanza grande, formano una progressione aritmetica del terzo ordine che ha per ragione}
\[
\pi+sn+\frac{s(s-3)}{2}.
\]

Questa progressione, prolungata per i valori più piccoli di \(s\), vale a definire in ogni caso l’\textit{espressione virtuale} del numero delle superficie aggiunte di dato ordine linearmente indipendenti, e in particolare il \textit{genere numerico} \(p_a\).

In questo senso si può dire che la progressione aritmetica anzidetta, che è d’altronde una conseguenza delle formule di postulazione, tiene luogo della conoscenza effettiva di queste formule\origfootnote{1}{s068R-1}{Indipendentemente dalla circostanza che il gruppo base assegnato alle \(\Phi\) sia costituito dalle singolarità di una superficie \(F\) rispetto a cui le \(\Phi\) debbono essere aggiunte, si può dimostrare che i numeri esprimenti la dimensione del sistema della superficie \(\Phi\) d’ordine abbastanza alto costrette a passare con certe molteplicità per i punti e per le curve di un gruppo base assegnato, danno luogo sempre ad analoga progressione aritmetica del terzo ordine, la cui conoscenza vale, in molti casi, a supplire la conoscenza effettiva delle formule di postulazione. Cfr. \textsc{G. Castelnuovo}, \textit{Alcuni risultati sui sistemi lineari di curve appartenenti ad una superficie algebrica.} Mem. Soc. It. delle scienze (detta dei XL), serie III, t. 10, 1896. Riportato nelle “Memorie scelte„ pag. 335.}. % [sic] della superficie -> delle superficie, the following participle "costrette" is plural.

Riprendiamo a considerare i sistemi delle superficie \(\Phi_{n-3+s}\) aggiunte alla \(F\).

Indichiamo con \(h\geq0\) il più piccolo numero tale che per \(s\geq h\) codesti sistemi seghino su un piano la totalità delle curve aggiunte dello stesso ordine; vuol dire che, se non è \(h=0\), il sistema delle \(\Phi_{n-3+h-1}\) (che è d’ordine maggiore od uguale ad \(n-3\)) segherà su un piano un sistema di curve aggiunte non completo, con la de%
%CONT>>
% scan: s069L
% !PRINTED 116
\unskip\weberpage{116}%
% head: CAPITOLO QUARTO
%<<CONT
% stamp: Dark scan border at the right edge and scattered scan specks.
ficienza \(\delta_{h-1}>0\). Per \(s<h-1\) i sistemi delle \(\Phi_{n-3+s}\) (\(s=0, 1, \ldotp\ldotp\ldotp\ldotp\ h-2\)) segheranno sullo stesso piano sistemi di curve aggiunte, completi o no, che a priori potranno supporsi avere le deficienze
\[
\delta_0\geq0,\quad \delta_1\geq0;\quad \ldotp\ldotp\ldotp\ldotp,\quad \delta_{h-2}\geq0. % ?? uncertain: punctuation after \delta_1\geq0 -- semicolon or comma with a stray dot.
\]

Ora abbiamo visto che, in corrispondenza alla deficienza \(\delta_{h-1}\), e in confronto della sua espressione virtuale, cresce appunto di \(\delta_{h-1}\) la dimensione del sistema delle superficie aggiunte d’ordine \(n-3+h-2\) e quella del sistema delle curve da esse segato su \(F\), che varrà
\[
p_a+\pi_{h-2}-1+\delta_{h-1}.
\]

Similmente se \(\delta_{h-2}>0\) crescerà, in senso analogo, di \(\delta_{h-2}\) la dimensione del sistema segato dalle \(\Phi_{n-3+h-3}\), che varrà
\[
p_a+\pi_{h-3}+\delta_{h-2}+\delta_{h-1}. % [sic] p_a+\pi_{h-3}+\delta_{h-2}+\delta_{h-1} -> p_a+\pi_{h-3}-1+\delta_{h-2}+\delta_{h-1}, the constant -1 in the preceding dimension formula is missing here.
\]

Infine la dimensione del sistema canonico, segato dalle \(\Phi_{n-4}\), risulterà
\[
p_g-1=p_a-1+\delta_0+\delta_1+\ldotp\ldotp\ldotp\ldotp+\delta_{h-1}.
\]

Cioè l’\textit{irregolarità \(p_g-p_a\) della superficie eguaglia la somma delle deficienze dei sistemi di curve segati sopra un piano dalle superficie aggiunte ad \(F\), d’ordine}
\[
n-3+s\qquad \text{per}\qquad s=1,2,\ldotp\ldotp\ldotp\ldotp
\]

Possiamo precisare alquanto gli enunciati che precedono, dimostrando anzitutto che \textit{se le superficie aggiunte di un ordine maggiore od uguale ad \(n-2\) segano sopra un piano il sistema completo delle curve aggiunte dello stesso ordine, lo stesso accade per tutti i sistemi lineari di superficie aggiunte, d’ordine superiore.} (Si dirà poi come l’enunciato si estenda per \(s=0\), cioè al caso in cui le superficie aggiunte d’ordine \(n-3\) segano su un piano il sistema completo delle aggiunte di quest’ordine).

Per dimostrare il teorema enunciato, potremo passare dai sistemi delle superficie \(\Phi\) ai sistemi di curve da essi segati sopra \(F\), e quindi alle serie che questi sistemi segano sopra una sezione piana \(C\) di \(F\). E basterà stabilire che: « se le curve del sistema \(|sC+C'|\), per \(s\geq1\), segano su \(C\) la serie completa, anche la serie segata su \(C\) dalle curve di \(|(s+1)C+C'|\) è certo completa ».

Supporremo dapprima che sia \(s>1\) e — per semplicità di discorso — facciamo \(s=2\). L’ipotesi è dunque che le curve di \(|2C+C'|\), sezioni delle \(\Phi_{n-1}\) aggiunte, seghino su \(C\) la serie completa \(g\) d’ordine \(2\pi-2+2n\) determinata dalle curve aggiunte \(\varphi_{n-1}\). Ed evidentemente basta mostrare che la somma minima di questa \(g\) e della serie segata dalle rette è necessariamente completa.

% scan: s069R
% !PRINTED 117
\weberpage{117}
% head: IL GENERE NUMERICO E IL TEOREMA DI RIEMANN-ROCH, ECC.
% stamp: Dark scan border at the left edge and scattered scan specks.
Si tratta insomma di stabilire (con \textsc{Castelnuovo}) il seguente:

\textit{Lemma sulla integrità della somma di due serie appartenenti ad una curva.}

Sopra una curva di genere \(\pi\) si sommi ad una serie completa e non speciale \(g_m\) di un certo ordine \(m\) una serie \(\infty^1\) almeno \(g_n\) senza punti fissi, che sia contenuta in \(g_m\) e dia come residua rispetto a questa una serie non speciale; la somma minima di \(g_m\) e \(g_n\) è necessariamente completa.

Infatti si sommino alla \(g_m\), ch’è di dimensione \(m-\pi\) due gruppi \(G\) e \(G'\) della \(g_n\), da ritenere come fissi: si avranno così due serie lineari d’ordine \(n+m\) e di dimensione \(m-\pi\), che hanno a comune la serie formata da \(2n\) punti fissi costituenti il gruppo \(G+G'\) e dai gruppi della \(|g_m-g_n|\) che è una serie d’ordine \(m-n\), supposta non speciale, e perciò di dimensione
\[
m-n-\pi.
\]

Pertanto le nostre due serie, che hanno come fissi i gruppi \(G\) e \(G'\), saranno contenute in una serie di dimensione minima
\[
m+n-\pi
\]
che risulta dunque completa, c. d. d.

Questo lemma si applica immediatamente al caso \(h\geq2\). Ma, con lieve modificazione, il ragionamento si estende al caso \(h=1\) (non però ad \(h=0\)). Si tratta di provare che sommando alla serie \(g_{2\pi-2+n}\) segata su \(C\) dalle curve aggiunte \(\varphi_{n-2}\) la serie segata dalle rette, si ottiene sempre, come somma minima, la serie completa segata dalle \(\varphi_{n-1}\).

Infatti si sommino alla \(g_{2\pi-2+n}\) due gruppi fissi \(G\) e \(G'\) segati da due rette \(a\) e \(a'\) del piano; si avranno due serie lineari di dimensione \(\pi-2+n\), che hanno a comune la \(g^{\pi-1}_{2\pi-2}\) canonica (non già di dimensione \(\pi-2\), ma di dimensione \(\pi-1\)): la somma minima delle due serie coi due gruppi fissi \(G\) e \(G'\) è ora una serie d’ordine \(2\pi-2+2n\) e di dimensione \(\pi-2+n-1\) anzichè \(\pi-2+n\). Questa serie, non completa, viene segata su \(C\) dal sistema di tutte le \(\varphi_{n-1}\) che passano per il punto comune alle rette \(a\) e \(a'\). Ora il minimo sistema lineare contenente entro di sè due sistemi di \(\varphi_{n-1}\) con un punto base diverso è il sistema completo di tutte le \(\varphi_{n-1}\). Risulta pertanto che la somma minima della serie \(g_{2\pi-2+n}\) e della serie segata dalle rette del piano riesce necessariamente la serie completa segnata dalle \(\varphi_{n-1}\).

\edsection{6. Nota sulle condizioni di regolarità di una superficie.}{6. Nota sulle condizioni di regolarità di una superficie.}{sec.IV.6}

Abbiamo veduto che, data una superficie \(F\) d’ordine \(n\) dello spazio ordinario, che per semplicità vogliamo supporre dotata di singolarità
%CONT>>
% scan: s070L
% !PRINTED 118
\weberpage{118}
% head: CAPITOLO QUARTO
%<<CONT
% stamp: Dark scan border at the right edge and an extraneous slanted stroke in the right margin beside "Infatti, per il".
normali, se le \(\Phi_{n-2}\) aggiunte a questa, d’ordine \(n-2\), segano su un piano il sistema completo delle curve aggiunte alla sezione di \(F\), lo stesso accade per le superficie \(\Phi_{n-2+r}\) aggiunte d’ordine \(>n-2\); in questo caso dunque la irregolarità \(p_g-p_a\) della superficie è uguale alla deficienza della serie (canonica) segata sulle sezioni piane \(C\) di \(F\) dalle \(\Phi_{n-3}\). In particolare, se le \(\Phi_{n-3}\) segano sulle dette \(C\) la serie canonica completa, l’irregolarità \(p_g-p_a=0\), ossia la superficie \(F\) è regolare.

Ma il criterio di regolarità, così stabilito da \textsc{Enriques}, può ridursi ad altro più semplice.

\textit{La superficie \(F\), d’ordine \(n\) dello spazio ordinario è regolare \((p_a=p_g)\) se le superficie \(\Phi_{n-3}\) d’ordine \(n-3\), aggiunte ad essa, segano su una sezione piana la serie canonica completa.}

Daremo un rapido cenno della dimostrazione che \textsc{Castelnuovo}\origfootnote{1}{s070L-1}{\textit{Alcune proprietà fondamentali dei sistemi lineari di curve tracciati sopra una superficie algebrica.} Annali di Mat., serie II, t. 25 (1897). Cfr. \textit{Memorie scelte.} XXIII.} ha recato di questo teorema. Conviene anzitutto stabilire il

\textit{Lemma.} – Se un gruppo di \(m<n\) punti della nostra superficie \(F\) scelti fra gli \(n\) punti segati da una retta \(a\), appartiene ad una serie speciale \(g_m^h\), di dimensione \(h\geq1\), sopra una sezione piana \(C\) per \(a\), lo stesso accade per ogni altra sezione di \(F\) con un qualsiasi piano per \(a\).

È una conseguenza immediata dell’ipotesi che le \(\Phi_{n-3}\) seghino sulle sezioni piane \(C\) di \(F\) la serie canonica completa. Infatti, per il teorema di Riemann-Roch, un gruppo della detta \(g_m^h\) su \(C\) offre precisamente \(m-h\) condizioni ai gruppi canonici di \(C\), ovvero anche alle \(C'\), che debbono contenerlo: ma il numero di tali condizioni è indipendente dalla scelta di una particolare \(C\) del fascio segato su \(F\) dai piani per \(a\).

Sulla base del lemma anzidetto siamo in grado di riconoscere che « se le \(\Phi_{n-3}\) segano la serie canonica completa sulle sezioni piane \(C\) di \(f\), la serie caratteristica del sistema completo \(|C|\) riesce completa ». % [sic] f -> F, the surface is otherwise denoted by capital F.

Sappiamo che il sistema completo \(|C|\) si può costruire su \(F\) mercè il sistema lineare di tutte le aggiunte \(\Phi_{n-2}\) che passano per la sezione \(C'\) di una \(\Phi_{n-3}\) (teorema del resto). Basterà quindi mostrare che codeste \(\Phi_{n-2}\) per una \(C'\) segano sopra un piano la serie completa delle curve \(\varphi_{n-2}\) d’ordine \(n-2\) aggiunte alla sezione di \(F\).

A tal uopo si cerchi di costruire una superficie \(\Phi_{n-2}\) aggiunta ad \(F\), come luogo di una curva \(\varphi_{n-2}\) dello stesso ordine \(n-2\) aggiunta alla sezione \(C\) di un piano \(\alpha\) variabile in un fascio, che abbia per asse una retta generica dello spazio, \(a\), non incidente a \(C'\). Alla \(\varphi_{n-2}\) s’imporrà:
%CONT>>
% scan: s070R
% !PRINTED 119
\weberpage{119}
% head: IL GENERE NUMERICO E IL TEOREMA DI RIEMANN-ROCH, ECC.
%<<CONT
% stamp: Dark scan border at the left edge and scattered scan specks.
1) di passare per i \(2\pi-2\) punti di un gruppo canonico, se zione di \(C'\); % [sic] se zione -> sezione, the word is split at a printed line end without a hyphen.

2) di passare per \(s-1\) punti, \(A_1, A_2, \ldotp\ldotp\ldotp\ldotp A_{s-1}\), scelti fra gli \(n\) punti comuni a \(C\) e alla retta \(a\), designando \(s\) la dimensione della serie caratteristica di \(C\) resa completa, cioè la dimensione della serie completa cui appartiene la \(g_n\) segata sulla \(C\) dalle rette del piano \(\alpha\);

3) e infine di toccare un piano \(\gamma\), passante per un punto \(A_1\) di codesto gruppo di \(s-1\) punti e non contenente \(a\).

Le condizioni 1) e 2) determinano, in un piano generico per \(a\), un fascio di curve \(\varphi_{n-2}\), cui appartiene la curva composta della retta \(a\) e della \(\varphi_{n-3}\) passante per il gruppo canonico di \(C\) sezione di \(C'\). Aggiungendo la condizione 3) resta determinata nel piano \textit{una} curva \(\varphi_{n-2}\); e, al variare di questo piano, la \(\varphi_{n-2}\) così costruita descriverà una superficie \(\Phi_{n-2+h}\), aggiunta ad \(F\) che a priori potrà essere d’un ordine \(n-2+h>n-2\), passando \(h\) (\(>0\)) volte per la retta \(a\). Per il nostro scopo occorre dimostrare che deve essere \(h=0\). Si supporrà dunque \(h>0\) e si ridurrà all’assurdo questa ipotesi.

A tal uopo si osservi che, essendo \(h>0\), si avranno, in un punto \(A_s\) di \(a\), \(h\) piani tangenti alla \(\Phi_{n-2+h}\), uno qualunque dei quali \(\bar{\alpha}\) segherà questa superficie secondo una curva composta della retta \(a\) contata \(h\) volte e di una \(\varphi_{n-2}\) passante per gli \(s\) punti \(A_1A_2\ldotp\ldotp\ldotp\ldotp A_{s-1}A_s\). Ora si può porre \(A_s\) in una delle \(n-(s-1)\) intersezioni della superficie \(F\) o della curva \(\bar{C}\) sezione di \(\bar{\alpha}\) con \(a\), fuori di \(A_1A_2\ldotp\ldotp\ldotp\ldotp A_{s-1}\): siccome la \(\bar{\varphi}_{n-3}\) sezione di \(\bar{\alpha}\) non passa per \(A_1\) (essendo \(a\) non incidente a \(C'\)), bisogna che il gruppo di punti \(A_1A_2\ldotp\ldotp\ldotp\ldotp A_s\) appartenga, non già ad una, ma ad un fascio di curve \(\varphi_{n-2}\) aggiunte alla sezione \(\bar{C}\) di \(F\) con \(\bar{\alpha}\), e quindi che codesto gruppo appartenga ad una \(g_s^1\) speciale, di dimensione (almeno) eguale ad 1. Ma questo non può accadere sopra un particolare piano \(\bar{\alpha}\) per \(a\), senza che accada per tutti i piani \(\alpha\). Pertanto si deduce che \(h=0\), e così riusciamo a costruire una superficie \(\Phi_{n-2}\) aggiunta ad \(F\) che soddisfi alle condizioni di:

1) passare per una curva \(C'\) sezione di \(f\) con una aggiunta \(\Phi_{n-3}\); % [sic] f -> F, the surface is otherwise denoted by capital F.

2) passare per \(s-1\) fra gli \(n\) punti intersezioni di \(F\) con una retta \(a\);

3) toccare un piano \(\gamma\) per uno di questi punti.

Variando i punti che figurano nella condizione 2) si possono ottenere \(\infty^{s+1}\) \(\Phi_{n-2}\) per \(C'\), le quali segano, su un piano generico \(\alpha\), \(\infty^s\) curve \(\varphi_{n-2}\), cioè la totalità delle curve aggiunte \(\varphi_{n-2}\) passanti per il gruppo canonico \(G_{2\pi-2}\) sezione di \(C'\). Ciò significa che le \(\Phi_{n-2}\) per \(C'\) segano su \(F\) un sistema lineare \(|C|\) di dimensione \(s+1\), che ha la serie caratteristica completa, c. d. d.

Per stabilire il teorema di Castelnuovo basta ormai rilevare che: « sopra una curva \(C\) di genere \(\pi(>0)\) la somma minima della serie
%CONT>>
% scan: s071L
% !PRINTED 120
\weberpage{120}
% head: CAPITOLO QUARTO
%<<CONT
% stamp: Scattered scan speckles and stray marginal marks.
canonica \(g_{2\pi-2}^{\pi-1}\) e d’un’altra serie completa \(g_n^s\), priva di punti fissi e non composta con una involuzione, è sempre completa ».

Ciò risulta dall’osservare che un gruppo \(G_n\) della seconda serie presenta \(s\) condizioni indipendenti ad un gruppo della stessa \(g_n^s\) che debba contenerlo e \(n-s\) condizioni ad un gruppo della serie canonica, quindi\origfootnote{1}{s071L-1}{Cfr. \textsc{Castelnuovo}, \textit{Sui multipli di una serie lineare di gruppi di punti appartenente ad una curva algebrica.} Rendic. Circolo di Palermo, 1893. \textit{Memorie scelte}, VIII.} almeno
\[
s+n-s-1=n-1
\]
condizioni ad un gruppo della serie somma minima \(g_{2\pi-2+n}=g_{2\pi-2}^{\pi-1}+g_n^s\); perciò la dimensione di questa supera di (almeno) \(n-1\) la dimensione della serie canonica, e vale dunque
\[
(n-1)+\pi-1=n+\pi-2:
\]
vuol dire che la serie minima \(g_{2\pi-2}^{\pi-1}+g_n^s\) ha la dimensione della serie completa \(g_{2\pi-2+n}^{\pi-2+n}\), ossia coincide con questa.

\textit{Nota.} – Il teorema di Castelnuovo che esprime le condizioni di regolarità della superficie, in relazione al sistema \(|C|\) delle sue sezioni piane, si estende facilmente ai sistemi lineari irriducibili \(|C|\) anche non semplici, come l’autore stesso ha avvertito. Avremo dunque:

\textit{La condizione necessaria e sufficiente perchè una superficie sia regolare \((p_a=p_g)\) è che il sistema \(|C'|\) aggiunto ad un sistema lineare irriducibile \(|C|\), di genere \(\pi>0\), \(\infty^2\) almeno, seghi sulla \(C\) generica la serie canonica completa, cioè che sia di dimensione \(p_g+\pi-1\).}

\textit{Ogni sistema lineare irriducibile completo appartenente ad una superficie regolare ha la serie caratteristica completa.}

Ritroveremo più avanti quest’ultimo teorema come caso particolare di una proprietà relativa alla deficienza della serie caratteristica dei sistemi lineari completi appartenenti a superficie irregolari.

\edsection{7. Esempi di superficie irregolari.}{7. Esempi di superficie irregolari.}{sec.IV.7}

Gli sviluppi che precedono vengono illustrati da qualche esempio. Per la superficie generale d’ordine \(n\), priva di singolarità, le superficie aggiunte \(\Phi_{n-4}\) costituiscono la totalità delle superficie d’ordine \(n-4\); il numero effettivo di queste, linearmente indipendenti, è uguale al numero virtuale
\[
p_g=p_a=\binom{n-1}{3}=\frac{(n-1)(n-2)(n-3)}{6}\text{\origfootmark{2}{s071L-2}}.
\]\origfoottext{2}{s071L-2}{Se si vogliono applicare le formule generali di postulazione, si deve ritenere che la curva doppia d’ordine \(n=0\) sia di genere \(\varrho=1\).}

% scan: s071R
% !PRINTED 121
\weberpage{121}
% head: IL GENERE NUMERICO E IL TEOREMA DI RIEMANN-ROCH, ECC.
% stamp: Dark vertical scan artefact along the left edge; scattered speckles.
Dunque, le superficie generali del proprio ordine sono regolari; in particolare i piani, le quadriche e le superficie cubiche che non siano coni, quindi anche tutte le superficie razionali hanno il genere superficiale
\[
p_g=p_a=0.
\]

Ora, se s’impone ad una superficie d’ordine \(n\) di possedere una curva doppia, d’ordine relativamente piccolo, si ottengono del pari superficie regolari; se invece si vuole che la superficie possegga una curva doppia d’ordine piuttosto alto, si trova a priori un numero di condizioni che eccede il numero dei parametri da cui dipende la superficie del dato ordine; pertanto riesce difficile, per questa via, ottenere esempi di superficie irregolari.

Ma si prenda a considerare un cono \(F\) d’un ordine \(n>2\), e di genere \(p>0\), e sia \(O\) il suo vertice (punto \(n\)-plo); qui mancano certo le superficie aggiunte d’ordine \(n-4\), che dovrebbero passare per il punto \(O\) con la molteplicità \(n-2>n-4\), sicchè risulta (come già sappiamo)
\[
p_g=0.
\]

Mancano anche per la stessa ragione le superficie \(\Phi_{n-3}\) aggiunte ad \(F\) che dovrebbero segnare su di essa \(p_a+p\) curve aggiunte alle sezioni piane linearmente indipendenti, sicchè
\[
p_a+p\leq 0,
\]
cioè
\[
p_a\leq -p.
\]

Ma esistono invece dei coni aggiunti \(\Phi_{n-2}\) (con \(O\) \((n-2)\)-plo) i quali segano sopra un piano il sistema completo delle curve d’ordine \(n-2\) aggiunte alla sezione di \(F\). Perciò i nostri risultati ci autorizzano a dire che il genere numerico
\[
p_a=-p.
\]

Si conclude che \textit{le rigate di genere \(p\) hanno il genere numerico \(-p\) e l’irregolarità}
\[
p_g-p_a=p.
\]

Questo teorema riesce dimostrato trasformando la rigata in un cono; ma anche senza eseguire questa trasformazione si può ripetere il ragionamento precedente, basato su ciò, che: \(1^\circ)\) non esistono sulla superficie curve aggiunte alle sezioni piane \(C\); e \(2^\circ)\) le curve di \(|C+C'|\) (sezioni delle aggiunte \(\Phi_{n-2}\)) segano su una \(C\) la serie completa.

L’esempio delle rigate è noto fino dai primi studi di A. \textsc{Cayley}\origfootnote{1}{s071R-1}{. \textit{On the deficiency of certain surfaces.} Math. Annalen, Bd. III.} % ?? uncertain: the small dot following the footnote label -- printed full stop (read), scan speck.
%CONT>>
% scan: s072L
% !PRINTED 122
\weberpage{122}
% head: CAPITOLO QUARTO
%<<CONT
% stamp: Dark vertical scan artefact along the right edge; scattered speckles and ink smudges.
(1871) e di M. \textsc{Noether}, che hanno riconosciuto l’espressione virtuale del genere diventare in questo caso negativa. Successivamente il primo esempio di superficie irregolare, per cui \(p_g>0\) è stato indicato da G. \textsc{Castelnuovo}\origfootnote{1}{s072L-1}{Cfr. G. \textsc{Castelnuovo}, \textit{Osservazioni intorno alla geometria sopra una superficie.} Rend. Ist. Lombardo, 1891. Nota I. “ Memorie scelte „ n. XVII.}.

Si imponga ad una superficie \(F_6\) d’ordine 6 di possedere 8 punti tripli \(A_1,A_2,\ldots,A_8\) in posizione arbitrariamente assegnata nello spazio. Le \(F_6\) dipendono da 83 costanti, l’imposizione di un punto triplo importa 10 condizioni lineari, quindi le \(F_6\) che passano triplamente per \(A_1,\mathinner{\ldotp\ldotp\ldotp\ldotp},A_8\) saranno \(\infty^3\), e si spezzeranno nelle terne di quadriche (appartenenti a un fascio) che passano per quegli 8 punti. Ma invece di prendere \(A_1,A_2,\mathinner{\ldotp\ldotp\ldotp\ldotp},A_8\) in posizione arbitraria, si supponga che essi costituiscano il gruppo base di una rete di quadriche \(\Phi_2\). In questo caso, combinando linearmente le terne di \(\Phi_2\), si otterrà un sistema lineare di \(F_6\) irriducibili, di dimensione 9\origfootnote{2}{s072L-2}{Si può dimostrare che il sistema delle \(F_6\) con gli otto punti tripli non può avere dimensione \(r>9\), valutando \(r\) per mezzo della dimensione del sistema che queste \(F_6\) segano sopra una quadrica della rete.}. % ?? uncertain: damaged letters in "superficie" and "appartenenti" -- i and r (read), apostrophe-like and dot-like marks.
% EDITORIAL D0069 RISOLTO: The damaged i in superficie retains its dot and remnants at the stem position. The r of appartenenti is severely broken, but its position and surviving upper stroke support the ordinary letters read in the transcription.

Ora le \(F_6\) così costruite risultano di genere \(p_g=3\), avendo come aggiunte d’ordine \(6-4=2\) le quadriche \(\Phi_2\) della rete considerata. Ma gli 8 punti base \(A_1,A_2,\mathinner{\ldotp\ldotp\ldotp\ldotp},A_8\) impongono alle \(\Phi_2\) sette condizioni lineari indipendenti invece di otto, che è il numero virtuale di queste condizioni. Perciò il genere numerico delle \(F_6\) vale % ?? uncertain: number of dots in this point-list ellipsis -- four (read), three.
% EDITORIAL D0070 RISOLTO: The magnified list distinctly contains four evenly spaced dots between the comma after A_2 and the comma before A_8.
\[
p_a=p_g-1=2.
\]

È interessante notare che quello costruito è il primo caso di una serie di esempi analoghi. Si possono costruire in generale superficie irriducibili d’ordine \(2n\) \((n\geq 3)\) che posseggono 8 punti \(n\)-pli nei punti base di una rete di quadriche: queste \(F_{2n}\) risultano irregolari di genere
\[
p_g=\frac{n(n-1)}{2}
\]
e
\[
p_a=n-1.
\]

Qui interessa osservare che le quadriche della rete segano sopra una \(F_{2n}\) un sistema lineare composto con quartiche ellittiche, e queste quartiche formano su \(F_{2n}\) un fascio irrazionale di genere
\[
p_g-p_a=\frac{(n-1)(n-2)}{2};
\]
infatti il genere del fascio eguaglia il genere del cono osculatore della \(F_{2n}\) in uno dei punti \(A\), che è punto base per il fascio delle quartiche stesse.

% scan: s072R
% !PRINTED 123
\weberpage{123}
% head: IL GENERE NUMERICO E IL TEOREMA DI RIEMANN-ROCH, ECC.
% stamp: Scattered scan speckles and marginal streaks, especially at the right edge.
Aggiungiamo che la serie degli anzidetti esempi, indicati da \textsc{Castelnuovo}, si lascia generalizzare come ha mostrato L. \textsc{Campedelli} (1936)\origfootnote{1}{s072R-1}{L. \textsc{Campedelli}, in Atti R. Accademia di Torino, 1936.}.

La costruzione di una \(F_{2n}\) irriducibile con 8 punti \(n\)-pli dipende invero da quella di un fascio lineare di curve ellittiche d’un ordine \(m\) divisore di \(n\) \((n=ms)\), che possegga 8 punti base \(s\)-pli sopra una quadrica: questo gruppo di punti base appartiene sempre ad una quartica ellittica, e dà su questa un gruppo proveniente dalla divisione per \(n\) della serie segata sulla curva da tutte le superficie d’ordine \(2n\). Il fascio delle curve ellittiche di cui si tratta si lascia proiettare sul piano da uno dei suoi punti base in un fascio di Halphen di curve \(C_{3m}\) con 9 punti base \(m\)-pli.

A partire dal fascio indicato sopra una quadrica si ottengono superficie \(F_{2n}\), con 8 punti \(n\)-pli, irregolari, che posseggono sempre un fascio irrazionale di genere \(p_g-p_a\).

\textit{Nota.} – In base al teorema del \S~6 che afferma l’integrità della serie caratteristica d’un sistema lineare completo sopra una superficie regolare, si riconosce in generale che « il possesso d’un fascio irrazionale (di genere \(\pi>0\)) di curve, porta che la superficie sia irregolare ».

Più precisamente troveremo che la superficie ha l’irregolarità \(p_g-p_a\geq\pi\). Ma di ciò più avanti.

\edsection{8. Il teorema di Riemann-Roch per le superficie: sistemi più ampi del sistema canonico.}{8. Il teorema di Riemann-Roch per le superficie: sistemi più ampi del sistema canonico.}{sec.IV.8}

Abbiamo riconosciuto che un sistema lineare di grado \(n\) e genere \(\pi\), che sia l’aggiunto di un altro sistema irriducibile sopra una superficie di genere numerico \(p_a\), ha la dimensione
\[
r\geq p_a+n-\pi+1,
\]
e questa disuguaglianza fondamentale, pur nelle condizioni restrittive che qui sono supposte, costituisce un’estensione alle superficie del noto teorema di Riemann-Roch relativo alle curve. Conviene anzitutto rendersi conto del significato della restrizione che \(|C|\) sia il sistema aggiunto di un altro sistema irriducibile. Per semplicità di discorso riteniamo che \(|C|\) sia privo di punti base sulla superficie \(F\), e lo stesso accada pel sistema \(|L|\) di cui esso è l’aggiunto. Designando con \(|K|\) il sistema canonico avremo
\[
|C|=|L+K|.
\]

% scan: s074L
% !PRINTED 124
\weberpage{124}
% head: CAPITOLO QUARTO
% stamp: Dark vertical scan artefact along the right edge; scattered speckles and short stray marks.
Così la condizione imposta a \(|C|\) di essere l’aggiunto di un altro sistema, si traduce in quella di contenere il sistema canonico, dando luogo ad un sistema residuo \(|L|\), d’ordine maggiore di zero, o, come diremo, di essere \textit{più ampio} del sistema canonico.

Un sistema \(|C|\) più ampio del sistema canonico \(|K|\) si potrà sempre ritenere come l’aggiunto del sistema
\[
|L|=|C-K|,
\]
ma non si può dire a priori che \(|L|\) sia irriducibile. Anzi vi sono dei casi in cui, pur accrescendo la dimensione di \(|C|\) col moltiplicarlo per un numero grande ad arbitrio, il sistema \(|L|\) sempre riesce riducibile, contenendo delle componenti fisse, non connesse con le componenti variabili. Ciò accade invero se \(|C|\) è il sistema delle sezioni piane di una superficie \(F\) la quale possegga un punto multiplo proprio influente sul genere, per esempio un punto multiplo isolato, d’ordine \(\nu\geq 3\); % ?? uncertain: punctuation after 3 -- semicolon (read), comma with a speck above it.
che rappresenta una curva fondamentale \(\theta\) di \(|C|\), sebbene qui sia facile riconoscere che la disuguaglianza di cui si tratta (dimensione di \(|C|\geq p_a+n-\pi+1\)) sussiste a più forte ragione. Infatti, designando con \(\bar{n}\) e \(\bar{\pi}\) i caratteri del sistema \(\bar{C}=|C-\theta|\), % [sic] \bar{C}=|C-\theta| -> |\bar{C}|=|C-\theta|, the left-hand system delimiters are absent in the print.
avremo che \(|\bar{C}|\) è l’aggiunto del sistema irriducibile \(|\bar{L}|=|L-\theta|\):
\[
|\bar{C}|=|\bar{L}'|,
\]
e quindi la dimensione di \(|\bar{C}|\) è
\[
\bar{r}\geq p_a+\bar{n}-\bar{\pi}+1;
\]
ora il grado di \(\theta\) è \(-\nu\) e il suo genere \(\varrho=\frac{(\nu-1)(\nu-2)}{2}\geq 1\), quindi si avrà
\[
\begin{aligned}
n&=\bar{n}-\nu\\ % ?? uncertain: small marks after nu -- scan speckles with no punctuation (read), comma or full stop.
\bar{\pi}&=\bar{\pi}+\varrho+\nu-1 % ?? uncertain: faint stroke over the left-hand pi -- overbar (read), stray scan mark.
\end{aligned}
\]
e di conseguenza la dimensione di \(|C|\), \(r>\bar{r}\), varrà
\[
r>p_a+n-\pi+1.
\]

Ma, per procedere rigorosamente con la massima generalità, dimostreremo che ogni sistema lineare \(|C|\) più ampio del sistema canonico \(|K|\) ha la dimensione \(r\geq p_a+n-\pi+1\), ricorrendo ad altre considerazioni che valgono ugualmente per sistemi \(|C|\) riducibili o irriducibili.

Data sopra la superficie \(F\) una curva \(C\) irriducibile o riducibile, possiamo sempre supporre che le superficie aggiunte ad \(F\), \(\Phi_{n-4+s}\) (d’ordine \(n-4+s\) abbastanza alto) che siano costrette a passare per \(C\) non passino di conseguenza per un’altra curva o per punti
%CONT>>
% scan: s074R
% !PRINTED 125
\weberpage{125}
% head: IL GENERE NUMERICO E IL TEOREMA DI RIEMANN-ROCH, ECC.
%<<CONT
% stamp: Scattered scan speckles and short stray marks.
fuori di \(C\) e della curva doppia\origfootnote{1}{s074R-1}{Ciò risulta dall’osservare che nelle formule di postulazione relative ad una curva composta (o ad una curva e a un gruppo di punti) ciascuna componente porta un numero effettivo di condizioni.}. Ciò posto la curva \(C\) sarà contenuta nel sistema \(|L'_s|\) segato su \(F\) dalle \(\Phi_{n-4+s}\), e darà rispetto a questo un sistema residuo irriducibile (senza punti base)
\[
|D|=|L'_s-C|.
\]

Ora, designando con \(r'_s\) la dimensione, con \(\pi'_s\) e \(n'_s\) il genere e il grado di \(|L'_s|\), abbiamo trovato
\[
r'_s\geq p_a+n'_s-\pi'_s+1
\]
(anzi in questo caso vale proprio il segno di eguaglianza); staccando una curva \(D\), si troverà che il sistema residuo \(|C|=|L'_s-D|\) deve avere la dimensione
\[
r\geq p_a+n-\pi+1,
\]
qualora si sappia che la serie segata da \(|L'_s|\) su \(D\) è non speciale, siccome può affermarsi nel caso che \(|C|\) sia più ampio del sistema canonico \(|K|\), le cui curve segano su \(D\) gruppi residui della serie caratteristica di \(D\).

Sviluppiamo il calcolo introducendo i caratteri genere \(\varrho\) e grado \(\nu\) di \(D\) e il numero \(t\) delle intersezioni di \(C\) e \(D\). La serie segata da \(|L'_s|\) su \(D\) avrà l’ordine \(\nu+t\), e (se sia completa) la dimensione \(\nu+t-\varrho\), quindi lo staccamento di \(D\) diminuisce di \(\nu+t-\varrho+1\) al più la dimensione \(r'_s\) di \(|L'_s|\); segue
\[
r\geq r'_s-(\nu+t-\varrho+1).
\]

D’altra parte
\[
\begin{aligned}
n'_s&=n+\nu+2t\\
\pi'_s&=\pi+\varrho+t-1
\end{aligned}
\]
e quindi
\[
n'_s-\pi'_s=n-\pi+\nu-\varrho+t+1;
\]
per conseguenza
\[
r\geq p_a+n-\pi+1.
\]

Abbiamo così dimostrato che: \textit{per i sistemi lineari irriducibili o riducibili di grado \(n\) e genere \(\pi\) più ampi del sistema canonico, la dimensione vale sempre}
\[
r\geq p_a+n-\pi+1.
\]

Questa diseguaglianza costituisce il cosidetto teorema di Riemann-Roch per le superficie, che si estenderà poi — nel modo che vedremo — ai sistemi lineari qualunque.

% scan: s075L
% !PRINTED 126
\weberpage{126}
% head: CAPITOLO QUARTO
% stamp: Dark vertical scan artefact along the right edge; scattered speckles and stray marks.
Intanto, riferendoci ai \textit{sistemi più ampi del sistema canonico} distingueremo:

1) i sistemi \textit{regolari} per cui vale l’eguaglianza
\[
r=p_a+n-\pi+1
\]

2) e i sistemi \textit{sovrabbondanti} per cui invece
\[
r>p_a+n-\pi+1.
\]

Abbiamo già veduto che esistono certi sistemi regolari, chè tali sono gli aggiunti ai multipli d’ordine assai elevato d’un sistema irriducibile \(|L|\).

Il risultato che abbiamo ottenuto porta subito conseguenze notevoli.

Sia \(|C|\) un \textit{sistema lineare irriducibile} di genere \(\pi\) e grado \(n\), \textit{più ampio del sistema canonico e regolare} sopra la superficie \(F\), e si sommi ad esso una curva irriducibile \(D\) di genere \(\varrho\) e grado \(\nu\), la quale intersechi le \(C\) in un certo numero \(m\) \((>0)\) di punti.

Il grado e il genere di \(|C+D|\) sono rispettivamente
\[
n+\nu+2m
\]
e
\[
\pi+\varrho+m-1;
\]
e quindi la loro differenza vale
\[
(n-\pi)+(\nu-\varrho)+m+1>n-\pi
\]
se sia
\[
m\geq\varrho-\nu.
\]

Si deduce che \textit{il sistema \(|C+D|\) ha, in tale ipotesi, dimensione maggiore di}
\[
p_a+n-\pi+1
\]
che è la dimensione di \(|C|\); e perciò è certo \textit{irriducibile}.

Inoltre questo sistema \(|C+D|\) sarà \textit{regolare} e dovrà segare su \(D\) una serie (non speciale) completa, altrimenti staccando \(D\) si otterrebbe un sistema \(|C|\) sovrabbondante.

Si avverta che questo ampliamento di \(|C|\) sarebbe impossibile per \(m=0\), cioè se la curva \(D\) fosse fondamentale per \(|C|\); in questo caso dovrà dunque aversi \(\nu-\varrho<0\), il che resulta a priori dall’osservare che la curva fondamentale \(D\) di genere \(\varrho\geq 0\), deve avere il grado negativo: \(\nu<0\), essendo \(-\nu\) il numero delle intersezioni di \(D\) colle sue curve residue (di \(|C-D|\)).

Se nell’osservazione precedente s’inverte l’ufficio di \(|C|\) e \(|D|\), si è condotti ad affermare: \textit{dato sopra la superficie \(F\) un sistema lineare} (virtualmente privo di punti base) \(|C|\), \textit{comunque riducibile,}
%CONT>>
% scan: s075R
% !PRINTED 127
\weberpage{127}
% head: IL GENERE NUMERICO E IL TEOREMA DI RIEMANN-ROCH, ECC.
%<<CONT
% stamp: Scattered scan speckles and stray marginal marks.
\textit{si può sempre sommargli un sistema lineare irriducibile regolare \(|D|\) in guisa da ottenere un sistema irriducibile \(|D+C|\) egualmente regolare.}

È sufficiente assumere come sistema \(|D|\) il multiplo \(|L_s|=|sL|\) d’un sistema regolare \(|L|\) privo di curve fondamentali, e perciò avente colle \(C\) (o colle componenti di esse) un numero d’intersezioni grande con \(s\). Nel caso che le curve \(C\) siano riducibili, p. es., \(C=C_1+C_2+\mathinner{\ldotp\ldotp\ldotp\ldotp}\), si sommerà anzitutto un multiplo di \(|L|\) a \(C_1\) e poi il sistema irriducibile (ampio quanto si vuole) così ottenuto a \(C_2\) e così di seguito, finchè si pervenga ad un sistema \(|sL+C|\) irriducibile.

\edsection{9. Eliminazione delle curve eccezionali.}{9. Eliminazione delle curve eccezionali.}{sec.IV.9}

Riprendiamo le considerazioni del paragrafo precedente in ordine alla somma di una curva ad un sistema lineare regolare \(|C|\), sopra una superficie \(F\). E prendiamo come curva sommanda a \(|C|\) una curva eccezionale \(\omega\) di genere \(\varrho=0\) e di grado \(\nu\geq -1\). Se essa non è fondamentale per \(|C|\), e quindi incontra le \(C\) in \(m>0\) punti, sommandola a \(|C|\) si otterrà sempre un sistema lineare di maggior dimensione, irriducibile e regolare come \(|C|\). E le curve di questo sistema somma avranno con essa \(m+\nu\) intersezioni. Perciò la \(\omega\) potrà sommarsi successivamente a \(|C|\), e si otterrà una serie illimitata di sistemi di dimensione crescente se è \(\nu\geq 0\), cioè nel caso in cui la \(\omega\) sia una curva eccezionale di seconda specie. Invece se la \(\omega\) è di prima specie \((\nu=-1)\), la serie dei sistemi ottenuta si arresterà al sistema \(|C+m\omega|\) per cui la \(\omega\) diventa fondamentale.

Ora si consideri, per ogni sistema lineare \(|C|\) di grado \(n\) e genere \(\pi\), appartenente alla superficie \(F\), la differenza \(2\pi-2-n\) (numero delle intersezioni delle sue curve colle curve canoniche virtuali) che dà un \textit{carattere addittivo} \(a(C)\) dei sistemi sulla stessa \(F\). Questa differenza diminuisce almeno di una unità, ogni volta che si somma \(\omega\) a \(|C|\) o ad uno dei sistemi successivamente così ottenuti, e quindi, se la serie di questi è illimitata, finisce certo per divenire negativa, sicchè si avranno sopra \(F\) sistemi di genere \(\pi\) e di grado \(n>2\pi-2\). % [sic] addittivo -> additivo, the printed word contains a doubled t.

Ciò posto noi vogliamo supporre che \textit{la superficie \(F\) non contenga sistemi per cui la differenza \(-a(C)=n-(2\pi-2)>0\).} In tale ipotesi essa non conterrà curve eccezionali di seconda specie e potrà contenere soltanto un numero finito \(s\) di curve eccezionali di prima specie.

Per semplicità discorreremo di queste ritenendo che siano curve eccezionali irriducibili, salvo ad avvertire poi ciò che deve dirsi nel caso escluso (che risponde a curve eccezionali immagini di punti infinitamente vicini).

% scan: s076L
% !PRINTED 128
\weberpage{128}
% head: CAPITOLO QUARTO
% stamp: Dark vertical scan border along the right edge; scattered scan specks.
Si assuma che il sistema (irriducibile regolare) \(|C|\) sia, non solo virtualmente sì anche effettivamente, privo di punti base su \(F\). Se come si è detto innanzi, le \(C\) segano la curva eccezionale \(\omega\) in \(m\) punti, la \(\omega\) sarà curva fondamentale (di grado \(-1\) e perciò diminuente di \(1\) il grado delle curve residue) per il sistema \(|C+m\omega|\). E trasformando la superficie \(F\) in un’altra \(F'\) per mezzo di questo sistema, la \(\omega\) si muterà in un punto semplice \(O\) di \(F'\).

Ora ognuna delle altre \(s-1\) curve eccezionali su \(F\) sarà senza connessione colla \(\omega\) e perciò darà su \(F'\) ancora una curva eccezionale di prima specie (ovvero un punto semplice della superficie); altrimenti si otterrebbero su \(F'\) curve eccezionali di grado \(\geq 0\) (ossia di seconda specie), e quindi si costruirebbero sulla superficie \(F'\) e anche su \(F\), sistemi \(|C|\) per cui il carattere
\[
-a(C)>0.
\]

Dunque la trasformazione di \(F\) in \(F'\) riesce a diminuire il numero delle curve eccezionali appartenenti alla superficie, sicchè dopo \(s\) trasformazioni otterremo una superficie \(\Phi\) affatto priva di curve eccezionali.

A vero dire si è semplificato il nostro discorso supponendo che le curve eccezionali su \(F\) siano tutte irriducibili. Ma se vi è su \(F\) una curva eccezionale di prima specie composta di \(h\) parti (date in un certo ordine \(\omega,\,\omega' \ldotp\ldotp\ldotp\ldotp \omega^{(h-1)}\)), l’ultima di queste componenti costituisce sempre una curva eccezionale di grado \(\nu=-1\), e quando si riduce questa curva ad un punto colla trasformazione che muta \(F\) in \(F'\), la componente che la precede diventa una curva eccezionale (di prima specie) di grado \(\nu=-1\). Così la curva eccezionale composta di \(h\) parti (con cui si formano \(h\) curve eccezionali di \(F\), una irriducibile ed \(h-1\) riducibili) diventa su \(F'\) una curva eccezionale composta di \(h-1\) parti. E finalmente scompare dopo \(h\) trasformazioni, riducendosi (su \(\Phi\)) ad un gruppo di \(h\) punti semplici infinitamente vicini. % ?? uncertain: the degree letter in both occurrences -- Greek \nu or italic v.
% EDITORIAL D0074 RISOLTO: Both magnified degree letters have the same short left entry stroke and pointed right branch as the nu used on pp. 86–87, supporting the retained Greek-letter encoding.

Concludiamo pertanto: \textit{Una superficie si può generalmente trasformare in un’altra priva di curve eccezionali. Questa trasformazione viene a mancare soltanto per le superficie} (di genere \(p_g=0\), bigenere \(P=0\) ecc.) \textit{sulle quali esistono sistemi lineari \(|C|\) per cui il carattere} \(-a(C)=n-(2\pi-2)>0\), cioè, come vedremo per \textit{le superficie che appartengono alla famiglia delle rigate.}

\edsection{10. Integrità della serie segata da un sistema lineare piccolo sulle curve di un sistema grande.}{10. Integrità della serie segata da un sistema lineare piccolo sulle curve di un sistema grande.}{sec.IV.10}

L’estensione del teorema di Riemann-Roch ai sistemi lineari qualunque (anche a quelli « speciali », che sono meno ampi del
%CONT>>
% scan: s076R
% !PRINTED 129
\weberpage{129}
% head: IL GENERE NUMERICO E IL TEOREMA DI RIEMANN-ROCH, ECC.
%<<CONT
% stamp: Scattered scan specks and a short stray stroke above “modo”.
sistema canonico) si darà nel modo più semplice fondandosi sulla integrità della serie che un sistema lineare « \textit{piccolo} » sega sopra la curva generica di un sistema « \textit{normalmente grande} ». « Piccolo » e « grande » sono termini relativi che si definiscono come segue: si considera « grande » rispetto ad un sistema lineare dato \(|C|\) (virtualmente privo di punti base sopra la superficie \(F\)) un sistema lineare \(|D|\) che contenga \(|C|\) per modo che la differenza \(|D-C|\) eguagli il multiplo \(|sL|\), secondo un intero \(s\) grande quanto si vuole, di un sistema lineare irriducibile \(|L|\) senza punti base, ovvero costituisca un sistema « più ampio » di \(|sL|\). E si dirà che \(|D|\) è \textit{normalmente} grande, rispetto a \(|C|\), se la condizione precedente sia soddisfatta in ordine ad un sistema \(|L|\) \textit{privo di curve fondamentali}, di modo che, scegliendo \(s\) sufficientemente grande si può ottenere che il sistema \(|B|=|D-C|\) sia irriducibile\origfootnote{1}{s076R-1}{Ciò non è più vero se \(|L|\) e quindi \(|L_1|=|D|\) possegga una curva fondamentale \(\theta\) che non sia fondamentale per \(|C|\), giacchè allora \(\theta\) si stacca come parte fissa da \(|D-C|\). Quindi anche il lemma fondamentale non sussiste più per \(|C|\) in confronto a questo sistema (\textit{non} normalmente) grande \(|D|\): per esempio il sistema canonico \(|K|\) non sega più, in generale, sopra una \(D\) la serie completa (segata da \(|K+\theta|\)).}. % ?? uncertain: the subscript in |L_1| in the footnote -- 1 or s.
% EDITORIAL D0075 RISOLTO: The footnote subscript is a clearly formed numeral 1 with a foot, not the curved s used elsewhere in the text.

Si vuol dimostrare il \textit{Lemma fondamentale.}

\textit{Un sistema lineare completo \(|C|\), irriducibile o riducibile, sega una serie completa sopra la curva irriducibile d’un sistema lineare \(|D|\), che sia, nei suoi confronti, normalmente grande.}

Per stabilire questo lemma, si assuma sopra una \(D\) un gruppo arbitrario \(G\) della serie completa \(CD\), segata da \(|C|\), e gli si associ un gruppo \(G'\), la cui esistenza sarà dimostrata nel seguito, che soddisfa alle due proprietà seguenti:

1) \(G+G'\) sia un gruppo della serie caratteristica (generalmente non completa) segata sulla \(D\) dalle altre curve di \(|D|\);

2) \(G'\) sia un gruppo della serie segata sulla \(D\) da una curva irriducibile \(B\) di \(|D-C|\).

Allora \(G+G'\) costituirà il gruppo base d’un fascio di curve \(D\) e rispetto a questo si avrà una curva fondamentale irriducibile \(B\) passante per \(G'\); di conseguenza al fascio suddetto apparterrà una curva \(D\) spezzata in \(B\) e in una \(C\) per \(G\); ogni gruppo \(G\) della serie completa \(CD\), su \(D\), sarà dunque gruppo sezione di una \(C\).

Occorre soltanto riconoscere che, essendo \(|D|\) assai grande rispetto a \(|C|\), esiste effettivamente per ogni gruppo \(G\), della serie completa \(CD\), un gruppo \(G'\) associabile ad esso che soddisfi alle proprietà 1) e 2). A tal uopo conviene mostrare che le due serie lineari, a priori non complete, \(DD-G\) e \(BD\), segate su \(D\) rispettivamente dalle \(D\) per \(G\) e dalle \(B\), e contenute nella medesima serie
% Gathering signature printed at the foot of the page.
% EDITORIAL: segnatura originale: \vadjust{\hbox to \linewidth{\textsc{Enriques F.} - \textit{Superficie algebriche.}\hfil 9}}
\printersignature{\textsc{Enriques F.} - \textit{Superficie algebriche.}\quad 9}
%CONT>>
% scan: s077L
% !PRINTED 130
\weberpage{130}
% head: CAPITOLO QUARTO
%<<CONT
% stamp: Dark vertical scan border along the right edge; a faint bracket-like mark between “almeno” and “un” near the bottom is treated as non-text.
completa, hanno certo a comune una effettiva serie lineare, descritta dai gruppi \(G'\). E perciò basta stabilire che le deficienze di codeste due serie sono inferiori alle loro dimensioni, perchè allora, come si vede facilmente, la somma delle dimensioni delle due serie è maggiore della dimensione delle serie completa. % [sic] delle serie completa -> della serie completa, the printed article does not agree with the singular noun and adjective.

Ora quest’ultima condizione è certamente soddisfatta poichè si riconosce che, per \(s\) grande tendente all’infinito, la dimensione della serie caratteristica di \(|D|\), su \(D\), e quindi anche la dimensione della serie \(DD-G\), diventa infinita del \(2^\circ\) ordine come \(\dfrac{s^2}{2}\,n\) (\(n\) grado di \(|L|\)), e lo stesso accade per la dimensione della serie \(BD\); invece il difetto della serie caratteristica \(DD\), e quindi anche quello della serie \(DD-G\), se pure possa crescere indefinitamente con \(s\), non potrà a priori diventare infinito d’ordine superiore al primo.

Ci limiteremo a svolgere questo semplice computo nel caso in cui sia
\[
|D|=|L_s| \qquad |C|=|L|.
\]

Designando con \(n\) e \(\pi\) grado e genere di \(|C|\) e con \(n_s\) e \(\pi_s\) grado e genere di \(|L_s|\), la dimensione \(r_s\) di \(|L_s|\) viene data dal teorema di Riemann-Roch:
\[
r_s\geq p_a+n_s-\pi_s+1,
\]
dove
\[
n_s=s^2n,\qquad \pi_s=s\pi+\frac{s(s-1)}{2}n-s+1,
\]
e quindi
\[
r_s\geq \frac{s^2}{2}n-s\pi-\frac{s}{2}n+s-1;
\]
dunque la dimensione effettiva della serie caratteristica \(|DD|=(L_sL_s)\) vale % The equality sign is printed both at the end and at the beginning of consecutive lines.
\[
r_s-1\geq \frac{s^2}{2}n+\ldotp\ldotp\ldotp\ldotp
\]
tralasciando i termini d’ordine \(s\). E poichè il gruppo \(G=CL_s\) comprende \(sn\) punti; anche la dimensione della serie effettiva \(DD-G\) sarà
\[
\frac{s^2}{2}n+\ldotp\ldotp\ldotp\ldotp,
\]
tralasciando i termini d’ordine inferiore al \(2^\circ\).

D’altra parte anche la dimensione della serie segata su \(D=L_s\) dalle curve \(B\) di \(|L_{s-1}|=|L_s-C|\) sarà
\[
\frac{s^2}{2}n+\ldotp\ldotp\ldotp\ldotp
\]

Ora si può valutare il difetto di codeste serie o almeno un limite superiore di esso che dipende dal difetto della serie caratte% ?? uncertain: the bracket-like mark between “almeno” and “un” -- a non-text mark or a printed right square bracket.
%CONT>>
% scan: s077R
% !PRINTED 131
\unskip\weberpage{131}%
% head: IL GENERE NUMERICO E IL TEOREMA DI RIEMANN-ROCH, ECC.
%<<CONT
% stamp: Thin scan trace at the left edge; scattered scan specks.
\unskip ristica \(DD\) e quindi dal suo indice di specialità. Se questa serie sia completa (di dimensione \(r_s-1=p_a+n_s-\pi_s\)) il suo indice di specialità sarà \(p_a\), se invece essa abbia il difetto \(d\), il suo indice di specialità sarà
\[
i=d+p_a.
\]

In ogni caso questo indice \(i\) sarà inferiore all’ordine della serie residua, cioè al numero delle intersezioni di \(D=L_s\) con una curva canonica \(K\):
\[
DK=s(2\pi-2-n),
\]
e di conseguenza il difetto della predetta serie \(DD\), e quindi anche i difetti di \(DD-G\) e di \(BD\), non potranno diventare infiniti con \(s\), d’ordine superiore al primo. \hfill c. d. d.

Lo stesso computo si estende, con lievi modificazioni, al caso in cui, essendo \(|L|\) un sistema privo di curve fondamentali, sia ancora \(|D|=|L_s|=|sL|\) ed \(|L|=|B+C|\), ovvero in cui sia
\[
|D|=|L_s+E|,
\]
e \(|C|\) coincida con \(|L|\) o sia contenuto parzialmente in esso:
\[
|L|=|B+C|.
\]

Si perviene egualmente alla conclusione che il sistema completo \(|C|\) sega sulla curva \(D\) d’un sistema normalmente grande una serie completa.

\edsection{11. Il teorema di Riemann-Roch per i sistemi lineari qualunque.}{11. Il teorema di Riemann-Roch per i sistemi lineari qualunque.}{sec.IV.11}

Al posto del sistema lineare \(|C|\), considerato nel precedente paragrafo, si assuma il sistema canonico \(|K|\), supponendo dunque \(p_g>0\), e si costruisca un sistema lineare irriducibile \(|D|\), che sia normalmente grande rispetto a \(|K|\). Potremo affermare che \(|K|\) sega, sopra una curva generica \(D\) di \(|D|\), la serie completa residua della serie caratteristica.

Questa affermazione vale anche per \(p_g=0\), nel senso che « sopra una superficie di genere \(p_g=0\), la \textit{serie caratteristica} di un sistema \(|D|\), normalmente grande, è una serie \textit{non speciale} ». Per dimostrarlo si assuma un sistema lineare ausiliario \(|C|\) e si consideri, \textit{insieme} ad esso, il suo sistema aggiunto \(|C'|\). Se \(|D|\) è abbastanza grande rispetto a \(|C|\) e a \(|C'|\), ambedue questi sistemi segano sopra una \(D\) generica serie complete: diciamo \(g\) e \(g'\). Se la serie caratteristica \(DD\) è speciale e quindi contenuta nella serie canonica \(DD'\), dove \(|D'|=|D+C'-C|\), la serie \(g\) sarà contenuta in \(g'\), e precisamente un gruppo \(G\) di \(g\) farà parte d’un gruppo \(G'=G+G_1\) di \(g'\). Si tratta di riconoscere che, in tale ipotesi, anche \(|C'|\) contiene \(|C|\) e quindi \(p_g>0\), contro l’ipotesi.

% scan: s078L
% !PRINTED 132
\weberpage{132}
% head: CAPITOLO QUARTO
% stamp: Dark vertical scan border along the right edge; scattered scan specks.
A tal uopo basta osservare che la curva \(C\) passante per \(G\) (e, nel caso di riducibilità, ogni componente di essa) ha con la \(D\) (normalmente grande) un numero d’intersezioni maggiore di \(CC'\), e quindi fa parte della \(C'\) per \(G+G_1\).

Così abbiamo stabilito che \textit{sopra una superficie di genere \(p_g\geq 0\), la serie caratteristica d’un sistema lineare normalmente grande} (in confronto del sistema canonico) \textit{ha l’indice di specialità \(p_g\).}

Ora si consideri sulla superficie \(F\) un sistema lineare \(|C|\) comunque speciale d’indice \(i\) \((\geq 0)\) ed un sistema \(|D|\) normalmente grande rispetto ad esso. Vogliamo riconoscere che « l’indice di specialità della serie segata su una \(D\) dal sistema \(|D+C|\) è uguale all’indice di specialità di \(|C|\), cioè al numero \(i\) delle curve canoniche \(K\) linearmente indipendenti di cui fa parte una \(C\) ».

Infatti, essendo completa la serie segata da \(|K|\) su \(D\), per un gruppo \(G\) della serie \(CD\) su una \(D\) passano tante curve canoniche \(K\), linearmente indipendenti, quanto è l’indice di specialità della serie segata sulla stessa \(D\) da \(|D+C|\); ciascuna di queste \(K\) ha con \(C\) (o, nel caso di riducibilità, con una qualsiasi componente di \(C\)) un numero d’intersezioni \(CD>CK\) e quindi contiene per intero la \(C\).

Ora siamo in grado di dedurre la dimensione \(r\) del sistema \(|C|\) di grado \(n\), genere \(\pi\) e indice di specialità \(i\), che si ottiene staccando una \(D\) dal sistema \(|D+C|\); perciò occorre togliere dalla dimensione di quest’ultimo sistema che è regolare, la dimensione della serie segata su \(D\) da \(|D+C|\), di cui si conosce l’indice di specialità \(i\), ma che può non essere completa. In questa maniera si trova un limite inferiore per \(r\): precisamente si trova, come pei sistemi più ampi del sistema canonico: \(r\geq p_a+n-\pi+1\) quando sia \(i=0\), cioè per tutti i sistemi \(|C|\) \textit{non speciali}, e in generale
\[
r\geq p_a+n-\pi+1-i
\]
\textit{per i sistemi speciali d’indice \(i\).}

\textit{Questa diseguaglianza costituisce l’espressione più generale del teorema di Riemann-Roch per le superficie.}

Essendo stata riconosciuta per i sistemi lineari completi irriducibili o riducibili \(|C|\), senza punti base sopra la superficie \(F\), si deduce che essa è valida per qualsiasi sistema lineare comunque dotato di punti base assegnati: infatti un punto base \(i\)-plo che s’imponga alle \(C\) ne diminuisce il grado di \(i^2\), il genere di \(\dfrac{i(i-1)}{2}\), e la dimensione di
\[
\frac{i(i+1)}{2}=i^2-\frac{i(i-1)}{2}
\]
al più.

% scan: s078R
% !PRINTED 133
\weberpage{133}
% head: IL GENERE NUMERICO E IL TEOREMA DI RIEMANN-ROCH, ECC.
% stamp: Thin scan trace along the left edge; scattered scan specks.
I sistemi lineari \(|C|\) per cui la dimensione \(r\) riesce eguale a \(r=p_a+n-\pi+1-i\) si diranno \textit{regolari}; \textit{sovrabbondanti} se sia
\[
r>p_a+n-\pi+1-i.
\]

Il teorema di Riemann-Roch per le curve di genere \(p\) ci dà la dimensione della serie completa speciale \(g_n^r\) espressa per la dimensione della serie residua di essa rispetto alla \(g_{2p-2}^{p-1}\) canonica. Codesto teorema ha come ovvio complemento l’osservazione che la serie residua della \(g_n^r\) ha l’ordine \(2p-2-n\). Anche per i sistemi lineari speciali appartenenti ad una superficie, si può completare il teorema di RiemannRoch, esprimendo i caratteri genere e grado del sistema residuo, per mezzo di quelli del sistema dato. Infatti siano \(|C_1|\) e \(|C_2|\) due sistemi lineari residui l’uno dell’altro rispetto al sistema canonico \(|K|\). Siano \(n_1\) e \(\pi_1\) i caratteri, genere e grado, del primo sistema; \(n_2\) e \(\pi_2\) i caratteri del secondo, e si designi con \(m\) il numero delle intersezioni di \(|C_1|\) e \(|C_2|\). Introducendo il genere lineare \(p^{(1)}\) della superficie, potremo esprimere il genere e il grado di \(|K|\) per mezzo delle formule % [sic] RiemannRoch -> Riemann-Roch, the internal hyphen is absent in this occurrence.
\[
\begin{aligned}
p^{(1)}&=\pi_1+\pi_2+m-1\\
p^{(1)}-1&=n_1+n_2+2m.
\end{aligned}
\]

Ma si può calcolare \(m\) tenendo conto che \(|K|\) sega su \(C_1\) un gruppo di
\[
2\pi_1-2-n_1=n_1+m
\]
punti, sicchè
\[
m=2\pi_1-2-2n_1.
\]
Avremo pertanto
\[
\begin{aligned}
p^{(1)}&=\pi_1+\pi_2+2\pi_1-2-2n_1-1=3\pi_1+\pi_2-2n_1-3\\
p^{(1)}-1&=n_1+n_2+4\pi_1-4-4n_1=4\pi_1-3n_1+n_2-4,
\end{aligned}
\]
da cui si ricava
\[
\begin{aligned}
\pi_2&=p^{(1)}-3\pi_1+2n_1+3\\
n_2&=p^{(1)}+3n_1-4\pi_1+3.
\end{aligned}
\]

\textit{Queste formule valgono ad esprimere i caratteri del sistema \(|C_2|\) residuo del sistema speciale \(|C_1|\) rispetto al sistema canonico, per mezzo dei caratteri di \(|C_1|\).}

Si noti che le dimensioni dei due sistemi residui uno dell’altro, \(|C_1|\) e \(|C_2|\), sono date da
\[
\begin{aligned}
r_1&=p_a+n_1-\pi_1-r_2+\omega_1\\
r_2&=p_a+n_2-\pi_2-r_1+\omega_2,
\end{aligned}
\]
%CONT>>
% scan: s079L
% !PRINTED 134
\weberpage{134}
% head: CAPITOLO QUARTO
%<<CONT
% stamp: Dark vertical scan border along the right edge; scattered scan specks.
designando \(\omega_1\) e \(\omega_2\) le rispettive \textit{sovrabbondanze}. Quindi dalle formule precedenti si ricava
\[
\omega_1=\omega_2.
\]
cioè: \textit{due sistemi speciali residui l’uno dell’altro rispetto al sistema canonico hanno la medesima sovrabbondanza.}

\edsection{12. Deficienza della serie caratteristica.}{12. Deficienza della serie caratteristica.}{sec.IV.12}

Sopra la superficie \(F\) di genere numerico \(p_a\) e genere geometrico \(p_g\geq p_a\) si consideri un sistema lineare irriducibile \(|C|\) di genere \(\pi\), di grado \(n\) e dimensione \(r\), che per semplicità supponiamo non speciale. Sopra la curva generica \(C\) abbiamo due serie lineari, residue l’una dell’altra rispetto alla serie canonica; sono: la serie caratteristica \(g_n^{r-1}\), segata dalle altre curve \(C\), e la serie d’ordine \(2\pi-2-n\) e di dimensione \(p_g-1\) segata dalle curve canoniche \(K\). Se queste serie hanno rispettivamente la deficienza \(d\) e \(d'\), esse sono contenute in due serie complete di dimensione
\[
r-1+d \qquad \text{e} \qquad p_g-1+d'.
\]

Applicando il teorema di Riemann-Roch sopra la curva si ha
\[
r-1+d=n-\pi+p_g-1+d'+1
\]
cioè
\[
r=n-\pi+p_g+d'-d+1.
\]
Quindi se \(d'=0\)
\[
r=p_g-d+n-\pi+1.
\]

Ora, se \(|C|\) è un sistema lineare regolare, normalmente grande rispetto al sistema canonico, la sua dimensione vale
\[
r=p_a+n-\pi+1
\]
e d’altra parte sappiamo che le curve canoniche \(K\) segano sopra \(C\) la serie completa, sicchè si ha \(d'=0\).

Si deduce che
\[
d=p_g-p_a.
\]

In parole: \textit{la serie caratteristica di un sistema lineare completo irriducibile, normalmente grande, che sia regolare}\origfootnote{1}{s079L-1}{Si vedrà più avanti che « ogni sistema lineare normalmente grande è regolare ».}, \textit{ha la deficienza}
\[
d=p_g-p_a.
\]

% scan: s079R
% !PRINTED 135
\weberpage{135}
% head: IL GENERE NUMERICO E IL TEOREMA DI RIEMANN-ROCH, ECC.
% stamp: Scattered scan specks and stray ink-like marks.
Ciò posto vogliamo dimostrare che la serie caratteristica di un qualunque sistema lineare irriducibile \(|C|\) appartenente alla superficie \(F\) ha sempre una deficienza
\[
d\leq p_g-p_a.
\]

Per ottenere questo risultato stabiliremo il seguente

\textit{Lemma}: Se un sistema lineare
\[
|C|=|C_1+C_2|
\]
somma di due sistemi irriducibili sega la serie completa su una curva del sistema \(|C_1|\), la deficienza \(d_1\) della serie caratteristica di \(|C_1|\) è minore od uguale di quella \(d_2\) di \(|C_2|\):
\[
d_1\leq d_2.
\]

Si designino con \(n_1,\ \pi_1,\ n_2,\ \pi_2\) i caratteri di \(|C_1|\) e \(|C_2|\), con \(r_1\) ed \(r_2\) le loro rispettive dimensioni, e con \(s\) il numero delle intersezioni di una \(C_1\) con una \(C_2\); infine sia \(r\) la dimensione di \(|C|\). Per ipotesi la serie segata da \(|C|\) sopra una \(C_1\) è completa. Quindi anche le \(C\) passanti per il gruppo \(G\) degli \(s\) punti comuni ad una \(C_1\) e ad una \(C_2\) segheranno su \(C_1\) una serie completa di dimensione \(r_1-1+d_1\).

D’altra parte la dimensione della serie completa segata da \(C\) su \(C_1\) è
\[
r-r_2-1
\]
e quindi la dimensione della serie segata su \(C_1\) dalle \(C\) per \(G\) sarà
\[
r-r_2-1-(s-\varepsilon),
\]
dove \(s-\varepsilon\) designa il numero delle condizioni linearmente indipendenti che il gruppo di \(s\) punti \(G\) presenta alle \(C\) che debbono contenerlo. Così avremo
\[
r_1-1+d_1=r-r_2-1-(s-\varepsilon),
\]
cioè
\[
d_1=r-(r_1+r_2)-(s-\varepsilon).
\]

Questa eguaglianza vale anche indipendentemente dall’ipotesi che \(|C|\) seghi su \(|C_1|\) la serie completa, purchè si sostituisca a \(d_1\) la deficienza \(d'_1\) della serie caratteristica di \(|C_1|\) rispetto a quella segata dalle \(C\) per \(G\) \((d'_1\leq d_1)\); applicando lo stesso ragionamento al sistema \(|C_2|\) si trova:
\[
d_2\geq d'_2=r-(r_1+r_2)-(s-\varepsilon).
\]

Dal confronto delle espressioni di \(d_1\) e \(d_2\) risulta in ogni caso
\[
d_1\leq d_2 \qquad\qquad \text{c. d. d.}
\]

% scan: s080L
% !PRINTED 136
\weberpage{136}
% head: CAPITOLO QUARTO
% stamp: Dark vertical scan border along the right edge; scattered scan specks.
Ciò posto, essendo dato sopra la superficie \(F\) un sistema lineare qualsiasi irriducibile \(\infty^1\) almeno \(|C|\) di grado \(n\) e genere \(\pi\), si costruisca su \(F\) un altro sistema lineare irriducibile \(|D|\), che sia regolare e abbastanza grande in modo che la deficienza della sua serie caratteristica valga \(p_g-p_a\).

\ifshownotes\clearmargin\fi\edFlush
Sappiamo (\S\ 8) che sommando \(|C|\) a \(|D|\) (per essere il numero delle intersezioni \((CD)\) maggiore di \(n-\pi\)) si ottiene ancora un sistema regolare \(|C+D|\) che sega sopra \(C\) la serie completa. Perciò la deficienza \(d\) della serie caratteristica di \(|C|\) dovrà essere
\[
d\leq p_g-p_a.
\]

In parole: \textit{la serie caratteristica di un sistema lineare irriducibile sopra una superficie di generi \(p_a\) e \(p_g\geq p_a\) ha la deficienza}
\[
d\leq p_g-p_a.
\]

Così \textit{la irregolarità} \(p_g-p_a\) d’una superficie viene anche \textit{definita} come il \textit{massimo valore della deficienza della serie caratteristica} per i sistemi lineari irriducibili che ad essa appartengono.

\textit{Sopra una superficie regolare} \((p_g=p_a)\) \textit{ogni sistema lineare irriducibile completo ha la serie caratteristica completa.} O, sotto l’aspetto proiettivo: una superficie regolare normale, in un certo spazio, ha come sezioni iperpiane curve normali. Se le sezioni d’una superficie normale non sono normali, si deduce
\[
p_a<p_g.
\]

Possiamo illustrare le cose dette coll’esempio delle \textit{superficie rigate.}

Sappiamo che le rigate (a sezioni) di genere \(p\) hanno il genere geometrico \(p_g=0\) e il genere numerico \(p_a=-p\). Pertanto le rigate normali d’ordine \(n\) dovranno appartenere in generale a spazi di dimensione \(r=n-2p+1\), ed invero \textsc{C. Segre}\origfootnote{1}{s080L-1}{\textsc{C. Segre}, \textit{Courbes et surfaces réglées.} Math. Annalen, 1892.} fino dal 1890, ha riconosciuto che una rigata non speciale d’ordine \(n\) e genere \(p\), che non sia un cono, appartiene precisamente ad uno spazio \(S_{n-2p+1}\); invece una rigata \textit{speciale}, le cui sezioni piane o iperpiane siano curve speciali d’indice \(i\), è normale in uno spazio di \(n-2p+1+i\) dimensioni (escluso sempre il caso del cono in cui il sistema delle sezioni risulta sovrabbondante); così la serie caratteristica del sistema delle sezioni iperpiane ha in ogni caso la deficienza \(p\).

\textit{Osservazione.} -- Il teorema sulla deficienza della serie caratteristica dei sistemi lineari completi, ha condotto \textsc{G. Castelnuovo} a riconoscere che: \textit{ogni superficie contenente un fascio irrazionale di genere \(p\) \((>0)\) di curve è irregolare di irregolarità}
\[
p_g-p_a\geq p. % ?? uncertain: a faint high mark after the first p -- a scan speck or a printed prime.
\]

% scan: s080R
% !PRINTED 137
\weberpage{137}
% head: IL GENERE NUMERICO E IL TEOREMA DI RIEMANN-ROCH, ECC.
% stamp: Scattered scan specks.
Questo teorema rientra in un altro più generale stabilito da \textsc{F. Enriques}\origfootnote{1}{s080R-1}{\textit{Una proprietà delle serie continue di curve appartenenti ad una superficie algebrica regolare.} Rendic. Circolo Mat. Palermo, t. XIII, 1899.}: \textit{una superficie che contenga un sistema algebrico irriducibile \(\{C\}\) di curve non equivalenti è irregolare}; più precisamente \textit{se \(\{C\}\) contiene \(\infty^d\) \((d>0)\) curve non equivalenti l’irregolarità della superficie vale}
\[
p_g-p_a\geq d.
\]

\textsc{Enriques} ha riconosciuto la deficienza della serie (canonica) segata sopra una \(C\) dal sistema aggiunto \(|C'|\). In quella occasione \textsc{Castelnuovo} --- cui egli aveva comunicato il risultato --- ha osservato la deficienza della serie caratteristica del sistema lineare completo a cui appartiene una \(C\). La dimostrazione assume la forma più semplice nell’esposizione di \textsc{F. Severi}\origfootnote{2}{s080R-2}{\textit{Osservazioni sui sistemi continui di curve appartenenti ad una superficie algebrica.} Atti R. Accad. Scienze di Torino, 1904.}. Invero basta notare che il sistema continuo completo \(\{C\}\) sarà formato da \(\infty^d\) sistemi lineari \(|C|\) aventi, in generale, una certa dimensione \(r\): quindi la serie caratteristica di \(|C|\) sopra una \(C\) generica avrà la dimensione \(r-1\) e sarà contenuta nella serie caratteristica del sistema continuo di dimensione \(r-1+d\), cioè nella serie segata su \(C\) dalle curve infinitamente vicine di \(\{C\}\).

Si avverta esplicitamente che il teorema relativo al fascio irrazionale rientra in quello così stabilito: basta notare che, sottraendo da un sistema lineare abbastanza ampio i gruppi di \(p\) curve d’un fascio di genere \(p\), si ottiene un sistema continuo formato da \(\infty^p\) sistemi lineari disequivalenti.

\edsection{13. Nota storica.}{13. Nota storica.}{sec.IV.13}

La considerazione del genere numerico d’una superficie si è presentata fino dai primi studi di \textsc{A. Cayley}, \textsc{H. G. Zeuthen} e \textsc{M. Noether}. Dopo che \textsc{Noether} ebbe definito il genere (geometrico) \(p\) d’una superficie \(F_n\) d’ordine \(n\), come numero delle superficie aggiunte \(\Phi_{n-4}\) linearmente indipendenti, era naturale di cercare una espressione di \(p\) per mezzo dei caratteri della curva doppia di \(F_n\); e all’uopo soccorrevano le formule di postulazione di \textsc{Cayley}. Questo stesso matematico\origfootnote{3}{s080R-3}{\textsc{A. Cayley}, \textit{On the deficiency of certain Surfaces.} Math. Annalen, Bd. III (1871).} ebbe a notare che, nel caso delle rigate di genere \(p\), il calcolo accennato conduce al valore negativo \(-p\), mentre il genere geometricamente definito non può, per sua natura, discendere al disotto di zero e, in questo caso, è precisamente nullo.

% scan: s081L
% !PRINTED 138
\weberpage{138}
% head: CAPITOLO QUARTO
% stamp: Dark vertical scan-edge band in the right margin; scattered specks.
Frattanto \textsc{Zeuthen}, da uno studio sulle corrispondenze fra superficie algebriche, era condotto a dimostrare (sia pure con qualche restrizione in ordine alle singolarità) che « l’espressione aritmetica (virtuale) del genere, secondo \textsc{Cayley}, rimane invariata per trasformazioni birazionali della superficie, e ciò indipendentemente dal suo (presunto) significato geometrico e dalla dimostrazione dell’invarianza del sistema canonico fornita da \textsc{Noether} ».

Questi risultati venivano interpretati da \textsc{Noether}\origfootnote{1}{s081L-1}{\textit{Mathematische Annalen.} Bd. VIII (1875).} nel senso che il genere d’una superficie $F_n$ d’ordine $n$ può essere definito in due modi: per via \textit{geometrica} come numero effettivo delle superficie aggiunte $\Phi_{n-4}$ linearmente indipendenti, e per via \textit{numerica} mediante le formule di postulazione; s’introducono così il \textit{genere geometrico} e il \textit{genere numerico} della superficie, ambedue a priori invarianti per trasformazioni birazionali, ma che si presumono coincidenti fra loro, salvochè il genere numerico può diventare negativo (per le rigate) e perciò sembra assorbire l’altro in un carattere di significato più generale.

Senonchè la presunzione di \textsc{Noether}, che debba essere, per $p_a$ non negativo, $p_a=p_g$, è venuta a cadere in seguito ai nuovi esempi di superficie non regolari aggiuntisi all’esempio delle rigate. Come si è detto il primo esempio di queste è stato recato da G. \textsc{Castelnuovo} in una Nota dell’Istituto lombardo del 1891. Altri esempi sono emersi dalle ricerche di E. \textsc{Picard} e di G. \textsc{Humbert} sulle superficie possedenti integrali di differenziali totali di prima specie, e in particolare sulle superficie iperellittiche, per cui $p_g=1$ e $p_a=-1$.

Successivi studi hanno messo in luce larghe classi di superficie irregolari, come si è indicato nell’Osservazione del § 7.

Frattanto le ricerche iniziate nel 1893 e proseguite nel 1896 (nella citata « Introduzione ») hanno condotto F. \textsc{Enriques} a riconoscere il significato geometrico della differenza $p_g-p_a$ in rapporto alla deficienza della serie (canonica) che il sistema $|C'|$ aggiunto ad un sistema lineare irriducibile $|C|$ sega sulla curva generica di $|C|$: donde risulta la dimostrazione affatto generale dell’invarianza del $p_a$, essendo già stabilita l’invarianza del $p_g$. Per \textsc{Enriques} l’irregolarità $p_g-p_a$ è il valore massimo della deficienza indicata, restando a priori possibile che esistano sistemi $|C|$ per cui la detta deficienza sia $<p_g-p_a$. In questo punto E. \textsc{Picard} ha completato la teoria, dimostrando (per via trascendente) che il sistema $|C'|$ aggiunto ad un sistema irriducibile semplice $|C|$ ($\infty^3$ almeno) è sempre regolare, e perciò $p_g-p_a$ si definisce semplicemente come deficienza della serie segata sulla sezione piana d’una
%CONT>>
% scan: s081R
% !PRINTED 139
\weberpage{139}
% head: IL GENERE NUMERICO E IL TEOREMA DI RIEMANN-ROCH, ECC.
%<<CONT
% stamp: Dark scan-edge marks at the upper right and along the right margin; scattered specks.
superficie $F_n$ dalle superficie aggiunte $\Phi_{n-3}$\origfootnote{1}{s081R-1}{E. \textsc{Picard}, \textit{Sur quelques questions se rattachant à la connexion linéaire dans la théorie des fonctions algébriques de deux variables indépendantes.} v. Crelle, Bd. 129, 1905. Cfr. \textsc{Picard} et \textsc{Simart}, \textit{Traité}, t. II, pag. 437. Una dimostrazione algebrico-geometrica (ed anche un’estensione delle condizioni di sussistenza) del teorema, è stata data dal \textsc{Severi}, nella Nota \textit{Sulla regolarità del sistema aggiunto} (Lincei, 8 nov. 1908), fondandosi però sul teorema della completezza della serie caratteristica di un sistema continuo}. In ordine a questo teorema che per $p_g=p_a$ comprende il teorema di Castelnuovo esposto nel § 6, vedansi gli sviluppi del seguente paragrafo.

L’estensione alle superficie del teorema di Riemann-Roch appare anzitutto accennata da M. \textsc{Noether} nel 1886\origfootnote{2}{s081R-2}{\textit{Extension du théorème de Riemann-Roch aux surfaces algébriques.} Comptes Rendus de l’Académie des Sciences, Paris, t. 103.}. L’autore considera un sistema lineare completo irriducibile $|C|$ di dimensione $r$, di grado $n$ e genere $\pi$, avente un certo indice di specialità $i\geq 0$, sopra una superficie di genere (geometrico) $p$. La serie caratteristica di $|C|$, segata da $|C|$ su una $C$ generica, è una $g_n^{r-1}$ residua della serie $g_{2\pi-2-n}^{p-i}$ segnata dal sistema canonico $|K|$; e \textsc{Noether}, senza dichiararlo esplicitamente, postula che le due serie (segate da sistemi completi) sieno complete; in tale ipotesi il teorema di Riemann-Roch relativo alla curva $C$ fornisce l’eguaglianza
\[
r=p+n-\pi+1-i.
\]

La questione viene ripresa da \textsc{Enriques} fino dalle « Ricerche » del 1893 nelle quali riesce a dimostrare che per le superficie \textit{regolari} di genere $p>1$, la serie caratteristica di un sistema completo è sempre completa, se tale sia la serie caratteristica del sistema canonico; in tale ipotesi (che risulterà poi portare una restrizione superflua) resta così dimostrata la diseguaglianza
\[
r\geq p+n-\pi+1,
\]
e l’autore avverte che essa non può convertirsi in una eguaglianza, cioè che esistono certo sistemi \textit{sovrabbondanti}, bastando all’uopo la presenza di curve fondamentali di genere $\varrho>0$. % ?? uncertain: the genus symbol before > 0 -- \varrho (transcribed) or \rho.
% EDITORIAL D0078 RISOLTO: The magnified glyph has the loop and visibly curled lower terminal of varrho, matching the genus symbol elsewhere in the supplied pages.

Nell’ « Introduzione » del 1896, colla nuova definizione del genere numerico, \textsc{Enriques} dava per tutte le superficie regolari e non regolari l’estensione del teorema di Riemann-Roch ai sistemi aggiunti, che vuol dire — sostanzialmente — per i sistemi più ampi del sistema canonico.

Per estendere il teorema ai sistemi lineari qualunque pareva necessario di ricercare il massimo della deficienza della serie caratteristica e dimostrare che esso è dato da $p_g-p_a$. Questo risultato
%CONT>>
% scan: s082L
% !PRINTED 140
\weberpage{140}
% head: CAPITOLO QUARTO
%<<CONT
% stamp: Dark vertical scan-edge band in the right margin; scattered specks.
appunto è stato conseguito da G. \textsc{Castelnuovo} in due memorie\origfootnote{1}{s082L-1}{\textit{Alcuni risultati sui sistemi lineari di curve appartenenti ad una superficie algebrica.} Memorie della Società Italiana delle Scienze, detta dei XL (1896, vol. X). \textit{Alcune proprietà fondamentali dei sistemi lineari di curve tracciati sopra una superficie algebrica.} Annali di Matematica, serie II, vol. XXV. Riprodotte in “ Memorie scelte „, XXII e XXIII.} degli anni 1896 e 1897.

Il procedimento usato a tal uopo si basa sopra una conveniente estensione dell’idea che conduce a trovare le più semplici condizioni di regolarità di una superficie, siccome abbiamo esposto nel § 6. Si tratta di metodi che sembrano ancora suscettibili di applicazione feconda, ma riescono alquanto laboriosi.

Una semplificazione notevole è stata recata in quest’ordine di questioni da F. \textsc{Severi}, in due Note del 1903 e del 1905\origfootnote{2}{s082L-2}{\textit{Sulla serie caratteristica di un sistema lineare di curve appartenente ad una superficie algebrica.} Rendiconti R. Accademia dei Lincei, serie V, vol. XII (1903). \textit{Sul teorema di Riemann-Roch e sulle serie continue di curve appartenenti ad una superficie algebrica.} Atti Accademia delle Scienze di Torino, vol. XL (1905).}; particolarmente nella seconda Nota l’autore riesce ad invertire l’ordine delle deduzioni, giustificando prima il teorema di Riemann-Roch e traendone poi il massimo della deficienza della serie caratteristica dei sistemi irriducibili.

Il punto di partenza di questi sviluppi è il noto teorema di Noether dell’$Af+B\varphi$ esteso alle superficie nella forma seguente:

Ogni superficie $f_m(xyz)=0$, d’ordine $m$, non contenente il piano $z=0$, che passi per i punti del gruppo $G$ comune al piano $z=0$ e alla curva $D_s$ d’ordine $ns$, intersezione completa della superficie $f_n$ d’ordine $n$, priva di punti multipli propri, con una superficie $f_s$ d’ordine $s$, ha un’equazione della forma
\[
f_m=Af_n+Bf_s+zf_{m-1}=0,
\]
essendo $f_{m-1}$ un polinomio d’ordine $m-1$ ed $A$ e $B$ polinomi in $x$, $y$, di grado $m-n$ e $m-s$ rispettivamente.

Assumendo addirittura che la superficie $f_n$ sia dotata di singolarità normali, da questo teorema segue che le superficie $f_m$ d’ordine qualunque, passanti per i punti doppi di una $D_s$, segano su $D_s$ la serie completa: invero si dimostra che la proposizione è vera per $m-1$, se è vera per $m$.

Da ciò \textsc{Severi} trae il

\textit{Lemma.} – Per $s>n-4$ le superficie $\Phi_{n-4}$ aggiunte alla $f_n$ segano sopra la curva $D_s$ (intersezione completa di $f_n$ e di $f_s$) una serie completa.

Infatti le superficie d’ordine $n-4$ passanti per i punti doppi
%CONT>>
% scan: s082R
% !PRINTED 141
\weberpage{141}
% head: IL GENERE NUMERICO E IL TEOREMA DI RIEMANN-ROCH, ECC.
%<<CONT
% stamp: Faint scan-edge marks in the right margin; scattered specks.
di $D_s$ vengono a contenere per intero la curva doppia di $f_n$ e quindi risultano aggiunte ad essa.

Come conseguenza del precedente lemma si ha che « le curve canoniche $K$, e poi anche le curve d’un sistema (speciale) completo $|C|$ contenuto in $|K|$, segano sulle curve $D_s$ la serie completa. E perciò, staccando $D_s$ dal sistema completo $|D_s+C|$, si ottiene per la dimensione di $|C|$ la diseguaglianza che esprime il teorema di Riemann-Roch. % [sic] opening « only -> closing » apparently omitted, the quotation is left open in print.

D’altra parte lo stesso lemma enunciato innanzi, quando si sia verificato che il sistema completo $|D_s|=|sD|$ (privo di curve fondamentali proprie) è regolare (cfr. il § 15), porta che la serie caratteristica di una $D_s$ ha la deficienza $p_g-p_a$.

Da ciò si deduce (come qui è fatto nel paragrafo precedente) che la serie caratteristica d’un qualsiasi sistema lineare irriducibile $|C|$, $\infty^1$ almeno, ha una deficienza $\leq p_g-p_a$.

Il ragionamento che conduce a tale risultato è quello stesso che \textsc{Enriques} aveva adoperato nelle « Ricerche » per passare dalla serie caratteristica del sistema canonico a quella d’un altro sistema lineare qualunque, sopra una superficie regolare, e che \textsc{Castelnuovo} aveva ripreso nella forma più generale, come un mezzo affatto secondario, per giustificare il teorema sulla deficienza della serie caratteristica nei casi particolari che sfuggivano alla sua dimostrazione diretta (sistemi non semplici).

La nuova dimostrazione del teorema di Riemann-Roch per i sistemi lineari qualunque, quale è offerta in questo trattato\origfootnote{1}{s082R-1}{Cfr. F. \textsc{Enriques}, \textit{Sur l’extension du théorème de Riemann-Roch aux systèmes linéaires de courbes appartenant à une surface algébrique.} Bull. des Sciences Mathématiques, 1940 (Cfr. in particolare il n. 5).}, è almeno tanto semplice nel concetto quanto quella del \textsc{Severi}. Essa ha, ai nostri occhi, il vantaggio di indicare la ragione profonda del fatto che le curve canoniche $K$ (o le curve d’un sistema speciale) segano sulle $D_s$ una serie completa: che non sta nell’essere le $D_s$ intersezioni complete della superficie $F_n$ data con una $F_s$, bensì nell’essere $|K|$ assai piccolo rispetto al sistema $|D'_s|$. Inoltre il nostro procedimento porge la dimostrazione più diretta del teorema che si ha in vista, rimanendo nello stesso ordine d’idee in cui si tratta il caso preliminare dei sistemi aggiunti o dei sistemi più ampi del sistema canonico, senza che faccia bisogno di ricorrere a qualche motivo nuovo, come sembrava necessario seguendo le vie di \textsc{Castelnuovo} e di \textsc{Severi}.

Ed ancora, poichè il teorema di Riemann-Roch per un sistema lineare $|C|$ viene ottenuto dalla considerazione della serie che un sistema $|L|=|C+D|$ sega sopra una curva $D$, che si distaccherà
%CONT>>
% scan: s083L
% !PRINTED 142
\weberpage{142}
% head: CAPITOLO QUARTO
%<<CONT
% stamp: Dark vertical scan-edge band in the right margin; scattered specks and marginal marks.
da $|L|$, esso resta giustificato \textit{per il sistema} $|C|=|L-D|$ comunque questo sia \textit{riducibile}.

Ora la considerazione relativa al distacco di una $D$ da $|L|$ si incontra già nelle « Ricerche » di \textsc{Enriques} (del 1893)\origfootnote{1}{s083L-1}{L. c., IV, 4.}, e perciò si può ritenere che \textsc{Enriques} stesso, e \textsc{Castelnuovo} con lui, hanno conosciuto assai presto, almeno implicitamente, la possibilità di estendere il teorema di Riemann-Roch ai sistemi riducibili. Secondo la loro mentalità induttiva essi dovevano rendere esplicita la conoscenza in occasione di casi concreti. L’estensione ai sistemi riducibili viene formulata in una loro memoria comune del 1900\origfootnote{2}{s083L-2}{\textsc{Castelnuovo-Enriques}, \textit{Sopra alcune questioni fondamentali nella teoria delle superficie algebriche.} Annali di Mat., t. VI, serie III, n. 4.}.

Il teorema di Riemann-Roch porge anche notevoli criteri per l’esistenza effettiva delle \textit{curve virtuali}, siccome diremo nel seguente paragrafo.

\edsection{14. Curve virtuali.}{14. Curve virtuali.}{sec.IV.14}

Teniamo presente che il calcolo della dimensione d’un sistema lineare $|C|=|A-B|$, residuo di una $B$ rispetto ad un sistema regolare e non speciale $|A|$, si è ottenuto valutando la dimensione della serie segata da $|A|$ su questa curva $B$; appare di qui che il teorema di Riemann-Roch, esprimente il risultato di codesto calcolo, deve porgere anche un criterio per l’esistenza effettiva delle curve virtuali di $|A-B|$. Ciò è anzitutto evidente nel caso in cui $|A|$ seghi sopra una $B$ una serie non speciale.

La formula che dà la dimensione del sistema aggiunto ad un altro di genere $\pi$:
\[
r\geq p_a+\pi-1
\]
è una prima applicazione di questo criterio, che si è presentata già nella « Introduzione » di \textsc{Enriques} del 1896. Una seconda applicazione è la formula che dà la dimensione del sistema $i$-canonico sopra una superficie di genere lineare $p^{(1)}$:
\[
P_i-1\geq p_a+\frac{i(i-1)}{2}(p^{(1)}-1)
\qquad (i>1)
\]
quale s’incontra in una Nota di \textsc{Castelnuovo} del 1897\origfootnote{3}{s083L-3}{\textit{Sul genere lineare di una superficie ecc.} Rend. R. Acc. dei Lincei, 1897; riprodotta in « Memorie scelte », XXIV.} e nella memoria di \textsc{Castelnuovo-Enriques} del 1900\origfootnote{4}{s083L-4}{\textit{Sopra alcune questioni fondamentali della teoria delle superficie.} Ann. di Matematica, 1900.}: infatti il sistema
%CONT>>
% scan: s083R
% !PRINTED 143
\weberpage{143}
% head: IL GENERE NUMERICO E IL TEOREMA DI RIEMANN-ROCH, ECC.
%<<CONT
% stamp: Dark scan-edge marks at the upper right and lower right; scattered specks.
$i$-canonico $|iK|$ si può dedurre da un sistema regolare $|A|=|B+iK|$, $i$-mo aggiunto d’un sistema regolare $|B|$, staccandone $|B|$, e $B+iK$ sega appunto sopra una $B$ una serie non speciale. % The two equals signs reproduce repetition at the printed line break.

Una terza applicazione del teorema di Riemann-Roch come criterio d’esistenza effettiva delle curve virtuali è data da \textsc{Severi} (1905) e contempla in generale il caso che $|A|$ seghi sopra una curva $B$ una serie comunque speciale: siccome diremo più avanti.

Volendo esaminare la questione d’esistenza delle curve d’un sistema differenza $|C|=|A-B|$, è lecito supporre che i due sistemi $|A|$ e $|B|$, che definiscono le curve virtuali $C$, sieno ambedue \textit{irriducibili, regolari e privi di curve fondamentali}, nonchè di punti base sopra una superficie $F$, priva essa stessa di singolarità; imperocchè ai detti sistemi $|A|$ e $|B|$ si può sempre sommare un altro sistema regolare e privo di curve fondamentali.

Ciò posto osserviamo che:

1) L’esistenza effettiva delle curve virtuali di $|C|=|A-B|$, esige anzitutto che il grado del sistema diminuendo sia maggiore od eguale a quello del diminutore
\[
A^2\geq B^2.
\]

Infatti se $A^2<B^2$, le curve $C$ debbono avere un numero negativo d’intersezioni colle curve del sistema $|A+B|$, essendo $(A+B)(A-B)=A^2-B^2$, e perciò dovrebbero essere componenti fisse di questo sistema che per le ipotesi ammesse, è irriducibile.

2) In secondo luogo perchè esistano curve di $|A-B|$ è necessario che il numero delle intersezioni delle curve $A$ e $B$ sia compreso fra i gradi dei due sistemi, $|A|$ e $|B|$:
\[
B^2\leq AB\leq A^2.
\]

Se così non fosse le curve $C$ avrebbero un numero negativo d’intersezioni colle curve $B$ o colle curve $A$, avendosi rispettivamente: per $AB<B^2$
\[
B(A-B)<0
\]
e per $AB>A^2$
\[
A(A-B)<0.
\]

Le condizioni 1) e 2) sono \textit{necessarie} sia per l’esistenza effettiva di $|C|=|A-B|$, sia per quella d’un multiplo $|i(A-B)|$. Anzi se
\[
A^2>B^2
\]
bisogna che sia
\[
B^2<AB<A^2,
\]
altrimenti le $C$ (o le $iC$) — che non sono in senso assoluto curve d’ordine zero — avrebbero zero intersezioni colle $B$ o colle $A$ e
%CONT>>
% scan: s084L
% !PRINTED 144
\weberpage{144}
% head: CAPITOLO QUARTO
%<<CONT
% stamp: Dark vertical scan-edge band in the right margin; scattered specks.
perciò sarebbero curve fondamentali per il relativo sistema (che abbiamo supposto non possedere curve siffatte).

Alle due condizioni necessarie sopra enunciate, fa riscontro la \textit{condizione sufficiente} per l’esistenza di un multiplo di $|C|$:

\textit{Se il grado virtuale di $|C|=|A-B|$ è positivo ed $A^2>AB>B^2$, le curve del sistema multiplo $|iC|$, per $i$ abbastanza grande, hanno sempre esistenza effettiva.} % The two consecutive > signs reproduce repetition at the printed line break.

Infatti, si consideri la serie segata dal sistema $|iA-(i-1)B|$ sopra una curva $B$ (supponendo che tale sistema esista effettivamente); questa serie è d’ordine
\[
AB+(i-1)(AB-B^2),
\]
che, per $i$ abbastanza grande. viene a superare l’ordine dei gruppi canonici di $B$; perciò la detta serie è non speciale, e si potrà calcolare la dimensione del sistema residuo % ?? uncertain: the mark after "grande" -- printed full stop (retained) or scan speck.
\[
|iA-(i-1)B-B|=|iA-iB|
\]
mediante la formula
\[
r_i\geq p_a+n_i-\pi_i+1,
\]
dove
\[
\begin{aligned}
n_i&=i^2n\\
\pi_i&=i\pi+\frac{i(i-1)}{2}n-i+1 % ?? uncertain: the faint mark after the final 1 -- scan speck (not transcribed) or full stop.
\end{aligned}
\]
($n=(A-B)^2=$ grado virtuale di $|C|$; $\pi=$ genere virtuale di $|C|$).

Dalla formula precedente risulta $r_i>0$ per valori sufficientemente alti di $i$ e quindi l’esistenza effettiva delle curve di $|iC|$.

Tuttavia nel discorso precedente si è supposta a priori l’esistenza del sistema $|iA-(i-1)B|$; questo supposto si giustifica per i valori di $i$ che rendono $r_i>0$, giacchè in tali condizioni si possono calcolare successivamente le dimensioni virtuali di $|iA-B|$, $|iA-B-B|=|iA-2B|$ ...., fino a $|iA-(i-1)B|$, che risultano tutte a fortiori maggiori di $r_i$: invero per i detti valori di $i$, la dimensione della serie segata da $|iA-(i-2)B|$ sopra $B$ sarà a fortiori superiore a quella della serie canonica, e quindi non speciale, ecc. % The four dots before the comma following |iA-2B| are printed.

Il criterio d’esistenza delle curve virtuali che così abbiamo dimostrato si applica al sistema pluricanonico di una superficie per cui il genere lineare sia $p^{(1)}>1$, e quindi il grado virtuale del sistema canonico
\[
p^{(2)}=p^{(1)}-1>0.
\]

Si trova così la formula già scritta innanzi
\[
P_i\geq p_a+\frac{i(i-1)}{2}(p^{(1)}-1)+1.
\]

% scan: s084R
% !PRINTED 145
\weberpage{145}
% head: IL GENERE NUMERICO E IL TEOREMA DI RIEMANN-ROCH, ECC.
% stamp: Thin dark scan-edge line at the left; scattered specks.
Si noti che il criterio non conduce ad alcuna risposta decisiva intorno all’esistenza del multiplo $|iC|$ di un sistema virtuale di grado $n=0$, e così non permette di affermare a priori l’esistenza di curve pluricanoniche per $p^{(1)}=1$; sebbene questa esistenza venga poi dimostrata mediante altre considerazioni per le superficie prive di curve eccezionali (non appartenenti alla famiglia delle rigate).

\edFlush
Il caso in cui il grado di $|C|$ sia $n<0$ resta a fortiori indeciso; non si può dire se i due sistemi $|A|$ e $|B|$, o i loro multipli $|iA|$ ed $|iB|$, per $i$ sufficientemente grande, diano luogo o meno ad una differenza effettiva.

A complemento di questa discussione vogliamo ancora considerare il caso in cui i sistemi $|A|$ e $|B|$ abbiano lo stesso grado:
\[
A^2=B^2.
\]

In questo caso la differenza $|A-B|$ o un suo multiplo, non può avere esistenza effettiva se il numero delle intersezioni delle $A$ e $B$ è diverso da $A^2$; si ha dunque la condizione necessaria d’esistenza
\[
AB=A^2=B^2,
\]
e il grado virtuale di $|C|=|A-B|$ vale
\[
n=(A-B)^2=0.
\]

Ora la curva $C$ non può avere esistenza effettiva se non come curva d’ordine zero, quando sia $A\equiv B$. Invece la curva di $|iC|$ avrà esistenza effettiva, ancora come curva d’ordine zero, quando sia
\[
iA\equiv iB.
\]

Qui vi è luogo ad osservare che qualora non sia abbia, per qualche valore di $i$, $iA\equiv iB$, si potrà costruire sulla superficie (che possiamo supporre avere per sezioni piane o iperpiane le curve $B$) una serie illimitata di sistemi $|iA-(i-1)B|$, tutti formati di curve dello stesso ordine, e per conseguenza contenuti in una o più serie continue di curve non equivalenti, ciò che porta che la superficie sia irregolare: % [sic] non sia abbia -> non si abbia, an extra a is printed in sia.
\[
p_g>p_a.
\]

Si può dunque affermare che:

\textit{Sopra una superficie regolare due sistemi lineari di curve $|A|$ e $|B|$ dello stesso grado, che s’intersechino in}
\[
AB=A^2=B^2
\]
\textit{punti, sono parti aliquote di uno stesso sistema multiplo:}
\[
|iA|=|iB|,
\]
% Printed signature line at the foot of the page.
% EDITORIAL: segnatura originale: {\footnotesize \textsc{Enriques F.} - \textit{Superficie algebriche.}\hfill 10}
\printersignature{\textsc{Enriques F.} - \textit{Superficie algebriche.}\quad 10}
%CONT>>
% scan: s085L
% !PRINTED 146
\weberpage{146}
% head: CAPITOLO QUARTO
%<<CONT
% stamp: Dark vertical scan-edge band in the right margin; scattered specks.
sicchè il multiplo secondo $i$ della loro differenza è la curva effettiva d’ordine zero: $iA-iA$.

Se si prescinde dalla regolarità della superficie, si può affermare in generale che: « se due curve $A$ e $B$ dello stesso grado $n$ si segano in $n$ punti, due multipli di esse appartengono ad una medesima serie continua di curve, o come si dice, sono \textit{algebricamente equivalenti} »\origfootnote{1}{s085L-1}{Cfr. F. \textsc{Severi}, Math. Annalen, 62 (1906), pag. 194. A. \textsc{Todd}, Proc. of the Cambridge Phil. Soc. 23 marzo, 1939. Per le relazioni con un teorema sulla base di W. V. D. \textsc{Hodge}, J. London Math. Soc. 12, 1937, cfr. B. \textsc{Segre}, Ann. di Mat., 1937 e J. \textsc{Bronowski}, J. London Math. Soc., 1938.}.

Il criterio d’esistenza effettiva di $|C|=|A-B|$, che abbiamo detto essere stato indicato da \textsc{Severi}\origfootnote{2}{s085L-2}{\textit{Sulle curve algebriche virtuali appartenenti ad una superficie algebrica.} Rendic. Ist. Lombardo II, 38 (1905).}, consiste nella semplice \textit{estensione formale del teorema di Riemann-Roch alle curve virtuali}, quale si ha definendo l’\textit{indice di specialità} $i$ di $|C|$, come « numero delle curve linearmente indipendenti dal sistema residuo di $|C|$ stesso rispetto al sistema canonico (impuro) $|K|$ », cioè come dimensione del sistema
\[
|K-C|=|B+K-A|
\]
aumentata di una unità.

Designando con $n$ e $\pi$ i caratteri virtuali di $|C|$, la diseguaglianza
\[
r\geq p_a+n-\pi+1-i,
\]
porge un limite inferiore per la dimensione $r$ di $|C|$ e assicura quindi dell’esistenza effettiva del sistema quando sia
\[
p_a+n-\pi+1-i\geq 0.
\]

Per giustificare ciò che si è detto nel nostro ordine d’idee (e limitandosi, per semplicità di discorso all’ipotesi $i>0$) basta tener presente il lemma che ci ha condotto al teorema di Riemann-Roch pei sistemi speciali, osservando che il sistema $|B+K-A|$ può ritenersi come un sistema piccolo rispetto al sistema normalmente grande $|sB|$ per $s$ assai elevato; perciò $|B+K-A|$ segherà sulla curva del detto sistema una serie completa di dimensione $i-1$, e quindi l’indice di specialità della serie segata sulla curva stessa da $|A+(s-1)B|$ avrà precisamente l’indice di specialità $i$: di qui si deduce la formula scritta innanzi che dà un limite inferiore della dimensione $r$ di % [sic] e quindi l’indice di specialità della serie -> e quindi la serie, the printed subject repeats the index that the predicate assigns.
\[
|A+(s-1)B-sB|=|A-B|.
\]

% scan: s085R
% !PRINTED 147
\weberpage{147}
% head: IL GENERE NUMERICO E IL TEOREMA DI RIEMANN-ROCH, ECC.
% stamp: Faint vertical scan-edge marks at the right; scattered specks.
Qui conviene aggiungere l’osservazione che l’indice di specialità $i$ di una curva virtuale $C$ può riuscire grande quanto si vuole, a differenza dell’indice di una curva effettiva che è sempre
\[
i\leq p_g.
\]

Infatti il residuo di $C$ rispetto a $|K|$,
\[
|B+K-A|,
\]
non è detto a priori che sia contenuto in $|K|$: anzi l’esservi contenuto esprime, in altro modo, la condizione necessaria e sufficiente perchè una curva virtuale speciale abbia esistenza effettiva.

\edsection{15. Sistemi lineari regolari e sovrabbondanti.}{15. Sistemi lineari regolari e sovrabbondanti.}{sec.IV.15}

Sopra una superficie di genere aritmetico $p_a$ si abbia un sistema lineare irriducibile completo, di genere $\pi$, grado $n$ e indice di specialità $i$: la sua dimensione soddisfa alla diseguaglianza che costituisce il teorema di Riemann-Roch
\[
r\geq p_a+n-\pi+1-i,
\]
ovvero
\[
r=p_a+n-\pi+1-i+\omega
\]
con
\[
\omega\geq 0;
\]
il sistema si è detto \textit{regolare} se $\omega=0$, e \textit{sovrabbondante} se $\omega>0$; in quest’ultimo caso $\omega$ ne designa la \textit{sovrabbondanza}.

Vi è luogo a ricercare se esistano dei criteri che permettano di valutare la sovrabbondanza d’un sistema $|C|$, o almeno di decidere se $|C|$ sia o meno regolare.

Il risultato più importante, in quest’ordine d’idee, concerne la regolarità del sistema aggiunto ad un sistema lineare irriducibile $|C|$, che noi abbiamo dimostrata per un multiplo abbastanza grande di $|C|$ ($\infty^3$ o $\infty^2$ almeno) e che si è detto essere stata riconosciuta da Picard (§ 13) per il sistema aggiunto a $|C|$ stesso.

Qui possiamo stabilire molto semplicemente il \textit{teorema di Picard} ove si aggiunga tuttavia una restrizione superflua:

\textit{Il sistema $|C'|$, aggiunto ad un sistema lineare irriducibile $|C|$ di dimensione $r\geq p_g-p_a$ è certo regolare} (di dimensione $p_a+\pi-1$).\origfootnote{1}{s085R-1}{Cfr. \textsc{Enriques}, Bull. de Sc. Math., l. c. 1940.}

La dimostrazione riposa sopra un

\textit{Lemma.} – Siano $|C|$ e $|D|$ due sistemi lineari irriducibili, diciamo per semplicità di dimensione maggiore d’uno. Si può ricono%
%CONT>>
% scan: s086L
% !PRINTED 148
\unskip\weberpage{148}%
% head: CAPITOLO QUARTO
%<<CONT
% stamp: Dark vertical scan-edge band at the right margin.
scere che « \(\lvert C\rvert\) sega sulla curva del sistema somma \(\lvert C+D\rvert\) una serie completa se la dimensione di \(\lvert D\rvert\) superi \(p_g-p_a\) ». Infatti si consideri sulla curva \(\bar{L}\) di \(\lvert L\rvert=\lvert C+D\rvert\) un gruppo \(G\) della serie completa contenente i gruppi intersezione delle \(C\). Siccome la serie caratteristica di \(\lvert L\rvert\) ha la deficienza \(\delta\leq p_g-p_a\), fra gli \(\infty^r\) gruppi \(LD\) ve ne saranno \(\infty^{r-(p_g-p_a)}\) che associati a \(G\) costituiscono gruppi della serie caratteristica di \(\lvert L\rvert\): sia \(G'\) uno di questi gruppi. Ora per \(G+G'\) passa un fascio di curve \(L\), e, rispetto a questo, la curva \(D\) per \(G'\) è fondamentale: staccandola resta una curva di \(\lvert C\rvert\) che passa per \(G\). Così dunque la serie completa dei gruppi \(G\) viene segata su \(\bar{L}\) dal sistema \(\lvert C\rvert\).

Ciò posto, si assuma sopra una superficie un sistema lineare irriducibile \(\lvert C\rvert\), di un certo genere \(\pi\) e di dimensione maggiore di \(p_g-p_a\); vogliamo provare che il suo aggiunto \(\lvert C+K\rvert\) (\(K\) designando il sistema canonico) è regolare, di dimensione \(p_a+\pi-1\).

A tal uopo si associ a \(\lvert C\rvert\) un sistema lineare irriducibile \(\lvert D\rvert\) di dimensione maggiore di \(p_g-p_a\), il cui aggiunto sia regolare; e si costruisca il sistema % ?? uncertain: "A tal" -- "A tal" or "A. tal"; a small intervening mark may be a scan speck.
\[
\lvert L\rvert=\lvert C+D+K\rvert.
\]

Sopra una curva \(L\), tanto il sistema \(\lvert C+K\rvert\) che il sistema \(\lvert D+K\rvert\), segheranno serie complete, le quali dovranno essere residue l’una dell’altra rispetto alla serie canonica. Siccome si conosce la dimensione della serie segata da \(\lvert D+K\rvert\), che è eguale alla dimensione del sistema \(\lvert D+K\rvert\), si dedurrà la dimensione della serie segata da \(\lvert C+K\rvert\), che è quella segata da questo stesso sistema.

Dicansi invero \(\varrho\) ed \(m\) il genere e il grado di \(\lvert D+K\rvert\), ed \(s\) il numero delle intersezioni di una \(C\) con una \(D+K\): la dimensione della serie segata da \(\lvert D+K\rvert\) su \(L\) vale
\[
p_a+m-\varrho+1
\]
e il suo ordine è \(m+s\) mentre il genere di \(L\) è
\[
\bar{\pi}=\pi+\varrho+s-1.
\]

Quindi il teorema di Riemann-Roch ci dà che la dimensione della serie segata sulla stessa \(L\) da \(\lvert C+K\rvert\) è
\[
p_a+m-\varrho+1-(m+s-\bar{\pi})-1=p_a+\pi-1.
\]

Un altro criterio di regolarità (che riesce utile nelle questioni di postulazione d’una superficie rispetto alle varietà \(V_{n-1}\) d’uno spazio \(S_n\) a \(n>3\) dimensioni) concerne i multipli d’un sistema lineare: \textit{Sono regolari i multipli d’ordine assai elevato d’un sistema irriducibile, \(\infty^2\) almeno, \(\lvert C\rvert\), privo di curve fondamentali.}

% scan: s086R
% !PRINTED 149
\weberpage{149}
% head: IL GENERE NUMERICO E IL TEOREMA DI RIEMANN-ROCH, ECC.
% stamp: Vertical scan-edge streaks and a dark mark at the lower right.
Si consideri un certo valore di \(h\) per cui il sistema \(\lvert hC\rvert\) contenga parzialmente il sistema canonico \(\lvert K\rvert\), di guisa che si abbia
\[
\lvert hC-K\rvert=\lvert D\rvert.
\]

Sappiamo che, per \(s\) abbastanza alto il sistema aggiunto ad \(\lvert sC\rvert\), cioè \(\lvert sC+K\rvert\), è regolare. Ora si avrà
\[
\lvert(s+h)C\rvert=\lvert sC+K+D\rvert,
\]
dove \(D\) sarà una curva composta, in generale, di più componenti, semplici o multiple, di grado positivo o negativo. Ciò posto si tratta di mostrare che il sistema regolare \(\lvert sC+K\rvert\) si amplia successivamente in sistemi regolari (secondo il criterio del § 8) quando si sommano ad esso le componenti irriducibili di \(D\).

Infatti, designando con \(L\) una componente irriducibile di \(D\), di caratteri \(\varrho\) e \(\nu\), ed essendo comunque \(\varrho\geq0\), si vede che il numero delle intersezioni di \(L\) con le curve sommande di \(\lvert sC+K\rvert\) (dove \(s\) è alto quanto si vuole) vale
\[
m=st+2\varrho-2-r>\varrho-\nu, % [sic] r -> \nu, the surrounding text denotes the degree by \nu.
\]
essendo \(t>0\) perchè \(L\) non è curva fondamentale per \(\lvert C\rvert\). L’analoga diseguaglianza sussiste sempre per il numero delle intersezioni delle curve di \(\lvert sC+K+L\rvert\) con una componente di \(D\), e così successivamente in rapporto alla somma di \(\lvert sC+K\rvert\) con tutte le componenti irriducibili di \(\lvert D\rvert\).

In ultima analisi il sistema \(\lvert sC+K+D\rvert\), che è il sistema multiplo \(\lvert(s+h)C\rvert\), per \(s+h\) assai elevato, risulta regolare, c. d. d.

Basterà spiegare la cosa riferendosi al primo caso significativo, in cui si tratti di sommare \(\lvert sC+K+L\rvert\) alla \(L\) stessa, che si supponga entrare come componente doppia in una curva \(D\) ed essere di genere \(\varrho\geq0\) e di grado negativo \(\nu=-\nu'\). Diciamo dunque che, in questo caso, il numero delle intersezioni di \(L\) con le curve di \(\lvert sC+K+L\rvert\) vale
\[
st+2\varrho-2+\nu'-\nu'=st+2\varrho-2>\varrho+\nu',
\]
per \(s\) abbastanza alto in confronto a \(\nu'=-\nu\).

Il teorema innanzi stabilito si può enunciare dicendo che:

\textit{Ogni sistema lineare normalmente grande} (rispetto al sistema canonico) \textit{è regolare.}

Giova rilevare esplicitamente che la \textit{proprietà non sussiste} per sistemi grandi che \textit{non} siano tali \textit{normalmente}.

L’esempio più semplice è dato dai sistemi lineari che posseggono curve fondamentali proprie:

\textit{Se il sistema lineare \(\lvert C\rvert\) ha una curva fondamentale propria \(\theta\)}
%CONT>>
% scan: s087L
% !PRINTED 150
\weberpage{150}
% head: CAPITOLO QUARTO
%<<CONT
% stamp: Dark vertical scan-edge band at the right margin.
\textit{di genere \(\varrho>0\), i multipli comunque grandi di \(\lvert C\rvert\) hanno la sovrabbondanza}
\[
\omega\geq\varrho\text{\origfootmark{1}{s087L-1}}.
\]\origfoottext{1}{s087L-1}{Cfr. \textsc{Enriques}, \textit{Ricerche ecc.}, 1893.}

Infatti il distacco di una \(\theta\), di grado \(-\nu\), da \(\lvert sC\rvert\) diminuisce di 1 la dimensione effettiva del sistema, di \(\nu\) il suo grado e di \(\nu+\varrho-1\) il suo genere, e quindi di \(\varrho-1\) la sua dimensione virtuale; vuol dire che la sovrabbondanza \(\omega\) di \(\lvert sC\rvert\) supera di \(\varrho\) quella \(\omega'\) di \(\lvert sC-\theta\rvert\): % ?? uncertain: "distacco di" -- "distacco di" or "distacco. di"; the intervening point may be a scan speck.
\[
\omega=\omega'+\varrho.
\]

Si noti che, in generale (per \(\varrho>1\)) anche \(\omega'>0\), e perciò il genere \(\varrho\) della curva fondamentale è soltanto un limite inferiore della sovrabbondanza di \(\lvert sC\rvert\). Qui riesce istruttivo il semplice esempio proiettivo di una superficie \(F\) d’ordine \(n\) (a sezioni piane \(C\)) dotata d’un punto multiplo isolato d’ordine \(i>3\), su cui le superficie \(\varphi_s\) (d’ordine \(s\) abbastanza alto) segano il sistema lineare \(\lvert sC\rvert\).

Si vede tosto, non soltanto che le \(\varphi_s\) generali segano su \(F\) un sistema sovrabbondante (che ha la curva fondamentale \(\theta\) rappresentata dall’intorno di \(O\)), sì anche che sono sovrabbondanti i sistemi segati su \(F\) dalle \(\varphi_s\) passanti per \(O\) colle molteplicità \(1,\,2,\,\ldotp\ldotp\ldotp\ldotp\,i-3\), mentre è regolare quello segato dalle \(\varphi_s\) aggiunte che passano \(i-2\) volte per \(O\). Donde si trae facilmente il calcolo della sovrabbondanza \(\omega\) di \(\lvert sC\rvert\). % ?? uncertain: accent in "sì anche" -- sì or si.
% EDITORIAL D0083 RISOLTO: The crop shows the separate grave stroke over the i, supporting sì rather than unaccented si.

\edsection{16. Precisazioni sulla regolarità del sistema aggiunto ad una curva\\ sopra una superficie regolare.}{16. Precisazioni sulla regolarità del sistema aggiunto ad una curva sopra una superficie regolare.}{sec.IV.16}

Vogliamo approfondire la questione della regolarità o della sovrabbondanza dei sistemi lineari più ampi del sistema canonico, riferendoci al caso semplice delle \textit{superficie regolari}, con \(p_g=p_a=p>0\).
% The equality sign is repeated at the printed line break.

Se \(\lvert C\rvert\) è un sistema lineare più ampio del sistema canonico \(\lvert K\rvert\), si ha
\[
\lvert C-K\rvert=\lvert L\rvert,
\]
dove \(\lvert L\rvert\) designa un sistema effettivo, e quindi
\[
\lvert C\rvert=\lvert L+K\rvert=\lvert L'\rvert.
\]

Così \(\lvert C\rvert\) si può considerare come sistema aggiunto ad una curva \(L\) effettiva. Ora \textit{se \(L\) è irriducibile}, di un genere \(\pi>0\), il \textit{suo sistema aggiunto è certo regolare}, di dimensione \(p+\pi-1\) giacchè esso sega sopra \(L\) una serie (completa) che non può superare la serie canonica \(g^{\pi-1}_{2\pi-2}\).

% scan: s087R
% !PRINTED 151
\weberpage{151}
% head: IL GENERE NUMERICO E IL TEOREMA DI RIEMANN-ROCH, ECC.
La stessa conclusione --- \textit{regolarità del sistema aggiunto ad \(L\)} --- si estende al caso in cui la curva \(L\) sia una \textit{curva connessa comunque spezzata in parti semplici}.

\textit{Curva connessa} formata di componenti semplici significa « curva che comunque divisa in due parti dà luogo a curve aventi qualche punto (di connessione) a comune », ovvero anche topologicamente, una « curva o superficie riemanniana su cui può darsi sempre un cammino continuo che ne congiunga due punti arbitrari ».

Per le curve connesse, formate di componenti semplici, che sono così definite, sussiste il

\textit{Lemma di composizione geometrica delle curve connesse.} Una curva connessa \(L\) formata di \(s\) componenti semplici si può sempre ottenere sommando ad una parte connessa formata di \(s-1\) componenti \(L_1+L_2+\ldotp\ldotp\ldotp\ldotp+L_{s-1}\) una \(s\)-ma componente \(L_s\) che abbia con questa qualche punto comune (di connessione).

Si dimostra per induzione completa: supponendo che la data \(L\) contenga una parte connessa, formata di \(h\) componenti irriducibili semplici con \(h<s\), questa parte deve essere connessa colla rimanente e perciò con una almeno delle sue componenti; sommandole tale componente si ottiene una parte connessa di \(L\), formata da \(h+1\) componenti.

Come \textit{Corollario}: data una curva connessa \(L\) formata di \(s\) componenti semplici, si può prendere queste in un ordine \(L_1,\,L_2,\,\ldotp\ldotp\ldotp\ldotp,\,L_s\) per modo che nella
\[
L=L_1+L_2+\ldotp\ldotp\ldotp\ldotp+L_s
\]
tutte le parti \(L_1+L_2+\ldotp\ldotp\ldotp\ldotp+L_h\) (\(h=1,\,2,\,\ldotp\ldotp\ldotp\ldotp,\,s\)) siano connesse.

In forza di questa proprietà è facile riconoscere che « il sistema aggiunto \(\lvert L'\rvert\) è regolare », come sopra si è enunciato. Infatti si designi con \(\pi_h\) il genere di \(L_1+\ldotp\ldotp\ldotp\ldotp+L_h\), con \(\bar{\pi}_h\) il genere di \(L_h\), e con \(i_h\) il numero delle intersezioni \((L_1+\ldotp\ldotp\ldotp\ldotp+L_{h-1})L_h\). Avremo che il genere di
\[
L=L_1+\ldotp\ldotp\ldotp\ldotp+L_s=(L_1+\ldotp\ldotp\ldotp\ldotp+L_{s-1})+L_s
\]
sarà
\[
\pi_s=\bar{\pi}_s+\pi_{s-1}+i_s-1,
\]
e la serie segata da \(\lvert L'\rvert\) su \(L_s\) sarà d’ordine
\[
n_s=2\bar{\pi}_s-2+i_s
\]
e di dimensione
\[
r_s\leq\bar{\pi}_s-2+i_s;
\]
quindi lo staccamento di \(L_s\) da \(\lvert L'\rvert\) importerà
\[
r_s+1\leq\bar{\pi}_s+i_s-1
\]
condizioni.

% scan: s088L
% !PRINTED 152
\weberpage{152}
% head: CAPITOLO QUARTO
% stamp: Dark vertical scan-edge band at the right margin.
Risulta quindi: se è regolare il sistema \(\lvert(L_1+\ldotp\ldotp\ldotp\ldotp+L_{s-1})'\rvert\) aggiunto ad \(\lvert L_1+\ldotp\ldotp\ldotp\ldotp+L_{s-1}\rvert\) si ha precisamente % ?? uncertain: point after "sistema" -- scan speck (reading adopted) or printed full stop.
% EDITORIAL D0084 RISOLTO: The isolated tiny dot is well below the text baseline, unlike printed punctuation. It is best treated as a scan/printing speck.
\[
r_s+1=\bar{\pi}_s+i_s-1
\]
ed è regolare anche il sistema \(\lvert L'\rvert\) aggiunto ad \(\lvert L\rvert=\lvert L_1+\ldotp\ldotp\ldotp\ldotp+L_s\rvert\). Il teorema enunciato si verifica per una \(L\) formata di \(s\) componenti (semplici e connesse fra loro) se si verifica per una curva formata di \(s-1\) componenti. E così può ritenersi dimostrato in generale, essendo regolare \(\lvert L'_1\rvert\).
% The plus sign is repeated at the printed line break.

Al contrario di ciò che accade per le curve connesse, non sono regolari, ma sovrabbondanti, i sistemi aggiunti a curve spezzate sconnesse. Più precisamente:

\textit{Il sistema \(\lvert L'\rvert\) aggiunto ad una curva \(L\) composta di \(s\) parti irriducibili semplici senza punti comuni ha la sovrabbondanza \(s-1\).}

Si dimostra collo stesso ragionamento induttivo, adoperato per le curve connesse. Si tratti per es. di una curva \(L\) (di genere \(\pi\)) composta di due parti \(L_1\) e \(L_2\), di genere \(\pi_1\) e \(\pi_2\), senza punti comuni (\(\pi=\pi_1+\pi_2-1\)). Allora il sistema aggiunto \(\lvert L'\rvert=\lvert(L_1+L_2)'\rvert\) sega su \(L_1\) la serie canonica completa \(g^{\pi_1-1}_{2\pi_1-2}\), ed il residuo di \(L_1\) rispetto ad esso --- che è \(\lvert L'_2\rvert\) --- sega su \(L_2\) la serie canonica \(g^{\pi_2-1}_{2\pi_2-2}\); quest’ultimo sistema \(\lvert L'_2\rvert\) è regolare, quindi \(\lvert L'\rvert\) è sovrabbondante, di dimensione
% The equality sign is repeated at the printed line break.
\[
p+\pi=p+(\pi_1+\pi_2-1)-1+\omega
\]
con
\[
\omega=1.
\]

Dopo ciò è anche facile valutare la sovrabbondanza del sistema \(\lvert(L+\Sigma s_i\theta_i)'\rvert\) aggiunto ad una curva sconnessa \(L+\Sigma s_i\theta_i\), formata di una curva irriducibile \(L\) di grado \(>0\), diciamo per semplicità di una curva variabile in un sistema lineare \(\infty^2\), e di più curve irriducibili \(\theta_i\) di grado \(\nu_i<0\), fondamentali per questo sistema \(\lvert L\rvert\), contate ciascuna un certo numero \(s_i\) di volte.

\edsection{17. Segue: il sistema aggiunto ad una curva connessa formata di\\ componenti multiple sopra una superficie regolare.}{17. Segue: il sistema aggiunto ad una curva connessa formata di componenti multiple sopra una superficie regolare.}{sec.IV.17}

I risultati indicati non esauriscono la questione della regolarità o sovrabbondanza del sistema aggiunto ad una curva \(L\) sopra una superficie regolare, quando la \(L\) sia comunque composta di parti semplici e multiple. La difficoltà che qui s’incontra tiene invero al concetto della « curva connessa », che occorre estendere appunto al caso di una \(L\) qualunque, composta come sopra. % ?? uncertain: "che occorre" -- "che occorre" or "che. occorre"; the small intervening point may be a scan speck.

Designamo come \textit{curva (aritmeticamente) connessa} una curva \(L\), comunque formata di componenti irriducibili semplici e multiple,
%CONT>>
% scan: s088R
% !PRINTED 153
\weberpage{153}
% head: IL GENERE NUMERICO E IL TEOREMA DI RIEMANN-ROCH, ECC.
%<<CONT
% stamp: Faint vertical scan-edge streak along the right margin.
quando due parti qualunque, \(L_1\) e \(L_2\), in cui essa venga divisa, hanno sempre un numero (virtuale) di punti comuni
\[
L_1L_2>0:
\]
non è escluso che qualche parte di \(L_1\) abbia un numero negativo \(-n\) di punti comuni con \(L_2\) (contenente almeno una componente multipla di \(L_2\), di grado negativo), purchè la parte restante abbia allora più che \(n\) punti comuni colla stessa \(L_2\).

La considerazione delle curve connesse, comunque formate con parti multiple, dà luogo a difficoltà che sono state esposte da A. \textsc{Franchetta}\origfootnote{1}{s088R-1}{\textit{Sulla caratterizzazione delle curve eccezionali riducibili di prima specie}, in Boll. dell’Unione Mat., Bologna, 1941.}, e che qui conviene segnalare.

Contrariamente a ciò che dà l’intuizione nel caso di parti semplici (aventi con altre un numero \(\geq0\) di punti comuni) \textit{non} si può dire in generale che aggiungendo ad una curva connessa \(L\) una componente \(C\) che la incontri in \(n>0\) punti si ottiene sempre una curva connessa \(L+C\).

Infatti \textsc{Franchetta} adduce \textit{esempi di curve sconnesse formate di \(s\) componenti, scomponibili in una curva connessa formata di \(s-1\) componenti e in una componente che l’incontra in un numero positivo di punti.}

In conseguenza di ciò cade anche la dimostrazione del lemma (\textit{di composizione geometrica}) che abbiamo stabilito nel paragrafo precedente per le curve connesse formate di componenti semplici: non si può dire in generale che una curva aritmeticamente connessa formata di \(s\) componenti, si ottenga sempre aggiungendo una componente irriducibile ad una curva connessa formata di \(s-1\) componenti.

Tuttavia sussistono per le curve aritmeticamente connesse alcune proprietà, messe in luce da \textsc{Franchetta}\origfootnote{2}{s088R-2}{l. c.}, che prendono il posto delle precedenti, quali vengono enunciate dai seguenti teoremi:

\textit{Teorema I.} -- Una curva \(L\) (formata comunque di componenti irriducibili semplici o multiple) è aritmeticamente connessa se ogni sua parte connessa è connessa colla rimanente (cioè ha con questa un numero virtuale di punti comuni \(>0\)).

Invero se la \(L=A+B\) non sia connessa, avendosi \(AB\leq0\), o \(A\) è connessa e si ha dunque una parte connessa di \(L\) che interseca la residua in un numero \(\leq0\) di punti, ovvero è possibile decomporre \(A\) in due parti \(A_1\) e \(A_2\) non connesse fra loro:
\[
A=A_1+A_2,\qquad A_1A_2\leq0;
\]
%CONT>>
% scan: s089L
% !PRINTED 154
\weberpage{154}
% head: CAPITOLO QUARTO
%<<CONT
% stamp: Thin dark vertical scan-edge line at the right margin.
mentre dall’essere \(AB\leq0\) segue che una almeno delle due curve, sia \(A_1\), un numero non positivo d’intersezioni con \(B\) ma in questo caso si ha % [sic] "sia A_1, un numero" -> "sia A_1, ha un numero", the clause appears to omit its verb.
\[
L=A_1+(A_2+B)
\]
con
\[
A_1(A_2+B)=A_1A_2+A_1B\leq0,
\]
ossia vi è una parte \(A_1\) di \(A\), entro \(L\), sconnessa colla parte residua. Ora sarà \(A_1\) connessa ovvero conterrà una parte \(A_{11}\) sconnessa colla curva residua \((L-A_{11})\) e così di seguito. Ma la serie delle parti successive \(A,\,A_1,\,A_{11},\,\ldotp\ldotp\ldotp\ldotp\), il cui numero di componenti va diminuendo, avrà termine con una curva connessa, formata di una o più componenti irriducibili: l’ultima parte così ottenuta sarebbe una curva connessa in sè e sconnessa colla residua.

\textit{Teorema II.} -- Da una curva \(L\) (aritmeticamente) connessa, si può togliere una parte \(B\) connessa (formata di una o più componenti) per modo che la parte residua \(L-B\) sia ancora connessa.

Fra tutte le possibili parti connesse di \(L\), se ne scelga una, \(B\), che abbia il minimo numero (positivo) di punti comuni colla parte residua \(A=L-B\); diciamo che \(A\) risulta pure connessa. Se così non fosse, si avrebbe una divisione in parti sconnesse
\[
\begin{aligned}
A&=A_1+A_2:\\
A_1A_2&\leq0,
\end{aligned}
\]
dove si può supporre che \(A_1\) sia, in sè, connessa (teor. I). Ma, in tal caso, la \(B\) dovrà avere con ciascuna delle parti \(A_1\) e \(A_2\) di \(A\) un numero positivo d’intersezioni, chè altrimenti sarebbe \((B+A_1)A_2\leq0\) o \((B+A_2)A_1\leq0\) e perciò la \(L\) sarebbe sconnessa; avremo dunque
\[
BA_1>0\qquad\text{e}\qquad BA_2>0,
\]
quindi
\[
BA_1<BA=B(A_1+A_2),
\]
e a fortiori
\[
(B+A_2)A_1<BA;
\]
così la \(A_1\), parte connessa di \(L\), avrebbe colla residua un numero d’intersezioni minore che la \(B\), contro la nostra ipotesi.

I teoremi di \textsc{Franchetta} tengono luogo, fino ad un certo punto, delle proprietà intuitive che caratterizzano le curve connesse formate di componenti semplici, siccome vedremo nel seguito.

Tuttavia non sembra facile dedurne, senza eccezione, la regolarità del sistema \(\lvert L'\rvert\) aggiunto ad una curva connessa qualsiasi \(L\) (con componenti multiple di grado negativo). Si può tentare una via alquanto indiretta, che è in stretto rapporto con i nostri sviluppi sul teorema di Riemann-Roch, ma che resta tuttavia subordinata ad un’ipotesi di lavoro che verrà messa in luce più sotto.

% scan: s089R
% !PRINTED 155
\weberpage{155}
% head: IL GENERE NUMERICO E IL TEOREMA DI RIEMANN-ROCH, ECC.
% stamp: Vertical scan-edge streak at the right margin; faint horizontal smudges below the dimension formula.
Si assuma sulla superficie \(F\) un sistema lineare regolare \(\lvert C\rvert\), che sia multiplo d’ordine assai elevato d’un sistema irriducibile \(\infty^2\) almeno, privo di curve fondamentali, sicchè risulti, in primo luogo, irriducibile \(\lvert D\rvert=\lvert C+L\rvert\): si può ritenere che tanto \(\lvert L\rvert\) come il suo sistema aggiunto \(\lvert L'\rvert=\lvert L+K\rvert\) seghino serie complete sopra una curva \(\bar{D}\) di questo sistema \(\lvert C+L\rvert\) normalmente grande. Ora si riconoscerà che anche \(\lvert C\rvert\) sega su tale curva \(\bar{D}\) una serie completa. Infatti, sia \(G\) un gruppo della serie completa \(CD\) e \(\Gamma\) un gruppo residuo di \(G\) rispetto alla serie caratteristica di \(\lvert D\rvert\), sezione di \(D\) con \(L\); \(G+\Gamma\) costituisce il gruppo base d’un fascio di curve \(D\), che non hanno alcuna intersezione colle componenti di \(L\), fuori di \(\bar{D}\).

Ora se s’impone ad una \(D\) di questo fascio di contenere un punto di \(L\) fuori di \(\bar{D}\), stante la connessione di \(L\), l’intera curva connessa si staccherà dalla \(D\), restando una \(C\) per \(G\), c. d. d.

La dimostrazione precedente si basa sul presupposto che essendo la \(L\) aritmeticamente connessa, quando s’imponga ad una \(D\) del fascio predetto, di contenere un punto di \(L\), fuori di \(\Gamma\), la \(L\) si stacchi completamente, ossia ogni sua componente si stacchi con la molteplicità con cui essa compare in \(L\).

Tale proprietà ci pare vera, ma non ne possediamo una dimostrazione esauriente. Noi la introdurremo come ipotesi di lavoro, restando così inteso che il risultato precedente e le conseguenze che ne trarremo, sono subordinati alla validità di tale ipotesi.

Ora, avendo riconosciuto che il sistema regolare \(\lvert C\rvert\) sega sulla curva \(D\) di \(\lvert D\rvert=\lvert C+L\rvert\) una serie completa \(g\), che è residua di quella segata da \(\lvert L'\rvert=\lvert L+K\rvert\), si potrà valutare la dimensione di quest’ultima (e quindi di \(\lvert L'\rvert\)) che eguaglia l’indice di specialità di \(g\) diminuito di una unità. Designando con \(n\) il grado di \(\lvert C\rvert\), con \(\varrho\) il suo genere, e con \(m\) il numero delle intersezioni \(CL\), l’ordine della serie \(g\) sarà
\[
n+m
\]
e la sua dimensione
\[
r=p+n-\varrho+1,
\]
mentre il genere di \(D\) vale
\[
\pi+\varrho+m-1.
\]

Quindi l’indice di specialità \(i\) di \(g\) è dato da:
\[
n+m-(\pi+\varrho+m-1)+i=p+n-\varrho+1,
\]
e di qui segue che la dimensione della serie \((CL')\) residua di \(CD\) è % [sic] (CL') -> (DL'), the preceding paragraph places this residual series on D.
\[
i-1=p+\pi-1.
\]

Dunque \(\lvert L'\rvert\) ha questa stessa dimensione ed è regolare.

% scan: s090L
% !PRINTED 156
\weberpage{156}
% head: CAPITOLO QUARTO
% stamp: Dark vertical scan-edge band at the right margin.
In conclusione: \textit{Sopra una superficie regolare di genere \(p>0\), il sistema aggiunto ad una curva aritmeticamente connessa è sempre regolare.}

\edsection{18. Sulla riducibilità del sistema canonico.}{18. Sulla riducibilità del sistema canonico.}{sec.IV.18}

Il teorema stabilito nel paragrafo precedente porta a notevoli applicazioni in ordine al sistema bicanonico e ai sistemi pluricanonici. Per comprenderne il senso conviene rifarci dalla questione della riducibilità del sistema canonico.

Noether riteneva di poter dimostrare che il sistema canonico di una superficie di genere \(p>0\) possa essere riducibile soltanto in due casi:

1) per la presenza di curve eccezionali (di genere \(\varrho=0\) e grado \(\nu=-1\)) che costituiscono parti fisse del sistema canonico impuro;

2) per essere il genere lineare (assoluto) \(p^{(1)}=1\), nel qual caso le curve canoniche pure sono in generale composte colle curve ellittiche d’un fascio.

Così per \(p^{(1)}>1\) il sistema canonico puro sarebbe sempre irriducibile. Ma la deduzione di Noether cade in errore, e si danno effettivi \textit{esempi di curve canoniche} (pure) \(K\) \textit{riducibili}, appartenenti a superficie di genere \(p>0\) e genere lineare \(p^{(1)}>1\):

1) Il primo esempio viene offerto dal piano doppio con curva di diramazione \(C_{2n}\) d’ordine \(2n\), dotata d’un punto \(O\) \((2n-6)\)-plo (\(n>4\)) e di due punti \([3,3]\), \(A_1\) e \(A_2\), allineati con \(O\) sopra una retta \(a\)
\[
(p=2n-7,\qquad p^{(1)}=4n-15).
\]

Le curve canoniche \(K\) hanno per immagini (doppie) le curve \(C_{n-3}\) con \(O\) \((n-4)\)-plo passanti semplicemente per \(A_1\) e \(A_2\), e perciò spezzate in \(C_{n-4}\) con \(O\) \((n-5)\)-plo e nella parte fissa \(a\). Ora questa retta rappresenta una curva \(\theta\) (irriducibile) componente fissa delle \(K\), che non è eccezionale, avendo il genere \(\varrho=0\) e il grado \(\nu=-2\): siccome appare anche dal fatto che essa ha due intersezioni colle residue componenti delle \(K\), perchè il numero di queste intersezioni vale:
\[
\begin{aligned}
&2\varrho-2-2\nu=2\\
&(\varrho=0,\qquad \nu=-2).
\end{aligned}
\]

2) Un secondo esempio è offerto dal piano doppio con curva di diramazione \(C_{10}\), d’ordine 10, dotata di 3 punti \([3,3]\) sopra una retta \(a\), che non fa parte di \(C_{10}\). Un semplice computo di costanti ci dà un sistema lineare di \(C_{10}\) con 3 punti 3-pli assegnati su \(a\) e altrettanti punti 3-pli vicini a questi in direzioni date, sistema di
%CONT>>
% scan: s090R
% !PRINTED 157
\weberpage{157}
% head: IL GENERE NUMERICO E IL TEOREMA DI RIEMANN-ROCH, ECC.
%<<CONT
% stamp: Short isolated horizontal mark in the right margin beside the line ending "Ora".
dimensione 29, che pertanto non si riduce a quello \(\infty^{26}\) costituito dalla \(a\) contata 2 volte e da curve \(C_8\) con 3 tacnodi e tangenti tacnodali date. Ora, sul piano doppio che abbiamo costruito, le immagini delle curve canoniche sono le coniche spezzate nella retta fissa \(a\) e in una retta variabile. La corrispondente parte fissa \(\theta\) del sistema canonico \(\lvert K\rvert\) è una curva di genere \(\varrho=1\) e grado \(\nu=-1\); i generi della superficie \(F\) valgono
\[
p=3,\qquad p^{(1)}=6.
\]

Si calcolerà il grado \(\nu=-1\) di \(\theta\) avvertendo che la \(\theta\) --- parte fissa di \(\lvert K\rvert\) --- ha 2 intersezioni colle curve residue, sicchè
\[
2\varrho-2-2\nu=-2\nu=2,\qquad \nu=-1.
\]

D’altronde la curva completa che risponde alla \(a\) su \(F\) è una curva di genere 4 composta di 3 curve ellittiche di grado \(-1\) rispondenti ai punti \([3,3]\) e della nostra \(\theta\), che ha un punto comune con ciascuna di esse: il grado complessivo di codesta curva completa è
\[
-3-1+2.3=2.
\]

3) Un terzo esempio è offerto dal piano doppio con curva di diramazione \(C_{12}\), d’ordine 12, con 4 punti \([3,3]\) appartenenti ad una retta \(a\) (che non fa parte di \(C_{12}\)); con un computo di costanti si dimostra l’esistenza di \(C_{12}\) irriducibili, che non contengono \(a\). Ora le curve canoniche \(K\) sono date sul piano doppio da cubiche, spezzate nella retta \(a\) e in una conica variabile. Alla retta \(a\) risponde una \(\theta\) componente fissa delle \(K\) che ha il genere \(\varrho=1\) e il grado \(\nu=-2\); infatti la \(\theta\) ha 4 intersezioni colle curve residue, sicchè % ?? uncertain: isolated dash to the right of "Ora" -- margin artefact (reading adopted) or printed dash.
% EDITORIAL D0086 RISOLTO: The isolated horizontal stroke lies well beyond the line-end word and below the ordinary hyphen/dash position, in the margin rather than the text; omission is warranted.
\[
2\varrho-2-2\nu=-2\nu=4,\qquad \nu=-2\text{\origfootmark{1}{s090R-1}}. % [sic] "rappresentato del fascio" -> "rappresentato dal fascio", the preposition appears to be a typographical error.
\]\origfoottext{1}{s090R-1}{Citiamo ancora l’esempio di una superficie irregolare con \(p_a=2\), \(p_g=3\), \(p^{(1)}=11\), quale è rappresentata sul piano doppio con \(C_{12}\) di diramazione composta di 3 quartiche \(C_4\) (d’un fascio) che si toccano in 8 punti d’una conica \(C_2\); le immagini delle curve canoniche sono spezzate nella conica \(C_2\), cui risponde una \(\theta\) fissa di genere 3, e in una retta variabile. La superficie contiene un fascio irrazionale (ellittico) di curve, rappresentato del fascio delle \(C_4\), d’accordo con un teorema generale sui piani doppi irregolari di M. \textsc{De Franchis}.}

La curva completa che risponde ad \(a\) su \(F\) è una curva di genere 5 composta di 4 curve ellittiche di grado \(-1\) rispondenti ai punti \([3,3]\) e della nostra \(\theta\), che ha un punto comune con ciascuna di esse: grado complessivo di codesta curva completa è
\[
-4-2+2.4=2.
\]

% scan: s091L
% !PRINTED 158
\weberpage{158}
% head: CAPITOLO QUARTO
% stamp: Right-edge scan line and scattered speckles.
\textit{Nota.} -- Nello stato attuale delle nostre conoscenze non si sa se esistono superficie con \(p^{(1)}>1\) il cui sistema canonico sia riducibile in modo che le sue parti variabili siano composte colle curve \(k\) d’un fascio\origfootnote{1}{s091L-1}{Esempi effettivi sono stati costruiti recentemente da \textsc{G. Pompilj}: \textit{Alcuni esempi di superficie algebriche a sistema canonico puro degenere}, Rend. Acc. Lincei, s. 8, v. IV, 1948.}. In ordine a tale possibilità possiamo soltanto dimostrare la seguente proposizione:

\textit{Se le curve canoniche pure d’una superficie di genere \(p>1\)} (all’infuori di parti fisse) \textit{sono composte colle curve \(k\) d’un fascio di genere \(\pi\) e grado \(n\), si ha in ogni caso}
\[
\pi\leq p^{(1)},\qquad n\leq p^{(2)}=p^{(1)}-1.
\]

Si designino con \(\theta\) le eventuali componenti fisse delle \(K\):
\[
|K|=\Sigma\theta+|sk|
\]
e sieno \(\varrho\) e \(\nu\) i caratteri di una \(\theta\). Esprimiamo il grado di \(|K|\) come somma dei numeri d’intersezione delle sue componenti \(k\) (irriducibili o meno) e \(\theta\); avremo
\[
s(2\pi-2-n)+\Sigma(2\varrho-2-\nu)=p^{(1)}-1.
\]

Ora, riferendoci ad una superficie trasformata priva di curve eccezionali sarà per ogni \(\theta\)
\[
2\varrho-2-\nu\geq0,
\]
quindi
\[
2\pi-2-n\leq p^{(1)}-1.
\]

D’altra parte una curva \(k\) ha, colla curva residua \((s-1)k+\Sigma\theta\), il numero virtuale d’intersezioni
\[
2\pi-2-2n\geq0.
\]

Di qui si deduce
\[
n\leq\pi-1,
\]
e
\[
\pi\leq p^{(1)},\qquad n\leq p^{(1)}-1,\qquad \text{c. d. d.}
\]

\edsection{19. Sulla regolarità del sistema bicanonico}{19. Sulla regolarità del sistema bicanonico}{sec.IV.19}

Le possibilità che il sistema canonico puro \(|K|\) d’una superficie, con \(p>0\) e \(p^{(1)}>1\), sia riducibile, costringe ad esaminare più profondamente la questione « se il sistema bicanonico e i sistemi pluricanonici siano regolari ». Per ciò occorre dimostrare che, comunque le \(K\) siano spezzate in parti semplici e multiple, sempre (per \(p^{(1)}>1\)) esse sono connesse; quindi è regolare il sistema aggiunto, % [sic] Le possibilità -> La possibilità, the following verb "costringe" is singular.
\[
|2K|=|K'|,
\]
%CONT>>
% scan: s091R
% !PRINTED 159
\weberpage{159}
% head: IL GENERE NUMERICO E IL TEOREMA DI RIEMANN-ROCH, ECC.
%<<CONT
% stamp: Scattered scan speckles.
e similmente, essendo pur connesse le curve di \(|2K|\), è regolare i sistema tricanonico % [sic] i sistema -> il sistema, the singular article is missing its final l.
\[
|3K|=|(2K)'|\quad\text{ecc.}
\]

In parole: \textit{Sopra una superficie regolare di genere \(p>0\) e di genere lineare \(p^{(1)}>1\), il bigenere vale esattamente}
\[
P=p+p^{(1)}
\]
\textit{e lo \(i\)-genere}
\[
P_i=p+\frac{i(i-1)}{2}p^{(1)}.
\]

Per dimostrare questo importante teorema conviene riconoscere successivamente che, per \(p>0\) e \(p^{(1)}>1\): % ?? uncertain: closing mark in the exponent of the second p -- closing parenthesis, comma-like damaged mark.

1) le parti variabili del sistema bicanonico sono irriducibili;

2) le curve canoniche \(K\), comunque riducibili, sono sempre connesse, di \textit{ordine di connessione} \(\geq2\) (cioè una parte qualsiasi di \(K\) ha almeno 2 intersezioni colla residua).

3) di conseguenza il sistema bicanonico puro è regolare di dimensione \(p+p^{(1)}-1\) e non maggiore.

Queste deduzioni si estendono a fortiori, ai sistemi pluricanonici d’ordine \(>2\).

1) La prima proposizione si dimostra facilmente. Anzitutto se \(p\geq2\) e \(p^{(1)}>1\), le parti variabili delle curve canoniche sono irriducibili o formate colle curve irriducibili \(C\), di genere \(\pi>1\), appartenenti ad un fascio; quindi le curve \(C'\) aggiunte ad una \(C\) formano un sistema lineare di dimensione \(p+\pi-1\) (almeno) secanti su \(C\) la serie canonica completa \(g_{2\pi-2}^{\pi-1}\); due \(C'\) hanno almeno \(2\pi-2\) intersezioni variabili e perciò le \(C'\) non sono composte colle curve d’un fascio; a fortiori il sistema bicanonico \(|2K|\) che contiene parzialmente \(|C'|\) non può avere parti variabili riducibili composte colle curve d’un fascio.

Il teorema si estende al caso \(p=1\). In questo caso, se le parti variabili delle curve bicanoniche sono composte colle curve irriducibili \(C\) di un fascio, anzitutto --- per essere \(p^{(1)}>1\) --- il genere di queste dovrà essere \(\pi>1\)\origfootnote{1}{s091R-1}{Infatti vedremo in seguito che se una superficie possiede un fascio di curve ellittiche ed il suo genere geometrico è positivo allora il suo genere lineare è necessariamente \(p^{(1)}=1\).}. Ciò posto, se le componenti \(C\) che entrano nelle curve di \(|2K|\) sono in numero di \(s\) (\(P=s+1\) designando il bigenere della superficie), per essere \(p=1\), vi saranno \(s\) curve \(C_{\frac{1}{2}}\) i cui doppi appartengono al fascio \(|C|\):
\[
K=\theta+sC_{\frac{1}{2}}.
\]

% scan: s092L
% !PRINTED 160
\weberpage{160}
% head: CAPITOLO QUARTO
% stamp: Right-edge scan line and scattered speckles.
Ora il sistema \(|C'_{\frac{1}{2}}|\) aggiunto ad una curva \(C_{\frac{1}{2}}\) e contenuto in \(|2K|\) dovrebbe avere come parti variabili le curve \(C\) di \(|C|\); ma ciò è impossibile perchè la sua dimensione vale non già uno, bensì (almeno) \(p+\pi-1=\pi>1\)\origfootnote{1}{s092L-1}{La prop. 1) si dimostra anche in altra maniera, facendo vedere anzitutto che deve essere \(\pi=2\) e quindi necessariamente \(p>1\): Se le parti variabili delle curve bicanoniche \(2K\) sono composte di \(s\) (\(>1\)) curve \(C\) d’un fascio lineare, completo, designando con \(\theta\) le parti fisse di \(|2K|\) scriveremo
\[
|2K|=\Sigma\theta+|sC|.
\]

Ora si indichino con \(n\) e \(\pi\) il grado e il genere di \(|C|\), e con \(\nu\) e \(\varrho\) quelli di una \(\theta\), e valutiamo il carattere additivo delle \(2K\) che esprime il numero delle loro intersezioni con una \(K\); avremo
\[
2K\cdot K=2p^{(1)}-2=\Sigma(2\varrho-2-\nu)+s(2\pi-2-n).
\]

Ma, riferendoci, com’è lecito, ad una superficie senza curve eccezionali:
\[
2\varrho-2-\nu\geq0,
\]
e perciò
\[
s(2\pi-2-n)\leq2p^{(1)}-2.
\]
Quindi, affinchè il fascio \(|C|\) sia completo, la sua dimensione dovrà essere \(r=1\); perciò, se \(|C|\) non è contenuto nel sistema canonico \(|K|\):
\[
p+n-\pi+1\leq1
\]
e
\[
n<\pi-1; % [sic] n<\pi-1 -> n\leq\pi-1, the preceding inequality gives only a non-strict bound when p=1.
\]
e questa diseguaglianza vale a fortiori se \(|C|\) è contenuto in \(|K|\), perchè in tal caso \(|2C|\) fa parte del sistema aggiunto
\[
|C'|=|C+K|.
\]

Ciò posto, se \(\pi>2\), si avrà
\[
\begin{aligned}
s(\pi-1)&\leq2p^{(1)}-2\\
\pi-1&\geq2\\
p^{(1)}-1&\geq s
\end{aligned}
\]
e di conseguenza la dimensione di \(|2K|\) vale
\[
\geq p+p^{(1)}-1>s,
\]
in contraddizione coll’ipotesi che \(|2K|\) sia formato dai gruppi di \(s\) curve del fascio \(|C|\) e quindi abbia la dimensione \(s\).

Dopo avere così stabilito la nostra proposizione nel caso \(\pi>2\), resta a discutere l’ipotesi \(\pi=2\).

In questo caso il fascio delle curve \(C\) di genere \(\pi=2\) permette di rappresentare la superficie sopra un piano doppio con curva di diramazione \(f_{2m}\) di un certo ordine pari \(2m\), dotata d’un punto \(O\) \((2m-6)\)-plo o \((2m-5)\)-plo e di altre singolarità elementari: punti 4-pli e punti \([3,3]\).

In questo piano le curve bicanoniche hanno per immagini curve \(\Phi_{2m-6}\) aggiunte ad \(f_{2m}\), che si compongono di una \(\Phi_{2m-6-s}\) fissa e di \(s\) rette per \(O\). Siccome fra le curve bicanoniche c’è almeno una curva canonica contata due volte, la \(\Phi_{2m-6-s}\) dovrà essere il doppio di una curva \(\Psi_{m-3-h}\) (\(s=2h\)) che, insieme ad \(h\) rette per \(O\), costituisce l’immagine di una curva canonica. Essendo \(h\geq1\), si deduce \(p>1\), caso già esaurito dalle prime considerazioni del testo!}.

% scan: s092R
% !PRINTED 161
\weberpage{161}
% head: IL GENERE NUMERICO E IL TEOREMA DI RIEMANN-ROCH, ECC.
% stamp: Scattered speckles and a curved stray ink mark crossing "connessa" near the middle of the page.
2) Per dimostrare la connessione delle curve canoniche \(K\), basterà dimostrare che sono connesse le curve bicanoniche di \(|2K|\), e per ciò si tratta di ridurre all’assurdo l’ipotesi che queste siano sconnesse. Ora il sistema bicanonico \(|2K|\) consterà di una parte fissa \(L\) (comunque composta) e di parti variabili irriducibili \(C\):
\[
|2K|=L+|C|.
\]

Se le curve del detto sistema sono sconnesse, si possono dividere in due parti aventi un numero \(\leq0\) di punti comuni: una di queste parti conterrà una \(C\) ed una parte (fissa) \(E\) di \(L\), mentre l’altra sarà una \(D\) fissa: % ?? uncertain: apparent dot after "dividere" -- scan speck (not transcribed), printed full stop.
\[
L=D+E,\qquad (C+E)D\leq0.
\]

Ora è lecito ritenere che la \(D\) sia connessa; se non lo fosse si lascerebbe scomporre in due parti \(D_1\) e \(D_2\) con un numero nullo o negativo di punti comuni, e una di queste --- sia p. es. \(D_1\) --- avrebbe
\[
D_1(C+E+D_2)\leq0
\]
intersezioni colla residua; e così seguitando si arriverebbe in ogni caso a trovare una parte connessa di \(L\), e quindi di \(C+L\), avente un numero \(\leq0\) d’intersezioni colla curva residua.

Supposto dunque che \(D\) sia connessa, si designi con \(\varrho\) il suo genere e con \(\nu\) il suo grado: per essere \(D\) connessa sarà
\[
\varrho\geq0.
\]\origfootnote{1}{s092R-1}{Questa proposizione si dimostra, con induzione completa, per le curve \(D\) composte di \(s\) componenti irriducibili ove si ritenga dimostrata per le curve composte di meno che \(s\) componenti. Infatti la \(D\) si può decomporre in due parti \(D_1\) e \(D_2\), ciascuna connessa in sè e connesse fra loro, con \(n>0\) punti comuni; quindi, designando con \(\varrho_1\) e \(\varrho_2\) i generi di queste curve, si avrà
\[
\begin{aligned}
\varrho&=\varrho_1+\varrho_2+n-1,\\
\varrho_1&\geq0,\quad \varrho_2\geq0,\quad \varrho\geq0;
\end{aligned}
\]
\hfill c. d. d.}

Scriviamo il numero delle intersezioni di \(D\) colla curva residua rispetto al sistema bicanonico \(|2K|\):
\[
D(2K-D)=D(C+E)=4\varrho-4-3\nu\leq0.
\]

Se ne deduce
\[
\nu\geq\frac{4}{3}(\varrho-1).
\]

Ora, se \(\varrho=0\), per essere \(\nu\) intero, sarà
\[
\nu\geq-1:
\]
% signature: Printed gathering signature.
% EDITORIAL: segnatura originale: \textsc{Enriques F.} - \textit{Superficie algebriche.}\hfill 11
\printersignature{\textsc{Enriques F.} - \textit{Superficie algebriche.}\quad 11}
%CONT>>
% scan: s093L
% !PRINTED 162
\weberpage{162}
% head: CAPITOLO QUARTO
%<<CONT
% stamp: Right-edge scan line and scattered speckles.
quindi o \(\nu\geq0\) ciò che porta \(p=0\), ovvero la \(D\) è una curva eccezionale della superficie (\(\varrho=0,\ \nu=-1\)). Si può supporre che la superficie sia trasformata in guisa da non contenere curve siffatte, e perciò è lecito scartare l’ipotesi \(\varrho=0\), e prendere
\[
\varrho\geq1.
\]

Ma se \(\varrho=1\), non potendo aversi \(\nu>0\) (altrimenti \(p=0\)) sarà \(\nu=0\), e allora si trova che le curve bicanoniche e pluricanoniche di \(|iK|\) avranno con \(D\) zero intersezioni, mentre la \(D\) stessa ha sempre zero intersezioni colle residue di \(|iK-D|\), \(|iK-2D|\) ecc.; da ciò segue che le curve di \(|iK|\) sono composte colle curve ellittiche d’un fascio e il genere lineare \(p^{(1)}=1\), contro la nostra ipotesi (\(p^{(1)}>1\)).

In conclusione, se le curve bicanoniche sono sconnesse, una parte \(D\) di esse, connessa in sè ed avente un numero \(\leq0\) d’intersezioni colle residue, dovrà ritenersi di genere \(\varrho>1\) e di grado
\[
\nu\geq\frac{4}{3}(\varrho-1)>\varrho-1.
\]

Ma questa conclusione è assurda, contrastando coll’ipotesi che \(D\) sia parte della parte fissa \(L\) delle curve di \(|2K|\), giacchè la dimensione del sistema completo \(|D|\) d’indice di specialità \(p\), secondo il teorema di Riemann-Roch, vale
\[
p+\nu-(\varrho-1)-p>0.
\]

Dunque le curve del sistema bicanonico \(|2K|\) e perciò anche le curve \(K\) del sistema canonico, sono connesse: una parte qualsiasi di una \(K\), di genere \(\varrho\) e grado \(\nu\), avrà colla residua un numero pari \(>0\) di punti comuni:
\[
2\varrho-2-2\nu\geq2,
\]
e così l’\textit{ordine di connessione delle \(K\)} è \(\geq2\).

3) Dall’essere connesse le curve canoniche pure \(K\) (per \(p>0\), \(p^{(1)}>1\)) si deduce che il sistema bicanonico \(|2K|=|K'|\) è regolare. Similmente resultano connesse tutte le curve pluricanoniche, e quindi ancora regolari i relativi sistemi
\[
|K_i|=|\{(i-1)K\}'|.
\]

\edsection{20. Nota sulla irriducibilità del sistema bicanonico.}{20. Nota sulla irriducibilità del sistema bicanonico.}{sec.IV.20}

Dopo avere dimostrato che sopra una superficie regolare con \(p>0\) e \(p^{(1)}>1\) \textit{il sistema bicanonico puro} è regolare, si pone la domanda se esso sia sempre irriducibile, e --- poichè abbiamo escluso la riducibilità delle sue parti variabili --- si tratta di \textit{dimostrare che non ha parti fisse} (fuori delle curve eccezionali).

% scan: s093R
% !PRINTED 163
\weberpage{163}
% head: IL GENERE NUMERICO E IL TEOREMA DI RIEMANN-ROCH, ECC.
% stamp: Faint vertical scan line in the right margin and scattered speckles.
Ma la dimostrazione urta in difficoltà d’ordine delicato, quando le curve canoniche si suppongono comunque formate di componenti semplici e multiple. Ci limiteremo a dimostrare il teorema nel caso in cui \textit{le curve canoniche} (pure) \(K\) \textit{siano formate di componenti semplici.}

In tal caso sappiamo che una \(K\) rimane connessa quando si tolga da essa una componente (irriducibile) \(K_1\) avente colla curva residua \(K_2\) il minimo numero \(n\) (\(\geq2\)) di punti comuni.

Ora, designando con \(\pi_1\) e \(\pi_2\) i generi di \(K_1\) e \(K_2\) rispettivamente, avremo il genere di \(K\):
\[
\pi=\pi_1+\pi_2+n-1.
\]

Quindi il distacco di \(K_1\) da \(|K|\) diminuisce di \(r_1+1=\pi_1+n-1>0\) la dimensione del sistema regolare aggiunto a \(|K|\), \(r_1\) designando la dimensione della serie segata da \(|K'|\) su \(K_1\). Ciò significa che la \(K_1\) non può essere parte fissa di \(|K'|=|K_1+K'_2|\), chè altrimenti togliendo la detta parte si otterrebbe il sistema \(|K'_2|\) colla medesima dimensione.
% The repeated plus signs at the printed line break are retained.

Ora --- tolta la componente \(K_1\) --- le curve di \(|K_2|\) restano connesse, ma non più (in generale) di ordine di connessione \(\geq2\): esse potranno avere dei \textit{manici}, cioè delle componenti aventi \textit{un} punto comune colla curva residua, le quali --- ove sieno di genere 0 --- costituiranno parti fisse di \(|K'_2|\). Tuttavia questi manici non sono, in ogni caso, parti fisse di \(|K'|\); e ciò perchè altrimenti la serie \(g_{2\pi_1-2+n}\) segata da \(|K'|\) su \(K_1\) avrebbe dei punti fissi e non sarebbe completa, mentre deve avere la dimensione
\[
r_1=\pi_1-2+n.
\]

Dunque si è riconosciuto che, per essere le \(K\) curve connesse d’ordine di connessione \(\geq2\), il sistema aggiunto \(|K'|\) non può contenere come parti fisse nè la \(K_1\), nè i manici della curva residua \(K_2\). Ora, spogliando una \(K_2\) di tali manici, si avrà una curva connessa \(\overline{K}_2\), d’ordine di connessione \(\geq2\), e si dimostrerà similmente che non sono parti fisse di \(|\overline{K}'_2|\), nè una componente toglibile senza rompere la connessione della curva, nè i manici della curva residua. Così proseguendo risulta che nessuna delle componenti di \(K\) può essere parte fissa di \(|2K|=|K'|\). \hfill c. d. d.

\edsection{21. Generi di una superficie spezzata.}{21. Generi di una superficie spezzata.}{sec.IV.21}

Agli sviluppi di questo capitolo vogliamo aggiungere qualche indicazione in ordine alle superficie spezzate. Se una superficie \(F\) di un certo ordine \(m\), che vogliamo supporre dotata di singolarità
%CONT>>
% scan: s094L
% !PRINTED 164
\weberpage{164}
% head: CAPITOLO QUARTO
%<<CONT
% stamp: Right-edge scan line and scattered speckles.
normali, viene a spezzarsi in due superficie \(F_1\) ed \(F_2\) rispettivamente d’ordine \(m_1\) ed \(m_2\), si domanda di determinare i caratteri numerici di \(F\) in funzione di quelli di \(F_1\) e \(F_2\) e della curva ad esse comune.

Conviene precisare che la superficie limite \(F_1+F_2\) avrà una curva doppia limite di quella di \(F\), che sarà composta di tre parti: una parte d’un certo ordine \(\delta_1\), doppia per \(F_1\) e non appartenente ad \(F_2\), una seconda parte di cui indicheremo l’ordine con \(\delta_2\), doppia per \(F_2\) e non appartenente ad \(F_1\), ed infine una terza parte --- diciamo d’ordine \(\delta\) --- appartenente ad \(F_1\) ed \(F_2\). All’infuori di questa si avrà una curva \(C\) di un certo genere \(\pi\), intersezione di \(F_1\) ed \(F_2\), che stabilisce la connessione fra le due superficie \(F_1\) ed \(F_2\), la cui somma viene pensata come limite di \(F\).

Ora si tratta di determinare i caratteri numerici della superficie spezzata \(F=F_1+F_2\) in funzione dei caratteri analoghi delle due superficie componenti, e di quelli della curva comune \(C\) attraverso cui sono connesse.\origfootnote{1}{s094L-1}{Che debba esistere una curva di connessione. d’ordine maggiore di zero, risulta dal principio di degenerazione applicato alle curve sezioni piane di \(F\). Cfr. \textsc{Enriques-Chisini}, Op. cit.}. % ?? uncertain: dots after "connesse" and "connessione" -- printed full stops, scan speckles.

Parliamo anzitutto del genere aritmetico. Indicheremo semplicemente con \(p\) (eguale \(p_a\)) il genere aritmetico di \(F\); con \(p_1\) e \(p_2\) rispettivamente i generi aritmetici di \(F_1\) ed \(F_2\), mentre si è detto \(\pi\) il genere della curva \(C\).

È chiaro che si può calcolare il \(p\) valendoci delle formule di postulazione relative alle curve doppie di \(F_1\) ed \(F_2\), o meglio alle curve d’ordine \(\delta_1\delta_2\) e \(\delta\) che costituiscono il limite della curva doppia di \(F\) quando questa degeneri in \(F_1\) ed \(F_2\). In questa maniera si può pervenire alla formula % [sic] \delta_1\delta_2 -> \delta_1,\delta_2, the three previously distinguished degrees are being enumerated.
\[
p=p_1+p_2+\pi,
\]
che costituisce l’estensione alle superficie della formula di Noether che dà il genere di una curva spezzata\origfootnote{2}{s094L-2}{In generale, per le varietà di qualunque dimensione: cfr. \textsc{F. Severi}, \textit{Fondamenti per la geometria sulle varietà algebriche.} Rendic. Circolo Mat. di Palermo, 1909.}.

Ma possiamo giustificare la formula senza sviluppare il calcolo sopra indicato.

Infatti le superficie \(\Phi_{m_1+m_2-4}\) d’ordine \(m_1+m_2-4\) aggiunte alla \(F_1+F_2\) (per cui la \(C\) non è da ritenere come curva doppia) segheranno sopra \(F_1\) il sistema aggiunto a \(|C|\) che ha la dimensione virtuale:
\[
p_1+\pi-1,
\]
e per ciascuna di queste curve si avranno in generale infinite superficie \(\Phi_{m_1+m_2-4}\); la dimensione virtuale di questo sistema di \(\Phi_{m_1+m_2-4}\)
%CONT>>
% scan: s094R
% !PRINTED 165
\weberpage{165}
% head: IL GENERE NUMERICO E IL TEOREMA DI RIEMANN-ROCH, ECC.
%<<CONT
% stamp: Scattered scan speckles.
sarà \(p_2\), essendovi \(\infty^{p_2-1}\) \(\Phi_{m_1+m_2-4}\) spezzate nella \(F_1\) e in una superficie \(\Phi_{m_2-4}\) aggiunta alla \(F_2\). In tal guisa si trova la dimensione del sistema delle \(\Phi_{m_1+m_2-4}\) aggiunte ad \(F_1+F_2\):
\[
p-1=p_1+p_2+\pi-1.
\]

A questo discorso si potrebbe obiettare che le superficie \(\Phi\) aggiunte ad \(F_1+F_2\), per il fatto di essere obbligate a passare per una curva d’ordine \(\delta_2\) non appartenente ad \(F_1\), potrebbero segare su \(F_1\) un sistema di dimensione inferiore a quella del sistema \(|C'|\) aggiunto a \(|C|\). Ma, trattandosi di dimensioni \textit{virtuali}, questa circostanza non ha effetti sul nostro calcolo, a cui potrebbe sostituirsi un calcolo analogo su superficie aggiunte d’ordine più elevato. % ?? uncertain: subscript on the second F in "F_1+F_2" -- 2, a dash-like damaged character.

Concludiamo pertanto che il \textit{genere numerico di una superficie spezzata le cui componenti di genere \(p_1\) e \(p_2\) siano connesse mediante una curva comune di genere \(\pi\), vale}
\[
p=p_1+p_2+\pi.
\]

Proponiamoci ora di valutare il genere lineare \(p^{(1)}\) di \(F_1+F_2\).

A tal uopo occorre conoscere i generi lineari \(p_1^{(1)}\) e \(p_2^{(1)}\) di \(F_1\) ed \(F_2\), il genere della curva di connessione \(C\) che già abbiamo designato con \(\pi\), e i gradi \(n_1\) ed \(n_2\) che questa curva possiede rispettivamente sopra \(F_1\) e sopra \(F_2\).

\textit{Il genere lineare della superficie spezzata vale}
\[
p_1^{(1)}+p_2^{(1)}+8\pi-n_1-n_2-9.
\]

Infatti, come abbiamo osservato, le superficie d’ordine \(m_1+m_2+4\) aggiunte ad \(F_1+F_2\) segano su \(F_1\) curve \(C'\) aggiunte a \(C\) che sono di genere % [sic] m_1+m_2+ +4 -> m_1+m_2-4, the same adjoint order is printed with -4 elsewhere on this page.
% The repeated plus signs at the printed line break are retained.
\[
p_1^{(1)}+3\pi-3-n_1;
\]
similmente le curve che le stesse \(\Phi\) segano su \(F_2\) sono di genere
\[
p_2^{(1)}+3\pi-3-n_2.
\]

Ora una curva del primo sistema su \(F_1\) ed una curva del secondo su \(F_2\) costituiranno insieme l’ intersezione di \(F_1+F_2\) con una \(\Phi_{m_1+m_2-4}\) allorchè abbiano a comune un gruppo canonico di \(2\pi-2\) punti costituenti le loro intersezioni con \(C\), sicchè il genere delle curve composte di queste due sarà
\[
p^{(1)}=p_1^{(1)}+p_2^{(1)}+8\pi-n_1-n_2-9;
\]
\hfill c. d. d.
% scan: s095L
% !PRINTED none
\weberblank{[s095L]}
% head: none
% blank page
% stamp: Right-edge scan line and scattered speckles.


\clearpage\thispagestyle{plain}
% scan: s095R
% !PRINTED none
\weberpage{[s095R]}
% head: none
% stamp: Handwritten "167" in the upper-right margin; scattered scan speckles.
\edchapter{V}{INVARIANTI NUMERICI E PIANI MULTIPLI}{Capitolo V. INVARIANTI NUMERICI E PIANI MULTIPLI}{INVARIANTI NUMERICI E PIANI MULTIPLI}

\edsection{1. L’invariante di Zeuthen-Segre.}{1. L’invariante di Zeuthen-Segre.}{sec.V.1}

In ciò che precede si sono introdotti gli invarianti delle superficie che vengono definiti in relazione al sistema canonico, seguendo un procedimento che appare naturale estensione di quello che conduce a definire il genere di una curva. Ma quest’ultimo procedimento può essere esteso in sensi diversi, e si riesce quindi a definire altri caratteri numerici delle superficie che hanno del pari un significato invariante e che possono essere adoperati con vantaggio negli sviluppi ulteriori della teoria, sebbene non siano caratteri essenzialmente nuovi, potendosi esprimere per mezzo del genere numerico \(p_a\) e del genere lineare \(p^{(1)}\), come avremo occasione di spiegare.

Sopra la superficie \(F\) si consideri un fascio lineare di curve irriducibili di genere \(\pi\geq0\), aventi un certo numero \(n\) di punti base effettivi (semplici o multipli), e si designi con \(\delta\) il numero delle curve del fascio dotate di un punto doppio, fuori dei punti base.

\textit{L’espressione}
\[
I=\delta-n-4\pi,
\]
\textit{formata col numero dei punti base, col genere e col numero delle curve dotate d’un punto doppio d’un fascio lineare, non dipende dalla scelta arbitraria del fascio, e costituisce perciò un invariante} (relativo) \textit{della superficie.}

Il significato invariante di questo carattere è stato riconosciuto anzitutto da \textsc{Zeuthen}\origfootnote{1}{s095R-1}{\textsc{H. G. Zeuthen}, \textit{Etudes géométriques etc.} Math. Ann., IV, 1871.} riferendosi ad un fascio di sezioni piane di una superficie. \textsc{C. Segre}\origfootnote{2}{s095R-2}{\textsc{C. Segre}, \textit{Intorno ad un carattere delle superficie ecc.} Atti Acc. Scienze di Torino, 1896.} ha dimostrato direttamente l’eguaglianza dei caratteri analoghi in ordine a due fasci qualunque appartenenti alla superficie. E perciò \(I\) viene designato di solito come \textit{invariante di Zeuthen-Segre.}

% scan: s096L
% !PRINTED 168
\weberpage{168}
% head: CAPITOLO QUINTO
% stamp: Scattered scan specks and a dark binding-edge artefact on the right.
Esponiamo la dimostrazione di Segre, confrontando due fasci in relazione affatto generica, \(|C_1|\) e \(|C_2|\), presi sulla superficie.

Il primo fascio avrà un certo numero \(n_1\) di punti base che vogliamo supporre semplici e distinti; il secondo avrà similmente un certo numero \(n_2\) di punti base, diversi dai primi e pure semplici e distinti. Designeremo rispettivamente con \(\pi_1\) e \(\pi_2\) i generi delle curve generiche dei due fasci, e con \(s\) il numero delle intersezioni \((C_1C_2)\). E considereremo la curva \(L\) luogo dei punti di contatto delle curve \(C_1\) e \(C_2\), cioè la jacobiana dei due fasci.

Una curva \(C_1\) incontra \(L\) nei punti del gruppo jacobiano della serie \(g_s^1\) segata da \(|C_2|\) su \(C_1\), e perciò in \(2s+2\pi_1-2\) punti, fuori degli \(n_1\) punti base che pure appartengono ad \(L\). Similmente una \(C_2\) sega \(L\) in \(2s+2\pi_2-2\) punti variabili.

Ora sulla curva \(L\), il cui genere designeremo con \(\pi\), avremo due serie \(\infty^1\), rispettivamente d’ordine \(2s+2\pi_1-2\) e \(2s+2\pi_2-2\), che designeremo con \(g_1\) e \(g_2\). Consideriamo il gruppo jacobiano di \(g_1\) che è formato di
\[
4s+4(\pi_1-1)+2\pi-2
\]
punti.

Si osservi che i punti doppi della serie \(g_1\) su \(L\) vengono dati da:

1) i \(\delta_1\) punti doppi del fascio \(|C_1|\); infatti uno di questi punti è doppio per il gruppo segato dalla relativa \(C_1\) su \(L\);

2) i \(\tau\) punti ove una \(C_1\) e una \(C_2\) hanno un contatto tripunto, perchè ivi si hanno due punti infinitamente vicini di una \(C_1\) comuni ad \(L\);

3) e infine gli \(n_2\) punti base del fascio \(C_2\); infatti uno di questi punti figura come punto fisso per la serie segata dalle \(C_2\) sulla \(C_1\) che passa per esso, e quindi conta due volte nello jacobiano di questa serie, sicchè costituisce un semplice contatto della \(C_1\) con la \(L\).

Ciò posto avremo: % ?? uncertain: punctuation after "avremo" -- colon or an isolated scan speck.
\[
4s+4(\pi_1-1)+2\pi-2=\delta_1+\tau+n_2. % ?? uncertain: the small mark above \pi -- scan speck or prime; read as a speck.
\]

Similmente, scambiando \(|C_1|\) con \(|C_2|\) si otterrà
\[
4s+4(\pi_2-1)+2\pi-2=\delta_2+\tau+n_1.
\]

Sottraendo si deduce
\[
\delta_1-n_1-4\pi_1=\delta_2-n_2-4\pi_2,
\]
\hfill c. d. d.

\textit{Osservazione.} -- Abbiamo supposto che i punti base di ciascuno dei due fasci \(|C_1|\) e \(|C_2|\) fossero semplici e distinti, ma possiamo affrancarci da queste restrizioni.

Pongasi per esempio che un punto \(O\) base di \(|C_2|\) sia un punto multiplo d’ordine \(r\) maggiore di uno, a tangenti variabili. Allora è
%CONT>>
% scan: s096R
% !PRINTED 169
\weberpage{169}
% head: INVARIANTI NUMERICI E PIANI MULTIPLI
%<<CONT
% stamp: Scattered scan specks and fine vertical streaks near the right edge.
facile mostrare che questo punto avrà la molteplicità \(2r-1\) per la jacobiana \(L\): si verifica infatti che la \(L\) completata con due curve \(\overline{C}_1\) e \(\overline{C}_2\), appartenenti rispettivamente ai due fasci, costituisce la curva jacobiana della rete cui appartengono i due fasci
\[
\overline{C}_2+|C_1|\qquad \text{e}\qquad \overline{C}_1+|C_2|,
\]
e perciò avrà in \(O\) la molteplicità \(3r-1\); togliendo la \(C_2\) resta che \(L\) ha in \(O\) la molteplicità \(2r-1\). Ora quando nel ragionamento precedente si va a cercare la molteplicità con cui \(O\) figura appartenere al gruppo di punti sezione di \(L\) con la \(C_1\) che passa per \(O\), si troverà che questa molteplicità è \(2r\) e quindi che \(O\) appartiene \(2r\) volte al gruppo jacobiano della \(g_s^1\) sopra la \(L\); di conseguenza il punto \(O\) viene sempre a figurare una volta nel computo dell’ordine della serie jacobiana di \(g_1\).

Per quel che concerne il caso in cui uno dei nostri fasci possegga punti base infinitamente vicini, si può dire che la nostra formula vale sempre conformemente ad una esigenza di continuità, purchè si tenga conto delle curve del fascio che vengono ad acquistare un punto doppio (diminuendo il genere \(\pi\)) in un punto base del fascio stesso, o in un punto ad esso infinitamente vicino. Basterà riferirci al caso in cui il fascio \(|C_1|\) abbia un punto base semplice \(O\) in cui le sue curve si tocchino: allora il fascio contiene una curva avente in \(O\) un punto doppio e una seconda curva (infinitamente vicina alla precedente) che ha un punto doppio infinitamente vicino ad \(O\): queste due curve debbono figurare nel computo di \(\delta\), che sarà dunque eguale al numero delle curve del fascio dotate di un punto doppio fuori dei punti base, aumentato di 2.

Dopo ciò possiamo valutare l’equivalenza di \textit{una curva \(C_1\) dotata di un punto \(r\)-plo}, fuori dei punti base, mostrando che essa \textit{tien luogo in generale di \((r-1)^2\) curve dotate di punto doppio.} Si confronti il fascio \(|C_1|\) a cui appartiene la curva dotata del punto \(r\)-plo con un secondo fascio generico \(|C_2|\). Anzitutto, come già si è osservato innanzi, il punto \(r\)-plo, \(P\), sarà \((r-1)\)-plo per la jacobiana dei due fasci \(|C_1|\) e \(|C_2|\). E sopra ogni ramo di questa jacobiana il punto \(P\) figurerà come \((r-1)\)-plo per la serie segata dalle \(C_2\); sicchè, ripetendo il ragionamento precedente, \(P\) verrà a contare \((r-1)^2\), volte nel computo dell’invariante \(I\). % [sic] "(r-1)^2, volte" -> "(r-1)^2 volte", the comma separates the multiplier from "volte".

In luogo di una curva dotata di un punto multiplo, può trovarsi in un fascio \(|C|\) una curva che contenga tutta una componente doppia o multipla: per esempio una
\[
C^{(0)}=2\chi+\psi,
\]
dove la \(\chi\) abbia un certo genere \(\varrho\geq0\) e la residua curva \(\psi\) (che può anche mancare) incontri la \(\chi\) in un certo numero \(i\geq0\) di punti.
%CONT>>
% scan: s097L
% !PRINTED 170
\weberpage{170}
% head: CAPITOLO QUINTO
%<<CONT
% stamp: Scattered scan specks and a vertical binding-edge line on the right.
Non è difficile valutare l’equivalenza di questa curva \(2\chi+\psi\), cioè il numero delle curve \(C\) dotate di punto doppio che vengono da essa assorbite.

Questo calcolo si trova sviluppato in una memoria di \textsc{Castelnuovo-Enriques}\origfootnote{1}{s097L-1}{\textsc{Castelnuovo-Enriques}, \textit{Sopra alcune questioni fondamentali ecc.} Ann. di Mat., 1901.}. Lo esporremo qui nel caso più semplice in cui non interviene la complicazione dei punti base infinitamente vicini, riferendoci perciò all’ipotesi che il fascio non possegga punti base, che è il caso in cui ci converrà in seguito applicare la formula in relazione ai fasci irrazionali.

Paragonando il fascio lineare \(|C|=|C_1|\) ad un altro fascio generico \(|C_2|\), si osserverà anzitutto che la \(\chi\) si stacca dalla jacobiana \(L\) dei due fasci, e la curva \(L_1\) parte residua di questa, passerà per gli \(i\) punti comuni a \(\chi\) e \(\psi\) che sono tripli per la curva \(C^{(0)}\). Seguendo un procedimento analogo a quello tenuto nel caso generale, si dovrà cercare l’ordine delle serie segata su \(L_1\) dalle curve dei fasci \(|C_1|\) e \(|C_2|\). % [sic] "serie segata" -> "serie segate", agreement with the plural "delle serie".

Intanto l’ordine della serie \(g_{m_1}^1\) segata dalle curve \(C_1\) sulla \(L_1\) è
\[
m_1=2s+2\pi_1-2\qquad \text{con}\qquad s=(C_1C_2),
\]
giacchè tante sono le \(C_2\) che toccano una \(C_1\) generica.

Per determinare l’ordine della serie \(g_{m_2}^1\) segata su \(L_1\) dalle \(C_2\), occorre tener presente che il gruppo jacobiano della serie \(g_s^1\) segata dalle \(C_1\) su una \(C_2\) generica comprende anche il gruppo dei punti d’intersezione della \(C_2\) stessa colla \(\chi\). Indicando con \(\sigma\) il numero di questi punti \((C_2\chi)\), si avrà che le curve \(C_1\) toccanti propriamente una \(C_2\) generica sono
\[
m_2=2s+2\pi_2-2-\sigma,
\]
e questo è anche il numero delle intersezioni \((C_2L)\). Ora, indicando con \(P\) il genere della jacobiana \(L_1\), il gruppo jacobiano della serie \(g_{m_1}^1\) su \(L_1\) sarà composto di \(2m_1+2P-2\) punti, e questo stesso gruppo è formato da: % [sic] "(C_2L)" -> "(C_2L_1)", the preceding count concerns the residual curve L_1.

1) i \(\bar{\delta}\) punti doppi delle restanti curve nodate che appartengono al fascio \(|C_1|\);

2) i \(\tau\) contatti tripunti di una \(C_1\) con una \(C_2\);

3) gli \(n_2\) punti base del fascio \(|C_2|\);

4) gli \(i\) punti comuni a \(\chi\) e \(\psi\), ciascuno dei quali, essendo triplo per il gruppo segato dalla \(C^{(0)}=2\chi+\psi\) su \(L_1\) deve essere contato due volte;

% scan: s097R
% !PRINTED 171
\weberpage{171}
% head: INVARIANTI NUMERICI E PIANI MULTIPLI
% stamp: Scattered scan specks and faint vertical streaks near the lower right edge.
5) i \(2\sigma+2\varrho-2\) punti doppi della serie segata dal fascio \(|C_2|\) sulla \(\chi\), attesochè la \(L_1\) passa per questi punti, che sono doppi per il gruppo segato dalla \(C^{(0)}\).

Si avrà dunque
\[
2m_1+2P-2=\bar{\delta}+\tau+n_2+2i+2\sigma+2\varrho-2.
\]

Invece, per la serie segata dalle curve \(C_2\) su \(L_1\) si ha semplicemente
\[
2m_2+2P-2=\delta_2+\tau,
\]
mancando i punti base per il fascio \(|C_1|\).

Eliminando da queste due relazioni \(2P-\tau-2\) e sostituendo i valori di \(m_1\) e \(m_2\) si ottiene
\[
\bar{\delta}-4\pi_1+2\varrho-2+2i=\delta_2-n_2-4\pi_2.
\]

Ma se il fascio \(|C_1|\) fosse dotato solo di curve aventi punti doppi isolati, e si indicasse con \(\delta_1\) il loro numero, si avrebbe
\[
I=\delta_1-4\pi_1.
\]

Dal confronto segue che \textit{una curva \(C^{(0)}\) dotata di una componente doppia di genere \(\varrho\) che intersechi in \(i\) punti la curva residua conta per}
\[
\delta_0=2\varrho-2+2i
\]
\textit{curve dotate di un punto doppio}, nel calcolo dell’invariante di Zeuthen-Segre, almeno nell’ipotesi che il \textit{fascio non sia dotato di punti base.}

Il ragionamento che precede si estende al caso in cui il fascio \(|C|=|C_1|\) contenga una curva
\[
C^{(0)}=h\chi+\psi,
\]
contenente una componente multipla secondo \(h\) (maggiore di 1), priva di punti doppi, la quale incontri in \(i\) punti la residua curva \(\psi\).

In questo caso si avrà
\[
\delta_0=(h-1)(2\varrho-2)+hi.
\]

Più in generale, si supponga che il fascio \(|C|\) contenga una curva spezzata
\[
C^{(0)}=h_1\chi_1+h_2\chi_2+\ldotp\ldotp\ldotp\ldotp+h_t\chi_t
\]
costituita di \(t\) componenti \(\chi\) semplici o multiple \((h_r\geq1)\), e che le curve \(\chi_r\) e \(\chi_s\) abbiano \(i_{rs}\) punti a comune. L’equivalenza di questa % ?? uncertain: the subscript on i -- rs or r,s.
% EDITORIAL D0093 RISOLTO: The two subscript letters are contiguous in the scan and no comma is present.
%CONT>>
% scan: s098L
% !PRINTED 172
\weberpage{172}
% head: CAPITOLO QUINTO
%<<CONT
% stamp: Scattered scan specks and a dark binding-edge artefact on the right.
curva ai fini del calcolo dell’invariante di Zeuthen-Segre è dato dalla \textit{formula di Godeaux}\origfootnote{1}{s098L-1}{L. \textsc{Godeaux}, \textit{Sur le calcul de l’invariant de Zeuthen-Segre.} Mémoires de la Société des Sciences du Hainaut, t. 66, 1920. Cfr. Bull. de l’Académie Royale de Belgique, 13 ott. 1934. V. anche L. \textsc{Campedelli}, \textit{Sul computo dell’invariante di Zeuthen-Segre per una superficie algebrica.} - \textit{Ancora sul computo ecc.} Rend. Acc. Lincei, serie 6, t. XIX, 1934.}. % [sic] "dato" -> "data", agreement with the subject "L’equivalenza".
\[
\delta_0=\Sigma(h_r-1)(2\varrho_r-2)+\Sigma i_{rs}(h_r+h_s-1) % ?? uncertain: the subscript on i -- rs or r,s.
% EDITORIAL D0094 RISOLTO: The displayed Godeaux formula visibly has the contiguous subscript rs, without a separating comma.
\]

Questa formula è suscettibile di corrispondere a tutti i casi, quando si tenga conto di due modificazioni:

1) Se \textit{una delle curve, \(\chi_1\), possiede un punto doppio} o multiplo, \(\delta_0\) deve essere debitamente accresciuto: se si tratta di un punto doppio e si assuma \(\varrho_1\) come \textit{genere effettivo}, a \(\delta_0\) si deve aggiungere \(2h_1-1\); se invece con \(\varrho_1\) si rappresenta \textit{il genere virtuale} di \(\chi_1\), il \(\delta_0\) va aumentato soltanto di un’unità.

2) Nel caso che il fascio \(|C|\) abbia punti base infinitamente vicini sopra una componente multipla \(\chi_1\) della \(C^{(0)}\), si possono eliminare codesti punti mediante una trasformazione della superficie che li muti in curve eccezionali: ciò equivale ad aggiungere alle \(\chi_r\) uno o più intorni di punti base. \textit{Con questa convenzione la formula di Godeaux risponde anche al caso di punti base infinitamente vicini.}

\edsection{2. Fascio irrazionale.}{2. Fascio irrazionale.}{sec.V.2}

L’invariante di Zeuthen-Segre relativo ad una superficie \(F\) può calcolarsi non soltanto in relazione ai fasci lineari tracciati su \(F\), sì anche a fasci irrazionali che ad \(F\) appartengono.

Precisamente l’\textit{invariante \(I\) si esprime in relazione ad un fascio di genere \(\varrho>0\) costituito di curve di genere \(\pi\) che contenga \(\Delta\) curve dotate di punto doppio mediante la formula}
\[
I=\Delta+4(\varrho-1)(\pi-1)-4\ \text{\origfootmark{2}{s098L-2}}.
\]\origfoottext{2}{s098L-2}{Cfr. \textsc{Castelnuovo-Enriques}, l. c.}

Questa formula si può dimostrare direttamente, riprendendo il confronto di due fasci \(|C_1|\) e \(|C_2|\) col sostituire ad uno dei due il fascio irrazionale \(\{C\}\) di cui qui si discorre. Ma si può anche ricondurre questa dimostrazione al caso di due fasci lineari, sostituendo a \(\{C\}\) un fascio lineare costituito dai gruppi di \(n\) curve \(C\) di una \(g_n^1\), dove \(n\) è un intero sufficientemente grande, per esempio uguale a \(\varrho+1\). Invero la circostanza che uno dei due fasci sia composto di curve riducibili non infirma i ragionamenti svolti nel precedente paragrafo.

% scan: s098R
% !PRINTED 173
\weberpage{173}
% head: INVARIANTI NUMERICI E PIANI MULTIPLI
% stamp: Scattered scan specks and vertical scan streaks near the right edge.
Adunque i gruppi di \(n\) curve \(C\) della \(g_n^1\) predetta formeranno un fascio lineare di curve, di genere
\[
\Pi=n\pi-n+1
\]
privo di punti base; e in questo fascio si avranno:

1) \(\Delta\) curve contenenti una \(C\) con punto doppio e

2) \(2n+2\varrho-2\) curve contenenti una \(C\) di genere \(\pi\) contata due volte, ciascuna delle quali equivale a \(2\pi-2\) curve dotate di punto doppio;

quindi si avrà
\[
I=\Delta+(\pi-2)(2n+2\varrho-2)-4\Pi=\Delta+4(\varrho-1)(\pi-1)-4, % [sic] "(\pi-2)" -> "(2\pi-2)", the preceding sentence gives the multiplicity as 2\pi-2.
\]
\hfill c. d. d.

\edsection{3. Espressione dell’invariante di Zeuthen-Segre per mezzo dei generi.}{3. Espressione dell’invariante di Zeuthen-Segre per mezzo dei generi.}{sec.V.3}

Si è detto che l’\textit{invariante \(I\) di Zeuthen-Segre è un invariante relativo della superficie \(F\)}; infatti esso \textit{cresce di uno quando un punto \(O\) di \(F\) si muti in una curva eccezionale.} L’asserzione si verifica in due modi: qualora si prenda un punto \(O\) di \(F\) che non sia punto base e sia semplice per una curva del fascio \(|C|\), mutando \(O\) in una curva eccezionale, crescerà di uno il numero \(\delta\) delle curve dotate di punto doppio; invece se \(O\) cade in un punto base del fascio \(|C|\), la trasformazione diminuisce di uno il numero dei punti base.

Osserviamo ora che una trasformazione della superficie \(F\) sostituente ad un punto una curva eccezionale fa variare in senso inverso l’invariante \(I\) e il genere lineare \(p^{(1)}\) della superficie, giacchè sappiamo che quest’ultimo decresce di una unità. Perciò la somma
\[
I+p^{(1)}
\]
costituirà un \textit{invariante assoluto} della superficie. Ma non si tratta di un invariante nuovo, poichè \textsc{Noether} ha dimostrato che esso si esprime per il genere numerico mediante la formula
\[
I+p^{(1)}=12p_a+9\ \text{\origfootmark{1}{s098R-1}}.
\]\origfoottext{1}{s098R-1}{\textsc{Noether}, op. c., Math. Ann., VIII.}

Questa relazione si collega ad un gruppo di formule che mettono in evidenza diversi caratteri numerici invariantivi delle superficie. La dimostreremo qui seguendo la via indicata da T. \textsc{Bonnesen}\origfootnote{2}{s098R-2}{P. \textsc{Bonnesen}, \textit{Sur les séries linéaires triplement infinies de courbes algébriques sur une surface algébrique.} Bull. de l’Académie Royal des Sciences et des Lettres de Danemarque, 1906, n. 4.}. % [sic] "P. Bonnesen" -> "T. Bonnesen", the main text gives the initial T.; "Académie Royal" -> "Académie Royale", adjective agreement in French.

% scan: s099L
% !PRINTED 174
\weberpage{174}
% head: CAPITOLO QUINTO
% stamp: Scattered scan specks and a vertical binding-edge line on the right.
Riferiamoci ad un sistema lineare irriducibile semplice \(\infty^3\) di curve \(|C|\), che possiamo ritenere sezioni piane di una superficie \(F\) dotata di una curva doppia \(D\); porremo per semplicità che la \(D\) sia irriducibile e dotata soltanto di \(t\) punti tripli che siano pure tripli per la \(F\). L’invariante \(I\) della superficie \(F\) si può calcolare in relazione ad un fascio di sezioni piane \(C\), e vien dato da
\weq{(1)}{I=\delta-n-4\pi,}
dove \(n\) è l’\textit{ordine} di \(F\), \(\pi\) il genere delle \(C\) e \(\delta\) designa la \textit{classe} della superficie \(F\). La classe di \(F\) è pur data dal grado effettivo del sistema \(|C_j|\) segato su \(F\) delle polari \(\Phi_{n-1}\); questo sistema possiede come punti base semplici i \(\tau\) punti cuspidali di \(F\) sopra \(D\), sicchè \(\delta+\tau\) eguaglia il grado di
\[
|C_j|=|3C+K|
\]
calcolato a prescindere dai \(\tau\) punti base suddetti. Introducendo il genere lineare \(p^{(1)}\) (genere delle curve canoniche impure \(K\)) avremo quindi
\weq{(2)}{\delta+\tau=3n+12(\pi-1)+p^{(1)}-1.}

Ora al sistema \(\infty^3\) \(|C|\) si legano alcuni caratteri (caratteri proiettivi della superficie \(F\)) e in particolare i seguenti:

\(\alpha\)) l’ordine \(d\) della curva doppia \(D\) che dipende dal grado \(n\) e dal genere \(\pi\) di \(|C|\):
\weq{(3)}{d=\frac{(n-1)(n-2)}{2}-\pi;}

\(\beta\)) il genere \(\varrho\) della curva doppia \(D\);

\(\gamma\)) il numero \(t\) dei punti tripli di \(D\), che sono anche tripli per la superficie \(F\); inoltre, chiamando \(\Delta\) la curva che corrisponde alla curva doppia \(D\) su un modello \(F'\) privo di singolarità della superficie \(F\):

\(\delta\)) il genere (effettivo) \(\Pi\) della curva delle coppie neutre di \(|C|\), curva \(\Delta\) che viene rappresentata sulla \(D\) doppia con \(\tau\) punti di diramazione nei punti cuspidali. Occorre notare che la curva \(\Delta\) delle coppie neutre appare possedere \(t\) terne di punti doppi in corrispondenza ai \(t\) punti tripli della \(D\), e perciò il \textit{genere virtuale} di \(\Delta\) è
\[
\Pi+3t.
\]

Fra i caratteri anzidetti oltre le (1) e (2) intercedono diverse relazioni da cui appare che oltre i due caratteri \(n\) e \(\pi\) del sistema \(|C|\), due soli caratteri della superficie sono suscettibili di assumere valori indipendenti e, fino a un certo punto, arbitrari.

Anzitutto osserviamo che la curva \(\Delta\) covariante del sistema \(\infty^3\) \(|C|\) risulta essere sopra la superficie \(F'\) (trasformata di \(F\) ove è
%CONT>>
% scan: s099R
% !PRINTED 175
\weberpage{175}
% head: INVARIANTI NUMERICI E PIANI MULTIPLI
%<<CONT
% stamp: Scattered scan specks and irregular vertical streaks near the right edge.
sciolta la curva doppia \(D\)) una curva parziale del sistema \(|(n-4)C|\); più precisamente, sommando a \(\Delta\) una curva canonica \(K\) si ottiene
\[
|\Delta+K|=|(n-4)C|.
\]

Perciò le intersezioni di \(\Delta\) con le curve del sistema aggiunto \(|\Delta+K|\) formeranno un gruppo canonico sopra la curva \(\Delta\) di genere virtuale \(\Pi+3t\), e si avrà quindi
\weq{(4)}{2(\Pi+3t)-2=2(n-4)d,}
dove \(d\) è espresso della formula (2). % [sic] "della formula (2)" -> "dalla formula (3)", the preposition is anomalous and the preceding formula for d is labelled (3).

D’altra parte, poichè la curva doppia \(D\) è in corrispondenza \((1,2)\) con la \(\Delta\) e si hanno su essa \(\tau\) punti di diramazione, avremo, per una nota formula di Zeuthen,
\weq{(5)}{2(2\varrho-2)+\tau=2\Pi-2.}

Finalmente possiamo procurarci una sesta relazione, valutando il numero delle intersezioni di una \(C_j\), jacobiana di una rete di \(|C|\), con la \(\Delta\), o, ciò che è lo stesso, con la \(D\). A tal uopo osserviamo che
\[
|C_j|=|3C+K|\qquad \text{e}\qquad |\Delta+K|=|(n-4)C|;
\]
per conseguenza il numero delle intersezioni \((C_j\Delta)=(C_jD)\) è dato da
\[
(3C+K)[(n-4)C-K]=2n^2-7n+15+2\pi(n-7). % [sic] "2n^2-7n+15+2\pi(n-7)" -> "2n^2-7n+15+2\pi(n-7)-p^{(1)}", the subsequent equation (6) retains the -p^{(1)} term.
\]

Ma questo stesso numero \((C_jD)\) si può calcolare direttamente sulla superficie \(F\), cercando il numero \(\alpha\) delle intersezioni di \(D\) con la seconda polare \(\Phi_{n-2}\) di un punto generico \(O\); infatti la \(C_j\) intersezione della polare di \(O\) incontra la curva doppia \(D\): nei \(\tau\) punti cuspidali e negli \(\alpha\) punti comuni a \(D\) e a \(\Phi_{n-2}\) (fuori dei punti tripli di \(D\)) in ciascuno dei quali il piano tangente ad una falda di \(F\) va a passare per \(O\). Ordunque avremo % ?? uncertain: the intersection-count symbol, transcribed as Latin a throughout this page -- a or \alpha.
% EDITORIAL D0095 RISOLTO: Enlargement shows the Greek alpha form consistently in both prose occurrences and both displayed formulas. The corresponding Latin-a transcription should be changed in these four places.
\[
(C_jD)=\alpha+\tau,
\]
dove
\[
\alpha=(n-2)d-3t.
\]

Per conseguenza possiamo scrivere
\weq{(6)}{2n^2+n+14+2\pi(n-7)-p^{(1)}=(n-2)d-3t.} % [sic] "n+14" -> apparently "-7n+15-\tau", comparison with the preceding expressions for (C_jD) and a gives a different linear and constant part.

Alle sei relazioni che abbiamo scritte conviene aggiungerne una settima che introduce il genere \(p_a\): si ha
\[
\begin{aligned}
p_a&=\frac{(n-1)(n-2)(n-3)}{6}-(n-4)\left[\frac{(n-1)(n-2)}{2}-\pi\right]+2t+\varrho-1=\\
&=\frac{1}{6}(n-1)(n-2)(9-2n)+\pi(n-4)+2t+\varrho-1.
\end{aligned}
\]

% scan: s100L
% !PRINTED 176
\weberpage{176}
% head: CAPITOLO QUINTO
% stamp: Scattered scan specks and a dark binding-edge artefact on the right.
Abbiamo trovato sette equazioni che legano gli undici caratteri: \(n,\pi,\delta,I,d,\varrho,\Pi,\tau,t,p_a,p^{(1)}\).

Fra questi \(n\) e \(\pi\) possono assumersi fino ad un certo punto in modo arbitrario, dipendendo dalla scelta di un fascio \(|C|\) sopra la superficie \(F\); pertanto è chiaro a priori che i sette caratteri
\[
\delta,\quad I,\quad \varrho,\quad \Pi,\quad \tau,\quad t % [sic] "\delta,I,\varrho,\Pi,\tau,t" -> "\delta,I,d,\varrho,\Pi,\tau,t", only six entries are printed after "i sette caratteri".
\]
potranno esprimersi in funzione dei detti \(n\) e \(\pi\) e degli invarianti \(p_a\) e \(p^{(1)}\) della superficie. In particolare si dovrà esprimere \(I\) come funzione di \(p_a\) e \(p^{(1)}\): a priori si potrebbe supporre che questa funzione dipenda altresì da \(n\) e \(\pi\), ma in fatto dovrà riuscire indipendente da questi, trattandosi di un carattere invariante di \(F\).

Non resta che eseguire il calcolo, e poichè le (1) e (2) contengono già una espressione esplicita di \(I\) e di \(d\), si tratta di risolvere cinque equazioni lineari, in cui figurano come parametri \(n\), \(\pi\), \(p_a\) e \(p^{(1)}\). Si ottiene così: % [sic] "(2)" -> "(3)", the expression for d on page 174 is labelled (3).
\[
\begin{aligned}
\delta&=n+4\pi-p^{(1)}+12p_a+9\\
\varrho&=\frac{1}{2}(n^2-7n+48)+(n-12)\pi+9p_a-2p^{(1)}\\
\Pi&=n^2-6n+2(n-10)\pi+12p_a-3p^{(1)}+36\\
t&=\frac{1}{6}(n^3-9n^2+26n-78)-(n-8)\pi-4p_a+p^{(1)}\\ % ?? uncertain: the exponent on the first n in the formula for t -- 3 or 2.
% EDITORIAL D0096 RISOLTO: The enlarged superscript has two rounded right lobes and no diagonal/base of a 2: the scan prints n cubed.
\tau&=2n+8\pi-12p_a+2p^{(1)}-22.
\end{aligned}
\]

Dalla prima di queste formule si ricava l’espressione di Noether dell’invariante di Zeuthen-Segre:
\[
I=12p_a-p^{(1)}+9.
\]

\textit{Nota.} -- Le quattro formule che danno \(\varrho\), \(\Pi\), \(t\) e \(\tau\), mettono in evidenza quattro invarianti relativi costruiti coi caratteri d’un sistema lineare \(\infty^3\) \(|C|\) appartenente alla superficie. In particolare è degna di nota l’ultima formula, che contiene linearmente \(n\) e \(\pi\); essa definisce l’invariante
\[
2n+8\pi-\tau=12p_a-2p^{(1)}+22.
\]

Combinando questa formula colla prima si otterrà un altro invariante relativo molto semplice \origfootnote{1}{s100L-1}{Cfr. \textsc{Bonnesen}, l. c.}, cioè
\[
2\delta-\tau=36p_a-4p^{(1)}+40.\ \text{\origfootmark{2}{s100L-2}}
\]\origfoottext{2}{s100L-2}{Cfr. \textsc{Bonnesen}, l. c.}

% scan: s100R
% !PRINTED 177
\weberpage{177}
% head: INVARIANTI NUMERICI E PIANI MULTIPLI
% stamp: Scattered scan specks and fine vertical streaks near the right edge.
Questo invariante è costruito a partire da due caratteri inerenti ad un sistema lineare \(\infty^3\) \(|C|\) tracciato sopra la superficie: \(\delta\) designa, in generale, il numero delle curve \(C\) di un fascio dotate di un punto doppio (fuori dei punti base) e \(\tau\) indica il numero dei fasci contenuti nel sistema \(\infty^3\) \(|C|\) che posseggono un punto base doppio (fuori dei punti base di \(|C|\)).

\textit{Osservazione.} -- L’espressione dell’invariante di Zeuthen-Segre per mezzo del genere numerico e del genere lineare conduce a ritrovare il \textit{genere lineare del piano e delle rigate}, che abbiamo calcolato nel Cap. II, \S\ 6.

Anzitutto si può calcolare l’invariante \(I\) del piano, riferendosi ad un fascio di coniche, per cui si ha:

il genere \(\pi=0\),

il numero dei punti base \(n=4\),

e il numero delle curve dotate di punto doppio \(\delta=3\).

Avremo
\[
I=\delta-n-4\pi=-1;
\]
quindi dalla formula
\[
I=12p_a-p^{(1)}+9
\]
essendo \(p_a=0\), si ricava
\[
p^{(1)}=10.
\]

Passiamo ora alle rigate irrazionali di genere \(p_a=-p\) \((p>0)\). Qui si può calcolare il valore di \(I\) in relazione al fascio delle rette generatrici, che è di genere \(p\) e per cui
\[
\pi=0,\quad \Delta=0,\quad n=0.
\]

Si avrà
\[
I=\Delta+4(p-1)(\pi-1)-4=-4p.
\]

Ora dalla formula noetheriana si deduce
\[
12p_a-p^{(1)}+9=-4p
\]
ed essendo \(p_a=-p\),
\[
p^{(1)}=-8(p-1)+1.
\]

Allo stesso risultato si perviene riferendosi invece che al fascio irrazionale delle generatrici ad un fascio lineare di sezioni piane.

\edsection{4. Curve cuspidate di una rete.}{4. Curve cuspidate di una rete.}{sec.V.4}

Riprendiamo il sistema lineare \(\infty^3|C|\) e i caratteri di esso considerati nel precedente paragrafo. Le formule stabilite ci permet%

% footer: Printed signature line.
% EDITORIAL: segnatura originale: {\small\noindent \textsc{Enriques F.} - \textit{Superficie algebriche.}\hfill 12\par}
\printersignature{\textsc{Enriques F.} - \textit{Superficie algebriche.}\quad 12}
%CONT>>
% scan: s101L
% !PRINTED 178
\unskip\weberpage{178}%
% head: CAPITOLO QUINTO
%<<CONT
% stamp: Dark scan-edge line in the right margin and scattered scan specks.
tono di determinare i caratteri appartenenti ad una rete di curve $|C|$, e in particolare il numero $\chi$ delle curve di una rete che sono dotate di una cuspide.

A tal uopo riferiamoci alla superficie $F$ che ha per sezioni piane le curve $C$. Il luogo dei punti di $F$ per cui la sezione piana tangente possiede una cuspide è la \textit{curva parabolica} $C_h$ di $F$, sezione della $F$ stessa con la superficie Hessiana $H$. E quindi le $\chi$ sezioni piane dotate di cuspide che passano per un punto $O$ dello spazio, corrisponderanno alle intersezioni di $C_h$ con la curva $C_j$, jacobiana della rete segata dai piani per $O$, fuori dei punti cuspidali di $F$, che appartengono semplicemente alla $C_j$, e, come vedremo, sono doppi per la curva parabolica $C_h$.

Per valutare $\chi$ occorre pertanto:

1) determinare il comportamento della hessiana nella curva doppia nodale della superficie $F$;

2) trovare ulteriormente le molteplicità che la curva $C_h$ possiede nei $\tau$ punti cuspidali di $F$.

Il primo calcolo si effettua semplicemente sostituendo alla data una superficie, che torniamo momentaneamente a indicare con $F$, che abbia nell’origine un punto doppio biplanare, coi piani tangenti $x$ e $y$; scriviamo in coordinate omogenee $x$, $y$, $z$, $u$:
\[
F=u^{n-2}xy+u^{n-3}\varphi_3(x\,y\,z)+\ldotp\ldotp\ldotp\ldotp+\varphi_n(x\,y\,z).
\]

Lo hessiano di $F$ è
\[
H=
\begin{vmatrix}
\dfrac{\partial^2 F}{\partial x^2} &
\dfrac{\partial^2 F}{\partial x\partial y} &
\dfrac{\partial^2 F}{\partial x\partial z} &
\dfrac{\partial^2 F}{\partial x\partial u} \\[6pt]
\dfrac{\partial^2 F}{\partial y\partial x} &
\dfrac{\partial^2 F}{\partial y^2} &
\dfrac{\partial^2 F}{\partial y\partial z} &
\dfrac{\partial^2 F'}{\partial y\partial u} \\[6pt] % ?? uncertain: the small mark after F in row 2, column 4 -- a prime; an ink blemish.
\dfrac{\partial^2 F}{\partial z\partial x} &
\dfrac{\partial^2 F}{\partial z\partial y} &
\dfrac{\partial^2 F}{\partial z^2} &
\dfrac{\partial^2 F}{\partial z\partial u} \\[6pt]
\dfrac{\partial^2 F}{\partial u\partial x} &
\dfrac{\partial^2 F}{\partial u\partial y} &
\dfrac{\partial^2 F}{\partial u\partial z} &
\dfrac{\partial^2 F}{\partial u^2}
\end{vmatrix}
\]

Eseguendo il calcolo si trova che il comportamento dell’hessiana nell’origine delle coordinate $(x=y=z=0,\ u=1)$ resta definito dalla hessiana della superficie approssimante del terzo ordine, e così $H$ possiede nel punto biplanare di $F$ un punto quadruplo. Ritornando ora alla nostra superficie $F$ dotata di curva nodale, si deduce che la sua hessiana $H$ d’ordine $4n-8$ passa quattro volte per la detta curva nodale. Per conseguenza la curva parabolica $C_h$ apparterrà al sistema lineare
\[
|4K+8C|,
\]
designando al solito con $K$ le curve canoniche della superficie.

% scan: s101R
% !PRINTED 179
\weberpage{179}
% head: INVARIANTI NUMERICI E PIANI MULTIPLI
% stamp: Faint vertical scan artefacts in the right margin and scattered specks.
Ora si vede che i $\tau$ punti cuspidali di $F$ (semplici per le $C_j$) sono doppi per la curva parabolica $C_h$. Infatti, riferendoci ad una superficie trasformata su cui venga sciolta la curva nodale di $F$, un punto cuspidale di $F$ darà su tale superficie un punto base doppio per un fascio di $|C|$, nel quale si avranno dunque \textit{due} curve $C$ dotate di cuspide.

Ciò posto, il numero delle curve cuspidate di una rete entro $|C|$ verrà dato da
\[
\chi=(C_j\cdot C_h)-2\tau.
\]

Ma essendo
\[
|C_j|=|3C+K|,\qquad |C_h|=|8C+4K|,
\]
si avrà
\[
(C_j\cdot C_h)=24n+20(2\pi-2-n)+4p^{(1)}-4,
\]
e ricordando l’espressione del $\tau$ dal paragrafo precedente
\[
\tau=2n+8\pi-12p_a+2p^{(1)}-22
\]
segue
\[
\chi=24(p_a+\pi).
\]

Si arriva così alla relazione scoperta da \textsc{Zeuthen}: \textit{il numero delle curve cuspidate di una rete di genere $\pi$ è in generale $24(p_a+\pi)$.}

A dir vero noi abbiamo dimostrato questa formula riferendoci ad una rete di curve $C$ che sia contenuta totalmente in un sistema $\infty^3$. Ma è facile estendere il ragionamento fatto riferendosi ad un sistema lineare $\infty^3$ $|L+C|$ che contenga parzialmente $|C|$. E così la formula verrà giustificata in ogni caso\origfootnote{1}{s101R-1}{La dimostrazione della formula data nel testo si può esporre anche in una forma poco diversa, dove non si fa più uso delle proprietà dell’Hessiana. A tal uopo conviene associare alla superficie $F$ la sua duale $F'$ che è trasformata di $F$ ottenuta facendo corrispondere alle sezioni piane le jacobiane $C_j$ delle reti di $|C|$. Se ora si costruisce la superficie duale di $F'$ si dovrà ricadere sopra $F$. Ciò significa che le curve jacobiane delle reti di $C_j$ sopra $F$, diciamo le $C_{jj}$, debbono identificarsi con le sezioni piane $C$, a meno di una componente fissa, che si riconosce essere la curva parabolica $C_h$.

Infatti si considerino le polari rispetto ad $F$ dei punti di un piano $\alpha$, le quali formano una rete. È chiaro anzitutto che la sezione piana $C$ di $\alpha$ appartiene alla jacobiana $C_{jj}$ di codesta rete, perchè la polare di un punto $P$ di $C$ sega $F$ in una curva che passa doppiamente per $P$. Inoltre, se si assume un punto $A$ sulla curva parabolica $C_h$, il piano tangente in $A$ incontra $\alpha$ secondo una retta $a$, e le polari dei punti di $a$ segano $F$ secondo le curve di un fascio che ha in $A$ un punto base con tangente fissa, fra le quali ce n’è una dotata di un punto doppio in $A$; perciò la curva $C_h$ fa parte della seconda jacobiana $C_{jj}$.

Dalla relazione
\[
|C_{jj}|=|C_h+C|
\]
segue quindi
\[
|C_h|=|8C+4K|.
\]
}.

% scan: s102L
% !PRINTED 180
\weberpage{180}
% head: CAPITOLO QUINTO
% stamp: Dark scan-edge line in the right margin and scattered specks.
\edsection{5. Notizia storica.}{5. Notizia storica.}{sec.V.5}

La relazione del numero delle curve cuspidate di una rete col genere numerico è stata scoperta da \textsc{H. G. Zeuthen}\origfootnote{1}{s102L-1}{\textsc{H. G. Zeuthen}, \textit{Études géométriques etc.} Math. Ann., Bd. IV, 1871. Cfr. la memoria, pure citata, di \textsc{Noether}, dei Math. Ann., VIII.}. Supponendo nota l’invarianza del genere numerico (che \textsc{Zeuthen} e \textsc{Noether} riconoscevano, sia pure con qualche restrizione relativa alle singolarità della superficie), \textsc{Zeuthen} veniva così a provare che il numero $\chi-24\pi$, calcolato a partire da una rete (di sezioni piane), risulta indipendente da questa, e perciò costituisce un carattere intrinseco della superficie. \textsc{F. Severi} nel 1902 ha dato una dimostrazione diretta di questo teorema, con un procedimento affatto elementare\origfootnote{2}{s102L-2}{\textsc{F. Severi}, \textit{Il genere aritmetico ed il genere lineare in relazione alle reti di curve tracciate sopra una superficie algebrica.} Atti Acc. di Torino, vol. XXXVII, 1902.}.

La dimostrazione svolta nel testo si trova in \textsc{Bonnesen}, l. c.\origfootnote{3}{s102L-3}{Cfr. anche \textsc{M. Pannelli}, \textit{Sui sistemi lineari triplamente infiniti di curve tracciati sopra una superficie algebrica.} Rend. Circ. Mat. Palermo, t. XX, 1905. La considerazione relativa alla jacobiana di una rete di jacobiane, che abbiamo indicato in una nota precedente, s’incontra anche in \textsc{E. Caporali}.}.

Non possiamo chiudere questo cenno storico che completa le notizie date nei precedenti paragrafi, senza ricordare un ordine d’idee in cui si ricerca l’interpretazione geometrica funzionale delle formule numerative di cui innanzi si è discorso. È noto come questa veduta si sia introdotta nella teoria delle curve, col teorema di \textsc{Painlevè-Castelnuovo} (interpretazione della formula di Zeuthen relativa alle corrispondenze $[1,n]$) e colla considerazione della serie jacobiana di una serie lineare, in cui \textsc{Enriques} ha riconosciuto il modo più semplice di definire la serie canonica e dimostrarne l’invarianza. % ?? uncertain: the accent in Painlevè -- grave (Painlevè); acute (Painlevé).
% EDITORIAL D0098 RISOLTO: The enlarged small-cap E carries a grave accent (sloping down to the right), so retain the printed spelling Painlevè.

Per trovare una interpretazione analoga della geometria numerativa sopra le superficie, \textsc{Severi}, in successive memorie a partire dal 1932, ha cercato di costruire una teoria delle serie di gruppi di punti, che chiama \textit{di equivalenza}, la quale dovrebbe costituire una naturale estensione della teoria delle serie lineari sopra una curva.

Per un’ampia bibliografia sull’argomento vedi il libro di \textsc{F. Severi}: « Serie, sistemi di equivalenza e corrispondenze algebriche sulle varietà algebriche », Roma, ed. Cremonese, 1942.

\edsection{6. Caratteri di una rete e piani multipli.}{6. Caratteri di una rete e piani multipli.}{sec.V.6}

Le formule stabilite nel precedente paragrafo permettono anche di esprimere i caratteri di una rete di curve $|C|$, per mezzo di due
%CONT>>
% scan: s102R
% !PRINTED 181
\weberpage{181}
% head: INVARIANTI NUMERICI E PIANI MULTIPLI
%<<CONT
% stamp: Scattered scan specks.
di essi, p. es. del genere $\pi$ e del grado $n$, e degli invarianti della superficie: $p_a$ e $p^{(1)}$.

Sia $|C|$ una rete di curve irriducibili (di genere $\pi$ e grado $n$) priva di punti base sopra la superficie $F$; ad essa apparterranno in genere i seguenti caratteri numerici:

1) il numero $N$ delle \textit{intersezioni delle curve $C$ colla jacobiana della rete $C_j$}:
\[
N=2n+2\pi-2;
\]

2) il numero $\delta$ delle \textit{curve $C$ di un fascio dotate di punto doppio}, che sappiamo essere espresso da
\[
\delta=n+4\pi+12p_a-p^{(1)}+9;
\]

3) il \textit{genere} della detta \textit{jacobiana} $C_j$ $(=3C+K)$
\[
P=p^{(1)}+9\pi-9;
\]

4) il numero $k$ dei \textit{contatti tripunti} delle $C$ della rete (punti base d’un fascio di $C$ osculatrici);

5) il numero $i$ delle \textit{curve} della rete $|C|$ \textit{dotate d’una cuspide};

6) il numero $d$ dei \textit{doppi contatti} delle $C$ della rete (punti base d’un fascio di $C$ bitangenti);

7) il numero $t$ delle $C$ della rete dotate di due punti doppi.

Si può supporre che la rete delle $C$ sia segata su $F$ da una stella di piani, per un punto $O$, che sarà in generale multiplo per $F$ e rappresenterà una curva irriducibile $L$ (la quale sommata alle $C$ dà il sistema delle sezioni piane).

Allora la $F$ viene proiettata da $O$ sopra un \textit{piano multiplo $n$-plo}, i cui punti rispondono a gruppi di $n$ punti di $F$ appartenenti all’involuzione $I$ definita dalla rete $|C|$. E i caratteri anzidetti assumono il seguente significato:

1) $N$ è l’\textit{ordine della curva di diramazione $D$} del piano $n$-plo (curva luogo dei punti del piano cui rispondono gruppi di $I$ dotati d’una coincidenza);

2) $P$ è il \textit{genere della curva di diramazione $D$};

3) $\delta$ è la \textit{classe} della detta \textit{curva $D$};

4) $k$ è il numero delle rette per $O$ che hanno un contatto tripunto colla superficie $F$, e quindi il \textit{numero delle cuspidi} della $D$;

5) $i$ è il numero delle curve cuspidate, sezioni piane di $F$ per $O$, e quindi il \textit{numero dei flessi} della \textit{curva di diramazione $D$};

6) $d$ è il numero delle rette per $O$ bitangenti ad $F$ e quindi il \textit{numero dei nodi di $D$};

% scan: s103L
% !PRINTED 182
\weberpage{182}
% head: CAPITOLO QUINTO
% stamp: Dark scan-edge line in the right margin and scattered isolated spots.
7) $t$ è il numero dei piani per $O$ bitangenti ad $F$ e quindi il \textit{numero delle tangenti doppie} (con contatti distinti) \textit{della $D$}.

Questi caratteri sono legati dalle \textit{formule di Plücker} relative alla curva $D$, che in generale è irriducibile e non possiede altre singolarità, puntuali o tangenziali.

In tal guisa si perviene a stabilire le seguenti relazioni:
\weq{A)}{
\left\{
\begin{aligned}
N&=2n+2\pi-2\\
P&=p^{(1)}+9\pi-9\\
\delta&=n+4\pi+12p_a-p^{(1)}+9\\
k&=3(n+6\pi+p^{(1)}-4p_a-11)\\
i&=24(\pi+p_a)\\
d&=2[(n+\pi)^2-5n-17\pi-2p^{(1)}+6p_a+24]\\
t&=\frac{1}{2}[(n+4\pi+12p_a-p^{(1)})^2+15n-6\pi+132p_a-17p^{(1)}+74].
\end{aligned}
\right.
}

\textit{Queste formule esprimono in generale i caratteri di una rete irriducibile $|C|$ priva di punti base sopra la superficie $F$, ovvero i caratteri plückeriani della curva di diramazione $D$ del corrispondente piano multiplo, per mezzo del genere e del grado, $\pi$ ed $n$, di $|C|$ e degli invarianti $p_a$ e $p^{(1)}$ della superficie.}

Tali invarianti sono espressi per mezzo dei caratteri di $D$ dalle formule:
\weq{B)}{
\left\{
\begin{aligned}
p^{(1)}&=P-9\pi+9\\
p_a&=\frac{i}{24}-\pi\\
&\text{o anche}\\
4p_a&=n+P-3\pi-\frac{k}{3}-2\\
(\pi&=\frac{N}{2}-n+1).
\end{aligned}
\right.
}

\textit{Osservazione.} – Tra le formule $A)$ hanno particolare importanza la quarta e la quinta, che mettono in evidenza due caratteri invarianti della rete $|C|$ contenenti linearmente $n$ e $\pi$, cioè $3(n+6\pi-k)$ e $24\pi-i$. % [sic] 3(n+6\pi-k) -> 3(n+6\pi)-k, comparison with the fourth formula A).

Ma soprattutto è importante l’ultima di queste due formule che ci porge l’invariante assoluto $p_a$. Come si è detto questa relazione fra il $p_a$ e il numero delle curve cuspidate di una rete è stata scoperta da \textsc{G. H. Zeuthen}. Ed \textsc{F. Severi} ha dimostrato in modo diretto ed elementare che l’espressione $24\pi-i$ è eguale per tutte le reti prive di punti base appartenenti alla superficie $F$. Si ha così una nuova dimostrazione dell’invarianza del $p_a$. % [sic] G. H. Zeuthen -> H. G. Zeuthen, the initials appear in the opposite order on page 180.

% scan: s103R
% !PRINTED 183
\weberpage{183}
% head: INVARIANTI NUMERICI E PIANI MULTIPLI
% stamp: Scattered scan specks.
\edsection{7. Corrispondenza $(1,n)$ fra due superficie: piani doppi.}{7. Corrispondenza (1,n) fra due superficie: piani doppi.}{sec.V.7}

Si abbia una corrispondenza $(1,n)$ fra due superficie $f$ ed $F$: ad un punto di $F$ corrisponde \textit{un} punto su $f$, mentre ai punti di $f$ corrispondono su $F$ i gruppi di una \textit{involuzione} $I$ d’ordine $n$ (ogni punto di $F$ appartiene ad \textit{un} gruppo della $I$). Aggiungiamo per semplicità l’ipotesi che \textit{non} vi siano punti \textit{fondamentali}, cioè punti di $f$ o $F$ a cui risponda sull’altra superficie una curva.

Ad un sistema lineare irriducibile $|c|$, dato sopra $f$, risponderà su $F$ un sistema lineare $|C|$, appartenente all’involuzione $I$. Alla jacobiana $|c_j|$ di una rete contenuta in $|c|$, risponde su $F$ una curva che fa parte della jacobiana della rete delle $|C|$ omologhe. Ma per completare la jacobiana di questa occorre aggiungervi la curva delle coincidenze $D$, che corrisponde alla curva di diramazione sopra $f$; infatti, un punto di $D$ è un punto base di un fascio di $|C|$ dentro la suddetta rete, che ha una tangente fissa, e perciò contiene una $C$ dotata di punto doppio. Ciò posto, ricordando la definizione delle curve canoniche $(|K|=|C_j-3C|)$ si deduce il teorema:

\textit{Per effetto di una corrispondenza $(1,n)$ fra le superficie $f$ ed $F$, le trasformate delle curve canoniche di $f$, sommate alla curva omologa della curva di diramazione, costituiscono curve canoniche di $F$}\origfootnote{1}{s103R-1}{Cfr. \textsc{Enriques}, \textit{Ricerche ecc.}.}. % ?? uncertain: the punctuation at the end of the footnote -- ecc..; ecc.

Il teorema va modificato quando vi sia sopra $f$ \textit{un punto fondamentale} a cui risponda sopra $F$ una curva (ovvero una curva e un gruppo di punti); \textit{questa curva fondamentale va aggiunta alle immagini delle curve canoniche di $f$ e alla curva di coincidenza}, per formare una curva canonica di $F$.

Il teorema enunciato sussiste anche quando le \textit{curve canoniche} di $f$ siano \textit{virtuali}. Si può farne applicazione al problema delle curve canoniche e pluricanoniche di un \textbf{piano doppio}.

Si abbia una superficie $F$ contenente un’involuzione razionale $I$ di $2^\circ$ ordine, mercè cui viene rappresentata sul piano doppio
\[
z^2=f_{2n}(x,y)
\]
dove la curva di diramazione $f_{2n}$ d’ordine $2n$, potrà avere dei punti singolari, la cui molteplicità è da valutare secondo le convenzioni del Cap. III, § 7. Sul piano doppio si assume la curva $\varphi_m$ d’un sistema lineare che non abbia punti base fuori dei punti singolari di $f_{2n}$; a $\varphi_m$ risponde su $F$ una curva che ha per aggiunte le trasformate delle $\varphi_{m-3}$, aggiunte alla $\varphi_m$ nel piano, sommate alla curva di coincidenza omologa alla $f_{2n}$. Se, in luogo di una prima aggiunta, si considera una seconda aggiunta $\varphi_{m-6}$; alla trasformata di essa dovrà sommarsi la curva di coincidenza contata due volte, e così via.

% scan: s104L
% !PRINTED 184
\weberpage{184}
% head: CAPITOLO QUINTO
% stamp: Dark vertical scan-edge marks in the right margin and scattered specks.
Infatti sul piano le curve canoniche virtuali $K$ sono definite dalla sottrazione
\[
|K|=|\varphi_{m-3}-\varphi_m|,
\]
e perciò si otterranno curve canoniche sopra $F$ sommando le $K$ alla curva di coincidenza $f_{2n}$; sommando a queste le omologhe delle $\varphi_m$ si avranno dunque le aggiunte a queste ultime curve.

Applichiamo il risultato ottenuto assumendo al posto delle $\varphi_m$ la curva di diramazione $f_{2n}$ e le seconde aggiunte a questa: $f_{2n-6}$. Per avere su $F$ le seconde aggiunte alla curva di coincidenza contata due volte, che risponde alla $f_{2n}$, basterà sommare questa curva contata due volte alle trasformate delle $f_{2n-6}$. Sottraendo dal sistema così ottenuto la stessa curva di coincidenza contata due volte, si avranno le curve bicanoniche di $F$, fra cui sono quelle, appartenenti all’involuzione $I$, che hanno per immagini le $f_{2n-6}$.

Dunque: \textit{le curve seconde aggiunte alla curva di diramazione del piano doppio sono immagini di curve bicanoniche} appartenenti all’involuzione $I$. Questa deduzione è d’accordo con ciò che si è visto al Cap. III, § 7: che le curve canoniche della superficie sono rappresentate dalle $f_{n-3}$ appartenenti alla metà del sistema secondo aggiunto alla curva di diramazione.

Il ragionamento si estende alle curve pluricanoniche: \textit{si ottengono curve tricanoniche del piano doppio sommando le terze aggiunte della curva di diramazione a questa stessa curva}, e così via.

Ma ritorniamo alle curve bicanoniche, per osservare che, anche in altro modo, \textit{si ottengono curve bicanoniche del piano doppio sommando la curva di diramazione alle aggiunte delle curve $f_{n-3}$, immagini delle curve canoniche}; basta ricordare che il sistema bicanonico sega gruppi della serie canonica sulle curve del sistema canonico.

Ora si noti che il sistema bicanonico è trasformato in se stesso dall’involuzione $I$, mentre le sue curve non appartengono in generale a questa involuzione. Quindi dovremo considerare codesto sistema come uno spazio lineare di $P-1$ dimensioni ($P$ designando il bigenere di $F$) su cui la $I$ subordina un’omografia involutoria con due spazi di punti uniti $S_r$ e $S_t$, dove
\[
P-1=r+t+1.
\]

Uno dei due spazi, sia p. es. $S_r$, risponderà al sistema delle $f_{2n-6}$, seconde aggiunte della $f_{2n}$, mentre l’altro spazio $S_t$ corrisponderà al sistema formato dalla $f_{2n}$ di diramazione presa insieme con le $f_{n-6}$ aggiunte delle immagini delle curve canoniche. Quindi, designando con $\pi=t+1$ il genere delle $f_{n-3}$, risulterà
\[
P=r+\pi+1.
\]

\textit{Questa formula esprime il bigenere del piano doppio} in funzione del genere $(r+1)$ delle curve $f_{2n-3}$ aggiunte alla curva di diramazione
%CONT>>
% scan: s104R
% !PRINTED 185
\weberpage{185}
% head: INVARIANTI NUMERICI E PIANI MULTIPLI
%<<CONT
% stamp: Scattered scan specks.
e del genere $\pi$ delle curve metà delle biaggiunte $f_{2n-6}$, cioè del genere dell’involuzione di $2^\circ$ ordine che la $I$ subordina su una curva canonica. % ?? uncertain: the final vowel of delle before biaggiunte -- delle; dello.

Le cose dette si estendono facilmente ai plurigeneri. Ci limiteremo a rilevare che se i plurigeneri debbono essere tutti nulli, mancheranno tutti i successivi aggiunti alla curva di diramazione, e perciò questa curva si potrà trasformare birazionalmente in una curva d’ordine $2n$ con un punto $(2n-2)$-plo o in una quartica o in una sestica con due punti tripli infinitamente vicini, o in un caso particolare di queste\origfootnote{1}{s104R-1}{Teorema di \textsc{Castelnuovo-Enriques}, 1900. Cfr. \textsc{Enriques-Conforto}, \textit{Le superficie razionali}, § 27, pag. 489.}.

Pertanto si può affermare che:

\textit{Un piano doppio avente tutti i generi nulli è razionale ovvero riferibile ad una rigata iperellittica.}

\edsection{8. Formule di corrispondenza.}{8. Formule di corrispondenza.}{sec.V.8}

Il teorema che abbiamo dato nel paragrafo precedente in ordine alla trasformazione delle curve canoniche in una corrispondenza $(1,n)$ fra due superficie $f$ ed $F$, dice anzitutto che, in ogni caso, il genere $P_g$ di $F$ non può essere inferiore a quello di $f$:
\[
P_g\geq p_g.
\]

Inoltre lo stesso teorema conduce subito ad una relazione fra i generi lineari $P^{(1)}$ e $p^{(1)}$ delle due superficie.

Conviene premettere un’osservazione relativa alla curva di diramazione $D$ della corrispondenza sopra $f$. Essa sarà in generale dotata di un certo numero di nodi a cui rispondono sopra $F$ gruppi di $n$ punti con due coincidenze distinte. Inoltre la $D$ possiederà un certo numero di cuspidi; a ciascuna delle quali risponde un gruppo di $n$ punti di cui tre sono venuti a coincidere. Ciò posto, per un noto teorema di Zeuthen, ad una curva canonica $n$-pla su $f$ che incontri in $s$ punti la curva di diramazione $d$, risponde una curva di genere % [sic] D (twice) -> d, the branch curve is denoted d later in this paragraph, while D denotes the coincidence curve.
\[
n(p^{(1)}-1)+\frac{s}{2}+1.
\]

Questa curva incontra in $s$ punti la curva di coincidenza $D$, sicchè designando con $\varrho$ il genere effettivo\origfootnote{2}{s104R-2}{Si ammette che la curva $d$ non presenti altre singolarità che i nodi e le cuspidi che hanno il significato spiegato innanzi rispetto alla nostra corrispondenza.} di questa, si avrà
\[
P^{(1)}=n(p^{(1)}-1)+\varrho+\frac{3}{2}s.
\]

% scan: s105L
% !PRINTED 186
% [printed number read as 184 in the first pass; the spread sequence gives 186]
\weberpage{186}
% head: CAPITOLO QUINTO
% ?? uncertain: the page number -- 184 is visible in the overlapped area; the underlying page number is not securely readable.
% stamp: A large irregular white occlusion covers the upper-left text and several line beginnings. Displaced fragments "185", "pre" and "da", and a dark right-edge line, are visible.
\sourcewarning{La parte superiore sinistra della scansione è coperta. Nei tre capoversi seguenti si conservano soltanto i frammenti leggibili, senza integrare le lacune. Il richiamo della prima nota è collocato provvisoriamente. Il testo continuo riprende con «Si designi dunque».}
\textit{lega i generi lineari delle due superficie in cor-} In essa deve farsi $\varrho=1$ ed $s=0$ \textit{quando manchi} a una $f$ senza singolarità) \textit{la curva di dirama-}\origfootnote{1}{s105L-1}{\textsc{Enriques}, \textit{Ricerche ecc.}, 1893.} % ?? uncertain: covered beginnings, continuations after cor- and dirama-, and the first footnote mark -- no reliable reading of the covered portions is possible; only visible fragments are retained, and the note is attached here provisionally.

mente si può scrivere una seconda relazione numerici $P_a$ e $p_a$ di $F$ ed $f$\origfootnote{2}{s105L-2}{Cfr. \textsc{Severi}, \textit{Sulle relazioni che legano i caratteri invarianti di due superficie algebriche.} Rend. Ist. Lombardo, 1903.}. % ?? uncertain: the clipped suffix before si può -- mente; uente; the preceding text and the beginning of the line containing numerici are covered.

po si assumerà un fascio di curve $|c|$ sopra $f$ e il fascio urve corrispondenti $|C|$ sopra $F$. Per mezzo dei caratteri di esti due fasci si valuteranno gli invarianti di Zeuthen-Segre $i$ ed $I$ delle due superficie e si scriverà la relazione da cui essi sono legati. % ?? uncertain: clipped word beginnings -- po (possibly uopo); urve (possibly curve); esti (possibly questi); other covered text cannot be read.

Si designi dunque con $m$ il numero dei punti base di $|c|$, con $\pi$ il genere delle $c$, e con $\delta$ il numero delle $c$ dotate di punto doppio; avremo
\[
i=\delta-m-4\pi.
\]

Ora il numero $M$ dei punti base del fascio $|C|$, e il genere $\Pi$ delle $C$ sono dati rispettivamente da
\[
M=mn
\]
e
\[
\Pi=n(\pi-1)+1+\frac{t}{2}
\]
ove con $t$ si designa il numero delle intersezioni di una $c$ con la curva di diramazione $d$ su $f$. Quanto al numero $\Delta$ delle curve $C$ dotate di punto doppio, si troveranno anzitutto tra queste le $\delta$ curve che rispondono alle $c$ dotate di punto doppio, ognuna delle quali, avendo $n$ punti doppi, dovrà contare per $n$. Ma vi saranno in generale altre $\Delta-n\delta$ curve che si ottengono dalle $c$ che toccano la curva $d$, fatta eccezione per quelle che hanno un contatto improprio, passando per una cuspide della $d$, alla quale risponde su $F$ un punto in cui si riuniscono tre degli $n$ punti corrispondenti.

Pertanto, designando con $\varrho$ il genere effettivo della curva di diramazione $d$, e con $\chi$ il numero delle cuspidi di essa si avrà
\[
\Delta=n\delta+2t+2\varrho-2-\chi,
\]
quindi si avrà l’invariante di Zeuthen-Segre della $F$ espresso da
\[
I=\Delta-M-4\Pi=n\delta+2\varrho-2-\chi-nm-4n(\pi-1)-4.
\]

% scan: s105R
% !PRINTED 187
\weberpage{187}
% head: INVARIANTI NUMERICI E PIANI MULTIPLI
% stamp: Scattered scan specks.
Ciò posto, ricordando la relazione di Noether
\[
\begin{aligned}
i&=12p_a-p^{(1)}+9\\
I&=12P_a-P^{(1)}+9,
\end{aligned}
\]
si dedurrà
\[
24(P_a+1)=24n(p_a+1)+6\varrho-6+3s-2\chi.
\]

\textit{Questa relazione lega i generi numerici delle due superficie $f$ ed $F$ in corrispondenza} $(1,n)$, nell’ipotesi semplificativa che \textit{non} esistano \textit{punti fondamentali} della corrispondenza sopra $f$. Anche qui \textit{se manca} la curva di diramazione si assumerà
\[
\varrho=1,\qquad s=0.
\]

Le formule che legano i generi numerici e i generi lineari di due superficie $f$ ed $F$ in corrispondenza $(1,n)$ si possono anche scrivere in altra forma, introducendo al posto del numero $s$ che designa il numero delle intersezioni della curva di diramazione $d$ di $f$ con le curve canoniche $k$, il grado $\nu$ della curva $d$: grado virtuale che si esprime per mezzo di $s$ e del genere virtuale di $d$, $\varrho+\varepsilon+\chi$, dove con $\varepsilon$ si indica il numero dei nodi di $d$ e con $\chi$ il numero delle cuspidi. Si avrà
\[
s=2(\varrho+\varepsilon+\chi)-2-\nu
\]
e quindi \textit{le formule che legano i caratteri di $f$ e di $F$ diventeranno}:
\[
\begin{aligned}
P^{(1)}&=n(p^{(1)}-1)+4\varrho+3(\varepsilon+\chi)-\frac{3}{2}\nu-3\\
24(P_a+1)&=24n(p_a+1)+12\varrho-12+6\varepsilon+4\chi-3. % [sic] -3 -> -3\nu, substitution of the preceding expression for s leaves a factor \nu.
\end{aligned}
\]

La prima formula si può dedurre egualmente comparando i gradi del sistema canonico su $f$ e su $F$. Ricordando che si ottiene una curva canonica $K$ di $F$ sommando alla trasformata di una $k$ di $f$ la curva di coincidenza $D$ che la incontra in $s$ punti; indicando con $x$ il grado di $D$, avremo:
\[
P^{(1)}-1=n(p^{(1)}-1)+x+2s.
\]

Questa formula si accorda con la precedente, ove si calcoli i valore di $x$ che è % [sic] i valore -> il valore, the final l of the article is absent.
\[
x=\frac{\nu}{2}-(\varepsilon+\chi).
\]

In particolare, se si ha su $f$ una curva di diramazione o una componente di essa di grado $\nu$, che sia priva affatto di punti doppi, la curva di coincidenza che le corrisponde sopra $F$ ha il grado $\dfrac{\nu}{2}$. Ciò appare chiaro specialmente nel caso di una corrispondenza $(1,2)$, $n=2$, dove la curva di diramazione deve ritenersi in generale
%CONT>>
% scan: s106L
% !PRINTED 188
\weberpage{188}
% head: CAPITOLO QUINTO
%<<CONT
% stamp: Dark vertical scan/binding line along the right margin and scattered specks.
priva di punti doppi: a questa curva di grado \(\nu\) corrisponde una curva doppia di grado \(2\nu\), la cui componente ha dunque il grado \(\frac{\nu}{2}\).

È interessante confrontare le formule qui ottenute con quelle che porgono il genere numerico e il genere lineare dei piani multipli che abbiamo date nel \S~6.

Si debbono ritrovare queste formule tenendo conto che il genere lineare (relativo) del piano vale \(p^{(1)} = 10\)\origfootnote{1}{s106L-1}{Cfr. Cap. II, \S~6; Cap. V, \S~3.}, mentre il suo genere numerico è \(p_a = 0\).

A tal uopo si riprendano le notazioni adoperate nelle formule \(B)\) del \S~6, mettendo così \(P\) al posto di \(\varrho\), \(k\) al posto di \(\chi\), \(d\) al posto di \(\varepsilon\) e \(N^2\) al posto di \(\nu\), dove
\[
N = 2n + 2\pi - 2.
\]

Con queste posizioni, le nostre formule danno
\[
P^{(1)} = 9n + 4P + 3(d + k) - \frac{3}{2}N^2 - 3 % ?? uncertain: sign before 3/2 -- minus or division sign; a dot is visible above the stroke.
% EDITORIAL D0105 RISOLTO: Enlargement shows a minus bar and a small isolated upper speck, with no lower dot of a division sign. Preserve the minus reading.
\]
e poichè
\[
P = \frac{(N - 1)(N - 2)}{2} - (d + k),
\]
si deduce
\[
P^{(1)} = P - 9\pi + 9;
\]
che è la prima delle nostre formule \(B)\).

Similmente la formula che dà \(24(P_a + 1)\) diventa
\[
24(P_a + 1) = 24n + 12(P - 1) + 6d + 4k - 3N^2
\]
da cui
\[
4P_a = n + P - 3\pi - \frac{k}{3} - 2
\]
che è la seconda delle nostre formule \(B)\).

\textit{In questo modo si sono ritrovate le formule che danno i caratteri di un piano multiplo.}

\textit{Esempi.} -- Riprendiamo le notazioni di questo paragrafo, per applicare le formule ottenute ad alcuni esempi interessanti.

Si consideri la \textit{superficie di Kummer}, superficie singolare del complesso quadratico di rette dello spazio, che è definita come superficie del \(4^\circ\) ordine \(f_4\) dotata di 16 punti doppi\origfootnote{2}{s106L-2}{Cfr. \textsc{Enriques-Chisini}, \textit{Lezioni ecc.}, Libro VI, \S~40, vol. IV, pag. 246.}. È noto che la \(f_4\), che è superficie duale di sè stessa, possiede 16 piani tangenti doppi, che la toccano ciascuno lungo una conica, cui appartengono 6 punti doppi. Ciò risulta senz’altro dalla rappresentazione della \(f_4\) su un piano doppio con sestica di diramazione spezzata in 6 rette, quale si ottiene proiettando la \(f_4\) da un suo punto doppio. Ora si conside%
%CONT>>
% scan: s106R
% !PRINTED 189
\unskip\weberpage{189}%
% head: INVARIANTI NUMERICI E PIANI MULTIPLI
%<<CONT
% stamp: Scattered scan specks.
rino 3 piani tangenti doppi della \(f_4\) passanti per un punto doppio \(A\); questi s’intersecheranno a due a due secondo rette, che proiettano da \(A\) tre punti doppi, \(B\), \(C\), \(D\); quindi risulta che il tetraedro \(ABCD\) i cui vertici sono doppi per \(f_4\), contiene semplicemente i rimanenti 12 punti doppi della \(f_4\), tre su ciascuna faccia.

Ciò posto, designando con \(T(x, y, z) = 0\) l’equazione di codesto tetraedro, costruiamo nello spazio \(S_4(x\,y\,z\,u)\) la superficie \(F\) di equazione
\[
\begin{aligned}
&u^2 = T(x\,y\,z)\\
&(f_4(x, y, z) = 0),
\end{aligned}
\]
che è rappresentata sopra la \(f_4\) doppia.

Importa osservare che su questa superficie doppia non v’è curva di diramazione d’ordine maggiore di zero. Ma vi sono 16 curve di diramazione infinitesime, costituite dagli intorni dei 16 punti doppi, per cui \(T\) passa un numero dispari di volte. Di qui risulta che la nostra \(f_4\) doppia è irriducibile; e del resto ciò dipende anche dal criterio di Comessatti\origfootnote{1}{s106R-1}{Ricordiamo che in base al criterio di Comessatti la condizione necessaria e sufficiente affinchè una curva irriducibile \(C\), tangente ovunque l’incontra alla curva di diramazione \(D_{2m}\), sia immagine di una curva spezzata è che il gruppo dei punti di contatto appartenga alla serie segata su \(C\) dalle curve \(D_m\) metà della curva di diramazione.}, osservando che su ciascuna quartica sezione piana doppia di \(f_4\), vi sono 8 punti critici apparenti che cadono nei punti di contatto di un quadrilatero, e non appartengono ad una conica. La corrispondenza \((1, 2)\) che intercede tra \(f_4\) ed \(F\) conduce dal sistema delle sezioni piane di \(f_4\) che ha il genere \(\pi = 3\) e il grado \(n = 4\), ad un sistema lineare di genere 5 e grado \(8 = (2 \cdot 5 - 2)\) sopra \(F\). Si deduce che la \(F\) è, come la \(f_4\) doppia, di genere
\[
p_g = 1
\]
ed ha, come quella, curva canonica d’ordine zero, e quindi tutti i plurigeneri eguali ad uno.

A priori il genere lineare (assoluto) di \(F\) sarà dunque \(P^{(1)} = 1\). Ma valutiamolo con la nostra formula:
\[
P^{(1)} = n(p^{(1)} - 1) + 4\varrho + 3(\varepsilon + \chi) - \frac{3}{2}\nu - 3.
\]

A tal uopo si osservi che la curva di diramazione \(d\) su \(f_4\) è costituita dai 16 intorni dei 16 punti doppi, ciascuno dei quali è una curva infinitesima di genere zero e grado \(-2\) \((\varepsilon = \chi = 0)\). Troveremo dunque
\[
P^{(1)} = -4 \cdot 15 + \frac{3}{2} \cdot 32 - 3 = -15.
\]

% scan: s107L
% !PRINTED 190
\weberpage{190}
% head: CAPITOLO QUINTO
% stamp: Dark vertical scan/binding line along the right margin and scattered specks.
Questa conclusione paradossale riesce chiarita se si pensa che a ciascuna curva di diramazione di genere zero e grado \(-2\), risponde a priori una curva di genere zero e grado \(-1\), che è eccezionale, e che nel caso nostro appare come l’intorno di un punto semplice di \(F\). Per costruire la curva canonica di \(F\) noi abbiamo sommato alla trasformata della curva canonica d’ordine zero di \(f_4\), la trasformata della curva di diramazione, per conseguenza abbiamo valutato il \(P^{(1)}\) tenendo conto degli intorni di 16 punti semplici come se questi fossero curve eccezionali di \(F\) agli effetti del calcolo del genere lineare. In altre parole abbiamo valutato il genere lineare di una superficie trasformata di \(F\) con 16 curve eccezionali. Eliminando queste curve si trova, come deve essere:
\[
P^{(1)} = 1.
\]

Procedendo ancora a valutare il genere numerico di \(F\), si avrà
\[
24(P_a + 1) = 96 - 12 \cdot 15 - 12 + 3 \cdot 32 = 0
\]
ossia
\[
P_a = -1.
\]

\textit{Nota.} -- Come è stato dimostrato da \textsc{F. Klein}\origfootnote{1}{s107L-1}{Cfr. \textsc{Enriques-Chisini}, l. c.}, la superficie di Kummer \(f_4\) si può rappresentare come una involuzione sopra la superficie \(\Phi\) di cui gli elementi sono le coppie di punti di una curva \(C\) di genere 2; precisamente, si prendono come coniugate su \(C\) le coppie di punti che sono residue l’una dell’altra rispetto ad una \(g_4^2\). In tal guisa si definisce su \(\Phi\) un’involuzione \(I\) che ha 16 punti di coincidenza, in corrispondenza alle 16 coppie di \(C\) provenienti dalla bisezione della \(g_4^2\). Ora la \(I\) essendo rappresentata dalla \(f_4\), si ha che la \(\Phi\) viene rappresentata da questa \textit{superficie di Kummer doppia, con \(16\) punti di diramazione, che cadono nei suoi \(16\) punti doppi.} Pertanto si deduce che la superficie \(\Phi\) equivale alla \(F\) di cui innanzi abbiamo scritto l’equazione. E così \textit{la superficie iperellittica rappresentante l’insieme delle coppie di punti di una curva di genere \(2\), ha i generi}
\[
p_g = 1,\qquad p_a = -1
\]
\textit{e tutti i plurigeneri uguali ad \(1\).} Avremo occasione di ritornare più avanti su tale superficie\origfootnote{2}{s107L-2}{Per la costruzione indicata innanzi della superficie iperellittica come superficie di Kummer doppia, con curva di diramazione definita dal tetraedro \(T\), cfr. \textsc{Enriques-Severi}, \textit{Mémoire sur les surfaces hyperelliptiques.}, Acta Math., 1909.}. % ?? uncertain: punctuation after "hyperelliptiques" -- ".,", ",", or ".".
% EDITORIAL D0106 RISOLTO: Enlargement shows a full stop after the final s followed by a smaller descending comma before Acta. The existing period-plus-comma transcription is retained.

Un altro esempio interessante, in ordine alla nostra teoria si ha considerando la superficie di genere \(p_a = p_g = 0\) e di bigenere
%CONT>>
% scan: s107R
% !PRINTED 191
\weberpage{191}
% head: INVARIANTI NUMERICI E PIANI MULTIPLI
%<<CONT
% stamp: Scattered ink/scan marks, including a dark spot beside the final display.
\(P = 1\), costituita dalla sestica \(f_6\) che passa doppiamente per gli spigoli di un tetraedro \(T\) (cfr. Cap. II, \S~8). Costruiamo la superficie \(F\) definita dalle equazioni
\[
u^2 = T(x, y, z),\qquad f_6(x, y, z) = 0,
\]
che è rappresentata sopra la \(f_6\) doppia. Poichè \(T\) ed \(f_6\) hanno comuni i 6 spigoli del tetraedro \(T\) contati 2 volte, sulla \(f_6\) non c’è curva di diramazione d’ordine \(> 0\). Ma non ci sono nemmeno curve di diramazione infinitesime, cioè punti di diramazione: a priori questi potrebbero cadere soltanto nei vertici di \(T\); ma uno di questi vertici \(A\) proviene su \(f_6\) dalla sovrapposizione di 3 punti semplici, i cui intorni appartengono ai piani facce del tetraedro passanti per \(A\). Ora, poichè \(A\) è \(4\)-plo per la sezione di \(f_6\) con una faccia \(\alpha\), i punti infinitamente vicini ad \(A\) su \(\alpha\) non saranno punti di diramazione della nostra \(f_6\) doppia. Ciò posto osserviamo che le curve \(C\) sezioni piane di \(f_6\) sono curve doppie irriducibili (criterio di Comessatti) e quindi rappresentano su \(F\) curve \(C\) di genere \(\pi = 7\) e di grado \(N = 2 \cdot 6 = 2\pi - 2\). Siccome il doppio di \(|C|\) coincide col doppio del sistema aggiunto: % The equals sign is repeated at the printed line break.
\[
|2C'| = |2C|,
\]
si deduce che su \(F\)
\[
|C'| = |C|
\]
e perciò il genere geometrico di \(F\) vale
\[
P_g = 1.
\]

La \(F\) possiede una curva canonica d’ordine zero, quindi il suo genere lineare
\[
P^{(1)} = 1
\]
e tutti i suoi plurigeneri \(P_i = 1\).

Confrontiamo i caratteri di \(F\) con quelli della superficie doppia \(f_6\) che ha
\[
p_a = p_g = 0,\qquad p^{(1)} = 1
\]
(cfr. \S~II, 8). Le nostre formule ci danno % [sic] \S II, 8 -> Cap. II, \S 8, the reference at the top of this page has that form.
\[
\begin{aligned}
&n = 2,\qquad \varrho = 1,\qquad \nu = 0,\qquad \varepsilon = \chi = 0,\\
&\qquad P^{(1)} = 2(p^{(1)} - 1) + 4\varrho - 3 = 1
\end{aligned}
\]
e
\[
24(P_a + 1) = 2 : 24(p_a + 1) = 48 % [sic] 2 : 24 -> 2 \cdot 24, the displayed result is 48.
\]
donde
\[
P_a = 1.
\]

\textit{La sestica \(f_6\) passante doppiamente per gli spigoli di un tetraedro che --- come semplice --- è una superficie regolare di genere zero e di}
%CONT>>
% scan: s108L
% !PRINTED 192
\weberpage{192}
% head: CAPITOLO QUINTO
%<<CONT
% stamp: Dark vertical scan/binding line along the right margin and scattered specks.
\textit{bigenere uno con curva bicanonica d’ordine zero, si può ritenere come una superficie doppia coi generi eguali ad \(1\) e con curva canonica d’ordine zero}\origfootnote{1}{s108L-1}{Cfr. \textsc{F. Enriques}, \textit{Un’osservazione relativa alle superficie di bigenere uno.} Rendic. Acc. Bologna, 1908.}.

\edsection{9. Teorema d’esistenza.}{9. Teorema d’esistenza.}{sec.V.9}

Ritorniamo alla considerazione dei piani multipli. Quando una superficie \(F\) con dati generi \(p_a\) e \(p^{(1)}\) venga rappresentata sopra un piano multiplo di un certo ordine \(n\), in corrispondenza ad una sua rete di curve di genere \(\pi\) e grado \(n\), abbiamo appreso a valutare i caratteri della curva di diramazione del piano \(n\)-plo per mezzo di \(n\), \(\pi\), \(p_a\) e \(p^{(1)}\) (cfr. \S~8).

Le formule così ottenute indicano già che la curva di diramazione di un piano \(n\)-plo, per \(n > 2\), non può assumersi ad arbitrio. A dir vero, data una qualunque curva piana
\[
f(x, y) = 0
\]
si può costruire il piano \(n\)-plo (\textit{ciclico})
\[
z^n = f(xy)
\]
che ha appunto, come curva di diramazione, la \(f = 0\).

Ma la curva di diramazione di questo piano \(n\)-plo non è costituita dalla sola \(f\) irriducibile, bensì dalla \(f\) contata \(n - 1\) volte, alla quale si aggiunge ancora la retta all’infinito, contata una o più volte, quando l’ordine \(m\) di \(f\) non sia multiplo di \(n\)\origfootnote{2}{s108L-2}{Per vederlo basta riferirsi ad una retta \(n\)-pla o meglio ad una riemanniana ad \(n\) fogli, che si costruisca assumendo nel piano della variabile complessa \(m\) punti di diramazione, a ciascuno dei quali risponda la sostituzione ciclica fra i rami \((1\,2\,3\ldots n)\). Un giro che avvolga tutti gli \(m\) punti di diramazione dà luogo ad una permutazione fra i rami espressa da \((1\,2\,3\ldots n)^m\), e questa è la sostituzione corrispondente ad un giro intorno al punto all’infinito, che pertanto riesce di fatto un punto di diramazione, quando la detta sostituzione non sia identica, cioè \(m\) non sia multiplo di \(n\).}. Così dunque la curva di diramazione di un piano ciclico \(n\)-plo, non è affatto una curva irriducibile arbitrariamente data.

Poichè la curva di diramazione di un piano \(n\)-plo per \(n > 2\) non può darsi ad arbitrio, nasce il problema di « determinare le condizioni perchè una curva piana \(f(xy) = 0\) sia curva di diramazione di un piano \(n\)-plo ». Queste condizioni condurranno ad un \textit{teorema d’esistenza} per i piani multipli o per le funzioni algebriche di due variabili, teorema che presenta un senso analogo, ma sotto un certo
%CONT>>
% scan: s108R
% !PRINTED 193
\weberpage{193}
% head: INVARIANTI NUMERICI E PIANI MULTIPLI
%<<CONT
aspetto diverso, dal teorema d’esistenza per le rette \(n\)-ple o per le funzioni algebriche di una variabile, dato da \textsc{Riemann}\origfootnote{1}{s108R-1}{Cfr. p. es. \textsc{Enriques-Chisini}, \textit{Lezioni.} Libro V, cap. III, \S~33 (vol. III, pag. 355).}. Infatti, nel caso di Riemann il gruppo dei punti di diramazione si può dare ad arbitrio, ed in corrispondenza ad esso si ha un numero finito di classi di funzioni algebriche (o di rette \(n\)-ple), ognuna delle quali dipende dall’assegnare le sostituzioni fra i rami, in relazione ad un sistema di cappi avvolgenti i punti di diramazione; qui invece, come si è detto, si tratta di trovare le condizioni particolari a cui la curva di diramazione deve soddisfare per l’esistenza del piano \(n\)-plo e, almeno in generale, quando queste condizioni siano soddisfatte, si avrà soltanto una classe\origfootnote{2}{s108R-2}{Cfr.: \textsc{O. Chisini}, \textit{Sulla identità birazionale delle funzioni algebriche di due variabili dotate di una medesima curva di diramazione.} Rend. Ist. Lombardo di Scienze e Lettere, 1944.} di superficie rappresentata sul piano, ossia un piano \(n\)-plo unico, cui risponde un determinato sistema di sostituzioni in relazione ai cappi avvolgenti la curva di diramazione, di cui diciamo in appresso.

Tuttavia, il significato algebrico della questione resta analogo a quello della costruzione di una funzione algebrica d’una sola variabile di cui sono dati i punti di diramazione. Quando si proietti la superficie \(F(xyz) = 0\) dal punto all’infinito \(Z\) dell’asse delle \(z\) sopra il piano \(z = 0\), si ha in generale in questo piano una curva luogo dei punti critici della \(z(xy)\) che si rappresenta annullando il discriminante di questa funzione,
\[
\Delta(xy) = 0
\]
cioè il risultante delle due equazioni
% stamp: A handwritten underline appears beneath part of "risultante".
\[
F(xyz) = 0,\qquad \frac{\partial F}{\partial z} = 0.
\]

Però codesta curva \(\Delta = 0\) si decompone in generale in due parti,
\[
\Delta = \Delta_1 \cdot \Delta_2^2,
\]
una delle quali, \(\Delta_2 = 0\), è una curva critica apparente, proiezione della curva doppia di \(F\), mentre l’altra \(\Delta_1 = 0\) (la cosidetta \textit{parte principale del discriminante}) costituisce propriamente la curva di diramazione della funzione algebrica \(z(xy)\). Il nostro problema consiste nell’ « assegnare le condizioni perchè una curva piana sia curva di diramazione di una funzione algebrica di due variabili » e quindi nella « costruzione di questa funzione algebrica con data curva di diramazione ». La posizione del problema stesso si può precisare assumendo che la curva piana \(f\), che vuolsi essere curva di
% footer: Printer's signature.
% EDITORIAL: segnatura originale: \vadjust{\hbox to \linewidth{\textsc{Enriques F.} - \textit{Superficie algebriche.}\hfil 13}}
\printersignature{\textsc{Enriques F.} - \textit{Superficie algebriche.}\quad 13}
%CONT>>
% scan: s109L
% !PRINTED 194
\weberpage{194}
% head: CAPITOLO QUINTO
%<<CONT
% stamp: Dark vertical scan/binding line along the right margin and scattered specks.
diramazione di un piano \(n\)-plo, sia una curva di ordine pari \(2m\)\origfootnote{1}{s109L-1}{Se fosse di ordine dispari le si aggiungerebbe necessariamente la retta all’infinito.}, proiezione da un punto \(O\) del contorno apparente \(C\) di una superficie \(F\), così che la \(f\) possegga un certo numero \(d\) di nodi ed un certo numero di cuspidi, che rispondano a punti doppi apparenti della detta curva \(C\): oltre a questi punti doppi la \(f\) può possedere nodi e cuspidi rispondenti a punti doppi isolati (rispettivamente punti conici e punti biplanari) della \(F\); ma da questi punti doppi di \(f\) si può prescindere ritenendoli come virtualmente inesistenti. % ?? uncertain: mark above the first C -- unmarked C with a speck, C', or \bar{C}.

Ciò premesso, consideriamo nel piano della curva data \(f\) un fascio di rette, il cui centro, per semplicità di discorso, possiamo supporre non appartenente ad \(f\), e sia per esempio il fascio \(y = tx\) col centro nell’origine \(O\). Per un valore generico di \(t\), si ha una retta \(n\)-pla su cui sono segnati \(2m\) punti di diramazione semplici; nel piano rappresentativo della variabile complessa \(x\) si assumono \(2m\) cappi uscenti dal punto \(O\) e avvolgenti codesti \(2m\) punti; si possono assegnare ad arbitrio \(2m\) trasposizioni sui rami corrispondenti ai detti cappi, con la condizione che il gruppo da esse generato sia transitivo e che il loro prodotto sia l’identità: in tal guisa si trova un numero finito di curve irriducibili, in generale birazionalmente distinte, che vengono rappresentate sulla retta \(n\)-pla \(y = tx\) e che sono egualmente di genere
\[
\pi = \frac{m}{2} - n + 1; % [sic] m/2 -> m, the preceding text specifies 2m simple branch points.
\]
queste curve sono definite a meno di trasformazioni birazionali, cioè rispondono a classi determinate di curve o funzioni algebriche, generalmente distinte.

Il numero \(r\) delle dette curve \(C_1\,C_2\,\ldotp\ldotp\ldotp\ldotp\,C_r\) è stato determinato da \textsc{Hurwitz}, ma qui non occorre conoscerlo. Importa invece tener presente che, quando si faccia variare \(t\) descrivendo un cammino chiuso nel piano \(\tau\) di questa variabile complessa, le curve \(C_1,\ C_2,\ \ldotp\ldotp\ldotp\ldotp,\ C_r\) (o le classi di curve cui esse appartengono) verranno in generale permutate fra loro: ciò avviene in primo luogo perchè vengono fra loro permutati i \(2m\) punti di diramazione intersezioni della retta \(y = tx\) con \(f\), ed anche perchè il sistema dei cappi, coi quali si costruisce per es. la riemanniana ad \(n\) fogli di \(C_1\), può ridursi per continuità ad un nuovo sistema di cappi, che si lascia riportare all’antico mutando opportunamente le relative trasposizioni. Ora se la \(f\) è curva di diramazione di un piano \(n\)-plo, proiezione di una superficie \(F\) razionalmente determinata, bisogna che una delle dette curve, sia per es. \(C_1\), essendo sezione di \(F\), venga razionalmente
%CONT>>
% scan: s109R
% !PRINTED 195
\weberpage{195}
% head: INVARIANTI NUMERICI E PIANI MULTIPLI
%<<CONT
determinata e perciò la sua riemanniana non si scambi con quelle di \(C_2,\ \ldotp\ldotp\ldotp\ldotp\,C_r\), ma ritorni sempre in \(C_1\) per ogni giro chiuso percorso dal punto \(t\) nel piano di questa variabile complessa.

Si tratta in primo luogo di stabilire le \textit{condizioni d’invarianza della} (classe a cui appartiene) \(C_1\) o \textit{della sua riemanniana, rispetto a tutti i possibili cammini chiusi percorsi da \(t\) nel piano \(\tau\).} Queste condizioni saranno certo condizioni necessarie per l’esistenza d’un piano \(n\)-plo avente \(f\) come curva di diramazione.

Precisiamo il senso della nostra ricerca. Anzitutto si scelga un valore iniziale di \(t\), sia p. es. \(t = 0\), e nel piano della variabile complessa \(x\) si assegni un sistema ordinato di cappi uscenti da \(O\) ed avvolgenti i punti di diramazione \(A_1\,A_2\,\ldotp\ldotp\ldotp\ldotp\,A_{2m}\) (intersezioni di \(f\) colla retta \(y = 0\)); quindi, in ordine ad una data nomenclatura dei rami \(1,\ 2,\ \ldotp\ldotp\ldotp\ldotp\,n\) della funzione algebrica \(z(x)\), si prenda nota delle trasposizioni che rispondono ai detti cappi \(OA_i\), riuscendo così a definire la riemanniana ad \(n\) fogli rappresentativa della nostra \(C_1\) (quella di cui si ricercano le condizioni d’invarianza rispetto ai giri chiusi percorsi da \(t\) nel piano \(\tau\)).

Avendo definita la riemanniana iniziale, fra quelle che rappresentano le rette multiple \(y = tx\), faremo variare per continuità la \(t\) (a partire dal punto \(t = 0\)) descrivendo un cammino chiuso nel piano \(\tau\) di codesta variabile; mentre \(t\) varia in questo modo, anche la \(C_1\) varia per continuità e la sua riemanniana si può costruire con fogli sovrapposti sul medesimo piano della variabile complessa \(x\), poichè (nel linguaggio della geometria algebrica) alla retta multipla \(y = tx\) --- che ha per punti di diramazione le intersezioni con \(f\) --- si può sostituire in generale la sua proiezione sull’asse delle \(x\) fatta dal punto all’infinito dell’asse \(y\).

Ora, in corrispondenza ad una variazione continua di \(t\) nel piano \(\tau\), escludendo le posizioni di \(t\) cui rispondono due punti \(A_r\) e \(A_s\) coincidenti, vedremo nel piano della variabile complessa \(x\) i punti \(A_1\,A_2\,\ldotp\ldotp\ldotp\ldotp\,A_{2m}\) moversi con continuità senza mai coincidere, descrivendo linee qualsiasi (chiuse o aperte) ed è lecito far variare i cappi \(OA_i\) che li avvolgono, sempre con continuità, in modo che essi --- conservando l’ordine in cui si succedono --- rimangano linee aperte e sciolte, non attraversantisi fra loro.

Dopo ciò si consideri, nel piano \(\tau\), un giro chiuso percorso dal punto \(t\), a partire dall’origine \(t = 0\). Se questo giro non comprenda alcun \textit{punto critico} cui rispondono due punti \(A_i\) coincidenti, esso potrà deformarsi per continuità riducendosi ad un ciclo nullo, perciò in corrispondenza ad esso il sistema ordinato di cappi e le relative sostituzioni che definiscono la riemanniana \(C_1\) dovrà ritornare in sè stesso: insomma il detto ciclo non dà luogo ad alcuna condizione per l’invarianza della \(C_1\).

% scan: s110L
% !PRINTED 196
\weberpage{196}
% head: CAPITOLO QUINTO
% stamp: Dark vertical scan/binding line along the right margin and scattered specks.
Pertanto le condizioni d’invarianza della \(C_1\) si ridurranno a quelle che rispondono a giri elementari o cappi del piano \(\tau\) avvolgenti un punto critico \(t\), per cui coincidono due punti \(A_i\).

Ora siffatti giri possono appartenere a tre specie diverse:

1) giri che avvolgono un punto \(T\), cui risponde una tangente semplice della curva \(f\) condotta dal punto esterno \(O\);

2) giri che avvolgono un punto \(D\) cui risponde una retta per \(O\) che va ad un nodo della curva \(f\); e

3) giri che avvolgono un punto \(Q\) cui risponde una retta per \(O\) che va ad una cuspide di \(f\).

Per ciascuno di questi tre casi è facile determinare le condizioni d’invarianza della riemanniana \(C_1\), incominciando dall’ipotesi semplice in cui i due punti \(A_1A_2\) che vengono a coincidere nel punto critico (\(T\) o \(D\) o \(Q\)) appartengono a cappi \(OA_1\) e \(OA_2\) contigui ed \textit{onestamente vicini}: intendiamo con ciò che al limite (quando \(A_1\) e \(A_2\) si sovrappongono) le due linee \(OA_1\) e \(OA_2\), aventi comune il punto terminale \(A_1 = A_2\), delimitino una superficie in cui non cade alcuno degli altri punti \(A_i\) per modo che codeste linee possano ridursi infinitamente vicine senza attraversare altri cappi della riemanniana variabile e quindi anche senza mutare le sostituzioni sui rami di \(z(x)\) che loro corrispondono.

Esaminiamo successivamente i tre casi definiti innanzi:

1) Se \(A_1\) e \(A_2\) vanno a coincidere in un punto \(T\), sulla retta multipla limite \(OT\) i due punti di diramazione \(A_1\) e \(A_2\) si perdono, confondendosi in un punto critico apparente immagine di un punto doppio di \(C_1\) (e dopo il giro di \(t\) si scambiano l’uno coll’altro); ciò significa che le due trasposizioni relative ai due cappi \(OA_1\) e \(OA_2\) (onestamente vicini) debbono essere \textit{identiche}. Viceversa se questo accade, per esempio se ad \(OA_1\) e \(OA_2\) risponde la medesima trasposizione \((12)\), la riemanniana \(C_1\), in corrispondenza al cappio del piano \(\tau\) avvolgente \(T\), rimane invariante.

2) Se \(A_1\) e \(A_2\) vanno a coincidere in un punto \(D\), sulla retta multipla limite \(OD\) i due punti di diramazione \(A_1\) e \(A_2\) si sovrappongono in un punto che è immagine di due punti distinti di \(C_1\). Ciò significa che le trasposizioni relative ai due cappi \(OA_1\) e \(OA_2\) (onestamente vicini) debbono portare su coppie di rami distinte, ossia essere permutabili e in generale \textit{disgiunte} come \((12)\) e \((34)\)\origfootnote{1}{s110L-1}{Saranno invece identiche e perciò ancora permutabili se il punto \(D\) sia un nodo virtualmente inesistente di \(f\).}. Viceversa se questo accade, la riemanniana \(C_1\), in corrispondenza al cappio del piano \(\tau\) avvolgente \(D\), rimane invariante.

% scan: s110R
% !PRINTED 197
\weberpage{197}
% head: INVARIANTI NUMERICI E PIANI MULTIPLI
% stamp: A thin vertical scan/binding line is visible at the right edge; scattered specks.
3) Infine se \(A_1\) e \(A_2\) vanno a coincidere in un punto \(Q\), sulla retta multipla limite \(OQ\) i due punti di diramazione \(A_1\) e \(A_2\) si sovrappongono in \(Q\), e per effetto del giro di \(t\) si scambiano fra loro; ma il punto \(Q\) diventa immagine di un flesso di \(C_1\), per modo che le trasposizioni relative ad \(A_1\) e \(A_2\) dovranno dare come prodotto una sostituzione ciclica del \(3^\circ\) ordine, portando perciò su due coppie di rami aventi un ramo a comune, cioè essere \textit{concatenate} come \((12)\) e \((23)\).

Le condizioni 1), 2) e 3) assicurano l’invarianza della riemanniana \(C_1\) rispetto ai giri avvolgenti i punti \(T\), \(D\) e \(Q\), nel piano \(\tau\), quando i due punti \(A_rA_s\) che vanno a coincidere in codesti punti siano i termini di cappi \(OA_r\) e \(OA_s\) che diventano onestamente vicini. Il caso in cui \(OA_r\) e \(OA_s\) non diventano onestamente vicini, si lascia ricondurre al precedente: in questo caso infatti accade che le due linee anzidette che vengono ad avere l’estremo comune \(A_r = A_s\) formeranno una linea chiusa delimitante una superficie entro cui cade qualcun altro dei punti \(A_i\); ma si può sempre ridurre la linea \(OA_s\) infinitamente vicina alla \(OA_r\), facendola variare per continuità traversando i suddetti punti \(A_i\): vuol dire che per effetto di questa riduzione, cambiando l’ordine dei cappi, cambia anche la trasposizione corrispondente al cappio \(OA_i\)\origfootnote{1}{s110R-1}{Cfr. p. es. \textsc{Enriques-Chisini}, \textit{Lezioni}, Libro V, \S~3, vol. III, pag. 18.}. % ?? uncertain: subscript in the final OA -- i or 1.
% EDITORIAL D0108 RISOLTO: The enlarged final subscript matches the dotted i in the nearby A_i occurrences, rather than the digit 1.

Possiamo riassumere il risultato ottenuto dicendo che: \textit{le condizioni d’invarianza di una riemanniana \(C_1\), per riguardo alle sue variazioni nel fascio delle rette \(n\)-ple \(y = tx\), coi punti di diramazione sulla curva \(f\), sono relative ai giri che la \(t\) compie nel piano \(\tau\), avvolgendo rispettivamente i punti critici \(T\), \(D\) e \(Q\)} (punti che rispondono alle tangenti proprie e alle rette che vanno ai nodi e alle cuspidi della curva di diramazione \(f\)); \textit{esse si esprimono dicendo che le trasposizioni corrispondenti ai punti di diramazione \(A_r\) e \(A_s\) che vengono a coincidere in uno di quei punti critici, debbono essere rispettivamente identiche, o permutabili} (e in generale \textit{disgiunte}) \textit{o concatenate, quando il sistema dei cappi che definisce la riemanniana sia trasformato in guisa che i cappi avvolgenti i detti punti, \(OA_r\) e \(OA_s\), diventino onestamente vicini.}

Le condizioni sopra espresse sono \textit{condizioni necessarie}, perchè la curva \(f\), dotata di nodi e di cuspidi, sia curva di diramazione d’un piano \(n\)-plo \((n > 2)\). Vogliamo riconoscere che esse sono altresì \textit{sufficienti} per l’esistenza di tale piano \(n\)-plo.

A tal uopo si consideri il fascio di piani per l’asse \(z\), aventi per tracce, nel piano \((z = 0)\) della curva \(f(xy) = 0\), le rette \(y = tx\). In ciascuno di questi piani si vuol costruire una curva ben definita che, per proiezione dal punto all’infinito dell’asse \(z\), diciamo \(Z\), dia luogo alla retta \(n\)-pla \(C_1\) coi punti di diramazione in \(A_1\,A_2\ldotp\ldotp\ldotp\ldotp\,A_{2m}\)
%CONT>>
% scan: s111L
% !PRINTED 198
\weberpage{198}
% head: CAPITOLO QUINTO
%<<CONT
% stamp: Dark vertical scan line at the right edge and scattered specks.
intersezioni colla curva $f$. La $C_1$, come si è detto innanzi, è definita come riemanniana in modo invariante rispetto ai giri chiusi di $t$ nel piano; si può costruire quindi una immagine proiettiva irriducibile di $C_1$ cioè una curva $C^{(n+h)}$, d’un certo ordine $n+h$ abbastanza alto, che passi $h$ volte per il punto $Z$ e tocchi le rette $ZA_i$. Ad una siffatta curva $C^{(n+h)}$ si può imporre altresì di passare per $n$ punti distinti $O_1\,O_2\,\ldotp\ldotp\ldotp\ldotp\,O_n$, della retta $OZ$, ed inoltre di soddisfare un certo numero di condizioni lineari, dipendenti da $h$; si troverà così un numero finito di curve $C^{(n+h)}$ irriducibili, e fra queste un numero finito $s$ di curve rappresentate dalla riemanniana $C_1$, cui risponderanno $s$ funzioni algebriche egualmente diramate $z_1(x),\,z_2(x),\,\ldotp\ldotp\ldotp\ldotp\,z_s(x)$\origfootnote{1}{s111L-1}{Cfr. \textsc{Enriques-Chisini}, \textit{Lezioni}. Libro V, § 33, vol. III, pag. 355.}; allora la somma
\[
z(x)=z_1(x)+z_2(x)+\ldotp\ldotp\ldotp\ldotp+z_s(x),
\]
corrisponderà del pari ad una curva irriducibile $C$ rappresentata sulla riemanniana $C_1$ e razionalmente definita. Facciamo variare codesta $C$ col parametro $t$, e si otterrà una superficie $F$ rappresentata sul piano $n$-plo $z=0$ colla curva di diramazione $f$. È ovvio che la $f$ fa parte della curva di diramazione del piano $n$-plo ottenuto proiettando la $F$ dal punto $Z$; si potrebbe tuttavia dubitare che a questa curva si aggiunga un certo numero di rette di diramazione uscenti da $O$; ma se così fosse dovrebbe qualche piano per la retta $OZ$ toccare la $F$ lungo una curva, e ciò è impossibile se i punti $O_1O_2\ldots O_n$ sono distinti, come si è supposto.

In conclusione: \textit{Le condizioni d’esistenza di un piano $n$-plo $(n>2)$ di cui sia assegnata la curva di diramazione $f(xy)=0$, dotata di nodi e di cuspidi,} (e non passante per l’origine $O=[00]$) \textit{si esprimono mercè le condizioni d’invarianza di una delle riemanniane rappresentate sulla retta multipla $y=tx$, in relazione ai cammini chiusi che il parametro $t$ può compiere, a partire dall’origine $O$ nel piano della propria variabile complessa, avvolgendo i punti che rispondono alle tangenti condotte ad $f$ da $O$ ovvero alle rette congiungenti $O$ coi suoi nodi o colle sue cuspidi.}

È essenziale osservare che le condizioni richieste per l’esistenza di un piano multiplo avente una certa curva di diramazione $f$, sono \textit{condizioni qualitative}, che non possono variare quando la $f$ stessa varii con continuità. Quindi il teorema d’esistenza che abbiamo stabilito dà luogo al seguente:

\textit{Corollario. -- Dato nel piano un sistema continuo di curve $f$, dotate di nodi e di cuspidi, se una di codeste curve è curva di diramazione di un piano $n$-plo, altrettanto accade per tutte le curve del sistema.}

% scan: s111R
% !PRINTED 199
\weberpage{199}
% head: INVARIANTI NUMERICI E PIANI MULTIPLI
% stamp: Scattered scan specks.
\edsection{10. Complementi: forme limiti della curva di diramazione secondo Chisini.}{10. Complementi: forme limiti della curva di diramazione secondo Chisini.}{sec.V.10}

Il teorema d’esistenza pei piani multipli, col suo corollario, esposto nel precedente paragrafo, è stato stabilito da \textsc{F. Enriques} nella memoria « Sulla costruzione delle funzioni algebriche di due variabili possedenti una data curva di diramazione », inserita negli Annali di Matematica s. IV, t. I (1923).

A questo teorema si collegano poi interessanti sviluppi. Anzitutto \textsc{O. Zariski}\origfootnote{1}{s111R-1}{\textsc{O. Zariski}, \textit{On the problem of existence of algebraic functions of two variables possessing a given branch curve}. Am. J. Math., vol. 51 (1929). \textit{On the linear connection index of the algebraic surfaces $z^n=f(xy)$}. Proc. Mat. Accad. Sc. U. S. A., vol. 15 (1929). \textit{On the irregularity of cyclic multiple planes}. Ann. of Math., II s., vol. 32 (1931). Cfr. \textit{Algebraic surfaces}, Berlino, Springer, 1935.} ha trasformato le condizioni di esistenza di Enriques in una forma topologico-gruppale, dalla quale egli deduce in particolare che una curva priva di cuspidi non può essere curva di diramazione di un piano multiplo d’ordine $n>2$, tranne il caso della curva contata $n-1$ volte che risponde a un piano $n$-plo ciclico. % ?? uncertain: journal abbreviation and punctuation in footnote 1 -- "Mat." or "Nat."; comma or period after "Springer".
% EDITORIAL D0109 RISOLTO: The enlarged note clearly prints an M with a central valley (Mat.), and a comma after Springer. Preserve both printed readings; Mat. is likely a bibliographic misprint, recorded separately.

Ma, sia sotto la forma di Enriques che sotto quella di Zariski, il teorema di esistenza dei piani multipli va incontro a notevoli difficoltà nella sua pratica applicazione, dando luogo al problema ulteriore di determinare le curve di diramazione soddisfacenti alle condizioni poste. Un contributo importante a questo problema è stato dato da \textsc{O. Chisini}\origfootnote{2}{s111R-2}{\textsc{O. Chisini}, \textit{Un teorema di esistenza dei piani multipli}. Note I e II, Atti Acc. Lincei, VI, 19 (1934). \textit{Un più generale teorema di esistenza dei piani multipli}. Atti Acc. Lincei, VI, 27 (1938). \textit{Altre curve di diramazione dei piani $n$-pli}. Atti Acc. Lincei, VI, 29 (1939).}, con la ricerca delle \textit{forme limiti} contenute in un sistema continuo di curve di diramazione, dotate di un certo numero di nodi e di cuspidi.

Per comprendere la ricerca di \textsc{Chisini}, si consideri una superficie $F$ d’un certo ordine $m+h$, con un punto $h$-plo $Z$, la quale venga a spezzarsi in due parti $f_1$ ed $f_2$ intersecanti le rette per $Z$ rispettivamente in $n$ ed $s$ punti $(m=n+s)$. Qui la curva di diramazione del piano $m$-plo, cu sui $F$ viene proiettata da $Z$, si spezzerà nelle due curve di diramazione del piano $n$-plo e del piano $s$-plo corrispondenti ad $f_1$ e ad $f_2$, e nella curva doppia proiezione della curva comune ad $f_1$ ed $f_2$, che diventa curva critica apparente. Questa osservazione conduce a ricercare se si possa costruire la curva di diramazione di un piano $m$-plo, con $m=n+s$, sommando alle curve di diramazione dei due piani $n$-plo ed $s$-plo, una curva doppia, che dovrà avere un comportamento opportuno rispetto a quelle, % [sic] cu sui -> su cui, apparent transposition of the initial consonants.
%CONT>>
% scan: s112L
% !PRINTED 200
\weberpage{200}
% head: CAPITOLO QUINTO
%<<CONT
% stamp: Dark vertical scan line at the right edge and scattered specks.
per modo che la curva composta, costruita come si è detto, rappresenti la curva limite di un sistema continuo di curve di diramazione di piani $m$-pli, dotate di un certo numero di nodi e di cuspidi.

Ciò che si è detto vale a spiegare la genesi del pensiero di \textsc{Chisini}. Tuttavia la ricerca generale delle curve doppie da aggiungere alla somma delle curve di diramazione di un piano $n$-plo e di un piano $s$-plo, dà luogo ancora a notevoli difficoltà. Per superarle, \textsc{Chisini}, si restringe al caso in cui sia, per es., $s=1$, cioè si abbia soltanto la curva di diramazione di un piano $n$-plo, da cui, con l’aggiunta di una conveniente curva doppia, si voglia trarre la curva di diramazione di un piano $(n+1)$-plo. Basterà spiegare la costruzione di Chisini riferendosi ad un semplice esempio.

Proponiamoci di costruire il sistema delle curve di diramazione dei piani quadrupli che nascono per proiezione della superficie generale del $4^\circ$ ordine $F_4$. Queste curve di diramazione sono $C_{12}$ d’ordine 12 con $d=12$ nodi e $k=24$ cuspidi.

Se facciamo passare la superficie $F_4$ per il centro di proiezione $Z$, la curva di diramazione del piano quadruplo $C_{12}$, degenera nella curva di diramazione di un piano triplo $C_{10}$, dotata di $k'=18$ cuspidi (e non più di nodi), e in una retta $C_1$ tangente doppia della $C_{10}$, contata due volte. A sua volta la $C_{10}$ è suscettibile di degenerare in una sestica $C_6$, affatto generale, curva di diramazione di un piano doppio, e in una conica sestitangente ad essa, contata due volte, e ciò in corrispondenza alla condizione che si imponga alla $F_4$ di passare doppiamente per il punto $Z$. Infine se il punto doppio $Z$ della $F_4$ diventi triplo, la corrispondente curva di diramazione $C_6$ del nominato piano doppio degenera in una cubica contata due volte.

Giova ora esaminare che cosa diventino le singolarità della curva di diramazione del piano multiplo, in corrispondenza alle successive degenerazioni sopra indicate.

Anzitutto, quando la anzidetta $C_{12}$ degenera nella $C_{10}+2C_1$, ciascuno dei due punti di contatto $P_1$ e $P_2$ di $C_1$ e $C_{10}$ appare come limite di 3 cuspidi della $C_{12}$, e ciascuno degli altri 6 punti comuni a $C_1$ e $C_{10}$ risulta limite di due nodi di $C_{12}$. L’asserto si giustifica riferendosi alla superficie $F_4$ che viene a passare per il centro di proiezione $Z$; occorre tener conto che le cuspidi di $C_{12}$ si ottengono proiettando da $Z$ le intersezioni di $F_4$ con la $1^{\mathrm{a}}$ e la $2^{\mathrm{a}}$ polare del punto $Z$, e che i nodi di $C_{12}$ provengono dalle bitangenti alla $F_4$ passanti per $Z$, le quali diventano ora le sei tangenti altrove condotte per $Z$ alla quartica sezione della $F_4$ col piano tangente in $Z$, ciascuna contata due volte.

Ciò che si è osservato in ordine alla degenerazione delle singolarità della $C_{12}$ quando questa diventa $C_{10}+2C_1$, sussiste ancora nei ri%
%CONT>>
% scan: s112R
% !PRINTED 201
\unskip\weberpage{201}%
% head: INVARIANTI NUMERICI E PIANI MULTIPLI
%<<CONT
% stamp: Scattered scan specks.
guardi della $C_{10}$ degenerante in $C_6+2C_2$; vale a dire che i 6 punti di contatto della $C_2$ con la $C_6$ vengono ad assorbire ciascuno 3 cuspidi della $C_{10}$.

Abbiamo osservato in qual modo la curva di diramazione $C_{12}$ di un piano quadruplo rappresentativo di una $F_4$ possa farsi degenerare successivamente in una $C_{10}+2C_1$ e poi in una $C_6+2C_2+2C_1$, e infine anche, se si vuole, in una $2C_3+2C_2+2C_1$, dove la $C_{10}$ è curva di diramazione d’un piano triplo, bitangente alla retta $C_1$, e la $C_6$ (curva di diramazione di un piano doppio) è sestitangente alla conica $C_2$. Ora questo processo analitico di riduzione si lascia invertire conducendo quindi ad un processo sintetico costruttivo dei piani multipli, che costituisce appunto il \textit{teorema di esistenza di Chisini}. Infatti l’autore risale dalle forme limiti a forme più generali dei piani multipli ritrovando le terne di cuspidi e le coppie di punti doppi che sono assorbiti dalle singolarità limiti. Così, prendendo le mosse dal piano doppio con sestica $C_6$ di diramazione e particolarizzando la $C_6$ in guisa che ammetta una conica $C_2$ sestitangente, si considererà la curva $C_6+2C_2$ come caso particolare di una $C_{10}$, per la quale verranno soddisfatte le condizioni d’Enriques spiegate nel precedente paragrafo, e che perciò risulterà curva di diramazione d’un piano triplo. E poi, sommando ancora alla $C_{10}$ una bitangente $C_1$, si proverà similmente, che la $C_{10}+2C_1$ (o la $C_6+2C_2+2C_1$) è forma limite di una $C_{12}$ che costituisce la curva di diramazione d’un piano quadruplo, rappresentativo di una superficie $F_4$ a sezioni piane di genere 3. È inutile dire che il caso a cui ci siamo riferiti non è che un esempio particolare, sebbene caratteristico, di un metodo che conduce l’autore al riconoscimento delle condizioni d’esistenza per una più vasta famiglia di piani multipli. Non ci proponiamo di esporre qui gli sviluppi più generali di \textsc{Chisini}, ma ci limiteremo a chiarire il suo metodo costruttivo ed anche le questioni delicate cui esso dà luogo col riferirci ad una famiglia notevole di superficie che comprende in sè le quartiche $F_4$, sulla quale dovremo ritornare nel seguito di queste Lezioni: diciamo delle superficie coi generi $p_a=p_g=1$ e $p^{(1)}=1$ possedenti una curva canonica d’ordine zero. % The plus sign at a printed line break is repeated in the source.

D’accordo colle formule del § 6, si ottengono superficie appartenenti alla anzidetta famiglia, partendo dai piani multipli $n$-pli, con $n=2\pi-2$, definiti da una curva di diramazione $C_{3n}=C_{6\pi-6}$ d’ordine $2n+2\pi-2=3n=6\pi-6$, dotata di $k(\pi)=24(\pi-2)$ cuspidi e $d(\pi)=18\pi^2-78\pi+84$ nodi: infatti un piano $n$-plo siffatto (ove esista) ha i generi $p_a=p^{(1)}=1$ ed anche il $p_g=1$, possedendo una curva canonica d’ordine zero; e le immagini delle rette del piano appartengono ad un sistema lineare di genere $\pi$ e grado $n=2\pi-2$, e quindi di dimensione $\pi$. Così il piano $n$-plo
%CONT>>
% scan: s113L
% !PRINTED 202
\weberpage{202}
% head: CAPITOLO QUINTO
%<<CONT
% stamp: Short slanted mark in the left margin beside "In generale"; dark right-edge scan line and scattered specks.
si può far nascere per proiezione, da punti esterni, di una superficie $F_n$ d’ordine $n=2\pi-2$ dello $S_\pi$, a sezioni iperpiane canoniche di genere $\pi$.

Ora qui si tenterà di provare l’esistenza di piani $n$-pli rispondenti alle $C_{3n}$ di diramazione sopra indicate, e quindi delle $F_n=F_{2\pi-2}$ di $S_\pi$ a sezioni canoniche di genere $\pi$, per tutti i valori di $\pi=3,4,5\ldots$. Già nella discussione precedente abbiamo indicato come si pervenga a costruire la curva di diramazione $C_{12}$ di un piano quadruplo (rappresentativo di una $F_4$) corrispondente a $\pi=3$, a partire dal piano doppio con sestica di diramazione $C_6$, sestitangente ad una conica $C_2$. Ora si procederà similmente a costruire il piano sestuplo corrispondente a $\pi=4$, a partire da una $C_{12}$ sestitangente ad una conica $C_2$: ci domanderemo se $C_{12}+2C_2$ sia la forma degenere di una $C_{16}$, curva di diramazione di un piano quintuplo e se, sommando poi a questa una retta bitangente $C_1$, contata due volte, si pervenga ad una $C_{18}$, curva di diramazione d’un piano sestuplo. Se questo accade, mentre la $C_{12}$ sopra indicata possiede $k(\pi)=k(3)=24$ cuspidi e $d(\pi)=d(3)=12$ nodi, la $C_{16}$ avrà in più 18 cuspidi (3 cuspidi corrispondendo a ciascuno dei 6 contatti di $C_{12}$ con $C_2$) e 24 nodi (2 nodi nascendo da ciascuna delle 12 intersezioni semplici di $C_{12}$ con $C_2$), e così in tutto 42 cuspidi e 36 nodi; sommando poi alla $C_{16}$ la bitangente $C_1$ contata due volte, si otterrà una $C_{18}$ avente in più altre 6 cuspidi e altri 24 nodi, e così % The equality sign before 24 is repeated across a printed line break.
\[
k(4)=48\ \text{cuspidi}
\]
e
\[
d(4)=60\ \text{nodi}.
\]

In generale per proseguire questo procedimento costruttivo, in corrispondenza ai valori successivi di $\pi$, giova anticipare una nozione che verrà stabilita nel seguente paragrafo, cioè che le superficie $F_n$ (di generi $p_a=p_g=p=1$, $p^{(1)}=1$) dipendono da almeno $9p-2p^{(1)}+12=19$ parametri essenziali o moduli, invarianti rispetto alle trasformazioni birazionali. Per conseguenza, quando sia dato un piano multiplo definito da una curva di diramazione $C_{3n}=C_{6\pi-6}$ coi caratteri anzidetti, disponendo di un parametro, si potrà particolarizzare la $C_{3n}$ in guisa che possegga una conica sestitangente $C_2$. E per seguire il metodo di Chisini siamo indotti sempre a domandarci se la $C_{3n}+2C_2$ possa riguardarsi come forma limite di una $C_{3n+4}$, d’ordine $3n+4$, che possieda, oltre alle cuspidi e ai nodi della $C_{3n}$, 18 cuspidi e $12n-24=24(\pi-2)$ nodi, corrispondenti rispettivamente ai 6 punti di contatto della $C_2$ colla $C_{3n}$ e alle $6n-12$ intersezioni ulteriori di queste curve; e se le $C_{3n+4}$ così costruite risultino curve di diramazione di piani $(n+1)$-pli. % The equality sign between C_{3n} and C_{6\pi-6} is repeated across a printed line break.

Se ciò possa farsi (in guisa che vengano soddisfatte le condizioni
%CONT>>
% scan: s113R
% !PRINTED 203
\weberpage{203}
% head: INVARIANTI NUMERICI E PIANI MULTIPLI
%<<CONT
% stamp: Scattered scan specks and a short dark edge trace at the lower left.
d’invarianza del precedente paragrafo) le $C_{3n+4}$ così costruite risulteranno curve di diramazione di piani $(n+1)$-pli. E, come diremo, sommando ad esse una conveniente retta bitangente $C_1$, si otterrà poi la curva di diramazione di un piano $(n+2)$-plo, rappresentativo di una $F_{n+2}$, a sezioni canoniche di genere $\pi+1$ dello spazio $S_{\pi+1}$.

Ma la domanda sopra espressa non comporta sempre una risposta affermativa. Perchè la $C_{3n}+2C_2$ sia effettivamente forma limite di una $C_{3n+4}$ soddisfacente alle nostre richieste, bisogna che le singolarità di essa si risolvano in nodi e cuspidi nel modo indicato, e che per la curva $C_{3n+4}$ così costruita si verifichino le condizioni d’invarianza che permettono di ritenerla come curva di diramazione di un piano $(n+1)$-plo. Ora l’insieme di queste condizioni non viene soddisfatto per una qualsiasi conica sestitangente a $C_{3n}$; ma, riferendosi, per esempio al caso $n=4$, si può soddisfare riducendo la $C_{12}$ alla forma limite di Chisini $C_{12}=2C_2+2C_1+C_6$ e ammettendo poi che vi sia una seconda conica sestitangente a $C_6$ diversa da $C_2$. Da ciò appare come la difficoltà si possa superare per i più piccoli valori di $n$, ma non si vede che possa farsi per i più grandi.

Comunque, se si sia riusciti a vedere la curva $C_{3n}+2C_2$ come forma limite di una $C_{3n+4}$ che sia curva di diramazione di un piano $(n+1)$-plo, converrà ancora sommare a questa una retta bitangente $C_1$, scelta in guisa che la curva $C_{3n}+2C_2+2C_1$ sia forma limite della curva di diramazione di un piano $(n+2)$-plo. Ma qui (senza bisogno di ricorrere a considerazioni analoghe a quelle che, nel caso di $n=4$, ci hanno condotto alla scelta della $C_2$ sestitangente a $C_{12}$), la difficoltà della scelta di $C_1$ si supera subito in base a questa semplice osservazione: che il piano $(n+1)$-plo definito dalla curva di diramazione $C_{3n+4}$ ha i generi uno e perciò appare rappresentativo di una $F_{n+2}$ $(n=2\pi-2)$ a sezioni canoniche di genere $\pi+1$, proiettata da un suo punto semplice; a questo punto risponde la retta eccezionale, bitangente a $C_{3n+4}$, che è la richiesta $C_1$.

Ciò che abbiamo detto vale a spiegare il metodo costruttivo di Chisini, e in qualche modo a saggiarne il significato coll’esame di un problema concreto. Per quel che concerne questo problema --- relativo all’esistenza e alla classificazione delle superficie $F_n$, coi generi uno --- ne daremo la soluzione diretta nel § 3 del Cap. VII ove si troverà anche una notizia storica dell’ordine di queste ricerche.

Aggiungiamo alle precedenti alcune indicazioni bibliografiche. Mentre nelle citazioni fatte ci riferivamo a piani $n$-pli con $n>2$, si possono annoverare alcune ricerche speciali sul caso dei piani doppi: \textsc{F. Enriques}, « Sui piani doppi di genere 1 », Memorie della Società It. delle Scienze, detta dei XL, Roma, 1896; « Sui piani doppi di genere lineare $p^{(1)}=1$ ». R. Accademia dei Lincei, 1898.
%CONT>>
% scan: s114L
% !PRINTED 204
\weberpage{204}
% head: CAPITOLO QUINTO
%<<CONT
% stamp: Dark vertical scan line at the right edge and scattered specks.
\textsc{L. Campedelli}, « Sopra i piani doppi con tutti i generi uguali all’unità ». Rend. del Sem. mat. di Padova, vol. XI, 1940; « Le superficie con i generi uguali all’unità, rappresentabili in infiniti modi sul piano doppio », Rend. del Sem. Mat. della R. Università di Roma, s. V, vol. I, 1940 e « La classificazione dei piani doppi con tutti i generi uguali all’unità ». Atti del $2^\circ$ Congresso dell’Unione Mat. Ital., Bologna, 1940; « Sui piani doppi con curva di diramazione dell’ottavo ordine ». R. Accademia dei Lincei, 1932; « Sui piani doppi con curva di diramazione del $10^\circ$ ordine », ibidem, 1932.; « Sopra alcuni piani doppi notevoli con curva di diramazione del $10^\circ$ ordine », ibidem, 1932. % [sic] 1932.; -> 1932;, a period is printed immediately before the semicolon.

Al caso dei piani doppi si ravvicina in qualche modo lo studio di quei particolari piani quadrupli che dipendono da equazioni abeliane, cioè dall’estrazione di due radicali quadratici: \textsc{P. Libois}, Rend. Acc. Lincei, 1934.

A classi notevoli di piani multipli $n$-pli, con $n>2$, si riferiscono ancora i seguenti lavori: \textsc{B. Segre}, « Sulla caratterizzazione delle curve di diramazione dei piani multipli generali », Reale Acc. d’Italia, vol. I, 1930. --- \textsc{G. Pompilj}, « Sulla rappresentazione algebrica dei piani tripli », Rend. Sem. mat. della R. Università di Roma, 1939; « Osservazioni sui piani tripli ». R. Istituto Lombardo, 1941; « Sui piani multipli diedrici ». Rend. Sem. mat. della R. Università di Roma, 1939.

\edsection{11. I moduli di una classe di superficie algebriche.}{11. I moduli di una classe di superficie algebriche.}{sec.V.11}

La considerazione dei piani multipli con curva di diramazione semplice, su cui può generalmente rappresentarsi una superficie algebrica, e il relativo teorema di esistenza che abbiamo spiegato nel precedente paragrafo, permette di dire qualcosa in ordine al numero dei \textit{moduli} o parametri da cui dipende la determinazione di una \textit{classe} di superficie entro una famiglia continua cui rispondono certi valori dei caratteri $p$ e $p^{(1)}$. Assumiamo per semplicità che si tratti di superficie regolari $p=p_g=p_a$, aventi qualche plurigenere $P_i>0$, e perciò trasformabili in guisa da eliminare le curve eccezionali, sicchè il $p^{(1)}$ diventi per esse il \textit{genere lineare assoluto}.

Sopra una superficie $F$ soddisfacente alle condizioni anzidette, scegliamo ad arbitrio un sistema lineare irriducibile, senza punti base $|C|$, di genere $\pi$ e grado $n\leq2\pi-2$, che supporremo regolare e non speciale, quindi di dimensione
\[
r=p+n-\pi+1.
\]

Una rete generica contenuta entro $|C|$ permette di rappresentare la $F$ sopra un piano multiplo d’ordine $n$, dotato di una curva
%CONT>>
% scan: s114R
% !PRINTED 205
\weberpage{205}
% head: INVARIANTI NUMERICI E PIANI MULTIPLI
%<<CONT
% stamp: Dark speck in the left margin and scattered scan marks.
di diramazione $D$ d’ordine $2n+2\pi-2$, la quale possiederà un certo numero $d$ di nodi e un certo numero $k$ di cuspidi, dove questi numeri si esprimono per $p$ e $p^{(1)}$, mercè le formule del § 6. Ora, se facciamo variare la curva $D$, d’ordine $2n+2\pi-2$, entro un sistema continuo di curve aventi ancora $d$ nodi e $k$ cuspidi, abbiamo visto che questa curva è curva di diramazione per un piano $n$-plo, cui risponde una superficie avente gli stessi caratteri $p$ e $p^{(1)}$ di $F$. Però non tutte le curve di questo sistema, di dimensione
\[
s\geq(n+\pi-1)(2n+2\pi+1)-d-2k, % ?? uncertain: sign before d -- minus sign with a speck above it; division sign.
% EDITORIAL D0110 RISOLTO: The scan shows a minus bar with one tiny off-axis upper speck and no lower dot; it is not a complete division sign.
\]
corrisponderanno a superficie birazionalmente distinte dalla $F$; anzitutto le $\infty^8$ trasformate omografiche di $D$ conducono alla stessa $F$, e alla stessa rete di curve $|C|$ sopra di essa; in secondo luogo, vi sono nel sistema $\infty^r$ $|C|$, infinite reti, ciascuna delle quali dipende da $3r-6$ parametri. Avremo pertanto almeno
\[
(n+\pi-1)(2n+2\pi+1)-d-2k-8-(3r-6) % ?? uncertain: sign before 8 -- minus sign with a speck above it; division sign.
% EDITORIAL D0111 RISOLTO: Enlargement shows only an upper speck over the minus bar, not the paired dots of division. The minus reading is retained.
\]
superficie, le quali debbono ritenersi birazionalmente distinte dalla $F$. Infatti se $\infty^1$ curve $D$ non corrispondenti a reti contenute entro $|C|$, diano luogo a superficie birazionalmente identiche, ciò significa che la $F$ possiede una serie continua di reti di curve non equivalenti, in cui è contenuta una rete di $C$, e ciò contraddice l’ipotesi che la superficie sia regolare $(p_a=p_g)$.

A dir vero in confronto a questa asserzione si può sollevare il dubbio che una rete di $C$, diciamo $|C|_2$, possa presentarsi come limite di una rete di curve $|L|_2$, d’ordine più alto, dotata di punti base, così come la rete delle rette del piano appare limite d’una rete di coniche con tre punti base. Ma il dubbio espresso viene facilmente escluso per le superficie $F$ su cui il grado d’un sistema lineare di curve non può superare il doppio del genere meno due, e che perciò possono ritenersi prive di curve eccezionali, alle quali ci riferiamo. Invero se una rete $|L|_2$ di genere effettivo $\pi$ e di grado effettivo $n$ possieda qualche punto base d’ordine $i\geq1$, le $L$ avranno colle curve canoniche $K$ un numero virtuale d’intersezioni $2\pi-2-n-\Sigma i^2$, e al limite dovranno spezzarsi in curve $C$ aumentate d’una componente fissa $\theta$, la quale avrà colle $K$ il numero negativo d’intersezioni: $-\Sigma i^2$. Questa $\theta$ avrà dunque un certo genere $\varrho$ e un grado
\[
\nu=2\varrho-2+\Sigma i^2>2\varrho-2; % ?? uncertain: degree symbol \nu -- Greek nu; italic Latin v.
\]
e la sua presenza sulla superficie $F$ contraddice l’ipotesi fatta in base a cui essa si è trasformata in guisa da non possedere curve eccezionali. Qui per la giustezza del nostro computo (affermazione che si avranno $\infty^{3r-6+8}$ piani multipli identici), giova aggiungere che la nostra $F$ non può ammettere una serie continua di trasfor%
%CONT>>
% scan: s115L
% !PRINTED 206
\unskip\weberpage{206}%
% head: CAPITOLO QUINTO
%<<CONT
% stamp: Dark right-edge scan line, scattered specks, and a small curved mark beside "razionali" in the footnote.
mazioni birazionali in se stessa; giacchè ne risulterebbe: o che il sistema completo $|C|$ dovrebbe appartenere ad una più ampia serie di curve disequivalenti, ciò che contraddice all’ipotesi $p_a=p_g$, ovvero che $|C|$ sarebbe trasformato in se stesso, e ciò contraddice all’ipotesi che la superficie $F$ non appartenga alla famiglia delle rigate\origfootnote{1}{s115L-1}{Infatti le superficie che posseggono un’infinità continua di trasformazioni proiettive sono razionali o rigate. Teorema di \textsc{Enriques-Fano}. Cfr. \textsc{Enriques-Conforto}, \textit{Le superficie razionali}, Op. cit., cap. VI, § 37, pag. 496.}.

\edFlush
Siamo ora in grado di scrivere il numero $M$ dei moduli da cui dipende la classe delle superficie $F$, o almeno un limite inferiore di questo numero. Avremo infatti
\[
M\geq(n+\pi-1)(2n+2\pi+1)-d-2k-8-(3r-6),
\]
ossia
\[
M\geq(n+\pi-1)(2n+2\pi+1)-d-2k-3r-2.
\]

In questa formula sostituiamo a $d$ e $k$ le loro espressioni per $p$ e $p^{(1)}$ date dalle formule del § 6:
\[
\begin{aligned}
d&=2[(n+\pi)^2-5n-17\pi-2p^{(1)}+6p_a+24]\\
k&=3(n+6\pi+p^{(1)}-4p_a-11);
\end{aligned}
\]
inoltre sostituiamo ad $r$ il suo valore
\[
r=p+n-\pi+1.
\]

Toveremo così % [sic] Toveremo -> Troveremo, the r after T is missing in print.
\[
M\geq9p-2p^{(1)}+12,
\]
ovvero, più precisamente,
\[
M=9p-2p^{(1)}+12+\omega,
\]
dove $\omega\geq0$ designa la sovrabbondanza del sistema delle curve piane $D$; vuol dire che nel sistema continuo delle curve $D$ d’ordine $2n+2\pi-2$, dotate di $d+k$ punti doppi, le condizioni imposte che $k$ punti doppi diventino cuspidi, non sono $k$ condizioni indipendenti, ma dipendono soltanto da $k-\omega$ condizioni.

Concludiamo intanto:

\textit{Le superficie regolari con qualche plurigenere $P_i>0$, di genere superficiale $p$ e genere lineare assoluto $p^{(1)}$, dipendono, in generale, da}
\[
M=9p-2p^{(1)}+12+\omega
\]
\textit{moduli, con $\omega\geq0$.}

Per le superficie con $p_a<p_g$, il computo che precede si deve modificare in rapporto alla proprietà caratteristica di queste super%
%CONT>>
% scan: s115R
% !PRINTED 207
\unskip\weberpage{207}%
% head: INVARIANTI NUMERICI E PIANI MULTIPLI
%<<CONT
% stamp: Scattered scan specks.
ficie, che stabiliremo più avanti, cioè alla proprietà di contenere serie $\infty^{p_g-p_a}$ di sistemi lineari disequivalenti. Quindi per le \textit{superficie irregolari il numero dei moduli sarà}
\[
M=10p_a-p_g-2p^{(1)}+12+\omega
\]
con $\omega\geq0$\origfootnote{1}{s115R-1}{Salvo il significato di $\omega$ la formula vale a fortiori per le superficie che ammettono un gruppo continuo di trasformazioni in sè stesse.}.

Ritorniamo alle superficie regolari $(p_a=p_g=p)$. Per le superficie con curva canonica o bicanonica d’ordine zero ($p=1$, $p^{(1)}=1$, e $p=0$, $p^{(1)}=1$) si avrà rispettivamente
\[
9p-2p^{(1)}+12=19
\]
e
\[
9p-2p^{(1)}+12=10,
\]
e verificheremo che in effetto il numero dei moduli è appunto, per le prime
\[
M=19,
\]
e per le seconde
\[
M=10;
\]
dunque
\[
\omega=0.
\]

Ma in genere si troverà
\[
\omega>0.
\]

\edsection{12. Digressione sulla integrità della serie caratteristica d’un sistema completo di curve piane dotate di nodi e di cuspidi.}{12. Digressione sulla integrità della serie caratteristica d’un sistema completo di curve piane dotate di nodi e di cuspidi.}{sec.V.12}

Per esprimere il computo dei moduli di una superficie con una formula più significativa conviene cercare il significato del carattere $\omega$, che appare nella formula del paragrafo precedente. È agevole riconoscere che « la sovrabbondanza $\omega$ di un sistema continuo di curve piane $\{D\}$, d’un certo ordine $n$ dotate di $d$ nodi e di $k$ cuspidi, è minore o eguale all’indice di specialità della sua serie caratteristica sopra una generica $D$, cioè della serie segata su di essa dalle curve dello stesso ordine che passano semplicemente per i $d$ nodi e per le $k$ cuspidi, e toccano in questi punti le tangenti cuspidali ».

Infatti ricordiamo che, essendo data una serie continua $\{D\}$ di curve piane $D$, d’un certo ordine $n$ dotate di $d$ nodi e di $k$ cuspidi, le curve di $\{D\}$ infinitamente vicine ad una data $D=D_0$ hanno con questa 2 intersezioni assorbite in ciascun nodo e 3 intersezioni assorbite in ciascuna cuspide\origfootnote{2}{s115R-2}{Osservazione di De \textsc{Jonquières} e \textsc{Beck} (Math. Annalen, Bd. 14, 1878). Cfr. \textsc{Enriques-Chisini}, Libro II, cap. II (vol. I, pag. 331, cfr. vol. II, pag. 286).}. Vuol dire che la serie caratteristica,
%CONT>>
% scan: s116L
% !PRINTED 208
\weberpage{208}
% head: CAPITOLO QUINTO
%<<CONT
% stamp: Dark scan-edge band at the right and scattered scan specks.
segata su \(D_0\) dalle \(D\) infinitamente vicine è contenuta nella serie completa segata sulla detta \(D_0\) dalle curve dello stesso ordine passanti per i suoi \(d\) nodi e le sue \(k\) cuspidi e tangenti in questi punti alle tangenti cuspidali. Ora, se il sistema continuo \(\{D\}\) ha la sovrabbondanza \(\omega\), la dimensione di esso vale, non già
\[
r = \frac{n(n+3)}{2} - d - 2k,
\]
bensì
\[
r + \omega = \frac{n(n+3)}{2} - d - 2k + \omega,
\]
e quindi la dimensione della serie \(g_m\), d’ordine \(m = n - 2d - 3k\); segata sulla \(D_0\), di genere % [sic] m = n - 2d - 3k -> m = n^2 - 2d - 3k, apparent omitted square exponent. % ?? uncertain: punctuation after 3k -- semicolon or comma.
% EDITORIAL D0113 RISOLTO: Enlargement shows an upper dot and a lower comma-shaped stroke: the punctuation is a semicolon.
\[
\pi = \frac{n(n-3)}{2} + 1 - d - k,
\]
dalle \(D\) infinitamente vicine è \(\ge r + \omega = m - \pi + \omega\), cioè codesta serie caratteristica è speciale, d’indice di specialità almeno eguale ad \(\omega\). A prima vista si sarebbe tentati di dire che \(\omega\) è senz’altro eguale al detto indice di specialità, cioè che « la serie caratteristica del sistema completo \(\{D\}\), costituito delle curve piane d’ordine \(n\) con \(d\) nodi e \(k\) cuspidi, è completa ». Infatti è facile verificare che viene soddisfatta la seguente \textit{proprietà differenziale}: nel fascio determinato da una \(D\) e da una sua curva aggiunta \(\varphi_n\) che tocchi nelle \(k\) cuspidi le sue tangenti cuspidali, la curva infinitamente vicina alla \(D\), possiede \(d\) nodi e \(k\) cuspidi vicini a quelli della \(D\) stessa.

Ma questa proprietà differenziale non porta di conseguenza che ognuna delle anzidette curve con \(d\) nodi e \(k\) cuspidi, vicine a \(D\), faccia parte di una serie continua di curve dotate dello stesso numero di nodi e di cuspidi. Invero, se ci riferiamo alla varietà \(V\) rappresentativa degli elementi (curve) del sistema \(D\), questa \(V\) ci apparirà come intersezione di una varietà \(W\), dello spazio \(S_R\) a \(R = \frac{n(n+3)}{2}\) dimensioni, rappresentativa del sistema di tutte le curve piane d’ordine \(n\) dotate di \(d + k\) nodi, con \(k\) ipersuperficie o falde ipersuperficiali lineari del detto spazio \(S_R\); e potrebbe accadere p. es. che queste \(k\) falde abbiano in ogni punto comune lo stesso spazio tangente a \(R - (k - 1)\) dimensioni, pur intersecandosi secondo una varietà di \(R - k\) (e non di \(R - k + 1\)) dimensioni: in questo caso un punto di \(V\) (rappresentante una curva di \(\{D\}\)) avrebbe, non già \(\infty^{r-1}\), bensì \(\infty^r\) punti vicini appartenenti a \(V\) (cioè immagini di curve di \(\{D\}\)).

% scan: s116R
% !PRINTED 209
\weberpage{209}
% head: INVARIANTI NUMERICI E PIANI MULTIPLI
% stamp: Broken scan streaks at the right edge and scattered scan specks.
Sebbene taluni indizii e gli esempi che conosciamo tendano ad escludere questo dubbio, non siamo riusciti a dimostrare ciò che crediamo vero, e dobbiamo limitarci a introdurre una

\textit{Ipotesi di lavoro:} Il sistema continuo completo delle curve piane \(D\) dotate di \(d\) nodi e \(k\) cuspidi ha la serie caratteristica completa; perciò \textit{la sovrabbondanza} \(\omega\) di questo sistema \(\{D\}\) \textit{eguaglia l’indice di specialità} della serie stessa, segata su una \(D\) dalle curve infinitamente vicine.

Questa ipotesi conduce ad alcune conseguenze notevoli, ed ha un certo valore euristico, in rapporto agli sviluppi che seguono.

\edsection{13. I moduli delle superficie regolari di genere \(p > 3\).}{13. I moduli delle superficie regolari di genere p > 3.}{sec.V.13}

Abbiamo veduto come venga espresso in generale il numero dei moduli da cui dipende una classe di superficie algebriche (con qualche plurigenere maggiore di zero, cioè non appartenenti alla famiglia delle rigate), quando queste superficie vengano rappresentate sopra piani multipli con curve di diramazione \(D\) d’un certo, ordine \(n\), dotate di un certo numero di nodi e di cuspidi. La formula conseguita nel \S~11 contiene, oltre i caratteri, generi superficiali \(p_a\), \(p_g\) e genere lineare assoluto \(p^{(1)}\), anche un carattere numerico \(\omega\), non negativo, che è la sovrabbondanza del sistema continuo delle \(D\), e — secondo la nostra \textit{ipotesi di lavoro} (\S~12) — l’indice di specialità della serie caratteristica di \(\{D\}\). % ?? uncertain: mark after "certo" -- printed comma or scan speck.
% EDITORIAL D0114 RISOLTO: A comma-shaped mark sits at the normal baseline immediately after certo. Preserve this printed punctuation, despite the unusual syntax.

Ora il significato ammesso per \(\omega\) è suscettibile di un’altra interpretazione e c’induce a scrivere una formula più significativa, in ordine al numero dei moduli delle superficie regolari di genere \(p > 3\).

A tal uopo si consideri una superficie \(F = F_n\) dello spazio ordinario, di un certo ordine \(n\), dotata di singolarità normali, e la rappresentazione di \(F\) sopra un piano \(n\)-plo, che si ottiene per proiezione della \(F\) stessa da un punto esterno generico, \(O\). In questa rappresentazione la curva di diramazione \(D\) nasce per proiezione dalla curva \(C_j\) jacobiana della rete delle sezioni piane \(C\) con piani per \(O\). Ora sulla curva piana \(D\) la serie caratteristica completa del sistema continuo \(\{D\}\), costituito dalle curve dello stesso ordine collo stesso numero di nodi e di cuspidi, ha precisamente l’indice di specialità \(\omega\). Si tratta di ritrovare questa serie, di specialità \(\omega\), sulla curva jacobiana \(C_j\). A tal uopo si osservi che la \(C_j\) è sezione di \(F\) con la polare del punto \(O\), ossia (designando con \(K\) le curve canoniche di \(F\)) è una curva del sistema lineare \(|3C + K|\) segato sopra \(F\) dalle superficie \(\varphi_{n-1}\) aggiunte ad essa d’ordine \(n - 1\), che passano per i \(\tau\) punti cuspidali di \(F\), il cui gruppo, appartenente alla curva doppia di \(F\), verrà designato con \(G\). Le altre curve \(C_j\), passanti per i \(\tau\) punti cuspidali, segano sulla \(C_j\), fuori di essi, una serie lineare che,
% Printer's signature at the foot of the page.
% EDITORIAL: segnatura originale: \vadjust{\hbox to \linewidth{\small\textsc{Enriques F.} - \textit{Superficie algebriche.}\hfil 14}}
\printersignature{\textsc{Enriques F.} - \textit{Superficie algebriche.}\quad 14}
%CONT>>
% scan: s117L
% !PRINTED 210
\weberpage{210}
% head: CAPITOLO QUINTO
%<<CONT
% stamp: Dark scan-edge band at the right and scattered scan specks.
come subito vedremo, viene proiettata da \(O\) nella serie segata sulla curva di diramazione \(D\) dalle sue prime polari. Infatti, detta \(\overline{C}_j\) la sezione di \(F\) con la superficie polare di un punto \(\overline{O}\), il piano tangente ad \(F\) in un punto \(P\) comune a \(C_j\) e \(\overline{C}_j\) passa per la retta \(O\overline{O}\): quindi la tangente alla curva \(D\) nel punto \(P'\), proiezione di \(P\) da \(O\), passa per la traccia della retta \(O\overline{O}\) sul piano multiplo. % ?? uncertain: punctuation after the first \overline{C}_j -- comma or only the subscript descender.
% EDITORIAL D0115 RISOLTO: Enlargement shows only the descender of the subscript j; no separate comma follows it before la. Remove the transcription's extra comma.

Segue da ciò che la serie caratteristica di \(\{D\}\) sulla \(D\) proiezione di \(C_j\), nasce per proiezione dalla serie che viene segata su \(C_j\) dalle curve del sistema lineare \(|L| = |4C + K|\) passanti per \(G\) (ovverosia dalle superficie aggiunte \(\varphi_n\) passanti per i \(\tau\) punti del detto \(G\)). Dunque il carattere \(\omega\), che designa l’indice di specialità della predetta serie caratteristica di \(\{D\}\), si può anche interpretare come « indice di specialità della serie segata sopra la curva \(C_j\), che è una curva di \(|3C + K|\) per \(G\), dalle curve del sistema lineare \(|L| = |4C + K|\) passanti egualmente per \(G\) ». % The equals sign is printed both before and after the line break.

Questa \textit{interpretazione proiettiva del carattere \(\omega\)}, rispetto alla superficie \(F\), è valida egualmente per superficie regolari e irregolari. E conduce a valutare il numero dei moduli da cui dipende una classe di superficie \(F\) entro la più ampia famiglia continua di superficie che la contiene, in un modo che si accorda assai bene coll’intuizione.

Infatti si consideri il sistema continuo di tutte le superficie \(F\), aventi lo stesso ordine \(n\) e una curva doppia coi medesimi caratteri, su cui si trovi lo stesso numero \(\tau\) di punti cuspidali: le superficie di un tale sistema \(\{F_n\}\) cui appartenga \(F\), determinano su una di esse, e in particolare sulla data \(F\), un \textit{sistema lineare caratteristico}, costituito dalle intersezioni di \(F\) colle superficie del sistema infinitamente vicine, e questo sistema lineare (che a priori si potrebbe dubitare non completo) viene segato sulla \(F\) dalle superficie aggiunte \(\varphi_n\) d’ordine \(n\), tangenti nei punti del gruppo \(G\) ai piani cuspidali di \(F\).

Il nostro teorema d’esistenza cui si aggiunga l’ipotesi di lavoro del \S~12, ci assicura che il sistema continuo \(F_n\) segherà su una \(F\) il \textit{sistema lineare caratteristico completo}; e, reciprocamente, questo presupposto \textit{equivale alla detta ipotesi di lavoro}, e dà il modo di computare il numero dei moduli da cui dipende la classe di \(F\) entro \(\{F_n\}\). Naturalmente per avere superficie con \textit{moduli generali}, converrà partire da un modello proiettivo \(F\) di cui si possa assicurare a priori l’esistenza entro una famiglia continua di superficie coi dati caratteri; per esempio da una \(F\) che risponda ad un sistema lineare regolare di grado \(n\) e genere \(\pi\) convenientemente elevato, per cui sia soddisfatta la diseguaglianza \(n > \pi - 1\), ovvero anche — salvo il caso di genere piccolo o di qualche eccezione analoga — da una superficie canonica, avente per sezioni piane \(\infty^3\) curve \(K\). Vediamo a che conduca il computo dei moduli, riferendoci alle due scelte di
%CONT>>
% scan: s117R
% !PRINTED 211
\weberpage{211}
% head: INVARIANTI NUMERICI E PIANI MULTIPLI
%<<CONT
% stamp: Scattered specks; a stray slanting mark lies below the minus sign in the second displayed formula for S.
\(F\) sopra indicate e supponendo, per semplicità di discorso, che la superficie sia regolare.

Anzitutto colla prima scelta si ritrova la formula dei moduli del \S~11. Infatti si tratta di valutare, la dimensione \(R\) del sistema caratteristico di \(\{F_n\}\) e poi il numero \(S\) che indica quanti sono i parametri da cui dipendono le \(F_n\) del detto sistema continuo birazionalmente identiche ad una data \(F\); sarà quindi
\[
M = R - S.
\]

Per calcolare \(R\) valutiamo anzitutto la dimensione del sistema \(|L| = |4C + K|\) senza punti base, che è lecito supporre seghi su una \(C_j\) una serie completa e non speciale, e poi sottragghiamo da codesta dimensione \(\tau - \omega\), cioè il numero delle condizioni indipendenti che il gruppo dei \(\tau\) punti cuspidali di \(F\) impone alle curve \(L\) che debbano contenerlo. La dimensione di \(|L|\) varrà
\[
p + 4\pi + 6n - 4,
\]
e quindi la dimensione del sistema caratteristico di \(\{F_n\}\) sarà
\[
R - 1 = p + 4\pi + 6n - 4 - \tau + \omega,
\]
dove
\[
\tau = 2n + 8\pi - 12p + 2p^{(1)} - 22,
\]
sicchè
\[
R = 13p + 4n - 4\pi - 2p^{(1)} + 19 + \omega.
\]

D’altra parte si calcola \(S\) aggiungendo 15 (numero dei parametri di un’omografia spaziale) al numero dei parametri che definiscono un sistema lineare \(\infty^3\) entro il sistema lineare completo \(|C|\) di dimensione \(r\). Si ottiene così
\[
S = 4r - 12 + 15 = 4r + 3
\]
cioè
\[
S = 4p + 4n - 4\pi + 7.
\]
Pertanto si trova
\[
M = R - S = 9p - 2p^{(1)} + 12 + \omega,
\]
che è la formula data nel \S~11.

Ma, tenuto conto del significato di \(\omega\) (in accordo colla nostra ipotesi di lavoro del \S~12), si può dare una formula più significativa, mostrando che per \(p > 3\) si ha in generale
\[
\begin{aligned}
&\qquad\qquad \omega \ge p,\\
&M = 10p - 2p^{(1)} + 12 + \omega'\\
&\qquad\qquad\qquad\qquad\qquad (\omega' \ge 0).
\end{aligned}
\]

% scan: s118L
% !PRINTED 212
\weberpage{212}
% head: CAPITOLO QUINTO
% stamp: Dark scan-edge band at the right and scattered scan specks.
A tal uopo riprendiamo il calcolo precedente, riferendoci ad una superficie \(F_n\) d’ordine \(n = p^{(1)} - 1\) le cui sezioni piane siano curve canoniche \(K\). Qui occorrono le seguenti modificazioni:

1) le curve \(L\) di \(|L| = |5K|\) segano su una \(C_j = K_j\) la serie canonica completa, perciò il calcolo di \(R\), indicando ancora con \(\omega\) l’indice di specialità della serie segata sulla su \(K_j\) dalle curva di \(|5K|\) che passano per il gruppo \(G\), ci darà ancora % [sic] "sulla su" -> "sulla"; "dalle curva" -> "dalle curve", duplicated preposition and number disagreement.
\[
R - 1 = p + 4\pi + 6n - 4 - \tau + \omega
\]
da cui risulta che i \(\tau\) punti cuspidali impongono, non \(\tau - \omega\) ma \(\tau - (\omega - 1)\) condizioni indipendenti alle \(L\) che debbano contenerli:

2) il sistema \(|K|\), a differenza del sistema regolare non speciale \(|C|\) considerato innanzi, ha la dimensione
\[
\begin{aligned}
&\qquad p - 1 = p + n - \pi\\ % ?? uncertain: small mark after \pi -- scan speck (best reading) or full stop.
&(n = p^{(1)} - 1,\qquad \pi = p^{(1)}),
\end{aligned}
\]
anzichè
\[
p + n - \pi + 1;
\]
quindi il numero delle superficie \(F_n\) birazionalmente identiche ad una data risulta
\[
4p + 4n - 4\pi + 3 % ?? uncertain: small mark after 3 -- scan speck (best reading) or full stop.
\]
invece di
\[
4p + 4n - 4\pi + 7.
\]

Pertanto si troverà che il numero dei moduli di \(F_n\) viene dato da
\[
M = 9p - 2p^{(1)} + 16 + \theta
\]
dove \(\theta = \omega - 4\) designa l’indice di specialità della serie segata su una \(K_j\) dalle curve di \(|5K|\) passanti per il gruppo \(G\) dei punti cuspidali che appartengono alla superficie \(F\) immagine di un sistema lineare scelto entro \(|K|\).

Ora si può valutare \(\omega\), o almeno un suo limite inferiore, osservando che il gruppo \(G\) dei punti cuspidali è un gruppo di punti appartenenti alla curva \(K_j\) jacobiana di una rete di curve \(K\), ed intersezione parziale con curva una \(L\) del sistema % [sic] con curva una L -> con una curva L, transposed words.
\[
|L| = |5K|
\]
che è aggiunto a
\[
|K_j| = |4K|.
\]

Il gruppo \(G_\tau\) sarà dunque un gruppo speciale, ed imporrà alle curve di \(|5K|\) che debbono contenerlo, precisamente \(\tau - \theta\) condizioni, dove \(\theta\) sia la dimensione della serie completa cui appartiene il \(G_\tau\) sopra \(K_j\) (teorema di Riemann-Roch).

% scan: s118R
% !PRINTED 213
\weberpage{213}
% head: INVARIANTI NUMERICI E PIANI MULTIPLI
% stamp: Scattered scan specks.
Ma si ha un limite inferiore di \(\theta\), notando che la rete avente come jacobiana la \(K_j\) appartiene ad \(\infty^{p-4}\) sistemi lineari \(\infty^3\) di curve \(K\), contenuti entro il sistema canonico. Ciascuno di questi sistemi \(\infty^3\) definisce un gruppo \(G_\tau\) sopra \(K_j\), che è il gruppo dei punti doppi per \(\infty^1\) curve del sistema. Pertanto risulta
\[
\theta \ge p - 4
\]
e infine
\[
M = 10p - 2p^{(1)} + 12 + \omega' \qquad (\omega' \ge 0).
\]

Enunciamo il risultato ottenuto, tenendo presente che in esso gioca l’ipotesi di lavoro del \S~12, di cui occorrerebbe dare la dimostrazione:

\textit{Il numero dei moduli da cui dipende una classe di superficie regolari di genere \(p > 3\), con sistema canonico irriducibile semplice \((p^{(1)} > 5)\), è probabilmente}
\[
M \ge 10p - 2p^{(1)} + 12.
\]

Nel seguente paragrafo diremo come questa formula si estenda al caso delle superficie irregolari.

\edsection{14. Nota storica e complementi.}{14. Nota storica e complementi.}{sec.V.14}

Il primo tentativo di determinare il numero \(M\) dei moduli da cui dipende una classe di superficie algebriche è dovuto a M. \textsc{Noether}\origfootnote{1}{s118R-1}{M. \textsc{Noether}, \textit{Anzahl der Moduln einer Classe algebraischer Flächen.} Sitzungsberichte der Akademie, Berlin, 1888.} (1888) che ritiene valida, sotto larghe condizioni, la formula
\[
M = 10p - 2p^{(1)} + 12 \qquad (p = p_a = p_g),
\]
quella appunto che, secondo la nostra ipotesi di lavoro del \S~12, darebbe in generale un limite inferiore di \(M\) per le superficie regolari con \(p > 3\), corrispondendo ad \(\omega' = 0\). \textsc{Noether} è giunto a questa specie di divinazione mediante un intuito geniale. Egli è partito dalla formula che dà la postulazione di una curva gobba \(C\) (dotata di punti tripli) per le superficie \(F_n\) d’ordine \(n\), che debbano passare doppiamente per essa. Se il numero delle condizioni così imposte ad \(F_n\) è negativo, \(-R\), presume che \(R\) indichi il numero delle condizioni perchè la detta curva \(C\) sia doppia per una \(F_n\); quindi si troverà il numero dei moduli relativi alle \(F_n\) colla \(C\) doppia, valutando il numero \(N\) dei parametri da cui dipende la \(C\):
\[
M = N - R.
\]

% scan: s119L
% !PRINTED 214
\weberpage{214}
% head: CAPITOLO QUINTO
% stamp: Dark scan-edge band at the right and scattered scan specks.
La questione dei moduli delle superficie è stata ripresa da \textsc{Enriques} in due lavori, del 1908 e del 1912\origfootnote{1}{s119L-1}{F. \textsc{Enriques}, \textit{Sui moduli delle superficie algebriche.} Rend. Lincei, giugno 1908; \textit{Sui moduli di una classe di superficie algebriche e sul teorema d’esistenza per le funzioni algebriche di due variabili.} Att. R. Accademia di Torino, 1912. Cfr. Annali di Mat., 1923-24.}. La valutazione si fa considerando — nello spazio ordinario — il sistema continuo delle superficie \(F_n\) d’ordine \(n\) dotate di una curva doppia con dati caratteri, e ammettendo che il sistema caratteristico, segato su una \(F_n\) dalle \(F_n\) infinitamente vicine, sia completo: che in tale presupposto si contenga implicitamente un’ipotesi che non viene sufficientemente dimostrata accertando la corrispondente proprietà differenziale, non si accorse allora l’autore; la lacuna è stata segnalata soltanto da \textsc{Zariski} (1935)\origfootnote{2}{s119L-2}{O. \textsc{Zariski}, \textit{Algebraic Surfaces.} Berlino, 1935 (cfr. V, 8).} dopo che una lacuna analoga era stata trovata dal \textsc{Severi} (1921) nella dimostrazione del teorema fondamentale che esprime la proprietà caratteristica delle superficie irregolari (cfr. Cap. IX, \S~6).

Comunque, \textsc{Enriques}, basandosi sull’ipotesi implicita sopra accennata, trovava anzitutto la formula del \S~11:
\[
M = 10p_a - p_g - 2p^{(1)} + 12 + \omega \qquad (\omega \ge 0)
\]
valida per tutte le superficie regolari e irregolari che non appartengono alla famiglia delle rigate (cioè con qualche plurigenere \(P_i > 0\)); e in secondo luogo la formula più espressiva che dovrebbe valere, in generale, per le superficie regolari (con sistema canonico irriducibile semplice) il cui genere \(p > 3\) \((p^{(1)} \ge 6)\):
\[
M = 10p - 2p^{(1)} + 12 + \omega'; \qquad (\omega' \ge 0)
\]
questa formula, come abbiam detto, per \(\omega = 0\) si riduce a quella di Noether. % [sic] \omega = 0 -> \omega' = 0, the preceding formula uses \omega'.
Nella seconda Nota (del 1912) \textsc{Enriques} traeva dalla formula dei moduli una qualche veduta del teorema d’esistenza dei piani multipli, rilevando che le condizioni perchè una curva piana dotata di un certo numero di nodi e di cuspidi sia curva di diramazione d’un piano multiplo sono di natura puramente aritmetica, di guisachè si può dire che « se una curva \(D\) sia curva di diramazione d’un piano \(n\)-plo, lo stesso dovrà accadere per tutte le curve del più ampio sistema continuo che sono dotate dello stesso numero di nodi e di cuspidi ». L’autore avvertiva, in questa occasione, che siffatte condizioni aritmetiche potranno determinarsi collo studio del gruppo di monodromia, come ha realizzato di poi nella memoria degli Annali di Matematica del 1923 e qui si è esposto nel \S~9.

Ora è chiaro che, partendo — come si è fatto innanzi nel \S~11
%CONT>>
% scan: s119R
% !PRINTED 215
\weberpage{215}
% head: INVARIANTI NUMERICI E PIANI MULTIPLI
%<<CONT
% stamp: Scattered scan specks.
— dal teorema d’esistenza dei piani multipli, si riesce a dare la prima formula dei moduli che (esclusa la famiglia delle rigate) è assolutamente generale, e resta così immune dalla critica di Zariski. La quale tocca invece la seconda formula più espressiva (per \(p > 3\)), che in questa esposizione si è fatta dipendere dall’ipotesi di lavoro del \S~12.

La questione dei moduli delle superficie algebriche è stata ripresa da B. \textsc{Segre}\origfootnote{1}{s119R-1}{B. \textsc{Segre}, \textit{Sui moduli delle superficie algebriche irregolari.} Rendic. Acc. Lincei, vol. XIX (1934).} nel 1934, in ispecie coll’intento di trovare una formula più espressiva per le superficie irregolari. In questa ricerca però rimane il presupposto (dell’integrità del sistema caratteristico di un sistema continuo di superficie \(F_n\)) che noi abbiamo ricondotto all’ipotesi di lavoro del \S~12; giacchè il \textsc{Segre} non si era accorto della difficoltà messa in luce più tardi dalla critica di Zariski.

Qui notiamo che, sulla base del detto presupposto, il ragionamento svolto innanzi per il caso generale delle superficie regolari con \(p > 3\) si estende subito al caso delle superficie irregolari (tenendo presente la proprietà caratteristica di queste cui si è accennato). In tal guisa si perverrebbe alla formula
\[
M = 12p_a - 2p_g - 2p^{(1)} + 12 + \omega \qquad (\omega \ge 0)
\]

Ma il \textsc{Segre} tiene un’altra via che lo conduce ad un’altra formula più significativa, tantochè già nel caso delle superficie regolari è condotto a scrivere \(\omega \ge p - 1\), e quindi
\[
M \ge 10p - 2p^{(1)} + 11,
\]
formula che (salvo le restrizioni sopra accennate) dovrebbe valere per tutte le superficie con qualche plurigenere non nullo, senza che faccia bisogno di ammettere \(p > 3\), curve canoniche irriducibili, ecc.

Del resto questo risultato rientra nella formula più generale che dà i moduli delle superficie irregolari. Infatti B. \textsc{Segre} trova \(\omega \ge 2p_g - p_a - 1\), e quindi
\[
M = 9p_a + p_g - 2p^{(1)} + 11 + \theta % ?? uncertain: tiny mark after \theta -- scan speck (best reading) or full stop.
\]

Per chiudere riportiamo una formula comunicataci dal Prof. G. \textsc{Castelnuovo} il quale per il carattere \(\omega\) ha trovato la seguente disuguaglianza:
\[
\omega \ge 2p_a + p_g + p^{(1)} - I - 3.
\]
% scan: s120L
% !PRINTED none
\weberblank{[s120L]}
% head: none
% blank page
% stamp: Dark scan-edge band at the right and scattered scan specks.


\clearpage\thispagestyle{plain}
% scan: s120R
% !PRINTED none
\weberpage{[s120R]}
% head: none
% stamp: Scattered scan specks and a narrow scan-edge mark at the upper right.
\edchapter{VI}{SUPERFICIE REGOLARI: MINIMO DEI GENERI\\E CONDIZIONI DI RAZIONALITÀ}{Capitolo VI. SUPERFICIE REGOLARI: MINIMO DEI GENERI E CONDIZIONI DI RAZIONALITÀ}{SUPERFICIE REGOLARI: MINIMO DEI GENERI E CONDIZIONI DI RAZIONALITÀ}

In questo capitolo vogliamo considerare per semplicità superficie regolari di genere
\[
p_g = p_a = p,
\]
sebbene molti risultati che a queste si riferiscono si estenderanno senza difficoltà alle superficie irregolari; come avremo spesso occasione di avvertire.

\edsection{1. Limite inferiore del genere lineare.}{1. Limite inferiore del genere lineare.}{sec.VI.1}

Consideriamo una superficie \(F\) di genere \(p > 0\), che possiamo supporre priva di curve eccezionali (Cap. IV, \S~9). La \(F\) possiede per ipotesi delle curve canoniche che supporremo dapprima essere di ordine maggiore di zero. Vogliamo dimostrare che il genere lineare di \(F\) (che è qui il genere lineare assoluto) vale
\[
p^{(1)} > 0.
\]

Occorre riconoscere che l’ipotesi \(p^{(1)} \le 0\) conduce ad un assurdo. Sia dunque \(p^{(1)} \le 0\), e quindi il grado del sistema canonico \(|K|\)
\[
p^{(2)} = p^{(1)} - 1 < 0.
\]

Avremo così (almeno) una curva canonica \(K\) di grado negativo, la quale potrà essere costituita da più parti semplici o multiple:
\[
K = s_1K_1 + s_2K_2 + \mathinner{\ldotp\ldotp\ldotp\ldotp} + s_rK_r,
\]
dove le \(K_i\) non possono essere curve di genere virtuale \(\varrho_i = 0\) e di grado \(\nu_i = -1\), che sarebbero curve eccezionali di prima specie sopra \(F\). Intersechiamo una delle componenti di \(K\), sia per esempio \(K_1\), con la \(K\). Il numero delle intersezioni sarà % ?? uncertain: letter forms -- \varrho or \rho for genus, and \nu or Latin v for degree.
\[
2\varrho_1 - 2 - \nu_1. % ?? uncertain: letter forms -- \varrho or \rho, and \nu or Latin v.
\]

% scan: s121L
% !PRINTED 218
\weberpage{218}
% head: CAPITOLO SESTO
% stamp: Dark vertical scan artefact at the right edge; scattered specks.
Questo numero non può essere negativo. Infatti se
\[
2\varrho_1 - 2 - \nu_1 < 0,
\]
essendo \(\varrho_1 \geq 0\), segue
\[
\nu_1 \geq -1
\]
e quindi (essendo esclusi i valori \(\nu_1 = -1\), \(\varrho_1 = 0\)),
\[
\nu_1 \geq 0.
\]

Ma questo è assurdo perchè una componente \(K_1\) di grado \(\nu_1 \geq 0\) non può avere un numero negativo di intersezioni con la \(K = s_1K_1 + s_2K_2 + \ldots\). Ora, avendo provato che le componenti \(K_i\) di \(K\) debbono avere con \(K\) un numero non negativo d’intersezioni, si deduce facilmente che la \(K\) stessa deve avere un grado \(p^{(2)} \geq 0\). % The plus sign is repeated at a printed line break.

Infatti, designando con \(n_{ij}\) il numero delle intersezioni di due componenti \(K_i\) e \(K_j\) della \(K\), il grado della \(K\) sarà espresso da
\[
p^{(2)} = \Sigma s_i^2\nu_i + 2\Sigma s_i s_j n_{ij}.
\]

L’espressione del \(p^{(2)}\) ci appare ora come una forma quadratica nelle \(s_i\) e, tenendo presente il teorema di Eulero sulle funzioni omogenee, si riconosce che essa deve essere essenzialmente non negativa perchè le sue derivate parziali rappresentanti il doppio del numero delle intersezioni \((K \cdot K_j)\) sono
\[
\frac{\partial p^{(2)}}{\partial s_i} = 2(s_i\nu_i + \Sigma s_j n_{ij}) \geq 0.
\]

La diseguaglianza \(p^{(1)} > 0\) sussiste ancora nel caso che la superficie \(F\) possegga una sola curva canonica d’ordine zero \((p = 1)\). Infatti si assuma sopra la \(F\), priva di curve eccezionali, un sistema lineare irriducibile \(|C|\) senza punti base di genere \(\pi\) e grado \(n\); avremo in questo caso
\[
n = 2\pi - 2,
\]
essendo il numero delle intersezioni \((C \cdot K) = 2\pi - 2 - n = 0\). E, designando con \(|C'|\) il sistema aggiunto di \(|C|\), di genere \(\pi'\) e di grado \(n'\), sarà
\[
|C'| = |C|.
\]

Ora, valutando il genere \(\pi' = \pi\) di \(|C'|\), si trova
\[
\pi' = \pi = \pi + p^{(1)} + 2\pi - 2 - n - 1
\]
e quindi
\[
p^{(1)} = 1.
\]

Il teorema dimostrato \((p^{(1)} \geq 1)\) sussiste anche per le superficie di genere \(p = 0\) che abbiano un plurigenere \(P_i > 0\). Infatti, si
%CONT>>
% scan: s121R
% !PRINTED 219
\weberpage{219}
% head: SUPERFICIE REGOLARI : MINIMO DEI GENERI E CONDIZIONI, ECC.
%<<CONT
% stamp: Scattered scan specks.
può ripetere il ragionamento, ove si ponga al posto della curva canonica \(K\) la curva \(i\)-canonica \(iK\), la quale se fosse \(p^{(1)} \leq 0\) dovrebbe avere con ciascuna delle sue componenti (di genere \(\varrho\) e grado \(\nu\)) un numero non negativo d’intersezioni, espresso da
\[
i(2\varrho - 2 - \nu).
\]

Anche nel caso \((P_i = 1)\) in cui la curva \(i\)-canonica sia d’ordine zero, e quindi un sistema \(|C|\) coincida col suo \(i\)-esimo aggiunto, si conclude come nel caso \(p = 1\), che
\[
p^{(1)} = 1.
\]

Infine avvertiamo che la dimostrazione del nostro teorema è affatto indipendente dall’ipotesi che la \(F\) sia regolare. Enunciamo dunque, senza restrizioni, che:

\textit{Una superficie di genere \(p_g > 0\), o anche di genere \(p_g = 0\) che abbia un plurigenere \(P_i > 0\), ha il genere lineare assoluto}
\[
p^{(1)} \geq 1.
\]

Vogliamo offrire una seconda dimostrazione del teorema sopra enunciato, la quale dipende da un lemma che riceverà nel seguito importanti applicazioni. Abbiasi una superficie \(F\), che vogliamo supporre priva di curve eccezionali, di genere lineare \(p^{(1)} \leq 0\), e si assuma sopra di essa un sistema lineare irriducibile regolare \(|C|\), di genere \(\pi\) e di grado \(n \leq 2\pi - 2\). Le \(C\) incontrano una curva canonica virtuale \(K\) in
\[
2\pi - 2 - n \geq 0
\]
punti. Invece le curve \(C'\) del sistema aggiunto \((|C'| = |C + K|)\) incontrano le \(K\) in
\[
2\pi - 2 - n + p^{(1)} - 1 < 2\pi - 2 - n
\]
punti. Così, passando da un sistema lineare al suo aggiunto, il numero delle intersezioni delle curve del sistema con la curva canonica va diminuendo. Ma ciò porta che il processo dell’aggiunzione stessa debba avere un termine, ovvero che si arrivi a sistemi lineari di genere \(\pi\) e grado \(N > 2\pi - 2\); ma l’esistenza di un siffatto sistema \(|L|\) porta ancora di necessità che la serie degli aggiunti \(|L'|\), \(|L''|\), ecc. debba avere un termine, perchè il numero delle intersezioni \((L \cdot L) > (L' \cdot L) > (L'' \cdot L)\ldots\).

Dunque l’ipotesi \(p^{(1)} \leq 0\) porta che, a partire da un sistema lineare \(|C|\), si ottenga soltanto una serie limitata di successivi aggiunti, e perciò che un sistema \(|C|\) non possa essere contenuto in nessuno dei suoi sistemi aggiunti \(|C^{(i)}|\), ossia che tutti i plurigeneri della superficie siano \(P_i = 0\).

% scan: s122L
% !PRINTED 220
\weberpage{220}
% head: CAPITOLO SESTO
% stamp: Dark vertical scan artefact at the right edge and a short scuff in the upper-left margin.
Qui occorrono due osservazioni. La prima è che l’ipotesi di una superficie \(F\) priva di curve eccezionali si introduce soltanto a scopo di riduzione all’assurdo; se si rigetta tale ipotesi si deve ammettere senz’altro che la superficie \(F\) possegga un sistema con \(n > 2\pi - 2\), a partire dal quale il processo di aggiunzione si estingue.

La seconda osservazione si riferisce alla eventuale riducibilità del sistema \(|C'|\) aggiunto al sistema irriducibile \(|C|\) da cui siamo partiti. Se la superficie \(F\) è regolare e \(|C|\) privo di curve fondamentali, è chiaro che \(|C'|\) non può essere riducibile, e così in particolare non può ridursi al sistema \(|C|\) cui si aggiunga una curva canonica fissa. Ma, indipendentemente dalla irriducibilità dei sistemi aggiunti successivi a \(|C|\), basta in realtà pervenire ad un sistema \(|L|\) comunque riducibile, di genere \(\pi\) e di grado \(N > 2\pi - 2\); ogni multiplo di questo sistema darà sempre un sistema lineare che soddisfa a fortiori alla medesima diseguaglianza. Riuscirà quindi agevole costruire un sistema irriducibile che soddisfi del pari ad una diseguaglianza analoga, sommando \(|sL|\) ad un conveniente sistema lineare irriducibile, su di che non giova qui trattenersi ulteriormente.

Frattanto concludiamo affermando il risultato sopra enunciato.

\edsection{2. Il bigenere e i plurigeneri.}{2. Il bigenere e i plurigeneri.}{sec.VI.2}

Riferiamoci ad una superficie \(F\) tale che per ogni sistema lineare di genere \(\pi\), sia il grado \(n \leq 2\pi - 2\), la quale, come abbiamo visto (Cap. IV, \S~9) può trasformarsi in guisa da eliminare le sue curve eccezionali. A partire da un sistema lineare irriducibile \(|C|\), che può essere il sistema delle sezioni iperpiane di una \(F\) priva di singolarità, cerchiamo di costruire il sistema bicanonico di \(F\). Occorre perciò staccare una \(C\) dal sistema secondo aggiunto \(|C''|\). Ora se si suppone che sia la differenza
\[
d = 2\pi - 2 - n \geq 0,
\]
le \(C''\) segheranno sopra \(C\) una serie d’ordine \(2\pi - 2 + d\) che è certo non speciale e perciò lo staccamento di \(C\) da \(|C''|\) importa al più \(\pi - 1 + d\) condizioni. Quindi, calcolando con il teorema di Riemann-Roch la dimensione di \(|C''|\), si trova che quella del sistema bicanonico \(|C'' - C|\) sarà espressa, da % [sic] espressa, da -> espressa da, intrusive comma between the participle and its complement.
\[
P - 1 \geq p_a + p^{(1)} - 1;
\]
avremo dunque
\[
P \geq p_a + p^{(1)}.
\]

Questo risultato è indipendente dall’esistenza effettiva delle curve canoniche (d’ordine \(d \geq 0\)), cioè sussiste per \(p_g \geq 0\). In particolare \textit{per le superficie regolari \((p = p_a = p_g)\), prive di curve ec}%
%CONT>>
% scan: s122R
% !PRINTED 221
\weberpage{221}%
% head: SUPERFICIE REGOLARI : MINIMO DEI GENERI E CONDIZIONI, ECC.
%<<CONT
% stamp: Scattered scan specks and a short dark mark below the running head.
\textit{cezionali, con curve canoniche virtuali d’ordine \(d\) \((\geq 0)\) il bigenere vale sempre}
\weq{1)}{P \geq p + p^{(1)}.}

Se vi sono curve canoniche irriducibili di genere \(p^{(1)} > 1\), su cui allora il sistema bicanonico sega la serie canonica completa, si ha esattamente:
\weq{2)}{P = p + p^{(1)}.}

In luogo dell’eguaglianza 2) si ha \textit{in generale} la diseguaglianza
\weq{3)}{P \geq p + 1 \qquad \text{per} \qquad p^{(1)} = 1,}
giacchè il \textit{sistema canonico} di grado \(p^{(1)} - 1 = 0\) si spezza allora nelle \textit{curve ellittiche \(k\) d’un fascio} (lineare sulla \(F\) regolare): se una curva canonica è composta di \(p - 1\) curve \(k\) d’un tal fascio, le curve bicanoniche sono formate da \(2p - 2\) curve \(k\) e danno quindi un sistema lineare di dimensione
\[
P - 1 = 2p - 2;
\]
se alle \(p - 1\) curve \(k\) si aggiunge qualche curva fissa, il doppio di questa può costituire una curva \(k\) suscettibile di variare nel fascio di esse; e perciò può riuscire
\[
P - 1 > 2p - 2.
\]

I risultati precedenti subiscono qualche modificazione nel caso \(p = 1\), quando c’è \textit{una curva canonica effettiva d’ordine \(d = 0\)}: si ha allora in luogo della 1):
\[
P = 1
\]
e
\[
p^{(1)} = 1 \qquad (P < p + p^{(1)}).
\]

Infatti se è \(d = 0\), il sistema \(|C|\) considerato sopra la superficie \(F\) è di grado
\[
n = 2\pi - 2,
\]
ed essendo \(p = 1\) si ha
\[
|C| = |C'| = |C''|;
\]
quindi il sistema bicanonico coincide col \textit{sistema d’ordine zero} che costituisce il sistema canonico e che ha la dimensione
\[
p - 1 = P - 1 = 0;
\]
il genere di tale sistema \(|K|\) vale
\[
p^{(1)} = 1,
\]
%CONT>>
% scan: s123L
% !PRINTED 222
\weberpage{222}
% head: CAPITOLO SESTO
%<<CONT
% stamp: Dark vertical scan artefact at the right edge; scattered specks and short stray marks.
perchè sommando una \(K\) d’ordine zero ad una \(C\) si ha una curva composta di genere
\[
\pi = \pi + p^{(1)} + 0 - 1.
\]

Così \textit{le superficie regolari di genere \(p = 1\) che posseggono una curva canonica d’ordine zero si distinguono da quelle con curve canonica d’ordine maggiore di zero per avere: le prime} % [sic] curve canonica -> curva canonica, the printed noun and adjective disagree in number.
\[
P = 1,
\]
\textit{e le seconde}
\[
P \geq p + p^{(1)} \geq 2.
\]

Per \textit{una superficie regolare senza curve eccezionali \((n \leq 2\pi - 2)\) di genere \(p = 0\), con curva canonica virtuale d’ordine \(d = 0\) \((d = 2\pi - 2 - n)\) sussiste ancora, come per \(d > 0\),} la \textit{diseguaglianza} 1) \textit{cioè} % The minus sign is repeated at a printed line break.
\[
P \geq p^{(1)},
\]
\textit{avendosi}
\[
P = 1,\qquad p^{(1)} = 1.
\]

Infatti, essendo \(|C'|\) diverso da \(|C|\), anche \(|C''|\) sarà distinto da \(|C'|\) e segherà sopra una \(C\) una serie \(g_{2\pi-2}\) non speciale; di conseguenza lo staccamento di \(C\) da \(|C''|\) importerà \(\pi - 1\) condizioni lineari al più: si deduce che esso importa proprio \(\pi - 1\) condizioni, risultando
\[
|C''| = |C|,\qquad P = 1,\qquad p^{(1)} = 1.
\]

\textit{Osservazione.} -- È importante osservare che i risultati precedenti, in ispecie la \textit{diseguaglianza} fondamentale
\weq{1)}{P \geq p + p^{(1)},}
cadono assolutamente in \textit{difetto} per le superficie contenenti sistemi lineari \(|C|\) di genere \(\pi\) e grado \(n > 2\pi - 2\); ossia aventi curve canoniche virtuali d’ordine negativo
\[
-d = n - (2\pi - 2).
\]

Infatti il sistema \(|C''|\) segherà ora sopra una \(C\) una serie d’ordine \(2\pi - 2 - d\) che potrà essere speciale.

Valga ad esempio il caso del piano, in cui si assume il sistema lineare \(|C_7|\) formato dalle curve generali del \(7^\circ\) ordine. Le curve aggiunte sono le quartiche \(C_4\), e le seconde aggiunte le rette \(C_1\), secanti sopra una \(C_7\) una serie \(g_7^2\), serie speciale d’indice di specialità
\[
i = 10.
\]

Qui si può calcolare il genere (lineare \(p^{(1)}\)) del sistema canonico virtuale
\[
|K| = |C_1 - C_4|.
\]
%CONT>>
% scan: s123R
% !PRINTED 223
\weberpage{223}
% head: SUPERFICIE REGOLARI : MINIMO DEI GENERI E CONDIZIONI, ECC.
%<<CONT
% stamp: Scattered scan specks and short scuffs near the upper margin.
Siccome le \(K\) segano le \(C_1\) in \(-d = -3\) e le \(C_4\) (di genere 3) in \(-12\) punti, si avrà
\[
0 = p^{(1)} + 3 - 12 - 1
\]
e quindi ritroviamo che il \textit{genere lineare (relativo) del piano vale:}
\[
p^{(1)} = 10.
\]

Frattanto, non essendo possibile staccare \(|C|\) da \(|C_1|\) manca affatto il sistema bicanonico e perciò
\[
P = 0.
\]

Le formule che esprimono il bigenere per mezzo del genere superficiale e lineare, si estendono ai plurigeneri. Quando la superficie \(F\) priva di curve eccezionali possegga \textit{curve canoniche virtuali d’ordine \(d > 0\)} \((d = 2\pi - 2 - n)\) \textit{si ha}
\[
P_i \geq p_a + \frac{i(i - 1)}{2}(p^{(1)} - 1) + 1
\]
e per una superficie regolare \((p_a = p_g = p)\) di \textit{genere lineare \(p^{(1)} > 1\), con curve \((i - 1)\) canoniche irriducibili:}
\[
P_i = p + \frac{i(i - 1)}{2}(p^{(1)} - 1) + 1.
\]

\textit{L’eguaglianza diventa diseguaglianza nel caso \(p^{(1)} = 1\)}, dove le \textit{curve pluricanoniche} di grado \(i^2(p^{(1)} - 1) = 0\), si compongono tutte colle curve ellittiche di un medesimo fascio:
\[
P_i \geq p + 1.
\]

Fanno eccezione alla diseguaglianza fondamentale
\[
P_i \geq p + \frac{i(i - 1)}{2}(p^{(1)} - 1) + 1
\]
le superficie di bigenere \(P = P_2 = 1\) che posseggono una curva canonica virtuale d’ordine \(d = 0\) e una curva bicanonica effettiva, pure d’ordine zero. In questo caso la serie segata dal terzo aggiunto \(|C'''| = |C'|\) sopra una \(C\) è la serie canonica (di dimensione \(\pi - 1\)), quindi il suo staccamento da \(|C'''|\), importando \(\pi\) condizioni, non è possibile, sicchè
\[
P_3 = 0.
\]

Dunque \textit{le superficie regolari di genere zero e bigenere uno con curva bicanonica d’ordine zero si distinguono da quelle con curva bicanonica d’ordine maggiore di zero per avere:}

\textit{le prime il trigenere}
\[
P_3 = 0,
\]
%CONT>>
% scan: s124L
% !PRINTED 224
\weberpage{224}
% head: CAPITOLO SESTO
%<<CONT
% stamp: Dark vertical scan artefact at the right edge; scattered specks.
\textit{mentre le seconde hanno}
\[
P_3 \geq 1.
\]

Per le prime superficie si ha poi
\[
P_{2i} = 1,\qquad P_{2i+1} = 0,
\]
mentre per le seconde tutti i plurigeneri sono essenzialmente positivi. Anzi essi vengono a superare l’unità anche nel caso \(p^{(1)} = 1\). Infatti, combinando una curva bicanonica contata tre volte con una curva tricanonica contata due volte si ottiene un \textit{fascio di curve sesticanoniche}, sicchè
\[
P_6 \geq 2.
\]

\edsection{3. Esempi: piani doppi di genere lineare \(p^{(1)} = 1\).}{3. Esempi: piani doppi di genere lineare p\textasciicircum{}(1) = 1.}{sec.VI.3}

Una esemplificazione interessante delle formule del precedente paragrafo viene offerta dai piani doppi di genere lineare \(p^{(1)} = 1\)\origfootnote{1}{s124L-1}{Cfr. \textsc{F. Enriques}, \textit{Sui piani doppi di genere lineare \(p^{(1)} = 1\).} Rendic. R. Acc. Lincei, Aprile e Maggio 1898.}. Ivi si dimostra che ad esclusione del piano doppio con sestica di diramazione affatto generale
\[
z^2 = D_6(xy)
\]
che ha curva canonica d’ordine zero (genere e plurigeneri \(p = P_i = 1\)) gli altri piani doppi, per cui \(p^{(1)} = 1\) e qualcuno dei plurigeneri \(P_i > 1\), contengono un fascio (razionale o meno) di curve ellittiche \(L\), trasformato in sè dall’involuzione \(I\) le cui coppie rispondono ai punti del piano. E quindi si è condotti a due famiglie di superficie rappresentabili su tali piani doppi:

I) superficie con un fascio lineare di curve ellittiche \(L\) (irriducibili) trasformate in sè dall’involuzione \(I\) che subordina su di esse una involuzione lineare \(g_2^1\), le quali vengono rappresentate su piani doppi con curva di diramazione \(D_{2n}\) d’ordine \(2n\) \((n > 3)\) dotata d’un punto \((2n - 4)\)-plo o \((2n - 3)\)-plo \(O\), ed eventualmente di altre singolarità abbassanti il genere: punti quadrupli e punti \([3, 3]\); le rette per \(O\) sono le immagini sul piano doppio delle curve \(L\);

II) superficie con un fascio di curve ellittiche irriducibili \(L\), che possono essere curve d’un fascio razionale, trasformate in sè dall’involuzione \(I\) subordinante su di esse un’involuzione ellittica \(\gamma_2^1\), ovvero curve d’un fascio (razionale o meno) associate a due a due da \(I\) in un’involuzione razionale del fascio; queste superficie si rappresentano su piani doppi per modo che alle \(L\) rispondano le
%CONT>>
% scan: s124R
% !PRINTED 225
\weberpage{225}
% head: SUPERFICIE REGOLARI : MINIMO DEI GENERI E CONDIZIONI, ECC.
%<<CONT
% stamp: Scattered scan specks and short vertical edge marks.
curve ellittiche d’un fascio di Halphen\origfootnote{1}{s124R-1}{Cfr. p. es., \textsc{Enriques-Conforto}, \textit{Le superficie razionali.} (Zanichelli - Bologna).} (curve \(C_{3s}\) con 9 punti \(s\)-pli) e pertanto le curve di diramazione di essi (supposte prive di curve eccezionali) sono composte colle curve del detto fascio ed eventualmente anche colla cubica passante per i suoi nove punti base.

Fra i piani doppi della categoria I sono notevoli quelli con curva di diramazione \(D_{2n}\) dotata di coppie di punti tripli infinitamente vicini (punti \([3, 3]\)) allineati col punto \((2n - 4)\)-plo \(O\).

Per \(n = 4\) abbiamo il piano doppio con \(D_8\) di diramazione dotata di un punto 4-plo \(O\) e di due punti tripli infinitamente vicini \(A\) e \(A_1\), allineati con \(O\); la \(\mathrm{D}_8\) si compone di una \(D_7\) con \(O^3A^2A_1^2\) e della retta \(r = OA\). Il genere di questo piano doppio vale % ?? uncertain: typeface of the second D_8 -- upright D or italic D.
% EDITORIAL D0121 RISOLTO: At enlarged scale its vertical stem is upright, unlike the slanted D_8 and D_7 on adjacent lines; the supplied typeface is faithful.
\[
p_g = p_a = p = 1,
\]
avendosi \textit{una curva canonica} rappresentata dalla retta \(r = OA\) (soggetta a passare semplicemente per \(O\) e \(A\)). C’è qualche difficoltà a determinare il genere lineare \(p^{(1)}\), giacchè alla retta doppia \(r = OA\), sembra rispondere soltanto una coppia di curve razionali coincidenti sconnesse. Si scioglie la difficoltà riconoscendo che la superficie \(F\) rappresentata per proiezione sul piano doppio possiede una curva canonica impura, costituita da due curve eccezionali infinitamente vicine i cui punti rispondono ai punti generici di \(r\), e da una curva infinitesima rappresentata dall’intorno del punto \(A_1\); la quale, per una trasformazione di \(F\) in una superficie priva di singolarità, si muta in una curva d’ordine maggiore di zero e di genere \(p^{(1)} = 1\). D’accordo con questo riconoscimento la \(F\) avrà il bigenere
\[
P = p + p^{(1)} = 2,
\]
essendovi un fascio di curve bicanoniche ellittiche rappresentate sul piano doppio dalle rette per \(O\): aggiungendo le curve eccezionali si avranno come immagini delle curve bicanoniche impure le coppie di rette formate da \(r\) e da una retta variabile per \(O\). Ciò che si è detto si rende chiaro considerando il sistema lineare \(\infty^3\) delle coniche per \(AA_1\), che porta a trasformare la \(F\) in un cono quadrico doppio senza curve eccezionali: le immagini di codeste coniche doppie, con 10 punti di diramazione, sono curve iperellittiche di genere 4, formanti un sistema lineare di grado 4 \((< 2 \cdot 4 - 2)\); le quali vengono segate in 2 punti (coniugati) dalla curva canonica che risponde all’intorno di \(A_1\), e in 4 punti, costituiti da due coppie di punti coniugati nella \(g_2^1\), dalle curve bicanoniche; perciò queste ultime sono
% Printer's signature at the foot of the page.
% EDITORIAL: segnatura originale: \hbox to \textwidth{\textsc{Enriques F.} - \textit{Superficie algebriche.}\hfil 15}
\printersignature{\textsc{Enriques F.} - \textit{Superficie algebriche.}\quad 15}
%CONT>>
% scan: s125L
% !PRINTED 226
\weberpage{226}
% head: CAPITOLO SESTO
%<<CONT
% stamp: Dark vertical scan artefact at the right edge; scattered specks.
rappresentate sul cono dalle sezioni piane d’un fascio, e sul piano dalle rette per \(O\). Precisiamo che il detto cono quadrico doppio possiederà una curva di diramazione \(D_{10}\) del \(10^\circ\) ordine, di cui la \(D_7\) si può ritenere come proiezione da un punto triplo \(P\) a cui è infinitamente vicino un’altro punto triplo \(P_1\) che si proietta in \(O\). La retta \(PP_1\) giace sul piano tangente al cono lungo la generatrice \(k\) per \(P\) e i piani per essa segano sul cono le immagini delle curve bicanoniche. % [sic] un’altro -> un altro, apostrophe before a masculine form.

Accanto al piano doppio menzionato innanzi, che ha una curva di diramazione \(D_8 = rD_7\) con \(O^4A^3A_1^3\), si può considerare il piano doppio con \(D_8\) di diramazione passante per \(O^4\) e per due coppie di punti tripli infinitamente vicini allineati con \(O\): \(A^3A_1^3\) e \(B^3B_1^3\). Questa \(D_8\) è spezzata nelle due rette \(r = OA\) ed \(s = OB\) e in una \(D_6(O^2A^2A_1^2B^2B_1^2)\). Il genere di tale piano doppio è
\[
p = 0,
\]
e il bigenere
\[
P = 1,
\]
essendovi una curva bicanonica impura rappresentata dalla coppia delle rette \(r\) e \(s\). Questa curva si riduce ad una coppia di curve eccezionali contate due volte. Valutando il trigenere del piano doppio si troverebbe \(P_3 = 0\). Ciò risulta indirettamente dalle osservazioni che seguono.

Se si trasforma il piano doppio in una superficie \(F\) senza curve eccezionali, appare che la curva bicanonica di \(F\) è d’ordine zero: si può costruire \(F\) mercè il sistema lineare \(\infty^4\) delle cubiche per \(OAA_1BB_1\), che sono immagini di curve di genere 4 e grado 6; si è condotti in tal guisa ad una superficie cubica doppia (normale in \(S_4\)) con curva di diramazione del \(10^\circ\) ordine.

Vedremo nel seguito che la \(F\) si può trasformare anche nella superficie \(F_6\) del \(6^\circ\) ordine che passa doppiamente per gli spigoli d’un tetraedro.

Consideriamo ora i piani doppi (ancora della categoria I) con curva di diramazione \(D_{10}\) del \(10^\circ\) ordine, dotata d’un punto 6-plo \(O\) e di due coppie di punti tripli infinitamente vicini \(AA_1\) e \(BB_1\) allineati con \(O\): la \(D_{10}\) è composta delle due rette \(r = OA\) e \(s = OB\) e di una \(D_8(O^4A^2A_1^2B^2B_1^2)\). Qui si ha una curva canonica impura rappresentata dalle due rette eccezionali \(r\) ed \(s\), e da una curva residua (che diventa d’ordine \(> 0\) sopra una superficie trasformata senza singolarità) formata di due componenti ellittiche che rispondono agli intorni dei punti \(A_1\) e \(B_1\), sicchè
\[
p_g = p_a = p = 1,\qquad p^{(1)} = 1.
\]

% scan: s125R
% !PRINTED 227
\weberpage{227}
% head: SUPERFICIE REGOLARI : MINIMO DEI GENERI E CONDIZIONI, ECC.
% stamp: Scattered scan specks.
Il bigenere \(P\) deve superare l’unità e difatti si ha una rete di curve bicanoniche che hanno per immagini sul piano doppio le coppie di rette per \(O\); si conclude dunque che
\[
P = 3 \qquad (> p + p^{(1)}).
\]

Accanto al piano doppio menzionato si consideri quello con curva di diramazione \(D_{10}\) dotata di un punto 6-plo \(O\) e di 3 coppie di punti 3-pli infinitamente vicini \(AA_1\), \(BB_1\), \(EE_1\), allineati con \(O\). La superficie \(F\) così definita non possiede curva canonica, non essendovi alcuna conica passante doppiamente per \(O\) e semplicemente per \(A\), \(B\), \(E\); quindi % ?? uncertain: final vowel of “coppie” -- “coppie” or “coppia”.
% EDITORIAL D0122 RISOLTO: The final letter has the open right side and horizontal bar of e, clearly visible in the enlarged scan.
\[
p_a = p_g = p = 0.
\]

Invece si hanno curve bicanoniche ellittiche rappresentate dalle rette per \(O\), cui si aggiungono le immagini eccezionali delle rette di diramazione
\[
r = OA,\qquad s = OB,\qquad t = OE;
\]
perciò
\[
p^{(1)} = 1
\]
e
\[
P = 2 \qquad (> p^{(1)}).
\]

Fra codeste curve bicanoniche vi sono tre curve doppie (in corrispondenza alle tre rette \(r\), \(s\), \(t\) o meglio agli intorni di \(A_1\), \(B_1\) ed \(E_1\)), le quali prese insieme semplicemente costituiscono una curva tricanonica della superficie \(F\):
\[
P_3 = 1 \qquad (= p^{(1)}).
\]

Il doppio della detta curva tricanonica appartiene al triplo del sistema bicanonico; si ottiene così il sistema lineare delle curve sesticanoniche di \(F\), le quali hanno per immagini le terne di rette per \(O\). Si ha dunque:
\[
P_6 = 4.
\]

Nella categoria II troviamo pure esempi interessanti per la nostra teoria. Si assuma un piano doppio con curva di diramazione \(D_{6n}\) \((n > 1)\) costituita da una curva ellittica irriducibile con 9 punti \(2n\)-pli, appartenente --- com’è noto --- ad un fascio di Halphen\origfootnote{1}{s125R-1}{Ogni curva (doppia) del fascio di Halphen rappresenta due curve \(L\) distinte non appartenenti all’involuzione \(I\) giacchè ciò accade per la \(D_{6n}\); del resto ciò è d’accordo col criterio di \textsc{Comessatti}, poichè la \(D_{6n}\) può ridursi per continuità ad una \(\varphi_3^n\) doppia. Cfr. p. es. \textsc{Enriques-Chisini}, Libro V, cap. IV, \S~38 (vol. III, pag. 445).}. La superficie \(F\) così definita possiede una curva canonica avente per % ?? uncertain: superscript on \varphi_3 in the footnote -- n or m.
% EDITORIAL D0123 RISOLTO: The superscript has two stems joined by a single arch, consistent with n, not the extra arch of m. Mathematical context is consistent but is not used as sole evidence.
%CONT>>
% scan: s126L
% !PRINTED 228
\weberpage{228}
% head: CAPITOLO SESTO
%<<CONT
% stamp: Dark vertical scan-edge artefact at the right margin; scattered specks.
immagine una curva d’ordine \(3n-3\) soggetta a passare \(n-1\) volte per i 9 punti base del fascio, la quale si riduce alla cubica \(\varphi_3\) pei 9 punti contata \(n-1\) volte; diciamo alla \(\varphi_3^{n-1}\). Si ha pertanto
\[
p_a=p_g=p=1,\qquad p^{(1)}=1.
\]

Ora la curva di diramazione \(D_{6n}\), presa insieme con la \(\varphi_3^{n-2}\) aggiunta alla \(\varphi_3^{n-1}\) (da contare due volte) formerà una curva bicanonica di \(F\) (cfr. Cap. V, § 7), mentre un’altra curva bicanonica è data dalla \(\varphi_3^{2n-2}\) da contare due volte. Si deduce che le curve del fascio di Halphen definito dalla \(D_{6n}\), prese insieme colla \(\varphi_3^{n-2}\) fissa, rappresentano un fascio di curve bicanoniche, costituite dunque da curve ellittiche irriducibili \(L\)\origfootnote{1}{s126L-1}{Ogni curva \(L\) di questo fascio è immagine di \textit{due} curve ellittiche distinte in accordo col criterio di \textsc{Comessatti}, perchè il gruppo dei punti critici apparenti segato su \(L\) dalla \(D_{6n}\) viene segato egualmente dalla \(D_3^{2n}\), doppio della \(D_3^n\). Cfr. \textsc{Enriques-Chisini}, l. c.} (\textit{non} appartenenti all’involuzione \(I\)) e da una parte fissa, che è una curva ellittica multipla secondo \(n-2\), la cui componente appartiene all’involuzione \(I\) ed equivale ad una parte \((n-1)\)-ma di \(L\). Così avremo
\[
P=2\qquad (=p+p^{(1)}).
\]

Si trova poi facilmente che le curve tricanoniche sono date dalle coppie di curve \(L\), sicchè
\[
P_3=3.
\]

Accanto al precedente troviamo ancora nella categoria II il piano doppio con curva di diramazione \(D_{6n}=D_{6n-3}+D_3\), composta d’una curva irriducibile \(D_{6n-3}\) di Halphen, con 9 punti \((2n-1)\)-pli e dalla cubica \(D_3\) che passa per questi 9 punti.

La superficie \(F\) così definita possiede \textit{una} curva canonica rappresentata dalla \(D_3\) contata \(n-1\) volte, diciamo \(D_3^{n-1}\), sicchè
\[
p_a=p_g=p=1,\qquad p^{(1)}=1.
\]

Le curve bicanoniche trasformate in sè dall’involuzione \(I\) sono date: dalla \(D_3^{n-1}\) doppia, contata due volte, e dalla curva di diramazione \(D_{6n-3}+D_3\) presa insieme all’aggiunta dell’immagine della curva canonica, che è la \(D_3^{n-2}\) doppia. Quindi si ha un fascio di curve bicanoniche ellittiche rappresentate dalle curve d’ordine \(6n-3\) del fascio di Halphen semplici\origfootnote{2}{s126L-2}{Ved. nota precedente.}, prese insieme colla curva multipla rappresentata dalla \(D_3^{n-2}\) doppia:
\[
P=2.
\]

% scan: s126R
% !PRINTED 229
\weberpage{229}
% head: SUPERFICIE REGOLARI : MINIMO DEI GENERI E CONDIZIONI, ECC.
% stamp: Dark horizontal scan-edge artefact at the upper right margin; scattered specks.
Non ci soffermeremo sui piani doppi della categoria II aventi una curva di diramazione comunque composta colle curve d’un fascio di Halphen: i quali portano in generale a superficie irregolari con
\[
p_a=1,\qquad p_g>1,\qquad p^{(1)}=1. % ?? uncertain: The subscript on the first p is indistinct -- a; i.
\]

Ma vogliamo indicare un piano doppio con \(p^{(1)}=1,\ p=0,\ P=1\), dotato di curva bicanonica d’ordine maggiore di zero, che possiede un fascio di curve ellittiche (non autoconiugate in \(I\)) rappresentate dalle curve d’un fascio di Halphen, e nondimeno esce dalla categoria II per la presenza di una retta parte della curva di diramazione, immagine di una curva eccezionale. È un piano doppio già presentatosi ad \textsc{Enriques}, che ha una curva di diramazione \(D_{10}\) d’ordine 10 e s’incontra nella classificazione di questi piani di L. \textsc{Campedelli}\origfootnote{1}{s126R-1}{L. \textsc{Campedelli}, \textit{Sui piani doppi con curva di diramazione del 10\textsuperscript{o} ordine} e \textit{Sopra alcuni piani doppi notevoli con curva di diramazione del 10\textsuperscript{o} ordine}. R. Acc. Lincei, 1932.}: la \(D_{10}\) si compone di una \(D_9+r\), dove la \(D_9\) è una curva ellittica (irriducibile) di Halphen con 9 punti tripli, sei dei quali \(AA_1,\ BB_1,\ EE_1\) sono a coppie infinitamente vicini mentre gli altri tre (distinti) si trovano sulla retta \(r\).

La superficie \(F\) rappresentata su questo piano doppio non possiede curve canoniche, poichè una conica passante per i tre punti tripli di \(D_9\) su \(r\) e per i tre punti \(A,\ B,\) e \(E\) si spezzerebbe nella \textit{curva eccezionale} che risponde alla retta \(r\)\origfootnote{2}{s126R-2}{È facile vedere che la curva corrispondente a codesta retta di diramazione di genere 0 e grado \(-2\) ha sopra \(F\) il genere 0 e il grado \(-1\).} e in una retta per \(A,\ B\) e \(E\), che invece non sono allineati.

Si ha dunque
\[
p=0.
\]

D’altra parte una retta doppia per \(A,\ B\) e \(E\), se esistesse, avrebbe 4 punti di diramazione (tre nei detti tre punti e uno su \(r\)), quindi il genere lineare assoluto di \(F\) (genere delle curve canoniche private della curva eccezionale) vale
\[
p^{(1)}=1.
\]

Ora sulla \(F\) si ha \textit{una} curva bicanonica irriducibile\origfootnote{3}{s126R-3}{L’irriducibilità della cubica \(D_3\) doppia risulta dal criterio di \textsc{Comessatti}, tenuta presente la costruzione del fascio di Halphen cui appartiene la \(D_9\).}, che è rappresentata sul piano doppio dalla cubica \(D_3\) passante per i 9 punti tripli di \(D_9\) cui si aggiunge la curva eccezionale rappresentata dalla retta \(r\). Dunque il bigenere di \(F\) vale
\[
P=1.
\]

% scan: s127L
% !PRINTED 230
\weberpage{230}
% head: CAPITOLO SESTO
% stamp: Dark vertical scan-edge artefact at the right margin; scattered specks.
In conformità del nostro teorema si dovrà avere il trigenere \(P_3\ge1\) e il sestigenere \(P_6\ge2\).

Per verificarlo si osservi che (la \(D_3\) possedendo una curva aggiunta d’ordine 0) una curva tricanonica di \(F\) verrà data dalla curva di diramazione \(D_9\) cui si aggiunga la curva eccezionale \(r\). Codesta curva tricanonica è isolata, cioè non suscettibile di variare in un fascio, poichè altrimenti, non avendo punti comuni colla curva bicanonica rappresentata da \(D_3\), dovrebbe contenerla e si avrebbe \(p_g>0\). Vuol dire che le curve ellittiche doppie di 9\textsuperscript{o} ordine, del fascio di Halphen definito dalla \(D_9\), sono immagini di curve irriducibili (autoconiugate nell’involuzione \(I\) definite su \(F\) dal piano doppio). % [sic] definite -> definita, apparently modifying the singular "involuzione".
Questa deduzione si può anche confermare direttamente col criterio di Comessatti. Infine abbiamo \textit{una} curva tricanonica \(D_3\) e un fascio di curve sesticanoniche irriducibili rappresentate dalle curve del fascio di Halphen cui appartengono la curva tricanonica contata due volte e la curva bicanonica contata tre volte. Dunque si ha: % [sic] D_3 -> D_9, the tricanonical curve is identified above with the branch curve D_9.
\[
P_3=1,\qquad P_6=2.
\]

\edsection{4. Condizioni di razionalità di una superficie.}{4. Condizioni di razionalità di una superficie.}{sec.VI.4}

Una superficie razionale, cioè riferibile punto per punto al piano, deve avere come questo il genere \(p=p_a=p_g=0\) e insieme col genere debbono annullarsi tutti i suoi plurigeneri. Gli esempi che abbiamo addotto (in particolare la superficie del 6\textsuperscript{o} ordine che passa doppiamente per gli spigoli di un tetraedro) mostrano che una superficie regolare di genere \(p=0\) può avere il bigenere \(P>0\), sicchè l’annullamento del genere non basta a fornire le condizioni di razionalità di una superficie. Ora è estremamente interessante riconoscere che queste condizioni possono esprimersi annullando oltre al genere anche il bigenere, cioè scrivendo
\[
p_a=P=0
\]
(da cui segue necessariamente anche \(p_g=0\)). Questo teorema è dovuto a G. \textsc{Castelnuovo} che lo ha dato in una memoria del 1896\origfootnote{1}{s127L-1}{G. \textsc{Castelnuovo}, \textit{Sulle superficie di genere zero}. Soc. It. delle Scienze (detta dei XL), cfr. \textit{Memorie scelte}, pag. 307.}.

La dimostrazione del nostro teorema si lascia ricondurre ai criteri di razionalità di Noether e di Clebsch-Noether, relativi alle superficie con un fascio lineare di curve razionali e ai tipi di piani doppi da essi determinati. A tale scopo si parte da un sistema lineare irriducibile \(\lvert C\rvert\) sopra la superficie data e si costruiscono suc%
%CONT>>
% scan: s127R
% !PRINTED 231
\unskip\weberpage{231}%
% head: SUPERFICIE REGOLARI : MINIMO DEI GENERI E CONDIZIONI, ECC.
%<<CONT
% stamp: Scattered scan specks and short marks in the margins.
cessivamente il suo sistema aggiunto \(\lvert C'\rvert\), il secondo aggiunto \(\lvert C''\rvert\) e così di seguito. Si riconoscerà che, essendo \(p_a=P=0\) la serie di questi aggiunti successivi ha un termine, e l’ultimo sistema aggiunto \(\infty^1\) almeno riesce composto di curve razionali o ellittiche e conduce quindi ad uno dei casi di razionalità contemplati nei criteri di Noether e di Clebsch-Noether di cui si è discorso.

La discussione si può svolgere in forma assai semplice distinguendo due ipotesi: 1\textsuperscript{o}) la superficie data \(F\) contiene un sistema lineare irriducibile \(\lvert C\rvert\) di genere \(\pi\) e di grado \(n>2\pi-2\), che può supporsi essere il sistema delle sezioni piane o iperpiane di una \(F\), affatto priva di singolarità; 2\textsuperscript{o}) per ogni sistema lineare irriducibile dato sopra la superficie \(F\) sussiste la diseguaglianza \(n\le2\pi-2\), designando \(n\) e \(\pi\) il grado e il genere del sistema.

\textit{Ipotesi I.} -- È chiaro che costruendo i successivi aggiunti \(\lvert C'\rvert,\ \lvert C''\rvert\), ecc. a partire dal sistema \(\lvert C\rvert\) dato sopra \(F\), si ottiene una serie limitata di sistemi; in altre parole il procedimento d’aggiunzione si estingue. Infatti le curve \(C',\ C'',\ \ldots\ldotp\) avranno un ordine decrescente
\[
2\pi-2,\qquad 2\pi-2-d,\qquad 2\pi-2-2d,\qquad \ldots\ldotp
\]
essendo
\[
d=n-(2\pi-2).
\]

Pertanto vogliamo considerare l’ultimo sistema aggiunto \(\lvert C^{(t)}\rvert\) di dimensione \(r_t\ge1\).

Supponiamo dapprima che i sistemi aggiunti \(\lvert C'\rvert,\ \lvert C''\rvert,\ \ldots\ldotp,\ \lvert C^{(t)}\rvert\) della serie anzidetta siano irriducibili. Il sistema \(\lvert C^{(t)}\rvert\), primo della serie che non ha un aggiunto di dimensione maggior di zero, è costituito di curve di genere zero o uno. Se il sistema \(\lvert C^{(t)}\rvert\) è formato di curve di genere zero, il teorema di Noether\origfootnote{1}{s127R-1}{Cfr. p. e. \textsc{Enriques-Conforto}, Op. cit.} sulle superficie contenenti un fascio lineare di curve razionali permette di dedurre senz’altro che la \(F\) è razionale.

Se invece il sistema \(\lvert C^{(t)}\rvert\) è di genere 1, conviene distinguere l’ipotesi che la sua dimensione sia \(r_t\ge2\), ovvero \(r_t=1\). Se \(r_t\ge2\) il sistema completo di curve ellittiche \(\lvert C^{(t)}\rvert\), che, per la regolarità della superficie, ha serie caratteristica completa, avrà il grado \(n_t=r_t\), e perciò conterrà dentro di sè una rete di grado 2, mercè la quale la superficie \(F\) viene rappresentata sul piano doppio con quartica di diramazione: da ciò segue\origfootnote{2}{s127R-2}{Teorema di Clebsch. Cfr. \textsc{Enriques-Conforto}, Op. cit.} che la nostra superficie è razionale.

Se invece \(r_t=1\), il sistema \(\lvert C^{(t-1)}\rvert\) che precede \(\lvert C^{(t)}\rvert\) nella serie degli aggiunti, dovrà essere di genere due, e si potrà supporre
%CONT>>
% scan: s128L
% !PRINTED 232
\weberpage{232}
% head: CAPITOLO SESTO
%<<CONT
% stamp: Dark vertical scan-edge artefact at the right margin; scattered specks.
di dimensione \(r_{t-1}\ge2\), perchè altrimenti si potrebbe risalire ad un sistema precedente, ancora di genere due. Ora si può vedere che il sistema \(\lvert C^{(t-1)}\rvert\) non è semplice: tutte le \(C^{(t-1)}\) passanti per un punto \(P\) della superficie \(F\) passano di conseguenza per un punto \(P'\) che è coniugato a \(P\) nella \(g_2^1\) sopra di essa; infatti il coniugato di \(P\) sopra una \(C^{(t-1)}\) variabile per \(P\) dovrà essere fisso, altrimenti descriverebbe una curva \(C^{(t)}\) o una componente di \(C^{(t)}\) di genere zero, contro le nostre ipotesi.

Ciò posto, essendo la \(F\) una superficie regolare, il sistema \(C^{(t-1)}\) aggiunto a \(C^{(t-2)}\) sarà un sistema regolare, di dimensione \(p_{t-2}-1\), designando \(p_{t-2}\) il genere di \(\lvert C^{(t-2)}\rvert\), ed avrà la serie caratteristica completa e non speciale. Perciò questa serie, che abbiamo visto essere composta con la \(g_2^1\), sarà precisamente il doppio di questa \(g_2^1\). Quindi il sistema \(\lvert C^{(t-1)}\rvert\) sarà un sistema \(\infty^3\) di curve di genere 2, appartenente ad un’involuzione di coppie di punti, mercè cui la nostra superficie \(F\) viene rappresentata sopra una quadrica doppia \(Q\), con curva di diramazione del sesto ordine \(D_6\). Siccome le generatrici di \(Q\) non possono possedere un numero dispari di punti di diramazione, si deduce che la \(Q\) è un cono, il cui vertice costituisce un punto di diramazione da aggiungere alla \(D_6\), secante le generatrici in tre punti. Il cono doppio così definito si lascia rappresentare sul piano doppio di Noether con sestica di diramazione dotata di due punti tripli infinitamente vicini, e da ciò segue che il piano doppio è razionale\origfootnote{1}{s128L-1}{Teorema di Noether. Cfr. \textsc{Enriques-Conforto}, Op. cit.}.

Abbiamo ammesso che la serie dei successivi sistemi aggiunti \(\lvert C'\rvert,\ \lvert C''\rvert,\ \ldots\ldotp\) sia formata di sistemi irriducibili. Se percorrendo questa serie s’incontra un sistema \(C^{(h)}\) dotato di una curva fissa \(\theta\) e di cui la parte variabile sia irriducibile, appartenente ad un sistema di dimensione \(r_h\ge2\), si potrà in generale proseguire la costruzione della nostra serie di sistemi aggiunti ponendo al posto di \(\lvert C^{(h)}\rvert\) il sistema irriducibile \(\lvert C^{(h)}-\theta\rvert\), spogliato della parte fissa \(\theta\). E se questo sistema sia di genere 0 o 1, e quindi (per l’integrità della serie caratteristica) di grado \(n_h=r_h-1\) ovvero \(n_h=r_h\), si concluderà (come innanzi) che la superficie \(F\) è razionale.

Resta da approfondire l’esame del caso in cui il sistema \(\lvert C^h-\theta\rvert\), spogliato della parte fissa \(\theta\), sia un fascio ovvero sia composto colle curve \(L\) d’un fascio che, per la regolarità della superficie, sarà \textit{lineare}.

In primo luogo se le curve \(L\) del fascio sono razionali il noto teorema di Noether ci permette di affermare la razionalità della superficie.

% scan: s128R
% !PRINTED 233
\weberpage{233}
% head: SUPERFICIE REGOLARI : MINIMO DEI GENERI E CONDIZIONI, ECC.
% stamp: Scattered scan specks.
Si assuma invece che le curve \(L\) abbiano il genere \(\varrho>0\). Osserviamo intanto che la \(\theta\), parte fissa di \(\lvert C^h\rvert\) dovrà essere curva fondamentale per il sistema irriducibile \(\lvert C^{h-1}\rvert\), giacchè \(\lvert C^h\rvert\) sega su una \(C^{h-1}\) la serie canonica completa, priva di punti fissi.

Ciò posto, le curve \(L\) incontreranno le \(C^{h-1}\) in \(\nu>1\) punti variabili (poichè dall’ipotesi \(\nu=1\) si dedurrebbe che le \(L\) sono razionali), e la serie canonica completa segata da \(\lvert C^h-\theta\rvert\) su una \(C^{h-1}\) sarà composta coi gruppi dell’involuzione d’ordine \(\nu\) segata dal fascio delle \(L\); di qui si deduce che le \(C^{h-1}\) sono curve iperellittiche essendo \(\nu=2\): il fascio delle \(L\), con cui si compongono le curve di \(\lvert C^h-\theta\rvert\), è un fascio lineare di bisecanti.

Ora ricordiamo che il sistema completo \(\lvert C^{h-1}\rvert\) aggiunto al sistema irriducibile \(\lvert C^{h-2}\rvert\) (sopra la superficie \(F\) regolare di genere \(p=0\)) è un sistema regolare (di dimensione \(\pi_{h-2}-1\)) e perciò ha serie caratteristica completa e non speciale. Questa serie è composta colla \(g_2^1\) perchè altrimenti il coniugato d’un punto \(P\) sopra una \(C^{h-1}\) per \(P\) descriverebbe una curva razionale, formante la \(L\) o una componente della \(L\) per \(P\). Per conseguenza il sistema \(\lvert C^{h-1}\rvert\) non è semplice, ma composto con una involuzione di 2\textsuperscript{o} ordine e il grado \(n_{h-1}\) di codesto sistema di curve iperellittiche sarà il doppio della dimensione diminuita di una unità:
\[
n_{h-1}=2(r_{h-1}-1);
\]
mentre il suo genere \(\pi_{h-1}\) varrà
\[
\pi_{h-1}=r_{h-1}-1.
\]

Possiamo arrivare a una determinazione più precisa di questi caratteri, mercè la considerazione seguente. Se si stacca dal sistema \(\lvert C^{h-1}\rvert\) una curva \(L\) (su cui le \(C^{h-1}\) segnano le coppie di una \(g_2^1\)) si ottiene un sistema lineare di dimensione \(r_{h-1}-2\), che deve avere come sistema aggiunto, il sistema aggiunto di \(\lvert C^{h-1}\rvert\), diminuito della \(L\), e perciò il sistema costituito da tutti i gruppi di \(\pi_{h-1}-2\) curve \(L\); quindi lo staccamento della \(L\) da \(\lvert C^{h-1}\rvert\) ne diminuisce il genere di una sola unità. Ma d’altra parte le \(C^{h-1}\) che vengono a spezzarsi nella \(L\) e in una curva residua, debbono essere connesse, e quindi la \(L\) deve incontrare le curve residue in \(i\ge1\) punti; così, essendo \(\pi_{h-1}-1\) il genere di queste curve, si avrà
\[
\pi_{h-1}=(\pi_{h-1}-1)+\varrho+i-1
\]
da cui (essendo \(\varrho>0\)) segue
\[
\varrho=1,\qquad i=1.
\]

% scan: s129L
% !PRINTED 234
\weberpage{234}
% head: CAPITOLO SESTO
% stamp: Dark vertical scan-edge artefact at the right margin; scattered specks.
Dunque, le curve residue di una \(L\) rispetto a \(\lvert C^{h-1}\rvert\) segano la \(L\) in un punto che è necessariamente fisso, e comune alle \(L\) del nostro fascio, e ciò porta che queste curve residue si spezzino in tante \(L\) del fascio stesso.

Ora, mercè il sistema \(\lvert C^{h-1}\rvert\), trasformiamo la superficie \(F\) in una superficie doppia d’ordine \(m=r_{h-1}-1\) nello spazio ad \(m+1=r_{h-1}\) dimensioni. Questa superficie, su cui le curve \(L\) sono rappresentate da rette, sarà un cono avente un certo vertice \(O\), e dovrà possedere una curva di diramazione \(D\), d’ordine % The plus sign in m+ +1 is repeated at the printed line break.
\[
2+2\pi_{h-1}=2m+2.
\]

Se questa \(D\) deve incontrare le generatrici della superficie in almeno 3 punti variabili (affinchè il genere delle \(L\) sia \(\varrho>0\)), deve essere \(m=2\); infatti un iperpiano condotto per \(O\) sega il nostro cono in \(m\) generatrici, e perciò si deve avere
\[
3m\le2m+2.
\]

In tal guisa siamo pervenuti a concludere che il sistema \(\lvert C^{h-1}\rvert\) conduce a rappresentare la nostra superficie \(F\) sopra un cono quadrico doppio, con sestica di diramazione, che, come già abbiamo ricordato, è razionale (Noether).

\textit{Ipotesi II.} -- Dobbiamo ancora esaminare la seconda ipotesi: ogni sistema lineare irriducibile tracciato sopra la superficie \(F\) ha genere \(\pi\) e grado \(n\), soddisfacenti alla diseguaglianza
\[
n\le2\pi-2.
\]

Assumiamo questa ipotesi che ridurremo all’assurdo. Secondo l’assunto, la superficie \(F\) si può supporre trasformata in guisa da non contenere curve eccezionali, e, se si vuole, priva di singolarità in un conveniente spazio. Le sezioni iperpiane di essa costituiscono un sistema lineare \(\lvert C\rvert\), di cui designiamo con \(n\) il grado, e con \(\pi\) il genere, essendo, come si è detto
\[
n\le2\pi-2.
\]

Allora il bigenere \(P\) della \(F\), si può valutare aumentando di una unità la dimensione del sistema bicanonico, il quale si ottiene staccando una \(C\) dal sistema secondo aggiunto \(\lvert C''\rvert\), e, come abbiamo visto, si trova che per una superficie regolare di genere superficiale \(p=p_a=p_g=0\), e di genere lineare \(p^{(1)}\), il bigenere vale
\[
P\ge p^{(1)}.
\]

Per conseguenza, se \(P=0\), sarà
\[
p^{(1)}\le0.
\]

% scan: s129R
% !PRINTED 235
\weberpage{235}
% head: SUPERFICIE REGOLARI : MINIMO DEI GENERI E CONDIZIONI, ECC.
% stamp: Scattered scan specks.
Allora si consideri la serie dei sistemi aggiunti successivamente a \(\lvert C\rvert,\ \lvert C'\rvert,\ \lvert C''\rvert,\ \ldots\ldotp\); il numero delle intersezioni \((2\pi-2-n)\) che \(\lvert C\rvert\) ha con le curve canoniche virtuali, cresce passando al sistema aggiunto di \(p^{(1)}-1\), cioè in realtà diminuisce almeno di una unità. Pertanto dovrà avverarsi uno almeno dei due fatti che seguono:

1\textsuperscript{o}) o la differenza \(n-(2\pi-2)\) costruita con i caratteri di un sistema \(\lvert C\rvert\) diventa positiva per qualcuno degli aggiunti successivi, ciò che contraddice direttamente la nostra ipotesi;

2\textsuperscript{o}) ovvero la serie dei detti aggiunti è limitata.

In tal caso, così come accadeva nella Ipotesi I), si arriva ad un ultimo sistema lineare, composto con le curve razionali di un fascio lineare, oppure ad un fascio di curve ellittiche che si presenta come aggiunto alle curve di genere 2 iperellittiche di un sistema lineare \(\infty^2\) almeno di grado 4. % [sic] Ipotesi I) -> Ipotesi I, the closing parenthesis has no matching opening parenthesis.
Con lo stesso ragionamento svolto nella Ipotesi I, si deduce di qui la razionalità della superficie \(F\). Ma non importa nemmeno ripetere quel ragionamento; basta avvertire che, in contraddizione alla nostra Ipotesi II, siamo pervenuti a costruire sopra \(F\) un sistema lineare di genere \(\pi\) e grado \(n>2\pi-2\), sicchè l’ipotesi stessa viene ridotta all’assurdo. Qui è opportuno osservare che quando si abbia sopra \(F\) un sistema \(\lvert L\rvert\) comunque riducibile che non contenga componenti eccezionali, di genere \(\pi\) e grado \(n>2\pi-2\), sommando un multiplo di esso ad un sistema regolare, si può sempre ottenere un sistema lineare irriducibile i cui caratteri soddisfino alla medesima diseguaglianza. Concluderemo enunciando il teorema: \textit{le condizioni necessarie e sufficienti perchè una superficie sia razionale sono espresse dall’annullarsi del genere numerico e del bigenere:}
\[
p_a=P=0;
\]
dove la condizione \(P=0\) porta a priori che sia anche \(p_g=0\). Così in altre parole può dirsi che le \textit{superficie razionali} vengono caratterizzate come « \textit{superficie regolari di bigenere zero} ».

Questo teorema appartiene a G. \textsc{Castelnuovo} che lo ha dimostrato sostanzialmente col metodo qui tenuto: dove si sono introdotte soltanto semplificazioni di particolari.
% scan: s130L
% !PRINTED none
\weberblank{[s130L]}
% head: none
% blank page
% stamp: Dark vertical scan-edge artefact at the right margin; scattered specks.


\clearpage\thispagestyle{plain}
% scan: s130R
% !PRINTED none
\weberpage{[s130R]}
% head: none
% stamp: Short horizontal scan-edge artefact near the upper right margin; scattered specks.
\edchapter{VII}{CLASSIFICAZIONE DELLE SUPERFICIE\\DI GENERE LINEARE \(p^{(1)}=1\).}{Capitolo VII. CLASSIFICAZIONE DELLE SUPERFICIE DI GENERE LINEARE p\textasciicircum{}(1)=1.}{CLASSIFICAZIONE DELLE SUPERFICIE DI GENERE LINEARE \(p^{(1)}=1\).}

\edsection{1. Superficie di genere zero e bigenere 1, con curva bicanonica d’ordine zero.}{1. Superficie di genere zero e bigenere 1, con curva bicanonica d’ordine zero.}{sec.VII.1}

Abbiamo veduto (Cap. VI, § 2) che tra le superficie regolari di genere \(p=0\) e di bigenere \(P=1\) quelle dotate di curva bicanonica d’ordine zero, sono caratterizzate, di fronte alle trasformazioni birazionali, dal valore del trigenere, che è per esse
\[
P_3=0,
\]
ovvero del sestigenere, che è
\[
P_6=1
\]
(anzichè \(P_6>1\)).

Ora vogliamo approfondire lo studio di queste superficie \(F\), riconducendole ad un tipo proiettivamente definito, che sarà il piano doppio con curva di diramazione \(D_8\), formata da una sestica con due tacnodi \(M\) ed \(N\) e con un punto doppio nell’intersezione delle loro tangenti tacnodali-\(p\) e \(q\), cui si aggiungono queste due rette \(p\) e \(q\). Riconosceremo poi che questo piano doppio si può trasformare in una superficie del 6\textsuperscript{o} ordine \(F_6\), passante doppiamente per gli spigoli di un tetraedro, che perciò può ugualmente prendersi come tipo della nostra famiglia di superficie (\(p=0,\ P=1,\ P_3=0\)). % ?? uncertain: The short mark after "tacnodali" may be a hyphen or an artefact -- tacnodali-p; tacnodali p.

Ricordiamo anzitutto le proprietà fondamentali delle nostre superficie \(F\), riferendoci ad un modello di esse privo di curve eccezionali:

1\textsuperscript{o}) ogni curva \(C\) di genere \(\pi\) tracciata su \(F\) ha il grado
\[
n=2\pi-2;
\]

2\textsuperscript{o}) la dimensione del sistema lineare completo \(\lvert C\rvert\) (virtualmente privo di punti base), è
\[
r\ge\pi-1;
\]
%CONT>>
% scan: s131L
% !PRINTED 238
\weberpage{238}
% head: CAPITOLO SETTIMO
% stamp: Dark vertical scan band at the right edge; scattered specks.
\(3^\circ\)) per ogni sistema irriducibile \(|C|\), di genere \(\pi > 1\), si ha esattamente
\[
r = \pi - 1
\]
cioè il sistema è regolare;

\(4^\circ\)) il sistema lineare \(|C|\) ha come aggiunto un sistema \(|C'|\), distinto da esso, coi medesimi caratteri; la relazione fra i due sistemi è reciproca, essendo anche \(|C|\) l’aggiunto di \(|C'|\): i due sistemi sono due diverse metà del medesimo sistema doppio
\[
|L| = |2C| = |2C'|;
\]

\(5^\circ\)) alla superficie \(F\) possono appartenere \textit{curve ellittiche}, cioè curve irriducibili di genere virtuale \(\pi = 1\), e quindi di grado \(n = 0\). Queste curve possono essere \textit{isolate}, cioè costituenti un sistema regolare, di dimensione \(r = 0\), ovvero appartenenti ad un fascio (sistema sovrabbondante, di dimensione \(r = 1\)). Se \(|C|\) è un fascio di curve ellittiche sopra \(F\), il sistema aggiunto \(|C'|\) è costituito da due curve ellittiche isolate \(\dfrac{C}{2}\) e \(\left(\dfrac{C}{2}\right)'\), aggiunte l’una dell’altra, i cui doppi appartengono al fascio. Infatti la curva aggiunta \(C'\), non avendo intersezioni variabili con le \(C\) del fascio, ed essendo distinta da una \(C\), deve constare di parti di curve \(C\), e d’altra parte \(|2C'|\) deve equivalere a una coppia di curve \(C\).

Reciprocamente se nel fascio \(|C|\) si ha una curva metà \(\dfrac{C}{2}\) anche la sua aggiunta \(\left(\dfrac{C}{2}\right)'\) sarà una curva metà del fascio, e le curve aggiunte a \(\dfrac{C}{2} + \left(\dfrac{C}{2}\right)'\) saranno le \(2\dfrac{C}{2} = C\). E quindi per la proprietà \(4^\circ\)) a un fascio di curve ellittiche \(|C|\) non può appartenere che una sola coppia di curve ellittiche doppie, formate con due metà del fascio, che sono aggiunte l’una dell’altra;

\(6^\circ\)) una curva \(\theta\), che non sia componente delle curve ellittiche di un fascio \(|C|\), sega necessariamente le \(C\) in un numero pari di punti. Infatti questo numero è doppio del numero delle intersezioni di \(\theta\) con le curve \(\dfrac{C}{2}\);

\(7^\circ\)) una curva \(\theta\) irriducibile di genere (virtuale) zero e quindi di grado \(-2\), sopra la superficie \(F\), non può possedere alcuna curva aggiunta.

Infatti, se esistesse un’aggiunta \(\theta'\), questa dovrebbe incontrare la \(\theta\) in \(-2\) punti e quindi coinciderebbe con \(\theta\); ma da \(\theta = \theta'\) seguirebbe che il genere \(p = 1\), contro l’ipotesi.

In altre parole l’esistenza di \(\theta'\) porta che la \(\theta\) sia di genere \(> 0\);

% scan: s131R
% !PRINTED 239
\weberpage{239}
% head: CLASSIFICAZIONE DELLE SUPERFICIE DI GENERE LINEARE ECC.
% stamp: Scattered scan specks and faint marginal marks.
\(8^\circ\)) sopra la superficie \(F\), dato un fascio di curve ellittiche \(|C|\), non può esistere una curva irriducibile \(\theta\) di genere zero che le bisechi.

Infatti il sistema \(C + \theta\), che sarebbe un fascio di curve di genere virtuale \(\varrho = 2\), dovrebbe ammettere un sistema aggiunto \(|L|\) di dimensione 1, contenente la curva composta \(\theta + \dfrac{C}{2} + \left(\dfrac{C}{2}\right)'\); invece si può riconoscere che questa curva non può variare in un fascio. A tal uopo si osservi che la \(\dfrac{C}{2}\) deve incontrare la \(L\) in un solo punto, essendo
\[
\frac{C}{2}L = \frac{C}{2}\theta + \frac{C}{2}\frac{C}{2} + \frac{C}{2}\left(\frac{C}{2}\right)'
\]
e
\[
\frac{C}{2}\frac{C}{2} = 0;
\]
quindi le \(L\) segano \(\dfrac{C}{2}\), e analogamente \(\left(\dfrac{C}{2}\right)'\), in un punto fisso e perciò — se ve ne sono \(\infty^1\) che non si riducono alla \(\theta\) stessa aumentata della \(C' \doteq \dfrac{C}{2} + \left(\dfrac{C}{2}\right)'\) — risultano prive d’intersezioni colle \(C\) e quindi composte colle dette \(C\) e con una curva \(\theta'\) aggiunta a \(\theta\): ma l’esistenza di \(\theta'\) porterebbe che il genere di \(\theta\) sia \(> 0\), contro l’ipotesi. % ?? uncertain: relation after C' -- \doteq or =; a small mark is visible above the upper stroke.
% EDITORIAL D0126 RISOLTO: The enlarged scan has a small, dark, nearly circular dot centered above the upper equality stroke. This supports the supplied dotted equality rather than an ordinary equality.

Teniamo presenti le proposizioni 1) .... 8).

Essendo \(|C|\) un sistema lineare irriducibile, di dimensione \(\pi - 1\), si consideri il sistema doppio \(|L|\) che ha il grado \(4n = 8\pi - 8\), il genere \(4\pi - 3\), e la dimensione \(4\pi - 4\). Entro il sistema \(|L|\) si hanno due sistemi \(\infty^{2\pi-2}\), diciamo \(\Sigma\) e \(\Sigma'\), costituiti: il primo dalle coppie di curve \(C\), e il secondo dalle coppie di curve \(C'\), ciascuno dei quali ha l’indice
\[
i = \frac{1}{2}\binom{2\pi - 2}{\pi - 1}.
\]

Indicando con \(S_{4\pi-4}\) lo spazio i cui punti rappresentano gli elementi (curve) del sistema \(|L|\), si avranno in esso due varietà, rappresentative di \(\Sigma\) e \(\Sigma'\), di dimensione \(2\pi - 2\), e d’ordine \(i\), le quali avranno a comune \(i^2\), ovvero una infinità di punti. Così si trovano delle curve \(L\) che sono contemporaneamente decomponibili in due curve \(C\) e in due curve \(C'\); queste \(L\) sono necessariamente spezzate, e permettono in generale di costruire entro un sistema \(|C|\) di genere \(\pi > 1\), un sistema parzialmente contenuto in esso, di genere minore di \(\pi\).

Precisiamo che si troverà una curva composta di due \(C\):
\[
C = C_1 + C_2,\qquad C = C_3 + C_4
\]
%CONT>>
% scan: s132L
% !PRINTED 240
\weberpage{240}
% head: CAPITOLO SETTIMO
%<<CONT
% stamp: Dark vertical scan band at the right edge; stray marginal marks.
decomposta in modo che una parte della prima, insieme ad una parte della seconda, dovrà dare una curva \(C'\), sia per es.
\[
C_1 + C_3 = C',\qquad C_2 + C_4 = C',
\]
sicchè, essendo
\[
C_1 + C_3 = (C_1 + C_2)'
\]
si dedurrà
\[
C_3 = C'_2
\]
e analogamente
\[
C_4 = C'_1.
\]

Ciò importa che le curve \(C_1\), \(C_2\) e similmente le \(C_3\) e \(C_4\) siano di genere maggiore o uguale a uno.

Distingueremo due casi:

\textit{a}) fra le due curve \(C_1\) e \(C_2\) ve n’è una, per es. \(C_1\), di genere \(\pi_1 > 1\).

\textit{b}) le due curve \(C_1\) e \(C_2\) sono di genere 1.

Nel caso \textit{a}) si può affermare che la \(C_1\) appartiene a un sistema lineare completo \(|C_1|\) di dimensione \(\pi_1 - 1 \geq 1\). Siccome questo sistema è contenuto parzialmente in \(|C|\), la sua dimensione \(\pi_1 - 1 < \pi - 1\), e quindi il genere \(\pi_1 < \pi\). % The less-than sign is repeated on the two sides of a printed line break.

Inoltre, le parti variabili delle curve \(C_1\) dovranno essere irriducibili, e ancora di genere \(\varrho > 1\) (\(\varrho < \pi\)). Invero, pongasi che codeste parti siano composte con \(s\) curve \(K\) di un fascio lineare (\(s > 1\)). Se le \(K\) fossero di genere maggiore di 1, il sistema completo multiplo di \(|K|\) secondo \(s\), risulterebbe di dimensione maggiore di \(s\), e quindi irriducibile.

Sia invece il genere delle \(K\) uguale ad 1. In tale ipotesi i gruppi di \(sK\) sono soltanto parti delle curve \(C_1\):
\[
|C_1| = \theta + |sK|.
\]

E affinchè il genere di \(C_1\) sia \(\pi_1 > 1\) bisogna che la \(\theta\), o una componente \(\theta_1\) di essa, incontri in qualche punto le \(K\) (una curva non intersecante le \(K\) resultando parte d’una \(K\) e quindi di genere minore o uguale 1). Il numero di queste intersezioni essendo pari, e perciò maggiore o ugual due, l’addizione di \(\theta_1\) ad \(|sK|\) riesce certo ad ampliare il sistema \(|sK|\) se \(\theta_1\) è di genere 1, ed anche se \(\theta_1\) è di genere 0, poichè in tale ipotesi il numero delle intersezioni con le \(K\) deve essere maggiore di 2 (cfr. prop. 8) e quindi maggiore o uguale a 4.

In conclusione, il sistema delle \(C_1\), ovvero un sistema completo in esso contenuto, costituisce un sistema lineare irriducibile, di genere maggiore d’uno.

Pertanto, l’applicazione reiterata del nostro processo di riduzione. ove non s’incontri il caso \textit{b}), conduce a trovare un sistema lineare parzialmente contenuto in \(|C|\) e di genere due. % [sic] "riduzione." -> "riduzione,", the following lower-case subordinate clause continues the sentence.

% scan: s132R
% !PRINTED 241
\weberpage{241}
% head: CLASSIFICAZIONE DELLE SUPERFICIE DI GENERE LINEARE ECC.
% stamp: Scattered scan specks and a faint dark line at the right edge.
Mettiamoci invece nel caso \textit{b}), supponendo dunque che fra le curve di un sistema \(|C|\) di genere \(\pi\) si trovi qualche curva spezzata in due curve di genere 1,
\[
C = C_1 + C_2.
\]

Siccome \(C_1\) e \(C_2\) sono curve di genere 1 e di grado 0, i loro doppi \(2C_1\) e \(2C_2\) daranno luogo a due fasci di curve ellittiche, le cui parti variabili designeremo con \(K_1\) e \(K_2\), e, come innanzi, è facile riconoscere che nessuno di questi due fasci può possedere parti fisse che non siano componenti delle curve del fascio stesso. Quindi le curve dei due fasci \(|K_1|\) e \(|K_2|\) dovranno intersecarsi in un numero di punti maggiore di zero, affinchè il grado di \(|C_1 + C_2|\) risulti maggiore di zero. Ora due curve metà di \(K_1\) e \(K_2\), diciamo \(\dfrac{K_1}{2}\) e \(\dfrac{K_2}{2}\) costituiscono due curve ellittiche isolate, aventi un certo numero \(s\) d’intersezioni, tale che
\[
1 \leq s \leq \pi - 1.
\]

In conclusione si può ritenere che, ove il nostro procedimento di riduzione non conduca a costruire un fascio di genere 2, si arriverà sempre a determinare due curve ellittiche isolate, che torniamo a chiamare \(C_1\) e \(C_2\), con un certo numero \(s \leq \pi - 1\) di punti a comune (\(s > 0\)). A rigore non si può affermare che queste \(C_1\) e \(C_2\) sieno senza parti comuni; potrebbe accadere per esempio che la \(C_1\) sia spezzata in due curve di genere zero con due punti comuni, \(C_1 = \theta + \theta_1\) e similmente sia \(C_2 = \theta + \theta_2\); allora \(C_1\) e \(C_2\) avrebbero a comune la componente \(\theta\), ma nelle deduzioni fondate sull’uso dei caratteri virtuali, è « come se » le due curve \(\theta + \theta_1\) e \(\theta + \theta_2\) sieno irriducibili con
\[
s = (\theta + \theta_1,\ \theta + \theta_2)
\]
punti comuni.

Il procedimento di riduzione che abbiamo messo in opera, a partire da un sistema lineare irriducibile, sopra la superficie \(F\), conduce a costruire su \(F\): o un sistema lineare \(\infty^1\) \(|C|\) di genere 2, ovvero una coppia di curve ellittiche \(C_1\) e \(C_2\) con \(s(\geq 1)\) punti comuni.

Nel primo caso le curve \(C\) debbono ritenersi irriducibili, perchè sono a priori irriducibili e di genere 2 le loro parti variabili costituenti un fascio: se invero queste fossero curve \(K\) di genere 1, \(|C|\) possederebbe una parte fissa \(\theta\), secante le \(K\) soltanto in 2 punti comuni, che è impossibile.

Pertanto si troveranno entro \(|C|\) delle curve spezzate in due componenti ellittiche isolate \(C_1\) e \(C_2\), con \textit{un} punto comune.

Nel secondo caso, quando si abbiano su \(F\) due curve ellittiche isolate \(C_1\) e \(C_2\) con \(s > 1\) punti comuni, si può trovare su \(F\) una coppia di curve ellittiche secantisi in \(s' < s\) (\(s' \geq 1\)) punti, e quindi anche
% Printed gathering signature:
% EDITORIAL: segnatura originale: \footnotesize \textsc{Enriques F.} - \textit{Superficie algebriche.}\hfill 16
\printersignature{\textsc{Enriques F.} - \textit{Superficie algebriche.}\quad 16}
%CONT>>
% scan: s134L
% !PRINTED 242
\weberpage{242}
% head: CAPITOLO SETTIMO
%<<CONT
% stamp: Dark right-edge band, marginal specks, and an oblique non-text stroke after "un’altra".
una coppia di curve ellittiche con \textit{un}-punto comune. La possibilità di questa nuova riduzione si dimostra come segue. % ?? uncertain: short mark between "un" and "punto" -- printed hyphen or ink/scanning artefact.

Se \(s > 1\) il sistema lineare \(|C| = |C_1 + C_2|\) è di genere \(\pi = s + 1 > 2\). Quindi fra le coppie di curve \(2C\), si avranno in generale \(\dfrac{1}{4}\dbinom{2\pi - 2}{\pi - 1}^{2} \geq 9\) curve spezzate. Se queste contengono delle componenti di genere maggiore o uguale due, si arriva a costruire col procedimento già esposto una coppia di curve ellittiche con un numero d’intersezioni \(s' < s\). Se invece tutte le dette curve spezzate sono formate da componenti ellittiche occorre svolgere le considerazioni seguenti. % The equality sign is repeated on the two sides of a printed line break.

Precisamente si deve ammettere che oltre alla coppia di curve spezzate che costituisce
\[
C_1 + C_2 + C'_1 + C'_2,
\]
si trovi un’altra coppia di curve spezzate in componenti ellittiche
\[
C_3 + C_4 + C'_3 + C'_4,
\]
dove \(C_3\) e \(C_4\) sieno distinte da \(C_1\), \(C_2\), \(C'_1\) e \(C'_2\): l’esistenza di \(C_3 + C_4\) si deduce dall’ammettere che le due varietà \(\Sigma\) e \(\Sigma'\) rappresentative delle coppie di curve \(C\) e \(C'\) abbiano qualche altro punto a comune \textit{distinto} dai due punti che rispondono alle coppie
\[
C_1 + C_2 \qquad C'_1 + C'_2
\]
e
\[
C_1 + C'_2 \qquad C_2 + C'_1;
\]
ed invero vedremo fra poco che le dette coppie di curve ellittiche isolate \(C_1 + C_2\) e \(C'_1 + C'_2\), non possono rispondere a intersezioni di \(\Sigma\) e \(\Sigma'\) che cadano infinitamente vicine.

Ordunque, essendo \(C_3 + C_4\) una curva spezzata del sistema \(|C|\) diversa da \(C_1 + C_2\) e da \(C_1 + C'_2\), bisognerà che \(C_3\) e \(C_4\) seghino in qualche punto ciascuna delle due curve \(C_1\) e \(C_2\), chè altrimenti conterrebbero come parte la \(C_1\) stessa, ovvero una delle due curve \(C_2\) o \(C'_2\); segue che, per esempio, la \(C_3\) dovrà incontrare ciascuna delle \(C_1\) e \(C_2\) in un numero di punti minore di \(s\); sicchè abbiamo ottenuto, come si voleva, una coppia di curve ellittiche isolate \(C_1\) e \(C_3\) con \(s' < s\) punti a comune.

Resta da giustificare che il nostro sistema \(|C|\) non può contenere due coppie di curve ellittiche isolate infinitamente vicine \(C_1 + C_2\) e \(C_3 + C_4\), dove sia per es. \(C_3\) infinitamente vicina a \(C_1\) e \(C_4\) a \(C_2\). Infatti le curve \(2C_1\) e \(2C_3\) dovrebbero definire un fascio di curve ellittiche, di cui \(C_1\) risulterebbe componente fissa, sicchè la \(C_1\) stessa
%CONT>>
% scan: s134R
% !PRINTED 243
\weberpage{243}
% head: CLASSIFICAZIONE DELLE SUPERFICIE DI GENERE LINEARE ECC.
%<<CONT
% stamp: Dark vertical scan band at the left edge and scattered specks.
verrebbe a variare in un fascio di curve ellittiche, mentre si è supposta isolata.

Da ciò che precede si ricava la conclusione che la nostra superficie \(F\) possiede una coppia di curve ellittiche isolate \(C_1\) e \(C_2\) aventi un punto a comune, per modo che le loro aggiunte, \(C'_1\) e \(C'_2\), sono incidenti alla \(C_2\) e alla \(C_1\). La curva composta \(C_1 + C_2 + C'_2\) riesce pertanto di genere 3, e definisce una rete di curve \(|K|\) di questo genere: le \(K\) segano \(C_2\) e \(C'_2\) ciascuna in un sol punto, che deve essere fisso; quindi le \(K\) hanno due punti base semplici, che cadono precisamente nei punti comuni a \(C'_1\) e \(C_2\), e \(C_1\) e \(C'_2\), giacchè la \(C'_1\) fa parte di \(\infty^1\) curve \(K\) avendo come residuo rispetto a \(|K|\) il fascio delle curve ellittiche \(|2C_2|\): l’asserto risulta dall’osservare che le curve \(C'_1 + 2C_2\) sono seconde aggiunte alle \(K = C_1 + C_2 + C'_2\), e quindi sono equivalenti, appartenendo al medesimo sistema \(|K|\).

Pertanto le \(K\) risulteranno curve iperellittiche, segandosi a due a due in due punti variabili; queste \(K\) segano poi sulla \(C_1\) le coppie di una \(g^1_2\), essendovi una sola curva residua di \(C_1\) rispetto a \(|K|\), che è la \(C_2 + C'_2\). Si aggiunga che i punti base delle \(K\) sono, sopra una \(K\) qualunque, due punti doppi della relativa \(g^1_2\).

Per mezzo della rete di curve iperellittiche \(|K|\) di genere 3, si può ora rappresentare la superficie \(F\) sopra un piano doppio, che avrà una curva di diramazione \(D_8\) di ordine 8.

Di questa \(D_8\) faranno parte due rette di diramazione, \(p\) e \(q\), corrispondenti ai punti base di \(|K|\); e il punto \(O = pq\) sarà l’immagine della curva ellittica \(C'_1\), e quindi quadruplo per la \(D_8\), ovvero doppio per la sua componente \(D_6\) che rimane togliendo le rette \(p\) e \(q\). Le rette del fascio di centro \(O\) saranno, nel piano doppio, le immagini delle curve ellittiche del fascio \(|2C_2| = |2C'_2|\); alle \(C_2\) e \(C'_2\) corrisponderanno due tacnodi della \(D_6\) con le tangenti tacnodali \(p\) e \(q\), ossia due punti \([3,3]\) della curva di diramazione \(D_8\). Queste singolarità della \(D_8\) valgono effettivamente a caratterizzarla come curva di diramazione di un piano doppio, di genere \(p_a = p_g = 0\) e \(P = 1\), con curva bicanonica (pura) d’ordine zero, siccome si riconosce in base al Cap. V; \S\ 7.

Concludiamo pertanto che \textit{le superficie regolari di genere zero e bigenere \(1\), con curva bicanonica d’ordine zero — quali sono caratterizzate dal valore \(P_3 = 0\) del trigenere, ovvero \(P_6 = 1\) del sestigenere — formano una sola famiglia, che ha per tipo il piano doppio con curva di diramazione}
\[
D_8 = D_6 + p + q,
\]
\textit{costituita da una sestica con due tacnodi e con un punto doppio nel punto d’incontro delle loro tangenti tacnodali, cui si sommino queste stesse tangenti.}

% scan: s135L
% !PRINTED 244
\weberpage{244}
% head: CAPITOLO SETTIMO
% stamp: Dark right-edge scan band, scattered marginal marks, and a small raised blot after "intorni".
Questa \textit{famiglia} di superficie dipende da 10 \textit{moduli}. Infatti la costruzione della \(D_6\) con 4 punti doppi assegnati e quindi della \(D_8\) dipende da
\[
\frac{6 \cdot 9}{2} - 5 \cdot 3 = 12
\]
costanti, sicchè, aggiungendo le 6 costanti che dipendono dalla scelta dei punti singolari della \(D_6\), si hanno \(\infty^{18}\) curve \(D_8\). Dai 18 parametri che ad esse appartengono, bisogna toglierne 8, in corrispondenza alle \(\infty^8\) \(D_8\) omografiche, che definiscono, naturalmente, la stessa superficie \(F\).

Si consideri la superficie sestica \(F_6\), dello spazio ordinario, passante doppiamente per gli spigoli d’un tetraedro \(T\). Essendo la \(F_6\) una superficie regolare con \(p = 0\), \(P = 1\), \(P_3 = 0\), essa deve potersi rappresentare su un piano doppio con curva di diramazione \(D_8\), dotata delle singolarità dette innanzi. Si ottiene effettivamente una rappresentazione di \(F_6\) su un piano doppio mercè la proiezione sghemba da una coppia di spigoli opposti \(a\) e \(b\) del tetraedro \(T\), giacchè le rette incidenti ad \(a\) e \(b\) incontrano la \(F_6\) in due punti semplici. Ma questa costruzione conduce in generale ad un piano doppio con curva di diramazione \(D_{10}\), che soltanto con una trasformazione ulteriore del piano può ricondursi alla nostra \(D_8\). Giova perciò particolarizzare la detta costruzione, assumendo un piano rappresentativo che passi per uno spigolo \(c\) del tetraedro, incidente ad \(a\) e \(b\). Allora le rette del piano doppio appaiono immagini delle curve \(K\) segate su \(F_6\) dalle quadriche per \(a\), \(b\), \(c\). Queste \(K\) sono curve iperellittiche di genere 3, fra le quali si trova una curva spezzata negli intorni di 3 rette doppie del tetraedro \(T\): \(c\), \(d\), ed \(e\), essendo \(d\) ed \(e\) gli spigoli di \(T\) ulteriori sezioni di \(F_6\) con i piani \(ca\) e \(cb\). Si noti che le componenti di questa curva \(K\) spezzata, cioè gli intorni delle 3 rette doppie \(c\), \(d\), ed \(e\), sono curve ellittiche connesse come le \(C_1\), \(C_2\), \(C'_2\), di cui ci siamo valsi innanzi per costruire la rete delle \(K\) che ci ha condotto al piano doppio con \(D_8\) di diramazione. % ?? uncertain: raised mark after "intorni" before "delle" -- scan/ink speck (reading adopted) or printed asterisk.

Si riesce così a verificare nel modo più semplice le conclusioni che si traggono a priori per la \(F_6\), dalla proprietà di essere superficie regolare di genere \(p = 0\) e bigenere \(P = 1\), con curva bicanonica d’ordine zero.

Ora si osservi che vi sono nello spazio \(\infty^{13}\) sestiche \(F_6\) passanti doppiamente per gli spigoli di un dato tetraedro \(T\) (delle quali abbiamo scritto l’equazione nel \S\ 11 del Cap. III) e che vi sono \(\infty^3\) omografie coi punti uniti nei vertici di \(T\); segue da ciò che le \(F_6\) dipendono da 10 parametri o invarianti proiettivi, proprio come i piani doppi con \(D_8\) di diramazione, su cui abbiamo appreso a rappresentarle. Che questi invarianti proiettivi diano altrettanti moduli delle superficie \(F_6\), rispetto alle trasformazioni birazionali si deduce
%CONT>>
% scan: s135R
% !PRINTED 245
\weberpage{245}
% head: CLASSIFICAZIONE DELLE SUPERFICIE DI GENERE LINEARE ECC.
%<<CONT
% stamp: Dark vertical scan band at the left edge and scattered specks.
da ciò che il sistema delle sezioni piane è completo, ed essendo la superficie regolare, non può essere contenuto in una serie continua di sistemi analoghi, benchè invero possa aversi sopra \(F\) una serie discontinua di sistemi siffatti\origfootnote{1}{s135R-1}{L’esistenza effettiva di una tale serie risulta dalla proprietà di \(F_6\) di possedere un’infinità discontinua di trasformazioni birazionali in sè stessa. Queste si lasciano definire in rapporto alle \(\infty^1\) quartiche ellittiche \(K\) segate dalle quadriche passanti per 4 spigoli \(a\,b\,d\,e\) di \(T\) che formino un quadrilatero sghembo. Su ciascuna di queste curve si ottiene una trasformazione birazionale, sommando ad un punto \(P\) il gruppo virtuale d’ordine zero, differenza delle coppie \(G_d - G_e\) sezioni di due spigoli incidenti \(d\) ed \(e\): dove giova notare che queste coppie \(G_d\) e \(G_e\) non possono essere equivalenti, nè dar luogo a multipli equivalenti, giacchè ne risulterebbe che le stesse \(d\) ed \(e\), o i loro multipli riuscirebbero equivalenti, ciò che contradice al fatto che \(d\) ed \(e\) hanno un punto a comune. Pertanto la trasformazione definita sopra ogni \(K\) del nostro fascio, portante da un punto \(P\) al punto \(P' = P + G_d - G_e\) non potrà essere ciclica, e così darà luogo ad una serie infinita di potenze, che danno ancora trasformazioni birazionali della superficie in sè stessa.}.

Segue da ciò che la \(F_6\), al pari ed in luogo del corrispondente piano doppio con \(D_8\) di diramazione, si può assumere come tipo generale delle superficie regolari di genere \(p = 0\) e di bigenere \(P = 1\), con curva bicanonica d’ordine zero (\(P_3 = 0\)).

Il problema di costruire la trasformata \(F_6\) a partire dal piano doppio con \(D_8\) di diramazione, si può risolvere determinando in questo piano una cubica \(K_3\) passante per i 5 punti doppi di \(D_6\), che tocchi \(D_6\) in 4 punti e che sia immagine di una curva ellittica spezzata sopra la superficie \(F\) rappresentata dal nostro piano doppio. Che fra le cubiche \(K_3\) aggiunte a \(D_6\) e quadritangenti ad essa esistano effettivamente delle immagini di curve spezzate, si può dimostrare adoperando il criterio di Comessatti, già più volte richiamato. Si tratta di provare che c’è una \(K_3\) quadritangente a \(D_6\), i cui punti di contatto stanno sopra un’altra cubica parimenti aggiunta a \(D_6\). E per ciò si partirà da una cubica qualunque \(K_3\) e si sceglieranno su di essa due punti \(P\) e \(Q\) aventi lo stesso tangenziale \(O\); poi si sceglierà un gruppo di 4 punti \(A_1A_2A_3A_4\), sezione di \(K_3\) con un’altra cubica \(K'_3\) passante per \(O\) e tangente in \(P\) e \(Q\) alla nostra \(K_3\); si verifica quindi che vi sono \(\infty^5\) curve del \(6^\circ\) ordine \(D_6\) passanti doppiamente per \(O\) ed aventi in \(P\) e \(Q\) due tacnodi con le rispettive tangenti tacnodali \(p = PO\) e \(q = QO\) e tangenti a \(K_3\) in \(A_1\), \(A_2\), \(A_3\), \(A_4\): al variare della \(K_3\) e dei punti che abbiamo scelto sopra di essa per la costruzione precedente, la \(D_6\) viene a variare nel sistema completo delle sestiche dotate di due tacnodi e di un punto doppio nel punto d’incontro delle loro tangenti tacnodali \(p = PO\) e \(q = QO\) e tangenti a \(K_3\) in \(A_1\), \(A_2\), \(A_3\), \(A_4\). Si deduce che, data a priori una tale \(D_6\), esistono delle cubiche \(K_3\) aggiunte e quadritangenti ad essa,
%CONT>>
% scan: s136L
% !PRINTED 246
\weberpage{246}
% head: CAPITOLO SETTIMO
%<<CONT
% stamp: Dark right scan edge, scattered specks, and short non-text strokes.
che sono immagini di due curve ellittiche distinte sopra la superficie $F$.

Ciò posto, si considerino su $F$ le curve di genere 2 che hanno per immagine le rette per uno dei due tacnodi di $D_6$, sia per es. per $P$, che possiamo supporre distinto da $O$; sommando a codeste curve di genere due una curva ellittica bisecante corrispondente alla nostra $K_3$ si definirà un sistema lineare $\infty^3$ di genere 4 e grado 6 che conduce a rappresentare la $F$ sopra la sestica $F_6$ che passa doppiamente per gli spigoli d’un tetraedro. Le immagini delle sezioni piane di questa $F_6$ saranno sul piano doppio $\infty^3$ curve del $6^{\mathrm{o}}$ ordine $K_6$, aventi a comune con $D_6$ i punti doppi $O$, $P$, $Q$ e $P'$, $Q'$ infinitamente vicini a questi ultimi; aventi inoltre un punto doppio variabile e 8-tangenti alla $D_6$. % ?? uncertain: Exponent of infinity after "sul piano doppio" -- 2 or 3.
% EDITORIAL D0129 RISOLTO: The enlarged glyph has the two curved lobes and open left side of 3, and matches the nearby printed infinity-cubed expressions.

Il significato della costruzione operata sul piano doppio si può anche spiegare dicendo che: sopra la superficie $F$, che abbiamo riconosciuto contenere un quadrilatero semplice di curve ellittiche isolate, $C_1+C_2+C'_1+C'_2$, esistono altre due curve ellittiche isolate, $C_3$ e $C'_3$ (corrispondenti alla sopra nominata $K_3$), incidenti ai lati del quadrilatero. Tre curve ellittiche a due a due incidenti come $C_1$, $C_2$ e $C_3$, costituiscono una curva di genere 4 appartenente ad un sistema lineare $\infty^3$, che conduce a trasformare $F$ in una $F_6$ dove le dette curve hanno per immagini tre rette doppie, lati d’un trilatero e le loro aggiunte, $C'_1 C'_2 C'_3$, dan luogo del pari a tre rette doppie, spigoli d’un tetraedro di cui fa parte il trilatero, passanti per il vertice opposto alla faccia di esso.

Riassumendo il risultato ottenuto enunceremo il teorema\origfootnote{1}{s136L-1}{F. \textsc{Enriques}, \textit{Sopra le superficie algebriche di bigenere uno.} Memorie Soc. It. delle Scienze, detta dei XL, 1906.}:

\textit{Le superficie regolari di genere $p=0$ e $P=1$, con curva bicanonica d’ordine zero — quali sono caratterizzate dal valore $P_3=0$ del trigenere (o $P_6=1$ del sestigenere) — si lasciano trasformare in sestiche $F_6$ passanti doppiamente per gli spigoli d’un tetraedro.}

\textit{Osservazione.} – Giova osservare alcuni casi particolari, notevoli della sestica $F_6$. % ?? uncertain: Mark after "particolari" -- a comma or a non-text blemish.

Anzitutto può accadere che i 4 piani del tetraedro, i cui spigoli sono rette doppie di $F_6$, vengano a passare per un punto. Il tetraedro si riduce, quindi, ad un angoloide, il cui vertice diventa quadruplo per $F_6$. Questo caso è stato osservato dal \textsc{Castelnuovo}. Per proiezione del punto quadruplo la $F_6$ viene rappresentata sopra un piano doppio con curva di diramazione $D_{10}$, la quale, con una trasformazione quadratica si lascia ricondurre alla $D_8$ di cui innanzi si è definito il tipo.

% scan: s136R
% !PRINTED 247
\weberpage{247}
% head: CLASSIFICAZIONE DELLE SUPERFICIE DI GENERE LINEARE ECC.
% stamp: Dark left scan edge and scattered non-text specks.
Un’altra degenerazione possibile del tetraedro si ha immaginando che due spigoli opposti diventino rette sghembe infinitamente vicine, ed insieme le coppie di facce per codeste rette vengano pure a coincidere. Si ottiene così una $F_6$ che possiede una retta tripla e una retta doppia, infinitamente vicina sghemba con essa, nonchè altre due rette doppie fra loro sghembe, incidenti alla retta tripla. Questa $F_6$ corrisponde ad un piano doppio, con una curva di diramazione $D_8=p+q+D_6$, dove uno dei tacnodi della $D_6$ è infinitamente vicino al punto doppio $O$.

\edsection{2. Le superficie con tutti i generi uguali ad 1.}{2. Le superficie con tutti i generi uguali ad 1.}{sec.VII.2}

Abbiamo veduto che le superficie regolari di genere $p=1$, con curva canonica d’ordine zero, sono caratterizzate di fronte ad altre aventi egualmente il genere lineare $p^{(1)}=1$, dall’essere il bigenere $P=1$, anzichè $P>1$. Anche tutti i plurigeneri ulteriori della superficie risultano $P_i=1$.

Assunto come modello una superficie $F$ priva di curve eccezionali, ogni sistema lineare $|C|$ di genere $\pi$, senza punti base sopra di essa, ha il grado $n=2\pi-2$ ed è l’aggiunto di sè stesso, sicchè per un $|C|$ irriducibile la serie caratteristica è la serie canonica.

\textit{Sopra la superficie $F$ $(p=P=1)$ una curva di genere virtuale $\pi\geq0$, irriducibile o riducibile purchè connessa, appartiene ad un sistema lineare completo di dimensione $r=\pi$} (cfr. Cap. IV, § 17).

Invece una curva sconnessa appartiene, in generale, ad un sistema lineare di dimensione superiore al suo genere $\pi$; in particolare: \textit{i gruppi di $s$ $(>1)$ curve ellittiche di un fascio, che si trovi sopra $F$, appartengono ad un sistema lineare (riducibile) di genere 1 e di dimensione $s$.} Sono anche sovrabbondanti i sistemi lineari contenenti come parte fissa $\theta$ una curva di genere zero.

Notiamo ancora che: se ad un sistema lineare $|C|$ senza parti fisse, che si trovi sopra $F$, si sommi una curva fissa $\theta$, intersecante in più di un punto le componenti variabili delle $C$, il sistema lineare in generale si amplia. Il solo caso in cui questo ampliamento non ha luogo è quello in cui si sommi ai gruppi di $s$ curve ellittiche d’un fascio una curva di genere zero unisecante tali curve.

Per classificare le nostre superficie di genere 1 $(p=P=1)$, prenderemo in considerazione sopra una data $F$ i sistemi lineari irriducibili di genere minimo $\pi>1$, che ad essa appartengono. Potremo supporre anzitutto che il genere del detto sistema minimo $|C|$ sia $\pi=2$, e poi $\pi=3$, $\pi=4$, \ldots.

L’ipotesi $\pi=2$ porta che il grado del sistema sia $n=2$ e la sua dimensione $r=2$, sicchè la $F$ verrà rappresentata sopra un piano
%CONT>>
% scan: s137L
% !PRINTED 248
\weberpage{248}
% head: CAPITOLO SETTIMO
%<<CONT
% stamp: Dark right scan edge, scattered specks, and diagonal non-text strokes.
doppio con sestica di diramazione $D_6$. E di fatto il piano doppio con $D_6$ di diramazione ha il genere % ?? uncertain: Capital letter after "D_6." -- E or È; a nearby mark obscures its upper part.
% EDITORIAL D0131 RISOLTO: The enlarged capital has no separate accent above it. A blemish overlaps its upper-right/middle portion but does not form a grave accent.
\[
p=p_a=p_g=1,
\]
e il bigenere $P=1$, essendovi su di esso una sola curva canonica d’ordine zero.

Per $\pi=3$ si è condotti alla superficie generale del $4^{\mathrm{o}}$ ordine $F_4$, che ha veramente il genere $p=1$, e possiede una curva canonica d’ordine zero.

Per $\pi=4$ saremo condotti a considerare una superficie $F_6$ di $S_4$, quale si ottiene come intersezione completa di una quadrica $Q_2$ e di una varietà cubica $V_3$. Infatti le $\infty^{14}$ quadriche di $S_4$ segano su $F_6$ il sistema doppio delle sezioni iperpiane, che è di genere
\[
4+4+6-1=13,
\]
e quindi anche di dimensione 13; perciò vi è una quadrica $Q_2$ che contiene $F_6$. Similmente si prova che per la $F_6$ passano $\infty^5$ varietà cubiche $V_3$ irriducibili, sicchè la $F_6$ riesce definita come intersezione completa di una $Q_2$ con una $V_3$. D’altronde è chiaro che la superficie intersezione così definita ha come sezioni iperpiane curve canoniche, il cui sisema lineare risulta pertanto un sistema $\infty^4$ di genere 4, aggiunto di sè stesso, onde segue % [sic] sisema -> sistema, the printed word omits "t".
\[
p=p_a=p_g=1,\qquad P=1.
\]

Per $\pi=5$ si è condotti ad una superficie $F_8$ d’ordine 8 di $S_5$, che facilmente si vede essere la superficie base di una rete di quadriche.

È facile prolungare la serie delle superficie di genere 1 che abbiamo costruito per i primi valori di $\pi$, ad es. considerando le intersezioni di quadriche e di varietà razionali a tre dimensioni, a curve sezioni ellittiche. Ma, senza indugiarci su questi esempi, rispondiamo alla questione d’ordine generale che qui si presenta:

\textit{Le superficie regolari di genere $p=1$ e bigenere $P=1$} (superficie dotate di curva canonica d’ordine zero) \textit{si lasciano classificare in base al valore minimo del genere $\pi$ $(>1)$ delle curve irriducibili che ad esse appartengono, e pertanto si distribuiscono in infinite famiglie aventi come tipo superficie $F_n$ d’ordine $n=2\pi-2$, a sezioni canoniche di genere $\pi$, appartenenti allo spazio $S_\pi$.}

\textit{Esistono famiglie di superficie $F_n$ rispondenti a tutti i valori di $\pi=2,\,3,\,4,\,5,\,\ldots\ldotp$, ciascuna famiglia dipendendo esattamente da 19 invarianti proiettivi esprimenti altrettanti moduli di esse rispetto alle trasformazioni birazionali. Superficie corrispondenti a diversi valori di $n$ o di $\pi$ sono generalmente irriducibili per trasformazioni birazionali.}

% scan: s137R
% !PRINTED 249
\weberpage{249}
% head: CLASSIFICAZIONE DELLE SUPERFICIE DI GENERE LINEARE ECC.
% stamp: Blotch in the upper-right margin, scattered specks, and dark left scan edge.
L’esistenza effettiva di superficie $F_n$ per tutti i valori interi di $\pi$, o per i valori pari di $n=2\pi-2$, si riconosce osservando che esistono almeno $\infty^{18}$ superficie $F_n$ particolari le quali si ottengono come proiezione di una $F_{n+2}$ dotata di punto doppio. Per dimostrarlo basta osservare che esistono fra le $\infty^{19}$ $F_n$, $\infty^{18}$ $F_n$ particolari che contengono una conica. Per il momento giustifichiamo questa affermazione con un semplice computo di costanti: a priori l’esistenza di un iperpiano quadritangente ad $F_n$ importerà $x-\pi+4$ condizioni, qualora l’esistenza di uno di questi porti di conseguenza l’esistenza di $\infty^x$ iperpiani simili. Ma se la sezione di un iperpiano quadritangente alla $F_n$ si spezzi in una conica e in una residua curva d’ordine $n-2$ (che dovrà avere con la conica 4 punti a comune) l’esistenza dell’anzidetto iperpiano trarrà con sè quella di un sistema $\infty^{\pi-3}$ di iperpiani simili, quali sono gli iperpiani per il piano della conica, e si avrà quindi $x=\pi-3$. Pertanto, avendosi $x-\pi+4=1$, non vi saranno in generale iperpiani quadritangenti ad $F_n$, intersecanti la superficie in una curva spezzata di cui fa parte una conica, e le superficie per cui siffatti iperpiani esistono, cioè le $F_n$ contenenti una conica, dovranno soddisfare ad una condizione.

Ora, partendo da una $F_n$, contenente una conica $C$, si consideri sopra di essa il sistema somma di $C$ e delle sezioni iperpiane; questo sarà un sistema di genere $\pi+1$ e dimensione $\pi+1$, rispetto a cui la $C$ è curva fondamentale; pertanto questo sistema condurrà all’esistenza di una $F_{n+2}$ di $S_{\pi+1}$ a curve sezioni canoniche, dotata di un punto doppio $O$, che risponde a $C$, la quale viene proiettata da $O$ nella $F_n$ da cui siamo partiti, o in una superficie proiettiva ad essa\origfootnote{1}{s137R-1}{Come si vede, il ragionamento precedente viene suggerito dal tentativo di costruire induttivamente le $F_n$ con $n=2\pi-2$ (passando da $\pi$ a $\pi+1$) mediante il procedimento costruttivo dei piani multipli di \textsc{Chisini} (cfr. V, 10). Ma la questione delicata che s’incontra per quella via (in ordine alla scelta di una conveniente conica sestitangente alla curva di diramazione), si risolve qui coll’uso del teorema di \textsc{Riemann-Roch} per le superficie coi generi uno.}.

Ma la famiglia delle superficie $F_{n+2}$ dotate di punto doppio, dipendendo soltanto da 18 invarianti assoluti, sarà contenuta in una famiglia di $F_{n+2}$ dipendenti da 19 invarianti, e perciò prive in generale di punti doppi.

Invero l’esistenza di questa famiglia più generale di $F_n$, cioè il relativo computo dei moduli, si giustifica (come nel § 11 del cap. V) rappresentando la superficie stessa sopra un piano $n$-plo e calcolando la dimensione del sistema di curve (con un dato numero di nodi e di cuspidi) definito dalla curva di diramazione $C$. Qui, trattandosi di una $F_n$ dotata di punto doppio, la $C$ avrà un punto doppio da
%CONT>>
% scan: s138L
% !PRINTED 250
\weberpage{250}
% head: CAPITOLO SETTIMO
%<<CONT
% stamp: Dark right scan edge, scattered specks, and short non-text strokes.
ritenere virtualmente inesistente, che a priori si comprende debba scomparire per la curva più generale del sistema sopra indicato.

Tuttavia le nostre affermazioni danno origine a qualche dubbio critico d’ordine delicato. Occorre provare che:

$1^{\mathrm{o}}$) le $F_n$ le quali (secondo il § 11 Cap. V) posseggono 19 moduli almeno; posseggono 19 moduli e non più; % ?? uncertain: Punctuation after "almeno" -- semicolon or comma.
% EDITORIAL D0132 RISOLTO: Both the upper dot and the lower comma-shaped part are visible in the enlarged scan, confirming a semicolon despite the unusual syntax.

$2^{\mathrm{o}}$) le $F_n$ (o le $F_{n+2}$) contenenti una conica posseggono effettivamente 18 moduli e non più, cioè che la condizione richiesta per il possesso della conica non è identicamente soddisfatta per la $F_n$ più generale della sua famiglia.

Queste due affermazioni si giustificano in base al computo degli invarianti proiettivi che spettano ad una superficie razionale a curve sezioni canoniche di genere $\pi$, dotata di punto triplo, appartenente allo spazio $S_\pi$. Tali superficie si lasciano rappresentare sul piano mediante un sistema lineare $\infty^\pi$, di genere $\pi$ e di grado $2\pi-2$, possedente una cubica fondamentale, su cui debbono trovarsi 12 punti base del sistema costituenti un $G_{12}$\origfootnote{1}{s138L-1}{Il teorema a cui qui facciamo appello trovasi in P. \textsc{Du Val}, \textit{On rational surfaces whose prime sections are canonical curves.} (Proc. of the Math. Society of London, 1932). Esso si trova anche virtualmente stabilito nella esposizione di \textsc{Enriques-Conforto}, \textit{Le superficie razionali}, Libro I, § 41, pag. 184 e § 46, pag. 220. In questo ultimo luogo (cfr. anche § 46, pag. 223) si dimostra che un sistema lineare $|C|$ si può trasformare in guisa che non possegga curve fondamentali staccantisi dal sistema aggiunto, ovvero in un sistema lineare costituito da curve di un certo ordine $n$ dotato di due punti base di molteplicità $r$ ed $s$, con $n=r+s$; o invece di curve di un certo ordine $n$, con punto base $(n-r)$-plo, ed altri punti base $r$-pli ad esso infinitamente vicini in direzioni distinte. In questi ultimi due casi il sistema $|C|$, supposto di genere $\pi$, ha il grado $n>2\pi-2$; perciò i sistemi $|C|$ a serie caratteristica canonica ricadono tutti nell’enunciato di pag. 184.}; si hanno perciò esattamente 12 invarianti assoluti, in corrispondenza agli 11 punti indipendenti del $G_{12}$, e all’invariante assoluto della cubica. Siccome l’imposizione di un punto triplo ad una superficie $F_n$, importa 6 condizioni (e non più) si deduce che le $F_n$ di $S_\pi$ non possono formare una famiglia dipendente da più che 19 moduli. La stessa considerazione mostra che le $F_{n+2}$ generali, dipendenti da 19 invarianti, non possono possedere un punto doppio, chè, imponendo a questo punto di diventare triplo, si viene ad assoggettare la superficie a 6 condizioni e non più. Da ciò segue anche che le $F_n$ possedenti una conica, derivando dalla proiezione di una $F_{n+2}$ con punto doppio, dipendono da 18 e non da 19 moduli. % ?? uncertain: Second character of "sistema" in "punti base del sistema" -- damaged i or an apostrophe, giving "s’stema".
% EDITORIAL D0133 RISOLTO: Two vertically aligned ink components are visible: a dot and a shortened stem. This is a damaged i rather than a single apostrophe.

Resta da giustificare l’affermazione che le superficie $F_n$ $(n=2\pi-2)$ rispondenti a diversi valori di $\pi$ $(>1)$, sono generalmente irriducibili. Ciò può farsi induttivamente in base ad un semplice
% The equality sign in "(n= =2\pi-2)" is printed both before and after a line break.
%CONT>>
% scan: s138R
% !PRINTED 251
\weberpage{251}
% head: CLASSIFICAZIONE DELLE SUPERFICIE DI GENERE LINEARE ECC.
%<<CONT
% stamp: Scattered non-text specks and short marginal strokes.
computo di costanti, come quello di cui ci siano valsi innanzi per mostrare che una $F_n$ non possiede in generale una conica. Con un simile computo si riconosce invero che « \textit{sopra una superficie $F_n$ di $S_\pi$} (corrispondente al minimo valore $\pi$ del genere delle curve che ad essa appartengono) \textit{non vi sono in generale altri sistemi lineari di curve che il sistema delle sezioni iperpiane e i suoi multipli} ». % [sic] siano -> siamo, the printed verb is inconsistent with the first-person plural construction "ci ... valsi".

Anzitutto si domanda se fra gli $\infty^\pi$ iperpiani di $S_\pi$ possono aversi degli iperpiani $m$-tangenti ad $F_n$ tali che le loro sezioni con $F_n$ risultino spezzate in due curve di generi $\pi_1$ e $\pi_2$, con $m$ punti comuni: $\pi=\pi_1+\pi_2+m-1$. Se esiste un iperpiano siffatto, ne esistono $\infty^x$, con $x=\pi_1+\pi_2$; per ciò il numero dei parametri da cui dipendono gli iperpiani cercati sarà $\pi-m-x=-1$: vuol dire che le $F_n$ possedenti siffatte curve spezzate sono soggette ad una condizione. In modo simile si riconosce che la sezione di una $F_n$ con una ipersuperficie di un certo ordine qualsiasi non potrà in generale spezzarsi se non in due curve che siano alla loro volta intersezione completa di due ipersuperficie; in quest’ultimo caso il computo di costanti dà ancora una condizione da soddisfare per la superficie $F_n$, ma è ovvio che tale condizione risulta qui come conseguenza del fatto che la $F_n$ contiene il sistema delle sezioni iperpiane, e perciò viene identicamente soddisfatta.

Ora è necessario avvertire che, come già si è detto innanzi (trattando delle $F_n$ che contengono una conica), il risultato a cui conduce il nostro computo di costanti lascia sussistere un dubbio che già toglie valore alla dimostrazione di Noether che la superficie del $4^{\mathrm{o}}$ ordine dello spazio ordinario contiene in generale soltanto curve intersezioni complete.

Tuttavia la cosa può essere stabilita rigorosamente ed anche semplicemente, sulla base del principio di continuità, riducendo una $F_n$ al caso limite della superficie razionale dotata di punto triplo. Ci limiteremo ad indicare, in breve, come possa svolgersi questa dimostrazione nel caso elementare del teorema di Noether sulle superficie del $4^{\mathrm{o}}$ ordine: l’estensione del ragionamento alle $F_n$ di genere uno non presenta difficoltà.

Qui dunque si fa variare per continuità una superficie quartica $F_4$, finchè si riduca ad un generale monoide $\Phi_4$ possedente un punto triplo $O$: il monoide, che contiene 12 rette $a_1a_2a_3\ldots\ldotp a_{12}$ per $O$, si lascia rappresentare sul piano, per proiezione da $O$, mediante un sistema lineare di curve del $4^{\mathrm{o}}$ ordine $C_4$ passanti per i 12 punti base $A_1 A_2\ldots\ldotp A_{12}$ tracce delle rette $a_i$: i quali appartengono ad una cubica fondamentale $C_3$ di genere uno.

La dimostrazione che ci proponiamo di svolgere mira a stabilire che: mentre la $F_4$ tende per continuità al monoide generale $\Phi_4$, un sistema lineare irriducibile completo, di curve $C$, tracciate su
%CONT>>
% scan: s139L
% !PRINTED 252
\weberpage{252}
% head: CAPITOLO SETTIMO
%<<CONT
% stamp: Dark right scan edge, scattered specks, and short non-text strokes.
$F_4$, si riduce al limite ad un sistema lineare di curve $\overline{C}$, intersezioni complete di $\Phi_4$ con superficie d’un certo ordine $d$, che non ne contengono il punto triplo $O$. Allora, designando con $\pi$ il genere delle $\overline{C}$, e quindi anche delle $C$, ed essendo perciò $2\pi-2$ il grado e $\pi$ la dimensione di $|C|$, resulterà che le dette $C$ su $F_4$ debbono avere $2\pi-2$ punti comuni colle intersezioni delle superficie $F_d$ d’ordine $d$, le quali sono egualmente di genere $\pi$ e di grado $2\pi-2$; siccome una $C$ non può contenere una serie $g_{2\pi-2}$ di dimensione $\pi$, si deduce che il sistema $|C|$ su $F_4$ coincide col sistema lineare delle intersezioni complete delle $F_d$.

Per sviluppare la dimostrazione così indicata, osserviamo che, quando la $F_4$ si riduce per continuità a $\Phi_4$, a priori un sistema $|C|$ di genere $\pi$ e grado $2\pi-2$, appartenente ad $F_4$, si ridurrà ad un sistema $\infty^\pi$ di curve limiti \textit{connesse} formate dalle intersezioni di $\Phi_4$ con superficie passanti per $O$ un certo numero di volte, aumentate delle rette $a_i$, contate rispettivamente $h_i$ volte $(h_i\geq0)$ e dall’intorno di $O$ contato $s$ $(\geq0)$ volte. Quindi le immagini di tali curve sul piano rappresentativo di $\Phi_4$ saranno, a priori, curve d’un certo ordine $n$, appartenenti ad un sistema lineare completo
\[
|C_n|=C_m(A_1^{r_1} A_2^{r_2}\ldots A_{12}^{r_{12}})+sC_3+\Sigma h_iA_i,
\]
dove le $C_m$ (d’ordine $m\leq n$) appartengono ad un sistema lineare irriducibile $\infty^\pi$. E siccome le $C_m$ hanno il genere $\pi$ e il grado $2\pi-2$, sarà
\[
\begin{aligned}
n^2-\Sigma(r_i+s-h_i)^2&=2\pi-2\\
n(n-3)-\Sigma(r_i+s-h_i)(r_i+s-h_i-1)&=2\pi-2,
\end{aligned}
\]
e quindi
\[
3n=\Sigma(r_i+s-h_i):
\]

Ciò significa che per il sistema $|C_n|$ la cubica $C_3$ è \textit{virtualmente fondamentale}, cioè avente zero intersezioni con le $C_n$.

Conviene discutere anzitutto l’ipotesi:

1) la $C_3$ è anche effettivamente fondamentale per il sistema completo $|C_m|$, cioè % ?? uncertain: Small mark immediately after "1)" -- a scan speck or a comma.
% EDITORIAL D0134 RISOLTO: The extra mark is a tiny cross-shaped speck immediately below the parenthesis; it lacks a comma body or descending hook.
\[
|C_n|=|C_m|,\qquad s=0,\qquad h_i=0.
\]

In questa ipotesi dimostriamo facilmente che $|C_n|$ è un multiplo del sistema $C_4$ $(A_1A_2\ldots\ldotp A_{12})$ e perciò il corrispondente sistema su $\Phi_4$ è costituito dalle intersezioni complete di questo monoide colle superficie $F_d$ d’ordine $d=\dfrac{n}{4}$.

Infatti, designando con $G$ il gruppo sezione della cubica ellittica
%CONT>>
% scan: s139R
% !PRINTED 253
\weberpage{253}
% head: CLASSIFICAZIONE DELLE SUPERFICIE DI GENERE LINEARE ECC.
%<<CONT
% stamp: Scattered non-text specks and marginal marks.
$C_3$ con una retta, avremo su questa le seguenti equivalenze:
\[
\begin{aligned}
\Sigma r_iA_i&\equiv nG\\
\Sigma A_i&\equiv4G;
\end{aligned}
\]
moltiplicando la seconda relazione per $r_1$ e sottraendo la prima si trova:
\[
\mathop{\Sigma}\limits_{i=2}^{i=12}(r_i-r_1)A_i\equiv(n-4r_1)G\,;
\]
ma quest’ultima relazione non può sussistere, essendo arbitrariamente scelto il gruppo degli 11 punti $A_2A_3\ldots\ldotp A_{12}$, se non è $n=4r_1$ e
\[
r_1=r_2=\ldots\ldotp=r_{12}=r;
\]
di modo che le $C_n$ sono curve d’ordine $4r$ appartenenti al sistema $|rC_4|$, \hfill c. d. d.

Se, invece, l’ipotesi 1) non è soddisfatta, avremo due possibilità che si tratta di escludere.

2) Anzitutto che delle $C_n$ non faccia parte la $C_3$, ma soltanto l’intorno di qualche punto $A_i$:
\[
s=0,\qquad h_i>0.
\]

Ora se la $C_3$ è fondamentale per $|C_m|$, questo sistema lineare irriducibile è un multiplo, secondo un certo numero $r=r_1=r_2=\ldots\ldotp=r_{12}$ di $C$, e — se non deve ampliarsi quando gli si sommino gli intorni dei punti $A_i$ — bisogna che il-punto sommando sia uno solo, per esempio $A_1$, e che abbia per $C_m$ la molteplicità 1, quindi $r=1$. Ma è ovvio che una curva $C$ passante per i 12 punti $A_i$ sommata ad uno di questi, non può presentarsi come limite di una curva tracciata sopra $F_4$, quando questa superficie si riduca al monoide $\Phi_4$. % ?? uncertain: Mark between "il" and "punto" -- a printed hyphen or a non-text stroke.

Se, invece, si suppone che la $C_3$ non sia fondamentale per $|C_m|$, questo sistema lineare è regolare (e non sovrabbondante), quindi esso si amplia sempre se gli si sommi l’intorno di qualche punto $A_i$. E ciò va contro l’ipotesi che il $|C_n|$ sia completo.

3) Finalmente conviene discutere l’ipotesi che la cubica $C_3$ compaia un certo numero $s>0$ di volte come parte fissa delle $C_n$. In questo caso il $|C_m|$, che non possiede la $C_3$ come curva fondamentale, è regolare, e quindi si amplia se gli si sommi l’intorno di un punto $A_i$ che non sia punto base per esso: così potranno figurare tra i punti costituenti la somma $\Sigma h_iA_i$, soltanto un certo numero $t$ di punti $A_i$ e ciascuno $h_i\leq s$ volte. Designando con $l=C_mC_3$
%CONT>>
% scan: s140L
% !PRINTED 254
\weberpage{254}
% head: CAPITOLO SETTIMO
%<<CONT
% stamp: Scattered specks, short non-text strokes, and traces of a dark right scan edge.
il numero delle intersezioni variabili di una $C_m$ colla cubica $C_3$, si avrà ora
\[
l+ts-3s\geq0.
\]

Quindi si sarà condotti a distinguere due casi a priori possibili: o $l+t\geq4$, dalla quale diseguaglianza segue l’ampliarsi del sistema lineare $C_m+C_3+\Sigma A_i$; ovvero $l+t<4$, che porta $s\leq1$; in quest’ultimo caso la dimensione di $|C_m|$ risulterebbe inferiore di 1 al genere di $|C_n|$: e perciò non potrebbe $|C_n|$ essere il sistema limite di un sistema lineare appartenente alla $F_4$, la cui dimensione eguaglia il genere.

Escluse così le ipotesi 2) e 3) resta soltanto l’ipotesi 1) da cui, come abbiam detto, si trae che le curve appartenenti ad una generale $F$ sono soltanto le intersezioni complete di questa con altre superficie d’ordine $d=1,\,2,\,\ldots$.

\textit{Osservazione.} – Come si è detto il ragionamento che precede si estende al caso delle superficie $F$ di genere 1. Tuttavia si arriva così a riconoscere, non già che tutte le curve appartenenti ad una $F$ generale con 19 moduli sono intersezioni di essa con ipersuperficie, bensì che tutte queste curve e in particolare le sezioni iperpiane di $F$, ove non siano le curve irriducibili di genere minimo appartenenti ad $F$, costituiscono sistemi multipli d’un medesimo sistema di curve di genere $\varrho$ e di grado $\nu$; si avrà dunque
\[
n=s\nu=2\pi-2=s(2\varrho-2).
\]

Viceversa, risalendo da una superficie razionale con punto triplo ad una superficie di genere uno a sezioni iperpiane canoniche, si riconosce (almeno induttivamente) che quando
\[
n=s\nu\qquad \text{e}\qquad 2\pi-2=s(2\varrho-2),
\]
esistono almeno due famiglie di superficie $F_n$ di genere 1: una superficie multipla $s$-pla ed una superficie semplice, l’una e l’altra dipendenti da 19 moduli e fra loro irriducibili.

\edsection{3. Notizia storica.}{3. Notizia storica.}{sec.VII.3}

I primi esempi di superficie $F_n$ di genere 1 contenenti lo stesso numero di moduli delle superficie del $4^{\mathrm{o}}$ ordine, e ad esse irriducibili, sono stati indicati da \textsc{Castelnuovo}\origfootnote{1}{s140L-1}{Cfr. F. \textsc{Enriques}, \textit{Ricerche ecc.}, III, 6.} (1893). Nel 1896, \textsc{Enriques}\origfootnote{2}{s140L-2}{F. \textsc{Enriques}, \textit{Sui piani doppi di genere uno.} Memorie della Soc. It. d. Scienze (detta dei XL), 1896.} caratterizzava, dal punto di vista delle trasformazioni birazionali, la famiglia di queste superficie regolari di genere 1 con curva ca%
%CONT>>
% scan: s140R
% !PRINTED 255
\unskip\weberpage{255}%
% head: CLASSIFICAZIONE DELLE SUPERFICIE DI GENERE LINEARE ECC.
%<<CONT
% stamp: Dark left scan edge and scattered non-text specks.
nonica d’ordine zero, mediante il valore del bigenere $P=1$. Più tardi \textsc{Enriques}\origfootnote{1}{s140R-1}{F. \textsc{Enriques}, \textit{Le superficie di genere uno}, Rendic. Acc. di Bologna, 1908.} e \textsc{Severi}\origfootnote{2}{s140R-2}{F. \textsc{Severi}, \textit{Le superficie algebriche con curva canonica d’ordine zero.} Atti Ist. Veneto (vol. LXV, 2), 1909.}, indipendentemente l’uno dall’altro, riuscivano a dimostrare l’esistenza di famiglie di superficie $F_n$ di $S_\pi$, d’ordine pari $n=2\pi-2$, per tutti i valori di $\pi$. L’uno e l’altro prendevano come punto di partenza il computo dei moduli fatto da \textsc{Enriques} (1908 Cfr. Cap. V, § 11). Poichè le $F_n$ dipendono da 19 moduli (almeno), l’esistenza di una famiglia di $F_n$ con tanti moduli, risulta per \textsc{Enriques} dal fatto che si hanno particolari $F_n$, iperellittiche, con 3 moduli, dotate di 16 punti doppi: di qui anche si argomenta che il numero dei moduli appartenenti alle $F_n$ dovrà essere 19 e non superiore; la prova di ciò richiede di stabilire che il possesso di 16 punti doppi caratterizza le $F_n$ iperellittiche, siccome è noto accadere per $n=4$ (superficie di Kummer).

\textsc{Severi} riconosce per via algebrico-geometrica l’esistenza delle $F_n$ $(n=2\pi-2)$ per qualsiasi valore di $\pi$, costruendo le $F_n$ particolari con 4 punti doppi, che si ottengono come trasformate di quartiche $F_4$ possedenti 4 curve di genere (effettivo e virtuale) zero, senza punti comuni. Ammettendo che il possesso di tali curve dia luogo per le $F_4$ a 4 condizioni indipendenti (ciò che implica il teorema di Noether che « le curve appartenenti ad una $F_4$ generale sono intersezioni complete » e qualcosa di più), ne deduce che le $F_n$ generali non possono contenere più di 19 moduli, e quindi dipendono esattamente da 19 moduli.

Qui è opportuno notare che la costruzione di $F_n$ si può fare a partire da quartiche o da altre superficie particolari coi generi uno, in molti modi diversi, che richiedono però qualche attenzione come è accennato nelle « Lezioni » di Enriques-Campedelli (pag. 311). Per esempio si ottengono $F_n$ con $n=2\pi-2$, dove $\pi$ è della forma $\pi=2+3s$, a partire da un piano doppio con sestica di diramazione che sia l’inviluppo di una serie $\infty^1$ di cubiche: all’uopo basta sommare alle curve $C$, curve di genere due, immagini delle rette del piano, $s$ curve ellittiche $K$, immagini delle dette cubiche, costruendo così il sistema $C+sK$, che è appunto di genere $\pi=2+3s$\origfootnote{3}{s140R-3}{Se $\pi$ è della forma $\pi=3+3s$, si procederà similmente a partire da una $F_4$ contenente una retta; e se $\pi=4+3s$, si partirà invece da una $F_6$ particolare contenente una cubica piana.}.

Siccome si parte da particolari superficie aventi esattamente 18 moduli ed a queste rispondono tutte le $F_n$ contenenti una cubica piana, e perciò soggette ad \textit{una} condizione (e non più), si deduce così che le $F_n$ dipendono proprio da 19 moduli, e non da più.

% scan: s141L
% !PRINTED 256
\weberpage{256}
% head: CAPITOLO SETTIMO
% stamp: Dark right scan edge, scattered specks and stray marginal strokes.
Frattanto le questioni di cui si discorre sono state chiarite da \textsc{A. Franchetta}\origfootnote{1}{s141L-1}{\textsc{A. Franchetta}, « \textit{Sulle superficie regolari di genere uno con curva canonica d’ordine zero} ». Istituto Lombardo, 1942.}, nel modo che abbiamo esposto nel testo. L’autore avendo cercato di collegare la dimostrazione dell’esistenza della serie infinita delle $F_n$ al teorema d’esistenza di Chisini, è condotto a dare a codesta dimostrazione la forma più semplice, passando dalle $F_n$ particolari contenenti una conica alle $F_{n+2}$. \textsc{Franchetta} dimostra poi che le $F_n$ (per un dato valore di $n$) non possono dipendere da più che 19 moduli, ricorrendo alle $F_n$ razionali con un punto triplo: appunto come si è veduto innanzi.

Che le $F_n$, corrispondenti a valori diversi di $n$, siano generalmente irriducibili per trasformazioni birazionali dipende, come abbiam detto, dal teorema che « sopra una $F_n$ generale non si hanno altre curve che intersezioni complete con ipersuperficie ». Del fatto così enunciato si acquista facilmente la nozione in base al teorema di Noether, concernente le quartiche $F_4$, ed al computo di costanti su cui esso si fonda. Ma con ciò non si ha ancora una dimostrazione rigorosa. \textsc{Enriques} e \textsc{Severi} hanno pensato che una tale dimostrazione potrebbe darsi (per $n\geq 4$) facendo ridurre per continuità una $F_n$ generale ad una $F_n$ iperellittica, con moduli generici, valendosi della proprietà, che si riconosce per via trascendente\origfootnote{2}{s141L-2}{\textsc{Enriques} e \textsc{Severi}, \textit{Mémoire sur les surfaces hyperelliptiques.} Acta Math., vol. XXXII, n. 27.}, che, su di esse, non vi sono altre curve se non le intersezioni complete con ipersuperficie e le curve di contatto con ipersuperficie tangenti. E \textsc{Severi} ha realizzato questo disegno, in modo elegante (resta nel suo discorso una piccola lacuna cui si supplisce facilmente). La dimostrazione algebrico-geometrica che qui si è porta, riducendosi al caso limite delle $F_n$ razionali con punto triplo, appartiene a \textsc{A. Franchetta}, nella sua nota già innanzi citata.

Qui daremo qualche indicazione bibliografica intorno alle varie dimostrazioni del teorema di Noether su le superficie del $4^\circ$ ordine e su le superficie generali d’ordine superiore.

\textsc{M. Noether}, « Zur Grundlegung der algebraischen Raumcurven ». Abhandlungen der Akademie der Wissenschaften, Berlin, 1883.

\textsc{K. Rohn}, « Die Raumcurven auf der Flaechen vierter Ordnung ». Berichten der koenigl. Saechs. Gesellschaft der Wissenschaften, Leipzig, 1897. % [sic] auf der Flaechen -> auf den Flaechen, apparent German agreement error.

\textsc{F. Severi}, « Le superficie algebriche con curva canonica d’ordine zero ». Atti Ist. Veneto (LXV, 2), 1909.

\textsc{G. Fano}, « Sulle varietà algebriche che sono intersezioni complete di più forme ». R. Acc. di Torino, vol. XLIV, 1909.

% scan: s141R
% !PRINTED 257
\weberpage{257}
% head: CLASSIFICAZIONE DELLE SUPERFICIE DI GENERE LINEARE ECC.
% stamp: Scattered scan specks and stray marginal marks.
\textsc{S. Lefschetz}, « On certain numerical Invariants of algebraics Varieties. (Trans. of Am. Math. Society, vol. 22, 1921); « Concerning Noether’s theorem on curves traced on non singular surfaces ». Bull. Am. Math. Society, vol 29, 1923. % [sic] algebraics Varieties. -> algebraic Varieties »., apparent extra s and missing closing guillemet.

\textsc{A. Franchetta}, « Sulle superficie regolari di genere uno con curva canonica d’ordine zero ». Istituto lombardo, 18 giugno 1942.

\textsc{B. D’Orgeval}, « Sur les surfaces algébriques dont tous les genres sont 1 ». Gauthier-Villars, Paris, 1945.

A chi persegua una classificazione completa delle superficie coi generi 1, si presenta ora una nuova domanda, cioè se le superficie $F_n$ con $n=2\pi-2$, a curve sezioni canoniche in un $S_\pi$, formino, entro questo spazio, una sola famiglia, ovvero più famiglie diverse, ciascuna dipendente da 19 invarianti assoluti o moduli. A questo proposito, d’accordo coll’osservazione finale del paragrafo precedente, si può notare che già nello spazio $S_5$ si hanno 2 tipi di $F_8$, cioè la superficie semplice $F_8$ intersezione di tre quadriche, e la superficie di Veronese $F_4$ doppia, con curva di diramazione d’ordine 12, che è la trasformata di un piano doppio con sestica di diramazione. Analogamente si trovano nello spazio $S_9$ due tipi irriducibili di $F_{16}$, una delle quali si deduce da una superficie del $4^\circ$ ordine dello $S_3$. Di qui si è condotti a riconoscere in $S_6$ due tipi di $F_{10}$ particolari dotate di punto doppio, che da questo punto vengano proiettate rispettivamente nella $F_8$ e nella $F_4$ doppia sopra indicata. Ma \textsc{P. du Val}\origfootnote{1}{s141R-1}{\textsc{P. du Val}, \textit{Superficie di genere uno che non sono base per un sistema di quadriche} e \textit{Osservazioni sulle superficie di genere uno che non sono base per un sistema di quadriche.} Rend. R. Acc. dei Lincei, febbraio-marzo, 1932.} ha mostrato che queste due $F_{10}$ appartengono ad una medesima famiglia di $F_{10}$ generali, senza punti doppi; e questo in accordo con il risultato generale stabilito da \textsc{B. d’Orgeval}, nel suo volume già citato, il quale appunto ha dimostrato che le superficie a curve sezioni canoniche dello $S_\pi$ formano una sola famiglia. % ?? uncertain: second digit of the subscript in F_{16} after S_9 -- 6 or 8.
% EDITORIAL D0136 RISOLTO: The enlarged second subscript digit has the lower bowl and open upper-left stroke of 6; it is not an 8 with two closed bowls.

\edsection{4. Superficie di genere lineare $p^{(1)}=1$ con $P>1$.}{4. Superficie di genere lineare p\textasciicircum{}(1)=1 con P>1.}{sec.VII.4}

Si abbia una superficie $F$, di genere lineare (assoluto) $p^{(1)}=1$ e di genere superficiale $p>1$. Le sue $\infty^{p-1}$ curve canoniche $K$ debbono formare un sistema lineare di curve ellittiche, di grado $p^{(1)}-1=0$, che sarà dunque composto colle curve ellittiche d’un fascio (senza punti base). Lo stesso può dirsi delle curve bicanoniche o $i$-canoniche di $F$, che debbono formare un sistema lineare di genere
\[
\frac{i(i+1)}{2}(p^{(1)}-1)+1=1,
\]
% EDITORIAL: segnatura originale: \hbox to \linewidth{\textsc{Enriques F.} - \textit{Superficie algebriche.}\hfil 17}
\printersignature{\textsc{Enriques F.} - \textit{Superficie algebriche.}\quad 17}
%CONT>>
% scan: s142L
% !PRINTED 258
\weberpage{258}
% head: CAPITOLO SETTIMO
%<<CONT
% stamp: Dark right scan edge, scattered specks and a diagonal stray stroke through the middle portion of the page.
e di grado
\[
i^2(p^{(1)}-1)=0.
\]

Dunque: \textit{le curve canoniche o pluricanoniche d’una superficie di genere lineare $p^{(1)}=1$ sono in generale riducibili, componendosi colle curve ellittiche d’un fascio.}

L’esistenza d’un fascio di curve ellittiche, e d’un fascio lineare per le superficie regolari, deriva quindi come conseguenza dall’essere il $p^{(1)}=1$, ogni qualvolta il genere della superficie sia $p>1$ o uno dei suoi plurigeneri $P_i$ risulti
\[
P_i>1.
\]

Ma questa diseguaglianza sussisterà sempre per le superficie regolari, non razionali, che posseggano qualche curva canonica o bicanonica d’ordine maggiore di zero, giacchè in tale ipotesi si è visto (Cap. VI, \S\ 2) che è per $p>1$
\[
P\geq 2,
\]
e per $p=0$:
\[
P\geq 1,\qquad P_6\geq 2.
\]

Siccome, d’altra parte, si è pur trovato che le superficie per cui $p=0$ e $P=1$, con curva bicanonica d’ordine zero ($P_6=1$), posseggono sempre dei fasci lineari di curve ellittiche, così è lecito affermare che:

\textit{Le superficie regolari (non razionali), di genere lineare (assoluto) $p^{(1)}=1$; posseggono un fascio lineare di curve ellittiche, ad eccezione di quelle che han tutti i generi eguali ad $1$} (superficie $F_n$ con $n=2\pi-2$, a sezioni canoniche in $S_\pi$ caratterizzate da $p=P=1$). % ?? uncertain: punctuation after p^{(1)}=1 -- semicolon or comma.

Per classificare le superficie $F$ di genere lineare $p^{(1)}=1$, con $P>1$, contenenti un fascio lineare di curve ellittiche $C$, conviene rilevare anzitutto che ad esse appartengono \textit{due caratteri invarianti}, suscettibili di assumere \textit{valori interi arbitrari}, cioè:

1) il genere $p$ ($=p_a$) che è definito dal numero $S$ delle curve $C$ del fascio dotate di punto doppio, avendosi l’invariante di Zeuthen-Segre $I=S-4=12p+8\ (=12p_a-p^{(1)}+9)$,

2) e il \textit{determinante} $d$, che si definisce come il minimo numero di punti determinabile razionalmente sopra una $C$, in funzione del parametro da cui la $C$ stessa dipende.

Per comprendere il significato di quest’ultimo carattere conviene ricordare che in generale sopra una curva ellittica $C$ d’ordine $n$, del piano o d’un qualunque spazio $S_r$, non si possono costruire \textit{razionalmente} (in funzione dei coefficienti della sua equazione) altre serie che i multipli della $g_n$, segata dalle rette del piano o dagli iperpiani dello $S_r$. E così non è possibile in generale abbassare l’ordine
%CONT>>
% scan: s142R
% !PRINTED 259
\weberpage{259}
% head: CLASSIFICAZIONE DELLE SUPERFICIE DI GENERE LINEARE ECC.
%<<CONT
% stamp: Dark left scan edge and scattered specks.
$n$ di una $C$ ellittica variabile operando razionalmente sui parametri da cui dipende la variazione di essa\origfootnote{1}{s142R-1}{Cfr. \textsc{F. Enriques}, \textit{Sulle superficie algebriche con un fascio di curve ellittiche.} Rendic. Acc. Lincei, 1912. In questa nota la deduzione del n. 3 è infirmata da un errore di calcolo. Cfr. pure \textsc{Enriques-Chisini}, \textit{Lezioni}, vol. III, pag. 345.}.

\edFlush
Pertanto una serie $\infty^1$ di curve ellittiche $C$, d’ordine $n$, assunta nello spazio, formerà in generale una superficie $F$ con un fascio di $C$ la quale non sarà suscettibile di trasformazione che abbassi l’ordine delle $C$.

Ora possiam dire che: le superficie $F$, con un fascio di curve ellittiche $C$, di \textit{determinante} $d>1$ si possono trasformare in guisa che le $C$ diventino \textit{curve d’ordine} $d$, e così --- nell’ipotesi della regolarità --- in superficie $F_m$ d’un certo ordine $m$, dello spazio ordinario, con retta $(m-d)$-pla (per $d=2$ in piani doppi); e non è possibile una trasformazione di queste che abbassi l’ordine delle $C$.

Per giustificare l’asserto si noti che --- a tenore della definizione --- la $F$ conterrà una curva $L$ $d$-secante le curve $C$, e non altre curve che seghino le $C$ in un numero minore di punti. Ora, sommando una $C$ alla $L$, se ne aumenterà il grado di $2d$ e il genere di $d$, sicchè la dimensione di $|L+C|$ data dal teorema di Riemann-Roch, crescerà su quella di $|L|$ di $2d>d>1$, e quindi le curve $L+C$, $L+2C$, $\ldotp\ldotp\ldotp\ldotp\ L+sC$, saran contenute in più ampi sistemi lineari irriducibili, $d$-secanti le $C$. Per $d=2$ si ottiene un piano doppio, con una certa curva di d’ramazione $D_{2n}$, su cui le $C$ vengono rappresentate dalle rette per un punto $(2n-4)$-plo di $D_{2n}$. % ?? uncertain: character after d in d’ramazione -- apostrophe or a defective i.

Per approfondire lo studio delle nostre superficie di genere lineare $p^{(1)}=1$ (superficie regolari, non razionali, con $P>p\geq 1$), conviene studiare in particolare quelle di \textit{determinante} $d=1$.

Una superficie $F$ di questa famiglia contiene un fascio lineare di curve ellittiche $C$ ed una \textit{curva $\theta$ unisecante} le $C$. Mercè un sistema lineare $\infty^3$ contenuto in $|3\theta+sC|$ si può trasformare $F$ in una superficie $F_m$, d’un certo ordine $m$, su cui le $C$ siano cubiche segate dai piani per una retta $(m-3)$-pla $a$; queste $C$, variabili nel fascio delle dette sezioni piane, avranno un flesso razionalmente distinto, che si può supporre fisso, in un punto $A$ $(m-2)$-plo per $F_m$, su $a$; e si può anche supporre che $a$ sia la relativa tangente d’inflessione delle $C$. Quindi la $F_m$ si rappresenta, per proiezione da $A$, su un piano doppio con curva di diramazione $D_{2n}$, d’un certo ordine pari $2n$, dotata d’un punto $(2n-3)$-plo $O$ (traccia di $a$). La $D_{2n}$, potrà avere altri punti singolari abbassanti il genere del piano doppio: punti 4-pli o punti $[3,3]$. Ma la presenza di tali singolarità permette in generale di abbassare l’ordine di $D_{2n}$ con trasformazioni quadratiche
%CONT>>
% scan: s143L
% !PRINTED 260
\weberpage{260}
% head: CAPITOLO SETTIMO
%<<CONT
% stamp: Dark right scan edge, scattered specks and stray marginal strokes.
del piano. Se si vuole che la $D_{2n}$ \textit{sia d’ordine minimo}, si deve ammettere che essa appartenga ad uno dei due tipi seguenti:

I) $D_{2n}$ con un punto $O$ $(2n-3)$-plo, senza altre singolarità influenti sul genere; e

II) $D_{2n}$ con un punto $O$ $(2n-3)$-plo e un punto $B$ 4-plo, formata della retta $OB$ e di una $D_{2n-1}$ con $O$ $(2n-4)$-plo e $B$ 3-plo.

I piani doppi del tipo I sono di genere
\[
p=n-2,
\]
avendosi $\infty^{p-1}$ curve canoniche rappresentate dai gruppi di $p-1$ rette per $O$. Le curve $i$-canoniche sono date dai gruppi di $i(p-1)$ rette per $O$, quindi
\[
P_i=i(p-1)+1.
\]

Per i piani doppi del tipo II, si hanno $\infty^{n-4}$ curve canoniche rappresentate dai gruppi di $p-1=n-4$ rette per $O$, cui si aggiunge la retta di diramazione eccezionale\origfootnote{1}{s143L-1}{Sulla retta $OP$ vi sono tre punti singolari di $D_{2n}$, cioè $P$, $O$ e il punto infinitamente vicino ad $O$, quindi il grado della retta di diramazione sul piano, è $-2$, e la curva corrispondente sulla superficie si compone di due curve di grado $-1$ (e di genere 0).}, quindi il genere % [sic] OP and P -> OB and B, the point is called B in the preceding description of type II.
\[
p=n-3.
\]

Similmente lo $i$-genere vale
\[
P_i=i(p-1)+1.
\]

Concludiamo intanto:

\textit{Le superficie $F$ di genere lineare assoluto $p^{(1)}=1$ e di determinante $d=1$, con $P\geq 1$, si lasciano rappresentare su piani doppi del tipo I,} con curva di diramazione $D_{2n}$ dotata d’un punto $O$ $(2n-3)$-plo, \textit{ovvero del tipo II,} cioè con $D_{2n}$ di diramazione dotata d’un punto $O$ $(2n-3)$-plo e fuori di esso un punto 4-plo $B$: codesti piani doppi danno luogo a tante famiglie distinte secondo i valori del genere, che è pel tipo I:
\[
p=n-2
\]
e pel tipo II
\[
p=n-3,
\]
e per esse è sempre
\[
P_i=i(p-1)+1.
\]

Ogni famiglia delle nostre $F^1$ contiene infinite classi di superficie birazionalmente distinte, ciascuna delle quali è caratterizzata da
%CONT>>
% scan: s143R
% !PRINTED 261
\weberpage{261}
% head: CLASSIFICAZIONE DELLE SUPERFICIE DI GENERE LINEARE ECC.
%<<CONT
% stamp: Dark left scan edge and scattered specks.
un certo numero di parametri o moduli, che è facile valutare. La $D_{2n}$ del tipo I per cui sia fissato il punto $(2n-3)$-plo $O$, dipende da
\[
\frac{2n(2n+3)}{2}-\frac{(2n-3)(2n-2)}{2}=8n-3
\]
parametri, e --- poichè vi sono $\infty^6$ omografie che lascian fermo il punto $O$ --- possiede
\[
8n-9
\]
invarianti assoluti. Così le $F^1$ \textit{del tipo I}, per ogni valore del genere $p=n-2$ dipendono da
\[
8p+7
\]
moduli.

Inoltre le $D_{2n}$ con $O$ $(2n-3)$-plo e $B$ 4-plo dipendono da
\[
8n-13
\]
parametri e possiedono
\[
8n-13-6=8n-19
\]
invarianti rispetto alle $\infty^6$ trasformazioni quadratiche che mutano in sè i due fasci di rette $O$ e $B$. Quindi, essendo $p=n-3$, si hanno
\[
8p+5
\]
moduli.

Questi risultati contraddicono, a prima vista, ciò che si è stabilito nel \S\ 11 del Cap. V, che, in generale, una classe di superficie regolari di generi $p$ e $p^{(1)}$ dipende da almeno
\[
9p-2p^{(1)}+12
\]
moduli; poichè per $p^{(1)}=1$ si ha
\[
9p-2p^{(1)}+12=9p+10
\]
che supera $8p+7$ e $8p+5$.

Ma la spiegazione del paradosso sta in ciò, che \textit{le superficie $F^1$ di genere lineare $p^{(1)}=1$ e di determinante $1$, con $P>1$, di un dato genere $p$, rientrano in famiglie $F^2$, di determinante $2$, aventi lo stesso genere $p$.}

Infatti il piano doppio del tipo I rientra come caso particolare nel piano doppio $\mathrm{I}'$ con $D_{2n}$ di diramazione dotata d’un punto $(2n-4)$-plo $O$, che è egualmente di genere
\[
p=n-2,
\]
e dipende da
\[
10n-12
\]
invarianti, ossia da
\[
10p+8
\]
%CONT>>
% scan: s144L
% !PRINTED 262
\weberpage{262}
% head: CAPITOLO SETTIMO
%<<CONT
% stamp: Dark right scan edge, scattered specks and a diagonal stray stroke near the middle of the page.
moduli; e per
\[
n\geq 4,\qquad p>1,
\]
questo numero riesce sempre
\[
\geq 9p+10\qquad (=9p-2p^{(1)}+12).
\]

Similmente il piano doppio del tipo II rientra, come caso particolare, in quello $\mathrm{II}'$ con $D_{2n}$ di diramazione dotata d’un punto $(2n-4)$-plo $O$ e d’un punto 4-plo $B$, che è egualmente di genere
\[
p=n-3
\]
e dipende da
\[
10n-22
\]
invarianti rispetto ad $\infty^6$ trasformazioni quadratiche, ossia da
\[
10p+8
\]
moduli.

Se ora, dopo avere discorso delle superficie $F^1$ di genere lineare $p^{(1)}=1$ e di determinante $d=1$ ($P>1$), passiamo ad esaminare le $F^2$ di determinante $d=2$ (con un fascio lineare di curve ellittiche $C$), andiamo incontro a qualcosa d’inaspettato: la classificazione di tali superficie dipende, non soltanto dal genere $p$, sì anche da un nuovo carattere, che è il numero $s$ delle curve del fascio $|C|$ che si riducono a curve ellittiche doppie; d’altronde questo carattere si lascia esprimere per il bigenere $P$ mediante la formula
\[
P=2(p-1)+s-1. % [sic] +s-1 -> +s+1, the same relation is printed with +1 on page 263.
\]

Per classificare le superficie $F^2$, partiamo dall’osservazione che una superficie con un fascio lineare di curve ellittiche $C$, su cui esistono curve bisecanti le $C$, contiene un sistema lineare irriducibile del tipo $(\theta+hC)$ e quindi si lascia rappresentare (su una rigata razionale doppia e) sopra un piano doppio con curva di diramazione $D_{2n}$, d’un certo ordine $2n$, dotata d’un punto $O$ $(2n-4)$-plo. La $D_{2n}$ potrà avere, in generale, altre singolarità elementari abbassanti il genere $p$: punti 4-pli o punti $[3,3]$. Ma se si vuole che l’ordine della $D_{2n}$ non possa abbassarsi per trasformazioni quadratiche del piano, si avranno i seguenti \textit{tipi irriducibili}:

$\mathrm{I}'$) $D_{2n}$ dotata di un punto $(2n-4)$-plo senza altre singolarità abbassanti il genere:
\[
p=n-2;
\]

$\mathrm{II}'$) $D_{2n}$ dotata d’un punto $(2n-4)$-plo e di un punto 4-plo
\[
p=n-3;
\]

III) $D_{2n}$ dotata d’un punto $(2n-4)$-plo $O$ e di $s$ $(\geq 1)$ punti $[3,3]$: $A_1A_2\ldotp\ldotp\ldotp\ldotp A_s$. La $D_{2n}$ contiene $s$ rette per $O$ e una rimanente
%CONT>>
% scan: s144R
% !PRINTED 263
\weberpage{263}
% head: CLASSIFICAZIONE DELLE SUPERFICIE DI GENERE LINEARE ECC.
%<<CONT
% stamp: Dark left scan edge and scattered specks.
parte $D'_{2n-s}$ dotata di $s$ tacnodi per cui le nominate rette sono tangenti tacnodali:
\[
p=n-2-s.
\]

\textit{Le superficie $F$ di genere lineare $p^{(1)}=1$ e di determinante $d=2$} ($P>p\geq 1$) danno luogo, per ogni valore del genere $p$, \textit{non solo} ai \textit{piani doppi dei tipi} $\mathrm{I}'$ e $\mathrm{II}'$ \textit{per cui}
\[
p=n-2,\qquad n-3,
\]
e
\[
P=2(p-1), % [sic] P=2(p-1) -> P=2(p-1)+1, an apparent omitted +1 when compared with the plurigenus formulas.
\]
sì anche ad \textit{infinite famiglie di piani doppi del tipo III, dipendenti da un altro intero}
\[
\begin{aligned}
s&=1,\,2,\,3\ \ldotp\ldotp\ldotp\ldotp: \\
p&=n-2-s,
\end{aligned}
\]
le quali \textit{si lasciano distinguere dando}, accanto al genere $p$, \textit{il bigenere}
\[
P=2(p-1)+s+1.
\]

Infatti, su un piano doppio del tipo III, le curve canoniche verranno rappresentate da gruppi di $p-1=n-3$ rette variabili per $O$ (rette doppie immagini di curve ellittiche irriducibili $C$), cui van sommate le $s$ rette $OA_1\ldotp\ldotp\ldotp\ldotp OA_s$ (immagini di curve fisse $\dfrac{C}{2}$), o meglio gli intorni dei punti $[3,3]$: $A_1\ldotp\ldotp\ldotp\ldotp A_s$ (cfr. Cap. III, \S\ 8, 11 e Cap. V, \S\ 7); e le curve bicanoniche verranno date dai gruppi di $2(p-1)+s$ rette variabili per $O$. Aggiungasi che per un piano doppio del tipo III \textit{i successivi plurigeneri} sono espressi da % [sic] p-1=n-3 -> p-1=n-3-s, inconsistent with p=n-2-s printed above.
\[
P_{2i}=2i(p-1)+si+1,
\]
e
\[
P_{2i+1}=(2i+1)(p-1)+si+1.
\]

Si può anche valutare il numero $M$ dei \textit{moduli} da cui dipende una classe di piani doppi del tipo \textit{III}. Invero la costruzione di una $D'_{2n-s}$ con punto-$(2n-s-4)$-plo $O$ assegnato e dati tacnodi $A_1\ldotp\ldotp\ldotp\ldotp A_s$, con rette tacnodali $OA_1\ldotp\ldotp\ldotp\ldotp OA_s$, dipende da $5(2n-s)-6-6s$ parametri, % ?? uncertain: stroke between punto and (2n-s-4) -- printed hyphen or scan mark.
mentre la determinazione del gruppo degli $s+1$ punti $O$, $A_1$, $\ldotp\ldotp\ldotp\ldotp$, $A_s$ mette in evidenza $2s-3$ invarianti proiettivi. % [sic] 2s-3 -> 2(s-3), the contribution is written as 2(s-3) in the following total.
Dunque le nostre $D'_{2n-s}$ e le $D_{2n}$ dipenderanno da $5(2n-s)-6-6s+2(s-3)=10n-9s-12$ invarianti o moduli. E poichè % Both plus signs are printed, one at the end and one at the start of consecutive lines.
\[
p=n-2-s,
\]
avremo
\[
M=10p+s+8,
\]
ovvero, introducendo il bigenere $P$:
\[
M=8p+P+9.
\]

% scan: s145L
% !PRINTED 264
\weberpage{264}
% head: CAPITOLO SETTIMO
% stamp: Dark right scan edge, scattered specks and stray marginal strokes.
Per $p=0$, $P=1$, curva bicanonica d’ordine zero, si ritrovano i 10 moduli del piano doppio tipico già incontrato con $D_8$ di diramazione.

I risultati conseguiti in ordine alla classificazione delle superficie di genere lineare $p^{(1)}=1$ e di determinante 2, conducono induttivamente a enunciare, in generale, per le superficie $F^d$ di determinante $d$ qualunque, il seguente teorema:

\textit{Nella classificazione delle superficie regolari di genere lineare $p^{(1)}=1$ ($P>1$) e di determinante $d>1$ (con un fascio lineare di curve ellittiche $C$ che possono supporsi d’ordine $d$) si ha da tener conto, non soltanto del genere $p$, sì anche di altri caratteri interi in corrispondenza ai divisori $r$ di $d$ ($1<r\leq d$), ciascuno dei quali caratteri, $s_r$, indica il numero delle curve ellittiche $\dfrac{C}{r}$, che contate $r$ volte appartengono al fascio $|C|$.}

\textit{Le varie famiglie di superficie $F^d$}, che rispondono a questi caratteri, \textit{si lasciano anche distinguere coi valori dei plurigeneri}: in corrispondenza a ciascuna delle $s_r$ curve $\dfrac{C}{r}$ il plurigenere d’ordine $i$, $P_i$ cresce (su $i(p-1)+1$) di $\left[\dfrac{i(r-1)}{r}\right]$, massimo intero contenuto in $\dfrac{i(r-1)}{r}$, e quindi lo $r$-genere $P_r$ cresce di $r-1$, e, in generale, $P_{ir}$ di $i(r-1)$ ecc.

Per giustificare l’enunciato si osservi che, se sopra una $F^d$, al fascio delle $C$ appartiene una curva ellittica $\dfrac{C}{r}$ che --- presa come $r$-pla --- costituisce una $C$, questa non influisce sul numero delle $C$ dotate di punto doppio, ai fini del computo dell’invariante di Zeuthen-Segre, e perciò neppure sul genere numerico $p_a$ della superficie $F^d$, nè sul suo genere geometrico $p_g=p_a$; la $\dfrac{C}{r}$, contata $r-1$ volte, viene ad aggiungersi come componente fissa alle parti variabili delle curve canoniche $K$ di $F^d$ (composte colle $C$). Infatti si costruisca su $F^d$ una rete di curve $|L+C|$, contenente il fascio di curve spezzate $L+|C|$; è facile vedere che la jacobiana di questa rete contiene $r-1$ volte come parte fissa ogni $\dfrac{C}{r}$, e da ciò segue che la $\dfrac{(r-1)C}{r}$ figurerà, come componente fissa, nelle curve canoniche $K$.

Il problema della classificazione delle superficie di genere lineare $p^{(1)}=1$ si può ritenere sostanzialmente risoluto dal teorema precedente. Ma a questo conviene aggiungere una relazione interessante che lega le nostre superficie $F^d$, di determinante $d>1$, alle $F^1$ di
%CONT>>
% scan: s145R
% !PRINTED 265
\weberpage{265}
% head: CLASSIFICAZIONE DELLE SUPERFICIE DI GENERE LINEARE ECC.
%<<CONT
% stamp: Dark left scan edge, scattered specks and stray marginal marks.
determinante 1, ed offre un mezzo di studio delle $F^d$. Si consideri, sopra una $F^d$, il relativo fascio di curve ellittiche $C$ che può supporsi d’ordine $d$. Ad una $C$ risponde in generale una varietà rappresentativa delle $g_d^{d-1}$ che ad essa appartengono; e se si ritengono queste $g_d^{d-1}$ come elementi « punti » di una curva, si ha una curva ellittica $c$ birazionalmente identica alla $C$, sopra cui è dato razionalmente \textit{un} punto, in corrispondenza alla serie segata sulla $C$ dai piani o dagli iperpiani dello spazio che la contiene. Così alla superficie $F^d$, luogo delle curve del fascio $|C|$, si può associare una $F^1$, luogo delle curve $c$, che sono birazionalmente identiche alle $C$ omologhe. Questa asserzione sembrerebbe, a prima vista, contraddire al fatto che la $F^1$ e la $F^d$ sono superficie birazionalmente distinte; ma conviene notare che fra una $C$ e la $c$ omologa vien data non \textit{una} trasformazione birazionale, bensì un gruppo di $d$ trasformazioni, che non sono razionalmente distinguibili in funzione del parametro della $C$ stessa, e così fra $F^1$ e $F^d$ si avrà una corrispondenza algebrica $[d,d]$. Invero, preso su $c$ un punto $P$, si determina sulla $C$ una serie $g_d^{d-1}$, entro la quale si può fissare razionalmente un gruppo $G_d$ di $d$ punti. Ora se si associa a $P$ un punto $P'$, fra i $d$ del $G_d$, ($G_d=P'+G_{d-1}$) resta definita una corrispondenza biunivoca fra $C$ e $c$, dove ad un punto $A$ di $C$ risponde su $c$ un punto $A'$ per cui il gruppo $A'+G_{d-1}$ appartiene alla serie $g_d^{d-1}$ omologa ad $A$.

Diremo pertanto:

\textit{Ad ogni superficie $F^d$ di genere lineare $p^{(1)}=1$ e di determinante $d>1$ ($P>1$), contenente un fascio lineare di curve ellittiche $C$, si può associare una superficie $F^1$ di determinante $1$, con un fascio di curve birazionalmente identiche alle $C$, la quale si trova con $F^d$ in corrispondenza $[d,d]$.}

Alle curve $C$ di $F^d$ dotate di punto doppio risponderanno in particolare curve $c$ dotate di punto doppio su $F^1$, quindi \textit{le due superficie} $F^1$ e $F^d$ avranno lo stesso invariante di Zeuthen-Segre e lo \textit{stesso genere} $p_a$.

\textit{Notizia.} -- In questo paragrafo si sono esposti, chiariti, in qualche punto corretti e completati, i risultati sulla classificazione delle superficie di genere lineare $p^{(1)}=1$, ottenuti da \textsc{Enriques}, specialmente nelle Note degli anni 1898, 1914: « \textit{Sui piani doppi di genere lineare $p^{(1)}=1$} ». Rendic. Acc. Lincei, $1^\circ$ sem. 1898, « \textit{Sulla classificazione delle [superficie algebriche e particolarmente sulle superficie di genere lineare $p^{(1)}=1$} ». Rendic. Acc. Lincei, febbr. 1914. % ?? uncertain: mark before superficie in the second quoted title -- opening square bracket, opening parenthesis, or scan blemish.

Le cose dette si estendono facilmente alle superficie irregolari ($p_g>p_a$), tostochè sia stabilita la proprietà caratteristica (relativa alle serie continue non lineari di curve) che le concerne (Cap. IX, \S\ 2 e 3).

Precisamente, in una Nota « \textit{Intorno alle superficie algebriche di}
%CONT>>
% scan: s146L
% !PRINTED 266
\weberpage{266}
% head: CAPITOLO SETTIMO
%<<CONT
% stamp: Scattered scan specks and a dark right-edge scan artefact.
\textit{genere lineare \(p^{(1)}=1\)} »\origfootnote{1}{s146L-1}{R. Acc. di Bologna, Dic. 1906.}, \textsc{Enriques} ha stabilito che: « \textit{le superficie irregolari di genere lineare \(p^{(1)}=1\) e di genere numerico \(p_a\geq0\) \((p_g>p_a)\) contengono un fascio irrazionale, di genere \(p_g-p_a\), di curve ellittiche} ». Tali superficie si lasciano classificare nello stesso modo che le regolari, aggiungendosi per esse il carattere \(p_g-p_a\).

Fuori di queste rimangono soltanto superficie di genere lineare \(p^{(1)}=1\) irregolari, che hanno il genere numerico \(p_a=-1\), e si lasceranno classificare completamente come «superficie ellittiche ed iperellittiche, con \(\infty^1\) e \(\infty^2\) trasformazioni birazionali in sè stesse » (Cap. XI).

\clearpage\thispagestyle{plain}
% scan: s146R
% !PRINTED none
\weberpage{[s146R]}
% head: none
% stamp: Scattered scan specks and a dark left-edge scan artefact.
\edchapter{VIII}{SUPERFICIE REGOLARI CANONICHE E PLURICANONICHE}{Capitolo VIII. SUPERFICIE REGOLARI CANONICHE E PLURICANONICHE}{SUPERFICIE REGOLARI CANONICHE E PLURICANONICHE}

\edsection{\textbf{1. Introduzione.}}{1. Introduzione.}{sec.VIII.1}

Le superficie di genere \(p_g>3\) posseggono un sistema canonico (puro) \(|K|\) di dimensione \(p_g-1\geq3\), e, mercè questo si lasciano trasformare di regola in \textit{superficie canoniche} dello spazio a \(p_g-1\) dimensioni \(S_{p_g-1}\), aventi per sezioni iperpiane le curve \(K\). In questo capitolo esamineremo la costruzione di alcune superficie canoniche, riferendoci, al caso delle \textit{superficie regolari}: \(p_g=p_a=p\). % ?? uncertain: the low mark between "e" and "mercè" -- comma (transcribed) or scan speck.

Le superficie canoniche di \(S_{p-1}\), aventi il genere lineare (assoluto) \(p^{(1)}\) \((>1)\), saranno in generale d’ordine
\[
p^{(2)}=p^{(1)}-1
\]
e, nelle loro proprietà proiettive, rispecchieranno le proprietà invarianti dell’intera classe delle superficie trasformate.

Danno luogo tuttavia ad eccezione, in ordine alla costruzione delle superficie canoniche: anzitutto le superficie di \textit{genere lineare} \(p^{(1)}=1\), su cui le curve canoniche sono spezzate nelle curve ellittiche d’un fascio. È l’\textit{eccezione essenziale} che ci riporta alla famiglia di superficie caratterizzata dal possesso di un fascio di curve ellittiche, studiata nel precedente capitolo: qui vengono a mancare affatto, non solo la superficie canonica, sì anche le superficie pluricanoniche di cui diciamo in appresso, tutte le curve pluricanoniche essendo egualmente composte colle curve ellittiche del fascio.

Altre eccezioni o casi in cui il discorso generale subisce almeno qualche modificazione, sono:

1) le superficie per cui il genere \(p\leq3\);

2) quelle su cui il sistema canonico, pur essendo irriducibile \(\infty^3\) almeno \((p>3)\), appartiene ad una involuzione \(I_n\), d’ordine \(n>1\), sicchè la superficie canonica si riduce ad una superficie multipla (\(n\)-pla);

% scan: s147L
% !PRINTED 268
\weberpage{268}
% head: CAPITOLO OTTAVO
% stamp: Scattered scan specks, a dark right-edge artefact, and a diagonal mark across "cui sezioni".
3) le superficie, di genere lineare \(p^{(1)}>1\), per cui le \textit{curve canoniche} (pure) \(K\) siano comunque riducibili; come si è visto al Cap. IV, § 18 si danno esempi effettivi in cui le \(K\) contengono parti fisse (non eccezionali) di genere zero e \(>0\); resta dubbio che esistano superficie su cui le parti variabili delle \(K\) sieno composte colle curve irriducibili \(K\) d’un fascio.

Se sopra una superficie \(F\) il sistema canonico puro \(|K|\) possiede una componente fissa, eccezionale o no, ed ha il genere \(p>3\), avendo componenti variabili irriducibili, le \(\infty^{p-1}\) curve canoniche condurranno sempre ad una superficie canonica, ma questa sarà d’ordine eguale al grado effettivo del sistema canonico \(\bar{p}^{(2)}<p^{(2)}=p^{(1)}-1\). Se invece le parti variabili delle curve canoniche sono riducibili, non si avrà più una superficie canonica. Comunque, anche in questo caso come nei precedenti, per \(p^{(1)}>1\), si potrà ottenere un modello proiettivo della classe, ricorrendo alle \textit{superficie bicanoniche o tricanoniche}, ecc., le cui sezioni piane sono rispettivamente curve bicanoniche o tricanoniche. Ma di ciò più avanti.

Nei seguenti paragrafi ci riferiremo anzitutto a superficie regolari con curve canoniche (pure) irriducibili (senza punti base) di genere \(p>3\), formanti un sistema lineare semplice \(\infty^3\) almeno di genere \(p^{(1)}\) e grado \(p^{(1)}-1\). La prima domanda che si presenta è questa: esistono in un medesimo spazio \(S_{p-1}\), cioè per lo stesso valore del genere \(p\), superficie canoniche di diverso ordine, corrispondenti dunque a diversi valori del \(p^{(1)}\)? Ossia: i due caratteri \(p\) e \(p^{(1)}\) sono indipendenti fra loro? La risposta, data dal \textsc{Noether} fino dal 1870, è affermativa.

L’indipendenza di \(p^{(1)}\) e \(p\), si prova subito con effettivi esempi. Basta osservare che un punto triplo e un tacnodo, imposti alle superficie \(F\) di un dato ordine e con date singolarità, diminuiscono egualmente il \(p\) di una unità, ma diminuiscono il \(p^{(1)}\) rispettivamente di 3 e di 2.

\edsection{\textbf{2. Superficie canoniche di genere \(p=4\) e genere lineare \(p^{(1)}=6,5\).}}{2. Superficie canoniche di genere p=4 e genere lineare p\textasciicircum{}(1)=6,5.}{sec.VIII.2}

In questo paragrafo e nei primi che gli succedono, ci proponiamo di classificare le superficie regolari di genere superficiale \(p=4\) e di genere lineare \(p^{(1)}>1\), costruendo le relative superficie canoniche, appartenenti allo spazio ordinario \(S_3\). Supporremo \textit{di regola} che si tratti di superficie con \textit{curve canoniche pure \(K\) formanti un sistema semplice irriducibile} (privo di punti base sopra un modello senza curve eccezionali), salvo a dedicare qualche nota ai casi di riducibilità; e così cercheremo anzitutto le superficie canoniche semplici, che rispondono all’ipotesi d’un sistema \(|K|\) non appartenente ad una involuzione \(I_n\) d’ordine \(n>1\).

% scan: s147R
% !PRINTED 269
\weberpage{269}
% head: SUPERFICIE REGOLARI CANONICHE E PLURICANONICHE
% stamp: Scattered scan specks.
Per \(p=4\) e \(p^{(1)}>1\), una superficie canonica \(F\), supposta semplice, deve avere l’ordine
\[
p^{(2)}=p^{(1)}-1\geq5.
\]

Il tipo semplice corrispondente al più piccolo valore di \(p^{(1)}\) è dunque: \textit{la superficie generale del \(5^\circ\) ordine, \(F_5\), superficie canonica con}
\[
p=4,\qquad p^{(1)}=6.
\]

È chiaro che una superficie con \(p=4\) e \(p^{(1)}=6\), di cui le curve canoniche siano irriducibili può ridursi birazionalmente al modello proiettivo della \(F_5\). Come esempi citiamo le superficie indicate da Noether:

\textit{a}) La superficie del \(7^\circ\) ordine \(F_7\), con conica tripla e punto triplo fuori di essa. Con una trasformazione quadratica di prima specie la \(F_7\) si riconduce ad una particolare \(F_5\), sulla quale sono dati un punto \(O\) ed una conica; il punto \(O\) risponde alla retta eccezionale in cui il piano della conica tripla sega \(F_7\).

\textit{b}) La \(F_7\) con sestica doppia \(C_6\) di genere 3. Nella consueta rappresentazione piana di una superficie cubica \(F_3\) passante per \(C_6\), la \(C_6\) stessa appare avere per immagine una \(C'_6\) \((1^3\,2^2\,3^2\,4^2\,5^2\,6)\) il cui comportamento rispetto ai punti fondamentali \(1,2,3,4,5,6\), viene indicato dal simbolo precedente; quindi le intersezioni variabili residue di due \(F_3\) per \(C_6\) sono cubiche gobbe, rappresentate da \(C_3(2,3,4,5,6^2)\), formanti una rete omaloidica; si deduce che le stesse \(F_3\) formano il sistema trasformante di una trasformazione cremoniana dello spazio, che riduce la nostra \(F_7\) alla \(F_5\).

\textit{c}) La superficie \(F_7\) con quintica doppia \(C_5\) di genere 0. Anche qui le \(F_3\) per \(C_5\) formano un sistema omaloidico, che però ha come curva base la \(C_5\) stessa completata da una retta che ad essa si appoggia in 4 punti, che è retta eccezionale della \(F_7\); e mercè questo sistema trasformante di \(F_3\) la \(F_7\) si riduce cremonianamente alla \(F_5\).

Ora vogliamo dimostrare con tutto rigore e generalità che: \textit{ogni superficie \(F\) coi generi \(p=4\) e \(p^{(1)}=6\), tale che la parte variabile delle curve canoniche sia irriducibile, si lascia trasformare in una quintica canonica \(F_5\).}

Indichiamo con \(|k|\) la parte variabile del sistema canonico che per ipotesi non è composta con le curve di un fascio. Facendo l’immagine proiettiva di questo sistema si ottiene una superficie \(F_n\) il cui ordine \(n\) \((\leq5)\) eguaglia il grado effettivo di \(|k|\).

Se \(n=5\) la \(F_n\) è la quintica e risulta a posteriori
\[
|k|=|K|.
\]

Se è \(n<5\), non può aversi una superficie canonica semplice (che sarebbe di genere \(p\leq1\)) e neppure un piano multiplo (che por%
%CONT>>
% scan: s148L
% !PRINTED 270
\unskip\weberpage{270}%%
% head: CAPITOLO OTTAVO
%<<CONT
% stamp: Scattered scan specks, a dark right-edge artefact, and diagonal marks across the text.
terebbe \(p=3\)), ma soltanto resta possibile — a priori almeno — una quadrica \(F_2\) doppia: s’incontrerebbe di fatto questo caso se si avesse una componente fissa eccezionale delle curve canoniche pure, cioè un punto base di \(|K|\) sopra un modello di \(F\) privo di curve eccezionali. Le curve \(k\) essendo di genere \(\pi\leq p^{(1)}\) \((=6)\), la curva di diramazione sulla \(F_2\) sarà una \(C_{2\pi+2}\) d’ordine \(2\pi+2\), e si proverà anzitutto che non è certo
\[
\pi<5,
\]
e perciò sono possibili a priori soltanto i casi:
\[
\pi=6,5.
\]

Invero non può essere \(\pi<5\), ossia \(\pi\leq4\), giacchè le curve canoniche \(K\) dovrebbero segare le \(k\) di genere \(\pi\) in gruppi residui della serie caratteristica \(g_n^2=g_4^2\), e quindi in non più di 2 punti, mentre deve essere \(kK>kk=4\). Ma, se \(\pi=5\), si ha sulla \(F_2\) doppia una curva di diramazione d’ordine 12; e, se \(\pi=6\), una curva d’ordine 14. La prima ipotesi porta ad un piano doppio con curva di diramazione \(C_{12}\) dotata di due punti 6-pli, che ha come curve canoniche le coniche pei due punti singolari (completate dalla retta eccezionale che li congiunge) e quindi il genere lineare assoluto \(p^{(1)}=5\) anzichè \(p^{(1)}=6\). La seconda ipotesi conduce ad un piano doppio con \(C_{14}\) di diramazione dotata di un punto 8-plo \(O\) e d’un punto 6-plo \(O_1\), ovvero di due punti 7-pli infinitamente vicini \(OO_1\), che ha come curve canoniche le \(\infty^3\) cubiche per \(O^2\) e \(O_1\), (completate dalla retta eccezionale \(OO_1\)) e perciò è di genere \(p=5\) anzichè \(p=4\).

Abbiamo detto che \(p^{(1)}=6\) è il minimo valore che può assumere il genere lineare \(p^{(1)}\) per una superficie di genere \(p=4\), che ammetta come modello proiettivo una superficie canonica semplice. Ma non è più così se il sistema canonico \(|K|\) appartiene ad una involuzione \(I_2\) d’ordine 2, ammettendo quindi un punto base.

Si ha allora come \textit{superficie canonica} la \textit{quadrica doppia} con curva di diramazione d’ordine 12, che conduce, come si è detto, al \textit{piano doppio con curva di diramazione} \(C_{12}\) d’ordine 12 dotata di due punti 6-pli. Questo piano doppio con \(p=4\) e \(p^{(1)}=5\), risponde al \textit{minimo valore del \(p^{(1)}\) per \(p=4\)} (d’accordo colla diseguaglianza generale \(p^{(1)}\geq2p-3\), che stabiliremo più avanti). Ciò risulta dall’analisi precedente perchè un valore minore del \(p^{(1)}\) dovrebbe condurre ad una superficie canonica semplice d’ordine 3 o 2, che sarebbe di genere \(p=0\), ovvero ad un piano triplo che avrebbe il genere \(p=3\).

E dalla stessa analisi resta anche escluso che un’ipotetica riducibilità del sistema canonico puro (cioè la presenza di componenti fisse, eccezionali o no) porti ad altri tipi di superficie con \(p=4\) e \(p^{(1)}=5\).

% scan: s148R
% !PRINTED 271
\weberpage{271}
% head: SUPERFICIE REGOLARI CANONICHE E PLURICANONICHE
% stamp: Scattered scan specks.
Enunceremo dunque che: \textit{Ogni superficie con \(p=4\) di genere lineare minimo \(p^{(1)}=5\) si può trasformare nel piano doppio con \(C_{12}\) di diramazione, dotata di due punti sestupli.}

Come esempio si consideri la superficie \(F_6\) del \(6^\circ\) ordine con due rette doppie \(a\) e \(b\) sghembe fra loro: le quadriche aggiunte segano su una tale \(F_6\) le curve canoniche, \(K\), le quali appartengono all’involuzione segata dalle rette incidenti ad \(a\) e \(b\).

\edsection{\textbf{3. Superficie con \(p=4\) e \(p^{(1)}=7\).}}{3. Superficie con p=4 e p\textasciicircum{}(1)=7.}{sec.VIII.3}

Passiamo a costruire il modello canonico per le superficie di generi \(p=4\) e \(p^{(1)}=7\). Avremo, nelle ipotesi apparentemente più generali, superficie \(F_6\) d’ordine 6 dello spazio ordinario, le cui quadriche aggiunte debbono spezzarsi in un piano variabile e in un piano fisso, intersecante \(F_6\) secondo una cubica doppia.

Così \textit{per \(p=4\) e \(p^{(1)}=7\) si ha la superficie canonica del \(6^\circ\) ordine \(F_6\) con cubica piana doppia.}

Ogni superficie i cui generi valgono \(p=4\) e \(p^{(1)}=7\), con sistema canonico semplice, si lascia trasformare birazionalmente nella \(F_6\). Così, per esempio, le superficie del \(6^\circ\) ordine \(F_6\) con 6 punti tripli, le cui curve canoniche (intersezioni delle quadriche per i 6 punti) sono precisamente di genere 7 e grado 6.

\textit{Nota.} – Ma, se si vuole approfondire lo studio della famiglia delle superficie per cui \(p=4\) e \(p^{(1)}=7\), conviene esaminare quelle per cui il sistema canonico \(|K|\) (che si suppone irriducibile a prescindere da eventuali componenti fisse) non sia semplice, sicchè la superficie canonica si riduca ad una cubica \(F_3\) doppia, ovvero ad una quadrica \(F_2\), doppia o tripla. Il principale scopo di questo esame è di rispondere alla domanda « se tutte le superficie per cui \(p=4\) e \(p^{(1)}=7\) appartengono ad una medesima famiglia rappresentata dalla \(F_6\) canonica, ovvero se si abbiano più famiglie distinte di superficie coi detti caratteri ».

Per quanto concerne le cubiche \(F_3\) doppie è facile persuadersi che esse rientrano come caso particolare nella famiglia delle \(F_6\).

Anzitutto la \(F_3\) non può essere rigata, poichè non sarebbe normale e ne resulterebbe \(p=5\); quindi essa sarà a sezioni piane ellittiche e possiederà una curva di diramazione \(C_{12}\) d’ordine
\[
2p^{(1)}-2=12.
\]

Questa \(C_{12}\) deve essere intersezione completa della \(F_3\) con una superficie del \(4^\circ\) ordine, altrimenti si troverebbe su \(F_3\) un fascio di coniche secanti in 6 punti la \(C_{12}\), che sarebbero immagini di curve di genere 2, secanti le curve canoniche \(K\) in 2 punti, sicchè le \(K\)
%CONT>>
% scan: s149L
% !PRINTED 272
\weberpage{272}
% head: CAPITOLO OTTAVO
%<<CONT
% stamp: Scattered scan specks and a dark right-edge scan artefact.
stesse risulterebbero iperellittiche e le sezioni piane di \(F_3\) razionali, ciò che si è escluso.

Quindi la \(F_3\) doppia si lascerà rappresentare sul \textit{piano doppio con curva di diramazione \(C_{12}\) d’ordine 12 dotata di 6 punti 4-pli}, che — di fatto — ha i generi \(p=4\) e \(p^{(1)}=7\), e su cui le curve canoniche hanno per immagini le cubiche per i 6 punti contate due volte.

Dimostriamo che \textit{il detto piano doppio rientra come caso particolare nella famiglia delle \(F_6\) con cubica piana doppia}; considerando il fascio di superficie determinato dalla \(F_6\) e da una sua aggiunta \(F_3\) contata due volte:
\[
F_3^2+\lambda F_6=0,
\]
per \(\lambda=0\) si ottiene in questo fascio la superficie cubica doppia \(F_3^2\), con curva di diramazione segata da \(F_6=0\).

Per dimostrarlo è sufficiente riconoscere che una sezione piana \(f_6\) di \(F_6\), appartenente ad un fascio \(f_3^2+\lambda f_6=0\) si riduce per continuità (per \(\lambda=0\)) alla cubica \(f_3\) doppia coi punti di diramazione nei 12 punti base semplici del detto fascio. Ora — se supponiamo la \(f_6\) affatto generale (cioè senza punti doppi) — essa costituisce un inviluppo di classe 30 che, al limite per \(\lambda=0\), deve ridursi alla cubica inviluppo \(f_3=0\) contata due volte (che è di classe \(2\cdot6=12\)), aumentata dei 18 fasci che han per centri i punti di diramazione della detta cubica doppia; ed è evidente che i punti di contatto delle tangenti condotte da un qualunque punto \(P\) alla \(f_6\), per \(\lambda=0\), vanno a cadere nelle 18 intersezioni di \(f_6\) e \(f_3\)\origfootnote{1}{s149L-1}{Cfr. \textsc{Enriques-Chisini}, \textit{Lezioni}, vol. III.}. Se poi la \(f_6\) ha 3 punti doppi allineati, presi come base del fascio \(f_3^2+\lambda f_6=0\), i 3 fasci con quei centri contati 2 volte saranno da togliere dai 18 anzidetti, restando dunque un \(G_{12}\) di diramazione della \(f_3\) doppia.

Mentre le superficie canoniche costituite da cubiche doppie con \(C_{12}\) di diramazione rientrano come caso particolare nel tipo più generale delle \(F_6\) con curva piana doppia del \(3^\circ\) ordine, all’opposto quelle superficie canoniche che sono costituite da quadriche doppie formano una famiglia distinta.

Senza indugiarci in un esame più minuto, diciamo che vi è un tipo di superficie \(F\), con \(p=4\) e \(p^{(1)}=7\), su cui le curve canoniche sono iperellittiche con due punti base (presi fra i punti doppi della \(g_2^1\)), le quali conducono a quadriche doppie canoniche rappresentabili sopra il \textit{piano doppio con curva di diramazione \(C_{10}\) d’ordine 10, dotata di due punti \([3,3]\)}.

E questi piani doppi porgono effettivamente \textit{una seconda famiglia di superficie con \(p=4\) e \(p^{(1)}=7\), che non rientrano fra le \(F_6\) canoniche semplici}, perchè dipendono da 39 moduli (\(39=65
%CONT>>
% scan: s149R
% !PRINTED 273
\text{\weberpage{273}}
% head: SUPERFICIE REGOLARI CANONICHE E PLURICANONICHE
%<<CONT
% stamp: Scattered scan specks.
-18-8\)\edFlush{}), laddove abbiam visto che le \(F_6\) dipendono da 38 moduli.

Maggiori difficoltà presenta l’esame delle superficie canoniche, con \(p=4\) e \(p^{(1)}=7\), che si riducono a quadriche triple. Le superficie di questa famiglia vengono caratterizzate dalla proprietà che le curve canoniche pure \(K\) non hanno punti base e contengono una \(g_3^1\), sicchè \(|K|\) risulta appartenente ad un’involuzione razionale del \(3^\circ\) ordine \(I_3\). \textit{La quadrica canonica tripla \(Q_2^3\)}, studiata da \textsc{G. Pompilj}\origfootnote{1}{s149R-1}{\textsc{G. Pompilj}, \textit{Sulle superficie algebriche le cui curve canoniche posseggono una \(g_3^1\).} Rendiconti Istituto Lombardo, 1940.}, \textit{è un cono e possiede una curva di diramazione \(C_{18}\) d’ordine 18.}

Si riconosce che le \textit{quadriche triple canoniche \(Q_2^3\) non rientrano come casi particolari nella famiglia delle \(F_6\) semplici con cubica piana doppia, bensì costituiscono una terza famiglia di superficie con \(p=4\) e \(p^{(1)}=7\), distinta dalle precedenti.}

Per dimostrarlo basterà constatare che una curva piana \(C_6\) con 3 punti doppi in linea retta, sezione di \(F_6\), non può ridursi per continuità ad una conica tripla \(C_2^3\), sezione di \(Q_2^3\), dotata di 18 punti di diramazione. A tal uopo si osserverà anzitutto che il sistema \(\infty^{27}\) delle sestiche piane \(C_6\) non contiene entro di sè quello \(\infty^{29}\) delle coniche triple dello stesso genere con 24 punti di diramazione. Non si può escludere tuttavia che fra le \(\infty^{27}\) \(C_6\) compaiano particolari \(C_2^3\) con 24 punti critici; ma se si vuole trovare entro questa famiglia una \(C_2^3\) con 18 punti di diramazione distinti, che costituisca una \(C_6\) degenere con 3 punti doppi in linea retta, bisogna che 6 fra i 24 punti critici suddetti coincidano a coppie in 3 punti critici apparenti \textit{allineati}, e ciò è impossibile per una \textit{\(C_2^3\) irriducibile}.

\edsection{\textbf{4. Superficie con \(p=4\) e \(p^{(1)}=8\).}}{4. Superficie con p=4 e p\textasciicircum{}(1)=8.}{sec.VIII.4}

Passiamo al caso successivo: \(p=4,\ p^{(1)}=8\). La superficie canonica (semplice) \(F_7\) che risponde a questi caratteri deve possedere in generale una curva doppia \(C_7\), che, contata due volte, costituisce l’intersezione completa della \(F_7\) con la quadrica aggiunta \(\Phi_2\). La \(C_7\), per una formula già vista\origfootnote{2}{s149R-2}{Cfr. cap. V, § 3.}, possiede un punto triplo che risulta doppio per la \(\Phi_2\) aggiunta la quale pertanto è un cono. Questo è in accordo col fatto che la \(C_7\) è la curva di contatto della \(\Phi_2\) con una superficie d’ordine 7, e perciò appartiene al sistema metà del sistema lineare \(|C_{14}|\) costituito dalle intersezioni di \(\Phi_2\) con le superficie del \(7^\circ\) ordine dello spazio. Ma dalla rappresentazione piana della quadrica appare che il detto sistema \(|C_{14}|\) non è in generale dimezzabile; affinchè lo sia, bisogna che la quadrica sia
% Signature line at the foot of the page:
% EDITORIAL: segnatura originale: \textsc{Enriques F.} - \textit{Superficie algebriche.}\hfill 18
\printersignature{\textsc{Enriques F.} - \textit{Superficie algebriche.}\quad 18}
%CONT>>
% scan: s150L
% !PRINTED 274
\weberpage{274}
% head: CAPITOLO OTTAVO
%<<CONT
% stamp: Scattered scan specks, a dark right-edge artefact, and a diagonal mark near "la quale".
un cono, il cui sistema rappresentativo è costituito dalle coniche per due punti infinitamente vicini \(O,\ O_1\): le immagini delle \(C_{14}\) sono allora curve dello stesso ordine \(C'_{14}\) \((O^7O_1^7)\) e la metà del sistema \(|C'_{14}|\) è a priori un sistema \(C'_7(O^4O_1^3)\) o un sistema più particolare in esso contenuto. Ma le \(C'_7(O^4O_1^3)\) che non abbiano in \(O\) una molteplicità superiore a 4, sono immagini di curve \(C_7\) di \(\Phi_2\), curve di contatto della stessa con delle \(F_7\), che non possono essere doppie per una \(F_7\) irriducibile perchè le generatrici del cono \(\Phi_2\) avrebbero con tale \(F_7\) otto intersezioni. Bisogna dunque che la curva doppia di una delle nostre \(F_7\) appartenga al sistema lineare che nel piano rappresentativo di \(\Phi_2\) ha per immagine \(C_7\) \((O^5O_1^2)\) e quindi passi tre volte per il vertice del cono ed abbia il genere \(\varrho=4\). % [sic] C_7(O^5O_1^2) -> C'_7(O^5O_1^2), the preceding plane-image notation uses a prime.

Abbiamo dimostrato che la \(F_7\) canonica \((p=4,\ p^{(1)}=8)\) ove esista, possiede come curva doppia una \(C_7\) di genere \(\varrho=4\), dotata di punto triplo \(V\), la quale appartiene al sistema metà di quello segato su \(\Phi_2\) dalle superficie del \(7^\circ\) ordine che passano 3 volte per \(V\) (o si comportano come quelle che hanno ivi un punto triplo). Ora si tratta di riconoscere l’esistenza effettiva di tali \(F_7\). Questa si può inferire induttivamente dal computo dei moduli; se si assume la formula più espressiva del § 13, Cap. V:
\[
M\geq10p-2p^{(1)}+12=36,
\]
si è condotti a ritenere che ogni \(C_7\) (definita come innanzi sopra un cono quadrico di vertice \(V\) triplo per essa) sia curva doppia per \(\infty^{26}\) superficie canoniche \(F_7\): infatti le \(C_7\) sopra il cono quadrico sono \(\infty^{17}\) e così, rispetto alle \(\infty^7\) trasformazioni proiettive del cono danno luogo ad \(\infty^{10}\) curve proiettivamente distinte; d’altra parte invece le superficie \(F_7\) aventi una medesima curva doppia \(C_7\) sono \(\infty^{26}\), poichè staccando il cono si hanno \(\infty^{25}\) \(F_5\) residue, passanti per la \(C_7\).

La conclusione precedente è raggiunta soltanto in via induttiva, sia perchè la formula dei moduli adoperata dipende da un’ipotesi non rigorosamente accertata, sia perchè essa vale a stabilire l’estensione di una famiglia continua di superficie con dati \(p\) e \(p^{(1)}\), quando si presuppone data l’esistenza di qualche superficie della famiglia; così bisogna almeno riconoscere l’esistenza effettiva di una particolare superficie \(F_7\) (coi generi \(p=4,\ p^{(1)}=8\)) dotata di una \(C_7\) doppia: si arriva per esempio a costruire una simile \(F_7\) partendo da una sestica \(F_6\) con 5 punti tripli e un tacnodo, per cui è appunto \(p=4,\ p^{(1)}=8\). Ma più direttamente si può dimostrare che ogni \(C_7\) (appartenente alla serie considerata sopra un cono quadrico) è curva doppia di una superficie (canonica). \(F_7\), col metodo indicato % ?? uncertain: the dot after "(canonica)" -- printed full stop (transcribed) or scan speck.
%CONT>>
% scan: s150R
% !PRINTED 275
\weberpage{275}
% head: SUPERFICIE REGOLARI CANONICHE E PLURICANONICHE
%<<CONT
% stamp: Scattered scan specks and an extra short stroke below the dash before "essendovi".
da \textsc{A. Franchetta}\origfootnote{1}{s150R-1}{\textsc{A. Franchetta}, \textit{Su alcuni esempi di superficie canoniche}, Seminario Mat. di Roma, 1939.}, e ne risulta quindi confermato per le stesse \(F_7\) il computo dei moduli.

Il metodo a cui ricorriamo consiste in questo: si riconosce che una \(C_7\) è sempre curva doppia per una (ed anzi per \(\infty^6\)) superficie razionale del settimo ordine \(F'_7\) che possegga inoltre un punto quadruplo assegnato; combinando linearmente questa particolare \(F'_7\) con una \(F_7\) degenere passante doppiamente per \(C_7\), si ottengono le \(F_7\) canoniche richieste.

Per realizzare questo programma, si osservi anzitutto che esistono superficie razionali del \(7^\circ\) ordine \(F'_7\) dotate di una curva doppia del \(7^\circ\) ordine e di un punto quadruplo isolato: una superficie siffatta viene rappresentata sul piano dal sistema delle curve di undicesimo ordine \(C_{11}(A^5B_1^3\ldotp\ldotp\ldotp\ldotp B_9^3D_1\ldotp\ldotp\ldotp\ldotp D_8)\) dove i 18 punti \(A,\ B\) e \(D\) appartengono ad una curva fondamentale del settimo ordine \(K_7\) \((A^3B_1^2\ldotp\ldotp\ldotp\ldotp B_9^2D_1\ldotp\ldotp\ldotp\ldotp D_8)\); le \(C_{11}\) sono curve di genere 8 e la serie caratteristica del sistema \(\infty^3\) è autoresidua, essendo il sistema \(|C_{11}|\) equivalente alla somma della curva fondamentale \(K_7\) e del suo aggiunto; perciò appunto la superficie rappresentata è una \(F'_7\) a sezioni piane di genere 8, e quindi dotata di una curva doppia \(C_7\), la quale possiede inoltre un punto quadruplo corrispondente alla \(K_7\). % ?? uncertain: dot counts in the index ranges of C_{11} and K_7 -- four (transcribed) or three per range.
% EDITORIAL D0143 RISOLTO: The enlarged C_11 and K_7 expressions each visibly have four separate baseline dots between B_1 and B_9 and between D_1 and D_8.

Ora le \(F'_7\) così costruite dipendono da un certo numero di invarianti assoluti proiettivi che è facile valutare: questo numero è 19. Inoltre si può calcolare la dimensione del sistema delle \(F'_7\) che posseggono una stessa curva doppia \(C_7\) e uno stesso punto quadruplo \(O\) fuori di esso; questa dimensione supera di uno quella del sistema caratteristico, segato sopra una \(F'_7\) dalle altre \(F'_7\) dello stesso sistema, e perciò vale 6. % [sic] fuori di esso -> fuori di essa, apparent agreement error with "curva doppia".
Ciò posto, facendo variare il punto quadruplo, si avranno (al più) \(\infty^9\) \(F'_7\) con la medesima curva doppia \(C_7\), sicchè — essendovi \(\infty^{19}\) \(F'_7\) — si avranno (almeno) \(\infty^{10}\) curve \(C_7\) doppie per una \(F'_7\). Ma le \(C_7\) così ottenute sono proprio quelle curve del \(7^\circ\) ordine, e di genere 4, giacenti sopra un cono quadrico, e passanti triplamente per il vertice, che sopra abbiamo definito. Combinando linearmente una \(F'_7\) con le \(\infty^{25}\) \(F_7\) degeneri costituite dal cono quadrico contenente \(C_7\) e da una \(F_5\) per \(C_7\), si ottengono \(\infty^{26}\) superficie canoniche \(F_7\) che passano doppiamente per la \(C_7\) assegnata. Con ciò viene dimostrato che: \textit{esistono \(\infty^{36}\) superficie canoniche del \(7^\circ\) ordine \(F_7\), proiettivamente distinte, le cui curve doppie sono le \(\infty^{10}\) curve del \(7^\circ\) ordine e di genere 4 che si definiscono mercè la bisezione del sistema lineare segato sopra un cono quadrico dalle superficie del \(7^\circ\) ordine passanti triplamente per il vertice.}

\textit{Le curve \(C_7\) si possono caratterizzare come curve gobbe iperel}%
%CONT>>
% scan: s151L
% !PRINTED 276
\unskip\weberpage{276}%
% head: CAPITOLO OTTAVO
%<<CONT
% stamp: Dark band at the right scan edge and scattered marginal specks.
\textit{littiche di genere $4$, dotate di un punto triplo}: infatti le curve gobbe del $7^{\mathrm{o}}$ ordine di genere 4 iperellittiche sono $\infty^{26}$ e il possesso di un punto triplo costituisce per esse \textit{una} condizione, come appare a priori da ciò che le $g_7^3$ possedenti una terna neutra sopra una curva iperellittica di genere 4 sono definite dal contenere il doppio della $g_2^1$ e perciò dipendono da tre costanti arbitrarie. % ?? uncertain: The exponent of the first infinity symbol -- 26 or 28.
% EDITORIAL D0144 RISOLTO: The last digit has a lower closed bowl with an open upper stroke, agreeing with 6 rather than 8.

Abbiamo spiegato la costruzione generale delle superficie canoniche $F_7$, dotate di una curva doppia $C_7$: curva di genere 4 iperellittica, dotata di punto triplo. Ma se in luogo di chiedere la costruzione di una $F_7$ con $C_7$ doppia, si domandi quella di una $F_7$ canonica possedente qualche curva multipla di molteplicità superiore a 2, si presenta subito l’esempio semplicissimo della $F_7$ dotata di una retta tripla $a$ e di due coniche doppie in piani per essa. Quando siano assegnate la retta tripla e le due coniche doppie si hanno $\infty^{33} F_7$: infatti staccando i piani delle due coniche si diminuisce il numero dei parametri di 2, e resta la dimensione del sistema delle $F_5$ passanti semplicemente per una retta e per due coniche in piani per essa, che è 31.

Ora si osservi che la curva composta di una retta e di due coniche in piani per essa possiede un invariante assoluto, che è il birapporto dei 4 punti sulla retta. Si avranno quindi $\infty^{34}$ superficie canoniche $F_7$ proiettivamente distinte, che posseggono una retta tripla e due coniche doppie. Si affaccia la domanda: queste $F_7$ rientrano come superficie particolari nella famiglia delle $F_7$ canoniche con $C_7$ doppia, ovvero costituiscono una serie continua distinta di superficie, che dipende soltanto da 34 anzichè da 36 moduli? Se la formula che dà i moduli fosse stabilita senza eccezione nella sua forma più espressiva citata innanzi, la questione verrebbe risolta senz’altro a priori adottando la prima alternativa. Comunque si può riconoscere direttamente in qual guisa le $F_7$ con retta tripla e due coniche doppie sieno particolarizzazioni delle $F_7$ con $C_7$ doppia.

% The first sum below has a plus sign at both sides of the printed line break.
Si tratta di mostrare che le curve dello spazio ordinario $C_2 + C'_2 + 3a$, costituite di due coniche generiche e della retta d’intersezione dei loro piani, rientrano come caso particolare nella famiglia delle nostre $C_7$, e poichè le nominate $C_2 + C'_2 + 3a$ soddisfano alle tre condizioni di possedere un punto triplo e di essere iperellittiche (cioè di contenere una $g_2^1$), basta mostrare che codeste curve composte rientrano nella famiglia delle $C_7$ di genere 4. Ora le $C_7$ di genere 4 dello spazio ordinario si possono ottenere come proiezioni di curve $C_8$ dello spazio $S_4$ da un punto semplice $O$; e fra codeste $C_8$ si trovano delle curve spezzate in due quartiche $C_4$, $C'_4$ di $1^{\mathrm{a}}$ specie passanti doppiamente per $O$ ed aventi ivi lo stesso piano tangente $\alpha$, le quali da $O$ si proiettano in una curva composta della retta tripla $a$ traccia del piano $\alpha$ e delle due coniche $C_2$, $C'_2$ proie% ?? uncertain: The plane symbol after "tangente" and "piano" -- Greek alpha or italic a.
% EDITORIAL D0145 RISOLTO: Both plane glyphs have the loop and crossing/rightward tail of alpha; the neighboring italic a denoting the line has a different shape.
%CONT>>
% scan: s151R
% !PRINTED 277
\unskip\weberpage{277}%
% head: SUPERFICIE REGOLARI CANONICHE E PLURICANONICHE
%<<CONT
% stamp: Dark band at the left scan edge and scattered specks.
zioni di $C_4$, $C'_4$. Così appare che la curva $C_2 + C'_2 + 3a$ può ritenersi come una particolarizzazione delle nostre $C_7$ iperellittiche e dotate di un punto triplo. Si può vedere direttamente che le condizioni richieste per questa degenerazione della $C_7$, sono in numero di 9, e perciò tutte le curve composte $C_2 + C'_2 + 3a$ rientrano nella famiglia delle nostre $C_7$; come abbiamo già visto ciò porta che le $F_7$ con retta tripla e due coniche doppie figurino entro la famiglia delle generali superficie canoniche con $p = 4$ e $p^{(1)} = 8$, soddisfacendo a 2 condizioni particolari.

Le superficie canoniche $F_7$ che abbiamo imparato a costruire rispondono ai valori dei generi $p = 4$ e $p^{(1)} = 8$ ($p^{(1)}$ genere lineare \textit{assoluto}) quando le curve canoniche $K$ sopra una superficie trasformata $F$ senza curve eccezionali formino un sistema semplice irriducibile e privo di punti base. Se invece si suppone che il sistema $|K|$ abbia su $F$ un \textit{punto base} (semplice), questo sistema avrà il grado effettivo $p^{(2)} - 1 = p^{(1)} - 2 = 6$; e quindi la superficie canonica corrispondente si ridurrà ad una sestica $F_6$, sopra la quale il punto base del sistema canonico dovrà avere per immagine una retta eccezionale contata due volte: codesta retta fa parte semplicemente dal sistema canonico impuro e figura poi come componente fissa del sistema canonico puro. % [sic] dal -> del, the preposition appears erroneous after "fa parte".

Ora è facile costruire effettivamente una superficie del $6^{\mathrm{o}}$ ordine $F_6$ che possieda una conica doppia assegnata $C$ e nel piano $\gamma$ di questa un tacnodo $O$ il cui piano tacnodale $\alpha$, diverso da $\gamma$, sia parimente assegnato; si ottengono tutte le $F_6$ soddisfacenti alle condizioni indicate per combinazione lineare di due particolari $F_6$: la $F_6$ costituita da una quadrica per $C$ contata due volte, e dal piano $\alpha$ pure contato due volte, e una $F_6$ composta del piano $\gamma$ e di una $F_5$ che passi semplicemente per $C$ ed abbia in $O$ un punto biplanare, toccando con una delle sue falde il piano tacnodale assegnato $\alpha$.

La $F_6$ così costruita possiede di fatto una componente fissa delle curve canoniche, che è la retta $a$ lungo la quale il piano $\alpha$ tocca $F_6$ e che perciò figura due volte nelle intersezioni di $F_6$ colle $\infty^3$ quadriche aggiunte, costituite da un piano variabile e dal piano fisso $\alpha$. Designando con $\varrho = 0$ e $-\nu$, il genere e il grado di codesta retta $a$, il genere e il grado di $2a$ varranno rispettivamente: % ?? uncertain: The shape of the genus symbol -- \varrho or \rho.
% EDITORIAL D0146 RISOLTO: The genus glyph has a curved descender that turns left and back, matching varrho rather than the straight-descender rho form.
\[
2 \cdot 0 - \nu - 1 = -(\nu + 1) \qquad \text{e} \qquad -4\nu;
\]
si può quindi valutare il grado $\bar{p}^{(2)}$ e il genere (virtuale) $\bar{p}^{(1)}$ del sistema canonico impuro su $F_6$ (quest’ultimo è il suo genere relativo):
\[
\begin{aligned}
\bar{p}^{(2)} &= 6 - 4\nu + 4 = 10 - 4\nu\\
\bar{p}^{(1)} &= 8 - (\nu + 1) + 2 - 1 = 8 - \nu \qquad (p^{(1)} = 8);
\end{aligned}
\]
%CONT>>
% scan: s152L
% !PRINTED 278
\weberpage{278}
% head: CAPITOLO OTTAVO
%<<CONT
% stamp: Dark band at the right scan edge and scattered marginal specks.
e dalla relazione
\[
\bar{p}^{(2)} = \bar{p}^{(1)} - 1
\]
si ricava quindi
\[
\begin{aligned}
10 - 4\nu &= 7 - \nu,\\
\nu &= 1.
\end{aligned}
\]

Così resta dimostrato a posteriori che la retta $a$ è effettivamente retta eccezionale per $F_6$ ($-\nu = -1$), sicchè per la presenza di essa il genere lineare relativo
\[
\bar{p}^{(1)} = 7 = p^{(1)} - 1
\]
e il genere lineare assoluto $p^{(1)} = 8$; inoltre codesta retta, tolta una volta dalle curve canoniche impure, rimane ancora componente fissa delle curve canoniche pure, rispondendo dunque ad un punto base semplice di $|K|$ sopra una superficie trasformata $F$\origfootnote{1}{s152L-1}{Anche in altro modo si può verificare che la $F_6$ è una superficie di genere lineare assoluto $p^{(1)} = 8$ le cui curve canoniche pure possiedono come componente fissa la retta eccezionale. Basta invero trasformare quadraticamente la $F_6$ mediante le quadriche che passano per la conica $C$ e per un punto $A$ di essa; si ottiene così una superficie del $7^{\mathrm{o}}$ ordine con conica tripla, che possiede inoltre un punto doppio (di \textsc{Nöther}) $O$, con un tacnodo infinitamente vicino; codesto punto singolare (immagine di una curva ellittica fondamentale di grado $-1$) dà luogo ad un punto base semplice per le curve canoniche segate dalle quadriche, parti variabili delle superficie cubiche aggiunte.}.

Contiamo ora il numero dei parametri da cui dipendono le nostre $F_6$. Le $F_5$ passanti per la conica $C$ e aventi un punto biplanare $O$ con una falda tangente al piano $\alpha$ dipendono da 37 parametri; e perciò le $F_6$ con conica doppia, tacnodo e piano tacnodale assegnato saranno $\infty^{38}$.

Siccome le singolarità assegnate per le $F_6$ possono trasformarsi proiettivamente in altre singolarità comunque date nello spazio (non possedendo nessun invariante assoluto), ed anzi quelle singolarità sono trasformate in sè stesse da $\infty^3$ omologie, si vede che: \textit{le nostre $F_6$ superficie canoniche con $p = 4$ e $p^{(1)} = 8$, corrispondenti ad un sistema canonico dotato di punto base, dipendono da $35$ invarianti assoluti o moduli.}

Anche qui il numero dei moduli è minore di quello (36) che appartiene alle superficie canoniche $F_7$, sicchè si è indotti a chiedere se le dette $F_6$ costituiscano una particolarizzazione delle $F_7$, corrispondente alla degenerazione di questa in una $F_6$ e in un piano.

Ci limiteremo ad accennare in qual modo la cosa possa realizzarsi. Se alla nostra $F_6$ si somma il piano $\gamma$ si ottiene una superficie $F_7$ che possiede la conica tripla $C$ ed inoltre due rette doppie infinitamente vicine coincidenti con $a$; la stessa $F_7$ si può ritenere dotata
%CONT>>
% scan: s152R
% !PRINTED 279
\weberpage{279}
% head: SUPERFICIE REGOLARI CANONICHE E PLURICANONICHE
%<<CONT
% stamp: Scattered specks and faint vertical marks in the right margin.
soltanto della conica tripla $C$ e della retta doppia $a$, ritenendo come virtualmente inesistente la retta doppia infinitamente vicina ad $a$, che vogliamo considerare come linea di connessione fra la $F_6$ e il piano $\gamma$; così la $F_7$ composta avrà il genere (numerico) della $F_6$, che andrebbe aumentato di zero (genere del piano) e di zero (genere della linea di connessione). Si tratta di riconoscere che la $F_7$ degenere definita in tal guisa, tenuto conto della singolarità puntuale che la nostra costruzione viene ad assegnarle nel punto $O$, può considerarsi come caso particolare di una $F_7$ canonica con curva doppia $C_7$ d’ordine 7 e genere 4 (iperellittica e dotata di un punto triplo). A tale scopo gioverà considerare la detta $C_7$ come proiezione da un punto $P$ di una $C_8$ d’ordine 8 e genere 4 dello spazio $S_4$, e quindi far degenerare la $C_8$ in una curva $C_6$ di genere 3 appartenente ad un cono quadrico di vertice $P$ ed avente in $P$ un punto doppio, e in due rette $a$ ed $o$ incidenti fra loro, la prima delle quali sia una corda di $C_6$ e l’altra passi per $P$: la proiezione di questa $C_8$ degenere, da $P$ nello spazio ordinario, potrà ritenersi come curva limite di una $C_7$ variabile che venga a degenerare in una conica tripla $C$ e in una retta $a$ giacente nel piano di $C$; questa curva degenere va completata col sommarle una curva infinitesima costituita dall’intorno del punto $O$ traccia della retta $o$, in un piano per $a$. Si comprende in tal modo come anche il \textit{possesso del tacnodo della $F_6$ possa risultare come limite del possesso della curva doppia $C_7$ che definisce la famiglia delle $F_7$.}

\textit{Nota.} -- Per approfondire lo studio delle superficie con $p = 4$ e $p^{(1)} = 8$ occorre esaminare i casi in cui il sistema canonico $|K|$, irriducibile nelle sue componenti variabili, appartenga ad una involuzione, dando luogo ad una superficie canonica multipla. Si è condotti ad esaminare superficie canoniche costituite da cubiche $F_3$ doppie, ovvero da quadriche $F_2$ doppie o triple. Senza indugiarci su questo esame ci limiteremo a indicare alcune famiglie notevoli che in tal guisa si ottengono.

Anzitutto si trova una $F_3$ doppia che risponde al caso in cui le $K$, appartenenti ad un’involuzione razionale $I_2$ (considerate sopra una trasformata senza curve eccezionali), abbiano un punto base. La $F_3$, come nel precedente paragrafo, si riconosce a sezioni piane ellittiche (cioè non rigata) e la sua curva di diramazione risulta ora una curva d’ordine $2p^{(1)} - 2 = 14$. Rappresentando codesta $F_3$ sopra il piano si ottiene un \textit{piano doppio con curva di diramazione $C_{12}$ d’ordine $12$, dotata di $5$ punti $4$-pli e d’un punto $[3,3]$}. Effettivamente questo piano doppio ha come immagini (doppie) delle curve canoniche le cubiche per i 6 punti singolari di $C_{12}$ e perciò
\[
p = 4,\qquad p^{(1)} = 8.
\]

% scan: s153L
% !PRINTED 280
\weberpage{280}
% head: CAPITOLO OTTAVO
% stamp: Dark band at the right scan edge and scattered marginal specks.
Ora si può riconoscere che \textit{la famiglia delle superficie, coi generi $p = 4$ e $p^{(1)} = 8$, definita dai nominati piani doppi rientra in quella delle $F_6$ canoniche semplici}, che corrispondono al caso di un sistema canonico dotato di punto base. Infatti si considerino le cubiche aggiunte ad una delle dette $F_6$, cioè passanti per la sua conica doppia $C_2$ e per il suo tacnodo $O$ (nel piano di $C_2$); fra queste superficie se ne prenda una $F_3$ che tocchi in $O$ il piano tacnodale di $F_6$ e si costruisca il fascio
\[
F_3^2 + \lambda F_6 = 0.
\]

In questo fascio, per $\lambda = 0$, si vede la $F_6$ degenerare con continuità nella $F_3^2$, con curva di diramazione d’ordine 14, sezione delle $F_6$ fuori di $C_2$: della curva di diramazione viene a far parte la retta che, contata due volte, costituisce la sezione ulteriore di $F_6$ col piano di $C_2$.

Una seconda famiglia di superficie canoniche multiple è quella delle \textit{quadriche doppie} riferibili al tipo del \textit{piano doppio con curve di diramazione $C_{12} = r + C_{11}$ comprendente una retta eccezionale $r$ ed una $C_{11}$ d’ordine $11$ dotata di due punti $4$-pli infinitamente vicini $AA_1$ su $r$ e di due punti $[3,3]$ $B$ e $D$, uno dei quali $B$, appartenga ancora alla retta $r$} (essendo il punto triplo infinitamente vicino $B_1$ fuori di $r$). Le immagini doppie delle curve canoniche $K$ sono qui le coniche per $A$ e $D$, completate dalla retta eccezionale $AA_1D$; perciò
\[
p = 4 \qquad \text{e} \qquad p^{(1)} = 8,
\]
e sopra un modello della superficie senza curve eccezionali le curve canoniche avranno 3 punti base. La domanda se questi piani doppi rientrino nella famiglia definita dalle superficie canoniche $F_7$ o in quella delle $F_6$ (sistema canonico puro semplice con un punto base o una componente fissa eccezionale) sembra dovere ricevere risposta negativa, sebbene il numero dei moduli da cui codesti piani doppi dipendono sia 33, cioè inferiore a quello delle dette $F_7$ (che ne hanno 36) e delle $F_6$ (che ne hanno 35). Invero non si vede come una superficie composta di una quadrica doppia e di due piani possa possedere le singolarità caratteristiche della $F_6$.

Non c’indugeremo oltre sulle superficie canoniche multiple coi generi $p = 4$ e $p^{(1)} = 8$, lasciando il problema delle quadriche triple canoniche coi detti generi; diremo soltanto che una ricerca di \textsc{G. Pompilj} (che avremo occasione di citare nel seguito) riesce ad escludere l’esistenza di questo tipo.

Infine ci limiteremo a segnalare che, per una classificazione completa delle superficie con $p = 4$ e $p^{(1)} = 8$, converrebbe ancora esaminare le superficie canoniche multiple che potrebbero nascere in corrispondenza all’ipotesi di curve canoniche riducibili con una parte fissa, le cui componenti variabili sieno di genere $< p^{(1)}$.

% scan: s153R
% !PRINTED 281
\weberpage{281}
% head: SUPERFICIE REGOLARI CANONICHE E PLURICANONICHE
% stamp: Dark band at the left scan edge and scattered specks.
\edsection{5. Superficie con $p = 4$ e $p^{(1)} = 9$.}{5. Superficie con p = 4 e p\textasciicircum{}(1) = 9.}{sec.VIII.5}

Le superficie canoniche $F_8$ di genere $p = 4$ e genere lineare (assoluto) $p^{(1)} = 9$ con un sistema canonico irriducibile semplice, senza punti base, si lasciano costruire con lo stesso metodo adoperato innanzi per $p^{(1)} = 7,8$. La $F_8$ deve possedere, secondo le nostre formule
\[
t = 4
\]
punti tripli, ed una superficie aggiunta $\Phi_3$ avente, in questi, 4 punti doppi. Occorre considerare le $F_9$ che passano triplamente per i 4 punti doppi di $\Phi_3$ e dimezzare il sistema lineare delle curve da esse segate su $\Phi_3$.

Ricorrendo alla consueta rappresentazione della $\Phi_3$ sopra un piano, dove le sezioni piane hanno per immagini delle curve $C_3$ passanti per tre coppie di punti infinitamente vicini, $PP_1$, $QQ_1$ ed $RR_1$, sopra una conica, il sistema metà da noi ricercato vien dato dalle curve del $9^{\mathrm{o}}$ ordine,
\[
C_9(P^4P_1Q^4Q_1R^4R_1).
\]

Queste curve sono sulla $\Phi_3$ $\infty^{21}$ curve del $12^{\mathrm{o}}$ ordine $C_{12}$ di genere 10 e a priori costituiscono possibili curve doppie di superficie canoniche $F_8$. Ora osserviamo che una $C_{12}$ la quale sia curva doppia per una $F_8$ irriducibile sarà pure doppia per $\infty^{13} F_8$, giacchè staccando la $\Phi_3$ aggiunta restano $\infty^{12}$ superficie del $5^{\mathrm{o}}$ ordine $F_5$ passanti per la stessa $C_{12}$. Siccome le superficie cubiche con 4 punti doppi sono $\infty^{15}$, fra loro proiettivamente identiche, le nostre $C_{12}$ appartengono ad una famiglia che dipende da 21 invarianti assoluti. Se ammettiamo che ciascuna $C_{12}$ sia curva doppia per una (e quindi per $\infty^{13}$) $F_8$ irriducibile, si avranno $\infty^{34} F_8$ proiettivamente distinte, cioè il numero dei moduli delle nostre superficie canoniche verrà dato esattamente dalla formula
\[
M = 10p - 2p^{(1)} + 12 = 40 - 18 + 12 = 34.
\]

Le ragioni esposte nel Cap. V, \S\ 13, che tendono a stabilire l’anzidetta espressione di $M$ almeno come limite inferiore del numero dei moduli per le superficie di genere $p > 3$ (a sistema canonico semplice), inducono a ritenere che di fatto tutte le $C_{12}$ definite innanzi siano curve doppie per effettive $F_8$ irriducibili (ciascuna per $\infty^{13}$ di queste). Il risultato così induttivamente conseguito si può confermare col metodo di \textsc{Franchetta}, e così viene anche, almeno verosimilmente, confermato per questa classe di superficie il computo dei moduli, secondo la formula scritta innanzi. Anzitutto si riconosce che \textit{esistono effettivamente superficie canoniche irriducibili $F_8$ aventi}
%CONT>>
% scan: s154L
% !PRINTED 282
\weberpage{282}
% head: CAPITOLO OTTAVO
%<<CONT
% stamp: Dark right-edge band, dark patch at the lower left, and scattered specks; a diagonal non-print mark crosses M in the first display.
\textit{una $C_{12}$ doppia}, e ciò almeno quando questa $C_{12}$ sia curva doppia per una superficie razionale $F'_8$, dotata altresì di un punto quadruplo isolato: le $F_8$ che vengono associate in tal guisa ad una $F'_8$ razionale, si ottengono per combinazione lineare della $F'_8$ con le superficie composte della subaggiunta $\Phi_3$ e di una $F_5$ passante per la stessa $C_{12}$. Quanto alla esistenza effettiva delle $F'_8$ razionali, si dimostra ricorrendo al sistema rappresentativo di esse nel piano, che è un sistema lineare $\infty^3$ di curve di genere $p^{(1)} = 9$ con serie caratteristica autoresidua, quale può definirsi per esempio mediante le curve di undicesimo ordine $C_{11}\,(A_1^3\ldotp\ldotp\ldotp\ldotp A_{12}^3\ B_1\ldotp\ldotp\ldotp\ldotp B_5)$ dove i 17 punti $A$ e $B$ appartengono ad una curva fondamentale del settimo ordine $K_7\,(A_1^2\ldotp\ldotp\ldotp\ldotp A_{12}^2\ B_1\ldotp\ldotp\ldotp\ldotp B_5)$. % ?? uncertain: An oblique mark follows B in "17 punti A e B" -- a scan speck or a prime on B.

Dopo avere così dimostrato che esistono superficie canoniche irriducibili $F_8$, si è condotti altresì ad ammettere che il \textit{numero dei moduli} da cui esse dipendono è esattamente
\[
M = 10p - 2p^{(1)} + 12 = 34,
\]
e quindi che tutte le $C_{12}$ ottenute dal problema di bisezione spiegato innanzi sono effettivamente curve doppie di superficie canoniche $F_8$; questa affermazione si giustifica almeno ove si ammetta che l’imposizione di un punto quadruplo alle nostre $F_8$, porti per esse 17 condizioni e non meno; giacchè si ritrova in tal guisa il numero $(34 - 17 = 17)$ degli invarianti assoluti delle $F'_8$ razionali, quale si lascia valutare dalla rappresentazione piana.

\textit{Nota.} -- Anche per $p = 4$ e $p^{(1)} = 9$ s’incontrano superficie canoniche multiple, che possono dar luogo a diversi casi. Ci limiteremo a segnalare due casi. Anzitutto la $F_4$ doppia a sezioni piane di genere 2 (cioè dotata d’una retta doppia) con curva di diramazione d’ordine 12, la quale si lascia rappresentare sul \textit{piano doppio con curva di diramazione $C_{14}$ d’ordine $14$, dotata d’un punto $6$-plo $O$ e di $8$ punti $4$-pli, $A_1A_2\ldotp\ldotp\ldotp\ldotp A_8$}. Qui le curve canoniche hanno per immagini (doppie) le quartiche per $O^2A_1A_2\ldotp\ldotp\ldotp\ldotp A_8$, sicchè risulta appunto
\[
p = 4 \qquad \text{e} \qquad p^{(1)} = 9.
\]

È agevole riconoscere che il nominato piano doppio o la corrispondente $F_4^2$, rientrano, come \textit{caso particolare, nella famiglia delle superficie con $p = 4$ e $p^{(1)} = 9$} definita dalle $F_8$ semplici. Ciò risulta, con tutta sicurezza, dal computo dei moduli, perchè gli invarianti proiettivi della $C_{14}$ di diramazione del detto piano doppio sono in numero di
\[
119 - 19 - 64 - 8 = 28,
\]
e questo numero riesce inferiore al numero dei moduli della famiglia cui le $F_4^2$ appartengono, che è certo
\[
M \geq 9p - 2p^{(1)} + 12 = 30.
\]

% scan: s154R
% !PRINTED 283
\weberpage{283}
% head: SUPERFICIE REGOLARI CANONICHE E PLURICANONICHE.
% ?? uncertain: The dot following the running head -- a printed full stop or a scan speck.
% stamp: Scattered marginal specks.
Un secondo caso notevole in cui --- per $p = 4$ e $p^{(1)} = 9$ --- s’incontra una superficie canonica multipla è quello della cubica $F_3$ doppia, con curva di diramazione d’ordine 16, la quale si lascia rappresentare sul \textit{piano doppio con $C_{12}$ di diramazione dotata di $4$ punti $4$-pli, $A_1A_2A_3A_4$, e di $2$ punti $[3,3]$: $B_1$ e $B_2$}. Qui le immagini delle curve canoniche $K$ sono date dalle cubiche per $A_1A_2A_3A_4B_1B_2$; e si hanno due punti base di $|K|$ in corrispondenza a $B_1$ e $B_2$. I nostri piani doppi dipendono da 32 invarianti o moduli; non ci fermeremo ad investigare se essi rientrino o meno nella famiglia definita dalle $F_8$ semplici, cui appartengono $10p - 2p^{(1)} + 12 = 34$ moduli.

\edsection{6. Superficie con $p = 4$ e $p^{(1)} > 9$.}{6. Superficie con p = 4 e p\textasciicircum{}(1) > 9.}{sec.VIII.6}

Le difficoltà che già appariscono nella costruzione delle superficie canoniche $F_8$, con $p = 4$ e $p^{(1)} = 9$, vanno crescendo quando si passi ai valori superiori del $p^{(1)}$. Il metodo che abbiamo seguito porta anzitutto alla bisezione d’un sistema lineare sopra una superficie $\Phi_{n-5}$ ($n = p^{(1)} - 1$), problema che non si lascia più risolvere così semplicemente quando si tratta di superficie non razionali. E questo caso si presenta già per le $F_9$ con $p^{(1)} = 10$; poichè s’incontra qui una superficie aggiunta $\Phi_4$ d’ordine 4, con
\[
t = 10
\]
punti doppi, e si tratta di dimezzare il sistema lineare segato sulla $\Phi_4$ dalle superficie $F_9$ d’ordine 9 che passano triplamente per questi 10 punti doppi.

Il problema che così si presenta per $p^{(1)} = 10, 11,\ldotp\ldotp\ldotp\ldotp$ dà luogo alle seguenti osservazioni: % ?? uncertain: The number of dots after 11 -- four or three.
% EDITORIAL D0149 RISOLTO: At enlargement, four separate full-stop shapes follow the comma.

$1^{\mathrm{o}}$) Non si può in generale affermare che la bisezione del sistema segato sopra una $\Phi_{n-5}$ con $t$ punti doppi dalle $F_n$ passanti triplamente per essi sia possibile, e nemmeno che essa dia un sistema univocamente determinato.

$2^{\mathrm{o}}$) In ogni caso occorrerebbe ancora decidere se le ipotetiche curve $C$ del sistema metà anzidetto, tutte o in parte, siano effettivamente curve doppie per superficie $F_n$ irriducibili.

3) Dalle osservazioni $1^{\mathrm{a}}$) e $2^{\mathrm{a}}$) segue che le superficie canoniche (semplici) dello spazio ordinario ($p = 4$, $p^{(1)} = n + 1$), anzichè costituire un unico sistema continuo, potranno dar luogo per i valori di $n > 8$ a più famiglie distinte (forse possedenti lo stesso numero di moduli $M = 10p - 2p^{(1)} + 12 = 50 - 2n$). Questa cir-
% The continuation is not supplied; the following scan is printed page 286.
%CONT>>

\editiongap

% scan: s155L
% !PRINTED 286
\weberpage{286}
% head: CAPITOLO OTTAVO
% stamp: Dark band at the right scan edge and scattered marginal specks.
\edsection{8. Superficie con $p = 5$ e $p^{(1)} = 10$.}{8. Superficie con p = 5 e p\textasciicircum{}(1) = 10.}{sec.VIII.8}

Passando ora alle superficie $F$ con $p = 5$ e $p^{(1)} = 10$ (con sistema canonico irriducibile, semplice) daremo anzitutto un esempio in cui è facile costruire la relativa superficie canonica $F_9$, d’ordine 9 ($= p^{(2)} = p^{(1)} - 1$) nello spazio $S_4$. Invero si consideri in $S_3$ la $F = F_6$ del $6^{\mathrm{o}}$ ordine con 5 punti tripli $A_1A_2\ldotp\ldotp\ldotp\ldotp A_5$, per cui si ha appunto $p = 5$ e $p^{(1)} = 10$. Le quadriche dello $S_3$, passanti semplicemente pei 5 punti suddetti, sono le superficie $\Phi_2$ ($= \Phi_{n-4}$) aggiunte alla $F_6$ ($= F_n$), secanti su $F_6$ le curve canoniche. Ora il sistema di tali $\Phi_2$ rappresenta una varietà cubica $V_3^3$ dello spazio $S_4$, $V_3^3$ particolare (studiata da \textsc{Segre} e \textsc{Castelnuovo}) che è dotata di 10 punti doppi in corrispondenza ai 10 spigoli del pentaedro $A_1A_2\ldotp\ldotp\ldotp\ldotp A_5$. Su questa $V_3^3$ le varietà cubiche dello $S_4$ (che in generale non passano per alcuno dei punti doppi) segano superficie $F_9$ del $9^{\mathrm{o}}$ ordine, aventi per immagini in $S_3$ le $F_6$ che passano triplamente per $A_1A_2\ldotp\ldotp\ldotp\ldotp A_5$. Pertanto la $F_6$ di $S_3$ si lascia trasformare in una superficie canonica $F_9$ dello $S_4$, intersezione completa di due varietà cubiche di questo spazio, una delle quali è una $V_3^3$ dotata di 10 punti doppi.

Ma le $F_6$ con 5 punti tripli da cui siamo partiti costituiscono soltanto una famiglia particolare di superficie con $p = 5$ e $p^{(1)} = 10$, ed è chiaro che ogni altra superficie con tali caratteri darà luogo in generale ad una $F_9$ canonica dello $S_4$ intersezione completa di due varietà cubiche: imperocchè su tale $F_9$ si riconosce facilmente che le $\infty^{34}$ varietà cubiche $S_4$ segano un sistema lineare completo (aggiunto al sistema segato dalle quadriche) di dimensione 32. % [sic] varietà cubiche S_4 -> varietà cubiche di S_4, a preposition before the ambient-space symbol appears to be omitted.

D’altra parte tutte le $F_9$ di $S_4$, intersezioni complete di due varietà cubiche, formano un sistema continuo coi medesimi caratteri $p = 5$ e $p^{(1)} = 10$, sicchè si può dire che il \textit{tipo generale delle superficie semplici, corrispondenti ai generi $p = 5$ e $p^{(1)} = 10$, è la superficie $F_9$ di $S_4$ intersezione completa di due varietà cubiche}, la quale dipende da
\[
2 \cdot 34 - 2 - 24 = 42
\]
invarianti assoluti e perciò possiede $M$ moduli:
\[
M = 10p - 2p^{(1)} + 12 = 42.
\]

È interessante vedere come la nominata superficie canonica si lasci costruire col metodo spiegato nei precedenti paragrafi, a partire da una sua proiezione dello stesso ordine, che è una superficie canonica non normale $F'_9$ dello spazio $S_3$, dotata d’una curva doppia $C_{18}$ d’ordine 18, appartenente ad una superficie del $4^{\mathrm{o}}$ ordine $\Phi_4$. Perciò giova rilevare che la proiezione $F'_9$ di una $F_9$, fatta da un punto esterno $O$, deve possedere $t = 6$ punti tripli, $B_1B_2\ldotp\ldotp\ldotp\ldotp B_6$, i
%CONT>>
% scan: s155R
% !PRINTED 287
\weberpage{287}
% head: SUPERFICIE REGOLARI CANONICHE E PLURICANONICHE
%<<CONT
% stamp: Scattered specks and faint vertical marks in the right margin.
quali rispondono alle 6 rette per $O$ della varietà cubica che passa per $F_9$ e per $O$ stesso\origfootnote{1}{s155R-1}{È noto ed ovvio che sopra una $V_3^3$ di $S_4$, e per un punto $O$ di essa, vi sono 6 rette, le quali appartengono alla superficie sezione di $V_3^3$ collo spazio $S_3$ tangente in $O$, e perciò stanno sul cono quadrico osculatore a questa in $O$.}, e quindi appartengono ad una conica $C_2$\origfootnote{2}{s155R-2}{A priori si potrebbe credere che le superficie del $4^{\mathrm{o}}$ ordine con 6 punti doppi sopra una conica costituiscano un sistema contenuto in un altro sistema più ampio di $\Phi_4$ con 6 punti doppi, non più appartenenti ad una conica. In realtà le superficie del $4^{\mathrm{o}}$ ordine con 6 punti doppi si distribuiscono in più sistemi continui $\infty^{28}$, uno dei quali, non ampliabile, è costituito appunto dalle $\Phi_4$ con i 6 punti doppi sopra una conica. Che vi siano $\Phi_4$ con 6 punti doppi non appartenenti ad una conica si vede subito, fissando 5 punti doppi in posizione generica, $A_1$, $A_2,\ldotp\ldotp\ldotp,$ $A_5$, e imponendo quindi un punto doppio (in posizione non assegnata) alle $\Phi_4$ del sistema lineare di dimensioni $34 - 20 = 14$, così definito: si ottengono in tal guisa $\infty^{28}\Phi_4$ con 6 punti doppi, non appartenenti ad un piano. Ma, d’altra parte, se s’impone ad una $\Phi_4$ di possedere un piano tangente lungo una conica $C$, le $\infty^4$ sestiche segate sulla $\Phi_4$ dalle quadriche per $C$ (fra le quali si trovano $\infty^2$ quadriche spezzate nel piano di $C$ e in un altro piano variabile), incontreranno la $C$ in 6 punti fissi, che risultano doppi per $\Phi_4$; ora è facile vedere che l’imposizione di un piano tangente lungo una conica, porta appunto per $\Phi_4$ 6 condizioni (quali si esprimono esigendo che $\Phi_4$ contenga una conica autoresidua rispetto al sistema delle sezioni piane): si deduce quindi che le $\Phi_4$ con 6 punti doppi sopra una conica formano esse pure un sistema continuo $\infty^{28}$ di superficie del $4^{\mathrm{o}}$ ordine con 6 punti doppi. % ?? uncertain: The number of dots between A_2 and A_5 -- four or three.
% EDITORIAL D0150 RISOLTO: The enlarged footnote shows a comma immediately after A_2, then three dots, then a comma; the first comma must not be counted as a fourth dot.

Riteniamo che sopra una superficie $\Phi_4$, dotata di 6 punti doppi non appartenenti ad una conica, non sia generalmente possibile bisecare il sistema segato dalle superficie del $9^{\mathrm{o}}$ ordine che passano triplamente per i 6 punti: il problema andrebbe esaminato rappresentando la $\Phi_4$ sopra un piano doppio, per proiezione da uno dei suoi punti doppi.}.

Ora si vede che la $C_{18}$ doppia di $F'_9$ deve appartenere ad un sistema metà di quello segato sulla $\Phi_4$ aggiunta delle superficie di $9^{\mathrm{o}}$ ordine passanti triplamente per i 6 punti $B_1B_2\ldotp\ldotp\ldotp\ldotp B_6$, che sono doppi per la detta $\Phi_4$. La bisezione di questo sistema $|C_{36}|$ si ottiene facilmente osservando che le $C_{18}$ del richiesto sistema metà si lasciano definire come curve appartenenti al sistema somma del sistema $|C_{16}|$, segato su $\Phi_4$ dalle superficie generali del $4^{\mathrm{o}}$ ordine passanti per $B_1B_2\ldotp\ldotp\ldotp\ldotp B_6$, e della conica $C_2$ secondo la quale il piano di $C_2$ tocca $\Phi_4$. % [sic] delle -> dalle, the surfaces appear to be the agents of "segato".

La $\Phi_4$ è una superficie regolare di genere 1, con curva canonica d’ordine zero, e perciò ogni sistema lineare privo di punti base (in punti semplici) su di essa ha la dimensione eguale al genere $\pi$. Pertanto il sistema $|C_{16}|$, di grado $64 - 2 \cdot 6 = 52$, ha il genere eguale alla dimensione $26 + 1 = 27$. Sommando al $|C_{16}|$ la conica $C_2$ di genere 0, che sega le $C_{16}$ (fuori di $B_1B_2\ldotp\ldotp\ldotp\ldotp B_6$) in 2 punti, si ottengono $C_{18}$ di genere
\[
\pi = 27 + 0 + 2 - 1 = 28.
\]

% scan: s156L
% !PRINTED 288
\weberpage{288}
% head: CAPITOLO OTTAVO
% stamp: Dark scan-edge bands and scattered specks.
Quindi anche la dimensione del sistema $|C_{18}|$ vale 28, e supera di 1 quella del sistema $|C_{16}|$ cui si sommi la conica fissa $C_2$: ciò significa che il sistema $|C_{18}|$ è irriducibile. Ordunque abbiamo ottenuto sopra la $\Phi_4$ con 6 punti doppi $B_1B_2\ldotp\ldotp\ldotp\ldotp B_6$ appartenenti ad una conica, $\infty^{28}$ curve $C_{18}$, generalmente irriducibili, dotate di 6 punti tripli in $B_1B_2\ldotp\ldotp\ldotp\ldotp B_6$, che sono a priori « possibili » curve doppie per superficie canoniche non normali $F'_9$. Ammettiamo che, fra queste $C_{18}$, ve ne siano $\infty^x$ $(x\leq 28)$ per cui passino doppiamente effettive superficie $F'_9$ irriducibili. Ognuna di queste $C_{18}$ appartiene allora, come curva doppia, a $\infty^5 F'_9$ irriducibili, essendovi già una e quindi $\infty^4 F'_9$ spezzate nella $\Phi_4$ e in una $F_5$ per essa. Si avranno quindi $\infty^{x+5} F'_9$ con curva doppia $C_{18}$ appartenente ad una data $\Phi_4$ (dotata di 6 punti doppi sopra una conica). Ma vi sono $\infty^{14}\Phi_4$ cogli stessi 6 punti doppi e, mutando questi, si ottengono $\infty^{28}\Phi_4$, sicchè si avranno, nello spazio $S_3$, $\infty^{28+x+5}F'_9$; fra queste ve ne sono $\infty^{15}$ proiettivamente identiche ad una data, e $\infty^{19}$ birazionalmente identiche rispondenti a proiezioni di $F_9$ dallo $S_4$ da $\infty^4$ centri diversi; perciò il numero delle $F'_9$ e delle $F_9$, birazionalmente distinte, cioè il numero dei moduli delle dette superficie canoniche risulta:
\[
M=28+x+5-19=x+14,
\]
e poichè si sa d'altra parte che
\[
M=42,
\]
si deduce
\[
x=28.
\]

Vuol dire che « \textit{tutte le curve $C_{18}$ ottenute dalla bisezione indicata innanzi, sono curve doppie di superficie canoniche non normali $F'_9$, effettivamente esistenti e generalmente irriducibili} ».

Terminiamo proponendo allo studioso alcune

\textit{Domande.} -- Se fra le superficie a sistema canonico semplice, per $p=5$ e $p^{(1)}=10$, esistono o meno superficie su cui questo sistema possegga un punto base, conducenti perciò a superficie canoniche $F_8$ di $S_4$, o a $F'_8$ dello $S_3$, con curva doppia $C_{11}$ d'ordine 11. E se tali superficie rientrino come casi particolari nella famiglia delle $F_9$, o delle $F'_9$, con 42 moduli, che abbiamo descritta. Se, d'altra parte, rientrino come casi particolari nella famiglia delle $F_9$, o delle $F'_9$, superficie in cui il sistema canonico (possedente un punto base) è composto con un'involuzione razionale $I_2$, cioè $F_8$ che si riducono ad $F_4$ doppie, rappresentabili sul piano doppio con curva di diramazione $C_{12}$, d'ordine 12, dotata di 4 punti 4-pli e d'un punto $[3,3]$: si noti che questo piano doppio ha giusto i caratteri $p=5$ e $p^{(1)}=10$, e che dipende soltanto da 41 moduli.

% scan: s156R
% !PRINTED 289
\weberpage{289}
% head: SUPERFICIE REGOLARI CANONICHE E PLURICANONICHE
% stamp: Dark scan-edge marks and scattered specks.
\textit{Nota.} -- Passando al caso $p=5$ e $p^{(1)}=11$, si osservi ora il \textit{piano doppio con $C_{12}$ di diramazione dotata di 3 punti 4-pli e di 2 punti $[3,3]$, che ha i generi $p=5$ e $p^{(1)}=11$ e il sistema canonico (non semplice) con 2 punti base}: questo piano doppio, che conduce ad una quartica canonica doppia con $C_{18}$ di diramazione, dipende da 40 moduli. Se le superficie canoniche semplici $F_{10}$ di $S_4$ (con $p=5$ e $p^{(1)}=11$) dipendono egualmente da $M=10p-2p^{(1)}+12=40$ moduli, vuol dire che, come è da ritenere probabile, si avranno per $p=5$ e $p^{(1)}=11$ \textit{due famiglie distinte} di superficie: le $F_{10}$ e le $F_4$ doppie indicate. Qui giova avvertire che esistono certo $F_{10}$ canoniche semplici dello spazio $S_4$; si ottengono particolari $F_{10}$ (dotate di una retta doppia) partendo da una $F_6$ dello spazio ordinario con 4 punti tripli e 1 tacnodo. % ?? uncertain: the subscript on C after "quartica canonica doppia" -- 18 or 16.
% EDITORIAL D0151 RISOLTO: The final subscript has two closed lobes, supporting 8, and differs from the open upper form of the nearby 6 digits.

\edsection{9. Le superficie sopra una varietà a 3 o più dimensioni.}{9. Le superficie sopra una varietà a 3 o più dimensioni.}{sec.VIII.9}

Anzichè proseguire l'analisi dei precedenti paragrafi, preferiamo indicare come possano ottenersi notevoli esempi di superficie canoniche iperspaziali, con costruzione diretta, in rapporto alla geometria delle varietà a più di 2 dimensioni.

Quando si passa dalle superficie alle varietà a 3 o più dimensioni, s'incontrano tre ordini di proprietà:

$1^\circ$) proprietà di ordine generale, che si presentano come naturale estensione di quelle delle curve e delle superficie, e che quindi è lecito ritenere come note; alcune di queste sono già indicate da \textsc{M. Noether} nella sua memoria fondamentale dei Math. Ann. Bd. 8;

$2^\circ$) proprietà più o meno previste o prevedibili come estensione di quelle delle superficie, che tuttavia possono esigere dimostrazioni meno evidenti, e talvolta anche difficili;

$3^\circ$) proprietà affatto nuove ed inaspettate che non hanno alcun riscontro nella teoria delle superficie\origfootnote{1}{s156R-1}{Come esempio ricordiamo i risultati di \textsc{Fano} ed \textsc{Enriques} intorno alle questioni di razionalità; involuzioni dello spazio non razionali; diversi tipi di varietà contenenti una congruenza di curve razionali, ecc.

Fra i lavori sulla geometria delle varietà citiamo: \textsc{F. Severi}, \textit{Fondamenti per la geometria sulle varietà algebriche} (Rend. del Circ. Mat. di Palermo, vol. XXVIII, 1909). \textsc{A. Rosenblatt}, \textit{Varietà algebriche a tre e più dimensioni} (Atti del Congresso Internazionale dei Matematici, 1928) che è un rapporto completo dei risultati conseguiti in questo campo. In esso si trovano anche estese indicazioni bibliografiche. \textsc{B. Segre}, \textit{Quelques résultats nouveaux dans la géométrie sur une $V_3$ algébrique} (Mémoires de l'Académie R. de Belgique 2\textsuperscript{e} s., t. XIV, 1936).}.

% EDITORIAL: segnatura originale: \textsc{Enriques F.} - \textit{Superficie algebriche.}\hfill 19
\printersignature{\textsc{Enriques F.} - \textit{Superficie algebriche.}\quad 19}

% scan: s157L
% !PRINTED 290
\weberpage{290}
% head: CAPITOLO OTTAVO
% stamp: Dark scan-edge bands and scattered specks.
Prenderemo in considerazione soltanto il primo ordine di problemi concernenti la teoria generale dei sistemi lineari di superficie che appartengono ad una varietà a 3 dimensioni, i sistemi aggiunti, il sistema canonico, e gli invarianti a cui esso dà luogo.

Per semplicità di discorso ci riferiamo ad una varietà a 3 dimensioni $V$, priva di singolarità in un conveniente iperspazio, e consideriamo sopra di essa sistemi lineari di superficie privi di punti o linee base. Una superficie $F$ su $V$ appartiene in generale ad un sistema lineare completo $|F|$, e si può operare su tali sistemi --- come sui sistemi lineari di curve su una superficie --- per \textit{somma e sottrazione}.

Se il sistema $|F|$ ha la dimensione almeno eguale a tre, ogni sistema triplamente infinito contenuto in $F$ possiederà una \textit{superficie jacobiana} $F_j$, luogo dei punti doppi delle sue superficie; le $F_j$ relative ai vari sistemi contenuti totalmente in $|F|$ sono \textit{equivalenti} fra loro, e determinano un sistema lineare completo: il \textit{sistema jacobiano} di $|F|$, che indicheremo con $|F_j|$. % [sic] F after "contenuto in" -> |F|, apparently missing system bars in a sentence about subsystems.

Analogamente a ciò che accade per i sistemi di curve sopra una superficie sussiste il \textit{teorema fondamentale} espresso dalla relazione:
\[
|(F+\Phi)_j|=|F_j+4\Phi|=|4F+\Phi_j|,
\]
e la differenza
\[
|F_j-4F|
\]
definisce il \textit{sistema canonico} di $V$\origfootnote{1}{s157L-1}{È appena necessario avvertire che le superficie canoniche, così definite rispetto a trasformazioni birazionali sopra la varietà $V$, non si debbono confondere con le superficie \textit{proiettivamente canoniche}, che sono definite, in uno spazio $S_{p-1}$, dalla proprietà che le loro sezioni iperpiane sono curve canoniche (della stessa superficie).}, effettivo o virtuale, mentre il \textit{sistema aggiunto} ad $F$ è dato da
\[
|F'|=|F_j-3F|.
\]

Se la varietà $V_n$ d'ordine $n$ appartiene allo spazio a 4 dimensioni $S_4$, essa possiederà in generale una superficie doppia $\Gamma$, i cui punti rispondono alle \textit{coppie neutre} del sistema delle sezioni iperpiane sopra una varietà trasformata; e, salvo che si abbiano singolarità più complicate, il sistema canonico viene segato dalle \textit{varietà aggiunte} $\Phi_{n-5}$ d'ordine $n-5$, definite semplicemente dal passaggio per la superficie doppia $\Gamma$. Invece le aggiunte $\Phi_{n-4}$ segheranno sulla $V_n$ le superficie aggiunte alle sezioni iperpiane.

Quando la varietà $V_n$ di $S_4$ possegga singolarità più alte, cioè una superficie multipla d'ordine $i>2$ ovvero una curva $i$-pla $(i>2)$ o un punto $i$-plo isolato $(i>3)$, le varietà $\Phi$ \textit{aggiunte} a
%CONT>>
% scan: s157R
% !PRINTED 291
\weberpage{291}
% head: SUPERFICIE REGOLARI CANONICHE E PLURICANONICHE
%<<CONT
% stamp: Scan-edge marks and scattered specks.
$V_n$ saranno definite in generale dalle condizioni di passare $i-1$ volte per ogni superficie $i$-pla, $i-2$ volte per una curva $i$-pla e $i-3$ volte per un punto $i$-plo isolato.

Sopra una varietà $V$ a 3 dimensioni (appartenente o meno ad uno spazio $S_4$) le superficie $F'$ del sistema aggiunto ad un sistema $|F|$ segano una $F$ generica secondo curve canoniche; e, viceversa, questa proprietà caratterizza le superficie aggiunte $F'$, quando il sistema $|F|$ non abbia superficie fondamentali.

I caratteri del sistema canonico di una varietà porgono altrettanti \textit{caratteri invarianti} della $V$ stessa, già incontrati da \textsc{Noether}. Sono:

$1^\circ$) la dimensione del sistema canonico, aumentata di un'unità, cioè il numero delle superficie canoniche linearmente indipendenti, che è il genere \textit{geometrico principale} $P$ (o $P_g$) di $V$;

$2^\circ$) il genere geometrico $p_g$ e il genere lineare $p^{(1)}$ delle superficie canoniche che si definiscono rispettivamente come \textit{genere superficiale $p_g$ e genere lineare $p^{(1)}$} della varietà $V$.

Il genere $x$ della curva $K$ intersezione di due superficie canoniche, e il grado $y$ del sistema canonico (numero dei punti comuni a 3 superficie canoniche) non costituiscono dei nuovi caratteri di $V$, perchè si esprimono in funzione dei precedenti in base alle formule
\[
\begin{aligned}
2x+y-1&=p^{(1)}\\
4y&=p^{(1)}-1=p^{(2)},
\end{aligned}
\]
le quali traducono il fatto che $|2K|$ è il sistema canonico della superficie canonica. Si hanno dunque tre soli caratteri indipendenti, come già è stato notato da \textsc{Noether}.

Di fronte ai caratteri geometrici si possono definire per $V$ i caratteri numerici: il \textit{genere numerico principale} $P_a$ e il \textit{genere numerico superficiale} $p_a$\origfootnote{1}{s157R-1}{Cfr. \textsc{Severi}, l. c.}.

Dalle cose dette segue che, se si ha in $S_4$ una superficie $F$ intersezione completa di due varietà $V^m$ e $V^n$ degli ordini $m$ ed $n$, prive di singolarità sopra la $F$, il sistema canonico viene segato dalle varietà d'ordine $m+n-5$\origfootnote{2}{s157R-2}{Cfr. \textsc{Severi}, \textit{Su alcune questioni di postulazione.} Rend. Circ. Mat. di Palermo, vol. XVII, 1903.}; si ha qui l'estensione del noto teorema di Noether per le curve gobbe\origfootnote{3}{s157R-3}{Cfr. \textsc{Enriques-Chisini}, op. cit., Libro V, cap. V, \S~41 (vol. III, pag. 528).}.

Se la $F$ non è intersezione completa di $V^m$ e $V^n$, ma queste s'incontrano ulteriormente secondo un'altra superficie $F^\alpha$, allora sulla $F$ il sistema canonico viene segato dalle varietà d'ordine $m+n-5$ che passano per $F^\alpha$. % ?? uncertain: the superscript on both occurrences of F -- \alpha, a, or x.

% scan: s158L
% !PRINTED 292
\weberpage{292}
% head: CAPITOLO OTTAVO
% stamp: Dark scan-edge bands and scattered specks.
In modo analogo, se nello spazio ad $r$ dimensioni si hanno $r-2$ ipersuperficie (varietà ad $r-1$ dimensioni), degli ordini $n_1,n_2,\ldotp\ldotp\ldotp\ldotp,n_{r-2}$, prive di singolarità, che s'incontrino in una superficie composta $F+F^\alpha$, sulla $F$ le curve canoniche vengono segate dalle ipersuperficie d'ordine $n_1+n_2+\ldotp\ldotp\ldotp\ldotp+n_{r-2}-r-1$ passanti per $F^\alpha$. % ?? uncertain: the superscript on the two residual-surface symbols F -- \alpha, a, or x.

\edsection{10. Esempii di superficie canoniche iperspaziali.}{10. Esempii di superficie canoniche iperspaziali.}{sec.VIII.10}

Le nozioni precedenti conducono a semplici esempi di superficie proiettivamente canoniche negli iperspazi. Anzitutto si consideri nello spazio $S_4$ la superficie intersezione completa di due varietà $V^m$ e $V^n$ senza punti singolari. Questa sarà una superficie canonica se
\[
m+n-5=1,
\]
giacchè, per quanto si è detto, gli iperpiani segheranno su questa le curve canoniche. Qui si riottengono i due tipi già indicati innanzi: la superficie $F_8$ intersezione di una quadrica e di una varietà $V_3^4$, per cui $p=5$ e $p^{(1)}=9$; e la superficie $F_9$, intersezione di due varietà cubiche, per cui $p=5$ e $p^{(1)}=10$.

Ma, più in generale, saranno superficie canoniche dello spazio $S_r$, le intersezioni complete di $r-2$ ipersuperficie, degli ordini $n_1,n_2,\ldotp\ldotp\ldotp\ldotp,n_{r-2}$, quando questi ordini $(\geq 2)$ siano tali che si abbia
\[
n_1+n_2+\ldotp\ldotp\ldotp\ldotp+n_{r-2}-r-1=1,
\]
cioè
\[
n_1+n_2+\ldotp\ldotp\ldotp\ldotp+n_{r-2}=r+2.
\]

Questa relazione, tenuto conto che deve essere $n_i\geq 2$, fornisce alcune notevoli superficie canoniche, quali sono le seguenti.

Nello spazio $S_5$ si ha la superficie $F_{12}$ con $p=6$ e $p^{(1)}=13$, intersezione completa di due quadriche e di una varietà cubica. E nello spazio $S_6$ si ha la $F_{16}$ con $p=7$ e $p^{(1)}=17$, intersezione completa di 4 quadriche.

Questi esempi rientrano in famiglie più generali di superficie canoniche, quali si lasciano definire come intersezione di una varietà normale a 3 dimensioni di un iperspazio a curve sezioni ellittiche ovvero proiettivamente canoniche, con una ipersuperficie rispettivamente d'ordine 3 o 2.

Per le $V_3^n$ dello spazio $S_r$, con $n=2r-4$, a curve sezioni proiettivamente canoniche\origfootnote{1}{s158L-1}{Varietà di questo tipo sono i coni che proiettano una superficie con tutti i generi uguali ad 1 da un punto esterno allo spazio che la contiene. Per i primi valori dell'ordine questi coni sono contenuti in famiglie più ampie di varietà senza punti singolari. Comunque il ragionamento del testo sembra estendibile anche ai coni, con lievi modificazioni.}, che supponiamo prive di punti singolari,
%CONT>>
% scan: s158R
% !PRINTED 293
\weberpage{293}
% head: SUPERFICIE REGOLARI CANONICHE E PLURICANONICHE
%<<CONT
% stamp: Scan-edge marks and scattered specks.
si noterà che le loro sezioni iperpiane $F$ sono superficie di genere 1 con curva canonica d'ordine zero, e perciò il sistema $|F'|$ aggiunto ad $F$ è costituito dalla sola superficie d'ordine zero. Per conseguenza il sistema aggiunto a $|2F|$ è
\[
|(2F)'|=|F|;
\]
quindi le superficie $F$, ossia gli iperpiani dello spazio $S_r$ segano curve canoniche sopra le superficie di $|2F|$ ossia sulle intersezioni delle quadriche.

Quanto alle $V_3^n$ dello spazio $S_{n+1}$ $(n>3)$ a curve sezioni ellittiche\origfootnote{1}{s158R-1}{Le varietà a 3 dimensioni a curve sezioni ellittiche sono classificate nella nota di \textsc{F. Enriques}, \textit{Sui sistemi lineari di superficie algebriche le cui intersezioni variabili sono curve ellittiche.} Rendic. Acc. Lincei, maggio-giugno 1894. Cfr. Math. Annalen, Bd. 46.}, che sono razionali, avendo le sezioni normali, basterà notare che le intersezioni delle quadriche, e perciò le superficie del sistema $|2F|$, doppio del sistema $|F|$ delle sezioni iperpiane, sono superficie coi generi uguali ad 1, e con curva canonica d'ordine zero: infatti le superficie di $|2F|$ hanno come sezioni iperpiane curve d'ordine $2n$, di genere $n+1$, appartenenti a spazi ad $n$ dimensioni, e quindi curve canoniche. Segue di qui che il sistema $|2F|$ possiede come unica superficie aggiunta la superficie d'ordine zero, e perciò il sistema aggiunto a $|3F|$ è
\[
|(3F)'|=F, % [sic] F -> |F|, apparently missing system bars, unlike the preceding analogous display.
\]
sicchè le superficie $F$, ossia gli iperpiani dello spazio $S_{n+1}$ segano sulla intersezione di $V_3^n$ con una varietà cubica, le loro curve canoniche.

In conclusione, abbiamo ottenuto \textit{due serie di superficie canoniche, di genere $p$, nello spazio $S_{p-1}$}:

$1^\circ$) \textit{le superficie $F_{3p-6}$, d'ordine $3p-6$ $(r=p-1,n=p-1)$, con $p^{(1)}=3p-5$, intersezioni complete di una ipersuperficie cubica con una varietà razionale normale a curve sezioni ellittiche, $V_3^{p-2}$}: % [sic] n=p-1 -> n=p-2, inconsistent with the exponent p-2 printed at the end of the same item.

$2^\circ$) \textit{e le superficie $F_{4p-12}$, d'ordine $4p-12$, con $p^{(1)}=4p-11$ $(r=p-1,n=2p-6)$ intersezioni complete di una quadrica con una varietà $V_3^{2p-6}$, a curve sezioni proiettivamente canoniche}.

Vedremo più avanti altre classi di superficie canoniche definite sulla varietà a 3 dimensioni a curve sezioni razionali. Qui non ci indugeremo ulteriormente sulla costruzione delle superficie canoniche degli iperspazi, limitandoci a segnalare recenti lavori, che si riattaccano in qualche modo all'argomento, in cui questo esame viene proseguito, con la costruzione di altri esempi interessanti: \textsc{G. Fano}, « Superficie algebriche e varietà a 3 dimensioni a curve sezioni canoniche », Atti Acc. Lincei, 1936; « Su alcune varietà al%
%CONT>>
% scan: s159L
% !PRINTED 294
\unskip\weberpage{294}%
% head: CAPITOLO OTTAVO
%<<CONT
% stamp: Dark scan-edge bands and scattered specks.
gebriche a 3 dimensioni aventi curve sezioni canoniche » in « Scritti matematici offerti a L. Berzolari », pag. 329 (1936); « Sulle varietà algebriche a tre dimensioni a curve sezioni canoniche », Mem. Acc. d'Italia, 1937.

\textsc{L. Godeaux}, « Sur la construction d'une surface canonique », Bull. de la Soc. Math. de Liège 5, (1936); « Sur une surface canonique appartenant à la variété de Segre représentant les couples de points de deux plans », Bull. de l'Académie Belgique, 22, (1936); « Sur les surfaces algébriques dont les sections hyperplanes sont des courbes canoniques », I, II, III, Bull. de la Soc. Math. de Liège, 6, (1937), « Construction d'une surface canonique du septième ordre », Bull. de la Soc. des Sciences de Liège, 1944; « Construction d'une surface canonique du huitième ordre », Bull. des Sciences math., 1944; « Construction d'une surface canonique du neuvième ordre », Bull. de l'Acad. de Belgique, 1944.

\textsc{P. Bourniat}, « Note sur quelques surfaces canoniques et sur des surfaces de genre un », Bull. de l'Acad. de Belgique 5, 22, (1936); « Note sur quelques surfaces canoniques », Acad. de Belgique 3, 22, (1936); « Sur des hypersurfaces et variétés canoniques », Acad. de Belgique (5), 23, (1937).

\textsc{E. A. Maxwell}, « Regular canonical surfaces of genus three and four », Proc. Cambridge phil. Soc., 33, (1937.)

\edsection{11. Minimo valore del genere lineare rispetto al genere superficiale.}{11. Minimo valore del genere lineare rispetto al genere superficiale.}{sec.VIII.11}

Già dalla costruzione delle superficie canoniche appare chiaro che il genere lineare (assoluto) $p^{(1)}$ di una superficie non può scendere al di sotto di un certo valore minimo, in corrispondenza ad un dato genere superficiale $p$, (o $p_g$), almeno se le curve canoniche sono irriducibili, ciò che esclude, per $p>2$, le superficie con $p^{(1)}=1$. Ora ci proponiamo di dimostrare, senza eccezione alcuna, che \textit{per le superficie con $p^{(1)}>1$ il minimo valore del $p^{(1)}$ rispetto al $p$ o $p_g$ $(>2)$ è dato da}
\[
p^{(1)}=2p-3.
\]

In altre parole si ha
\[
p^{(1)}\geq 2p-3,
\]
o, ciò che è lo stesso,
\[
p\leq \frac{p^{(1)}+3}{2}:
\]
\textit{il massimo del $p$ rispetto al $p^{(1)}$ $(>1)$ è $\dfrac{p^{(1)}+3}{2}$.}

Questa diseguaglianza viene stabilita dal \textsc{Noether}, pel caso di curve canoniche (pure) irriducibili, nella sua memoria dei Math.
%CONT>>
% scan: s159R
% !PRINTED 295
\weberpage{295}
% head: SUPERFICIE REGOLARI CANONICHE E PLURICANONICHE
%<<CONT
% stamp: Dark scan-edge bands and scattered specks.
Ann. Bd. 8 (1875). Se si suppongono le curve canoniche $K$ irriducibili, essa è una conseguenza del teorema del Clifford per cui l'ordine di una serie speciale $g_n^r$ sopra una curva qualsiasi è $n\geq 2r$\origfootnote{1}{s159R-1}{Cfr. \textsc{Enriques-Chisini}, Op. c., Libro V, cap. I, \S~11 (vol. III, pag. 87).-}. Infatti per la serie caratteristica del sistema $|K|$ il grado vale $n=p^{(1)}-1$ e la dimensione $r=p-2$, sicchè % ?? uncertain: short mark following the first footnote's final period -- printed hyphen or scan speck.
\[
p^{(1)}-1\leq 2p-4. % [sic] \leq -> \geq, the preceding sentence states n \geq 2r.
\]

Ma la stessa diseguaglianza è valida comunque le $K$ canoniche si ammettano riducibili purchè la parte variabile sia spezzata in curve di genere $\pi\geq 3$. Anzitutto se le loro componenti variabili $k$ sono irriducibili il loro grado vale $n\leq p^{(1)}-1$ (Cap. IV, \S~18), mentre la dimensione di $|k|$ è quella stessa di $|K|$ cioè $p-1$; si ha quindi
\[
n\geq 2p-4
\]
e a fortiori
\[
\begin{aligned}
p^{(1)}-1&\geq 2p-4\\
p^{(1)}&\geq 2p-3.
\end{aligned}
\]
Se poi le $k$ variabili si suppongono composte con $p-1$ curve d'un fascio, di genere $\pi\geq 3$ il genere di $|k|$ $(\leq p^{(1)})$ sarà
\[
\geq (p-1)(\pi-1)+1;
\]
si deduce
\[
p^{(1)}\geq 2p-1.
\]
Ora \textit{se per una superficie il genere $p$ $(>2)$ raggiunge il massimo $p=\dfrac{p^{(1)}+3}{2}$ $(p^{(1)}>2)$, le curve canoniche sono irriducibili iperellittiche.}

Infatti, se le $K$ sono irriducibili, per la serie caratteristica $g_n^r$ su una $K$ $(n=p^{(1)}-1,r=p-2)$ deve aversi $n=2r$, e ciò porta appunto che la $K$ sia iperellittica\origfootnote{2}{s159R-2}{Cfr. \textsc{Enriques-Chisini}, Op. c., Libro V, \S~12, vol. III, pag. 97.}. Se poi le $K$ fossero riducibili, il ragionamento si ripeterebbe per le loro componenti variabili che sono certo, queste, irriducibili (altrimenti $p^{(1)}\geq 2p+1$): le $K$ risultano così iperellittiche e conducono ad una superficie canonica doppia, riferibile al piano doppio che indichiamo in appresso, per cui risulta a posteriori non esservi componenti fisse di $|K|$. % [sic] 2p+1 -> 2p-1, the preceding displayed bound for the composed case has a minus sign.

È notevole che, per le nostre superficie con $p$ massimo, risulti così escluso anche il caso che le curve canoniche pure $K$ contengano componenti fisse eccezionali, ossia $|K|$ possegga qualche punto base (sopra un modello privo di curve eccezionali): invero, se vi sono $r$ punti base risulta
\[
p^{(1)}\geq 2p-3+s. % [sic] s -> r, the base-point count is denoted r immediately above.
\]

% scan: s160L
% !PRINTED 296
\weberpage{296}
% head: CAPITOLO OTTAVO
% stamp: Dark scan-edge bands, scattered specks, and a small stray mark beside the first footnote.
La superficie canonica corrispondente ad un $p$ massimo $(p^{(1)}=2p-3)$ è una superficie normale a curve sezioni iperpiane razionali, d'ordine $p-2$, nello $S_{p-2}$ e perciò si lascia rappresentare sul piano mediante un sistema lineare di curve appartenente ad uno dei seguenti tipi\origfootnote{1}{s160L-1}{Cfr. \textsc{Enriques-Chisini}, Libro V, cap. II, \S~22, vol. III, pag. 194.}: % [sic] S_{p-2} -> S_{p-1}, the canonical space is denoted S_{p-1} in the preceding discussion.

I) rete delle rette,

II) sistema $\infty^5$ delle coniche,

III) sistema delle curve d'ordine $n$ con un punto base $O$ di molteplicità $n-1$, ed eventualmente anche un altro punto base semplice distinto da $O$, ovvero più punti base semplici, infinitamente vicini ad $O$. % ?? uncertain: small mark following the first O -- O or O' (the mark may be a scan blemish).
% EDITORIAL D0155 RISOLTO: The mark above-right of O is a fine cross-shaped vertical blemish with a separate lower speck, unlike the solid oblique primes in the surrounding type. It does not support adding a prime.

I piani doppi i cui sistemi canonici hanno per immagini tali sistemi di curve razionali sono rispettivamente:

I) il piano doppio con curva di diramazione $C_8$ dell'$8^\circ$ ordine: $p=3$, $p^{(1)}=3$;

II) il piano doppio con curva di diramazione $C_{10}$ del $10^\circ$ ordine; $p=7$, $p^{(1)}=9$; % [sic] p=7 -> p=6, inconsistent with p^{(1)}=2p-3 and the printed value p^{(1)}=9.

III) il piano doppio con curva di diramazione $C_{2n+6}$, d'ordine $2n+6$, avente un punto $2n$-plo $O$: $p=2n+1$, $p^{(1)}=4n-2$; e % [sic] p^{(1)}=4n-2 -> p^{(1)}=4n-1, inconsistent with the relation p^{(1)}=2p-3 stated on this page.

III$'$) casi particolari del tipo III) dove la curva di diramazione (che può supporsi d'ordine non abbassabile con trasformazioni quadratiche) possegga, oltre $O$, un punto 6-plo o 4-plo, essendo $p=2n-2$, $p^{(1)}=4n-8$ o rispettivamente $p=2n$, $p^{(1)}=4n-4$. % [sic] 4n-8 -> 4n-7 and 4n-4 -> 4n-3, both inconsistent with the relation p^{(1)}=2p-3 stated on this page.

Invero, come si è detto, sono da escludere quei tipi di piani doppi per cui la serie caratteristica del sistema canonico puro $|K|$ abbia qualche punto fisso (punto base di $|K|$): così p. es. un punto $[5,5]$ ovvero un punto $[3,3]$.

Enunciamo frattanto che: \textit{le superficie di genere superficiale $p>2$ con $p^{(1)}>1$, per cui il genere lineare assoluto $p^{(1)}$ raggiunge il minimo, $p^{(1)}=2p-3$, sono rappresentabili sopra un piano doppio appartenente ad uno dei tipi} I), II) III) e III$'$)\origfootnote{2}{s160L-2}{Cfr. \textsc{F. Enriques}, \textit{Sopra le superficie algebriche di cui le curve canoniche sono iperellittiche.} Rendic. Acc. Lincei, serie V, vol. V, pag. 191 (1896).}.

Alla diseguaglianza di Noether
\[
p\leq \frac{p^{(1)}+3}{2}
\]
si può sostituire un'altra diseguaglianza più espressiva che vale per le superficie con sistema canonico semplice $(p>3)$, e che è stata
%CONT>>
% scan: s160R
% !PRINTED 297
\weberpage{297}
% head: SUPERFICIE REGOLARI CANONICHE E PLURICANONICHE
%<<CONT
% stamp: Dark scan-edge bands, a dark patch at the upper right, and scattered specks.
messa in luce da \textsc{Castelnuovo}\origfootnote{1}{s160R-1}{\textsc{G. Castelnuovo}, \textit{Osservazioni intorno alla geometria sopra una superficie.} Nota II. (Rendic. Ist. Lomb., serie II, vol. XXIV, 1891). «Memorie Scelte», numero XVIII.}. \textit{Quando il sistema canonico è irriducibile e semplice, il genere superficiale ($p$ o $p_g>2$) non può superare il valore $\dfrac{p^{(1)}+6}{3}$}; si ha dunque
\[
p\leq \frac{p^{(1)}+6}{3},\qquad p^{(1)}\geq 3p-6.
\]

Per dimostrare questa diseguaglianza conviene partire da un lemma di Bertini concernente la serie $g_n^r$ autoresidue che possono appartenere ad una curva di genere $\pi$. % [sic] la serie -> le serie, the following adjective and verb are plural.
Sopra la curva che contiene una tale $g_n^r$ si assumano due gruppi $G$ e $G'$ della serie, che abbiano a comune un gruppo $\Gamma$ di $r-1$ punti generici; se la serie è semplice (cioè non composta con un'involuzione), $G$ e $G'$ non avranno altri punti in comune oltre agli $r-1$ punti di $\Gamma$, e il gruppo $\Gamma$ non potrà appartenere ad una serie lineare $g_{r-1}^s$, di dimensione $s>0$. Ciò posto, il teorema di Riemann-Roch dice. che il gruppo $G$ della $g_n^r$ presenta $n-r$ condizioni ai gruppi canonici che debbano contenerlo\origfootnote{2}{s160R-2}{Cfr. p. e. \textsc{Enriques-Chisini}, l. c., Libro V, cap. I, vol. III, pag. 85.}. % ?? uncertain: the mark after "dice" -- period, comma, or scan speck.
% EDITORIAL D0156 RISOLTO: A dark compact full stop is plainly visible at the baseline after dice, with no descending comma hook. Preserve the printed punctuation.
Inoltre il gruppo $G'$ presenterà $n-r-(r-1)=n-2r+1$ condizioni (a priori non indipendenti), ai gruppi canonici che ne contengono già $r-1$ punti, e perciò a quelli che siano già stati costretti a contenere il $G$. Quindi il numero delle condizioni perchè un gruppo canonico contenga $G+G'-\Gamma$ è minore o uguale di
\[
n-r+(n-2r+1)=2n-3r+1.
\]

Ora i gruppi canonici soddisfacenti a queste condizioni non possono essere più di uno, cioè coincidono col $G+G'$, per cui i punti di $\Gamma$ sono doppi; altrimenti il $\Gamma$ (che è costituito di $r-1$ punti generici della curva, con $r<\pi$), dovrebbe appartenere ad una serie $g_{r-1}^s$, con $s>0$. Si deduce che
\[
2n-3r+1\geq \pi-1,
\]
essendo $\pi-1$ il numero delle condizioni che valgono a determinare il $G+G'$ entro la serie canonica.

Ma poichè la $g_n^r$ è autoresidua,
\[
n=\pi-1,
\]
quindi
\[
\pi-3r\geq 0.
\]

% scan: s161L
% !PRINTED 298
\weberpage{298}
% head: CAPITOLO OTTAVO
% stamp: Scattered scan specks and marginal strokes; a diagonal stray stroke crosses the equality sign in p = 5.
Applichiamo questa formula alla serie caratteristica, semplice, del sistema canonico sopra una superficie di generi \(p\) e \(p^{(1)}\), ed avremo
\[
n=\pi-1=p^{(1)}-1;\qquad r=p-2,
\]
e quindi
\[
p^{(1)}-3p+6\geq 0,\qquad \text{c. d. d.}
\]

\textsc{Castelnuovo}, dopo avere scritta questa diseguaglianza, che deduce appunto dalla formula di Bertini, porge effettivi esempi di superficie con sistema canonico irriducibile semplice, per cui il genere \(p\) raggiunge il massimo rispetto al \(p^{(1)}\), ed anzi riesce a determinare l’\textit{intera famiglia delle superficie per cui codesto massimo è raggiunto}. I primi esempi sono:

1°) la superficie del 5° ordine di \(S_3\) per cui \(p=4\) e \(p^{(1)}=6\);

2°) la superficie canonica \(F_8\) dello spazio a quattro dimensioni, intersezione di una quadrica con una varietà del 4° ordine, per cui, come si è visto, \(p=5\) e \(p^{(1)}=9\); % ?? uncertain: the subscript of F in example 2 -- 6 / 8.
% EDITORIAL D0157 RISOLTO: The enlarged subscript has two closed loops, confirming 8. The source therefore agrees with the degree of a quadric/quartic intersection; this is a transcription correction, not a mathematical emendation.

3°) la superficie che si ottiene in \(S_6\) sopra il cono a tre dimensioni del 4° ordine, che proietta da un punto \(O\) una superficie di Veronese, intersecandolo con una ipersuperficie del 4° ordine che passi per un cono quadrico a due dimensioni, di vertice \(O\), giacente su di esso:
\[
p=7,\qquad p^{(1)}=15.
\]

Lasciamo allo studioso di verificare che la \(F_{14}\) così costruita è effettivamente una superficie canonica.

\textit{Per \(p>7\) le superficie canoniche semplici, d’ordine \(p^{(1)}-1\), il cui genere lineare ha il minimo valore \(p^{(1)}=3p-6\), si ottengono come intersezioni di una varietà semplicemente infinita di piani \(V_3^{p-3}\) d’ordine \(p-3\) nello \(S_{p-1}\) con una ipersuperficie del 4° ordine, \(V_{p-2}^4\), passante per \(p-5\) piani della \(V_3^{p-3}\).}

Ci limiteremo a accennare come si giustifichi che una superficie \(F\) d’ordine \(3p-7\) ottenuta in tal guisa, sia effettivamente una superficie canonica dello \(S_{p-1}\). A tal uopo conviene ricordare come la varietà \(V_3^{p-3}\) a sezioni iperpiane razionali\origfootnote{1}{s161L-1}{Le varietà a tre dimensioni a curve sezioni razionali sono: serie razionali semplicemente infinite di piani, ovvero coni del 4° ordine proiettanti da un punto le superficie di Veronese, quali vengono rappresentate sullo spazio ordinario da quadriche tangenti in un punto. Cfr. \textsc{F. Enriques}, \textit{Sui sistemi lineari di superficie algebriche ad intersezioni variabili iperellittiche}. Acc. Lincei, 1893. Cfr. Math. Ann., Bd. 46.} venga rappresentata nello spazio per proiezione da \(p-4\) dei suoi punti: il sistema rappresentativo è costituito da \(\infty^{p-1}\) superficie d’ordine \(p-3\) passanti \(p-4\) volte per una retta \(r\) e semplicemente per altre \(p-4\) rette \(t\) incidenti a questa, che sono le tracce dei piani per i centri di proiezione.

% scan: s161R
% !PRINTED 299
\weberpage{299}
% head: SUPERFICIE REGOLARI CANONICHE E PLURICANONICHE
% stamp: Edge shadow, scattered scan specks, and a diagonal stray stroke near “condotto”.
Si può supporre che la superficie \(F\) contenga i \(p-4\) centri di proiezione anzidetti, e quindi che essa abbia come immagine sopra lo \(S_3\) rappresentativo una superficie \(F_{2p-3}\) contenente la \(r\) come retta \((p-4)\)-pla e come rette triple le \(p-4\) rette \(t\), incidenti ad essa. Riesce quindi evidente che le superficie \(\Phi_{2p-7}\) aggiunte alla \(F_{2p-3}\) si spezzano nei \(p-4\) piani \(rt\) e in una residua superficie \(\Phi_{p-3}\) d’ordine \(p-3\) che passa semplicemente per le \(p-4\) rette \(t\). Così appare che le sezioni iperpiane della superficie normale \(F\) dello \(S_{p-7}\) sono le curve canoniche, c. d. d. % [sic] S_{p-7} -> S_{p-1}, the preceding page places this normal surface in S_{p-1}.

\edsection{12. Limite superiore del genere lineare \(p^{(1)}\) rispetto al genere superficiale.}{12. Limite superiore del genere lineare p\textasciicircum{}(1) rispetto al genere superficiale.}{sec.VIII.12}

La stessa costruzione delle superficie canoniche così come il computo dei moduli, suggerisce che il genere lineare \(p^{(1)}\) non possa superare un certo massimo rispetto al \(p\) (\(p_a\) o \(p_g\)). Se pure non si sia riusciti ancora a segnare il massimo effettivamente raggiunto dal \(p^{(1)}\), convien dire che una prima \textit{diseguaglianza} in questo senso è stata scritta da \textsc{A. Rosenblatt}\origfootnote{1}{s161R-1}{\textsc{A. Rosenblatt}, \textit{Sur quelques inégalités ecc.} Comptes rendus Ac. des Sciences, Paris, t. 154 (1912).}: % ?? uncertain: the place name in footnote 1 -- Peris / Paris.
% EDITORIAL D0158 RISOLTO: The second letter has the closed lower bowl and right stem of a, rather than the open right side and central bar of e.
\[
p^{(1)}\leq 16p_a+27.
\]

Più di recente \textsc{F. Severi}\origfootnote{2}{s161R-2}{Rendic. Circolo Mat. di Palermo, t. LVI, pag. 79 (1932).}, appoggiandosi ad un risultato trascendente di \textsc{Hodge}, è pervenuto alla \textit{diseguaglianza}:
\[
p^{(1)}\leq 2p_g+8p_a+11-\varrho,
\]
dove \(\varrho\) designa il numero base: numero dei sistemi continui di curve indipendenti tracciati sopra la superficie priva di curve eccezionali.

La questione è sempre aperta sotto l’aspetto geometrico, anche nel caso delle superficie regolari: riprendendo il calcolo dei moduli che (per \(p_a=p_g=p>3\) e per superficie canoniche semplici) ci ha condotto alla formula
\[
M=10p-2p^{(1)}+12+\omega,
\]
si ha da ricercare il limite superiore, e possibilmente il massimo effettivamente raggiunto, della sovrabbondanza \(\omega\). Si tratta invero di riconoscere il numero minimo delle condizioni che il gruppo dei punti cuspidali d’una superficie d’ordine \(n\) impone alle intersezioni delle superficie aggiunte dello stesso ordine, che debbano contenerli.

% scan: s162L
% !PRINTED 300
\weberpage{300}
% head: CAPITOLO OTTAVO
% stamp: Edge shadow and scattered scan specks and strokes, including a diagonal stroke between “ha” and “quindi”.
\edsection{13. Sistemi canonici appartenenti ad un’involuzione.}{13. Sistemi canonici appartenenti ad un’involuzione.}{sec.VIII.13}

Le costruzioni dei paragrafi precedenti mettono in rilievo che il sistema canonico di una superficie di genere \(p>3\), pur supposto irriducibile, può appartenere ad una involuzione \(I_n\) di un ordine \(n\geq 2\), in guisa che la corrispondente superficie canonica risulti multipla. Qui giova dire qualcosa intorno a tale eventualità.

Il caso più semplice, che abbiamo incontrato negli esempi precedenti, risponde all’ipotesi che la superficie contenga un’involuzione razionale di 2° ordine \(I_2\) cioè sia rappresentabile sul piano doppio con una qualunque curva di diramazione \(C_{2n}\) d’ordine pari \(2n\): in tal caso --- come già fu osservato da \textsc{Noether} --- le curve canoniche (aventi per immagini curve d’ordine \(n-3\) doppie e formanti la metà del sistema secondo aggiunto alla \(C_{2n}\)), appartengono sempre all’involuzione \(I_2\); si ha quindi come canonica una superficie razionale doppia (eventualmente riducibile essa stessa ad una superficie multipla).

È particolarmente notevole il fatto che il possesso di una involuzione segua necessariamente dall’ipotesi che il genere lineare \(p^{(1)}\) sia abbastanza piccolo rispetto al \(p\): \(p^{(1)}<3p-6\), giacchè si è riconosciuto che, se il sistema canonico è irriducibile semplice,
\[
p^{(1)}\geq 3p-6.
\]

Però non sussiste la proposizione inversa, giacchè esistono per esempio piani doppi con \(p^{(1)}>3p-6\); per \(n\) assai grande è tale il piano doppio con \(C_{2n}\) di diramazione affatto priva di singolarità, per cui
\[
p=\frac{n(n-3)}{2}+1,\qquad p^{(1)}=2(n-3)^2+1
\]
dove il rapporto \(\frac{p^{(1)}}{p}\) per \(n=\infty\) tende a 4.

Le superficie canoniche doppie rispondono sempre a superficie possedenti un’involuzione del secondo ordine \(I_2\), ma questa \textit{involuzione non è necessariamente razionale}. Alcuni esempi di superficie le cui curve canoniche appartengono ad un’involuzione \(I_2\) non razionale, sono stati già indicati da altri, di cui non ci sovviene la citazione. Nel nostro ordine d’idee possiamo costruire un esempio interessante di superficie canonica doppia, non riferibile ad un piano doppio, partendo dalla superficie canonica \(F_{10}\) con \(p=4\) e \(p^{(1)}=11\), la cui curva doppia possiede \(t=20\) punti tripli. Questa \(F_{10}\) ammetterà una superficie aggiunta del 5° ordine \(\Phi_5\), dotata di 20 punti conici. Nel fascio
\[
\lambda F_{10}+\Phi_5^2=0
\]
%CONT>>
% scan: s162R
% !PRINTED 301
\weberpage{301}
% head: SUPERFICIE REGOLARI CANONICHE E PLURICANONICHE
%<<CONT
% stamp: Scattered scan specks and a dark strip along the left edge.
si vedrà la \(F_{10}\) ridursi per continuità alla \(\Phi_5^2\), sulla quale non si avrà alcuna curva di diramazione propria, bensì \(t=20\) punti conici costituenti altrettante curve di diramazione infinitesime. La quintica \(\Phi_5\), di generi \(p=4\) e \(p^{(1)}=6\) costituisce così una \textit{superficie canonica doppia} coi generi \(p=4\) e \(p^{(1)}=11\). Ci basti di avere accennato a questo caso, che indichiamo come argomento di studio al lettore volonteroso.

I primi esempi di superficie il cui sistema canonico appartiene ad un’involuzione \(I_n\) d’ordine \(n>2\) sono stati segnalati da \textsc{B. Segre}\origfootnote{1}{s162R-1}{\textsc{B. Segre}, \textit{Sulle superficie algebriche aventi il sistema canonico composto con un’involuzione}. Rendic. Acc. Lincei, s. VI, vol. XVI, pag. 316 (1932).}. Il più semplice è dato dalla superficie del 6° ordine \(F_6\) possedente due oscnodi \(A\) e \(B\): \(p=4\), \(p^{(1)}=9\). Le \(\infty^3\) quadriche aggiunte alla \(F_6\) toccano in \(A\) e \(B\) i piani oscnodali, e perciò s’incontrano a due a due in coniche tangenti a questi stessi piani, che giacciono in piani per \(A\) e \(B\): le dette coniche formano, nello spazio, una congruenza del 1° ordine e segano sulla sestica \(F_6\) un’involuzione \(I_4\), con cui risulta composto il sistema canonico. La \textit{superficie canonica} corrispondente è una \textit{quadrica quadrupla}.

Ad una quadrica, e più precisamente, ad un cono quadrico 4-plo canonico, coi medesimi caratteri \(p=4\) e \(p^{(1)}=11\), si è pur condotti dal \textit{piano doppio con curva di diramazione} \(C_{18}\), d’ordine 18, con 8 punti 6-pli, giacchè qui le curve canoniche hanno per immagini doppie le sestiche di genere due passanti doppiamente per gli 8 punti, le quali appartengono, com’è noto, ad un’involuzione \(I_2\) (3° tipo di Bertini).

Lasciamo allo studioso di confrontare i due esempi e di riconoscere se essi rispondano a tipi distinti di superficie e se rientrino o meno in una medesima famiglia di superficie più ampia.

Il secondo esempio di \textsc{B. Segre} è la \textit{sestica \(F_6\) dotata di un oscnodo \(A\) e di un punto triplo \(B\) cui sia vicina una retta doppia infinitesima}: \(p=4\), \(p^{(1)}=7\). È chiaro che anche qui le quadriche aggiunte si segano a due a due nelle coniche di una congruenza del 1° ordine, le quali incontrano la sestica in terne di punti di un’involuzione \(I_3\), colla quale viene composto il sistema canonico. La sestica \(F_6\) conduce dunque ad una quadrica tripla canonica, che ha, come si è detto,
\[
p=4\qquad \text{e}\qquad p^{(1)}=7.
\]

Questo caso di \textsc{B. Segre} viene a realizzare un’osservazione di \textsc{Castelnuovo}, il quale, fin dal 1892, aveva avvertito che, qualora le curve canoniche di una superficie contengano una serie \(g_3^1\), la serie caratteristica del sistema dovrà essere composta con questa, e perciò il sistema canonico apparterrà ad un’involuzione razionale del 3° ordine \(I_3\).

% scan: s163L
% !PRINTED 302
\weberpage{302}
% head: CAPITOLO OTTAVO
% stamp: Dark edge strip, marginal strokes, and a diagonal stray stroke through “cubico”.
Il problema di « \textit{determinare tutte le superficie \((p>2)\) le cui curve canoniche} (irriducibili senza punti base) \textit{contengono una \(g_3^1\) e che perciò si trasformano in superficie canoniche triple, a sezioni razionali} » è stato risoluto di recente da \textsc{G. Pompilj} nella Nota\origfootnote{1}{s163L-1}{Rendic. Ist. Lombardo, 27 marzo 1941.} « Sulle superficie algebriche le cui curve canoniche posseggono una \(g_3^1\) ». Vi sono \textit{tre tipi di superficie canoniche} che rispondono a queste condizioni:

1) un piano triplo con curve di diramazione \(C_{12}\) d’ordine 12 (già segnalato nelle « Lezioni » di Enriques-Campedelli)\origfootnote{2}{s163L-2}{Ivi si nota che un tal piano triplo con \(p=3\) e \(p^{(1)}=4\) (che il computo dei moduli mostra essere caso particolare della famiglia) nasce dalla superficie del 5° ordine \(F_5\) dotata di tacnodo, per proiezione da questo punto singolare. Si può similmente partire da una superficie canonica \(F_6\) dello \(S_3\), che acquisti un punto triplo.

Questa considerazione si estende facilmente; così dalle superficie canoniche \(F_7\), \(F_8\ldots\) dello \(S_3\), che acquistino un punto triplo, si ottiene la costruzione di particolari piani multipli canonici, piani quadrupli o quintupli o in generale \((p^{(1)}-1)\)-pli, di genere \(p=3\) e \(p^{(1)}=5,6\ldots\)

Invece facendo acquistare alle dette superficie canoniche due punti tripli, si ottengono superficie con \(p=2\) di genere lineare \(p^{(1)}=2,3,\ldots\), dove il fascio delle curve canoniche possiede \(p^{(1)}-1\) punti base.}, che ha
\[
p=3\qquad \text{e}\qquad p^{(1)}=4;
\]

2) un cono quadrico triplo con curve di diramazione \(C_{18}\) d’ordine 18 per cui
\[
p=4\qquad \text{e}\qquad p^{(1)}=7,
\]
nel qual tipo rientra (come caso particolare?) l’esempio di \textsc{B. Segre};

3) un cono cubico, razionale normale in \(S_4\), triplo con curva di diramazione \(C_{24}\) d’ordine 24, per cui
\[
p=5\qquad \text{e}\qquad p^{(1)}=10.
\]

Questi tre tipi sono analiticamente rappresentati dalle seguenti equazioni:

1) piano triplo
\[
z^3+3p_4z+q_6=0,
\]
dove \(p_4(xy)\) e \(q_6(xy)\) eguagliati a zero danno curve generali d’ordine 4 e 6 rispettivamente: curva di diramazione
\[
4p_4^3+q_6^2=0;
\]

2) piano triplo
\[
z^3+3p_6z+q_9=0,
\]
%CONT>>
% scan: s163R
% !PRINTED 303
\weberpage{303}
% head: SUPERFICIE REGOLARI CANONICHE E PLURICANONICHE
%<<CONT
% stamp: Scattered scan specks and stray strokes, including one above “posseggono” in footnote 2.
\(p_6\) e \(q_9\) curve d’ordine 6 e 9, \(p_6\) passante per due punti tripli infinitamente vicini \(O^3\) e \(O_1^3\), e \(q_9\) per \(O^5\) e \(O_1^4\); curva di diramazione
\[
4p_6^3+q_9^2=0;
\]

3) piano triplo
\[
z^3+3p_8+q_{12}=0; % [sic] 3p_8 -> 3p_8z, the factor z present in the corresponding terms of the two preceding equations is absent.
\]
\(p_8\) e \(q_{12}\) curve d’ordine 8 e 12, passanti per un punto \(O\) e per due punti \(O_1\) e \(O_2\) prossimi ad esso come viene indicato dai simboli \(p_8(O^5O_1^2O_2^2)\) e \(q_{12}(O^8O_1^4O_2^4)\); curva di diramazione
\[
4p_8^3+q_{12}^2=0.
\]

Non c’indugeremo sull’analisi di \textsc{Pompilj} che si basa sulla teoria dei piani tripli da lui stesso stabilita\origfootnote{1}{s163R-1}{\textsc{G. Pompilj}, \textit{Sulla rappresentazione algebrica dei piani tripli}. Rendic. Seminario mat. Università di Roma, s. 4, vol. III. \textit{Osservazioni sui piani tripli}. Rendic. Istituto Lombardo, 27 febbraio 1941.} e che può anche condursi in altro modo\origfootnote{2}{s163R-2}{Cfr. \textsc{A. Franchetta}, «\textit{Sulle superficie le cui curve canoniche posseggono una \(g_3^1\)}». Boll. dell’U. M. I., 1942.}. Ma vogliamo indicare come si giustifichi l’osservazione di \textsc{Castelnuovo} che « se le curve canoniche di una superficie contengono una \(g_3^1\) la serie caratteristica è necessariamente composta con questa ».

Per dimostrare l’asserto basta considerare una curva di genere \(\pi>3\), la quale contenga una \(g_3^1\) ed una serie autoresidua di dimensione \(r\geq 1\), entro cui sarà lecito prendere una \(g_n^1\) con \(n=\pi-1\): si ridurrà all’assurdo l’ipotesi che la \(g_n^1\) non sia composta colla \(g_3^1\) (all’infuori di punti fissi). Invero se la \(g_n^1\) non è composta colla \(g_3^1\) si può trasformare la detta curva in una \(C_m\) d’ordine \(m=n+3\) su cui la \(g_3^1\) e la \(g_n^1\) siano segate rispettivamente dalle rette per due punti \(A\) e \(B\): \(A\) \(n\)-plo e \(B\) 3-plo. Ora una coppia di rette per \(B\) dovrebbe far parte di una curva d’ordine \(m-3\) aggiunta alla \(C_m\), e perciò dovrebbe esistere una curva d’ordine \(m-5\) passante \(n-1=m-4\) volte per \(A\): ciò che è assurdo.

\edsection{14. Superficie di genere lineare \(p^{(1)}=2\): primo caso \(p=2\).}{14. Superficie di genere lineare p\textasciicircum{}(1)=2: primo caso p=2.}{sec.VIII.14}

Quando per una superficie \(F\) il sistema canonico è di dimensione minore di 3, cioè il genere \(p<4\), ovvero se codesto sistema (puro) è riducibile, purchè sia il genere lineare \(p^{(1)}>1\), o anche se esso appartiene ad una involuzione, in tutti questi casi in cui non si ha più una superficie canonica semplice d’ordine \(p^{(1)}-1\), si può assumere utilmente come modello proiettivo della classe di \(F\) la \textit{superficie bicanonica} di \(S_{P-1}\) (\(P=p+p^{(1)}\)), le cui sezioni iperpiane sono le curve bicanoniche, o la \textit{superficie tricanonica}, ecc. Mostriamo il
%CONT>>
% scan: s164L
% !PRINTED 304
\weberpage{304}
% head: CAPITOLO OTTAVO
%<<CONT
% stamp: Scattered scan specks, marginal strokes, and a diagonal stray stroke below “corrispondenza”.
partito che si può trarre da questa considerazione, determinando in tal guisa le superficie regolari di genere lineare più basso: \(p^{(1)}=2,3\).

Pongasi che il genere lineare (assoluto) di \(F\) sia \(p^{(1)}=2\), e il suo genere superficiale sia \(p_a=p_g=p\). Sarà, d’accordo colla relazione che porge il massimo del \(p\) rispetto al \(p^{(1)}\):
\[
p\leq 2;
\]
invero le curve canoniche non possono formare una rete di curve irriducibili, che sarebbe di grado \(p^{(2)}=p^{(1)}-1=1\), e nemmeno possono essere riducibili contenendo come parti le curve d’un fascio, che dovrebbero essere di genere \(<2\) e perciò ellittiche, portando quindi \(p^{(1)}=1\). Avremo dunque da esaminare tre casi:
\[
p=2,\qquad p=1,\qquad p=0.
\]
E in corrispondenza ai tre valori indicati di \(p\), sarà il bigenere
\[
P=p+p^{(1)}=4,3,2,
\]
e il trigenere
\[
P_3=p+3p^{(1)}-2=6,5,4.
\]

\textit{Primo caso}: \(p^{(1)}=2\), \(p=2\), \(P=4\). Qui si ha un fascio (lineare) di curve canoniche \(K\), di genere 2, e un sistema lineare \(\infty^3\) di curve bicanoniche, di genere \(3p^{(1)}-2=4\), secanti sulle \(K\) la \(g_2^1\) e perciò appartenenti ad un’involuzione razionale \(I_2\). Si avrà dunque come modello proiettivo della classe delle date superficie \(F\), una superficie bicanonica, d’ordine \(4(p^{(1)}-1)=4\); ma questa si ridurrà ad una quadrica doppia con curva di diramazione \(C_{10}\) d’ordine 10, secante le coniche sezioni piane in 10 punti; dovendo quindi le generatrici possedere un numero pari di punti di diramazione, tale quadrica sarà un cono, sul quale si aggiunge alla \(C_{10}\), come punto di diramazione, il vertice \(V\): che, del resto, risponde al punto base del fascio \(|K|\) su \(F\), il quale (per essere autoresiduo) deve cadere, per ogni \(K\), in un punto doppio della \(g_2^1\).

Riassumendo:

\textit{Le superficie di generi \(p=2\), \(p^{(1)}=2\), ammettono come modello la superficie bicanonica costituita dal cono quadrico doppio con curva di diramazione d’ordine 10, intersezione completa d’una superficie del 5° ordine che non passa per il vertice, cui si aggiunge come punto di diramazione questo stesso vertice.} Il detto cono si lascia anche rappresentare sul \textit{piano doppio con curva di diramazione d’ordine 10 dotata di due punti 5-pli infinitamente vicini}.

\textit{Osservazione.} -- All’analisi precedente si può obiettare che essa suppone le curve canoniche \(K\) irriducibili. A priori le \(K\) potrebbero essere composte colle curve \(k\) variabili in un fascio (lineare) e con
%CONT>>
% scan: s164R
% !PRINTED 305
\weberpage{305}
% head: SUPERFICIE REGOLARI CANONICHE E PLURICANONICHE
%<<CONT
% stamp: Dark edge strip and scattered scan specks and marginal strokes.
altre curve fisse \(\theta\): \(K=nk+\Sigma\theta\). Ma, ad ogni modo, le \(k\) non potrebbero essere di genere \(\pi=0\) o \(\pi=1\), giacchè ne seguirebbe che la superficie è razionale (\(p=0\)) ovvero che ha il \(p^{(1)}=1\). Le \(k\), il cui genere non può superare \(p^{(1)}\) (cfr. Cap. IV, \S~18), saranno dunque di genere 2.

Ora il sistema \(\infty^3\) \(|k'|\) aggiunto a \(|k|\) --- che è contenuto in \(|2K|\) --- avrà il grado 4 e apparterrà ad una involuzione razionale \(I_2\), conducendo ad un modello della superficie costituito da un cono quadrico doppio con curva di diramazione \(C_{10}\) d’ordine 10: questo è il modello generale delle superficie bicanoniche (\(p=1\), \(p^{(1)}=2\)) definite innanzi, e risulta a posteriori che non esistono parti fisse del fascio canonico \(|K|\), nè del sistema bicanonico. % [sic] p=1 -> p=2, this refers to the p=2 case just described.

\edsection{15. Superficie con \(p^{(1)}=2\) e \(p=1\) (\(P=3\), \(P_3=5\)).}{15. Superficie con p\textasciicircum{}(1)=2 e p=1 (P=3, P\_3=5).}{sec.VIII.15}

Qui c’è una curva canonica \(K\) di genere \(p^{(1)}=2\) è \(\infty^2\) curve bicanoniche di genere \(3p^{(1)}-2=4\) secantisi a due a due in \(4(p^{(1)}-1)=4\) punti. % ?? uncertain: the conjunction after p^{(1)}=2 -- è / e; the apparent accent may be a stray mark.
% The minus sign is printed both at the end of one line and at the beginning of the next; both signs are retained.
Se si suppone la \(K\) irriducibile, saranno pure irriducibili le curve bicanoniche di \(|2K|\), e non potranno essere iperellittiche, altrimenti la serie caratteristica della rete, formata dalle coppie di una \(g_2^1\), non sarebbe completa, siccome esige la regolarità della data superficie \(F\). Ciò posto avremo su \(F\) un sistema lineare tricanonico \(\infty^4\) \(|3K|\), formato pure di curve irriducibili, che --- segando su una curva bicanonica la serie canonica ed essendo \(p>0\) --- sarà necessariamente semplice. Esso conduce ad una superficie tricanonica \(F_9\) d’ordine \(9(p^{(1)}-1)=9\) dello spazio \(S_4\), a sezioni iperpiane di genere \(6p^{(1)}-5=7\).

La costruzione diretta della \(F_9\) di \(S_4\) sembra offrire qualche difficoltà. Invece riesce facile caratterizzare una proiezione \(F_6\) di essa nello spazio ordinario. A tal uopo si osservi che la \(F_9\) possiede una retta tripla, corrispondente alla curva canonica \(K\), su cui le curve tricanoniche segano una \(g_3^1\). Da un punto di codesta retta tripla la \(F_9\) viene proiettata in una superficie del 6° ordine \(F_6\) dello \(S_3\) che a priori deve contenere 3 rette eccezionali. La \(F_6\) possiederà un punto doppio singolare, \(O\), traccia della retta tripla di \(F_9\), e tale che l’intorno di \(O\) risponda alla curva \(K\) di genere 2. Ora le sezioni piane di \(F_6\) saranno di genere 7, cioè sestiche dotate di 3 punti doppi, le sezioni piane per \(O\) di genere 4, la singolarità del punto doppio \(O\) corrispondente ad una curva fondamentale di genere 2 abbasserà di 3 il genere delle sezioni piane, e perciò queste non avranno punti doppi lontani da \(O\). Si deduce che la \(F_6\) possiede 3 rette doppie per \(O\) (non confondibili in una retta tripla, che porterebbe un fascio di sezioni piane ellittiche e quindi \(p^{(1)}=1\)); e codeste rette \(a\), \(b\) e \(c\)
% Printer's signature line at the foot of the page.
% EDITORIAL: segnatura originale: \makebox[\linewidth]{\textsc{Enriques F.} - \textit{Superficie algebriche.}\hfill 20}
\printersignature{\textsc{Enriques F.} - \textit{Superficie algebriche.}\quad 20}
%CONT>>
% scan: s165L
% !PRINTED 306
\weberpage{306}
% head: CAPITOLO OTTAVO
%<<CONT
% stamp: Dark edge strip, marginal strokes, and a diagonal stray stroke across the equality sign in p^{(1)}=2.
stanno nel piano\origfootnote{1}{s165L-1}{Altrimenti il punto \(O\) risulterebbe triplo.} tangente in \(O\) alla \(F_6\) (piano che contato due volte ne costituisce il cono osculatore): la singolarità su ogni sezione piana di \(F_6\) per \(O\) consta di due rami che hanno un contatto del 5° ordine (contatto sestipunto).

Una \(F_6\), dotata delle singolarità indicate, possiede una quadrica aggiunta che si spezza nel piano delle tre rette doppie contato due volte; affinchè vi siano tre rette eccezionali, dovranno le nominate rette doppie essere cuspidali.

La \(F_6\) che abbiamo costruito ha il genere geometrico \(p_g=1\), possedendo una quadrica aggiunta costituita dal piano \(a\,b\,c\) contato due volte; essa ha anche il genere numerico \(p_a=1\), ossia è regolare; infatti per questa superficie, le cui sezioni piane sono di genere 7, esistono \(\infty^7\) (\(7=p_a+\pi-1\)) superficie aggiunte del 3° ordine, quali sono definite dalle condizioni di passare per le rette \(a\), \(b\), \(c\), ed avere un contatto del 4° ordine colle falde di \(F_6\) in \(O\); ed infine ha il genere lineare assoluto \(p^{(1)}=2\) perchè il suo bigenere vale \(P=p+p^{(1)}=3\), essendovi \(\infty^2\) curve bicanoniche segate dai piani per \(O\) cui va sommato il piano \(a\,b\,c\) contato tre volte: il valore del \(p^{(1)}\) risulta anche dall’ordine della \(F_6\), a sezioni piane tricanoniche, tenuto conto delle 3 rette eccezionali:
\[
6=9(p^{(1)}-1)-3=9-3\ \text{\origfootmark{2}{s165L-2}}.
\]\origfoottext{2}{s165L-2}{La classificazione delle superficie, con \(p^{(1)}=2\), \(p=1,2\), è data nella Nota di \textsc{F. Enriques}, \textit{Le superficie algebriche di genere lineare \(p^{(1)}=2\)}. Rendic. Acc. Lincei, Febbraio 1897.}
Si può costruire una \(F_6\) con le singolarità sopra indicate partendo da una superficie cubica \(F_3\) che contenga tre rette \(a\), \(b\), \(c\) giacenti in un piano e passanti per un punto \(O\). Esistono \(\infty^4\) superficie cubiche \(F'_3\) che toccano la \(F_3\) in \(O\), dando luogo sopra ogni sezione piana ad un contatto sestipunto: queste \(F'_3\) toccano di conseguenza \(F_3\) secondo le tre rette \(a\), \(b\), \(c\), come appare dalla rappresentazione piana, e segano ulteriormente \(F_3\) secondo una cubica piana, risultando caratterizzate da queste proprietà. Ora una superficie composta \(F_3+F'_3\) fornisce una particolare sestica che possiede le \(a\), \(b\), \(c\), come rette tacnodali, e presenta in \(O\) la singolarità sopra indicata (sopra ogni sezione piana per \(O\) due rami con contatto sestipunto). Combiniamo linearmente la \(F_3+F'_2\) con la sestica degenere costituita dal piano \(a\,b\,c\) contato tre volte, e da un cono cubico di vertice \(O\); in tal guisa si otterrà una \(F_6\) che avrà in \(O\) la singolarità richiesta, e per cui le rette \(a\), \(b\), \(c\), risultano cuspidali. % [sic] F'_2 -> F'_3, the same composite surface is written F_3+F'_3 immediately above.

Si verifica tosto che la \(F_6\) così costruita possiede un sistema \(\infty^7\) di superficie cubiche aggiunte \(\Phi_3\), definite dalle condizioni di passare
%CONT>>
% scan: s165R
% !PRINTED 307
\weberpage{307}
% head: SUPERFICIE REGOLARI CANONICHE E PLURICANONICHE
%<<CONT
% stamp: Scattered scan specks, a dark left edge, and fine marginal marks.
per le tre rette \(a\), \(b\), \(c\), e di toccare semplicemente in \(O\) la \(F_3\), per modo che la sezione di \(\Phi_3\) con un piano per \(O\) venga a contenere 5 fra i 6 punti doppi successivi della sezione di \(F_6\). E da ciò segue che la superficie \(F_6\), le cui sezioni piane sono di genere 7, è regolare, giacchè le dette \(\Phi_3\) segano, sopra una sezione piana di \(F_6\), il sistema completo delle sue curve aggiunte.

Frattanto riassumiamo il risultato ottenuto enunciando che: \textit{Le superficie regolari coi generi \(p=1\) e \(p^{(1)}=2\) si lasciano trasformare in sestiche \(F_6\) dotate di tre rette cuspidali giacenti in un piano e passanti per un punto doppio singolare \(O\), tali che sopra ogni sezione piana per esso si abbiano due rami con contatto del 5° ordine.}

\textit{Osservazione.} -- Nell’analisi precedente si è supposto che la curva canonica \(K\) sia irriducibile, e quindi che la rete delle curve bicanoniche \(|2K|\) non abbia parti fisse. Accenniamo rapidamente come possiamo emanciparci da questa ipotesi restrittiva, e perciò esaminiamo se per \(p=1\), \(p^{(1)}=2\) sia possibile un sistema bicanonico
\[
|2K|=|k_2|+\Sigma\theta
\]
avente qualche componente fissa \(\theta\), le cui parti variabili siano curve \(k_2\) irriducibili (cfr. Cap. IV, \S~19). Si designano con \(\pi\) ed \(n\) genere e grado di una \(k_2\), con \(\varrho\) e \(\nu\) genere e grado di una \(\theta\). Allora il carattere addittivo che esprime il numero virtuale delle intersezioni delle curve bicanoniche colle canoniche viene dato da % [sic] addittivo -> additivo, apparent doubled t.
\[
2p^{(1)}-2=2=2\pi-2-n+\Sigma(2\varrho-2-\nu).
\]
E, supponendo eliminate le curve eccezionali della superficie (e quindi per \(\varrho=0\), \(\nu\leq -2\)), si avrà per tutte le componenti \(\theta\):
\[
2\varrho-2-\nu\geq 0,
\]
e quindi
\[
2\pi-2-n\leq 2.
\]
Ora se esiste qualche \(\theta\), per essere la \(K\), e perciò anche le curve di \(|2K|\), connesse d’ordine di connessione \(\geq 2\), sarà il numero delle intersezioni di una \(k_2\) colla \(\Sigma\theta\) (residua rispetto al sistema bicanonico)
\[
k_2\cdot\Sigma\theta=4\pi-4-3n\geq 2.
\]

Segue di qui
\[
\begin{aligned}
4\pi-4-2n&=(4\pi-4-3n)+n\leq 4,\\
n&\leq 2\qquad \text{e perciò}\qquad n=2(n>1),\\ % ?? uncertain: faint mark after the first 2 in this line -- no punctuation / a full stop.
% EDITORIAL D0160 RISOLTO: The faint mark is a short horizontal stroke at mid-height to the right of 2, not a compact baseline period. It is consistent with a scan blemish.
2\pi-2&\leq n+2=4;
\end{aligned}
\]
cioè la rete \(|k_2|\) è costituita di curve di genere \(\pi\leq 3\), e di grado \(n=2\). Ma questa conclusione è assurda perchè (nell’ipotesi più
%CONT>>
% scan: s166L
% !PRINTED 308
\weberpage{308}
% head: CAPITOLO OTTAVO
%<<CONT
% stamp: Scattered scan specks, short marginal marks and a dark right scan edge.
favorevole: \(\pi = 3\)) porta come conseguenza la possibilità di rappresentare la data superficie \(F\) sopra un piano doppio con curva di diramazione \(C_8\) d’ordine 8, che --- per essere \(p = 1\) --- deve possedere due singolarità elementari, punti quadrupli o punti \([3, 3]\), onde risulta
\[
p^{(1)} = 1.
\]

Dunque il sistema bicanonico d’una \(F\) di genere \(p = 1\) con \(p^{(1)} = 2\), non possiede parti fisse, e la classificazione di codeste superficie riesce perfetta.

\edsection{16. Superficie con \(p^{(1)} = 2\) e \(p = 0\) (\(P = 2,\ P_3 = 4\)).}{16. Superficie con p\textasciicircum{}(1) = 2 e p = 0 (P = 2, P\_3 = 4).}{sec.VIII.16}

Si è condotti ad una \textit{superficie tricanonica} dello spazio ordinario \(S_3\): superficie \(F_9\) d’ordine 9, a sezioni piane di genere 7 e perciò \textit{dotata d’una curva doppia \(C_{21}\) d’ordine \(21\)}, la quale non appartiene ad alcuna superficie \(\Phi_5\), ma è egualmente doppia per \(\infty^4\) superficie biaggiunte \(\Phi_{10}\) d’ordine 10, che segano su di essa le sestiche gobbe di genere 4 d’un fascio, curve bicanoniche.

Non è facile costruire direttamente la \(F_9\) e definirne precisamente le singolarità, ma si può almeno dimostrarne l’\textit{esistenza}, partendo da particolari superficie della famiglia di cui essa porge il modello proiettivo. A tal uopo si può partire dal piano doppio di \textsc{Campedelli}, dotato di curve di diramazione \(C_{10}\) d’ordine 10 con un punto 4-plo \(A\) e 5 punti \([3, 3]\), per il quale
\[
p = 0 \qquad \text{e} \qquad p^{(1)} = 2,
\]
avendosi un fascio di curve bicanoniche di genere \(3p^{(1)} - 2 = 4\), rappresentate dalle quartiche (doppie) che passano doppiamente per \(A\) e semplicemente per le 5 coppie di punti tripli infinitamente vicini di \(C_{10}\), formanti le singolarità \([3, 3]\). Il piano doppio così definito rappresenta una superficie \(F\) su cui le curve bicanoniche non sono iperellittiche e perciò le curve tricanoniche formano un sistema semplice (vedi l’Osservazione più avanti). Mercè questo sistema la \(F\) si lascia trasformare nella \(F_9\) tricanonica sopraindicata, o meglio, in un caso particolare di essa.

Invero i nostri piani doppi dipendono da 4 moduli, essendovi 4 invarianti proiettivi della \(C_{10}\) di diramazione:
\[
4 = 65 - 8 - 5 \cdot 9 - 8;
\]
invece il numero dei moduli, che spettano alla famiglia della superficie \(F_8\), con \(p = 0\) e \(p^{(1)} = 2\), è certo % ?? uncertain: subscript of F after "famiglia della superficie" -- 8 or 9.
% EDITORIAL D0161 RISOLTO: The enlarged final subscript has two closed bowls. Preserve the printed F_8; the apparent mathematical inconsistency is recorded separately.
\[
M \geq 9p - 2p^{(1)} + 12 = 8.
\]

% scan: s166R
% !PRINTED 309
\weberpage{309}
% head: SUPERFICIE REGOLARI CANONICHE E PLURICANONICHE
% stamp: Scattered scan specks, a dark left scan edge and thin marks near the right margin.
\textit{Osservazione.} -- Aggiungeremo che la tricanonica corrispondente ad una superficie \(F\) regolare con \(p = 0\) e \(p^{(1)} = 2\) riesce certo semplice, almeno se le curve bicanoniche di \(F\) sono irriducibili.

Accenniamo rapidamente come possa giustificarsi l’asserto.

Anzitutto conviene stabilire che dall’essere irriducibili le curve bicanoniche \(K_2\) segue che sono irriducibili e senza punti base le tricanoniche \(K_3\) del sistema aggiunto \(|K'_2| = |K_3|\). Siccome le \(K'_2\) segano sopra una \(K_2\) la serie canonica completa, se esse sono riducibili dovranno contenere una parte fissa \(\theta\) (irriducibile) fondamentale pel fascio \(|K_2|\), e allora il sistema tricanonico spogliato della \(\theta\) costituirà l’aggiunto alla curva \(K_2 - \theta\), sicchè \((K_2 - |\theta|)'\) avrebbe egual dimensione di \(|K'_2|\), cioè dimensione 3: si tratta quindi di riconoscere che ciò è impossibile perchè la \(K_2 - \theta\) è una curva connessa di genere \(< 4\) (Cfr. Cap. IV, § 17). A tale scopo si realizzeranno le \(K_2\) come sestiche canoniche dello spazio ordinario, proiettate sul piano in sestiche con due punti tripli, e se ne studieranno le possibili degenerazioni: così appunto si vede che il distacco di una \(\theta\) irriducibile da \(|K_2|\) conduce necessariamente ad una curva connessa di genere \(\leq 2\). % [sic] (K_2 - |\theta|)' -> |K_2 - \theta|', the printed bars enclose only theta rather than the residual system named in the preceding clause.

Una analisi simile vale ad escludere che il \(|K_3|\) irriducibile abbia qualche punto base sopra la data superficie \(F\): in primo luogo non può esservi un punto base che sia semplice per la curva \(K_3\) che vi passa, ecc. Ciò posto, essendo \(|K_3|\) un sistema lineare irriducibile senza punti base, di grado 9, ove esso appartenga ad una involuzione \(I_n\) d’ordine \(n\), dovrà essere \(n\) divisore di 9 e perciò \(n = 3\). In tale ipotesi si è condotti ad esaminare una superficie cubica \(F_3\) razionale a sezioni ellittiche, tripla, con curve di diramazione \(C_{12}\) d’ordine 12, che dovrebbe costituire la tricanonica della \(F\) di genere \(p_a = p_g = 0\) e di genere lineare \(p^{(1)} = 2\). Ma è facile persuadersi che una tale cubica \(F_3\) tripla, di generi \(p = 0\) e \(p^{(1)} = 2\) \textit{non esiste}. Diciamo anzitutto che la nostra \(F_3\), che contata 3 volte costituirebbe la tricanonica di una superficie \(F\) con \(p = 0\) e \(p^{(1)} = 2\), deve possedere almeno un punto doppio (singolare). Invero, se si suppone la \(F_3\) priva di punti doppi si avranno su di essa due reti omaloidiche di cubiche gobbe complementari (rispetto al sistema delle sestiche segato dalle quadriche): \(|C_3|\) e \(|L_3|\); e si vede tosto che queste debbono rappresentare due reti di curve, rispettivamente \(|C|\) e \(|L|\), di grado 3 e di genere 4. Se così non fosse, poichè le \(C_3 + L_3\) segano la \(C_{12}\) di diramazione della \(F_3\) tripla in 24 punti, sarebbe una delle due reti, per esempio \(|C|\), di genere \(\pi < 4\); ma, se \(\pi < 3\), tutti i plurigeneri della \(F\) risultano nulli; se invece \(\pi = 3\), le curve bicanoniche di \(F\), che --- per ipotesi --- formano un fascio, segherebbero su una \(C\) le coppie di una \(g_2^1\), ciò che è incompatibile col fatto che le \(C\) contengono una \(g_3^1\) senza punti fissi.

% scan: s167L
% !PRINTED 310
\weberpage{310}
% head: CAPITOLO OTTAVO
% stamp: Scattered scan specks, marginal marks, a short diagonal mark near "fissato" and dark scan edges.
Ciò posto, essendo le \(C\) e le \(L\) curve di genere 4, si riconosce che su di esse la \(g_3^1\) caratteristica è autoresidua (rispetto alla serie canonica). Se così non fosse, ad es. per la \(C\), si prova che il genere delle curve canoniche virtuali \(K\), che hanno 3 punti comuni colle \(C\), deve valere \(p^{(1)} = 1\), anzichè \(p^{(1)} = 2\). Invero si assuma come modello di una \(F\) una \(F_n\) semplice, d’un certo ordine \(n\) assai alto, su cui le \(C\) vengano segate dai piani per un punto \((n - 3)\)-plo \(O\): le curve del sistema aggiunto \(|C'| = |C + K|\) appartenenti ad \(F\) saranno proiettate da \(O\) in curve piane del \(6^{\mathrm{o}}\) ordine e, in particolare le \(C'\) per un gruppo della \(g_3^1\) caratteristica di \(C\)(che sono \(\infty^1\)) in sestiche con un punto triplo: ciò significa che il genere delle \(C\) vale % [sic] C -> C', the preceding sentence describes the adjoint curves C'.
\[
7 = 4 + p^{(1)} + 3 - 1,
\]
e quindi
\[
p^{(1)} = 1.
\]

Resta dunque fissato che ciascuna delle due reti di curve \(|C|\) e \(|L|\) su \(F\) possiede una serie caratteristica autoresidua. Di conseguenza le curve bicanoniche \(K_2\) (di \(|2K|\)) segano su una \(C\) gruppi della serie doppia della detta \(g_3^1\) e quindi sono contenute parzialmente nel sistema doppio \(|2C|\), e così similmente in \(|2L|\).

Siccome il \(|2C|\) completo ha la dimensione 6 e il \(|2C|\) minimo la dimensione 5, vi è \textit{una} curva del fascio \(|K_2|\) appartenente al \(|2C|\) minimo, e questa è caratterizzata dall’appartenere alla involuzione \(I_3\) le cui terne rispondono ai punti della \(F_3\), e perciò deve appartenere altresì al minimo sistema \(|2L|\).

Ora si rappresenti la cubica \(F_3\) punto per punto su un piano \(\alpha\): le sezioni piane di \(F_3\) avranno per immagini le cubiche per 6 punti \(A_1A_2\ldotp\ldotp\ldotp\ldotp A_6\), non appartenenti ad una conica; le curve del minimo sistema \(|2C|\) potranno supporsi rappresentate dalle coniche di \(\alpha\); e quelle del minimo sistema \(|2L|\) da \(C_{10}\) d’ordine 10 passanti 4 volte per i 6 punti \(A_i\). Quindi dovrebbe esservi nel piano \(\alpha\) una conica, passante per qualche punto \(A_i\), che faccia parte insieme di una \(C_{10}\) e incontri le \(C_{10}\) medesime in 2 punti variabili: ma una tale conica non esiste. Così è rifiutata l’ipotesi di una \(F_3\) (immagine tripla di una \(F\) tricanonica con \(p = 0\) e \(p^{(1)} = 2\)) che non possegga punti doppi.

Supponiamo d’altra parte che esista una \(F_3\) dotata di un punto doppio \(O\) (comunque singolare), le cui sezioni piane siano immagini delle curve tricanoniche \(K_3\) di \(F\).

Le sezioni piane \(C_3\) di \(F_3\) per \(O\) (come si è detto prima per le eventuali cubiche gobbe appartenenti a questa superficie) saranno immagini di curve \(C\) del genere 4, e le curve bicanoniche \(K_2\) di \(F\) avranno per immagini su \(F_3\) delle sestiche semplici, incontranti le dette \(C_3\) in 6 punti variabili, e perciò non passanti per \(O\). Di conse%
%CONT>>
% scan: s167R
% !PRINTED 311
\unskip\weberpage{311}%%
% head: SUPERFICIE REGOLARI CANONICHE E PLURICANONICHE
%<<CONT
% stamp: Scattered scan specks, short marginal marks and a dark left scan edge.
guenza il sistema quadricanonico aggiunto a \(|K_3|\) cioè
\[
|2K_2| = |K'_3| = |4K|
\]
sarà costituito in generale di curve non passanti per \(O\); e ad \(O\) risponderà su \(F\) una curva \(\theta\) fondamentale per il detto \(|K_3|\). Ma da ciò si deduce che togliendo la \(\theta\) da tale sistema si ha l’aggiunto a \(|C|\):
\[
\begin{aligned}
|C'| &= |K'_3 - \theta|\\
(|K_3 - \theta| &= |C|),
\end{aligned}
\]
e quindi il genere delle \(C\) deve risultare eguale a quello delle \(K_3\) diminuito di una unità:
\[
7 - 1 = 6;
\]
conseguenza assurda perchè il genere delle \(C\) vale, come si è detto,
\[
\pi = 4!
\]

Così è esaurita la discussione: \textit{le \(F_3\) triple bicanoniche di genere \(p = 0\) e \(p^{(1)} = 2\) non esistono.} \hfill c. d. d. % [sic] bicanoniche -> tricanoniche, the preceding discussion concerns the tricanonical cubic model.

\edsection{17. Superficie di genere lineare \(p^{(1)} = 3\): primo caso \(p = 3\)\origfootnote{1}{s167R-1}{La classificazione delle superficie di genere lineare \(p^{(1)} = 3\), per \(p = 1, 2, 3\), che viene qui in varii sensi approfondita ed estesa al caso \(p = 0\), è data da \textsc{F. Enriques} nella Nota \textit{Sulle superficie algebriche di genere lineare \(p^{(1)} = 3\)}. Rendic. Lincei, maggio 1897. Manca soltanto il tipo semplice per \(p^{(1)} = 3\) e \(p = 2\), del § 18.}.}{17. Superficie di genere lineare p\textasciicircum{}(1) = 3: primo caso p = 3.}{sec.VIII.17}

Esaminiamo le superficie (regolari) di genere lineare \(p^{(1)} = 3\). Il loro genere superficiale
\[
p \leq \frac{p^{(1)} + 3}{2}
\]
non potrà superare il valore 3 (cfr. § 11). Avremo dunque a priori 4 casi:
\[
p = 3,\, 2,\, 1,\, 0.
\]

\textit{Primo caso:} \(p^{(1)} = 3,\ p = 3\). Il sistema canonico \(|K|\), supposto irriducibile, è una rete di grado \(p^{(2)} = p^{(1)} - 1 = 2\), che conduce a rappresentare la superficie sul \textit{piano doppio con curva di diramazione \(C_8\) d’ordine \(8\)}.

L’ipotesi della riducibilità di \(|K|\) in una componente fissa \(\Sigma\theta\) e in una parte variabile \(|k|\) non conduce ad alcun caso nuovo. Si designino con \(\pi\) ed \(n\) i caratteri delle \(k\), e con \(\varrho\) e \(\nu\) quelli di una \(\theta\). Calcolando il carattere additivo che indica il numero delle intersezioni delle curve canoniche con queste stesse curve, avremo:
\[
2\pi - 2 - n + \Sigma(2\varrho - 2 - \nu) = 2.
\]

% scan: s168L
% !PRINTED 312
\weberpage{312}
% head: CAPITOLO OTTAVO
% stamp: Scattered scan specks, marginal marks, a short diagonal mark below the equalities and a dark right scan edge.
Ora il numero delle intersezioni di una \(k\) (irriducibile) col resto delle curve canoniche, \(\Sigma\theta\), è
\[
2\pi - 2 - 2n \geq 0,
\]
ed il grado del sistema \(|k|\) è \(n \geq 2\), quindi
\[
2\pi - 2 - n \geq 2.
\]

D’altra parte, se si ammette di avere eliminato dalla superficie ogni curva eccezionale, sarà per ciascuna \(\theta\):
\[
2\varrho - 2 - \nu \geq 0;
\]
per conseguenza
\[
\begin{aligned}
2\pi - 2 - n = 2,\qquad &2\varrho - 2 - \nu = 0,\\
n = 2 \qquad \text{e}\qquad &\pi = 3.
\end{aligned}
\]

Pertanto le \(k\) formano una rete che conduce a trasformare la superficie data in un piano doppio con \(C_8\) di diramazione, e le \(\theta\) risultano a posteriori non esistere.

\edsection{18. Superficie con \(p^{(1)} = 3\) e \(p = 2\) (\(P = 5\)).}{18. Superficie con p\textasciicircum{}(1) = 3 e p = 2 (P = 5).}{sec.VIII.18}

Supposto che il fascio canonico \(|K|\) sia irriducibile a prescindere dalle componenti eccezionali fisse che rispondono ai suoi punti base, le \(\infty^4\) curve bicanoniche, secanti su ogni \(K\) la serie canonica completa \(g_4^2\), formeranno pure un sistema irriducibile, di genere \(3p^{(1)} - 2 = 7\) e grado \(4(p^{(1)} - 1) = 8\), privo di punti base. E questo sistema dovrà essere semplice, se le curve canoniche del fascio \(|K|\) non sono iperellittiche, come conviene in primo luogo supporre. Quindi si è condotti ad una \(F_8\) bicanonica semplice dello \(S_4\), a sezioni iperpiane di genere \(3p^{(1)} - 2 = 7\), su cui le \(K\) sono quartiche piane aventi due punti comuni e perciò giacenti in piani per una retta \(a\), la quale non sta su \(F_8\). Gli iperpiani per \(a\) segano \(F_8\) secondo coppie di quartiche \(K\); segue che i piani delle \(K\) formano un cono quadrico di \(2^{\mathrm{a}}\) specie \(Q\), che ha per asse la \(a\). Ora si proietti \(F_8\) sullo spazio \(S_3\) da un punto generico \(P\) di \(Q\): si avrà in \(S_3\) una superficie \(F'_8\) con due rette quadruple \(o\) e \(o'\) infinitamente vicine, uscenti da un punto \(O\). Ciò significa che la \(F_8\) è intersezione completa del cono quadrico \(Q\) e d’una ipersuperficie del \(4^{\mathrm{o}}\) ordine. Ma (come già abbiamo veduto) l’intersezione di quadrica e quartica in \(S_4\) è in genere una superficie canonica di generi \(p = 5\) e \(p^{(1)} = 9\): per avere una \(F_8\) bicanonica (\(p = 2,\ p^{(1)} = 3\)) bisogna imporre all’intersezione suddetta convenienti singolarità che ne abbassino il genere, e precisamente --- poichè il genere 9 delle sezioni iperpiane deve abbassarsi % ?? uncertain: capital E before "questo sistema" -- E or È.
%CONT>>
% scan: s168R
% !PRINTED 313
\weberpage{313}
% head: SUPERFICIE REGOLARI CANONICHE E PLURICANONICHE
%<<CONT
% stamp: Scattered scan specks, dark left-edge artefacts and a short horizontal mark near the upper right edge.
di 2 --- converrà imporre alla \(F_8\) una linea doppia \(C_2\) del \(2^{\mathrm{o}}\) ordine. Ora, riferendosi alla \(F'_8\) proiezione di \(F_8\) in \(S_3\) da un punto del cono quadrico che la contiene, si vede che la \(C'_2\) proiezione di \(C_2\) non può constare di due rette, incidenti o meno, che non stiano in un piano per \(o\) e quindi deve essere una conica giacente in un piano per \(o\). In conclusione il \textit{tipo della superficie bicanonica semplice, per \(p = 2\) e \(p^{(1)} = 3\), è una \(F_8\) d’ordine \(8\) dello spazio \(S_4\) appartenente a un cono quadrico \(Q\) e avente come proiezione da un punto di \(Q\) nello spazio ordinario una \(F_8\) dotata di due rette quadruple infinitamente vicine \(o\) \(o'\) incidenti e di una conica doppia, giacente in un altro piano per \(o\).}

Si costruisce la \(F'_8\) combinando linearmente: una quartica con due rette doppie \(o\) e \(o'\) passante per la conica \(C'_2\), contata due volte, ed una \(F'_8\) degenere composta del piano \(\omega = o\,o'\) contato 4 volte e del piano \(\gamma\) della conica \(C'\) contato due volte, oltrechè di una generica quadrica. E si verifica tosto che la superficie \(F'_8\) dotata delle singolarità che ne risultano è proprio la nostra superficie bicanonica, non normale. Infatti le superficie aggiunte ad essa si ridurranno: 1) per effetto della presenza delle 2 rette 4-ple \(o\) e \(o'\): a coni tangenti lungo \(o\) (cui si aggiunge il piano \(\omega = oo'\) contato due volte); e 2) per effetto della conica doppia \(C'_2\) (il cui piano \(\gamma\) si stacca dai detti coni) a piani per \(o\) secanti su \(F'_8\) le quartiche \(K\); onde segue \(p = 2,\ p^{(1)} = 3\). Inoltre c’è una \(\Phi_7\) biaggiunta, \(\Phi_7 = \omega^5\gamma^2\), che sega \(F'_8\) soltanto secondo le curve singolari, dimodochè le sezioni piane di \(F'_8\) risultano curve bicanoniche.

\textit{Ogni superficie con \(p = 2\) e \(p^{(1)} = 3\), con curve canoniche} (irriducibili) \textit{non iperellittiche, si può trasformare nella \(F_8\) bicanonica semplice di \(S_4\).}

Un esempio particolare è dato dal \textit{piano doppio con curva di diramazione \(C_{12}\) d’ordine \(12\), dotata di \(8\) punti \(4\)-pli \(A_1A_2\ldotp\ldotp\ldotp\ldotp A_8\)}. Qui le curve canoniche sono rappresentate dalle cubiche \(C_3\) per \(A_1A_2\ldotp\ldotp\ldotp\ldotp A_8\), le quali --- prese come doppie, coi punti di diramazione negli incontri con \(C_{12}\) --- costituiscono, in generale, curve di genere 3 non iperellittiche. In effetto ci sono \(\infty^4\) curve bicanoniche aventi per immagini sul piano curve d’ordine 12 passanti 4 volte per \(A_1A_2\ldotp\ldotp\ldotp\ldotp A_8\) e 8-tangenti alla \(C_{12}\) di diramazione; e queste rappresentano un sistema lineare completo \(\infty^4\) semplice. Entro il detto sistema completo vi sono \(\infty^3\) curve bicanoniche che hanno per immagini le sestiche \(C_6\) passanti doppiamente per i punti 4-pli di \(C_{12}\) e formano un sistema lineare, non completo, appartenente ad una involuzione del \(4^{\mathrm{o}}\) ordine.

Nell’esempio precedente si vede una particolare \(F_8\) bicanonica che --- da un certo punto esterno --- viene proiettata in un cono del secondo ordine quadruplo. Un altro esempio notevole è dato da una \(F_8\) che può proiettarsi da un punto esterno \(O\) in una \textit{quartica \(F_4\)
%CONT>>
% scan: s169L
% !PRINTED 314
\weberpage{314}
% head: CAPITOLO OTTAVO
%<<CONT
% stamp: Scattered scan specks, short marginal and diagonal marks and a dark right scan edge.
di genere \(1\) doppia con curva di diramazione piana del \(4^{\mathrm{o}}\) ordine \(F_4\) dotata altresì di un piano tangente lungo una conica, sulla quale si trovano \(6\) punti conici di diramazione.}

Invero se la \(F_8\) è trasformata in sè da una involuzione \(I_2\), subordinata da un’omologia armonica di centro \(O\), essa avrà per proiezione da \(O\) una \(F_4\), sulla quale le curve canoniche \(K\) hanno per immagini le sezioni piane di genere 3 d’un fascio, e che perciò deve essere priva di curva doppia: siccome poi la \(I_2\) coniuga a coppie le curve canoniche del fascio \(|K|\) su \(F_8\), in questo fascio si debbono trovare due \(K\) unite, una delle quali si proietta nella curva di diramazione \(C_4\) di \(F_4^2\), l’altra (iperellittica) nella conica \(C\) sezione d’un piano tangente \(\alpha\), e su questa si trovano (oltre le due intersezioni con \(C_4\)) 6 punti di diramazione, che sono i punti conici di \(F_4\) su \(\alpha\). Reciprocamente, si assuma una \(F_4\) con piano \(\alpha\) tangente lungo una conica \(C\) e si costruisca la \(F = F_4^2\) estraendo su \(F_4\) la \(\sqrt{\alpha\beta}\), in guisa da possedere una quartica \(C_4\) di diramazione \(\beta = 0\), cui si aggiunge la conica critica apparente \(C\). È facile persuadersi che le \(\infty^1\) sezioni dei piani passanti per la retta comune ai piani \(\alpha\) e \(\beta\) costituiscono quartiche doppie, immagini di curve spezzate come la sezione del piano \(\beta\); e perciò --- prese semplicemente --- rappresentano \(\infty^1\) curve canoniche di \(F\), sicchè per la detta \(F\) risulta \(p_g = 2\). Si può anche verificare che \(p_a = 2\), sia mostrando a priori la regolarità della superficie, sia calcolando il \(p_a\) mediante le formule stabilite per la corrispondenza \([1, 2]\) fra \(F_4\) ed \(F\): dove occorre tener conto che i 6 punti doppi di \(F_4\) sulla conica \(C\) costituiscono 6 curve infinitesime di diramazione (da aggiungere alla quartica sezione del piano \(\beta\)) cui rispondono curve eccezionali o punti semplici della superficie \(F\).

\textit{Ma un secondo tipo affatto distinto}, che non rientra in quello delle \(F_8\) bicanoniche semplici (\(p = 2,\ p^{(1)} = 3\)) è dato dalla \textit{superficie bicanonica costituita da una superficie \(F_4\) di Segre}\origfootnote{1}{s169L-1}{Cfr. \textsc{F. Enriques-F. Conforto}: « \textit{Le superficie razionali} », pag. 144.} (con due punti doppi) \textit{doppia}, ovvero dal \textit{piano doppio con curva di diramazione \(C_{10}\) d’ordine \(10\), dotata di \(2\) punti \(4\)-pli e di \(2\) punti \([3, 3]\)}\origfootnote{2}{s169L-2}{O, se piace meglio, di una \(C_{10}\) costituita da una retta \(r\) e da una \(C_9\) con tre punti 3-pli su \(r\) e, fuori di essa, due punti 3-pli infinitamente vicini.}.

Qui le curve canoniche sono rappresentate da \(\infty^1\) coniche d’un fascio (ciascuna da prendersi come doppia con \(6 + 2 = 8\) punti di diramazione) e le bicanoniche da \(\infty^4\) quartiche con due punti base doppi e due contatti:
\[
p = 2,\qquad p^{(1)} = 3.
\]

Si dimostra facilmente che \textit{ogni superficie \(F\) con \(p = 2\) e \(p^{(1)} = 3\) che conduca in \(S_4\) ad una bicanonica doppia è a curve canoniche iper%
%CONT>>
% scan: s169R
% !PRINTED 315
\unskip\weberpage{315}%%
% head: SUPERFICIE REGOLARI CANONICHE E PLURICANONICHE
%<<CONT
% stamp: Scattered scan specks, a speckled mark in the upper right margin and a dark left scan edge.
ellittiche e si lascia rappresentare sulla \(F_4^2\) di \(S_4\) o sul piano doppio anzidetto.}

Anzitutto, se il sistema bicanonico \(|2K|\) di \(F\) non è semplice, vuol dire che le curve di esso passanti per un punto \(A\) debbono passare di conseguenza per (almeno) un altro punto \(A'\); ma questo non può essere fuori della \(K\) che passa per \(A\), e se sta sulla \(K\) (supposta irriducibile) vuol dire che la serie canonica \(g_4^2\) segnata da \(|2K|\) su \(K\) è composta e la \(K\) è iperellittica.

Ora se le \(K\) sono iperellittiche la superficie bicanonica trasformata di \(F\) si ridurrà ad una \(F_4\) doppia, e perchè la \(F_4\) sia normale in \(S_4\) dovrà essere a sezioni iperpiane ellittiche, ossia una superficie di Segre, intersezione completa di due quadriche di \(S_4\). Sulla \(F_4^2\), le cui sezioni iperpiane doppie sono di genere \(3p^{(1)} - 2 = 7\), si avrà una curva di diramazione \(C_{12}\) d’ordine 12. Aggiungasi che la nostra \(F_4\) di \(S_4\) non è la più generale intersezione di due quadriche, bensì è dotata di due punti doppi, i quali cadono nei punti base del fascio delle coniche canoniche (immagini di due punti doppi della \(g_2^1\) sulle \(K\)). La rappresentazione punto per punto della \(F_4\) sul piano conduce ad un sistema rappresentativo delle sezioni iperpiane costituito da curve del \(4^{\mathrm{o}}\) ordine con due punti doppi e due punti base semplici di contatto: queste sono le immagini delle curve bicanoniche d’un piano doppio, con \(p = 2\) e \(p^{(1)} = 3\), di cui la curva di diramazione \(C_{10}\) possiede due punti 4-pli e due punti \([3, 3]\).

Abbiamo affermato che le \(F_4\) di Segre doppie costituiscono un \textit{secondo tipo} di superficie con \(p = 2\) e \(p^{(1)} = 3\), non rientrante come caso particolare nella famiglia che ha per tipo la \(F_8\) bicanonica semplice. Per dimostrarlo si può ricorrere ad un computo di moduli: tanto la \(F_8\) bicanonica semplice come la \(F_4^2\) dipendono da \(24 = 9p - 2p^{(1)} + 12\) moduli, e perciò formano due famiglie distinte. Ma ciò appare anche in modo diretto, perchè se si cerca di ridurre per continuità la \(F'_8\), proiezione della \(F_8\) in \(S_3\), alla \(F'_4\) proiezione della \(F_4\) di Segre, contata due volte, si perviene, non già ad una quartica dotata di conica doppia irriducibile (siccome accade di regola pel nostro secondo tipo), bensì ad una quartica con due rette doppie infinitamente vicine: pertanto si hanno veramente \textit{due famiglie distinte di superficie con \(p = 2\) e \(p^{(1)} = 3\)}, aventi a comune una particolare sottofamiglia.
% The equality sign is printed at both the end of one line and the beginning of the next; both occurrences are retained.

\textit{Osservazione.} -- Anche qui si può completare l’analisi, esaminando l’ipotesi che il \textit{sistema canonico \(|K|\) sia riducibile.}

Anzitutto è lecito affermare che dovranno essere, in ogni caso, irriducibili le curve \(k\) componenti variabili delle \(K\), perchè altrimenti, essendo esse composte colle curve d’un fascio lineare (regolarità della superficie) si avrebbero almeno \(\infty^2\) \(K\) e \(p > 2\).

% scan: s170L
% !PRINTED 316
\weberpage{316}
% head: CAPITOLO OTTAVO
% stamp: Scattered scan specks, short marginal marks, a diagonal mark below an inequality and a dark right scan edge.
Ora si potrà supporre che le \(K\) posseggano, oltre le \(k\), delle componenti fisse \(\theta\), non eccezionali
\[
|K| = \Sigma\theta + |k|.
\]

Designando con \(\pi\) e \(n\) i caratteri di \(|k|\) e con \(\varrho\) e \(\nu\) quelli d’una \(\theta\), avremo
\[
2\pi - 2 - n + \Sigma(2\varrho - 2 - \nu) = p^{(1)} - 1 = 2,
\]
mentre \(\pi \geq 2\).

Inoltre, riferendoci ad una superficie senza curve eccezionali, sarà per ogni \(\theta\)
\[
2\varrho - 2 - \nu \geq 0.
\]

Siccome poi una \(k\) ha \(2\pi - 2 - 2n\) intersezioni colla residua curva \(\Sigma\theta\):
\[
2\pi - 2 - 2n \geq 0.
\]

Da queste diseguaglianze risultano pei valori di \(\pi\) ed \(n\) tre possibilità:
\weq{1)}{\pi = 2,\ n = 0;}
\weq{2)}{\pi = 2,\ n = 1;}
\weq{3)}{\pi = 3,\ n = 2.}

Ma l’ipotesi 1) si scarta subito. Il sistema \(|k'|\) aggiunto a \(|k|\) sarebbe un sistema \(\infty^3\) di curve del grado
\[
0 + 2 \cdot 2 + 2 = 6,
\]
formato di curve iperellittiche (bisecanti le \(k\)) e perciò condurrebbe ad una rigata cubica doppia, la quale non è normale, sicchè \(|k'|\) (che ha la dimensione \(p + 2 - 1 = 3\)) dovrebbe essere contenuto in un sistema completo di dimensione \(> 3\): ciò che è impossibile.

Anche l’ipotesi 2) (\(\pi = 2,\ n = 1\)) si scarta facilmente.

Perchè il sistema aggiunto \(|k'|\) che è un sistema di curve iperellittiche, di genere 4 e grado 7, dotato d’un punto base semplice, condurrebbe ancora ad una rigata cubica doppia non normale, onde, la sua dimensione sarebbe \(> 3\), ciò che è impossibile.

Resta l’ipotesi 3) (\(\pi = 3;\ n = 2\)). Ma in questa ipotesi il sistema \(|2k|\) ci conduce proprio alla superficie \(F_4\) doppia cui conduceva \(|K'| = |2K|\) nel supposto delle \(K\) irriducibili. Si ricade quindi sullo stesso piano doppio e, a posteriori, si riconosce che le \(K\) sono irriducibili: \(|K| = |k|\).

\edsection{19. Superficie con \(p^{(1)} = 3\) e \(p = 1\) (\(P = 4\)).}{19. Superficie con p\textasciicircum{}(1) = 3 e p = 1 (P = 4).}{sec.VIII.19}

Supposto che la curva canonica \(K\) sia irriducibile, il sistema bicanonico \(|2K|\), che è \(\infty^3\), sarà pure irriducibile e privo di punti base su
%CONT>>
% scan: s170R
% !PRINTED 317
\weberpage{317}
% head: SUPERFICIE REGOLARI CANONICHE E PLURICANONICHE
%<<CONT
% stamp: Scattered scan specks and faint edge artefacts.
\(K\). Se esso è semplice conduce a trasformare la superficie data in una \(F_8\) d’ordine 8, dello spazio \(S_3\), a sezioni piane di genere \(3p^{(1)} - 2 = 7\), e perciò dotata d’una curva doppia \(C_{14}\) d’ordine 14. Questa \(C_{14}\) appartiene ad \textit{una} superficie biaggiunta \(\Phi_7\), d’ordine \(2 \cdot 8 - 9 = 7\), che insieme ad un piano costituisce una \(\Phi_8\) secante su \(F_8\) una curva bicanonica; anzi la \(C_{14}\) è intersezione completa delle \(F_8\) e \(\Phi_7\left(\dfrac{7 \cdot 8}{4} = 14\right)\). D’altra parte deve esistere \textit{una} quartica aggiunta \(\Phi_4\), che sega su \(F_8\), fuori di \(C_{14}\), la \(K\) canonica d’ordine 4 (\(4 \cdot 8 - 2 \cdot 14 = 4\)). Ciò posto è facile caratterizzare la curva \(C_{14}\): essa è, in generale, la curva doppia della superficie \(\Phi_7\) a sezioni ellittiche, superficie possedente una quartica aggiunta \(\Phi_4\) e perciò razionale.
% The minus signs at the three printed line breaks occur on both adjoining lines; all occurrences are retained.

Pertanto \textit{la superficie bicanonica semplice \(F_8\) di generi \(p = 1\) e \(p^{(1)} = 3\) è definita, in generale, dal possesso di una curva doppia \(C_{14}\), d’ordine \(14\), che è la curva doppia di una superficie razionale \(\Phi_7\) del \(7^{\mathrm{o}}\) ordine a sezioni ellittiche}\origfootnote{1}{s170R-1}{Per le superficie a sezioni ellittiche cfr. p. es. \textsc{Enriques - Conforto}. \textit{Le superficie razionali}. Libro II, cap. III.}. Per costruire la \(F_8\) basterà combinare linearmente la quartica \(\Phi_4\) aggiunta a \(\Phi_7\), contata due volte, e la \(\Phi_7\) stessa, presa insieme con un piano generico.

Si può valutare il numero dei moduli da cui dipende la \(F_8\): si trova che la \(C_{14}\) (o la \(\Phi_7\)) possiede 12 invarianti proiettivi ed è doppia per \(\infty^4 F_8\), sicchè vi sono \(\infty^{16} F_8\) proiettivamente distinte: il numero dei moduli di \(F_8\) vale:
\[
M = 16 = 10p - 2p^{(1)} + 12.
\]

Ora nella famiglia delle \(F_8\) bicanoniche semplici con \(p = 1\) e \(p^{(1)} = 3\) si vedranno rientrare, come casi particolari, le superficie bicanoniche multiple, che --- almeno nell’ipotesi di una curva canonica irriducibile --- sono soltanto quartiche doppie \(F_4^2\), e danno luogo a due tipi di queste.

Se sopra una superficie \(F\) con \(p = 1\) e \(p^{(1)} = 3\) si ha una curva canonica \(K\) irriducibile, e il sistema bicanonico \(|2K|\), che sega su \(K\) la \(g_4^2\) completa, non sia semplice, bisognerà che esso appartenga ad un’involuzione del \(2^{\mathrm{o}}\) ordine \(I_2\) per cui la \(K\) sia: 1) curva doppia, ovvero: 2) curva iperellittica trasformata in se stessa. E siccome in nessun caso \(|2K|\) può avere dei punti base, si deduce che « le sole superficie bicanoniche multiple per \(p = 1\) e \(p^{(1)} = 3\), sono quartiche doppie \(F_4^2\) ».

Studiamo successivamente i due casi cui conducono le due ipotesi sopra indicate 1) e 2) in ordine all’appartenenza di \(|2K|\) ad una involuzione \(I_2\).

% scan: s171L
% !PRINTED 318
\weberpage{318}
% head: CAPITOLO OTTAVO
% stamp: Dark right scan edge, scattered specks, and oblique scan marks in the text and margins.
Anzitutto: se si assume l’ipotesi 1) la superficie bicanonica trasformata della data \(F\) sarà una quartica \(F_4\) doppia con quartica piana di diramazione \(K\). E perchè la \(K\) abbia il genere \(p^{(1)} = 3\), e anche perchè la \(F_4^2\) resulti di genere \(p = 1\), bisognerà che le sezioni piane di \(F_4\) non abbiano punti doppi, e che la \(F_4\) stessa sia di genere 1. Si può costruire una \(F_4\) di genere 1 doppia, rappresentante una \(F\) con \(p = 1\) e \(p^{(1)} = 3\), facendo fascio di una \(F_8\) bicanonica (coi detti generi) e della sua aggiunta \(F_4\) contata due volte; in questo fascio
\[
\lambda F_8 + F_4^2 = 0
\]
si vede la superficie generica (che è una \(F_8\)) ridursi per continuità (per \(\lambda = 0\)) alla \(F_4\) doppia. La \(F_4\) possiede come punti doppi i punti tripli di \(F_8\), e questi si aggiungono come punti (curve infinitesime) di diramazione alla curva di diramazione propria che è la la quartica \(K\) intersezione semplice di \(F_8\); la rimanente parte che è la \(C_{14}\), curva di contatto di \(F_4\) e di una \(\Phi_7\), costituisce una curva critica apparente. % [sic] la la -> la, the article is repeated across two printed lines.

Ora il numero dei punti tripli della \(C_{14}\), curva doppia di \(F_8\) (e anche di \(\Phi_7\)), si può valutare con la nota formula, che dà
\[
t = 10.
\]

Questo numero si ritrova anche in altro modo dalla formula di corrispondenza che lega i generi numerici delle due superficie di genere 1, \(F_4\) e \(F\), in corrispondenza \([1, 2]\).

Da ciò che precede risulta esistere, come caso particolare della \(F_8\) bicanonica con \(p = 1\) e \(p^{(1)} = 3\), una \(F_4\) di genere 2 doppia possedente come elementi di diramazione una quartica sezione piana e 10 punti conici. % [sic] genere 2 -> genere 1, this quartic is described as genus 1 above and on the following page.
Tuttavia si potrebbe pensare che le \(F_4\) doppie innanzi costruite rappresentino soltanto una sottofamiglia di un tipo più generale di superficie bicanoniche doppie con \(p = 1\) e \(p^{(1)} = 3\), le quali non rientrino nella famiglia delle \(F_8\). Per risolvere questo dubbio, bisogna approfondire lo studio delle nostre \(F_4^2\). Esse possono definirsi estraendo su \(F_4\) la radice quadrata d’un polinomio di \(8^{\mathrm{o}}\) ordine, il quale può ridursi ad \(\alpha\Phi_7\), essendo \(\alpha = 0\) il piano della \(K\) di diramazione e \(\Phi_7\) la superficie del \(7^{\mathrm{o}}\) ordine che ha come curva doppia la \(C_{14}\) (curva critica apparente). A codesto polinomio \(\alpha\Phi_7\) si può sostituire un altro polinomio \(\alpha F_7\), facendo variare liberamente la superficie entro il sistema delle \(F_7\), d’ordine 7, che passano triplamente per i 10 punti doppi di \(F_4\) e la toccano lungo una curva del sistema completo \(\lvert C_{14}\rvert\). Anzi si può ampliare questo sistema sommandovi due volte gli intorni dei 10 punti conici, prendendo dunque al posto di una superficie \(F_7\) che passi semplicemente per codesti 10 punti e tocchi \(F_4\) lungo la curva critica apparente del \(14^{\mathrm{o}}\) % ?? uncertain: the apostrophe in “gl’intorni” -- “gl’intorni” or “glintorni”.
% EDITORIAL D0163 RISOLTO: The enlarged image visibly contains a complete lowercase i after gl, followed by a space and the initial i of intorni. The printed reading is neither of the alternatives offered in the original doubt.
%CONT>>
% scan: s171R
% !PRINTED 319
\weberpage{319}
% head: SUPERFICIE REGOLARI CANONICHE E PLURICANONICHE
%<<CONT
% stamp: Dark left scan edge, scattered specks, and a faint vertical dotted scan mark near the right margin.
ordine (che torniamo a chiamare \(C_{14}\)). Ed è poi lecito sommare o sottrarre ad una tale \(F_7\), una superficie d’ordine pari \(2n\), che tocchi la \(F_4\) secondo una curva \(C_{4n}\), sostituendo così alla \(C_{14}\) una curva critica apparente \(C_{14} \pm C_{4n}\). Più precisamente vogliamo dimostrare che si può togliere dal sistema \(\lvert C_{14}\rvert\), segato su \(F_4\) dalle \(F_7\) tangenti che passano per i suoi 10 punti conici, due quartiche sezioni piane (ciascuna contata due volte), in guisa da definire la \(F_4^2\) mediante l’estrazione di una \(\sqrt{\alpha F_3}\); dove \(F_3\) eguagliato a zero rappresenti una cubica tangente a \(F_4\) lungo una sestica \(C_6\) di genere 3; quest’ultima potrà dunque assumersi come curva critica apparente della \(F_4^2\).

A tal uopo osserviamo che il sistema lineare delle \(C_{28}\) segate su \(F_4\) dalle \(F_7\) per 10 punti conici ha il grado
\[
4 \cdot 7^2 - 2 \cdot 10,
\]
e quindi il sistema metà delle \(C_{14}\) segate dalle \(F_7\) tangenti è di grado
\[
7^2 - 5 = 44,
\]
e perciò ha il genere e la dimensione \(\dfrac{44 + 2}{2} = 23\). Ora il sistema \(\lvert C_{14}\rvert\)-sega sopra una quartica piana \(C_4\) una serie \(g_{14}^{11}\), sicchè togliendo la \(C_4\) da \(\lvert C_{14}\rvert\) si otterrà un sistema \(\lvert C_{10}\rvert\) di dimensione e genere \(23 - 12 = 11\). E poichè questo \(\lvert C_{10}\rvert\) sega su un’altra quartica sezione piana una \(g_{10}^7\), togliendo anche questa curva si avrà un sistema lineare \(\infty^3\) di sestiche \(C_6\) di genere 3, curve di contatto di cubiche passanti per 10 punti conici di \(F_4\). % [sic] \lvert C_{14}\rvert-sega -> \lvert C_{14}\rvert sega, an extraneous hyphen separates the subject from the verb.

In conclusione: \textit{la bicanonica con \(p = 1\) e \(p^{(1)} = 3\), rappresentata sopra una superficie del \(4^{\mathrm{o}}\) ordine \(F_4\) di genere 1 doppia, è definita in relazione ad una \(F_4\) che sia toccata lungo una sestica \(C_6\) di genere 3 da una superficie del \(3^{\mathrm{o}}\) ordine \(F_3\) ed abbia quindi su questa \(C_6\) 10 punti conici; la \(F_4^2\) si ottiene estraendo su \(F_4\) la \(\sqrt{\alpha F_3}\)}, dove \(\alpha\) designa il piano della quartica di diramazione.

Resta da verificare che la \textit{detta \(F_4^2\) rientra come caso particolare nella famiglia delle bicanoniche \(F_8\) (\(p = 1,\ p^{(1)} = 3\)) semplici.}

Perciò basta calcolare il numero degli invarianti o moduli da cui dipendono le nostre \(F_4^2\). Anzitutto si osserverà che « l’esistenza di una sestica di genere 3, \(C_6\), autoresidua rispetto al sistema segato dalle superficie cubiche, porta per la \(F_4\) 10 condizioni e il possesso di 10 punti doppi (in generale conici) sopra la \(C_6\) ». Infatti l’esistenza di una \(C_6\) di genere 3 sopra \(F_4\) porta intanto \textit{una} condizione per \(F_4\); soddisfatta questa condizione si avranno \(\infty^6\) cubiche \(F_3\) secanti su \(F_4\) due sestiche \(C_6\) dello stesso genere 3, residua l’una dell’altra; e due sestiche siffatte --- formando insieme una curva \(C_{12}\) di genere 19 --- avranno 14 punti comuni; ma perchè le \(C_6\) residue siano equi%
%CONT>>
% scan: s172L
% !PRINTED 320
\unskip\weberpage{320}%
% head: CAPITOLO OTTAVO
%<<CONT
% stamp: Dark right scan edge, scattered specks, and an oblique scan mark across “proiettivi”.
% The word “equivalenti” is joined across the page boundary without retaining the printed line-end hyphen.
\unskip valenti bisogna che fra i 14 punti suddetti vi siano 4 punti variabili e 10 punti base, che cadranno necessariamente in 10 punti doppi di \(F_4\); e l’esistenza di un punto base, punto fisso per la serie \(g_{14}^{11}\) segata dalle \(C_6\) di un sistema su una \(C_6\) residua, porta \textit{una} condizione, ma l’esistenza di 10 punti base aventi come serie residua la \(g_4^2\) canonica, porta soltanto 9 condizioni. In tutto, dunque, le condizioni cui la \(F_4\) deve soddisfare per possedere una \(C_6\) autoresidua di genere 3 sono in numero di 10 e portano di conseguenza 10 punti doppi di \(F_4\) su \(C_6\).

% The minus sign in 34 - - 10 is printed at both sides of a line break and is retained twice.
Ciò posto le \(F_4\) possedenti una \(C_6\) autoresidua dipendono da \(34 - 10 = 24\) parametri; a questi vanno aggiunti i 3 parametri da cui dipende la scelta di una sezione piana \(K\), e tolti i 15 parametri di una omografia spaziale: resta che le nostre \(F_4^2\) dipendono da. % ?? uncertain: punctuation after “da” -- a full stop or a scan speck.
\[
M = 24 + 3 - 15 = 12
\]
invarianti proiettivi o moduli. E poichè
\[
M < 9p - 2p^{(1)} + 12 = 15;
\]
si deduce che la famiglia delle dette \(F_4^2\) è contenuta in una più ampia famiglia di \(F_8\) bicanoniche semplici, e precisamente in quella che sopra abbiamo definita.

C’è un altro tipo di bicanonica \(F_4^2\) rispondente ad una \(F\) coi generi \(p = 1\) e \(p^{(1)} = 3\), e con curva canonica (irriducibile) iperellittica, e si ottiene prendendo come doppia una quartica \(F_4\) a sezioni di genere 2 (cioè dotata di retta doppia) con una curva di diramazione \(C_{12}\) d’ordine 12; ed anche questa \(F_4^2\) rientra come caso particolare nella famiglia delle \(F_8\) bicanoniche semplici.

Rappresentando la \(F_4\) sopra un piano, si ottiene un \textit{piano doppio con curva di diramazione \(C_{10}\) d’ordine 10, dotata d’un punto 4-plo \(O\) e di 4 punti \([3, 3]\): \(A_1A_2A_3A_4\)}. Difatti su questo piano doppio c’è \textit{una} curva canonica \(K\) di genere \(p^{(1)} = 3\), che ha per immagine la conica per \(OA_1A_2A_3A_4\), e un sistema di curve bicanoniche, rappresentate doppiamente da \(\infty^3\) quartiche \(C_4\) con \(O\) doppio e tangenti in \(A_1A_2A_3A_4\), le quali formano il sistema rappresentativo delle sezioni piane della \(F_4\).

Si riconosce che \textit{questo piano doppio rientra come caso particolare nella famiglia delle bicanoniche semplici \(F_8\)} (con \(p = 1\) e \(p^{(1)} = 3\)) perchè la sua \(C_{10}\) di diramazione dipende da
\[
65 - 8 - 36 - 8 = 13
\]
invarianti proiettivi o moduli, ed è
\[
13 < 9p - 2p^{(1)} + 12 = 15.
\]

Possiamo riassumere i risultati della discussione fatta enunciando che: \textit{tutte le superficie con \(p = 1\) e \(p^{(1)} = 3\)}, almeno se la curva ca%
%CONT>>
% scan: s172R
% !PRINTED 321
\unskip\weberpage{321}%
% head: SUPERFICIE REGOLARI CANONICHE E PLURICANONICHE
%<<CONT
% stamp: Dark left scan edge and scattered specks.
% The word “canonica” is joined across the page boundary without retaining the printed line-end hyphen.
\unskip nonica è irriducibile, \textit{rientrano nella famiglia più generale delle bicanoniche semplici \(F_8\) dello spazio ordinario.} Sembra anche che l’eventuale riducibilità della curva canonica non dia luogo ad alcun caso nuovo. Ma non c’indugeremo in un esame minuto della questione.

\edsection{20. Superficie con \(p = 0\) e \(p^{(1)} = 3\).}{20. Superficie con p = 0 e p\textasciicircum{}(1) = 3.}{sec.VIII.20}

% The equals sign and the minus sign repeated at printed line breaks in this paragraph are both retained.
Qui le curve bicanoniche formanti una rete di grado \(4(p^{(1)} - 1) = 8\), conducono in generale ad un piano multiplo d’ordine 8, perciò conviene cercare il tipo delle date superficie ricorrendo alla \textit{superficie tricanonica, la quale è una \(F_{18}\) d’ordine 18 dello spazio a \(3p^{(1)} - 2 = 7\) dimensioni.}

% The minus sign in the sectional-genus expression is printed at both sides of a line break and is retained twice.
I caratteri elevati di questa superficie rendono difficile di costruirla direttamente nello \(S_7\), o di definirne precisamente una proiezione \(F'_{18}\) nello \(S_3\): la quale dovrà avere sezioni piane di genere \(6p^{(1)} - 2 = 16\), e perciò possiederà, in generale, una curva doppia \(C_{120}\) d’ordine 120. Ma possiamo almeno dimostrare l’\textit{esistenza della \(F_{18}\)}, perchè si conosce un caso particolare della famiglia, quale è un piano doppio di Campedelli, definito da una curva di diramazione \(C_{10}\) d’ordine 10 con 6 punti \([3, 3]\), non appartenenti ad una conica.

È facile rendersi conto che il detto piano doppio conduce ad una superficie tricanonica \(F_{18}\) semplice. Ma questa non è la \(F_{18}\) più generale, perchè la \(C_{10}\) di diramazione sopra nominata dipende da 3 invarianti proiettivi o moduli, e si ha
\[
3 < 9p - 2p^{(1)} + 12 = 6.
\]

Aggiungasi che, almeno se la rete delle curve bicanoniche di \(F\) (superficie regolare) è irriducibile, la \textit{tricanonica corrispondente (\(p = 0\) e \(p^{(1)} = 3\)) riesce certo semplice.} Infatti una superficie tricanonica multipla (d’ordine \(\dfrac{18}{2} = 9\)) potrebbe nascere soltanto nell’ipotesi che le curve bicanoniche \(K_2\) di \(F\) siano iperellittiche, risultando allora composta la serie canonica (completa) che le curve tricanoniche segano sopra una \(K_2\). In questa ipotesi la serie caratteristica della rete bicanonica su una generica curva \(K_2\), sarà una \(g_8^2\) composta colla \(g_2^1\) che appartiene alla \(K_2\) stessa, e perciò contenuta in una \(g_8^4\) completa; ma (se la superficie deve essere regolare) ciò orta \(P = 6\) anzichè \(P = 3\). % [sic] orta -> porta, the initial p is absent from the visible word.

\edsection{21. Superficie pluricanoniche semplici e multiple.}{21. Superficie pluricanoniche semplici e multiple.}{sec.VIII.21}

Nei precedenti paragrafi abbiamo incontrato diversi tipi di superficie che dan luogo a bicanoniche doppie, in ispecie i tre tipi generali seguenti:

% EDITORIAL: segnatura originale: \noindent\textsc{Enriques F.} - \textit{Superficie algebriche.}\hfill 21
\printersignature{\textsc{Enriques F.} - \textit{Superficie algebriche.}\quad 21}

% scan: s173L
% !PRINTED 322
\weberpage{322}
% head: CAPITOLO OTTAVO
% stamp: Dark right scan edge, scattered specks, and oblique scan marks across “avverta” and in the left margin.
I) Per \(p = 3\) e \(p^{(1)} = 3\), il piano doppio con curva di diramazione \(C_8\) d’ordine 8;

II) Per \(p = 2\) e \(p^{(1)} = 2\), il piano doppio con curva di diramazione \(C_{10}\) d’ordine 10 dotata di due punti quintupli infinitamente vicini;

III) Per \(p = 2\) e \(p^{(1)} = 3\), il piano doppio con curva di diramazione \(C_{12}\) d’ordine 12, dotata di due punti quadrupli e di due punti \([3, 3]\).

Per i tipi I) e II) non soltanto la bicanonica sì anche la tricanonica risulta una superficie doppia, in quanto le curve bicanoniche (di cui le tricanoniche sono le aggiunte) sono curve iperellittiche. % ?? uncertain: accent on “si” -- “si” or “sì”.
% EDITORIAL D0165 RISOLTO: The enlarged scan shows a grave-accent stroke above the i in sì, distinguishable from the nearby ordinary i dot. The conjunction is printed with the accent.

Ora ci proponiamo di dimostrare che questi tipi I) e II) sono i soli tipi di superficie, con \(p > 0\) e \(p^{(1)} > 1\), le cui tricanoniche siano superficie multiple.

A tale scopo si avverta che, essendo le curve bicanoniche o le parti variabili di esse curve irriducibili (cfr. \S~19, Cap. IV), le aggiunte e quindi le tricanoniche, o le loro parti variabili, formeranno certo un sistema semplice se le dette bicanoniche non sono iperellittiche (e quindi appartenenti ad una involuzione razionale \(I_2\)); se dunque una superficie tricanonica regolare (per \(p > 0\) e \(p^{(1)} > 1\)) non è semplice, essa deve essere rappresentabile su un piano doppio per modo che le bicanoniche (iperellittiche) abbiano per immagini le curve razionali di un sistema lineare \(\infty^2\) almeno, cui vada sommata eventualmente qualche componente fissa.

Conviene esaminare successivamente i tre casi in cui le immagini delle dette bicanoniche si riducono a:

1) le \(\infty^2\) rette del piano;

2) le \(\infty^5\) coniche del piano;

3) o, infine, alle curve \(C_n\) d’un certo ordine \(n\) (\(> 1\)) formanti un sistema lineare di dimensione \(2n - s\) con un punto base \((n - 1)\)-plo \(O\) e senz’altri punti base (\(s = 0\)); ovvero con un altro punto base semplice \(O_1\) (\(s = 1\)), o con \(s\) (\(> 1\)) punti base semplici \(O_1O_2\ldots O_s\) infinitamente vicini ad \(O\) (in direzioni distinte).

L’ipotesi 1) --- che dà \(P = p + p^{(1)} = 3\), \(p = 1\) e \(p^{(1)} = 2\) --- è impossibile come appare dall’esame delle superficie con tali caratteri che abbiamo svolto nel \S~15.

L’ipotesi 2) conduce al piano doppio con \(C_8\) di diramazione, cioè al tipo I) con \(p = 3\), \(p^{(1)} = 3\), incontrato nel \S~17.

L’ipotesi 3) si realizza per \(n = 2\), quando si assuma come curva di diramazione del piano doppio una \(C_8\) d’ordine 8 con due punti tripli infinitamente vicini \(O\) e \(O_1\): le immagini delle curve bicanoniche sono le coniche per \(O\) e \(O_1\).

% scan: s173R
% !PRINTED 323
\weberpage{323}
% head: SUPERFICIE REGOLARI CANONICHE E PLURICANONICHE
% stamp: Dark left scan edge, scattered specks, and faint vertical scan marks near the right margin.
Ma il piano doppio così definito rientra nel tipo più generale dove si assuma, come curva di diramazione una \(C_{10}\) d’ordine 10 con due punti 5-pli infinitamente vicini, che è il tipo II) per \(p = 2\) e \(p^{(1)} = 2\): quando la \(C_{10}\) abbia altresì un punto doppio, si riconduce quadraticamente ad una \(C_8\) con punto \([3, 3]\).

Per discutere in generale le possibilità cui dà luogo l’ipotesi 3), pongasi che le bicanoniche del nostro piano doppio contengano una parte fissa d’ordine \(d\), da sommarsi alle \(C_n\); affinchè l’addizione di questa parte non ampli il sistema \(\lvert C_n\rvert\) bisogna che essa consti di rette fondamentali per \(C_n\), contate rispettivamente \(h_1h_2\ldots h_s\) volte, cioè di rette \(OO_i\) e quindi che passi colla molteplicità \(d\) per \(O\) e colla molteplicità \(h_i\) per \(O_i\). Così la curva di diramazione del piano doppio sarà una \(C_{2m}\) d’ordine \(2m\) che ha le curve \(C_n + \Sigma h_i \cdot OO_i\), d’ordine \(n + d = n + \Sigma h_i\), come seconde aggiunte; ed avremo
\[
2m = n + d + 6.
\]

Oltre a ciò basterà ritenere che la \(C_{2m}\) deve possedere in \(O\) la molteplicità \(n - 1 + d + 2 = n + d + 1 = 2m - 5\). Ora sopra un piano doppio con curva di diramazione \(C_{2m}\) dotata di un punto \(O\) \((2m - 5)\)-plo, le immagini delle curve bicanoniche sono in generale curve d’ordine \(2m - 6 = n + d\) passanti per \(O\) colla molteplicità
\[
2m - 8 = n + d - 2;
\]
affinchè questa molteplicità cresca di 1, diventando \(n + d - 1\) (come accade per le \(C_n + \Sigma h_i \cdot OO_i\)), bisogna che vicino ad \(O\) si trovi un altro punto \((2m - 5)\)-plo di \(C_{2m}\); ma la somma
\[
2m - 5 + 2m - 5 \leq 2m,
\]
e quindi
\[
m = 5,\qquad 2m = 10
\]
o (in particolare)
\[
m = 4,\qquad 2m = 8;
\]
così si ritrova il tipo II).

In conclusione enunceremo il teorema:

\textit{La tricanonica di una superficie regolare di generi \(p > 0\) e \(p^{(1)} > 1\) è sempre una superficie semplice, d’ordine \(\leq 9(p^{(1)} - 1)\), fatta eccezione per due tipi I) e II) di superficie riferibili al piano doppio, che rispondono ai valori \(p = 3\) e \(p^{(1)} = 3\), e \(p = 2\) e \(p^{(1)} = 2\).}

Per questi due tipi I) e II) risultano doppie tanto la bicanonica che la tricanonica: la bicanonica è per la I) una superficie di Veronese (\(F_4\) di \(S_5\)) doppia e per la II) un cono quadrico doppio dello \(S_3\); la tricanonica è, nel primo caso, una superficie doppia \(F_9\), d’ordine 9 a sezioni ellittiche in \(S_9\), nel secondo caso una superficie doppia \(F_6\), d’ordine 6 a sezioni ellittiche in \(S_6\).

% scan: s174L
% !PRINTED 324
\weberpage{324}
% head: CAPITOLO OTTAVO
% stamp: Dark right scan edge, scattered specks, and oblique marks in the margins and the lower blank area.
Siccome dunque le curve tricanoniche delle nostre superficie non sono, in nessun caso, iperellittiche, ne segue (per \(p > 0\)) che le loro aggiunte non appartengono ad una involuzione, cioè: \textit{sopra una superficie regolare per \(p > 0\) e \(p^{(1)} > 1\), le curve quadricanoniche} (o le componenti variabili di esse), \textit{e quindi anche le \(i\)-canoniche per \(i > 4\) formano sempre un sistema semplice.}

Questa proposizione affermante la semplicità delle superficie pluricanoniche (\(p^{(1)} > 1\)) per \(i > 3\), si estenderà verosimilmente al caso \(p = 0\); ma occorre per ciò una analisi più difficile, in cui non c’indugeremo (cfr. per \(p^{(1)} = 2\) la discussione del \S~16). % ?? uncertain: the small mark following “cfr.” -- a scan speck or a second printed stop.

\clearpage\thispagestyle{plain}
% scan: s174R
% !PRINTED none
\weberpage{[s174R]}
% head: none
% stamp: Handwritten “325” in the upper-right margin; scattered specks and a thin scan line near the lower-right edge.
\edchapter{IX}{SUPERFICIE IRREGOLARI E SISTEMI CONTINUI\\DI CURVE DISEQUIVALENTI}{Capitolo IX. SUPERFICIE IRREGOLARI E SISTEMI CONTINUI DI CURVE DISEQUIVALENTI}{SUPERFICIE IRREGOLARI E SISTEMI CONTINUI DI CURVE DISEQUIVALENTI}

\edsection{1. Introduzione.}{1. Introduzione.}{sec.IX.1}

Richiamiamo alcune osservazioni che abbiamo fatto nel Cap. IV, \S~12.

E anzitutto avvertiamo che, parlando di curve e di sistemi lineari o continui di curve \(C\), sopra una superficie \(F\), sottintenderemo di regola che si tratti di \textit{curve irriducibili prive di punti multipli}\origfootnote{1}{s174R-1}{Se si vuole considerare anche curve \(C\) con punti multipli, conviene ritenere che questi punti siano punti base assegnati per i sistemi lineari e continui a cui ci si riferisce.}.

Abbiamo veduto che, sopra una superficie regolare, un sistema lineare completo di curve \(\lvert C\rvert\), ha la serie caratteristica completa. Ed una superficie cui appartenga un sistema continuo di curve \(C\), formato di \(\infty^d\) sistemi lineari distinti, cioè un sistema \(\{C\}\) contenente \(\infty^d\), con \(d > 0\), curve disequivalenti, è irregolare, d’irregolarità
\[
p_g - p_a \geq d.
\]

L’affermazione si giustifica senz’altro per sistemi continui \(\{C\}\) di curve irriducibili \(\infty^{r+d}\), formati di \(\infty^d\) sistemi lineari che hanno una certa dimensione (generica) \(r\): giacchè la serie caratteristica segata su una \(C\) generica dalle curve infinitamente vicine di \(\{C\}\) avrà la dimensione
\[
r + d - 1,
\]
e quindi la serie caratteristica del sistema lineare \(\lvert C\rvert\) avrà la deficienza (almeno) eguale a \(d\); basta quindi invocare il teorema di Castelnuovo (Cap. IV, \S~12) per cui
\[
d \leq p_g - p_a.
\]

Aggiungasi che questa diseguaglianza vale anche per sistemi continui \(\{C\}\) formati di curve riducibili; perchè, in primo luogo,
%CONT>>
% scan: s175L
% !PRINTED 326
\weberpage{326}
% head: CAPITOLO NONO
%<<CONT
% stamp: Dark right scan edge, scattered specks, and oblique scan marks near “abbordare” and in the left margin.
si può prescindere da eventuali componenti fisse che facciano parte delle \(C\), ed in secondo luogo, se le \(C\) variabili siano riducibili, sottraendole da un sistema lineare abbastanza grande \(\lvert L\rvert\) si potrà ottenere un sistema continuo \(\{L - C\}\) che --- all’infuori di componenti fisse --- verrà costituito di curve irriducibili e conterrà almeno --- come \(\{C\}\) --- \(\infty^d\) curve disequivalenti.

Scopo principale di questo Capitolo è di esporre le argomentazioni che portano ad invertire il teorema precedente, nel senso che: sopra una superficie irregolare, di generi \(p_a\) e \(p_g > p_a\), e quindi d’irregolarità \(q = p_g - p_a\), esistono sempre sistemi continui di curve formati di \(\infty^q\) sistemi lineari distinti.

Il possesso di sistemi continui di curve disequivalenti costituirà così la \textit{proprietà caratteristica delle superficie irregolari.} E il teorema che l’enuncia e la precisa come sopra, in ordine al valore di \(p_g - p_a\), si potrà chiamare il \textit{teorema fondamentale} nella teoria delle superficie irregolari.

Ma prima di abbordare la dimostrazione del teorema fondamentale, in vista delle difficoltà d’ordine delicato che essa presenta, conviene mettersi nelle ipotesi più generali in cui se ne faccia astrazione, ed esaminare dunque che cosa possa dirsi delle superficie contenenti sistemi continui \(\infty^q\) di curve disequivalenti, quando si sappia soltanto che
\[
0 \leq q \leq p_g - p_a,
\]
come si è riconosciuto innanzi. A tal uopo distingueremo due irregolarità della superficie, a priori non sempre eguali: \textit{l’irregolarità numerica}
\[
p_g - p_a
\]
e \textit{l’irregolarità geometrica} \(q\), cioè la massima dimensione d’una serie continua di curve disequivalenti che si trovi sulla superficie:
\[
q \leq p_g - p_a.
\]

Questo numero \(q\) designa, in ogni caso, un carattere invariante della superficie stessa: tutti gli esempi noti di superficie irregolari (superficie con un fascio irrazionale di curve, superficie delle coppie di punti d’una curva ecc.) danno \(q > 0\) ed anzi \(q = p_g - p_a\).

\edsection{2. Condizione aritmetica perchè una curva, sopra una superficie d’irregolarità geometrica \(q\), appartenga ad una serie continua \(\infty^q\) di curve disequivalenti.}{2. Condizione aritmetica perchè una curva, sopra una superficie d’irregolarità geometrica q, appartenga ad una serie continua ∞\textasciicircum{}q di curve disequivalenti.}{sec.IX.2}

Sia \(q > 0\) l’irregolarità geometrica d’una superficie \(F\): vuol dire che ad \(F\) appartiene una serie continua \(\infty^q\) \(\{D\}\) e non una serie più ampia, di curve disequivalenti.

% scan: s175R
% !PRINTED 327
\weberpage{327}
% head: SUPERFICIE IRREGOLARI E SISTEMI CONTINUI DI CURVE, ECC.
% stamp: Dark left scan edge, scattered specks, and faint vertical scan marks near the right margin.
Si domanda: può darsi un criterio aritmetico perchè una curva \(C\), su \(F\), appartenga similmente ad una serie \(\infty^q\) di curve disequivalenti, ossia ad una serie formata di \(\infty^q\) sistemi lineari completi distinti?

La risposta è fornita dal calcolo della dimensione virtuale del sistema individuato dalla curva \(C\) (di genere \(\pi\), grado \(n\), e indice di specialità \(i\)):
\[
r = p_a + n - \pi + 1 - i.
\]

Invero sia \(\lvert C\rvert\) un sistema lineare di curve, irriducibili o riducibili, di grado \(n\), di genere \(\pi\), e di indice di specialità \(i\), \textit{aritmeticamente effettivo}, cioè tale che la sua dimensione virtuale
\[
p_a + n - \pi + 1 - i \geq 0;
\]
ad esso associeremo un altro sistema lineare regolare \(\lvert D\rvert\) (irriducibile ecc.) abbastanza grande, per modo che anche \(\lvert C + D\rvert\) sia regolare. Il sistema \(\lvert D\rvert\) potrà supporsi contenuto in un sistema continuo \(\{D\}\) formato di \(\infty^q\) sistemi lineari disequivalenti, uno dei quali --- diverso da \(\lvert D\rvert\) --- vogliamo designare con \(\lvert\overline{D}\rvert\).

Ora il sistema \(\lvert C\rvert\) si lascia costruire staccando una curva \(D\) da \(\lvert C + D\rvert\), e la sua dimensione risulta \(\geq p_a + n - \pi + 1 - i\), giacchè \(\lvert C + D\rvert\) sega sulla detta \(D\) una serie (completa o meno) che ha precisamente l’indice di specialità \(i\) (cfr. Cap. IV, \S~11).

% The minus sign before the final barred D is printed at both sides of a line break and is retained twice.
Consideriamo la serie segata dallo stesso sistema \(\lvert C + D\rvert\) su una \(\overline{D}\). Non è possibile che, per una scelta generica di \(\overline{D}\) in \(\{D\}\), la detta serie abbia un indice di specialità superiore ad \(i\), perchè --- variando la \(\overline{D}\) con continuità --- si può ridurla alla \(D\) e, in questo passaggio al limite, l’indice di specialità della serie non può diminuire. Pertanto la dimensione del sistema lineare \(\lvert\overline{C}\rvert = \lvert C + D - \overline{D}\rvert\) risulterà (secondo il noto calcolo che porge il teorema di Riemann-Roch) \(\geq p_a + n - \pi + 1 - i\): così, al variare di \(\overline{D}\), si otterranno \(\infty^q\) sistemi lineari completi distinti, formanti una serie irriducibile \(\{C\}\).

Enunciamo il teorema così conseguito:

\textit{Sopra una superficie di irregolarità geometrica \(q > 0\), un sistema lineare aritmeticamente effettivo, cioè di grado \(n\), genere \(\pi\) e indice di specialità \(i\), tale che}
\[
p_a + n - \pi + 1 - i \geq 0,
\]
\textit{appartiene sempre ad una serie \(\infty^q\) di sistemi lineari disequivalenti.}

Si vede qui la più semplice definizione della irregolarità geometrica (e poi della differenza \(p_g - p_a\) in rapporto al teorema fondamentale che porge la proprietà caratteristica delle superficie irregolari).

% scan: s176L
% !PRINTED 328
\weberpage{328}
% head: CAPITOLO NONO
% stamp: Dark vertical scan-edge band at the right and scattered scan specks.
\edsection{3. La varietà di Picard corrispondente ad una superficie irregolare.}{3. La varietà di Picard corrispondente ad una superficie irregolare.}{sec.IX.3}

Nel precedente paragrafo si è visto come, sopra una superficie d’irregolarità geometrica \(q\ (>0)\), gli \(\infty^q\) sistemi lineari formanti un certo sistema continuo \(\{|C|\}\) o \(\{C\}\) siano legati a quello di un altro sistema continuo analogo \(\{D\}\): vi è una corrispondenza biunivoca fra gli elementi (sistemi lineari) di \(\{C\}\) e quelli di \(\{D\}\), essendo omologhi i sistemi lineari residui l’uno dell’altro rispetto a un \(|C+D|\). Questa osservazione (che appartiene a \textsc{Castelnuovo}) si può precisare dicendo: \textit{Gli elementi (sistemi lineari) delle serie continue \(\infty^q\) appartenenti ad una superficie di irregolarità geometrica \(q\ (>0)\), si possono ritenere come i punti di una medesima varietà abeliana \(V\)}, che \textsc{Castelnuovo} ha voluto intitolare al nome del geometra francese, come \textit{varietà di Picard}, corrispondente alla superficie.

Invero l’operazione \(+D-\overline{D}\) per cui si passa da un sistema lineare \(|C|\) ad un altro sistema lineare \(|\overline{C}|\) contenuto in \(\{C\}\), si può interpretare come una trasformazione puntuale della varietà \(V\) in sè stessa; è ovvio che tutte le trasformazioni analoghe sono permutabili e formano un gruppo \(\infty^q\) semplicemente transitivo senza eccezioni entro la detta varietà \(V=V_q\) (a \(q\) dimensioni): che appunto perciò viene caratterizzata quale \textit{varietà abeliana}, non degenere.

Il significato del teorema conseguito viene chiarito dalle considerazioni seguenti.

Si ponga il problema di determinare tutte le curve d’un dato ordine \(n\) appartenenti ad una superficie algebrica \(F\).

A priori queste curve che — per appartenere alla \(F\) — vengono assoggettate a condizioni algebriche, si ripartiranno in un certo numero di famiglie e perciò costituiranno, in generale, dei sistemi continui algebrici.

Tali sistemi saranno sistemi lineari se la \(F\) è regolare. Se, invece, la \(F\) è di irregolarità geometrica \(q\) \((p_g>p_a)\) codesti sistemi — riguardati come varietà aventi per elementi dei sistemi lineari distinti — saranno, non già varietà algebriche di tipo qualsiasi, ma varietà abeliane \(V_q\) o, almeno, formeranno coi loro multipli, delle varietà abeliane.

A meglio comprendere che cosa ciò importi varranno alcuni richiami che gioverà fare appunto sulle varietà abeliane (cfr. § 9).

\edsection{4. Proprietà caratteristica delle superficie algebriche irregolari: teorema fondamentale per \(p_g=0\).}{4. Proprietà caratteristica delle superficie algebriche irregolari: teorema fondamentale per p\_g=0.}{sec.IX.4}

La dimostrazione del teorema fondamentale sulle superficie irregolari, che abbiamo enunciato nel § 2, si dà nel modo più sem%
%CONT>>
% scan: s176R
% !PRINTED 329
\unskip\weberpage{329}%
% head: SUPERFICIE IRREGOLARI E SISTEMI CONTINUI DI CURVE, ECC.
%<<CONT
% stamp: Scattered scan specks and faint vertical marks in the outer right margin.
plice per le superficie di genere \(p_g=0\), e richiede invece considerazioni più delicate per \(p_g>0\). Perciò conviene distinguere questi due casi.

Supponiamo dapprima che sia \(p_g=0\) e quindi \(p_a=-p<0\).

Si consideri, sopra la nostra superficie \(F\), un sistema lineare completo \(|C|\), che sia regolare, non speciale, di genere \(\pi\) e di grado \(n\), e quindi di dimensione
\[
r=-p+n-\pi+1.
\]

E ad esso si associ un altro sistema lineare \(|D|\), che sia pure regolare, di genere \(\varrho\) e di grado \(\nu\), e quindi di dimensione
\[
r'=-p+\nu-\varrho+1,
\]
e tale che anche la somma \(|C+D|\) sia un sistema regolare (cfr. Cap. IV, § 8).

Designando con \(m\) il numero delle intersezioni di \(C\) e \(D\), scriveremo il genere e il grado di \(|C+D|=|E|\):
\[
\begin{aligned}
\Pi &= \pi+\varrho+m-1\\
N &= n+\nu+2m,
\end{aligned}
\]
e quindi la sua dimensione:
\[
R=-p+N-\Pi+1=-p+(n-\pi)+(\nu-\varrho)+m+2,
\]
cioè
\[
R=r+r'+p+m.
\]

Ricordiamo che, in base al \textit{principio di degenerazione o di spezzamento}\origfootnote{1}{s176R-1}{Cfr. \textsc{Enriques-Chisini}, L. V, Cap. III, 35 (vol. III, pag. 405).}, una curva irriducibile, dotata di un certo numero \(m\) di punti doppi, non può ridursi, per continuità, ad una curva spezzata le cui componenti, ove siano prive di punti multipli, non abbiano almeno \(m+1\) punti comuni (cioè con un punto doppio di più, attraverso a cui vengano connesse le due componenti della curva spezzata). Da questo principio si deduce che, imponendo alle curve \(E\) di possedere \(m\) punti doppi, non si potrà ottenere un sistema continuo \(\overline{E}\) di dimensione \(r+r'+p>r+r'\) formato di curve irriducibili, che contenga entro di sè il sistema delle curve di \(|E|=|C+D|\) spezzate in una curva di \(|C|\) e in una curva di \(|D|\), ma si verrà a definire un sistema di cui una parte almeno (della stessa dimensione) contenga entro di sè un sistema continuo di curve spezzate in una curva dell’ordine di \(C\) e in una curva dell’ordine di \(D\), queste due componenti descrivendo dunque due sistemi continui non lineari, diciamo \(\{C\}\) e \(\{D\}\), complementari.

% scan: s177L
% !PRINTED 330
\weberpage{330}
% head: CAPITOLO NONO
% stamp: Dark right scan edge and scattered specks; a diagonal stray mark appears near "da una".
Ad ogni curva \(C\) del più ampio sistema continuo \(\{C\}\) corrisponderà, come residuo rispetto ad \(|E|\), un sistema lineare di curve \(|D|\) e viceversa; sicchè è facile valutare la dimensione del sistema continuo \(\{C\}\), che dovrà essere (almeno) \(r+p\), e similmente quella di \(\{D\}\), che dovrà essere \(r'+p\). Anzi si può dire che \(\{C\}\) consterà di una serie (almeno) \(\infty^p\) di sistemi lineari disequivalenti, aventi la dimensione eguale ad \(r\), perchè la dimensione dei sistemi lineari completi formanti \(\{C\}\), che è a priori \(\ge r\), non può essere \(>r\), visto che codesti sistemi lineari si riducono per continuità al \(|C|\) da cui siamo partiti, che ha proprio la dimensione \(r\).

Ma nel caso che il genere geometrico di \(F\) sia \(p_g=0\), sappiamo che la \(F\) non può contenere un sistema continuo di curve formato di \(\infty^d\) sistemi lineari distinti, dove sia \(d>p_g-p_a=p\); pertanto si conclude che, sopra le superficie irregolari di genere geometrico \(p_g=0\), esistono sistemi continui di \(\infty^p\) curve disequivalenti, e non sistemi formati da una serie più ampia di curve siffatte.

Riassumeremo i risultati ottenuti enunciando il \textit{Teorema fondamentale ristretto: Una superficie di genere numerico negativo \(p_a=-p\) contiene certo sistemi continui formati di \(\infty^p\) curve disequivalenti, e, se il genere geometrico \(p_g=0\), non possiede sistemi più ampi di curve disequivalenti; in altri termini l’irregolarità geometrica vale}
\[
q\ge -p % [sic] q >= -p -> q >= p, the immediately preceding statement asserts a family of dimension p with p_a=-p.
\]
e per \(p_g=0\) \textit{eguaglia l’irregolarità numerica}
\[
p_g-p_a=p.
\]

\edsection{5. Teorema fondamentale per \(p_g>0\).}{5. Teorema fondamentale per p\_g>0.}{sec.IX.5}

Cerchiamo di estendere il teorema fondamentale alle superficie di genere geometrico \(p_g>0\) e di genere numerico qualunque \(p_a<p_g\). A tal uopo riprendiamo a considerare, sopra la superficie \(F\) (supposta di generi \(p_a\) e \(p_g\)) due sistemi lineari non speciali regolari \(|C|\) e \(|D|\) coi caratteri \(n,\pi\) e \(\nu,\varrho\) tali che anche la somma \(|E|=|C+D|\) sia regolare, segando su una \(D\) una serie completa e non speciale \(g_{m+\nu}\).
% line-break repetition: The equals sign after |E| is printed at both the end and the beginning of adjacent lines.

Designando con \(N\) e \(\Pi\) i caratteri di \(|E|\), si vede (come innanzi nel caso \(p_g=0\)) che il sistema lineare \(|E|\) avrà la dimensione
\[
\begin{aligned}
& R=p_a+N-\Pi+1=p_a+(n-\pi)+(\nu-\varrho)+m+2\\
& \qquad\qquad\qquad\qquad(m=CD),
\end{aligned}
\]
ossia
\[
R=r+r'+m-p_a,
\]
%CONT>>
% scan: s177R
% !PRINTED 331
\weberpage{331}
% head: SUPERFICIE IRREGOLARI E SISTEMI CONTINUI DI CURVE, ECC.
%<<CONT
% stamp: Dark left scan edge and scattered scan specks.
dove \(r\) e \(r'\) sono le dimensioni di \(|C|\) e \(|D|\). Ora, se il sistema lineare \(|C|\) deve essere contenuto in un sistema continuo \(\{C\}\) di dimensione \(r+q=r+p_g-p_a\) (formato di \(\infty^q\) sistemi lineari distinti), bisogna che l’imposizione alle \(E\) sopra \(F\) di \textit{\(m\) punti doppi}, onde risulta \textit{lo spezzamento di tali curve} in \(C\) di \(\{C\}\) e \(D\) del sistema residuo \(\{D\}\), porti — non già \(m\) condizioni indipendenti (o \(m-p_a\) come apparirebbe dall’esistenza di \(|C|\) e di \(|D|\)) — ma \textit{soltanto \(m-p_g\) condizioni}. Se ciò accade, il sistema continuo \(\{C\}\), di dimensione \(p_g+n-\pi+1\), avrà, su una \(C=C_0\), la serie caratteristica completa, e sarà quindi non ampliabile.

Così, per mostrare che « l’imposizione alle \(E\) su \(F\) di \(m\) punti doppi, onde risulta lo spezzamento di tali curve in curve di \(\{C\}\) e \(\{D\}\), non può portare più di \(m-p_g\) condizioni », e quindi l’irregolarità geometrica \(q\) di \(F\) eguaglia l’irregolarità numerica \((q=p_g-p_a)\), si è condotti ad esaminare la serie caratteristica del sistema continuo \(\{C\}\); pare che si raggiungerà la dimostrazione cercata ove si riconosca che codesta serie è completa.
% line-break repetition: The minus sign between p_g and p_a is printed at both the end and the beginning of adjacent lines.

E a tal uopo basterà considerare le curve \(E\) di \(|E|=|C+D|\) che passano pel gruppo \(G\) degli \(m\) punti comuni alle due componenti \(C_0\) e \(D_0\) di una \(E\) spezzata, \(E_0=C_0+D_0\); queste \(E\) segano su \(C_0\) la serie caratteristica completa, ed una di esse, insieme ad \(E_0\), determina un fascio lineare di curve che si toccano nei punti del \(G\): segue da ciò che anche la curva del fascio infinitamente vicina ad \(E_0\) possiede \(m\) punti doppi vicini a quelli del \(G\)\origfootnote{1}{s177R-1}{Cfr. p. es. \textsc{Enriques-Chisini}, \textit{Lezioni}, L. II, § 5, vol. I, pag. 183 e L. III, Cap. II, § 20, vol. II, pag. 167.}, e perciò, se anche a queste curve infinitamente vicine sia applicabile il principio di degenerazione, essa deve spezzarsi in due curve \(C_1\) e \(D_1\), la prima delle quali — vicina a \(C_0\) — sega su \(C_0\) un gruppo (arbitrariamente assegnabile) della serie caratteristica.

Abbiamo così raggiunto la dimostrazione (per \(p_g\) qualunque) del teorema fondamentale?

No, l’apparenza della dimostrazione non ci deve ingannare: il teorema di cui si tratta concerne una \textit{proprietà integrale}, e noi abbiamo stabilito soltanto una \textit{proprietà differenziale} che ne è una conseguenza, ma che viceversa non implica necessariamente la prima. Invero si è dimostrato tutt’al più che « la serie caratteristica di \(C_0\) viene segata completamente da curve \(C_1\), infinitamente vicine ad essa », ma non che le \(C_1\) appartengano veramente, come curve tendenti al limite, ad un sistema continuo di curve proprie \(\{C\}\).

C’è qui un punto critico della dimostrazione, di cui rileveremo più avanti il significato storico (§ 6): la condizione per cui « \(m-p_g\) punti doppi imposti alle \(E\) debbono portare di conseguenza altri
%CONT>>
% scan: s178L
% !PRINTED 332
\weberpage{332}
% head: CAPITOLO NONO
%<<CONT
% stamp: Dark right scan edge, scattered specks, and a diagonal stray stroke near "risulterà".
\(p_g\) punti doppi » si è verificata non « in grande » ma « in piccolo », a meno d’infinitesimi d’ordine superiore, per le curve infinitamente vicine ad \(E_0\).

La proprietà differenziale « in piccolo » consegue, come si è detto, dalla proprietà integrale o « in grande » che vuolsi stabilire, ma questa non è a sua volta conseguenza di quella.

Per chiarire meglio la cosa si consideri una varietà \(V_r\) di dimensione \(r\), la quale sia definita in uno spazio \(S_R\) \((R>r)\) come intersezione di un certo numero di ipersuperficie (o di falde lineari differenzialmente distinte di una stessa ipersuperficie algebrica): per valutare la dimensione \(r\) si è condotti a ricercare il numero \(s\) dei punti linearmente indipendenti nell’intorno del prim’ordine di un punto generico \(P\) di \(V_r\); ma è chiaro che le varietà di cui \(V_r\) è definita come intersezione possono avere in ogni punto di questa un certo contatto, ed in tal caso il numero dei punti vicini a \(P\), di cui si è detto, risulterà \(s>r\).

Per esempio, per \(r=1\) la \(V_1\) può essere la curva di contatto di due o più superficie dello spazio \(S_3\), curva da contare due volte; allora il numero dei punti indipendenti di essa, vicini ad un punto \(P\), non è più \(r=1\), bensì \(s=2\).

Da questo chiarimento risulta anche l’esigenza a cui si deve soddisfare per colmare la lacuna che rimane nella nostra dimostrazione del teorema fondamentale generale per \(p_g>0\), passando dalla proprietà riconosciuta in piccolo alla proprietà in grande, di cui si è detto sopra.

Si consideri la \(V_r\) dello \(S_R\), definita come intersezione di \(m\) falde ipersuperficiali, che rispondono alle curve di \(|E|=|C+D|\) dotate di \(m\) punti doppi (costituenti il più ampio sistema continuo cui appartengono le \(E\) spezzate in una curva \(C\) dell’ordine di \(C_0\) e in una \(D\) dell’ordine di \(D_0\)): le dette \(m\) falde non debbono toccarsi tutte lungo la \(V_r\) o in un punto generico \(P\) di essa. Vuol dire che non può aversi un ramo lineare di curva uscente da \(P\) che abbia un contatto d’ordine finito \(s\) \((\ge 1)\) colle fade suddette e non con \(V_r\); in altri termini: se un tal ramo ha \(h<s+1\) \((>0)\) punti successivi comuni con \(V_r\), avrà comune con questa varietà anche il successivo punto \((h+1)\)—mo. % [sic] fade -> falde, the surrounding text refers to the same hypersurface sheets.

Ovvero, ritornando dalla \(V_r\) alla nostra superficie \(F\), occorre mostrare che per ogni curva \(C_1\) infinitamente vicina a \(C_0\) nell’intorno del prim’ordine, esiste una serie di curve infinitamente vicine \(C_1C_2\ldotp\ldotp\ldotp\ldotp\) negli intorni d’ordine successivo \(1,\ 2,\ 3\,\ldotp\ldotp\ldotp\ldotp\); serie \textit{regolare}, (cioè assimilabile ad un ramo lineare). La \(C_1\) fa parte di una curva spezzata \(E_1=C_1+D_1\) vicina a \(C_0+D_0\); la \(C_2\) farà parte di una curva successiva \(C_2+D_2\), e così via. % ?? uncertain: number of dots immediately following C_1C_2 -- five or four.
% EDITORIAL D0167 RISOLTO: The enlarged scan shows four separate baseline dots after C_2; the transcription has five.

Per assolvere l’esigenza così posta, e quindi completare la di%
%CONT>>
% scan: s178R
% !PRINTED 333
\unskip\weberpage{333}%
% head: SUPERFICIE IRREGOLARI E SISTEMI CONTINUI DI CURVE, ECC.
%<<CONT
% stamp: Dark left scan edge and scattered scan specks.
mostrazione generale del teorema fondamentale si affaccia assai semplicemente un’idea\origfootnote{1}{s178R-1}{La parte che segue è stata lasciata incompleta dall’Autore e quindi le argomentazioni ivi svolte presentano numerose lacune; tuttavia si è ritenuto opportuno riportarla perchè contiene idee che forse, opportunamente completate, potranno fornire lo spunto per una sistemazione della teoria. (Nota dei revisori).}: poichè su \(F\) sono date due curve infinitamente vicine non equivalenti, \(C_0\) e \(C_1\), l’operazione \(C_1-C_0\) applicata successivamente ad una curva, p. es. a \(C_1\), varrà a definire la serie delle curve successive \(C_2C_3\,\ldotp\ldotp\ldotp\ldotp\), o quella dei sistemi lineari da essa determinati. Qui giuoca, e dir vero, un’intuizione ardita, cioè viene assunto che « possa operarsi sulla superficie \(F\), per somma e sottrazione di curve, anche quando si tratti di curve infinitamente vicine ». Si vuole giustificare questa intuizione. % [sic] e dir vero -> a dir vero, apparent misprint in the idiomatic phrase.

Prendiamo le mosse da un questione in qualche modo analoga, che si riferisce alle curve. Se, sopra una curva \(L\) di genere \(p>0\), si considerano due gruppi di \(p\) punti, \(G\) e \(G'\), disequivalenti, l’operazione \(+G'-G\) applicata successivamente a partire da \(G\) conduce ad una serie (generalmente) illimitata di gruppi % [sic] un questione -> una questione, missing feminine ending before questione.
\[
G,\quad G',\quad G'',\quad G''',\quad \ldotp\ldotp\ldotp\ldotp,
\]
che rispondono agli omologhi di \(G\) nelle potenze di una trasformazione di prima specie \(T\) della varietà di Jacobi \(V\), relativa ad \(L\). Perciò la serie dei gruppi successivi, \(G'',\,G'''\ldotp\ldotp\ldotp\ldotp\), non cessa di essere definita\origfootnote{2}{s178R-2}{Nel senso preciso delle condizioni differenziali che definiscono i punti infinitamente vicini. Cfr. \textsc{Enriques-Chisini}, \textit{Lezioni}, L. IV, vol. II.} quando \(G'\) diventi infinitamente vicino a \(G\): in questo caso la trasformazione \(T\) diventa una trasformazione infinitesima generatrice del gruppo continuo della \(V\), secondo \textsc{S. Lie}, e definisce una serie analitica \(t\), traiettoria del gruppo di Jacobi di \(L\).

Per passare di qui alla considerazione delle curve infinitamente vicine sopra una superficie, conviene giustificare anzitutto il principio di spezzamento per la curva infinitamente vicina ad una \(C_0+D_0\), che abbia (come nella costruzione indicata innanzi) \(m=C_0D_0\) punti doppi vicini ai punti comuni a \(C_0\) e \(D_0\); si tratta dunque di mostrare che tale curva è una \(C_1+D_1\), formata di due parti rispettivamente vicine a \(C_0\) e \(D_0\).

A tal uopo si può supporre che \(C_0\) e \(D_0\) siano curve normali non speciali appartenenti a spazi di dimensione assai elevata rispetto al loro genere: a questo caso si può invero ridursi mercè una trasformazione della superficie data \(F\). Allora è facile verificare che, nello spazio di \(F\), le curve, irriducibili \(E\) di ordine \(N\) e genere \(\Pi\), che costituiscono la famiglia definita dalle \(C+D\) spezzate, formano una varietà avente una certa dimensione \(R\), entro cui si ha una varietà
%CONT>>
% scan: s179L
% !PRINTED 334
\weberpage{334}
% head: CAPITOLO NONO
%<<CONT
% stamp: Dark right scan edge and scattered specks; a diagonal stray stroke crosses the plus sign after 2n in item 1).
di dimensione \(R-m\) costituita dalle \(E=\overline{E}\) (spezzate) con \(m\) punti doppi comprendenti la famiglia delle \(C+D\), per modo che lo spezzamento delle \(E\) \textit{nello spazio} di \(F\) (non diciamo sulla superficie) dipenda proprio da \(m\) condizioni indipendenti, e non da meno.

Per dimostrarlo, si proietti la curva spezzata \(C_0+D_0\) con \(m\) punti doppi (comuni a \(C_0\) e \(D_0\)) da un conveniente spazio lineare sopra una retta multipla \(N\)-pla (riducibile)con \(2N+2\Pi-2\) punti critici. % [sic] (riducibile)con -> (riducibile) con, no separating space is printed.

Questa retta \(N\)-pla si può ritenere come una riemanniana ad \(N\) fogli, che — per una certa scelta di cappi, presi in un dato ordine e per una opportuna nomenclatura dei rami — si lascia definire da un sistema di trasposizioni sui rami, simile a quella che risponde al teorema di Lüroth-Clebsch\origfootnote{1}{s179L-1}{Cfr. \textsc{Enriques-Chisini}, \textit{Lezioni}, L. V, Cap. I, § 3, vol. III, pag. 26; Cap. III § 33, vol. III, pag. 359-376.}; così (essendo \(n,\pi\) e \(\nu,\varrho\) i caratteri rispettivamente di \(C_0\) e \(D_0\)) si avranno:

1) anzitutto \(2n+2\pi-2\) cappi a ciascuno dei quali risponde un numero pari di trasposizioni simili \((12)\ (12)\,\ldotp\ldotp\ldotp\ldotp\,((n-1)n),\ ((n-1)n)\);

2) poi \(m\) cappi (che vanno agli \(m\) punti critici apparenti immagini dei punti comuni a \(C_0\) e \(D_0\)) a ciascuno dei quali risponde il prodotto di due trasposizioni
\[
(n(n+1))\cdot(n(n+1))=1;
\]

3) e infine \(2\nu+2\varrho-2\) cappi a ciascuno dei quali risponde un numero pari di trasposizioni simili
\[
\begin{aligned}
((n+1)(n+2)),&\quad ((n+1)(n+2))\,\ldotp\ldotp\ldotp\ldotp\\
((n+\nu-1)(n+\nu)),&\quad ((n+\nu-1)(n+\nu)).
\end{aligned}
\]

Ora, pel teorema d’esistenza\origfootnote{2}{s179L-2}{Cfr. \textsc{Enriques-Chisini}, \textit{Lezioni}, L. V, Cap. III, § 33 l. c.}, si può far variare per continuità la retta multipla così definita, staccando i due punti di diramazione che confluiscono in uno dei punti critici apparenti. Appare quindi che la curva spezzata \(\overline{E}=C_0+D_0\), nello spazio di \(F\), è suscettibile di variare in una certa varietà \(\infty^R\ V\) di curve irriducibili \(E\), ed entro a questa si trova una varietà \(\infty^{R-m}\ \overline{V}\) di curve spezzate con \(m\) punti doppi, per modo che le curve infinitamente vicine alla \(C_0+D_0\) con \(m\) punti doppi costituiscono precisamente l’intorno lineare di questo elemento-curva su \(\overline{V}\). % ?? uncertain: first letter in the two infinity exponents -- R or n; R is the best reading.
% EDITORIAL D0168 RISOLTO: Enlarged image shows the same italic capital R in both exponents; the bowl and diagonal leg are visible.

Infine dunque si è dimostrato il \textit{lemma fondamentale}: « ogni curva \(E\) con \(m\) punti doppi, infinitamente vicina a \(C_0+D_0\) (nello spazio di \(F\)) appartiene ad una serie (regolare) \(\infty^1\) di curve spezzate che comprende la \(C_0+D_0\), e perciò deve ritenersi effettivamente
%CONT>>
% scan: s179R
% !PRINTED 335
\weberpage{335}
% head: SUPERFICIE IRREGOLARI E SISTEMI CONTINUI DI CURVE, ECC.
%<<CONT
% stamp: Dark left scan edge and upper-right horizontal streak; a thin horizontal marginal mark extends from "giustificare".
spezzata in due curve improprie \(C_1\) e \(D_1\) rispettivamente dell’ordine di \(C_0\) e \(D_0\), vicine a queste curve ».

Il lemma fondamentale così stabilito significa che « la curva infinitamente vicina a \(C_0\) che fa parte di una curva vicina a \(C_0+D_0\) con \(m\) punti doppi, quale si è imparato a costruire sopra la superficie \(F\), \textit{dipende} soltanto da \(C_0\) e \textit{non} dalla scelta \textit{di una \(D_0\)}, che in infiniti modi diversi può associarsi ad essa, entro un sistema regolare \(|D|\) ».

Su questa base vogliamo riconoscere che successivamente alla detta curva \(C_1\), si ha ancora, nell’intorno di \(2^\circ\) ordine di \(C_0\), una curva \(C_2\), e poi una curva \(C_3\) ecc.

Sia lecito qui parlare di \textit{curve infinitamente vicine}, che si presumono definite da condizioni differenziali allo stesso modo dei punti sopra una curva, come di \textit{curve proprie}; riservandoci di giustificare più avanti l’uso che facciamo di questo linguaggio.

Sulla superficie \(F\), che contiene i sistemi lineari \(|C|\) e \(|D|\) considerati innanzi, si costruisca un sistema lineare normalmente grande \(|L|\) di un certo genere \(P\); per modo che \(|C|\) e \(|D|\) seghino sopra una \(L\) generica, ed anche sopra ogni \(L\) (irriducibile) dotata di un punto doppio, \textit{serie complete} d’ordine inferiore a \(P\)\origfootnote{1}{s179R-1}{Cfr. Cap. IV, § 10.}. Consideriamo in \(|L|\) un fascio generico di curve \(L\), avente un certo gruppo \(\Gamma\) di punti base: diciamo anzitutto che « le curve infinitamente vicine disequivalenti \(C_0\) e \(C_1\) segano sulle \(L\) del fascio, gruppi disequivalenti ».

Questa affermazione consegue da ciò che \(|C_0|\) sega su ogni \(L\) una serie completa: così — essendo \(M=CL>C_0^2\) — tante sono le curve equivalenti infinitamente vicine a \(C_0\) quanti sono su \(L\) i gruppi equivalenti a \(C_0L\) vicini ad esso; e se \(C_1\) segasse su una \(L\) un gruppo equivalente a \(C_0L\), dovrebbe pure essere \(C_1\) equivalente a \(C_0\).

Ora si possono modificare lievemente e precisare le nostre assunzioni, assumendo che \(C_0\) appartenga, entro \(\{C\}\), ad un sistema regolare di dimensione \(r=p+n-\pi+1\) sicchè la curva stessa sia individuata da \(r\) punti base semplici che formino un sottogruppo \(A\) di \(\Gamma\) (e similmente, ove occorra, che \(|D_0|\) sia pure di dimensione \(r'=p_a+\nu-\varrho+1\) sicchè la curva \(D_0\) resti individuata imponendo \(r'\) punti base semplici che formino un sottogruppo \(B\) dello stesso \(\Gamma\)). Allora nel gruppo base \(\Gamma\) del fascio delle \(L\) si sceglieranno fuori di \(A\) e \(B\), altri \(P-M\) punti (riducibili ad un solo punto da contare \(P-M\) volte), formanti un certo gruppo \(S\), che sommati al gruppo di \(M\) punti sezioni di \(C_0\) con \(L\) daranno su questa un gruppo \(G\) di \(P\) punti, non speciale, e sommati al gruppo sezione di \(C_1\), daranno similmente un gruppo di \(P\) punti \(G_1\), infinitamente vicino % [sic] r=p+n-\pi+1 -> r=p_a+n-\pi+1, the parallel formula for r' uses p_a.
%CONT>>
% scan: s180L
% !PRINTED 336
\weberpage{336}
% head: CAPITOLO NONO
%<<CONT
% stamp: Dark right scan edge and scattered marginal specks and strokes; a diagonal stray stroke crosses "perciò".
e disequivalente a \(G\). La corrispondenza fra \(G\) e \(G_1\) (trasformazione infinitesima del gruppo della varietà di Jacobi di \(L\)) genera sulla \(L\) del fascio una serie analitica \(t\) di gruppi di \(P\) punti, cui appartiene un gruppo \(G_2\) infinitamente vicino a \(G\) nell’intorno di \(2^\circ\) ordine di \(G\), cioè successivo a \(G_1\). Al variare di \(L\) nel suo fascio, il gruppo \(G_2\), che riesce determinato algebricamente e razionalmente\origfootnote{1}{s180L-1}{L’espressione differenziale che definisce il \(G_1\), da cui può ricavarsi questa di \(G_2\), trovasi calcolata da \textsc{Chisini}. Cfr. \textsc{Enriques-Chisini}, \textit{Lezioni}, L. VI, Cap. II, § 28, vol. IV, pag. 143.}, descriverà sulla superficie \(F\) una curva \(C_2\) infinitamente vicina alla \(C_0\), nel suo intorno di \(2^\circ\) ordine, e successiva alla \(C_1\), cui saranno da aggiungere le curve infinitesime, con esse non connesse, che circondano i punti del gruppo \(A\).

Per giustificare l’asserzione fatta, che al variare di \(L\) nel fascio \(G_2\) generi una curva \(C_2\) vicina a \(C_1\) e dello stesso ordine, occorre riconoscere che il luogo di \(G_2\), definito a meno di curve irriducibili \(L\) del fascio, non contiene qualcuno dei punti del suo gruppo base \(\Gamma\) fuori di \(A\), e perciò che sopra una \(L\) \textit{qualunque} del nostro fascio (scelto in modo generico entro \(|L|\)) il gruppo di punti \(G_2\) (vicino a \(G\) nell’intorno di \(2^\circ\) ordine) non può essere speciale.

A tal uopo osserviamo che:

1) sopra una curva \(L\), di genere \(P\), il gruppo \(G_2\) non può essere speciale senza che siano speciali anche i gruppi \(G\) e \(G_1\); imperocchè un gruppo infinitamente vicino ad un gruppo non speciale è certo non speciale;

2) poichè il gruppo di \(M\) punti \(LC_0\) (tolto \(A\)) non appartiene ad una serie lineare di dimensione \(>0\), sommando ad esso un gruppo generico di \(P-M\) punti (il gruppo \(S\)) si ottiene un gruppo non speciale\origfootnote{2}{s180L-2}{C’è qui una semplice conseguenza del teorema di Riemann-Roch. Cfr. p. es. \textsc{Enriques-Chisini}, \textit{Lezioni}. L. V, Cap. I, vol. III.}: per avere un gruppo di \(P\) punti speciale bisogna assoggettare \(S\) ad \textit{una} condizione, e almeno ad \textit{un’altra condizione} conviene assoggettare la scelta di \(S\) sopra \(L\) affinchè anche il gruppo \(G=S+LC\), sia speciale; così dunque la scelta di un gruppo \(S\) di \(P-M\) punti di \(L\) che, sommato a \(LC_0\) e a \(LC_1\), dia luogo a due gruppi speciali \(G\) e \(G_1\), importa almeno \textit{due} condizioni a cui \(S\) deve soddisfare.

Ciò posto riprendiamo la costruzione precedentemente disegnata di un fascio di curve \(L\) in \(|L|\): si assumerà anzitutto una particolare curva \(\overline{L}\) di \(|L|\) (irriducibile di genere \(P\)) e su questa un gruppo di \(r\) punti imposti come punti base a \(|C|\), e poi un gruppo \(S\) di \(P-M\) punti da aggiungere a quelli del gruppo \(LC_0\) (da cui è tolto \(A\)) e infine, sulla medesima \(\overline{L}\), un gruppo \(\Gamma\) della serie carat%
%CONT>>
% scan: s180R
% !PRINTED 337
\unskip\weberpage{337}%
% head: SUPERFICIE IRREGOLARI E SISTEMI CONTINUI DI CURVE, ECC.
%<<CONT
% stamp: Dark left scan edge, a thin vertical line in the outer right margin, and scattered specks.
teristica \(g_N^{R-r}\) di \(\overline{L}\), contenente \(A\). Se il gruppo di \(P\) punti \(G=S+LC_0\) deve essere speciale su \(\overline{L}\) e deve avere, sulla stessa curva, un gruppo vicino \(G_1=S+LC_1\) pure speciale, bisogna che la scelta del gruppo \(S\) soddisfi, come abbiam detto, ad almeno due condizioni; e se, invece, le analoghe condizioni debbono essere soddisfatte sopra una curva \(L\) del fascio che ha per gruppo base \(\Gamma\), disponendosi qui di un parametro, bisognerà tuttavia che la scelta di \(S\) e quindi quella di \(\Gamma\) soddisfi sopra \(\overline{L}\) ad una condizione almeno, e così avremo, un fascio di \(L\) che non è scelto in modo generico entro \(|L|\). % ?? uncertain: superscript of g -- R-r or n-r; R-r is the best reading.
% EDITORIAL D0169 RISOLTO: Enlarged image supports italic capital R followed by -r, consistent with the earlier dimension notation.
% line-break repetition: The plus sign after S in G=S+ +LC_0 is printed at both the end and the beginning of adjacent lines.

In tal guisa resta provato l’asserto: il gruppo \(G_2\) costruito, come innanzi, sulla \(L\) di un fascio generico di \(|L|\) descrive una curva vicina a \(C_1\) dello stesso ordine (non passante per qualcuno dei punti base del \(\Gamma\)). E giova avvertire che il nostro ragionamento non cade in difetto per la circostanza che entro un fascio di \(L\) si abbiano, in generale, curve (irriducibili) di genere effettivo \(P-1\), dotate di un punto doppio \(X\); imperocchè è lecito ritenere codeste curve di genere virtuale \(P\) (assumendo il punto doppio \(X\) come virtualmente inesistente). Vuol dire che un gruppo \(S\) di \((P-1)+1\) punti sopra una \(L\) siffatta dovrà ritenersi non speciale ove non appartenga ad una serie lineare \(\infty^1\) che abbia una coppia neutra \(X_1X_2\) in corrispondenza ad \(X\). Così, anche in questo caso, la specialità di \(G\) e \(G_1\) porta due condizioni almeno cui la scelta del gruppo \(S\) deve soddisfare.

Nel discorso precedente si parla delle curve \(C_1\ C_2\,\ldotp\ldotp\ldotp\ldotp\) infinitamente vicine a \(C_0\), come se fossero curve proprie. Per giustificare l’uso di questo linguaggio, conviene mostrare che le dette curve possono veramente definirsi mediante un \textit{passaggio al limite} di curve proprie. Ciò accade anzitutto per la curva \(C_1\) infinitamente vicina a \(C_0\) nell’intorno del prim’ordine: giacchè abbiam veduto che \(C_1\) fa parte della curva riducibile vicina a \(C_0\) entro un fascio limite per \(\lambda=0\) di \(C_0+D_0+\lambda E\) \((|E|=|C+D|)\) che ha come punti base gli \(m\) punti del gruppo \(C_0D_0\). È ben vero che la curva di questo fascio è, per \(\lambda\ne 0\), irriducibile e tende alla \(C_0+D_0\); ma, secondo il lemma fondamentale stabilito innanzi, per \(\lambda\) infinitesimo questa curva acquista \(m\) punti base e si spezza in \(C_1+D_1\): vuol dire che la serie analitica \(t\ \infty^1\) dei gruppi di punti determinata da \(G\) e \(G_1\) sulle \(L\) del nominato fascio, non dipende affatto dalla scelta di \(D_0\) o del sistema \(|D|\) che si è associato a \(|C|\).

Ora si può costruire una curva propria la cui variazione definisce \(C_2\), considerando sopra una \(L\) generica del detto fascio:

1) il gruppo \(G\) sezione di \(C_0\) cui si aggiungono i punti di \(A\);

2) il gruppo \(G'\) costituito (oltre che da \(A\)) dai punti di \(C_0+D_0+\lambda E\) vicini a quelli di \(G\);
% line-break repetition: The plus sign after C_0 is printed at both the end and the beginning of adjacent lines.

3) il gruppo di \(P\) punti \(G''\), che risponde a \(G'\), nella trasformazione \((GG')\) della varietà di Jacobi di \(L\).

% EDITORIAL: segnatura originale: \textsc{Enriques F.} - \textit{Superficie algebriche.}\hfill 22
\printersignature{\textsc{Enriques F.} - \textit{Superficie algebriche.}\quad 22}

% scan: s181L
% !PRINTED 338
\weberpage{338}
% head: CAPITOLO NONO
% stamp: Scattered scan specks, a dark right-edge band, and short marginal strokes.
Il luogo descritto da \(G''\), al variare di \(L\) nel fascio, vale appunto a definire la curva \(C_2\) successiva a \(C\) e \(C_1\).

A dir vero, finchè \(\lambda\ne 0\), il detto luogo è una curva irriducibile che tende a una curva più ampia di \(C_0\), cioè a \(C_0+D_0\); infatti un punto di \(C_0+D_0+\lambda E\), vicino a \(C_0\), col variare di \(L\) nel fascio, è suscettibile di scambiarsi con un punto vicino a \(D_0\), sicchè la curva che vien definita dal \(G''\) variabile sarà generata dalla somma di \(G''\) e di un \(\bar{G}''\) analogo definito rispetto a \(D_0\) e vicino a questa \(D_0\) e la somma \(G''+\bar{G}''\) descriverà una curva \(E_2\), secante sulle \(L\) gruppi della serie lineare segata da \(|E|=|C+D|\), o (a priori) un multiplo di questa serie, il cui limite — per \(\lambda=0\) — è \(C_0+D_0\)\origfootnote{1}{s181L-1}{La giustificazione dell’asserto si riduce a questo: si abbiano su una curva \(L\) di genere \(P\) due gruppi di \(P\) punti \(G\) e \(\bar{G}\), e ancora altri due \(G'\) e \(\bar{G}'\) tali che \(G+\bar{G}=G'+\bar{G}'\) e si costruiscano i gruppi \(G''=G'+|G'-G|\) e \(\bar{G}''=\bar{G}'+|\bar{G}'-\bar{G}|\), allora \(G''+\bar{G}''=G'+\bar{G}'=G+\bar{G}\).}.

Ma, sopra \(L\), i gruppi \(G\,G'\,G''\) appartengono ad una serie analitica \(\infty^1\) di gruppi di \(P\) punti che — per \(\lambda\to 0\) — quando essi diventano infinitamente vicini, dipende soltanto da \(C_0\) e \(C_1\) e \textit{non da} \(D_0\) o da \(|D|\), e perciò vale veramente a definire, nell’intorno di \(2^\circ\) ordine di \(C_0\), la curva \(C_2\), parte irriducibile di \(E_2\) per \(\lambda\to 0\), di cui volevasi stabilire l’esistenza.

Abbiamo costruito su \(F\) una curva \(C_2\) vicina a \(C_1\) e successiva ad essa, nell’intorno di \(2^\circ\) ordine di \(C_0\), la quale fa parte di una \(E_2=C_2+D_2\) spezzata, del sistema lineare \(|E|=|C+D|\). Similmente, a partire da \(C_2\), si costruirà una successiva curva \(C_3\), vicina a \(C_0\) nell’intorno di \(3^\circ\) ordine ecc. Così viene assolta l’esigenza dimostrativa dichiarata innanzi, cioè riesce stabilito che vicino alla \(E_0=C_0+D_0\) spezzata, vi sono, successivamente in tutti gl’intorni d’ordine finito, curve del sistema \(|E|=|C+D|\) spezzate, formanti una serie regolare (assimilabile ad un ramo lineare di curva); e quindi risulta dimostrato in generale il teorema fondamentale.

Sopra una superficie di genere geometrico \(p_g\geq 0\) e di genere numerico \(p_a<p_g\), ogni sistema lineare regolare non speciale di curve irriducibili \(|C|\) è contenuto in un sistema continuo \(\{C\}\) che sega su una \(C\) la serie caratteristica completa e perciò riesce formato da \(\infty^q\) (\(q=p_g-p_a\)) sistemi lineari distinti. Quindi: \textit{l’irregolarità geometrica d’una superficie eguaglia sempre l’irregolarità numerica:}
\[
q=p_g-p_a.
\]

E (\S~2) \textit{sopra una superficie di irregolarità \(q\) ogni curva \(C\), coi caratteri \(n\), \(\pi\) e \(i\), che sia aritmeticamente effettiva, cioè tale che}
\[
p_a+n-\pi+1-i\geq 0,
\]
\textit{appartiene ad una serie continua \(\infty^q\) di curve disequivalenti.}

% scan: s181R
% !PRINTED 339
\weberpage{339}
% head: SUPERFICIE IRREGOLARI E SISTEMI CONTINUI DI CURVE, ECC.
% stamp: Dark left-edge scan band and scattered specks.
\edsection{6. Storia della teoria dei sistemi continui.}{6. Storia della teoria dei sistemi continui.}{sec.IX.6}

La teoria dei sistemi continui di curve appartenenti alle superficie irregolari ha un significato fondamentale, non soltanto per lo studio algebrico-geometrico delle superficie, sì anche per la dottrina (trascendente) degli integrali semplici (di differenziali totali) che ad esse appartengono, valendo a ricollegare questi due aspetti diversi della scienza delle superficie. Pertanto, mettendoci dal punto di vista generale che risponde alla sua genesi, vogliamo approfondire la comprensione della teoria generale di questi sistemi continui offrendone una prospettiva storica, dove si veda pure il significato degli errori per cui la teoria stessa è dovuta passare: cosa tanto più essenziale se si ritenga che lo sviluppo della teoria non abbia ancora toccato ad una sistemazione definitiva.

La memoria di \textsc{Castelnuovo-Enriques} del 1900 « Sopra alcune questioni fondamentali nella teoria delle superficie algebriche »\origfootnote{1}{s181R-1}{\textit{Annali di Matematica}, t. VI, s. II, pag. 165.}, mentre da una parte conchiudeva fruttuosamente una serie di sforzi, dall’altra poneva problemi nuovi, come quello della determinazione delle superficie coi plurigeneri nulli (e quindi della superficie di genere \(p_g=0\) e di genere numerico negativo). Per approfondire tali problemi e in generale per far progredire la conoscenza delle superficie, anche in rapporto alla teoria trascendente della scuola francese, si presentava come istrumento lo studio dei sistemi continui di curve disequivalenti.

La considerazione generale di sistemi siffatti risale ad una Nota \textsc{G. Humbert} del 1893\origfootnote{2}{s181R-2}{\textit{Sur une propriété d’une classe de surfaces algébriques.} (Comptes rendus, 1893).}: l’autore dimostra che il possesso di un tale sistema di curve porta l’esistenza di integrali semplici di prima specie (integrali di Picard di differenziali totali) annessi alla superficie. Su questa Nota richiamava l’attenzione dei colleghi \textsc{C. Segre} e poco dopo l’argomento diventava di studio per la scuola geometrica italiana. Già nel 1896 \textsc{Enriques} vi portava un contributo indiretto\origfootnote{3}{s181R-3}{\textit{Un’osservazione relativa alla rappresentazione parametrica delle curve algebriche.} Rendic. Circolo mat. di Palermo, 1895.} mostrando che « sopra una superficie una serie razionale di curve è sempre contenuta in un sistema lineare ».

Nel 1899 \textsc{Enriques} inizia lo studio geometrico delle superficie contenenti sistemi continui \(\{C\}\) di curve disequivalenti\origfootnote{4}{s181R-4}{\textit{Una proprietà delle serie continue di curve appartenenti ad una superficie algebrica regolare.} Rendic. Circolo mat. di Palermo, 1899.}, stabilendo che esse sono irregolari. L’autore mostra che il sistema aggiunto
%CONT>>
% scan: s182L
% !PRINTED 340
\weberpage{340}
% head: CAPITOLO NONO
%<<CONT
% stamp: Long black streak in the top margin, a right-edge band, and scattered marks; a diagonal stroke crosses a mathematical relation.
ad un sistema lineare \(|C|\) contenuto in \(\{C\}\) non può segare sopra una \(C\) la serie canonica completa. Avendo egli comunicato, prima della pubblicazione, il risultato al \textsc{Castelnuovo}, questi gli rispondeva offrendo una seconda dimostrazione del teorema, ove si fa vedere che la serie caratteristica di \(|C|\) è necessariamente deficiente. E nella Nota da lui pubblicata \textsc{Enriques}, riporta, dopo la propria, la dimostrazione del Collega.

Successivamente (nel 1901) \textsc{Enriques} stesso — avendo di mira di riconoscere il legame fra le superficie irregolari geometricamente definite (\(p_g-p_a>0\)) e quelle che posseggono integrali semplici di prima specie — dimostrava che « le superficie possedenti \(q\) integrali semplici di prima specie con \(2q\) periodi, contengono sistemi continui di curve disequivalenti\origfootnote{1}{s182L-1}{\textit{Sur les surfaces algébriques admettant des intégrales de différentielles totales de première espèce.} Annales de Toulouse (2), 3, pag. 77.} ».

D’altra parte \textsc{Castelnuovo} riusciva a dimostrare che « una superficie di genere \(p_g=0\) e \(p_a<0\), contenente un sistema continuo di curve disequivalenti, possiede di conseguenza anche un fascio irrazionale di curve » (\S~11). Era una proprietà il cui significato dipendeva da ulteriori eventuali progressi della teoria, e \textsc{Castelnuovo} si limitava a comunicarla ad \textsc{Enriques}, che la renderà pubblica più tardi (1904), quando potrà dedurne che « ogni superficie con \(p_g=0\) e \(p_a<0\) possiede sempre un tal fascio irrazionale ». % ?? uncertain: relation after the first p_g -- printed equality crossed by a later stroke, or printed \ne.

Frattanto, nel 1902, \textsc{F. Severi}, da poco laureato nell’Università di Torino, per consiglio del suo maestro \textsc{C. Segre}, veniva come assistente alla cattedra di geometria dell’Università di Bologna, a l’ora tenuta dall’\textsc{Enriques}, allo scopo di entrare in più stretta relazione coi cultori della teoria delle superficie. \textsc{Enriques} attirava subito l’attenzione del giovane geometra sull’argomento dei sistemi continui, e in ispecie gli suggeriva lo studio di una particolare famiglia di superficie contenente sistemi siffatti, cioè delle superficie che rappresentano le coppie di punti di una curva, alle quali \textsc{Severi} dedicava appunto un lavoro, pubblicato nel 1903\origfootnote{2}{s182L-2}{\textit{Sulle superficie che rappresentano le coppie di punti di una curva algebrica.} Atti della R. Accademia di Torino, Gennaio 1903.

Convien dire che, indipendentemente dal \textsc{Severi}, anche \textsc{M. De Franchis} pubblicava, nel frattempo, uno studio \textit{Sulla varietà \(\infty^2\) delle coppie di punti di due curve o di una curva algebrica.} Rendic. Circolo mat. di Palermo, 1903.}. % ?? uncertain: spacing and letters in "a l’ora" -- "a l’ora" or "all’ora".

Lo studio delle superficie rappresentative delle coppie di punti d’una curva offriva esempi appropriati ad illuminare questioni d’ordine generale: ivi, in particolare, \textsc{Severi} incontrava il problema della base, che più tardi costituirà oggetto di qualcuno dei più im%
%CONT>>
% scan: s182R
% !PRINTED 341
\unskip\weberpage{341}%
% head: SUPERFICIE IRREGOLARI E SISTEMI CONTINUI DI CURVE, ECC.
%<<CONT
% stamp: Dark edge bands and scattered specks; a round black blot follows "all’espressione".
portanti suoi studi. Frattanto esso dava occasione all’introduzione del concetto più esteso della serie caratteristica.

L’idea era in germe nella dimostrazione che abbiam detto: \textsc{Castelnuovo} aveva comunicato ad \textsc{Enriques}, circa la deficienza della serie caratteristica di un \(|C|\) contenuto in un \(\{C\}\) di curve disequivalenti, perchè ivi si vede la serie deficiente di \(|C|\) ampliarsi col limite della serie caratteristica di un sistema lineare infinitamente vicino. Ma più precisamente \textsc{Castelnuovo} rilevava in maniera esplicita questo concetto riferendosi alle superficie (iperellittiche) che rappresentano le coppie di punti di una curva \(L\) di genere \(2\); su tali superficie si ha un sistema continuo di curve \(C\) di genere due, che è del grado due e contiene entro di sè le immagini delle serie di coppie di punti di \(L\) con un punto fisso: orbene una curva \(C\) di tale sistema è isolata, sicchè manca il sistema lineare che dovrebbe segare su di essa la serie caratteristica, la quale (per essere \(p_g=1\) e la curva canonica d’ordine zero) sarebbe la \(g_2^1\) canonica; ma sulla detta \(C\) la \(g_2^1\) stessa viene segata dalle curve di \(\{C\}\) infinitamente vicine.

Il \textsc{Castelnuovo} avendo comunicato questa osservazione al \textsc{Severi}, questi fu tratto a svilupparla in generale, avvertendo anzitutto che la serie caratteristica di una \(C\) può essere definita, indipendentemente dall’esistenza di un sistema lineare più ampio che contenga \(C\), come differenza
\[
(C+D)\,C-DC,
\]
dove \(|D|\) sia un sistema lineare ausiliario qualsiasi (tale che \(|C+D|\) seghi su \(C\) una serie completa). Infatti questa serie differenza è indipendente dalla scelta del sistema ausiliario \(|D|\)\origfootnote{1}{s182R-1}{\textsc{Severi}, \textit{Osservazioni sui sistemi continui di curve appartenenti ad una superficie algebrica.} Atti della R. Accademia di Torino, Febbraio, 1904.}. Come si vede c’è qui l’interpretazione geometrica o funzionale del procedimento con cui \textsc{Enriques} aveva già definito il grado virtuale della curva \(C\), riguardando appunto \(|C|\) come differenza di due sistemi lineari di curve.

Definita in tal guisa la serie caratteristica di una curva \(C\), indipendentemente dall’esistenza di un più ampio sistema lineare \(|C|\), non vi è difficoltà a riconoscere che « se \(C\) appartiene comunque ad un sistema continuo \(\{C\}\), le curve di questo infinitamente vicine a \(C\) segano su di essa gruppi della serie caratteristica ». E quindi la dimostrazione dei teoremi di Enriques, che sopra abbiamo ricordato, si riduceva all’espressione più semplice.

Fino da quando si era riconosciuto che il possesso di un sistema continuo di curve disequivalenti porta la irregolarità della super%
%CONT>>
% scan: s183L
% !PRINTED 342
\unskip\weberpage{342}%
% head: CAPITOLO NONO
%<<CONT
% stamp: Scattered marginal marks and a black blot in the first footnote near "anzichè".
ficie (\(p_a<p_g\)), si affacciava naturalmente la richiesta di invertire il teorema, stabilendo dunque che il possesso di sistemi continui siffatti costituisce la proprietà caratteristica delle superficie irregolari. Anzi tale richiesta assumeva una forma più precisa: infatti dei geometri che solo per forza di esempi si erano staccati dall’ipotesi implicitamente ammessa da \textsc{Noether} circa l’integrità della serie caratteristica dei sistemi lineari completi, dovevano domandarsi se non sarebbe invece completa la serie caratteristica di un più ampio sistema continuo. Ma il punto era di trovare un mezzo per discutere la questione; a tal uopo occorreva infatti un criterio che permettesse in qualche modo di contare le condizioni di spezzamento di una curva contenuta in un sistema lineare. \textsc{Enriques} (1904) trovò un criterio siffatto nel \textit{principio di degenerazione o di spezzamento}: se una curva irriducibile \(C\), variando con continuità, venga a spezzarsi in due componenti, queste dovranno avere in comune almeno un punto di connessione, che non sia limite di un punto doppio variabile di \(C\).

Sotto l’aspetto topologico c’è qui una verità evidente: perchè in forza della continuità la riemanniana di \(C\), che è connessa, deve rimanere connessa in ogni passaggio al limite, conducente ad una \(C\) spezzata\origfootnote{1}{s183L-1}{Solo di recente è stata segnalata un’osservazione di \textsc{Noether} che precorre al detto principio. Nella memoria \textit{Ueber die reductiblen algebraischen Curven} (Acta Mathematica, 1896), questi, avendo svolto la teoria delle serie di gruppi di punti (segnate da sistemi lineari di curve o superficie) sopra una curva spezzata, piana o gobba, avverte che per le curve gobbe (!) la teoria non rientra in quella relativa alle curve irriducibili, se le componenti non abbiano almeno un punto comune.

Il principio di spezzamento, anzichè da un punto di vista topologico, può anche essere stabilito per via algebrico-geometrica. Cfr. p. es., \textsc{Enriques-Chisini}, Libro V, cap. III, \S~36, vol. III, pag. 405.}.

Coll’aiuto del principio di spezzamento \textsc{Enriques} credette di poter dimostrare che « i sistemi continui completi hanno la serie caratteristica completa »\origfootnote{2}{s183L-2}{\textit{Sulla proprietà caratteristica delle superficie algebriche irregolari.} Rendic. Acc. di Bologna, dicembre 1904.} e così che un sistema lineare regolare non speciale (di caratteri \(n\) e \(\pi\)), sopra una superficie d’irregolarità \(q=p_g-p_a\), appartiene ad un sistema continuo di dimensione \(p_g+n-\pi+1\), composto dunque di \(\infty^q\) sistemi lineari disequivalenti. Egli è ricorso a tal uopo alla rappresentazione della superficie sopra un piano multiplo che — ai fini della dimostrazione stessa — non è affatto essenziale, tantochè il ragionamento si può ripetere in maniera sostanzialmente identica sopra la superficie medesima, come ha osservato \textsc{Severi}\origfootnote{3}{s183L-3}{\textit{Intorno alla costruzione dei sistemi completi non lineari.} Rendic. Circolo mat. di Palermo, 1905.}.

% scan: s183R
% !PRINTED 343
\weberpage{343}
% head: SUPERFICIE IRREGOLARI E SISTEMI CONTINUI DI CURVE, ECC.
% stamp: Dark left-edge band and scattered marginal specks and strokes.
Conviene dir subito che il ragionamento di \textsc{Enriques} conteneva un errore, come viene spiegato negli sviluppi del \S~5: l’autore non si era accorto del punto critico che qui s’incontra e passa, senza accorgersene, da una proprietà differenziale ad una proprietà integrale. La difficoltà, ossia l’insufficienza del principio di spezzamento per \(p_g>0\), è stata anche illustrata nel detto \S~5 rilevando che il vero significato di codesto principio consiste in ciò che « lo spezzamento della curva di un sistema lineare \(|C+D|\) in due curve dell’ordine di \(C\) e di \(D\) con \(m\) punti comuni importa (non più di) \(m\) condizioni »; invece per \(p_g>0\), occorre riconoscere che queste condizioni sono, in generale, soltanto \(m-p_g\). Dimostrare che questo accade per le curve infinitamente vicine ad una curva spezzata \(C+D\) del nostro sistema lineare significa appunto stabilire una proprietà differenziale che non si è autorizzati a prendere in un senso integrale. La deduzione riesce giusta soltanto per \(p_g=0\) (e \(p_a<0\)), nel qual caso invero non fa d’uopo affatto di ricorrere alla considerazione della serie caratteristica.

Ma l’errore o la lacuna del ragionamento di \textsc{Enriques} non fu avvertito, per quindici anni, dai geometri, tantochè bisogna arrivare al 1921 per trovare la critica formulata da \textsc{Severi}; e a questa epoca il teorema fondamentale, concernente l’esistenza di serie \(\infty^q\) di curve disequivalenti, era stato dimostrato, per altra via, con procedimento trascendente da \textsc{Poincaré} (1910)\origfootnote{1}{s183R-1}{\textsc{H. Poincaré}, \textit{Sur les courbes tracées sur les surfaces algébriques.} Annales de l’Ecole Normale Sup., pag. 3 (27).}. Nel 1904 il ragionamento di \textsc{Enriques} pareva convalidato anche dall’esposizione che ne faceva lo stesso \textsc{Severi}, conservandone — come abbiamo accennato — lo schema essenziale\origfootnote{2}{s183R-2}{Cfr. la citata memoria: \textit{Intorno alla costruzione dei sistemi completi non lineari che appartengono ad una superficie irregolare.} Rendic. del Circolo mat. di Palermo, 1905.

Più tardi nella Nota \textit{Nuovi contributi alla teoria dei sistemi continui di curve appartenenti ad una superficie algebrica} (Rendic. Lincei, 2 e 16 aprile 1916), al n. 1, \textsc{Severi} riesponeva la dimostrazione di \textsc{Enriques}, riprendendo anzi la forma originale (rappresentazione della superficie sul piano multiplo).}.

E pertanto il teorema, così accettato dai geometri, diveniva il punto di partenza di importanti sviluppi.

Questi procedono essenzialmente in due direzioni:

1) \textsc{Enriques} si attacca anzitutto alla classificazione delle superficie di genere \(p_g=0\), e più tardi ai problemi più generali di classificazione delle superficie, che costituiranno l’argomento dei due capitoli seguenti X e XI. Fortunatamente questi sviluppi si
%CONT>>
% scan: s184L
% !PRINTED 344
\weberpage{344}
% head: CAPITOLO NONO
%<<CONT
% stamp: Right-edge scan band and diagonal marks near "abbiamo" and "un fatto".
basano o si possono basare per la maggior parte sul teorema fondamentale ristretto.

Ciò è chiaro a priori per il Cap. X relativo alle superficie di genere \(p_g=0\), ma vale anche per i risultati sulle superficie di genere numerico \(p_a<0\), che esponiamo nel Cap. XI, dove tuttavia occorre modificare in qualche punto le dimostrazioni. Fanno eccezione i teoremi concernenti le superficie con \(p_a=0\) e \(p_g=1\): che codeste superficie non possono avere la curva canonica d’ordine zero, e che posseggono sempre un fascio di curve ellittiche. % ?? uncertain: damaged initial in "d’ordine" -- d or l.

2) \textsc{Castelnuovo} e \textsc{Severi} si volgevano invece a sviluppare le conseguenze del teorema fondamentale in ordine alla teoria trascendente: dove si trattava di precisare dal punto di vista quantitativo il legame che il detto teorema poneva in luce fra le superficie irregolari e quelle possedenti integrali semplici di prima specie. In proposito abbiamo già accennato al primo tentativo da noi fatto per invertire l’osservazione di \textsc{Humbert}, cioè al teorema che le superficie possedenti \(q\) integrali semplici di prima specie con \(2q\) periodi contengono sistemi continui di curve disequivalenti e quindi sono irregolari. Questo teorema assumerebbe il più vasto significato ove la teoria trascendente di Picard permettesse a priori di conoscere che il numero dei detti integrali di prima specie, annessi ad una superficie qualunque, è sempre eguale alla metà del numero dei periodi.

In mancanza di ciò, già prima della Nota d’\textsc{Enriques} « Sulla proprietà caratteristica.... », \textsc{Severi} aveva investigato in altro modo la questione delle superficie possedenti integrali semplici di prima specie, dimostrando che « il possesso di integrali di seconda specie (e perciò di prima) porta che la serie caratteristica di un sistema lineare completo sia deficiente e quindi anche che la superficie sia irregolare »\origfootnote{1}{s184L-1}{\textit{Sulle superficie algebriche che posseggono integrali di Picard della 2\textsuperscript{a} specie.} Rendic. Lincei, 1904. Cfr. Math. Annalen, 1905.}. Sebbene questa dimostrazione valga a stabilire un fatto fondamentale e contenga un’analisi interessante relativa alle singolarità sulla curva polare dell’integrale, non sembra che essa possa venir completata fino a trovare il numero degli integrali di prima specie in funzione dell’irregolarità, senza l’aggiunta di qualche elemento essenziale\origfootnote{2}{s184L-2}{L’analisi della questione conduce \textsc{Severi} ad un risultato \textit{Sulla-differenza fra i numeri degli integrali di Picard della prima e della seconda specie appartenenti ad una superficie irregolare.} % ?? uncertain: mark between "Sulla" and "differenza" -- printed hyphen or scan streak.
% EDITORIAL D0173 RISOLTO: The enlarged image shows a horizontal connecting hyphen between Sulla and differenza, at the printed baseline/midline; retain it as printed.
Atti R. Accad. Torino, 1905. Al medesimo risultato giungeva, col suo procedimento, \textsc{Picard} (Comptes, rendus, 1905): Cfr. il \textit{Traité} di \textsc{Picard} e \textsc{Simart}, II, pag. 417.}. % [sic] Comptes, rendus -> Comptes rendus, apparently extraneous comma in the journal title.

Invece la considerazione dei sistemi continui di \(\infty^q\) curve dise%
%CONT>>
% scan: s184R
% !PRINTED 345
\unskip\weberpage{345}%
% head: SUPERFICIE IRREGOLARI E SISTEMI CONTINUI DI CURVE, ECC.
%<<CONT
% stamp: Scattered scan specks and small marginal strokes.
quivalenti permette a \textsc{Castelnuovo}\origfootnote{1}{s184R-1}{\textit{Sugli integrali semplici appartenenti ad una superficie irregolare.} Rendic. Lincei, Maggio-Giugno, 1905.} e a \textsc{Severi}\origfootnote{2}{s184R-2}{\textit{Il teorema d’Abel sulle superficie algebriche.} Annali di Mat., agosto, 1905.}, indipendentemente l’uno dall’altro, di precisare sotto l’aspetto quantitativo l’osservazione di \textsc{Humbert}, dimostrando che « le superficie irregolari, di irregolarità \(q=p_g-p_a\), posseggono proprio \(q\) integrali semplici di prima specie, con \(2q\) periodi ». La dimostrazione di questo teorema capitale, cui non arriva la considerazione dell’equazione lineare di Fuchs adoperata da \textsc{Picard}, si vedrà poi svolta per altra via da \textsc{Poincaré}, nella memoria citata, in cui ritrova, come abbiam detto, anche l’esistenza di sistemi continui con \(\infty^q\) curve disequivalenti.

Per la dimostrazione del teorema sugli integrali di cui si è detto innanzi, \textsc{Castelnuovo} introduce l’importante concetto di quella che egli ha chiamato la \textit{varietà di Picard} annessa ad una superficie irregolare (cfr. \S~3). Questa nozione permette in generale di costruire il sistema continuo definito da una curva \(C\), a partire da un altro qualunque sistema continuo \(\{D\}\) appartenente alla superficie. \textsc{Enriques} ha osservato che la costruzione si può precisare, in base al teorema di Riemann-Roch, traendone che « un sistema lineare di curve di caratteri \(n\), \(\pi\) ed \(i\), \textit{aritmeticamente effettivo}, cioè tale che \(p_a+n-\pi+1-i\geq 0\), appartiene sempre ad una serie \(\infty^q\) di sistemi lineari disequivalenti » (cfr. \S~2). L’osservazione fu dall’autore comunicata al \textsc{Severi}, e poichè questi era riuscito proprio allora a semplificare notevolmente la dimostrazione del teorema di Riemann-Roch, a lui lasciava l’\textsc{Enriques} di pubblicarla\origfootnote{3}{s184R-3}{\textsc{Severi}, \textit{Sul teorema di Riemann-Roch e sulle serie continue di curve appartenenti ad una superficie algebrica.} R. Accad. Torino, 1905.}.

Veniamo ora alla critica che ha scoperto il punto debole nella dimostrazione del teorema di Enriques sui sistemi continui. Abbiam detto che questa è stata mossa dal \textsc{Severi} nel 1921. Si trova in una serie di Note « Sulla teoria degli integrali semplici di prima specie appartenenti ad una superficie algebrica »\origfootnote{4}{s184R-4}{Rendic. Lincei, 1921.}, in cui viene esposta, semplificandola in qualche punto, la dimostrazione di \textsc{Poincaré} del 1910, comprendente insieme l’esistenza dei sistemi continui e degli integrali di prima specie\origfootnote{5}{s184R-5}{La dimostrazione di \textsc{Poincaré}, ulteriormente semplificata da \textsc{Lefschetz}, viene esposta da \textsc{O. Zariski}, \textit{Algebraic Surfaces} (Berlino, Springer, 1935). Cfr. VII, 4, 5 (pag. 128, 133).}. L’autore osserva che (grazie a \textsc{Poincaré}) almeno la parte essenziale del teorema (il teorema fondamentale come è enunciato qui nel \S~5) resta salvo. Più tardi si riprenderanno i tentativi per estendere questo teorema dimostrando
%CONT>>
% scan: s185L
% !PRINTED 346
\weberpage{346}
% head: CAPITOLO NONO
%<<CONT
% stamp: Dark right-edge band, scattered marginal marks, and a diagonal stroke near "particolare".
come sia sempre completa la serie caratteristica di un sistema continuo completo qualsiasi; ma di ciò più avanti (\S~7).

\textsc{Enriques}, in più Note del 1936-1937, su le « Curve infinitamente vicine sopra una superficie algebrica »\origfootnote{1}{s185L-1}{Rendic. Lincei, aprile 1936 e Rendic. Seminario Matematico della R. Università di Roma, 1936.} e su « La proprietà caratteristica delle superficie algebriche irregolari e le curve infinitamente vicine »\origfootnote{2}{s185L-2}{Rendic. Lincei, ottobre 1937.} ha tentato di colmare la lacuna segnalata nella dimostrazione del teorema fondamentale, cercando di provare che ad una curva \(C\) (quando la serie caratteristica di \(C\) abbia la deficienza \(q=p_g-p_a\)) sono infinitamente vicine, non solo \(\infty^q\) curve disequivalenti nell’intorno del prim’ordine, sì anche curve successive a queste nell’intorno del secondo ordine ecc. L’idea fondamentale di tale dimostrazione (che già appare nelle « Lezioni » di \textsc{Enriques} e \textsc{Campedelli}), consiste in ciò che « sopra una varietà abeliana si definisce, con \textsc{Lie}, una trasformazione infinitesima generatrice di un gruppo \(\infty^1\), a partire da due punti infinitamente vicini ».

Così, in particolare, nella varietà di Jacobi costituita dai gruppi di \(p\) punti di una curva del genere \(p\), si riesce a definire una serie \(\infty^1\) analitica di elementi (gruppi di punti) a partire da due gruppi infinitamente vicini. Questa costruzione si estende ai sistemi continui di curve disequivalenti su una superficie, quando si presupponga l’esistenza di un sistema i cui elementi formino la varietà (abeliana) di Picard. Le considerazioni sviluppate da \textsc{Enriques} mirano a giustificare tale costruzione, passando da curve infinitamente vicine nell’intorno del prim’ordine a curve nell’intorno di secondo ordine e così via, indipendentemente dal detto presupposto; anzi collo scopo di pervenire così a giustificarlo. Il punto di partenza è dunque che, vicino ad una curva spezzata \(C+D\), vi sono, nell’intorno del primo ordine, curve \(\bar{C}+\bar{D}\) spezzate in due componenti, una dell’ordine di \(C\) e un’altra dell’ordine di \(D\), aventi lo stesso numero \(m\) di punti comuni. Qui, avverte l’autore, conviene anzitutto giustificare l’applicazione implicita che si fa del principio di spezzamento nell’infinitesimo; ma ciò non offre difficoltà, in quanto è dato riconoscere che, uscendo dalla superficie, si può costruire, nello spazio, una serie di curve spezzate che contiene \(C+D\) e \(\bar{C}+\bar{D}\). % ?? uncertain: dot-like mark after "(gruppi di punti)" -- scan speck or printed full stop.

Giustificata così la \textit{realtà} dello spezzamento per le curve infinitamente vicine, e quindi l’esistenza effettiva di una \(\bar{C}\) disequivalente infinitamente vicina alla \(C\), la costruzione di \textsc{Enriques} porge una serie di operazioni su \(C+D\) e \(\bar{C}+\bar{D}\), che viene a dipendere soltanto da \(C\) e \(\bar{C}\), e conduce appunto a prolungare nei successivi intorni di \(C\) la serie di curve definita dalla \(\bar{C}\).

% scan: s185R
% !PRINTED 347
\weberpage{347}
% head: SUPERFICIE IRREGOLARI E SISTEMI CONTINUI DI CURVE, ECC.
% stamp: Scattered specks and small marginal marks; a small mark lies above the last minus sign.
In occasione della Nota del 1937, questa dimostrazione fu contestata da \textsc{B. Segre}.

Tuttavia \textsc{Enriques} (ribattendo le obiezioni del \textsc{Segre}) concedeva volentieri la preferenza ad un’altra dimostrazione basata semplicemente sul principio di continuità, che veniva proposta dal \textsc{Segre} stesso, salvo a discutere del suo buon fondamento: di questo interessante tentativo e degli sviluppi che esso ha ricevuto da parte di altri geometri, discorriamo espressamente nel seguente \S~7.

Soltanto dopo il fallimento del tentativo anzidetto abbiamo ripreso, com’è naturale, la dimostrazione criticata, che — con qualche precisazione e chiarimento — si è esposta nel precedente \S~5.

\edsection{7. Su diversi tentativi di dimostrazione e d’estensione del teorema\\fondamentale.}{7. Su diversi tentativi di dimostrazione e d’estensione del teorema fondamentale.}{sec.IX.7}

La dimostrazione del teorema fondamentale, che esprime la proprietà caratteristica delle superficie irregolari, è stata raggiunta in generale (per \(p_g\geq 0\)) mediante una considerazione di infinitesimi che supplisce ad un calcolo differenziale non eseguito. Considerazioni di tal natura, anche restando nel campo dell’analisi algebrica, sono d’ordine delicato e sollevano la diffidenza dei matematici che non hanno con esse sufficiente familiarità, o quanto meno lasciano l’impressione che si sia introdotto qualche elemento intuitivo, in contrasto colle esigenze del rigore logico.

Questa impressione era tanto più spiegabile di fronte a qualche passaggio un po’ rapido o meno giustificato delle Note citate: la cui esposizione\origfootnote{1}{s185R-1}{In ispecie seguendo la via indicata nelle Note del 1936.} si è ripresa qui cercando di precisare e chiarire i punti più delicati. Per tali motivi appariva e può apparire ancora desiderabile di fornire altre dimostrazioni del teorema medesimo, che s’impongano a tutti, senza dubbi ed esitazioni.

Abbiamo già detto che una dimostrazione siffatta, coll’uso degli integrali appartenenti alla superficie, è stata data dal \textsc{poincaré} (1910), ed oggi si può vedere, per esempio, nell’esposizione, semplificata del \textsc{Lefschetz}, riferita da \textsc{Zariski}. Ma, nell’ambito della teoria algebrico-geometrica, merita speciale menzione il tentativo che abbiamo accennato essere stato intrapreso da \textsc{B. Segre}\origfootnote{2}{s185R-2}{\textit{Un teorema fondamentale della geometria sulle superficie algebriche e il principio di spezzamento.} Annali di Mat., 1938.}. % ?? uncertain: initial of "poincaré" -- small capital p or full capital P.
% EDITORIAL D0175 RISOLTO: The initial p has the same small-cap height as the following letters; it is not a full-height capital. Retain lowercase input to textsc.

L’idea direttiva di questo autore si può esprimere come segue. Si consideri il sistema lineare \(|E|=|C+D|\) che contiene la curva spezzata \(E_0=C_0+D_0\), dotata di \(m\) punti doppi che formano il gruppo \(C_0D_0\); e s’impongano alle curve \(E\), vicine ad \(E_0\), \(m-p_g\)
%CONT>>
% scan: s186L
% !PRINTED 348
\weberpage{348}
% head: CAPITOLO NONO
%<<CONT
% stamp: Dark scan-border lines at the upper left and right edge; scattered marginal specks and short strokes.
punti doppi (vicini ad $m-p_g$ punti del detto gruppo); si vuol dimostrare che le curve $\bar{E}$ così definite sono necessariamente spezzate in una curva dell’ordine di $C_0$ e in una curva dell’ordine di $D_0$: e a tal uopo conviene ridurre all’assurdo l’ipotesi che le dette $\bar{E}$ siano irriducibili. Perciò si costruisca su una $\bar{E}$ con $m-p_g$ punti doppi assegnati, la serie caratteristica $g$ del sistema continuo $\{\bar{E}\}$ resa completa, e si passi al limite per $\bar{E}\to E_0$; la serie $g$ ha la dimensione
\[
2p_g+(n-\pi)+(\nu-\varrho)+1, % ?? uncertain: the symbols in the final parenthesis -- Greek \nu or Latin v; \varrho or \rho.
\]
ma, al limite, essa appare ridursi alla serie $g_c+g_d$ formata coi gruppi delle serie caratteristiche di $C_0$ e di $D_0$ (rese complete) cui si aggiungano $p_g$ punti fissi del gruppo $C_0D_0$ (i nuovi punti doppi acquistati dalla $E_0$ spezzata); orbene la dimensione della serie limite $g_c+g_d$ è
\[
2p_g+(n-\pi)+(\nu-\varrho), % ?? uncertain: the symbols in the final parenthesis -- Greek \nu or Latin v; \varrho or \rho.
\]
di una unità inferiore a quella di $g$!

L’autore affermava che questa diminuzione della dimensione di una serie è assurda, contrastando al principio di continuità. Ma \textsc{Enriques}, avendo fatto per conto suo un tentativo analogo (nel caso più semplice $p_g=1$), si era arrestato di fronte ad un imbarazzante paradosso, ed esponeva quindi la difficoltà incontrata\origfootnote{1}{s186L-1}{\textsc{Enriques}, \textit{Sulla proprietà caratteristica delle superficie algebriche irregolari.} Rendic. Lincei, Giugno, 1938.}. Si consideri, nel piano, il sistema delle curve $C_8$ d’ordine 8 con 15 punti doppi; fra queste $C_8$ vi sono le curve spezzate in due quartiche $C_4+\bar{C}_4$, che hanno comuni 15 punti $A_1\ldotp\ldotp\ldotp\ldotp A_{15}$ limiti dei punti doppi di $C_8$ e un altro sedicesimo punto $O$. Ora sopra una $C_8$ variabile, la serie canonica $g=g_{10}^5$ è segata dalle quintiche $C_5$ per i 15 punti doppi; al limite per $C_8\to C_4+\bar{C}_4$ la detta serie canonica si riduce alla serie composta dei gruppi delle due serie canoniche $g_4^2$ su $C_4$ e su $\bar{C}_4$, che viene segata dalle $C_5$ per $A_1\ldotp\ldotp\ldotp\ldotp A_{15}$, aventi di conseguenza il punto fisso comune $O$; così sembra che al limite la serie $g$ anzidetta si riduca ad una serie di dimensione inferiore. Nasce il dubbio che la conservazione della dimensione, che pareva conseguire dal principio di continuità, si trovi in difetto!

È merito del \textsc{Segre} di avere chiarito la difficoltà, spiegando il paradosso. Il principio di conservazione della dimensione d’una serie lineare consegue veramente da una esigenza di continuità che, spiegata in maniera esplicita, risulta irrefutabile; ma quando si cerca il limite della serie $g$, per $C_8\to C_4+\bar{C}_4$, bisogna considerare, non soltanto ciò che diviene la detta $g$ mentre una $C_5$ aggiunta a $C_8$ tende ad una $C_5$ generica per $A_1\ldotp\ldotp\ldotp\ldotp A_{15}$, bensì anche il limite di
%CONT>>
% scan: s186R
% !PRINTED 349
\weberpage{349}
% head: SUPERFICIE IRREGOLARI E SISTEMI CONTINUI DI CURVE, ECC.
%<<CONT
% stamp: Dark scan-edge marks and scattered marginal specks.
essa in corrispondenza ad una variazione per cui $C_8$ e $C_5$ tendano insieme a due curve degeneri contenenti come parte $C_4$ (o $\bar{C}_4$), dove i gruppi sezioni di $C_8$ e $C_5$ diventano indeterminati. In altre parole: l’esempio paradossale mette in luce l’esigenza che nel passaggio al limite della serie $g$ si tenga conto dell’\textit{intero limite}!

B. \textsc{Segre} spiega bene il valore di questa esigenza coll’esempio di un’iperbole che, variando, si riduca, al limite, alla coppia dei suoi asintoti: la $g_2^1$ segata sull’iperbole dalle rette per il centro, si riduce al punto fisso che è il centro, ma anche alle due simmetrie rispetto al centro stesso che si hanno sugli asintoti. In modo simile, nell’esempio di \textsc{Enriques}, tenuto conto che le curve $C_8$ passanti per $A_1\ldots A_{15}$ su $C_4$ non passano di conseguenza per $O$, il limite della $g=g_{10}^5$ viene dato da due serie di dimensione 5 che si ottengono associando gruppi di $C_4$ e di $\bar{C}_4$: per dire soltanto di una di queste serie, essa si costruisce associando ai gruppi canonici di $\bar{C}_4$ i gruppi di una $G_6^3$ su $C_4$, che non ha più il punto fisso $O$, ma si definisce ponendo $g_6^3=g_4^2+2O$. % [sic] C_8 -> C_5, the cutting curves are called quintics C_5 on p. 348.

Spiegato così l’imbarazzante paradosso, B. \textsc{Segre} riconosceva nella propria dimostrazione del teorema fondamentale una lacuna, e si accingeva a colmarla colla ricerca di un criterio atto a determinare il limite dell’intersezione di due curve, che, variando su una superficie, vengano ad acquistare una parte comune. A questo criterio infinitesimale egli accennava in una Nota « Sur un théorème fondamental de géométrie sur les surfaces algébriques » del 9 gennaio 1939\origfootnote{1}{s186R-1}{\textit{Comptes-rendus}, 9 Janvier 1939.}, cui si accompagnava un’altra Nota di \textsc{Enriques}\origfootnote{2}{s186R-2}{\textit{Sur la propriété caractéristique des surfaces algébriques irrégulières.} Comptes-rendus, 3 Janvier 1939.} che con piacere salutava il successo del Collega.

Frattanto B. \textsc{Segre} lasciava l’Italia e non consta che sia uscita fino ad ora una sua pubblicazione più diffusa sull’argomento.

Successivamente il \textsc{Severi} in una memoria dell’Accademia d’Italia\origfootnote{3}{s186R-3}{\textit{La teoria generale dei sistemi continui di curve sopra una superficie algebrica.} Memorie dell’Accademia d’Italia, vol. XIX, 1941}, offre una revisione critica dell’intera teoria dei sistemi continui di curve sopra le superficie algebriche. Disgraziatamente questa esposizione è affetta, da alcuni errori che, in parte, sono stati segnalati da \textsc{Enriques} in una Nota « Sui sistemi continui di curve appartenenti ad una superficie algebrica », e che qui occorre riprendere in esame. % ?? uncertain: the mark after “affetta” -- a comma or a scan speck.
% EDITORIAL D0178 RISOLTO: The enlarged scan shows a distinct comma-shaped mark below the baseline following affetta. Retain the printed punctuation.

Per soddisfare all’esigenza dimostrativa che si riferisce al ragionamento del \textsc{Segre}, come si è spiegato innanzi, il \textsc{Severi} riprende da questi — in una forma alquanto modificata — un
%CONT>>
% scan: s187L
% !PRINTED 350
\weberpage{350}
% head: CAPITOLO NONO
%<<CONT
% stamp: Dark vertical scan border at the right; scattered specks and short marginal strokes.
lemma di geometria differenziale, da cui dovrebbe risultare che la serie caratteristica $g$ del sistema continuo $\{\bar{E}\}$ su $\bar{E}$ (resa completa), passando al limite per $\bar{E}\to E_0=C_0+D_0$, avrebbe necessariamente come punti fissi i $p_g$ punti $P$ del gruppo $C_0D_0$ che non sono limiti degli $m-p_g$ punti doppi di $\bar{E}$, e ciò in qualunque modo sia fatto il passaggio al limite, anche quando le curve secanti su $\bar{E}$ la $g$ tendano a contenere come parte la $C_0$, di guisa che il $\mathit{lim}\,g$ risulti indeterminato. Il lemma di Segre-Severi che permetterebbe di dedurre questa conclusione, esprime una proprietà di geometria differenziale, per cui i detti punti $P$ sarebbero limiti di punti appartenenti a gruppi della serie $G$. Questa proprietà differenziale discenderebbe a sua volta da quella, riconosciuta da \textsc{Enriques} nella dimostrazione originale del teorema fondamentale, che « le curve infinitamente vicine alla $C_0$ segano su di essa, la serie caratteristica completa »; in tal guisa sarebbe colmata la lacuna nella dimostrazione di quel teorema.

Ora \textsc{Enriques}, che già aveva trovato qualche difficoltà a ricostruire il criterio infinitesimale del \textsc{Segre} sulle indicazioni fornitegli dall’A., non riuscì a comprendere questo passaggio e perciò dichiarò il detto lemma « oscuro e non convincente »\origfootnote{1}{s187L-1}{Commentarii Mathematici Helvetici: entrata in redazione il 20 Marzo 1942.}. Ma credette che si potesse tuttavia riuscire allo scopo facendo segare la serie $\bar{g}$ sulla $\bar{E}$ variabile da curve che si mantengono irriducibili nel passaggio al limite di cui è questione.

In ciò egli s’ingannava e \textsc{Severi} ha giustamente denunziato il suo errore\origfootnote{2}{s187L-2}{\textit{Intorno ai sistemi continui di curve sopra una superficie algebrica.} Commentarii Mathematici Helvetici (30 Novembre 1942). A questo articolo la direzione dei Commentari non consentì ad \textsc{Enriques} di rispondere, e perciò questi non potè allora nè riconoscere il proprio errore, nè dichiarare gli errori in cui, da parte sua, il suo critico era caduto.}. Senonchè il riconoscimento di questo getta un dubbio su tutta la dimostrazione proposta da \textsc{Segre} e \textsc{Severi}. Non si tratta soltanto di correggere il lemma denunciato nella Memoria di \textsc{Severi}, come ha fatto l’A. stesso\origfootnote{3}{s187L-3}{Nota presentata alla R. Acc. d’Italia, 16 Gennaio 1942.} in seguito ad un esempio di \textsc{Bompiani}; ma piuttosto di chiedersi se una proprietà differenziale quale viene espressa nel lemma anzidetto valga veramente ad escludere che il limite di $g$ possa dare origine su $C_0$ ad una serie senza punti fissi. Chi approfondisca l’analisi della questione si persuaderà agevolmente che questo dubbio non può essere eliminato, anzi che c’è qui una radicale insufficienza del lemma, cioè che « questo lemma di Segre-Severi non risponde allo scopo per cui viene introdotto ». Non insistiamo su ciò, avendo ragione di ritenere che il \textsc{Severi} stesso abbia ormai riconosciuto l’insufficienza del suo
%CONT>>
% scan: s187R
% !PRINTED 351
\weberpage{351}
% head: SUPERFICIE IRREGOLARI E SISTEMI CONTINUI DI CURVE, ECC.
%<<CONT
% stamp: A dark scan-edge line at the left and scattered marginal specks.
ragionamento. Ad ogni modo conviene avere accennato qui ad una idea, che — se non oggi — potrebbe avere domani un’applicazione fruttuosa.

Abbiamo accennato che il \textsc{Severi} (e prima di lui B. \textsc{Segre}) ha tentato di estendere il teorema fondamentale sulle superficie irregolari, stabilendo che, sopra una superficie qualsiasi « la serie caratteristica di un sistema continuo irriducibile senza punti multipli (comunque speciale o sovrabbondante) è sempre completa ».

Comunque oggi l’A. ha riconosciuto, traverso esempii, che il risultato più esteso che aveva in vista patisce qualche eccezione\origfootnote{1}{s187R-1}{Ed anche patisce eccezione \textit{l’unicità} del sistema continuo cui appartiene una curva irriducibile, che l’A. aveva creduto di dimostrare colle considerazioni svolte a fine di pag. 378 e al principio di pag. 379 nella memoria su \textit{La teoria generale ecc.}, del 1941.}.

Ma si restringa quanto occorre la scelta dei sistemi continui di curve irriducibili $\{C\}$, in guisa che la proprietà sussista \textit{per ipotesi}: e sia dunque $\{C\}$ un sistema continuo $\infty^{d+1}$ che abbia su una $C$ la serie caratteristica completa di dimensione $d$; conviene precisare diverse circostanze che possono presentarsi.

Il sistema $\{C\}$ conterrà un’infinità di sistemi lineari di curve disequivalenti, se $d+1$ superi la dimensione del sistema lineare completo $|C|$, cioè se questo sistema di dimensione
\[
r=p_a+n-\pi+1-i+\omega \qquad (\text{con }\omega\geq0)
\]
abbia la serie caratteristica deficiente, di deficienza
\[
q-\varepsilon \qquad (q=p_g-p_a).
\]

Ora giova distinguere i \textit{casi elementari} abnormi, che, colle loro combinazioni, danno luogo a tutti i casi possibili. A priori può accadere:

1) che il sistema $|C|$ sia regolare e speciale d’indice di specialità $i$:
\[
0<i<p_g,
\]
colle notazioni precedenti:
\[
\omega=\varepsilon=0;
\]

2) che $|C|$ sia non speciale e sovrabbondante, di sovrabbondanza $\omega>0$, mentre la sua serie caratteristica abbia la deficienza massima $q=p_g-p_a$: $i=0$, $\varepsilon=0$.

3) che $|C|$ sia non speciale e sovrabbondante di sovrabbondanza $\omega=\varepsilon$, per avere serie caratteristica di deficienza $q-\varepsilon$: $i=0$, $\varepsilon>0$.

% scan: s188L
% !PRINTED 352
\weberpage{352}
% head: CAPITOLO NONO
% stamp: Dark horizontal mark in the upper-left margin, dark right scan border, and scattered marginal specks and strokes.
Nel \textit{caso} 1) (ammesso possibile) la dimensione di $|C|$ vale
\[
r=p_a+n-\pi+1-i,
\]
e quella di $\{C\}$ è $r+q$, quindi \textit{il sistema continuo è formato di $\infty^q$ sistemi lineari} (in generale) \textit{collo stesso indice di specialità $i$.}

Nel \textit{caso} 2) la dimensione di $|C|$ essendo $r=p_a+n-\pi+1+\omega$, quella di $\{C\}$ risulta ancora $r+q$, quindi $\{C\}$ \textit{è formato di $\infty^q$ sistemi lineari colla medesima sovrabbondanza.} % The plus sign is repeated on both sides of the printed line break.

Ma nel \textit{caso} 3), la dimensione di $|C|$ essendo $r=p_a+n-\pi+1+\varepsilon$, quella di $\{C\}$ risulta $r+q-\varepsilon$, e non si può dire a priori se, per esempio, $\{C\}$ \textit{sia formato di $\infty^q$ sistemi lineari regolari aventi generalmente la dimensione $r-\varepsilon$, fra cui si trovi il $|C|$ sovrabbondante di dimensione $r$, ovvero da una serie $\infty^{q-\varepsilon}$ di sistemi lineari tutti colla medesima sovrabbondanza $\omega=\varepsilon$.} % The plus sign is repeated on both sides of the printed line break.

Abbiamo così nominato soltanto i due casi estremi fra quelli che nell’ipotesi 3) appariscono a priori possibili. La possibilità effettiva che si presentino queste diverse ipotesi, e in particolare la prima (che un errore assai riposto aveva condotto a negare) è stata riconosciuta nella citata Nota critica di \textsc{Enriques}\origfootnote{1}{s188L-1}{Commentarii Mathematici Helvetici, 20 Marzo 1942.}, ove trovasi approfondito lo studio dell’esempio che qui viene brevemente accennato.

Si consideri la superficie rigata $F_5$ del $5^\circ$ ordine che ha una retta tripla e una retta doppia: la $F_5$ ha il genere geometrico $p_g=0$, il genere numerico $p_a=-p=-2$ e l’irregolarità $q=2$.

Sopra $F_5$ il sistema $|C|$ delle sezioni piane è sovrabbondante di sovrabbondanza $\varepsilon=1$, avendo il grado $n=5$, il genere $\pi=2$, e la dimensione $r=p_a+n-\pi+1+\varepsilon=3$, mentre la deficienza della serie caratteristica è $q-\varepsilon=1$.

Il sistema lineare $|C|$ è contenuto in un sistema continuo $\{C\}$ di dimensione 4, e si può vedere che questo è formato non già di $\infty^1$ sistemi lineari $\infty^3$, bensì di $\infty^2$ sistemi lineari regolari $\infty^2$, fra cui si trova il detto $|C|$ di dimensione 3. Per verificare tale asserzione si sommerà al sistema $|C|$ una coppia di generatrici $a\ b$ e mercè il sistema lineare così ottenuto (che è $\infty^6$) si costruirà una superficie $F_9$ d’ordine 9, rigata normale nello spazio $S_6$, di cui la $F_5$ appare come proiezione da una coppia di generatrici incidenti $a'\ b'$. La $F_9$ ha un punto doppio $O$ nel punto d’incrocio di $a'$ e $b'$, ma non possiede una linea doppia; perciò non vi sono per essa infinite coppie di generatrici incidenti, e così gli $\infty^2$ sistemi lineari costituenti $\{C\}$ cioè i sistemi segati su $F_9$ dagli iperpiani per una coppia di generatrici, sono, in generale, di dimensione 2, \hfill c. d. d.

% scan: s188R
% !PRINTED 353
\weberpage{353}
% head: SUPERFICIE IRREGOLARI E SISTEMI CONTINUI DI CURVE, ECC.
% stamp: Vertical scan lines near the right edge and scattered marginal specks.
D’altra parte si può indicare su $F_9$ un sistema continuo $\{L\}$ formato di $\infty^1$ sistemi lineari sovrabbondanti, tutti di dimensione $3\ (=2+1)$: tale è il sistema continuo segato dagli iperpiani passanti per il punto doppio $O$ e per una generatrice della rigata. E conviene avvertire che questo esempio deve ritenersi rientrare nel caso 3), e non nel caso 2), perchè la serie caratteristica di $\{L\}$, è una serie non speciale $g_5^3$.

\textit{Osservazione.} -- Gli $\infty^1$ sistemi lineari $|L|$ formanti $\{L\}$, sono — sulla rigata $F_5$ — i residui delle generatrici rispetto al sistema lineare $\infty^5$
\[
|C+a+b-A-B|,
\]
che ha due punti base, $A$ e $B$, rispettivamente su $a$ e $b$, i quali costituiscono una coppia neutra per il sistema
\[
|C+a+b|.
\]

Ora, se si cambiano $a$ e $b$ in altre due generatrici $\bar{a}$ e $\bar{b}$, la coppia $A\ B$ non sarà più neutra per $|C+\bar{a}+\bar{b}|$, e quindi la sottrazione di una generatrice da questi sistemi lineari condurrà ad $\infty^2$ sistemi lineari $|\bar{L}|=|C+\bar{a}+\bar{b}-A-B|$.

I due sistemi, $\{L\}$ formato di $\infty^1$ sistemi lineari ed $\{\bar{L}\}$ formato di $\infty^2$ sistemi lineari, sono distinti, cioè $\{L\}+\{\bar{L}\}$ è riducibile in due varietà $\infty^4$ come sistema di curve, ed invece dovrà ritenersi irriducibile come varietà $\infty^2$ di sistemi lineari $|\bar{L}|$ che contiene entro di sè la varietà $\infty^1$ degli $|L|$. Si esprime questa circostanza dicendo che il detto sistema $\{L\}$ è \textit{esorbitante}. Il fenomeno della esorbitanza (di un sistema lineare o di una serie continua di sistemi lineari) è stato osservato da \textsc{Rosenblatt} e \textsc{Albanese} e poi da \textsc{Severi}\origfootnote{1}{s188R-1}{\textit{Nuovi contributi ecc.} Lincei, 1916.}.

Altri esempi di sistemi esorbitanti sono dati dal sistema delle sezioni piane d’un cono, affatto generale, d’ordine $n\geq3$.

Per $n=3$ il sistema lineare $\infty^3$ delle sezioni piane è contenuto nella serie degli $\infty^1$ sistemi lineari ($\infty^2$) formati dalle terne di generatrici del cono, cui va sommato l’intorno del vertice.

Per $n=4$ il sistema lineare $\infty^3$ delle sezioni piane del cono è contenuto nella serie degli $\infty^3$ sistemi lineari ($\infty^1$) formati dalle quaterne di generatrici del cono, cui va sommato l’intorno del vertice.

Per $n\geq5$ il sistema lineare $\infty^3$ delle sezioni piane del cono è contenuto nella serie degli $\infty^n$ sistemi lineari ($\infty^0$) formati dai gruppi di $n$ generatrici del cono, cui si sommi sempre l’intorno del vertice.

% EDITORIAL: segnatura originale: \noindent\textsc{Enriques F.} - \textit{Superficie algebriche.}\hfill 23
\printersignature{\textsc{Enriques F.} - \textit{Superficie algebriche.}\quad 23}

% scan: s189L
% !PRINTED 354
\weberpage{354}
% head: CAPITOLO NONO
% stamp: Dark vertical scan border at the right; scattered specks and short marginal strokes.
\edsection{8. Il sistema paracanonico.}{8. Il sistema paracanonico.}{sec.IX.8}

Secondo il teorema fondamentale (\S~4) e la condizione numerica del \S~2, il sistema canonico $|K|$ di una superficie $f$, di genere numerico $p_a\geq0$ e di genere geometrico $p_g>p_a$, sarà sempre contenuto in un sistema continuo $\infty^q$ ($q=p_g-p_a$) di sistemi lineari distinti, di dimensione
\[
r\geq p_a+p^{(1)}-1-p^{(1)}+1=p_a
\]

Si è creduto di poter precisare la dimensione di questo \textit{sistema paracanonico} $\{K\}$, determinando la deficienza della serie caratteristica del sistema canonico $|K|$, almeno nell’ipotesi che \textit{le curve canoniche pure sieno irriducibili}, e senza punti multipli, di genere $p^{(1)}>1$\origfootnote{1}{s189L-1}{\textsc{Enriques}, \textit{Des courbes paracanoniques appartenant à une surface algébrique irrégulière.} Soc. Royale des Sciences de Liège (19 Ottobre, 1939).}.

La detta deficienza sarebbe $q=p_g-p_a$ e quindi il sistema paracanonico (ammettendo che abbia su una $K$ la serie caratteristica completa) avrebbe la dimensione $2p_g-p_a-1$. Ma la \textit{deduzione} è \textit{erronea}\origfootnote{2}{s189L-2}{Pare che cada sostanzialmente nel medesimo errore anche la dimostrazione dello stesso teorema che viene offerta da \textsc{Severi} nella citata memoria su \textit{La teoria generale dei sistemi continui ecc}, del 1941.} siccome si riconosce riprendendo il ragionamento dell’A. e mettendo in luce un’obiezione d’ordine delicato che viene ad infirmarlo.

A tal uopo si assuma sulla superficie $F$ un sistema lineare regolare e non speciale $|C|$, di grado $n$ e genere $\pi$, il cui aggiunto $|C'| = |C+K|$ avrà la dimensione % The equality sign is repeated on both sides of the printed line break.
\[
p_a+\pi-1
\]
e segherà su una $C$ una serie canonica deficiente, di dimensione % ?? uncertain: the damaged n in “canonica” -- canonica or canoica.
\[
\pi-1-q.
\]

Indicando con $G=CK$ il gruppo di $2\pi-2-n$ punti comuni ad una $C$ e ad una $K$, si vuol calcolare il numero delle condizioni che il $G$ offre alle curve di $|C'|$ che debbono contenerlo. E per ciò si assume che le $C'$ contenenti $G$ seghino su $C$ la serie caratteristica di $|C|$ resa completa, ossia la serie caratteristica $g$ di $\{C\}$, segata sulla detta $C$ dalle $\bar{C}$ (in generale disequivalenti) ad essa infinitamente vicine. L’assunzione si vorrebbe giustificare dicendo che, fra le curve residue di una $\bar{C}$ rispetto a $|C'|$, vi sarebbe almeno una curva $\bar{K}$ vicina a $K$, sicchè $C+K$ e $\bar{C}+\bar{K}$ determinerebbero un
%CONT>>
% scan: s189R
% !PRINTED 355
\weberpage{355}
% head: SUPERFICIE IRREGOLARI E SISTEMI CONTINUI DI CURVE, ECC.
%<<CONT
% stamp: Scattered specks and short marginal strokes.
fascio di curve $C'$ per $G$, passanti per un gruppo di $g$ arbitrariamente dato.

Ma qui appunto sta l’errore! Se i sistemi lineari paracanonici hanno una dimensione inferiore a quella, $p_g-1$, del sistema canonico $|K|$, un sistema paracanonico $|\bar{K}|$ vicino a $K$ non possiede una curva $\bar{K}$ vicina ad una $K$ qualsiasi: il limite di un $|\bar{K}|$ è un sistema lineare minore, contenuto in quest’ultimo.

Abbiamo spiegato l’errore che infirma il ragionamento precedente, da cui seguirebbe la deficienza della serie canonica di $|K|$ essere $q=p_g-p_a$. Ora si può vedere da effettivi esempi che la conclusione stessa è lungi dall’essere verificata. Anzi sembra \textit{probabile} che « sopra una superficie d’irregolarità $q=p_g-p_a>1$, ($p_a\geq0$) con curve canoniche $K$ irriducibili, senza punti multipli, \textit{gli $\infty^q$ sistemi lineari paracanonici siano, in generale, regolari} e quindi il sistema continuo $\{K\}$ sia $\infty^{p_g}$: la \textit{serie} caratteristica di $|K|$, (almeno se questo sistema non è esorbitante) avrebbe quindi la \textit{deficienza} 1 ».

Vediamo come ciò si verifichi per la superficie rappresentativa delle coppie di punti d’una quartica $L$ di genere $q=3$.

Questa superficie $F$ ha il genere geometrico $p_g=3$ e il genere numerico $p_a=0$\origfootnote{1}{s189R-1}{Cfr. F. \textsc{Severi}, \textit{Sulle superficie che rappresentano le coppie di punti di una curva algebrica.} Atti R. Acc. di Torino, 25 Gennaio 1903.}: le curve canoniche di $F$ sono date dalla serie delle coppie di punti di $L$ estratte dalle quaterne delle $g_4^1$ speciali, e perciò sono $\infty^2$; invece le curve paracanoniche sono date dalle serie di coppie di punti estratte dalle quaterne delle $g_4^1$ non speciali, e perciò formano un sistema continuo $\infty^3$, costituito di $\infty^3$ curve disequivalenti, entro cui si trova il sistema sovrabbondante $|K|$.

Si può verificare facilmente che i sistemi paracanonici $|K|$ sono, in generale, di dimensione $p_a=0$ e non $>0$; altrimenti la serie caratteristica $g_6^x$ di $\{K\}$ su una $K$ (di genere $p^{(1)}=7$) avrebbe la dimensione
\[
x\geq3,
\]
mentre una curva di genere superiore a 4, non iperellittica, quale è la $K$, non può contenere una $g_6^3$.

Ciò che si è dimostrato per la superficie rappresentativa delle coppie di punti d’una quartica $L$ di genere $q=3$, sembra estendersi alle superficie $F$ rappresentative delle coppie di punti d’una curva $L$ di genere $q=4,\,5\ldots$.

Per $q=4$, la $L$ venendo realizzata da una sestica di $S_3$, le curve canoniche $K$ di $F$ saranno date dalle rigate intersezioni della congruenza delle corde di $L$ coi complessi lineari di $S_3$ e perciò si avrà
\[
p_g=6;
\]
%CONT>>
% scan: s190L
% !PRINTED 356
\weberpage{356}
% head: CAPITOLO NONO
%<<CONT
% stamp: Dark right scan border, scattered specks, and short marginal strokes.
invece si avranno curve paracanoniche $\bar{K}$ di $F$ in corrispondenza alle serie delle coppie di punti di $L$ estratte dalle sue $g_6^1$ non speciali. Se, come crediamo, si verifichi che alle $g_6^1$ di una $g_6^2$ non speciale rispondono le curve $\bar{K}$ d’un sistema lineare $|\bar{K}|$ completo, il sistema continuo paracanonico $\{\bar{K}\}$ di $F$ avrà la dimensione
\[
p_g=p_g-p_a+p_a=4+2=6
\]
contenendo in sè $\infty^4$ sistemi lineari regolari di dimensione
\[
p_a=2.
\]

Però il sistema canonico $\infty^5$ non sarà contenuto in codesto sistema paracanonico totale $\{\bar{K}\}$, ma dovrà essere \textit{esorbitante}, perchè la rigata delle corde di $L$ che appartengono ad un complesso lineare speciale non può essere limite di una rigata le cui generatrici congiungano le coppie di una $g_6^1$ (non speciale) della $L$ stessa.

Anche per $q=5$, la $L$ essendo realizzata in $S_4$ da una curva d’ordine 8 (base di una rete di quadriche) sembra che la superficie $F$, rappresentativa delle coppie di punti di $L$, possegga un sistema paracanonico $\{K\}$ formato di $\infty^q=\infty^5$ sistemi lineari regolari di dimensione
\[
p_a=5 \qquad (p_g=10),
\]
se si verifichi che le $g_8^1$ estratte da $g_8^3$ non speciali formino un sistema (non lineare) contenuto in un sistema lineare $\infty^5$ completo.

Gli esempi citati, di cui occorre approfondire lo studio, suggeriscono diversi problemi intorno alle superficie irregolari, in ispecie « se il sistema paracanonico $\{K\}$ di una superficie con $p_a\geq0$ e $p_g>p_a$, abbia la dimensione $p_g$ e consti, in generale, di $\infty^{p_g-p_a}$ sistemi lineari regolari $\infty^{p_a}$ ».

Ma almeno il caso in cui il genere lineare assoluto $p^{(1)}=1$, e quindi le curve canoniche, per $p_g>2$, sono riducibili, componendosi colle curve ellittiche d’un fascio, dà luogo ad una eccezione.

In questo caso si ha da fare con superficie $F$ di irregolarità $q=p_g-p_a$, possedenti un fascio di curve ellittiche di genere $q$; infatti si può provare che il sistema continuo $\{C\}$, a cui appartiene una qualunque curva $C$, sega sulle $k$ (componenti delle curve canoniche $K$) una serie di gruppi equivalenti: se così non fosse si troverebbe oltre il fascio $|k|$ un secondo fascio ellittico di curve e si dedurrebbe che le $k$ sono curve di modulo costante e che $p_a=-1$ (cfr. Cap. X). % The minus sign is repeated on both sides of the printed line break.

Ciò posto, se $p_a>0$, le curve canoniche di $F$ si comporranno di $p_g-1>q-1$ curve $k$, formanti entro il fascio una serie non speciale: i sistemi paracanonici saranno dunque egualmente di dimensione $p_g-1$, dando luogo ad un sistema continuo sovrabbondante di dimensione % ?? uncertain: the blotched sign in p_g-1>q-1 -- minus or plus.
% EDITORIAL D0180 RISOLTO: The enlarged sign has a horizontal bar with an irregular blot below it, not a centred vertical plus stroke. The reading is minus.
\[
2p_g-p_a-1.
\]

% scan: s190R
% !PRINTED 357
\weberpage{357}
% head: SUPERFICIE IRREGOLARI E SISTEMI CONTINUI DI CURVE, ECC.
% stamp: Scattered specks and marginal strokes, with faint vertical scan lines near the lower right edge.
Invece, per $p_a=0$, il sistema canonico $|K|$ sarà il sistema lineare $\infty^{p_g-1}$ formato dalle $2q-2=2p_g-2$ curve dei gruppi canonici entro il fascio $|k|$, ed i sistemi lineari paracanonici avranno la dimensione $p_g-2$, dando luogo ad un sistema continuo $\{K\}$ di dimensione
\[
2p_g-p_a-2.
\]

\textit{Osservazione.} -- Qui conviene aggiungere che per $p^{(1)}=1$ e $p_a=0$, non si presenta il caso di superficie irregolari (di genere $p_g=1$) con curva canonica pura d’ordine zero. Infatti \textit{per una superficie $F$ di genere $p_g=1$ con curva canonica d’ordine zero, non si può avere $p_a=0$}\origfootnote{1}{s190R-1}{Cfr. \textsc{Enriques}, \textit{Intorno alle superficie algebriche di genere lineare $p^{(1)}=1$.} Rendic. Accademia Bologna, 9 dicembre 1906 (n. 1).}. Giacchè un sistema lineare $|C|=|C'|$ di genere $\pi$ e grado $n=2\pi-2$, sopra la $F$, darebbe luogo ad una serie continua $\infty^1$ di sistemi lineari di dimensione $\pi-1$, e uno di questi diversi da $|C|$ e non contenente $|C|$, dovrebbe segare sopra una $C$ una serie non speciale di grado $2\pi-2$ e di dimensione $\pi-1$.

\edsection{9. Digressione sulle varietà abeliane.}{9. Digressione sulle varietà abeliane.}{sec.IX.9}

Dicesi varietà abeliana (di rango 1) a $q$ dimensioni, $V_q$, una varietà (algebrica) che ammette una rappresentazione parametrica mediante funzioni (abeliane, cioè) $2q$ volte periodiche di $q$ variabili indipendenti. rappresentazione univoca nel prismatoide dei periodi: se queste funzioni, come si suppone, non sono degeneri, la $V_q$ ammette un gruppo $\infty^q$ di trasformazioni permutabili semplicemente transitivo, senza eccezione: vuol dire che dati due punti (semplici) \textit{qualunque} di $V_q$, esiste \textit{una} trasformazione birazionale di $V_q$ che muta il primo punto nel secondo. % [sic] indipendenti. rappresentazione -> indipendenti, rappresentazione, the printed full stop is followed by a lower-case continuation.
Il più semplice esempio di varietà abeliana è la varietà dei gruppi di $p(=q)$ punti, $G_p$; di una curva $C$ di genere $p$, ossia la \textit{varietà di Jacobi} corrispondente alla curva\origfootnote{2}{s190R-2}{Cfr., p. es., \textsc{Enriques-Chisini}, \textit{Lezioni}, Libro VI, cap. III, \S~36. vol. IV, pag. 223.}. % ?? uncertain: punctuation after G_p -- semicolon or comma with an ink speck.
% EDITORIAL D0181 RISOLTO: Enlargement shows a separate round dot above a comma tail: a semicolon, as transcribed.

Questa $V_p$ ammette un gruppo misto semplicemente transitivo di trasformazioni birazionali in sè stessa che viene generato dalle $\infty^p$ trasformazioni involutorie di seconda specie definite dalle $g_{2p}^p$ appartenenti alla curva e dalle $\infty^p$ trasformazioni permutabili di prima specie, costituenti un gruppo continuo $\Gamma_p$, definite come prodotti di un numero pari di trasformazioni di seconda specie. La detta varietà di Jacobi, $V_p$, contiene una \textit{varietà eccezionale}, $V_{p-1}$, corrispondente alla serie dei gruppi $G_p$ speciali sopra la curva; tale $V_{p-1}$ viene trasformata in una $V_{p-2}$ di dimensione $p-2$ dalle
%CONT>>
% scan: s191L
% !PRINTED 358
\weberpage{358}
% head: CAPITOLO NONO
%<<CONT
% stamp: Dark right scan border and scattered specks.
trasformazioni di prima o di seconda specie, per esempio dalla involuzione fra gruppi \(G_p\) definita su \(C\) da una \(g^p_{2p}\)\origfootnote{1}{s191L-1}{Cfr. \textsc{Enriques-Chisini}, \textit{Lezioni}, Libro VI, cap. II, § 27 (vol. IV, pag. 143).}.

Quando fra i \(p\) integrali indipendenti di prima specie della curva \(C\) di genere \(p\) ve ne sono \(q<p\) (linearmente indipendenti) riducibili con \(2q\) periodi, entro la \(V_p\) di Jacobi che corrisponde alla curva si trova una schiera \(\infty^{p-q}\) di varietà abeliane \(V_q\) costituenti un sistema d’imprimitività per il \(\Gamma_p\), ciascuna ammettente un gruppo algebrico \(\infty^q\) semplicemente transitivo di trasformazioni birazionali in sè ed anche, per il teorema di Picard-Poincaré, una schiera analoga \(\infty^q\) di varietà abeliane \(V_{p-q}\)\origfootnote{2}{s191L-2}{Cfr. p. es., \textsc{Enriques-Chisini}, \textit{Lezioni}, Libro VI. Interpretazione geometrica di Castelnuovo. Cap. III, 4, 38 (vol. IV, pag. 234).}.

\textit{Varietà abeliane} \(V_q\) di questo tipo, contenute imprimitivamente entro una \(V_p\) di Jacobi, sono le \textit{varietà di Picard} annesse a superficie algebriche irregolari\origfootnote{3}{s191L-3}{Cfr. \textsc{R. Torelli}, \textit{Sulle serie algebriche semplicemente infinite di gruppi di punti appartenenti a una curva algebrica}, Rend. Circ. Mat. di Palermo, t. XXXVII, 1914 (§ II, n. 11).}: infatti, sopra una superficie \(F\) irregolare, di irregolarità \(q\), le curve disequivalenti di un sistema continuo \(\infty^q\) segano, sopra una curva di genere \(p\) (\(>q\)) una serie formata di \(\infty^q\) gruppi di punti disequivalenti o di \(\infty^q\) serie lineari, che costituisce appunto una varietà abeliana \(V_q\) entro la \(V_p\) di Jacobi corrispondente alla curva.

La precedente osservazione è di portata generale perchè si dimostra\origfootnote{4}{s191L-4}{Cfr. \textsc{G. Castelnuovo - F. Enriques}, \textit{Sur les intégrales simples....}, Ann. de l’école normale sup., s. 3, t. XXIII,, 1906.}, ma con procedimenti trascendenti, che ogni varietà abeliana (di rango \(1\)) è una varietà di Picard annessa ad una qualche superficie algebrica irregolare; di modo che si può subito enunciare il teorema di \textsc{Torelli}: \textit{Esistono sempre varietà jacobiane contenenti imprimitivamente una data varietà abeliana.} % ?? uncertain: dots after "simples" -- four or five.
% EDITORIAL D0182 RISOLTO: Four baseline dots followed by the surrounding comma are visible. The transcription already preserves the four dots. % ?? uncertain: punctuation after "XXIII" -- two commas, period and comma, or comma with a scan speck.
% EDITORIAL D0183 RISOLTO: The enlargement shows two separate comma tails, not a period followed by a comma. Retain the printing and its sic proposal. % [sic] "XXIII,," -> "XXIII,", apparently duplicated comma after the volume number.

Se però non si vuole far ricorso al risultato, ottenuto per via trascendente, di \textsc{Castelnuovo-Enriques}, si possono svolgere le considerazioni seguenti che portano ad un teorema meno preciso di quello di \textsc{Torelli}, ma sufficiente per i successivi sviluppi della nostra esposizione.

Ponendoci dal punto di vista geometrico più generale \textit{definiremo} come \textit{varietà abeliana} a \(q\) dimensioni \(V_q\), una \(V_q\) che possegga un \textit{gruppo \(\infty^q\) \(\Gamma\) di trasformazioni birazionali in sè, semplicemente transitivo senza eccezione}, tale dunque che esista \textit{una} trasformazione non degenere di \(\Gamma\) in cui si corrispondono \textit{due punti} (semplici) \textit{qualunque di \(V_q\)}.

% scan: s191R
% !PRINTED 359
\weberpage{359}
% head: SUPERFICIE IRREGOLARI E SISTEMI CONTINUI DI CURVE, ECC.
% stamp: Scattered scan specks.
La proprietà del gruppo \(\Gamma\), di essere semplicemente transitivo senza eccezione sulla \(V_q\), porta di conseguenza che: \textit{non esiste sulla \(V_q\) alcun sistema lineare di dimensione \(>0\) di varietà \(V_{q-1}\), che sia trasformato transitivamente in sè stesso da \(\infty^r\) (\(r>0\)) trasformazioni di \(\Gamma\).} In altre parole: una varietà \(V_{q-1}\) su \(V_q\), la quale non sia trasformata in sè stessa da un sottogruppo continuo del \(\Gamma\), non può essere trasformata dal \(\Gamma\) nelle varietà di un sistema lineare \(\lvert V_{q-1}\rvert\).

Si abbia, sulla varietà abeliana \(V_q\), un sistema lineare \(\lvert V_{q-1}\rvert\) di varietà (irriducibili), di dimensione \(r>0\), che sia trasformato transitivamente in sè stesso da un sottogruppo \(\infty^r\) (algebrico) del \(\Gamma\). Si può rappresentare il \(\lvert V_{q-1}\rvert\) mercè uno spazio lineare \(S_r\), in cui vien dato un gruppo continuo transitivo \(\infty^r\) di omografie. Ora, entro questo gruppo, si troverà un sottogruppo continuo \(\Gamma'\), di dimensione \(h\geq1\), che lascia invariata una retta \(a\) di \(S_r\), permutandone i punti in \(\infty^1\) modi. Quindi si avrà nel \(\Gamma\) di \(V_q\) un sottogruppo \(\Gamma_h\) di dimensione \(h\), che scambia fra loro in \(\infty^1\) modi le varietà \(\overline{V}_{q-1}\) di un fascio lineare entro \(\lvert V_{q-1}\rvert\) e porta un punto generico di una \(\overline{V}_{q-1}\) nei punti di una \(\overline{V}_h\) formata da \(\infty^1\overline{V}_{h-1}\) appartenenti alle \(\overline{V}_{q-1}\) del detto fascio. Ma in questo fascio c’è (almeno) una \(\overline{V}_{q-1}\) invariante per il \(\Gamma_h\) (che risponde ad un punto unito sulla retta unita \(a\) di \(S_r\)) e quindi la \(\overline{V}_h\) dovrà contenere una \(\overline{V}_{h-1}\) invariante per il \(\Gamma_h\): associando due punti di questa \(\overline{V}_{h-1}\) si avranno non una, ma \(\infty^1\) trasformazioni del \(\Gamma_h\) e così si contraddice al supposto che il \(\Gamma\) di \(V_q\) sia semplicemente transitivo senza eccezioni.

Dunque sulla \(V_q\) non si hanno sistemi lineari \(\lvert V_{q-1}\rvert\) invarianti per un sottogruppo continuo di trasformazioni del \(\Gamma\), ma si possono avere \(\lvert V_{q-1}\rvert\) invarianti per un gruppo finito, sottogruppo del \(\Gamma\) stesso.

Ora, sopra la \(V_q\) si può scegliere una varietà \(V_{q-1}\), a \(q-1\) dimensioni, che non sia trasformata in sè stessa da alcuna trasformazione (non identica) del gruppo \(\Gamma\). Infatti, essendo la \(V_q\) proiettivamente realizzata in un conveniente spazio \(S_r\), una sezione iperpiana generica di essa non verrà certo mutata in sè da alcuna trasformazione del \(\Gamma\).

Ciò posto la serie \(\infty^q\) delle trasformate della detta \(V_{q-1}\), diciamo \(\{V_{q-1}\}\), costituirà un ente algebrico identico al gruppo \(\Gamma\) e perciò anche alla varietà \(V_q\). In generale nella detta serie potrà trovarsi un numero finito \(s\) di \(V_{q-1}\) equivalenti a quella da cui siamo partiti. Quindi segando le \(V_{q-1}\) della detta serie con una curva algebrica generica tracciata su \(V_q\), si otterrà una serie comprendente \(\infty^q\) gruppi di punti disequivalenti, che si lascia rappresentare con una varietà abeliana \(W_q\) (contenuta nella varietà di Picard della data \(V_q\)), la quale, per \(s=1\), riesce identica alla schiera \(\{V_{q-1}\}\) e alla stessa \(V_q\). Invece se \(s>1\) la schiera indicata, e la \(V_q\), vengono rap%
%CONT>>
% scan: s192L
% !PRINTED 360
\unskip\weberpage{360}%
% head: CAPITOLO NONO
%<<CONT
% stamp: Dark right scan border and scattered specks.
presentate sopra la \(W_q\) multipla secondo il numero \(s\), giacchè ad \(sV_{q-1}\) equivalenti risponderà un solo punto di \(W_q\). Riferendoci a questo secondo caso, \(s>1\), fra la \(W_q\) e la \(V_q\) si avrà dunque una corrispondenza \((1,s)\), per effetto della quale viene definita su \(V_q\) una involuzione d’ordine \(s\), \(I'\), invariante per le trasformazioni di \(\Gamma\) e perciò priva di punti doppi, il cui luogo darebbe origine ad eccezione per la semplice transitività di \(\Gamma\).

In conclusione: \textit{una varietà abeliana \(V_q\) si può sempre rappresentare sopra un’altra varietà abeliana, \(W_q\), semplice o multipla senza punti di diramazione, che sia contenuta imprimitivamente nella varietà di Jacobi corrispondente ad una curva di genere \(p(\geq q)\).}

Viceversa: \textit{ogni varietà \(V_q\) irriducibile, che si lasci rappresentare sopra una varietà abeliana \(W_q\), semplice o multipla senza punti di diramazione, contenuta imprimitivamente in una varietà di Jacobi, è essa stessa una varietà abeliana.}

Per dimostrare questo teorema si consideri la corrispondenza \((s,s)\), in cui ad un gruppo \((A,B\,\ldotp\ldotp\ldotp\ldotp)\) dell’involuzione \(I'\) risponde un altro gruppo di punti coniugati \((A',B'\,\ldotp\ldotp\ldotp\ldotp)\); corrispondenza definita dalla trasformazione del gruppo abeliano di \(W_q\) che porta uno nell’altro i punti omologhi di essa: \(\overline{A}\) e \(\overline{A}'\). Si tratta di dimostrare che codesta corrispondenza \((s,s)\) si scinde in \(s\) trasformazioni birazionali, e così, per esempio, che il punto \(A'\) dipende da \(A\), non solo algebricamente, ma anche univocamente, e quindi razionalmente. A tal fine faremo variare per continuità il punto \(A\) su \(V_q\) descrivendo un ciclo chiuso \(\gamma\), e riconosceremo che \(A'\) si muove pure, per continuità, descrivendo del pari un ciclo chiuso e perciò non può scambiarsi p. es. con \(B'\), ma deve tornare immancabilmente in \(A'\). % ?? uncertain: mark after "A" in "A tal fine" -- scan mark or apostrophe.

Se il ciclo descritto da \(A\) è un ciclo \(\gamma\) nullo, la nostra asserzione è ovvia, poichè esso non può contenere entro di sè alcun punto di diramazione della corrispondenza \((s,s)\), essendo l’involuzione \(I_s\) priva di punti doppi. Ma la cosa si può provare anche per un ciclo \(\gamma\) non nullo, sopra lo spazio riemanniano che risponde alla varietà \(V_q\). E la prova risulta da un’esigenza di continuità, ove si tenga conto che la corrispondenza \((A\,A')\) fra il ciclo \(\gamma\) e il cammino \(\gamma'\) descritto senza ambiguità dal punto omologo \(A'\) è suscettibile di variare con continuità fino a ridursi all’identità ove si faccia variare \(A'\) fino a coincidere con \(A\), e al limite il cammino \(\gamma'\) supposto aperto, non può ridursi ad un cammino chiuso, i punti coniugati nell’involuzione \(I_s\) restando sempre distinti.

Risulta così che la varietà \(V_q\) possiede un gruppo semplicemente transitivo di trasformazioni birazionali in sè, essendovi appunto \textit{una} trasformazione di essa che porta un punto \(A\) in un altro punto \(A'\). In particolare vi sarà una trasformazione birazionale che porta un punto \(A\) in un punto \(B\) coniugato ad esso nell’involuzione \(I_s\),
%CONT>>
% scan: s192R
% !PRINTED 361
\weberpage{361}
% head: SUPERFICIE IRREGOLARI E SISTEMI CONTINUI DI CURVE, ECC.
%<<CONT
% stamp: Scattered specks and a dark streak along the right scan edge.
e così la detta involuzione sarà generata da un gruppo finito di \(s\) trasformazioni: \(G_s\). % ?? uncertain: initial word -- e or è; the mark above it may be scan noise.

Ora delle trasformazioni birazionali \((A\,A')\) che abbiamo trovato su \(V_q\) (e che rispondono alle trasformazioni permutabili della varietà \(s\)-pla \(W_q\)) non si può dire, a priori, che sono permutabili. Siccome esse trasformano in sè l’involuzione generata dal gruppo \(G_s\), i cui gruppi di punti rispondono ai punti della \(W_q\), si può dire soltanto che la trasformata di una di esse, \(\Pi\), per mezzo di un’altra \(T\),
\[
T\Pi T^{-1}=\Pi'
\]
differisce da \(\Pi\) per una trasformazione ciclica \(\Omega\) del \(G_s\):
\[
\Pi'=\Pi\Omega.
\]

Però la \(\Omega\) deve restare immutata quando la \(T\) varia nel gruppo continuo \(\Gamma\) che contiene l’identità, e perciò
\[
\Omega=1.
\]

Si deduce dunque che la \(V_q\) possiede un gruppo \(\infty^q\) semplicemente transitivo di trasformazioni permutabili in sè, ossia è una varietà abeliana.

\textit{Nota.} -- La teoria trascendente delle varietà abeliane \(V_q\) risulta dai classici lavori di \textsc{Jacobi}, \textsc{Riemann} e poi di \textsc{Poincarè}, \textsc{Picard}; per le superficie (iperellittiche) --- caso \(q=2\) --- lo studio è stato approfondito da \textsc{G. Humbert} e poi anche dalla scuola italiana. % ?? uncertain: accent on the final e of "Poincarè" -- grave or acute.
% EDITORIAL D0186 RISOLTO: The enlarged accent runs from upper left to lower right: a grave accent. Retain the historical spelling as printed.

\edsection{10. Segue: il genere delle varietà abeliane.}{10. Segue: il genere delle varietà abeliane.}{sec.IX.10}

Vogliamo completare le cose dette intorno alle varietà abeliane, valutandone il genere (geometrico, principale).

Incominceremo dalle varietà di Jacobi, dimostrando che: \textit{La varietà di Jacobi, \(V_p=J_p\), corrispondente a una curva di genere \(p\), ha il genere uno.}

Proveremo che sulla \(J_q\) delle \(q\)-ple di punti di una curva \(C\) di genere \(p\) (\(q\leq p\)) la varietà \(K\) delle \(q\)-ple tolte dai singoli gruppi di una \(g^{q-1}_{2p-2}\) contenuta nella serie canonica di \(C\), è una varietà canonica. Si prenda dapprima \(q=2\) (\(p\geq2\)). Sulla superficie \(J'_2\) delle coppie (non ordinate) dei punti della curva \(C\), si avranno \(\infty^1\) curve \(M\) irriducibili e birazionalmente identiche alla \(C\), ognuna delle quali risponde alla serie delle coppie con un punto fisso; sicchè le \(M\) costituiscono, sulla superficie, un sistema d’indice due. % ?? uncertain: mark beside J_2 in "Sulla superficie" -- J_2 or J'_2.
% EDITORIAL D0187 RISOLTO: The raised mark immediately after J is a distinct dark apostrophe/prime aligned at superscript height. The scan supports a printed prime; its mathematical appropriateness is a separate issue.

Prendiamo una \(g^1_{2p-2}\) contenuta nella serie canonica di \(C\) e sia \(K\) la curva immagine delle coppie tolte da codesta serie \(\infty^1\); deter%
%CONT>>
% scan: s193L
% !PRINTED 362
\unskip\weberpage{362}%
% head: CAPITOLO NONO
%<<CONT
% stamp: Dark right scan border and scattered specks.
miniamo le intersezioni di \(K\) con una curva \(M\). A tal uopo si consideri la \(M\) che risponde alla serie delle coppie con un punto fisso \(A\) (punto generico di \(C\)) e sia
\[
A,\quad A_1,\quad A_2\quad \ldotp\ldotp\ldotp\ldotp\quad A_{2p-3}
\]
il gruppo della \(g^1_{2p-2}\) che contiene \(A\). I punti corrispondenti alle coppie
\[
(A\,A)\,(A\,A_1)\,\ldotp\ldotp\ldotp\ldotp\,(A\,A_{2p-3})
\]
appartengono alla \(M\) e ne costituiscono un gruppo canonico, e d’altra parte tutti questi punti, all’infuori del primo \((A\,A)\), appartengono anche alla \(K\). Ora il gruppo di questi punti, comuni a \(K\) ed \(M\), è residuo della serie caratteristica di \(M\) rispetto alla sua serie canonica. Infatti si consideri, accanto alla \(M\), la curva \(M'\) rappresentativa delle coppie di \(C\) che hanno comune un altro punto \(B\): le \(M\) e \(M'\) s’intersecano nel punto omologo alla coppia \((A\,B)\), e questo tende al punto \((A\,A)\) quando \(B\) si avvicini infinitamente ad \(A\), sicchè \((A\,A)\) costituirà un gruppo della serie caratteristica sulla curva \(M\). Ora, poichè la \(K\) sega sopra ogni curva del sistema \(M\) un gruppo residuo della serie caratteristica, per il \textit{terzo criterio d’equivalenza} (cfr. Cap. III, § 13) la \(K\) stessa sarà una curva canonica della \(J_q\), \hfill c. d. d.

Dopo aver così dimostrato il nostro teorema per \(q=2\), lo dimostreremo per \(q=3\) e così successivamente per induzione completa. Infatti, supposto vero il teorema per la varietà dei gruppi di \(q-1\) punti di \(C\) (\(q\leq p\)) si riconosce vero anche per la \(J_q\) delle \(q\)-ple.

A tal uopo si consideri sulla \(J_q\) la serie delle \(\infty^1\) varietà \(M_{q-1}\), irriducibili e birazionalmente identiche, ognuna delle quali rappresenta le \(q\)-ple con un punto fisso: questa serie è d’indice \(q\). Sia \(K_{q-1}\) la varietà delle \(q\)-ple tolte dai singoli gruppi di una \(g^{q-1}_{2p-2}\) contenuta nella serie canonica di \(C\), e determiniamo l’intersezione della \(K_{q-1}\) colla \(M_{q-1}\). Sia \(g^{q-2}_{2p-3}\) la serie residua del punto \(A\) rispetto alla serie considerata. Aggiungendo ai gruppi di \(q-1\) punti estratti dai gruppi della \(g^{q-2}_{2p-3}\) il punto fisso \(A\), si ottiene una varietà \(K_{q-2}\) che è comune alla \(M_{q-1}\) rispondente al punto \(A\) e alla \(M_{q-1}\) immagine della \(g^{q-1}_{2p-2}\). % ?? uncertain: prime on M in "e alla M'_{q-1} immagine" -- M'_{q-1} or M_{q-1}.
% EDITORIAL D0188 RISOLTO: The enlarged M has no prime; the area between its upper right serif and immagine is blank. Remove the transcribed prime only; do not substitute the mathematically expected K. % [sic] M'_{q-1} -> K_{q-1}, the image of the series was denoted K_{q-1} earlier in this paragraph.
Ora se, sulla \(M_{q-1}\), alla \(K_{q-2}\) sommiamo la varietà \(M_{q-2}\) delle \((q-1)\)-ple aventi il punto fisso \(A\) (cioè la varietà dei gruppi di \(q\) punti, dei quali fa parte due volte il punto \(A\)) si otterrà una varietà \textit{canonica} (impura) della \(M_{q-1}\); e ciò per il teorema ipoteticamente assunto come base della nostra induzione, perchè la detta varietà risponde alla serie dei gruppi di \(q-1\) punti estratti dai singoli gruppi della serie
\[
g^{q-2}_{2p-3}+A,
\]
serie col punto fisso \(A\), contenuta nella serie canonica di \(C\).

% scan: s193R
% !PRINTED 363
\weberpage{363}
% head: SUPERFICIE IRREGOLARI E SISTEMI CONTINUI DI CURVE, ECC.
% stamp: Thin dark left scan border and scattered specks.
Dunque la \(K_{q-1}\) sega sulla \(M_{q-2}\) una varietà residua di una varietà canonica rispetto alla \(M_{q-2}\); quindi per il terzo criterio di equivalenza (che facilmente si estende dal caso delle curve a quello delle varietà a qualsiasi dimensione) basterà provare che la \(M_{q-2}\) costituisce una varietà del sistema caratteristico della \(M_{q-1}\). Ed invero la \(M_{q-1}\) relativa al punto fisso \(A\) e la \(M'_{q-1}\) relativa al punto fisso \(B\), hanno in comune la varietà \(\overline{M}_{q-2}\) delle \(q\)-ple con i due punti fissi \(A\) e \(B\), e quando \(B\) tende ad \(A\), la \(M'_{q-1}\) diviene infinitamente vicina alla \(M_{q-1}\) e la \(\overline{M}_{q-2}\) tende appunto alla \(M_{q-2}\). % [sic] first M_{q-2} -> M_{q-1}, the preceding paragraph places this intersection on M_{q-1}.

Così essendo stabilito il teorema enunciato, si faccia \(p=q\); avremo: \textit{sulla varietà di Jacobi \(J_p\), rappresentativa dei gruppi di \(p\) punti di una curva di genere \(p\), la varietà \(K_{p-1}\) delle \(p\)-ple speciali è una varietà canonica} (a priori impura). Ma essa è precisamente una \textit{varietà eccezionale}, che si lascia trasformare in una varietà di dimensione \(p-2\) sulla varietà, birazionalmente identica alla \(J_p\), rappresentativa dei gruppi di \(p\) punti disequivalenti della curva \(C\) di genere \(p\). Si ottiene d’altronde questa trasformazione usufruendo della corrispondenza involutoria che intercede, su \(C\), fra i gruppi di \(p\) punti residui di una \(g^p_{2p}\).

Pertanto ne segue che la \textit{varietà di Jacobi \(J_p\) ha il genere uno.}

Il teorema stabilito per le varietà di Jacobi \(J_p=V_p\) si estende ora, facilmente, alle varietà abeliane più generali, e basterà indicarne in breve la dimostrazione. Anzitutto, riferendoci ad una \(V_q\) semplice contenuta imprimitivamente in una \(V_p\) (\(q<p\)), questa \(V_q\) farà parte di una schiera di varietà della stessa dimensione, schiera d’indice uno; e quindi la varietà canonica (impura) della \(V_p\) segherà su ciascuna di tali \(V_q\) \textit{una} varietà canonica. In secondo luogo il teorema si estende ancora al caso di una \(V_q\) che venga rappresentata sopra una varietà \(W_q\) multipla, contenuta in \(V_p\), senza varietà di diramazione (cfr. § 9). Infatti la trasformazione \((1,s)\) che fa passare da \(W_q\) a \(V_q\) muta la varietà canonica di \(W_q\) in una varietà canonica di \(V_q\)\origfootnote{1}{s193R-1}{Teorema di Painlevé-Castelnuovo stabilito per \(q=1\) in \textsc{Enriques-Chisini}, \textit{Lezioni}. Libro V, cap. I, § 9 (vol. III, pag. 72). Per \(q=2\) cfr. il Cap. V, § 7 di questo trattato. Il teorema si estende senza difficoltà per tutti i valori di \(q\).}.

In conclusione possiamo enunciare che: \textit{Le varietà abeliane \(V_q\) (\(q\geq1\)) sono sempre di genere uno, e con varietà canonica pura d’ordine zero.}

\textit{Notizia.} -- La valutazione del genere delle varietà di Jacobi \(V_p\), e in generale delle varietà abeliane, riesce immediata sotto l’aspetto trascendente, poichè si costruisce subito un integrale \(p\)-plo di prima specie annesso a \(V_p\). La dimostrazione geometrica del
%CONT>>
% scan: s194L
% !PRINTED 364
\weberpage{364}
% head: CAPITOLO NONO
%<<CONT
% stamp: Dark scan borders at the upper left and right, with scattered specks.
teorema che abbiamo sopra esposta, è stata data dal \textsc{Severi} in una Nota del 1911\origfootnote{1}{s194L-1}{\textit{Sulle superficie e varietà algebriche irregolari di genere geometrico nullo.} R. Acc. Lincei, 1911.}.

Un’altra dimostrazione geometrica dello stesso teorema è stata indicata da \textsc{B. Segre}\origfootnote{2}{s194L-2}{\textit{Determinazione di certi gruppi covarianti di due o più serie lineari.} Rendic. Circolo Mat. Palermo, 1932.} e consiste nel mostrare che le serie lineari \(g^{p-1}_{2p}\) contenute in una serie \(g^p_{2p}\) sopra una curva \(C\) di genere \(p\), hanno per immagini sulla corrispondente varietà di Jacobi, le varietà \(V_{p-1}\) di un sistema lineare il cui sistema caratteristico, sommato all’intersezione coll’immagine della varietà dei gruppi speciali, costituisce il sistema canonico su una \(V_{p-1}\). Ci limiteremo a spiegare la cosa nel caso più semplice \(p=2\), avvertendo che l’estensione del ragionamento per \(p>2\) riesce invero meno semplice di quanto si potrebbe desiderare. % ?? uncertain: punctuation after "Rendic" in the footnote -- period or hyphen.
% EDITORIAL D0189 RISOLTO: The short flattened mark is at full-stop baseline height, substantially below normal hyphen height. Retain Rendic. as a damaged full stop.

Consideriamo dunque una curva \(C\) di genere \(p=2\), e sopra di essa una serie lineare generica \(g^2_4\): le \(g^1_4\) contenute in questa serie danno luogo, sopra la superficie di Jacobi \(V_2\), alle curve \(L\) di una rete, priva di curve fondamentali, che avranno un’intersezione variabile \(A\) colla curva eccezionale \(V_1\) immagine della \(g^1_2\) canonica di \(C\), o un punto base \(B=V_0\) nel punto trasformato di essa mercè l’involuzione definita dalla detta \(g^2_4\). Le curve canoniche \(K\) di \(V_2\) sono definite, a priori, sulla superficie, dalla proprietà di segare su una generica \(L\) gruppi residui della serie caratteristica di \(\lvert L\rvert\), aumentata dalla intersezione \(A\) della \(L\) colla curva eccezionale \(V_1\) o del punto base \(B\) (trasformato della \(V_1\)).

Si tratta dunque di dimostrare che la serie \(A+B+g^2_6\) su una \(L\) è contenuta (totalmente) nella serie canonica \(g^4_8\). E la tesi deriva in sostanza dalla interpretazione geometrica della formula di Zeuthen che lega le due serie canoniche delle curve \(C\) e \(L\) in corrispondenza \((2,3)\)\origfootnote{3}{s194L-3}{Questo passaggio riesce assai semplice (per \(p=2\)) ove s’introduca la curva ausiliaria \(M\) che rappresenta le coppie di punti di \(L\) complementari rispetto ad una \(g^1_4\), contenuta nella \(g^2_4\).}.

Così dunque la rete delle curve \(L\) su \(V_2\), o meglio su una trasformata di questa senza curve eccezionali, risulta essere di genere \(\pi=5\), di grado effettivo \(6\) e di grado virtuale \(n=8\), (\(=2\pi-2\)), avendo due punti base semplici, con serie caratteristica canonica. Il possesso di una tal rete \(\lvert L\rvert\), senza curve fondamentali proprie, porta appunto che la \(V_2\) sia una superficie di genere \(p_g=1\) con curva canonica d’ordine zero, \hfill c. d. d. % ?? uncertain: first letter of the concluding abbreviation -- c or o.
% EDITORIAL D0190 RISOLTO: The initial letter has a clearly open right side: c.

% scan: s194R
% !PRINTED 365
\weberpage{365}
% head: SUPERFICIE IRREGOLARI E SISTEMI CONTINUI DI CURVE, ECC.
% stamp: Scattered specks and a dark right scan edge.
\edsection{11. Fascio irrazionale sopra le superficie di genere geometrico nullo.}{11. Fascio irrazionale sopra le superficie di genere geometrico nullo.}{sec.IX.11}

Dalla circostanza che una varietà abeliana \(V_q\) di dimensione \(q\geq2\) ha il genere geometrico eguale ad uno, si deduce che «\textit{ogni superficie \(F\) non eccezionale appartenente ad una varietà abeliana \(V_q\), ha il genere geometrico \(p_g>0\)}».

Per \(q=2\) la cosa è evidente, la \(F\) coincidendo con \(V_2\). Per provare in generale il teorema occorre osservare che per \(q>2\) «la \(F\) su \(V_q\) fa parte di un sistema continuo di superficie analoghe, di dimensione \(q-2\) almeno».

Infatti, applicando alla \(F\) le \(\infty^q\) trasformazioni del gruppo \(\Gamma\) della \(V_q\) si avranno, in generale, \(\infty^q\) superficie analoghe ad \(F\): fa eccezione il caso in cui la \(F\) ammetta una serie continua di trasformazioni in sè; ma, in ogni modo (per la semplice transitività del gruppo \(\Gamma\)) le trasformazioni della \(F\) saranno al più \(\infty^2\), e quindi si avranno almeno \(\infty^{q-2}\) superficie trasformate di \(F\).

Ciò posto la proprietà di \(F\) di essere di genere geometrico non nullo (come abbiamo enunciato) è una conseguenza della proprietà generale che: Sopra una varietà \(V_q\) di dimensione \(q\), la quale abbia il genere geometrico \(P>0\), ogni superficie \(F\) che appartenga ad una schiera continua di dimensione \(q-2\) almeno, ha il genere \(p_g>0\).

Dimostriamo questa proposizione riferendoci dapprima al caso \(q=3\), e poi al caso \(q=4\), ecc.

Supposto \(q=3\), la superficie canonica che esiste per ipotesi sulla \(V_q\) (anche se costituita da una sola superficie, d’ordine zero) interseca la \(F\) lungo una curva che appartiene al sistema residuo della serie del sistema caratteristico di \(F\) rispetto al suo sistema canonico\origfootnote{1}{s194R-1}{Cfr. VIII, § 9.}; ciò basta per provare l’esistenza di quest’ultimo, cioè che il genere di \(F\) vale \(p_g>0\).

Se \(q_g=4\), le superficie \(F\), che sono almeno \(\infty^2\), si possono distribuire in almeno \(\infty^1\) sistemi continui semplicemente infiniti, che dan luogo dunque ad una serie \(\infty^1\) almeno di varietà \(V_3\): ora, per essere il genere di \(V_4\) maggior di zero, una varietà canonica di \(V_4\) segherà su ciascuna delle suddette \(V_3\) superficie speciali, residue del sistema caratteristico di \(V_3\), rispetto al sistema canonico, e quindi codeste \(V_3\) saranno di genere \(P>0\). Da ciò segue che le \(\infty^1\) superficie che le costituiscono sono pure di genere \(p_g>0\). % ?? uncertain: small mark after q in "Se q_g=4" -- subscript g, comma, or scan mark.
% EDITORIAL D0191 RISOLTO: A separately set subscript g is visible beside q; it is neither a comma nor an isolated speck. Retain the printed q_g and the proposed mathematical correction only in the apparatus. % [sic] q_g = 4 -> q = 4, this case is introduced as q = 4 in the preceding paragraph.

Con procedimento analogo, passando da \(q\) a \(q+1\), si prova la cosa per una varietà \(V_q\) di dimensione qualunque.

Il teorema stabilito che le superficie appartenenti ad una varietà abeliana sono di genere \(p_g>0\), ci permette ora di precisare, per le
%CONT>>
% scan: s195L
% !PRINTED 366
\weberpage{366}
% head: CAPITOLO NONO
%<<CONT
% stamp: Heavy scan speckling and a dark right border; an ink-like diagonal stroke crosses the i in "delle i curve", and a vertical mark follows "-p".
superficie di genere \(p_g=0\), la proprietà caratteristica delle superficie-irregolari, dimostrando che: \textit{Le superficie di genere \(p_g=0\), irregolari, posseggono un fascio irrazionale di genere \(-p_a=p\) di curve. Viceversa: ogni superficie irregolare di genere \(p_a=0\), che possegga un fascio irrazionale di genere \(p\) (\(>0\)), è una superficie irregolare di genere numerico \(-p\).} % ?? uncertain: hyphen in "superficie-irregolari" -- printed hyphen or scan streak.
% EDITORIAL D0192 RISOLTO: A regular horizontal hyphen is visible after the continued word superficie; retain the printed hyphen. % [sic] p_a = 0 -> p_g = 0, the converse concerns the geometric-genus condition stated immediately before.

Sia \(F\) una superficie irregolare e --- senza fare per ora alcuna ipotesi circa i valori dei suoi caratteri --- si consideri sopra di essa un sistema continuo semplicemente infinito di curve \(\{C\}\), \textit{d’indice} \(i>2\), che possiamo supporre \textit{semplice}, cioè tale che le \(i\) curve \(C\) passanti per un punto generico della superficie non abbiano di conseguenza a comune altri punti.

Si indichi con
\[
C_i=\Sigma C
\]
la curva somma delle \(i\) curve di \(\{C\}\) che escono da un punto della superficie \(F\).

Si possono presentare tre casi:

\textit{a}) le \(\infty^2\) curve \(C_i\) sono tutte equivalenti sopra la superficie \(F\);

\textit{b}) le dette \(C_i\) si distribuiscono in \(\infty^1\) sistemi semplicemente infiniti di curve equivalenti;

\textit{c}) le \(C_i\) sono tutte disequivalenti fra loro, ovvero ciascuna di esse riesce equivalente ad un numero finito, \(s\geq1\), di curve.

Ora proveremo che \textit{nel caso a) tutte le curve \(C\) sono equivalenti fra loro, cioè la serie \(\{C\}\) è contenuta in un sistema lineare.}

A tal uopo basta provare che le \(\infty^1\) curve \(C\) segano una serie di gruppi equivalenti sopra una qualsiasi curva \(M\) tracciata sopra \(F\), e invocare, per esempio, il primo criterio d’equivalenza. Invero la proprietà affermata per la serie segata sulla curva \(M\) deriva da un noto criterio d’equivalenza di Severi-Castelnuovo relativo alla serie \(\infty^1\) di gruppi di punti appartenenti ad una curva\origfootnote{1}{s195L-1}{Cfr. III, 13. Vedi pure \textsc{Enriques-Chisini}, \textit{Lezioni}. Libro V, cap. IV, § 42 (vol. III, pag. 483).}. % ?? uncertain: mark after "affermata" -- scan noise or comma.

In secondo luogo proveremo che \textit{nel caso b) si può costruire sopra la superficie \(F\) un fascio irrazionale di curve.}

Infatti si faccia muovere su \(F\) un punto \(A\) per modo che le curve \(C_i\) uscenti da esso si mantengano equivalenti; siccome la condizione imposta è algebrica, il punto \(A\) descriverà una curva algebrica \(K\), ed è chiaro che le curve \(K\) così definite su \(F\), non potendo segarsi a due a due, dovranno formare un fascio (privo di punti base semplici per la superficie). Supponendo, per semplicità di discorso, che le \(K\) siano irriducibili (chè altrimenti basterebbe considerare al
%CONT>>
% scan: s195R
% !PRINTED 367
\weberpage{367}
% head: SUPERFICIE IRREGOLARI E SISTEMI CONTINUI DI CURVE, ECC.
%<<CONT
% stamp: Heavy scan speckling, dark edge streaks, and a large speckled patch in the lower right margin.
posto delle \(K\) le loro componenti irriducibili), dimostreremo ora che il fascio \((K)\) è irrazionale.

Invero, la stessa considerazione fatta innanzi mostra che le curve \(C\) di \(\{C\}\) debbono segare, sopra una qualunque \(K\), \(\infty^1\) gruppi di punti equivalenti. Quindi, per il \textit{secondo criterio d’equivalenza}\origfootnote{1}{s195R-1}{Cfr. III, 13.}, le \(C\) saranno equivalenti a meno di curve del fascio \((K)\), cioè --- designando con \(C\) e \(C_0\) due curve qualsiasi di \(\{C\}\) --- si avrà
\[
C+\Sigma K=C_0+\Sigma'K,
\] % ?? uncertain: relation sign between the two sums -- = or \equiv.
% EDITORIAL D0194 RISOLTO: Enlargement shows two printed horizontal strokes; nearby dots do not form a third bar. Equality as transcribed.
dove le due sommatorie \(\Sigma\) e \(\Sigma'\) contengono uno stesso numero di addendi. Di conseguenza, se le curve \(C\) non sono tutte equivalenti, il fascio \((K)\) non può essere lineare, chè altrimenti sommando e sottraendo uno stesso numero di curve \(K\) si passerebbe da una \(C\) ad una \(C_0\) equivalente.

A questo ragionamento, fondato sul secondo criterio d’equivalenza, si può muovere un’obiezione, nell’ipotesi che il fascio \((K)\) contenga delle curve spezzate, come
\[
K=K_1+K_2\qquad \text{e}\qquad \overline{K}=\overline{K}_1+\overline{K}_2;
\]
giacchè allora potrebbe aversi, per esempio,
\[
C_0=C+K_1-\overline{K}_1+\ldotp\ldotp\ldotp\ldotp
\] % ?? uncertain: relation sign after C_0 -- = or \equiv.
% EDITORIAL D0195 RISOLTO: The relation has two clear parallel bars: equality.

Ma il dubbio che qui si solleva viene eliminato dalla considerazione che le curve \(C_0\) e \(C\) appartengono ad un medesimo sistema continuo \(\{C\}\) e al variare di \(C_0\) in \(\{C\}\) le componenti delle \(K\) spezzate, essendo in numero finito, rimangono fisse; di qui, riducendo per continuità la \(C_0\) alla \(C\), risulta necessariamente
\[
C_0\to C,\qquad K_1=\overline{K}_1.
\] % ?? uncertain: relation sign between K_1 and the barred K_1 -- = or \equiv.
% EDITORIAL D0196 RISOLTO: The relation has two clear parallel bars; the overbar belongs to the K on the right.

Infine, ponendosi nel \textit{caso c)}, si prova subito che dev’essere il \textit{genere della superficie \(p_g\geq1\).} Rappresentando infatti le curve di \(\{C\}\) con i punti di una curva \(L\), le \(\infty^2\) curve \(C_i\) definiscono sulla \(L\) stessa una serie \(\gamma'_i\) che si può considerare come una varietà doppiamente infinita, semplice o multipla, contenuta nella varietà di Jacobi che corrisponde alla curva \(L\), rappresentando le serie lineari disequivalenti della \(L\) stessa. % ?? uncertain: prime on gamma -- \gamma'_i or \gamma_i.

La discussione precedente conduce a riconoscere che, per una superficie \(F\) di genere \(p_g=0\), irregolare, ogni serie \(\{C\}\) d’indice \(i\), costituita di curve non equivalenti, che si scelga sopra di essa, dà luogo al caso \textit{b)}, sicchè si trova sulla superficie un fascio irrazionale di curve \(K\).

% scan: s196L
% !PRINTED 368
% [printed number read as 352 in the first pass; the spread sequence gives 368]
\weberpage{368}
% head: CAPITOLO NONO % ?? uncertain: Only the ending "NO" is exposed; the full head is read from the matching page in scan s197L -- CAPITOLO NONO / another obscured head ending in NO.
% stamp: Folded or overlapping printed leaves obscure the upper-left portion. The visible folios are 352 and, on the oblique fold, 353. Overlaid fragments include "Nel caso 1) (g" (with "caso" italic), "e qu", "TI DI CURVE, ECC.", "continuo {L}", "dimensione", "ni pas-", and "e F". % ?? uncertain: Several overlaid fragments are cropped -- g / p; qu / q; TI / UI; L / C; e F / F.
\sourcewarning{La scansione presenta fogli sovrapposti. Per i due capoversi seguenti la trascrizione dichiara di avere recuperato le parti coperte da una scansione corrispondente, s197L, non inclusa nel materiale fornito. Si conserva il testo ricevuto, ma queste letture non sono verificabili qui. La limitazione termina prima di «Notizia».}
Si riconosce quindi che il genere di \((K)\) eguaglia il genere numerico mutato di segno: \(p=-p_a\). Perchè, se fosse \(p<-p_a\), in forza del secondo criterio d’equivalenza, si riuscirebbe a costruire su \(F\) una serie \(\{C\}\) di curve non equivalenti, secanti su ciascuna \(K\) gruppi di punti non equivalenti, e quindi (avverandosi il caso \(b\)), un secondo fascio irrazionale di curve; ma il possesso di due fasci irrazionali su \(F\), \((K)\) e \((M)\), porterebbe che questa superficie debba avere il genere geometrico \(p_g>0\): la cosa appare evidente se i due fasci sono formati di curve unisecantisi\origfootnote{1}{s196L-1}{Infatti si ottiene una curva canonica di \(F\) sommando alle curve di \((M)\) che passano per i punti di un gruppo canonico di una \(K\) le curve di \((K)\) che passano per i punti di un gruppo canonico di una \(M\).}; mentre il caso che le curve dei due fasci si seghino in un certo numero \(s>1\) di punti, conduce a rappresentare la \(F\) sopra una superficie multipla (\(s\)-pla) con due fasci di curve unisecantisi, donde risulta, a fortiori\origfootnote{2}{s196L-2}{Cfr. V, 7.}, che il suo genere \(p_g>0\). % ?? uncertain: The covered portions of this paragraph, including its opening, are restored solely from the matching exposed text in scan s197L -- identical wording / wording differing beneath the fold.

In conclusione la proprietà caratteristica delle superficie irregolari di genere \(p_g=0\), che importa il possesso di un fascio irrazionale di genere \(p=-p_a\), viene dimostrata, secondo l’enunciato precedente. % ?? uncertain: The covered beginning and portions of the next line are read from the matching paragraph in scan s197L -- identical wording / wording differing beneath the fold.

\textit{Notizia.} – Il teorema che concerne l’esistenza di un fascio irrazionale sopra le superficie irregolari di genere \(p_g=0\), è stato enunciato da \textsc{Enriques}, come conseguenza della proprietà caratteristica delle superficie irregolari, e ciò nella sua citata nota dei Rendic. dell’Accademia di Bologna nel 1904. Invero, fino dal 1900 il \textsc{Castelnuovo} aveva avuto luogo di osservare e comunicare al collega che « dall’esistenza di un sistema continuo di curve disequivalenti sopra una superficie di genere \(p_g=0\), si deduce l’esistenza di un fascio irrazionale ». Questa proprietà, che l’\textsc{Enriques} richiama nella sua Nota ricordandone l’autore, era giustificata per via trascendente, come segue.

Sia \(\{C\}\) un sistema continuo \(\infty^1\) di curve non equivalenti, sopra la superficie \(F\) di genere \(p_g=0\). Allora, come fu rilevato per la prima volta da G. \textsc{Humbert}\origfootnote{1}{s196L-3}{Comptes rendus, 1893; Journal del Math., s. 4, t. X, 1894.}% [sic] In-text (1), footnote (3) -> in-text (3), the printed call and the number printed beside its footnote disagree.
, si costruiscono su \(F\) degli integrali semplici (di differenziali totali) di prima specie, sommando i valori che gli integrali abeliani di prima specie pertinenti alla serie \(\{C\}\) assumono in corrispondenza alle \(i\ (>1)\) curve \(C\) uscenti da uno stesso punto \(A\) di \(F\). E questi integrali semplici \(I_1\,I_2\,\ldots.\) saranno necessariamente funzioni l’uno dell’altro, altrimenti — come è stato
%CONT>>
% scan: s196R
% !PRINTED 369
\weberpage{369}
% head: SUPERFICIE IRREGOLARI E SISTEMI CONTINUI DI CURVE, ECC.
%<<CONT
% stamp: Dark scan borders and scattered specks.
indicato da \textsc{Noether}\origfootnote{1}{s196R-1}{\textit{Ueber die totalen algebraischen differentialausdrücke.} Math. Annalen, Bd. 29, 1887. Confr. \textsc{Picard} et \textsc{Simart}, op. c., vol. I, Cap. V, \S\ 15.} — si dedurrebbe l’esistenza di un integrale doppio di prima specie, e quindi \(p_g>0\).

Ciò posto il \textsc{Castelnuovo} osserva che le linee \(I=\mathrm{cost.}\) saranno composte colle curve algebriche d’un fascio irrazionale che vengano definite facendo muovere \(A\) sopra la superficie in guisa che le \(i\) curve \(C\) uscenti da un punto descrivano gruppi di curve equivalenti. Questo ragionamento, specie nella forma che assume traverso le considerazioni della Nota di \textsc{Castelnuovo} « Sugli integrali semplici appartenenti ad una superficie irregolare »\origfootnote{2}{s196R-2}{R. Acc. Lincei, 21 Maggio, 4 e 18 Giugno, 1905. Cfr. Memorie scelte, pag. 473 e segg. (v. in ispecie il n. 10, pag. 490).}, è stato tradotto in quello geometrico che sopra si è esposto, da F. \textsc{Severi}\origfootnote{3}{s196R-3}{\textit{Sulle superficie e varietà algebriche irregolari di genere geometrico nullo.} Rendic. Lincei, s. V, vol. XX, 1911.}.

\edsection{12. Nota sulle superficie di irregolarità uno.}{12. Nota sulle superficie di irregolarità uno.}{sec.IX.12}

In vista di possibili applicazioni conviene rilevare che le considerazioni svolte nel precedente paragrafo conducono anche a caratterizzare le superficie di irregolarità \(q=p_g-p_a=1\), per le quali sussiste il teorema: \textit{Le superficie \(F\) di irregolarità \(q=1\) posseggono un fascio ellittico di curve} (di genere \(\pi\geq0\)). \textit{Sopra di esse ogni sistema continuo di curve si ottiene da un sistema lineare sommandovi la differenza di due curve del detto fascio.}

Osserviamo che la varietà di Picard di \(F\) è una curva ellittica. Si assuma ora sopra \(F\) un qualsiasi sistema \(\infty^1\) di curve disequivalenti, \(\{C\}\), d’indice \(\nu>1\): per ogni punto \(P\) di \(F\) passeranno \(\nu\) curve \(C\), e, al variare di \(P\), si otterranno \(\infty^2\) gruppi di \(\nu\) curve; quindi, a partire da uno particolare di questi punti, \(P=\overline{P}\), vi saranno \(\infty^1\) punti \(P\) per cui passa un gruppo di \(\nu\) curve \(C\) equivalente al gruppo delle curve passanti per \(\overline{P}\). Il luogo di questi punti \(P\) sarà una curva \(K\) per \(\overline{P}\), e tutte le \(K\) analoghe formeranno su \(F\) un fascio ellittico di curve \((K)\). Stabilita in tal guisa la prima parte del teorema, per dimostrare la seconda, basterà avvertire che le curve \(C\) del sistema continuo \(\{C\}\) segheranno su ogni \(K\) gruppi equivalenti, e di conseguenza (per il secondo criterio d’equivalenza) due \(C\), non equivalenti, differiranno fra loro soltanto per la differenza di due \(K\).

% EDITORIAL: segnatura originale: \textsc{Enriques F.} - \textit{Superficie algebriche.}\hfill 24
\printersignature{\textsc{Enriques F.} - \textit{Superficie algebriche.}\quad 24}
% scan: s198L
% !PRINTED none
\weberblank{[s198L]}
% head: none
% blank page
% stamp: Handwritten "370" in the upper-right margin; scattered specks and a dark vertical scan border at the right.


\clearpage\thispagestyle{plain}
% scan: s198R
% !PRINTED none
\weberpage{[s198R]}
% head: none
% stamp: Handwritten "371" in the upper-right margin; scattered specks and vertical scan lines along the right margin.
\edchapter{X}{LE SUPERFICIE DI GENERE GEOMETRICO NULLO}{Capitolo X. LE SUPERFICIE DI GENERE GEOMETRICO NULLO}{LE SUPERFICIE DI GENERE GEOMETRICO NULLO}

\edsection{1. Introduzione.}{1. Introduzione.}{sec.X.1}

La proprietà caratteristica delle superficie irregolari di genere \(p_g=0\), che abbiamo stabilito nei \S\S\ 4 e 11 del precedente capitolo, permette, come vedremo, di caratterizzare — in modo completo — le superficie di genere \(p_g=0\), e — in particolare — di assegnare, mediante l’annullamento dei plurigeneri, le condizioni necessarie e sufficienti che definiscono la famiglia delle superficie rigate.

Lo sviluppo delle nostre deduzioni procederà nell’ordine seguente.

Anzitutto tratteremo delle superficie che contengono sistemi di curve di genere \(\pi\) e di grado \(n>2\pi-2\), dimostrando che esse sono trasformabili in rigate, razionali o irrazionali; il caso razionale risponde (com’è noto) alla regolarità della superficie.

In secondo luogo esporremo un lemma concernente le curve spezzate di un fascio irrazionale, che ci permetterà di precisare in una forma più semplice il secondo criterio d’equivalenza quale occorre adoperare negli sviluppi che seguono.

Di qui trarremo intanto che le superficie di genere \(p_g=0\) e \(p_a<-1\) sono senz’altro riferibili a rigate, di genere \(p_a=-p\). Restano quindi a studiare le superficie coi generi \(p_g=0\) e \(p_a=-1\), le quali danno luogo a:

1) superficie trasformabili in rigate ellittiche, e

2) superficie contenenti un fascio ellittico di curve di genere \(\pi\geq1\).

Tratteremo separatamente i due casi che si presentano, per \(\pi>1\) e \(\pi=1\), dimostrando che, nell’uno e nell’altro, la superficie contiene un secondo fascio lineare di curve ellittiche, birazionalmente identiche.

Infine approfondiremo lo studio della famiglia delle superficie (\textit{ellittiche}) definite dalle anzidette proprietà, ne assegneremo la co%
%CONT>>
% scan: s199L
% !PRINTED 372
\unskip\weberpage{372}%
% head: CAPITOLO DECIMO
%<<CONT
% stamp: Scattered scan specks, a dark right-hand edge, and a diagonal blemish near "formeranno" and "dell’aggiunzione".
struzione, ne valuteremo i plurigeneri e riconosceremo che esse ammettono un gruppo ellittico, semplicemente infinito, di trasformazioni birazionali in se stesse, ecc. Termineremo il capitolo caratterizzando le superficie coi plurigeneri nulli e quelle coi plurigeneri eguali a zero o ad uno.

\edsection{2. Le superficie contenenti un sistema di curve di genere \(\pi\) e di grado \(n>2\pi-2\).}{2. Le superficie contenenti un sistema di curve di genere π e di grado n>2π-2.}{sec.X.2}

Sia \(F\) una superficie la quale contenga un sistema lineare di curve \(|C|\), \(\infty^2\) almeno, irriducibile per cui \(n>2\pi-2\). Le curve \(C'\), \(C''\) ecc., successive aggiunte a \(|C|\), segheranno le \(C\) in un numero decrescente di punti:
\[
2\pi-2,\qquad 4\pi-4-n,\qquad 6\pi-6-2n\,\ldots.
\]
e perciò formeranno una serie limitata; in altri termini, l’operazione dell’aggiunzione, a partire dal sistema lineare \(|C|\), dovrà estinguersi dopo un numero finito di passaggi, dando luogo ad un ultimo sistema aggiunto, \(|C^t|\), di dimensione \(\geq0\), che non possiede più alcuna curva aggiunta, d’ordine \(>0\). % ?? uncertain: A diagonal blemish immediately follows "formeranno" -- "formeranno una" / "formeranno, una".

Conosciamo già un caso in cui l’aggiunzione si estingue in siffatta guisa: è il caso delle superficie regolari di genere \(p_g=p_a=0\), con bigenere \(P=0\), in cui è noto che la superficie \(F\) risulta razionale\origfootnote{1}{s199L-1}{Teorema di Castelnuovo cfr. Cap. V, \S\ 4.}. Se si prescinde da questo caso, conviene esaminare l’ipotesi che la superficie su cui si trova il nostro sistema \(|C|\) sia irregolare, di genere \(p_g=0\) e \(p_a=-p\leq-1\), e perciò contenga un fascio irrazionale, di genere \(p\), di curve \(K\). Se le \(K\) sono di genere \(\varrho=0\), la \(F\) risulta senz’altro trasformabile in una rigata di genere \(p\)\origfootnote{2}{s199L-2}{Teorema di Enriques cfr. p. es. \textsc{Enriques-Conforto}, \textit{Le superficie razionali}, Libro II, cap. I, \S\ 4-6 (pag. 253).}. Occorre esaminare e ridurre all’assurdo l’ipotesi che il genere delle \(K\) sia invece \(\varrho\geq1\). E gioverà distinguere i due casi:
\weq{1)}{p>1}
\weq{2)}{p=1.}

\textit{Primo caso:} \(p>1\).

Si designi con \(s\) il numero delle intersezioni delle \(C\) colle \(K\), e si osservi che le \(C+pK-p\overline{K}\), dan luogo ad una schiera \(\infty^p\) di sistemi lineari distinti secanti sulle \(K\) serie equivalenti; da ciò segue
\[
s>1.
\]

Ora, siccome le curve canoniche (virtuali) di \(F\) segano le \(K\), di genere \(\varrho\), in \(2\varrho-2\) punti, le \(C'\) segheranno le stesse \(K\) in \(s+2\varrho
%CONT>>
% scan: s199R
% !PRINTED 373
\text{\weberpage{373}}
% head: LA SUPERFICIE DI GENERE GEOMETRICO NULLO
%<<CONT
% stamp: Scattered scan specks and marginal edge traces.
-2\)\edFlush{} punti, le \(C''\) le segheranno in \(s+2\,(2\varrho-2)\) punti, e così via. Infine le curve dell’ultimo sistema aggiunto \(|C^t|\) segheranno le \(K\) in un numero di punti che sarà, a fortiori, \(m\geq2\). Ma questa conclusione apparisce subito assurda, se le \(C^t\) sono irriducibili, perchè, possedendo esse una involuzione irrazionale \(\gamma_m^1\) di genere \(p\), segata dalle \(K\), dovranno avere il genere (effettivo e virtuale) maggiore o uguale a \(m(p-1)+1>p\), e quindi possedere un sistema aggiunto di dimensione \(>0\).

La conclusione (che è una riduzione all’assurdo della ipotesi da cui siamo partiti) si estende al caso in cui la \(C^t\) sia formata di curve riducibili, poichè in ogni modo una componente di \(C^t\) dovrà segare le \(K\) in almeno \(2\) punti: questa affermazione consegue dall’osservare che le \(C^t\) appartengono ad una serie continua di curve disequivalenti, che si ottiene sommando ad esse la differenza di due curve del fascio \((K)\): \(K-\overline{K}\).

\textit{Secondo caso:} \(p=1\).

Si proceda come nel caso precedente: a partire dal sistema \(|C|\) coi caratteri \(\pi\) ed \(n>2\pi-2\), si arriverà ad un ultimo sistema aggiunto \(|C^t|\), costituito di curve secanti le \(K\) in \(m\geq2\) punti; ma qui pare che il sistema aggiunto a \(C^t\) possa effettivamente mancare, se il genere (effettivo e virtuale) di \(C^t\) sia eguale ad \(1\), come quello del fascio delle \(K\). Però la nostra ipotesi porta che le curve \(C^t\) supposte irriducibili (o le componenti di esse) appartengano ad una schiera ellittica, \(\infty^1\), di curve ellittiche disequivalenti; e siccome le curve di questa schiera non possono formare un secondo fascio irrazionale (diverso da \(|K|\)), così dobbiamo ammettere che le curve \(C^t+K-\overline{K}\) diano luogo ad un sistema continuo \(\{C^t\}\) di curve di grado \(\nu\geq1\).

Ma di qui si trae che la superficie \(F\) deve possedere \(\infty^1\) curve razionali e quindi deve essere riferibile ad una rigata ellittica, in contraddizione colle nostre ipotesi. Infatti, se \(\nu=1\), la schiera ellittica \(\infty^1\) di curve \(C^t\) unisecantisi conterrà \(\infty^1\) coppie di curve equivalenti, e (per il teorema di Lüroth)\origfootnote{1}{s199R-1}{Cfr. \textsc{Enriques-Chisini}, \textit{Lezioni.} Libro II, cap. I, \S\ 3 (vol. I, pag. 169) e Libro V, cap. I, \S\ 7 (vol. III, pag. 48).} i loro punti d’intersezione genereranno su \(F\) \(\infty^1\) curve razionali. Se sia invece \(\nu>1\), le curve di \(C^t\) (irriducibili) con \(\nu-1\) punti fissi daran luogo ad una schiera \(\infty^1\) di curve ellittiche unisecantisi, sicchè si ricade nel caso precedente.

Così rimane stabilito il teorema fondamentale:

\textit{Le superficie contenenti un sistema lineare irriducibile, \(\infty^2\) almeno, di genere \(\pi\) e grado \(n>2\pi-2\), sono trasformabili in rigate, razionali o irrazionali.}

% scan: s200L
% !PRINTED 374
\weberpage{374}
% head: CAPITOLO DECIMO
% stamp: Scattered scan specks, a dark right-hand border, and a diagonal blemish near "qualcuno" and "certo".
E da ciò, come abbiamo annunciato nel \S\ 9 del Cap. IV, segue che:

\textit{Le superficie appartenenti alla famiglia delle rigate sono le sole superficie contenenti infinite curve eccezionali; e tutte le superficie, fuori di questa famiglia, possono trasformarsi in superficie prive di curve eccezionali.}

Quanto al \textit{teorema fondamentale}, conviene aggiungere che esso \textit{sussiste} nelle condizioni più ampie, \textit{qualunque sia la dimensione del sistema \(|C|\) fra i cui caratteri interceda la diseguaglianza \(n>2\pi-2\), purchè le \(C\) stesse non contengano come componenti delle curve eccezionali.}

Per dimostrarlo si osservi che il carattere \(a(C)=n-(2\pi-2)\) è un carattere additivo dei sistemi lineari a cui si riferisce, e perciò non può essere positivo per la somma di due sistemi senza essere positivo per qualcuno degli addendi. Così basterà dimostrare che:

\textit{Appartiene certo alla famiglia delle rigate qualunque superficie \(F\) su cui si trovi una curva \(C\), irriducibile e non eccezionale, di genere \(\pi\) e grado \(n>2\pi-2\).}

Infatti il multiplo della detta curva \(C\) secondo un numero \(s\) abbastanza grande avrà il genere e il grado
\[
\pi_s=s\pi+\frac{s(s-1)}{2}n-s+1,\qquad n_s=s^2n;
\]
e quindi avrà la dimensione
\[
r\geq p_a+n_s-\pi_s+1>0,
\]
costituendo un sistema irriducibile (per \(n>0\)) per cui
\[
n_s>2\pi_s-2.
\]

E si avverta esplicitamente che ciò che si è detto vale anche per \(\pi=0\), perchè, essendo escluso che \(C\) sia curva eccezionale e quindi che sia \(n=-1\), sarà
\[
n\geq0;
\]
ed invero, nel caso \(n=0\) che solo occorre esaminare, il multiplo \(|sC|\), per \(s\) assai grande, risulterà sempre composto colle curve razionali d’un fascio, donde segue che la \(F\) è trasformabile in una rigata.

\textit{Notizia.} – Il teorema che le superficie su cui l’aggiunzione si estingue, e perciò quelle che contengono sistemi lineari coi caratteri \(n\) e \(\pi\) per cui \(n>2\pi-2\), appartengono alla famiglia delle rigate; costituisce il risultato fondamentale della Memoria di \textsc{Castelnuovo-Enriques}, « Sopra alcune questioni fondamentali nella teoria delle superficie algebriche », pubblicata negli Annali di Matematica, 1901. % ?? uncertain: The punctuation after "rigate" resembles a semicolon -- semicolon / comma.
% EDITORIAL D0203 RISOLTO: The enlarged scan shows a small upper dot and a distinct lower comma tail: a semicolon.

% scan: s200R
% !PRINTED 375
\weberpage{375}
% head: LA SUPERFICIE DI GENERE GEOMETRICO NULLO
% stamp: Scattered scan specks and marginal traces; a curved extraneous mark touches "circostanze".
La dimostrazione originale di questo teorema non faceva uso della proprietà delle superficie irregolari di genere \(p_g=0\), di contenere un fascio irrazionale di curve (proprietà scoperta, come si è detto, nel 1904). In luogo di questa interveniva una diseguaglianza che deriva dal confronto dell’espressione noetheriana dell’invariante di Zeuthen-Segre col teorema di Riemann-Roch. Accenniamo in breve al concetto di questo ragionamento dimostrativo.

Se (sopra una superficie \(F\) che deve essere di genere \(p_g=0\)), a partire da un sistema lineare \(|C|\), l’aggiunzione si estingue, pervenendosi ad un ultimo sistema aggiunto \(|C^t|\), il genere di questo sistema dovrà essere \(\pi\leq-p_a=p\).

Supponiamo, per semplicità, che, \(|C^t|\) sia irriducibile, \(\infty^1\) almeno; designando con \(n\) il suo grado, avremo che l’invariante di Zeuthen-Segre della superficie vale % ?? uncertain: The mark after "che" is read as a comma -- comma / scan blemish.
\[
\delta-n-4\pi=-12p-\dot{p}^{(1)}+9,\qquad (p=-p_a)\text{\origfootmark{1}{s200R-1}} % ?? uncertain: A small mark above p in this formula and in its footnote is unclear -- p^{(1)} / \dot{p}^{(1)}.
% EDITORIAL D0205 RISOLTO: Both occurrences have a distinct centred dot above p. Its repetition at exactly the same typographical position in the formula and in the footnote settles the reading as dotted p.
\]\origfoottext{1}{s200R-1}{Qui \(\dot{p}^{(1)}\) designa il genere lineare relativo della superficie.}
mentre il teorema di Riemann-Roch ci dà la sua dimensione
\[
r\geq n-\pi+1-p.
\]
Di qui si deduce
\[
5\pi+r\geq11p+p^{(1)}-8+\delta\geq11p+p^{(1)}-8.
\]

Quindi una discussione appropriata, in cui occorre distinguere l’ipotesi \(p^{(1)}\leq1\) da quella \(p^{(1)}>1\), conduce a trovare sopra la superficie:

1) o un sistema lineare di curve di genere \(\pi\) e di dimensione
\[
r\geq3\pi-5,
\]
onde risulta che la superficie è razionale o riferibile ad una rigata di genere \(\pi\)\origfootnote{2}{s200R-2}{Questa è una conseguenza del ragionamento adoperato da \textsc{Enriques} nella Nota \textit{Sulla massima dimensione dei sistemi lineari di curve di dato genere appartenenti ad una superficie algebrica.} Atti Acc. di Torino, 1894.};

2) ovvero un sistema lineare di curve spezzate in curve razionali, onde risulta egualmente che la superficie è riferibile ad una rigata.

Questo cenno dà appena un’idea della via seguita: occorre poi un’analisi assai minuta, in rapporto a diverse circostanze complicatrici, e un esame speciale dei casi in cui \(p<3\).

La dimostrazione notevolmente più semplice del nostro teorema fondamentale, basata sul fascio irrazionale appartenente ad una
%CONT>>
% scan: s201L
% !PRINTED 376
\weberpage{376}
% head: CAPITOLO DECIMO
%<<CONT
% stamp: Dark scanner border at the right edge; scattered specks and marginal marks.
superficie irregolare di genere \(p_g=0\), è stata disegnata nelle « Lezioni » di \textsc{Enriques-Campedelli}\origfootnote{1}{s201L-1}{Cfr. Cap. III.} pubblicate dai Rendiconti del Seminario Matematico di Roma (1934). Codesta dimostrazione è stata qui sviluppata in una forma più semplice, evitando di fare appello, come ivi si faceva, al teorema di Castelnuovo-Enriques sui piani doppi che hanno i plurigeneri nulli. Supplisce a ciò la considerazione elementare adoperata innanzi, che una serie \(\infty^1\) di coppie di elementi equivalenti, sopra un ente ellittico \(\infty^1\), è una \(g_2^1\) e perciò è sicuramente razionale. % ?? uncertain: small mark before "Cap." in the footnote -- scan speck or an additional full stop.

Diremo infine che l'\textit{estensione del teorema fondamentale} al caso di \textit{una sola curva}, non eccezionale, di genere \(\pi\ (\geq 0)\) e grado \(n>2\pi-2\), è stata messa in luce da L. \textsc{Campedelli} nella Nota « Intorno alle superficie algebriche su cui esistono curve di genere \(\pi\) e di grado \(n>2\pi-2\) » (Rendic. Lincei, \(1933_2\)), e l'A. stesso ha rilevato come l'uso esplicito di tale proprietà valga a semplificare in più punti la nostra teoria. % ?? uncertain: small mark after "esplicito" -- scan speck or a printed full stop.

\textit{Nota.} – Merita speciale rilievo anche la proposizione di cui ci siamo serviti per dimostrare il teorema fondamentale di questo paragrafo. \textit{Le superficie su cui l'aggiunzione si estingue appartengono alla famiglia delle rigate.}

E pertanto, richiamando gli sviluppi dati in fine del § 1 del Cap. VI, potremo affermare che: \textit{Le superficie non appartenenti alla famiglia delle rigate, e perciò riducibili a un modello senza curve eccezionali, hanno il genere lineare assoluto}
\[
p^{(1)}\geq 1.
\]

Nel Cap. VI, § 1, si era riconosciuto soltanto che per \(p^{(1)}\leq 0\), estinguendosi l'aggiunzione, dovrebbero essere tutti i plurigeneri \(P_i=0\); qui si vede di più che la superficie dovrebbe essere riferibile ad una rigata (in contraddizione coll'ipotesi della eliminabilità delle curve eccezionali).

\edsection{3. Lemma sulle curve riducibili di un fascio.}{3. Lemma sulle curve riducibili di un fascio.}{sec.X.3}

In vista degli sviluppi che seguono occorre precisare quale sia la \textit{valenza} di una curva comunque spezzata, appartenente ad un fascio razionale o irrazionale, ai fini del computo delle curve dotate di punto doppio, il cui numero figura nell'espressione dell'invariante di Zeuthen-Segre (Cap. V, § 2, 3). Ci riferiremo, per semplicità, a superficie senza curve eccezionali, non appartenenti dunque alla famiglia delle rigate, e dimostreremo che: \textit{In ordine al computo dell'invariante}
%CONT>>
% scan: s201R
% !PRINTED 377
\weberpage{377}
% head: LA SUPERFICIE DI GENERE GEOMETRICO NULLO
%<<CONT
% stamp: Scattered scan specks and marginal marks.
\textit{di Zeuthen-Segre relativo ad una superficie che possegga un fascio \(|C|\), razionale o irrazionale, di curve di genere \(\pi>0\), irriducibili, privo di punti base, ogni curva spezzata del fascio ha la valenza \(\geq 0\).}

Di conseguenza il numero, \(\delta\) o \(\Delta\), delle curve del fascio dotate di punto doppio deve ritenersi sempre \(\geq 0\). \textit{Se \(\delta=0\) o \(\Delta=0\), il fascio non può contenere alcuna curva spezzata, tranne curve ellittiche multiple (di grado zero), nel caso che le curve del fascio siano ellittiche \((\pi=1)\).}

Per giustificare queste asserzioni, conviene richiamare la \textit{formula di Godeaux} che (Cap. V, § 2) abbiamo detto esprimere la valenza di una particolare curva \(\bar{C}\) del fascio, la quale sia comunque spezzata in parti multiple. Designamo con \(\bar{\delta}\) la valenza di una particolare \(C=\bar{C}\) che supporremo spezzata in componenti multiple di genere \(\varrho_1,\ \varrho_2\,\ldots.\):
\[
\bar{C}=t_1\theta_1+t_2\theta_2+\ldots.
\]
incontrantisi fra loro in
\[
s_{hk}=\theta_h\theta_k
\]
punti. Si avrà:
\[
\bar{\delta}\geq \Sigma(t_i-1)(2\varrho_i-2)+\Sigma s_{hk}(t_h+t_k-1),
\]
il segno \(>\) riferendosi a possibili complicazioni, per esempio al caso in cui qualcuna delle componenti \(\theta_i\) possegga un punto doppio, o a quello in cui alcuni dei punti di connessione comuni a due fra le componenti, \(\theta_h\) e \(\theta_k\), vengano a coincidere in un punto comune ad \(r>2\) componenti\origfootnote{1}{s201R-1}{La diseguaglianza \(\bar{\delta}\geq 0\) è stata riconosciuta da \textsc{Castelnuovo-Enriques} (Annali di Mat., 1901) nel caso più semplice in cui la \(\bar{C}\) sia spezzata in una curva irriducibile multipla \(\theta_1\) di genere \(\varrho_1\) e in un'altra componente semplice \(\theta_2\) che la incontri in \(s_{12}=\theta_1\theta_2\ (>0)\) punti. La formula generale sopra scritta viene pòrta, estendendo il ragionamento usato in quel caso, da L. \textsc{Godeaux}, \textit{Sur le calcul de l'invariant de Zeuthen-Segre.} (Mém. de la Soc. des Sciences de Hainart, 1920) e poi anche da L. \textsc{Campedelli}, \textit{Sul computo dell'invariante di Zeuthen-Segre; Ancora sul computo ecc.} Lincei, 1934. Quest'ultimo autore ne ha tratto la dimostrazione che \(\bar{\delta}\geq 0\) e più precisamente può aversi \(\bar{\delta}=0\) solo quando il fascio \((C)\) sia di genere \(\pi=1\) e la \(C\) sia una curva ellittica multipla di genere \(\varrho=1\): % ?? uncertain: place name in the journal citation -- Hainart or Hainaut.
% EDITORIAL D0208 RISOLTO: The enlarged print unambiguously shows an r, not a u, before t. Retain Hainart, with any suggested spelling correction confined to the apparatus.
\[
\bar{\delta}=(t+1)(2\varrho-2)=0. % [sic] (t+1) -> (t-1), the preceding general formula has the factor t_i-1.
\]

È la proprietà enunciata nel nostro lemma che figura pure nelle \textit{Lezioni} di \textsc{Enriques-Campedelli} (Sem. Mat. di Roma, 1934) e che ora riceve una dimostrazione notevolmente più semplice dovuta ad A. \textsc{Franchetta}.}.

Osserviamo che, essendo la superficie priva di curve eccezionali, si ha, per ogni \(i\),
\[
s_{ii}=\theta_i\theta_i\leq 2\varrho_i-2;
\]
%CONT>>
% scan: s202L
% !PRINTED 378
\weberpage{378}
% head: CAPITOLO DECIMO
%<<CONT
% stamp: Dark scanner border at the right edge; scattered specks and small marginal strokes.
quindi
\[
\bar{\delta}\geq \Sigma s_{ii}(t_i-1)+\Sigma s_{hk}(t_h+t_k-1);
\]
di più, poichè per \(h\neq k\) è \(s_{hk}\geq 0\),
\[
\begin{aligned}
\Sigma s_{hk}(t_h+t_k-1)&=\Sigma s_{hk}(t_h-1)+\Sigma s_{hk}(t_k-1)+\Sigma s_{hk}\geq\\
&\geq\Sigma s_{hk}(t_h-1)+\Sigma s_{hk}(t_k-1).
\end{aligned}
\]

Si ha dunque
\[
\bar{\delta}\geq\Sigma s_{ii}(t_i-1)+\Sigma s_{hk}(t_h-1)+\Sigma s_{hk}(t_k-1).
\]

Introduciamo la curva
\[
C^*=\theta_1+\theta_2+\ldots
\]
e indichiamo con \(\varrho\) e \(\nu\) i suoi caratteri virtuali. L'espressione a secondo membro della disuguaglianza precedente, uguaglia il numero delle intersezioni virtuali \(C^*(\bar{C}-C^*)\); pertanto il teorema sarà dimostrato se faremo vedere che \(C^*(\bar{C}-C^*)\geq 0\), ossia se ridurremo all'assurdo l'ipotesi
\[
C^*(\bar{C}-C^*)=-i\qquad(i>0).
\]

Notando che \(C^*\bar{C}=C^*C=0\), l'ipotesi adottata porterebbe
\[
\nu=C^*C^*=i>0;
\]
quindi la curva \(rC^*\) individuerebbe un sistema lineare \(|rC^*|\) di dimensione \(d_r\) che, per \(r\) abbastanza grande risulterebbe
\[
d_r\geq p_a+r^2i-r\varrho-\frac{r(r-1)}{2}i-r+1=p_a+\frac{r^2i}{2}-\frac{r}{2}(2\varrho-2-i)+1, % [sic] -r+1 -> +r+1, the printed right-hand expression has the opposite sign for this r term.
\]
ossia infinita del second'ordine al tendere di \(r\) all'infinito.

D'altra parte la curva \(r\bar{C}\) individua un sistema lineare \(|r\bar{C}|\), composto con le curve del fascio; la sua dimensione \(D_r\) è data, per \(r\) abbastanza grande, da \(D_r=r-p\), ove \(p\) è il genere del fascio; ossia è infinita del prim'ordine, al tendere di \(r\) all'infinito. Ma è chiaro che il sistema \(|rC^*|\) è contenuto parzialmente nel sistema \(|r\bar{C}|\), è quindi che \(d_r\leq D_r\). L'ipotesi da cui siamo partiti è dunque assurda. % [sic] è quindi -> e quindi, the accented verb interrupts the coordination of the two clauses introduced by "che".

Risulta inoltre che può essere \(\bar{\delta}=0\) solo se \(\Sigma s_{hk}=0\), e ciò è possibile solo se \(C^*\) è formata da una sola componente \(\theta\), e se \(\theta\cdot\theta=2\varrho-2=0\) ossia se \(\varrho=1\).

\edsection{4. Superficie di genere \(p_g=0\) con \(p_a<-1\).}{4. Superficie di genere p\_g=0 con p\_a<-1.}{sec.X.4}

Sia \(F\) una superficie di genere geometrico \(p_g=0\) e di genere numerico negativo \(p_a=-p<-1\). Sappiamo (Cap. IX, § 11) che essa
%CONT>>
% scan: s202R
% !PRINTED 379
\weberpage{379}
% head: LA SUPERFICIE DI GENERE GEOMETRICO NULLO
%<<CONT
% stamp: Dark scanner borders at the top and right; scattered specks.
deve contenere un fascio di genere \(p\) di curve \(K\), aventi il genere \(\pi\geq 0\). Proviamo che dev'essere \(\pi=0\) e perciò la \(F\) riferibile ad una rigata di genere \(p\), riducendo all'assurdo l'ipotesi \(\pi>0\).

Si assuma dunque come ipotesi che sia \(\pi>0\) e quindi la \(F\) (non riferibile a rigata) possa ridursi ad una superficie senza curve eccezionali; il suo genere lineare assoluto sarà (§ 2, Nota)
\[
p^{(1)}\geq 1.
\]
Scriviamo ora la formula che dà l'invariante di Zeuthen-Segre di \(F\) in relazione al fascio irrazionale \(K\). Avremo (Cap. V, § 2)
\[
\Delta-4+4(p-1)(\pi-1)=12p_a-p^{(1)}+9
\]
e quindi, per \(\pi>0\),
\[
\Delta\leq 13-12p-p^{(1)}.
\]
Ma da ciò segue, per \(p>1\), \(\Delta<0\) che è in contraddizione col lemma del nostro § 3.

Si conclude dunque \(\pi=0\). Così: \textit{Le superficie di genere \(p_g=0\) e \(p_a=-p<-1\), sono riferibili a rigate di genere \(p\).}

\edsection{5. Superficie \(F\) con \(p_g=0\) e \(p_a=-1\), possedenti un fascio ellittico di curve \(K\) di genere \(\pi>1\): lemma I.}{5. Superficie F con p\_g=0 e p\_a=-1, possedenti un fascio ellittico di curve K di genere π>1: lemma I.}{sec.X.5}

L'analisi delle superficie di generi \(p_g=0\) e \(p_a=-p=-1\) porta ad esaminare superficie \(F\) che posseggono un fascio ellittico di curve \(K\) di genere \(\pi\geq 1\); e conviene distinguere i due casi che si presentano possibili, quando la \(F\) non appartenga alla famiglia – delle rigate cioè: % ?? uncertain: dash after "famiglia" -- printed dash or stray scan mark.
% EDITORIAL D0209 RISOLTO: A regular horizontal dash follows famiglia at the end of the line; it is present in the printing.
\[
\pi>1\qquad\text{e}\qquad\pi=1.
\]

Discutiamo qui il caso \(\pi>1\).

Anzitutto si osservi che la nostra superficie \(F\), che possiamo assumere senza curve eccezionali, ha il genere lineare assoluto \(p^{(1)}=1\). Infatti, calcolando l'invariante di Zeuthen-Segre per il fascio ellittico \((K)\)\origfootnote{1}{s202R-1}{Cfr. V, 2 e 3.} si trova
\[
I=\Delta-4+4(p-1)(\pi-1)=\Delta-4=-12-p^{(1)}+9,
\]
sicchè, essendo \(\Delta\geq 0\), si deduce \(p^{(1)}\leq 1\), e quindi (§ 2, Nota)
\[
p^{(1)}=1.
\]

Giova anche notare che per il nostro fascio \((K)\) si ha di conseguenza
\[
\Delta=0,
\]
%CONT>>
% scan: s203L
% !PRINTED 380
\weberpage{380}
% head: CAPITOLO DECIMO
%<<CONT
% stamp: Dark scanner border at the right; scattered specks and a curved stray stroke between the first display and the following line.
e perciò (§ 3) il fascio stesso (formato di curve di genere maggiore di uno) non contiene alcuna curva \(K\) dotata di punto doppio, o comunque spezzata.

Ciò posto procediamo a stabilire una serie di lemmi che ci condurranno poi ad un teorema definitivo

\textsc{Lemma I.} – \textit{La superficie \(F\) possiede almeno una curva ellittica parabicanonica, secante le curve \(K\) del fascio in gruppi della serie bicanonica \(g_{4\pi-4}^{3\pi-4}\)}; e similmente \textit{una curva ellittica para-\(i\)-canonica}, per \(i>2\).

Infatti si consideri il sistema lineare \(|K''|\), secondo aggiunto ad una curva \(K\) del nostro fascio. Essendo \(p^{(1)}=1\), e le curve canoniche virtuali segando le \(K\) in gruppi di \(2\pi-2\) punti, il genere delle \(K'\) vale
\[
\pi'=\pi+1+2\pi-2-1=3\pi-2
\]
e quindi la dimensione di \(|K''|\) è
\[
r\geq-p+\pi'-1=3\pi-4,
\]
cioè almeno uguale a quella della serie bicanonica di \(K\). Ora, se la serie segata dal \(|K''|\) su \(K\) non è completa, per un gruppo di essa passa almeno un fascio di \(K''\), e si deduce l'esistenza di una curva \(K''\) spezzata nella \(K\) e in una curva bicanonica, che è di genere
\[
3p^{(1)}-2=1:
\]
questa costituisce una curva ellittica irriducibile, ovvero una curva composta di curve ellittiche senza punti comuni.

Nel caso che \(|K''|\) seghi su \(K\) la serie bicanonica completa, si consideri il sistema continuo \(\{K''\}\) (secondo paraggiunto), formato da \(\infty^1\) sistemi lineari, che si ottengono a partire da un particolare \(|K''|\), aggiungendovi la differenza di due curve \(K_1-K_2\) del fascio \(K\).

Notiamo che tutti i sistemi lineari formanti \(\{K''\}\) segheranno sopra una \(K\) la medesima serie bicanonica, e similmente riusciranno sempre equivalenti i gruppi segati da due curve comunque disequivalenti di un sistema paraggiunto (d'ordine \(1,\ 2,\ 3,\ \ldots.\)).

Segue di qui che si avranno (almeno) \(\infty^1\) curve di \(\{K''\}\) passanti per un medesimo gruppo bicanonico di \(K\), e quindi si avrà una curva di questo sistema spezzata in una \(K\) e in una residua curva parabicanonica: quest'ultima, come la curva bicanonica virtuale, sarà di genere
\[
3p^{(1)}-2=1
\]
e perciò sarà costituita da una o più curve ellittiche irriducibili senza punti comuni.

% scan: s203R
% !PRINTED 381
\weberpage{381}
% head: LA SUPERFICIE DI GENERE GEOMETRICO NULLO
% stamp: Scattered scan specks and marks along the right margin.
In modo perfettamente analogo si dimostra così che il sistema paraggiunto d'ordine \(i=3,\ 4\ldots.\), \(\{K^i\}\), di dimensione \((2i-1)\pi-2i\), contiene (almeno) \(\infty^1\) curve passanti per un gruppo della serie \(i\)-canonica scelto su una \(K\), e perciò anche una curva spezzata nella \(K\) e in una \textit{curva para-\(i\)-canonica} ellittica o formata di curve ellittiche irriducibili, senza punti comuni.

Ma non si può affermare che tutte queste curve ellittiche, per esempio la curva parabicanonica e la paratricanonica siano formate da componenti irriducibili \textit{distinte}: poichè tutte potrebbero ridursi ad una curva paracanonica multipla, cioè ad una curva ellittica il cui \textit{ordine} (numero delle intersezioni colle \(K\)) valga \(2\pi-2\).

\edsection{6. Lemma II: disegno della dimostrazione.}{6. Lemma II: disegno della dimostrazione.}{sec.X.6}

\textit{Sopra la superficie \(F\)} (di generi \(p_g=0\) e \(p_a=-1\), con un fascio di curve di genere \(\pi>1\)) \textit{esistono almeno due curve ellittiche, non formate dalle stesse componenti, i cui multipli costituiscono due curve parapluricanoniche}, secanti le \(K\) in gruppi di punti equivalenti.

Abbiamo già costruito una curva parabicanonica (o anche una parapluricanonica d'ordine superiore), diciamo \(C\), che è una curva di genere uno, formata di componenti irriducibili, semplici o multiple, senza punti comuni. Vediamo ora di costruire una seconda curva, \(D\), dello stesso ordine \(4\pi-4\), che non sia parabicanonica, ma seghi su ogni \(K\) una serie non bicanonica, il cui multiplo secondo \(s\) sia la serie \(2s\)-canonica, segata dalla curva \(sC\).

Più precisamente, nel caso in cui il genere \(\pi\) delle \(K\) sia pari, avendo dunque il divisore \(s=2\), ci proponiamo di costruire una curva ellittica \(D\) che, pur avendo lo stesso ordine \(4\pi-4\) di \(C\), seghi su ogni \(K\) un gruppo, non bicanonico, bensì appartenente ad una serie semiquadricanonica di questa curva. Ed a ciò riusciremo nel caso più semplice quando vengano soddisfatte le seguenti due ipotesi semplificative che discuteremo nei paragrafi successivi: I) su ogni curva \(K\) si può determinare razionalmente una serie \(\bar{g}\) semibicanonica (ma non canonica); II) entro la predetta \(\bar{g}\) si può determinare razionalmente un gruppo.

In queste ipotesi potremo anzitutto determinare razionalmente su una \(K\) una serie semibicanonica \(\bar{g}=g_{2\pi-2}^{\pi-2}\) (diversa dalla serie canonica \(g=g_{2\pi-2}^{\pi-1}\)) e quindi determinare in questa un gruppo \(G_{2\pi-2}\) e così costruire un sistema lineare di curve \(|L|\), abbastanza ampio, che seghi su ogni \(K\) la relativa serie \(\bar{g}\). Quindi il sistema lineare \(|L'|\), aggiunto ad \(|L|\), segherà su ciascuna \(K\) la serie semiquadricanonica \(\bar{g}+g\). Confrontando ora \(|2L'|\) e \(|2K''|\) (\(|K''|\) secondo aggiunto di una \(K\)), in ordine al secondo criterio d'equivalenza (qui precisato col lemma del § 3), riconosceremo che
%CONT>>
% scan: s204L
% !PRINTED 382
\weberpage{382}
% head: CAPITOLO DECIMO
%<<CONT
% stamp: Dark scanner border at the right; scattered specks and a stray mark before "(non verificata per F)".
essi differiscono soltanto per un certo numero \(r\) di curve \(K_1\,\ldots.\,K_r\) del fascio \((K)\), da contare due volte:
\[
|2L'|=|2K''+2(K_1+\ldots.+K_r)|,
\]
e ne dedurremo che il sistema \(|L'-(K_1+\ldots.+K_r)|\) ha gli stessi caratteri (grado, genere e quindi dimensione virtuale \(3\pi-4\)) di \(|K''|\). Perciò esso dà luogo ad un sistema continuo di dimensione \(3\pi-3\) (di una unità superiore a quella della serie segata su una \(K\)), da cui è possibile staccare una \(K\): la curva residua è la \(D\) semiparaquadricanonica che ci proponevamo di costruire: \(C\) e \(D\) sono certo disequivalenti, e non sono nemmeno paraequivalenti, mentre \(2C\) e \(2D\) sono paraequivalenti.

Qualora non si avverino le ipotesi II e 1 dette innanzi, il ragionamento si ripete \textit{mutatis mutandis} in questo senso: % [sic] II e 1 -> II e I, the hypotheses are labelled I and II.

1) Quando cada l'ipotesi II (sussistendo la I), in luogo di una \(\bar{g}\), serie semibicanonica (non canonica), si costruirà su \(K\) una serie \(\bar{\bar{g}}\) multipla della detta \(\bar{g}\) secondo un numero dispari \(i\) per cui la detta ipotesi sia soddisfatta. Quindi al posto del sistema \(|L|\) o meglio del suo aggiunto \(|L'|\), di cui si è discorso innanzi, si riuscirà a costruire un sistema lineare \(|M|\), che seghi su ogni \(K\) la serie \(\bar{\bar{g}}\). Allora il confronto di \(|2M|\) e di \(|2K^i|\) (dove \(|K^i|\) designa lo \(i\)-mo aggiunto di una \(K\)) ci farà riconoscere che il primo sistema, diminuito di un certo numero di curve \(K_1\,\ldots.\,K_r\) del fascio \((K)\) contate due volte, riesce paraequivalente al secondo. Da ciò si deduce che il sistema \(|M-(K_1+\ldots.+K_r)|\) ha gli stessi caratteri di \(|K^i|\); e quindi (§ 5) che appartiene ad un sistema continuo in cui si trova una curva spezzata in una \(K\) e in una residua curva ellittica \(D\).

La \(D\) ha il medesimo ordine della \(C=|K^i-\bar{K}|\) ma le due curve non sono equivalenti nè paraequivalenti, bensì hanno doppi paraequivalenti, che segano sulle \(K\) gruppi equivalenti.

2) In difetto dell'ipotesi I, si passerà dalla superficie \(F\) ad una \(F^*\) cogli stessi caratteri, rappresentata sulla \(F\) \(m\)-pla \((m>1)\) senza curva di diramazione, in modo che l'ipotesi semplificativa (non verificata per \(F\)) si avveri per \(F^*\). Si costruirà quindi su \(F^*\) una curva ellittica \(D\) (d'ordine \(4\pi-4\) o \(i(2\pi-2)\)) semiparapluricanonica, a cui risponderà sulla \(F\) una curva \(D\) parimente ellittica, (d'ordine \(m(4\pi-4)\), ovvero \(im(2\pi-2)\)). Questa curva \(D\) non è paraequivalente alla curva parapluricanonica \(C\) del medesimo ordine (cioè alla parabicanonica \(C=|K''-\bar{K}|\) che si costruisce su \(F\) secondo il § 5); ma i doppi \(2D\) e \(2C\) sono paraequivalenti, segando sulle \(K\) gruppi equivalenti.

Nel caso in cui \(\pi\) sia un numero dispari, indicando con \(s\) un suo divisore primo (dispari), ci proponiamo di costruire su \(F\) una curva
%CONT>>
% scan: s204R
% !PRINTED 383
\weberpage{383}
% head: LA SUPERFICIE DI GENERE GEOMETRICO NULLO
%<<CONT
% stamp: Dark scanner borders at the top and right; scattered specks.
ellittica \(D\) d'ordine \(4\pi-4\), che seghi su ogni \(K\) un gruppo, non bicanonico, appartenente ad una serie \(s\)-ma parte delle serie \(s\)-bicanonica. Ed a ciò riusciremo nel caso più semplice in cui vengano soddisfatte due ipotesi I e II, analoghe a quelle che si presentano nel primo caso di \(\pi\) pari ed \(s=2\). % [sic] delle serie -> della serie, the following adjective "s-bicanonica" is singular.

Invero, essendo data su ogni \(K\) secondo tali ipotesi una serie non speciale \(\bar{g}=g_{2\pi-2}^{\pi-2}\), \(s\)-ma parte della serie \(s\)-canonica, e in questa razionalmente un gruppo di punti, si costruirà un sistema lineare di curve abbastanza ampio \(|L|\), che seghi su ciascuna \(K\) la detta serie \(\bar{g}\). Quindi il sistema aggiunto \(|L'|\) (d'ordine \(4\pi-4\)) segherà sulla \(K\) la serie \(s\)-ma parte della serie \(s\)-bicanonica \(\bar{g}+g\). E si riconoscerà che \(|sL'|\) differisce da \(|sK''|\) (\(|K''|\) secondo aggiunto di una \(K\)) per un certo numero \(r\) di curve \(K_1\,\ldots.\,K_r\) del fascio \((K)\) da contare \(s\) volte:
\[
|sL'|=|sK''+s(K_1+\ldots.+K_r)|.
\]

In conseguenza di questa relazione, il sistema \(|L'-(K_1+\ldots.+K_r)|\) avrà i caratteri (grado, genere e dimensione \(3\pi-4\)) di \(|K''|\) e sarà contenuto in un sistema continuo \(\infty^{3\pi-3}\), cui appartiene una curva spezzata in \(K\) e in una \(D\): quest'ultima è una curva di genere 1 e grado 0, il cui \(s\)-plo riesce paraequivalente alla curva \(s\)-pla della parabicanonica \(C\).

Ma, come nel primo caso (\(\pi\) pari ed \(s=2\)) converrà esaminare in qual modo il ragionamento si modifichi ove non si verifichino le ipotesi I e II, sopra accennate.

1) Quando non si avveri l'ipotesi II, in luogo di una \(\bar{g}\), d'ordine \(2\pi-2\), diversa dalla serie canonica \(g\), ma \(s\)-ma parte della serie \(s\)-canonica, si prenderà in considerazione su \(K\) una serie \(\bar{\bar{g}}=(2\pi-2)\bar{g}-(\pi-2)g\), d'ordine \(\pi(2\pi-2)\), non pluricanonica, che dipende razionalmente dalla \(\bar{g}\) supposta data, secondo l'ipotesi I; della quale \(\bar{\bar{g}}\) si riuscirà a determinare \textit{un} gruppo. In tal guisa si costruirà, in luogo di \(|L|\) o meglio di \(|L'|\), un sistema lineare \(|M|\) che seghi su ciascuna \(K\) la relativa serie \(\bar{\bar{g}}\). Quindi si proverà che
\[
|sM|=|sK+s(K_1+\ldots.+K_r)| % [sic] sK -> sK^\pi, the following sentence compares the same system with |K^\pi|.
\]
e perciò
\[
|M-(K_1+\ldots.+K_r)|
\]
ha gli stessi caratteri di \(|K^\pi|\), e contiene una curva spezzata in una \(K\) e in una residua curva ellittica \(D\): la \(D\) ha lo stesso ordine \(\pi(2\pi-2)\) della curva para-\(\pi\)-canonica \(C=|K^\pi-\bar{K}|\), ma le due curve non sono equivalenti nè paraequivalenti; soltanto i multipli di esse, \(sC\) e \(sD\), costituiscono curve para-\(s\)-canoniche, secanti sulle \(K\) gruppi equivalenti. % [sic] para-s-canoniche -> para-s\pi-canoniche, the preceding text calls C para-\pi-canonica and then takes its s-fold multiple.

% scan: s205L
% !PRINTED 384
\weberpage{384}
% head: CAPITOLO DECIMO
% stamp: Dark scanner border at the right; scattered specks and short diagonal stray strokes near "anzitutto" and "annunciato".
2) Se poi non si avveri l'ipotesi I, si riuscirà ad ogni modo a costruire razionalmente su \(K\) non una, ma un gruppo di \(m>1\) serie \(\bar{g}\) (o \(\bar{\bar{g}}\)) e si otterrà quindi una superficie \(F^*\), cogli stessi caratteri di \(F\), rappresentata sopra la \(F\) \(m\)-pla senza curva di diramazione; sulla quale \(F^*\) le curve ellittiche \(K^*\), corrispondenti alle \(K\), contengono ciascuna una determinata serie \(\bar{g}\) (o \(\bar{\bar{g}}\)). Quindi (essendo verificata per \(F^*\) l'ipotesi I) si riuscirà a costruire su \(F^*\) una curva ellittica \(D^*\), parte aliquota d'una curva parapluricanonica, cui risponderà su \(F\) la cercata curva ellittica \(D\).

Abbiamo spiegato, in tal guisa, il disegno generale della dimostrazione del nostro lemma II; svolgiamo ora, nei suoi particolari, questa dimostrazione, distinguendo i due casi che si presentano: \(\pi\) pari ed \(s=2\), \(\pi\) dispari ed \(s\) suo divisore primo (diverso da 2).

\edsection{7. Primo caso: il genere \(\pi\) delle curve \(K\) sia pari.}{7. Primo caso: il genere π delle curve K sia pari.}{sec.X.7}

Ricordiamo anzitutto che la bisezione di una serie lineare \(g_{2n}^{2n-\pi}\) con \(n>\pi\), sopra una curva di genere \(\pi\), conduce a \(2^{2\pi}\) serie sempre distinte\origfootnote{1}{s205L-1}{Cfr. p. es. \textsc{Enriques-Chisini}, \textit{Lezioni}, Libro V, cap. III, § 35 (vol. III, pag. 394). Cfr. anche Libro VI, cap. III, § 34 (vol. IV, pag. 199).}.

In particolare la serie doppia della serie canonica sopra una \(K\), cioè la serie bicanonica \(g_{4\pi-4}^{3\pi-4}\), darà luogo a \(2^{2\pi}\) serie metà, una delle quali è la serie canonica \(g=g_{2\pi-2}^{\pi-1}\) e le altre sono \(2^{2\pi}-1\) serie semibicanoniche \(\bar{g}=g_{2\pi-2}^{\pi-2}\), distinte dalla \(g\) e distinte fra loro, che qui prendiamo a considerare. Siccome la \(K\), variando nel suo fascio, rimane irriducibile e non acquista mai punti doppi, le dette serie \(\bar{g}\) non vengono mai a coincidere; tuttavia esse potranno scambiarsi fra loro, per una circolazione della curva \(K\) nel proprio fascio (irrazionale) e converrà supporre in generale che sopra una \(K\) si possa definire razionalmente — non \textit{una} serie — ma un gruppo di \(m(\geq 1)\) serie \(\bar{g}=g_{2\pi-2}^{\pi-2}\) distinte fra loro.

Qui, come abbiamo annunciato, introdurremo due ipotesi semplificative, che saranno poi giustificate relativamente alle conseguenze che se ne deducono.

\gesperrt{Ipotesi I.} – Che sia \(m=1\), cioè che per ogni \(K\) sia dato determinare univocamente, e perciò \textit{razionalmente, una serie} \(\bar{g}\) semibicanonica (e non canonica).

\gesperrt{Ipotesi II.} – Che in una serie \(\bar{g}\) che si assume come razionalmente data, si possa determinare univocamente, e perciò \textit{razionalmente, un gruppo}, che, al variare di \(K\) su \(F\), darà luogo ad una
%CONT>>
% scan: s205R
% !PRINTED 385
\weberpage{385}
% head: LA SUPERFICIE DI GENERE GEOMETRICO NULLO
%<<CONT
% stamp: Scattered scan specks and small marginal marks.
curva \(L\) secante le \(K\) in gruppi semibicanonici di \(\bar{g}\). Questa seconda ipotesi potrebbe sembrare evidente, poichè la conoscenza della serie \(\bar{g}\) su \(K\), conduce in generale a trasformare la \(K\) in una \(K_{2\pi-2}\) dello spazio \(S_{\pi-2}\) proiettivamente definita, su cui gli iperpiani segano la serie \(\bar{g}\), e così questa serie viene data \textit{linearmente}, in guisa che un iperpiano di \(S_{\pi-2}\) vale a staccarne \textit{un} gruppo. Ma a priori non è detto che una serie \(\bar{g}=g^{\pi-2}\) razionalmente data su \(K\) possa venire linearmente definita in modo univoco. Giacchè fra le infinite curve \(K_{2\pi-2}\) immagini della serie \(\bar{g}\), fra loro proiettive, la scelta di una \(K_{2\pi-2}\) particolare, quale si riesce a determinare univocamente quando si parta da un gruppo della \(\bar{g}\) stessa, dipenderà in generale da qualche irrazionalità numerica che converrà aggiungere al campo di razionalità dato-con \(K\) e con \(\bar{g}\), e tali irrazionalità daranno poi — al variare di \(K\) — irrazionalità algebriche\origfootnote{1}{s205R-1}{Per chiarire la difficoltà si consideri una quartica piana \(K_4\) definita come inviluppo di una serie \(\infty^1\) d'indice 2 di coniche quadritangenti \(K_2\); su \(K_4\) resta così definita razionalmente la \(g_4^1\) delle quaterne di punti di contatto, che però non è definita linearmente mercè un determinato sistema lineare \(\infty^1\) di coniche con 4 punti base. Astrattamente la \(g_4^1\) può ritenersi come una curva razionale rappresentabile sopra una retta; ma per determinare la rappresentazione occorre dare un punto della detta curva razionale, che importa l'aggiunta di una irrazionalità quadratica. (Teorema di Noether, cfr. \textsc{Enriques-Chisini}, \textit{Lezioni.} Libro V, cap. III, § 32 (vol. III, pag. 341). % [sic] 341). -> 341))., the parenthesis opened before "Teorema" is not closed in print.
}. % [sic] dato-con -> dato con, an extraneous hyphen joins these words.

In base alle ipotesi I e II si costruisca intanto, sulla \(F\) una curva \(L\) che seghi sopra ciascuna \(K\) del nostro fascio ellittico un gruppo di \(2\pi-2\) punti della serie semibicanonica \(\bar{g}\) (che si è assunta come data). Le curve \(L'\) aggiunte ad \(L\) segheranno sulle \(K\) gruppi di una serie semiquadricanonica, \(\bar{g}+g\), diversa dalla bicanonica e formeranno un sistema lineare \(|L'|\) che potrà ritenersi ampio quanto si vuole, sommandovi — ove occorra — un certo numero di curve \(K\). Ora si confronti il sistema lineare \(|2L'|\) a \(|2K''|\) (doppio del secondo aggiunto ad una \(K\)). Siccome i due sistemi segano su ogni \(K\) gruppi equivalenti, per il secondo criterio d'equivalenza, precisato dal lemma del § 3, essi dovranno differire soltanto per curve del fascio \((K)\). Ma sopra questo fascio ellittico la differenza di due gruppi di elementi (il primo dei quali assai grande) riuscirà equivalente ad un gruppo di elementi da prendere positivamente, sicchè si potrà scrivere
\[
|2L'|=|2K''+K_1+K_2+\ldots.|. % ?? uncertain: small marks above K_1 and K_2 in this display -- scan specks or overbars.
% EDITORIAL D0210 RISOLTO: Magnified comparison shows detached irregular specks, not horizontal bars; retain both unbarred K symbols.
\]

Distinguiamo due casi che si presentano a priori possibili, cioè che il numero delle \(K_1\ K_2\,\ldots.\) sia un numero pari, \(2r\), ovvero che sia un numero dispari \(2r+1\).

% EDITORIAL: segnatura originale: \textsc{Enriques F.} - \textit{Superficie algebriche.}\hfill 25
\printersignature{\textsc{Enriques F.} - \textit{Superficie algebriche.}\quad 25}

% scan: s206L
% !PRINTED 386
\weberpage{386}
% head: CAPITOLO DECIMO
% stamp: Scattered specks and stray strokes near the formulas; a dark scanning band along the right edge.
Nel primo caso, facendo variare il gruppo delle curve \(K_1\,K_2\,\ldotp\ldotp\ldotp\ldotp\) entro la serie lineare \(g_{2r}^{2r-1}\) che esso determina nel fascio ellittico si può fare in modo che il gruppo stesso si riduca ad un gruppo contato due volte:
\[
\begin{aligned}
&2K_1+2K_2+\ldotp\ldotp\ldotp\ldotp+2K_r\\
&(K_1=K_{r+1},\quad K_2=K_{r+2}\,\ldotp\ldotp\ldotp\ldotp).
\end{aligned}
\]

Nel secondo caso si può ridurre il gruppo suddetto ad un altro di \(r+1\) curve, fra cui \(r\) curve contate due volte:
\[
2K_1+2K_2+\ldotp\ldotp\ldotp\ldotp+2K_r+K_0.
\]

Quindi si dovrà avere
\[
|2L'-2K''-2(K+K_1+\ldotp\ldotp\ldotp\ldotp+K_r)|=|K_0|; % [sic] 2(K+K_1+\ldotp\ldotp\ldotp\ldotp+K_r) -> 2(K_1+\ldotp\ldotp\ldotp\ldotp+K_r), the following virtual curve omits the unindexed K.
\]
e perciò i caratteri (\(\pi\) e \(n=0\)) di \(K_0\) si esprimono mediante quelli (genere \(\varrho\) e grado \(\nu\)) della curva virtuale
\[
L'-K''-(K_1+\ldotp\ldotp\ldotp\ldotp+K_r)
\]
il cui doppio equivale a \(K_0\).

Si avrà dunque
\[
\begin{aligned}
n&=4\nu=0\\
\pi&=2\varrho+\nu-1=2\varrho-1,
\end{aligned}
\]
e per conseguenza \(\pi\) deve essere dispari. Essendosi supposto \(\pi\) pari, il secondo caso esaminato come possibile non può presentarsi e si avrà
\[
|2L'-2(K_1+\ldotp\ldotp\ldotp\ldotp+K_r)|=|2K''|.
\]

Ora il sistema lineare \(|2L'-2(K_1+\ldotp\ldotp\ldotp\ldotp+K_r)|\) ammette due metà disequivalenti:
\[
|L'-(K_1+\ldotp\ldotp\ldotp\ldotp+K_r)|
\]
e
\[
|K''|.
\]

Pertanto il primo sistema avrà gli stessi caratteri, genere e grado, del secondo, che sono:
\[
4\pi-4+\pi+1-1=5\pi-4
\]
e
\[
2\cdot(4\pi-4)=8\pi-8.
\]

La sua dimensione sarà dunque
\[
3\pi-4,
\]
che è la dimensione della serie completa da esso segata su una \(K\). Per conseguenza un gruppo di \(3\pi-4\) punti di questa serie, apparterrà ad \(\infty^1\) curve del sistema continuo, \(\{L'-(K_1+\ldotp\ldotp\ldotp\ldotp+K_r)\}\)
%CONT>>
% scan: s206R
% !PRINTED 387
\weberpage{387}
% head: LA SUPERFICIE DI GENERE GEOMETRICO NULLO
%<<CONT
% stamp: Scattered specks and faint vertical scanning streaks at the right edge.
formato da \(\infty^1\) sistemi lineari disequivalenti, e vi sarà almeno una curva di questo sistema continuo, spezzata in una \(K\) e in una curva residua \(D\): il genere di \(D\) risulta eguale a quello di \(C=|K''-\bar K|\), cioè eguale ad \textit{uno}. E — come si richiedeva — la curva ellittica \(D\), dello stesso ordine della parabicanonica \(C\), non è equivalente (bensì semiequivalente) ad essa e perciò \textit{essenzialmente distinta} da quella, cioè contiene certo una qualche componente (ellittica) distinta. % ?? uncertain: punctuation after "disequivalenti" -- a comma or a scanning speck.
% EDITORIAL D0211 RISOLTO: A curved descender below the baseline marks a comma after disequivalenti.

Ma conviene discutere le ipotesi I e II che figurano come presupposto della precedente deduzione. Cominciamo dalla

\begin{center}
\textit{Discussione dell’ipotesi II.}
\end{center}

L’ipotesi II, verificata per \(\pi=2\)\origfootnote{1}{s206R-1}{In questo caso la serie \(\bar g\) si riduce ad una coppia unica, che perciò è data colla \(\bar g\).}, non è più vera a rigore per i valori superiori del numero pari \(\pi\) \((\pi=4,\,6\,\ldotp\ldotp\ldotp\ldotp)\). Ma, se non si può, in generale, staccare su \(K\) \textit{un} gruppo \(G\) della \(\bar g\) o della serie semibicanonica \(\bar g+g\) (che si ottiene dalla \(\bar g\) supposta data razionalmente) dimostriamo tuttavia che si riesce a determinare un gruppo formato da \(i=\pi-1\) gruppi \(G\) (\(i\) dispari), il quale appartiene ad una serie — diciamo \(\bar{\bar g}=i\bar g\) — non pluricanonica, ma metà della serie \(2i\)-canonica. Quindi si potrà costruire su \(F\) una curva \(M\) che seghi ogni \(K\) in un gruppo della corrispondente \(\bar{\bar g}\), e operando su \(|M|\) come prima si operava su \(|L|\) o meglio su \(|L'|\), cioè staccando un certo numero di curve \(K\), si otterrà una curva \(D\), semipara-\(2i\)-canonica ellittica (essenzialmente distinta dalla \(iC\)) il cui doppio è paraequivalente alla curva bicanonica \(C\) contata \(i\) volte.

Dunque il nostro ragionamento fondato sull’ipotesi II si ripete con lievi modificazioni (sostituendo la serie \(\bar{\bar g}\) a \(\bar g\) o a \(\bar g+g\)), ove si dimostri quello che sopra è enunciato: potersi determinare \textit{un} gruppo formato da \(i=\pi-1\) (numero dispari) gruppi \(G\) di \(\bar g\) (e quindi linearmente \(\bar{\bar g}\)). A tale scopo si consideri nello spazio \(S_{\pi-2}\) la famiglia di tutte le curve proiettive fra loro, \(K_{2\pi-2}\), che sono immagini della serie semibicanonica \(\bar g=g_{2\pi-2}^{\pi-2}\) data su \(K\). E si cerchi di costruire due quadriche \(Q\) e \(Q'\), o un fascio di quadriche, in codesto spazio, che sia (proiettivamente) covariante di \(K_{2\pi-2}\): allora resterà determinato come covariante del fascio lo \(i\)-gono, con \(i=\pi-1\), autoconiugato comune a \(Q\) e \(Q'\), le cui facce segheranno su \(K_{2\pi-2}\) \(\pi-1\) gruppi della \(\bar g\).

La costruzione del fascio anzidetto di \(Q\) e \(Q'\) si fa univocamente e perciò razionalmente a partire dalla \(K\) su cui è data \(\bar g\): ciò appare subito ove \(K_{2\pi-2}\) non appartenga ad alcuna quadrica di \(S_{\pi-2}\), perchè la minima serie doppia della serie \(\bar g\), che è una serie bicanonica
%CONT>>
% scan: s207L
% !PRINTED 388
\weberpage{388}
% head: CAPITOLO DECIMO
%<<CONT
% stamp: A slanting stray stroke in the middle text, scattered marginal specks, and a dark right-edge scanning band.
(completa o no) segata su \(K_{2\pi-2}\) dalle quadriche del suo spazio, viene definita linearmente sulla relativa curva bicanonica \(K_{4\pi-4}\) dello \(S_{3\pi-4}\) (pur essa linearmente data), mediante gli iperpiani d’un sistema lineare \(\alpha\) ben determinato dalla curva \(K_{4\pi-4}\). Così un fascio di quadriche dello \(S_{\pi-2}\) riesce determinato razionalmente in funzione di \(\bar g\), da un fascio d’iperpiani dello \(S_{3\pi-4}\) scelto in \(\alpha\) imponendo il passaggio per un conveniente gruppo di punti, dello \(S_{3\pi-4}\), fissati indipsndentemente dalla \(K_{4\pi-4}\): ossia da una serie \(g_{4\pi-4}^1\) contenuta nella bicanonica \(g_{4\pi-4}^{3\pi-4}\) (e più precisamente nella minima serie doppia di \(\bar g\)). % [sic] indipsndentemente -> indipendentemente, apparent substitution of s for e.

Ma se la \(K_{2\pi-2}\) appartenga ad \(\infty^r\) quadriche di \(S_{\pi-2}\) (\(r\geq1\)) il ragionamento precedente cade in difetto, perchè ad un gruppo della serie minima doppia di \(\bar g\), definito come sezione d’un iperpiano di \(\alpha\) sulla \(K_{4\pi-4}\), non risponde più una quadrica dello \(S_{\pi-2}\) bensì un sistema lineare \(\infty^{r+1}\) di quadriche; e similmente ad un sistema lineare \(\infty^t\) di iperpiani dello \(S_{3\pi-4}\) scelto in \(\alpha\), risponde un sistema lineare \(\infty^{r+t+1}\) di quadriche dello \(S_{\pi-2}\), secante su \(K_{\pi-2}\) una serie di dimensione minore \(g_{4\pi-4}^t\). Tuttavia, facendo % [sic] K_{\pi-2} -> K_{2\pi-2}, the curve has the latter subscript earlier in this paragraph.
\[
t=3\pi-4-2=3\pi-6,
\]
si avrà nello \(S_{\pi-2}\) un sistema lineare di quadriche \(|Q|\), di 2 unità inferiore al sistema totale: ora c’è una schiera armonica associata\origfootnote{1}{s207L-1}{Date in uno spazio lineare a quante si vogliano dimensioni una quadrica-luogo \(a_x^2=0\) e una quadrica inviluppo \(\alpha_u^2=0\), c’è (secondo \textsc{Battaglini} e \textsc{Rosanes}) un invariante lineare simultaneo di esse, l’\textit{armonizzante} \(a_\alpha^2\), che col suo annullamento definisce la relazione di armonia fra una quadrica luogo e il sistema lineare associato di quadriche inviluppo. Cfr. \textsc{Enriques-Chisini}, \textit{Lezioni}. Libro III cap. I, \S~9 (vol. I, pag. 61).} a \(|Q|\) di quadriche inviluppo, che definisce un \(i\)-gono (\(i=\pi-1\)) coniugato e quindi un gruppo di \(i\) gruppi di \(\bar g\) \hfill c. d. d.

\begin{center}
\textit{Discussione dell’ipotesi I.}
\end{center}

Giustificata l’ipotesi II (almeno nel senso che se essa non sia verificata per la serie \(\bar g\), lo sarà per \(\bar{\bar g}\)) vediamo ora come si giustifichi relativamente anche l’ipotesi I, in questo senso che — se essa non si avveri per la nostra superficie \(F\) — si avvererà per una \(F^*\) cogli stessi caratteri rappresentata sulla \(F\) contata \(m\) \((>1)\) volte, senza curva di diramazione, sicchè sarà dato costruire su \(F^*\) una curva ellittica \(D^*\), cui risponderà su \(F\) una \(D\) semiparapluricanonica.

Pongasi che sopra ciascuna \(K\) di \(F\) sia dato segare univocamente non \textit{una} serie semibicanonica \(\bar g\), ma un gruppo di \(m\) \((>1)\) serie \(\bar g\) siffatte, scambiabili fra loro per una circolazione di \(K\) nel
%CONT>>
% scan: s207R
% !PRINTED 389
\weberpage{389}
% head: LA SUPERFICIE DI GENERE GEOMETRICO NULLO
%<<CONT
% stamp: Scattered marginal specks and short stray strokes.
fascio \((K)\). E, per semplicità di discorso, si supponga \(\pi>2\) e quindi \(\pi\geq4\), e si accetti in rapporto ad ognuna di queste \(\bar g\) l’ipotesi II che, in funzione razionale di essa, possa costruirsene una immagine lineare, cioè si riesca a definire \textit{una} curva \(K_{\pi-2}\) dello \(S_{\pi-2}\), scelta fra le infinite curve proiettive corrispondenti. % [sic] K_{\pi-2} -> K_{2\pi-2}, the corresponding curve is given the latter subscript below.

La varietà algebrica \(\infty^1\) che ha per elementi le nostre serie \(g\) sulle \(K\) (associate a queste stesse \(K\)) si può ritenere astrattamente come una serie \(\infty^1\Sigma\) di spazi \(S_{\pi-2}\), contenuta in uno spazio di più dimensioni \(S_h\). Siccome due fra le \(m\) serie \(\bar g\) su una \(K\) non vengono mai a coincidere, la detta serie, \(\Sigma\), contiene una involuzione ellittica \(\gamma_m^1\) di spazi, senza elementi doppi, e perciò è essa stessa ellittica. % [sic] g -> \bar g, the same series carry an overbar later in the paragraph.

Ora in ciascuno spazio \(S_{\pi-2}\) di \(\Sigma\) si può costruire (per ipotesi) \textit{una} curva \(K^*=K_{2\pi-2}\), immagine della \(\bar g\) omologa scelta sulla relativa \(K\): il luogo delle curve \(K^*\) (birazionalmente identiche alle \(K\)) sarà una superficie \(F^*\), rappresentata sulla \(F\) \(m\)-pla (priva di curva di diramazione) in guisa che ad ogni \(K\) di \(F\) risponderà su \(F^*\) un gruppo di \(m\) curve \(K^*\), sempre distinte fra loro. Sulla \(F^*\) le curve \(K^*\), birazionalmente identiche alle \(K\), formano \textit{un fascio ellittico} \((K^*)\).

È facile vedere che la superficie \(F^*\) ha, come \(F\), il genere numerico \(p_a=-1\) e il genere lineare \(p^{(1)}=1\), d’accordo colla formula che dà l’invariante di Zeuthen-Segre di \(F^*\) in relazione al fascio ellittico delle \(K\), prive di punti doppi.

Ciò posto si può dire per la \(F^*\) quel che — nell’ipotesi I — si diceva per la \(F\), giacchè sulla \(F^*\), per ogni curva \(K^*\) (di genere \(\pi\)), è data \textit{una} serie semibicanonica \(\bar g\), razionalmente e linearmente definita dagli iperpiani di \(S_h\). Si deduce che la \(F^*\) contiene una curva semiparabicanonica ellittica \(D^*\), cui risponderà su \(F\) una curva \(D\), egualmente ellittica, semipara-\(2m\)-canonica.

\textit{Osservazione.} -- Nel discorso che precede si è esclusa l’ipotesi \(\pi=2\), in cui, tuttavia, il discorso stesso potrebbe ripetersi con qualche modificazione. Ma basti per questo caso avvertire che, se le \(K\) sono di genere \(\pi=2\), su ciascuna di esse viene determinato il gruppo dei 6 punti doppi della \(g_2^1\), sempre distinti fra loro, che — al variare di \(K\) — descriveranno una curva ellittica semiparasesticanonica

\edsection{8. Secondo caso: il genere \(\pi\) delle curve \(K\) sia dispari.}{8. Secondo caso: il genere π delle curve K sia dispari.}{sec.X.8}

Designeremo con \(s\) un numero primo (dispari) divisore di \(\pi\), e considereremo su una \(K\) qualsiasi una fra le serie non speciali \(\bar g = g_{2\pi-2}^{\pi-2}\) che si ottengono dalla divisione per \(s\) della serie \(s\)-canonica. Introduciamo, provvisoriamente, due ipotesi semplificative analoghe alle I e II del paragrafo precedente.
% The equality sign is repeated at the printed line break.

% scan: s208L
% !PRINTED 390
\weberpage{390}
% head: CAPITOLO DECIMO
% stamp: A dark right-edge scanning band and scattered spots and stray strokes.
\textit{Ipotesi I}: che si possa determinare univocamente, cioè razionalmente, sulla \(K\) variabile del fascio \((K)\) una serie \(\bar g\).

\textit{Ipotesi II}: che, data razionalmente sopra ogni \(K\) una \(\bar g\), si possa determinare razionalmente \textit{un} gruppo, che così venga definito univocamente al variare di \(K\) in \((K)\); in altre parole che la \(\bar g\), data razionalmente, possa definirsi linearmente (dando luogo ad \textit{una} \(K_{\pi-2}\) in \(S_\pi\)). % [sic] K_{\pi-2} in S_\pi -> K_{2\pi-2} in S_{\pi-2}, comparison with the same construction on the preceding pages.

In base a queste ipotesi si può costruire su \(F\) una curva \(L\) che seghi ciascuna \(K\) in un gruppo di \(\bar g\); ed \(|L|\), che è definito a meno di addendi eguali a \(K\), costituirà un sistema lineare ampio quanto sivuole. % [sic] sivuole (printed si- / vuole) -> si vuole, an anomalous line-end hyphen joins two separate words.

Indichiamo con \(|L'|\) il sistema lineare, d’ordine \(4\pi-4\) aggiunto di \(|L|\), e con \(|K''|\) il secondo aggiunto ad una \(K\), che è del medesimo ordine.

Confrontiamo, rispetto al solito criterio d’equivalenza, i due sistemi \(|sL'|\) e \(|sK''|\); essi differiranno per un certo numero di curve del fascio \((K)\) da sommare al sistema meno ampio \(|sK''|\); anzi è lecito assumere che la curva sommanda a questo sistema minore sia formata da un certo numero di curve \(s\)-ple
\[
sK_1,\quad sK_2,\quad \ldotp\ldotp\ldotp\ldotp\quad sK_r
\]
ed eventualmente anche da una curva \(tK_0\), \(t\)-pla, del detto fascio, dove \(t<s\):
\[
|sL'|=|sK''+s(K_1+\ldotp\ldotp\ldotp\ldotp+K_r)+tK_0|.
\]

Ma si prova che
\[
t=0,
\]
perchè la curva virtuale
\[
\left(\frac{t}{s}K_0\right)=|L'-K''-(K_1+\ldotp\ldotp\ldotp\ldotp+K_r)|
\]
dovrebbe avere un genere \(\varrho\) intero, mentre il genere di \(tK_0\) (di grado 0) vale
\[
t\pi-t+1=s\varrho-s+1
\]
sicchè
\[
t(\pi-1)=s(\varrho-1):
\]
questa relazione non può sussistere con \(\varrho\) intero, perchè \(s\) (divisore primo di \(\pi\)) è primo con \(\pi-1\) e con \(t\).

Si avrà dunque
\[
|sL'|=|sK''+s(K_1+\ldotp\ldotp\ldotp\ldotp+K_r)|
\]
e
\[
|sL'-s(K_1+\ldotp\ldotp\ldotp\ldotp+K_r)|=|sK''|.
\]

% scan: s208R
% !PRINTED 391
\weberpage{391}
% head: LA SUPERFICIE DI GENERE GEOMETRICO NULLO
% stamp: Scattered specks and short stray strokes around the text.
Quindi
\[
|L'-(K_1+\ldotp\ldotp\ldotp\ldotp+K_r)|\quad \text{e}\quad |K''|
\]
avranno i medesimi caratteri (grado \(8\pi-8\) e genere \(5\pi-4\)) e perciò la stessa dimensione virtuale \(3\pi-4\), che è la dimensione della serie da essi segata su una \(K\). Si deduce che, nel sistema continuo \(\infty^{3\pi-3}\{L'-(K_1+\ldotp\ldotp\ldotp\ldotp+K_r)\}\) vi è una curva spezzata in una \(K\) e in una residua curva \(D\), che (come la parabicanonica \(C = |K''-\bar K|\)) sarà ellittica, d’ordine \(4\pi-4\); e la \(sD\), come la \(sC\), sarà una curva para-\(2s\)-canonica, le due curve segando le \(K\) in gruppi (\(2s\)-canonici) equivalenti. % ?? uncertain: the prime on K in the expression for C -- one prime or two primes.
% EDITORIAL D0212 RISOLTO: Two adjacent prime strokes are visible at upper right of K.
% The equality sign in the expression for C is repeated at the printed line break.

Quando non si avverino le ipotesi II e I il ragionamento precedente può tuttavia ripetersi con qualche modificazione.

\textit{Se non si avveri l’ipotesi II}, cioè non si possa determinare univocamente un gruppo della serie non speciale \(\bar g=g_{2\pi-2}^{\pi-2}\), \(s\)-ma parte della serie \(s\)-canonica, che si suppone razionalmente data secondo l’ipotesi I, converrà considerare su \(K\), al posto della serie non speciale \(\bar g\), \(s\)-ma parte della serie \(s\)-canonica, la serie
\[
\bar{\bar g}=(2\pi-2)\bar g-(\pi-2)g,
\]
di cui si definisce razionalmente un gruppo prendendo la somma dei gruppi \(\bar G_{2\pi-2}\) di \(\bar g\) che hanno come \((\pi-2)\)-plo un punto di un gruppo (canonico) \(G_{2\pi-2}\) di \(g\) e staccandone poi il \(G_{2\pi-2}\) contato \(\pi-2\) volte: è essenziale notare che la \(\bar{\bar g}\), il cui \(s\)-plo è una serie \(s\pi\)-canonica, non è una serie \(\pi\)-canonica, perchè altrimenti dalle relazioni
\[
(2\pi-2)\bar g-(\pi-2)g=\pi g
\]
ossia
\[
(2\pi-2)\bar g=(2\pi-2)g
\]
e
\[
s\bar g=sg,
\]
essendo \(2\pi-2\) ed \(s\) numeri primi fra loro, si ricaverebbe
\[
\bar g=g\origfootmark{1}{s208R-1}.
\]\origfoottext{1}{s208R-1}{Basta risolvere l’equazione d’analisi indeterminata \((2\pi-2)x-sy=1\); si dedurrà quindi
\[
[(2\pi-2)x-sy]\bar g=[(2\pi-2)x-sy]g.
\]
}

Ora, al posto di \(|L|\) o meglio di \(|L'|\), costruiremo un sistema lineare \(|M|\), abbastanza ampio, che seghi sopra ciascuna \(K\) la serie \(\bar{\bar g}\) (costruita in funzione della \(\bar g\) che si suppone razionalmente data). Confrontiamo \(|sM|\) col sistema lineare \(|sK^\pi|\) multiplo di \(|K^\pi|\),
%CONT>>
% scan: s209L
% !PRINTED 392
\weberpage{392}
% head: CAPITOLO DECIMO
%<<CONT
% stamp: Scattered specks and stray strokes, with a dark scanning band along the right edge.
\(\pi\)-mo aggiunto ad una \(K\); designando con \(K_0,\,K_1\,\ldotp\ldotp\ldotp\ldotp\,K_r\) curve particolari del fascio \((K)\), avremo
\[
|sM|=|sK^\pi+s(K_1+\ldotp\ldotp\ldotp\ldotp+K_r)+tK_0|,
\]
dove \(t=0\) e perciò
\[
|sM|=|sK^\pi+s(K_1+\ldotp\ldotp\ldotp\ldotp+K_r)|.
\]

Di qui si deduce che il sistema lineare
\[
|M-(K_1+\ldotp\ldotp\ldotp\ldotp+K_r)|
\]
ha gli stessi caratteri — grado, genere e dimensione virtuale — di \(|K^\pi|\) e perciò appartiene ad un sistema continuo la cui dimensione supera di 1 quella della serie \(\bar{\bar g}\), entro cui si troverà dunque una curva spezzata in una \(K\) e in una residua curva ellittica \(D\), d’ordine \(\pi(2\pi-2)\), il cui multiplo secondo \(s\) costituirà una curva para-\(s\pi\)-canonica.
% The minus sign in the expression for the order is repeated at the printed line break.

Infine discutiamo ciò che può farsi se per la superficie \textit{non si avveri l’ipotesi I}: vuol dire che sopra una curva \(K\) variabile non ci è dato determinare razionalmente \textit{una serie non speciale} \(\bar g\), fra quelle il cui \(s\)-plo costituisce la serie \(s\)-canonica, ma soltanto un gruppo di \(m\) \((>1)\) serie \(\bar g\), scambiabili fra loro per una circolazione di \(K\) nel fascio \((K)\). Come già abbiam visto nel primo caso (\(\pi\) pari ed \(s=2\)) conviene allora passare dalla superficie \(F\) ad una \(F^*\), cogli stessi caratteri, rappresentata sulla \(F\) \(m\)-pla senza curva di diramazione, per la quale \(F^*\) si verifichi l’ipotesi II. Quindi si costruirà su \(F^*\) una curva ellittica \(D^*\), d’un certo ordine \((2\pi-2)\), o più in generale \(\pi(2\pi-2)\), a cui risponderà su \(F\) una \(D\), d’ordine \(m\pi(2\pi-2)\) parimente ellittica: e la \(D\), \(s\)-ma parte di una curva para-\(sm\pi\)-canonica, conterrà certo qualche componente distinta da quelle di una \(C\) para-\(m\pi\)-canonica, costruita su \(F\) secondo il \S~5.

\edsection{9. Lemma III.}{9. Lemma III.}{sec.X.9}

\textit{Sopra una superficie \(F\), di generi \(p_g=0\) e \(p_a=-1\), possedente un fascio ellittico di curve di genere \(\pi>1\), si abbia una curva \(E\), la quale seghi le \(K\) in gruppi di punti equivalenti a quelli segati da una curva para-\(i\)-canonica \(C\); se, a prescindere da eventuali componenti comuni colla \(C\), la \(E\) non ha punti comuni colla \(C\), essa è come la \(C\) una curva para-\(i\)-canonica ellittica.}

Si confrontino la \(E\) e la \(C\) in ordine al secondo criterio d’equivalenza (qui precisato col lemma del \S~3): esse saranno curve equi%
% The word "equivalenti" continues across the page break.
%CONT>>
% scan: s209R
% !PRINTED 393
\unskip\weberpage{393}%
% head: LA SUPERFICIE DI GENERE GEOMETRICO NULLO
%<<CONT
% stamp: Scattered marginal specks and short stray strokes.
valenti a meno della differenza \(K_1-K_2\) di due curve del fascio \((K)\):
\[
|E|=|C+K_1-K_2|.
\]

Quindi la \(E\) avrà lo stesso genere \textit{uno} e lo stesso grado \textit{zero} della \(C\) e sarà, com’essa, una curva para-\(i\)-canonica.

\edsection{10. Conclusione: le superficie \(F\) posseggono anche un fascio lineare di curve ellittiche.}{10. Conclusione: le superficie F posseggono anche un fascio lineare di curve ellittiche.}{sec.X.10}

I lemmi I, II e III, stabiliti nei precedenti paragrafi ci conducono ora ad un teorema che conchiude l’analisi delle nostre superficie \(F\).

\textit{Le superficie \(F\), di genere \(p_g=0\) e \(p_a=-1\), aventi un fascio ellittico di curve \(K\) di genere \(\pi>1\), posseggono altresì un fascio lineare, senza punti base, di curve ellittiche.}

Invero si costruiscano su \(F\) (lemma II) due curve ellittiche para-\(i\)-canoniche \(C\) e \(D\), secanti le \(K\) in gruppi equivalenti di \(n=i(2\pi-2)\) punti. Su ciascuna \(K\) le dette curve determinano una serie lineare \(g_n^1\), la quale possiederà un certo numero
\[
m\leq2n+2\pi-2
\]
di gruppi \(G_n\) dotati d’un punto doppio, all’infuori di quelli che possono essere assorbiti dai gruppi sezioni di \(C\) e di \(D\). Il luogo dei detti gruppi \(G_n\) è una curva \(E\) che ha (rispetto alle \(K\)) lo stesso ordine \(mn\) delle \(mC\) e \(mD\), e sega le \(K\) secondo gruppi equivalenti alle sezioni di codeste curve; inoltre, al variare di \(K\) nel fascio \((K)\), non può mai accadere che un gruppo \(G_n\) venga a coincidere col gruppo sezione di \(C\), perchè ciò importerebbe l’esistenza di punti doppi per l’involuzione ellittica segata dalle \(K\) su una componente ellittica della detta \(C\). Quindi la curva \(E\) ha zero intersezioni colla \(mC\), e di conseguenza è anch’essa una curva para-\(mi\)-canonica come la \(mC\) e \(mD\).

Ciò posto si rappresenti la \(F\) su una rigata ellittica \(n\)-pla, \(\Phi\), in modo che alle \(K\) rispondano le generatrici, ognuna delle quali rappresenti la \(g_n^1\) definita, come sopra, dalle sezioni di \(C\) e \(D\). La \(\Phi\) si potrà realizzare come rigata di un certo spazio \(S_r\), su cui si abbiano due direttrici minime rispondenti a \(C\) e \(D\), appartenenti a due spazi complementari \(\gamma_1\) e \(\gamma_2\), di \(s\) e \(r-s-1\) dimensioni; per esempio si può riferire \(\Phi\) ad un cono il cui vertice risponda a \(C\) e di cui \(D\) sia una sezione iperpiana. E sulla \(\Phi\) si avrà una curva ellittica \(E\), non avente punti comuni con \(\gamma_1\) e \(\gamma_2\). Ora le \(\infty^1\) omografie dello \(S_r\) che hanno come spazi di punti uniti \(\gamma_1\) e \(\gamma_2\), trasformeranno l’immagine di \(E\) in un sistema razionale \(\infty^1\) di curve ellittiche di grado 0, e perciò in un fascio lineare di curve ellittiche, senza punti base;
%CONT>>
% scan: s210L
% !PRINTED 394
\weberpage{394}
% head: CAPITOLO DECIMO
%<<CONT
% stamp: A dark right-edge scanning band, a lower-left blot, and scattered specks and stray strokes, including a stroke above the section heading.
a questo fascio della \(\Phi\) \(n\)-pla risponderà similmente su \(F\) un fascio lineare di curve ellittiche \(C\), senza punti base, di cui in tal guisa viene provata l’esistenza, \hfill c. d. d.

Non è escluso che le \(C\) sieno riducibili: in ogni caso le loro componenti variabili \textit{irriducibili} formeranno un \textit{fascio lineare, che torniamo a indicare con} \(|C|\).

Aggiungasi che il calcolo dell’invariante di Zeuthen-Segre per questo fascio lineare irriducibile \(|C|\), ci dà
\[
\delta-4=12p_a-p^{(1)}+9=-4,
\]
e quindi \(\delta=0\): vuol dire che le curve ellittiche \(C\), che variando in \(|C|\) non acquistano mai un punto doppio, hanno tutte lo stesso modulo, cioè sono birazionalmente identiche.

\edsection{11. Superficie di generi \(p_g=0\) e \(p_a=-1\) con un fascio ellittico di curve di genere \(\pi=1\).}{11. Superficie di generi p\_g=0 e p\_a=-1 con un fascio ellittico di curve di genere π=1.}{sec.X.11}

Il teorema che afferma l’esistenza di un fascio lineare di curve ellittiche sopra le superficie con \(p_g=0\) e \(p_a=-1\) possedenti un fascio ellittico di curve \(K\) di genere \(\pi>1\), si estende ora alle superficie \(F\) coi medesimi caratteri (\(p_g=0,\ p_a=-1\) e \(p^{(1)}=1\)) che posseggono un fascio ellittico di curve \(K\) di genere \(\pi=1\).

Giova partire dall’osservazione che, per essere l’invariante di Zeuthen-Segre
\[
\Delta-4=12p_a-p^{(1)}+9=-4,
\]
deve aversi
\[
\Delta=0;
\]
e quindi le \(K\) ellittiche, che variando in \((K)\) non acquistano mai un punto doppio, sono fra loro birazionalmente identiche.

Si trasformi la \(F\), su cui le \(K\) siano realizzate proiettivamente come curve d’ordine \(m\), in una superficie d’un conveniente iperspazio, che torniamo a chiamare \(F\), in guisa che le \(K\) diventino curve ellittiche normali \(K_m\), dello stesso ordine \(m\) appartenenti ciascuna ad uno spazio \(S_{m-1}\): ogni \(K\) della primitiva \(F\) sarà, in generale, proiezione della corrispondente \(K_m\). Ora due \(K_m\), curve ellittiche normali collo stesso modulo, saranno fra loro proiettive in un certo numero \(n\) di modi diversi: sarà precisamente \(n=2m^2\) se le dette curve sono di modulo generale, ed invece \(n=4m^2\) o risp. \(n=6m^2\), se esse sono armoniche o equianarmoniche. Così fra le due \(K_m\) suddette viene definito razionalmente l’insieme \(\Gamma_n\) di \(n\) omografie iperspaziali; tenendo fissa una delle due \(K_m\) e sopra di essa un punto \(P\), resta quindi definita una curva \(C\) che passa per \(P\) ed incontra l’altra \(K_m\)
%CONT>>
% scan: s210R
% !PRINTED 395
\weberpage{395}
% head: LA SUPERFICIE DI GENERE GEOMETRICO NULLO
%<<CONT
% stamp: Scattered marginal specks and small stray strokes.
variabile negli \(n\) punti corrispondenti a \(P\) rispetto alle \(n\) omografie del \(\Gamma_n\).

Sembra, a prima vista, che — le \(K_m\) essendo fra loro identiche — le suddette omografie si mantengano sempre distinte fra loro e facciano corrispondere ad un punto \(P\), fissato come sopra, \(n\) punti pure distinti; così la curva \(C\), contenendo un’involuzione ellittica \(\gamma_n^1\), priva di punti doppi, risulterebbe senz’altro ellittica e, al variare di \(P\), ci darebbe un fascio lineare \(|C|\) di curve ellittiche, direttrici del fascio ellittico.

Queste curve potranno essere riducibili; in tal caso torneremo a indicare con \(C\) le loro componenti irriducibili, che sono ancora ellittiche senza punti comuni, e torneremo pure a indicare con \(n\) il numero delle loro intersezioni colle \(K=K_m\). Avremo così, sopra la superficie d’irregolarità \(p_g-p_a=1\), possedente il fascio ellittico di curve ellittiche \((K)\), un altro fascio \(|C|\) di curve ellittiche direttrici delle \(K\), fascio necessariamente razionale\origfootnote{1}{s210R-1}{Cfr. cap. IX, \S~11.}.

Ma al discorso fatto per dimostrare che le \(C\) sono di genere uno, si può muovere un’obiezione che cade soltanto di fronte ad un esame approfondito, in cui si tenga conto che trattiamo di superficie di genere \(p_a=-1\). Infatti non è escluso a priori che le omografie intercedenti fra due \(K_m\) variabili vengano a degenerare, e così — supponendo per esempio \(n=6\) — che si abbia un fascio di sestiche \(K_6\), normali, in \(S_5\), entro a cui si trovino sestiche degeneri in cubiche piane doppie \(K_3^2\): allora tra una \(K_6\) generica e la \(K_3\) intercederanno, non più 36, ma 18 omografie degeneri, che portano fra le due curve altrettante corrispondenze \([1,\,2]\). % [sic] n=6 -> m=6, the order of the curves K_m is denoted m on the preceding page.

A priori non è nemmeno escluso che fra le \(K_m\) d’un fascio ellittico, fra loro generalmente proiettive, vi siano delle curve dotate di cuspide o spezzate, in più parti semplici o multiple; ma questo caso non può incontrarsi sulla nostra superficie di genere \(p_a=-1\), (e \(p^{(1)}=1\)) dove abbiamo visto che il numero delle \(K_m\) dotate di punto doppio deve essere
\[
\Delta=13+12p_a-p^{(1)}=0.
\]

Essendo \(p_a=-1\), si tratta dunque di rimuovere soltanto l’eccezione in cui il fascio \((K_m)\) contenga delle curve ellittiche multiple che, in ordine all’invariante di Zeuthen Segre, hanno la valenza zero nel computo di \(\Delta\) (cfr. \S~3).

A tal uopo si osservi che, in rapporto ai due fasci \((K)\) e \(|C|\), la nostra superficie \(F\) si può rappresentare sopra un cilindro cubico \(n\)-plo, \(\Phi\), dove le \(C\) rispondano alle sezioni piane normali \(c\) di \(\Phi\) e le \(K\) alle rette sue generatrici: si avrà su \(\Phi\) una curva di dirama%
%CONT>>
% scan: s211L
% !PRINTED 396
\unskip\weberpage{396}%
% head: CAPITOLO DECIMO
%<<CONT
% stamp: Dark right scan edge, lower-left blotch, scattered speckles, and stray diagonal and marginal marks.
zione formata di un certo numero di sezioni piane \(c\) rispondenti a curve multiple \(C=sC_s\) e (eventualmente) anche da un certo numero di generatrici \(k\), rispondenti a \(K\) degeneri in curve ellittiche multiple \(K=rK_r\). Se \(p_a=-1\), e quindi (per le \(K\)) \(\Delta=0\), il cilindro multiplo \(\Phi\) offre la rappresentazione d’una \(F\), che può supporsi priva di singolarità, su cui una \(C_s\) e una \(K_r\) s’incontrano in \(\frac{n}{rs}\) punti (punti semplici della superficie), che rispondono al punto \(A=ck\) di \(\Phi\): invero se qualcuno, \(A'\), fra i detti punti intersezione di \(C_s\) e \(K_r\), fosse doppio o multiplo per \(F\), la curva degenere limite di una \(K\) sarebbe, non già (come si è supposto) \(rK_r\), bensì \(rK_r\) aumentata dell’intorno di \(A'\), sicchè (\S\ 3) ne seguirebbe \(\Delta>0\). Ma ora si dimostra all’opposto che, se al punto \(A\) del cilindro \(n\)-plo \(\Phi\) risponda un punto della superficie \(F\), questo deve essere un punto singolare (doppio o multiplo) di \(F\), sicchè l’ipotesi di una \(K\) degenere in \(rK_r\) viene ridotta all’assurdo. Invero si consideri un piano \(\bar{\alpha}\), obliquo alle generatrici di \(\Phi\), che passi per \(A\), e un piano variabile \(\alpha\) che non passi per \(A\) ma tenda ad \(\bar{\alpha}\). La sezione di \(\Phi\) col piano \(\alpha\) è una curva multipla \(l\), \(n\)-pla, su cui si hanno due punti di diramazione \(A_1\) e \(A_2\) intersezioni di \(c\) e \(k\). Sopra la riemanniana della \(l\) multipla si descriva un cammino chiuso \(\nu\) che avvolga \(A_1\) e \(A_2\), il quale può ritenersi somma di due cappi elementari avvolgenti \(A_1\) e \(A_2\). Al primo cappio risponde una sostituzione sui rami equivalente al prodotto di % ?? uncertain: the first closed-path symbol -- \nu, a faint \gamma, or Latin v.
\[
(s-1)\frac{n}{s}
\]
trasposizioni, e al secondo una sostituzione equivalente a
\[
(r-1)\frac{n}{r}
\]
trasposizioni. Quindi al cammino chiuso \(\gamma\), che al limite --- per \(\alpha\to\bar{\alpha}\) --- diventa un cappio avvolgente \(A\), risponde una sostituzione formata di almeno \(n\) trasposizioni: vuol dire che la curva \(L\) omologa alla \(l\) non può avere punti semplici in corrispondenza ad \(A\), giacchè nel caso più sfavorevole questa curva darebbe luogo ad un ramo d’ordine \(n\), cui risponderebbe sulla riemanniana \(l\) ad \(n\) fogli, una sostituzione sui rami formata da \(n-1\) trasposizioni!

Concludiamo pertanto:

\textit{Le superficie \(F\) di generi \(p_g=0\) e \(p_a=-1\), possedenti un fascio ellittico di curve ellittiche \(K\), possiedono altresì un fascio lineare di curve ellittiche \(C\), direttrici delle \(K\).}

\textit{Nota.} -- Nel ragionamento svolto innanzi si è fatto uso del fecondo concetto delle singolarità delle curve esposto da \textsc{O. Chisini}
%CONT>>
% scan: s211R
% !PRINTED 397
\weberpage{397}
% head: LA SUPERFICIE DI GENERE GEOMETRICO NULLO
%<<CONT
% stamp: Scattered scan speckles and stray marginal marks, especially at right.
nella Memoria « Le singolarità di un ramo superlineare di curva piana definite mediante un prodotto di sostituzioni »\origfootnote{1}{s211R-1}{Atti del R. Istituto Veneto di Scienze, Lettere ed Arti, t. XXX, parte \(2^{\mathrm{a}}\), 1921. Questa memoria posteriore alla pubblicazione del secondo volume delle \textit{Lezioni} di \textsc{Enriques-Chisini}, reca alla teoria ivi svolta un ulteriore contributo importante.}.

Giova rilevare esplicitamente che, se si lascia cadere l’ipotesi \(p_a=-1\), si hanno di fatto superficie possedenti un fascio ellittico di curve ellittiche \(K\), cui appartengono curve spezzate: diciamo superficie \(F\) per cui \(p_a\geq0\) e \(p_g=p_a+1\), contenenti un fascio lineare di curve \(C\), direttrici di \((K)\), di genere \(\pi>1\). Se alle superficie con \(p_g=0\) e \(p_a=-1\) studiate nei precedenti paragrafi diamo il nome di \textit{ellittiche} (per un motivo che verrà esposto più avanti) queste nuove \(F\) che presentano con esse alcune notevoli analogie, potranno denominarsi \textit{paraellittiche}. La principale analogia è che sulle superficie paraellittiche, come sulle ellittiche con \(p_a=-1\), si hanno \textit{due fasci di curve birazionalmente identiche}: \textit{un fascio ellittico} di curve ellittiche \((K)\) e \textit{un fascio lineare} di curve, che qui non sono ellittiche ma di genere \(\pi>1\). Però, mentre le superficie ellittiche ammettono --- come vedremo --- un gruppo \(\infty^1\) di trasformazioni birazionali in sè stesse, questa proprietà non sussiste per le superficie paraellittiche.

Ci limitiamo qui a indicare due esempi di superficie paraellittiche, costruite a partire da un cilindro cubico multiplo \(f(xy)=0\):
\weq{1)}{\begin{aligned}
&u^2=(z-a)(z-b)(z-c)(z-d)\,\psi(xy)\\
&\qquad f(xy)=0\quad(p_a=3,\quad p_g=4);
\end{aligned}}
\(\psi=0\) rappresenta un cilindro quadrico parallelo al \(f\), che sega \(f=0\) in 6 generatrici di diramazione, da aggiungere alle 4 cubiche di diramazione
\[
z=a,\quad z=b,\quad z=c,\quad z=d;
\]

\weq{2)}{\begin{aligned}
&u^3=(z-a)(z-b)(z-c)\,\psi(xy)\\
&f(xy)=0\quad(p_a=0,\quad p_g=1);
\end{aligned}}
\(\psi=0\) designa un cilindro cubico che ha con \(f\) un contratto tripunto lungo tre generatrici, e lo tocca secondo un’altra generatrice \(k\) segandolo ulteriormente in una generatrice \(k'\): la curva di diramazione del cilindro triplo \(f\) consta dunque delle tre cubiche % [sic] contratto -> contatto, an extra r is printed in the contact term.
\[
z=a,\quad z=b,\quad z=c,
\]
e delle due generatrici
\[
k\qquad\text{e}\qquad k'.
\]
Infine, segnaliamo ai giovani ricercatori questo interessante oggetto
%CONT>>
% scan: s212L
% !PRINTED 398
\weberpage{398}
% head: CAPITOLO DECIMO
%<<CONT
% stamp: Dark right scan edge, lower-left blotch, and stray marginal marks and speckles.
di studio: \textit{determinare e classificare tutte le superficie paraellittiche}; e rimandiamo alle Note a pag. 123 e 134 delle « Lezioni » di \textsc{Enriques-Campedelli} nei Rendiconti del Seminario Matematico della R. Università di Roma (1934): dalle quali il lettore potrà raccogliere qualche osservazione e suggestione utile.

\edsection{12. Superficie ellittiche.}{12. Superficie ellittiche.}{sec.X.12}

Riassumiamo i risultati ottenuti nei precedenti paragrafi: \textit{Ogni superficie irregolare \(F\) di genere \(p_g=0\) che non appartenga alla famiglia delle rigate} (razionali o meno) \textit{ha il genere numerico \(p_a=-1\) e contiene un fascio ellittico \((K)\) di curve di genere \(\pi\geq1\), e un fascio lineare di curve ellittiche \(|C|\)}; quindi essa può rappresentarsi sopra un \textit{cilindro cubico \(\Phi\) multiplo secondo il numero \(n=CK\), con una curva di diramazione costituita di sezioni piane normali a \(\Phi\).}

Viceversa un \textit{cilindro cubico \(n\)-plo \(\Phi\)}, con curva di diramazione costituita di sezioni piane normali (fissate in guisa da soddisfare alle condizioni d’esistenza di una superficie \(F\) irriducibile), \textit{non} rappresenta \textit{sempre una superficie di genere \(p_g=0\)}, perchè le immagini delle sezioni piane di \(\Phi\) possono essere curve ellittiche riducibili le cui componenti costituiscono un fascio \((C)\), non lineare, \textit{di genere \(\varrho\geq1\)}.

Qui appare che le nostre superficie, con \(p_g=0\) e \(p_a=-1\), non riferibili a rigate, rientrano in una \textit{famiglia più generale di superficie \(F\)}, con \(p_a=-1\), che posseggono un \textit{fascio ellittico di curve \(K\)} di genere \(\pi\geq1\) e un \textit{fascio di genere \(\varrho\geq0\) di curve ellittiche \(C\)}, per le quali \(F\) \textit{il genere geometrico sarà \(p_g=\varrho\)}.

Invero l’esistenza dei due fasci \((C)\) e \((K)\) che caratterizzano le dette \(F\) più generali, essendo \(p_a=-1\), porta che l’irregolarità di \(F\) sia
\[
q\geq\varrho+1
\]
e quindi
\[
p_g\geq\varrho;
\]
ma, d’altra parte, le curve canoniche di \(F\), formate di curve \(C\), debbono segare ogni \(K\) in gruppi della \(g_{2\pi-2}^{\pi-1}\) composti con l’involuzione di genere \(\varrho\) segata da \((C)\), priva di punti doppi\origfootnote{1}{s212L-1}{Teorema di Painlevé-Castelnuovo. Cfr. \textsc{Enriques-Chisini}, \textit{Lezioni}. Libro V, cap. I, \S\ 9 (vol. III, pag. 72).}; e perciò debbono appartenere alla serie canonica \(g_{2\varrho-2}^{\varrho-1}\) entro il fascio \((C)\). Resulterà dunque:
\[
p_g\leq\varrho
\]
e pertanto
\[
p_g=\varrho,
\]
\hfill c. d. d.

% scan: s212R
% !PRINTED 399
\weberpage{399}
% head: LA SUPERFICIE DI GENERE GEOMETRICO NULLO
% stamp: Scattered scan speckles and stray marginal marks.
Le superficie \(F\) più generali, caratterizzate --- come sopra è detto --- dall’avere il genere numerico \(p_a=-1\) e dal possedere due fasci \((K)\) e \((C)\) di curve irriducibili, cioè

1) un fascio ellittico di curve \(K\) di un certo genere \(\pi\geq0\) (per cui dall’essere \(p_a=-1\) segue \(\Delta=0\));

2) e un fascio di genere \(p_g\) di curve ellittiche \(C\), secantisi in \(n=CK\geq1\) punti, si diranno \textit{superficie ellittiche}: e, più precisamente superficie ellittiche \textit{proprie} quelle, per \(\pi\geq1\), che non sono riferibili a rigate, e \textit{improprie} quelle, per \(\pi=0\), che si riducono a rigate di genere 1.

Il detto nome « ellittiche » deriva da una rappresentazione parametrica delle superficie nominate che è in rapporto col gruppo \(\infty^1\) delle loro trasformazioni in sè di cui discorriamo più avanti, e che è stata indicata da \textsc{P. Painlevé}\origfootnote{1}{s212R-1}{\textsc{P. Painlevé}, \textit{Leçons sur la théorie analytique des équations différentielles}. (Paris, Hermann, 1897).}.

Per riguardo a questa rappresentazione parametrica il numero \(n=CK\) riceve il nome di \textit{determinante} della superficie.

La rappresentazione parametrica sopra accennata si traduce nella seguente proprietà geometrica: \textit{Ogni superficie ellittica \(F\) ammette un gruppo ellittico \(\infty^1\) di trasformazioni birazionali in se stessa, che ha come traiettorie le curve del fascio \((C)\).}

Dimostriamo il teorema, facendo vedere che esiste una trasformazione birazionale di \(F\) --- che lascia invariate tutte le curve \(C\) --- in cui si corrispondono due punti \(P_o\) e \(P'_o\) arbitrariamente scelti sopra una di tali curve, \(C_o\). % ?? uncertain: the subscripts in P_o, P'_o, and C_o -- italic o or zero.
% EDITORIAL D0214 RISOLTO: The small round subscripts are slanted italic letter o, rather than the upright numeric zero used elsewhere on the page.

Invero si consideri un’altra curva \(C\) dello stesso fascio \((C)\) e su di essa si scelga un punto \(P\). Se le \(C\), fra loro birazionalmente identiche, sono curve ellittiche di modulo generale, vi sono --- com’è noto --- \textit{due} trasformazioni di \(C_o\) in \(C\), che mutano \(P_o\) in \(P\)\origfootnote{2}{s212R-2}{Cfr. p. es. \textsc{Enriques-Chisini}, \textit{Lezioni}. Libro V, cap. III, \S\ 27 (vol. III, pag. 262).}; e tra queste trasformazioni una sola\origfootnote{3}{s212R-3}{Cfr. \textsc{Enriques-Chisini}, l. c., cap. IV, \S\ 39 (vol. III, pag. 446).} (razionalmente determinata) fa corrispondere i gruppi sezioni delle curve \(K\), che formano su ciascuna una involuzione ellittica \(\gamma_n^1\). Pertanto, essendo data su \(C_o\) la trasformazione di prima specie \((P_oP'_o)\), viene anche determinata razionalmente con essa \textit{una} trasformazione di prima specie della \(C\), in cui ad un punto \(P\) risponderà un punto \(P'\), per modo che la \((PP')\) si ridurrà alla \((P_oP'_o)\) quando la \(C\) si riduca a \(C_o\). % ?? uncertain: the repeated subscripts o -- italic o or zero.
% EDITORIAL D0215 RISOLTO: Magnified repeated occurrences, especially (P_o P' _o), show the same slanted italic o used above.

La dimostrazione si estende senza difficoltà ai casi in cui le \(C\) sieno curve ellittiche armoniche o equianarmoniche.

% scan: s213L
% !PRINTED 400
\weberpage{400}
% head: CAPITOLO DECIMO
% stamp: Dark right scan edge and scattered marginal marks and speckles.
Giova anche osservare che la proprietà stabilita per le superficie ellittiche è caratteristica:

\textit{Ogni superficie \(F\) che ammetta un gruppo ellittico \(\infty^1\) \(\Gamma\) di trasformazioni birazionali in se stessa, è una superficie ellittica, di genere numerico \(p_a=-1\), propria o impropria, possedente un fascio ellittico di curve \(K\) di genere \(\pi(\geq0)\) e un fascio di genere \(p_g\) di curve ellittiche \(C\).}

Infatti la \(F\) dovrà possedere in primo luogo un fascio \((C)\), di un certo genere \(\varrho\geq0\), formato dalle traiettorie del gruppo \(\Gamma\), che saranno curve ellittiche identiche allo stesso gruppo \(\Gamma\). Ma, in secondo luogo, si costruirà su \(F\) un fascio ellittico \((K)\), colla costruzione seguente: si consideri sopra \(F\) una curva (irriducibile) \(L\) e la serie ellittica \(\infty^1\) \(\{L\}\) costituita dalle sue trasformate mercè le trasformazioni del gruppo \(\Gamma\); le curve \(L\) della detta serie saranno certo disequivalenti, dovendo segare su una \(K\) gruppi di punti disequivalenti, che si corrispondono in una trasformazione di prima specie; però fra i gruppi di curve \(L\) uscenti da un punto \(P\) variabile sopra \(F\), ve ne saranno \(\infty^1\) fra loro equivalenti: ciò significa che vi è su \(F\) un fascio ellittico di curve su cui le \(L\) della serie \(\{L\}\) segano gruppi equivalenti. (Cfr. Cap. IX, \S\ 11).

Dopo avere costruiti in tal guisa su \(F\) i due fasci \((C)\) e \((K)\), si determinano agevolmente i caratteri della superficie. Anzitutto il calcolo dell’invariante di Zeuthen-Segre per \((K)\) ci dà
\[
\Delta=13-12p_a-p^{(1)}\geq0 % [sic] 13-12p_a-p^{(1)} -> 13+12p_a-p^{(1)}, the following lower bound on p_a requires a plus sign.
\]
e quindi
\[
p_a\geq-1\qquad\qquad(p^{(1)}\geq1).
\]

D’altra parte la presenza di un fascio ellittico \((K)\) e d’un fascio \((C)\) di genere \(\varrho\), porta che (esista su \(F\) una serie di \(\infty^{\varrho+1}\) curve disequivalenti, e quindi che sia)
\[
\begin{aligned}
&p_g-p_a\geq\varrho+1\\
&p_g\geq\varrho.
\end{aligned}
\]

Infine una considerazione già fatta innanzi vale a stabilire che le curve canoniche di \(F\) --- che per \(\varrho>0\) sono formate con \(2\varrho-2\) curve \(C\) --- debbono dare entro il fascio \((C)\) una serie contenuta nella serie canonica \(g_{2\varrho-2}^{\varrho-1}\); dal confronto segue
\[
p_g=\varrho,\qquad p_a=-1.
\]
\hfill c. d. d.

\textit{Osservazione.} -- Da ciò che abbiamo detto si deduce che sopra una superficie ellittica \(F\) di determinante \(n>1\) si ha una \textit{involuzione, \(I_n\), formata dai gruppi di \(n\) punti comuni alle \(C\) e \(K\), che viene}
%CONT>>
% scan: s213R
% !PRINTED 401
\weberpage{401}
% head: LA SUPERFICIE DI GENERE GEOMETRICO NULLO
%<<CONT
% stamp: Scattered marginal marks and scan speckles.
\textit{generata da un gruppo abeliano \(\Gamma_n\) di trasformazioni birazionali} (permutabili), precisamente dalle trasformazioni del gruppo \(\Gamma\) della superficie che lasciano ferme le \(C\) e le \(K\). Se i gruppi di punti della \(I_n\) si assumono come elementi (punti) di una nuova superficie \(f\), si ottiene una superficie ellittica di determinante 1, sopra la quale la \(F\) viene rappresentata in modo multiplo, secondo \(n\); quando sia \(p_g=0\), questa superficie ellittica impropria si riduce ad una rigata ellittica, e se si vuole ad un cilindro cubico \(f(xy)=0\).

Qui giova rilevare che le \textit{superficie ellittiche \(f\) di determinante 1} si definiscono in generale come rappresentative della \textit{varietà \(\infty^2\) delle coppie di punti appartenenti ad una curva ellittica \(C\) e ad una curva \(K\) di genere \(\pi\) \((\geq0)\).}

A partire da queste superficie, mediante l’introduzione di irrazionalità algebriche atte a risolvere il gruppo abeliano \(\Gamma_n\), si potranno costruire le superficie ellittiche di determinante \(n>1\). A tal uopo si può dire, a priori, che occorreranno \textit{un solo radicale} ovvero \textit{due radicali non sovrapposti}, secondo che il detto gruppo \(\Gamma_n\) sia \textit{ciclico ovvero abeliano a base due}. Perchè il \(\Gamma_n\), che è un gruppo finito di trasformazioni in sè di una curva ellittica \(C\), non può avere una base \(>2\)\origfootnote{1}{s213R-1}{Cfr. \textsc{Enriques-Chisini}, \textit{Lezioni}. Libro V, cap. III, \S\ 28 (vol. III, pag. 271) e cap. IV, \S\ 39 (vol. III, pag. 446-451).}. Nel paragrafo seguente, limitandoci in particolare all’ipotesi \(p_g=0\), e poi ancora, più precisamente, nel \S\ 16, spieghiamo il senso di questa costruzione, che vale a stabilire l’\textit{esistenza} delle superficie ellittiche, e le condizioni che vi si riferiscono.

\edsection{13. Costruzione delle superficie ellittiche di genere \(p_g=0\).}{13. Costruzione delle superficie ellittiche di genere p\_g=0.}{sec.X.13}

Per costruire le superficie ellittiche (proprie) \(F\) di genere \(p_g=0\), aventi un qualsiasi determinante \(n>1\), prendiamo le mosse dal cilindro cubico \(f(xy)=0\), che realizza --- come si è detto --- la superficie ellittica (impropria) di genere \(p_g=0\) e di determinante 1.

Mediante i due fasci di curve irriducibili \(|C|\) e \((K)\), che la caratterizzano, la superficie \(F\) si lascerà rappresentare sopra un cilindro \(n\)-plo \(f\), con curva di diramazione composta di un certo numero di sezioni piane normali \(c\), per modo che le \(K\) abbiano come immagini le generatrici \(k\) del cilindro, e le \(C\) (traiettorie del gruppo \(\infty^1\) di trasformazioni della \(F\)) rispondano alle sezioni piane normali \(c\) dello stesso cilindro. La curva di diramazione del cilindro \(n\)-plo sarà costituita da un certo numero di curve \(c\), cui risponderanno curve multiple del tipo \(C=sC_s\).

% EDITORIAL: segnatura originale: \textsc{Enriques F.} - \textit{Superficie algebriche.}\hfill 26
\printersignature{\textsc{Enriques F.} - \textit{Superficie algebriche.}\quad 26}

% scan: s214L
% !PRINTED 402
\weberpage{402}
% head: CAPITOLO DECIMO
% stamp: Dark right scan edge, lower-left blotch, marginal marks, and a curved stroke beneath the first displayed formula.
Quindi la rappresentazione della superficie ellittica \(F\) sul cilindro \(n\)-plo \(f\) deve soddisfare alle seguenti condizioni:

1) che la curva \(K\) definita da una \(k\) \(n\)-pla sia irriducibile e contenga una involuzione \(g_n\) del tipo ciclico ovvero abeliano (propriamente detto) a base 2;

2) che le curve \(C\) corrispondenti alle sezioni piane normali \(c\) di \(f\), curve \(n\)-ple senza punti di diramazione, siano generalmente irriducibili.

Osserviamo invero che, se questa condizione 2) non sia soddisfatta, la superficie \(F\), rappresentata dal cilindro \(n\)-plo, non sarà più una superficie (ellittica) di determinante \(n\), ma avrà un determinante \(<n\), e riuscirà, in generale, di genere \(p_g>0\). Così accade, per esempio, per la superficie
\[
\left\{
\begin{aligned}
u^2&=(z-a)(z-b)(z-e)(z-d)\,\psi^2(xy)\\
f(xy)&=0
\end{aligned}
\right.
\]
che è una superficie irriducibile \(F\) rappresentata sul cilindro doppio, con curva di diramazione composta dalle sezioni piane \(z=a\), \(z=b\), \(z=d\), \(z=e\).

Infatti ad una sezione normale generica \(c\):
\[
z=c\qquad(c=\text{cost.})
\]
risponde sulla detta \(F\) una coppia di curve ellittiche distinte
\[
u=\pm\sqrt{(c-a)(c-b)(c-d)(c-e)}\,\psi,
\]
sicchè le curve \(C\) irriducibili che compongono le immagini delle \(c\) formano ora un fascio (non più lineare, ma ellittico) di curve unisecanti le \(K\); e pertanto il genere geometrico di \(F\) vale:
\[
p_g=1.
\]

Qui conviene richiamare i criteri perchè una curva ellittica \(n\)-pla, senza punti di diramazione, sia irriducibile\origfootnote{1}{s214L-1}{Cfr. \textsc{Enriques-Chisini}, \textit{Lezioni}, Libro V, cap. IV, \S\ 39 (vol. III, pag. 446).}. Vi sono due casi possibili: il caso \textit{ciclico} e il caso \textit{abeliano} (base 2). Nel primo caso la curva \(n\)-pla si costruisce, a partire p. es. da una cubica \(f(xy)=0\), estraendo su di essa un radicale \(n\)-mo che porta sopra un polinomio \(\varphi(xy)=0\),
\[
\left\{
\begin{aligned}
u^n&=\varphi(xy)\\
f(xy)&=0;
\end{aligned}
\right.
\]
affinchè la curva multipla non abbia punti di diramazione, bisogna che i \textit{punti critici}
\[
f=0,\qquad\varphi=0,
\]
%CONT>>
% scan: s214R
% !PRINTED 403
\weberpage{403}
% head: LA SUPERFICIE DI GENERE GEOMETRICO NULLO
%<<CONT
% stamp: Scattered scan speckles and marginal marks.
siano \textit{apparenti}, e perciò che la \(\varphi\) presa p. es. d’ordine \(n\), abbia tre contatti \(n\)-punti con \(f\). Quindi l’\textit{irriducibilità} della detta curva multipla dipende da ciò che «la serie \(g_3^2\) a cui appartiene la terna dei punti critici apparenti sia diversa dalla serie segata su \(f\) dalle rette del piano ».

Nel caso abeliano (in cui \(n=pq\)) la costruzione della curva \(n\)-pla richiede l’estrazione di due radicali d’indice \(p\) e \(q\) non sovrapposti, portanti rispettivamente su due polinomi \(\varphi(xy)\) e \(\psi(xy)\), che è lecito supporre d’ordine \(p\) e \(q\): affinchè la curva multipla non abbia punti di diramazione, bisogna che le curve \(\varphi\) e \(\psi\) abbiano con \(f\) tre punti di contatto rispettivamente d’ordine \(p-1\) e \(q-1\), i quali formano così due terne di punti critici apparenti; l’irriducibilità della detta curva multipla importa che le dette terne non appartengono alla \(g_3^2\) segata su \(f\) dalle rette del piano, ed inoltre che esse soddisfino ad una certa condizione di dissomiglianza.

Sviluppiamo la costruzione indicata nel \textit{caso ciclico}, che s’incontra in particolare quando il determinante \(n\) sia un numero primo. Abbiamo dunque che sopra la superficie \(F\) l’involuzione \(I_n\) costituita dai gruppi \(CK\), viene generata da un gruppo ciclico di trasformazioni \(\Gamma_n\), d’ordine \(n\).

A partire dal cilindro cubico \(n\)-plo
\[
f(xy)=0,
\]
fissiamo le sezioni piane normali di diramazione
\[
z=a_1,\quad z=a_2,\quad \ldotp\ldotp\ldotp\ldotp\quad z=a_t;
\]
gli \(n\) punti del gruppo \(G_n\) che risponde ad un punto generico di \(f\) potranno darsi mediante i valori di un radicale
\[
u=\sqrt[n]{(z-a_1)^{r_1}(z-a_2)^{r_2}\,\ldotp\ldotp\ldotp\ldotp\,(z-a_t)^{r_t}\cdot\varphi(xy)},
\]
soddisfacendo alle condizioni seguenti:

1) Per gli esponenti \(r_i<n\) si avrà
\[
r_1+r_2+\ldotp\ldotp\ldotp\ldotp+r_t\equiv0\qquad(\mathrm{mod.}\ n),
\]
altrimenti nascerebbe sulle generatrici del cilindro \(f\) un punto di diramazione all’infinito.

2) Se \(h_i\) è il massimo comune divisore di \(n\) ed \(r_i\), sicchè
\[
n=h_i s_i,\qquad r_i=h_i\varrho_i. % ?? uncertain: punctuation after h_i\varrho_i -- a full stop, a faint comma, or a scan speck.
\]
nella espressione di \(u\) appare il fattore
\[
\sqrt[n]{(z-a_i)^{r_i}}=\sqrt[s_i]{\varepsilon_i(z-a_i)^{\varrho_i}},
\]
dove \(\varepsilon_i\) è una radice \(h_i\)-ma dell’unità.

% scan: s215L
% !PRINTED 404
\weberpage{404}
% head: CAPITOLO DECIMO
% stamp: Dark right scan edge, scattered speckles, and diagonal marginal marks.
Così sulla generatrice \(n\)-pla del cilindro \(f\), il punto
\[
z=a_i
\]
è un punto di diramazione a cui corrisponde sulla \(K\) omologa un gruppo \(G_n\) (della \(g_n\) segata dalle \(C\)) costituito da \(h_i\) gruppi di \(s_i\) punti coincidenti. In altre parole, il \textit{punto \(z=a_i\)} è per la \(k\) multipla un punto di \textit{diramazione d’ordine \(s_i\)}, con
\[
s_i=\frac{n}{h_i},
\]
e si ha quindi
\[
\sum\frac{n}{s}(s-1)=2n+2\pi-2
\]
(sommatoria estesa ai valori di \(s=s_1,\ s_2\,\ldotp\ldotp\ldotp\ldotp\,s_t\)) designando \(\pi\) \((\geq1)\) il genere delle \(K\).

3) Ciascuno dei numeri \(s_i(i=1,2,\ \ldotp\ldotp\ldotp\ldotp\ t)\) è un divisore del minimo comune multiplo dei rimanenti.

Questa proprietà risulta in generale dal carattere abeliano del \(\Gamma_n\). Infatti, sulla retta multipla \(k\), non essendovi diramazioni all’infinito, il prodotto delle sostituzioni cicliche sui rami della funzione algebrica \(u(xyz)\), in ordine ai punti di diramazione \(a_i\), equivale all’identità:
\[
S_t\,\ldotp\ldotp\ldotp\ldotp\,S_2S_1=1
\]
sicchè, per esempio,
\[
S_1^{-1}=S_t\,\ldotp\ldotp\ldotp\ldotp\,S_3S_2,
\]
ossia
\[
S_1^{s_1-1}=S_t\,\ldotp\ldotp\ldotp\ldotp\,S_3S_2.
\]
Ma la sostituzione \(S_1^{s_1-1}\) è ciclica di periodo \(s_1\) mentre il prodotto \(S_t\,\ldotp\ldotp\ldotp\ldotp\,S_3S_2\) è ciclico con periodo divisore del minimo comune multiplo di \(s_2,s_3,\ldotp\ldotp\ldotp\ldotp,s_t\)\origfootnote{1}{s215L-1}{Cfr. p. es. \textsc{L. Bianchi}, \textit{Lezioni sulla teoria dei gruppi di sostituzioni e delle equazioni algebriche secondo Galois} (Pisa, 1900), cap. III, \S\ 30.}, avendosi
\[
S_t^{s_t}\,\ldotp\ldotp\ldotp\ldotp\,S_3^{s_3}S_2^{s_2}
=(S_t\,\ldotp\ldotp\ldotp\ldotp\,S_3S_2)^{s_2s_3\,\ldotp\ldotp\ldotp\ldotp\,s_t}=1.
\]

4) Poichè il cilindro \(n\)-plo \(f\) non possiede generatrici di diramazione, il polinomio \(\varphi(xy)\) che figura sotto il radicale, nell’espressione di \(u\), deve rappresentare un cilindro avente con \(f\) un contatto \(n\)-punto lungo ogni generatrice comune. E perchè le curve \(C\) rappresentate dalle \(z=\text{cost.}\) siano irriducibili, bisognerà che il gruppo delle generatrici critiche apparenti, non sia equivalente ad un multiplo delle terne di generatrici sezioni dei piani
\[
\alpha x+\beta y+\gamma=0.
\]

% scan: s215R
% !PRINTED 405
\weberpage{405}
% head: LA SUPERFICIE DI GENERE GEOMETRICO NULLO
% stamp: Scattered scan speckles and marginal marks.
Potremo anzi prendere \(\varphi\) di ordine \(n\), in guisa che \(\varphi=0\) tocchi \(f=0\) lungo tre generatrici (non giacenti in un piano) con contatto \(n\)-punto. Invero si riconosce facilmente che, facendo variare \(\varphi\) con continuità, la superficie \(F\) si conserva birazionalmente identica a se stessa, conservandosi le curve \(C\) e \(K\) rappresentate rispettivamente dalle sezioni piane \(z=\text{cost.}\) e dalle generatrici di \(f\). Quindi, se si assume al posto del cilindro \(\varphi\) un cilindro \(\psi\) d’ordine \(nh\), con \(h>1\), potremo sostituire a \(\psi(xy)\) un polinomio spezzato in un fattore semplice \(\varphi\) d’ordine \(n\) e in un fattore \(n\)-plo \(\varphi_1^n\) con \(\varphi_1\) d’ordine \(h-1\):
\[
\psi=\varphi_1^n\varphi;
\]
così risulterà
\[
u=\varphi_1\sqrt[n]{\varphi(xy)\cdot(z-a_1)^{r_1}\,\ldotp\ldotp\ldotp\ldotp\,(z-a_t)^{r_t}},
\]
da cui --- cambiando \(u\) in \(u\varphi_1\) --- si ottiene appunto l’espressione che volevasi stabilire.

In conclusione: \textit{come tipo delle superficie ellittiche di determinante \(n\), appartenenti al caso ciclico, si ha la superficie}
\[
\left\{
\begin{aligned}
u&=\sqrt[n]{(z-a)^{r_1}\,\ldotp\ldotp\ldotp\ldotp\,(z-a_t)^{r_t}\cdot\varphi(xy)}\\ % [sic] (z-a)^{r_1} -> (z-a_1)^{r_1}, the preceding formula indexes the first branch value by 1.
f(xy)&=0
\end{aligned}
\right.
\]
\textit{dove gli esponenti \(r_i\) soddisfano alle condizioni dette innanzi e \(\varphi(xy)=0\) è un cilindro d’ordine \(n\) che tocca \(f\), con contatto \(n\)-punto, secondo tre generatrici non appartenenti ad un piano.}

La costruzione delle superficie \(F\), a partire da un cilindro cubico \(n\)-plo \(f(xy)=0\), si ripete con poche modificazioni nel caso in cui il gruppo sia \textit{abeliano} (a base 2). Senonchè, in luogo di un solo radicale, si avranno ora nell’\textit{espressione di \(u\) due radicali non sovrapposti}. Non c’indugeremo su questa costruzione, che più avanti (\S\ 16) sviluppiamo, riferendoci alla famiglia delle superficie \(F\) per cui le \(K\) come le \(C\) siano curve ellittiche. Qui basti avvertire che ancora, come nel caso ciclico, si avrà l’eguaglianza
\[
\sum\frac{n(s-1)}{s}=2n-2\pi-2, % [sic] 2n-2\pi-2 -> 2n+2\pi-2, compare the corresponding equality on printed p. 404.
\]
e sempre \textit{ciascuno dei numeri \(s_i\) dovrà dividere il minimo comune multiplo dei rimanenti.}

\edsection{14. Superficie coi plurigeneri nulli: caratterizzazione delle rigate.}{14. Superficie coi plurigeneri nulli: caratterizzazione delle rigate.}{sec.X.14}

Abbiamo dimostrato che «le superficie irregolari di genere geometrico \(p_g=0\), non appartenenti alla famiglia delle rigate, sono superficie ellittiche, di genere numerico \(p_a=-1\)». Vogliamo ora determinare i plurigeneri di queste superficie ellittiche e verificare che codesti
%CONT>>
% scan: s216L
% !PRINTED 406
\weberpage{406}
% head: CAPITOLO DECIMO
%<<CONT
% stamp: Scattered specks, stray strokes, and a dark vertical scanning artefact at the right edge.
caratteri non possono essere tutti nulli, come accade per le rigate (razionali o no). In tal guisa, ricordando le condizioni di razionalità di Castelnuovo, date dall’annullamento del genere numerico e del bigenere (Cap. VI, § 4): $p_a=P_2=0$, la famiglia delle rigate verrà definita dall’annullamento dei plurigeneri, e più precisamente dal seguente teorema: \textit{Condizione necessaria e sufficiente perchè una superficie sia riferibile ad una rigata è che si annullino per essa il quadrigenere e il sestigenere:}
\[
P_4=P_6=0.
\]

Consideriamo una superficie (propriamente) ellittica $F$, di genere $p_g=0$ e di determinante $n\geq2$, rappresentata sul cilindro cubico $n$-plo $f(xy)=0$, nel modo innanzi definito, in rapporto ai due fasci caratteristici di essa $(K)$ e $|C|$. Si può determinare la curva canonica virtuale di $F$ sommando alla trasformata della curva canonica del cilindro $f$ la corrispondente della curva di diramazione data su questo cilindro multiplo (cfr. Cap. V, § 7).

Ora la curva canonica virtuale di $f$ si ottiene sottraendo il sistema lineare $2c$; costituito da due sezioni piane normali al cilindro, dal suo sistema aggiunto, che consta delle tre rette sezioni col piano all’infinito; e poichè queste rette sono immagini di curve eccezionali, sopra la $F$, si avrà soltanto da sommare alla trasformata della curva di diramazione la curva virtuale $-2C$ che risponde a $-2c$. Pertanto, designando con $z=a_i$ le componenti della curva di diramazione a ciascuna delle quali risponde una curva multipla $C=s_iC_i$, la curva canonica pura di $F$ verrà data da % ?? uncertain: punctuation after 2c -- semicolon or comma.
% EDITORIAL D0217 RISOLTO: Magnified punctuation has both an upper point and a lower curved descender: semicolon.
\[
D-2C=\Sigma(s-1)C_s-2C.
\]

Essendo le curve $C$ omologhe alle $c$, irriducibili, si vede tosto che la sottrazione indicata dalla formula precedente non è possibile, e così si conferma che l’anzidetta irriducibilità porta che il genere di $F$ sia $p_g=0$.

Infatti il sistema lineare $|D|$, definito dalla $D=\Sigma(s-1)C_s$, ha la dimensione $x=0$ e perciò non può contenere entro di sè il sistema $|C|$ e tanto meno $|2C|$. Per giustificare questa asserzione si avverta che la curva $D$ essendo formata di componenti ellittiche di grado $0$ dovrà definire a priori un sistema lineare composto colle curve del fascio $|C|$ e con componenti parti aliquote di tali curve:
\[
|D|=|\Sigma(s-1)C_s|=|xC+\Sigma h_sC_s|.
\]

Di qui si ricava
\[
|xC|=|\Sigma k_sC_s|,
\]
con
\[
0\leq k_s<s-1,
\]
%CONT>>
% scan: s216R
% !PRINTED 407
\weberpage{407}
% head: LA SUPERFICIE DI GENERE GEOMETRICO NULLO
%<<CONT
% stamp: Scattered scanning specks and stray marginal marks.
ed è chiaro che una tale relazione non può sussistere per $x>0$ (essendo le $C$ generiche irriducibili) perchè si contraddirebbe al \textit{principio di spezzamento}, in quanto una curva variabile di $(C)$ (irriducibile e perciò connessa) verrebbe a degenerare in parti $C_s$ non connesse fra loro.

Ora dalla formula scritta innanzi passiamo all’espressione delle curve pluricanoniche e quindi al calcolo dei plurigeneri di $F$.

In generale le curve $m$-canoniche saranno date da
\[
|m\Sigma(s-1)C_s-2mC|=|m(t-2)C-m\Sigma C_s|,
\]
dove la sommatoria si estende ai diversi valori di $s$ che corrispondono alle $t$ curve multiple del fascio $|C|$:
\[
s=s_1,\ s_2;\ \ldotp\ldotp\ldotp\ldotp,\ s_t. % ?? uncertain: punctuation after s_2 -- semicolon or comma.
% EDITORIAL D0218 RISOLTO: A clear upper point and lower curved tail form a semicolon.
\]

Se si designa con $x=P_m-1$ la dimensione del sistema $m$-canonico completo, questo sistema si comporrà di $x$ curve $C$ (irriducibili) variabili e di componenti fisse $C_s$, ciascuna da contare un certo numero di volte
\[
h<s.
\]

Avremo dunque
\[
|m(t-2)C-m\Sigma C_s|=|xC+\Sigma hC_s|,
\]
con $0\leq h<s$, la sommatoria essendo estesa ai valori
\[
s=s_1,\ s_2,\ \ldotp\ldotp\ldotp\ldotp,\ s_t
\]
e ai corrispondenti valori di
\[
h=h_1,\ h_2,\ \ldotp\ldotp\ldotp\ldotp,\ h_t.
\]

Ora dalla precedente relazione, essendo $sC_s=C$, si ricava
\[
|(m(t-2)-x)C|=\left|\Sigma\frac{m+h}{s}C\right|:
\]
vuol dire che le frazioni $\dfrac{m+h}{s}$ equivalgono a numeri interi:
\[
\frac{m+h}{s}=\varrho,
\]
ciascuno di questi interi ($\varrho=\varrho_1,\ \varrho_2,\ \ldotp\ldotp\ldotp\ldotp,\ \varrho_t$) esprimendo il valore per eccesso della frazione $\dfrac{m}{s}$, ovvero il massimo intero
\[
\left[\frac{m+s-1}{s}\right]
\]
contenuto in $\dfrac{m+s-1}{s}$.

% scan: s217L
% !PRINTED 408
\weberpage{408}
% head: CAPITOLO DECIMO
% stamp: Scattered specks, stray diagonal strokes, and a dark right-edge scanning artefact.
In conclusione avremo
\[
x=m(t-2)-\Sigma\left[\frac{m+s-1}{s}\right],
\]
e quindi \textit{il plurigenere d’ordine $m(\geq2)$ è dato dalla formula.} % ?? uncertain: punctuation after "formula" -- a full stop or a scanning speck.
\[
P_m=1+m(t-2)-\Sigma\left[\frac{m+s-1}{s}\right],
\]
dove si dovrà prendere $P_m=0$ quando la sua espressione risulti negativa.

Facciamo qui $m=2$, tenendo conto che — per $s=s_1,s_2,\ldotp\ldotp\ldotp\ldotp,s_t$ — si ha sempre
\[
s\geq2;
\]
avremo
\[
\left[\frac{s+1}{s}\right]=1,\qquad
\Sigma\left[\frac{s+1}{s}\right]=t,
\]
e quindi \textit{il bigenere vale:}
\[
P_2=t-3.
\]

Ora, per $t>3$, si ha certo $P_2>0$ e perciò anche $P_4>0$, $P_6>0$ ecc. Se invece $P_2=0$, essendo esclusa l’ipotesi $t=2$ che porta ad una $F$ riferibile a rigata, si avrà
\[
t=3.
\]

Ma, designando con $n$ il determinante di $F$, si ha l’eguaglianza
\[
\Sigma n(s-1)=2n+2\pi-2, % [sic] \Sigma n(s-1) -> \Sigma\frac{n}{s}(s-1), the corresponding equality on page 410 contains the factor n/s.
\]
da cui (per essere $s_i\geq2$, $\pi\geq1$)
\[
\Sigma\left(1-\frac{1}{s}\right)\geq1, % [sic] \geq1 -> \geq2, the equality with 2n gives the stronger bound used below for t=3.
\]
ossia
\[
\begin{aligned}
&t-\Sigma\frac{1}{s}\geq1\\ % [sic] \geq1 -> \geq2, the apparently intended bound is the one used in the following t=3 calculation.
&\Sigma\frac{1}{s}\leq t-1. % [sic] t-1 -> t-2, the next displayed calculation gives a bound of 1 when t=3.
\end{aligned}
\]

Per $t=3$ si avrà dunque
\[
\Sigma\frac{1}{s}=\frac{1}{s_1}+\frac{1}{s_2}+\frac{1}{s_3}\leq1.
\]

Poniamo che i numeri (interi e positivi) $s_1$, $s_2$, $s_3$ siano in ordine non decrescente:
\[
s_3\geq s_2\geq s_1;
\]
%CONT>>
% scan: s217R
% !PRINTED 409
\weberpage{409}
% head: LA SUPERFICIE DI GENERE GEOMETRICO NULLO
%<<CONT
% stamp: Scattered scanning specks and stray marginal marks.
la precedente diseguaglianza ci dà
\[
s_1=2,\qquad s_2=3,\qquad s_3\geq6
\]
oppure
\[
s_1=2\qquad s_3\geq s_2\geq4,
\]
ovvero
\[
s_1\geq3,\qquad s_3\geq s_2\geq3, % [sic] s_1\geq3 -> s_1=3, the following fourth case separately treats s_1\geq4.
\]
o infine
\[
s_3\geq s_2\geq s_1\geq4.
\]

\edFlush
Nel primo e nel secondo caso:
\[
\left[\frac{s_1+2}{s_1}\right]=2,\qquad
\left[\frac{s_2+2}{s_2}\right]=\left[\frac{s_3+2}{s_3}\right]=1
\]
e quindi
\[
P_3=0,
\]
mentre nel terzo e quarto caso:
\[
\left[\frac{s_1+2}{s_1}\right]=\left[\frac{s_2+2}{s_2}\right]=\left[\frac{s_3+2}{s_3}\right]=1,
\]
e
\[
P_3=1.
\]

Invece, calcolando il quadrigenere
\[
P_4=5-\Sigma\left[\frac{s+3}{s}\right], % ?? uncertain: delimiters around (s+3)/s -- square brackets with faint horizontal strokes or vertical bars.
% EDITORIAL D0220 RISOLTO: The bracket corners are faint but visible, particularly the inward stroke at the top left and bottom right.
\]
avremo:

nel primo e terzo caso
\[
P_4=0,
\]

nel secondo caso
\[
P_4=1
\]

e nel quarto
\[
P_4=2.
\]

E per il sestigenere
\[
P_6=7-\Sigma\left[\frac{s+5}{s}\right]
\]
si troverà:

nel primo e terzo caso
\[
P_6=1,
\]

nel secondo caso
\[
P_6=0
\]

e nel quarto
\[
P_6=2.
\]

% scan: s218L
% !PRINTED 410
\weberpage{410}
% head: CAPITOLO DECIMO
% stamp: Scattered specks, stray diagonal strokes, and a dark vertical scanning artefact at the right edge.
Per il nostro scopo è sufficiente rilevare che quando $P_2=P_4=0$ (primo e terzo caso della discussione precedente) risulta
\[
P_6\geq1.
\]

In conclusione per una superficie ellittica, non riferibile a rigata, uno almeno dei due plurigeneri $P_4$ o $P_6$ sarà diverso da zero:
\[
P_4+P_6\geq1,\qquad\qquad\text{c. d. d.}
\]

\textit{Osservazione. -- Le condizioni perchè una superficie sia riferibile ad una rigata si possono anche esprimere annullando il genere d’ordine $12$:}
\[
P_{12}=0.
\]

Infatti per le superficie ellittiche, aventi il $P_4\geq1$ o il $P_6\geq1$, risulta certo
\[
P_{12}\geq1.
\]

Giova anche rilevare che si avrà
\[
P_{12}=1
\]
per le \textit{superficie ellittiche} (con curva canonica virtuale d’ordine zero) possedenti un \textit{fascio ellittico di curve $K$ ellittiche}, ed invece
\[
P_{12}>1
\]
per le \textit{superficie ellittiche} possedenti un \textit{fascio ellittico di curve $K$ di genere $\pi>1$}, le cui curve canoniche virtuali segano le $K$ in $2\pi-2$ punti.

Invero la differenza fra le due specie di superficie $F$ consiste in ciò che:

nel caso delle $K$ di genere $\pi=1$ si ha l’\textit{eguaglianza}
\[
2n=\Sigma\frac{s-1}{s}n
\]
cioè
\weq{1)}{t-2=\Sigma\frac{1}{s},}
mentre nel caso delle $K$ di genere $\pi>1$ sussiste la
\[
2n+2\pi-2=\Sigma\frac{s-1}{s}n
\]
che porta la \textit{diseguaglianza}
\weq{2)}{t-2>\Sigma\frac{1}{s}.}

Che nel caso 1) si abbia
\[
P_{12}=1
\]
e non $P_{12}>1$, risulta subito, non soltanto dalla formula del $P_{12}$,
%CONT>>
% scan: s218R
% !PRINTED 411
\weberpage{411}
% head: LA SUPERFICIE DI GENERE GEOMETRICO NULLO
%<<CONT
% stamp: Scattered scanning specks and a dark artefact near the upper right edge.
sì anche dall’osservare che sulla superficie $F$ possedente due fasci di curve ellittiche non possono aversi curve pluricanoniche (pure) d’ordine $>0$.

Dimostriamo ora che, se sussiste la diseguaglianza 2), si ha necessariamente
\[
P_{12}>1
\]
e non $P_{12}=1$. Questa dimostrazione non è una semplice deduzione dalla formula del $P_{12}$, ma esige altresì che si tenga conto di una condizione d’esistenza delle superficie ellittiche (attinente al carattere abeliano del gruppo $\Gamma_n$) già avvertita nel § 13, cioè che « ciascuno dei numeri interi $s_i$ deve dividere il minimo comune multiplo dei rimanenti »\origfootnote{1}{s218R-1}{Così appunto la diseguaglianza $P_{12}>1$ viene stabilita nel § 15 delle Lezioni di \textsc{Enriques-Campedelli} (Seminario mat. di Roma, 1934). Ma nella Nota originale di \textsc{Enriques}; \textit{Sulla classificazione delle superficie algebriche ecc.} (Rendic. Lincei, febbraio 1914) la cosa non è chiarita, per modo che è naturale ritenere trattarsi di una semplice deduzione dalla formula generale che dà il $P_{12}$. Ora questa deduzione sarebbe erronea, come ha osservato \textsc{H. Geppert} nel suo Bericht \textit{Die Klassification der algebraischen Flaechen} (Jahresbericht der Mathematische Vereinigung, 1931); invero, facendo, per esempio $s=2$, $s_2=s_3=5$; si trova $P_{12}=1$. % [sic] s=2 -> s_1=2, the other two multiplicities in this example are indexed 2 and 3.

Perciò il \textsc{Geppert} propone di considerare il $P_{24}$ al posto del $P_{12}$. Ma in realtà per le superficie ellittiche con curve pluricanoniche d’ordine $>0$ si ha già $P_{12}>1$. Il caso di eccezione segnalato non risponde ad un’effettiva superficie ellittica, perchè il $2$ non è divisore del $5$!}.

Scriviamo la formula
\[
P_{12}=12t-23-\Sigma\frac{s+11}{s}, % [sic] \Sigma\frac{s+11}{s} -> \Sigma\left[\frac{s+11}{s}\right], the integer-part brackets used in the general formula and subsequent calculations are absent here.
\]
e discutiamo i casi che possono presentarsi per $t>4$, per $t=4$ e per $t=3$.

Per $t>4$ cioè $t\geq5$, si ha
\[
\left[\frac{s_i+11}{s_i}\right]\leq6
\]
e perciò è certo
\[
\begin{aligned}
&\Sigma\left[\frac{s_i+11}{s_i}\right]\leq24,\\ % [sic] 24 -> 6t, the preceding bound is 6 for each of t summands.
&P_{12}\geq12t-47\geq12(t-4)+1>1. % [sic] P_{12}\geq12t-47\geq12(t-4)+1>1 -> P_{12}\geq6t-23=6(t-4)+1>1, using the preceding per-term bound with t summands.
\end{aligned}
\]

Per $t=4$ il valore minimo del $P_{12}$:
\[
P_{12}=12(t-4)+1=1,
\]
%CONT>>
% scan: s219L
% !PRINTED 412
\weberpage{412}
% head: CAPITOLO DECIMO
%<<CONT
% stamp: Scattered specks, stray diagonal strokes, and a dark vertical scanning artefact at the right edge.
si raggiunge soltanto quando tutti e quattro i numeri $s_i$ siano eguali a $2$:
\[
s_1=s_2=s_3=s_4=2,
\]
sicchè sussista l’eguaglianza 1), perchè solo in questo caso sarà per ciascuno di essi
\[
\left[\frac{s_i+11}{s_i}\right]=6;
\]
se invece uno di codesti numeri supera $2$, come è richiesto dalla diseguaglianza 2), per esempio se
\[
s_4\geq3,
\]
si ha
\[
\left[\frac{s_4+11}{s_4}\right]\leq4
\]
e quindi
\[
P_{12}\geq3.
\]

Sia infine $t=3$, e si scrivano tutte le soluzioni intere della diseguaglianza
\[
\begin{aligned}
&\frac{1}{s_1}+\frac{1}{s_2}+\frac{1}{s_3}<1\\
&(s_1\leq s_2\leq s_3),
\end{aligned}
\]
soddisfacenti alla condizione che ciascuno degli $s_i$ divida il comune multiplo dei rimanenti. Si avrà
\[
\left\{
\begin{aligned}
s_1&=2,&\qquad s_2&=4,&\qquad s_3&=8,\\
s_1&=2,& s_2&=5,& s_3&=10,\\
s_1&=2,& s_2&\geq6,& s_3&\geq s_2,\\
s_1&=3,& s_2&\geq3,& s_3&>3,\\
s_1&>4,& s_2&\geq s_1,& s_3&\geq s_2. % [sic] s_1>4 -> s_1\geq4, the displayed enumeration otherwise omits s_1=4.
\end{aligned}
\right.
\]

In tutti questi casi si ha
\[
\begin{aligned}
&\Sigma\left[\frac{s+11}{s}\right]<12\\
&P_{12}>1,
\end{aligned}
\]
\hfill c. d. d.

\edsection{15. Nota storica.}{15. Nota storica.}{sec.X.15}

La classificazione delle superficie irregolari di genere $p_g=0$ ($p_a<0$), la dimostrazione che quelle che non siano riferibili a rigate sono superficie ellittiche, contenenti un fascio ellittico di curve
%CONT>>
% scan: s219R
% !PRINTED 413
\weberpage{413}
% head: LA SUPERFICIE DI GENERE GEOMETRICO NULLO
%<<CONT
% stamp: Scattered scanning specks and stray marginal marks.
ammettenti un gruppo $\infty^1$ di trasformazioni in se stesse, infine $K$ di genere $\pi\geq1$ e un secondo fascio lineare di curve ellittiche $C$, la caratterizzazione delle rigate coll’annullamento del quadrigenere e del sestigenere, costituiscono l’oggetto della memoria di \textsc{F. Enriques} « Sulle superficie algebriche di genere geometrico zero », pubblicata nei Rendiconti del Circolo matematico di Palermo, del 1905. Tuttavia (a prescindere da qualche imprecisione minore) la dimostrazione di questi risultati lasciava una lacuna, per quel che si riferisce al lemma II dei §§ 6, 7, 8; giacchè l’esistenza di una seconda curva ellittica essenzialmente distinta dalla parabicanonica $C$, veniva inferita semplicemente da ciò che le curve del sistema $\{K''\}$, secondo paraggiunto ad una $K$, passanti per un gruppo bicanonico di questa, debbono formare un sistema $\infty^1$ d’indice $>1$: infatti non è escluso che le curve spezzate di codesto sistema che contengono come parte la detta $K$, si riducano ad una sola da contarsi più volte. % ?? uncertain: primes on K in the system in braces -- two primes or one.
% EDITORIAL D0221 RISOLTO: The magnification shows two adjacent prime strokes.

Per colmare questa lacuna lo stesso autore ha escogitato il ragionamento che si appoggia sulla costruzione di curve secanti le $K$ in gruppi della serie semiquadricanonica o, in generale, parte aliquota di un multiplo della serie canonica. Questo procedimento è stato esposto da lui in una Nota dell’Accademia dei Lincei, del 1930\origfootnote{1}{s219R-1}{\textsc{Enriques}, \textit{Sopra le superficie algebriche trasformabili in rigate.} (Luglio, 1930).} e quindi da \textsc{Enriques-Campedelli}, nelle « Lezioni sulla classificazione delle superficie algebriche particolarmente di genere zero ». (Seminario matematico di Roma, 1934).

Ma anche nella nuova dimostrazione rimane un punto debole d’ordine delicato, perchè non è lecito ritenere a priori, come si fà provvisoriamente coll’ipotesi II dei §§ 6, 7 e 8, che, essendo data o aggiunta al campo di razionalità di una curva $K$, una serie lineare, si possa determinarne razionalmente \textit{un} gruppo. Fortunatamente ci è riuscito di superare la difficoltà riconoscendo che l’ipotesi II è verificata per qualche multiplo della serie suddetta, in guisa da giustificare il procedimento costruttivo di curve parapluricanoniche essenzialmente distinte.

Dobbiamo ancora aggiungere un rilievo concernente il caso in cui le nostre superficie posseggano un fascio ellittico di curve $K$ ellittiche (birazionalmente identiche); per essere rigorosa la costruzione delle traiettorie $C$, o meglio la dimostrazione che anch’esse sono ellittiche, si deve escludere il dubbio che talune fra codeste $K$ possano degenerare in curve multiple (dimostrazione che, per $p_a=-1$, non possono aversi superficie paraellittiche): questa
%CONT>>
% scan: s220L
% !PRINTED 414
\weberpage{414}
% head: CAPITOLO DECIMO
%<<CONT
% stamp: Scattered specks, stray diagonal strokes, and a dark vertical scanning artefact at the right edge.
precisazione — che si basa sull’analisi della singolarità nascente dall’incrocio di due curve di diramazione (\textsc{Chisini}) trovasi nelle « Lezioni » di \textsc{Enriques-Campedelli} del 1934 (§ 12).

Infine per quel che concerne la costruzione effettiva delle superficie ellittiche, la Memoria originale di \textsc{Enriques} del 1905 contemplava soltanto il caso ciclico. I primi esempi relativi al caso abeliano (a base due) sono stati messi in luce da \textsc{G. Bagnera} e \textsc{M. De Franchis} nel loro studio « Sopra le superficie algebriche che danno le coordinate del punto generico esprimibili con funzioni meromorfe quadruplamente periodiche di due parametri »\origfootnote{1}{s220L-1}{R. Accademia dei Lincei, $1907_1$.}.

Infine, la costruzione generale di tutte le superficie ellittiche (che si riconduce alla risoluzione di equazioni abeliane con gruppo a base due) è stata sviluppata da \textsc{O. Chisini} nelle Note su « Le superficie ellittiche il cui determinante è un numero composto » del 1921\origfootnote{2}{s220L-2}{R. Accademia dei Lincei, $1921_2$.}.

\edsection{16. Superficie con curve pluricanoniche d’ordine zero: tipi con\\curve ellittiche $K$ di modulo generale.}{16. Superficie con curve pluricanoniche d’ordine zero: tipi con curve ellittiche K di modulo generale.}{sec.X.16}

Nel § 14, e in particolare nell’Osservazione che lo chiude, abbiamo dimostrato che:

\textit{Le superficie propriamente ellittiche di generi $p_g=0$ e $p_a=-1$, si distinguono dalle rigate per avere il 12-genere.}
\[
P_{12}>0,
\]
\textit{anzichè}
\[
P_{12}=0.
\]

\textit{Più precisamente le superficie ellittiche con $p_g=0$ e $p_a=-1$ danno luogo a due famiglie:}

1) \textit{superficie con curva canonica virtuale d’ordine zero, caratterizzate dal possesso di un fascio ellittico di curve $K$ di genere $\pi=1$;}

2) \textit{superficie con curva canonica virtuale d’ordine $2\pi-2$, che contengono un fascio ellittico di curve $K$ di genere $\pi>1$.}

\textit{Le prime} (su cui ogni sistema lineare puro di genere $p$ ha il grado $n=2p-2$) \textit{hanno il}
\[
P_{12}=1,
\]
%CONT>>
% scan: s220R
% !PRINTED 415
\weberpage{415}
% head: LA SUPERFICIE DI GENERE GEOMETRICO NULLO
%<<CONT
% stamp: Scattered scanning specks and stray marginal marks.
\textit{mentre le seconde} (per cui $n<2p-2$) \textit{posseggono sistemi lineari di dimensione $r>0$ di curve pluricanonice} (di genere $p^{(1)}=1$) \textit{ed hanno il} % [sic] pluricanonice -> pluricanoniche, the h is absent in print.
\[
P_{12}\geq2.
\]

Ora vogliamo stringere più da vicino la classificazione delle superficie 1) per cui $P_{12}=1$, mostrando che esse danno luogo a pochi tipi ben determinati, di cui possiamo scrivere le equazioni.

Riferiamoci alle superficie ellittiche $F$ della famiglia 1); conviene distinguere i casi in cui:

\textit{Le curve ellittiche $K$} (fra loro birazionalmente identiche) \textit{hanno modulo generale}, ovvero \textit{le $K$ sono armoniche}, o \textit{sono equianarmoniche}.

Incominciamo qui dal \textit{caso generale}: proveremo che la superficie $F$ ha il \textit{determinante} $n=2$ (caso ciclico) oppure $n=4$ (caso abeliano).

A tal uopo occorre considerare l’involuzione $g_n^1$ segata sopra una curva ellittica $K$ dalle traiettorie $C$. Questa involuzione è generata da un gruppo abeliano $\Gamma_n$ e, per essere razionale (anzichè di genere $1$), dovrà contenere qualche trasformazione involutoria $I$ di seconda specie ($g_2^1$). Ora, se il gruppo $\Gamma_n$ non si riduce al $\Gamma_2$ costituito da una $I$ e dall’identità (siccome i prodotti di due $I$ sono trasformazioni $\pi$ di prima specie) il $\Gamma_n$ dovrà contenere tante trasformazioni $\pi$ quante $I$, diciamo $m$ delle une e delle altre: $n=2m$. Più precisamente le $\pi$ permutabili con una $I$ saranno involutorie, e quindi il gruppo abeliano $\Gamma_n$ generato dalle dette $\pi$ ed $I$ sarà il $\Gamma_4$ diedrico formato da due involuzioni $I$, una involuzione $\pi$ e l’identità.

Dunque la nostra superficie $F$ con un fascio di \textit{curve ellittiche $K$ di modulo generale sarà di determinante $n=2$ o $n=4$, e la relativa $g_n^1$ definita sopra una $K$ apparterrà ad uno dei due tipi seguenti:}

\textit{a)} $n=2,\quad g_n^1=g_2^1,$

\textit{b)} $n=4,\quad g_n^1=g_4^1$ \textit{abeliana e diedrica, generata da due trasformazioni di seconda specie permutabili, il cui prodotto è una $\gamma_2^1$ di prima serie.} % [sic] prima serie -> prima specie, the preceding discussion calls these transformations "di prima specie".

Nel primo caso si può assumere come \textit{tipo della superficie $F$ di determinante $n=2$, quello che è definito dalle equazioni:}
\weq{\mathrm{I}_a}{
\left\{
\begin{aligned}
u&=\sqrt{(z-a_1)(z-a_2)(z-a_3)(z-a_4)\varphi(xy)}\\ % ?? uncertain: extent of the radical bar -- all of \varphi(xy) under the radical or the terminal factor outside it.
% EDITORIAL D0222 RISOLTO: The overbar continues across the phi and the parenthesized variables; the end is faint but does not terminate before the factor.
f(xy)&=0,
\end{aligned}
\right.
}
dove $f$ è un cilindro cubico e $\varphi$ un cilindro quadrico che ha con esso un contatto bipunto lungo tre generatrici, non giacenti in un piano.

Nel secondo caso si tratta di costruire, nel piano $(u\,z)$, la curva
%CONT>>
% scan: s221L
% !PRINTED 416
\weberpage{416}
% head: CAPITOLO DECIMO
%<<CONT
% stamp: Scattered specks, a diagonal non-text mark beneath "quadrupla", and a dark vertical line at the right edge.
ellittica \(K\) contenente una \(g_4^1\) diedrica, a partire dalla corrispondente rappresentazione sopra la retta quadrupla \(k\) (\(u=0\)).

Questa costruzione si può spiegare in rapporto alla teoria generale delle equazioni abeliane\origfootnote{1}{s221L-1}{Cfr. p. es. L. \textsc{Bianchi}, op. cit., cap. IV, \S~76.}, ovvero riferendosi più specialmente alle nozioni che concernono le curve ellittiche\origfootnote{2}{s221L-2}{Cfr. \textsc{Enriques-Chisini}, \textit{Lezioni.} Libro V, cap. IV, \S~39 (vol. III).}, come qui faremo.

Osserviamo che la nostra \(g_4^1\) diedrica (su \(K\)) possiede 8 punti doppi che si distribuiscono in due quaterne: una quaterna di punti doppi della prima \(g_2^1\) costituita da due coppie di punti coniugati della seconda
\[
A_1\,A_2 \qquad \text{e} \qquad B_1\,B_2,
\]
ed una quaterna di punti doppi della seconda \(g_2^1\) costituita da due coppie della prima
\[
D_1\,D_2 \qquad \text{e} \qquad E_1\,E_2.
\]

Quindi la retta quadrupla \(k\) avrà 4 punti di diramazione (doppi):
\[
a\,b\,d\,e,
\]
e il campo di razionalità della \(K\) si potrà definire ponendo
\[
u=\sqrt{\varphi\cdot(z-a)(x-b)}\pm\sqrt{\psi\cdot(z-d)(z-e)}, % [sic] (x-b) -> (z-b), the other factors and the subsequent equations use z.
\]
(con \(\varphi\) e \(\psi\) costanti rispetto a \(z\) ed \(u\)) in guisa che ai due punti di diramazione \(a\) e \(b\) rispondano le coppie di punti doppi \(A_1\,A_2\) e \(B_1\,B_2\), e ai punti \(d\) ed \(e\) rispondano le coppie di punti doppi \(D_1\,D_2\) e \(E_1\,E_2\).

Che, effettivamente, l’equazione sopra scritta rappresenti la \(K\) ellittica rispondente alla nostra \(k\) quadrupla, si prova da ciò che la curva rappresentata possiede una \(g_4^1\) con 8 punti doppi, e d’altronde --- liberando l’equazione stessa dai radicali --- si ha un’equazione di \(4^{\mathrm{o}}\) grado in \(z\) ed \(u\), che è di \(2^{\mathrm{o}}\) grado in \(z\), e definisce una quartica avente un punto doppio nel punto all’infinito dell’asse \(z\) ed un altro punto doppio infinitamente vicino. Aggiungasi che le due trasformazioni di seconda specie, cioè le \(g_2^1\) dotate ciascuna di 4 punti doppi, da cui è formata la \(g_4^1\) diedrica, si lasciano rappresentare mediante equazioni
\[
\begin{aligned}
u'&=\frac{\varphi\cdot(z-a)(z-b)-\psi\cdot(z-e)(z-d)}{u}\\
u''&=\frac{\psi\cdot(z-d)(z-e)-\varphi\cdot(z-a)(z-b)}{u},
\end{aligned}
\]
%CONT>>
% scan: s221R
% !PRINTED 417
\weberpage{417}
% head: LA SUPERFICIE DI GENERE GEOMETRICO NULLO
%<<CONT
% stamp: Scattered scan specks and a thin vertical line at the right edge.
come risulta dall’identità:
\[
\begin{aligned}
&\frac{\varphi\cdot(z-a)(z-b)-\psi\cdot(z-d)(z-e)}
{\sqrt{\varphi\cdot(z-a)(z-b)}+\sqrt{\psi\cdot(z-d)(z-e)}}\\
&=\sqrt{\varphi\cdot(z-a)(z-b)}-\sqrt{\psi\cdot(z-d)(z-e)}.
\end{aligned}
\]

Pertanto le equazioni della \textit{superficie ellittica \(F\) di determinante \(n=4\), con \(K\) di modulo generale, si potranno ridurre al tipo:}
\weq{\mathrm{I}_b}{
\left\{
\begin{aligned}
u&=\sqrt{(z-a)(z-b)\cdot\varphi(xy)}+\sqrt{(z-d)(z-e)\cdot\psi(xy)}\\
f(xy)&=0,
\end{aligned}
\right.
}
dove \(f\) è un cilindro cubico e \(\varphi\) e \(\psi\) sono due cilindri quadrici che toccano \(f\) secondo due terne di generatrici, non giacenti in un piano e non equivalenti.

\edsection{17. Segue: Caso armonico.}{17. Segue: Caso armonico.}{sec.X.17}

Nel caso armonico oltre alla \(g_2^1\) e alla \(g_4^1\) del caso generale, si potrà avere sopra la curva ellittica \(K\) una \(g_n^1\) ciclica o abeliana, i cui gruppi di \(n\) punti siano trasformati in sè da una trasformazione singolare del quart’ordine. Infatti le curve ellittiche armoniche posseggono, oltre le trasformazioni ordinarie di prima e seconda specie, \(\pi\) ed \(I\), due serie continue \(\infty^1\) di trasformazioni singolari del \(4^{\mathrm{o}}\) ordine, \(\omega\) e \(\tau\), le quattro specie di trasformazioni essendo legate fra loro dalle relazioni seguenti\origfootnote{1}{s221R-1}{Cfr. p. es. \textsc{Enriques-Chisini}, \textit{Lezioni}, Libro V, cap. III, \S~27 (vol. III, pag. 260).}: % ?? uncertain: punctuation after "Lezioni" in the footnote -- comma or period.
% EDITORIAL D0223 RISOLTO: The punctuation has a curved descender below the baseline: comma.
\[
\begin{aligned}
\pi\omega&=\omega_1 & \omega_1\omega&=I & I\omega&=\tau & \tau\omega&=\pi\\
\pi I&=I_1 & I_1I&=\pi & \omega I&=\tau & \tau I&=\omega\\
\pi\tau&=\tau_1 & \tau_1\tau&=I & I\tau&=\omega & \omega\tau&=\pi\\
\pi_1\pi&=\pi_2 & I\pi&=I_1 & \omega\pi&=\omega_1 & \tau\pi&=\tau_1.
\end{aligned}
\]

Si possono classificare i varii tipi di \(g_n^1\) ciclica o abeliana appartenenti ad una curva ellittica armonica \(K\), partendo dall’osservazione che, se il relativo gruppo \(\Gamma_n\) non è generato da sole trasformazioni ordinarie (corrispondenti a \(K\) di modulo generale), esso dovrà contenere un sottogruppo invariante \(\Gamma_4\) generato da una trasformazione singolare \(\omega\) o \(\tau\). Ora questo \(\Gamma_4\) dà luogo ad una \(g_4^1\) con due punti quadrupli \(A\) e \(B\) e a una coppia di punti doppi \(D_1\) e \(D_2\) (formanti un gruppo della \(g_4^1\)), la quale conduce a rappresentare la \(K\) sopra una retta quadrupla \(k\) (\(u=0\)) con due punti di diramazione quadrupli: \(z=a\) e \(z=b\), e un punto di diramazione doppio: \(z=d\).
% EDITORIAL: segnatura originale: \vadjust{\hbox to \linewidth{\textsc{Enriques F.} - \textit{Superficie algebriche.}\hfil 27}}
\printersignature{\textsc{Enriques F.} - \textit{Superficie algebriche.}\quad 27}
%CONT>>
% scan: s222L
% !PRINTED 418
\weberpage{418}
% head: CAPITOLO DECIMO
%<<CONT
% stamp: Scattered marks, a diagonal non-text stroke beside K, and a dark vertical line at the right edge.
Se la \(g_n^1\) non è esaurita dalla detta \(g_4^1\) (cioè se \(n>4\)), il gruppo \(\Gamma_n'\) dovrà contenere le trasformazioni che rispondono alle proiettività della retta quadrupla trasformanti in sè il punto \(d\) e permutanti i due punti \(A\) e \(B\): e queste trasformazioni generano un \(\Gamma_8\) abeliano, formato da 4 trasformazioni singolari del \(4^{\mathrm{o}}\) ordine, due \(\omega\) e due \(\tau\), da 2 involuzioni di seconda specie \(I\), e da 2 trasformazioni \(\pi\) di prima specie: una involuzione \(\gamma_2^1\) e l’identità. % ?? uncertain: prime on the first \Gamma_n -- \Gamma_n' or \Gamma_n with a stray mark.
% EDITORIAL D0224 RISOLTO: The raised mark is a distinct narrow oblique prime in the expected position, not a diffuse speck.

In conclusione la \(g_n^1\) non generata da sole trasformazioni ordinarie, sulla \textit{curva ellittica armonica \(K\)}, apparterrà ad uno dei due tipi seguenti:

\textit{a}) \(n=4\), \(g_4^1\) ciclica generata da una trasformazione singolare del \(4^{\mathrm{o}}\) ordine: \(\omega\) o \(\tau\); e

\textit{b}) \(n=8\), \(g_8^1\) generata da un gruppo abeliano \(\Gamma_8\) contenente 4 trasformazioni singolari cicliche del \(4^{\mathrm{o}}\) ordine (due \(\omega\) e due \(\tau\)) e 4 trasformazioni ordinarie: due \(I\) e due \(\pi\), (cioè una \(\gamma_2^1\) e l’identità).

Nel \textit{primo caso} la \(K\) ellittica rappresentata sulla \(k\) quadrupla, coi punti di diramazione \(a\), \(b\) e \(d\), si realizza nel piano (\(z\,u\)) scrivendo, per esempio:
\[
u=\sqrt[4]{\varphi\cdot(z-a)(z-b)(z-d)^2}
\]
con \(\varphi\) costante (rispetto a \(z\) ed \(u\)).

L’equazione scritta risponde ad una certa scelta delle sostituzioni sui rami della funzione algebrica \(u(z)\), in rapporto alle curve di diramazione \(z=a,\ b,\ d\); se invece si assumono le sostituzioni inverse avremo il \textit{tipo simile} ma birazionalmente distinto:
\[
u=\sqrt[4]{\varphi\cdot(z-a)^3(z-b)^3(z-d)^2}.
\]

Quindi \textit{le superficie ellittiche armoniche di determinante \(n=4\), caso ciclico, possono ridursi al tipo rappresentato dalle equazioni:}
\weq{\mathrm{II}_a}{
\left\{
\begin{aligned}
u&=\sqrt[4]{(z-a)(z-b)(z-d)^2\cdot\varphi(xy)}\\
f(xy)&=0,
\end{aligned}
\right.
}
o ad un \textit{tipo simile}: dove \(f\) è un cilindro cubico e \(\varphi\) un cilindro del quarto ordine toccante \(f\) con contatto quadripunto lungo tre generatrici, non giacenti in un piano. I tipi simili che diversificano, come si è detto, per la scelta delle sostituzioni sui rami di \(u\), pur essendo birazionalmente distinti, hanno però i medesimi caratteri.

Nel caso propriamente abeliano possiamo ancora realizzare nel piano (\(uz\)) un tipo della curva ellittica \(K\), che contenga una \(g_8^1\) con due gruppi costituiti ciascuno di due punti quadrupli e un gruppo costituito di quattro punti doppi: per ciò si tratta di costruire una
%CONT>>
% scan: s222R
% !PRINTED 419
\weberpage{419}
% head: LA SUPERFICIE DI GENERE GEOMETRICO NULLO
%<<CONT
% stamp: Scattered scan specks and several vertical lines at the right edge.
retta 8-pla con tre punti di diramazione
\[
z=a,\qquad z=b,\qquad z=d
\]
in modo che ai primi due risponda una sostituzione sui rami ciclica del \(4^{\mathrm{o}}\) ordine, mentre al terzo risponda una semplice sostituzione di periodo 2.

Questa costruzione si effettua secondo i principii generali della teoria delle equazioni abeliane\origfootnote{1}{s222R-1}{Cfr. p. es., L. \textsc{Bianchi}, \textit{Lezioni sulla teoria dei gruppi di sostituzioni e delle equazioni algebriche secondo Galois.} (Pisa, 1900) cap. III, \S~30.}, e d’altronde la formula che così si ottiene si giustifica a posteriori\origfootnote{2}{s222R-2}{Cfr. O. \textsc{Chisini}, \textit{Le superficie ellittiche il cui determinante è un numero composto.} Rendic. Lincei, s. V, vol. XXX, \(1921_2\).} riconoscendo le sostituzioni sui rami che corrispondono ai punti di diramazione, di cui si è detto innanzi.

Diciamo che la curva \(K\) può ridursi al tipo rappresentato dalle equazioni
\[
u=\sqrt{(z-b)(z-d)\cdot\varphi}+\sqrt[4]{(z-a)(z-b)^3\cdot\psi}
\]
o ad un \textit{tipo simile} (che ne differisce per una diversa scelta delle sostituzioni anzidette): dove si designano con \(\varphi\) e \(\psi\) delle costanti (rispetto ad \(u\) e \(z\)).

Per giustificare l’asserto, si distinguano gli 8 rami di \(u\) con due indici:
\[
\begin{aligned}
&u_{11}\qquad u_{12}\qquad u_{13}\qquad u_{14}\\
&u_{21}\qquad u_{22}\qquad u_{23}\qquad u_{24},
\end{aligned}
\]
facendo corrispondere il primo indice ai due valori \(+\sqrt{}\) e \(-\sqrt{}\), del radicale quadratico che entra nell’espressione di \(u\), e il secondo ai quattro valori \(+\sqrt[4]{},\ i\sqrt[4]{},\ -\sqrt[4]{},\ -i\sqrt[4]{}\) del radicale del \(4^{\mathrm{o}}\) ordine.

Al punto di diramazione \(z=a\) risponderà la sostituzione sui rami
\[
S_a=(u_{11}\,u_{12}\,u_{13}\,u_{14})(u_{21}\,u_{22}\,u_{23}\,u_{24}).
\]

E al punto \(z=b\) risponderà egualmente una sostituzione ciclica del \(4^{\mathrm{o}}\) ordine, cioè la
\[
S_b=(u_{11}\,u_{24}\,u_{13}\,u_{22})(u_{12}\,u_{21}\,u_{14}\,u_{23}).
\]

Invece al punto \(z=d\) risponderà la sostituzione a periodo 2:
\[
S_d=(u_{11}\,u_{21})(u_{12}\,u_{22})(u_{13}\,u_{23})(u_{14}\,u_{24}).
\]

Si ottengono così tre sostituzioni che, per moltiplicazione, generano un gruppo abeliano d’ordine 8, avente la struttura del nostro \(\Gamma_8\) sopra indicato.

% scan: s223L
% !PRINTED 420
\weberpage{420}
% head: CAPITOLO DECIMO
% stamp: Scattered scan specks and a dark irregular vertical line at the right edge.
Si aggiunga che le sostituzioni anzidette soddisfano alla condizione d’esistenza\origfootnote{1}{s223L-1}{Cfr. p. es., \textsc{Enriques-Chisini}, \textit{Lezioni.} Libro V, cap. IV, \S~38 (vol. III, pag. 429).} della funzione algebrica \(u(z)\), (non diramata nel punto \(z=\infty\)), cioè che
\[
S_d\,S_b\,S_a=1;
\]
appunto in vista di tale condizione si è dovuto attribuire l’esponente 3 (anzichè 1) al termine \(z-b\) che figura sotto il radicale del \(4^{\mathrm{o}}\) ordine.

In tal guisa la funzione \(u(z)\) realizza la richiesta curva ellittica \(K\) con \(g_8^1\) abeliana, e ne porge anzi il \textit{tipo}, sebbene possa aversi anche un altro \textit{tipo simile}, irriducibile per trasformazioni birazionali, quale è dato da
\[
u=\sqrt{(z-a)(z-b)^3\cdot\varphi}+\sqrt[4]{(z-a)(z-d)\cdot\psi}\,\text{\origfootmark{2}{s223L-2}}:
\]\origfoottext{2}{s223L-2}{Dal punto di vista proiettivo la curva \(K\) rappresentativa della funzione \(u(z)\) avrà: in corrispondenza al punto di diramazione \(z=a\) una tangente (parallela all’asse \(u\)) con due contatti quadripunti, e in corrispondenza a \(z=b\) un tacnodo in cui si toccano due rami lineari con contatto quadripunto (che così importa \textit{due} punti quadrupli della \(g_8^1\)), infine corrispondentemente a \(z=d\) una tangente quadrupla, con quattro contatti semplici.}
altri tipi che apparentemente si presentano, si riducono a quelli qui indicati, con semplici scambi dei nomi delle curve di diramazione.

Conviene aggiungere che le due sostituzioni sui rami della \(u\) nell’intorno dei punti di diramazione \(z=a\) e \(z=b\) si estendono in trasformazioni birazionali, dello stesso ordine 4, fra i punti della curva \(K\) la cui equazione
\[
\theta(zu)=0
\]
si ottiene eliminando i radicali che figurano nell’espressione di \(u\); e analogamente si dica per la sostituzione di periodo 2 relativa al punto \(z=d\).

Infatti si può provare, per esempio, che \(u_{12}\) è funzione razionale di \(u_{11}\) e di \(z\); ciò segue dal considerare le due equazioni soddisfatte per un medesimo valore di \(z\):
\[
\theta(uz)=0,\qquad\theta(u'z)=0
\]
e la loro risultante
\[
R(uu')=0:
\]
questa, insieme alla \(\theta(uz)=0\), definisce \(u'(uz)\), come funzione razionale.

% scan: s223R
% !PRINTED 421
\weberpage{421}
% head: LA SUPERFICIE DI GENERE GEOMETRICO NULLO
% stamp: Scattered scan specks and short non-text marks in the margins.
Ora possiamo scrivere le \textit{equazioni di una superficie ellittica, tipo armonico, di determinante \(n=8\):}
\weq{\mathrm{II}_b}{
\left\{
\begin{aligned}
u&=\sqrt[2]{(z-a)(z-b)^3\cdot\varphi(xy)}+\sqrt[4]{(z-b)(z-d)\cdot\psi(xy)},\\
f(xy)&=0
\end{aligned}
\right.
}
dove \(f\) rappresenta un cilindro cubico e \(\varphi\) e \(\psi\) possono supporsi due cilindri, l’uno del \(2^{\mathrm{o}}\) e l’altro del \(4^{\mathrm{o}}\) ordine, il primo tangente ad \(f\) secondo 3 generatrici e il secondo tangente ad esso con contatto quadripunto secondo altre 3 generatrici. Perchè la superficie \(F\) risulti proprio di determinante \(n=8\) (irriducibilità delle curve \(C\) che rispondono su di essa a \(z=\text{cost.}\)) converrà supporre che le due terne di generatrici di contatto di \(f\) e \(\varphi\) e di \(f\) e \(\psi\) non giacciano in piani, sicchè per \(z=\text{cost.}\) risultino irriducibili le due curve
\[
u=\sqrt[2]{(z-a)(z-b)^3\cdot\varphi(xy)},\qquad f(xy)=0
\]
ed
\[
u=\sqrt[4]{(z-b)(z-d)\cdot\psi(xy)},\qquad f(xy)=0,
\]
e ancora che le due curve
\[
u=\sqrt[2]{(z-a)(z-b)^3\cdot\varphi(xy)}
\qquad\text{ed}\qquad
u=\sqrt[4]{(z-b)(z-d)\cdot\psi(xy)}
\]
siano birazionalmente distinte\origfootnote{1}{s223R-1}{Cfr. \textsc{Enriques-Chisini}, \textit{Lezioni.} Libro V, cap. IV, \S~33 (vol. III, pag. 452).}: ciò importa che la terna \(g_1\,g_2\,g_3\) delle generatrici di contatto quadripunto di \(\psi\) con \(f\) sia \textit{dissimile} dalla terna \(g'_1\,g'_2\,g'_3\) delle generatrici di contatto di \(\varphi\) ed \(f\), cioè contata due volte, costituisca una sestina non equivalente a quella che si ottiene sommando alla \(g'_1\,g'_2\,g'_3\) tre generatrici di \(f\) che giacciano in un piano.

\edsection{18. Caso equianarmonico.}{18. Caso equianarmonico.}{sec.X.18}

Nel caso equianarmonico, se la curva ellittica \(K\) contiene una \(g_n^1\) che non sia generata da semplici trasformazioni ordinarie (\(n=2,\ 4\)), essa dovrà venir generata da un gruppo \(\Gamma_n\) di trasformazioni, ciclico o abeliano (a base due), che contenga qualche trasformazione singolare del \(3^{\mathrm{o}}\) o del \(6^{\mathrm{o}}\) ordine. Infatti le curve ellittiche equianarmoniche posseggono, oltre le trasformazioni ordinarie di prima e di seconda specie, \(\pi\) ed \(I\), due serie continue \(\infty^1\) di trasformazioni singolari cicliche del \(3^{\mathrm{o}}\) ordine, \(\omega\) e \(\tau\), e altre due serie di trasformazioni singolari cicliche del \(6^{\mathrm{o}}\) ordine, \(\sigma\) e \(\varrho\)\origfootnote{2}{s223R-2}{\textsc{Enriques-Chisini}, \textit{Lezioni.} Libro V, cap. III, \S~27 (vol. III, pag. 260-63).}. E le trasfor%
% The two equality signs in n = = 2, 4 occur on consecutive printed lines.
%CONT>>
% scan: s224L
% !PRINTED 422
\unskip\weberpage{422}%
% head: CAPITOLO DECIMO
%<<CONT
% stamp: Scattered scan specks, short marginal marks, and a dark vertical line at the right edge.
mazioni dei vari tipi sono legate fra loro dalle seguenti relazioni:
\[
\left\{
\begin{aligned}
\pi\sigma&=\sigma_1 & \sigma_1\sigma&=\tau & \tau\sigma&=I & I\sigma&=\omega & \omega\sigma&=\varrho & \varrho\sigma&=\pi\\
\pi\tau&=\tau_1 & \tau_1\tau&=\omega & \omega\tau&=\pi & \sigma\tau&=I & I\tau&=\varrho & \varrho\tau&=\sigma\\
\pi I&=I_1 & I_1I&=\pi & \sigma I&=\omega & \omega I&=\sigma & \tau I&=\varrho & \varrho I&=\tau.\\ % ?? uncertain: mark after the final \tau in this row -- period, short dash, or scan speck.
\pi\omega&=\omega_1 & \omega_1\omega&=\tau & \pi\omega&=\pi & \sigma\omega&=\varrho & \varrho\omega&=I & I\omega&=\sigma\\ % [sic] \pi\omega=\pi -> \tau\omega=\pi, the same left-hand side already equals \omega_1 in this row and the reverse composition is listed above.
\pi\varrho&=\varrho_1 & \varrho_1\varrho&=\omega & \omega\varrho&=I & I\varrho&=\tau & \tau\varrho&=\sigma & \sigma\varrho&=\pi.
\end{aligned}
\right.
\]

Ora per la \(g_n\) abeliana non generata da trasformazioni ordinarie sulla \textit{curva \(K\) equianarmonica}, saranno possibili \textit{tre casi}:

\textit{a}) \(n=3\), \(g_3^1\) ciclica generata da una trasformazione singolare del \(3^{\mathrm{o}}\) ordine, \(\omega\) o \(\tau\);

\textit{b}) \(n=6\), \(g_6^1\) ciclica, generata da una trasformazione singolare del \(6^{\mathrm{o}}\) ordine, \(\sigma\) o \(\varrho\);

\textit{c}) \(n=9\), \(g_9^1\) abeliana (a base due) il cui gruppo contiene tre trasformazioni singolari cicliche del \(3^{\mathrm{o}}\) ordine, fra loro permutabili, diciamo \(\omega\,\omega_1\,\omega_2\), altre tre trasformazioni cicliche c. s. \(\tau=\omega^2,\ \tau_1=\omega_1^2,\ \tau_2=\omega_2^2\), ed infine tre trasformazioni ordinarie di prima specie \(\pi,\ \pi^2,\ \pi^3=1\).

Questa classificazione può essere facilmente giustificata a partire dall’osservazione che il gruppo \(\Gamma_n\) deve contenere un sottogruppo ciclico invariante \(\Gamma_3\), generato da una trasformazione singolare \(\omega\) o \(\tau\). Siccome il \(\Gamma_3\) genera una \(g_3^1\) ciclica con tre punti tripli, siamo condotti ad esaminare le trasformazioni di \(K\) che rispondono alle proiettività mutanti in sè una retta tripla del tipo
\[
u^3=(z-a)(z-b)(z-d);
\]
le quali sono:

1) le involuzioni, p. es. \((ab)(dd)\), che danno come doppio un punto della terna \((abd)\) e che scambiano fra loro gli altri due punti; e

2) le proiettività cicliche \((abd)\) o \((adb)\).

Ora l’aggiunta al \(\Gamma_3\) di una trasformazione del tipo 1) riesce ad ampliare il \(\Gamma_3\) stesso in un \(\Gamma_6\) ciclico (con un punto sestuplo, due punti tripli e tre punti doppi) che non è contenuto in un gruppo abeliano più ampio; mentre l’aggiunta di una trasformazione \((abd)\) del tipo 2) porta ad un \(\Gamma_9\) abeliano, contenente tre trasformazioni singolari del terz’ordine, il quale non è a sua volta contenuto in un gruppo abeliano più ampio.

Possiamo costruire facilmente le curve equianarmoniche \(K\) che rispondono alle anzidette \(g_n^1\) abeliane, per
\[
n=3,\qquad 6,\qquad 9,
\]
e quindi le corrispondenti superficie ellittiche.

% scan: s224R
% !PRINTED 423
\weberpage{423}
% head: LA SUPERFICIE DI GENERE GEOMETRICO NULLO
% stamp: Scattered scan specks and a thin vertical line at the right edge.
Avremo come \textit{tipo delle superficie ellittiche equianarmoniche di determinante \(n=3\)} (caso ciclico) \textit{le superficie}:
\weq{\mathrm{III}_a}{
\left\{
\begin{aligned}
u&=\sqrt[3]{(z-a)(z-b)(z-d)\cdot\varphi(xy)}\\
f(xy)&=0,
\end{aligned}
\right.
}
dove \(f\) è un cilindro cubico e \(\varphi\) un altro cilindro dello stesso ordine, che tocchi \(f\) con contatto tripunto secondo tre generatrici non giacenti in un piano.

E avremo come \textit{tipo delle superficie ellittiche equianarmoniche di determinante \(n=6\)} (caso ciclico) \textit{le superficie}:
\weq{\mathrm{III}_b}{
\left\{
\begin{aligned}
u&=\sqrt[6]{(z-a)(z-b)^2(z-d)^3\cdot\varphi(xy)}\\
f(xy)&=0,
\end{aligned}
\right.
}
dove \(f\) designa un cilindro cubico e \(\varphi\) un altro cilindro del \(6^{\mathrm{o}}\) ordine che tocca \(f\) con contatto sestipunto, secondo tre generatrici non appartenenti ad un piano.

Nel caso propriamente abeliano, in cui si tratta delle superficie ellittiche equianarmoniche \(F\) di determinante \(n=9\), si può riuscire alla costruzione di un tipo della \(F\), in modo analogo a quello tenuto pel caso armonico (\(n=8\)).

A tal uopo occorre costruire la curva ellittica equianarmonica \(K\) del piano (\(uz\)) che viene rappresentata sopra una retta multipla d’ordine 9 (\(u=0\)), immagine della nostra \(g_9^1\) abeliana, generata da due trasformazioni singolari del \(3^{\mathrm{o}}\) ordine, fra loro permutabili. Come tipo si potrà assumere la curva
\[
u=\sqrt[3]{(z-a)^2(z-d)\cdot\varphi}+\sqrt[3]{(z-b)(z-d)^2\cdot\psi}
\]
con \(\varphi\) e \(\psi\) costanti.

Per dimostrare che la \(u(z)\) così definita risponde effettivamente ad una curva ellittica equianarmonica contenente la sopra indicata \(g_9^1\) abeliana, converrà calcolare le sostituzioni sui rami della \(u\) in ordine ai punti di diramazione
\[
z=a,\qquad z=b,\qquad z=d.
\]

A tal uopo si designeranno i rami con due indici, il primo dei quali assumerà i valori 1, 2, 3 in relazione ai valori del primo radicale cubico
\[
\sqrt[3]{},\qquad \varepsilon\sqrt[3]{},\qquad \varepsilon^2\sqrt[3]{}
\qquad
\left(\varepsilon=e^{\frac{2\pi i}{3}}\right)
\]
che figura nell’espressione di \(u(z)\), mentre il secondo assumerà similmente tre valori in rapporto a quelli del secondo radicale cubico.
%CONT>>
% scan: s225L
% !PRINTED 424
\weberpage{424}
% head: CAPITOLO DECIMO
%<<CONT
% stamp: Heavy mottling, strongest along the right edge, with a dark vertical edge band and scattered marginal marks.
Avremo, in ordine ai tre punti di diramazione, le sostituzioni seguenti:

per \(z=a\)
\[
S_a=(u_{11}\,u_{31}\,u_{21})(u_{12}\,u_{32}\,u_{22})(u_{13}\,u_{33}\,u_{23});
\]

per \(z=b\)
\[
S_b=(u_{11}\,u_{12}\,u_{13})(u_{31}\,u_{32}\,u_{33})(u_{21}\,u_{22}\,u_{23})
\]

e per \(z=d\)
\[
S_d=(u_{11}\,u_{23}\,u_{32})(u_{12}\,u_{21}\,u_{33})(u_{13}\,u_{22}\,u_{31}).
\]

La terza sostituzione è l’inversa del prodotto delle altre due, di guisa che viene soddisfatta la condizione d’esistenza della retta multipla \(u=0\), non diramata nel punto all’infinito \(z=\infty\)\origfootnote{1}{s225L-1}{Cfr. \textsc{Enriques-Chisini}, \textit{Lezioni.} Libro V, cap. IV, \S~38 (vol. III, pag. 429).}: \(S_dS_bS_a=1\).

Quindi si riconosce che le dette sostituzioni generano un gruppo abeliano \(\Gamma_9\), che dà luogo ad una \(g_9^1\), egualmente abeliana, sopra una curva ellittica equianarmonica.

Così si è condotti a dimostrare che le \textit{superficie ellittiche equianarmoniche di determinante \(n=9\)} (caso abeliano) \textit{danno luogo al tipo}
\weq{\mathrm{III}_c}{
\left\{
\begin{aligned}
u&=\sqrt[3]{(z-a)^2(z-d)\cdot\varphi(xy)}+\sqrt[3]{(z-b)(z-d)^2\psi(xy)}\\
f(xy)&=0,
\end{aligned}
\right.
}
dove \(f\) è un cilindro cubico e \(\varphi\) e \(\psi\) possono supporsi del pari cilindri cubici osculanti \(f\) secondo due terne di generatrici non equivalenti, che non giacciano in uno stesso piano.

\edsection{19. Riassunto.}{19. Riassunto.}{sec.X.19}

Possiamo riassumere l’analisi delle superficie ellittiche di genere \(p_g=0\) con curva canonica virtuale d’ordine zero, enunciando il seguente teorema:

\textit{Le superficie ellittiche di genere \(p_g=0\) e \(p_a=-1\) con curva canonica virtuale d’ordine zero, caratterizzate dal valore del \(P_{12}=1\),} cioè le superficie con \(p_a=-1\) possedenti un fascio ellittico di curve ellittiche \(K\) e un secondo fascio lineare di curve ellittiche \(C\), \textit{si distribuiscono in sette famiglie di determinante \(n=2,\ 4;\ 4,\ 8;\ 3,\ 6,\ 9\), rappresentate da uno dei tipi d’equazioni \(I_a I_b II_a II_b III_a III_b III_c\), ovvero da un tipo simile.}

I tipi \(I_a\) e \(I_b\) (risp. tipo \textit{ciclico} e tipo \textit{abeliano}) rispondono a curve ellittiche \(C\) di \textit{modulo generale}; \(II_a\) e \(II_b\) (tipo \textit{ciclico} e tipo \textit{abeliano})
%CONT>>
% scan: s225R
% !PRINTED 425
\weberpage{425}
% head: LA SUPERFICIE DI GENERE GEOMETRICO NULLO
%<<CONT
% stamp: Heavy mottling and scattered scan specks, especially along the left side.
a curve \(C\) \textit{armoniche}, ed infine \(III_a\) e \(III_b\) a un tipo \textit{ciclico} con \(C\) \textit{equianarmoniche} e \(III_c\) \textit{a un tipo equianarmonico abeliano.}

Le superficie ellittiche predette si distinguono anche, fino ad un certo punto, mediante i valori dei plurigeneri inferiori al \(P_{12}\), calcolabili colla formula del \S~14. Le superficie \(I_a\) e \(I_b\) (determinante 2 e 4) hanno i caratteri:
\[
P_2=1,\qquad P_4=1,\qquad P_{12}=1;
\]
le superficie di tipo armonico \(II_a\) e \(II_b\) (determinante 4 e 8) hanno:
\[
(P_2=0),\qquad P_4=1,\qquad P_6=0,\qquad P_{12}=1;
\]
ed infine le superficie del tipo equianarmonico \(III_a\) (determinante 3) hanno
\[
P_2=0,\qquad P_3=1,\qquad P_4=1,\qquad P_{12}=1, % ?? uncertain: punctuation after P_{12}=1 -- comma, possibly followed by a period.
\]
mentre le \(III_b\) e \(III_c\) (determinante 6 e 9) hanno
\[
P_2=0,\qquad P_3=0,\qquad P_4=1,\qquad P_{12}=1.
\]

\textit{Notizia.} -- Le superficie ellittiche con curva canonica virtuale d’ordine zero, già osservate da \textsc{Enriques} nel suo studio sulle superficie di genere zero (1905), si presentano come particolari superficie iperellittiche (più precisamente come involuzioni sopra la superficie di Jacobi che corrisponde ad una curva degenere di genere due), ed in questo senso sono state incontrate da G. \textsc{Bagnera} e M. \textsc{De Franchis} nella memoria del 1907 « Sopra le superficie algebriche che hanno le coordinate del punto generico esprimibili con funzioni meromorfe quadruplamente periodiche di due parametri »\origfootnote{1}{s225R-1}{Rendic. Lincei, s. V, vol. XVI, \(1907_1\).}.

I detti autori scoprono che tali superficie danno luogo a un numero finito di tipi e riescono a classificarle assegnandone le rappresentazioni parametriche iperellittiche, che rispondono alle nostre 7 equazioni tipiche.

Indipendentemente da \textsc{Bagnera} e \textsc{De Franchis}, ma dopo di loro, anche \textsc{Enriques} e \textsc{Severi} hanno ritrovato in altro modo codesto risultato, sempre partendo dalla rappresentazione parametrica delle nostre superficie mediante funzioni iperellittiche\origfootnote{2}{s225R-2}{F. \textsc{Enriques} e F. \textsc{Severi}, \textit{Intorno alle superficie iperellittiche irregolari.} Rendic. Lincei, s. V, vol. XVII, \(1908_1\) (cfr. degli stessi autori ibidem, vol. XVI, \(1907_1\)). \textit{Mémoire sur les surfaces hyperelliptiques.} Acta Mathem., t. 32 e 33 (1909). Vedansi in ispecie i \S\S~56 e 57.}.

In codesta trattazione ci sembra particolarmente notevole il concetto di costruire la superficie iperellittica o abeliana propria, come superficie ellittica multipla priva di curve di diramazione, calcolandone i caratteri che valgono a definirla (\(p_a=-1,\ p_g=1,\)
%CONT>>
% scan: s226L
% !PRINTED 426
\weberpage{426}
% head: CAPITOLO DECIMO
%<<CONT
% stamp: Scattered scan specks and a dark vertical scan-edge line at the right.
\(P_4=1\) secondo un teorema di Picard-Enriques). Per esempio, riferendoci alla superficie ellittica più semplice, del tipo \(I_a\), si avrà una \(F_m\) d'un certo ordine \(m\) in \(S_3\), possedente una superficie biaggiunta \(\varphi_{2n-8}\) e si otterrà una superficie iperellittica propria (\(p_a=-1\), \(p_g=1\)) estraendo la \(\sqrt[2]{\varphi_{2n-8}}\) sopra la \(F_m\). % ?? uncertain: the letter following 2 in both subscripts of \varphi -- n or m.

La presente trattazione del problema, cioè la classificazione algebrica delle nostre superficie ellittiche (che, nella parte costruttiva, utilizza il citato lavoro di O. \textsc{Chisini}), ha formato oggetto di una Nota di \textsc{Enriques} « Sulle superficie ellittiche di genere zero »% EDITORIAL: numero originale fissato esplicitamente.
\origfootnote{3}{s226L-1}{Rendic. Lincei, s. VI, vol. XIX, 18 febbraio 1934.} del 1934, e si trova poi esposta nelle « Lezioni » di \textsc{Enriques-Campedelli} del Seminario Matematico di Roma (1934).

\clearpage\thispagestyle{plain}
% scan: s226R
% !PRINTED none
\weberpage{[s226R]}
% head: none
% stamp: Scattered scan specks.
\edchapter{XI}{CLASSIFICAZIONE GENERALE DELLE SUPERFICIE}{Capitolo XI. CLASSIFICAZIONE GENERALE DELLE SUPERFICIE}{CLASSIFICAZIONE GENERALE DELLE SUPERFICIE}

\edsection{1. Introduzione.}{1. Introduzione.}{sec.XI.1}

Il lettore che abbia seguito gli sviluppi di questo trattato, se — come spesso accade fra i matematici — porta amore soprattutto alle verità generali, può averne ritratto l'impressione che l'autore abbia dato troppo posto ad esempi e casi particolari, lasciandosi in qualche modo guidare dal sentimento di curiosità del naturalista che raccoglie in un museo i più diversi tipi di animali o di piante o di minerali.

Ma come il museo riesce a dare un'idea della ricchezza di forme della vita e conduce quindi a problemi generali della biologia, anche la raccolta di esempi, in questo campo delle matematiche, assume un significato essenziale sotto l'aspetto euristico o storico-costruttivo della scienza. Possiamo illustrarne il valore ripetendo le parole con cui G. \textsc{Castelnuovo} rendeva conto dei nostri sforzi, fatti in comune, per dissipare le oscurità che incontravamo agl'inizi della teoria delle superficie% EDITORIAL: numero originale fissato esplicitamente.
\origfootnote{1}{s226R-1}{\textit{La geometria algebrica e la scuola italiana.} Conferenza tenuta al Congresso Internazionale dei Matematici, Bologna, settembre 1928.}.

« Avevamo costruito, in senso astratto s'intende, un gran numero di modelli di superficie del nostro spazio o di spazi superiori; e questi modelli avevamo distribuito, per dir così, in due vetrine. Una conteneva le superficie regolari per le quali tutto procedeva come nel migliore dei mondi possibili; l'analogia permetteva di trasportare ad esse le proprietà più salienti delle curve piane. Ma quando cercavamo di verificare queste proprietà sulle superficie dell'altra vetrina, le irregolari, cominciavano i guai, e si presentavano eccezioni d'ogni specie. Alla fine lo studio assiduo dei nostri modelli ci aveva condotto a divinare alcune proprietà che dovevano sussistere, con modificazioni opportune, per le superficie di ambedue
%CONT>>
% scan: s227L
% !PRINTED 428
\weberpage{428}
% head: CAPITOLO UNDECIMO
%<<CONT
% stamp: Scattered scan specks, stray diagonal marks in the text, and a dark vertical scan-edge line at the right.
le vetrine; mettevamo poi a cimento queste proprietà colla costruzione di nuovi modelli. Se resistevano alla prova ne cercavamo, ultima fase, la giustificazione logica ».

Ma ora possiamo dire di più: le determinazioni di particolari famiglie o classi di superficie su cui ci siamo indugiati, non danno semplici esempi, il cui significato si lasci superare nel processo induttivo della scienza; poichè esse si ravvicinano in una sintesi che riesce a ciò che può dirsi scopo principale di questa teoria algebrico-geometrica; cioè alla classificazione generale delle superficie algebriche.

A tale proposito conviene notare che la geometria algebrica è, in sostanza, un'aritmetica superiore, dove il criterio stesso della generalità ha soltanto un significato relativo. C'è qui un'osservazione assai profonda, che risale al \textsc{Poncelet} e al \textsc{Noether}: da un punto di vista puramente logico non può dirsi, per esempio, che la famiglia delle superficie non riferibili a rigate sia più ampia di quella delle rigate. Se si procede a classificare le superficie (dello spazio ordinario o di un iperspazio) secondo il loro ordine \(n\), si presenta come più generale il caso di superficie non riferibili a rigate le cui sezioni iperpiane sono di genere \(\pi\) con
\[
n\leq 2\pi-2;
\]
ma se invece si procede a classificare le superficie secondo il genere \(\pi\) delle loro sezioni iperpiane, apparirà più generale il caso in cui l'ordine
\[
n>2\pi-2,
\]
cioè il caso che risponde a superficie riferibili a rigate.

Da ciò che si è detto non vogliamo trarre conclusioni paradossali. Ma se continueremo a ritenere lo studio delle superficie non appartenenti alla famiglia delle rigate come \textit{più generale} di quello delle rigate, questo giudizio, piuttosto che un'affermazione di fatto, sarà un \textit{giudizio di valore}, in cui si soppesano — per così dire — i problemi che si riferiscono all'uno e all'altro tipo, la ricchezza di casi cui danno luogo, la possibilità di subordinare l'una all'altra famiglia da punti di vista più significativi. Da ciò conviene trarre questo insegnamento: che ogni famiglia o classe di superficie, la quale presenti caratteri propri che la distinguono dalle rimanenti, costituisce un oggetto di studio altrettanto degno di essere perseguito, e che la definizione di una tale famiglia in ordine a caratteri interi invarianti importa spesso un interesse d'ordine generale. Per esempio la definizione della famiglia delle rigate mediante l'annullamento del quadrigenere e del sestigenere costituisce, non tanto un risultato particolare interessante codesta famiglia, quanto un teorema significativo della
%CONT>>
% scan: s227R
% !PRINTED 429
\weberpage{429}
% head: CLASSIFICAZIONE GENERALE DELLE SUPERFICIE
%<<CONT
% stamp: Scattered scan specks.
teoria generale; che c'insegna l'esistenza di curve canoniche o pluricanoniche (d'ordine \(\geq 0\)) sopra qualsiasi superficie non riferibile a rigata.

La nostra classificazione, raccogliendo in una sintesi i risultati delle analisi precedenti, riesce a distinguere 4 famiglie di superficie, che hanno proprietà essenzialmente diverse e che si definiscono col valore del \(P_{12}\) e — subordinatamente — del \(p^{(1)}\). Sono:

\textit{A}) le rigate (esistenza di sistemi lineari di genere \(\pi\) e grado \(n>2\pi-2\), infinite curve eccezionali non eliminabili, ed anche — come vedremo — serie continue di trasformazioni birazionali non formanti un gruppo di dimensione finita);

\textit{B}) le superficie (prive di curve eccezionali) con curve pluricanoniche d'ordine zero (sistemi lineari puri di grado \(n=2\pi-2\));

\textit{C}) le superficie (\(n<2\pi-2\) per \(\pi>1\)) aventi curve canoniche o pluricanoniche composte colle curve ellittiche d'un fascio (\(p^{(1)}=1\));

\textit{D}) le superficie (\(n<2\pi-2\)) con \(p^{(1)}>1\), che danno luogo a superficie canoniche o pluricanoniche porgenti un modello le cui proprietà proiettive rispecchiano le proprietà invariantive della classe di superficie.

Ognuna di queste famiglie costituisce un oggetto di studio interessante di per sè, ma il rilievo delle proprietà caratteristiche che la distinguono dalle altre ha, come abbiam detto, un valore in ordine alla teoria generale.

Ora, perchè la classificazione assuma il suo proprio significato, conviene estendere alle superficie di genere \(p_g>0\) alcune osservazioni fatte per il caso \(p_g=0\); e per questo scopo giova anzitutto approfondire lo studio delle superficie che ammettono una serie continua di trasformazioni birazionali in se stesse.

\edsection{2. Superficie che ammettono una serie continua di trasformazioni\\birazionali in se stesse: casi che conducono alle rigate.}{2. Superficie che ammettono una serie continua di trasformazioni birazionali in se stesse: casi che conducono alle rigate.}{sec.XI.2}

Sappiamo che le superficie riferibili a rigate posseggono serie continue di trasformazioni birazionali in se stesse; in particolare sopra una superficie razionale, ovvero sul piano, si conoscono% EDITORIAL: numero originale fissato esplicitamente.
\origfootnote{1}{s227R-1}{\textsc{Enriques}, 1893. Cfr. \textsc{Enriques-Chisini}, \textit{Lezioni.} Libro V; cap. II, § 28 (vol. III, pag. 200).} i tipi di gruppi continui di trasformazioni cremoniane, e d'altra parte si hanno pure serie continue di trasformazioni, come sono le quadratiche, che non generano per moltiplicazione gruppi continui di dimensioni finite. E serie simili si possono anche riconoscere sopra una superficie rigata di genere \(p>0\); dove tuttavia si avranno serie
%CONT>>
% scan: s228L
% !PRINTED 430
\weberpage{430}
% head: CAPITOLO UNDECIMO
%<<CONT
% stamp: Scattered scan specks and a dark vertical scan-edge line at the right.
di trasformazioni operanti sulla superficie in modo non transitivo, se sia \(p>1\).

Con qualche restrizione aggiuntiva, le osservazioni precedenti si invertono, cioè il possesso di un'infinità continua di trasformazioni birazionali permette di caratterizzare la famiglia delle rigate. Ciò può farsi da quattro punti di vista:

1) appartiene alla \textit{famiglia delle rigate} ogni superficie che possegga un \textit{gruppo razionale} \(\infty^1\) \textit{di trasformazioni in sè}.

2) appartiene alla \textit{famiglia delle rigate} ogni superficie che possegga un \textit{gruppo continuo} (algebrico o trascendente) \textit{di trasformazioni in sè}, che lasci transitivamente% EDITORIAL: numero originale fissato esplicitamente.
\origfootnote{1}{s228L-1}{S'intende che permuti in infiniti modi gli elementi (curve) di \(|L|\).} invariato un \textit{sistema lineare}, \(|L|\), \(\infty^1\) almeno, di curve.

3) appartiene alla \textit{famiglia delle rigate} ogni superficie \(F\) che possegga un \textit{gruppo continuo} di trasformazioni in sè, \textit{di dimensione} \(r\geq 2\), che operi \textit{intransitivamente} sui punti di \(F\), ovvero di \textit{dimensione} \(r\geq 3\) che operi \textit{transitivamente} su di essa.

4) appartiene alla \textit{famiglia delle rigate} ogni superficie che possegga una \textit{serie continua} di trasformazioni in sè, \textit{non generante un gruppo} di dimensione finita.

Il teorema I) discende direttamente dal teorema di Noether-Enriques (1898)\origfootnote{2}{s228L-2}{Cfr. p. es. \textsc{Conforto}, \textit{Le superficie razionali.} Libro II, cap. I, §§ 4, 5, 6 (pag. 245, 260, 262).} concernente le superficie con un fascio di genere \(p\geq 0\) di curve razionali; % ?? uncertain: the theorem label -- I) or 1).
% EDITORIAL D0228 RISOLTO: The glyph has horizontal top and bottom serifs and a vertical stem: capital I, distinct from the nearby numeric 1.
perchè le traiettorie del gruppo razionale \(\infty^1\) formano. appunto un fascio di curve razionali. % ?? uncertain: the mark after "formano" -- period, comma, or scan speck.
Ma convien dire che il risultato di cui si tratta è stato conseguito da P. \textsc{Painlevé} prima che \textsc{Enriques} avesse esteso al caso \(p>0\) il teorema dato da \textsc{Noether} per \(p=0\): infatti la circostanza che le curve razionali \(C\) d'un fascio sono traiettorie d'un gruppo \(\infty^1\) permette, agevolmente, di costruire una curva direttrice unisecante le \(C\), mercè cui le \(C\) stesse si trasformano razionalmente in rette.

Il teorema 2) enuncia, in sostanza, che « una superficie trasformata in sè da un gruppo continuo di omografie, è razionale o riferibile a rigata », e sotto questa forma proiettiva fu dato con qualche restrizione, da \textsc{Enriques} (1893), ed ha ricevuto la più completa dimostrazione da \textsc{Fano} (1895, 96, 97)\origfootnote{3}{s228L-3}{Cfr. \textsc{Conforto}, \textit{Le superficie razionali.} Op. cit. Libro II, cap. II, § 37 (pagina 496).}.

Tuttavia conviene porgere una nuova dimostrazione del teorema stesso, che contempli anche il caso di un gruppo \(\Gamma\) per cui resti invariato un sistema lineare \(|L|\) (\(\infty^2\) almeno) appartenente ad una
%CONT>>
% scan: s228R
% !PRINTED 431
\weberpage{431}
% head: CLASSIFICAZIONE GENERALE DELLE SUPERFICIE
%<<CONT
% stamp: Scattered scan specks.
involuzione, e il caso di un \(\Gamma\) che lasci invariato un fascio lineare \(|L|\); anzi la dimostrazione che proponiamo ricondurrà il teorema generale proprio a quest'ultimo caso.

Prendiamo le mosse dall'osservazione (di \textsc{Picard}) che « un gruppo continuo di trasformazioni birazionali di una superficie algebrica in sè è sempre contenuto in un \textit{gruppo continuo algebrico} »% EDITORIAL: numero originale fissato esplicitamente.
\origfootnote{1}{s228R-1}{Invero le condizioni perchè una superficie \(F\) venga trasformata in sè da trasformazioni birazionali di un dato ordine, sono algebriche, e quindi valgono a definire un gruppo algebrico, continuo o misto. In questa seconda ipotesi il gruppo misto contiene entro di sè un gruppo continuo.}.

Se la superficie \(F\) ammette un gruppo continuo algebrico di trasformazioni in sè, che lasci invariato transitivamente un fascio lineare di curve \(|L|\), si può supporre che le curve \(L\) sieno permutate da \(\Gamma\) secondo un gruppo \(\infty^1\), poichè altrimenti basterebbe porre al posto di \(\Gamma\) il sottogruppo delle sue trasformazioni che lascin ferme una ovvero due particolari curve \(L\).

Ciò posto converrà distinguere le seguenti ipotesi:

\textit{a}) il gruppo \(\Gamma\) è \(\infty^1\); allora le sue traiettorie, ammettendo \(\infty^1\) trasformazioni in sè, sono curve \(K\) (a priori razionali o ellittiche)\origfootnote{2}{s228R-2}{Teorema di \textsc{Schwartz}. Cfr. \textsc{Enriques-Chisini}, \textit{Lezioni.} Libro V, cap. III, § 31 (vol. III, pag. 309).} su cui le \(L\) segano un'involuzione razionale invariante e perciò sono razionali\origfootnote{3}{s228R-3}{Le curve che ammettono infinite trasformazioni in sè che lascino invariata una \(g^1\), si riducono mercè un multiplo di questa serie a curve (di un certo spazio) con infinite trasformazioni proiettive e perciò sono razionali. Cfr. p. es. \textsc{Enriques-Chisini}, op. cit. Libro V, cap. III.}: segue che la \(F\) è riferibile ad una rigata. % [sic] Schwartz -> Schwarz, the running text on this page spells the same name Schwarz.

\textit{b}) il gruppo \(\Gamma\) è \(\infty^1\) almeno, ed opera transitivamente sui punti della superficie. Allora le curve del fascio \(|L|\), trasformate in sè da \(\infty^1\) trasformazioni del \(\Gamma\), sono razionali o ellittiche.

Se sono razionali la \(F\) si lascia trasformare in una rigata. Se invece le \(L\) sono ellittiche, si consideri la serie \(\infty^1\) ellittica \(\{D\}\) che è costituita dalle trasformate di una curva \(D\), scelta in modo affatto generale, sopra la \(F\): per ogni punto \(P\) di \(f\) vi è un certo numero \(i\) di curve \(D\) di \(\{D\}\), e \(P\) è suscettibile di variare sopra una curva \(K\) in modo che i gruppi di \(i\,D\) uscenti da esso risultino equivalenti; perciò le \(K\) così definite debbono formare un fascio ellittico \((K)\), senza punti base\origfootnote{4}{s228R-4}{Cfr. Cap. IX, 12.}. % [sic] f -> F, the surface is otherwise denoted by capital F in this paragraph.

Quindi, fra le trasformazioni del nostro \(\Gamma\) permutanti gli elementi (curve) di \((K)\) vi saranno \(\infty^1\) trasformazioni, formanti un sottogruppo \(\Gamma'\), che avranno come traiettorie le \(K\). A priori, per il teorema di \textsc{Schwarz}, queste \(K\) risulterebbero razionali o ellittiche; ma il se% The word "secondo" is divided across the page break; the parts are joined without a hyphen.
%CONT>>
% scan: s229L
% !PRINTED 432
\unskip\weberpage{432}%
% head: CAPITOLO UNDECIMO
%<<CONT
% stamp: Scattered scan specks, stray diagonal marks in the text, and a dark vertical scan-edge line at the right.
\unskip condo caso si esclude, osservando che le trasformazioni di \(\Gamma'\) debbono lasciare invariata la serie lineare \(g_n^1\) segata su una \(K\) dalle \(L\) di \(|L|\). Si deduce che la \(F\) è riferibile ad una rigata, \hfill c. d. d.

Dopo avere così stabilito il teor. 2) pel caso in cui il sistema lineare \(|L|\), invariante pel gruppo \(\Gamma\), sia un fascio, è facile estendere la dimostrazione al caso in cui \(|L|\) sia un sistema lineare, \(\infty^r\), con \(r>1\). Infatti le trasformazioni di \(\Gamma\) (che può supporsi gruppo continuo algebrico) operano sugli elementi (curve) di \(|L|\) come omografie di un \(S_r\), formanti del pari un gruppo algebrico \(\Gamma'\) e saranno pure algebrici i sottogruppi di dimensione \(\geq 1\) formati da tutte le trasformazioni di \(\Gamma'\) permutabili con una data \(\omega\). Ora un sottogruppo siffatto, diciamo \(\overline{\Gamma}'\), lascia fermi gli stessi punti uniti di \(\omega\) e quindi anche una retta unita, non tutta costituita di punti uniti: alla quale corrisponde un fascio lineare di curve \(L\), transitivamente invariante per \(\overline{\Gamma}'\). Si ricade così nel caso trattato innanzi.

Il teorema 3)% EDITORIAL: numero originale fissato esplicitamente.
\origfootnote{1}{s229L-1}{Cfr. \textsc{Castelnuovo-Enriques}, \textit{Sur les surfaces algébriques admettant un groupe continu de transformations birationelles en elles-mêmes.} (Comptes rendus, Paris, 29 Luglio 1895).} si dimostra come segue. % [sic] birationelles -> birationnelles, apparent missing n in the French adjective.

Pongasi dapprima che il gruppo \(\Gamma\), \(\infty^r\), delle trasformazioni di \(F\) operi intransitivamente sulla superficie, vale a dire trasformi in sè le curve \(C\) d'un fascio sopra di essa. Se il \(\Gamma\) è (come può supporsi) un gruppo algebrico \(\infty^r\) con \(r\geq 2\), si otterrà un suo sottogruppo algebrico \(\infty^1\), \(\Gamma'\), imponendo alle trasformazioni di \(\Gamma\) di lasciar fermi \(r-1\) punti generici di \(F\): diciamo \(O, O'\ldots\).

Ora il gruppo \(\Gamma'\) — serie algebrica \(\infty^1\) di elementi con infinite trasformazioni in sè — sarà, a priori, razionale o ellittico. Ma non può essere ellittico perchè contiene una serie \(g_n^1\) invariante, d'un certo ordine \(n>1\), che risponde alla serie dei punti di \(F\) infinitamente vicini ad \(O\). Dunque è razionale e (conforme al teor. 1) la \(F\) è riferibile ad una rigata.

A dir vero il ragionamento precedente darebbe luogo ad un'eccezione se le trasformazioni di \(\Gamma'\) lasciassero fermi tutti i punti di \(F\) infinitamente vicini ad \(O\); ma in tal caso si ripeterà lo stesso discorso per riguardo all'intorno successivo di un punto \(O_1\) vicino ad \(O\) nell'intorno del prim'ordine, e similmente — ove occorra — per il primo intorno di un punto \(O_s\) vicino ad \(O\) nell'intorno d'ordine \(s\), che non sia tutto costituito di punti uniti pel \(\Gamma'\).

In modo analogo si dimostra il secondo caso del teorema, dove si ha un gruppo algebrico \(\Gamma\), \(\infty^3\) almeno, che opera transitivamente sui punti di \(F\). Riferendoci all'ipotesi più semplice, in cui il \(\Gamma\) abbia la dimensione 3, si otterrà un sottogruppo algebrico \(\infty^1\), \(\Gamma'\), imponendo alle trasformazioni di \(\Gamma\) di lasciar fermo un punto \(O\), e
%CONT>>
% scan: s229R
% !PRINTED 433
\weberpage{433}
% head: CLASSIFICAZIONE GENERALE DELLE SUPERFICIE
%<<CONT
% stamp: Scattered scan specks and faint vertical scan-edge marks at the right.
— come nel caso precedente — si riconoscerà che il \(\Gamma'\) è un gruppo \textit{razionale}.

Se invece il \(\Gamma\) abbia una dimensione \(r>3\), si otterrà similmente un suo sottogruppo \(\Gamma'\), razionale \(\infty^1\), imponendo alle trasformazioni di \(\Gamma\) di lasciar fermo un punto \(O\) di \(F\), e poi ancora un punto \(O_1\), infinitamente vicino ad \(O\) ecc.

Finalmente dimostriamo il teorema 4) che risponde ad una questione più riposta, sollevata da E. \textsc{Picard}% EDITORIAL: numero originale fissato esplicitamente.
\origfootnote{1}{s229R-1}{\textsc{Castelnuovo-Enriques}, \textit{Sopra alcune questioni fondamentali nella teoria delle superficie algebriche.} Annali di Mat., 1901 (n. 19).}.

Se la superficie \(F\) ammette una serie continua \(\Gamma\) — poniamo \(\infty^1\) — di trasformazioni in sè, che non generi un gruppo (di dimensione finita), vuol dire che i prodotti delle trasformazioni di \(\Gamma\), prese \(r\) ad \(r\), costituiranno in generale le trasformazioni di una serie \(\infty^r\), \(\Gamma_r\), dove \(r\) è suscettibile di diventare grande ad arbitrio. Quindi le trasformazioni di \(\Gamma_r\) muteranno una curva \(C\), scelta nel modo più generale su \(F\), nelle curve \(C_r\) di un sistema continuo \(\{C_r\}\) di dimensione \(r\). Se la curva \(C\) appartiene ad un sistema lineare privo di punti base, di genere \(\pi\) e grado \(n\), il \(\{C_r\}\) sarà formato di sistemi lineari \(|C_r|\) dello stesso genere e grado, aventi un certo numero di punti base di molteplicità \(i\) (\(i=1,\,2\,\ldots.\)); e il grado e il genere virtuali di un \(|C_r|\) saranno dati rispettivamente da % ?? uncertain: the number of dots after 2 in the parenthetical enumeration -- four or three.
% EDITORIAL D0230 RISOLTO: The enlarged enumeration shows four discrete baseline dots before the closing parenthesis.
\[
N_r=n+\Sigma h_i^2
\]
e
\[
\Pi_r=\pi+\Sigma\frac{h_i(h_i-1)}{2}
\]

Pertanto si avrà
\[
N_r-2\Pi_r=n-2\pi+\Sigma h_i.
\]

Ma, siccome il grado \(N_r\) di \(\{C_r\}\) (che a priori è \(\geq r-1\)) cresce con \(r\) quanto si vuole, si potrà scegliere \(r\) in modo che sia
\[
\Sigma h_i^2>(2\pi-n)^2
\]
e quindi
\[
\Sigma h_i\geq\sqrt{\Sigma h_i^2}>2\pi-n.
\]

Così appare che la superficie \(F\) possiede curve di genere \(\Pi_r\) e di grado \(N_r>2\Pi_r-2\) e perciò\origfootnote{2}{s229R-2}{Cfr. Cap. X, § 2.} è riferibile ad una rigata. \hfill c. d. d.

\edsection{3. Superficie ellittiche e iperellittiche.}{3. Superficie ellittiche e iperellittiche.}{sec.XI.3}

L'analisi del precedente paragrafo ci permette ora di «determinare tutte le superficie \(F\), non riferibili a rigate, che ammettono un gruppo continuo di trasformazioni in se stesse ».

% EDITORIAL: segnatura originale: \textsc{Enriques F.} - \textit{\gesperrt{Superficie algebriche.}}\hfill 28 
\printersignature{\textsc{Enriques F.} - \textit{\gesperrt{Superficie algebriche.}}\quad 28} % ?? uncertain: the footer title's spacing -- letter-spaced italic or ordinary italic.

% scan: s230L
% !PRINTED 434
\weberpage{434}
% head: CAPITOLO UNDECIMO
% stamp: Scattered scan specks and a dark vertical scan-edge line at the right.
Anzitutto si può supporre che questo gruppo \(\Gamma\) (non contenuto in un gruppo più ampio) sia \textit{algebrico}; quindi il gruppo darà luogo a uno dei casi seguenti:

1) \(\Gamma\) opera intransitivamente sui punti della superficie \(F\), ed è allora (§ 2) un gruppo \(\infty^1\) ellittico che muta in sè le curve ellittiche \(C\) d'un fascio, avente un certo genere \(\varrho\geq 0\);

2) ovvero \(\Gamma\) è un gruppo \(\infty^2\) di trasformazioni, che opera transitivamente sui punti di \(F\).

\textit{Primo caso.} – Nel primo caso \textit{la superficie \(F\) che ammette un gruppo algebrico \(\infty^1\) di trasformazioni birazionali è una superficie ellittica, ed ha genere numerico \(p_a=-1\), il genere geometrico \(p_g=\varrho\), e il genere lineare \(p^{(1)}=1\).}

Infatti le \(\infty^1\) trasformazioni di \(\Gamma\) portano una curva \(L\), scelta in modo generale su \(F\), nelle \(\infty^1\) curve di una serie ellittica \(\{L\}\); per ogni punto \(P\) vi sarà un certo numero \(i\) di curve \(L\) e \(P\) potrà muoversi in \(F\) sopra una curva \(K\) in modo che il gruppo delle \(i\,L\) suddette si mantenga equivalente a se stesso; quindi le curve \(K\) formeranno su \(F\) un fascio ellittico \((K)\), sulle curve del quale le \(L\) segano gruppi equivalenti: per questo motivo le curve \(K\) sono essenzialmente distinte dalle \(C\), traiettorie del \(\Gamma\).

Ora la superficie \(F\), che contiene un fascio di genere \(\varrho\) di curve \(C\) ellittiche e un fascio ellittico di curve \(K\) (di genere \(>0\)) è ellittica (Cap. X, § 12) è possiede i caratteri anzidetti, che si lasciano facilmente valutare. % [sic] è possiede -> e possiede, the conjunction is printed with an accent.

A tale scopo si calcoli il numero \(\Delta\) delle curve di \((K)\) dotate di punto doppio e si esprima per mezzo dell'invariante di Zeuthen-Segre. Poichè \((K)\) è un fascio ellittico si troverà
\[
\Delta-4=12p_a-p^{(1)}+9.
\]

E poichè le curve di \((K)\) sono trasformate l'una nell'altra dalle trasformazioni di \(\Gamma\) che non ne lasciano ferma alcuna:
\[
\Delta=0;
\]
quindi
\[
\begin{aligned}
&12\,p_a-p^{(1)}+13=0,\\
&\quad p_a=-1,\qquad p^{(1)}=1.
\end{aligned}
\]

Infine, se \(\varrho=0\) segue \(p_g=0\), e se \(\varrho>0\) i gruppi della \(g_{2\varrho-2}^{\varrho-1}\) entro il fascio \((C)\) costituiranno le curve canoniche di \(F\) (Cap. X, § 12); perciò sarà, in ogni caso
\[
p_g=\varrho.
\]

\textit{Nota.} – Il più semplice esempio di superficie ellittiche di genere \(p_g>0\) è quello delle superficie \(F\) di determinante 1, rappresentative delle coppie di punti di una curva ellittica e di una curva di genere
%CONT>>
% scan: s230R
% !PRINTED 435
\weberpage{435}
% head: CLASSIFICAZIONE GENERALE DELLE SUPERFICIE
%<<CONT
% stamp: Scattered scan specks.
\(p_g\), le quali si lasciano d'altra parte rappresentare sopra un cilindro cubico multiplo \(\Phi\) (cfr. Cap. X, § 13). Le superficie ellittiche di genere \(p_g>0\) e di determinante \(n>1\) si possono costruire partendo dalla \(F\) di determinante 1 multipla, o anche da un cilindro cubico multiplo. In sostanza si ripetono con poche modifiche le considerazioni svolte pel caso \(p_g=0\).

Appare pertanto che \textit{le superficie ellittiche per \(p_g>0\) danno luogo ad infinite famiglie corrispondenti ai seguenti caratteri interi}, fino ad un certo punto indipendenti:

\textit{a}) il \textit{genere geometrico} \(p_g\), genere del fascio di curve ellittiche \((K)\);

\textit{b}) il \textit{determinante} \(n\), numero delle intersezioni di una curva \(K\) con una curva del fascio ellittico \((C)\) (costituito dalle curve direttrici delle \(K\));

\textit{c}) il numero delle \textit{curve ellittiche multiple} (\(K=sK_s\)) che appartengono al \textit{fascio} \((K)\) e i relativi ordini di molteplicità \((s_i)\): come per \(p_g=0\) si può mettere in relazione questi numeri coi plurigeneri della superficie% EDITORIAL: numero originale fissato esplicitamente.
\origfootnote{1}{s230R-1}{Qui si presentano interessanti problemi classificatori: per esempio « determinare le superficie ellittiche corrispondenti ai primi valori di \(\pi\) ». Per \(\pi=2\), \(p_g=0\), citiamo la Memoria di L. \textsc{Campedelli}, \textit{Intorno alle superficie ellittiche con un fascio di curve di genere due.} Rendic. Seminario Mat. di Padova, 1935.}.

\textit{d}) il genere \(\pi\) delle curve del fascio ellittico \((C)\): \origfootmark{1}{s230R-1}
% The same printed footnote marker (1) occurs again here.
\[
\pi=1+n(p_g-1)+\Sigma\frac{(s_i-1)n}{2s_i};
\]

\textit{Secondo caso.} – Dopo aver esaurito la discussione del caso 1), esaminiamo ora il caso 2), ove si tratta delle superficie che ammettono un gruppo continuo, \(\infty^2\), \(\Gamma\) di trasformazioni in sè.

Anzitutto rileviamo che, ove le trasformazioni di \(\Gamma\) non siano tutte permutabili fra loro, quelle trasformazioni che sono permutabili con una data formeranno un sottogruppo algebrico \(\infty^1\) di \(\Gamma\), e quindi si sarà ricondotti a particolarizzazioni del caso 1).

Pertanto si può supporre che il \(\Gamma\) sia un gruppo \(\infty^2\), transitivo su \(F\), costituito da trasformazioni permutabili. Ed ancora è lecito aggiungere che esso sia transitivo \textit{senza eccezio. i} sopra la superficie \(F\) priva di curve eccezionali: altrimenti si avrebbe su \(F\) almeno un punto unito \(O\) o una curva unita \(\chi\) e, fissando un punto vicino ad \(O\) ovvero un punto della curva \(\chi\), si otterrebbe entro \(\Gamma\) un sottogruppo algebrico \(\infty^1\) con un punto unito, onde la \(F\) potrebbe trasformarsi in una rigata (§ 2). % [sic] eccezio. i -> eccezioni, the printed word has a point and a gap where n is expected.

Infine è agevole riconoscere che il nostro \(\Gamma\) opera in modo \textit{semplicemente} transitivo sui punti della \(F\), cioè che non può aversi un
%CONT>>
% scan: s231L
% !PRINTED 436
\weberpage{436}
% head: CAPITOLO UNDECIMO
%<<CONT
% stamp: Scattered scan speckling and a dark band along the right edge.
numero finito $n>1$ di trasformazioni di $\Gamma$ in cui ad un punto $A$ corrisponda un medesimo punto $A'$ della superficie. Invero questa ipotesi si riduce all’assurdo come segue.

Ammesso che l’ipotesi sia verificata, si consideri, accanto alla $F$, la superficie $F'$ (che offre la rappresentazione \textit{parametrica} del gruppo $\Gamma$ secondo S. \textsc{Lie}) i cui punti rispondono, senza eccezione, agli elementi (trasformazioni) di $\Gamma$; le trasformazioni di $\Gamma$ operando le une sulle altre per moltiplicazione, il $\Gamma$ opererà su $F'$ in modo semplicemente transitivo: la $F'$ è dunque una superficie iperellittica, cioè abeliana (Cap. IX, \S~9). D’altronde vi è fra $F$ ed $F'$ una corrispondenza $(1,n)$ in cui ai punti di $F$ rispondono su $F'$ i gruppi d’una involuzione $I_n$; e così la $F'$ è rappresentata sulla $F$ $n$-pla senza punti di diramazione. Segue di qui che l’involuzione $I_n$ su $F'$ è generata da un gruppo finito abeliano $G_n$ di trasformazioni del $\Gamma$, invariante per le $\infty^2$ trasformazioni del $\Gamma$ stesso; e pertanto alle $n$ trasformazioni del $\Gamma$ che fanno corrispondere due gruppi dell’involuzione $I_n$, risponderà una sola trasformazione di $F$ in cui si corrispondono i punti omologhi $A$ ed $A'$, contro l’ipotesi, da cui siamo partiti, che avessero $n$ trasformazioni del $\Gamma$ portanti $A$ in $A'$.

In conclusione possiam dire che \textit{le superficie $F$, non riferibili a rigate, che ammettono un gruppo $\infty^2$ di trasformazioni birazionali} (caso 2) sono \textit{superficie iperellittiche} cioè varietà a due dimensioni \textit{abeliane}, secondo la definizione del Cap. IX, \S~9.

Di queste superficie abbiamo già valutato il genere (Cap. IX, \S~10) che è
\[
p_g=1,
\]
ed è agevole trovare gli altri caratteri.

Anzitutto la \textit{superficie iperellittica $F$} (priva di curve eccezionali) \textit{non può possedere curve canoniche $K$ d’ordine $>0$.}

Infatti il gruppo $\infty^2$ delle trasformazioni di $F$, lasciando ferma la $K$, non porterebbe più i punti di essa in altri punti della superficie, avendosi così un’eccezione alla transitività del $\Gamma$. Qui giova anche avvertire che, a priori, la $F'$ non può possedere un sistema canonico di dimensione $p_g-1>0$, perchè, questo sistema lineare essendo trasformato in sè dalle trasformazioni del $\Gamma$, la $F$ stessa resulterebbe riducibile ad una rigata (\S~2).

Ora, la superficie $F$ possedendo una curva canonica d’ordine zero, il suo genere lineare (assoluto) sarà
\[
p^{(1)}=1.
\]

Resta da valutarne il genere numerico $p_a$. A tal uopo converrà dimostrare che l’irregolarità di $F$ vale
\[
q=p_g-p_a=2,
\]
sicchè
\[
p_a=-1.
\]

% scan: s231R
% !PRINTED 437
\weberpage{437}
% head: CLASSIFICAZIONE GENERALE DELLE SUPERFICIE
% stamp: Scattered scan speckling and faint ink smears.
Cominciamo dal dimostrare che « la varietà di Jacobi $V_p$, corrispondente ad una curva $C$ di genere $p$ che appare subito avere una \textit{irregolarità} $\geq 0$, ha precisamente l’irregolarità \textit{$p$ e non superiore}, cioè che non può aversi in $V_p$ una serie $\infty^r$ con $r>p$ di varietà $V_{p-1}$ disequivalenti ».

Questa proposizione consegue dal terzo criterio d’equivalenza (Cap. III, \S~13) che dalle superficie si estende agevolmente alle varietà a più dimensioni. In base a questo criterio è lecito affermare che due varietà $V_{p-1}$ appartenenti alla $V_p$, che sieno disequivalenti, dovranno segare $V_{p-2}$ disequivalenti sopra la $V_{p-1}$ che risponde alla serie dei gruppi di $p$ punti di $C$ con un punto fisso $A_1$, e quindi ancora $V_{p-3}$ disequivalenti sulla $V_{p-2}$ che risponde ai gruppi di $p$ punti di $C$ con due punti fissi $A_1$ e $A_2$ ecc.

Pertanto si dimostrerà la proposizione enunciata ove si riconosca che la $V_2$, rispondente alla serie dei gruppi di $p$ punti di $C$ con $p-2$ punti fissi $A_1 A_2 \ldotp\ldotp\ldotp\ldotp A_{p-2}$, cioè la superficie rappresentativa delle coppie di punti della curva $C$ di genere $p$, ha proprio l’irregolarità $p$ e non $>p$.

E per ciò basta osservare che una serie $\infty^r$ di curve disequivalenti su tale superficie, dà luogo sopra la $C$ ad una serie $\infty^r$ di gruppi di punti, disequivalenti, sicchè si ha certo
\[
r\leq p \qquad\qquad \text{c. d. d.}
\]

Una volta stabilita la proposizione sulla irregolarità di una varietà di Jacobi $V_p$, si deduce che « una varietà abeliana a $q$ dimensioni, $V_q$, contenuta imprimitivamente nella varietà di Jacobi $V_p$ $(q<p)$, cioè trasformata in sè da un sottogruppo $\infty^p$ del gruppo $\infty^q$ della $V_p$, ha pure l’irregolarità $q$ ». % [sic] subgroup \infty^p and group \infty^q -> subgroup \infty^q and group \infty^p, the dimensions are reversed relative to the surrounding notation.

Siccome una $V_{q-1}$ di $V_q$ è trasformata dal relativo gruppo abeliano $\Gamma_q$ nelle $V_{q-1}$ di una schiera $\infty^q$, l’irregolarità di $V_q$ è certo $\geq q$, e basterà dimostrare che non può essere $>q$.

Invero, se una $V_p$ contiene una siffatta $V_q$, sicchè il gruppo $\Gamma_p$ delle $\infty^p$ trasformazioni permutabili della $V_p$ contenga un sottogruppo algebrico $\Gamma_q$, $\infty^q$, che trasformi in sè le varietà $V_q$ di una schiera $\infty^{p-q}$ d’imprimitività, il $\Gamma_p$ conterrà altresì un sottogruppo algebrico $\Gamma_{p-q}$, $\infty^{p-q}$, che lascerà invariate le $V_{p-q}$ di una schiera $\infty^q$\origfootnote{1}{s231R-1}{Questo teorema (già enunciato in Cap. IX, \S~9) deriva, come si è detto, dall’interpretazione geometrica di un teorema di \textsc{Picard-Poincaré} relativo agli integrali abeliani riducibili sopra una curva, interpretazione dovuta a \textsc{Castelnuovo}. Cfr. \textsc{Enriques-Chisini}, \textit{Lezioni}. L. VI. Cap. III. \S~38 (vol. IV, pag. 237-241).}.

Ora, se una delle anzidette $V_q$ entro $V_p$ possiede un sistema continuo $\infty^r$ di $V_{q-1}$ disequivalenti, queste, mediante le $\infty^{p-q}$ trasfor%
%CONT>>
% scan: s232L
% !PRINTED 438
\unskip\weberpage{438}%
% head: CAPITOLO UNDECIMO
%<<CONT
% stamp: Scattered scan speckling and a dark band along the right edge.
mazioni del $\Gamma_{p-q}$, genereranno un sistema continuo di dimensione $r+p-q$ di $V_{p-1}$ disequivalenti, sicchè si avrà
\[
\begin{aligned}
r+p-q&\leq p,\\
r&\leq q.
\end{aligned}
\]

Dunque: la $V_q$, contenuta in $V_p$, ha l’irregolarità $q$ e non $>q$.

In particolare per $q=2$, una $V_2$ abeliana, ossia una superficie iperellittica $F$ che appartenga imprimitivamente ad una $V_p$ di Jacobi $(p\geq 2)$, essendo trasformata in sè transitivamente da $\infty^2$ trasformazioni del relativo $\Gamma_p$, avrà l’irregolarità
\[
p_g-p_a=2
\]
e quindi
\[
p_a=-1.
\]

Ma (Cap. X, \S~9) qualunque superficie iperellittica $F'$ si può sempre ridurre ad una tale $F$, semplice o multipla, contenuta imprimitivamente in $V_p$, e se si ha una $F$ multipla secondo un certo numero $n>1$, questa è, in ogni caso, priva di curva di diramazione; pertanto si può valutare il genere numerico di $F'$ adoprando la formula di Severi che lega i generi numerici $p_a=-1$ e $p'_a$ delle due superficie $F$ e $F'$ in corrispondenza $(1,n)$ (Cap. V; \S~28); mancando la curva di diramazione della $F$ multipla, codesta formula si scrive % ?? uncertain: punctuation after “Cap. V” -- semicolon or comma.
% EDITORIAL D0232 RISOLTO: The punctuation has an upper point and lower comma tail: semicolon.
\[
24(p'_a+1)=24n(p_a+1),
\]
e dà quindi
\[
p'_a=p_a=-1.
\]

In conclusione \textit{una superficie iperellittica ha i caratteri}
\[
p_g=1,\qquad p^{(1)}=1,\qquad p_a=-1.
\]

Inoltre, possedendo soltanto una curva canonica d’ordine zero i suoi \textit{plurigeneri saranno tutti eguali ad uno.}

\edsection{4. Caratterizzazione delle superficie ellittiche e iperellittiche mediante i valori dei generi.}{4. Caratterizzazione delle superficie ellittiche e iperellittiche mediante i valori dei generi.}{sec.XI.4}

I valori dei generi $(p_a,p_g,p^{(1)},P_4)$, che abbiamo trovato appartenere alle superficie ellittiche e iperellittiche, permettono reciprocamente di definire queste superficie (non riferibili a rigate) e quindi l’intera famiglia delle superficie possedenti un gruppo continuo di trasformazioni birazionali in sè stesse.

Dimostriamo dapprima che \textit{una superficie $F$ coi generi}
\[
p_a=-1,\qquad p_g>1,\qquad p^{(1)}=1,
\]
\textit{è una superficie ellittica.}

% scan: s232R
% !PRINTED 439
\weberpage{439}
% head: CLASSIFICAZIONE GENERALE DELLE SUPERFICIE
% stamp: Scattered scan speckling and a narrow dark mark at the right edge.
Infatti le $\infty^{p_g-1}$ curve canoniche di $F$ saranno composte colle curve ellittiche $K$ d’un fascio $(K)$, avente un certo genere $\varrho$; e a priori le $\infty^q$ $(q=p_g-p_a)$ curve $L$ disequivalenti d’un sistema completo $\infty^q$, costruito su $F$, segheranno sulle $K$ $\infty^q$ gruppi equivalenti ovvero appartenenti ad una serie formata di $\infty^1$ serie lineari: nel primo caso le $L$ di $\{L\}$ si ottengono l’una dall’altra per somma e sottrazione di curve $K$ di $(K)$ e si deduce
\[
\varrho=q=p_g+1;
\]
nel secondo caso il sistema $\{L\}$ viene formato da $\infty^1$ sistemi continui $\infty^{q-1}$, ciascuno dei quali si deduce da una sua curva sommando e sottraendo curve $K$ di $(K)$, sicchè
\[
\varrho=p_g;
\]
e oltre $(K)$ --- come più volte abbiamo fatto innanzi --- si costruisce su $F$ un secondo fascio irrazionale, ellittico, di curve direttrici delle $K$. Di conseguenza la superficie $F$ risulta essere ellittica (Cap. X, \S~12).

Ora vediamo che proprio si avvera questo secondo caso e non il primo.

A tal uopo calcoliamo l’invariante di Zeuthen-Segre della $F$ in ordine al fascio di curve ellittiche $K$: si avrà
\[
\Delta-4=12p_a-p^{(1)}+9=-4,
\]
e perciò
\[
\Delta=0;
\]
onde si trae che le $K$ di $(K)$, non acquistando mai punti doppi, sono fra loro birazionalmente identiche; in $(K)$ si troverà un certo numero di curve ellittiche multiple collo stesso modulo
\[
s_1K_1,\quad s_2K_2,\quad \ldotp\ldotp\ldotp\ldotp
\]
dove
\[
s_i>1.
\]

Ciò posto, il procedimento spiegato nel \S~11 del Cap. X (pel caso d’un fascio di genere $\varrho=1$) permette di costruire un fascio (a priori razionale o ellittico) di curve direttrici $C$, di genere $\pi>\varrho$, senza punti base. Sopra una $C$ le curve canoniche di $F$ segheranno la serie canonica $g_{2\pi-2}$ (in generale non completa). Siccome le $K$ segano su una $C$ un’involuzione di genere $\varrho$, che ha per punti $s_i$-pli le intersezioni colle $K_i$, le curve canoniche di $F$ saranno formate dai gruppi della $g^{\varrho-1}_{2\varrho-2}$ entro il fascio $(K)$ di genere $\varrho$, cui si sommino le curve $(s_i-1)K_i$. E da ciò risulta il genere di $F$
\[
p_g=\varrho \qquad\qquad \text{c. d. d.}
\]

% scan: s233L
% !PRINTED 440
\weberpage{440}
% head: CAPITOLO UNDECIMO
% stamp: An illegible inverted ink fragment above “si distribuiscono”; scattered speckling and vertical edge marks.
Anche per $p_g=0$ e $p_a=-1$, $p^{(1)}=1$, come per $p_g>1$ e $p^{(1)}=1$, la superficie $F$ risulta ellittica (Cap. X, \S~12) salvo che essa può essere ellittica impropria, riducendosi ad una rigata ellittica $(P_{12} = 0)$\origfootnote{1}{s233L-1}{Cfr. X, 14.}. % The equality sign is printed on both sides of the original line break.

Ma il caso $p_g=1$ dà luogo ad una eccezione: le superficie i cui generi valgono
\[
p_a=-1,\qquad p_g=1,\qquad p^{(1)}=1,
\]
si distribuiscono in due famiglie:

1) superficie con curva canonica ellittica d’ordine maggiore di zero;

2) superficie con curva canonica pura d’ordine zero.

Le superficie della prima famiglia sono ellittiche, cioè ammettono un gruppo ellittico $\infty^1$ di trasformazioni in sè, mentre le superficie della seconda famiglia sono iperellittiche, e non in generale ellittiche: vuol dire che esse ammettono un gruppo abeliano $\infty^2$, $\Gamma$, di trasformazioni in sè, ma soltanto per moduli particolari esiste in $\Gamma$ un sottogruppo algebrico (ellittico) $\infty^1$.

Per giustificare ciò che si è detto, intorno ai casi 1) e 2), conviene stabilire un

\textit{Lemma sulla superficie di irregolarità $q=2$:} « una superficie $F$ di irregolarità $q=p_g-p_a=2$, non contenente un fascio di genere due di curve, si lascia rappresentare sopra una superficie iperellittica $\Phi$, semplice o multipla ».

Per dimostrare questo lemma si consideri su $F$ un sistema continuo (semplice) $\{C\}$, formato di $\infty^2$ curve disequivalenti, e si ritengano gli elementi (curve) di $\{C\}$ come « punti » di una superficie $\Phi$: la quale (come varietà di Picard della $F$) sarà una superficie iperellittica. Fra le superficie $F$ e $\Phi$ intercederà una corrispondenza, in cui ai punti di $F$ (ossia ai sistemi $\infty^1$ di curve $C$ passanti per essi) risponderanno le curve $L$ di un sistema $\infty^2$ $\{L\}$ sopra $F$. Diciamo che le $\infty^2$ $L$ su $\Phi$ non possono essere tra loro equivalenti. Altrimenti le $\infty^1$ $C$ di $\{C\}$ passanti per un punto $O$ di $F$ segherebbero sopra un’altra $C=\bar{C}$ gruppi equivalenti e, per il terzo criterio d’equivalenza (Cap. III, \S~13), risulterebbero fra loro equivalenti. % [sic] sopra F -> sopra \Phi, the following sentence explicitly places the curves L on \Phi.

Per giustificare l’asserto conviene considerare sulla superficie $F$ la serie dei gruppi di punti $G$ sezioni di una particolare $C=\bar{C}$, non passante per $O$, colle $\infty^1$ $C$ per $O$; se si designa con $i$ l’indice di codesta serie, si avranno $i$ gruppi $G$ aventi comune un punto $P$ di $\bar{C}$ e $i$ gruppi $G$ aventi a comune un altro punto $P'$ della stessa curva; e in virtù di un criterio di \textsc{Severi-Castelnuovo}\origfootnote{2}{s233L-2}{Cfr. \textsc{Enriques-Chisini}, \textit{Lezioni}. Libro V, cap. IV, \S~42 (vol. III, pag. 483)-}, occorre % ?? uncertain: final punctuation in the second footnote -- hyphen or damaged full stop.
%CONT>>
% scan: s233R
% !PRINTED 441
\weberpage{441}
% head: CLASSIFICAZIONE GENERALE DELLE SUPERFICIE
%<<CONT
% stamp: Scattered scan speckling and small diagonal ink marks among the first lines.
e basta dimostrare che i due insiemi di gruppi $G$, $P(G)$ e $P'(G)$, così definiti sulla $\bar{C}$ sono equivalenti. A tal uopo si consideri su $\Phi$ la curva $\bar{L}$ che risponde al punto $O$ (o meglio alla serie $C(O)$ delle $C$ per $O$) e il punto $A$ che risponde alla curva $\bar{C}$: se, come supponiamo, le $\infty^1$ $L$ per $A$ sono equivalenti su $\Phi$, esse segheranno su $\bar{L}$ gruppi di punti equivalenti e così in ispecie saranno equivalenti, sopra la detta curva, i gruppi sezioni delle due $L$ --- diciamo $L_P$ e $L_{I'}$ --- corrispondenti ai punti $P$ e $P'$ di $F$. Ma --- ritornando alla superficie $F$ --- ciò significa che i due gruppi di curve $C$ per $O$ e $P$ e per $O$ e $P'$ secanti su $\bar{C}$ i gruppi di punti $P(G)$ e $P'(G)$, sono equivalenti \textit{entro la serie delle $C$ per $O$} e così fanno parte di una serie razionale $\infty^1$ di gruppi di curve $C$; ricordando che una serie razionale di gruppi di punti sopra una curva è sempre costituita di gruppi equivalenti\origfootnote{1}{s233R-1}{Cfr. \textsc{Enriques-Chisini}, \textit{Lezioni}. Libro V, cap. I, \S~10, Nota (vol. III, pag. 78).}, si deduce che $P(G)$ e $P'(G)$, sezioni di $\bar{C}$ con codesti gruppi di $C$, sono equivalenti \textit{sopra la detta $\bar{C}$.} % ?? uncertain: subscript of the second L after “diciamo” -- P' or I'.
% EDITORIAL D0234 RISOLTO: The subscript is an I with a stem and serifs, with no bowl of P, followed by a prime.

Ora, avendo dimostrato che le $\infty^2$ curve $L$ di $\Phi$ (corrispondenti ai punti di $F$) non possono essere tutte equivalenti, dobbiamo ancora esaminare l’ipotesi che ogni $L$ di $\{L\}$ faccia parte di una serie $\infty^1$ di curve $L$ equivalenti, e riconoscere che, in tale ipotesi, la superficie $F$ contiene un fascio di genere due. Invero nella detta ipotesi si troverà su $F$ un fascio di curve $K$ ai cui punti rispondono curve $L$ equivalenti, e si può stabilire che le $C$ di una serie $C(O)$, passanti per un punto qualunque $O$ della superficie, e quindi tutte le $C$, segano su una $K$ gruppi equivalenti: a tal uopo basta ripetere il ragionamento fatto innanzi, mostrando che sono equivalenti (entro la serie $C(O)$) i gruppi di curve comuni a $C(O)\ C(P)$ e $C(O)\ C(P')$, e quindi sono anche equivalenti, sulla curva $K$, i gruppi di punti $P(G)$ e $P'(G)$ sezioni della $K$ colle $C$ di $C(O)$.

In tal guisa si verifica che le $\infty^2$ curve $C$ di $F$, fra loro disequivalenti, segano su ogni $K$ gruppi equivalenti; quindi il sistema continuo $\infty^2$ $\{C\}$ si costruisce a partire da un sistema lineare $|C|$, o da una curva $C$, sommando e sottraendo curve $K$ del fascio $(K)$: il quale risulta perciò di genere ($\varrho\geq 2$ e quindi) $p_g-p_a=2$.

Così è giustificato l’enunciato \textit{Lemma}.

Ora si può aggiungere che: \textit{le superficie d’irregolarità $q=2$, con $p_g>0$, contenenti un fascio di genere due di curve} (di genere $\pi>0$) \textit{hanno il $p_a\geq 0$.}

Infatti se sia data una superficie $F$, di genere numerico $p_a=-1$, la quale contenga un fascio di genere due di curve $K$, dimostriamo che si ha per essa: $p_g>1$.

% scan: s234L
% !PRINTED 442
\weberpage{442}
% head: CAPITOLO UNDECIMO
% stamp: Scattered scan speckling and a long dark vertical mark at the right edge.
Calcolando l’invariante di Zeuthen-Segre in relazione al fascio $(K)$, si trova
\[
\Delta-4+8(\pi-1)=12p_a-p^{(1)}+9=-3-p^{(1)}
\]
e quindi
\[
p^{(1)}=1,\qquad \Delta=0,\qquad \pi=1;
\]
perciò le curve ellittiche del fascio $(K)$ risultano fra loro birazionalmente identiche. Allora (come in Cap. X, \S~11) si può costruire un fascio (che a priori sarebbe razionale o ellittico e risulterà poi ellittico) di curve $C$, senza punti base, direttrici delle $K$: quindi\origfootnote{1}{s234L-1}{. Cfr. l’osservazione al \S~8 del Cap. IX.} le curve canoniche di $F$ saranno composte coi gruppi canonici delle $K$ di $(K)$ ed eventualmente anche da parti fisse (in modo più preciso da $(s-1)C_s$ dove $sC_s=C$ sia una curva multipla $s$-pla del fascio $(C)$); si deduce % ?? uncertain: dot immediately after the footnote number -- printed full stop or scan speck.
\[
p_g\geq 2 \qquad\qquad \text{c. d. d.}
\]

In base al lemma sulle superficie d’irregolarità $q=2$, e all’osservazione che le superficie con $q=2$ contenenti un fascio di genere due hanno il $p_g\geq 2$, e quindi $p_a\geq 0$, possiamo dire che, una superficie $F$ con $p_a=-1$ e $p_g=1$ si lascia rappresentare sopra una superficie iperellittica $\Phi$, semplice o multipla, che avrà una certa curva di diramazione $D$. Alla $D$ risponde su $F$ la curva canonica; perciò, se il genere lineare di $F$ vale $p^{(1)}=1$, si avverrà uno dei due casi seguenti, che ci riporteranno ai casi 1) e 2) menzionati innanzi:

1) la $D$ è una curva ellittica d’ordine maggiore di zero, ovvero

2) la $D$ è una curva d’ordine zero.

1) Nel primo caso le $\infty^2$ trasformazioni in sè della superficie iperellittica $\Phi$, appartenenti al gruppo continuo $\Gamma$, mutano la $D$ (di grado zero) nelle curve d’un fascio, e ciascuna $D$ di questo (o ciascuna componente irriducibile della $D$) è invariante per un sottogruppo ellittico $\infty^1$ del $\Gamma$. Segue di qui che la superficie $\Phi$ multipla ammette un gruppo ellittico $\infty^1$ di trasformazioni, ossia è una superficie ellittica, c. d. d.

2) Se invece sopra la superficie multipla $\Phi$ (privata di curve eccezionali) non vi è curva di diramazione, la $F$ è, come la $\Phi$, una superficie abeliana, ossia iperellittica (Cap. IX, \S~9); e si avvera dunque il caso 2).

Si possono distinguere i due casi 1) e 2) osservando che, nel caso 1) (per essere $p_g=1$ e non $p_g>1$) alle curve ellittiche del fascio $(D)$ di $\Phi$, o meglio alle loro componenti irriducibili, rispondono su
%CONT>>
% scan: s234R
% !PRINTED 443
\weberpage{443}
% head: CLASSIFICAZIONE GENERALE DELLE SUPERFICIE
%<<CONT
% stamp: Scattered scan speckling and a short dark mark at the right edge.
$F$ le curve $K$ d’un fascio ellittico di curve ellittiche, in cui il numero delle curve dotate di punto doppio vale $\Delta=0$.

Pertanto in questo fascio $(K)$, formato di curve birazionalmente identiche, vi sarà tutt’al più un certo numero di curve multiple $D_i$ da contare ciascuna $s_i$ $(>1)$ volte, e la particolare curva $D$ che costituisce la curva di diramazione della $\Phi$ multipla sarà composta di un certo numero $r$ di codeste curve:
\[
D=\Sigma(s_i-1)D_i.
\]

Segue di qui che nel fascio $(K)$ sopra $F$ la curva canonica, rispondente alla $D$, è
\[
\Sigma\frac{s_i-1}{s_i}K.
\]

Quindi, se $r>1$, si avrà su $F$ un sistema bicanonico di dimensione $\geq r-1$:
\[
|rK+\Sigma\frac{s_i-2}{s_i}K|,
\]
aventi come componenti fisse le curve che rispondono alle $D_i$ di molteplicità $s_i>2$, contate $s_i-2$ volte; e, in ogni caso, ove sia $r=1$, si avrà su $F$ un sistema quadricanonico $\infty^1$ almeno, di cui le componenti variabili sono coppie di curve $K$, formanti una $g^1_2$ entro il $(K)$ ellittico. % [sic] aventi -> avente, apparent agreement error with the singular “un sistema bicanonico”.

In conclusione, nel caso 1), la superficie ellittica $F$ $(p_a=-1,\ p_g=1,\ p^{(1)}=1)$ ha il quadrigenere
\[
P_4>1,
\]
mentre la superficie iperellittica (cogli stessi $p_a=-1$, $p_g=1$ e $p^{(1)}=1$, ma con curva canonica d’ordine zero) avrà tutti i plurigeneri eguali ad uno e, in particolare il \textit{quadrigenere}:
\[
P_4=1.
\]

Riassumiamo i risultati ottenuti:

\textit{Le superficie, non riferibili a rigate} ($P_4+P_6>0$ o $P_{12}>0$), \textit{che ammettono un gruppo continuo di trasformazioni birazionali in sè stesse, sono caratterizzate globalmente dai valori dei generi:}
\[
p_a=-1,\qquad p^{(1)}=1;
\]
\textit{esse sono sempre o superficie ellittiche (con un gruppo $\infty^1$) e $p_g>1$ ovvero $p_g=1$ e $P_4>1$; o superficie iperellittiche (con un gruppo $\infty^2$) se}
\[
p_g=P_4=1.
\]

% scan: s235L
% !PRINTED 444
\weberpage{444}
% head: CAPITOLO UNDECIMO
% stamp: Scattered scan speckling, diagonal ink marks, and a long dark vertical mark at the right edge.
Che le superficie iperellittiche non siano in generale anche ellittiche, cioè che il loro gruppo $\infty^2$ non contenga generalmente alcun sottogruppo $\infty^1$ algebrico (ellittico), si può vedere riferendosi alla superficie di Jacobi che risponde ad una curva $C$ di genere due (superficie $F$ rappresentativa delle coppie di punti di $C$), o similmente alle involuzioni abeliane sopra $F$ o alle $F$ multiple senza curva di diramazione.

La $F$ iperellittica così costruita non contiene curve ellittiche, ove la $C$ non contenga una involuzione ellittica $\gamma_n^1$, d’un certo ordine $n>1$, ovvero degeneri in una coppia di curve ellittiche; nel caso particolare in cui la detta $F$ contenga una curva ellittica $C$, questa sarà portata dalle trasformazioni del gruppo $\Gamma$ di $F$ nelle curve d’un fascio ellittico $(C)$ e (d’accordo col teorema sugli integrali riducibili di Picard-Poincaré, interpretato geometricamente da \textsc{Castelnuovo}) vi sarà su $F$ un secondo fascio ellittico $(K)$ pure di curve ellittiche: le-quali si costruiscono come direttrici del fascio $(C)$ nel modo indicato nel \S~11 del Cap. X. Pertanto la $F$ si potrà ritenere in due modi diversi come superficie ellittica, ammettendo due gruppi ellittici $\infty^1$ di trasformazioni, sottogruppi del $\Gamma$ $\infty^2$, che hanno come traiettorie le $C$ e le $K$. Infine aggiungiamo che la detta $F$, due volte ellittica, di determinante $n>1$ (ove le $C$ e le $K$ si seghino in $n$ punti), si lascia rappresentare sopra la particolare superficie iperellittica (e due volte ellittica di determinante \textit{uno}) che è la superficie di Jacobi corrispondente alla curva di genere due degenere in due curve ellittiche, superficie da prendersi come multipla secondo $n$. % ?? uncertain: mark between “le” and “quali” -- printed hyphen or scan artefact; mark after “si” -- ink artefact or apostrophe.

\textit{Osservazione.} -- Le condizioni con cui abbiamo definito la famiglia delle superficie ellittiche ed iperellittiche ($p_a=-1\quad p^{(1)}=1$) contengono qualcosa di sovrabbondante:

\begin{center}
\textit{la condizione $p^{(1)}=1$ è conseguenza dalla $p_a=-1$.} % [sic] dalla -> della, apparent preposition error after “conseguenza”.
\end{center}

Dimostreremo ciò, per le superficie ellittiche con $p_g$ qualunque, nel \S~8; qui l’affermazione verrà giustificata per quanto concerne le superficie col $p_g=1$, cioè essenzialmente pel caso iperellittico.

Partiamo dalla conclusione acquisita innanzi che: « le superficie $F$ coi generi $p_a=-1$ e $p_g=1$ si lasciano rappresentare sopra una superficie iperellittica $\Phi$, semplice o multipla ». Se si ha una $\Phi$ multipla, vi sarà su questa una curva di diramazione $D$ di genere $p^{(1)}$ eguale al genere lineare di $F$, e si tratta di accertare che dall’ipotesi $p^{(1)}>1$ segue
\[
p_a\geq 0.
\]

Dimostriamo la cosa nell’ipotesi più semplice in cui la curva di diramazione $D$ della $\Phi$ sia una curva semplice, di genere $\varrho=p^{(1)}$, dotata di un certo numero $\delta$ di nodi e di un certo numero $\tau$ di cu%
%CONT>>
% scan: s235R
% !PRINTED 445
\unskip\weberpage{445}%
% head: CLASSIFICAZIONE GENERALE DELLE SUPERFICIE
%<<CONT
% stamp: Scattered scan speckling and small marginal ink marks.
spidi, caso a cui si riferisce la formula di Severi\origfootnote{1}{s235R-1}{Per completare la dimostrazione occorre dunque estendere questa formula al caso in cui la detta curva di diramazione sia comunque composta di parti multiple.

Che in tal guisa si pervenga al risultato enunciato siamo accertati \textit{a priori} da una considerazione (che riposa in parte sulla teoria degli integrali semplici di prima specie appartenenti ad una superficie), quale ci viene gentilmente comunicata da G. \textsc{Castelnuovo}. Ed è che, per le superficie irregolari non riferibili a rigate, designando con $I$ l’invariante di Zeuthen-Segre, si ha
\[
I+4=13+12p_a-p^{(1)}\geq 0.
\]

Infatti per $p_a<0$ (cioè $p_a=-1$) segue
\[
p^{(1)}=1.
\]
} che lega i generi numerici $p'_a$ e $p_a$ delle due superficie $\Phi$ ed $F$ in corrispondenza $(1,n)$ (cfr. Cap. V, \S~8). La detta formula ci dà
\[
24(p_a+1)\geq 24n(p'_a+1)+6(\varrho-1)-2\tau,\origfootmark{2}{s235R-2}
\]\origfoottext{2}{s235R-2}{Il caso in cui vale la diseguaglianza risponde all’ipotesi in cui si abbiano sù $D$ dei punti fondamentali per la corrispondenza $(1,n)$.}
dove
\[
p'_a=-1,\qquad \varrho>1,
\]
cioè
\[
24(p_a+1)\geq 6(\varrho-1)-2\tau.
\]

\edFlush
Ma, sopra la superficie iperellittica $\Phi$ la $D$ appartiene almeno ad un sistema $\infty^2$, e possiede una serie caratteristica (segata dalle curve infinitamente vicine) di grado % ?? uncertain: mark after “grado” -- scan speck or full stop.
\[
2\delta+3\tau\leq 2(\varrho+\delta+\tau-1).
\]

Pertanto si avrà
\[
\tau\leq 2(\varrho-1), % ?? uncertain: punctuation after 2(\varrho-1) -- comma or scan speck.
\]
e
\[
24(p_a+1)>0,
\]
cioè
\[
p_a\geq 0. \qquad\qquad \text{c. d. d.}
\]

Riassumendo: \textit{Le superficie coi generi}
\[
p_a=-1\qquad \text{e}\qquad p_g=1
\]
\textit{hanno il genere lineare}
\[
p^{(1)}=1
\]
\textit{e, dove posseggano una curva canonica d’ordine zero (ciò che accade per $P_4=1$) sono superficie iperellittiche.}

\textit{Notizia.} -- L’analisi delle superficie, non riferibili a rigate, che ammettono un gruppo continuo di trasformazioni in se stesse, s’ini% % ?? uncertain: accent on “se” in “in se stesse” -- se or sè.
% EDITORIAL D0239 RISOLTO: No accent is attached to e; the isolated higher speck is too far above the vowel and irregular in form.
%CONT>>
% scan: s236L
% !PRINTED 446
\unskip\weberpage{446}%
% head: CAPITOLO UNDECIMO
%<<CONT
% stamp: Dark vertical scan-edge line at the right; scattered specks and short marginal marks.
zia collo studio delle superficie iperellittiche fatto da \textsc{Picard} sotto l’aspetto trascendente\origfootnote{1}{s236L-1}{E. \textsc{Picard}, \textit{Mémoire sur la théorie des fonctions algébriques de deux variables indépendantes.} Journal de Mathématiques, 1889. (Cap. III). Cfr. il t. II del Traité di \textsc{Picard} e \textsc{Simart}, Parigi 1906.}.

Una brillante applicazione degli integrali semplici di differenziali totali che egli stesso ha introdotto nella scienza, conduce l’autore a caratterizzare le superficie, della specie indicata, che posseggono un gruppo \(\infty^2\) di trasformazioni permutabili, mercè:

1) l’esistenza di due integrali semplici di prima specie;

2) l’esistenza di un integrale doppio di prima specie (che, secondo \textsc{Clebsch} e \textsc{Noether}, vuol dire \(p_g=1\));

3) e la mancanza di una curva canonica (pura) d’ordine maggiore di zero.

La condizione 3) (che il \textsc{Picard}, ha avvertito ma, ingannato da un errore di \textsc{Noether}, riteneva superflua) porta l’invertibilità dei due integrali semplici e quindi la rappresentazione parametrica della superficie con funzioni (iperellittiche) quattro volte periodiche di due variabili indipendenti. % [sic] Picard, ha -> Picard ha, apparently intrusive comma between subject and verb.

Da questo risultato \textsc{Enriques}\origfootnote{2}{s236L-2}{\textit{Sulle superficie algebriche che ammettono un gruppo continuo di trasformazioni birazionali in se stesse.} Rendic. Circolo Mat. di Palermo, 1905.} ha tratto la caratterizzazione delle superficie iperellittiche mediante i valori dei generi \(p_a=-1\), \(p_g=P_4=1\).

Più tardi \textsc{Severi}\origfootnote{3}{s236L-3}{\textit{Sulle superficie algebriche che ammettono un gruppo continuo permutabile a due parametri di trasformazioni birazionali.} Atti Istituto Veneto, 1908.} ha tradotto la dimostrazione del teorema di Picard in forma geometrica, deducendo dirittamente l’esistenza del gruppo \(\infty^2\) di trasformazioni dalle condizioni \(p_a=-1\), \(p_g=1\) e curva canonica d’ordine zero, come anche qui si è fatto.

Il caso delle superficie con un gruppo ellittico \(\infty^1\), già incontrato da \textsc{Picard}, viene esaurito, sotto l’aspetto trascendente, dall’analisi di \textsc{Painlevé}\origfootnote{4}{s236L-4}{P. \textsc{Painlevé}, \textit{Leçons sur la théorie analytique des équations différentielles.} Stoccolma, 1895 (pag. 270).} che definisce le superficie ellittiche mediante una rappresentazione parametrica con funzioni ellittiche d’un parametro e algebriche d’un altro, e riesce quindi a classificare tutte le superficie, non riferibili a rigate, possedenti un gruppo continuo di trasformazioni in sè. La traduzione geometrica di questo risultato e la classificazione dell’intera famiglia di superficie mediante i valori dei generi, è stata conseguita da \textsc{Enriques} nella citata Memoria del 1905\origfootnote{5}{s236L-5}{\textit{Sulle superficie algebriche che ammettono un gruppo continuo ecc.} l. c. In parte già nella precedente Memoria, pubblicata egualmente nei Rendiconti del Circolo Mat. di Palermo: \textit{Sulle superficie di genere zero} (1905).}.

% scan: s236R
% !PRINTED 447
\weberpage{447}
% head: CLASSIFICAZIONE GENERALE DELLE SUPERFICIE
% stamp: Scattered scan specks and short marginal marks.
\edsection{5. Superficie coi generi \(p_g=1\) e \(p^{(1)}=1\).}{5. Superficie coi generi p\_g=1 e p\textasciicircum{}(1)=1.}{sec.XI.5}

I risultati stabiliti nei precedenti paragrafi ci permettono ora di classificare le superficie di genere geometrico \(p_g=1\), aventi il genere lineare \(p^{(1)}=1\). Queste superficie si possono distribuire in tre famiglie:

1) Superficie regolari di genere \(p_a=p_g=1\) con curva canonica d’ordine zero (caratterizzate dal valore del bigenere \(P_2=1\)), che danno luogo a infinite classi di superficie (con 19 moduli) a sezioni iperpiane canoniche di genere \(\pi\), d’ordine \(n=2\pi-2\), nello spazio \(S_\pi(\pi=2,\ 3\ldotp\ldotp\ldotp\ldotp)\) (Cfr. Cap. VII, \S~2).

2) Superficie iperellittiche coi generi \(p_a=-1\) e \(p_g=1\), possedenti del pari una curva canonica d’ordine zero (caratterizzate dal valore del quadrigenere \(P_4=1\)), che danno luogo ancora ad infinite classi di superficie normali a curve sezioni canoniche di genere \(\pi\), d’ordine \(n=2\pi-2\) nello spazio \(S_{\pi-2}\) (queste classi dipendendo, come diremo, da tre moduli).

3) Superficie possedenti un fascio (razionale o ellittico) di curve ellittiche: questa proprietà appartiene alle superficie 1) e 2) soltanto per valori particolari dei loro moduli.

Per dimostrare che « \textit{le tre famiglie di superficie} 1), 2), 3) \textit{esauriscono le superficie aventi i generi \(p_g=1\) e \(p^{(1)}=1\)} » conviene in parte richiamare osservazioni già fatte innanzi in questo trattato, in parte riconoscere alcuni fatti non ancora osservati. Si tratta delle proposizioni seguenti:

\textit{a}) le superficie di genere geometrico \(p_g=1\) hanno il genere numerico
\[
p_a\ge -1,
\]
cioè
\[
p_a=1 \qquad \text{o} \qquad p_a=0 \qquad \text{o} \qquad p_a=-1;
\]

\textit{b}) le superficie (regolari) per cui \(p_g=p_a=1\) con curva canonica d’ordine zero, sono definite dal valore del bigenere \(P_2=1\) (Cap. VII, \S~2); invece le superficie di questa famiglia che hanno il genere lineare \(p^{(1)}=1\), e posseggono una curva canonica (pura) d’ordine maggiore di zero, avendo il \(P_2>1\), conterranno anche un fascio (lineare) di curve ellittiche, bicanoniche o componenti delle bicanoniche;

\textit{c}) le superficie per cui \(p_g=1\), con curva canonica d’ordine zero hanno il genere numerico \(p_a=1\) ovvero \(p_a=-1\) (Cap. IX, \S~8) e, se \(p_a=-1\) hanno pure il genere lineare \(p^{(1)}=1\), e sono
%CONT>>
% scan: s237L
% !PRINTED 448
\weberpage{448}
% head: CAPITOLO UNDECIMO
%<<CONT
% stamp: Dark vertical scan-edge line at the right; scattered specks and diagonal marks.
superficie iperellittiche o superficie ellittiche, con un fascio di curve ellittiche, secondochè \(P_4=1\) o \(P_4>1\) (\S~4);

\textit{d}) le superficie per cui \(p_g=1\), \(p_a=0\) e \(p^{(1)}=1\) (che per la \textit{b}) debbono contenere una curva canonica d’ordine maggiore di zero) posseggono un fascio ellittico di curve paracanoniche, ed hanno quindi il bigenere \(P_2>1\).

Fra queste proposizioni la \textit{b}) e \textit{c}) richiamano conoscenze già acquisite, perciò basterà dimostrare le \textit{a}) e \textit{d}) che portano qualcosa di nuovo.

La prop. \textit{a}) essendo già stabilita per le superficie ellittiche e iperellittiche, basterà dimostrarla nell’ipotesi che la superficie \(F\) (di genere \(p_g=1\)) non possegga un gruppo continuo di trasformazioni in sè (più avanti, \S~7, questa proposizione apparirà come caso particolare d’un teorema più generale per cui il genere numerico delle superficie non riferibili a rigate vale sempre \(p_a\ge-1\)).

Sia \(F\) una superficie di genere \(p_g=1\), che non ammetta un gruppo continuo di trasformazioni, e sia \(q=p_g-p_a\) la sua irregolarità: si vuol ridurre all’assurdo l’ipotesi
\[
q>2,
\]
cioè
\[
p_a<-1.
\]

A tal uopo, assumendo \(q>2\), si consideri la varietà di Picard \(V_q\) corrispondente ad un sistema continuo \(\infty^q\) di curve \(C\) disequivalenti, preso sulla nostra \(F\). Come abbiamo veduto (Cap. IX, \S~11), ove la \(F\) non contenga un fascio irrazionale di curve \(K\), una serie \(\infty^1\) di curve \(C\) di \(\{C\}\) conduce a costruire una superficie, semplice o multipla, identica ad \(F\) entro la detta \(V_q\). Anzi, in vista di mostrare che \(p_g\ge q\), si può ritenere che si abbia in \(V_q\) una superficie semplice identica ad \(F\): giacchè nel caso di una superficie multipla, richiamando l’\textit{osservazione} del \S~4, si vede che se questa superficie ammette un gruppo continuo di trasformazioni in sè, mentre la \(F\) non possiede un tal gruppo, il genere numerico della \(F\) risulta \(p_a\ge0\), e di conseguenza per la \(F\) medesima, \(p_g\ge q\).

Ordunque si abbia in \(V_q\) una superficie identica ad \(F\) (che torniamo a indicare con \(F\)) la quale non ammetta un gruppo continuo di trasformazioni in sè; mediante le \(\infty^q\) trasformazioni della \(V_q\) la nostra \(F\) viene trasformata in una serie \(\infty^q\) di superficie analoghe, diciamo \(\{F\}\). Per semplicità supponiamo dapprima \(q=3\). Siccome la \(V_q=V_3\) ha il genere geometrico uno, il sistema lineare \(\infty^2\) caratteristico di \(F\), segato sopra una \(F\) dalle superficie infinitamente vicine, sarà costituito di curve speciali (residue dell’intersezione colla varietà canonica di \(V_3\) e, quindi, canoniche), sicchè risulterà per la \(F\)
\[
p_g\ge3. % ?? uncertain: small mark above the p in this display -- scan speck with no prime, or p'_g.
\]

% scan: s237R
% !PRINTED 449
\weberpage{449}
% head: CLASSIFICAZIONE GENERALE DELLE SUPERFICIE
% stamp: Scattered scan specks and marginal marks.
Il ragionamento si estende al caso \(q>3\). In questo caso si assumerà entro la serie \(\{F\}\) una serie \(\infty^{q-3}\) di superficie \(F\), che costituisce una varietà \(V_{q-1}\), e si supporrà in generale che tale varietà sia contenuta in un sistema lineare \(|V_{q-1}|\) di dimensione \(r\) (grande quanto si vuole); quindi si costruirà il sistema continuo \(\infty^{r+q}\) \(\{V_{q-1}\}\) generato da \(|V_{q-1}|\) mediante le \(\infty^q\) trasformazioni della \(V_q\). Questo sistema continuo ha un sistema lineare caratteristico \(\infty^{r+q-1}\) di varietà speciali o canoniche, segato sopra una \(V_{q-1}\) dalle varietà infinitamente vicine; perciò le dette \(V_{q-1}\) hanno il genere
\[
P\ge r+q.
\]

Ora si valuti l’infinità delle \(V_{q-1}\) di \(\{V_{q-1}\}\) che contengono una \(F\): vi sono \(\infty^q F\), e \(\infty^{r+q}V_{q-1}\) contenenti ciascuna \(\infty^{q-3}F\), sicchè la dimensione del sistema \(\{F\}\) risulta (nell’ipotesi più sfavorevole):
\[
r+2q-3-x=q
\]
dove \(x\) indica l’infinità del sistema delle \(V_{q-1}\) contenenti una \(F\): \(x=r+q-3\). Segue di qui che il sistema caratteristico di \(\{V_{q-1}\}\) in una \(V_q\) sega sopra una \(F\) in essa contenuta, un sistema lineare (speciale o canonico) la cui dimensione è almeno:
\[
r+q-1-(x-1)=3.
\]

E pertanto si deduce che il genere di \(F\) vale
\[
p_g\ge4.
\]

In conclusione abbiamo dimostrato che una superficie \(F\) di genere \(p_g<3\), che non contenga un fascio irrazionale di curve, non può avere l’irregolarità
\[
q>2;
\]
in particolare una \(F\) di genere \(p_g=1\) avrà il genere numerico
\[
p_a\ge-1.
\]

La conclusione si estende all’ipotesi che la \(F\) contenga un fascio irrazionale di curve \((C)\), perchè il computo dell’invariante di Zeuthen-Segre dà, in rapporto a questo fascio,
\[
\begin{aligned}
\Delta-4&\ge12p_a-p^{(1)}+9,\\
12(p_a+1)-p^{(1)}&\le0,\\ % [sic] \le0 -> \ge0, apparently reversed inequality in view of the bound asserted on the next line.
p_a&\ge-1
\end{aligned}
\]
e per \(p_a=1\), \(p^{(1)}=1\). % [sic] p_a=1 -> p_a=-1, a minus sign apparently omitted in the limiting case.

Il teorema \textit{d}) si dimostra come segue.

Si consideri sopra \(F\) un sistema lineare (irriducibile) \(|C|\), di un certo genere \(\pi\).

Il suo sistema aggiunto \(|C'|\) avrà la dimensione
\[
p_a+\pi-1=\pi-1,
\]
% EDITORIAL: segnatura originale: \makebox[\textwidth][l]{\textsc{Enriques F.} - \textit{Superficie algebriche.}\hfill 29}
\printersignature{\textsc{Enriques F.} - \textit{Superficie algebriche.}\quad 29}
%CONT>>
% scan: s239L
% !PRINTED 450
\weberpage{450}
% head: CAPITOLO UNDECIMO
%<<CONT
% stamp: Dark vertical scan-edge line at the right; scattered specks and a diagonal stroke near the inline p_g=1.
e segherà sopra una \(C\) una serie \(g^{\pi-2}_{2\pi-2}\), di deficienza uno, contenuta nella serie canonica. Inoltre il detto \(|C'|\) sarà contenuto in una serie \(\infty^1\) di sistemi lineari paraggiunti \(|\overline{C}'|\), secanti egualmente su \(C\) una serie \(g^{\pi-2}_{2\pi-2}\). Pertanto vi sarà una curva paracanonica ellittica definita come differenza
\[
\overline{C}'-C,
\]
e le curve paracanoniche analoghe \(K\) formeranno su \(F\) un fascio ellittico: aggiungasi che ogni curva \(2K\) apparterrà quindi ad un sistema lineare \(\infty^1\) (almeno) e perciò \(P_2>1\), \hfill c. d. d.

\textit{Osservazione.} -- Un esempio di superficie avente i generi \(p_a=0\), \(p_g=1\), \(p^{(1)}=1\), e possedente un fascio ellittico di curve ellittiche, è dato dalle più semplici superficie paraellittiche, che abbiamo incontrato nel \S~11 del Cap. X.

Può essere interessante di studiare e classificare le superficie con \(p_a=0\) e \(p_g=1\), aventi il genere lineare \(p^{(1)}=1\), o anche \(p^{(1)}=2\ldotp\ldotp\ldotp\ldotp\)

I risultati conseguiti si lasciano riassumere dicendo: \textit{Le superficie, non riferibili a rigate \((P_{12}>0)\), che posseggono una curva canonica o pluricanonica d’ordine zero} (e quindi sistemi lineari senza punti base di genere \(\pi\) e grado \(n=2\pi-2\)) \textit{sono definite dalla condizione}
\[
P_{12}=1,
\]
\textit{che porta}
\[
p^{(1)}=1
\]
\textit{e dà luogo alle seguenti famiglie:}

1) \textit{\(p_g=0\) e \(p_a=0\), superficie regolari con curva bicanonica d’ordine zero \((P_2=1)\), rappresentate dalla superficie del \(6^{\mathrm{o}}\) ordine che passa doppiamente per gli spigoli d’un tetraedro} (Cap. VII, \S~1);

2) \textit{\(p_g=0\) e \(p_a=-1\), superficie ellittica dei tipi} \(\mathrm{I}_a\ \mathrm{I}_b\ \mathrm{II}_a\ \mathrm{II}_b\ \mathrm{III}_a\ \mathrm{III}_b\ \mathrm{III}_e\) (Cap. X, \S\S~16-19); % ?? uncertain: final subscript in the list of types -- c, e.
% EDITORIAL D0241 RISOLTO: The subscript has the central crossbar and enclosed upper part of e, rather than the open curve of c.

3) \textit{\(p_g=1\) e \(p_a=1\) \((P_2=1)\), superficie regolari coi generi uno} (Cap. VII, \S~2);

4) \textit{\(p_g=1\) e \(p_a=-1\) \((P_4=1)\) superficie iperellittiche} (\S\S~4 e 6 di questo capitolo).

\edsection{6. Nota sulla teoria geometrica delle superficie iperellittiche.}{6. Nota sulla teoria geometrica delle superficie iperellittiche.}{sec.XI.6}

Fra le famiglie di superficie col \(P_{12}=1\), contrassegnate coi nn. 1) 2) 3) 4) alla fine del precedente paragrafo, conviene distinguere: da una parte le superficie di genere \(p_g=1\), che sono le 3) e 4) (\(p_g=1\),
%CONT>>
% scan: s239R
% !PRINTED 451
\weberpage{451}
% head: CLASSIFICAZIONE GENERALE DELLE SUPERFICIE
%<<CONT
% stamp: Scattered scan specks and marginal marks.
\(p_a=1\) e \(p_a=-1\)) e dall’altra le superficie di genere \(p_g=0\) (regolari con \(P_2=1\), ovvero ellittiche) che sono le 1) e 2). Queste ultime costituiscono un numero finito di famiglie a ciascuna delle quali appartiene un certo numero di moduli: precisamente: 10 moduli alle 1), 2 moduli o 1 modulo alle 2). Invece le superficie 3) e 4) danno luogo a \textit{infinite famiglie} distinte dai valori di (almeno) un carattere aritmetico, a ciascuna delle quali appartiene un certo numero di moduli.

Ciò che qui si asserisce è stato dimostrato per le superficie con tutti i generi uno (le 3)) nel Cap. XII, \S~2; ma sussiste anche per le superficie iperellittiche 4), come occorre qui spiegare. % [sic] Cap. XII -> Cap. VII, apparently inconsistent with the reference for the same family on the preceding page.

La \textit{classificazione delle superficie iperellittiche} risulta anzitutto, sotto l’aspetto trascendente, dalla rappresentazione parametrica di esse mediante funzioni 4 volte periodiche di due variabili indipendenti\origfootnote{1}{s239R-1}{Cfr. p. es. \textsc{Enriques-Chisini}, \textit{Lezioni.} Libro VI, cap. III, \S~39 (vol. IV, pag. 242).}. I periodi normali primitivi di esse si possono ridurre alla tavola
\[
\begin{tabular}{cccc}
\(1\) & \(0\) & \(g\) & \(h\) \\
\(0\) & \(\dfrac{1}{\delta}\) & \(h\) & \(g',\)
\end{tabular}
\]
dove \(\delta\) è un carattere intero (il cosidetto \textit{divisore} delle corrispondenti superficie). Ora per \(\delta=1\) si ha la \textit{superficie tipo}, che è la superficie di Jacobi, \(F_j\), rappresentativa della varietà delle \textit{coppie di punti della curva di genere due}; invece per \(\delta>1\) si ha una \(F_j\) multipla secondo \(\delta\) senza curva di diramazione, ovvero un’involuzione d’ordine \(\delta\) sulla stessa superficie \(F_j\).

La classificazione delle superficie iperellittiche così ottenuta ha un significato algebrico-geometrico, e si pone quindi il problema di giustificarla sotto questo aspetto. Diciamo in breve a quale ordini di considerazioni si venga condotti in tal guisa. % [sic] quale ordini -> quali ordini, adjective and plural noun do not agree.

Anzitutto si possono determinare i caratteri della superficie di Jacobi \(F_j\):
\[
p_g=1,\qquad p_a=-1
\]
e curva canonica d’ordine zero (quindi \(P_4=1\)). Ciò si ottiene rappresentando la \(F_j\) sopra la superficie doppia di Kummer, del \(4^{\mathrm{o}}\) ordine con 16 punti doppi, come si è visto nel \S~8 del Cap. V\origfootnote{2}{s239R-2}{Cfr. anche \textsc{Enriques-Chisini}, \textit{Lezioni.} L. V. Cap. III, \S~40 (Vol. IV pag. 24).}. Questa rappresentazione si basa sulla conoscenza delle trasformazioni (involutorie) di \(2^{\mathrm{a}}\) specie della varietà delle coppie di punti della curva \(C\) di genere due, quali si ottengono associando le cop%
%CONT>>
% scan: s240L
% !PRINTED 452
\unskip\weberpage{452}%
% head: CAPITOLO UNDECIMO
%<<CONT
% stamp: Dark vertical scan-edge line at the right; scattered specks and short marginal marks.
\unskip pie residue l’una all’altra rispetto ad una \(g^2_4\) presa su \(C\): un’involuzione così definita è birazionalmente identica ad una superficie di Kummer, del \(4^{\mathrm{o}}\) ordine, avente 16 punti doppi che corrispondono alle 16 coppie di punti della \(C\) provenienti dalla bisezione della \(g^2_4\).

Ora sulla superficie di Jacobi \(F_j\) i prodotti delle dette involuzioni \(I\) danno luogo alle trasformazioni di un gruppo continuo abeliano \(\Gamma\); e vi sono trasformazioni cicliche del \(\Gamma\) che generano gruppi finiti \(G_\delta\) in esso contenuti: un tale \(G_\delta\) (sia p. es. un \(G_\delta\) ciclico) dà luogo su \(F_j\) ad una involuzione \(I_\delta\) d’ordine \(\delta\), priva di punti doppi, cui risponde in generale una superficie iperellittica \(F_\delta\) diversa dalla \(F_j\). Si ottengono così le superficie iperellittiche di divisore \(\delta>1\). E si può vedere che esse possono anche prodursi a partire dalla \(F_j\) multipla secondo \(\delta\), senza curva di diramazione; a tale scopo occorre riconoscere che coi gruppi di un’involuzione \(I_\delta\), si può costruire una involuzione d’ordine superiore, diciamo \(r\delta\), che riesca birazionalmente identica alla \(F_j\).

In ciò che si è detto viene esposto il disegno di una \textit{teoria geometrica delle superficie iperellittiche}, basata sullo studio delle trasformazioni cicliche appartenenti al gruppo \(\Gamma\) della superficie di Jacobi \(F_j\): un tale studio si presenta come una estensione dello studio delle involuzioni cicliche sulle curve ellittiche, che si trova sviluppato nelle « Lezioni » di \textsc{Enriques-Chisini}\origfootnote{1}{s240L-1}{Libro V, cap. III, \S\S~27-28 (vol. III, pagg. 251-298).}.

Ma la considerazione delle superficie iperellittiche sotto l’aspetto geometrico ci mette di fronte anche ad altri problemi interessanti.

Invero noi abbiamo dimostrato (Cap. IX, \S~9) che una varietà abeliana \(V_q\) si può ritenere identica ad una varietà semplice o multipla contenuta imprimitivamente entro la varietà di Jacobi \(V_p\) che risponde ad una curva di genere \(p\ge q\); ma non sussiste in generale che possa prendersi sempre \(p=q\), perchè il numero dei moduli da cui dipende la \(V_q\) è maggiore di quello, \(3q-3\), delle curve di genere \(q\)\origfootnote{2}{s240L-2}{Cfr. \textsc{Enriques-Chisini}, \textit{Lezioni.} Libro V, cap. III, \S~33 (vol. III, pag. 355) e Libro VI, cap. III, \S~35 (vol. IV, pag. 216).}. Ora la teoria trascendente ci mostra che le superficie iperellittiche (con un dato divisore) dipendono da \textit{tre moduli}, e perciò si riducono tutte al tipo della superficie di Jacobi \(F_j\) semplice o multipla, e, in quest’ultimo caso, anche ad involuzioni (prive di punti doppi) sopra la \(F_j\).

La cosa si può riconoscere sotto l’aspetto geometrico, nel modo che brevemente accenniamo. Sopra una superficie iperellittica \(F\) (per cui \(p_a=-1\)) una curva di genere \(\pi\) appartiene, in generale, ad un sistema lineare di dimensione \(\pi-2\), in cui sarà contenuto
%CONT>>
% scan: s240R
% !PRINTED 453
\weberpage{453}
% head: CLASSIFICAZIONE GENERALE DELLE SUPERFICIE
%<<CONT
% stamp: Fingerprint-like smudge in the lower right margin; scattered scan specks and marginal marks.
un numero finito di curve \(C\), dotate di \(\pi-2\) punti doppi, e perciò di genere due.

Mercè le trasformazioni di prima specie della \(F\), appartenenti al gruppo \(\Gamma\), si otterrà quindi su \(F\) un sistema continuo \(\infty^2\) \(\{C\}\) di curve di genere due, fra loro birazionalmente identiche, i cui \textit{tre moduli} forniranno i moduli della \(F\). Ora una \(C\) di codesto sistema è invariante per \textit{una} trasformazione (involutoria) di seconda specie\origfootnote{1}{s240R-1}{Riferendosi alla rappresentazione parametrica del gruppo \(\Gamma\) una trasformazione di seconda specie si può fare rispondere a quella in cui una \(\pi\) di \(\Gamma\) si muta nell’inversa \(\pi^{-1}\). È facile persuadersi che, anche nel caso in cui la curva \(C\) di genere due possegga trasformazioni singolari in sè, non vi possono essere più trasformazioni di \(2^{\mathrm{a}}\) specie di \(F\) che la lascino invariata.} della \(F\), che subordina sulla \(C\) la sua \(g^1_2\); quindi il gruppo \(\Gamma\) è birazionalmente identico, sia alla superficie \(F\) (concepita come luogo di punti) sia al sistema \(\{C\}\) ritenuto come varietà \(\infty^2\) di elementi (curve \(C\)). Segue di qui che la \(F\) riesce birazionalmente identica alla serie \(\infty^2\) dei gruppi di punti segati dalle curve di \(\{C\}\) sopra una \(C\) particolare: se i gruppi di questa serie sono tutti disequivalenti la \(F\) risulta identica alla superficie di Jacobi che risponde alla detta \(C\); se invece per ogni gruppo della detta serie vi è un certo numero \(\delta\) di gruppi equivalenti, la \textit{superficie} iperellittica \(F\) risulta birazionalmente \textit{identica ad una involuzione \(I_\delta\) d ordine \(\delta\) sopra la superficie di Jacobi.} % [sic] d ordine -> d’ordine, apparently missing elision apostrophe.

Il teorema stabilito assume un significato geometrico interessante in relazione a diverse costruzioni che possono darsi delle superficie iperellittiche. Invero si può costruire una superficie iperellittica a partire dalla varietà di Jacobi corrispondente ad una curva \(C_p\) di genere \(p\). Quando su questa \(C_p\) si trovino due integrali abeliani riducibili a 4 periodi (secondo l’interpretazione geometrica del teorema di Picard-Poincaré dovuta a \textsc{Castelnuovo})\origfootnote{2}{s240R-2}{Cfr. \textsc{Enriques-Chisini}, \textit{Lezioni.} Libro VI, cap. III, \S~38 (vol. IV, pag. 234).} % ?? uncertain: final punctuation of this footnote -- period, short dash.
la relativa varietà di Jacobi \(V_p\) conterrà un sottogruppo algebrico \(\infty^2\) del gruppo \(\Gamma\) delle trasformazioni di prima specie, e si avrà quindi una varietà abeliana (cioè una superficie iperellittica) contenuta imprimitivamente in \(V_p\). Il teorema precedente ci dice che le superficie iperellittiche costruite in tal guisa si possono anche ritrovare come superficie di Jacobi semplici o multiple, ovvero come involuzioni sulle superficie di Jacobi, corrispondenti alla curva di genere \(p=2\).

Giova illustrare questa proprietà riferendosi al caso più semplice in cui si tratti delle superficie iperellittiche contenute imprimitivamente entro la varietà di Jacobi \(V_3\), che risponde ad una curva \(C_3\) di genere tre. Qui il gruppo \(\Gamma\) delle trasformazioni di
%CONT>>
% scan: s241L
% !PRINTED 454
\weberpage{454}
% head: CAPITOLO UNDECIMO
%<<CONT
% stamp: Dark vertical scan-edge line at the right margin, scattered specks, and short diagonal marks.
prima specie della \(V_3\), conterrà due sottogruppi algebrici complementari: un \(\Gamma_2\ \infty^2\) che dà luogo ad una schiera d’imprimitività costituita da \(\infty^1\) superficie iperellittiche \(F_2\), e un \(\Gamma_1\ \infty^1\) che dà luogo ad una schiera d’imprimitività costituita da \(\infty^2\) curve ellittiche \(C_1\), le \(F_2\) e le \(C_1\) incontrandosi in un certo numero di punti. Ed è agevole dimostrare che ciascuna \(C_1\) è una trasformata unirazionale della \(C_3\), cioè corrisponde ad una involuzione ellittica \(\gamma_n^1\) (di un certo ordine \(n\)) appartenente alla \(C_3\).

Pertanto si potrà costruire una superficie iperellittica \(F_2\), contenuta imprimitivamente in una varietà di Jacobi \(V_3\), a partire da una curva \(C_3\) di genere tre, cui appartenga una involuzione ellittica \(\gamma_n^1\); a tal uopo basterà costruire la \(C_3\) che risponde ad una cubica \(K\) multipla secondo il numero \(n\), sulla quale si assumano ad arbitrio 4 punti di diramazione colle relative sostituzioni fra i rami\origfootnote{1}{s241L-1}{Cfr. il teorema d’esistenza per le curve multiple in \textsc{Enriques-Chisini}, \textit{Lezioni.} Libro V, cap. IV, \S~38 (vol. III, pag. 427).}. Si noterà che le \(C_3\) contenenti una \(\gamma_n^1\) ellittica dipenderanno da 4 moduli: l’invariante assoluto della \(K\) e le \textit{distanze} di 3 punti di diramazione del rimanente, sulla stessa curva ellittica \(K\).

Tolto l’invariante assoluto di \(K\), che risponde a quello delle \(C_1\) sulla \(V_3\), restano 3 moduli appartenenti alle superficie complementari \(F_2\).

Aggiungasi che la \(V_3\) di Jacobi costruita innanzi conterrà, in generale, una involuzione costituita dai gruppi di punti sezioni delle \(F_2\) e delle \(C_1\): la varietà \(\overline{V}_3\) rappresentativa di codesti gruppi, che è trasformata unirazionale della \(V_3\), sarà la varietà di Jacobi di una curva di genere 3 degenere in una curva ellittica e in una curva di genere due; così dunque (d’accordo col teorema sopra esposto), le superficie iperellittiche \(F_2\), che formano un sistema d’imprimitività per il gruppo delle trasformazioni di prima specie della \(V_3\), risultano birazionalmente identiche ad una involuzione su una delle analoghe \(\overline{F}_2\), che è la superficie di Jacobi corrispondente alla predetta curva di genere due.

I cenni dati innanzi bastino a segnalare l’interesse, la bellezza e la ricchezza delle proprietà, che appartengono alle superficie iperellittiche.

Aggiungeremo soltanto che lo studio delle involuzioni irregolari iperellittiche sopra una superficie iperellittica \(F\) (e in ispecie sulla superficie di Jacobi \(F_j\)) si prolunga naturalmente in quello di altre involuzioni irregolari ellittiche o regolari coi generi uno, che possono appartenere alla \(F\) e conducono a superficie cui spetta ancora una
%CONT>>
% scan: s241R
% !PRINTED 455
\weberpage{455}
% head: CLASSIFICAZIONE GENERALE DELLE SUPERFICIE
%<<CONT
% stamp: Scattered scan specks in the margins.
rappresentazione parametrica con funzioni iperellittiche \textit{di rango} \(r>1\)\origfootnote{1}{s241R-1}{Dove ad ogni punto della superficie rispondono \(r\) punti entro il prismatoide dei periodi.}.

Le involuzioni irregolari ellittiche (che appartengono ad \(F\) con moduli particolari) danno luogo alle superficie ellittiche con \(p_a=-1\), \(p_g=0\) e \(P_{12}=1\), possedenti una curva canonica virtuale d’ordine zero, come già si è accennato nei \S\S~16-19 del Cap. X.

Le involuzioni regolari d’ordine 2 appartenenti ad una superficie iperellittica di divisore \(\delta\geq1\) conducono in generale a superficie (iperellittiche in senso esteso) di rango \(r=2\), e precisamente: per \(\delta=1\) alle superficie di Kummer del \(4^{\mathrm{o}}\) ordine con 16 punti doppi (Cap. V, \S~8); e per \(\delta>1\) ad una serie di superficie analoghe d’ordine \(n=2\pi-2\) a sezioni iperpiane canoniche di genere \(\pi\) nello spazio \(S_\pi\) \((\pi=4,5,6,\ldotp\ldotp\ldotp\ldotp)\), dotate di 16 punti doppi. Ma per moduli particolari si presentano pure sopra una \(F\) \textit{involuzioni regolari singolari}, le quali vengono generate da un gruppo finito di trasformazioni birazionali della detta \(F\); e in base a questa generazione (dimostrata da \textsc{Enriques-Severi}) sono state classificate nella loro integrità dalle memorie già citate, che vicendevolmente si completano, di \textsc{Enriques} e \textsc{Severi}\origfootnote{2}{s241R-2}{\textit{Mémoire sur les surfaces hyperellyptiques.} Acta Mathemat., t. 32 e 33 (1908).} % [sic] "hyperellyptiques" -> "hyperelliptiques", apparent French spelling error.
e di \textsc{Bagnera} e \textsc{De Franchis}\origfootnote{3}{s241R-3}{\textit{Le superficie algebriche le quali ammettono una rappresentazione parametrica mediante funzioni iperellittiche di due argomenti.} Memorie della Soc. It. d. Scienze (dotta dei XL), 1908.}. % ?? uncertain: "dotta" -- "dotta" or "detta"; the vowel is indistinct.
% EDITORIAL D0243 RISOLTO: The vowel is a closed round o, with no crossbar of e. Preserve the printed dotta.

Infine ricorderemo che questo ordine di considerazioni si estende, dalle superficie iperellittiche alle superficie regolari con tutti i generi uno, sulle quali è dato pure di trovare involuzioni regolari dello stesso genere \(p_g=1\) (o con \(p_g=0\) e bigenere \(P_2=1\)), che vengono generate da un gruppo finito di trasformazioni birazionali. In proposito basterà menzionare i lavori interessanti della scuola belga, e particolarmente di \textsc{Godeaux}.

\edsection{7. Classificazione generale delle superficie algebriche\origfootnote{4}{s241R-4}{Cfr. \textsc{Enriques}, \textit{Sulla classificazione delle superficie algebriche in particolare di genere lineare \(p^{(1)}=1\).} Rendic. Acc. Lincei, Febbraio 1914.}.}{7. Classificazione generale delle superficie algebriche.}{sec.XI.7}

Siamo ora in grado di raccogliere in una sintesi i risultati principali della teoria delle superficie, esponendo una classificazione generale delle superficie algebriche, secondo il valore del dodici-genere \(P_{12}\): nel seguito questa classificazione potrà stringersi più da vicino, come vedremo nei \S\S~8 e 9.

% scan: s242L
% !PRINTED 456
\weberpage{456}
% head: CAPITOLO UNDECIMO
% stamp: Broken dark scan-edge lines at the right margin and a diagonal mark between "lineari" and "di".
Distinguiamo le superficie algebriche secondo i valori del \(P_{12}\) e — subordinatamente — del genere lineare (assoluto) \(p^{(1)}\).

Avremo \textit{quattro grandi famiglie di superficie}:

\(A)\) \textit{per \(P_{12}=0\) le superficie riferibili a rigate} (razionali o meno);

\(B)\) \textit{per \(P_{12}=1\) la superficie con curva canonica}, effettiva o virtuale, \textit{d’ordine zero};

\(C)\) \textit{per \(P_{12}>1\) e \(p^{(1)}=1\) le superficie che contengono sistemi lineari infiniti di curve canoniche o pluricanoniche composte colle curve ellittiche d’un fascio};

\(D)\) \textit{per \(P_{12}>1\) e \(p^{(1)}>1\), le superficie che ammettono come modello proiettivo superficie canoniche o pluricanoniche.}

Conviene richiamare e spiegare le \textit{proprietà fondamentali e caratteristiche} che distinguono queste quattro famiglie di superficie:

\(A)\) \(P_{12}=0\). Superficie riferibili a rigate (razionali o no): posseggono sistemi lineari di curve di genere \(\pi\) e grado \(n>2\pi-2\), curve eccezionali in numero infinito non eliminabili, gruppi razionali \(\infty^1\) di trasformazioni e serie continue di trasformazioni birazionali in sè non formanti gruppo. Si dividono in

\(A')\) superficie razionali (regolari) caratterizzate dalle condizioni
\[
p_a=P_2=0,
\]
e

\(A'')\) rigate irrazionali, per cui
\[
p_a<0;
\]
le rigate di genere \(p=-p_a>1\) sono caratterizzate (come vedremo nel \S~8) dalla semplice condizione
\[
p_a<-1,
\]
mentre le rigate di genere \(p=1\) si distinguono dalle superficie propriamente ellittiche con \(p_a=-1\) per avere \(P_4=P_6=P_{12}=0\)). % [sic] "0)." -> "0.", unmatched closing parenthesis.

\(B)\) \(P_{12}=1,\ p^{(1)}=1\) (\(p_a\geq-1\)). Superficie possedenti una curva canonica o pluricanonica d’ordine zero. Per ogni sistema senza punti base di genere \(\pi\) e grado \(n\), sopra la superficie privata di curve eccezionali, si ha
\[
n=2\pi-2,
\]
sicchè si ottengono, in generale, come immagini proiettive di questo tipo, superficie \(F_n\) a sezioni piane o iperpiane di genere \(\pi\) d’ordine \(n=2\pi-2\), in uno spazio ad \(r=\pi,\ \pi-1\) o \(\pi-2\) dimensioni (i tre valori di \(r\) rispondendo ai tre casi: \(p_a=p_g=1\), \(p_a=p_g=0\), e \(p_a=-1\)).

Le superficie \(F_n\) della famiglia \(B)\) si dividono in due categorie:

\(B')\) superficie a sezioni (piane o iperpiane) canoniche, che hanno il genere geometrico \(p_g=1\) e il genere numerico \(p_a=1\) ovvero
%CONT>>
% scan: s242R
% !PRINTED 457
\weberpage{457}
% head: CLASSIFICAZIONE GENERALE DELLE SUPERFICIE
%<<CONT
% stamp: Scattered scan specks and small marginal strokes.
\(p_a=-1\), cioè superficie regolari con tutti i generi uno (Cap. VII, \S~2) e superficie iperellittiche;

\(B'')\) superficie le cui sezioni sono curve non canoniche ma rappresentative di serie lineari che provengono dalla divisione per \(s\) della serie \(s\)-canonica (\(s=2,3,4,6\)).

Queste sono le superficie regolari di genere \(p_a=p_g=0\) e \(P_2=1\) (\(P_3=0\)) riducibili alla sestica \(F_6\) che passa doppiamente per gli spigoli d’un tetraedro (Cap. VII, \S~1) e le superficie ellittiche con \(p_a=-1\) e \(p_g=0\) contenenti un fascio lineare di curve ellittiche, che abbiamo distinti in sette tipi \(I_a\,I_b\,II_a\,II_b\,III_a\,III_b\,III_c\) nei \S\S~16-19 del Cap. X. % [sic] "distinti" -> "distinte", apparent agreement error referring to "superficie".

Fra le superficie \(B'\) e le \(B''\) intercede questa differenza: che, mentre le \(B''\) danno luogo ad un numero finito di tipi, le \(B'\) danno luogo a un numero infinito di famiglie di superficie, dipendenti (almeno) da un carattere intero, suscettibile di assumere infiniti valori.

Ciò è stato dimostrato per le superficie regolari con tutti i generi uno nel \S~2 del Cap. VII; (dove si è pur visto che ad ogni famiglia appartengono 19 moduli); ma è pur vero per le superficie iperellittiche, dove, accanto alla superficie di Jacobi \(F_j\) corrispondente alla curva di genere due, vi sono le infinite altre che rispondono ad involuzioni d’ordine \(\delta\) su \(F_j\) ovvero a \(F_j\) multiple secondo \(\delta\), senza curva di diramazione: ogni classe essendo definita entro la rispettiva famiglia da 3 moduli (cfr. \S~6).

Invece le superficie \(B')\) danno luogo, come si è detto, a un numero finito, e precisamente ad 8 famiglie, nettamente caratterizzate e contenenti: la prima (superficie regolari) 10 moduli (Cap. VII, \S~1) e le altre (ellittiche) 2 moduli (le \(I_a\,I_b\)) o 1 modulo (le rimanenti del tipo armonico o equianarmonico). % [sic] B') -> B''), the preceding distinction assigns the finite families to B''.

Aggiungasi che le superficie \(B')\) contengono sempre fasci di curve ellittiche e perciò si ravvicinano alle superficie \(C)\), in confronto alle \(B'')\) che non contengono in generale simili fasci. % [sic] first B') and final B'') -> first B'') and final B'), these labels reverse the distinction stated above.

\(C)\) \(P_{12}>1\) e \(p^{(1)}=1\). Superficie possedenti un fascio (razionale o irrazionale) di curve ellittiche \(K\), componenti di curve canoniche o pluricanoniche. Insieme alle \(B'')\) queste superficie formano la vasta famiglia delle \textit{superficie che contengono un fascio di curve ellittiche.}

Queste superficie si lasciano classificare, come abbiam visto nel caso regolare (Cap. VII, \S~4), secondo alcuni caratteri interi, cioè:

1) il genere numerico \(p_a\);

2) il determinante \(n\), cioè l’ordine del più piccolo gruppo di punti che può determinarsi razionalmente su una \(K\);

% scan: s243L
% !PRINTED 458
\weberpage{458}
% head: CAPITOLO UNDECIMO
% stamp: Dark vertical scan-edge line at the right margin and diagonal marks near "mediante".
3) il numero delle curve del fascio che si riducono ad una curva multipla; cui si aggiunge nel caso irregolare

4) il genere \(\varrho\) del fascio che vale, in generale (\textit{Osservazione}, Cap. IX, \S~8) % ?? uncertain: the rho glyph in this paragraph and its displays -- \varrho (read here) or \rho.
% EDITORIAL D0244 RISOLTO: The loop has a curved lower return rather than the straight descender of rho; the same form appears in both equations.
\[
\begin{aligned}
&\varrho=p_g-p_a\\
&\quad(p_a\geq0)
\end{aligned}
\]
e, invece,
\[
\varrho=p_g-p_a-1
\]
nel caso in cui la superficie contenga oltre \((K)\) un secondo fascio ellittico di curve, caso che porta \(p_a=-1\) e conduce alle superficie ellittiche.

Per le superficie \(C)\) regolari abbiamo riconosciuto che le infinite famiglie di esse, dipendenti dai caratteri sopra indicati si lasciano distinguere mediante i valori dei plurigeneri (Cap. VII, \S~4).

Questi risultati sembrano, in gran parte, estendersi anche alle \textit{superficie irregolari}: \(p_a<p_g\). Ad ogni modo vi è qui un \textit{oggetto di studio}, che segnaliamo all’attenzione dei giovani cultori della geometria.

\(D)\) \(P_{12}>1,\ p^{(1)}>1\). Queste superficie, che debbonsi ritenere costituenti il \textit{caso più generale}, a differenza delle precedenti \(A)\ B)\) e \(C)\) danno luogo, per dati valori dei caratteri \(p_a,\ p_g,\ p^{(1)}\), ad un \textit{numero finito di famiglie}, ciascuna delle quali contiene un sistema continuo di classi, dipendenti da un \textit{certo numero} di \textit{parametri} o \textit{moduli}. Infatti alle superficie \(D)\) risponde un tipo di superficie canonica o pluricanonica, il cui ordine è \(p^{(1)}-1\) o risp. un multiplo di \(p^{(1)}-1\), la quale offre un modello proiettivo delle superficie stesse, riguardate di fronte alle trasformazioni birazionali: la classificazione delle superficie di un dato ordine, con alcuni caratteri dati, non può condurre che ad un numero finito di tipi aritmeticamente distinti.

Abbiamo già avuto occasione di studiare il \textit{problema generale della classificazione delle superficie con \(p^{(1)}>1\) nel caso regolare} (Cap. VIII). Le considerazioni svolte si estendono, in gran parte, al caso in cui si lasci cadere l’ipotesi \(p_a=p_g\); in ogni modo vi è qui un \textit{oggetto di ricerca} per gli studiosi. È degno di nota che quando si procede a costruire le superficie canoniche o pluricanoniche con dati \(p_g\) e \(p^{(1)}\) s’incontrino di solito superficie regolari e solo in via d’eccezione appaia la possibilità di superficie irregolari; così riescono regolari le superficie canoniche per \(p_g\geq4\), che si costruiscono nei primi casi \(p^{(1)}=5,6,7\ldotp\ldotp\ldotp\ldotp\), (Cap. VIII), e similmente sembra risultino necessariamente regolari anche le superficie tricanoniche per \(p^{(1)}=2\) e le bicanoniche per \(p^{(1)}=3\) che si costruiscono in
%CONT>>
% scan: s243R
% !PRINTED 459
\weberpage{459}
% head: CLASSIFICAZIONE GENERALE DELLE SUPERFICIE
%<<CONT
% stamp: Dark scan-edge line at the left margin, scattered specks, and small marks near the upper and lower right corners.
corrispondenza ad un \(p_a>0\), secondo il procedimento dei \S\S~14-20 del Cap. VIII,\origfootnote{1}{s243R-1}{Considerazioni che si completano coi teoremi sugli integrali semplici appartenenti ad una superficie conducono in generale ad affermare che le superficie, il cui sistema canonico non è composto con le curve di un fascio, di irregolarità \(q\geq3\) col \(p_a\geq0\) hanno il genere lineare \(p^{(1)}\geq6p_a+4\); se invece \(q=2\) e \(p_a\geq0\), allora \(p^{(1)}\geq3p_a+4\).

Perciò resterebbero da esaminare soltanto le superficie di irregolarità \(q=1\) possedenti un fascio ellittico di curve (Comunicazione gentilmente fattaci da G. \textsc{Castelnuovo}, Novembre 1942).}. % ?? uncertain: the mark after "VIII" before the footnote -- a full stop (read here), a comma, or a scan speck.
% EDITORIAL D0245 RISOLTO: The enlarged mark has an upper head and a tail descending leftward, visibly a comma rather than a period.

Qui conviene aggiungere che in forza di una diseguaglianza di cui discorriamo nel seguente paragrafo, le superficie \(D)\) (con \(p^{(1)}>1\)) hanno tutte il genere numerico: \(p_a\geq0\).

\edsection{8. Superficie di genere numerico negativo.}{8. Superficie di genere numerico negativo.}{sec.XI.8}

La classificazione che abbiamo esposta nel precedente paragrafo può essere stretta più da vicino sulla base di una diseguaglianza fra l’irregolarità e il genere geometrico, che fino ad oggi non si è riusciti a dimostrare in modo algebrico-geometrico, e che perciò occorre attingere alla teoria trascendente (cfr. Notizia, pag. 463). Sussiste il teorema: \textit{Per ogni superficie di genere \(p_g\) e irregolarità \(q=p_g-p_a\), che non contenga un fascio irrazionale di curve, si ha} % [sic] "pag. 463" -> "pag. 461", the Notizia reproduced in this packet begins on printed page 461.
\[
p_g>2(q-2)
\]
ossia
\[
p_g<2p_a+4.
\]
Dunque: \textit{se \(p_g\geq2p_a+4\), la superficie contiene un fascio irrazionale di curve \((C)\)}.

La precedente diseguaglianza è certo soddisfatta se il genere numerico
\[
p_a<-1\quad(p_g\geq0).
\]

Pertanto si può affermare che le superficie \(F\) di genere numerico \(p_a<-1\) contengono un fascio irrazionale \((C)\): fascio di genere \(p>0\) formato di curve \(C\) di un certo genere \(\pi\). Ma si riconosce agevolmente che \(\pi=0\) e quindi la \(F\) si può trasformare in una rigata (di genere \(p=-p_a\)).

A tal uopo ragioniamo per assurdo: se \(\pi>0\), la \(F\) si può ridurre ad una superficie priva di curve eccezionali, di genere lineare assoluto \(p^{(1)}\geq1\); quindi — calcolando il suo invariante di Zeuthen-Segre in relazione a \((C)\) — si trova (come nel Cap. X, \S~4)
\[
\begin{aligned}
&\Delta-4=12p_a-p^{(1)}+9,\\
&\Delta\geq0\\
&p_a\geq-1\ !
\end{aligned}
\]

% scan: s244L
% !PRINTED 460
\weberpage{460}
% head: CAPITOLO UNDECIMO
% stamp: Dark vertical scan-edge line at the right margin; diagonal strokes near the second line of the Zeuthen-Segre formula and small marginal dashes.
Concludiamo dunque che \textit{le superficie di genere numerico \(p_a<-1\) sono tutte riferibili a rigate} (di genere \(p=-p_a\)): nel teorema (Cap. X, \S~4) che assegna a tal uopo le condizioni \(p_a=-p\) e \(p_g=0\), la condizione \(p_g=0\) è superflua, risultando come conseguenza dall’altra.

Facciamo \(p_a=-1\). Se la superficie \(F\) ha il genere geometrico
\[
p_g>1,
\]
la diseguaglianza
\[
p_g\geq2p_a+4
\]
è ancora soddisfatta e quindi la \(F\) possiede un fascio irrazionale di curve \(C\) (di genere \(\pi>0\)). Ma, in relazione a questo fascio si trova l’invariante di Zeuthen-Segre
\[
\begin{aligned}
&\Delta-4=12p_a-p^{(1)}+9\\
&\quad p_a=-1,\quad\Delta\geq0
\end{aligned}
\]
e quindi
\[
p^{(1)}=1.
\]

Così la condizione \(p^{(1)}=1\), che — in aggiunta alle \(p_a=-1\) e \(p_g>1\) — caratterizza le superficie ellittiche (Cap. XI, \S~4) è una conseguenza delle altre due condizioni. E, siccome per \(p_g=0\) e \(p_a=-1\) si ha ancora \(p^{(1)}=1\) e la \(F\) risulta una superficie ellittica (propria o impropria), possiamo enunciare che: \textit{le superficie di genere numerico \(p_a=-1\) e di genere geometrico \(p_g\ne1\) sono ellittiche} (superficie ellittiche proprie o rigate di genere uno).

Infine esaminiamo il caso:
\[
p_a=-1,\quad p_g=1.
\]

In questo caso abbiamo veduto (Cap. X, \S~4) che la \(F\) è iperellittica (con \(P_4=1\)) ovvero si può ridurre ad una superficie iperellittica \(\Phi\), multipla con curva di diramazione ellittica ed allora è, essa stessa, una superficie ellittica (con \(P_4>1\)).

Così possiamo completare i teoremi stabiliti nel \S~4 del Cap. X, nel senso già ivi annunciato dicendo che:

\textit{Le superficie ellittiche ed iperellittiche sono caratterizzate globalmente dalla condizione}
\[
p_a=-1. % ?? uncertain: the small mark after "p_a=-1" -- a full stop (read here) or a scan speck.
\]
(congiunta a quelle che escludono le superficie ellittiche improprie, riferibili a rigate: \(P_4+P_6=0\), o \(P_{12}=0\)).

\textit{Le superficie iperellittiche} (che solo per moduli particolari sono anche — in doppio modo — ellittiche) \textit{vengono distinte dalle ellittiche dalle condizioni}
\[
p_g=P_4=1
\]
(anzichè
\[
p_g=1,\quad P_4>1).
\]

% scan: s244R
% !PRINTED 461
\weberpage{461}
% head: CLASSIFICAZIONE GENERALE DELLE SUPERFICIE
% stamp: Dark scan-edge marks at the upper right and right margin, with scattered specks.
Si possono riassumere i risultati ottenuti enunciando il teorema:

\textit{Le superficie di genere numerico negativo} (\(p_a<0\)) \textit{sono}:

1) \textit{riferibili a rigate} (irrazionali).

2) \textit{ovvero superficie ellittiche- o iperellittiche} (in quest’ultimo caso: \(p_g=1\)). % ?? uncertain: the short stroke after "ellittiche" -- a printed hyphen (read here) or a scan mark.

Od anche:

\textit{Il genere numerico negativo vale a definire l’intera famiglia delle superficie, non razionali, che ammettono un gruppo continuo di trasformazioni birazionali.}

\textit{Notizia.} – La diseguaglianza fondamentale \(p_g<2p_a+4\), per le superficie che non contengono un fascio irrazionale di curve, deriva — come si è detto — da considerazioni d’ordine trascendente. Premettiamo che il genere \(p_g\) di una superficie \(F\) è il numero degli integrali doppi di prima specie, linearmente indipendenti, che appartengono ad \(F\), e d’altra parte che l’irregolarità \(q\) è il numero degli integrali semplici di prima specie della medesima \(F\), e ricordiamo che \textsc{Noether} ha rilevato come si possa « costruire un integrale doppio di prima specie a partire da due integrali semplici che non siano funzioni l’uno dell’altro ».

Ora il fatto che la superficie \(F\) possegga due integrali semplici, \(I_1\) e \(I_2\), funzioni l’uno dell’altro, significa che la \(F\) contiene un fascio irrazionale di curve algebriche \(C\), su cui \(I_1\) e \(I_2\) si mantengono costanti.

Questa osservazione si è presentata per la prima volta a \textsc{Castelnuovo}, nei riguardi delle superficie di genere \(p_g=0\) (cfr. Cap. IX, \S~6); più tardi (1905), per le superficie di genere \(p_g\) qualunque, a \textsc{De Franchis}\origfootnote{1}{s244R-1}{\textit{Sulle superficie algebriche le quali contengono un fascio irrazionale di curve,} Rendic. Circolo Mat. di Palermo (25 Aprile 1905).}, \textsc{Castelnuovo}\origfootnote{2}{s244R-2}{\textit{Sulle superficie aventi il genere aritmetico negativo.} Ib., 14 Maggio 1905 (Memorie scelte, XXVII, pag. 501).} ed \textsc{Enriques}\origfootnote{3}{s244R-3}{\textit{Sulle superficie algebriche che ammettono un gruppo continuo di trasformazioni ecc.} (Ibid., 14 Maggio 1905).} i quali hanno scorto qui il mezzo di studio appropriato per classificare le superficie di genere \(p_a<0\).

La relazione fra il genere \(p_g\) e l’irregolarità \(q\), o fra il \(p_g\) e il \(p_a\), è stata precisata e dimostrata nel modo più semplice da \textsc{Castelnuovo} colla diseguaglianza anzidetta
\[
p_g<2p_a+4,
\]
valida per le superficie \(F\) che non contengono un fascio irrazionale di curve. Frattanto, sulla base di codesta relazione, \textsc{De Franchis}
%CONT>>
% scan: s245L
% !PRINTED 462
\weberpage{462}
% head: CAPITOLO UNDECIMO
%<<CONT
% stamp: Dark scan-edge band at the right margin; diagonal marks near "Comessatti" and "problema", and scattered specks.
enunciava che le \(F\) di genere \(p_a<-1\) posseggono un fascio irrazionale di curve\origfootnote{1}{s245L-1}{Aggiungasi che \textsc{De Franchis} (precisando un risultato anteriore, del 1904) stabilisce che « \textit{ogni piano doppio irregolare possiede un fascio ellittico o iperellittico di curve} ».} mentre \textsc{Castelnuovo} ed \textsc{Enriques} — facendo uso di un ragionamento già adoperato nell’ipotesi \(p_g=0\) (Cap. X, \S~4) — riconoscevano che codeste \textit{superficie sono riferibili a rigate, di genere \(p=-p_a\) \((>1)\)}.

Infine \textsc{Enriques} approfondiva l’analisi delle superficie per cui \(p_a=-1\), riuscendo a riconoscere che esse sono superficie ellittiche o iperellittiche.

Gli sviluppi di cui si è discorso hanno avuto un seguito negli importanti lavori di A. \textsc{Rosenblatt}\origfootnote{2}{s245L-2}{\textit{Sur les surfaces irrégulières ecc.} Rendic. Circolo Mat. di Palermo, 1913.} e A. \textsc{Comessatti}\origfootnote{3}{s245L-3}{Rendic. Acc. Lincei, 1919. e \textit{Intorno alle superficie algebriche irregolari ecc.} Rendic. Circolo Mat. di Palermo, t. XLVI, 1922.} che portano la classificazione delle superficie per cui
\[
p_g\geq2p_a+4.
\]

Già \textsc{Rosenblatt} (1913) riconosce che esse possono ridursi a pochi tipi. \textsc{Comessatti} (1919 e 1922) riprende l’analisi di \textsc{Castelnuovo} trattando il problema più generale di « determinare i legami in termini finiti che passano tra più funzioni di due variabili quando i loro mutui Jacobiani soddisfino a un numero sufficiente di equazioni lineari omogenee a coefficienti costanti ».

Di qui trae, in particolare, la classificazione delle superficie algebriche per cui
\[
p_g\geq2p_a+4,
\]
dimostrando che esse appartengono ai seguenti tipi:

1) rigate di genere \(p>1\);

2) superficie ellittiche con \(p_g>1\) e \(p_a=-1\);

3) superficie rappresentative delle coppie di punti di due curve di genere 2 e \(\pi\): \(p_g=2\pi,\ p_a=\pi-2\);

4) particolari superficie di genere lineare \(p^{(1)}=1\) col \(p_a\geq0\).

Anche di questi brillanti risultati, come della diseguaglianza di \textsc{Castelnuovo}, è desiderabile si riesca a dare una dimostrazione algebrico-geometrica.

% scan: s245R
% !PRINTED 463
\weberpage{463}
% head: CLASSIFICAZIONE GENERALE DELLE SUPERFICIE
% stamp: Scan-edge shading at the left and lower corners; scattered specks and a faint horizontal mark across the classification table.
\edsection{9. Riassunto della classificazione precisata.}{9. Riassunto della classificazione precisata.}{sec.XI.9}

I teoremi esposti sulle superficie di genere numerico negativo, valgono a stringere più da vicino la classificazione generale delle superficie, in particolare portando che per \(p^{(1)}>1\) deve essere \(p_a\geq0\).

Questa deduzione ha un significato d’ordine generale: per le superficie \(D\) (\(P_{12}>1,\ p^{(1)}>1\)), essendo \(p_a\geq0\), si ottiene un sistema bicanonico di dimensione \(P_2-1=p_a+p^{(1)}-1\), e quindi — per \(p^{(1)}>3\) — riesce \(P_2-1\geq3\). Così, a prescindere da un dubbio che deve essere criticamente esaminato\origfootnote{1}{s245R-1}{Come si è fatto per le superficie regolari nei \S\S~18-20 del cap. IV.} circa la riducibilità delle curve bicanoniche, le superficie per cui \(p^{(1)}>3\) ammettono come modello proiettivo una superficie bicanonica (che, in casi particolari, potrà essere multipla).

Ed appare quindi l’interesse speciale di possedere una classificazione completa delle superficie per cui \(p^{(1)}=2\) e \(p^{(1)}=3\), cui si riferiscono gli studi da noi esposti nei \S\S~14-20 del Cap. VIII; tanto più se avvenga di trovare che per tali superficie (a prescindere forse da qualche eccezione) si avveri a posteriori l’ipotesi della regolarità, che nei nostri studi è stata ammessa a priori\origfootnote{2}{s245R-2}{Cfr. la nota a piè di pagina in fine al \S~7 di questo Capitolo.}.

Riassumiamo i risultati più precisi conseguiti nella classificazione generale delle superficie algebriche, valendoci anche di qualche formula il cui significato è stato spiegato innanzi, nel \S~7.

\[
\begin{aligned}
&\begin{tabular}{@{}l@{}l@{}}
\begin{tabular}{@{}>{\raggedright\arraybackslash}p{0.33\linewidth}@{}}
\(P_{12}=0\): rigate \(n>2\pi-2\), curve eccezionali non eliminabili, schiera, continua di trasformazioni non formanti gruppo. % ?? uncertain: the small mark after "schiera" -- a comma (read here) or a scan speck.
\end{tabular}
&
\(\left\{
\begin{tabular}{@{}>{\raggedright\arraybackslash}p{0.61\linewidth}@{}}
\(p_a=0\) superficie razionali \((p_a=P_2=0)\)\\
\(p_a=-1\) rigate ellittiche \((p_a=-1,\)\newline \(\qquad P_4=P_6=0)\)\\
\(p_a<-1\) rigate di genere : \(p=-p_a>1\). % ?? uncertain: punctuation after "genere" -- a colon (read here) or no printed punctuation.
\end{tabular}
\right.\)
\end{tabular}
\end{aligned}
\]

\[
\begin{aligned}
&\begin{tabular}{@{}l@{}l@{}}
\begin{tabular}{@{}>{\raggedright\arraybackslash}p{0.20\linewidth}@{}}
\(P_{12}=1\): curva canonica virtuale d’ord. 0,\\
\(\qquad n=2\pi-2\)\\
\((p^{(1)}=1)\)
\end{tabular}
&
\(\left\{
\begin{tabular}{@{}l@{}l@{}}
\(p_g=1\) &
\(\left\{
\begin{tabular}{@{}>{\raggedright\arraybackslash}p{0.62\linewidth}@{}}
\(p_a=1,\) infinite famiglie-di superficie \((p_a=P_2=1)\)\\ % ?? uncertain: "famiglie-di" -- printed hyphen or scan stroke; the subscript in the parenthetical "p_a" -- a (read here) or g.
% EDITORIAL D0250 RISOLTO: The hyphen is visibly between famiglie and di; the subscript has the closed bowl and short right stroke of a, not the descender of g.
\(p_a=-1,\) infinite famiglie di superficie iperellittiche di divisore \(\delta=1,\ 2,\ \ldotp\ldotp\ldotp\ldotp\)\newline \(\qquad(p_a=-1,\ p_g=P_4=1).\)
\end{tabular}
\right.\)
\\
\(p_g=0\) &
\(\left\{
\begin{tabular}{@{}>{\raggedright\arraybackslash}p{0.62\linewidth}@{}}
\(p_a=0,\) superficie del \(6^{\mathrm{o}}\) ordine passanti doppiamente per gli spigoli d’un tetraedro \((p_a=P_3=0,\ P_2=1).\)\\
\(p_a=-1,\) superficie ellittiche \(I_a\,I_b\,II_a\,II_b\)\newline \(\qquad III_a\,III_b\,III_c\) \((P_2=0,1,\ P_4=0,1,\)\newline \(\qquad P_3=0,1,\ P_4=0,1).\) % [sic] final "P_4=0,1" -> "P_6=0,1", P_4 is repeated in the list of plurigenera.
\end{tabular}
\right.\)
\end{tabular}
\right.\)
\end{tabular}
\end{aligned}
\]

\Needspace{110mm}
% scan: s246L
% !PRINTED 464
\weberpage{464}
% head: CAPITOLO UNDECIMO
% stamp: Scattered scan speckles, dark right-edge shading, and a diagonal stray mark between "un’ascensione" and "alpina".
\[
\begin{aligned}
\begin{gathered}
P_{12}>1\\
(n<2\pi-2)
\end{gathered}
\;
\left\{
\begin{aligned}
p^{(1)}=1\ &\left\{
\begin{aligned}
&p_a=-1,\ \text{superficie ellittiche con un fascio}\\
&\quad\text{di genere }p_g\ \text{di curve ellittiche}\\
&\quad(p_a=-1,\ P_4\ne1).\\[1ex]
&p_a\ge0,\ \text{superficie le cui curve canoniche e}\\
&\quad\text{pluricanoniche sono composte colle curve}\\
&\quad\text{ellittiche d’un fascio di genere}\\
&\quad p_g-p_a:\ \text{dipendono da più caratteri interi}\\
&\quad\text{arbitrarii.}
\end{aligned}
\right.\\[1ex]
p^{(1)}>1\ &\left\{
\begin{aligned}
&p_a\ge0,\ \text{superficie che ammettono un modello}\\
&\quad\text{proiettivo canonico }(p_g>3)\ \text{o}\\
&\quad\text{bicanonico }(p^{(1)}>3)\ \text{ecc.: un numero}\\
&\quad\text{finito di tipi per un dato valore del}\\
&\quad p^{(1)}.
\end{aligned}
\right.
\end{aligned}
\right.
\end{aligned}
\]

A questo punto ci sia consentito fermarci un istante, come in un’ascensione alpina si ama sostare sul picco conseguito e di là contemplare lo spettacolo della Natura che si offre alla vista.

Cinquant’anni or sono s’iniziava in Italia lo studio di queste teorie, appena abbozzate dal genio di un precursore (\textsc{Max Noether}); allora, scherzando sulle difficoltà e le eccezioni che s’incontravano da ogni parte, si soleva dire che, mentre le curve algebriche (già composte in una teoria armonica) sono create da Dio, le superficie invece sono opera del Demonio.

Ora si palesa invece che piacque a Dio di creare per le superficie un ordine di armonie più riposte ove rifulge una meravigliosa bellezza, e ch’Ei volle in esse — diciamo col Poeta —
\begin{verse}
del creator suo spirito\\
più vasta orma stampar.
\end{verse}

La ricchezza delle proprietà e la bellezza, lungamente nascosta, che qui si palesano, non debbono costituire ragione di vano orgoglio per la scuola geometrica italiana o per i matematici stranieri che hanno collaborato a scoprirle, ma piuttosto debbono suscitare un senso di reverenza per quell’ordine meraviglioso degli enti matematici, che il pensiero trova innanzi a sè e quasi raccoglie, al pari delle specie viventi, dalla Natura Madre; e così alimentare la fede dei giovani ricercatori che dietro alle difficoltà, alle eccezioni, alle apparenti incongruenze, c’è realmente in questo mondo di enti, una divina armonia, che gli sforzi concordi degli studiosi riusciranno sempre meglio a mettere in luce. % ?? uncertain: possible accent over the first o in "piuttosto" -- piuttosto / piuttósto.
% scan: s246R
% !PRINTED none
\weberblank{[s246R]}
% head: none
% blank page
% stamp: Dark scan borders at the left and right edges, a short horizontal mark near the upper right, and scattered speckles.

\par\edFlush\clearpage\end{document}
